\documentclass[a4paper,10pt,openany]{ltjsbook}
\usepackage[haranoaji]{luatexja-preset}
\usepackage[margin=22mm]{geometry}
\usepackage{amsmath,amssymb,longtable,booktabs,tcolorbox}
\usepackage[backend=biber,style=numeric,sorting=none]{biblatex}
\addbibresource{references.bib}
\usepackage[hidelinks]{hyperref}
\hypersetup{pdftitle={技術士第一次試験 過去問体系化 作業版},pdfauthor={技術士第一次試験対策プロジェクト}}
\setcounter{tocdepth}{1}
\setcounter{secnumdepth}{1}
\title{技術士第一次試験\\過去問体系化\par\large 作業版：登録・分類・一部解説初稿}
\author{基礎・適性・応用理学／平成23年度〜令和7年度・令和元年度再試験}
\date{2026年10月8日}
\begin{document}
\maketitle
\frontmatter
\chapter*{この作業版の状態}
全1,270問の原典と公式正答を登録し、分類体系の修正と一部問題の独自解説を進めている。
未校正・未執筆の項目はその状態を表示する。本文中の解説例もAstra査読前の初稿である。
このPDFは完成教材の認定や、全問のOCR・解答の確認を意味しない。
\tableofcontents
\mainmatter
\chapter{対象資料と試験回の整理}
対象は平成23〜30年度、令和元〜7年度、令和元年度再試験の16試験回である。
基礎470問、適性240問、応用理学560問を対象とする。その他部門の保存PDFは教材対象に含めない。
問題・正答の根拠は日本技術士会の公式公開資料\cite{official_questions,official_answers}とし、各問題にPDF物理ページを示す。
\section{年度と試験回}
令和元年度の通常試験と再試験を別回として扱う。試験回を分母にする集計は16、年度を分母にする集計は15となる。
同年度に出題があったかを調べる際に、再試験を別年度として加算しない。
\section{原典へ戻る方法}
問題IDは年度、通常／再試験、科目、番号を組み合わせた識別子である。
PDFの物理ページは1から数える。紙面の印刷ページと表紙・空白ページの有無により一致しない場合がある。
\section{公式正答の例外}
複数の選択肢が正答として受理された問題と、全員へ得点が与えられた問題を区別する。
ALLは全選択肢が内容上正しいという意味ではない。訂正された正答は現行公式値を保持し、訂正前の値と注記を別に追跡する。

\chapter{物理専攻者の学習導線}
数学・物理の既知知識を、工学と地球科学の観測・設計・判断へ接続する。
新規用語の数、推論段階、計算負荷、図表の読解を区別して、自分にとって必要な準備を見積もる。
\section{公式問題群と学問体系}
基礎の公式5群は出題上の区分として保持する。一方、勉強するときは線形代数、振動、材料力学、数値解析などの学習分野へ分けて考える。
「固有」という共通語があっても、行列の性質と構造物の固有振動は同じ学習項目にはまとめない。
\section{初めて学ぶ分野の読み方}
情報分野では、状態と操作を明示することが物理の状態変数と時間発展の考え方に対応する。
化学分野では、反応式だけでなく溶液条件と、操作によって残る成分を追う。
地球科学では、一般理論と、実際の自然試料・観測データからの解釈を分ける。
制度・倫理では、出題当時の記述、公式正答、現行資料の説明をそれぞれ記録する。
\section{問題を解いた後}
問題幹の否定、単位、図の境界条件、選択肢番号の対応まで戻って確認する。
同じ解法の問題でも数値や選択肢順が変われば、正答番号を記憶する方法は通用しない。
未確認の類似候補を再出題と確定せず、差分を比較して承認した関係だけを復習に使う。

\chapter{基礎科目の構造}
公式群と学習分野を別軸にする。情報・論理では論理・集合、数値表現、データ構造、アルゴリズム、情報理論、計算機構成、ネットワーク、セキュリティを整理する。
解析では微積分・数値解析、線形代数・ベクトル解析、連続体・構造力学、振動・運動、計算力学、回路を区別する。
OS・DB・AIへの接続は、出題実績のある概念と補足を分けて示す。

\chapter{適性科目の読み方}
法令の義務・責務、職業倫理、組織の説明責任、安全、環境、研究・情報倫理を整理する。
問題の正答番号と、すべての選択肢の現在の妥当性は同じ確認作業ではない。
法改正や統計の更新、出題文の省略・不適切な表現がある場合は、当時資料と現在資料を並べて疑義を残す。
社会的責任規格のようなガイダンスを、法律や認証要求事項と混同しない。

\chapter{応用理学の領域をつなぐ}
物理及び化学、地球物理及び地球化学、地質の3領域を主分類とする。
一般の同位体の定義と、自然試料の同位体比・起源・年代の解釈を区別する。
一般波動、自然地震、地下構造の探査は共通の波動概念を持つが、問いの対象に応じて主分類を選ぶ。
結晶の格子・結合と、鉱物の同定・産状にも同様の区別を置く。
この作業版では領域全体の解説は未完成であり、校正済みの例から本文作成を進める。

\chapter{設計・計画}
\section{構造設計は力の流れと性能を結び付ける}
橋、機械の軸、容器や支持柱は、与えられた力を接合部から支持部へ伝える。
設計では材質を強くするだけでなく、荷重がどう流れ、変形がどこに集まり、支持機構がどの条件で失われるかを調べる。
同じ荷重でも細い柱は材料の破断前に座屈し、強い梁でも過大なたわみで使えなくなる。
本節群は構造強度の初稿である。他の設計・計画領域は後続の原典確認で加筆する。

\section{強度・剛性・安定を分ける}
\label{theory:structural:modes}
強度は破損・降伏等に耐える能力、剛性は荷重に対する変形の小ささ、安定は小さな擾乱が増幅せず平衡を維持する性質である。
ヤング率 \(E\) が同じでも降伏強度は異なり得る。逆に強度の高い材料が必ず高剛性であるとはいえない。
静的な荷重支持機構の喪失、圧縮による座屈、周期外力による共振を区別する。
共振の比較対象は外力周波数と固有振動数であり、応力と振動数を大小比較することは次元的にも誤る。
短く太い柱では圧縮降伏・圧壊が支配し得るが、引張破断強度をそのまま圧縮破損の閾値としない。

\section{引張試験の測定値から定義を作る}
\label{theory:structural:tensile}
試験機が直接取得するのは荷重 \(F\) と標点間の伸び \(\Delta L\) である。
初期断面積 \(A_0\)、初期標点距離 \(L_0\) を用いて
\[
 \sigma_{\mathrm n}=\frac{F}{A_0},\qquad
 \varepsilon_{\mathrm n}=\frac{\Delta L}{L_0}
\]
と規格化すると、試験片の初期寸法の影響を分けて比較できる。
除荷で回復する変形が弾性、残る変形が塑性である。弾性という語だけでは線形比例まで保証しない。
線形範囲では \(\sigma=E\varepsilon\) である\cite{struct_mlit_glossary}。
延性は破断まで塑性変形を許す能力であり、塑性と同じ語ではない。
単調引張の最大公称応力
\[
 \sigma_{\mathrm B}=\frac{\max F}{A_0}
\]
が引張強さである。くびれ後の瞬間断面積 \(A(t)\) を使う真応力 \(F/A(t)\)、破断時の公称応力、
繰返し負荷での疲労限度を混同しない。平成29年度の「応力」の語には、初期断面積を用いる前提を明示して読む。

\section{許容応力と安全率}
\label{theory:structural:allowable}
材料の限界応力を \(\sigma_{\mathrm{str}}\)、安全率を \(n>1\) とした単純なモデルでは
\[
 \sigma_{\mathrm{allow}}=\frac{\sigma_{\mathrm{str}}}{n},
 \qquad \sigma_{\mathrm{working}}\le \sigma_{\mathrm{allow}}
\]
である。安全側へ余裕を取るため許容値は限界値より小さくなる。
ヤング率も応力単位を持つため、強度をヤング率で割った比を許容応力と呼ぶことはできない。
実設計では何を限界応力にするか、温度、荷重種類、疲労、座屈等を分ける必要がある。
安全率を掛けて許容値を大きくする誤りと、材料強度だけで全ての破損形態を代表させる誤りを避ける。

\section{限界状態は一つではない}
\label{theory:structural:limit}
限界状態は要求された性能を満たせなくなる境界である。
通常の使用に必要な変位・振動等を扱う使用限界、最大耐荷性能を扱う終局限界、
繰返し負荷の損傷を扱う疲労限界を分ける。
破壊していなくても利用できなくなる状態があるため、終局だけを照査しても十分ではない。
国土交通省の空港構造設計要領は修復限界も扱い、疲労を終局に含める場合があることを明記しつつ別に分類する\cite{struct_mlit_limit}。
これはすべての構造種別で分類名と係数が同一という意味ではない。過去問で指定された対応と規準の粒度を保存する。

\section{不静定構造の余力とその限界}
\label{theory:structural:redundancy}
静定構造の反力・断面力は適切な平衡式から求められる。不静定構造では余剰な未知力があり、
変形の適合条件と材料の構成関係も用いる。
延性的な部材が局所降伏した後、他の経路へ荷重を再配分できると構造全体の余力が生まれる。
この説明には接合部の能力、十分な変形能力、荷重経路の存在が必要である。
塑性化による材料非線形、たわみによる幾何非線形を扱わず、常に線形比例を仮定して損傷後を解析することはできない。
\begin{tcolorbox}[title=部分損傷から全体安全へ飛躍しない]
不静定次数が高いことだけでは安全性の保証にならない。少数の塑性ヒンジによる機構形成、
全体座屈や接合部破壊で支持能力が失われる場合がある。
平成24年度の初回・公式5は保持するが、「すべての部材が破壊するまで安全」という一般命題へ拡張しない。
\end{tcolorbox}

\section{梁の曲げ：断面二次モーメントと断面係数}
\label{theory:structural:beam}
細長い均質梁の小変形・線形弾性では、平面断面が変形後も平面を保つと近似する。
中立軸からの距離を \(y\)、曲率を \(\kappa\) とすると軸ひずみは \(-y\kappa\)、応力は \(-Ey\kappa\) となる。
断面にわたりモーメントを積分して
\[
 I=\int_A y^2\,dA,\qquad |M|=EI|\kappa|,\qquad
 \sigma_{\max}=\frac{|M|c}{I}=\frac{|M|}{Z},\qquad Z=\frac{I}{c}
\]
を得る。\(c\) は外縁までの距離である。\(I\) の単位は長さの4乗、\(Z\) は3乗である。
断面を中立軸から離して配置すると曲げへの抵抗が増す。これは材料そのものの強度を変更する操作ではない。
梁理論の仮定、断面力、境界条件は大学の公開講義と対応する\cite{struct_mit_beam}。

\section{単純支持梁の等分布荷重を解く}
一定の荷重密度 \(q\)、長さ \(L\) の梁の両端を単純支持する。対称性と平衡から反力は両端 \(qL/2\) である。
左端からの位置 \(x\) でのモーメントは
\[
 M(x)=\frac q2 x(L-x),\qquad M_{\max}=\frac{qL^2}{8}.
\]
下向きたわみを正とする符号規約で \(EI v''=-M\) とし、\(v(0)=v(L)=0\) を用いて積分すると
\[
 v(x)=\frac{q}{24EI}\bigl(x^4-2Lx^3+L^3x\bigr),\qquad
 v_{\max}=v(L/2)=\frac{5qL^4}{384EI}
\]
を得る。最大応力は材料と断面係数、最大たわみは \(EI\) と荷重・長さに依存する。
大たわみ、短く太い梁のせん断変形、断面の局所座屈をこの式に含めたとはしない。

\section{寸法のべき則は固定条件を含む}
\label{theory:structural:scaling}
円断面の直径 \(d\)、矩形断面の幅 \(b\) と荷重方向の高さ \(h\) に対して
\[
 I_{\mathrm{circle}}=\frac{\pi d^4}{64},\quad
 Z_{\mathrm{circle}}=\frac{\pi d^3}{32},\qquad
 I_{\mathrm{rectangle}}=\frac{bh^3}{12},\quad
 Z_{\mathrm{rectangle}}=\frac{bh^2}{6}.
\]
\(q,L,E\) を固定して円梁の径だけを変えると \(\sigma_{\max}\propto d^{-3}\)、\(v_{\max}\propto d^{-4}\) である。
矩形の幅も固定すれば \(\sigma_{\max}\propto h^{-2}\) となる。断面二次モーメントをそのまま応力の分母へ入れると、
外縁距離の1乗を見落とす。
\begin{tcolorbox}[title=自重を固定荷重として扱わない]
自重だけなら \(q=\rho Ag\) である\cite{struct_mit_beam}。円梁の \(A\propto d^2\) も同時に変わるため、
同じ \(\rho,L,E\) では自重応力は \(d^{-1}\)、たわみは \(d^{-2}\) となる。
令和4・6年度の解は通常の固定 \(q\) 比較として示し、寸法を変えれば全ての荷重が一定であるとは主張しない。
\end{tcolorbox}

\section{座屈は圧縮平衡の安定性喪失である}
\label{theory:structural:buckling}
引張は細い定規を直線へ戻す方向に働くが、圧縮は小さな横たわみを増幅し得る。
理想的な等断面柱で横たわみを \(v(x)\)、圧縮力を \(P\) とすると、両端ピンの微小たわみモデルでは
\[
 EI v''+Pv=0,\qquad v(0)=v(L)=0
\]
となる。非自明解は \(v\propto\sin(m\pi x/L)\) であり、最小の固有値は
\[
 P_{\mathrm{cr}}=\frac{\pi^2EI}{(KL)^2}
\]
である。\(K\) は端部拘束に応じた有効長係数である。両端ピン1、固定・自由2、両端固定0.5という理想化を用いる。
同じ拘束と断面なら長さの逆2乗、円断面なら直径の4乗に依存する。
弾性式に降伏強度は直接出ないが、臨界応力 \(P_{\mathrm{cr}}/A\) が弾性域にあることが適用条件である。
細長比が小さい領域の降伏・圧壊をオイラー式だけで判定しない\cite{struct_mit_buckling}。
初期曲がりや荷重偏心のある実柱には、理想臨界前からたわみがある。図の突然の横変形を共振や剥離と呼ばない。

\section{疲労は振幅と平均応力を分ける}
\label{theory:structural:fatigue}
繰返しの荷重・ひずみが損傷を蓄積し、静的破壊強度より小さい負荷でも破損することがある。
応力の最大・最小から
\[
 \sigma_a=\frac{\sigma_{\max}-\sigma_{\min}}{2},\qquad
 \sigma_m=\frac{\sigma_{\max}+\sigma_{\min}}{2}
\]
を定義する。振幅と平均を混ぜず、同じ振幅の比較では通常の金属疲労モデルで平均引張応力が不利となる。
修正Goodman線はその影響を扱う設計近似であるが、材料・残留応力・塑性・多軸条件を問わず同一の寿命式にはならない。
データの温度、応力比、ひずみ制御か荷重制御か、対象繰返し数を読む\cite{struct_nims_data,struct_nasa_mean}。
疲労限度が存在するという理想化もすべての金属や超高サイクル域へ一般化しない。

\section{クリープは時間と温度の条件を持つ}
\label{theory:structural:creep}
金属材料が高温で荷重を受け続けると、一定の負荷でも塑性変形が時間とともに進む。
瞬間的なヤング率や単調引張強度だけでは長期の変形・破断時間を決められない。
NIMSの資料は高温機器の設計に、温度・荷重・経過時間を揃えた長時間試験が必要なことを示す\cite{struct_nims_data}。
高強度という語を高温という条件へ置換しない。また低温から常に全材料の時間依存変形がゼロとなるという定理にはしない。

\section{応力集中は局所と公称の差である}
\label{theory:structural:concentration}
穴、切欠き、段差等の幾何変化は応力の流れを曲げる。局所最大と遠方の公称応力の比で集中を評価する。
半径を緩やかにする等の設計で集中を減らすことがあるが、その値は形状・荷重・境界条件に依存する。
一軸応力ではvon Mises相当応力は軸応力の絶対値に一致するが、一般の多軸場は
\[
 \sigma_{\mathrm{VM}}=
 \sqrt{\frac{(\sigma_1-\sigma_2)^2+(\sigma_2-\sigma_3)^2+(\sigma_3-\sigma_1)^2}{2}}
\]
で評価する。主応力 \(\sigma_i\) の静水圧部分はこの値へ寄与しない。
断面を急変させたという言葉だけで、どの荷重・どの点の相当応力も必ず高くなるとは断言しない。

\section{比強度と価格当たり性能}
\label{theory:structural:specific}
軽量な荷重支持を考えるとき、比強度 \(\sigma/\rho\) は同じ質量でどれだけ引張荷重を支えられるかの指標となる。
価格が質量当たり \(p\) なら、単位比強度当たりの価格は \(p/(\sigma/\rho)\) である。
二つの材料で比強度が何倍かを求めてから、価格比をその倍率で割る。
令和5年度の値ではCFRPの比強度が鋼の24.6875倍、価格が266.6667倍なので、
単位比強度当たりの価格は約10.80倍となる。強度優位と費用優位は別である。
この価格は設問所与であり現在の市場価格ではない。繊維方向、圧縮・接合、加工、疲労等の性能をこの一指標へ圧縮しない。

\section{強度と荷重の確率には同時分布が必要である}
\label{theory:structural:reliability}
強度 \(R\) と作用 \(S\) にばらつきがあるモデルでは、破壊事象を \(S>R\) とする。
一般に
\[
 P_{\mathrm f}=\iint_{s>r} f_{R,S}(r,s)\,ds\,dr
\]
であり、独立なら \(f_{R,S}=f_R f_S\) と因数分解できる。
平成29年度の \(R\sim U[18,23]\)、\(S\sim U[10,20]\) について独立を補うと、
破壊領域は一辺2の直角三角形、支持矩形の面積は50なので \(P_{\mathrm f}=2/50=0.04\) となる。
\begin{tcolorbox}[title=周辺分布だけでは決まらない]
原文には独立性の明記がない。同じ \(U\sim U[0,1]\) から
\(R=18+5U,\ S=10+10U\) とすると破壊確率0である。
一方 \(R=18+5U,\ S=20-10U\) とすると破壊条件は \(U<2/15\) で、確率2/15となる。
どちらも所与の周辺一様分布を持つ。公式0.04を保持しつつ、独立モデルを用いた条件付き解であることを残す。
\end{tcolorbox}

\section{安全率の費用最小化と規範的判断を分ける}
\label{theory:structural:optimization}
対策を増やすと直接費用は増し、事故等の損害費用は減るという所与のモデルを考える。
\[
 C(x)=\frac{a}{1+bx}+cx,\quad
 C'(x)=-\frac{ab}{(1+bx)^2}+c,\quad
 C''(x)=\frac{2ab^2}{(1+bx)^3}.
\]
\(a,b,c>0\) かつ \(1+bx>0\) の域では凸であり、停留点が定義域内なら一意の最小である。
根が域外なら許容端点を比較し、開区間では端点に最小が実現しない場合もある。
令和元年度再試験の \((a,b,c)=(90,4,10)\)、\(x\ge1\) では \(x=1.25\)、
令和4年度の \((9,1,1)\)、\(x>1\) では \(x=2\) である。全費用の単位を億円にそろえる。
この数学的最小化は与式の範囲内の答えであり、法的な最低要求、重大危害、ALARPの著しい不均衡を
単純な費用関数だけで判断できるという主張ではない。

\section{典型的な問われ方と確認順序}
語の穴埋めでは測定値と規格化した量、材料物性と形状の効果、荷重形態と破損形態を分ける。
寸法の比例問題では断面二次モーメントと断面係数を先に書き、固定した荷重を確認する。
最適化では定義域と単位、確率では独立性と領域を先に読む。
同じ定規や柱の題材でも年度によって候補の順が違うため、正答番号だけを移さない。
個別問題から以上の共通理論へ参照し、全問の原典確認と分類を終えるまでは頻度や将来の出題を断定しない。

\section{信頼性を確率モデルとして設計する背景}
\label{theory:dependability:background}
強度を満たしても、制御、電源、通信など一つの必要機能が止まれば装置全体が使えなくなる。
一方、代替経路を設ければ一部の停止を許せる。この違いを、各部の成功又は故障という事象の論理関係へ写すのが信頼性設計である。
対象機能、環境条件、観測期間を先に定め、その期間内に要求機能を維持する確率を信頼度とする。
部品の配置図は電気回路そのものとは限らず、機能が成立する条件を表す。

\section{信頼度、寿命分布、故障率}
\label{theory:dependability:definition}
最初の故障までの寿命を非負確率変数\(T\)とし、
\[
 R(t)=P(T>t),\qquad F(t)=1-R(t),\qquad
 h(t)=\frac{f(t)}{R(t)}
\]
と定める。最後の式は密度\(f\)が存在し\(R(t)>0\)の範囲での条件付き故障率である。
信頼度は無次元、故障率は時間の逆数の次元を持つ。一定故障率\(\lambda\)を仮定した場合に限り\(R(t)=e^{-\lambda t}\)となる。
直並列の論理を扱う設問は時点又は期間を固定した信頼度を与えるため、指数寿命分布を無断で追加する必要はない。
\cite{reliability_nist_series}

\section{直列系：全ての必要機能が成功する}
\label{theory:dependability:series}
要素\(i\)が成功する事象を\(S_i\)とする。全てを必要とする系の成功は\(\bigcap_i S_i\)である。
一般には条件付き確率の積であり、独立を仮定したときだけ
\[
 R_{\mathrm{series}}=P\left(\bigcap_i S_i\right)=\prod_i r_i,
 \qquad r_i=P(S_i)
\]
となる。同じ信頼度\(r\)の要素を増やすと\(r^n\)であり、\(0<r<1\)では低下する。
最初の一要素故障が系を止める非修復モデルの条件である。\cite{reliability_nist_series}

\section{並列系：少なくとも一つの経路が成功する}
\label{theory:dependability:parallel}
系の停止は全経路の停止である。補集合を使うと、独立な要素について
\[
 R_{\mathrm{parallel}}
 =1-P\left(\bigcap_i S_i^{\mathrm c}\right)
 =1-\prod_i(1-r_i)
\]
となる。同一要素\(n\)個なら\(1-(1-r)^n\)である。
「並列」という語だけで成功率を足すと、二経路が同時成功する重複を数えて1を超え得る。
待機予備の切替故障、共通電源、負荷再配分による故障率変化は、全要素が独立して動くこの模型から分ける。\cite{reliability_nist_parallel}

\section{独立性がない場合と共通原因故障}
\label{theory:dependability:independence}
二要素では、周辺成功率\(p,q\)が分かっても
\[
 \max(0,p+q-1)\le P(S_1\cap S_2)\le\min(p,q)
\]
の幅が残る。下限は和事象の確率が1以下、上限は共通事象が各事象に含まれることから出る。
同じ地震、冷却系、電源又はソフトウェアの誤りで停止すれば、機器の台数が二つでも独立とは限らない。
例えば信頼度0.9の二要素が完全に同じ成功事象を持つなら並列系も0.9であり、独立模型の0.99にはならない。
原典が独立を明示しない計算問題では、公式の想定模型と条件不足を併記する。

\section{混合系の縮約と未知信頼度の逆算}
\label{theory:dependability:mixed}
独立な別要素からなる部分系を内側から一つの成功事象へまとめ、次の接続を読む。
縮約後の部分系同士が要素を共有すると、元の部品が独立でも部分系は独立ではない。
未知信頼度\(x\)を逆算する場合はまず\(0\le x\le1\)を課し、系の信頼度との区別を保つ。
二系の等式から\(x\)を割る操作は\(x>0\)が確認された場合に限る。
原典に非零条件がないときは、零根を都合よく消さない。全体の信頼度を問う候補へ単体\(x\)を入れない。

\section{要求信頼度を満たす最小冗長数}
\label{theory:dependability:redundancy}
独立並列の要素信頼度\(r\)、その外に直列で必要な一要素の信頼度\(s\)を置けば
\[
 R_n=s\{1-(1-r)^n\}.
\]
要求\(R_n\ge R_*\)を変形し、\(0<r<1\)かつ\(R_*<s\)なら
\[
 (1-r)^n\le1-\frac{R_*}{s},\qquad
 n\ge\frac{\log(1-R_*/s)}{\log(1-r)}.
\]
負の分母で割るため不等号の向きに注意する。最小整数の前後を直接代入して検算する。
有限の\(n\)で\(R_n<s\)なので、外の直列要素が達成可能な上限を作る。
これを改善せず並列部だけ増やしても要求を満たせない場合がある。

\section{多数決冗長系}
\label{theory:dependability:voting}
\(n\)個のうち任意の\(k\)個以上が成功すればよい場合、独立・同一信頼度\(r\)の模型では
\[
 R_{k/n}=\sum_{j=k}^{n}\binom{n}{j}r^j(1-r)^{n-j}.
\]
二項係数は、成功する要素の組合せを数えるものである。成功個数が異なる事象は排反なので和を取れる。
2/3系は\(3r^2(1-r)+r^3\)であり、全並列の1/3系とも全直列の3/3系とも異なる。
判定器自身、同じ故障出力を返す共通原因、異なる単体信頼度は別の模型を要する。\cite{reliability_nist_voting}

\section{故障の木：原因事象から上位事象へ}
\label{theory:dependability:fta}
故障の木解析（FTA）は、望ましくない上位事象がどの下位事象の論理結合で起きるかを表す。
独立な下位故障確率を\(p_i\)とした場合、ANDは\(\prod_i p_i\)、ORは\(1-\prod_i(1-p_i)\)となる。
信頼性ブロック図の成功事象とFTAの故障事象を区別する。形が似ていても、どちらの事象を入力としているかを先に読む。
末端から段階的に縮約し、同じ原因が複数枝へ再出現する場合は事象を重複した独立変数にしない。
独立でない二事象のORは常に\(p_1+p_2-P(E_1\cap E_2)\)であり、最後の項を勝手に\(p_1p_2\)へ置換できない。

\section{修復可能な系の稼働率}
\label{theory:dependability:availability}
信頼度は途中で故障しない確率であるが、稼働率は故障後の修復も含めて所要時点又は期間に使える割合である。
平均運転時間を\(m_{\mathrm u}\)、平均修復を含む非稼働時間を\(m_{\mathrm d}\)とする。
交互の運転・停止周期が安定し有限平均を持つ模型で、長期の時間割合は
\[
 \mathcal A=\frac{m_{\mathrm u}}{m_{\mathrm u}+m_{\mathrm d}}.
\]
多周期の運転時間総和を周期時間総和で割ると、平均の比としてこの式が出る。指数分布は必要条件ではない。
試験でMTBFを平均運転時間、MTTRを平均修復時間とする場合はMTBF/(MTBF+MTTR)となる。
同じ名称でも暦時間の故障間隔を含める定義を代入すると分母を二重計上する。
物流待ち、点検、計画停止を含む運用稼働率には、その時間範囲に対応する分母を使う。

\section{保全：故障前後と処置時期の判断基準}
\label{theory:dependability:maintenance}
予防保全は故障を避ける目的で故障前に処置し、事後保全は故障後に機能を回復する。
時間計画保全は暦又は運転時間等から時期を予定し、状態監視保全は測定や点検で把握する劣化状態から処置を決める。
点検そのものが定期的でも、処置の判断が状態に基づくなら時間計画の交換とは区別される。
国交省の設備保全資料は、重大な機能影響と状態監視の可能性で三種の使い分けを示す。\cite{maintenance_mlit_2013}
全設備へ同じ交換周期を当てはめず、故障の帰結、兆候の観測可能性、保全の費用を検討する。
JIS・IECの版ごとの全用語一致はここでは未確認である。

\section{安全側、縮退、誤操作対策、耐故障性}
\label{theory:dependability:handling}
故障時に安全を維持するフェールセーフ、影響を受ける機能を優先順位付きで終了するフェールソフト、
人為的な誤りを重大な失敗へ結び付けないフールプルーフ、想定する故障があっても機能するフォールトトレランスを分ける。
例えば機械の不意の再起動を防ぐ仕組みは故障時の安全条件を定め、開いたガードで起動しない仕組みは誤操作への備えにもなる。\cite{design_mhlw_failsafe,design_mhlw_foolproof}
耐故障性には検出・隔離・回復や冗長性が関わり、影響を外へ見せないフォールトマスキングはその一方式である。
航空宇宙のソフトウェアでも故障の予防と発生後の影響緩和を分けて設計する。\cite{design_nasa_tolerance}
安全に停止できる設備と機能継続が安全上必要な設備では要求が異なるため、「止めれば常に安全」とはしない。

\section{安全率の値を比べる前に定義を比べる}
\label{theory:dependability:factor}
安全率\(n\)は、基準強さ\(\sigma_{\mathrm{ref}}\)から許容応力を
\(\sigma_{\mathrm{allow}}=\sigma_{\mathrm{ref}}/n\)とする模型で、無次元の余裕を与える。
基準が引張強さか降伏強さか疲労限度かを先に指定する。同じ材料でも荷重履歴や破損形態が違えば値を転用しない。
設問の2以上又は1.3～2.0という代表値だけで、全設計に共通する法定下限を作ることはできない。
\begin{tcolorbox}[title=対象と資料時点を限定した例]
2011年の米国航空機規則25.303では、外力としての限界荷重に原則1.5を掛ける。ロケット全般の規定ではない。
厚労省のクレーン等安全規則213条では、玉掛けワイヤロープの切断荷重/最大荷重を6以上とする。
食品安全委員会の2008年ADI説明は、残留農薬等の無毒性量を通常100で割る。治療薬全般の一律係数ではない。
国交省の平成25年改正反映告示表では、対象条件を満たすエレベーター主索の常時使用時4.0、設置時5.0であり端部の値は別である。
この告示の2011当時版及び現行までの全改正履歴は未照合である。
\end{tcolorbox}
\cite{design_faa_2011,design_mhlw_crane,design_fsc_adi,design_mlit_elevator}
異なる分母と対象を持つ係数の大小を、危険又は安全の単純順位へ変換しない。

\section{設計管理と製造物責任法の安全性}
\label{theory:dependability:liability}
大量生産品では、利用者が製造過程の過失を立証することが難しい。この背景で製造物の欠陥に基づく責任を定め、被害者の保護を図るのが製造物責任法の目的である。
第2条の欠陥は、特性、通常予見される使用形態、引渡時期などを考慮した通常有すべき安全性の不足を意味する。
機能性の不足、故障、破損だけで直ちに全て法定の欠陥になるわけではなく、絶対安全の要求でもない。
消費者庁は製造上、設計上、指示・警告上の三側面を説明している。\cite{design_caa_pl1,design_caa_pl2}
設計者は必要機能だけでなく、合理的に予見する使用、危険の除去、残留する危険の伝達を一体に検討する。
\begin{tcolorbox}[title=確認した法令時点]
ここは各出題原典の第1条・第2条の引用と平成30年9月の消費者庁逐条解説を比較した説明である。
e-Gov現行本文の取得は未完了で、時効等を含む出題年から現行までの全改正履歴を確認した説明ではない。
個別問題ではこの確認範囲を部分校正として追跡する。
\end{tcolorbox}

\section{単体法、双対、多目的設計}
\label{theory:dependability:optimization}
線形計画は線形の目的と制約で資源配分を表す。単体法は基底を交換し可行な端点間を移る。
退化があると単純な「毎回必ず厳密改善する」という説明は成立せず、循環を避ける選択則が必要になる。
例えば主問題を\(\max c^{\mathsf T}x,\ Ax\le b,\ x\ge0\)と置けば、
双対は\(\min b^{\mathsf T}y,\ A^{\mathsf T}y\ge c,\ y\ge0\)である。
両方の可行解について
\[
 c^{\mathsf T}x\le y^{\mathsf T}Ax\le y^{\mathsf T}b
\]
なので双対は主目的の上界を与える。可行で有限な最適値のある線形計画で両最適値は一致する。\cite{design_mit_lp}
逆問題は観測から原因を推定する別の問題設定であり、双対の呼称の代わりにはならない。
多目的最小化では、費用、重量、性能、環境負荷などの改善が相反するトレードオフが生じる。
パレート最適とは、全目的で悪化せず少なくとも一目的を厳密に改善する他の可行解が存在しないことである。
「全目的を同時に厳密改善できない」だけは弱パレートの条件である。
可行目的値\(\{(0,1),(0,0)\}\)で(0,1)は弱パレートだが(0,0)に支配されるため通常のパレートではない。

\section{信頼性・設計安全の典型的な読み方}
回路では成功経路、FTAでは故障原因を読み、同じ図形でも事象の向きを取り違えない。
積を使う前に独立性を、根を除く前に確率の範囲と非零条件を確認する。
保全では故障前後と処置時期の判断基準、法令では目的と定義、安全率では分子分母と資料時点を明示する。
過去問の組合せ表は候補の順を年度ごとに確認し、同じ概念を学ぶことと同じ候補番号を記憶することを分ける。

\section{品質を工程と判断の関係で考える}
\label{theory:quality:background}
品質管理で問うのは、ばらつきの小ささだけではない。何を満たすべきかという要求、
どのように作るかという工程、どの情報で異常を判断するかという観測を結ぶことである。
良い製品を作ること、工程が安定していること、出荷ロットを受け入れることは判断対象が違う。
目標を設定し、工程を観測して改善する流れの中で、確率統計は限られた測定から判断する道具となる。

\section{平均・分散・共分散から合成する}
\label{theory:quality:variance}
有限な二次モーメントを持つ変数を列ベクトル \(X\)、係数を \(a\) とする。
線形結合 \(Y=a^{\mathsf T}X\) の平均は \(E[Y]=a^{\mathsf T}E[X]\) である。
平均からの偏差を二乗して展開すると
\[
 \operatorname{Var}(Y)
 =\sum_i a_i^2\sigma_i^2+2\sum_{i<j}a_i a_j\operatorname{Cov}(X_i,X_j)
 =a^{\mathsf T}\Sigma a
\]
となる。\(\sigma_i^2\) は成分の分散、\(\Sigma\) は共分散行列である。
平均の線形性には独立性は不要である。分散から共分散項を消すには無相関が必要で、
独立性はその十分条件となる。独立な差 \(X-Y\) でも係数 \(-1\) が二乗されるため
分散は和となる。共通加工機の温度変動などがあれば独立とは限らず、共分散を捨てない。
単位も二乗され、長さの分散は長さの二乗、標準偏差は長さである。\cite{stat_psu_linear}

\section{正規分布の閉性と上側確率}
\label{theory:quality:normal}
\(X\sim N(\mu,\sigma^2)\) の第二引数は分散とする。
独立な正規変数の線形結合も正規となる。これは特性関数
\(\exp(i\mu t-\sigma^2t^2/2)\) の積を取れば、平均は係数付き和、
分散は係数の二乗付き和になることから分かる。
一般には各成分が正規であるだけでは足りず、同時分布の条件も必要である。
標準化 \(Z=(X-\mu)/\sigma\) により
\[
 \Pr(X\ge x)=1-\Phi\!\left(\frac{x-\mu}{\sigma}\right)
\]
と表せる。\(\Phi\) は標準正規累積分布であり、右側の面積を求める。
試験表が「上側確率」を示すか、左側の累積確率を示すかを先に読む。
差の平均に符号を付け、分散を合成し、平方根で標準偏差を出してから標準化する。\cite{stat_nist_normal}

\section{標本平均と総量は別の変数である}
\label{theory:quality:sample_mean}
同じ平均 \(\mu\)、分散 \(\sigma^2\) を持つ独立な \(n\) 個について、
総量の期待値は \(n\mu\)、分散は \(n\sigma^2\) である。
標本平均は総量を \(n\) で割るため、期待値は \(\mu\)、分散は
\(\sigma^2/n\)、標準誤差は \(\sigma/\sqrt n\) となる。
物理量の平均を取っても期待値を標本数で割って小さくするわけではない。
原変数が正規なら有限の \(n\) でもこの平均は正規である。
中心極限定理による大標本の近似とは別の根拠である。
相関のある標本では共分散項が残り、標本数を増やすだけで同じ割合で不確かさが減るとは限らない。

\section{応力・強度から故障確率へ}
\label{theory:quality:stress_strength}
強度 \(R\) と作用応力 \(S\) の単位を揃え、限界状態 \(G=R-S\) を定める。
独立正規モデルでは
\[
 b=\frac{\mu_R-\mu_S}{\sqrt{\sigma_R^2+\sigma_S^2}},\qquad
 \Pr(G<0)=\Phi(-b)
\]
となる。\(b\) は平均の余裕をばらつきで規格化した無次元の指標である。
平均強度が大きくてもばらつきが大きければ故障確率は大きくなり得る。
安全率の比とこの確率を同一視しない。
非正規の裾や強度と応力の依存があると、この式の直接適用はできない。
接続部のすき間も内径から外径を引く限界状態であり、平均の大小は全個体の符号を保証しない。

\section{どの量を数えるかで分布を選ぶ}
\label{theory:quality:models}
同じ「ばらつき」でも、測る量の範囲が違う。公平なサイコロは有限集合の離散一様である。
一定の成功確率 \(p\) の独立試行を \(n\) 回行った成功数は二項分布で、
\[
 \Pr(K=k)={n\choose k}p^k(1-p)^{n-k},\qquad
 E[K]=np,\quad \operatorname{Var}(K)=np(1-p)
\]
である。製造数を固定して不適合品を数えるモデルもこれに対応する。
コインが公平でなくても、一定 \(p\) と独立性があれば二項である。
品質の共通原因による連発や、試行ごとの確率変化があればこの条件を満たさない。
\cite{stat_nist_binomial}

\section{ポアソンの回数と指数の待ち時間}
\label{theory:quality:poisson}
二項分布で機会を増やし、小確率にして \(np=\lambda\) を有限に保つと、
個数は \(\Pr(K=k)=e^{-\lambda}\lambda^k/k!\) のポアソンへ近づく。
平均と分散はともに \(\lambda\) である。
一定率 \(r\) のポアソン過程なら時間 \(t\) の回数の平均は \(rt\) で、
一度も生じない確率が \(e^{-rt}\) となる。
したがって次の事象までの待ち時間 \(T\) は
\(\Pr(T>t)=e^{-rt}\) の指数分布となり、平均は \(1/r\) である。
率 \(r\) と区間内平均 \(\lambda=rt\) の単位を区別する。
災害や交通事故の実測に対する適合は別に検証する必要がある。
大きな \(\lambda\) ではポアソンの歪度 \(1/\sqrt{\lambda}\) が小さくなり、
平均 \(\lambda\)、分散 \(\lambda\) の正規近似が有効になり得る。
「回数だから正規近似は常に不可能」とはしない。\cite{stat_nist_poisson,stat_nist_exponential}

\section{要求品質を設計特性へ展開する}
\label{theory:quality:qfd}
利用者の言葉である要求品質を、設計で調整する品質要素へ変換するのが品質機能展開である。
試験の品質表では要求 \(i\) の重要度 \(w_i\) と要素 \(j\) との関係係数 \(c_{ij}\) を使い、
\(v_j=\sum_i w_i c_{ij}\) を求める。単なる丸の数ではなく、重い要求との対応が寄与する。
係数尺度は判断の結果に影響するため、指定された◎・○・△の値を守る。
この加重得点は実測の成功確率や安全性そのものではない。
関連する特性間の相反もあり、得点最大の列だけを設計の唯一の目標にはしない。
\cite{quality_qfd_scale}

\section{検査と上流の品質保証}
\label{theory:quality:assurance}
検査は要求への適合を判定する活動である。検査で不適合を見つけても、
それだけで不適合を作る原因は除かれない。源流管理は仕様や工程など上流にさかのぼって
発生原因を抑える考え方である。
標準化で手順と判断を明確にし、維持管理と改善を結ぶ。
工程の標準は経験的な固定ルールではなく、要求・観測・変更理由を追跡できるものとする。
QC工程表は工程と管理項目などを結ぶ計画の道具であり、統計量を時間順に並べる管理図と違う。
\cite{quality_process_jsqc}

\section{PDCAをフィードバックとして読む}
\label{theory:quality:pdca}
Planは目標と達成するプロセスを定め、Doは実施、Checkは方針・目標・要求に対する
結果を監視・測定し、Actは性能を改善する行動を取る段階である。
目標なしの測定、測定なしの改善では循環の根拠が失われる。
CをChangeやChoose、AをAdjustやAchieveという近い一般語へ置き換えない。
Actの結果を次の計画と標準へ戻すことに技術的意味がある。
\cite{quality_iso_process}

\section{2008年版を扱う問題の資料時点}
\label{theory:quality:historical}
H27の設問はISO 9001:2008/JIS Q 9001:2008を指定する。
ISOの当時の実施ガイダンスN836でもPlan-Do-Check-Actの利用が示される。
後版である2015年版の公式解説は、プロセスの管理とリスクに基づく思考を併せて説明する。
ここでは一般の四段階を両資料で照合したのであり、2008年版の全条文を取得・照合した
とは扱わない。現行規格の詳細を求める場合はその版・追補を別途確認する。
\cite{quality_iso2008_guidance,quality_iso_process}

\section{工程能力と規格限界}
\label{theory:quality:capability}
工程能力は、安定した工程が作る特性のばらつきと規格の許容幅を比較する。
正規モデルで平均 \(\mu\)、標準偏差 \(\sigma\)、規格上限 \(U\)、下限 \(L\) とすると
\[
 C_p=\frac{U-L}{6\sigma},\qquad
 C_{pk}=\min\left(\frac{U-\mu}{3\sigma},\frac{\mu-L}{3\sigma}\right)
\]
である。\(C_p\) は幅の比較だけなので平均の片寄りを直接反映しない。
\(C_{pk}\) は近い規格限界までの余裕を見る。
両者は無次元であり、標準偏差そのものではない。
非正規、工程の漂移、異なる集団の混合では、指数から不適合率への読み替えに注意する。
一定の数値を全産業共通の合格基準として勝手に設定しない。\cite{quality_nist_capability}

\section{管理図・ヒストグラム・散布図}
\label{theory:quality:charts}
管理図は工程の特性又は統計量を時間・標本順に並べ、安定工程の変動を基準に異常を探す。
中心線と管理限界は観測モデルに基づき、利用者の規格限界とは目的が違う。
限界内に全点があっても、片側への連続や系統的変化があれば安定とは限らない。
ヒストグラムは一変数の測定値を区間に分けて分布を示す。
区間幅が異なるなら高さを度数密度にし、長方形の面積を度数に対応させる。
散布図は対になった二変数の関連を示す。
分布、時間変動、二変数の関連という軸を分ければ、図の名を暗記するだけでなく用途が分かる。
\cite{quality_nist_control,quality_nist_histogram}

\section{相関は線形関連を測る}
\label{theory:quality:correlation}
有限かつ非零の分散について相関は
\(\rho=\operatorname{Cov}(X,Y)/(\sigma_X\sigma_Y)\) であり、
Cauchy--Schwarzから \(-1\le\rho\le1\) となる。
端点は完全な線形関係で達成されるため除外しない。
独立なら無相関だが、無相関だけで独立とは限らない。
対称な \(X\) と \(Y=X^2\) はその反例で、三次モーメントが零でも関数の依存が残る。
相関の絶対値は線形関連の強さを表し、因果の方向や共通原因を判別しない。
正の相関が1へ近づく説明を、任意のデータが完全な直線上にあるという断定へ変えない。

\section{回帰と決定係数の条件}
\label{theory:quality:regression}
回帰は説明変数から応答を予測するモデルを作り、残差などで適合を評価する。
切片を含む一説明変数の通常の最小二乗法では、偏差平方和を
回帰が説明する部分と残差の部分に分けられ、決定係数は \(R^2=r^2\) である。
相関の絶対値 \(|r|\) とは違う。
切片なし、多変数、異なるモデルで同じ式を機械的に使わない。
主成分分析は複数変数を合成して分散を捉える次元縮約であり、指定された応答を予測する
回帰と目的が違う。相関と回帰のどちらも観測だけから因果の証明にはならない。
\cite{stat_psu_correlation}

\section{ロットと抜取検査の二つの危険}
\label{theory:quality:acceptance}
全数を検査できない場合、サンプルからロットを受け入れるか判断する。
計数値では不適合品などを数え、計量値では長さなどの連続量を使う。
良いロットを拒否する生産者危険と、悪いロットを受け入れる消費者危険を分ける。
母集団の不適合率 \(p\) を決めた条件付き受入確率を扱うので、
「受け入れたロットが良い確率」という逆向きの条件付き確率とは違う。
\cite{quality_nist_sampling}

\section{OC曲線から閾値の効果を導く}
\label{theory:quality:oc}
大きなロットから近似的に独立に \(n\) 個を抽出し、不適合数 \(K\le c\) なら合格とする。
受入確率は二項の累積で
\[
 P_a(p)=\sum_{k=0}^c {n\choose k}p^k(1-p)^{n-k}
\]
である。良い品質点を \(p_{\mathrm g}\)、悪い品質点を \(p_{\mathrm b}\) とすれば、
\(\alpha=1-P_a(p_{\mathrm g})\)、\(\beta=P_a(p_{\mathrm b})\) となる。
\(n\) を固定して \(c\) を増やすと受入集合が広がるため、\(\alpha\) は下がり
\(\beta\) は上がる。この包含関係がトレードオフの根拠である。
\((n,c)\) の両方を固定すれば危険も固定であり、原典の「手順固定」の範囲は注意を要する。
小さい有限ロットからの非復元抽出は超幾何となり、二項の独立近似を無条件に使わない。
\cite{quality_nist_oc}

\section{ロット間の切換えと複数回の抽出}
\label{theory:quality:switching}
先行ロットの成績で後続の検査をなみ検査・きつい検査へ切り換える考え方と、
一つのロットで段階的に標本を追加する多回抜取を区別する。
後者は標本を追加したうえで合否を決める手順である。
切換えの回数や復帰条件は適用する規格に従うものであり、今回取得していない
出題時のJIS全文の数値規則を補わない。\cite{quality_nist_sampling}

\section{仮説検定は誤りを管理する判断である}
\label{theory:quality:hypothesis}
帰無仮説 \(H_0\) を基準とし、対立仮説に対してどれほど不整合な観測かを判断する。
真の \(H_0\) を棄却するのが第1種、偽の \(H_0\) を棄却しないのが第2種の誤りである。
有意水準 \(\alpha\) は第1種の誤りを管理する事前の基準であり、
観測後に \(H_0\) が真である確率ではない。
棄却されなかったことも \(H_0\) の真の証明にならない。
第2種の確率 \(\beta\) は具体的な対立側の値に依存し、検出力は \(1-\beta\) である。
同じ二種類の誤りが、品質の受入判断では生産者と消費者の危険に対応する。
\cite{quality_nist_sampling}

\section{t統計量と信頼区間の被覆}
\label{theory:quality:t_interval}
独立な正規標本で母分散が未知なら、不偏分散
\(s^2=\sum_i(X_i-\bar X)^2/(n-1)\) を使う。
平均を推定した残差は和が零という一つの制約を持つため自由度は \(n-1\) となる。
\[
 T=\frac{\bar X-\mu_0}{s/\sqrt n},\qquad
 \bar X\pm t_{1-\alpha/2,n-1}\frac{s}{\sqrt n}
\]
は、順に \(H_0\) 下での一標本t統計量と母平均の両側信頼区間である。
信頼係数は \(1-\alpha\) であって \(\alpha\) の逆数ではない。
反復標本で作る区間のうち母平均を覆う割合を表すので、観測後の一つの固定区間について
母平均がそこに存在する事後確率とはしない。\cite{stat_nist_interval}

\section{典型的な問われ方}
寸法・力の計算は変数の和差と分散の合成を先に決め、正規表で候補の範囲を判定する。
品質表は指定された関係尺度を行の重要度と掛け、列を合計する。
知識問題は要求適合、工程監視、工程能力、受入検査の判断対象を分ける。
各年度の組合せ表の列と番号は原典で照合し、同じ理論でも候補番号が同じと仮定しない。
ここで扱う例は今回確認した問題群であり、全試験回の頻度の確定は別の集計に従う。


\section{先行関係と並行作業：機械と部品の二つの順序}
\label{theory:discrete:precedence}
複数の機械で複数の部品を加工するとき、同じ機械を同時に使わない資源条件と、同じ部品の加工順の条件を別々に書く。機械AとBが部品p,q,rを各p→q→r順に処理し、各部品がA→Bへ進むなら、A列の順、B列の順、各部品のAからBへの順が全て必要である。

アロー図の矢印は作業、節点はその完了・開始の条件を表す。ある節点へ複数の矢印が入るなら、その全てが完了してから次へ進む。ダミー作業は実際の加工を増やすのでなく、必要な先行関係を追加する。例えばAqとBpの両方が終わってからBqを始めたい場合、Bpの終了節点へAqの終了からダミーを入れればよい。

候補の割当ては、見た目の左右で判断せず、各機械の順序と各部品の順序を一つずつ照合する。不要な待ちを追加するだけの場合と、必要な順を保証しない場合も異なる。並行できるという記述を全作業同時開始と考えず、共有する機械・未完了の前工程の有無を調べる。


\section{過去問と個別の解法}

\section{M/M/1は待ちと処理を分ける}\label{theory:design:queue}
到着がPoisson過程、処理時間が指数分布、窓口1、十分な待ち場所の定常モデルを考える。到着率\(\lambda\)、サービス率\(\mu\)は同じ時間単位へそろえる。平均処理時間がtなら\(\mu=1/t\)である。利用率\(\rho=\lambda/\mu<1\)が定常分布の条件である。出生死滅の釣合いから\(p_n=(1-\rho)\rho^n\)となり、
\[L=\sum_{n=0}^{\infty}np_n=\frac{\rho}{1-\rho},\qquad W=\frac L\lambda=\frac1{\mu-\lambda}.\]
待ちだけなら\(W_q=W-1/\mu\)、\(L_q=\lambda W_q=\rho^2/(1-\rho)\)。系内人数Lと行列人数Lqを混同しない。旧原典はLを待ち行列長と称するが、補助計算\(L/\mu\)はWqに等しく、処理時間を足す手順は正しい。名称と数式の検証を分ける。\cite{design48_queue_ref}
\section{PERTの最長経路と総余裕}\label{theory:design:pert}
先行作業が複数あれば全て終了してから後続を始める。結合点の最早時刻は\(T_E(j)=\max_i[T_E(i)+t_{ij}]\)、最遅時刻は終点から\(T_L(i)=\min_j[T_L(j)-t_{ij}]\)で求める。ダミー作業は所要0で先行関係を伝える。総余裕\(TF_{ij}=T_L(j)-T_E(i)-t_{ij}\)が0の作業はクリティカルである。
全工期は最長経路長であり、最も安い作業を無条件に縮めても工期は下がらない。複数の最長経路があれば全てを短くする。平成30年度はa3とb2が同じ22日で二つの最長枝を作るため、その作業和集合が管理対象になる。短縮前後の最長経路を再計算し、別経路への支配交代も確かめる。
\section{線形計画の実行可能性と頂点}\label{theory:design:lp}
生産量を変数にし、資源上限を不等式、利益を目的関数へ置き換える。利益係数が最大の製品だけを増やす前に、全材料と時間の制約を確認する。二変数なら制約境界の交点を求め、実行可能領域の頂点の利益を比較する。台数が整数なら連続最適解をそのまま使わず整数候補を確認する。
令和2年度は\(3x+2y\le24\)、\(x+3y\le15\)の交点\((6,3)\)が当初の最適である。利益係数変更後もこの頂点を最適にするには、隣接頂点\((0,5),(8,0)\)と比較して
\[21+6\Delta c\ge15,\qquad21+6\Delta c\ge16+8\Delta c,\qquad-1\le\Delta c\le5/2.\]
端点では最適解が複数になる。元の点が最適であることと唯一であることを区別する。全組合せ最適化に指数時間下界が証明済みと主張せず、解ける部分クラスと近似品質の条件を別に扱う。
\section{ユニバーサルデザインの利用者と場面}\label{theory:design:universal}
UDは可能な限り多くの人が特別な変更なしに使える製品・環境を目指す。1997年作業班の七原則は、公平な利用、柔軟な利用、単純で直観的な利用、認知できる情報、失敗への寛容、少ない身体的努力、接近・利用に必要な大きさと空間である。原則名を対象例へ対応させ、自動ドアの省力、速度変更の柔軟性、多形式情報、低い位置からの到達を比較する。安全のため特定利用者を一律に排除する機能を、そのままUDの例にはしない。\cite{design48_ud}
\section{分割と機能独立性}\label{theory:design:principles}
多様な要求の調和を図り、複雑系を適切な機能単位へ分割して管理する。複雑であるほど優れた設計という結論は出ない。機能間の結合を減らすことは変更影響を局所化するが、インタフェース・全体相互作用を無視できるわけではない。ばらつきは設計対象であり理想値だけを実機の保証としない。
\section{角法と図面情報}\label{theory:design:drawing}
第三角では見る側と同じ側へ図を置き、平面図は上、右側面図は右。第一角では反対側となり平面図は下、左側面図は右になる。上と下の一語が異なる候補を全体の印象で正にしない。ISOは国際標準化機構で国別のイタリア規格ではない。公開規格には第一角・第三角の節があるが有料節の全図まで読了したとはしない。\cite{design48_iso}
\section{要求性能から構造詳細と照査へ}\label{theory:design:performance}
合理的構造体を実現するため、要求性能を設定し、構造計画、構造詳細の設定、性能照査へ進む。構造詳細は部材寸法・材料等の照査に必要な情報であり、詳細設計という一般工程名と異なる。供用期間の作用と劣化を踏まえて、設定した要求が満たされるかを評価する。使用性能という一つの性能だけに全要求を限定しない。\cite{design48_mlit}

\section{利得表は他者の行動ごとに比較する}
\label{theory:design:game}
各参加者が自分の利得を最大にするゲームでは、相手の行動を固定して自分の候補を比較する。相手のどの行動でも優る戦略が支配戦略である。両者の最適反応が同時に満たされる組合せは、片方だけ変更しても利得を増せないナッシュ均衡となる。行プレイヤーの利得を列ごと、列プレイヤーの利得を行ごとに比較し、表を転置して読む誤りを避ける。両者の合計利得の最大化は協調の評価で、独立の最適反応とは別である。

\section{期待値と決定ノードの後向き評価}
\label{theory:design:decision}
確率が既知の状態\(s\)について期待便益は\(\sum_s p_s B_s\)。一定費用\(C\)を引けば期待純便益になる。決定後に状態を観測して対応を選べる木では、末端から機会ノードを期待値、決定ノードを最大便益又は最小損失で置き換える。選べる二案を半々に平均してはならない。先に選ぶ案と、状態を見てから選べる案では利用できる情報が違う。

\section{限界利益と検査の期待費用}
\label{theory:design:economics}
一人又は一個当たり売上を\(p\)、変動費を\(v\)、期間固定費を\(F\)、目標利益を\(G\)とすると、必要数量は\((F+G)/(p-v)\)。\(p-v>0\)が前提である。最後に営業日数・設備数で割って、求める単位へ変換する。
\label{theory:design:inspection}
検査回数\(x\)で不具合確率が\(p(x)\)、一回費用\(c\)、損害額\(D\)なら期待総費用は\(cx+Dp(x)\)。検査を増やす費用と確率低下の利益を比較する。連続最適点を微分で求めた後、実際に許される整数値と端点を調べる。推定確率のモデルをそのまま現場の保証値とはしない。

\section{凸性は最適解の意味を強める}
\label{theory:design:convex}
線形計画は目的と制約がともに線形である。凸集合上で凸関数を最小化する凸計画では局所最適点が大域最適点となる\autocite{finish56_convex}。目的だけ凸でも実行可能集合が非凸ならその性質は保証されない。整数制約は実行可能集合を離散化し、連続問題の解を丸めるだけでは制約を満たさない場合がある。ヒューリスティックは候補解を得る手段で、最適性証明や誤差保証は別に扱う。

\section{経済的発注量は二つの費用を釣り合わせる}
\label{theory:design:eoq}
一定需要\(d\)、即時一括補充、欠品なしのモデルでは、発注量\(q\)に対し発注回数\(d/q\)、平均在庫\(q/2\)。一回費用\(k\)、年間単位在庫費\(h\)なら
\begin{equation}C(q)=\frac{kd}{q}+\frac{hq}{2},\qquad q_*=\sqrt{\frac{2kd}{h}}.\end{equation}
\(q>0\)で\(C''=2kd/q^3>0\)だから停留点は最小である\autocite{finish56_eoq}。上限、ロット単位、需要変動、納期、数量割引があれば再計算する。全需要に比例して一定の購入費は、このモデルでは最適発注量に影響しない。

\section{三つの解析公式の成立条件}
\label{theory:design:mathformulas}
Taylor展開の中心が0ならMaclaurin展開で\(f(x)=\sum_{k\ge0}f^{(k)}(0)x^k/k!\)。無限回微分可能だけでは級数と元の関数の一致は保証されず、剰余の収束を含む条件が必要である。複素指数の\(e^{i\theta}=\cos\theta+i\sin\theta\)から\(\theta=\pi\)で\(e^{i\pi}+1=0\)となる。
ロピタルの定理は\(0/0\)又は適切な無限大比等の不定形で、微分可能性、分母微分非零、微分比の極限等を確認して使う。任意の比の分子・分母を微分する規則ではない。

\section{不等式の向きと等号条件}
\label{theory:design:inequality}
非負の数の幾何平均は算術平均以下で、全数が同じなら等号を取る。\([0,\pi/2]\)上の\(\sin\theta\)の凹性から弦との比較で\(\sin\theta\ge2\theta/\pi\)。\(0<\theta\le\pi/2\)なら\(\sin\theta/\theta\ge2/\pi\)、右端で等号である。
一方、区間上で\(f''>0\)なら厳密凸で、異なる二点に対し\(f((x_1+x_2)/2)<(f(x_1)+f(x_2))/2\)。端点を含む不等式と厳密凸の厳密不等号を読み分ける。

\section{歯車の速度比とトルク}
\label{theory:design:gear}
回転動力は\(P=T\omega\)。速度比\(r=\omega_{\mathrm{in}}/\omega_{\mathrm{out}}\)、効率\(\eta\)なら\(T_{\mathrm{out}}=\eta rT_{\mathrm{in}}\)。減速で角速度が下がると、同じ動力に対応するトルクは増す。比の定義を入力/出力かその逆かで確認し、回転方向と大きさの比較を分ける。

\section{平面応力と平面ひずみは異なる近似}
\label{theory:design:plane}
薄板の自由表面に対する平面応力では\(\sigma_z,\tau_{xz},\tau_{yz}\simeq0\)で、面外ひずみは一般に0でない。軸方向変形を拘束した一様な長い構造の内部に対する平面ひずみでは\(\varepsilon_z,\gamma_{xz},\gamma_{yz}\simeq0\)で、軸応力は一般に残る。等方線形弾性なら後者の\(\sigma_z=\nu(\sigma_x+\sigma_y)\)。端部、軸方向の自由伸び、異方性を含む場合は境界条件からモデルを選び直す。

\subsection{平成23年度 問1-1}
\label{q:h23-reg-basic-1-1}

\noindent\textbf{原典：} sources/H23/basic.pdf、PDF物理1頁。

\textbf{公式正答：} 2（single）。

\par\textit{原典本文の校正と独自解説は未作成。分類のみの登録項目。}

\subsection{平成23年度 問1-2}
\label{q:h23-reg-basic-1-2}

\noindent\textbf{原典：} sources/H23/basic.pdf、PDF物理1頁。

\textbf{公式正答：} 4（single）。

\par\textit{原典本文の校正と独自解説は未作成。分類のみの登録項目。}

\subsection{平成23年度 問1-3}
\label{q:h23-reg-basic-1-3}

\noindent\textbf{原典：} sources/H23/basic.pdf、PDF物理2頁。

\textbf{公式正答：} 3（single）。

\par\textit{原典本文の校正と独自解説は未作成。分類のみの登録項目。}

\subsection{平成23年度 問1-4}
\label{q:h23-reg-basic-1-4}

\noindent\textbf{原典：} sources/H23/basic.pdf、PDF物理3頁。

\textbf{公式正答：} 1（single）。

\par\textit{原典本文の校正と独自解説は未作成。分類のみの登録項目。}

\subsection{平成23年度 問1-5}
\label{q:h23-reg-basic-1-5}

\noindent\textbf{原典：} sources/H23/basic.pdf、PDF物理4頁。

\textbf{公式正答：} 4（single）。

\par\textit{原典本文の校正と独自解説は未作成。分類のみの登録項目。}

\subsection{平成24年度 問1-1}
\label{q:h24-reg-basic-1-1}

\noindent\textbf{原典：} sources/H24/basic.pdf、PDF物理1頁。

\textbf{公式正答：} 4（single）。

\par\textit{原典本文の校正と独自解説は未作成。分類のみの登録項目。}

\subsection{平成24年度 問1-2}
\label{q:h24-reg-basic-1-2}

\noindent\textbf{原典：} sources/H24/basic.pdf、PDF物理2頁。

\textbf{公式正答：} 5（single）。

\par\textit{原典本文の校正と独自解説は未作成。分類のみの登録項目。}

\subsection{平成24年度 問1-3}
\label{q:h24-reg-basic-1-3}

\noindent\textbf{原典：} sources/H24/basic.pdf、PDF物理3頁。

\textbf{公式正答：} 3（single）。

\par\textit{原典本文の校正と独自解説は未作成。分類のみの登録項目。}

\subsection{平成24年度 問1-4}
\label{q:h24-reg-basic-1-4}

\noindent\textbf{原典：} sources/H24/basic.pdf、PDF物理4頁。

\textbf{公式正答：} 5（single）。

\par\textit{原典本文の校正と独自解説は未作成。分類のみの登録項目。}

\subsection{平成24年度 問1-5}
\label{q:h24-reg-basic-1-5}

\noindent\textbf{原典：} sources/H24/basic.pdf、PDF物理4頁。

\textbf{公式正答：} 5（single）。

\par\textit{原典本文の校正と独自解説は未作成。分類のみの登録項目。}

\subsection{平成25年度 問1-1}
\label{q:h25-reg-basic-1-1}

\noindent\textbf{原典：} sources/H25/basic.pdf、PDF物理1頁。

\textbf{公式正答：} 3（single）。

\par\textit{原典本文の校正と独自解説は未作成。分類のみの登録項目。}

\subsection{平成25年度 問1-2}
\label{q:h25-reg-basic-1-2}

\noindent\textbf{原典：} sources/H25/basic.pdf、PDF物理2頁。

\textbf{公式正答：} 5（single）。

\par\textit{原典本文の校正と独自解説は未作成。分類のみの登録項目。}

\subsection{平成25年度 問1-3}
\label{q:h25-reg-basic-1-3}

\noindent\textbf{原典：} sources/H25/basic.pdf、PDF物理3頁。

\textbf{公式正答：} 2（single）。

\par\textit{原典本文の校正と独自解説は未作成。分類のみの登録項目。}

\subsection{平成25年度 問1-4}
\label{q:h25-reg-basic-1-4}

\noindent\textbf{原典：} sources/H25/basic.pdf、PDF物理4頁。

\textbf{公式正答：} 4（single）。

\par\textit{原典本文の校正と独自解説は未作成。分類のみの登録項目。}

\subsection{平成25年度 問1-5}
\label{q:h25-reg-basic-1-5}

\noindent\textbf{原典：} sources/H25/basic.pdf、PDF物理5頁。

\textbf{公式正答：} 4（single）。

\par\textit{原典本文の校正と独自解説は未作成。分類のみの登録項目。}

\subsection{平成25年度 問1-6}
\label{q:h25-reg-basic-1-6}

\noindent\textbf{原典：} sources/H25/basic.pdf、PDF物理5頁。

\textbf{公式正答：} 3（single）。

\par\textit{原典本文の校正と独自解説は未作成。分類のみの登録項目。}

\subsection{平成26年度 問1-1}
\label{q:h26-reg-basic-1-1}

\noindent\textbf{原典：} sources/H26/basic.pdf、PDF物理1頁。

\textbf{公式正答：} 5（single）。

\par\textit{原典本文の校正と独自解説は未作成。分類のみの登録項目。}

\subsection{平成26年度 問1-2}
\label{q:h26-reg-basic-1-2}

\noindent\textbf{原典：} sources/H26/basic.pdf、PDF物理1頁。

\textbf{公式正答：} 4（single）。

\par\textit{原典本文の校正と独自解説は未作成。分類のみの登録項目。}

\subsection{平成26年度 問1-3}
\label{q:h26-reg-basic-1-3}

\noindent\textbf{原典：} sources/H26/basic.pdf、PDF物理2頁。

\textbf{公式正答：} 1（single）。

\par\textit{原典本文の校正と独自解説は未作成。分類のみの登録項目。}

\subsection{平成26年度 問1-4}
\label{q:h26-reg-basic-1-4}

\noindent\textbf{原典：} sources/H26/basic.pdf、PDF物理3頁。

\textbf{公式正答：} 3（single）。

\par\textit{原典本文の校正と独自解説は未作成。分類のみの登録項目。}

\subsection{平成26年度 問1-5}
\label{q:h26-reg-basic-1-5}

\noindent\textbf{原典：} sources/H26/basic.pdf、PDF物理3頁。

\textbf{公式正答：} 1（single）。

\par\textit{原典本文の校正と独自解説は未作成。分類のみの登録項目。}

\subsection{平成26年度 問1-6}
\label{q:h26-reg-basic-1-6}

\noindent\textbf{原典：} sources/H26/basic.pdf、PDF物理4頁。

\textbf{公式正答：} 3（single）。

\par\textit{原典本文の校正と独自解説は未作成。分類のみの登録項目。}

\subsection{平成27年度 問1-1}
\label{q:h27-reg-basic-1-1}

\noindent\textbf{原典：} sources/H27/basic.pdf、PDF物理1頁。

\textbf{公式正答：} 5（single）。

\par\textit{原典本文の校正と独自解説は未作成。分類のみの登録項目。}

\subsection{平成27年度 問1-2}
\label{q:h27-reg-basic-1-2}

\noindent\textbf{原典：} sources/H27/basic.pdf、PDF物理2頁。

\textbf{公式正答：} 4（single）。

\par\textit{原典本文の校正と独自解説は未作成。分類のみの登録項目。}

\subsection{平成27年度 問1-3}
\label{q:h27-reg-basic-1-3}

\noindent\textbf{原典：} sources/H27/basic.pdf、PDF物理2頁。

\textbf{公式正答：} 5（single）。

\par\textit{原典本文の校正と独自解説は未作成。分類のみの登録項目。}

\subsection{平成27年度 問1-4}
\label{q:h27-reg-basic-1-4}

\noindent\textbf{原典：} sources/H27/basic.pdf、PDF物理3頁。

\textbf{公式正答：} 3（single）。

\par\textit{原典本文の校正と独自解説は未作成。分類のみの登録項目。}

\subsection{平成27年度 問1-5}
\label{q:h27-reg-basic-1-5}

\noindent\textbf{原典：} sources/H27/basic.pdf、PDF物理4頁。

\textbf{公式正答：} 1（single）。

\par\textit{原典本文の校正と独自解説は未作成。分類のみの登録項目。}

\subsection{平成27年度 問1-6}
\label{q:h27-reg-basic-1-6}

\noindent\textbf{原典：} sources/H27/basic.pdf、PDF物理4頁。

\textbf{公式正答：} 5（single）。

\par\textit{原典本文の校正と独自解説は未作成。分類のみの登録項目。}

\subsection{平成28年度 問1-1}
\label{q:h28-reg-basic-1-1}

\noindent\textbf{原典：} sources/H28/basic.pdf、PDF物理1頁。

\textbf{公式正答：} 4（single）。

\par\textit{原典本文の校正と独自解説は未作成。分類のみの登録項目。}

\subsection{平成28年度 問1-2}
\label{q:h28-reg-basic-1-2}

\noindent\textbf{原典：} sources/H28/basic.pdf、PDF物理2頁。

\textbf{公式正答：} 4（single）。

\par\textit{原典本文の校正と独自解説は未作成。分類のみの登録項目。}

\subsection{平成28年度 問1-3}
\label{q:h28-reg-basic-1-3}

\noindent\textbf{原典：} sources/H28/basic.pdf、PDF物理3頁。

\textbf{公式正答：} 4（single）。

\par\textit{原典本文の校正と独自解説は未作成。分類のみの登録項目。}

\subsection{平成28年度 問1-4}
\label{q:h28-reg-basic-1-4}

\noindent\textbf{原典：} sources/H28/basic.pdf、PDF物理4頁。

\textbf{公式正答：} 4（single）。

\par\textit{原典本文の校正と独自解説は未作成。分類のみの登録項目。}

\subsection{平成28年度 問1-5}
\label{q:h28-reg-basic-1-5}

\noindent\textbf{原典：} sources/H28/basic.pdf、PDF物理5頁。

\textbf{公式正答：} 5（single）。

\par\textit{原典本文の校正と独自解説は未作成。分類のみの登録項目。}

\subsection{平成28年度 問1-6}
\label{q:h28-reg-basic-1-6}

\noindent\textbf{原典：} sources/H28/basic.pdf、PDF物理6頁。

\textbf{公式正答：} 1（single）。

\par\textit{原典本文の校正と独自解説は未作成。分類のみの登録項目。}

\subsection{平成29年度 問1-1}
\label{q:h29-reg-basic-1-1}

\noindent\textbf{原典：} sources/H29/basic.pdf、PDF物理1頁。

\textbf{公式正答：} 5（single）。

\par\textit{原典本文の校正と独自解説は未作成。分類のみの登録項目。}

\subsection{平成29年度 問1-2}
\label{q:h29-reg-basic-1-2}

\noindent\textbf{原典：} sources/H29/basic.pdf、PDF物理1頁。

\textbf{公式正答：} 5（single）。

\par\textit{原典本文の校正と独自解説は未作成。分類のみの登録項目。}

\subsection{平成29年度 問1-3}
\label{q:h29-reg-basic-1-3}

\noindent\textbf{原典：} sources/H29/basic.pdf、PDF物理2頁。

\textbf{公式正答：} 2（single）。

\par\textit{原典本文の校正と独自解説は未作成。分類のみの登録項目。}

\subsection{平成29年度 問1-4}
\label{q:h29-reg-basic-1-4}

\noindent\textbf{原典：} sources/H29/basic.pdf、PDF物理3頁。

\textbf{公式正答：} 5（single）。

\par\textit{原典本文の校正と独自解説は未作成。分類のみの登録項目。}

\subsection{平成29年度 問1-5}
\label{q:h29-reg-basic-1-5}

\noindent\textbf{原典：} sources/H29/basic.pdf、PDF物理4頁。

\textbf{公式正答：} 1（single）。

\par\textit{原典本文の校正と独自解説は未作成。分類のみの登録項目。}

\subsection{平成29年度 問1-6}
\label{q:h29-reg-basic-1-6}

\noindent\textbf{原典：} sources/H29/basic.pdf、PDF物理5頁。

\textbf{公式正答：} 2（single）。

\par\textit{原典本文の校正と独自解説は未作成。分類のみの登録項目。}

\subsection{平成30年度 問1-1}
\label{q:h30-reg-basic-1-1}

\noindent\textbf{原典：} sources/H30/basic.pdf、PDF物理1頁。

\textbf{公式正答：} 3（single）。

\par\textit{原典本文の校正と独自解説は未作成。分類のみの登録項目。}

\subsection{平成30年度 問1-2}
\label{q:h30-reg-basic-1-2}

\noindent\textbf{原典：} sources/H30/basic.pdf、PDF物理2頁。

\textbf{公式正答：} 1（single）。

\par\textit{原典本文の校正と独自解説は未作成。分類のみの登録項目。}

\subsection{平成30年度 問1-3}
\label{q:h30-reg-basic-1-3}

\noindent\textbf{原典：} sources/H30/basic.pdf、PDF物理3頁。

\textbf{公式正答：} 2（single）。

\par\textit{原典本文の校正と独自解説は未作成。分類のみの登録項目。}

\subsection{平成30年度 問1-4}
\label{q:h30-reg-basic-1-4}

\noindent\textbf{原典：} sources/H30/basic.pdf、PDF物理4頁。

\textbf{公式正答：} 2（single）。

\par\textit{原典本文の校正と独自解説は未作成。分類のみの登録項目。}

\subsection{平成30年度 問1-5}
\label{q:h30-reg-basic-1-5}

\noindent\textbf{原典：} sources/H30/basic.pdf、PDF物理5頁。

\textbf{公式正答：} 3（single）。

\par\textit{原典本文の校正と独自解説は未作成。分類のみの登録項目。}

\subsection{平成30年度 問1-6}
\label{q:h30-reg-basic-1-6}

\noindent\textbf{原典：} sources/H30/basic.pdf、PDF物理5頁。

\textbf{公式正答：} 4（single）。

\par\textit{原典本文の校正と独自解説は未作成。分類のみの登録項目。}

\subsection{令和元年度 問1-1}
\label{q:r01-reg-basic-1-1}

\noindent\textbf{原典：} sources/R01/basic.pdf、PDF物理1頁。

\textbf{公式正答：} 5（single）。

\par\textit{原典本文の校正と独自解説は未作成。分類のみの登録項目。}

\subsection{令和元年度 問1-2}
\label{q:r01-reg-basic-1-2}

\noindent\textbf{原典：} sources/R01/basic.pdf、PDF物理2頁。

\textbf{公式正答：} 2（single）。

\par\textit{原典本文の校正と独自解説は未作成。分類のみの登録項目。}

\subsection{令和元年度 問1-3}
\label{q:r01-reg-basic-1-3}

\noindent\textbf{原典：} sources/R01/basic.pdf、PDF物理3頁。

\textbf{公式正答：} 5（single）。

\par\textit{原典本文の校正と独自解説は未作成。分類のみの登録項目。}

\subsection{令和元年度 問1-4}
\label{q:r01-reg-basic-1-4}

\noindent\textbf{原典：} sources/R01/basic.pdf、PDF物理4頁。

\textbf{公式正答：} 5（single）。

\par\textit{原典本文の校正と独自解説は未作成。分類のみの登録項目。}

\subsection{令和元年度 問1-5}
\label{q:r01-reg-basic-1-5}

\noindent\textbf{原典：} sources/R01/basic.pdf、PDF物理5頁。

\textbf{公式正答：} 2（single）。

\par\textit{原典本文の校正と独自解説は未作成。分類のみの登録項目。}

\subsection{令和元年度 問1-6}
\label{q:r01-reg-basic-1-6}

\noindent\textbf{原典：} sources/R01/basic.pdf、PDF物理6頁。

\textbf{公式正答：} 3（single）。

\par\textit{原典本文の校正と独自解説は未作成。分類のみの登録項目。}

\subsection{令和元年度再試験 問1-1}
\label{q:r01-ret-basic-1-1}

\noindent\textbf{原典：} sources/R01\_retest/basic.pdf、PDF物理1頁。

\textbf{公式正答：} 5（single）。

\par\textit{原典本文の校正と独自解説は未作成。分類のみの登録項目。}

\subsection{令和元年度再試験 問1-2}
\label{q:r01-ret-basic-1-2}

\noindent\textbf{原典：} sources/R01\_retest/basic.pdf、PDF物理2頁。

\textbf{公式正答：} 3（single）。

\par\textit{原典本文の校正と独自解説は未作成。分類のみの登録項目。}

\subsection{令和元年度再試験 問1-3}
\label{q:r01-ret-basic-1-3}

\noindent\textbf{原典：} sources/R01\_retest/basic.pdf、PDF物理3頁。

\textbf{公式正答：} 1（single）。

\par\textit{原典本文の校正と独自解説は未作成。分類のみの登録項目。}

\subsection{令和元年度再試験 問1-4}
\label{q:r01-ret-basic-1-4}

\noindent\textbf{原典：} sources/R01\_retest/basic.pdf、PDF物理4頁。

\textbf{公式正答：} 2（single）。

\par\textit{原典本文の校正と独自解説は未作成。分類のみの登録項目。}

\subsection{令和元年度再試験 問1-5}
\label{q:r01-ret-basic-1-5}

\noindent\textbf{原典：} sources/R01\_retest/basic.pdf、PDF物理5頁。

\textbf{公式正答：} 2（single）。

\par\textit{原典本文の校正と独自解説は未作成。分類のみの登録項目。}

\subsection{令和元年度再試験 問1-6}
\label{q:r01-ret-basic-1-6}

\noindent\textbf{原典：} sources/R01\_retest/basic.pdf、PDF物理5頁。

\textbf{公式正答：} 5（single）。

\par\textit{原典本文の校正と独自解説は未作成。分類のみの登録項目。}

\subsection{令和2年度 問1-1}
\label{q:r02-reg-basic-1-1}

\noindent\textbf{原典：} sources/R02/basic.pdf、PDF物理1頁。

\textbf{公式正答：} 3（single）。

\par\textit{原典本文の校正と独自解説は未作成。分類のみの登録項目。}

\subsection{令和2年度 問1-2}
\label{q:r02-reg-basic-1-2}

\noindent\textbf{原典：} sources/R02/basic.pdf、PDF物理2頁。

\textbf{公式正答：} 4（single）。

\par\textit{原典本文の校正と独自解説は未作成。分類のみの登録項目。}

\subsection{令和2年度 問1-3}
\label{q:r02-reg-basic-1-3}

\noindent\textbf{原典：} sources/R02/basic.pdf、PDF物理3頁。

\textbf{公式正答：} 3（single）。

\par\textit{原典本文の校正と独自解説は未作成。分類のみの登録項目。}

\subsection{令和2年度 問1-4}
\label{q:r02-reg-basic-1-4}

\noindent\textbf{原典：} sources/R02/basic.pdf、PDF物理4頁。

\textbf{公式正答：} 5（single）。

\par\textit{原典本文の校正と独自解説は未作成。分類のみの登録項目。}

\subsection{令和2年度 問1-5}
\label{q:r02-reg-basic-1-5}

\noindent\textbf{原典：} sources/R02/basic.pdf、PDF物理5頁。

\textbf{公式正答：} 5（single）。

\par\textit{原典本文の校正と独自解説は未作成。分類のみの登録項目。}

\subsection{令和2年度 問1-6}
\label{q:r02-reg-basic-1-6}

\noindent\textbf{原典：} sources/R02/basic.pdf、PDF物理6頁。

\textbf{公式正答：} 3（single）。

\par\textit{原典本文の校正と独自解説は未作成。分類のみの登録項目。}

\subsection{令和3年度 問1-1}
\label{q:r03-reg-basic-1-1}

\noindent\textbf{原典：} sources/R03/basic.pdf、PDF物理1頁。

\textbf{公式正答：} 4（single）。

\par\textit{原典本文の校正と独自解説は未作成。分類のみの登録項目。}

\subsection{令和3年度 問1-2}
\label{q:r03-reg-basic-1-2}

\noindent\textbf{原典：} sources/R03/basic.pdf、PDF物理2頁。

\textbf{公式正答：} 2（single）。

\par\textit{原典本文の校正と独自解説は未作成。分類のみの登録項目。}

\subsection{令和3年度 問1-3}
\label{q:r03-reg-basic-1-3}

\noindent\textbf{原典：} sources/R03/basic.pdf、PDF物理3頁。

\textbf{公式正答：} 2（single）。

\par\textit{原典本文の校正と独自解説は未作成。分類のみの登録項目。}

\subsection{令和3年度 問1-4}
\label{q:r03-reg-basic-1-4}

\noindent\textbf{原典：} sources/R03/basic.pdf、PDF物理3頁。

\textbf{公式正答：} 3（single）。

\par\textit{原典本文の校正と独自解説は未作成。分類のみの登録項目。}

\subsection{令和3年度 問1-5}
\label{q:r03-reg-basic-1-5}

\noindent\textbf{原典：} sources/R03/basic.pdf、PDF物理4頁。

\textbf{公式正答：} 5（single）。

\par\textit{原典本文の校正と独自解説は未作成。分類のみの登録項目。}

\subsection{令和3年度 問1-6}
\label{q:r03-reg-basic-1-6}

\noindent\textbf{原典：} sources/R03/basic.pdf、PDF物理5頁。

\textbf{公式正答：} 3（single）。

\par\textit{原典本文の校正と独自解説は未作成。分類のみの登録項目。}

\subsection{令和4年度 問1-1}
\label{q:r04-reg-basic-1-1}

\noindent\textbf{原典：} sources/R04/basic.pdf、PDF物理1頁。

\textbf{公式正答：} 1（single）。

\par\textit{原典本文の校正と独自解説は未作成。分類のみの登録項目。}

\subsection{令和4年度 問1-2}
\label{q:r04-reg-basic-1-2}

\noindent\textbf{原典：} sources/R04/basic.pdf、PDF物理2頁。

\textbf{公式正答：} 4（single）。

\par\textit{原典本文の校正と独自解説は未作成。分類のみの登録項目。}

\subsection{令和4年度 問1-3}
\label{q:r04-reg-basic-1-3}

\noindent\textbf{原典：} sources/R04/basic.pdf、PDF物理2頁。

\textbf{公式正答：} 3（single）。

\par\textit{原典本文の校正と独自解説は未作成。分類のみの登録項目。}

\subsection{令和4年度 問1-4}
\label{q:r04-reg-basic-1-4}

\noindent\textbf{原典：} sources/R04/basic.pdf、PDF物理3頁。

\textbf{公式正答：} 1（single）。

\par\textit{原典本文の校正と独自解説は未作成。分類のみの登録項目。}

\subsection{令和4年度 問1-5}
\label{q:r04-reg-basic-1-5}

\noindent\textbf{原典：} sources/R04/basic.pdf、PDF物理3頁。

\textbf{公式正答：} 3（single）。

\par\textit{原典本文の校正と独自解説は未作成。分類のみの登録項目。}

\subsection{令和4年度 問1-6}
\label{q:r04-reg-basic-1-6}

\noindent\textbf{原典：} sources/R04/basic.pdf、PDF物理4頁。

\textbf{公式正答：} 2（single）。

\par\textit{原典本文の校正と独自解説は未作成。分類のみの登録項目。}

\subsection{令和5年度 問1-1}
\label{q:r05-reg-basic-1-1}

\noindent\textbf{原典：} sources/R05/basic.pdf、PDF物理1頁。

\textbf{公式正答：} 5（single）。

\par\textit{原典本文の校正と独自解説は未作成。分類のみの登録項目。}

\subsection{令和5年度 問1-2}
\label{q:r05-reg-basic-1-2}

\noindent\textbf{原典：} sources/R05/basic.pdf、PDF物理2頁。

\textbf{公式正答：} 4（single）。

\par\textit{原典本文の校正と独自解説は未作成。分類のみの登録項目。}

\subsection{令和5年度 問1-3}
\label{q:r05-reg-basic-1-3}

\noindent\textbf{原典：} sources/R05/basic.pdf、PDF物理3頁。

\textbf{公式正答：} 3（single）。

\par\textit{原典本文の校正と独自解説は未作成。分類のみの登録項目。}

\subsection{令和5年度 問1-4}
\label{q:r05-reg-basic-1-4}

\noindent\textbf{原典：} sources/R05/basic.pdf、PDF物理3頁。

\textbf{公式正答：} 4（single）。

\par\textit{原典本文の校正と独自解説は未作成。分類のみの登録項目。}

\subsection{令和5年度 問1-5}
\label{q:r05-reg-basic-1-5}

\noindent\textbf{原典：} sources/R05/basic.pdf、PDF物理4頁。

\textbf{公式正答：} 1（single）。

\par\textit{原典本文の校正と独自解説は未作成。分類のみの登録項目。}

\subsection{令和5年度 問1-6}
\label{q:r05-reg-basic-1-6}

\noindent\textbf{原典：} sources/R05/basic.pdf、PDF物理4頁。

\textbf{公式正答：} 2（single）。

\par\textit{原典本文の校正と独自解説は未作成。分類のみの登録項目。}

\subsection{令和6年度 問1-1}
\label{q:r06-reg-basic-1-1}

\noindent\textbf{原典：} sources/R06/basic.pdf、PDF物理1頁。

\textbf{公式正答：} 1（single）。

\par\textit{原典本文の校正と独自解説は未作成。分類のみの登録項目。}

\subsection{令和6年度 問1-2}
\label{q:r06-reg-basic-1-2}

\noindent\textbf{原典：} sources/R06/basic.pdf、PDF物理2頁。

\textbf{公式正答：} 2（single）。

\par\textit{原典本文の校正と独自解説は未作成。分類のみの登録項目。}

\subsection{令和6年度 問1-3}
\label{q:r06-reg-basic-1-3}

\noindent\textbf{原典：} sources/R06/basic.pdf、PDF物理2頁。

\textbf{公式正答：} 1（single）。

\par\textit{原典本文の校正と独自解説は未作成。分類のみの登録項目。}

\subsection{令和6年度 問1-4}
\label{q:r06-reg-basic-1-4}

\noindent\textbf{原典：} sources/R06/basic.pdf、PDF物理3頁。

\textbf{公式正答：} 3（single）。

\par\textit{原典本文の校正と独自解説は未作成。分類のみの登録項目。}

\subsection{令和6年度 問1-5}
\label{q:r06-reg-basic-1-5}

\noindent\textbf{原典：} sources/R06/basic.pdf、PDF物理4頁。

\textbf{公式正答：} 4（single）。

\par\textit{原典本文の校正と独自解説は未作成。分類のみの登録項目。}

\subsection{令和6年度 問1-6}
\label{q:r06-reg-basic-1-6}

\noindent\textbf{原典：} sources/R06/basic.pdf、PDF物理5頁。

\textbf{公式正答：} 3（single）。

\par\textit{原典本文の校正と独自解説は未作成。分類のみの登録項目。}

\subsection{令和7年度 問1-1}
\label{q:r07-reg-basic-1-1}

\noindent\textbf{原典：} sources/R07/basic.pdf、PDF物理1頁。

\textbf{公式正答：} 3（single）。

\par\textit{原典本文の校正と独自解説は未作成。分類のみの登録項目。}

\subsection{令和7年度 問1-2}
\label{q:r07-reg-basic-1-2}

\noindent\textbf{原典：} sources/R07/basic.pdf、PDF物理2頁。

\textbf{公式正答：} 4（single）。

\par\textit{原典本文の校正と独自解説は未作成。分類のみの登録項目。}

\subsection{令和7年度 問1-3}
\label{q:r07-reg-basic-1-3}

\noindent\textbf{原典：} sources/R07/basic.pdf、PDF物理2頁。

\textbf{公式正答：} 2（single）。

\par\textit{原典本文の校正と独自解説は未作成。分類のみの登録項目。}

\subsection{令和7年度 問1-4}
\label{q:r07-reg-basic-1-4}

\noindent\textbf{原典：} sources/R07/basic.pdf、PDF物理3頁。

\textbf{公式正答：} 1（single）。

\par\textit{原典本文の校正と独自解説は未作成。分類のみの登録項目。}

\subsection{令和7年度 問1-5}
\label{q:r07-reg-basic-1-5}

\noindent\textbf{原典：} sources/R07/basic.pdf、PDF物理3頁。

\textbf{公式正答：} 5（single）。

\par\textit{原典本文の校正と独自解説は未作成。分類のみの登録項目。}

\subsection{令和7年度 問1-6}
\label{q:r07-reg-basic-1-6}

\noindent\textbf{原典：} sources/R07/basic.pdf、PDF物理4頁。

\textbf{公式正答：} 3（single）。

\par\textit{原典本文の校正と独自解説は未作成。分類のみの登録項目。}


\chapter{情報・論理}
\noindent 分野解説は執筆中。以下は登録問題と校正済み例の作業用一覧。
\section{スタック：状態と操作を分ける}
\subsection{背景と基本概念}
データ構造は値の集合だけでなく、許された操作と順序を定める。スタックでは最後に入れた要素を先に取り出す。
底から上へ$S=(s_1,\ldots,s_n)$と書けば、PUSH $x$は$S'=(s_1,\ldots,s_n,x)$、POPは$s_n$を返し$S'=(s_1,\ldots,s_{n-1})$を残す。
空のスタックへのPOPは別途の扱いを要するが、ここで校正した例の操作列では発生しない。
\subsection{過去問の追跡}
取り出された値と、操作後に残った状態を別に記録する。複数回のPOPがあるときは、最後の戻り値を問われているかを確認する。
\begin{tcolorbox}[title=情報工学への接続]
関数呼出しの管理、深さ優先探索、式の解析などへ接続できる。これは学習上の補足であり、各用途の出題頻度をこの例から推定しない。
\end{tcolorbox}

\section{アッカーマン関数と入れ子再帰}

\subsection{令和7年度 問2-6}
\label{q:r07-reg-basic-2-6}

\noindent\textbf{原典：} sources/R07/basic.pdf、PDF物理7頁。

\textbf{公式正答：} 1（single）。

\paragraph{問題の要旨} アッカーマン関数の二段の再帰を展開してA(2,2)を求める。

\Needspace{4\baselineskip}\paragraph{確認した条件・記述（独自要約）}\mbox{}\par

\noindent\textbf{基底：} A(0,n)=\allowbreak{}n+1\par

\noindent\textbf{零引数：} A(m+1,0)=\allowbreak{}A(m,1)\par

\noindent\textbf{再帰：} A(m+1,n+1)=\allowbreak{}A(m,A(m+1,n))\par

\paragraph{問われる核心} アッカーマン関数の二段の再帰を展開してA(2,2)を求める。

\paragraph{必要知識} 非負整数の再帰基底と入れ子の引数評価。

\paragraph{考え方・正答を選ぶ理由（査読待ち初稿）} A(1,0)=\allowbreak{}A(0,1)=\allowbreak{}2、A(1,n+1)=\allowbreak{}A(0,A(1,n))=\allowbreak{}A(1,n)+1なのでA(1,n)=\allowbreak{}n+2。A(2,0)=\allowbreak{}A(1,1)=\allowbreak{}3、A(2,n+1)=\allowbreak{}A(1,A(2,n))=\allowbreak{}A(2,n)+2。したがってA(2,1)=\allowbreak{}5、A(2,2)=\allowbreak{}7、1である。内側再帰の結果を外側の引数へ渡す。

\Needspace{4\baselineskip}\paragraph{選択肢の検討（査読待ち初稿）}\mbox{}\par

\noindent\textbf{1：適切} 7。A(1,n)=\allowbreak{}n+2、A(2,n)=\allowbreak{}2n+3より7である。\par

\noindent\textbf{2：不適切} 8。各再帰段の増分2と初期3を適用した7ではない。\par

\noindent\textbf{3：不適切} 9。n=\allowbreak{}3のA(2,3)=\allowbreak{}9でありnを一つ多くしている。\par

\noindent\textbf{4：不適切} 10。再帰の初期値又は増分の誤りである。\par

\noindent\textbf{5：不適切} 11。A(2,4)=\allowbreak{}11であり二段多く展開している。\par

\paragraph{再利用する考え方} 定義した対象量、演算順序、保存条件を固定し、途中状態と候補の違いを比較する。

\paragraph{共通理論} \hyperref[theory:discrete:recursion]{第\ref*{theory:discrete:recursion}節}

\par 難易度 3（校準初稿）、優先度未評価。

\section{走査の不変条件・分岐と素数判定}

\subsection{平成24年度 問2-5}
\label{q:h24-reg-basic-2-5}

\noindent\textbf{原典：} sources/H24/basic.pdf、PDF物理7頁。

\textbf{公式正答：} 2（single）。

\paragraph{問題の要旨} 最小値と最大値の更新をelseで結んだ走査の不変条件と失敗例を判定する。

\Needspace{4\baselineskip}\paragraph{確認した条件・記述（独自要約）}\mbox{}\par

\noindent\textbf{判断対象：} 0<a\_i<1000、min初期1000、max初期0、elseによる排他的更新。\par

\noindent\textbf{原典確認：} 図・式・全5候補を29保存画像で確認。rawの数式欠落と演算子誤認は別校正JSONで扱う。\par

\paragraph{問われる核心} 最小値と最大値の更新をelseで結んだ走査の不変条件と失敗例を判定する。

\paragraph{必要知識} 0<a\_i<1000、min初期1000、max初期0、elseによる排他的更新。

\paragraph{考え方・正答を選ぶ理由（査読待ち初稿）} min=\allowbreak{}1000から、より小さい値を毎回取り込むため、minは走査済みの最小値である。max更新は最小更新をしなかったときだけ行う。例えば入力2,1では最小更新を続けてmax=\allowbreak{}0となるが、1,2ならmax=\allowbreak{}2となる。minは常に最小、maxは失敗する場合があり2である。

\Needspace{4\baselineskip}\paragraph{選択肢の検討（査読待ち初稿）}\mbox{}\par

\noindent\textbf{1：不適切} min常に最小・max常に最大でない。1,2の順ならmax=\allowbreak{}2となるので常に失敗ではない。\par

\noindent\textbf{2：適切} min常に最小・max失敗する場合あり。minの更新は常に行われ、maxはelseで初回等を飛ばす場合がある。\par

\noindent\textbf{3：不適切} min失敗の場合あり・max常に最大でない。minは各接頭列の最小を保持する。maxも成功例がある。\par

\noindent\textbf{4：不適切} minもmaxも失敗の場合あり。minは常に正しいので前半が違う。\par

\noindent\textbf{5：不適切} min失敗の場合あり・max常に最大。minは正しくmaxは失敗例があるので両方違う。\par

\paragraph{再利用する考え方} 定義した対象量、演算順序、保存条件を固定し、途中状態と候補の違いを比較する。

\paragraph{共通理論} \hyperref[theory:discrete:invariants]{第\ref*{theory:discrete:invariants}節}

\paragraph{関連問題} \hyperref[q:h29-reg-basic-2-3]{H29-REG-BASIC-2-3}

\par 難易度 2（校準初稿）、優先度未評価。

\subsection{平成29年度 問2-3}
\label{q:h29-reg-basic-2-3}

\noindent\textbf{原典：} sources/H29/basic.pdf、PDF物理8頁。

\textbf{公式正答：} 5（single）。

\paragraph{問題の要旨} 試し割りで素数を判定する流れ図の継続条件と割り切れ判定を決める。

\Needspace{4\baselineskip}\paragraph{確認した条件・記述（独自要約）}\mbox{}\par

\noindent\textbf{判断対象：} 因子対の平方根境界、継続判定、割り切れの向き。\par

\noindent\textbf{原典確認：} 図・式・全5候補を29保存画像で確認。rawの数式欠落と演算子誤認は別校正JSONで扱う。\par

\paragraph{問われる核心} 試し割りで素数を判定する流れ図の継続条件と割り切れ判定を決める。

\paragraph{必要知識} 因子対の平方根境界、継続判定、割り切れの向き。

\paragraph{考え方・正答を選ぶ理由（査読待ち初稿）} 合成数N=\allowbreak{}uvでは片方の因子が√N以下である。I=\allowbreak{}2からI≤√Nの間にN\%I=\allowbreak{}0を調べ、真なら非素数、偽ならIを増やす。範囲を越えたら素数である。5となる。平方根が整数のときも含めるため≤が必要である。

\Needspace{4\baselineskip}\paragraph{選択肢の検討（査読待ち初稿）}\mbox{}\par

\noindent\textbf{1：不適切} I≥√N・IがNで割り切れる。継続範囲が逆で割り算の向きも逆である。\par

\noindent\textbf{2：不適切} I≥√N・NがIで割り切れない。範囲と非素数へ進む判定が逆である。\par

\noindent\textbf{3：不適切} I≥√N・NがIで割り切れる。範囲が逆でN=\allowbreak{}9等を正しく調べられない。\par

\noindent\textbf{4：不適切} I≤√N・NがIで割り切れない。割り切れない時に非素数と出すので逆である。\par

\noindent\textbf{5：適切} I≤√N・NがIで割り切れる。2から平方根まで調べ、割り切れるなら非素数である。\par

\paragraph{再利用する考え方} 定義した対象量、演算順序、保存条件を固定し、途中状態と候補の違いを比較する。

\paragraph{共通理論} \hyperref[theory:discrete:invariants]{第\ref*{theory:discrete:invariants}節}

\paragraph{関連問題} \hyperref[q:h24-reg-basic-2-5]{H24-REG-BASIC-2-5}

\par 難易度 2（校準初稿）、優先度未評価。

\section{漸近上界と増大の階層}

\subsection{令和3年度 問2-6}
\label{q:r03-reg-basic-2-6}

\noindent\textbf{原典：} sources/R03/basic.pdf、PDF物理10頁。

\textbf{公式正答：} 3（single）。

\paragraph{問題の要旨} 漸近上界の定義から多項式・対数・指数・二重指数の四記述を判定する。

\Needspace{4\baselineskip}\paragraph{確認した条件・記述（独自要約）}\mbox{}\par

\noindent\textbf{ア：} 5n\textasciicircum{}3+1=\allowbreak{}O(n\textasciicircum{}3)\par

\noindent\textbf{イ：} n log2 n=\allowbreak{}O(n\textasciicircum{}1.5)\par

\noindent\textbf{ウ：} n\textasciicircum{}3 3\textasciicircum{}n=\allowbreak{}O(4\textasciicircum{}n)\par

\noindent\textbf{エ：} 2\textasciicircum{}(2\textasciicircum{}n)=\allowbreak{}O(10\textasciicircum{}(n\textasciicircum{}100))\par

\paragraph{問われる核心} 漸近上界の定義から多項式・対数・指数・二重指数の四記述を判定する。

\paragraph{必要知識} Oは十分大きいnで定数倍に抑えられる上界である。対数・多項式・指数の比を調べる。

\paragraph{考え方・正答を選ぶ理由（査読待ち初稿）} アは(5nの3乗+1)/nの3乗≤6。イのlog2(n)/√nは0へ、ウのnの3乗×(3/4)のn乗も0へ収束する。エは2の(2のn乗)と10の(nの100乗)の比較で、比の対数は2のn乗ln2−nの100乗ln10となり正無限大へ発散する。従って正正正誤、3である。

\Needspace{4\baselineskip}\paragraph{選択肢の検討（査読待ち初稿）}\mbox{}\par

\noindent\textbf{1：不適切} 正・誤・誤・誤。対数と多項式、3/4の指数減衰を不適切に否定する。\par

\noindent\textbf{2：不適切} 正・正・誤・正。ウは正でエは誤なので二か所が逆である。\par

\noindent\textbf{3：適切} 正・正・正・誤。ア・イ・ウの比は有界で、エは発散する。\par

\noindent\textbf{4：不適切} 正・誤・正・誤。イも正なので対数の成長を過大評価する。\par

\noindent\textbf{5：不適切} 誤・誤・誤・正。四記述の判断が逆である。\par

\paragraph{再利用する考え方} 定義した対象量、演算順序、保存条件を固定し、途中状態と候補の違いを比較する。

\paragraph{共通理論} \hyperref[theory:discrete:asymptotic]{第\ref*{theory:discrete:asymptotic}節}

\par 参考：\cite{mit_discrete2015}。

\par 難易度 3（校準初稿）、優先度未評価。

\section{二分探索木・挿入順の祖先制約}

\subsection{令和元年度 問2-2}
\label{q:r01-reg-basic-2-2}

\noindent\textbf{原典：} sources/R01/basic.pdf、PDF物理8頁。

\textbf{公式正答：} 1（single）。

\paragraph{問題の要旨} 同じ二分探索木を作る挿入順を、部分木の祖先が先に入る条件から選ぶ。

\Needspace{4\baselineskip}\paragraph{確認した条件・記述（独自要約）}\mbox{}\par

\noindent\textbf{判断対象：} 二分探索木の大小分岐と親子の挿入先行。\par

\noindent\textbf{原典確認：} 図・式・全5候補を29保存画像で確認。rawの数式欠落と演算子誤認は別校正JSONで扱う。\par

\paragraph{問われる核心} 同じ二分探索木を作る挿入順を、部分木の祖先が先に入る条件から選ぶ。

\paragraph{必要知識} 二分探索木の大小分岐と親子の挿入先行。

\paragraph{考え方・正答を選ぶ理由（査読待ち初稿）} 原典の根8、左5、右12、5の左右3/7、7の左6、12の左10を保つ。挿入順は親を子より先にし、別枝は前後してよい。1だけが12より前に10を置かず、7より前に6を置かない。

\Needspace{4\baselineskip}\paragraph{選択肢の検討（査読待ち初稿）}\mbox{}\par

\noindent\textbf{1：適切} 8,5,7,12,3,10,6。12を10より、7を6より先に入れ、根8と左根5も維持する。\par

\noindent\textbf{2：不適切} 8,5,7,10,3,12,6。10が12より先なので右部分木の根が10になる。\par

\noindent\textbf{3：不適切} 8,5,6,12,3,10,7。6が7より先なので5の右部分木の根が6になる。\par

\noindent\textbf{4：不適切} 8,5,3,10,7,12,6。10が12より先なので右部分木が変わる。\par

\noindent\textbf{5：不適切} 8,5,3,12,6,10,7。6が7より先なので左側部分木が変わる。\par

\paragraph{再利用する考え方} 定義した対象量、演算順序、保存条件を固定し、途中状態と候補の違いを比較する。

\paragraph{共通理論} \hyperref[theory:discrete:bst]{第\ref*{theory:discrete:bst}節}

\par 参考：\cite{princeton_bst}。

\par 難易度 2（校準初稿）、優先度未評価。

\section{論理式・命題}

\subsection{平成23年度 問2-3}
\label{q:h23-reg-basic-2-3}

\noindent\textbf{原典：} sources/H23/basic.pdf、PDF物理6頁。

\textbf{公式正答：} 2（single）。

\paragraph{問題の要旨} 否定上線の範囲を守り、与えられた論理式と等価な候補を選ぶ。

\Needspace{4\baselineskip}\paragraph{確認した条件・記述（独自要約）}\mbox{}\par

\noindent\textbf{対象式：} \(X=(\overline A+B)\overline{(A+\overline B)}\)\par

\noindent\textbf{演算：} ＋は論理和、積は論理積。第2候補は和全体と第2項の二重の否定線をもつ。\par

\paragraph{問われる核心} 否定の作用範囲と論理等価

\paragraph{必要知識} ド・モルガン則、吸収律、二重否定、真理値表。

\paragraph{考え方（査読待ち初稿）} \(\overline{A+\overline B}=\overline A B\)なので対象は\((\overline A+B)\overline A B=\overline A B\)である。第2候補も\(\overline{A\overline B+(A+\overline B)}=\overline{A+\overline B}=\overline A B\)となる。上線を個々の項に移す前に、和全体を覆う線を確定する。

\paragraph{正答を選ぶ理由} 対象と第2候補はともにA=\allowbreak{}0、B=\allowbreak{}1のときだけ1となる。

\Needspace{4\baselineskip}\paragraph{選択肢の検討（査読待ち初稿）}\mbox{}\par

\noindent\textbf{1：条件に一致しない候補} \((\overline A+B)+(\overline A B)\)。吸収律で否定A又はBとなり、A=\allowbreak{}B=\allowbreak{}0で1を返す。対象式は0である。\par

\noindent\textbf{2：条件に一致する候補} \(\overline{A\overline B+\overline{\overline A B}}\)。二重否定とド・モルガン則により否定AかつBへ簡約され、対象式と一致する。\par

\noindent\textbf{3：条件に一致しない候補} \((A\overline B)(\overline A B)\)。Aと否定Aの積が含まれ常に0である。A=\allowbreak{}0、B=\allowbreak{}1の対象値1を再現しない。\par

\noindent\textbf{4：条件に一致しない候補} \(\overline{(\overline A+B)+(\overline A+B)}\)。同一項の和の否定なのでAかつ否定Bとなり、対象と逆の入力で1になる。\par

\noindent\textbf{5：条件に一致しない候補} \((A+\overline B)(\overline A+\overline B)\)。分配・吸収で否定Bに簡約される。A=\allowbreak{}0、B=\allowbreak{}0で1だが対象は0である。\par

\paragraph{再利用する考え方} 定義と条件から否定の作用範囲と論理等価を構成し、成立の証明と誤った候補の反例を区別する。

\paragraph{共通理論} \hyperref[theory:information:boolean]{第\ref*{theory:information:boolean}節}

\par 難易度 2（校準初稿）、優先度未評価。

\subsection{平成25年度 問2-3}
\label{q:h25-reg-basic-2-3}

\noindent\textbf{原典：} sources/H25/basic.pdf、PDF物理8頁。

\textbf{公式正答：} 1（single）。

\paragraph{問題の要旨} 2つの積項にかかる否定を変形し、同じ論理関数を選ぶ。

\Needspace{4\baselineskip}\paragraph{確認した条件・記述（独自要約）}\mbox{}\par

\noindent\textbf{対象式：} \(X=\overline{(\overline A\overline B)}\,\overline{(\overline A+B)}\)\par

\paragraph{問われる核心} ド・モルガン変形と吸収

\paragraph{必要知識} 積と和の否定、変数の否定、相補律、吸収律。

\paragraph{考え方（査読待ち初稿）} \(\overline{\overline A\overline B}=A+B\)、\(\overline{\overline A+B}=A\overline B\)である。したがって\(X=(A+B)(A\overline B)=A\overline B\)。第1候補が最初の変形式と一致する。

\paragraph{正答を選ぶ理由} 対象はA=\allowbreak{}1、B=\allowbreak{}0でのみ1となり、第1候補と一致する。

\Needspace{4\baselineskip}\paragraph{選択肢の検討（査読待ち初稿）}\mbox{}\par

\noindent\textbf{1：条件に一致する候補} \((A+B)(A\overline B)\)。対象の2つの否定をド・モルガン則で置き換えた式であり一致する。\par

\noindent\textbf{2：条件に一致しない候補} \((A+B)(\overline A B)\)。否定AかつBとなり、A=\allowbreak{}1、B=\allowbreak{}0では0である。対象は1となる。\par

\noindent\textbf{3：条件に一致しない候補} \((\overline A+B)(A\overline B)\)。展開するとAと否定A又はBと否定Bが同時に必要となり、常に0である。\par

\noindent\textbf{4：条件に一致しない候補} \(\overline A\overline B+(A+\overline B)\)。A又は否定Bへ簡約され、A=\allowbreak{}B=\allowbreak{}0でも1となる。対象は0である。\par

\noindent\textbf{5：条件に一致しない候補} \(\overline A\overline B+(\overline A+B)\)。否定A又はBへ簡約され、A=\allowbreak{}1、B=\allowbreak{}0で0になる。対象は1である。\par

\paragraph{再利用する考え方} 定義と条件からド・モルガン変形と吸収を構成し、成立の証明と誤った候補の反例を区別する。

\paragraph{共通理論} \hyperref[theory:information:boolean]{第\ref*{theory:information:boolean}節}

\par 難易度 2（校準初稿）、優先度未評価。

\subsection{平成25年度 問2-4}
\label{q:h25-reg-basic-2-4}

\noindent\textbf{原典：} sources/H25/basic.pdf、PDF物理8頁。

\textbf{公式正答：} 4（single）。

\paragraph{問題の要旨} 全100万件中の2語の出現数から、一方のみを含む件数の可能範囲を求める。

\Needspace{4\baselineskip}\paragraph{確認した条件・記述（独自要約）}\mbox{}\par

\noindent\textbf{全集合：} 100万件\par

\noindent\textbf{L：} 論理を含む65万件\par

\noindent\textbf{I：} 情報を含む55万件\par

\noindent\textbf{求める量：} Lを含みIを含まない件数k。\par

\paragraph{問われる核心} 集合の重なりの上下限

\paragraph{必要知識} 包含排除、共通部分は小さい方以下、和集合は全集合以下。

\paragraph{考え方（査読待ち初稿）} 重なりを\(t=|L\cap I|\)とおく。和集合が100万を超えないので\(t\ge65+55-100=20\)万、一方\(t\le55\)万である。\(k=65-t\)より\(10\le k\le45\)万となる。

\paragraph{正答を選ぶ理由} 包含排除から重なりの両端20万・55万を求めるとkは10〜45万となる。

\Needspace{4\baselineskip}\paragraph{選択肢の検討（査読待ち初稿）}\mbox{}\par

\noindent\textbf{1：条件に一致しない候補} 0〜35万。重なりは最大55万なのでkの最小値は10万であり、0まで許す範囲は不正確である。\par

\noindent\textbf{2：条件に一致しない候補} 0〜45万。最大値45万は合うが、最小値は10万である。\par

\noindent\textbf{3：条件に一致しない候補} 0〜65万。最小値も最大値も実際の可能範囲より広い。\par

\noindent\textbf{4：条件に一致する候補} 10〜45万。重なり20〜55万に対して65万から差し引いた範囲10〜45万に一致する。\par

\noindent\textbf{5：条件に一致しない候補} 10〜65万。最小値は合うが、重なりは少なくとも20万あり、最大値は45万となる。\par

\paragraph{再利用する考え方} 定義と条件から集合の重なりの上下限を構成し、成立の証明と誤った候補の反例を区別する。

\paragraph{共通理論} \hyperref[theory:information:sets]{第\ref*{theory:information:sets}節}

\paragraph{関連問題} \hyperref[q:r01-ret-basic-2-5]{R01-RET-BASIC-2-5}

\par 難易度 2（校準初稿）、優先度未評価。

\subsection{平成27年度 問2-2}
\label{q:h27-reg-basic-2-2}

\noindent\textbf{原典：} sources/H27/basic.pdf、PDF物理5頁。

\textbf{公式正答：} 5（single）。

\paragraph{問題の要旨} 住民CとDの賛否規則を論理式へ変換し、必ず成り立つ記述を選ぶ。

\Needspace{4\baselineskip}\paragraph{確認した条件・記述（独自要約）}\mbox{}\par

\noindent\textbf{値：} 賛成1、反対0。住民は必ずどちらかを決める。\par

\noindent\textbf{C：} AとBの両方が賛成のときだけ反対。\par

\noindent\textbf{D：} AとCの両方が賛成のときだけ反対。\par

\paragraph{問われる核心} 自然言語のNANDと含意

\paragraph{必要知識} 論理積の否定、代入、吸収律、十分条件。

\paragraph{考え方（査読待ち初稿）} \(C=\overline{AB}\)、\(D=\overline{AC}\)である。代入して\(D=\overline A+AB=\overline A+B\)となる。A,Bの入力00,01,10,11に対するDは1,1,0,1であり、B=\allowbreak{}1の2行では必ずD=\allowbreak{}1となる。

\paragraph{正答を選ぶ理由} B=\allowbreak{}1を代入するとD=否定A又は1=\allowbreak{}1となる。

\Needspace{4\baselineskip}\paragraph{選択肢の検討（査読待ち初稿）}\mbox{}\par

\noindent\textbf{1：条件に一致しない候補} Cの賛成はAとBがともに賛成のときだけ。A=\allowbreak{}B=\allowbreak{}1ならC=\allowbreak{}0であり規則と逆である。\par

\noindent\textbf{2：条件に一致しない候補} Cの賛成はAとBの賛否が異なるときだけ。A=\allowbreak{}B=\allowbreak{}0でもC=\allowbreak{}1なので、異なるときだけという限定は誤りである。\par

\noindent\textbf{3：条件に一致しない候補} Dの賛成はAとBがともに賛成のときだけ。A=\allowbreak{}B=\allowbreak{}0でもD=\allowbreak{}1となり、両者賛成のときだけではない。\par

\noindent\textbf{4：条件に一致しない候補} Dの賛成はAとBの賛否が異なるときだけ。A=\allowbreak{}B=\allowbreak{}1でもD=\allowbreak{}1であり、異なるときだけではない。\par

\noindent\textbf{5：条件に一致する候補} Bが賛成すれば必ずDも賛成。Dは否定A又はBに簡約されるので、B=\allowbreak{}1ならAにかかわらずD=\allowbreak{}1となる。\par

\paragraph{再利用する考え方} 定義と条件から自然言語のNANDと含意を構成し、成立の証明と誤った候補の反例を区別する。

\paragraph{共通理論} \hyperref[theory:information:boolean]{第\ref*{theory:information:boolean}節}

\par 難易度 2（校準初稿）、優先度未評価。

\subsection{平成28年度 問2-2}
\label{q:h28-reg-basic-2-2}

\noindent\textbf{原典：} sources/H28/basic.pdf、PDF物理8頁。

\textbf{公式正答：} 2（single）。

\paragraph{問題の要旨} 2入力がともに1のときだけ0を出す真理値表に合う論理式を選ぶ。

\Needspace{4\baselineskip}\paragraph{確認した条件・記述（独自要約）}\mbox{}\par

\noindent\textbf{入力順：} (A,B)=\allowbreak{}(0,0),(0,1),(1,0),(1,1)\par

\noindent\textbf{出力：} 1,1,1,0\par

\paragraph{問われる核心} 真理値表とNAND

\paragraph{必要知識} AND・OR・NOT・XORの表、ド・モルガン則。

\paragraph{考え方（査読待ち初稿）} 表の0の行が\(AB=1\)の唯一の行なので、対象は\(\overline{AB}=\overline A+\overline B\)である。第2候補を選ぶ。

\paragraph{正答を選ぶ理由} 第2候補の出力1110が4行すべてで一致する。

\Needspace{4\baselineskip}\paragraph{選択肢の検討（査読待ち初稿）}\mbox{}\par

\noindent\textbf{1：条件に一致しない候補} \(A+B\)。ORは00で0となり指定表と異なる。\par

\noindent\textbf{2：条件に一致する候補} \(\overline A+\overline B\)。ド・モルガン則でANDの否定となり、11以外の全行で1となる。\par

\noindent\textbf{3：条件に一致しない候補} \(AB\)。ANDは11でのみ1となり、指定の否定と逆である。\par

\noindent\textbf{4：条件に一致しない候補} \(\overline A B\)。01でしか1にならず、00や10の出力を再現しない。\par

\noindent\textbf{5：条件に一致しない候補} \(A\overline B+\overline A B\)。XORは00で0となり、その行が指定表と異なる。\par

\paragraph{再利用する考え方} 定義と条件から真理値表とNANDを構成し、成立の証明と誤った候補の反例を区別する。

\paragraph{共通理論} \hyperref[theory:information:boolean]{第\ref*{theory:information:boolean}節}

\par 難易度 2（校準初稿）、優先度未評価。

\subsection{平成29年度 問2-4}
\label{q:h29-reg-basic-2-4}

\noindent\textbf{原典：} sources/H29/basic.pdf、PDF物理9頁。

\textbf{公式正答：} 3（single）。

\paragraph{問題の要旨} 4・100・400によるうるう年規則を、条件と動作の決定表へ対応させる。

\Needspace{4\baselineskip}\paragraph{確認した条件・記述（独自要約）}\mbox{}\par

\noindent\textbf{規則：} 4で割り切れ100で割り切れない、又は400で割り切れる年がうるう年。\par

\noindent\textbf{表条件3行：} 4、100、400でそれぞれ割り切れるか。\par

\noindent\textbf{4列：} (N,−,−)、(Y,N,−)、(Y,Y,N)、(Y,Y,Y)。全候補の条件部は一致。\par

\noindent\textbf{動作行：} うるう年／平年の排他的2行。候補値は表のうるう年実行をY、平年実行をNへ符号化したもの。\par

\noindent\textbf{表記：} 原典のX=実行、−=不実行。条件部の−は不要又は不可能。\par

\paragraph{問われる核心} 含意制約のある決定表

\paragraph{必要知識} 400の倍数⇒100の倍数⇒4の倍数、AND・OR・NOT、条件の不要欄、動作の排他。

\paragraph{考え方（査読待ち初稿）} 400の倍数なら100と4の倍数でもある。実現する4条件列は「4で不可」「4で可・100で不可」「100で可・400で不可」「400で可」であり、判定は平年、うるう年、平年、うるう年の順になる。論理式は\(L=F_4\overline{F_{100}}+F_{400}\)である。4列すべてをこの順に照合する。

\paragraph{正答を選ぶ理由} 4条件列の判定NYNYをすべて満たす表を選ぶ。

\Needspace{4\baselineskip}\paragraph{選択肢の検討（査読待ち初稿）}\mbox{}\par

\noindent\textbf{1：条件に一致しない候補} 4条件列でのうるう年判定：NYYY（Y=うるう年、N=平年）。1900年は平年であるがうるう年としている。\par

\noindent\textbf{2：条件に一致しない候補} 4条件列でのうるう年判定：NNYY（Y=うるう年、N=平年）。2024年はうるう年であるが平年としている。1900年は平年であるがうるう年としている。\par

\noindent\textbf{3：条件に一致する候補} 4条件列でのうるう年判定：NYNY（Y=うるう年、N=平年）。4つの有効な条件列がすべて規則に一致する。\par

\noindent\textbf{4：条件に一致しない候補} 4条件列でのうるう年判定：NYNN（Y=うるう年、N=平年）。2000年はうるう年であるが平年としている。\par

\noindent\textbf{5：条件に一致しない候補} 4条件列でのうるう年判定：NNNY（Y=うるう年、N=平年）。2024年はうるう年であるが平年としている。\par

\paragraph{再利用する考え方} 定義と条件から含意制約のある決定表を構成し、成立の証明と誤った候補の反例を区別する。

\paragraph{共通理論} \hyperref[theory:information:decision]{第\ref*{theory:information:decision}節}

\paragraph{関連問題} \hyperref[q:r03-reg-basic-2-4]{R03-REG-BASIC-2-4}

\par 難易度 2（校準初稿）、優先度未評価。

\subsection{平成30年度 問2-4}
\label{q:h30-reg-basic-2-4}

\noindent\textbf{原典：} sources/H30/basic.pdf、PDF物理8頁。

\textbf{公式正答：} 5（single）。

\paragraph{問題の要旨} 入力一致の検出を否定した論理式を、等価な排他的論理和へ変形する。

\Needspace{4\baselineskip}\paragraph{確認した条件・記述（独自要約）}\mbox{}\par

\noindent\textbf{対象式：} \(X=\overline{\overline A\overline B+AB}\)\par

\paragraph{問われる核心} XNORの否定とXOR

\paragraph{必要知識} 同値検出、排他的論理和、ド・モルガン則、相補律。

\paragraph{考え方（査読待ち初稿）} 対象は\(\overline{\overline A\overline B+AB}\)である。括弧内は入力が等しいとき1なので、その否定はXORである。ド・モルガン則で\((A+B)(\overline A+\overline B)=(A+B)\overline{AB}\)とも表せる。

\paragraph{正答を選ぶ理由} 第5候補の出力0110が対象のXORと一致する。

\Needspace{4\baselineskip}\paragraph{選択肢の検討（査読待ち初稿）}\mbox{}\par

\noindent\textbf{1：条件に一致しない候補} \((A+B)\overline{(A+B)}\)。任意の命題とその否定の積なので常に0である。\par

\noindent\textbf{2：条件に一致しない候補} \((A+B)(\overline A B)\)。否定AかつBのみを表し、A=\allowbreak{}1、B=\allowbreak{}0のXOR出力1を再現しない。\par

\noindent\textbf{3：条件に一致しない候補} \((AB)(\overline A\overline B)\)。Aと否定Aを同時に要求し、常に0である。\par

\noindent\textbf{4：条件に一致しない候補} \((AB)\overline{(AB)}\)。命題ABとその否定の積で常に0である。\par

\noindent\textbf{5：条件に一致する候補} \((A+B)\overline{(AB)}\)。ORが1かつANDが0という条件なので、入力が異なる2行でのみ1になり一致する。\par

\paragraph{再利用する考え方} 定義と条件からXNORの否定とXORを構成し、成立の証明と誤った候補の反例を区別する。

\paragraph{共通理論} \hyperref[theory:information:boolean]{第\ref*{theory:information:boolean}節}

\paragraph{関連問題} \hyperref[q:r03-reg-basic-2-2]{R03-REG-BASIC-2-2}

\par 難易度 2（校準初稿）、優先度未評価。

\subsection{令和元年度再試験 問2-5}
\label{q:r01-ret-basic-2-5}

\noindent\textbf{原典：} sources/R01\_retest/basic.pdf、PDF物理10頁。

\textbf{公式正答：} 1（single）。

\paragraph{問題の要旨} 情報の出現数と情報・論理の共通部分から、論理を含まない件数の範囲を求める。

\Needspace{4\baselineskip}\paragraph{確認した条件・記述（独自要約）}\mbox{}\par

\noindent\textbf{全集合：} 100万件\par

\noindent\textbf{I：} 情報を含む65万件\par

\noindent\textbf{共通部分：} 情報と論理の両方を含む55万件\par

\noindent\textbf{求める量：} 論理を含まない件数k。\par

\paragraph{問われる核心} 集合の部分情報と補集合の範囲

\paragraph{必要知識} 集合の差、既知共通部分、未知領域の上下限。

\paragraph{考え方（査読待ち初稿）} 情報を含むが論理を含まない部分は65-55=\allowbreak{}10万件である。情報外の35万件は、すべて論理を含む場合から一切含まない場合までありうる。よってkは10万から10+35=\allowbreak{}45万までである。H25の2語それぞれの出現数の条件と取り違えない。

\paragraph{正答を選ぶ理由} 必ず論理を含まない10万件と、残りの可変0〜35万件の和である。

\Needspace{4\baselineskip}\paragraph{選択肢の検討（査読待ち初稿）}\mbox{}\par

\noindent\textbf{1：条件に一致する候補} 10〜45万。情報を含み論理を含まない10万件に、情報を含まない35万件のうち0〜35万件が加わり、正確な範囲である。\par

\noindent\textbf{2：条件に一致しない候補} 10〜55万。上限55万は、論理を含む55万件以上が既にあることに反する。\par

\noindent\textbf{3：条件に一致しない候補} 10〜65万。上限65万は、論理を含む件数の既知下限55万を満たさない。\par

\noindent\textbf{4：条件に一致しない候補} 45〜65万。下限45万は誤りであり、情報外の全35万件が論理を含むならk=\allowbreak{}10万である。\par

\noindent\textbf{5：条件に一致しない候補} 45〜90万。上限90万と下限45万はいずれも既知の共通部分と情報のみの10万件から得る範囲に合わない。\par

\paragraph{再利用する考え方} 定義と条件から集合の部分情報と補集合の範囲を構成し、成立の証明と誤った候補の反例を区別する。

\paragraph{共通理論} \hyperref[theory:information:sets]{第\ref*{theory:information:sets}節}

\paragraph{関連問題} \hyperref[q:h25-reg-basic-2-4]{H25-REG-BASIC-2-4}

\par 難易度 2（校準初稿）、優先度未評価。

\subsection{令和元年度再試験 問2-6}
\label{q:r01-ret-basic-2-6}

\noindent\textbf{原典：} sources/R01\_retest/basic.pdf、PDF物理10頁。

\textbf{公式正答：} 4（single）。

\paragraph{問題の要旨} 4要素と2要素の直積から2要素集合へ送る写像の総数を求める。

\Needspace{4\baselineskip}\paragraph{確認した条件・記述（独自要約）}\mbox{}\par

\noindent\textbf{集合：} A=\{a,b,c,d\}、B=\{α,β\}、C=\{0,1\}。\par

\noindent\textbf{対象：} 任意の写像f:A×B→C。単射・全射などの追加制約なし。\par

\paragraph{問われる核心} 直積と写像の個数

\paragraph{必要知識} 積の法則、定義域の濃度、写像・単射・全射の区別。

\paragraph{考え方（査読待ち初稿）} 定義域A×Bの要素数は4×2=\allowbreak{}8である。各入力の像を0か1から独立に選べるので、写像は2\textsuperscript{8}=256通りである。同じ像が何回使われてもよく、像として両方を使う義務もない。

\paragraph{正答を選ぶ理由} 8入力について各2通りの選択を積の法則で数え、256となる。

\Needspace{4\baselineskip}\paragraph{選択肢の検討（査読待ち初稿）}\mbox{}\par

\noindent\textbf{1：条件に一致しない候補} 32。2\textsuperscript{5}であり、独立に像を選ぶ8要素を5要素と誤って数えている。\par

\noindent\textbf{2：条件に一致しない候補} 64。2\textsuperscript{6}であり、直積を4+2=\allowbreak{}6とする和の計数と混同している。\par

\noindent\textbf{3：条件に一致しない候補} 128。2\textsuperscript{7}であり、直積の8個すべての像の選択を数えていない。\par

\noindent\textbf{4：条件に一致する候補} 256。直積には4×2=\allowbreak{}8個あり、各要素に2通りの像を独立に選べるので2\textsuperscript{8}である。\par

\noindent\textbf{5：条件に一致しない候補} 512。2\textsuperscript{9}であり、独立の入力を1個多く数えている。\par

\paragraph{再利用する考え方} 定義と条件から直積と写像の個数を構成し、成立の証明と誤った候補の反例を区別する。

\paragraph{共通理論} \hyperref[theory:information:sets]{第\ref*{theory:information:sets}節}

\par 難易度 2（校準初稿）、優先度未評価。

\subsection{令和2年度 問2-2}
\label{q:r02-reg-basic-2-2}

\noindent\textbf{原典：} sources/R02/basic.pdf、PDF物理7頁。

\textbf{公式正答：} 3（single）。

\paragraph{問題の要旨} 3入力の真理値表を満たす論理式を選ぶ。

\Needspace{4\baselineskip}\paragraph{確認した条件・記述（独自要約）}\mbox{}\par

\noindent\textbf{入力順：} 000,001,010,011,100,101,110,111\par

\noindent\textbf{出力：} 0,1,0,0,0,1,1,1\par

\paragraph{問われる核心} 条件による論理関数の分解

\paragraph{必要知識} 真理値表、場合分け、選択回路、論理積の和。

\paragraph{考え方（査読待ち初稿）} y=\allowbreak{}0の4行を見ると出力はzであり、y=\allowbreak{}1の4行では出力はxである。場合分けを\(f=\overline y z+yx\)と書けば、第3候補となる。選択入力yが0ならz、1ならxを選ぶ2入力選択回路と解釈できる。

\paragraph{正答を選ぶ理由} 第3候補はyの両値でそれぞれ指定表を満たし、全8行に一致する。

\Needspace{4\baselineskip}\paragraph{選択肢の検討（査読待ち初稿）}\mbox{}\par

\noindent\textbf{1：条件に一致しない候補} \(xy+z\)。011でz=\allowbreak{}1により1となるが、指定表は0である。\par

\noindent\textbf{2：条件に一致しない候補} \(\overline x y+\overline y\overline z\)。000で否定yかつ否定zが1になり、指定表の0と異なる。\par

\noindent\textbf{3：条件に一致する候補} \(xy+\overline y z\)。y=\allowbreak{}0ではz、y=\allowbreak{}1ではxを返す。全8行の指定表を再現する。\par

\noindent\textbf{4：条件に一致しない候補} \(xy+\overline{xy}\)。命題又はその否定なので常に1であり、000等の0を再現しない。\par

\noindent\textbf{5：条件に一致しない候補} \(xy+\overline x z\)。011で否定xかつzが1となるほか、101で0となり、指定表と一致しない。\par

\paragraph{再利用する考え方} 定義と条件から条件による論理関数の分解を構成し、成立の証明と誤った候補の反例を区別する。

\paragraph{共通理論} \hyperref[theory:information:boolean]{第\ref*{theory:information:boolean}節}

\par 難易度 2（校準初稿）、優先度未評価。

\subsection{令和3年度 問2-2}
\label{q:r03-reg-basic-2-2}

\noindent\textbf{原典：} sources/R03/basic.pdf、PDF物理6頁。

\textbf{公式正答：} 2（single）。

\paragraph{問題の要旨} 入力が等しい2行の和を否定した式に等価な候補を選ぶ。

\Needspace{4\baselineskip}\paragraph{確認した条件・記述（独自要約）}\mbox{}\par

\noindent\textbf{対象式：} \(\overline{\overline A\overline B+AB}\)\par

\paragraph{問われる核心} 否定の作用範囲とXOR

\paragraph{必要知識} ド・モルガン則、相補律、二重否定。

\paragraph{考え方（査読待ち初稿）} 対象は\(\overline{\overline A\overline B+AB}\)である。括弧内は入力が等しいとき1なので、その否定はXORである。ド・モルガン則で\((A+B)(\overline A+\overline B)=(A+B)\overline{AB}\)とも表せる。

\paragraph{正答を選ぶ理由} 第2候補だけが入力01・10で1、00・11で0となる。

\Needspace{4\baselineskip}\paragraph{選択肢の検討（査読待ち初稿）}\mbox{}\par

\noindent\textbf{1：条件に一致しない候補} \((A+B)\overline{(A+B)}\)。命題とその否定の積で常に0である。\par

\noindent\textbf{2：条件に一致する候補} \((A+B)(\overline A+\overline B)\)。入力の少なくとも一方が1、かつ少なくとも一方が0という条件でありXORを表す。\par

\noindent\textbf{3：条件に一致しない候補} \((AB)(\overline A\overline B)\)。Aと否定Aが同時に必要となり常に0である。\par

\noindent\textbf{4：条件に一致しない候補} \((AB)\overline{(AB)}\)。ABとその否定の積で常に0である。\par

\noindent\textbf{5：条件に一致しない候補} \((A+B)+(\overline A+\overline B)\)。A又は否定Aが含まれ常に1となり、00・11での指定出力0と異なる。\par

\paragraph{再利用する考え方} 定義と条件から否定の作用範囲とXORを構成し、成立の証明と誤った候補の反例を区別する。

\paragraph{共通理論} \hyperref[theory:information:boolean]{第\ref*{theory:information:boolean}節}

\paragraph{関連問題} \hyperref[q:h30-reg-basic-2-4]{H30-REG-BASIC-2-4}

\par 難易度 2（校準初稿）、優先度未評価。

\subsection{令和3年度 問2-4}
\label{q:r03-reg-basic-2-4}

\noindent\textbf{原典：} sources/R03/basic.pdf、PDF物理8頁。

\textbf{公式正答：} 2（single）。

\paragraph{問題の要旨} 4・100・400によるうるう年規則を、条件と動作の決定表へ対応させる。

\Needspace{4\baselineskip}\paragraph{確認した条件・記述（独自要約）}\mbox{}\par

\noindent\textbf{規則：} 4で割り切れ100で割り切れない、又は400で割り切れる年がうるう年。\par

\noindent\textbf{表条件3行：} 4、100、400でそれぞれ割り切れるか。\par

\noindent\textbf{4列：} (N,−,−)、(Y,N,−)、(Y,Y,N)、(Y,Y,Y)。全候補の条件部は一致。\par

\noindent\textbf{動作行：} うるう年／平年の排他的2行。候補値は表のうるう年実行をY、平年実行をNへ符号化したもの。\par

\noindent\textbf{表記：} 原典のX=実行、−=不実行。条件部の−は不要又は不可能。\par

\paragraph{問われる核心} 含意制約のある決定表

\paragraph{必要知識} 400の倍数⇒100の倍数⇒4の倍数、AND・OR・NOT、条件の不要欄、動作の排他。

\paragraph{考え方（査読待ち初稿）} 400の倍数なら100と4の倍数でもある。実現する4条件列は「4で不可」「4で可・100で不可」「100で可・400で不可」「400で可」であり、判定は平年、うるう年、平年、うるう年の順になる。論理式は\(L=F_4\overline{F_{100}}+F_{400}\)である。4列すべてをこの順に照合する。

\paragraph{正答を選ぶ理由} 4条件列の判定NYNYをすべて満たす表を選ぶ。

\Needspace{4\baselineskip}\paragraph{選択肢の検討（査読待ち初稿）}\mbox{}\par

\noindent\textbf{1：条件に一致しない候補} 4条件列でのうるう年判定：NYYY（Y=うるう年、N=平年）。1900年は平年であるがうるう年としている。\par

\noindent\textbf{2：条件に一致する候補} 4条件列でのうるう年判定：NYNY（Y=うるう年、N=平年）。4つの有効な条件列がすべて規則に一致する。\par

\noindent\textbf{3：条件に一致しない候補} 4条件列でのうるう年判定：NNYY（Y=うるう年、N=平年）。2024年はうるう年であるが平年としている。1900年は平年であるがうるう年としている。\par

\noindent\textbf{4：条件に一致しない候補} 4条件列でのうるう年判定：NYNN（Y=うるう年、N=平年）。2000年はうるう年であるが平年としている。\par

\noindent\textbf{5：条件に一致しない候補} 4条件列でのうるう年判定：NNNY（Y=うるう年、N=平年）。2024年はうるう年であるが平年としている。\par

\paragraph{再利用する考え方} 定義と条件から含意制約のある決定表を構成し、成立の証明と誤った候補の反例を区別する。

\paragraph{共通理論} \hyperref[theory:information:decision]{第\ref*{theory:information:decision}節}

\paragraph{関連問題} \hyperref[q:h29-reg-basic-2-4]{H29-REG-BASIC-2-4}

\par 難易度 2（校準初稿）、優先度未評価。

\subsection{令和5年度 問2-5}
\label{q:r05-reg-basic-2-5}

\noindent\textbf{原典：} sources/R05/basic.pdf、PDF物理8頁。

\textbf{公式正答：} 1（single）。

\paragraph{問題の要旨} 2つの8ビット列を含む反復XORと最後のANDの結果を求める。

\Needspace{4\baselineskip}\paragraph{確認した条件・記述（独自要約）}\mbox{}\par

\noindent\textbf{A：} 01011101\par

\noindent\textbf{B：} 10101101\par

\noindent\textbf{演算：} \(C=(((A\oplus B)\oplus B)\oplus A)\cdot A\)\par

\noindent\textbf{定義：} XORとANDは各桁で独立に作用する。\par

\paragraph{問われる核心} XORの自己逆性とビット列演算

\paragraph{必要知識} XORの交換・結合律、自己相殺、ANDとゼロ。

\paragraph{考え方（査読待ち初稿）} XORの結合律と\(X\oplus X=0\)により、括弧内は\(A\oplus A\oplus B\oplus B=0\)である。最後に\(0\cdot A=0\)なので8ビットとも0になる。直接計算しても11110000→01011101→00000000→00000000となる。

\paragraph{正答を選ぶ理由} 同じビット列を2回XORすると消え、最後のANDでもゼロ列が保たれる。

\Needspace{4\baselineskip}\paragraph{選択肢の検討（査読待ち初稿）}\mbox{}\par

\noindent\textbf{1：条件に一致する候補} 00000000。XORで同じBと同じAがそれぞれ相殺され、ゼロ列とAのANDもゼロ列となる。\par

\noindent\textbf{2：条件に一致しない候補} 11111111。全ビット1ではなく、括弧内はゼロ列となる。\par

\noindent\textbf{3：条件に一致しない候補} 10101101。入力Bを残す候補だが、Bは2回XORされて消える。\par

\noindent\textbf{4：条件に一致しない候補} 01011101。入力Aを残す候補だが、括弧内でAも2回XORされて消える。\par

\noindent\textbf{5：条件に一致しない候補} 11110000。A XOR Bの途中結果であり、その後のBとAのXORを省略している。\par

\paragraph{再利用する考え方} 定義と条件からXORの自己逆性とビット列演算を構成し、成立の証明と誤った候補の反例を区別する。

\paragraph{共通理論} \hyperref[theory:information:boolean]{第\ref*{theory:information:boolean}節}

\par 難易度 2（校準初稿）、優先度未評価。

\subsection{令和5年度 問2-6}
\label{q:r05-reg-basic-2-6}

\noindent\textbf{原典：} sources/R05/basic.pdf、PDF物理8頁。

\textbf{公式正答：} 4（single）。

\paragraph{問題の要旨} 3部分集合の交差数と和集合の外側の数から全集合の元の個数を求める。

\Needspace{4\baselineskip}\paragraph{確認した条件・記述（独自要約）}\mbox{}\par

\noindent\textbf{各集合：} |A|=300、|B|=180、|C|=120。\par

\noindent\textbf{2集合交差：} |A∩B|=60、|A∩C|=40、|B∩C|=20。\par

\noindent\textbf{3集合交差：} |A∩B∩C|=10。\par

\noindent\textbf{外側：} A∪B∪Cの補集合の元は400個。\par

\paragraph{問われる核心} 3集合の包含排除と補集合

\paragraph{必要知識} 重複計数、3重交差の符号、全集合の分割。

\paragraph{考え方（査読待ち初稿）} 3集合の和集合は\(|A\cup B\cup C|=300+180+120-60-40-20+10=490\)である。3重交差の元は最初3回足し、2重交差として3回引いているので最後に1回加える。外側400を加えて全集合は890個となる。

\paragraph{正答を選ぶ理由} 包含排除で和集合490、補集合400と互いに排反な和をとり890となる。

\Needspace{4\baselineskip}\paragraph{選択肢の検討（査読待ち初稿）}\mbox{}\par

\noindent\textbf{1：条件に一致しない候補} 600。300+180+120の単純和であり、交差と外側400を扱っていない。\par

\noindent\textbf{2：条件に一致しない候補} 720。包含排除と外側を加えた890に一致しない。各領域の計数が必要である。\par

\noindent\textbf{3：条件に一致しない候補} 730。包含排除の和集合490と外側400の和に一致しない。\par

\noindent\textbf{4：条件に一致する候補} 890。600-(60+40+20)+10=\allowbreak{}490に外側400を加えた890に一致する。\par

\noindent\textbf{5：条件に一致しない候補} 1000。各集合の単純和600と外側400を加え、重複110を除いていない。\par

\paragraph{再利用する考え方} 定義と条件から3集合の包含排除と補集合を構成し、成立の証明と誤った候補の反例を区別する。

\paragraph{共通理論} \hyperref[theory:information:sets]{第\ref*{theory:information:sets}節}

\par 難易度 2（校準初稿）、優先度未評価。

\subsection{令和6年度 問2-2}
\label{q:r06-reg-basic-2-2}

\noindent\textbf{原典：} sources/R06/basic.pdf、PDF物理6頁。

\textbf{公式正答：} 2（single）。

\paragraph{問題の要旨} 実数の等式と不等式の2組について必要条件・十分条件を判定する。

\Needspace{4\baselineskip}\paragraph{確認した条件・記述（独自要約）}\mbox{}\par

\noindent\textbf{領域：} x,yは実数。\par

\noindent\textbf{ア：} x\textsuperscript{2}+y\textsuperscript{2}=0はx=\allowbreak{}y=\allowbreak{}0であるための条件。\par

\noindent\textbf{イ：} x+y>2はx>1又はy>1であるための条件。\par

\paragraph{問われる核心} 必要条件・十分条件と対偶

\paragraph{必要知識} 実数平方の非負性、含意の方向、対偶、反例。

\paragraph{考え方（査読待ち初稿）} 実数平方は非負なので和が0なら両者0、逆も成り立つ。イでx≤1かつy≤1ならx+y≤2なので、対偶によりx+y>2ならx>1又はy>1となる。逆はx=\allowbreak{}2、y=-100で成立しない。したがってアは必要十分、イは十分だが必要でない。

\paragraph{正答を選ぶ理由} アは同値、イは和が2より大きい条件から一方向に含意し逆向きに反例がある。

\Needspace{4\baselineskip}\paragraph{選択肢の検討（査読待ち初稿）}\mbox{}\par

\noindent\textbf{1：条件に一致しない候補} ア必要十分／イ必要だが十分でない。アは合うが、イの含意の向きが逆である。\par

\noindent\textbf{2：条件に一致する候補} ア必要十分／イ十分だが必要でない。実数平方の非負性でアは双方向に成立し、イは一方向のみ成立する。\par

\noindent\textbf{3：条件に一致しない候補} ア十分だが必要でない／イ必要十分。アも双方向に成立し、イは逆向きに成立しないので両欄が異なる。\par

\noindent\textbf{4：条件に一致しない候補} ア十分だが必要でない／イ必要だが十分でない。アは必要でもあり、イは十分側であるため両欄が異なる。\par

\noindent\textbf{5：条件に一致しない候補} ア必要だが十分でない／イ十分だが必要でない。イは合うが、アは十分でもある。\par

\paragraph{再利用する考え方} 定義と条件から必要条件・十分条件と対偶を構成し、成立の証明と誤った候補の反例を区別する。

\paragraph{共通理論} \hyperref[theory:information:implication]{第\ref*{theory:information:implication}節}

\par 難易度 2（校準初稿）、優先度未評価。

\section{キャッシュ・CPU}

\subsection{平成28年度 問2-4}
\label{q:h28-reg-basic-2-4}

\noindent\textbf{原典：} sources/H28/basic.pdf、PDF物理10頁。

\textbf{公式正答：} 4（single）。

\paragraph{問題の要旨} ヒット率95\%のキャッシュと主記憶の平均アクセス時間に最近の値を求める。

\Needspace{4\baselineskip}\paragraph{確認した条件・記述（独自要約）}\mbox{}\par

\noindent\textbf{ヒット総時間：} 1.00ns\par

\noindent\textbf{主記憶時間：} 100ns\par

\noindent\textbf{ヒット率：} 0.95\par

\paragraph{問われる核心} キャッシュの条件付き時間と確率重み付き平均

\paragraph{必要知識} キャッシュの条件付き時間と確率重み付き平均。ヒット確率hに対する期待時間は\(E[T]=hT_{\rm hit}+(1-h)T_{\rm miss}\)である。与えられた1ns/100nsをヒット/ミス総時間とする簡略モデルで5.95nsとなる。主記憶100nsを照合後の追加時間とする逐次モデルなら1+0.05×100=\allowbreak{}6nsであり、それでも候補4が最近。どちらの式も定義を付けて使い、実装一般の厳密時間を5.95nsに固定しない。

\paragraph{考え方・正答を選ぶ理由（査読待ち初稿）} ヒット確率hに対する期待時間は\(E[T]=hT_{\rm hit}+(1-h)T_{\rm miss}\)である。与えられた1ns/100nsをヒット/ミス総時間とする簡略モデルで5.95nsとなる。主記憶100nsを照合後の追加時間とする逐次モデルなら1+0.05×100=\allowbreak{}6nsであり、それでも候補4が最近。どちらの式も定義を付けて使い、実装一般の厳密時間を5.95nsに固定しない。

\Needspace{4\baselineskip}\paragraph{選択肢の検討（査読待ち初稿）}\mbox{}\par

\noindent\textbf{1：不適切} 0.05ns。ミス率0.05という無次元数だけであり、アクセス時間になっていない。\par

\noindent\textbf{2：不適切} 0.95ns。ヒット側0.95×1だけでミス側を忘れている。\par

\noindent\textbf{3：不適切} 5.00ns。ミス側0.05×100だけでヒット側を忘れている。\par

\noindent\textbf{4：適切} 5.95ns。ヒット/ミスの総時間を1/100nsとするモデルでは0.95+5=\allowbreak{}5.95ns。\par

\noindent\textbf{5：不適切} 95.0ns。ヒット率を遅い主記憶に掛けており条件付き時間の対応が逆。\par

\paragraph{再利用する考え方} 符号化の目的、情報量と物理的データ長、条件付き時間と平均時間を区別する。キャッシュの条件付き時間と確率重み付き平均の定義と適用条件を確認する。

\paragraph{共通理論} \hyperref[theory:information:cache]{第\ref*{theory:information:cache}節}

\par 参考：\cite{info_cache}。

\paragraph{関連問題} \hyperref[q:r02-reg-basic-2-6]{R02-REG-BASIC-2-6}

\par 難易度 2（校準初稿）、優先度未評価。

\subsection{令和2年度 問2-6}
\label{q:r02-reg-basic-2-6}

\noindent\textbf{原典：} sources/R02/basic.pdf、PDF物理10頁。

\textbf{公式正答：} 5（single）。

\paragraph{問題の要旨} キャッシュ平均時間と主記憶だけに対する速度比を組として求める。

\Needspace{4\baselineskip}\paragraph{確認した条件・記述（独自要約）}\mbox{}\par

\noindent\textbf{キャッシュ：} 50ns\par

\noindent\textbf{主記憶：} 450ns\par

\noindent\textbf{ヒット率：} 0.90\par

\paragraph{問われる核心} キャッシュの平均時間と逆数としての速度比

\paragraph{必要知識} キャッシュの平均時間と逆数としての速度比。ヒット/ミス総時間を50/450nsとする簡略モデルでは平均90ns。同じデータを呼ぶ速さは時間に反比例するので、主記憶単独450nsとの比は450/90=\allowbreak{}5である。ミス時に50ns照合後450nsを追加するモデルなら95ns・約4.74倍となり、それでも組の候補5が最も近い。主記憶時間を総時間か追加ペナルティか先に定義する。

\paragraph{考え方・正答を選ぶ理由（査読待ち初稿）} ヒット/ミス総時間を50/450nsとする簡略モデルでは平均90ns。同じデータを呼ぶ速さは時間に反比例するので、主記憶単独450nsとの比は450/90=\allowbreak{}5である。ミス時に50ns照合後450nsを追加するモデルなら95ns・約4.74倍となり、それでも組の候補5が最も近い。主記憶時間を総時間か追加ペナルティか先に定義する。

\Needspace{4\baselineskip}\paragraph{選択肢の検討（査読待ち初稿）}\mbox{}\par

\noindent\textbf{1：不適切} 45ns/2倍。45nsはヒット側0.9×50だけで、ミス45nsが抜ける。速度比も450/90とは違う。\par

\noindent\textbf{2：不適切} 60ns/2倍。ヒット/ミス重み付き平均90nsとは一致せず、450/60=\allowbreak{}7.5で2倍でもない。\par

\noindent\textbf{3：不適切} 60ns/5倍。60nsは平均に合わず、450/60=\allowbreak{}7.5なので5倍とも整合しない。\par

\noindent\textbf{4：不適切} 90ns/2倍。時間90nsは簡略モデルで合うが速度比は450/90=\allowbreak{}5、2ではない。\par

\noindent\textbf{5：適切} 90ns/5倍。0.9×50+0.1×450=\allowbreak{}90ns、450/90=\allowbreak{}5倍の両方が一致する。\par

\paragraph{再利用する考え方} 符号化の目的、情報量と物理的データ長、条件付き時間と平均時間を区別する。キャッシュの平均時間と逆数としての速度比の定義と適用条件を確認する。

\paragraph{共通理論} \hyperref[theory:information:cache]{第\ref*{theory:information:cache}節}

\par 参考：\cite{info_cache}。

\paragraph{関連問題} \hyperref[q:h28-reg-basic-2-4]{H28-REG-BASIC-2-4}

\par 難易度 2（校準初稿）、優先度未評価。

\section{数値表現}

\subsection{平成23年度 問2-1}
\label{q:h23-reg-basic-2-1}

\noindent\textbf{原典：} sources/H23/basic.pdf、PDF物理5頁。

\textbf{公式正答：} 4（single）。

\paragraph{問題の要旨} 32ビットのIPv4アドレスを上位から8ビットずつ10進表示へ変換する。

\Needspace{4\baselineskip}\paragraph{確認した条件・記述（独自要約）}\mbox{}\par

\noindent\textbf{ビット列：} 11000000101010000001111110101100\par

\noindent\textbf{区切り：} 11000000|10101000|00011111|10101100\par

\paragraph{問われる核心} 32ビットのIPv4アドレスを上位から8ビットずつ10進表示へ変換する。

\paragraph{必要知識} 2進位取りと8ビット塊の順序。

\paragraph{考え方・正答を選ぶ理由（査読待ち初稿）} 11000000=\allowbreak{}192、10101000=\allowbreak{}168、00011111=\allowbreak{}31、10101100=\allowbreak{}172である。原典32ビットをこの順に分割し192.168.31.172、4となる。ビットの順と表示基数を同時に変えない。

\Needspace{4\baselineskip}\paragraph{選択肢の検討（査読待ち初稿）}\mbox{}\par

\noindent\textbf{1：不適切} 172.31.192.168。8ビット塊の前後を入れ替えている。\par

\noindent\textbf{2：不適切} 172.31.168.192。前後の塊と中央の順序が違う。\par

\noindent\textbf{3：不適切} 172.31.255.255。全ビット1の塊はなく255にはならない。\par

\noindent\textbf{4：適切} 192.168.31.172。192、168、31、172を上位から並べている。\par

\noindent\textbf{5：不適切} 192.168.172.31。最後の31と172が逆である。\par

\paragraph{再利用する考え方} 定義した対象量、演算順序、保存条件を固定し、途中状態と候補の違いを比較する。

\paragraph{共通理論} \hyperref[theory:information:radix]{第\ref*{theory:information:radix}節}, \hyperref[theory:discrete:bitcount]{第\ref*{theory:discrete:bitcount}節}

\par 難易度 2（校準初稿）、優先度未評価。

\subsection{平成24年度 問2-1}
\label{q:h24-reg-basic-2-1}

\noindent\textbf{原典：} sources/H24/basic.pdf、PDF物理5頁。

\textbf{公式正答：} 4（single）。

\paragraph{問題の要旨} 数値表現と誤差に関する5記述から、最も適切なものを選ぶ。

\Needspace{4\baselineskip}\paragraph{確認した条件・記述（独自要約）}\mbox{}\par

\noindent\textbf{仮定：} 選択肢5は10進7桁精度で100000.0と0.01を加算する。\par

\paragraph{問われる核心} 有限精度計算で誤差を生む原因

\paragraph{必要知識} 有効桁数、丸め、桁落ち、情報落ち、指数範囲。

\paragraph{考え方・正答を選ぶ理由（査読待ち初稿）} 有限精度は固定・浮動のどちらでも誤差を生じ得る。浮動小数点には有限の指数範囲もある。4の「改善できることが多い」は仮数精度を増す効果として適切である。5は大小の加算で小さい量が失われる例である。

\Needspace{4\baselineskip}\paragraph{選択肢の検討（査読待ち初稿）}\mbox{}\par

\noindent\textbf{1：不適切な解答候補} 浮動小数点では桁あふれ・桁不足はない。指数の範囲は有限であり、上限・下限の範囲超過がある。\par

\noindent\textbf{2：不適切な解答候補} 桁落ちは乗算又は除算で生じる。近い数の減算又は逆符号の加算で有効桁が相殺される現象である。\par

\noindent\textbf{3：不適切な解答候補} 固定小数点なら誤差はない。小数桁数が有限なら表現できない値があり、演算後の丸めもある。\par

\noindent\textbf{4：適切な解答候補} 単精度の丸め誤差は倍精度により改善できることが多い。通常は仮数の桁数が増え、同じ計算の表現・丸め誤差を小さくできる。ただし常に完全に解消するという主張ではない。\par

\noindent\textbf{5：不適切な解答候補} 10進7桁なら100000.0+0.01に誤差はない。正確な100000.01には8桁の有効数字が必要で、7桁では0.01が保持されない。\par

\paragraph{再利用する考え方} 表現・演算の丸め、相殺による桁落ち、大小の加算による情報落ち、範囲超過を区別する。

\paragraph{共通理論} \hyperref[theory:information:error]{第\ref*{theory:information:error}節}

\par 参考：\cite{goldberg1991}。

\par 難易度 2（校準初稿）、優先度未評価。

\subsection{平成25年度 問2-5}
\label{q:h25-reg-basic-2-5}

\noindent\textbf{原典：} sources/H25/basic.pdf、PDF物理9頁。

\textbf{公式正答：} 2（single）。

\paragraph{問題の要旨} 0.85を二進小数4桁で切り捨て、その値を半分にして再度4桁で切り捨てる。

\Needspace{4\baselineskip}\paragraph{確認した条件・記述（独自要約）}\mbox{}\par

\noindent\textbf{初期値：} 10進0.85\par

\noindent\textbf{演算：} 二進小数点以下5桁目以降を各段階で切捨て。初回の格納値を2で除する。\par

\noindent\textbf{空欄順：} 初回の二進表現／除算後の二進表現／除算後の十進値\par

\paragraph{問われる核心} 有限二進小数の量子化と演算順

\paragraph{必要知識} 二進小数の位取り、切捨てによる量子化。

\paragraph{考え方・正答を選ぶ理由（査読待ち初稿）} 初回は \(\lfloor16\times0.85\rfloor/16=13/16=(0.1101)_2\)。これを半分にすると13/32=\allowbreak{}(0.01101)\textsubscript{2}である。4桁に切り捨てると6/16=\allowbreak{}(0.0110)\textsubscript{2}=0.375になる。

\Needspace{4\baselineskip}\paragraph{選択肢の検討（査読待ち初稿）}\mbox{}\par

\noindent\textbf{1：不適切な解答候補} 0.1101／0.0101／0.375。除算後の0110を0101としている。0101の値も0.3125であり提示の0.375に合わない。\par

\noindent\textbf{2：適切な解答候補} 0.1101／0.0110／0.375。13/16を半分にし、6/16へ切り捨てる両段階と一致する。\par

\noindent\textbf{3：不適切な解答候補} 0.1101／0.0110／0.425。0.425は元の0.85を正確に半分にした値であり、段階ごとの切捨てを無視している。\par

\noindent\textbf{4：不適切な解答候補} 0.1010／0.0101／0.375。初回は1101であり1010ではない。除算後の0101も誤りである。\par

\noindent\textbf{5：不適切な解答候補} 0.1010／0.0101／0.425。初回と除算後のビットが誤りで、0.425も切捨て後の値ではない。\par

\paragraph{再利用する考え方} 位取りの基数と、小数の有限表現が可能かどうかを区別する。小数は2倍操作で上位桁から求める。

\paragraph{共通理論} \hyperref[theory:information:radix]{第\ref*{theory:information:radix}節}

\par 難易度 2（校準初稿）、優先度未評価。

\subsection{平成27年度 問2-1}
\label{q:h27-reg-basic-2-1}

\noindent\textbf{原典：} sources/H27/basic.pdf、PDF物理5頁。

\textbf{公式正答：} 5（single）。

\paragraph{問題の要旨} 二進計算機で0.01を100回加え、正しいプログラムでも結果が1と一致しない原因を選ぶ。

\Needspace{4\baselineskip}\paragraph{確認した条件・記述（独自要約）}\mbox{}\par

\noindent\textbf{入力：} 10進0.01\par

\noindent\textbf{繰返し：} 100回加算。プログラムの誤りはない。\par

\paragraph{問われる核心} 進数依存の表現誤差と累積

\paragraph{必要知識} 有限二進小数の条件、入力と演算の丸め。

\paragraph{考え方・正答を選ぶ理由（査読待ち初稿）} 0.01=\allowbreak{}1/100の既約分母には5が残り、二進有限小数にはならない。格納値は近似であり、加算の結果も有限桁へ丸められるので、100回の累積値が数学的な1と一致しないことがある。誤差の符号や実際の値は形式・演算順に依存する。

\Needspace{4\baselineskip}\paragraph{選択肢の検討（査読待ち初稿）}\mbox{}\par

\noindent\textbf{1：不適切な解答候補} 桁落ち。ほぼ等しい量を減算して有効桁が失われる状況ではない。\par

\noindent\textbf{2：不適切な解答候補} 情報落ち。著しく桁の異なる量を加算して小さい量が消える原因を主題としていない。\par

\noindent\textbf{3：不適切な解答候補} オーバーフロー。1付近の合計が通常の浮動小数点の上限を超える状況ではない。\par

\noindent\textbf{4：不適切な解答候補} アンダーフロー。0.01は通常の浮動小数点の最小正数より十分大きい。\par

\noindent\textbf{5：適切な解答候補} 丸め誤差。0.01は二進有限小数で表せず、入力表現と各加算の丸めが蓄積し得る。\par

\paragraph{再利用する考え方} 表現・演算の丸め、相殺による桁落ち、大小の加算による情報落ち、範囲超過を区別する。

\paragraph{共通理論} \hyperref[theory:information:error]{第\ref*{theory:information:error}節}

\par 参考：\cite{goldberg1991}。

\par 難易度 2（校準初稿）、優先度未評価。

\subsection{平成27年度 問2-3}
\label{q:h27-reg-basic-2-3}

\noindent\textbf{原典：} sources/H27/basic.pdf、PDF物理6頁。

\textbf{公式正答：} 3（single）。

\paragraph{問題の要旨} 10進0.1を二進小数で小数点以下16桁まで表す。

\Needspace{4\baselineskip}\paragraph{確認した条件・記述（独自要約）}\mbox{}\par

\noindent\textbf{入力：} 10進0.1\par

\noindent\textbf{表示桁：} 小数点以下16桁\par

\paragraph{問われる核心} 循環二進小数の生成

\paragraph{必要知識} 小数の反復2倍法、分母と有限展開。

\paragraph{考え方・正答を選ぶ理由（査読待ち初稿）} 小数を順に2倍し整数部を取り出す。0.1→0.2→0.4→0.8→1.6で最初の4桁は0001である。残り0.6から1.2、0.4、0.8、1.6と循環し、0011の反復が現れる。16桁は0001100110011001となる。

\Needspace{4\baselineskip}\paragraph{選択肢の検討（査読待ち初稿）}\mbox{}\par

\noindent\textbf{1：不適切な解答候補} 0.0000011001100110。先頭の1が6桁目にあり、0.1の4桁目と合わない。\par

\noindent\textbf{2：不適切な解答候補} 0.0000110011001100。先頭の1が5桁目で、0.1の表現を一桁右へずらした形である。\par

\noindent\textbf{3：適切な解答候補} 0.0001100110011001。反復2倍の出力0001100110011001に一致する。\par

\noindent\textbf{4：不適切な解答候補} 0.0011001100110011。先頭の1が3桁目で、0.1の表現を一桁左へずらした形である。\par

\noindent\textbf{5：不適切な解答候補} 0.0110011001100110。先頭の1が2桁目で、0.1より大きい値である。\par

\paragraph{再利用する考え方} 位取りの基数と、小数の有限表現が可能かどうかを区別する。小数は2倍操作で上位桁から求める。

\paragraph{共通理論} \hyperref[theory:information:radix]{第\ref*{theory:information:radix}節}

\paragraph{関連問題} \hyperref[q:r06-reg-basic-2-1]{R06-REG-BASIC-2-1}

\par 難易度 2（校準初稿）、優先度未評価。

\subsection{平成27年度 問2-4}
\label{q:h27-reg-basic-2-4}

\noindent\textbf{原典：} sources/H27/basic.pdf、PDF物理6頁。

\textbf{公式正答：} 1（single）。

\paragraph{問題の要旨} 10進表記で2 TBの記憶容量を、問題が定義する2進表記のTBへ換算する。

\Needspace{4\baselineskip}\paragraph{確認した条件・記述（独自要約）}\mbox{}\par

\noindent\textbf{容量：} 2×10\textasciicircum{}12 B\par

\noindent\textbf{単位：} 設問は10進1 TB=\allowbreak{}10\textasciicircum{}12 B、2進1 TB=\allowbreak{}2\textasciicircum{}40 Bと定義する。現行の明確な名称では後者は1 TiBである。\par

\paragraph{問われる核心} 情報容量の接頭語と単位換算

\paragraph{必要知識} TBとTiBの単位量、情報量保存。

\paragraph{考え方・正答を選ぶ理由（査読待ち初稿）} 容量をバイトへ統一してから新単位で割る。\(2\times10^{12}/2^{40}=2(1000/1024)^4\simeq1.819\)であり、約1.8となる。原典の「2進のTB」を現在の用語で明確に区別すればTiBである。

\Needspace{4\baselineskip}\paragraph{選択肢の検討（査読待ち初稿）}\mbox{}\par

\noindent\textbf{1：適切な解答候補} 1.8 TB。2×10\textasciicircum{}12/2\textasciicircum{}40≈1.819を小数1桁へ概算した値に一致する。\par

\noindent\textbf{2：不適切な解答候補} 2.0 TB。2.0 TBは2×10\textasciicircum{}12/2\textasciicircum{}40≈1.819の概算と合わない。単位当たりの容量が大きくなるので数値は2より小さくなる。\par

\noindent\textbf{3：不適切な解答候補} 2.1 TB。2.1 TBは2×10\textasciicircum{}12/2\textasciicircum{}40≈1.819の概算と合わない。単位当たりの容量が大きくなるので数値は2より小さくなる。\par

\noindent\textbf{4：不適切な解答候補} 2.2 TB。2.2 TBは2×10\textasciicircum{}12/2\textasciicircum{}40≈1.819の概算と合わない。単位当たりの容量が大きくなるので数値は2より小さくなる。\par

\noindent\textbf{5：不適切な解答候補} 2.3 TB。2.3 TBは2×10\textasciicircum{}12/2\textasciicircum{}40≈1.819の概算と合わない。単位当たりの容量が大きくなるので数値は2より小さくなる。\par

\paragraph{再利用する考え方} 単位あたりのバイト数を先に定義し、容量の数値だけを比べない。

\paragraph{共通理論} \hyperref[theory:information:capacity]{第\ref*{theory:information:capacity}節}

\par 参考：\cite{nist_binary_prefix}。

\paragraph{関連問題} \hyperref[q:r01-ret-basic-2-3]{R01-RET-BASIC-2-3}

\par 難易度 2（校準初稿）、優先度未評価。

\subsection{平成28年度 問2-3}
\label{q:h28-reg-basic-2-3}

\noindent\textbf{原典：} sources/H28/basic.pdf、PDF物理9頁。

\textbf{公式正答：} 1（single）。

\paragraph{問題の要旨} 11を例に、2で割った余りを配列へ格納する二進変換流れ図の分岐と更新を選ぶ。

\Needspace{4\baselineskip}\paragraph{確認した条件・記述（独自要約）}\mbox{}\par

\noindent\textbf{初期化：} I=-1、Nは自然数。\par

\noindent\textbf{操作：} Nと0を比較。継続側でIを1増やしてA[I]へNを2で割った余りを保存し、Nを更新。終了側でA[I]からA[0]を逆順表示。\par

\noindent\textbf{空欄順：} アは継続条件、イは終了条件、ウはNの更新。\par

\paragraph{問われる核心} 整数商と余りによる基数変換

\paragraph{必要知識} 除法の等式N=\allowbreak{}2q+r、0≤r<2、終端条件、配列の順序。

\paragraph{考え方・正答を選ぶ理由（査読待ち初稿）} Nを11→5→2→1→0と整数商で更新し、余りは1、1、0、1の順で得る。格納順は下位桁からなので逆順表示により1011となる。正のNで継続し、N=\allowbreak{}0で停止する。N=\allowbreak{}0を入力として許す実装ではゼロ表示の別処理が必要だが、例示の正の入力の空欄判定には影響しない。

\Needspace{4\baselineskip}\paragraph{選択肢の検討（査読待ち初稿）}\mbox{}\par

\noindent\textbf{1：適切な解答候補} N>0／N=\allowbreak{}0／Nを2で割った商。正のNについて余りを保存し整数商へ更新、商0で終了する。\par

\noindent\textbf{2：不適切な解答候補} N>0／N=\allowbreak{}0／Nを2で割った余り。余りをNへ戻すと桁の取り出しが進まず、余り1では停止しない。\par

\noindent\textbf{3：不適切な解答候補} N=\allowbreak{}0／N>0／Nを2で割った商。初期値11が正の段階で終了し、桁を生成できない。\par

\noindent\textbf{4：不適切な解答候補} N=\allowbreak{}0／N>0／Nを2で割った余り。分岐が逆であり、更新も桁を消去する商ではない。\par

\noindent\textbf{5：不適切な解答候補} N>0／N=\allowbreak{}0／2N。Nを増加させるため正の入力で終了に近づかない。\par

\paragraph{再利用する考え方} 位取りの基数と、小数の有限表現が可能かどうかを区別する。小数は2倍操作で上位桁から求める。

\paragraph{共通理論} \hyperref[theory:information:radix]{第\ref*{theory:information:radix}節}

\par 難易度 2（校準初稿）、優先度未評価。

\subsection{平成29年度 問2-2}
\label{q:h29-reg-basic-2-2}

\noindent\textbf{原典：} sources/H29/basic.pdf、PDF物理7頁。

\textbf{公式正答：} 4（single）。

\paragraph{問題の要旨} 13.0を、符号1ビット・指数8ビット・仮数23ビットの指定浮動小数点形式へ変換する。

\Needspace{4\baselineskip}\paragraph{確認した条件・記述（独自要約）}\mbox{}\par

\noindent\textbf{数値：} 13.0\par

\noindent\textbf{表現：} 正規化値±2\textasciicircum{}a(1+x)、0≤x<1。指数はa+127、仮数はxの二進小数23桁。\par

\noindent\textbf{原典の導入範囲：} 絶対値2\textasciicircum{}-126以上2\textasciicircum{}128未満。IEEE形式全体の表現可能範囲と完全には同一でないが、13.0には影響しない。\par

\noindent\textbf{原典例：} -6.5=-2\textasciicircum{}2(1+0.625)、符号1、指数129。\par

\noindent\textbf{空欄順：} 符号／指数／仮数\par

\paragraph{問われる核心} 正規化と指数バイアス

\paragraph{必要知識} 単精度の3欄、指数バイアス127、正規化先頭1の省略。

\paragraph{考え方・正答を選ぶ理由（査読待ち初稿）} \(13=8(1+0.625)=2^3(1+0.101_2)\)である。正なので符号0、指数欄は3+127=\allowbreak{}130=\allowbreak{}(10000010)\textsubscript{2}。仮数欄は先頭の暗黙の1を保存せず、0.625の小数桁101に0を20個加えた23ビットとなる。

\Needspace{4\baselineskip}\paragraph{選択肢の検討（査読待ち初稿）}\mbox{}\par

\noindent\textbf{1：不適切な解答候補} 1／10000001／10010000000000000000000。13.0は正なので符号は0である。指数3へ127を加えると130で、10000010になる。0.625=\allowbreak{}(0.101)\textsubscript{2}なので仮数部は101の後に0が20個である。\par

\noindent\textbf{2：不適切な解答候補} 1／10000010／10100000000000000000000。13.0は正なので符号は0である。\par

\noindent\textbf{3：不適切な解答候補} 0／10000001／10100000000000000000000。指数3へ127を加えると130で、10000010になる。\par

\noindent\textbf{4：適切な解答候補} 0／10000010／10100000000000000000000。符号・指数・仮数の3欄がすべて一致する。\par

\noindent\textbf{5：不適切な解答候補} 0／10000001／10100000000000000000000。指数3へ127を加えると130で、10000010になる。\par

\paragraph{再利用する考え方} 先頭の1を除く仮数部と、バイアスを加えた指数部を別々に作る。符号を最後まで追跡する。

\paragraph{共通理論} \hyperref[theory:information:floating]{第\ref*{theory:information:floating}節}

\par 参考：\cite{oracle_single}。

\paragraph{関連問題} \hyperref[q:r01-ret-basic-2-4]{R01-RET-BASIC-2-4}

\par 難易度 2（校準初稿）、優先度未評価。

\subsection{平成30年度 問2-3}
\label{q:h30-reg-basic-2-3}

\noindent\textbf{原典：} sources/H30/basic.pdf、PDF物理8頁。

\textbf{公式正答：} 3（single）。

\paragraph{問題の要旨} 01011の5桁を固定して、2の補数と1の補数を求める。

\Needspace{4\baselineskip}\paragraph{確認した条件・記述（独自要約）}\mbox{}\par

\noindent\textbf{対象二進数：} 01011\par

\noindent\textbf{桁数：} 5\par

\noindent\textbf{定義：} 2の補数は2\textasciicircum{}k-X、1の補数は(2\textasciicircum{}k-1)-X。\par

\noindent\textbf{空欄順：} 2の補数／1の補数\par

\paragraph{問われる核心} 桁数を固定した補数と負数表現

\paragraph{必要知識} 基数の補数と基数から1を引いた補数、固定桁幅。

\paragraph{考え方・正答を選ぶ理由（査読待ち初稿）} 01011の十進値は11である。2の補数は32-11=\allowbreak{}21、10101となる。1の補数は31-11=\allowbreak{}20、10100となる。後者は各ビットを反転し、前者はそこへ1を加えても得られる。

\Needspace{4\baselineskip}\paragraph{選択肢の検討（査読待ち初稿）}\mbox{}\par

\noindent\textbf{1：不適切な解答候補} 11011／10100。2の補数は10101であり、11011ではない。\par

\noindent\textbf{2：不適切な解答候補} 10101／11011。1の補数は10100であり、11011ではない。\par

\noindent\textbf{3：適切な解答候補} 10101／10100。2の補数と1の補数がともに一致する。\par

\noindent\textbf{4：不適切な解答候補} 10100／10101。2の補数は10101であり、10100ではない。1の補数は10100であり、10101ではない。\par

\noindent\textbf{5：不適切な解答候補} 11011／11011。2の補数は10101であり、11011ではない。1の補数は10100であり、11011ではない。\par

\paragraph{再利用する考え方} 桁数を固定してから補数を計算する。1の補数はビット反転、2の補数はさらに1を加える。

\paragraph{共通理論} \hyperref[theory:information:complement]{第\ref*{theory:information:complement}節}

\paragraph{関連問題} \hyperref[q:r02-reg-basic-2-4]{R02-REG-BASIC-2-4}

\par 難易度 2（校準初稿）、優先度未評価。

\subsection{令和元年度 問2-1}
\label{q:r01-reg-basic-2-1}

\noindent\textbf{原典：} sources/R01/basic.pdf、PDF物理7頁。

\textbf{公式正答：} 1（single）。

\paragraph{問題の要旨} 10進11.5を二進数と16進数へ変換する。

\Needspace{4\baselineskip}\paragraph{確認した条件・記述（独自要約）}\mbox{}\par

\noindent\textbf{入力：} 11.5\par

\noindent\textbf{空欄順：} 二進数／16進数\par

\paragraph{問われる核心} 二進と16進の位取り対応

\paragraph{必要知識} 基数の位取り、4ビットと16進1桁の対応。

\paragraph{考え方・正答を選ぶ理由（査読待ち初稿）} 11=\allowbreak{}2\textsuperscript{3}+2+1より整数部1011、0.5=\allowbreak{}2\textsuperscript{-1}より小数部.1となる。16進整数11はBで、小数0.5を16倍すると8なのでB.8である。二進小数.1を.1000と4ビットにそろえて8へ変換してもよい。

\Needspace{4\baselineskip}\paragraph{選択肢の検討（査読待ち初稿）}\mbox{}\par

\noindent\textbf{1：適切な解答候補} 1011.1／B.8。11=\allowbreak{}8+2+1、0.5=\allowbreak{}1/2であり、16進では11=\allowbreak{}B、0.5=\allowbreak{}8/16になる。\par

\noindent\textbf{2：不適切な解答候補} 1011.0／C.8。整数12を表すCと、小数0を表す.0が誤りである。\par

\noindent\textbf{3：不適切な解答候補} 1011.1／B.5。16進.5は5/16であり、十進0.5とは異なる。\par

\noindent\textbf{4：不適切な解答候補} 1011.0／B.8。二進.0は小数0であり、0.5を失っている。\par

\noindent\textbf{5：不適切な解答候補} 1011.1／C.5。整数C=\allowbreak{}12、小数.5=\allowbreak{}5/16の両方が異なる。\par

\paragraph{再利用する考え方} 位取りの基数と、小数の有限表現が可能かどうかを区別する。小数は2倍操作で上位桁から求める。

\paragraph{共通理論} \hyperref[theory:information:radix]{第\ref*{theory:information:radix}節}

\par 難易度 2（校準初稿）、優先度未評価。

\subsection{令和元年度再試験 問2-3}
\label{q:r01-ret-basic-2-3}

\noindent\textbf{原典：} sources/R01\_retest/basic.pdf、PDF物理8頁。

\textbf{公式正答：} 2（single）。

\paragraph{問題の要旨} 10進表記で2 TBの記憶容量を、問題が定義する2進表記のTBへ換算する。

\Needspace{4\baselineskip}\paragraph{確認した条件・記述（独自要約）}\mbox{}\par

\noindent\textbf{容量：} 2×10\textasciicircum{}12 B\par

\noindent\textbf{単位：} 設問は10進1 TB=\allowbreak{}10\textasciicircum{}12 B、2進1 TB=\allowbreak{}2\textasciicircum{}40 Bと定義する。現行の明確な名称では後者は1 TiBである。\par

\paragraph{問われる核心} 情報容量の接頭語と単位換算

\paragraph{必要知識} TBとTiBの単位量、情報量保存。

\paragraph{考え方・正答を選ぶ理由（査読待ち初稿）} 容量をバイトへ統一してから新単位で割る。\(2\times10^{12}/2^{40}=2(1000/1024)^4\simeq1.819\)であり、約1.8となる。原典の「2進のTB」を現在の用語で明確に区別すればTiBである。

\Needspace{4\baselineskip}\paragraph{選択肢の検討（査読待ち初稿）}\mbox{}\par

\noindent\textbf{1：不適切な解答候補} 1.6 TB。1.6 TBは2×10\textasciicircum{}12/2\textasciicircum{}40≈1.819の概算と合わない。単位当たりの容量が大きくなるので数値は2より小さくなる。\par

\noindent\textbf{2：適切な解答候補} 1.8 TB。2×10\textasciicircum{}12/2\textasciicircum{}40≈1.819を小数1桁へ概算した値に一致する。\par

\noindent\textbf{3：不適切な解答候補} 2.0 TB。2.0 TBは2×10\textasciicircum{}12/2\textasciicircum{}40≈1.819の概算と合わない。単位当たりの容量が大きくなるので数値は2より小さくなる。\par

\noindent\textbf{4：不適切な解答候補} 2.2 TB。2.2 TBは2×10\textasciicircum{}12/2\textasciicircum{}40≈1.819の概算と合わない。単位当たりの容量が大きくなるので数値は2より小さくなる。\par

\noindent\textbf{5：不適切な解答候補} 2.4 TB。2.4 TBは2×10\textasciicircum{}12/2\textasciicircum{}40≈1.819の概算と合わない。単位当たりの容量が大きくなるので数値は2より小さくなる。\par

\paragraph{再利用する考え方} 単位あたりのバイト数を先に定義し、容量の数値だけを比べない。

\paragraph{共通理論} \hyperref[theory:information:capacity]{第\ref*{theory:information:capacity}節}

\par 参考：\cite{nist_binary_prefix}。

\paragraph{関連問題} \hyperref[q:h27-reg-basic-2-4]{H27-REG-BASIC-2-4}

\par 難易度 2（校準初稿）、優先度未評価。

\subsection{令和元年度再試験 問2-4}
\label{q:r01-ret-basic-2-4}

\noindent\textbf{原典：} sources/R01\_retest/basic.pdf、PDF物理9頁。

\textbf{公式正答：} 5（single）。

\paragraph{問題の要旨} 13.0を、符号1ビット・指数8ビット・仮数23ビットの指定浮動小数点形式へ変換する。

\Needspace{4\baselineskip}\paragraph{確認した条件・記述（独自要約）}\mbox{}\par

\noindent\textbf{数値：} 13.0\par

\noindent\textbf{表現：} 正規化値±2\textasciicircum{}a(1+x)、0≤x<1。指数はa+127、仮数はxの二進小数23桁。\par

\noindent\textbf{原典の導入範囲：} 絶対値2\textasciicircum{}-126以上2\textasciicircum{}128未満。IEEE形式全体の表現可能範囲と完全には同一でないが、13.0には影響しない。\par

\noindent\textbf{原典例：} -6.5=-2\textasciicircum{}2(1+0.625)、符号1、指数129。\par

\noindent\textbf{空欄順：} 符号／指数／仮数\par

\paragraph{問われる核心} 正規化と指数バイアス

\paragraph{必要知識} 単精度の3欄、指数バイアス127、正規化先頭1の省略。

\paragraph{考え方・正答を選ぶ理由（査読待ち初稿）} \(13=8(1+0.625)=2^3(1+0.101_2)\)である。正なので符号0、指数欄は3+127=\allowbreak{}130=\allowbreak{}(10000010)\textsubscript{2}。仮数欄は先頭の暗黙の1を保存せず、0.625の小数桁101に0を20個加えた23ビットとなる。

\Needspace{4\baselineskip}\paragraph{選択肢の検討（査読待ち初稿）}\mbox{}\par

\noindent\textbf{1：不適切な解答候補} 1／10000010／10100000000000000000000。13.0は正なので符号は0である。\par

\noindent\textbf{2：不適切な解答候補} 1／10000001／10100000000000000000000。13.0は正なので符号は0である。指数3へ127を加えると130で、10000010になる。\par

\noindent\textbf{3：不適切な解答候補} 0／10000001／10100000000000000000000。指数3へ127を加えると130で、10000010になる。\par

\noindent\textbf{4：不適切な解答候補} 0／10000001／10100000000000000000000。指数3へ127を加えると130で、10000010になる。\par

\noindent\textbf{5：適切な解答候補} 0／10000010／10100000000000000000000。符号・指数・仮数の3欄がすべて一致する。\par

\paragraph{再利用する考え方} 先頭の1を除く仮数部と、バイアスを加えた指数部を別々に作る。符号を最後まで追跡する。

\paragraph{共通理論} \hyperref[theory:information:floating]{第\ref*{theory:information:floating}節}

\par 参考：\cite{oracle_single}。

\paragraph{関連問題} \hyperref[q:h29-reg-basic-2-2]{H29-REG-BASIC-2-2}

\par 難易度 2（校準初稿）、優先度未評価。

\subsection{令和2年度 問2-4}
\label{q:r02-reg-basic-2-4}

\noindent\textbf{原典：} sources/R02/basic.pdf、PDF物理8頁。

\textbf{公式正答：} 2（single）。

\paragraph{問題の要旨} 100110の6桁を固定して、2の補数と1の補数を求める。

\Needspace{4\baselineskip}\paragraph{確認した条件・記述（独自要約）}\mbox{}\par

\noindent\textbf{対象二進数：} 100110\par

\noindent\textbf{桁数：} 6\par

\noindent\textbf{定義：} 2の補数は2\textasciicircum{}k-X、1の補数は(2\textasciicircum{}k-1)-X。\par

\noindent\textbf{空欄順：} 2の補数／1の補数\par

\paragraph{問われる核心} 桁数を固定した補数と負数表現

\paragraph{必要知識} 基数の補数と基数から1を引いた補数、固定桁幅。

\paragraph{考え方・正答を選ぶ理由（査読待ち初稿）} 100110の十進値は38である。2の補数は64-38=\allowbreak{}26、011010となる。1の補数は63-38=\allowbreak{}25、011001となる。後者は各ビットを反転し、前者はそこへ1を加えても得られる。

\Needspace{4\baselineskip}\paragraph{選択肢の検討（査読待ち初稿）}\mbox{}\par

\noindent\textbf{1：不適切な解答候補} 000110／000101。2の補数は011010であり、000110ではない。1の補数は011001であり、000101ではない。\par

\noindent\textbf{2：適切な解答候補} 011010／011001。2の補数と1の補数がともに一致する。\par

\noindent\textbf{3：不適切な解答候補} 000111／000110。2の補数は011010であり、000111ではない。1の補数は011001であり、000110ではない。\par

\noindent\textbf{4：不適切な解答候補} 011001／011010。2の補数は011010であり、011001ではない。1の補数は011001であり、011010ではない。\par

\noindent\textbf{5：不適切な解答候補} 011000／011001。2の補数は011010であり、011000ではない。\par

\paragraph{再利用する考え方} 桁数を固定してから補数を計算する。1の補数はビット反転、2の補数はさらに1を加える。

\paragraph{共通理論} \hyperref[theory:information:complement]{第\ref*{theory:information:complement}節}

\paragraph{関連問題} \hyperref[q:h30-reg-basic-2-3]{H30-REG-BASIC-2-3}

\par 難易度 2（校準初稿）、優先度未評価。

\subsection{令和2年度 問2-5}
\label{q:r02-reg-basic-2-5}

\noindent\textbf{原典：} sources/R02/basic.pdf、PDF物理9頁。

\textbf{公式正答：} 5（single）。

\paragraph{問題の要旨} 11010101を上位桁から二進→十進変換する流れ図の途中値を求める。

\Needspace{4\baselineskip}\paragraph{確認した条件・記述（独自要約）}\mbox{}\par

\noindent\textbf{入力二進数：} 11010101\par

\noindent\textbf{初期値：} s←最上位ビット、i←次のビット位置。\par

\noindent\textbf{反復：} i≥0ならs←2s+a\_i、i←i-1。i<0で終了。\par

\noindent\textbf{空欄：} 上位4桁までを処理した後から最終更新直前までの値。\par

\paragraph{問われる核心} 桁評価の漸化式とループ不変条件

\paragraph{必要知識} 位置記数法、s←2s+a\_i、上位からの評価、添字の下降。

\paragraph{考え方・正答を選ぶ理由（査読待ち初稿）} 上位から1ビット読むたび、既読の値を2倍して新しいビットを加える。更新列は1→3→6→13→26→53→106→213であり、空欄は26／53／106となる。最後だけでなく全段階の添字とビットを対応させる。

\Needspace{4\baselineskip}\paragraph{選択肢の検討（査読待ち初稿）}\mbox{}\par

\noindent\textbf{1：不適切な解答候補} 25／52／105。アは26であり25ではない。イは53であり52ではない。ウは106であり105ではない。\par

\noindent\textbf{2：不適切な解答候補} 25／52／106。アは26であり25ではない。イは53であり52ではない。\par

\noindent\textbf{3：不適切な解答候補} 26／52／105。イは53であり52ではない。ウは106であり105ではない。\par

\noindent\textbf{4：不適切な解答候補} 26／53／105。ウは106であり105ではない。\par

\noindent\textbf{5：適切な解答候補} 26／53／106。全更新値が漸化式s←2s+a\textsubscript{i}に一致する。\par

\paragraph{再利用する考え方} 各更新後の値が読了したビット列の値であるという不変条件を用い、最後の一致だけで途中を省略しない。

\paragraph{共通理論} \hyperref[theory:information:horner]{第\ref*{theory:information:horner}節}

\paragraph{関連問題} \hyperref[q:r04-reg-basic-2-5]{R04-REG-BASIC-2-5}

\par 難易度 2（校準初稿）、優先度未評価。

\subsection{令和4年度 問2-5}
\label{q:r04-reg-basic-2-5}

\noindent\textbf{原典：} sources/R04/basic.pdf、PDF物理8頁。

\textbf{公式正答：} 5（single）。

\paragraph{問題の要旨} 11001011を上位桁から二進→十進変換する流れ図の途中値を求める。

\Needspace{4\baselineskip}\paragraph{確認した条件・記述（独自要約）}\mbox{}\par

\noindent\textbf{入力二進数：} 11001011\par

\noindent\textbf{初期値：} s←最上位ビット、i←次のビット位置。\par

\noindent\textbf{反復：} i≥0ならs←2s+a\_i、i←i-1。i<0で終了。\par

\noindent\textbf{空欄：} 上位3桁までを処理した後から最終更新直前までの値。\par

\paragraph{問われる核心} 桁評価の漸化式とループ不変条件

\paragraph{必要知識} 位置記数法、s←2s+a\_i、上位からの評価、添字の下降。

\paragraph{考え方・正答を選ぶ理由（査読待ち初稿）} 上位から1ビット読むたび、既読の値を2倍して新しいビットを加える。更新列は1→3→6→12→25→50→101→203であり、空欄は12／25／50／101となる。最後だけでなく全段階の添字とビットを対応させる。

\Needspace{4\baselineskip}\paragraph{選択肢の検討（査読待ち初稿）}\mbox{}\par

\noindent\textbf{1：不適切な解答候補} 12／25／51／102。ウは50であり51ではない。エは101であり102ではない。\par

\noindent\textbf{2：不適切な解答候補} 13／26／50／102。アは12であり13ではない。イは25であり26ではない。エは101であり102ではない。\par

\noindent\textbf{3：不適切な解答候補} 13／26／52／101。アは12であり13ではない。イは25であり26ではない。ウは50であり52ではない。\par

\noindent\textbf{4：不適切な解答候補} 13／25／50／101。アは12であり13ではない。\par

\noindent\textbf{5：適切な解答候補} 12／25／50／101。全更新値が漸化式s←2s+a\textsubscript{i}に一致する。\par

\paragraph{再利用する考え方} 各更新後の値が読了したビット列の値であるという不変条件を用い、最後の一致だけで途中を省略しない。

\paragraph{共通理論} \hyperref[theory:information:horner]{第\ref*{theory:information:horner}節}

\paragraph{関連問題} \hyperref[q:r02-reg-basic-2-5]{R02-REG-BASIC-2-5}

\par 難易度 2（校準初稿）、優先度未評価。

\subsection{令和6年度 問2-1}
\label{q:r06-reg-basic-2-1}

\noindent\textbf{原典：} sources/R06/basic.pdf、PDF物理6頁。

\textbf{公式正答：} 1（single）。

\paragraph{問題の要旨} 10進0.6を二進小数で小数点以下16桁まで表す。

\Needspace{4\baselineskip}\paragraph{確認した条件・記述（独自要約）}\mbox{}\par

\noindent\textbf{入力：} 10進0.6\par

\noindent\textbf{表示桁：} 小数点以下16桁\par

\paragraph{問われる核心} 循環二進小数と表示精度

\paragraph{必要知識} 反復2倍法、剰余の循環。

\paragraph{考え方・正答を選ぶ理由（査読待ち初稿）} 0.6を2倍して整数1、残り0.2から0、残り0.4から0、残り0.8から1、残り0.6へ戻る。したがって1001を反復し、小数16桁は1001100110011001となる。分母5は2の累乗ではないので展開は有限ではない。

\Needspace{4\baselineskip}\paragraph{選択肢の検討（査読待ち初稿）}\mbox{}\par

\noindent\textbf{1：適切な解答候補} 0.1001100110011001。0.6の反復2倍から得る1001100110011001に一致する。\par

\noindent\textbf{2：不適切な解答候補} 0.1011001100110011。最初の3桁を101としているが、0.6では100となる。\par

\noindent\textbf{3：不適切な解答候補} 0.1100000000000000。0.75の有限二進表現であり、0.6と異なる。\par

\noindent\textbf{4：不適切な解答候補} 0.1100110011001100。先頭2桁が11なら0.75以上となり、0.6と合わない。\par

\noindent\textbf{5：不適切な解答候補} 0.1110011001100110。先頭3桁が111なら0.875以上となり、0.6と合わない。\par

\paragraph{再利用する考え方} 位取りの基数と、小数の有限表現が可能かどうかを区別する。小数は2倍操作で上位桁から求める。

\paragraph{共通理論} \hyperref[theory:information:radix]{第\ref*{theory:information:radix}節}

\paragraph{関連問題} \hyperref[q:h27-reg-basic-2-3]{H27-REG-BASIC-2-3}

\par 難易度 2（校準初稿）、優先度未評価。

\section{アルゴリズム・再帰}

\subsection{平成24年度 問2-2}
\label{q:h24-reg-basic-2-2}

\noindent\textbf{原典：} sources/H24/basic.pdf、PDF物理5頁。

\textbf{公式正答：} 2（single）。

\paragraph{問題の要旨} 再帰的な生成規則が表す6文字の二進記号列を選ぶ。

\Needspace{4\baselineskip}\paragraph{確認した条件・記述（独自要約）}\mbox{}\par

\noindent\textbf{規則：} S ::= 01 | 0S1。|は生成規則の選択。\par

\noindent\textbf{対象：} 整数への数値変換ではなく、0と1の列の生成。\par

\paragraph{問われる核心} 再帰規則の不変量と形式言語

\paragraph{必要知識} 基底ケース、再帰で両端に付ける操作、記号の個数・順序。

\paragraph{考え方（査読待ち初稿）} 基底は01である。再帰適用1回で0011、2回で000111となる。生成されるのは0をn個続けた後に1をn個続ける列（n≥1）であり、6文字ならn=\allowbreak{}3なので000111である。

\paragraph{正答を選ぶ理由} 長さ6で生成できる列は0\textsuperscript{3}1\textsuperscript{3}の000111のみである。

\Needspace{4\baselineskip}\paragraph{選択肢の検討（査読待ち初稿）}\mbox{}\par

\noindent\textbf{1：条件に一致しない候補} 000000。規則の基底01に1があり、1を含まない列は生成できない。\par

\noindent\textbf{2：条件に一致する候補} 000111。S→0S1→00S11→000111と生成できる。\par

\noindent\textbf{3：条件に一致しない候補} 001011。0と1は同数だが、1の後に0が出るため規則が作る順序に反する。\par

\noindent\textbf{4：条件に一致しない候補} 010101。01の単なる繰返しは、0で囲んだ内側が再びSになる規則では作れない。\par

\noindent\textbf{5：条件に一致しない候補} 011111。0と1の個数が異なり、両端へ1個ずつ付ける規則の不変量に反する。\par

\paragraph{再利用する考え方} 定義と条件から再帰規則の不変量と形式言語を構成し、成立の証明と誤った候補の反例を区別する。

\paragraph{共通理論} \hyperref[theory:information:grammar]{第\ref*{theory:information:grammar}節}

\par 難易度 2（校準初稿）、優先度未評価。

\subsection{平成26年度 問2-2}
\label{q:h26-reg-basic-2-2}

\noindent\textbf{原典：} sources/H26/basic.pdf、PDF物理6頁。

\textbf{公式正答：} 2（single）。

\paragraph{問題の要旨} ユークリッド互除法の流れ図で、継続・終了条件と出力変数を選ぶ。

\Needspace{4\baselineskip}\paragraph{確認した条件・記述（独自要約）}\mbox{}\par

\noindent\textbf{入力：} A,Bは正の自然数。\par

\noindent\textbf{反復：} R←AをBで割った余り。非零ならA←B、B←Rとして反復。\par

\noindent\textbf{空欄：} アは出力への横分岐、イはA←B、B←Rへの下分岐。\par

\noindent\textbf{候補順：} ア／イ／ウ\par

\paragraph{問われる核心} 最大公約数を保存する反復と停止

\paragraph{必要知識} 除法の等式、公約数の保存、余りの減少、終了時の出力。

\paragraph{考え方（査読待ち初稿）} \(A=QB+R\)より共通の約数は組(A,B)と(B,R)で変わらない。余りが非零なら後者へ更新し、余り0ならBが最大公約数である。終了前にB←0と更新してはいけない。原典矢印上のア・イの位置を個別に読み、候補番号を年度間で流用しない。

\paragraph{正答を選ぶ理由} 余り非零で更新、余り零で終了、出力はBという3条件を満たす。

\Needspace{4\baselineskip}\paragraph{選択肢の検討（査読待ち初稿）}\mbox{}\par

\noindent\textbf{1：条件に一致しない候補} R=\allowbreak{}0／R≠0／A。終了時の最大公約数はBであり、Aではない。AはBの倍数で、Rは0である。\par

\noindent\textbf{2：条件に一致する候補} R=\allowbreak{}0／R≠0／B。非零で更新、零で終了しBを出力する正しい対応である。\par

\noindent\textbf{3：条件に一致しない候補} R=\allowbreak{}0／R≠0／R。終了時の最大公約数はBであり、Rではない。AはBの倍数で、Rは0である。\par

\noindent\textbf{4：条件に一致しない候補} R≠0／R=\allowbreak{}0／A。分岐が逆で、余り0で更新しB=\allowbreak{}0を作るか、非零の余りで途中終了する。終了時の最大公約数はBであり、Aではない。AはBの倍数で、Rは0である。\par

\noindent\textbf{5：条件に一致しない候補} R≠0／R=\allowbreak{}0／B。分岐が逆で、余り0で更新しB=\allowbreak{}0を作るか、非零の余りで途中終了する。\par

\paragraph{再利用する考え方} 定義と条件から最大公約数を保存する反復と停止を構成し、成立の証明と誤った候補の反例を区別する。

\paragraph{共通理論} \hyperref[theory:information:euclid]{第\ref*{theory:information:euclid}節}

\paragraph{関連問題} \hyperref[q:r05-reg-basic-2-2]{R05-REG-BASIC-2-2}

\par 難易度 2（校準初稿）、優先度未評価。

\subsection{令和5年度 問2-2}
\label{q:r05-reg-basic-2-2}

\noindent\textbf{原典：} sources/R05/basic.pdf、PDF物理6頁。

\textbf{公式正答：} 4（single）。

\paragraph{問題の要旨} ユークリッド互除法の流れ図で、継続・終了条件と出力変数を選ぶ。

\Needspace{4\baselineskip}\paragraph{確認した条件・記述（独自要約）}\mbox{}\par

\noindent\textbf{入力：} A,Bは正の自然数。\par

\noindent\textbf{反復：} R←AをBで割った余り。非零ならA←B、B←Rとして反復。\par

\noindent\textbf{空欄：} アは更新への下分岐、イは出力への横分岐。H26とはアとイの位置が逆。\par

\noindent\textbf{候補順：} ア／イ／ウ\par

\paragraph{問われる核心} 最大公約数を保存する反復と停止

\paragraph{必要知識} 除法の等式、公約数の保存、余りの減少、終了時の出力。

\paragraph{考え方（査読待ち初稿）} \(A=QB+R\)より共通の約数は組(A,B)と(B,R)で変わらない。余りが非零なら後者へ更新し、余り0ならBが最大公約数である。終了前にB←0と更新してはいけない。原典矢印上のア・イの位置を個別に読み、候補番号を年度間で流用しない。

\paragraph{正答を選ぶ理由} 余り非零で更新、余り零で終了、出力はBという3条件を満たす。

\Needspace{4\baselineskip}\paragraph{選択肢の検討（査読待ち初稿）}\mbox{}\par

\noindent\textbf{1：条件に一致しない候補} R=\allowbreak{}0／R≠0／A。分岐が逆で、余り0で更新しB=\allowbreak{}0を作るか、非零の余りで途中終了する。終了時の最大公約数はBであり、Aではない。AはBの倍数で、Rは0である。\par

\noindent\textbf{2：条件に一致しない候補} R≠0／R=\allowbreak{}0／A。終了時の最大公約数はBであり、Aではない。AはBの倍数で、Rは0である。\par

\noindent\textbf{3：条件に一致しない候補} R=\allowbreak{}0／R≠0／B。分岐が逆で、余り0で更新しB=\allowbreak{}0を作るか、非零の余りで途中終了する。\par

\noindent\textbf{4：条件に一致する候補} R≠0／R=\allowbreak{}0／B。非零で更新、零で終了しBを出力する正しい対応である。\par

\noindent\textbf{5：条件に一致しない候補} R≠0／R=\allowbreak{}0／R。終了時の最大公約数はBであり、Rではない。AはBの倍数で、Rは0である。\par

\paragraph{再利用する考え方} 定義と条件から最大公約数を保存する反復と停止を構成し、成立の証明と誤った候補の反例を区別する。

\paragraph{共通理論} \hyperref[theory:information:euclid]{第\ref*{theory:information:euclid}節}

\paragraph{関連問題} \hyperref[q:h26-reg-basic-2-2]{H26-REG-BASIC-2-2}

\par 難易度 2（校準初稿）、優先度未評価。

\section{形式言語・木}

\subsection{平成26年度 問2-5}
\label{q:h26-reg-basic-2-5}

\noindent\textbf{原典：} sources/H26/basic.pdf、PDF物理8頁。

\textbf{公式正答：} 5（single）。

\paragraph{問題の要旨} 構文図の迂回と繰返し経路から生成できない文字列を選ぶ。

\Needspace{4\baselineskip}\paragraph{確認した条件・記述（独自要約）}\mbox{}\par

\noindent\textbf{判断対象：} 構文図の順方向・戻り・迂回を反復と空反復へ対応させる。\par

\noindent\textbf{原典確認：} 図・式・全5候補を29保存画像で確認。rawの数式欠落と演算子誤認は別校正JSONで扱う。\par

\paragraph{問われる核心} 構文図の迂回と繰返し経路から生成できない文字列を選ぶ。

\paragraph{必要知識} 構文図の順方向・戻り・迂回を反復と空反復へ対応させる。

\paragraph{考え方・正答を選ぶ理由（査読待ち初稿）} Lを英字、Dを数字と書くと、下側の反復はLの後にDを1回以上通る単位、最後にはLを1回通る。上側で反復を飛ばすこともでき、言語は(LD+)の0回以上反復の後Lである。5は英字同士の連続を含みこの構文で生成できない。

\Needspace{4\baselineskip}\paragraph{選択肢の検討（査読待ち初稿）}\mbox{}\par

\noindent\textbf{1：不適切} a2b3c。a2とb3の二反復の後にcで終わる経路がある。\par

\noindent\textbf{2：不適切} x98y。xの後に数字を二回通りyで終われる。\par

\noindent\textbf{3：不適切} w。上側の迂回で最後の英字だけを通れる。\par

\noindent\textbf{4：不適切} p5q。p5の一反復と最終qで生成できる。\par

\noindent\textbf{5：適切} abc45fg。中間の連続英字ab等を通る経路がなく生成できない。\par

\paragraph{再利用する考え方} 定義した対象量、演算順序、保存条件を固定し、途中状態と候補の違いを比較する。

\paragraph{共通理論} \hyperref[theory:information:grammar]{第\ref*{theory:information:grammar}節}, \hyperref[theory:discrete:syntax]{第\ref*{theory:discrete:syntax}節}

\par 難易度 2（校準初稿）、優先度未評価。

\subsection{令和元年度 問2-4}
\label{q:r01-reg-basic-2-4}

\noindent\textbf{原典：} sources/R01/basic.pdf、PDF物理9頁。

\textbf{公式正答：} 3（single）。

\paragraph{問題の要旨} 形式文法で生成できる数値列を、符号と整数部分の規則から判断する。

\Needspace{4\baselineskip}\paragraph{確認した条件・記述（独自要約）}\mbox{}\par

\noindent\textbf{生成規則：} 数値→整数/小数/整数小数、整数→数字列/符号数字列、小数→点数字列、数字列は非空\par

\paragraph{問われる核心} 生成規則における符号の所属と非空数字列

\paragraph{必要知識} 生成規則における符号の所属と非空数字列。根の数値は整数、小数、整数小数の3通りである。整数は数字列又は符号数字列、小数は点数字列であり、数字列は1桁から再帰的に長くできる。符号は単独で小数に付かないため-.37だけ生成不能で、-0.37なら生成できる。日常の数値表記の許容ではなく、提示規則だけを使う。

\paragraph{考え方・正答を選ぶ理由（査読待ち初稿）} 根の数値は整数、小数、整数小数の3通りである。整数は数字列又は符号数字列、小数は点数字列であり、数字列は1桁から再帰的に長くできる。符号は単独で小数に付かないため-.37だけ生成不能で、-0.37なら生成できる。日常の数値表記の許容ではなく、提示規則だけを使う。

\Needspace{4\baselineskip}\paragraph{選択肢の検討（査読待ち初稿）}\mbox{}\par

\noindent\textbf{1：適切} -19.1。符号-と数字列19が整数を生成し、小数.1と結合できる。\par

\noindent\textbf{2：適切} .52。小数点と非空数字列52で小数だけを生成できる。\par

\noindent\textbf{3：不適切} -.37。符号は整数の生成規則にだけ現れ、整数には必ず数字列が要る。整数桁のない-.37は生成不能。\par

\noindent\textbf{4：適切} 4.35。整数4と小数.35の結合で生成できる。\par

\noindent\textbf{5：適切} -125。符号-と数字列125で整数を生成できる。\par

\paragraph{再利用する考え方} 符号化の目的、情報量と物理的データ長、条件付き時間と平均時間を区別する。生成規則における符号の所属と非空数字列の定義と適用条件を確認する。

\paragraph{共通理論} \hyperref[theory:information:grammar]{第\ref*{theory:information:grammar}節}

\par 難易度 2（校準初稿）、優先度未評価。

\section{暗号の安全性・再送防止・秘匿証明}

\subsection{令和6年度 問2-6}
\label{q:r06-reg-basic-2-6}

\noindent\textbf{原典：} sources/R06/basic.pdf、PDF物理9頁。

\textbf{公式正答：} 4（single）。

\par\textbf{検討未確定：} 肢2の候補の有力性は個別方式・時点に依存する。肢3は秘密鍵非漏洩・真正公開鍵の条件、肢5は値の拘束・範囲・権限等を省き等式を取引正当性全体と括る。公式4・初回4を保持して限定前提を査読対象へ残す。

\paragraph{問題の要旨} ハッシュ・耐量子暗号・署名・再送攻撃・ゼロ知識証明の保証を比較する。

\Needspace{4\baselineskip}\paragraph{確認した条件・記述（独自要約）}\mbox{}\par

\noindent\textbf{方向：} 最も不適切なもの。\par

\noindent\textbf{時点：} R06年度。現在の推奨を出題当時へ無条件に転用しない。\par

\paragraph{問われる核心} 暗号の機能と成立条件・計算量的安全性

\paragraph{必要知識} ハッシュ族、耐衝突性、公開鍵署名、多要素、nonce、再送耐性、非対話ゼロ知識・範囲証明。

\paragraph{考え方・正答を選ぶ理由（査読待ち初稿）} 4は認証要素の複数化とリプレイ耐性を同一視している。NISTは再送防止をnonce等による新鮮さの証明として説明する。固定秘密を繰り返す方式に別要素を加えたというだけで過去メッセージの再受付が不可能になるわけではない。5の秘匿等式は証明可能な関係と取引全体の正当性を分け、2も候補族と個別方式の破綻を分ける。

\Needspace{4\baselineskip}\paragraph{選択肢の検討（査読待ち初稿）}\mbox{}\par

\noindent\textbf{1：標準的定義として適切} 固定出力ハッシュの安全性をパラメータ族の漸近性で定義。。実用の固定nビット出力と、パラメータkで変わる関数族の安全性を区別する。衝突自体は鳩の巣原理で存在し、効率的に見つけにくいという計算量の性質を要求する。\par

\noindent\textbf{2：候補族として理解・時点留意} 格子・符号・多変数多項式等は耐量子暗号の候補。。問題族の研究候補としての説明であり全方式の安全を保証しない。2022年にRainbowには実用攻撃が報告され、2024年の初期標準化は格子・ハッシュ系である。多変数方式全てが破られたとする逆の全称も避ける。\par

\noindent\textbf{3：適切（鍵保護等の前提）} 秘密鍵側だけが署名生成し誰でも検証。。署名者の秘密鍵が適切に保護され、検証者が真正な公開鍵と手続きを用いる前提で成立する。漏洩した秘密鍵での偽署名や鍵の置換を本人の意思の保証としない。\par

\noindent\textbf{4：不適切} 多要素認証ならリプレイ攻撃を無効にする。。要素数とメッセージの新鮮さは別である。使い捨て値の一回受付・nonce/challenge・タイムスタンプ等を正しく検証して再利用を排除する必要がある。\par

\noindent\textbf{5：等式証明は可能・全正当性に条件不足} X/Y/Zを秘密にしてZ=\allowbreak{}Y+Xの関係を証明。。コミットメントと適切な証明体系なら値を明かさず関係成立を示せる。ただし等式だけで取引の所有権・非負残高・承認等を全保証するものではなく範囲・権限などの別条件も必要である。\par

\paragraph{再利用する考え方} 目的量・脅威・成立条件を定めてから手段を選ぶ。秘匿、改ざん検出、主体認証、再送防止、復旧は別の性質であり、方式名だけで全てを保証しない。

\paragraph{共通理論} \hyperref[theory:network:hash]{第\ref*{theory:network:hash}節}, \hyperref[theory:network:keys]{第\ref*{theory:network:keys}節}, \hyperref[theory:network:factors]{第\ref*{theory:network:factors}節}, \hyperref[theory:network:replay]{第\ref*{theory:network:replay}節}, \hyperref[theory:network:zkp]{第\ref*{theory:network:zkp}節}

\par 参考：\cite{net_hash,net_nist2017,net_signature,net_pq2024,net_pq2022,net_zkp2021}。

\paragraph{関連問題} \hyperref[q:r03-reg-basic-2-1]{R03-REG-BASIC-2-1}

\par 難易度 3（校準初稿）、優先度未評価。

\section{データ容量・圧縮率・通信速度の単位換算}

\subsection{平成24年度 問2-4}
\label{q:h24-reg-basic-2-4}

\noindent\textbf{原典：} sources/H24/basic.pdf、PDF物理6頁。

\textbf{公式正答：} 3（single）。

\paragraph{問題の要旨} 2バイト文字の新聞を半分に圧縮したとき、指定容量へ格納できる頁数を求める。

\Needspace{4\baselineskip}\paragraph{確認した条件・記述（独自要約）}\mbox{}\par

\noindent\textbf{文字数/頁：} 10240\par

\noindent\textbf{文字幅：} 2バイト\par

\noindent\textbf{圧縮後/前：} 0.5\par

\noindent\textbf{容量：} 800Mバイト、M=\allowbreak{}1024\textsuperscript{2}\par

\paragraph{問われる核心} 指定単位と圧縮後データ長による容量計算

\paragraph{必要知識} 指定単位と圧縮後データ長による容量計算。圧縮前1頁は10240×2=\allowbreak{}20480バイト、圧縮後は10240バイト。Mは明示的に1024\textsuperscript{2}なので全容量は838860800バイト。割ると81920頁となり3が最近。現行SIのM=\allowbreak{}10\textasciicircum{}6や実際の文字コードの可変幅を、この問題の指定へ持ち込まない。

\paragraph{考え方・正答を選ぶ理由（査読待ち初稿）} 圧縮前1頁は10240×2=\allowbreak{}20480バイト、圧縮後は10240バイト。Mは明示的に1024\textsuperscript{2}なので全容量は838860800バイト。割ると81920頁となり3が最近。現行SIのM=\allowbreak{}10\textasciicircum{}6や実際の文字コードの可変幅を、この問題の指定へ持ち込まない。

\Needspace{4\baselineskip}\paragraph{選択肢の検討（査読待ち初稿）}\mbox{}\par

\noindent\textbf{1：不適切} 約1万頁。約81920頁に対して約8分の1であり、ビットとバイトの8倍を誤って適用した規模である。\par

\noindent\textbf{2：不適切} 約4万頁。圧縮後は1頁10240バイト。圧縮を無視すると40960頁だが指定条件に反する。\par

\noindent\textbf{3：適切} 約8万頁。800×1024\textsuperscript{2}/10240=\allowbreak{}81920頁で、この候補が最近。\par

\noindent\textbf{4：不適切} 約16万頁。81920頁の約2倍で、文字の2バイトと50\%圧縮を二重に消している。\par

\noindent\textbf{5：不適切} 約32万頁。81920頁の約4倍で、与えられた容量/1頁容量の比とは合わない。\par

\paragraph{再利用する考え方} 符号化の目的、情報量と物理的データ長、条件付き時間と平均時間を区別する。指定単位と圧縮後データ長による容量計算の定義と適用条件を確認する。

\paragraph{共通理論} \hyperref[theory:information:transfer]{第\ref*{theory:information:transfer}節}

\par 参考：\cite{info_deflate}。

\par 難易度 2（校準初稿）、優先度未評価。

\subsection{令和3年度 問2-3}
\label{q:r03-reg-basic-2-3}

\noindent\textbf{原典：} sources/R03/basic.pdf、PDF物理7頁。

\textbf{公式正答：} 2（single）。

\paragraph{問題の要旨} 圧縮後データ量と有効通信速度から伝送秒数を求める。

\Needspace{4\baselineskip}\paragraph{確認した条件・記述（独自要約）}\mbox{}\par

\noindent\textbf{元量：} 5GB、G=\allowbreak{}10\textasciicircum{}9\par

\noindent\textbf{圧縮後/前：} 0.5\par

\noindent\textbf{回線速度：} 200Mbps\par

\noindent\textbf{利用率：} 0.70\par

\paragraph{問われる核心} 伝送データ長と有効速度の単位換算

\paragraph{必要知識} 伝送データ長と有効速度の単位換算。圧縮後は2.5×10\textasciicircum{}9バイト=20×10\textasciicircum{}9ビット。有効速度は200×10\textasciicircum{}6×0.70=\allowbreak{}140×10\textasciicircum{}6ビット/秒。比は約142.857秒、最近は143秒の2。G/Mは原典指定の10\textasciicircum{}9/10\textasciicircum{}6であり、1024の累乗ではない。

\paragraph{考え方・正答を選ぶ理由（査読待ち初稿）} 圧縮後は2.5×10\textasciicircum{}9バイト=20×10\textasciicircum{}9ビット。有効速度は200×10\textasciicircum{}6×0.70=\allowbreak{}140×10\textasciicircum{}6ビット/秒。比は約142.857秒、最近は143秒の2。G/Mは原典指定の10\textasciicircum{}9/10\textasciicircum{}6であり、1024の累乗ではない。

\Needspace{4\baselineskip}\paragraph{選択肢の検討（査読待ち初稿）}\mbox{}\par

\noindent\textbf{1：不適切} 286秒。5GBを圧縮せず計算すると約285.7秒になるが半分への圧縮を無視する。\par

\noindent\textbf{2：適切} 143秒。20×10\textasciicircum{}9bit/(140×10\textasciicircum{}6bit/s)=\allowbreak{}142.857秒で最近。\par

\noindent\textbf{3：不適切} 100秒。利用率70\%を忘れた20×10\textasciicircum{}9/(200×10\textasciicircum{}6)=\allowbreak{}100秒。\par

\noindent\textbf{4：不適切} 18秒。圧縮後バイト量をビット量へ8倍しないと約17.86秒になる。\par

\noindent\textbf{5：不適切} 13秒。利用率も8倍換算も忘れると12.5秒でこの規模になる。\par

\paragraph{再利用する考え方} 符号化の目的、情報量と物理的データ長、条件付き時間と平均時間を区別する。伝送データ長と有効速度の単位換算の定義と適用条件を確認する。

\paragraph{共通理論} \hyperref[theory:information:transfer]{第\ref*{theory:information:transfer}節}

\par 参考：\cite{info_deflate}。

\paragraph{関連問題} \hyperref[q:r07-reg-basic-2-5]{R07-REG-BASIC-2-5}

\par 難易度 2（校準初稿）、優先度未評価。

\subsection{令和7年度 問2-5}
\label{q:r07-reg-basic-2-5}

\noindent\textbf{原典：} sources/R07/basic.pdf、PDF物理7頁。

\textbf{公式正答：} 4（single）。

\paragraph{問題の要旨} 既知の圧縮率、通信速度と伝送時間から回線利用率を逆算する。

\Needspace{4\baselineskip}\paragraph{確認した条件・記述（独自要約）}\mbox{}\par

\noindent\textbf{元量：} 250MB、M=\allowbreak{}10\textasciicircum{}6\par

\noindent\textbf{圧縮後/前：} 0.8\par

\noindent\textbf{回線速度：} 100Mbps\par

\noindent\textbf{時間：} 25秒\par

\paragraph{問われる核心} 通信時間式の逆算と圧縮率/利用率の区別

\paragraph{必要知識} 通信時間式の逆算と圧縮率/利用率の区別。圧縮後は200×10\textasciicircum{}6バイト=1.6×10\textasciicircum{}9ビット。\(t=D/(v\eta)\)を\(\eta=D/(vt)\)へ解くと1.6×10\textasciicircum{}9/(100×10\textasciicircum{}6×25)=\allowbreak{}0.64となる。比は無次元で百分率64\%。0.8は圧縮率であり利用率そのものではない。

\paragraph{考え方・正答を選ぶ理由（査読待ち初稿）} 圧縮後は200×10\textasciicircum{}6バイト=1.6×10\textasciicircum{}9ビット。\(t=D/(v\eta)\)を\(\eta=D/(vt)\)へ解くと1.6×10\textasciicircum{}9/(100×10\textasciicircum{}6×25)=\allowbreak{}0.64となる。比は無次元で百分率64\%。0.8は圧縮率であり利用率そのものではない。

\Needspace{4\baselineskip}\paragraph{選択肢の検討（査読待ち初稿）}\mbox{}\par

\noindent\textbf{1：不適切} 8\%。200Mバイトをそのまま100Mbit/sと25秒で割ると0.08となり、バイト/ビット換算が抜ける。\par

\noindent\textbf{2：不適切} 10\%。与えられた圧縮後ビット量1.6Gbitに対し、250Mbitしか25秒で送れない値。\par

\noindent\textbf{3：不適切} 16\%。1.6Gbitの4分の1しか25秒で送れず、時間と整合しない。\par

\noindent\textbf{4：適切} 64\%。250×10\textasciicircum{}6×0.8×8/(100×10\textasciicircum{}6×25)=\allowbreak{}0.64で整合。\par

\noindent\textbf{5：不適切} 80\%。圧縮前データを用いると0.80となるが、実際は0.8倍に圧縮して送る。\par

\paragraph{再利用する考え方} 符号化の目的、情報量と物理的データ長、条件付き時間と平均時間を区別する。通信時間式の逆算と圧縮率/利用率の区別の定義と適用条件を確認する。

\paragraph{共通理論} \hyperref[theory:information:transfer]{第\ref*{theory:information:transfer}節}

\par 参考：\cite{info_deflate}。

\paragraph{関連問題} \hyperref[q:r03-reg-basic-2-3]{R03-REG-BASIC-2-3}

\par 難易度 2（校準初稿）、優先度未評価。

\section{ビット列・積の法則と写像数}

\subsection{平成25年度 問2-1}
\label{q:h25-reg-basic-2-1}

\noindent\textbf{原典：} sources/H25/basic.pdf、PDF物理6頁。

\textbf{公式正答：} 4（single）。

\paragraph{問題の要旨} 二値ランプの個数が4から8へ増えたときの状態数の倍率を求める。

\Needspace{4\baselineskip}\paragraph{確認した条件・記述（独自要約）}\mbox{}\par

\noindent\textbf{判断対象：} 積の法則と二値状態の全組合せ。\par

\noindent\textbf{原典確認：} 図・式・全5候補を29保存画像で確認。rawの数式欠落と演算子誤認は別校正JSONで扱う。\par

\paragraph{問われる核心} 二値ランプの個数が4から8へ増えたときの状態数の倍率を求める。

\paragraph{必要知識} 積の法則と二値状態の全組合せ。

\paragraph{考え方・正答を選ぶ理由（査読待ち初稿）} 独立な二値の選択をn個並べると2のn乗通りである。8個なら256、4個なら16なので256/16=\allowbreak{}16、4である。ランプの種類が複数あっても設問の各位置の状態はON/OFFの二値である。

\Needspace{4\baselineskip}\paragraph{選択肢の検討（査読待ち初稿）}\mbox{}\par

\noindent\textbf{1：不適切} 2倍。ランプ個数の比であって状態数の比ではない。\par

\noindent\textbf{2：不適切} 4倍。2個増加の倍率であり今回は4個増える。\par

\noindent\textbf{3：不適切} 8倍。3個増加の倍率である。\par

\noindent\textbf{4：適切} 16倍。2の8乗を2の4乗で割ると16である。\par

\noindent\textbf{5：不適切} 32倍。5個増加の倍率である。\par

\paragraph{再利用する考え方} 定義した対象量、演算順序、保存条件を固定し、途中状態と候補の違いを比較する。

\paragraph{共通理論} \hyperref[theory:discrete:bitcount]{第\ref*{theory:discrete:bitcount}節}

\paragraph{関連問題} \hyperref[q:r07-reg-basic-2-1]{R07-REG-BASIC-2-1}

\par 難易度 2（校準初稿）、優先度未評価。

\subsection{平成28年度 問2-6}
\label{q:h28-reg-basic-2-6}

\noindent\textbf{原典：} sources/H28/basic.pdf、PDF物理11頁。

\textbf{公式正答：} 5（single）。

\paragraph{問題の要旨} IPv4とIPv6のビット長から表現可能なアドレス総数の倍率を求める。

\Needspace{4\baselineskip}\paragraph{確認した条件・記述（独自要約）}\mbox{}\par

\noindent\textbf{判断対象：} nビットの全列2のn乗と、指数の差による状態数の比。\par

\noindent\textbf{原典確認：} 図・式・全5候補を29保存画像で確認。rawの数式欠落と演算子誤認は別校正JSONで扱う。\par

\paragraph{問われる核心} IPv4とIPv6のビット長から表現可能なアドレス総数の倍率を求める。

\paragraph{必要知識} nビットの全列2のn乗と、指数の差による状態数の比。

\paragraph{考え方・正答を選ぶ理由（査読待ち初稿）} IPv4は4×8=\allowbreak{}32ビット、IPv6は8×16=\allowbreak{}128ビット。全ビット列数の比は2の128乗/2の32乗=2の96乗である。予約・用途による割当の割合でなく表現可能な全列を数える。

\Needspace{4\baselineskip}\paragraph{選択肢の検討（査読待ち初稿）}\mbox{}\par

\noindent\textbf{1：不適切} 2の8乗倍。IPv6とIPv4のビット長の差は96であり、この指数ではない。\par

\noindent\textbf{2：不適切} 2の16乗倍。IPv6とIPv4のビット長の差は96であり、この指数ではない。\par

\noindent\textbf{3：不適切} 2の32乗倍。IPv6とIPv4のビット長の差は96であり、この指数ではない。\par

\noindent\textbf{4：不適切} 2の64乗倍。IPv6とIPv4のビット長の差は96であり、この指数ではない。\par

\noindent\textbf{5：適切} 2の96乗倍。128−32=\allowbreak{}96ビットの差に対応する。\par

\paragraph{再利用する考え方} 定義した対象量、演算順序、保存条件を固定し、途中状態と候補の違いを比較する。

\paragraph{共通理論} \hyperref[theory:discrete:bitcount]{第\ref*{theory:discrete:bitcount}節}

\par 参考：\cite{rfc4291_addresses}。

\paragraph{関連問題} \hyperref[q:r04-reg-basic-2-6]{R04-REG-BASIC-2-6}

\par 難易度 2（校準初稿）、優先度未評価。

\subsection{令和4年度 問2-6}
\label{q:r04-reg-basic-2-6}

\noindent\textbf{原典：} sources/R04/basic.pdf、PDF物理9頁。

\textbf{公式正答：} 4（single）。

\paragraph{問題の要旨} IPv4とIPv6のビット長から表現可能なアドレス総数の倍率を求める。

\Needspace{4\baselineskip}\paragraph{確認した条件・記述（独自要約）}\mbox{}\par

\noindent\textbf{判断対象：} nビットの全列2のn乗と、指数の差による状態数の比。\par

\noindent\textbf{原典確認：} 図・式・全5候補を29保存画像で確認。rawの数式欠落と演算子誤認は別校正JSONで扱う。\par

\paragraph{問われる核心} IPv4とIPv6のビット長から表現可能なアドレス総数の倍率を求める。

\paragraph{必要知識} nビットの全列2のn乗と、指数の差による状態数の比。

\paragraph{考え方・正答を選ぶ理由（査読待ち初稿）} IPv4は4×8=\allowbreak{}32ビット、IPv6は8×16=\allowbreak{}128ビット。全ビット列数の比は2の128乗/2の32乗=2の96乗である。予約・用途による割当の割合でなく表現可能な全列を数える。

\Needspace{4\baselineskip}\paragraph{選択肢の検討（査読待ち初稿）}\mbox{}\par

\noindent\textbf{1：不適切} 2の4乗倍。IPv6とIPv4のビット長の差は96であり、この指数ではない。\par

\noindent\textbf{2：不適切} 2の16乗倍。IPv6とIPv4のビット長の差は96であり、この指数ではない。\par

\noindent\textbf{3：不適切} 2の32乗倍。IPv6とIPv4のビット長の差は96であり、この指数ではない。\par

\noindent\textbf{4：適切} 2の96乗倍。128−32=\allowbreak{}96ビットの差に対応する。\par

\noindent\textbf{5：不適切} 2の128乗倍。IPv6とIPv4のビット長の差は96であり、この指数ではない。\par

\paragraph{再利用する考え方} 定義した対象量、演算順序、保存条件を固定し、途中状態と候補の違いを比較する。

\paragraph{共通理論} \hyperref[theory:discrete:bitcount]{第\ref*{theory:discrete:bitcount}節}

\par 参考：\cite{rfc4291_addresses}。

\paragraph{関連問題} \hyperref[q:h28-reg-basic-2-6]{H28-REG-BASIC-2-6}

\par 難易度 2（校準初稿）、優先度未評価。

\subsection{令和7年度 問2-1}
\label{q:r07-reg-basic-2-1}

\noindent\textbf{原典：} sources/R07/basic.pdf、PDF物理5頁。

\textbf{公式正答：} 3（single）。

\paragraph{問題の要旨} 直積集合を定義域とする全写像の数を積の法則で求める。

\Needspace{4\baselineskip}\paragraph{確認した条件・記述（独自要約）}\mbox{}\par

\noindent\textbf{判断対象：} 直積の要素数と有限集合間の全写像数。\par

\noindent\textbf{原典確認：} 図・式・全5候補を29保存画像で確認。rawの数式欠落と演算子誤認は別校正JSONで扱う。\par

\paragraph{問われる核心} 直積集合を定義域とする全写像の数を積の法則で求める。

\paragraph{必要知識} 直積の要素数と有限集合間の全写像数。

\paragraph{考え方・正答を選ぶ理由（査読待ち初稿）} |A×B|=3×2=\allowbreak{}6である。各順序対の像をCの2要素から選ぶので、全写像は2の6乗=64、3である。単射・全射という条件はなく同じ像を複数の要素へ割り当ててよい。

\Needspace{4\baselineskip}\paragraph{選択肢の検討（査読待ち初稿）}\mbox{}\par

\noindent\textbf{1：不適切} 12。6×2は定義域と像の組の総数であり写像数ではない。\par

\noindent\textbf{2：不適切} 36。6の2乗は方向を逆にした数え方である。\par

\noindent\textbf{3：適切} 64。6要素の各像を独立に2択から選ぶ2の6乗である。\par

\noindent\textbf{4：不適切} 81。3の4乗は今回の定義域と値域の要素数に合わない。\par

\noindent\textbf{5：不適切} 128。2の7乗であり定義域要素を一つ多く数えている。\par

\paragraph{再利用する考え方} 定義した対象量、演算順序、保存条件を固定し、途中状態と候補の違いを比較する。

\paragraph{共通理論} \hyperref[theory:discrete:mappings]{第\ref*{theory:discrete:mappings}節}

\paragraph{関連問題} \hyperref[q:h25-reg-basic-2-1]{H25-REG-BASIC-2-1}

\par 難易度 2（校準初稿）、優先度未評価。

\section{文書特徴ベクトルのコサイン類似度}

\subsection{令和元年度 問2-3}
\label{q:r01-reg-basic-2-3}

\noindent\textbf{原典：} sources/R01/basic.pdf、PDF物理9頁。

\textbf{公式正答：} 3（single）。

\paragraph{問題の要旨} 単語出現回数のベクトルを長さで正規化し、文書Aへ最も近い文書を選ぶ。

\Needspace{4\baselineskip}\paragraph{確認した条件・記述（独自要約）}\mbox{}\par

\noindent\textbf{定義：} \(d(X,Y)=1-\frac{X\cdot Y}{|X||Y|}\)\par

\noindent\textbf{全18セル：} A(7,3,2)、B(2,3,0)、C(70,3,2)、D(21,9,6)、E(1,2,3)、F(7,30,20)。\par

\noindent\textbf{対象：} 3単語だけの頻度ベクトル。原典は非零。\par

\paragraph{問われる核心} 特徴ベクトルの方向とコサイン距離

\paragraph{必要知識} 内積、ユークリッドノルム、正の比例、Cauchy-Schwarz、ゼロベクトルでは式未定義。

\paragraph{考え方・正答を選ぶ理由（査読待ち初稿）} 距離は\(1-\frac{A\cdot X}{|A||X|}\)である。D=\allowbreak{}3Aなので内積186、ノルム積\(\sqrt{62}\sqrt{558}=186\)、距離0。他は正比例でなく距離が正となる。単語総数や成分差の大きさだけを比較しない。

\Needspace{4\baselineskip}\paragraph{選択肢の検討（査読待ち初稿）}\mbox{}\par

\noindent\textbf{1：不一致} B=\allowbreak{}(2,3,0)。原典の定義による距離は0.189860。非零のAと同じ向きではなく、Dの距離0より大きい。\par

\noindent\textbf{2：不一致} C=\allowbreak{}(70,3,2)。原典の定義による距離は0.088622。非零のAと同じ向きではなく、Dの距離0より大きい。\par

\noindent\textbf{3：適切} D=\allowbreak{}(21,9,6)。原典の定義による距離は0.000000。D=\allowbreak{}3Aで正の比例となり、内積/ノルム積が1になる。\par

\noindent\textbf{4：不一致} E=\allowbreak{}(1,2,3)。原典の定義による距離は0.355098。非零のAと同じ向きではなく、Dの距離0より大きい。\par

\noindent\textbf{5：不一致} F=\allowbreak{}(7,30,20)。原典の定義による距離は0.381056。非零のAと同じ向きではなく、Dの距離0より大きい。\par

\paragraph{再利用する考え方} 目的量・脅威・成立条件を定めてから手段を選ぶ。秘匿、改ざん検出、主体認証、再送防止、復旧は別の性質であり、方式名だけで全てを保証しない。

\paragraph{共通理論} \hyperref[theory:network:cosine]{第\ref*{theory:network:cosine}節}

\par 参考：\cite{net_cosine}。

\par 難易度 1（校準初稿）、優先度未評価。

\section{算術木・演算優先と後置記法}

\subsection{平成27年度 問2-6}
\label{q:h27-reg-basic-2-6}

\noindent\textbf{原典：} sources/H27/basic.pdf、PDF物理7頁。

\textbf{公式正答：} 4（single）。

\paragraph{問題の要旨} 左右の部分木と根の演算から算術木を中置式へ読み戻す。

\Needspace{4\baselineskip}\paragraph{確認した条件・記述（独自要約）}\mbox{}\par

\noindent\textbf{判断対象：} 算術木の左右順と演算子の優先・結合。\par

\noindent\textbf{原典確認：} 図・式・全5候補を29保存画像で確認。rawの数式欠落と演算子誤認は別校正JSONで扱う。\par

\paragraph{問われる核心} 左右の部分木と根の演算から算術木を中置式へ読み戻す。

\paragraph{必要知識} 算術木の左右順と演算子の優先・結合。

\paragraph{考え方・正答を選ぶ理由（査読待ち初稿）} 根は減算、左の子は乗算3×2、右の子は減算5−1である。よって3×2−(5−1)、4。葉だけを左から並べると括弧で指定する構造を失う。

\Needspace{4\baselineskip}\paragraph{選択肢の検討（査読待ち初稿）}\mbox{}\par

\noindent\textbf{1：不適切} (3−2)×(5−1)。根を掛算とする式で原典の減算根と違う。\par

\noindent\textbf{2：不適切} 5−1−3×2。左と右の部分式が逆である。\par

\noindent\textbf{3：不適切} 3×2−5−1。減算の左結合で(3×2−5)−1となり右部分木が違う。\par

\noindent\textbf{4：適切} 3×2−(5−1)。左部分木3×2から右部分木5−1を引く構造に合う。\par

\noindent\textbf{5：不適切} 3×2−5×1。右部分木は乗算でなく減算である。\par

\paragraph{再利用する考え方} 定義した対象量、演算順序、保存条件を固定し、途中状態と候補の違いを比較する。

\paragraph{共通理論} \hyperref[theory:discrete:expression]{第\ref*{theory:discrete:expression}節}

\paragraph{関連問題} \hyperref[q:r03-reg-basic-2-5]{R03-REG-BASIC-2-5}

\par 難易度 2（校準初稿）、優先度未評価。

\subsection{平成30年度 問2-5}
\label{q:h30-reg-basic-2-5}

\noindent\textbf{原典：} sources/H30/basic.pdf、PDF物理9頁。

\textbf{公式正答：} 1（single）。

\paragraph{問題の要旨} 乗除を加算より先に行う中置式を後置記法へ変換する。

\Needspace{4\baselineskip}\paragraph{確認した条件・記述（独自要約）}\mbox{}\par

\noindent\textbf{判断対象：} 算術木の左・右・根の後行順走査と演算優先。\par

\noindent\textbf{原典確認：} 図・式・全5候補を29保存画像で確認。rawの数式欠落と演算子誤認は別校正JSONで扱う。\par

\paragraph{問われる核心} 乗除を加算より先に行う中置式を後置記法へ変換する。

\paragraph{必要知識} 算術木の左・右・根の後行順走査と演算優先。

\paragraph{考え方・正答を選ぶ理由（査読待ち初稿）} 対象はa×b+c÷dである。左部分式a×bはab×、右部分式c÷dはcd÷となる。両者を加える根を最後に置けばab×cd÷+、1である。

\Needspace{4\baselineskip}\paragraph{選択肢の検討（査読待ち初稿）}\mbox{}\par

\noindent\textbf{1：適切} ab×cd÷+。左右の部分式ab×とcd÷を終えて根の+を出す。\par

\noindent\textbf{2：不適切} ab×c÷d+。(a×b)÷c+dとなり除算の対象が違う。\par

\noindent\textbf{3：不適切} abc×÷d+。a÷(b×c)+dとなり括りが違う。\par

\noindent\textbf{4：不適切} abc+d÷×。a×((b+c)÷d)となり加算対象が違う。\par

\noindent\textbf{5：不適切} abcd×÷+。a+b÷(c×d)となり根以下の構造が違う。\par

\paragraph{再利用する考え方} 定義した対象量、演算順序、保存条件を固定し、途中状態と候補の違いを比較する。

\paragraph{共通理論} \hyperref[theory:discrete:expression]{第\ref*{theory:discrete:expression}節}

\paragraph{関連問題} \hyperref[q:r03-reg-basic-2-5]{R03-REG-BASIC-2-5}

\par 難易度 2（校準初稿）、優先度未評価。

\subsection{令和3年度 問2-5}
\label{q:r03-reg-basic-2-5}

\noindent\textbf{原典：} sources/R03/basic.pdf、PDF物理9頁。

\textbf{公式正答：} 1（single）。

\paragraph{問題の要旨} 除算の優先を保って括弧内の中置式を逆ポーランド記法へ変換する。

\Needspace{4\baselineskip}\paragraph{確認した条件・記述（独自要約）}\mbox{}\par

\noindent\textbf{原典式：} (A+B÷C)×(D−F)\par

\noindent\textbf{OCR差：} 生OCRの(A+B−C)は原典に一致しない。÷を別校正へ記録。\par

\paragraph{問われる核心} 除算の優先を保って括弧内の中置式を逆ポーランド記法へ変換する。

\paragraph{必要知識} 演算子の優先順位と後行順走査。

\paragraph{考え方・正答を選ぶ理由（査読待ち初稿）} 原典は(A+B÷C)×(D−F)である。左部分木はAとB÷Cを加えるのでABC÷+、右部分木はDF−。根の乗算を最後にしてABC÷+DF−×、1。OCRの減算らしい記号を採用せず原典の÷を保持する。

\Needspace{4\baselineskip}\paragraph{選択肢の検討（査読待ち初稿）}\mbox{}\par

\noindent\textbf{1：適切} ABC÷+DF−×。B÷Cの後にAを足し、D−Fの後に掛ける順に合う。\par

\noindent\textbf{2：不適切} AB+C÷DF−×。(A+B)÷Cへ括りを変えている。\par

\noindent\textbf{3：不適切} ABC÷+D×F−。左式×Dの後でFを引いてしまう。\par

\noindent\textbf{4：不適切} ×+A÷BC−DF。先頭の演算子は後置ではない。\par

\noindent\textbf{5：不適切} AB+C÷D×F−。左の除算対象と最後の減算の括りが違う。\par

\paragraph{再利用する考え方} 定義した対象量、演算順序、保存条件を固定し、途中状態と候補の違いを比較する。

\paragraph{共通理論} \hyperref[theory:discrete:expression]{第\ref*{theory:discrete:expression}節}

\paragraph{関連問題} \hyperref[q:h27-reg-basic-2-6]{H27-REG-BASIC-2-6}, \hyperref[q:h30-reg-basic-2-5]{H30-REG-BASIC-2-5}

\par 難易度 2（校準初稿）、優先度未評価。

\section{拡張互除法の係数と行列積}

\subsection{令和元年度再試験 問2-2}
\label{q:r01-ret-basic-2-2}

\noindent\textbf{原典：} sources/R01\_retest/basic.pdf、PDF物理7頁。

\textbf{公式正答：} 4（single）。

\paragraph{問題の要旨} 拡張ユークリッド互除法の行列積から最大公約数と指定係数を求める。

\Needspace{4\baselineskip}\paragraph{確認した条件・記述（独自要約）}\mbox{}\par

\noindent\textbf{判断対象：} 割り算の余り、(0,1;1,−q)の行変換、Bézout係数。\par

\noindent\textbf{原典確認：} 図・式・全5候補を29保存画像で確認。rawの数式欠落と演算子誤認は別校正JSONで扱う。\par

\paragraph{問われる核心} 拡張ユークリッド互除法の行列積から最大公約数と指定係数を求める。

\paragraph{必要知識} 割り算の余り、(0,1;1,−q)の行変換、Bézout係数。

\paragraph{考え方・正答を選ぶ理由（査読待ち初稿）} 余り51,6,3,0からgcdは3。商1,1,8の変換行列を順に左から積むと(−1,2;9,−17)となる。指定第二行は9,−17で108×9−57×17=\allowbreak{}3、4である。5も線形和3を満たす別のBézout解だが、指定行列の第二行ではない。

\Needspace{4\baselineskip}\paragraph{選択肢の検討（査読待ち初稿）}\mbox{}\par

\noindent\textbf{1：不適切} (6, -1, 2)。指定する係数行又は最大公約数が違う。線形和は6となるが、gとの一致だけで指定行列の行とは限らない。\par

\noindent\textbf{2：不適切} (6, 1, -2)。指定する係数行又は最大公約数が違う。線形和は-6となるが、gとの一致だけで指定行列の行とは限らない。\par

\noindent\textbf{3：不適切} (6, 1, 2)。指定する係数行又は最大公約数が違う。線形和は222となるが、gとの一致だけで指定行列の行とは限らない。\par

\noindent\textbf{4：適切} (3, 9, -17)。指定行列の第二行と最大公約数が全て一致する。\par

\noindent\textbf{5：不適切} (3, -10, 19)。指定する係数行又は最大公約数が違う。線形和は3となるが、gとの一致だけで指定行列の行とは限らない。\par

\paragraph{再利用する考え方} 定義した対象量、演算順序、保存条件を固定し、途中状態と候補の違いを比較する。

\paragraph{共通理論} \hyperref[theory:information:euclid]{第\ref*{theory:information:euclid}節}, \hyperref[theory:discrete:extended]{第\ref*{theory:discrete:extended}節}

\par 参考：\cite{mit_discrete2015}。

\paragraph{関連問題} \hyperref[q:r06-reg-basic-2-5]{R06-REG-BASIC-2-5}

\par 難易度 3（校準初稿）、優先度未評価。

\subsection{令和6年度 問2-5}
\label{q:r06-reg-basic-2-5}

\noindent\textbf{原典：} sources/R06/basic.pdf、PDF物理8頁。

\textbf{公式正答：} 4（single）。

\paragraph{問題の要旨} 拡張ユークリッド互除法の行列積から最大公約数と指定係数を求める。

\Needspace{4\baselineskip}\paragraph{確認した条件・記述（独自要約）}\mbox{}\par

\noindent\textbf{判断対象：} 割り算の余り、(0,1;1,−q)の行変換、Bézout係数。\par

\noindent\textbf{原典確認：} 図・式・全5候補を29保存画像で確認。rawの数式欠落と演算子誤認は別校正JSONで扱う。\par

\paragraph{問われる核心} 拡張ユークリッド互除法の行列積から最大公約数と指定係数を求める。

\paragraph{必要知識} 割り算の余り、(0,1;1,−q)の行変換、Bézout係数。

\paragraph{考え方・正答を選ぶ理由（査読待ち初稿）} 余り39,26,13,0からgcdは13。商1,1,1の変換T=\allowbreak{}(0,1;1,−1)に対しTの3乗=(−1,2;2,−3)。指定第二行2,−3は104×2−65×3=\allowbreak{}13で4である。5の−3,5も線形和13だが指定行ではない。

\Needspace{4\baselineskip}\paragraph{選択肢の検討（査読待ち初稿）}\mbox{}\par

\noindent\textbf{1：不適切} (5, 2, -3)。指定する係数行又は最大公約数が違う。線形和は13となるが、gとの一致だけで指定行列の行とは限らない。\par

\noindent\textbf{2：不適切} (5, -3, 5)。指定する係数行又は最大公約数が違う。線形和は13となるが、gとの一致だけで指定行列の行とは限らない。\par

\noindent\textbf{3：不適切} (8, 3, -3)。指定する係数行又は最大公約数が違う。線形和は117となるが、gとの一致だけで指定行列の行とは限らない。\par

\noindent\textbf{4：適切} (13, 2, -3)。指定行列の第二行と最大公約数が全て一致する。\par

\noindent\textbf{5：不適切} (13, -3, 5)。指定する係数行又は最大公約数が違う。線形和は13となるが、gとの一致だけで指定行列の行とは限らない。\par

\paragraph{再利用する考え方} 定義した対象量、演算順序、保存条件を固定し、途中状態と候補の違いを比較する。

\paragraph{共通理論} \hyperref[theory:information:euclid]{第\ref*{theory:information:euclid}節}, \hyperref[theory:discrete:extended]{第\ref*{theory:discrete:extended}節}

\par 参考：\cite{mit_discrete2015}。

\paragraph{関連問題} \hyperref[q:r01-ret-basic-2-2]{R01-RET-BASIC-2-2}

\par 難易度 3（校準初稿）、優先度未評価。

\section{二段キャッシュの実効アクセス時間}

\subsection{平成23年度 問2-4}
\label{q:h23-reg-basic-2-4}

\noindent\textbf{原典：} sources/H23/basic.pdf、PDF物理6頁。

\textbf{公式正答：} 1（single）。

\par\textbf{検討未確定：} 原典はアクセス時間に上位段の探索時間を含めるか明示しない。候補に合う経路総時間モデルと逐次加算モデルを区別し、査読を要する。

\paragraph{問題の要旨} 二段キャッシュの実効アクセス時間。

\Needspace{4\baselineskip}\paragraph{確認した条件・記述（独自要約）}\mbox{}\par

\noindent\textbf{一次：} 1ns、95\%。\par

\noindent\textbf{二次：} 10ns、一次失敗条件下90\%。\par

\noindent\textbf{主記憶：} 100ns。\par

\paragraph{問われる核心} 二段キャッシュの実効アクセス時間。

\paragraph{必要知識} 条件付き確率、排他的経路、ヒット率、累積時間のモデル。

\paragraph{考え方・正答を選ぶ理由（査読待ち初稿）} 一次成功、一次失敗かつ二次成功、両方失敗を排他的に分ける。二次の90\%は一次失敗時の条件付きヒット率と読む。0.95×1+0.05×0.9×10+0.05×0.1×100=\allowbreak{}1.9 nsの1。二次や主記憶の時間を各経路の総所要時間として扱う試験モデル。各段を逐次探索し、その時間を別々に加えるモデルなら1+0.05×(10+0.1×100)=\allowbreak{}2 nsとなる。

\Needspace{4\baselineskip}\paragraph{選択肢の検討（査読待ち初稿）}\mbox{}\par

\noindent\textbf{1：条件内で適切} 1.9。解法の条件と結果を照合する。 計算結果と一致。\par

\noindent\textbf{2：不適切} 6.45。解法の条件と結果を照合する。 計算結果に一致しない。\par

\noindent\textbf{3：不適切} 11.45。解法の条件と結果を照合する。 計算結果に一致しない。\par

\noindent\textbf{4：不適切} 15.4。解法の条件と結果を照合する。 計算結果に一致しない。\par

\noindent\textbf{5：不適切} 19.95。解法の条件と結果を照合する。 計算結果に一致しない。\par

\paragraph{再利用する考え方} 条件付き確率、排他的経路、ヒット率、累積時間のモデル。 条件を式又は判定にし、全候補を照合してその年度の番号へ戻す。

\paragraph{共通理論} \hyperref[theory:information:cache]{第\ref*{theory:information:cache}節}

\par 参考：\cite{info_cache}。

\par 難易度 3（校準初稿）、優先度未評価。

\section{命令数とCPIからCPU実行時間を求める}

\subsection{平成25年度 問2-6}
\label{q:h25-reg-basic-2-6}

\noindent\textbf{原典：} sources/H25/basic.pdf、PDF物理9頁。

\textbf{公式正答：} 3（single）。

\paragraph{問題の要旨} 命令数とCPIからCPU実行時間を求める。

\Needspace{4\baselineskip}\paragraph{確認した条件・記述（独自要約）}\mbox{}\par

\noindent\textbf{命令：} 転送3500×CPI6、算術5000×5、分岐1500×4。\par

\noindent\textbf{周波数：} 2.0GHz。\par

\paragraph{問われる核心} 命令数とCPIからCPU実行時間を求める。

\paragraph{必要知識} CPI、クロック周期、命令構成による重み付き平均、SI接頭辞。

\paragraph{考え方・正答を選ぶ理由（査読待ち初稿）} 必要サイクル数は3500×6+5000×5+1500×4=\allowbreak{}52000。2.0GHzは2.0×10\textasciicircum{}9サイクル/秒なので時間は52000/(2.0×10\textasciicircum{}9)=\allowbreak{}26×10\textasciicircum{}-6秒=26μs。命令数10000だけを周波数で割らず、CPIで重み付けする。

\Needspace{4\baselineskip}\paragraph{選択肢の検討（査読待ち初稿）}\mbox{}\par

\noindent\textbf{1：不適切} 26ns。各文又は定義の判定を先に行う。並びが判定 26μs と異なる。\par

\noindent\textbf{2：不適切} 25μs。各文又は定義の判定を先に行う。並びが判定 26μs と異なる。\par

\noindent\textbf{3：条件内で適切} 26μs。各文又は定義の判定を先に行う。全項目が一致する。\par

\noindent\textbf{4：不適切} 25ms。各文又は定義の判定を先に行う。並びが判定 26μs と異なる。\par

\noindent\textbf{5：不適切} 26ms。各文又は定義の判定を先に行う。並びが判定 26μs と異なる。\par

\paragraph{再利用する考え方} CPI、クロック周期、命令構成による重み付き平均、SI接頭辞。 条件を式又は判定にし、全候補を照合してその年度の番号へ戻す。

\paragraph{共通理論} \hyperref[theory:information:cpu]{第\ref*{theory:information:cpu}節}

\par 難易度 2（校準初稿）、優先度未評価。

\subsection{平成29年度 問2-6}
\label{q:h29-reg-basic-2-6}

\noindent\textbf{原典：} sources/H29/basic.pdf、PDF物理10頁。

\textbf{公式正答：} 2（single）。

\paragraph{問題の要旨} 命令数とCPIからCPU実行時間を求める。

\Needspace{4\baselineskip}\paragraph{確認した条件・記述（独自要約）}\mbox{}\par

\noindent\textbf{命令：} 転送3500×CPI6、算術5000×5、分岐1500×4。\par

\noindent\textbf{周波数：} 2.0GHz。\par

\paragraph{問われる核心} 命令数とCPIからCPU実行時間を求める。

\paragraph{必要知識} CPI、クロック周期、命令構成による重み付き平均、SI接頭辞。

\paragraph{考え方・正答を選ぶ理由（査読待ち初稿）} 必要サイクル数は3500×6+5000×5+1500×4=\allowbreak{}52000。2.0GHzは2.0×10\textasciicircum{}9サイクル/秒なので時間は52000/(2.0×10\textasciicircum{}9)=\allowbreak{}26×10\textasciicircum{}-6秒=26μs。命令数10000だけを周波数で割らず、CPIで重み付けする。

\Needspace{4\baselineskip}\paragraph{選択肢の検討（査読待ち初稿）}\mbox{}\par

\noindent\textbf{1：不適切} 260ns。各文又は定義の判定を先に行う。並びが判定 26μs と異なる。\par

\noindent\textbf{2：条件内で適切} 26μs。各文又は定義の判定を先に行う。全項目が一致する。\par

\noindent\textbf{3：不適切} 260μs。各文又は定義の判定を先に行う。並びが判定 26μs と異なる。\par

\noindent\textbf{4：不適切} 26ms。各文又は定義の判定を先に行う。並びが判定 26μs と異なる。\par

\noindent\textbf{5：不適切} 260ms。各文又は定義の判定を先に行う。並びが判定 26μs と異なる。\par

\paragraph{再利用する考え方} CPI、クロック周期、命令構成による重み付き平均、SI接頭辞。 条件を式又は判定にし、全候補を照合してその年度の番号へ戻す。

\paragraph{共通理論} \hyperref[theory:information:cpu]{第\ref*{theory:information:cpu}節}

\par 難易度 2（校準初稿）、優先度未評価。

\section{再帰文法の生成列の不変条件}

\subsection{平成29年度 問2-5}
\label{q:h29-reg-basic-2-5}

\noindent\textbf{原典：} sources/H29/basic.pdf、PDF物理10頁。

\textbf{公式正答：} 4（single）。

\paragraph{問題の要旨} 再帰文法の生成列の不変条件。

\Needspace{4\baselineskip}\paragraph{確認した条件・記述（独自要約）}\mbox{}\par

\noindent\textbf{規則：} S::=01 | 0 S 1。\par

\paragraph{問われる核心} 再帰文法の生成列の不変条件。

\paragraph{必要知識} BNF、基底と再帰、生成列の構造、不変条件。

\paragraph{考え方・正答を選ぶ理由（査読待ち初稿）} 基底は01。再帰規則は生成列の外へ0と1を一つずつ付け、0011、000111となる。一般形0\textasciicircum{}n1\textasciicircum{}n(n≥1)なので4のみ生成可能。数が等しいだけでなく全0が全1の前に並ぶ必要がある。

\Needspace{4\baselineskip}\paragraph{選択肢の検討（査読待ち初稿）}\mbox{}\par

\noindent\textbf{1：不適切} 111110。各文又は定義の判定を先に行う。並びが判定 000111 と異なる。\par

\noindent\textbf{2：不適切} 111000。各文又は定義の判定を先に行う。並びが判定 000111 と異なる。\par

\noindent\textbf{3：不適切} 101010。各文又は定義の判定を先に行う。並びが判定 000111 と異なる。\par

\noindent\textbf{4：条件内で適切} 000111。各文又は定義の判定を先に行う。全項目が一致する。\par

\noindent\textbf{5：不適切} 000001。各文又は定義の判定を先に行う。並びが判定 000111 と異なる。\par

\paragraph{再利用する考え方} BNF、基底と再帰、生成列の構造、不変条件。 条件を式又は判定にし、全候補を照合してその年度の番号へ戻す。

\paragraph{共通理論} \hyperref[theory:information:grammar]{第\ref*{theory:information:grammar}節}

\par 難易度 2（校準初稿）、優先度未評価。

\section{ハミング距離と一ビット誤り訂正}

\subsection{令和元年度 問2-5}
\label{q:r01-reg-basic-2-5}

\noindent\textbf{原典：} sources/R01/basic.pdf、PDF物理10頁。

\textbf{公式正答：} 5（single）。

\paragraph{問題の要旨} ハミング距離と一ビット誤り訂正。

\Needspace{4\baselineskip}\paragraph{確認した条件・記述（独自要約）}\mbox{}\par

\noindent\textbf{比較：} 1110001 / 0001110\par

\noindent\textbf{受信：} 1001010\par

\noindent\textbf{検査式：} x5=\allowbreak{}x2+x3+x4、x6=\allowbreak{}x1+x3+x4、x7=\allowbreak{}x1+x2+x4、全てmod2。\par

\noindent\textbf{仮定：} 通信誤りは高々一ビット。\par

\paragraph{問われる核心} ハミング距離と一ビット誤り訂正。

\paragraph{必要知識} XOR、ハミング距離、mod2の検査式、最小距離と訂正能力。

\paragraph{考え方・正答を選ぶ理由（査読待ち初稿）} 1110001と0001110の各位置を比較すると距離7。受信1001010について無反転と七通りの一ビット反転を調べ、三つの検査ビットが全て再計算値に合う1011010を選ぶ。R01では第3位置を反転し、情報1011のパリティは010。 最小距離3なので、誤りが高々1という条件下で訂正が一意になる。

\Needspace{4\baselineskip}\paragraph{選択肢の検討（査読待ち初稿）}\mbox{}\par

\noindent\textbf{1：不適切} 5/1001010。各文又は定義の判定を先に行う。並びが判定 7/1011010 と異なる。\par

\noindent\textbf{2：不適切} 5/0001010。各文又は定義の判定を先に行う。並びが判定 7/1011010 と異なる。\par

\noindent\textbf{3：不適切} 5/1101010。各文又は定義の判定を先に行う。並びが判定 7/1011010 と異なる。\par

\noindent\textbf{4：不適切} 7/1001010。各文又は定義の判定を先に行う。並びが判定 7/1011010 と異なる。\par

\noindent\textbf{5：条件内で適切} 7/1011010。各文又は定義の判定を先に行う。全項目が一致する。\par

\paragraph{再利用する考え方} XOR、ハミング距離、mod2の検査式、最小距離と訂正能力。 条件を式又は判定にし、全候補を照合してその年度の番号へ戻す。

\paragraph{共通理論} \hyperref[theory:information:hamming]{第\ref*{theory:information:hamming}節}

\par 難易度 2（校準初稿）、優先度未評価。

\subsection{令和4年度 問2-4}
\label{q:r04-reg-basic-2-4}

\noindent\textbf{原典：} sources/R04/basic.pdf、PDF物理7頁。

\textbf{公式正答：} 5（single）。

\paragraph{問題の要旨} ハミング距離と一ビット誤り訂正。

\Needspace{4\baselineskip}\paragraph{確認した条件・記述（独自要約）}\mbox{}\par

\noindent\textbf{比較：} 1110101 / 1001111\par

\noindent\textbf{受信：} 1000010\par

\noindent\textbf{検査式：} x5=\allowbreak{}x2+x3+x4、x6=\allowbreak{}x1+x3+x4、x7=\allowbreak{}x1+x2+x4、全てmod2。\par

\noindent\textbf{仮定：} 通信誤りは高々一ビット。\par

\paragraph{問われる核心} ハミング距離と一ビット誤り訂正。

\paragraph{必要知識} XOR、ハミング距離、mod2の検査式、最小距離と訂正能力。

\paragraph{考え方・正答を選ぶ理由（査読待ち初稿）} 1110101と1001111の各位置を比較すると距離4。受信1000010について無反転と七通りの一ビット反転を調べ、三つの検査ビットが全て再計算値に合う1000011を選ぶ。R04では第7位置を反転し、情報1000のパリティは011。 最小距離3なので、誤りが高々1という条件下で訂正が一意になる。

\Needspace{4\baselineskip}\paragraph{選択肢の検討（査読待ち初稿）}\mbox{}\par

\noindent\textbf{1：不適切} 4/0000010。各文又は定義の判定を先に行う。並びが判定 4/1000011 と異なる。\par

\noindent\textbf{2：不適切} 5/1100010。各文又は定義の判定を先に行う。並びが判定 4/1000011 と異なる。\par

\noindent\textbf{3：不適切} 4/1001010。各文又は定義の判定を先に行う。並びが判定 4/1000011 と異なる。\par

\noindent\textbf{4：不適切} 5/1000110。各文又は定義の判定を先に行う。並びが判定 4/1000011 と異なる。\par

\noindent\textbf{5：条件内で適切} 4/1000011。各文又は定義の判定を先に行う。全項目が一致する。\par

\paragraph{再利用する考え方} XOR、ハミング距離、mod2の検査式、最小距離と訂正能力。 条件を式又は判定にし、全候補を照合してその年度の番号へ戻す。

\paragraph{共通理論} \hyperref[theory:information:hamming]{第\ref*{theory:information:hamming}節}

\par 難易度 2（校準初稿）、優先度未評価。

\section{LRUで参照履歴を更新する}

\subsection{令和4年度 問2-3}
\label{q:r04-reg-basic-2-3}

\noindent\textbf{原典：} sources/R04/basic.pdf、PDF物理6頁。

\textbf{公式正答：} 3（single）。

\paragraph{問題の要旨} LRUで参照履歴を更新する。

\Needspace{4\baselineskip}\paragraph{確認した条件・記述（独自要約）}\mbox{}\par

\noindent\textbf{容量：} 3ページ、初期空。\par

\noindent\textbf{参照：} 2,1,1,2,3,4,1,3,4。\par

\noindent\textbf{時間：} ヒットH、外部M。原典例はミスに追加Hを計上しない。\par

\paragraph{問われる核心} LRUで参照履歴を更新する。

\paragraph{必要知識} LRU、ページフォールト、ヒット時更新、アクセス時間の原典モデル。

\paragraph{考え方・正答を選ぶ理由（査読待ち初稿）} 古い→新しいの順で保持。2(M):[2]、1(M):[2,1]、1(H):[2,1]、2(H):[1,2]、3(M):[1,2,3]、4(M):1を追出し[2,3,4]、1(M):2を追出し[3,4,1]、3(H):[4,1,3]、4(H):[1,3,4]。外部5回・主記憶4回で5M+4Hの3。挿入の古さを使うFIFOとは違いヒットでも順序を更新する。

\Needspace{4\baselineskip}\paragraph{選択肢の検討（査読待ち初稿）}\mbox{}\par

\noindent\textbf{1：不適切} 7M+2H。各文又は定義の判定を先に行う。並びが判定 5M+4H と異なる。\par

\noindent\textbf{2：不適切} 6M+3H。各文又は定義の判定を先に行う。並びが判定 5M+4H と異なる。\par

\noindent\textbf{3：条件内で適切} 5M+4H。各文又は定義の判定を先に行う。全項目が一致する。\par

\noindent\textbf{4：不適切} 4M+5H。各文又は定義の判定を先に行う。並びが判定 5M+4H と異なる。\par

\noindent\textbf{5：不適切} 3M+6H。各文又は定義の判定を先に行う。並びが判定 5M+4H と異なる。\par

\paragraph{再利用する考え方} LRU、ページフォールト、ヒット時更新、アクセス時間の原典モデル。 条件を式又は判定にし、全候補を照合してその年度の番号へ戻す。

\paragraph{共通理論} \hyperref[theory:information:lru]{第\ref*{theory:information:lru}節}

\par 難易度 2（校準初稿）、優先度未評価。

\section{エレベータの状態遷移のガード条件}

\subsection{平成30年度 問2-2}
\label{q:h30-reg-basic-2-2}

\noindent\textbf{原典：} sources/H30/basic.pdf、PDF物理7頁。

\textbf{公式正答：} 3（single）。

\paragraph{問題の要旨} エレベータの状態遷移のガード条件。

\Needspace{4\baselineskip}\paragraph{確認した条件・記述（独自要約）}\mbox{}\par

\noindent\textbf{状態：} 停止・上昇・下降のみ。扉・故障・別階呼出を除外。\par

\noindent\textbf{a：} 上昇自己遷移。\par

\noindent\textbf{b：} 上昇→停止。\par

\noindent\textbf{c：} 下降→停止。\par

\noindent\textbf{d：} 下降自己遷移。\par

\noindent\textbf{e：} 停止→下降。\par

\paragraph{問われる核心} エレベータの状態遷移のガード条件。

\paragraph{必要知識} 状態、イベント、ガード、状態不変条件。

\paragraph{考え方・正答を選ぶ理由（査読待ち初稿）} 上昇を継続するaは現在階<目的階、停止へ戻るb,cはともに=。下降を続けるdと停止から下降を始めるeは現在階>目的階。a,b,c,d,e=<,=,=,>,>の3。矢印の始点・終点と比較を合わせて読む。

\Needspace{4\baselineskip}\paragraph{選択肢の検討（査読待ち初稿）}\mbox{}\par

\noindent\textbf{1：不適切} = = = = =。各文又は定義の判定を先に行う。並びが判定 < = = > > と異なる。\par

\noindent\textbf{2：不適切} = > < = =。各文又は定義の判定を先に行う。並びが判定 < = = > > と異なる。\par

\noindent\textbf{3：条件内で適切} < = = > >。各文又は定義の判定を先に行う。全項目が一致する。\par

\noindent\textbf{4：不適切} = < > = =。各文又は定義の判定を先に行う。並びが判定 < = = > > と異なる。\par

\noindent\textbf{5：不適切} > = = < >。各文又は定義の判定を先に行う。並びが判定 < = = > > と異なる。\par

\paragraph{再利用する考え方} 状態、イベント、ガード、状態不変条件。 条件を式又は判定にし、全候補を照合してその年度の番号へ戻す。

\paragraph{共通理論} \hyperref[theory:information:state]{第\ref*{theory:information:state}節}

\par 難易度 2（校準初稿）、優先度未評価。

\section{情報理論・符号・誤り検出}

\subsection{平成26年度 問2-3}
\label{q:h26-reg-basic-2-3}

\noindent\textbf{原典：} sources/H26/basic.pdf、PDF物理7頁。

\textbf{公式正答：} 4（single）。

\paragraph{問題の要旨} ISBN-10の重み付き剰余条件を保つ数字の変更を選ぶ。

\Needspace{4\baselineskip}\paragraph{確認した条件・記述（独自要約）}\mbox{}\par

\noindent\textbf{方向：} 正しい番号\par

\paragraph{問われる核心} 重み付き検査数字と剰余による誤り検出

\paragraph{必要知識} 重み付き検査数字と剰余による誤り検出。先頭9桁の和171に検査数字5を加えると176=\allowbreak{}16×11。変更後の全和を毎回計算するより、変更差が11の倍数か判定するとよい。候補差は45,15,1,11,10で、4だけ有効。これは検査数字規則に適合するという意味であり、番号が実際に出版物へ割り当て済みという保証ではない。

\paragraph{考え方・正答を選ぶ理由（査読待ち初稿）} 先頭9桁の和171に検査数字5を加えると176=\allowbreak{}16×11。変更後の全和を毎回計算するより、変更差が11の倍数か判定するとよい。候補差は45,15,1,11,10で、4だけ有効。これは検査数字規則に適合するという意味であり、番号が実際に出版物へ割り当て済みという保証ではない。

\Needspace{4\baselineskip}\paragraph{選択肢の検討（査読待ち初稿）}\mbox{}\par

\noindent\textbf{1：不適切} 第2桁0→5：95-70002-34-5。第2桁の重み9、差5で和が45増える。元和176に足した221は11の倍数でない。\par

\noindent\textbf{2：不適切} 第8桁3→8：90-70002-84-5。重み3の桁が5増えるので和が15増える。191は11の倍数でない。OCRの04は画像の84へ校正する。\par

\noindent\textbf{3：不適切} 検査数字5→6：90-70002-34-6。検査桁の重み1で和が177となり、余り1なので無効。\par

\noindent\textbf{4：適切} 第8/9桁34→82：90-70002-82-5。重み3と2の変化は3×5+2×(-2)=\allowbreak{}11。全和187=\allowbreak{}17×11となり有効。\par

\noindent\textbf{5：不適切} 第8/9桁34→56：90-70002-56-5。変化は3×2+2×2=\allowbreak{}10。186は余り10で無効。\par

\paragraph{再利用する考え方} 符号化の目的、情報量と物理的データ長、条件付き時間と平均時間を区別する。重み付き検査数字と剰余による誤り検出の定義と適用条件を確認する。

\paragraph{共通理論} \hyperref[theory:information:checksum]{第\ref*{theory:information:checksum}節}

\par 参考：\cite{info_errors}。

\paragraph{関連問題} \hyperref[q:r05-reg-basic-2-3]{R05-REG-BASIC-2-3}

\par 難易度 2（校準初稿）、優先度未評価。

\subsection{平成28年度 問2-5}
\label{q:h28-reg-basic-2-5}

\noindent\textbf{原典：} sources/H28/basic.pdf、PDF物理10頁。

\textbf{公式正答：} 4（single）。

\paragraph{問題の要旨} 偶数パリティ1ビットの付加による検出能力と、位置特定の可否を判断する。

\Needspace{4\baselineskip}\paragraph{確認した条件・記述（独自要約）}\mbox{}\par

\noindent\textbf{元データ例：} 1010100\par

\noindent\textbf{付加ビット：} 1の個数偶数なら0、奇数なら1\par

\noindent\textbf{送信列例：} 10101001\par

\paragraph{問われる核心} 偶数パリティの検出条件と訂正能力の区別

\paragraph{必要知識} 偶数パリティの検出条件と訂正能力の区別。パリティを\(b=\sum_ix_i\pmod2\)と置くと、データとbの全和は偶数になる。誤りベクトルeを足した受信列の検査値は\(\sum_ie_i\pmod2\)であり、奇数個反転だけ検出できる。異なる2ビット反転は検出不能で、単一検査値から反転位置も決まらない。元1010100の1は3個、付加1で10101001の1が4個になることも画像確認した。

\paragraph{考え方・正答を選ぶ理由（査読待ち初稿）} パリティを\(b=\sum_ix_i\pmod2\)と置くと、データとbの全和は偶数になる。誤りベクトルeを足した受信列の検査値は\(\sum_ie_i\pmod2\)であり、奇数個反転だけ検出できる。異なる2ビット反転は検出不能で、単一検査値から反転位置も決まらない。元1010100の1は3個、付加1で10101001の1が4個になることも画像確認した。

\Needspace{4\baselineskip}\paragraph{選択肢の検討（査読待ち初稿）}\mbox{}\par

\noindent\textbf{1：不適切} 1ビット反転検出には元データ8ビット以下が必要。。長さに関係なく1回の反転で全体の1の個数の偶奇が変わる。8ビットという上限はない。\par

\noindent\textbf{2：不適切} 1ビット反転検出には元データ2ビット以上が必要。。1ビットのデータにも付加でき、反転は検出できる。2ビットという下限はない。\par

\noindent\textbf{3：不適切} 8ビットデータの1ビット反転位置が特定できる。。受信全体の偶奇だけでは、データ中の各位置やパリティ位置のどれが変わったか区別できない。\par

\noindent\textbf{4：適切} 2ビット反転では誤りを検出できない。。異なる2ビットの反転は偶奇を2回反転するため元の偶数に戻る。検査値0でも元データ一致の保証はない。\par

\noindent\textbf{5：不適切} データによって付加ビットを決められない。。任意の有限ビット列で1の個数は偶数か奇数のいずれかなので、一意に付加値が決まる。\par

\paragraph{再利用する考え方} 符号化の目的、情報量と物理的データ長、条件付き時間と平均時間を区別する。偶数パリティの検出条件と訂正能力の区別の定義と適用条件を確認する。

\paragraph{共通理論} \hyperref[theory:information:parity]{第\ref*{theory:information:parity}節}

\par 参考：\cite{info_errors}。

\paragraph{関連問題} \hyperref[q:r06-reg-basic-2-4]{R06-REG-BASIC-2-4}

\par 難易度 2（校準初稿）、優先度未評価。

\subsection{令和5年度 問2-3}
\label{q:r05-reg-basic-2-3}

\noindent\textbf{原典：} sources/R05/basic.pdf、PDF物理7頁。

\textbf{公式正答：} 5（single）。

\paragraph{問題の要旨} ISBN-13の交互重み1/3の合同式から欠けた末尾検査数字を求める。

\Needspace{4\baselineskip}\paragraph{確認した条件・記述（独自要約）}\mbox{}\par

\noindent\textbf{番号：} 978-4-103-34194-X\par

\noindent\textbf{左から重み：} 1,3交互\par

\noindent\textbf{合同式：} 全和mod10=\allowbreak{}0\par

\paragraph{問われる核心} ISBN-13の検査数字と交互重みの合同式

\paragraph{必要知識} ISBN-13の検査数字と交互重みの合同式。原典の番号は978-4-103-34194-Xで、Xは末尾桁である。左から1/3交互の重みを掛けると9+21+8+12+1+0+3+9+4+3+9+12=\allowbreak{}101。\(101+X\equiv0\pmod{10}\)を満たす0〜9は9であり5が適切。ISBN-10の重み10〜1、法11とは異なる規則である。

\paragraph{考え方・正答を選ぶ理由（査読待ち初稿）} 原典の番号は978-4-103-34194-Xで、Xは末尾桁である。左から1/3交互の重みを掛けると9+21+8+12+1+0+3+9+4+3+9+12=\allowbreak{}101。\(101+X\equiv0\pmod{10}\)を満たす0〜9は9であり5が適切。ISBN-10の重み10〜1、法11とは異なる規則である。

\Needspace{4\baselineskip}\paragraph{選択肢の検討（査読待ち初稿）}\mbox{}\par

\noindent\textbf{1：不適切} 1。101+1=\allowbreak{}102はmod10で2。\par

\noindent\textbf{2：不適切} 3。101+3=\allowbreak{}104はmod10で4。\par

\noindent\textbf{3：不適切} 5。101+5=\allowbreak{}106はmod10で6。\par

\noindent\textbf{4：不適切} 7。101+7=\allowbreak{}108はmod10で8。\par

\noindent\textbf{5：適切} 9。101+9=\allowbreak{}110はmod10で0となる。\par

\paragraph{再利用する考え方} 符号化の目的、情報量と物理的データ長、条件付き時間と平均時間を区別する。ISBN-13の検査数字と交互重みの合同式の定義と適用条件を確認する。

\paragraph{共通理論} \hyperref[theory:information:checksum]{第\ref*{theory:information:checksum}節}

\par 参考：\cite{info_isbn}。

\paragraph{関連問題} \hyperref[q:h26-reg-basic-2-3]{H26-REG-BASIC-2-3}

\par 難易度 2（校準初稿）、優先度未評価。

\subsection{令和6年度 問2-4}
\label{q:r06-reg-basic-2-4}

\noindent\textbf{原典：} sources/R06/basic.pdf、PDF物理7頁。

\textbf{公式正答：} 5（single）。

\paragraph{問題の要旨} 偶数パリティ1ビットの付加による検出能力と、位置特定の可否を判断する。

\Needspace{4\baselineskip}\paragraph{確認した条件・記述（独自要約）}\mbox{}\par

\noindent\textbf{元データ例：} 1010100\par

\noindent\textbf{付加ビット：} 1の個数偶数なら0、奇数なら1\par

\noindent\textbf{送信列例：} 10101001\par

\paragraph{問われる核心} 偶数パリティの検出条件と訂正能力の区別

\paragraph{必要知識} 偶数パリティの検出条件と訂正能力の区別。パリティを\(b=\sum_ix_i\pmod2\)と置くと、データとbの全和は偶数になる。誤りベクトルeを足した受信列の検査値は\(\sum_ie_i\pmod2\)であり、奇数個反転だけ検出できる。異なる2ビット反転は検出不能で、単一検査値から反転位置も決まらない。元1010100の1は3個、付加1で10101001の1が4個になることも画像確認した。

\paragraph{考え方・正答を選ぶ理由（査読待ち初稿）} パリティを\(b=\sum_ix_i\pmod2\)と置くと、データとbの全和は偶数になる。誤りベクトルeを足した受信列の検査値は\(\sum_ie_i\pmod2\)であり、奇数個反転だけ検出できる。異なる2ビット反転は検出不能で、単一検査値から反転位置も決まらない。元1010100の1は3個、付加1で10101001の1が4個になることも画像確認した。

\Needspace{4\baselineskip}\paragraph{選択肢の検討（査読待ち初稿）}\mbox{}\par

\noindent\textbf{1：不適切} 1ビット反転検出には元データ8ビット以下が必要。。長さに関係なく1回の反転で全体の1の個数の偶奇が変わる。8ビットという上限はない。\par

\noindent\textbf{2：不適切} 1ビット反転検出には元データ2ビット以上が必要。。1ビットのデータにも付加でき、反転は検出できる。2ビットという下限はない。\par

\noindent\textbf{3：不適切} 8ビットデータの1ビット反転位置が特定できる。。受信全体の偶奇だけでは、データ中の各位置やパリティ位置のどれが変わったか区別できない。\par

\noindent\textbf{4：不適切} データによって付加ビットを決められない。。任意の有限ビット列で1の個数は偶数か奇数のいずれかなので、一意に付加値が決まる。\par

\noindent\textbf{5：適切} 2ビット反転では誤りを検出できない。。異なる2ビットの反転は偶奇を2回反転するため元の偶数に戻る。検査値0でも元データ一致の保証はない。\par

\paragraph{再利用する考え方} 符号化の目的、情報量と物理的データ長、条件付き時間と平均時間を区別する。偶数パリティの検出条件と訂正能力の区別の定義と適用条件を確認する。

\paragraph{共通理論} \hyperref[theory:information:parity]{第\ref*{theory:information:parity}節}

\par 参考：\cite{info_errors}。

\paragraph{関連問題} \hyperref[q:h28-reg-basic-2-5]{H28-REG-BASIC-2-5}

\par 難易度 2（校準初稿）、優先度未評価。

\section{有限マルコフ連鎖の遷移と定常分布}

\subsection{平成24年度 問2-3}
\label{q:h24-reg-basic-2-3}

\noindent\textbf{原典：} sources/H24/basic.pdf、PDF物理6頁。

\textbf{公式正答：} 4（single）。

\paragraph{問題の要旨} 3状態の天気遷移図から2日後の条件付き確率と長期極限を判定する。

\Needspace{4\baselineskip}\paragraph{確認した条件・記述（独自要約）}\mbox{}\par

\noindent\textbf{方向：} 最も不適切\par

\paragraph{問われる核心} 有限マルコフ連鎖の経路和・行列累乗・定常分布

\paragraph{必要知識} 有限マルコフ連鎖の経路和・行列累乗・定常分布。状態を晴・曇・雨と並べると遷移行列は対角1/2、非対角1/4である。2日後は\((P^2)_{ij}=\sum_kP_{ik}P_{kj}\)から自己3/8、異なる状態5/16。定常分布は(1/3,1/3,1/3)であり、全要素正のため初期天気によらず収束する。肢4の1/4が誤りである。

\paragraph{考え方・正答を選ぶ理由（査読待ち初稿）} 状態を晴・曇・雨と並べると遷移行列は対角1/2、非対角1/4である。2日後は\((P^2)_{ij}=\sum_kP_{ik}P_{kj}\)から自己3/8、異なる状態5/16。定常分布は(1/3,1/3,1/3)であり、全要素正のため初期天気によらず収束する。肢4の1/4が誤りである。

\Needspace{4\baselineskip}\paragraph{選択肢の検討（査読待ち初稿）}\mbox{}\par

\noindent\textbf{1：適切} 晴→2日後雨5/16。晴→晴→雨、晴→曇→雨、晴→雨→雨の排他的3経路を足すと1/8+1/16+1/8=\allowbreak{}5/16。\par

\noindent\textbf{2：適切} 曇→2日後曇3/8。曇→曇→曇、曇→雨→曇、曇→晴→曇を足すと1/4+1/16+1/16=\allowbreak{}3/8。\par

\noindent\textbf{3：適切} 雨→2日後曇5/16。自己遷移を一方の時刻に含む2経路と、晴を経る経路で5/16。\par

\noindent\textbf{4：不適切} 雨→2日後雨1/4。自己遷移を2回行う経路だけなら1/4だが、他の2状態を経て戻る1/16ずつを忘れている。合計3/8。\par

\noindent\textbf{5：適切} 初日雨から遠い将来の雨確率1/3。全ての1日遷移が正で、対称な行列の定常分布は一様。有限既約・非周期の連鎖でこの分布へ収束する。\par

\paragraph{再利用する考え方} 符号化の目的、情報量と物理的データ長、条件付き時間と平均時間を区別する。有限マルコフ連鎖の経路和・行列累乗・定常分布の定義と適用条件を確認する。

\paragraph{共通理論} \hyperref[theory:information:markov]{第\ref*{theory:information:markov}節}

\par 参考：\cite{info_markov}。

\paragraph{関連問題} \hyperref[q:h28-reg-basic-2-1]{H28-REG-BASIC-2-1}

\par 難易度 2（校準初稿）、優先度未評価。

\subsection{平成28年度 問2-1}
\label{q:h28-reg-basic-2-1}

\noindent\textbf{原典：} sources/H28/basic.pdf、PDF物理7頁。

\textbf{公式正答：} 4（single）。

\paragraph{問題の要旨} 同じ3状態の天気遷移図で、2日後確率の異なる記述を検証する。

\Needspace{4\baselineskip}\paragraph{確認した条件・記述（独自要約）}\mbox{}\par

\noindent\textbf{方向：} 最も不適切\par

\paragraph{問われる核心} 有限マルコフ連鎖の経路和・行列累乗・定常分布

\paragraph{必要知識} 有限マルコフ連鎖の経路和・行列累乗・定常分布。状態を晴・曇・雨と並べると遷移行列は対角1/2、非対角1/4である。2日後は\((P^2)_{ij}=\sum_kP_{ik}P_{kj}\)から自己3/8、異なる状態5/16。定常分布は(1/3,1/3,1/3)であり、全要素正のため初期天気によらず収束する。肢4の曇→晴3/16が不適切である。

\paragraph{考え方・正答を選ぶ理由（査読待ち初稿）} 状態を晴・曇・雨と並べると遷移行列は対角1/2、非対角1/4である。2日後は\((P^2)_{ij}=\sum_kP_{ik}P_{kj}\)から自己3/8、異なる状態5/16。定常分布は(1/3,1/3,1/3)であり、全要素正のため初期天気によらず収束する。肢4の曇→晴3/16が不適切である。

\Needspace{4\baselineskip}\paragraph{選択肢の検討（査読待ち初稿）}\mbox{}\par

\noindent\textbf{1：適切} 雨→雨3/8。雨→雨→雨1/4、雨→曇→雨1/16、雨→晴→雨1/16の合計3/8。\par

\noindent\textbf{2：適切} 晴→2日後雨5/16。晴→晴→雨、晴→曇→雨、晴→雨→雨の排他的3経路を足すと1/8+1/16+1/8=\allowbreak{}5/16。\par

\noindent\textbf{3：適切} 曇→2日後曇3/8。曇→曇→曇、曇→雨→曇、曇→晴→曇を足すと1/4+1/16+1/16=\allowbreak{}3/8。\par

\noindent\textbf{4：不適切} 曇→晴3/16。曇→曇→晴1/8、曇→雨→晴1/16、曇→晴→晴1/8の和は5/16。3/16ではない。\par

\noindent\textbf{5：適切} 初日雨から遠い将来の雨確率1/3。全ての1日遷移が正で、対称な行列の定常分布は一様。有限既約・非周期の連鎖でこの分布へ収束する。\par

\paragraph{再利用する考え方} 符号化の目的、情報量と物理的データ長、条件付き時間と平均時間を区別する。有限マルコフ連鎖の経路和・行列累乗・定常分布の定義と適用条件を確認する。

\paragraph{共通理論} \hyperref[theory:information:markov]{第\ref*{theory:information:markov}節}

\par 参考：\cite{info_markov}。

\paragraph{関連問題} \hyperref[q:h24-reg-basic-2-3]{H24-REG-BASIC-2-3}

\par 難易度 2（校準初稿）、優先度未評価。

\section{最短経路とボトルネック経路}

\subsection{平成23年度 問2-2}
\label{q:h23-reg-basic-2-2}

\noindent\textbf{原典：} sources/H23/basic.pdf、PDF物理5頁。

\textbf{公式正答：} 2（single）。

\paragraph{問題の要旨} 通信ネットワークの辺遅延を足し、AからFへの最短経路の所要値を求める。

\Needspace{4\baselineskip}\paragraph{確認した条件・記述（独自要約）}\mbox{}\par

\noindent\textbf{方向：} 最短遅延。\par

\noindent\textbf{原典の全辺：} AB=\allowbreak{}1、AC=\allowbreak{}2、BC=\allowbreak{}3、BD=\allowbreak{}1、BE=\allowbreak{}2、CE=\allowbreak{}4、CF=\allowbreak{}4、DF=\allowbreak{}4、EF=\allowbreak{}2。図の接続方向に従う。\par

\noindent\textbf{単位：} 原典が指定しない時間単位をms等と補わない。\par

\paragraph{問われる核心} 非負重みグラフの最短距離

\paragraph{必要知識} 辺重みの加算、緩和、Dijkstra、経由点の接続を図から読む。

\paragraph{考え方・正答を選ぶ理由（査読待ち初稿）} 非負の辺遅延を加算する。AからB=\allowbreak{}1、C=\allowbreak{}2、D=\allowbreak{}2、E=\allowbreak{}3、F=\allowbreak{}5と確定し、A-B-E-F=\allowbreak{}1+2+2を得る。各辺の幅の最小を評価するボトルネック問題とは異なる。

\Needspace{4\baselineskip}\paragraph{選択肢の検討（査読待ち初稿）}\mbox{}\par

\noindent\textbf{1：不一致} 1。最短でもA-B-E-Fの5を要する。単一の辺ABの1を総時間と取り違えている。\par

\noindent\textbf{2：適切} 5。A-B-E-Fの遅延は1+2+2=\allowbreak{}5。距離の確定でこれより短い経路がないと検証できる。\par

\noindent\textbf{3：不一致} 6。A-C-F=\allowbreak{}2+4やA-B-D-F=\allowbreak{}1+1+4は6だが、経路が存在することは最短性を保証しない。\par

\noindent\textbf{4：不一致} 8。より短い5の経路がある。図の辺を全て通る問題でも最長経路でもない。\par

\noindent\textbf{5：不一致} 10。5の経路と比較して大きい。総辺数や全遅延の合計を使わない。\par

\paragraph{再利用する考え方} 目的量・脅威・成立条件を定めてから手段を選ぶ。秘匿、改ざん検出、主体認証、再送防止、復旧は別の性質であり、方式名だけで全てを保証しない。

\paragraph{共通理論} \hyperref[theory:network:shortest]{第\ref*{theory:network:shortest}節}, \hyperref[theory:network:bottleneck]{第\ref*{theory:network:bottleneck}節}

\par 参考：\cite{net_dijkstra}。

\paragraph{関連問題} \hyperref[q:h26-reg-basic-2-6]{H26-REG-BASIC-2-6}

\par 難易度 2（校準初稿）、優先度未評価。

\subsection{平成25年度 問2-2}
\label{q:h25-reg-basic-2-2}

\noindent\textbf{原典：} sources/H25/basic.pdf、PDF物理7頁。

\textbf{公式正答：} 5（single）。

\paragraph{問題の要旨} 道路距離の最短施設への登録を、同距離で複数登録する条件も含めて数える。

\Needspace{4\baselineskip}\paragraph{確認した条件・記述（独自要約）}\mbox{}\par

\noindent\textbf{全辺：} Au2、uB2、uv3、AB1、Ax3、Bx3、Bw4、xw3、vw2、vC6、wC4、xC3。\par

\noindent\textbf{利用者：} u,v,w,x各1人。同距離なら全施設へ登録。\par

\paragraph{問われる核心} 道路距離の最短施設への登録を、同距離で複数登録する条件も含めて数える。

\paragraph{必要知識} 重み付き道路網の最短距離と同点の複数登録。

\paragraph{考え方・正答を選ぶ理由（査読待ち初稿）} 各利用者から施設A,B,Cへの最短距離を求める。uは(2,2,8)、vは(5,5,6)、wは(5,4,4)、xは(3,3,3)である。したがってu,vはA/B、wはB/C、xは全施設へ登録され、a=\allowbreak{}3,b=\allowbreak{}4,c=\allowbreak{}2。5である。最短距離の同点を一施設だけへ割り当てない。

\Needspace{4\baselineskip}\paragraph{選択肢の検討（査読待ち初稿）}\mbox{}\par

\noindent\textbf{1：不適切} a>b>c。Bの人数がAより多いので逆である。\par

\noindent\textbf{2：不適切} a=\allowbreak{}b>c。A3とB4は等しくない。\par

\noindent\textbf{3：不適切} a=\allowbreak{}b=\allowbreak{}c。三施設の人数は3,4,2で異なる。\par

\noindent\textbf{4：不適切} b>a=\allowbreak{}c。A3とC2は等しくない。\par

\noindent\textbf{5：適切} b>a>c。B4>A3>C2である。\par

\paragraph{再利用する考え方} 定義した対象量、演算順序、保存条件を固定し、途中状態と候補の違いを比較する。

\paragraph{共通理論} \hyperref[theory:network:shortest]{第\ref*{theory:network:shortest}節}

\par 難易度 3（校準初稿）、優先度未評価。

\subsection{平成26年度 問2-6}
\label{q:h26-reg-basic-2-6}

\noindent\textbf{原典：} sources/H26/basic.pdf、PDF物理8頁。

\textbf{公式正答：} 3（single）。

\paragraph{問題の要旨} 道路網で通行できる車高の最大値を、各経路の最小車高制限から求める。

\Needspace{4\baselineskip}\paragraph{確認した条件・記述（独自要約）}\mbox{}\par

\noindent\textbf{方向：} 車高最大。\par

\noindent\textbf{位置名：} 以下の節点名は解説の独自呼称。A、U=左上、V=上中央、W=右上、R=右中、L=左中、P=内左下、D=左下、E=下中央、C=中央、Q=内右下、B。\par

\noindent\textbf{全辺の車高制限：} AU8、AL9、UV4、UP3、VW9、WR7、RB6、LP2、LD6、PC6、PE6、DE7、CV5、CR3、CQ9、QB4、EQ5。\par

\noindent\textbf{実現経路：} A-L-D-E-P-C-V-W-R-Bの最小5。\par

\paragraph{問われる核心} 経路の最小車高制限を最大化するボトルネック

\paragraph{必要知識} max-min経路、閾値での到達可能性、図の交点と道路接続。

\paragraph{考え方・正答を選ぶ理由（査読待ち初稿）} 経路の通行可能車高は車高制限の最小、答えはそれを経路全体について最大化した値である。車高5の接続経路を示し、車高6以上の辺だけ残すと分断されることを確認するので最大値5。遅延の総和最小化を使わない。

\Needspace{4\baselineskip}\paragraph{選択肢の検討（査読待ち初稿）}\mbox{}\par

\noindent\textbf{1：最大でない} 3。車高5まで通る経路があり最大でない。単一の車高制限3の道路を避ける経路も存在する。\par

\noindent\textbf{2：最大でない} 4。車高制限4以下の道を通らなくても車高5の経路がある。最初に見つけた経路の車高制限を最大としない。\par

\noindent\textbf{3：適切} 5。図のAから左中・左下・下中央・内左下・中央・上中央・上右・右中・Bへの経路の最小車高制限は5である。\par

\noindent\textbf{4：通行経路なし} 6。車高6以上だけを残すとA側とB側が接続しない。経路全体の最小車高制限が車高を制限する。\par

\noindent\textbf{5：通行経路なし} 7。車高7以上の経路がないことは6以上で既に分断されることからも分かる。最大制限の一辺だけでは判断できない。\par

\paragraph{再利用する考え方} 目的量・脅威・成立条件を定めてから手段を選ぶ。秘匿、改ざん検出、主体認証、再送防止、復旧は別の性質であり、方式名だけで全てを保証しない。

\paragraph{共通理論} \hyperref[theory:network:bottleneck]{第\ref*{theory:network:bottleneck}節}, \hyperref[theory:network:shortest]{第\ref*{theory:network:shortest}節}

\par 参考：\cite{net_dijkstra}。

\paragraph{関連問題} \hyperref[q:h23-reg-basic-2-2]{H23-REG-BASIC-2-2}

\par 難易度 2（校準初稿）、優先度未評価。

\section{動的経路制御の情報共有と次ホップ}

\subsection{令和6年度 問2-3}
\label{q:r06-reg-basic-2-3}

\noindent\textbf{原典：} sources/R06/basic.pdf、PDF物理7頁。

\textbf{公式正答：} 2（single）。

\paragraph{問題の要旨} リンクステートと距離ベクトルの保持情報・経路計算を比較する。

\Needspace{4\baselineskip}\paragraph{確認した条件・記述（独自要約）}\mbox{}\par

\noindent\textbf{方向：} 最も不適切なもの。\par

\noindent\textbf{時点：} R06年度。現在の推奨を出題当時へ無条件に転用しない。\par

\paragraph{問われる核心} トポロジ共有と距離の分散更新

\paragraph{必要知識} 次ホップ、経路表、SPF、Bellman-Ford、OSPF、RIP、適用範囲。

\paragraph{考え方・正答を選ぶ理由（査読待ち初稿）} 2が保持情報を同一視している。リンクステートは範囲内のリンク情報、距離ベクトルは宛先までの隣接経由距離を更新する。\(d_i(t)=\min_j\{c_{ij}+d_j(t)\}\)で次ホップを選べるので全接続図が必須ではない。

\Needspace{4\baselineskip}\paragraph{選択肢の検討（査読待ち初稿）}\mbox{}\par

\noindent\textbf{1：適切（記述の条件内）} 経路表を参照して次ホップへ転送。。動的ルーティングが構築した経路表を用いた次ホップ転送の説明として妥当。実装には静的経路や転送用FIBもあるが、問題が説明する動的経路制御を否定しない。\par

\noindent\textbf{2：不適切} どちらもネットワークトポロジを把握。。距離ベクトルの節点は隣接節点の宛先距離を利用するのでネットワーク全体の詳細トポロジを必ず保持するわけではない。\par

\noindent\textbf{3：適切（記述の条件内）} リンク情報を広告し最短経路計算。。リンクステートは領域内のリンク状態を共有してSPF等を計算する。全インターネットの全接続を各ルータが保持する意味ではない。\par

\noindent\textbf{4：適切（記述の条件内）} 宛先距離を伝搬し最短となる次ホップ決定。。隣接経由コストと隣接節点からの距離を足して比較する。Bellman-Ford型の分散更新である。\par

\noindent\textbf{5：適切（記述の条件内）} OSPF/IS-ISはリンク状態、RIPは距離ベクトル。。代表例の対応は妥当。BGPのパスベクトル等もあり、全プロトコルを厳密に2種だけへ封じる分類ではない。\par

\paragraph{再利用する考え方} 目的量・脅威・成立条件を定めてから手段を選ぶ。秘匿、改ざん検出、主体認証、再送防止、復旧は別の性質であり、方式名だけで全てを保証しない。

\paragraph{共通理論} \hyperref[theory:network:routing]{第\ref*{theory:network:routing}節}, \hyperref[theory:network:shortest]{第\ref*{theory:network:shortest}節}

\par 参考：\cite{net_ospf,net_rip}。

\par 難易度 2（校準初稿）、優先度未評価。

\section{パスワード運用・端末・認証要素}

\subsection{平成23年度 問2-5}
\label{q:h23-reg-basic-2-5}

\noindent\textbf{原典：} sources/H23/basic.pdf、PDF物理7頁。

\textbf{公式正答：} 2（single）。

\par\textbf{検討未確定：} 選択肢2は公開鍵の真正性確認・認証済み保護通信という条件を省く。方式を用いるだけで全漏洩が防げると一般化しない。初回2・公式2を保持。

\paragraph{問題の要旨} パスワードの伝送・保管・端末入力・推測攻撃の対策を比較する。

\Needspace{4\baselineskip}\paragraph{確認した条件・記述（独自要約）}\mbox{}\par

\noindent\textbf{方向：} 最も適切なもの。\par

\noindent\textbf{時点：} H23年度。現在の推奨を出題当時へ無条件に転用しない。\par

\paragraph{問われる核心} 認証秘密の各ライフサイクルの脅威

\paragraph{必要知識} 安全な暗号通信、信頼した公開鍵、平文保存、辞書攻撃、キーロガー、WEP。

\paragraph{考え方・正答を選ぶ理由（査読待ち初稿）} 2は安全に認証した暗号通信でパスワード伝送を保護する意図として選ぶ。保管の平文露出、端末の入力監視、予測可能な列、WEPの弱点は別に残る。公開鍵の認証とプロトコル条件を本文で補い、原典の方式名だけの保証は未確定として保持する。

\Needspace{4\baselineskip}\paragraph{選択肢の検討（査読待ち初稿）}\mbox{}\par

\noindent\textbf{1：不適切} WEP無線LANなら覚えやすい単語パスワードを安全に利用。。WEPには既知の弱点があり安全を保証しない。無線区間の暗号と端末・Webサービスの安全性も別である。\par

\noindent\textbf{2：意図として適切・条件不足} 公開鍵暗号でパスワード伝送時の漏洩を防ぐ。。適切に鍵を認証して構成した暗号通信なら伝送中の秘匿に寄与する。ただし公開鍵という方式名だけで偽サーバへの入力や中間者攻撃を防ぐ保証はない。\par

\noindent\textbf{3：不適切} 忘れたとき通知できるようサーバにパスワードをそのまま保存。。ファイルを読める者へ認証秘密が露出する。サーバ側は適切なソルトと計算コストを伴うパスワードハッシュ等で検証する。\par

\noindent\textbf{4：不適切} 長いパスワードなら公共端末で安全。。共有端末のキーロガー等は入力そのものを読むため文字数だけでは防げない。長さは推測耐性の要素である。\par

\noindent\textbf{5：不適切} asdfghjkなら辞書攻撃に強い。。キーボード配列上の連続列は推測候補になり得る。辞書攻撃の候補は自然言語の辞書語だけとは限らない。\par

\paragraph{再利用する考え方} 目的量・脅威・成立条件を定めてから手段を選ぶ。秘匿、改ざん検出、主体認証、再送防止、復旧は別の性質であり、方式名だけで全てを保証しない。

\paragraph{共通理論} \hyperref[theory:network:password]{第\ref*{theory:network:password}節}, \hyperref[theory:network:tls]{第\ref*{theory:network:tls}節}, \hyperref[theory:network:wireless]{第\ref*{theory:network:wireless}節}

\par 参考：\cite{net_nist2017,net_nist2025,net_wep}。

\paragraph{関連問題} \hyperref[q:r05-reg-basic-2-1]{R05-REG-BASIC-2-1}

\par 難易度 2（校準初稿）、優先度未評価。

\subsection{平成29年度 問2-1}
\label{q:h29-reg-basic-2-1}

\noindent\textbf{原典：} sources/H29/basic.pdf、PDF物理6頁。

\textbf{公式正答：} 1（single）。

\par\textbf{検討未確定：} 肢4の機械的な定期変更推奨は2017年6月NIST指針と一致しない。公式1・初回1を保持し、当時の国内出題根拠と排他性を査読対象にする。

\paragraph{問題の要旨} 差出人・更新・パスワード・外部媒体に関する基本防御を判断する。

\Needspace{4\baselineskip}\paragraph{確認した条件・記述（独自要約）}\mbox{}\par

\noindent\textbf{方向：} 最も不適切なもの。\par

\noindent\textbf{時点：} H29年度。現在の推奨を出題当時へ無条件に転用しない。\par

\paragraph{問われる核心} 送信者を装う攻撃と認証秘密の運用

\paragraph{必要知識} なりすまし、アカウント侵害、パッチ、マルウェア対策、定期変更と侵害時変更。

\paragraph{考え方・正答を選ぶ理由（査読待ち初稿）} 差出人アドレスで添付安全を決める1が不適切。2・3・5は脅威低減の基本策である。4の周期変更について当時の推奨を追加査読する。2017年に既に存在するNIST指針を踏まえ、他肢全てを無条件正しいとは断定しない。

\Needspace{4\baselineskip}\paragraph{選択肢の検討（査読待ち初稿）}\mbox{}\par

\noindent\textbf{1：不適切} 知人のアドレスからなら添付を直ちに開く。。表示アドレスの偽装や知人アカウントの侵害があり、本人らしさだけでは添付の安全性を判断できない。\par

\noindent\textbf{2：適切} 端末とサーバへウイルス対策。。適切に運用・更新した対策は感染リスク低減に寄与する。対策ソフトだけで全感染や損害を消せるわけではない。\par

\noindent\textbf{3：適切} 脆弱性更新を確認しパッチ適用。。脆弱性の修正を適時に適用して既知の侵入経路を減らす。業務環境では影響確認と運用管理も必要である。\par

\noindent\textbf{4：再利用回避は適切・周期変更に疑義} 定期的に変更し以前のパスワードは再利用しない。。再利用の回避は妥当だが、機械的な周期変更の一般推奨には疑義がある。NISTの2017年6月版は侵害証拠がない恣意的な定期変更を推奨しない。出題時の資料の採用範囲は未確定。\par

\noindent\textbf{5：適切} 出所不明のプログラム・USBを使用しない。。由来不明の実行コードや媒体による感染を避ける対策である。USBという規格自体が常に悪性という意味ではない。\par

\paragraph{再利用する考え方} 目的量・脅威・成立条件を定めてから手段を選ぶ。秘匿、改ざん検出、主体認証、再送防止、復旧は別の性質であり、方式名だけで全てを保証しない。

\paragraph{共通理論} \hyperref[theory:network:password]{第\ref*{theory:network:password}節}, \hyperref[theory:network:malware]{第\ref*{theory:network:malware}節}, \hyperref[theory:network:targeted]{第\ref*{theory:network:targeted}節}

\par 参考：\cite{net_nist2017,net_target2020,net_ipa_account}。

\paragraph{関連問題} \hyperref[q:h30-reg-basic-2-1]{H30-REG-BASIC-2-1}, \hyperref[q:r01-ret-basic-2-1]{R01-RET-BASIC-2-1}

\par 難易度 2（校準初稿）、優先度未評価。

\subsection{平成30年度 問2-1}
\label{q:h30-reg-basic-2-1}

\noindent\textbf{原典：} sources/H30/basic.pdf、PDF物理6頁。

\textbf{公式正答：} 4（single）。

\par\textbf{検討未確定：} 原典2はSSL/TLSと一括表記し版・暗号設定を省く。SSLv3禁止は出題前の2015年であり、SSLを安全として一般推奨できない。総称としての出題意図と字義を区別し公式4を保持。

\paragraph{問題の要旨} 境界防御・保護通信・IoT・パスワード再利用・無線接続先を評価する。

\Needspace{4\baselineskip}\paragraph{確認した条件・記述（独自要約）}\mbox{}\par

\noindent\textbf{方向：} 最も不適切なもの。\par

\noindent\textbf{時点：} H30年度。現在の推奨を出題当時へ無条件に転用しない。\par

\paragraph{問われる核心} 攻撃面とパスワード漏洩の横展開

\paragraph{必要知識} 境界防御、認証したTLS、IoT接続、資格情報再利用、無線AP。

\paragraph{考え方・正答を選ぶ理由（査読待ち初稿）} 4は認証情報の漏洩を横展開するため不適切。各防御は対応する脅威を限定して評価する。原典2のSSLという名称はそのまま校正条件へ記録し、安全なTLSの意味と廃止済みSSLを区別する。

\Needspace{4\baselineskip}\paragraph{選択肢の検討（査読待ち初稿）}\mbox{}\par

\noindent\textbf{1：適切（限定防御）} ファイアウォールを利用。。許可通信を制御して不正アクセス等を減らす。許可通信内の攻撃や端末からの持出しまで単独で全防止するものではない。\par

\noindent\textbf{2：TLS意図で適切・SSL版に疑義} SSL/TLSのような暗号通信を確認。。適切なTLSと相手認証が必要である。原典のSSL/TLSを時代の総称として保持するが、SSLv3は2015年RFC7568で禁止されており、SSLを現行でも安全と読み替えない。\par

\noindent\textbf{3：適切（停止可能な機器）} 常時不要なIoTは使用後電源オフ。。常時使わず停止してよい機器について接続時間を減らす策である。継続監視や安全機能を必要とする装置を一律停止する指示ではない。\par

\noindent\textbf{4：不適切} 忘れないため複数サービスで同じパスワード。。一つのサービスの漏洩情報が他サービスの認証へ悪用される。サービスごとの独自の強いパスワードが必要である。\par

\noindent\textbf{5：適切} 意図しない無線AP接続を確認。。偽・無許可APへの自動接続を避ける。接続先確認に加え適切な暗号通信等も必要である。\par

\paragraph{再利用する考え方} 目的量・脅威・成立条件を定めてから手段を選ぶ。秘匿、改ざん検出、主体認証、再送防止、復旧は別の性質であり、方式名だけで全てを保証しない。

\paragraph{共通理論} \hyperref[theory:network:password]{第\ref*{theory:network:password}節}, \hyperref[theory:network:tls]{第\ref*{theory:network:tls}節}, \hyperref[theory:network:wireless]{第\ref*{theory:network:wireless}節}

\par 参考：\cite{net_ssl2015,net_ipa_account,net_wep}。

\paragraph{関連問題} \hyperref[q:h29-reg-basic-2-1]{H29-REG-BASIC-2-1}

\par 難易度 2（校準初稿）、優先度未評価。

\subsection{令和5年度 問2-1}
\label{q:r05-reg-basic-2-1}

\noindent\textbf{原典：} sources/R05/basic.pdf、PDF物理5頁。

\textbf{公式正答：} 1（single）。

\paragraph{問題の要旨} パスワード変更、PIN、生体、段階数と要素数、旧無線暗号を評価する。

\Needspace{4\baselineskip}\paragraph{確認した条件・記述（独自要約）}\mbox{}\par

\noindent\textbf{方向：} 最も適切なもの。\par

\noindent\textbf{時点：} R05年度。現在の推奨を出題当時へ無条件に転用しない。\par

\paragraph{問われる核心} 認証要素・試行制限・秘密の更新

\paragraph{必要知識} 知識/所持/生体、二段階と多要素、PINのローカル検証、レート制限、漏洩時変更。

\paragraph{考え方・正答を選ぶ理由（査読待ち初稿）} 1が漏洩時変更と恣意的な定期変更を区別する。2〜5は用途・試行制限、確率的な生体照合、認証要素、既知の暗号弱点を無視している。NIST2017と現行2025の表現強度は時点を明記して分ける。

\Needspace{4\baselineskip}\paragraph{選択肢の検討（査読待ち初稿）}\mbox{}\par

\noindent\textbf{1：適切} 複雑で独自のパスワードなら侵害事案なしの定期変更は基本不要。。機械的な周期変更と侵害時の変更を区別する。漏洩の疑いがあれば迅速に変更し、使い回しを避ける。サービス規程や長さ・予測可能性も重要。\par

\noindent\textbf{2：不適切} 4〜6桁PINは総当たりされるので使うべきでない。。PINを端末内で検証し試行を制限する場合等がある。桁数だけでオンライン無制限攻撃と同一視し全用途を禁止しない。\par

\noindent\textbf{3：不適切} 生体認証は本人がいなければ必ず失敗。。生体照合は確率的で、他人受入れや提示攻撃もある。固有の情報だから本人の存在を数学的に保証するわけではない。\par

\noindent\textbf{4：不適切} 二段階一要素は一段階二要素より安全。。段階は手続の回数、要素は知識・所持・生体等の種類であり違う。知識を2回入力しても2要素にはならない。段階数のみで安全性の順序は決まらない。\par

\noindent\textbf{5：不適切} 古いWEP対応機器は買換えまで安全。。WEPに既知の弱点があり、古い機器という理由で安全性が復活しない。買換えまでの運用も適切な代替・停止等を検討する。\par

\paragraph{再利用する考え方} 目的量・脅威・成立条件を定めてから手段を選ぶ。秘匿、改ざん検出、主体認証、再送防止、復旧は別の性質であり、方式名だけで全てを保証しない。

\paragraph{共通理論} \hyperref[theory:network:password]{第\ref*{theory:network:password}節}, \hyperref[theory:network:factors]{第\ref*{theory:network:factors}節}, \hyperref[theory:network:wireless]{第\ref*{theory:network:wireless}節}

\par 参考：\cite{net_nist2017,net_nist2025,net_wep}。

\paragraph{関連問題} \hyperref[q:h23-reg-basic-2-5]{H23-REG-BASIC-2-5}

\par 難易度 2（校準初稿）、優先度未評価。

\section{公開鍵暗号・署名・PKIの鍵の役割}

\subsection{平成27年度 問2-5}
\label{q:h27-reg-basic-2-5}

\noindent\textbf{原典：} sources/H27/basic.pdf、PDF物理6頁。

\textbf{公式正答：} 2（single）。

\paragraph{問題の要旨} 公開鍵基盤・電子署名・端末・無線暗号の役割を区別する。

\Needspace{4\baselineskip}\paragraph{確認した条件・記述（独自要約）}\mbox{}\par

\noindent\textbf{方向：} 最も適切なもの。\par

\noindent\textbf{時点：} H27年度。現在の推奨を出題当時へ無条件に転用しない。\par

\paragraph{問われる核心} 公開鍵の役割と証明書の信頼

\paragraph{必要知識} 公開鍵暗号、秘密鍵、署名生成と検証、CA証明書検証、WEP。

\paragraph{考え方・正答を選ぶ理由（査読待ち初稿）} 2がPKIの主体と公開鍵の結び付けを説明する。公開鍵で秘密を隠す暗号化と、秘密鍵で生成し公開鍵で検証する署名は目的・鍵の向きが違う。端末や無線の絶対安全という断定を排除する。

\Needspace{4\baselineskip}\paragraph{選択肢の検討（査読待ち初稿）}\mbox{}\par

\noindent\textbf{1：不適切} 公開鍵だけが必要。。対となる秘密鍵が復号や署名生成に必要である。公開する鍵と保護する鍵を分ける。\par

\noindent\textbf{2：適切} 認証局で公開鍵の所有者を保証。。CAが主体と公開鍵の対応へ署名し、利用者が証明書の署名・有効期限・信頼経路等を検証する仕組みである。保証範囲は証明書の検証条件内である。\par

\noindent\textbf{3：不適切} スマートフォンは感染しない。。OSの違いによって無条件免疫にはならない。不正アプリ、脆弱性、権限悪用等への対策が必要である。\par

\noindent\textbf{4：不適切} 公開鍵でダイジェストを暗号化して署名。。署名生成には秘密鍵を用い公開鍵で検証する。全署名方式を単純な「秘密鍵で暗号化」と表すのも不正確である。\par

\noindent\textbf{5：不適切} WEPなら完全に盗聴を防ぐ。。WEPは既知の攻撃がある。暗号を導入した事実だけで完全な秘匿を保証しない。\par

\paragraph{再利用する考え方} 目的量・脅威・成立条件を定めてから手段を選ぶ。秘匿、改ざん検出、主体認証、再送防止、復旧は別の性質であり、方式名だけで全てを保証しない。

\paragraph{共通理論} \hyperref[theory:network:keys]{第\ref*{theory:network:keys}節}, \hyperref[theory:network:pki]{第\ref*{theory:network:pki}節}, \hyperref[theory:network:wireless]{第\ref*{theory:network:wireless}節}

\par 参考：\cite{net_pki,net_signature,net_wep}。

\paragraph{関連問題} \hyperref[q:r03-reg-basic-2-1]{R03-REG-BASIC-2-1}

\par 難易度 2（校準初稿）、優先度未評価。

\subsection{令和3年度 問2-1}
\label{q:r03-reg-basic-2-1}

\noindent\textbf{原典：} sources/R03/basic.pdf、PDF物理6頁。

\textbf{公式正答：} 1（single）。

\paragraph{問題の要旨} 公開鍵暗号・証明書・署名の鍵の役割と無線の推奨を比較する。

\Needspace{4\baselineskip}\paragraph{確認した条件・記述（独自要約）}\mbox{}\par

\noindent\textbf{方向：} 最も適切なもの。\par

\noindent\textbf{時点：} R03年度。現在の推奨を出題当時へ無条件に転用しない。\par

\paragraph{問われる核心} 暗号化・署名・証明書の分離

\paragraph{必要知識} 公開鍵暗号の復号者、署名者と検証者、CA、WEP、端末感染。

\paragraph{考え方・正答を選ぶ理由（査読待ち初稿）} 1が秘匿のための公開鍵暗号を説明する。署名は秘密鍵で作り公開鍵で検証し、CAの証明は公開鍵の対応である。同じ鍵対でも暗号化と署名を一つの手順へ統合しない。

\Needspace{4\baselineskip}\paragraph{選択肢の検討（査読待ち初稿）}\mbox{}\par

\noindent\textbf{1：適切} 公開鍵で暗号化し秘密鍵で復号。。公開鍵暗号で受信者へ秘匿する基本の鍵の向きである。実用通信では鍵確立と共通鍵暗号を組み合わせる場合もある。\par

\noindent\textbf{2：不適切} CAはユーザと秘密鍵を結び付ける。。証明書が公開鍵と主体の対応を結び付ける。秘密鍵を第三者へ公開・配付して所有者を証明する仕組みではない。\par

\noindent\textbf{3：不適切} スマートフォン感染の報告はない。。不正アプリ・端末脆弱性等に対する対策を不要にはできない。報告がないという原典の前提も保持できない。\par

\noindent\textbf{4：不適切} 公開鍵で署名生成し秘密鍵で検証。。逆である。秘密鍵で生成し対応公開鍵で検証する。検証者に秘密鍵を配らない。\par

\noindent\textbf{5：不適切} 現在の無線LANではWEPを推奨。。WEPは既知の弱い無線暗号であり出題時の安全な推奨としない。現在でもWEPを安全と教えない。\par

\paragraph{再利用する考え方} 目的量・脅威・成立条件を定めてから手段を選ぶ。秘匿、改ざん検出、主体認証、再送防止、復旧は別の性質であり、方式名だけで全てを保証しない。

\paragraph{共通理論} \hyperref[theory:network:keys]{第\ref*{theory:network:keys}節}, \hyperref[theory:network:pki]{第\ref*{theory:network:pki}節}, \hyperref[theory:network:wireless]{第\ref*{theory:network:wireless}節}

\par 参考：\cite{net_pki,net_signature,net_wep}。

\paragraph{関連問題} \hyperref[q:h27-reg-basic-2-5]{H27-REG-BASIC-2-5}, \hyperref[q:r06-reg-basic-2-6]{R06-REG-BASIC-2-6}

\par 難易度 2（校準初稿）、優先度未評価。

\section{包除原理・補集合と多重共通部分}

\subsection{平成30年度 問2-6}
\label{q:h30-reg-basic-2-6}

\noindent\textbf{原典：} sources/H30/basic.pdf、PDF物理10頁。

\textbf{公式正答：} 2（single）。

\paragraph{問題の要旨} 三集合の包除原理から和集合の補集合の要素数を求める。

\Needspace{4\baselineskip}\paragraph{確認した条件・記述（独自要約）}\mbox{}\par

\noindent\textbf{数値：} |U|900、|A|300、|B|180、|C|128、|A∩B|60、|A∩C|43、|B∩C|26、三重9。\par

\noindent\textbf{対象：} 上線はA∪B∪C全体の補集合である。\par

\paragraph{問われる核心} 三集合の包除原理から和集合の補集合の要素数を求める。

\paragraph{必要知識} 三集合の重複補正と補集合の分母。

\paragraph{考え方・正答を選ぶ理由（査読待ち初稿）} 和集合の数は300+180+128−60−43−26+9=\allowbreak{}488である。原典では和集合全体へ上線が付いており、その補集合を問う。全体900から引いて412、2となる。

\Needspace{4\baselineskip}\paragraph{選択肢の検討（査読待ち初稿）}\mbox{}\par

\noindent\textbf{1：不適切} 385。与えた交差数の包除と全体900からの差引きに合わない。\par

\noindent\textbf{2：適切} 412。和集合488の補集合は900−488=\allowbreak{}412である。\par

\noindent\textbf{3：不適切} 420。三重交差等の補正を正しく扱った412とは異なる。\par

\noindent\textbf{4：不適切} 480。三重交差の加算を落とした等の誤りであり412ではない。\par

\noindent\textbf{5：不適切} 488。488は和集合の数であり問う補集合ではない。\par

\paragraph{再利用する考え方} 定義した対象量、演算順序、保存条件を固定し、途中状態と候補の違いを比較する。

\paragraph{共通理論} \hyperref[theory:information:sets]{第\ref*{theory:information:sets}節}, \hyperref[theory:discrete:inclusion]{第\ref*{theory:discrete:inclusion}節}

\par 参考：\cite{mit_discrete2015}。

\paragraph{関連問題} \hyperref[q:r04-reg-basic-2-2]{R04-REG-BASIC-2-2}

\par 難易度 2（校準初稿）、優先度未評価。

\subsection{令和4年度 問2-2}
\label{q:r04-reg-basic-2-2}

\noindent\textbf{原典：} sources/R04/basic.pdf、PDF物理5頁。

\textbf{公式正答：} 3（single）。

\paragraph{問題の要旨} 四集合の包除原理から四重共通部分の要素数を逆算する。

\Needspace{4\baselineskip}\paragraph{確認した条件・記述（独自要約）}\mbox{}\par

\noindent\textbf{判断対象：} 四集合の包除と二重6・三重4の数え上げ。\par

\noindent\textbf{原典確認：} 図・式・全5候補を29保存画像で確認。rawの数式欠落と演算子誤認は別校正JSONで扱う。\par

\paragraph{問われる核心} 四集合の包除原理から四重共通部分の要素数を逆算する。

\paragraph{必要知識} 四集合の包除と二重6・三重4の数え上げ。

\paragraph{考え方・正答を選ぶ理由（査読待ち初稿）} 単集合は4通り、二重は6通り、三重は4通りなので、和集合16=\allowbreak{}4×11−6×7+4×4−x=\allowbreak{}18−x。x=\allowbreak{}2、3である。各交差が同じ値でも、一つだけ引いたり足したりせず組合せ数を掛ける。

\Needspace{4\baselineskip}\paragraph{選択肢の検討（査読待ち初稿）}\mbox{}\par

\noindent\textbf{1：不適切} 8。三重共通部分の一つ4より大きく包含条件にも反する。\par

\noindent\textbf{2：不適切} 4。包除式では和集合14となり条件16と違う。\par

\noindent\textbf{3：適切} 2。44−42+16−2=\allowbreak{}16で条件に合う。\par

\noindent\textbf{4：不適切} 1。包除式で和集合17となる。\par

\noindent\textbf{5：不適切} 0。包除式で和集合18となる。\par

\paragraph{再利用する考え方} 定義した対象量、演算順序、保存条件を固定し、途中状態と候補の違いを比較する。

\paragraph{共通理論} \hyperref[theory:discrete:inclusion]{第\ref*{theory:discrete:inclusion}節}

\par 参考：\cite{mit_discrete2015}。

\paragraph{関連問題} \hyperref[q:h30-reg-basic-2-6]{H30-REG-BASIC-2-6}

\par 難易度 2（校準初稿）、優先度未評価。

\section{可逆性・圧縮限界・普遍性と情報源通信路分離}

\subsection{令和2年度 問2-1}
\label{q:r02-reg-basic-2-1}

\noindent\textbf{原典：} sources/R02/basic.pdf、PDF物理7頁。

\textbf{公式正答：} 4（single）。

\paragraph{問題の要旨} 可逆/非可逆圧縮、代表形式、未知情報源への圧縮可能性と圧縮限界を比較する。

\Needspace{4\baselineskip}\paragraph{確認した条件・記述（独自要約）}\mbox{}\par

\noindent\textbf{方向：} 最も不適切\par

\paragraph{問われる核心} 可逆性・平均符号長の限界と普遍的/適応的圧縮

\paragraph{必要知識} 可逆性・平均符号長の限界と普遍的/適応的圧縮。事前分布の既知性と、観測データから規則を構成することを区別する。DEFLATEの反復列への参照や動的符号表は、入力から圧縮表現を作れる具体例である。未知だから圧縮不能という肢が不適切。限界は一意復号や平均長の条件付きであり、特定の短い列にエントロピー平均下限を機械的に適用しない。

\paragraph{考え方・正答を選ぶ理由（査読待ち初稿）} 事前分布の既知性と、観測データから規則を構成することを区別する。DEFLATEの反復列への参照や動的符号表は、入力から圧縮表現を作れる具体例である。未知だから圧縮不能という肢が不適切。限界は一意復号や平均長の条件付きであり、特定の短い列にエントロピー平均下限を機械的に適用しない。

\Needspace{4\baselineskip}\paragraph{選択肢の検討（査読待ち初稿）}\mbox{}\par

\noindent\textbf{1：適切} 可逆圧縮は完全復元しテキスト等で多用。。元ビット列を一意に復元する性質が可逆の定義である。文字1つの変更が意味を変えるテキストでは重要。\par

\noindent\textbf{2：適切} 非可逆圧縮は完全復元しない場合があり音声・映像等で多用。。知覚上許せる情報損失を利用してデータ量を減らす。全ての音声/映像が非可逆という断定ではない。\par

\noindent\textbf{3：適切} JPEGは静止画、MPEGは動画の代表的圧縮方式。。代表用途として正しい。JPEG/MPEGは複数の方式・規格を持つ名称であり、全方式を一律に可逆/非可逆とは扱わない。\par

\noindent\textbf{4：不適切} 情報源の事前知識がなければデータ圧縮できない。。データから頻度を学ぶ適応方式や繰り返しを辞書へ写す方式が反例。真の生起確率を事前に与える必要はない。\par

\noindent\textbf{5：適切} 可逆圧縮には方式では超えられない限界がある。。一意復号可能な符号の平均長はエントロピーで制約され、全ての入力を必ず短くする単射も作れない。個別列が必ず同じ下限長という意味ではない。\par

\paragraph{再利用する考え方} 符号化の目的、情報量と物理的データ長、条件付き時間と平均時間を区別する。可逆性・平均符号長の限界と普遍的/適応的圧縮の定義と適用条件を確認する。

\paragraph{共通理論} \hyperref[theory:information:compression]{第\ref*{theory:information:compression}節}, \hyperref[theory:information:entropy]{第\ref*{theory:information:entropy}節}, \hyperref[theory:information:universal]{第\ref*{theory:information:universal}節}

\par 参考：\cite{info_deflate,info_lossless,info_universal,info_jpeg,info_mpeg}。

\paragraph{関連問題} \hyperref[q:r05-reg-basic-2-4]{R05-REG-BASIC-2-4}

\par 難易度 2（校準初稿）、優先度未評価。

\subsection{令和5年度 問2-4}
\label{q:r05-reg-basic-2-4}

\noindent\textbf{原典：} sources/R05/basic.pdf、PDF物理7頁。

\textbf{公式正答：} 1（single）。

\paragraph{問題の要旨} 可逆/非可逆圧縮、代表形式、未知情報源への圧縮可能性と圧縮限界を比較する。

\Needspace{4\baselineskip}\paragraph{確認した条件・記述（独自要約）}\mbox{}\par

\noindent\textbf{方向：} 最も不適切\par

\paragraph{問われる核心} 可逆性・平均符号長の限界と普遍的/適応的圧縮

\paragraph{必要知識} 可逆性・平均符号長の限界と普遍的/適応的圧縮。事前分布の既知性と、観測データから規則を構成することを区別する。DEFLATEの反復列への参照や動的符号表は、入力から圧縮表現を作れる具体例である。未知だから圧縮不能という肢が不適切。限界は一意復号や平均長の条件付きであり、特定の短い列にエントロピー平均下限を機械的に適用しない。

\paragraph{考え方・正答を選ぶ理由（査読待ち初稿）} 事前分布の既知性と、観測データから規則を構成することを区別する。DEFLATEの反復列への参照や動的符号表は、入力から圧縮表現を作れる具体例である。未知だから圧縮不能という肢が不適切。限界は一意復号や平均長の条件付きであり、特定の短い列にエントロピー平均下限を機械的に適用しない。

\Needspace{4\baselineskip}\paragraph{選択肢の検討（査読待ち初稿）}\mbox{}\par

\noindent\textbf{1：不適切} 情報源の事前知識がなければデータ圧縮できない。。データから頻度を学ぶ適応方式や繰り返しを辞書へ写す方式が反例。真の生起確率を事前に与える必要はない。\par

\noindent\textbf{2：適切} 可逆圧縮には方式では超えられない限界がある。。一意復号可能な符号の平均長はエントロピーで制約され、全ての入力を必ず短くする単射も作れない。個別列が必ず同じ下限長という意味ではない。\par

\noindent\textbf{3：適切} 非可逆圧縮は完全復元しない場合があり音声・映像等で多用。。知覚上許せる情報損失を利用してデータ量を減らす。全ての音声/映像が非可逆という断定ではない。\par

\noindent\textbf{4：適切} 可逆圧縮は完全復元しテキスト等で多用。。元ビット列を一意に復元する性質が可逆の定義である。文字1つの変更が意味を変えるテキストでは重要。\par

\noindent\textbf{5：適切} JPEGは静止画、MPEGは動画の代表的圧縮方式。。代表用途として正しい。JPEG/MPEGは複数の方式・規格を持つ名称であり、全方式を一律に可逆/非可逆とは扱わない。\par

\paragraph{再利用する考え方} 符号化の目的、情報量と物理的データ長、条件付き時間と平均時間を区別する。可逆性・平均符号長の限界と普遍的/適応的圧縮の定義と適用条件を確認する。

\paragraph{共通理論} \hyperref[theory:information:compression]{第\ref*{theory:information:compression}節}, \hyperref[theory:information:entropy]{第\ref*{theory:information:entropy}節}, \hyperref[theory:information:universal]{第\ref*{theory:information:universal}節}

\par 参考：\cite{info_deflate,info_lossless,info_universal,info_jpeg,info_mpeg}。

\paragraph{関連問題} \hyperref[q:r02-reg-basic-2-1]{R02-REG-BASIC-2-1}

\par 難易度 2（校準初稿）、優先度未評価。

\subsection{令和7年度 問2-2}
\label{q:r07-reg-basic-2-2}

\noindent\textbf{原典：} sources/R07/basic.pdf、PDF物理5頁。

\textbf{公式正答：} 4（single）。

\paragraph{問題の要旨} 情報源符号化、通信路符号化、分離原理、普遍符号化、多元情報理論を区別する。

\Needspace{4\baselineskip}\paragraph{確認した条件・記述（独自要約）}\mbox{}\par

\noindent\textbf{方向：} 最も不適切\par

\noindent\textbf{モデルの範囲：} 単一情報源・単一入力単一出力通信路と多元系を区別\par

\paragraph{問われる核心} 情報源/通信路符号化の目的と分離/普遍性の条件

\paragraph{必要知識} 情報源/通信路符号化の目的と分離/普遍性の条件。圧縮は情報源の冗長性を減らし、通信路符号化は雑音に耐える制御された冗長性を加える。定常無記憶な点対点モデルで十分長いブロックを使うなどの条件では両設計を分離できる。一方、普遍符号化は事前に真の情報源分布を知らない場合に対応する考えである。4はこの名称を取り違える。通信路容量との整合や有限遅延での性能まで無視できるわけではない。

\paragraph{考え方・正答を選ぶ理由（査読待ち初稿）} 圧縮は情報源の冗長性を減らし、通信路符号化は雑音に耐える制御された冗長性を加える。定常無記憶な点対点モデルで十分長いブロックを使うなどの条件では両設計を分離できる。一方、普遍符号化は事前に真の情報源分布を知らない場合に対応する考えである。4はこの名称を取り違える。通信路容量との整合や有限遅延での性能まで無視できるわけではない。

\Needspace{4\baselineskip}\paragraph{選択肢の検討（査読待ち初稿）}\mbox{}\par

\noindent\textbf{1：適切} 情報源系列から通信路入力への変換が符号化、受信から元系列等への変換が復号化。。通信系における符号器/復号器の基本的役割を述べる。非可逆では元に近い系列へ復元する場合もある。\par

\noindent\textbf{2：適切} 雑音なしで情報源長を固定し通信路入力長を小さくする課題が情報源符号化。。既定の情報を短い符号列へ表す圧縮の課題である。\par

\noindent\textbf{3：適切} 一様な無記憶情報源で通信路入力長を固定し情報源長/伝送レートを増す課題が通信路符号化。。固定長の通信路利用で送れるメッセージ数を増し、誤り確率を抑える通信路符号化の課題である。\par

\noindent\textbf{4：不適切} 情報源符号化で通信路の確率特性を逐一考慮しない扱いがユニバーサル符号化。。標準的な情報源と通信路の分離原理の話である。普遍符号化は情報源分布などが未知なクラスに対応する概念であり、通信路特性を考慮しないという定義ではない。\par

\noindent\textbf{5：適切} 相関した複数情報源や複数通信路の圧縮/伝送を多元情報理論で扱う。。複数利用者・相関情報源のモデルを扱う一般的な分野説明である。分離が多元モデルで常に最適という追加主張はない。\par

\paragraph{再利用する考え方} 符号化の目的、情報量と物理的データ長、条件付き時間と平均時間を区別する。情報源/通信路符号化の目的と分離/普遍性の条件の定義と適用条件を確認する。

\paragraph{共通理論} \hyperref[theory:information:separation]{第\ref*{theory:information:separation}節}, \hyperref[theory:information:universal]{第\ref*{theory:information:universal}節}, \hyperref[theory:information:entropy]{第\ref*{theory:information:entropy}節}

\par 参考：\cite{info_separation,info_universal,info_lossless}。

\par 難易度 2（校準初稿）、優先度未評価。

\section{スタック・データ構造}

\subsection{平成26年度 問2-4}
\label{q:h26-reg-basic-2-4}

\noindent\textbf{原典：} sources/H26/basic.pdf、PDF物理7頁。

\textbf{公式正答：} 4（single）。

\paragraph{問題の要旨} 空スタックへ指定された8操作を実行し、最後のPOPの値を求める。

\Needspace{4\baselineskip}\paragraph{確認した条件・記述（独自要約）}\mbox{}\par

\noindent\textbf{initial\_stack：} empty\par

\noindent\textbf{operations：} PUSH 1；PUSH 2；PUSH 3；POP；PUSH 4；PUSH 5；POP；POP\par

\paragraph{問われる核心} データ構造・形式言語

\paragraph{必要知識} LIFO、PUSH/POP、状態追跡

\paragraph{考え方・正答を選ぶ理由（査読待ち初稿）} スタックを底から上へ記録する。操作後は [1]→[1,2]→[1,2,3]→[1,2]→[1,2,4]→[1,2,4,5]→[1,2,4]→[1,2]。最後のPOPで取り出す値は4なので選択肢④。

\Needspace{4\baselineskip}\paragraph{選択肢の検討（査読待ち初稿）}\mbox{}\par

\noindent\textbf{1：不適切} 1。最終POP後も底に残る。\par

\noindent\textbf{2：不適切} 2。最終POP後もスタックに残る。\par

\noindent\textbf{3：不適切} 3。先のPOPですでに取り出される。\par

\noindent\textbf{4：適切・選択対象} 4。5をPOPした後の頂上4を最後にPOPする。\par

\noindent\textbf{5：不適切} 5。最後より一つ前のPOPで取り出される。\par

\paragraph{再利用する考え方} LIFO、PUSH/POP、状態追跡

\paragraph{共通理論} \hyperref[theory:discrete:stack]{第\ref*{theory:discrete:stack}節}

\par 難易度 2（校準初稿）、優先度未評価。

\subsection{令和元年度 問2-6}
\label{q:r01-reg-basic-2-6}

\noindent\textbf{原典：} sources/R01/basic.pdf、PDF物理11頁。

\textbf{公式正答：} 2（single）。

\paragraph{問題の要旨} スタック操作の途中状態を追って最後のPOPの戻り値を求める。

\Needspace{4\baselineskip}\paragraph{確認した条件・記述（独自要約）}\mbox{}\par

\noindent\textbf{判断対象：} 後入れ先出し、PUSH/POPの状態と戻り値。\par

\noindent\textbf{原典確認：} 図・式・全5候補を29保存画像で確認。rawの数式欠落と演算子誤認は別校正JSONで扱う。\par

\paragraph{問われる核心} スタック操作の途中状態を追って最後のPOPの戻り値を求める。

\paragraph{必要知識} 後入れ先出し、PUSH/POPの状態と戻り値。

\paragraph{考え方・正答を選ぶ理由（査読待ち初稿）} 底から[1,2,3,4]となりPOPで4と3を除く。[1,2]へ5をPUSHし、POPで5、最後のPOPで2が返る。残る1と最後に返った2を区別し2を選ぶ。

\Needspace{4\baselineskip}\paragraph{選択肢の検討（査読待ち初稿）}\mbox{}\par

\noindent\textbf{1：不適切} 1。1は最後にスタックへ残る底の値である。\par

\noindent\textbf{2：適切} 2。取り出し列4,3,5,2の最後である。\par

\noindent\textbf{3：不適切} 3。2回目のPOPで既に取り出されている。\par

\noindent\textbf{4：不適切} 4。1回目のPOPの値である。\par

\noindent\textbf{5：不適切} 5。最後から2回目のPOPの値である。\par

\paragraph{再利用する考え方} 定義した対象量、演算順序、保存条件を固定し、途中状態と候補の違いを比較する。

\paragraph{共通理論} \hyperref[theory:discrete:stack]{第\ref*{theory:discrete:stack}節}

\par 難易度 2（校準初稿）、優先度未評価。

\section{標的型攻撃の多層防御・報告・復旧}

\subsection{令和元年度再試験 問2-1}
\label{q:r01-ret-basic-2-1}

\noindent\textbf{原典：} sources/R01\_retest/basic.pdf、PDF物理6頁。

\textbf{公式正答：} 5（single）。

\paragraph{問題の要旨} 感染予防と、感染後に改ざんされたシステムの復旧を区別する。

\Needspace{4\baselineskip}\paragraph{確認した条件・記述（独自要約）}\mbox{}\par

\noindent\textbf{方向：} 最も適切なもの。\par

\noindent\textbf{時点：} R01年度。現在の推奨を出題当時へ無条件に転用しない。\par

\paragraph{問われる核心} 検知・除去・復旧は別の工程

\paragraph{必要知識} マルウェア、最小権限・攻撃面、信頼したバックアップ、なりすまし。

\paragraph{考え方・正答を選ぶ理由（査読待ち初稿）} 5が感染除去と既発生の改ざんからの回復を区別する。検知ソフトを導入したという事実だけでファイルの原状復帰を保証しない。1〜4は予防に逆行する。

\Needspace{4\baselineskip}\paragraph{選択肢の検討（査読待ち初稿）}\mbox{}\par

\noindent\textbf{1：不適切} 短い単純なパスワード。。推測が容易となる。記憶の容易さはパスワード管理等で扱い、短く単純という選択を安全策にしない。\par

\noindent\textbf{2：不適切} 性能のためウイルス対策不要。。性能だけを理由に必要な防御を除外するのは不適切。適切な製品・設定・更新と多層防御を検討する。\par

\noindent\textbf{3：不適切} 実在の企業名なら添付を信用。。表示名や企業名は偽装でき、正規アカウントも侵害され得る。内容・送信経路を確認する。\par

\noindent\textbf{4：不適切} 使わないポートも開く。。不要な待受サービスを増やすと攻撃面が広がる。必要なポートと送信・受信を限定する。\par

\noindent\textbf{5：適切} システム改ざんによって完全修復が難しい感染もある。。マルウェアの検出・除去と破壊・改ざん済みデータの回復は別である。信頼できるバックアップや再構築を要することがある。\par

\paragraph{再利用する考え方} 目的量・脅威・成立条件を定めてから手段を選ぶ。秘匿、改ざん検出、主体認証、再送防止、復旧は別の性質であり、方式名だけで全てを保証しない。

\paragraph{共通理論} \hyperref[theory:network:malware]{第\ref*{theory:network:malware}節}, \hyperref[theory:network:targeted]{第\ref*{theory:network:targeted}節}, \hyperref[theory:network:password]{第\ref*{theory:network:password}節}

\par 参考：\cite{net_target2020,net_ipa_account}。

\paragraph{関連問題} \hyperref[q:h29-reg-basic-2-1]{H29-REG-BASIC-2-1}, \hyperref[q:r07-reg-basic-2-4]{R07-REG-BASIC-2-4}

\par 難易度 2（校準初稿）、優先度未評価。

\subsection{令和2年度 問2-3}
\label{q:r02-reg-basic-2-3}

\noindent\textbf{原典：} sources/R02/basic.pdf、PDF物理8頁。

\textbf{公式正答：} 1（single）。

\paragraph{問題の要旨} 標的型攻撃の入口・多層防御・連絡・実行制御・訓練を判断する。

\Needspace{4\baselineskip}\paragraph{確認した条件・記述（独自要約）}\mbox{}\par

\noindent\textbf{方向：} 最も不適切なもの。\par

\noindent\textbf{時点：} R02年度。現在の推奨を出題当時へ無条件に転用しない。\par

\paragraph{問われる核心} 正規サービスの悪用と多層防御

\paragraph{必要知識} URLと実ファイルの違い、入口・内部・出口、実行許可、組織的報告。

\paragraph{考え方・正答を選ぶ理由（査読待ち初稿）} 1はサービス名の正規性と格納ファイルの安全性を混同している。IPA白書2020は正規オンラインストレージのURLからウイルスを取得させる攻撃を記録する。メール境界検査に加えて端末・実行制限・連絡・教育を組み合わせる。

\Needspace{4\baselineskip}\paragraph{選択肢の検討（査読待ち初稿）}\mbox{}\par

\noindent\textbf{1：不適切} 正規オンラインストレージを確認しメールゲートウェイで検知。。正規のサービスにも攻撃者は悪性ファイルを置ける。本文にURLだけ載る場合、メール内容の検査だけでリンク先の実体を常に検査・検知できない。\par

\noindent\textbf{2：適切（記述の条件内）} 複数対策の多層防御。。入口・端末・ネットワーク・運用等を組み合わせ、単独策が破られた場合も被害低減を図る。\par

\noindent\textbf{3：適切（記述の条件内）} 組織内の連絡窓口を周知。。受信・誤操作時に早期報告できる既知の窓口は調査・封じ込めへつながる。\par

\noindent\textbf{4：適切（記述の条件内）} 実行可能ファイルの制限。。許可した実行対象へ限定することで悪性実行の経路を減らす。万能検知ではなく業務許可の設計・更新が必要。\par

\noindent\textbf{5：適切（記述の条件内）} 模擬メールの定期訓練。。利用者が不審メールを認識・報告する練習となる。組織の承認・個人の扱い・連絡手順を含む適切な訓練として行う。\par

\paragraph{再利用する考え方} 目的量・脅威・成立条件を定めてから手段を選ぶ。秘匿、改ざん検出、主体認証、再送防止、復旧は別の性質であり、方式名だけで全てを保証しない。

\paragraph{共通理論} \hyperref[theory:network:targeted]{第\ref*{theory:network:targeted}節}, \hyperref[theory:network:malware]{第\ref*{theory:network:malware}節}, \hyperref[theory:network:sandbox]{第\ref*{theory:network:sandbox}節}

\par 参考：\cite{net_target2020,net_target2023}。

\paragraph{関連問題} \hyperref[q:r07-reg-basic-2-4]{R07-REG-BASIC-2-4}

\par 難易度 2（校準初稿）、優先度未評価。

\subsection{令和4年度 問2-1}
\label{q:r04-reg-basic-2-1}

\noindent\textbf{原典：} sources/R04/basic.pdf、PDF物理5頁。

\textbf{公式正答：} 1（single）。

\par\textbf{検討未確定：} 肢3は比較する端末管理・運用・利用者条件を省き「オフィスより危険性が高まる」とする。一般的な留意事項と普遍的比較命題を区別し、公式1・初回1を保持。

\paragraph{問題の要旨} テレワーク時の会議参加制御・VPN・人的攻撃・連絡体制を判断する。

\Needspace{4\baselineskip}\paragraph{確認した条件・記述（独自要約）}\mbox{}\par

\noindent\textbf{方向：} 最も不適切なもの。\par

\noindent\textbf{時点：} R04年度。現在の推奨を出題当時へ無条件に転用しない。\par

\paragraph{問われる核心} 通信路の秘匿と会議参加認証の違い

\paragraph{必要知識} 認可・認証、待機室、脆弱性管理、ショルダーサーフィン、連絡手順。

\paragraph{考え方・正答を選ぶ理由（査読待ち初稿）} 1は案内の秘匿と参加認証の混同である。IPAの2020年Web会議資料は安全配付と参加者確認・認証を併記する。3は運用条件に依存するリスク説明としてのみ理解し、無条件の相対リスク命題は査読対象へ残す。

\Needspace{4\baselineskip}\paragraph{選択肢の検討（査読待ち初稿）}\mbox{}\par

\noindent\textbf{1：不適切} 安全な経路で会議URLを送れば意図しない参加を防げる。。案内の安全な配付だけでなく参加者認証・待機室等が必要。正規参加者からの流出や共有などもありURL送付だけで十分とはいえない。\par

\noindent\textbf{2：適切} VPN製品の脆弱性情報を常に収集。。対策適用・接続管理のための情報収集が重要である。収集だけで修正済みとはならない。\par

\noindent\textbf{3：条件付きの傾向・全称未確定} オフィスよりフィッシング被害リスクが高まる。。相談・監視・端末管理が弱まる運用では危険性が増し得る。一律に全テレワークの危険度が全オフィスより高いとはいえず、原典に比較条件がない。\par

\noindent\textbf{4：適切} 第三者の多い場所での業務を避ける。。のぞき見・盗聴等の機会を減らす。接続暗号化は周囲からの視聴や会話聴取を防ぐものではない。\par

\noindent\textbf{5：適切} 連絡先不明はインシデントの被害拡大へつながる。。遅れた報告・封じ込めにより攻撃継続や拡大の可能性がある。事前に連絡先と手順を確認する。\par

\paragraph{再利用する考え方} 目的量・脅威・成立条件を定めてから手段を選ぶ。秘匿、改ざん検出、主体認証、再送防止、復旧は別の性質であり、方式名だけで全てを保証しない。

\paragraph{共通理論} \hyperref[theory:network:telework]{第\ref*{theory:network:telework}節}, \hyperref[theory:network:tls]{第\ref*{theory:network:tls}節}, \hyperref[theory:network:targeted]{第\ref*{theory:network:targeted}節}

\par 参考：\cite{net_meeting2020,net_telework}。

\paragraph{関連問題} \hyperref[q:r07-reg-basic-2-4]{R07-REG-BASIC-2-4}

\par 難易度 2（校準初稿）、優先度未評価。

\subsection{令和7年度 問2-4}
\label{q:r07-reg-basic-2-4}

\noindent\textbf{原典：} sources/R07/basic.pdf、PDF物理6頁。

\textbf{公式正答：} 5（single）。

\paragraph{問題の要旨} 多層防御・連絡・実行制御・知見共有と、外部解析への送信リスクを比較する。

\Needspace{4\baselineskip}\paragraph{確認した条件・記述（独自要約）}\mbox{}\par

\noindent\textbf{方向：} 最も不適切なもの。\par

\noindent\textbf{時点：} R07年度。現在の推奨を出題当時へ無条件に転用しない。\par

\paragraph{問われる核心} 標的型防御と検査先への情報開示

\paragraph{必要知識} 多層防御、許可リスト、窓口、攻撃指標、外部解析の情報取扱い。

\paragraph{考え方・正答を選ぶ理由（査読待ち初稿）} 5は検知のための送信それ自体が機密流出の経路となる。IPA白書2023の公開該当記述もオンライン検査・サンドボックスへのアップロードによる意図しない漏洩を警告する。承認された閉域解析等と無差別な外部投入を区別する。組織間の4の手口共有は、機密を管理した共有である。

\Needspace{4\baselineskip}\paragraph{選択肢の検討（査読待ち初稿）}\mbox{}\par

\noindent\textbf{1：適切（記述の条件内）} 複数策を多層的に組合せ。。異なる地点・仕組みの防御を組み合わせて侵入・実行・被害拡大を減らす。多数の同じ製品を置くだけでは独立した防御にならない。\par

\noindent\textbf{2：適切（記述の条件内）} 不審メールの組織内窓口を周知。。早期の情報収集と対応へつなぐ。利用者が勝手に危険ファイルを他者へ転送する指示ではない。\par

\noindent\textbf{3：適切（記述の条件内）} 実行可能ファイルを制限。。業務に必要な許可対象へ限定することで悪性コード実行を減らす。署名・パス・ハッシュ等の運用条件を管理する。\par

\noindent\textbf{4：適切（記述の条件内）} 攻撃手口を組織間で共有。。攻撃指標・手口の適切な共有は他組織の検知へ役立つ。機密内容や個人情報をそのまま公開する意味ではない。\par

\noindent\textbf{5：不適切} オンラインの検査・サンドボックスへ添付を積極投入。。公開の解析サービスは提出ファイルを保持・共有する場合があり、機密添付を送ると情報漏洩を起こし得る。利用規約・契約・組織の承認・対象情報の性質を確認せず積極投入するのは不適切。\par

\paragraph{再利用する考え方} 目的量・脅威・成立条件を定めてから手段を選ぶ。秘匿、改ざん検出、主体認証、再送防止、復旧は別の性質であり、方式名だけで全てを保証しない。

\paragraph{共通理論} \hyperref[theory:network:targeted]{第\ref*{theory:network:targeted}節}, \hyperref[theory:network:sandbox]{第\ref*{theory:network:sandbox}節}

\par 参考：\cite{net_target2020,net_target2023}。

\paragraph{関連問題} \hyperref[q:r01-ret-basic-2-1]{R01-RET-BASIC-2-1}, \hyperref[q:r02-reg-basic-2-3]{R02-REG-BASIC-2-3}, \hyperref[q:r04-reg-basic-2-1]{R04-REG-BASIC-2-1}

\par 難易度 2（校準初稿）、優先度未評価。

\section{TCPの窓・ACK・輻輳制御}

\subsection{令和7年度 問2-3}
\label{q:r07-reg-basic-2-3}

\noindent\textbf{原典：} sources/R07/basic.pdf、PDF物理6頁。

\textbf{公式正答：} 3（single）。

\par\textbf{検討未確定：} 原典1は出題前2001年のECN、2はRFC5681/6928の初期窓、5はFlightSizeと閾値下限を省く。旧来の1MSS基本模型の意図と2025年TCP全体の命題を区別する。初回3・公式3を保持、他肢の無条件適切性は未確定。

\paragraph{問題の要旨} TCPのACK単位・RTT単位の窓増加と損失時制御を比較する。

\Needspace{4\baselineskip}\paragraph{確認した条件・記述（独自要約）}\mbox{}\par

\noindent\textbf{方向：} 最も不適切なもの。\par

\noindent\textbf{時点：} R07年度。現在の推奨を出題当時へ無条件に転用しない。\par

\paragraph{問われる核心} 窓増加の単位・ECN・損失窓

\paragraph{必要知識} cwndとrwnd、MSS、RTT、ACK、Reno、ECN、FlightSize、初期窓と損失窓。

\paragraph{考え方・正答を選ぶ理由（査読待ち初稿）} 3の時間尺度が誤っている。1RTTで窓W分のセグメントがACKされ各ACKで1MSS増える理想模型なら\(W_{n+1}\simeq2W_n\)である。輻輳回避は\(W_{n+1}\simeq W_n+\mathrm{MSS}\)の加算増加と区別する。遅延ACK・累積ACK・ABC等で係数は変わる。1・2・5の普遍性は出題時点に既存RFCとも合わないため、公式3を保持しながら原典疑義を明示する。

\Needspace{4\baselineskip}\paragraph{選択肢の検討（査読待ち初稿）}\mbox{}\par

\noindent\textbf{1：旧模型限定・ECNに疑義} ネットワーク層以下から混雑情報は直接得られない。。損失やACKを用いる旧来の説明に限れば理解できるが、ECNではIPヘッダのCEとTCPのECE等で明示的に混雑を通知する。直接情報なしを現行TCP全般へ広げない。\par

\noindent\textbf{2：旧模型限定・初期窓に疑義} 開始時cwnd=\allowbreak{}1MSS。。基本模型の初期値としては使えるが、RFC5681はMSSに応じ2〜4MSSの初期上限を許し、実験的RFC6928は任意のIW10も扱う。常に1MSSではない。\par

\noindent\textbf{3：不適切} ACKごと+MSSなのでスロースタートは線形。。ACKが1セグメントごと返る模型では窓内のACK数も増えるので1RTTあたり概ね倍増となる。ACK回数に対する線形と時間に対する指数的増加を取り違えている。\par

\noindent\textbf{4：Reno型として適切} 輻輳回避は概ね1RTTあたり+1MSS。。Reno型の加算増加におけるRTT単位の説明として妥当。全ての現行TCP輻輳アルゴリズムをこの式へ固定しない。\par

\noindent\textbf{5：模型意図・FlightSizeと下限に疑義} タイムアウトで閾値=cwnd/2、cwnd=\allowbreak{}1MSS。。基本模型では意図を理解できるがRFC5681の閾値上限はmax(FlightSize/2,2SMSS)。未確認データ量とcwndは別で、下限もある。タイムアウトで損失窓へ戻す点は妥当。\par

\paragraph{再利用する考え方} 目的量・脅威・成立条件を定めてから手段を選ぶ。秘匿、改ざん検出、主体認証、再送防止、復旧は別の性質であり、方式名だけで全てを保証しない。

\paragraph{共通理論} \hyperref[theory:network:tcp]{第\ref*{theory:network:tcp}節}, \hyperref[theory:network:ecn]{第\ref*{theory:network:ecn}節}

\par 参考：\cite{net_tcp2009,net_ecn2001,net_iw2013}。

\par 難易度 3（校準初稿）、優先度未評価。


\chapter{解析}
\section{背景：物理モデルを計算できる形へ移す}
解析では、実際の現象を場・ベクトル・方程式として表し、解の性質や近似の誤差を調べる。流れの保存則には発散、局所的な循環には回転、温度やポテンシャルの変化には勾配を用いる。座標変換は形の複雑な領域を計算しやすくし、行列は連立方程式や変換をまとめる。数値解析では連続な微分・積分を有限個の値から計算する。

物理専攻の読者には数学そのものより、\emph{何を入力して何を返す演算か、条件はどこで使うか}の確認が重要である。勾配はベクトル、勾配の長さと発散はスカラー、三次元の回転はベクトルである。微分前に点を代入しない。変換の向きと行・列の順序を、問題の式に合わせて読む。

\section{内積・外積と直交射影}
\label{theory:analysis:projection}
正規直交基底では$a\cdot b=\sum_i a_i b_i$であり、内積はスカラーとなる。仕事$F\cdot\Delta x$は変位方向の力成分を取り出す例である。三次元の外積$a\times b$は両者に垂直で、向きは右手系で定まり、$a\times b=-b\times a$である。平行な場合はゼロだが、それだけを根拠に交換的とはいえない。

$a$を非零の$b$に平行な$p$と垂直な$q$に分けるとき、$p=tb$、$q=a-tb$とおく。$q\cdot b=0$から
\[
 t=\frac{a\cdot b}{b\cdot b},\qquad
 p=\frac{a\cdot b}{b\cdot b}b,\qquad q=a-p
\]
となる。$b$が単位ベクトルなら分母は1である。$q\cdot b=0$だけで候補を選ぶと、$a-q$が$b$の倍数でない候補も残る。二条件を同時に満たす必要がある。

原点を通る直線への垂足の\emph{位置}は$p$であり、点$a$から垂足へ向かう\emph{変位}は$p-a$である。始点と終点を逆にすると符号が反転する。これは、点から面・線への最短距離や力の成分分解にも使える。

\section{勾配と最急上昇率}
\label{theory:analysis:gradient}
スカラー場$f$の小さな変化は$df=\nabla f\cdot dx$で表せる。単位方向$e$への方向微分は$\nabla f\cdot e$であり、コーシー・シュワルツの不等式から
\[
 \nabla f\cdot e\le |\nabla f|\,|e|=|\nabla f|
\]
となる。$\nabla f\neq0$なら$e=\nabla f/|\nabla f|$で等号が成立し、勾配方向が最急上昇、勾配のノルムがその上昇率である。勾配と方向微分の対応は大学講義の整理とも一致する\parencite{mit_chain2007}。

二変数なら$\nabla f=(f_x,f_y)$、$|\nabla f|=\sqrt{f_x^2+f_y^2}$である。求める量が「長さ」なら成分の組を答えず平方根まで取る。$\nabla f=0$の点では一次の方向微分が全方向で0であり、極値の性質には二階微分などが必要となる。

\section{発散と回転：保存と循環}
\label{theory:analysis:divcurl}
右手系の直交座標でベクトル場$V=(V_x,V_y,V_z)$に対し
\[
 \nabla\cdot V=\frac{\partial V_x}{\partial x}
 +\frac{\partial V_y}{\partial y}+\frac{\partial V_z}{\partial z},
\]
\[
 \nabla\times V=
 \left(\frac{\partial V_z}{\partial y}-\frac{\partial V_y}{\partial z},
 \frac{\partial V_x}{\partial z}-\frac{\partial V_z}{\partial x},
 \frac{\partial V_y}{\partial x}-\frac{\partial V_x}{\partial y}\right).
\]
発散は成分と同じ軸の微分を\emph{足す}。回転は別の軸の微分を\emph{引く}。これらの出力の型と座標依存性は材料工学の数学講義でも明示される\parencite{mit_vectorcalc2005}。上式は直交するデカルト座標の成分式であり、円柱・球座標へ無変更で使わない。

発散定理により、小体積から外向きに出る流束の総量を体積で割った極限が発散になる。流速を$u$とすると質量保存は
\[
 \partial_t\rho+\nabla\cdot(\rho u)=0.
\]
密度が流れに沿って不変な非圧縮性条件では$\nabla\cdot u=0$となる。「平面の全ての点で成立」という設問では、式がある曲線上でだけゼロになる候補を除く。

二次元の$V=(u,v)$を三次元へ$(u,v,0)$と埋め込むと回転の面外成分は$\partial_xv-\partial_yu$となる。過去問がこのスカラーをrotと定義した場合は、その定義に従う。三次元での回転ベクトルや発散と取り違えない。

\section{連鎖律：ヤコビ行列の向きを式から決める}
\label{theory:analysis:chain}
$x=x(r,s),y=y(r,s)$と$f=f(x,y)$に対し、全微分へ代入すると
\[
 df=f_x(x_rdr+x_sds)+f_y(y_rdr+y_sds)
\]
である。$dr,ds$の係数を比べれば
\[
 \begin{pmatrix}f_r\\f_s\end{pmatrix}
 =\begin{pmatrix}x_r&y_r\\x_s&y_s\end{pmatrix}
 \begin{pmatrix}f_x\\f_y\end{pmatrix}.
\]
一方、$f=f(x,y),g=g(x,y)$を変数とする$u(f,g)$なら
\[
 \begin{pmatrix}u_x\\u_y\end{pmatrix}
 =\begin{pmatrix}f_x&g_x\\f_y&g_y\end{pmatrix}
 \begin{pmatrix}u_f\\u_g\end{pmatrix}.
\]
ヤコビ行列という語だけから配置を決めず、各行がどの微分を返すかを確認する。連鎖律は微分記号を分数として無条件に約分する操作ではない\parencite{mit_chain2007}。

関係式$f'=1+f^2$と$f(0)$から高階導関数を求める場合も、連鎖律と積の微分を順に使う。解$f$を全区間で求めることと、一点の導関数を求めることは別の仕事である。

\section{行列式と面積要素の変数変換}
\label{theory:analysis:jacobian}
局所的に微分可能で正則な変換$(u,v)\mapsto(x,y)$では、小長方形が平行四辺形へ写り、その面積倍率は行列式の\emph{絶対値}になる。従って
\[
 dxdy=\left|\det\frac{\partial(x,y)}{\partial(u,v)}\right|dudv.
\]
向き付きの行列式と、負にならない面積の倍率を分ける。MITの変数変換講義でも面積で絶対値が必要な理由をこの幾何から説明する\parencite{mit_jacobian2007}。

例えば$x=u+v,y=uv$の微分行列は$\left(\begin{smallmatrix}1&1\\v&u\end{smallmatrix}\right)$で、行列式は$u-v$、面積倍率は$|u-v|$である。また$u,v$を交換すると同じ$x,y$になるため、積分の置換には一対一となる枝を選ぶか重複を数える処理が必要である。平成28年度の設問では符号付き行列式の候補はあるが、面積要素へ用いる符号・領域条件が不足するため疑義を保持する。

累次積分では領域も先に読む。$0\le x\le1,0\le y\le\sqrt{1-x^2}$は四分の一円である。被積分関数が$x$なら、まず$y$について積分して$\int_0^1x\sqrt{1-x^2}\,dx$とできる。被積分関数1の領域面積と、$x$の重みを持つ一次モーメントは異なる。

\section{逆行列：積で確かめ、三角構造を使う}
\label{theory:analysis:inverse}
逆行列は$AA^{-1}=A^{-1}A=I$を満たす行列である。小さい行列の随伴行列による計算と、一般の連立方程式としての意味は大学講義に整理される\parencite{mit_inverse2007}。2行2列なら
\[
 \begin{pmatrix}a&b\\c&d\end{pmatrix}^{-1}
 =\frac{1}{ad-bc}\begin{pmatrix}d&-b\\-c&a\end{pmatrix},
 \qquad ad-bc\neq0.
\]
非対角成分の符号だけでなく、分母の符号も確認する。

単位下三角行列では逆も同じ形とおき、積の非対角成分を順に0へ合わせれば
\[
 \begin{pmatrix}1&0&0\\a&1&0\\b&c&1\end{pmatrix}^{-1}
 =\begin{pmatrix}1&0&0\\-a&1&0\\ac-b&-c&1\end{pmatrix}.
\]
第3行第1列には二段の結合項$ac$が現れる。上三角の問題でも同じように、掛け算の全経路を足す。元の非対角成分を一律に符号反転するだけでは逆にならない。

実際の大規模な線形方程式$Ax=b$では、逆行列を明示的に作る代わりに分解と代入で$x$を求めることが多い。ここでの小行列の計算は写像と検算を理解するために位置付ける。

\section{数値微分：テイラー展開で次数と誤差を読む}
\label{theory:analysis:difference}
等間隔の$x_i$、幅$h>0$に対して、滑らかな関数の展開は
\[
 f_{i\pm1}=f_i\pm hf_i'+\frac{h^2}{2}f_i''
 \pm\frac{h^3}{6}f_i'''+\frac{h^4}{24}f_i^{(4)}+\cdots.
\]
差を取ると一階中心差分、和から中央を引くと二階中心差分となる。
\[
 f_i'=\frac{f_{i+1}-f_{i-1}}{2h}+O(h^2),\qquad
 f_i''=\frac{f_{i+1}-2f_i+f_{i-1}}{h^2}+O(h^2).
\]
前進・後退の一階差分は通常$O(h)$、三点の片側一階差分$(3f_i-4f_{i-1}+f_{i-2})/(2h)$は$O(h^2)$である。名称を知らなくても定数・一次・二次関数を入れれば、多くの誤った候補を除ける。二階微分の分母は長さの二乗であり、次元も検算になる。

幅を小さくすると打切り誤差は減るが、近い値の差を小さい幅で割るため丸め誤差の影響は強くなり得る。数値微分の「正しい差分式」は滑らかな関数に対する近似であり、有限の幅で任意の関数へ恒等的に正確という意味ではない。

\section{数値積分：補間とモーメント条件}
\label{theory:analysis:quadrature}
求積式を$\sum_jw_jf(x_j)$とすると、多項式$1,x,x^2,\ldots$へ正確となるための条件は
\[
 \sum_jw_jx_j^k=\int_a^b x^k\,dx
\]
である。定数関数には重み和が区間幅となる必要がある。対称区間では奇関数部分は相殺されるため、偶数次数の条件を確認しやすい。これは式の係数を暗記するだけでなく、似た誤答式を検証する方法となる。

$a,b,c$が等間隔$h$の三点なら、二次補間多項式を積分するシンプソン法は
\[
 \int_a^cf(x)\,dx\approx\frac h3[f(a)+4f(b)+f(c)].
\]
全幅を$H=c-a=2h$と書けば係数は$H/6$である。記号の定義を取り違えない。滑らかな関数の一パネル誤差は$-h^5f^{(4)}(\xi)/90$であり、三次以下の多項式には正確となる\parencite{nist_quadrature}。

等間隔に縛られず、区間$[-1,1]$の二点を$\pm t$、等しい重みを$w$とすると、定数と二次のモーメント条件から$2w=2$、$2wt^2=2/3$となる。従って$w=1,t=1/\sqrt3$である。奇数次数は対称性で相殺されるので三次以下に正確な二点ガウス求積
\[
 \int_{-1}^1f(x)\,dx=f(-1/\sqrt3)+f(1/\sqrt3)
 \quad (\deg f\le3)
\]
が得られる。評価点が正しくても重みが違えば別の式である。

近似手法が指定された問題では、その手法で計算した値を答える。例えば$f(x)=1/(x+3)$を$[-1,1]$で一パネルのシンプソン法にかけると$25/36\approx0.69444$であり、厳密値$\ln2\approx0.693147$と違う。両者に最も近い候補が同じでも、導出を置き換えない。

\section{固有振動：復元力のモデル化}
\subsection{背景}
支持条件や接続を変えると、同じ材質・同じ部品でも振動特性は変わる。まず変位に対する復元力を求め、慣性との組合せを考える。
\subsection{理論と重要式}
一自由度の線形系では$m\ddot{x}+k_{\rm eff}x=0$から$\omega_0^2=k_{\rm eff}/m$。
質量が同じなら、有効剛性だけで固有振動数の大小を比較できる。
直列では同じ力を受けるばねの変形が加算される。並列では同じ変位を受けるばねの力が加算される。
\subsection{技術的意味}
左右の固定壁につながる質点を右へ変位させると、左のばねの伸びと右のばねの圧縮はともに左向きの復元力を生む。
図の左右という見かけから力を相殺してしまわない。
\subsection{一般の多自由度系への接続}
$M\ddot{q}+Kq=0$から$(K-\omega^2M)\phi=0$が得られる。
ただし過去問の主題がばね・はりの振動比較なら主分類は振動・力学とし、一般化固有値問題は横断参照とする。


\section{数値解析：誤差を発生箇所で分ける}
\label{theory:numerics:errors}
物理モデルを計算する際には、材料や境界条件を単純化するモデル誤差、連続な場を格子や多項式で表す離散化誤差、方程式を途中で解くのを止める反復誤差、有限の桁で演算する丸め誤差がある。実験と計算の差をすべてメッシュの粗さへ帰してはいけない。同じ評価量を固定し、どの操作で結果が変わるかを切り分ける。

例えば形状を変えず格子幅を半分にして応力が変わるなら離散化の影響が疑われる。格子を固定して停止閾値を厳しくすると値が変わるなら反復が不足している。倍精度へ変える操作は仮数の精度を増すが、モデルの誤りや切り捨てた高次項を復元しない。反復回数が多いから閾値を緩めるという操作は、速度と引換えに許容誤差を増す操作である。

\subsection{有限精度：桁落ちと情報落ち}
\label{theory:numerics:floating}
浮動小数点数は指数と有限桁の仮数で値を表す。大きな数へ小さい数を足すと、指数を揃えた小さい側の桁が保持範囲外へ落ちることがある。近い二数の差では共通の大きい部分が消え、各入力に既に含まれた誤差が差の値に対して大きくなる。これは異なる現象である。正確に表された二数の減算が常に新しい大誤差を作るわけではない\parencite{goldberg1991}。

実数の結合則は浮動小数点演算の同じ結果を保証しない。通常の倍精度で$10^{16}+1$の1が失われる場合、$(10^{16}+1)-10^{16}$と$(10^{16}-10^{16})+1$は異なる。桁落ちを避ける式変形の例は
\[
 \sqrt{x+1}-\sqrt{x}=\frac{1}{\sqrt{x+1}+\sqrt{x}}\quad(x\ge0)
\]
である。大きい$x$で左辺を近い数の差として直接計算するより、右辺は差による増幅を避けられる。ただしオーバーフローや入力精度など別の限界は残る。

\subsection{差分式：整合性と精度を別に調べる}
\label{theory:numerics:diff}
前節\ref{theory:analysis:difference}のTaylor展開から、同じ一階導関数を近似しても前進・後退差分の誤差は$O(h)$、中心差分は$O(h^2)$となる。二次精度の後退差分も
\[
 \frac{3f_i-4f_{i-1}+f_{i-2}}{2h}=f'_i+O(h^2)
\]
と導ける。定数項の係数和は$3-4+1=0$、一次項は$2hf'_i$、二次項は消える。補間多項式の微分から差分と余項を得る整理もある\parencite{nist_differentiation}。

式が目的の導関数に近づくことが整合性であり、誤差が減る次数が精度である。$(f_{i+1}-2f_i+f_{i-1})/h$は$hf_i''+O(h^3)$であって、一階微分に整合しない。分母を$h^2$に直しても二階微分である。検算には$f=1,x,x^2$を代入するとよい。格子を過度に細かくすると近い値の差の丸め誤差が増し、打切り誤差の低下だけから最終精度を断定できない。

\subsection{補間：与えられた点と元の関数を区別する}
\label{theory:numerics:interpolation}
相異なる$n+1$点を通る次数$n$以下の多項式は一意である。二つあれば差の多項式は$n+1$個の根を持つが次数は$n$以下であるため、差は零となる。Lagrange形式では
\[
 p(x)=\sum_{j=0}^{n}f(x_j)\prod_{k\ne j}\frac{x-x_k}{x_j-x_k}.
\]
各基底は自分の節点で1、他の節点で0となる\parencite{nist_interpolation}。三点なら$p=ax^2+bx+c$へ代入して解く方法も使える。補間は与えられた点を正確に通すが、元の関数全体を正確に復元するとは限らない。区間外の外挿では、補間値の一致だけで信頼性を判断しない。

\subsection{積分領域：三角形の境界を先に書く}
\label{theory:numerics:integration}
頂点$(0,0),(a,0),(0,a)$の三角形は$0\le x\le a,\ 0\le y\le a-x$である。長方形$[0,a]^2$へ置き換えると領域を増やす。一次関数の例では
\[
 \int_0^a\int_0^{a-x}(x+y)\,dy\,dx
 =\int_0^a\left[x(a-x)+\frac{(a-x)^2}{2}\right]dx
 =\frac{a^3}{3}.
\]
独立の幾何検算として、一次関数の三角形内の平均値は重心での値$2a/3$に等しく、面積$a^2/2$との積も$a^3/3$となる。単位も長さの三乗で確認できる。

\subsection{求積：重み和と奇偶の検算}
\label{theory:numerics:quadrature}
前節\ref{theory:analysis:quadrature}のモーメント条件を使うと、$[-1,1]$の一次関数$ax+b$に正確な公式は重み和2かつ$\sum_jw_jx_j=0$を満たす。端点の台形則、中央点則、三点のSimpson則はいずれも満たす。台形則の一般関数に対する余項には二階導関数が現れる\parencite{nist_quadrature}。

例えば端点重みを$1/4$ずつ、中央重みを1とすると重み和は$3/2$であり、定数$b$を$3b/2$へ誤る。$b=0$なら偶然ゼロへ一致するが、全ての一次関数に正確な一般式ではない。特定の値で一致することと、ある関数族に正確な公式であることを分ける。

\subsection{Newton法：更新の符号と終了条件}
\label{theory:numerics:newton}
現在値$x_n$で接線を引き、その零点を次の値とする。$f(x_n)+f'(x_n)(x_{n+1}-x_n)=0$から
\[
 x_{n+1}=x_n-\frac{f(x_n)}{f'(x_n)}.
\]
分母は導関数であり積ではない。単根近傍の条件が満たされると二次収束するが、初期値を問わず収束する保証はなく、$f'(x_n)=0$ではこの更新を使えない\parencite{nist_newton}。

フローチャートが$x\leftarrow x-\Delta x$なら$\Delta x=f/f'$、$x\leftarrow x+\Delta x$なら$\Delta x=-f/f'$である。小さいステップで終了する分岐には$|\Delta x|<\varepsilon$を合わせる。実際の計算ではステップが小さくても残差が十分小さいとは限らないため、$|f(x)|$、相対ステップ、最大反復数、導関数の小ささを別に確認する。試験の簡略図に合わせることと、実務の停止条件が十分であることを分ける。

\section{有限要素法：場を節点値で近似する}
\label{theory:fem:background}
連続体の変位を有限個の節点値で表し、材料則とつり合いを離散方程式へ移す。変位を$u_h=\sum_iN_i u_i$、微分演算を$L$とすると、ひずみは$\varepsilon=Lu_h=Bq$、線形弾性の応力は$\sigma=D\varepsilon$である。形状関数の微分を含む$B$を用い、要素剛性は$K_e=\int_{V_e}B^TDB\,dV$となる。この補間から剛性行列への関係はMIT講義でも整理される\parencite{mit_roy_fea1999}。以下の具体的な形状関数と幾何計算は独自に導出する。

\subsection{一次三角形：一定ひずみの理由}
\label{theory:fem:triangle}
非退化の三頂点を持つ一次三角形では、各変位成分を$u=a+bx+cy$と表す。微分は$u_x=b,u_y=c$であるため、微小ひずみは要素内で一定となる。均質な線形弾性なら応力も一定である。6節点の二次三角形には$x^2,xy,y^2$の項が加わり、微分は一般に一次関数となる。

これは任意の節点変位に対する要素の性質である。二次三角形でも二次項がゼロとなる特別な場ならひずみは一定になり得る。「一定ひずみ要素」を選ぶ設問で、その特殊例だけを理由に高次要素まで選ばない。退化三角形や不均質な材料では、一定ひずみから一定応力への移行にも条件が要る。

\subsection{双一次四角形：交差項が残る}
\label{theory:fem:q4}
自然座標の四隅を$(-1,-1),(1,-1),(1,1),(-1,1)$とすると
\[
 (N_1,N_2,N_3,N_4)=\frac14\bigl((1-\xi)(1-\eta),(1+\xi)(1-\eta),
 (1+\xi)(1+\eta),(1-\xi)(1+\eta)\bigr).
\]
各節点で対応する関数だけ1となり、和は1である。展開した$\xi,\eta,\xi\eta$の係数ベクトルは順に
\[
 \tfrac14(-1,1,1,-1),\quad\tfrac14(-1,-1,1,1),\quad
 \tfrac14(1,-1,1,-1).
\]
非定数の各係数の和が0という検算もできる。

長方形の物理座標なら$u=a+bx+cy+dxy$となり、$u_x=b+dy,u_y=c+dx$には位置依存が残る。一般四角形では座標写像のJacobianを介して微分し、ひずみはさらに形状に依存する。四隅4節点というだけで要素内一定応力と判断しない。

\subsection{変位の連続と応力の跳び}
\label{theory:fem:continuity}
適合する変位法では隣接要素が同じ辺の節点変位を共有する。一次の辺補間が同じなので辺上の変位も一致する。しかし法線方向の微分までは一般に一致せず、微分から作るひずみや応力に跳びが生じ得る。一様ひずみの場では隣接応力が一致する例もあるため、「常に不連続」という全称は強すぎる。

表示用に節点応力を平均すると滑らかな色分布になることがあるが、平均処理が近似精度の保証を与えるわけではない。変位、積分点応力、平均応力のどれを比較したかを明示する。

\subsection{メッシュ：局所応力に上下保証はない}
\label{theory:fem:mesh}
細分化するh改良は空間解像度を増し、高次化するp改良は近似関数の表現力を増す。滑らかな解、健全な形状、適切な積分や境界条件の下で有効となる。高勾配域へ要素を配置する考え方はMITの例とも対応する\parencite{mit_roy_fea1999}。

線形弾性の適合変位法における全体のエネルギーの傾向から、任意の点の応力が上界であるとはいえない。粗いメッシュでは応力集中を過小評価することもある。「変形が小さいので応力は常に高く安全側」という結論は成立しない。異なる荷重では重要な勾配の位置も変わり、一つの荷重で妥当だった分割を無検証で使い回さない。

複数のメッシュで同じ変位、反力、位置の応力を比較する。尖った角や点荷重に理想化した特異点ではピーク応力が細分ごとに増す場合があり、有限値への収束を当然視しない。要素の歪みや退化、丸め誤差、材料モデルの誤差は別に確認する。

\subsection{面積座標：重心・内心・外心を分ける}
\label{theory:fem:area}
三角形ABC内の点Pに対し、$S_A=S_{PBC},S_B=S_{PCA},S_C=S_{PAB}$とおき、$\lambda_i=S_i/S$と定める。非負で和は1となる。頂点Aでは$(1,0,0)$、重心では$(1/3,1/3,1/3)$である。面積座標は一次三角形の形状関数にもなる。

内心は三辺からの距離$r$が等しく、$S_A=r|BC|/2$などとなる。従って座標は向かい辺の長さの比である。$AB:BC:CA=3:4:5$なら$(\lambda_A,\lambda_B,\lambda_C)=(4,5,3)/12$となる。Bが直角で、外心は斜辺ACの中点なので$(1/2,0,1/2)$である。重心を内心や外心と混同せず、図の見た目より辺長比を優先する。三角形外の点に拡張する場合には符号付き面積が必要となる。

\subsection{二層の熱伝導：並列と直列の加算量}
\label{theory:fem:thermal}
定常伝導の平板では$Q=kA\Delta T/d$であり、熱流は温度勾配に比例する\parencite{openstax_thermal_conduction}。接触熱抵抗がなく、厚さ$l_1,l_2$の層に平行な流れでは温度勾配が共通であり、各断面の熱流を足す。層を横切る流れでは熱流が共通であり、各層の温度差を足す。$L=l_1+l_2$として
\[
 k_\parallel=\frac{k_1l_1+k_2l_2}{L},\qquad
 k_\perp=\frac{L}{l_1/k_1+l_2/k_2}.
\]
前者は厚さで重み付けした算術平均、後者は調和平均である。正の熱伝導率では両者とも各材料値の間に入り、$k_\parallel\ge k_\perp$となる。界面熱抵抗や過渡現象があれば、この簡単な二式のモデルをそのまま適用しない。

厚さを材料に逆に対応させた誤答式は等厚の特殊値で正しい式と一致する。試験では一般式を導く目的と読みつつ、その一致と時間条件の省略を解説の疑義へ残す。

\section{軸変形は釣合と適合を同時に満たす}
\label{theory:analysis:fixedbar}
\label{theory:analysis:compositebar}
\label{theory:analysis:parallelbar}
\label{theory:analysis:extension}
一様な棒の軸剛性は \(EA/l\)、伸びは \(Nl/(EA)\) である。力の釣合だけで決まらない反力には、固定端間の総変位がゼロという適合条件を加える。途中の荷重点を動かすと、左の伸びと右の縮みは同じ変位となり、両区間は荷重点に対して復元剛性を並列に提供する。右向きの外力に対する支持反力を右向き正で表せば、いずれも負になる。

二本の等長棒を回転しない剛体端板で結ぶ場合は伸び、従ってひずみが共通である。応力比は弾性係数比、軸力比は \(EA\) の比となる。端板が回転する条件や長さが違う場合はこの結論をそのまま使わない。計算は N と mm なら \(E\) を N/mm\(^2\) にそろえる。\cite{mit_material_elastic2000}

\section{横収縮と非圧縮極限}
\label{theory:analysis:hooke}
\label{theory:analysis:incompressible}
\label{theory:analysis:thermal}
等方線形弾性体の構成則は
\[
 \varepsilon_x=\frac{\sigma_x-\nu(\sigma_y+\sigma_z)}{E},\qquad
 \varepsilon_v=\frac{1-2\nu}{E}(\sigma_x+\sigma_y+\sigma_z).
\]
引張の横収縮は負符号で入る。任意の応力状態でも体積を変えない極限は \(\nu\to1/2\) である。このとき \(K=E/[3(1-2\nu)]\) は発散し、通常の有限圧縮性の逆構成則へそのまま代入する操作とは区別する。熱を受ける完全固定棒では全ひずみ \(\sigma/E+\alpha\Delta T=0\) から温度上昇に対する圧縮応力を得る。部分拘束では支持の変形も加えて解く。\cite{mit_material_elastic2000}

\section{弾性エネルギーの基準と安定性}
\label{theory:analysis:axialenergy}
\label{theory:analysis:potential}
\label{theory:analysis:verticalenergy}
線形ばねの蓄積仕事は力変位線図の三角形面積である。
\[
 U=\frac12 P\delta=\frac{P^2l}{2EA},\qquad
 U_s=\frac12ku^2.
\]
ばね自身の \(u\) は自然長から測る。鉛直ばねで平衡からの変位を \(x\) とすると \(u=mg/k+x\) であり、ばねエネルギーと重力エネルギーを合計して平衡値を引いた振動エネルギーは \(kx^2/2\) となる。ばねだけの全量とは違う。全ポテンシャル \(\Pi=ku^2/2-mgu\) は平衡で停留し、二階微分が正なので安定な最小となる。\cite{analysis57_shm}

\section{曲げ方向の高さと片持ち撓み}
\label{theory:analysis:section}
\label{theory:analysis:cantilever}
矩形断面の曲げ軸に関する断面二次モーメントは \(I=bh^3/12\)。断面積が等しくても高さ \(h\) が違えば曲げ剛性は違う。同じ \(E\) と荷重に対し小さい \(I\) が大きい撓みを与える。先端荷重を受ける細長い片持ち梁では
\[
 \delta=\frac{PL^3}{3EI}.
\]
他の支持や分布荷重なら係数も式も変わる。寸法比を求めるときは各梁の長さと断面を別々に代入し、無次元比で比較する。\cite{struct_mit_beam,analysis57_beam}

\section{孔縁の自由表面にも接線応力が残る}
\label{theory:analysis:inglis}
自由表面は応力テンソル全体がゼロという条件ではなく、法線 \(n\) に対する表面力 \(\sigma n=0\) である。楕円孔の水平端では法線が水平なので法線垂直応力とせん断応力が消え、接線垂直応力は残る。十分広い等方線形弾性板で遠方荷重が楕円長軸に直交する場合、孔端の応力集中係数は \(1+2a/b\)。円孔の3へ戻ることを検算に使える。有限幅、塑性や鋭い亀裂先端には別のモデルが必要となる。\cite{analysis57_inglis}

\section{小振動の剛性を配置から組み立てる}
\label{theory:analysis:spring}
\label{theory:analysis:springgravity}
\label{theory:analysis:stringmass}
運動を平衡からの変位で表すと \(m\ddot x+k_{\rm eff}x=0\) なので \(\omega=\sqrt{k_{\rm eff}/m}\)、\(f=\omega/(2\pi)\)。ばね並列は剛性の和、直列はコンプライアンスの和になる。質点を左右の壁へつなぐ二ばねは、一方が伸び他方が縮んでも復元力が同じ方向へ働く。重力の一定成分は線形ばねの平衡位置をずらすが剛性を変えない。

張った弦の中央に質点を置く場合は横変位に対する張力成分を小角近似で加える。各半長を \(L\) とすると \(k_{\rm eff}=2T/L\)。全長と半長を取り違えず、質量の変更も二乗振動数の比へ含める。\cite{analysis57_shm}

\section{梁振動と柱座屈の境界条件}
\label{theory:analysis:beamfreq}
\label{theory:analysis:buckling}
一様Euler--Bernoulli梁の曲げ固有角振動数は \(\omega_1=\beta_1^2\sqrt{EI/(\rho A)}/L^2\)。片持ちの特性式は \(\cosh\beta\cos\beta=-1\)、両端固定は \(=1\) で、それぞれ非零の最小根は約1.875と4.730である。単純支持では \(\beta_1=\pi\)。剛体モードのゼロ根を曲げ基本振動と混同しない。\cite{analysis57_beam}

Euler座屈は \(P_{\rm cr}=\pi^2EI/(KL)^2\) に支持を有効長係数 \(K\) として含める。ピン・ピンは1、固定・自由は2、固定・ピンは約0.699。後者の特性根 \(z\) は \(\tan z=z\)、\(K=\pi/z\) である。これは真直な細長柱の弾性理想モデルで、初期不整や短柱圧壊の設計耐力を保証しない。

\section{節点の力から伸びへ進む}
\label{theory:analysis:truss}
\label{theory:analysis:trussforce}
ピン節点の二力部材は軸方向だけに力を伝える。節点の水平・鉛直力を釣り合わせて部材力 \(N_i\) を求め、その後に \(\delta_i=N_il_i/(EA_i)\) を使う。軸力比と伸び比は長さ・断面の因子が違う。角度が水平基準か鉛直基準かを先に定め、正負によって引張・圧縮を分ける。\cite{mit_material_elastic2000}

\section{回転慣性と動力の二つの関係}
\label{theory:analysis:inertia}
\label{theory:analysis:torque}
\label{theory:analysis:geartorque}
\label{theory:analysis:pendulum}
回転運動は \(\tau=I\alpha\)、回転動力は \(P=\tau\omega\)。角速度を時間で変える問題と減速比の問題では、同じ式だけでは解けない。入力速度が出力の \(n\) 倍、効率1なら入力トルクは出力の \(1/n\)。換算慣性 \(I/n^2\) を使うときの角加速度も入力側へ換算する。

一様薄板の面外重心軸には垂直軸定理 \(I_z=I_x+I_y\) が使える。面外に厚さがある物体には同じ式を無条件に使わない。物理振子の微小角運動は \(I\ddot\theta+Mgd\theta=0\)、従って周期は \(2\pi\sqrt{I/(Mgd)}\)。支点から重心までの距離 \(d\) と物体の慣性モーメントを別々に求める。

\section{終端・定常の意味を保存則で読む}
\label{theory:analysis:terminal}
\label{theory:analysis:bernoulli}
\label{theory:analysis:motional}
線形抵抗落下は \(m\dot v=mg-kv\)。終端では \(\dot v=0\) で、速度依存が二乗抵抗とは違う。水平な定常・非圧縮・非粘性流れには \(A_av_a=A_bv_b\) と \(p+\rho v^2/2=\text{一定}\) を併用する。断面積比が速度二乗へ入ること、粘性損失や高さの差を原典条件から外せるかを確認する。

運動起電力の棒では \(\mathcal E=BLv\)、\(I=\mathcal E/R\)、制動力 \(IBL\)。定常で駆動力と制動が釣り合う。抵抗が大きいほど電流が小さくなるため制動が弱まるという極限も検算に使える。

\section{抵抗網は線の形でなく節点の接続で解く}
\label{theory:analysis:network}
\label{theory:analysis:tetra}
同じ二節点を結ぶ枝は並列、途中で分岐のない経路が直列である。交差線の接続点を確認し、対称性で同電位と証明できる節点をまとめる。任意網には内部節点で
\[
 \sum_j\frac{V_i-V_j}{R_{ij}}=0
\]
を解き、端子電圧を1とすれば \(R_{\rm eq}=1/I_{\rm terminal}\)。一辺の中点を端子にする場合はその辺を半抵抗二本に分ける。四面体の頂点・対辺中点・対辺の両中点は異なる端子条件であり、同じ正四面体だから同抵抗とはならない。\cite{analysis57_resistors}

\section{拡散の減衰と幾何の重心}
\label{theory:analysis:heat}
\label{theory:analysis:centroid}
正の拡散係数の通常の熱伝導は、固定温度又は断熱の安定境界下で空間固有値 \(\lambda\ge0\) に対し \(e^{-\kappa\lambda t}\) の応答を持つ。時間一階微分という形式だけでは非振動性を証明できず、連立一階方程式が回転運動を表す例もある。微分次数、演算子、境界条件を一組で扱う。

面積重心は \(\bar x=A^{-1}\int x\,dA\)、\(\bar y=A^{-1}\int y\,dA\)。一様面、線、体積の重心では積分の測度が違う。対称性で二座標が等しい場合も、値は面積一次モーメントから求める。四分円面と円弧を同じ式で扱わない。

\subsection{平成23年度 問3-1}
\label{q:h23-reg-basic-3-1}

\noindent\textbf{原典：} sources/H23/basic.pdf、PDF物理8頁。

\textbf{公式正答：} 2（single）。

\par\textit{原典本文の校正と独自解説は未作成。分類のみの登録項目。}

\subsection{平成23年度 問3-2}
\label{q:h23-reg-basic-3-2}

\noindent\textbf{原典：} sources/H23/basic.pdf、PDF物理9頁。

\textbf{公式正答：} 1（single）。

\par\textit{原典本文の校正と独自解説は未作成。分類のみの登録項目。}

\subsection{平成23年度 問3-3}
\label{q:h23-reg-basic-3-3}

\noindent\textbf{原典：} sources/H23/basic.pdf、PDF物理10頁。

\textbf{公式正答：} 5（single）。

\par\textit{原典本文の校正と独自解説は未作成。分類のみの登録項目。}

\subsection{平成23年度 問3-4}
\label{q:h23-reg-basic-3-4}

\noindent\textbf{原典：} sources/H23/basic.pdf、PDF物理11頁。

\textbf{公式正答：} 3（single）。

\par\textit{原典本文の校正と独自解説は未作成。分類のみの登録項目。}

\subsection{平成23年度 問3-5}
\label{q:h23-reg-basic-3-5}

\noindent\textbf{原典：} sources/H23/basic.pdf、PDF物理11頁。

\textbf{公式正答：} 1（single）。

\par\textit{原典本文の校正と独自解説は未作成。分類のみの登録項目。}

\subsection{平成24年度 問3-1}
\label{q:h24-reg-basic-3-1}

\noindent\textbf{原典：} sources/H24/basic.pdf、PDF物理8頁。

\textbf{公式正答：} 1（single）。

\par\textit{原典本文の校正と独自解説は未作成。分類のみの登録項目。}

\subsection{平成24年度 問3-2}
\label{q:h24-reg-basic-3-2}

\noindent\textbf{原典：} sources/H24/basic.pdf、PDF物理8頁。

\textbf{公式正答：} 2（single）。

\paragraph{問題の要旨} ばね接続から有効剛性を求め、3系の固有振動数を比較する。

\paragraph{必要知識} ばねの直列・並列合成、微小変位の復元力、固有角振動数

\paragraph{考え方（査読待ち初稿）} aでは復元力が2kxなので有効剛性2k。bでは変形量が加算され1/keff=1/k+1/k、よってkeff=k/2。cも左右の復元力が同方向に合計され2kx。質量が同じため固有角振動数はsqrt(keff/m)で比較でき、bのみ最小。正答②。

\paragraph{他の選択肢（査読待ち初稿）} ①aと③cはいずれも有効剛性2k。④は最小でないaを、⑤は最小でないcを含む。cの左右の力は打ち消し合わず、変位に対して同じ復元方向に働く。

\paragraph{関連問題} \hyperref[q:h26-reg-basic-3-5]{H26-REG-BASIC-3-5}

\par 難易度 2（校準初稿）、優先度未評価。

\subsection{平成24年度 問3-3}
\label{q:h24-reg-basic-3-3}

\noindent\textbf{原典：} sources/H24/basic.pdf、PDF物理9頁。

\textbf{公式正答：} 1（single）。

\par\textit{原典本文の校正と独自解説は未作成。分類のみの登録項目。}

\subsection{平成24年度 問3-4}
\label{q:h24-reg-basic-3-4}

\noindent\textbf{原典：} sources/H24/basic.pdf、PDF物理10頁。

\textbf{公式正答：} 4（single）。

\par\textit{原典本文の校正と独自解説は未作成。分類のみの登録項目。}

\subsection{平成24年度 問3-5}
\label{q:h24-reg-basic-3-5}

\noindent\textbf{原典：} sources/H24/basic.pdf、PDF物理11頁。

\textbf{公式正答：} 4（single）。

\par\textit{原典本文の校正と独自解説は未作成。分類のみの登録項目。}

\subsection{平成25年度 問3-1}
\label{q:h25-reg-basic-3-1}

\noindent\textbf{原典：} sources/H25/basic.pdf、PDF物理10頁。

\textbf{公式正答：} 5（single）。

\par\textit{原典本文の校正と独自解説は未作成。分類のみの登録項目。}

\subsection{平成25年度 問3-2}
\label{q:h25-reg-basic-3-2}

\noindent\textbf{原典：} sources/H25/basic.pdf、PDF物理11頁。

\textbf{公式正答：} 3（single）。

\par\textit{原典本文の校正と独自解説は未作成。分類のみの登録項目。}

\subsection{平成25年度 問3-3}
\label{q:h25-reg-basic-3-3}

\noindent\textbf{原典：} sources/H25/basic.pdf、PDF物理12頁。

\textbf{公式正答：} 1（single）。

\par\textit{原典本文の校正と独自解説は未作成。分類のみの登録項目。}

\subsection{平成25年度 問3-4}
\label{q:h25-reg-basic-3-4}

\noindent\textbf{原典：} sources/H25/basic.pdf、PDF物理13頁。

\textbf{公式正答：} 5（single）。

\par\textit{原典本文の校正と独自解説は未作成。分類のみの登録項目。}

\subsection{平成25年度 問3-5}
\label{q:h25-reg-basic-3-5}

\noindent\textbf{原典：} sources/H25/basic.pdf、PDF物理13頁。

\textbf{公式正答：} 2（single）。

\par\textit{原典本文の校正と独自解説は未作成。分類のみの登録項目。}

\subsection{平成25年度 問3-6}
\label{q:h25-reg-basic-3-6}

\noindent\textbf{原典：} sources/H25/basic.pdf、PDF物理14頁。

\textbf{公式正答：} 4（single）。

\par\textit{原典本文の校正と独自解説は未作成。分類のみの登録項目。}

\subsection{平成26年度 問3-1}
\label{q:h26-reg-basic-3-1}

\noindent\textbf{原典：} sources/H26/basic.pdf、PDF物理9頁。

\textbf{公式正答：} 2（single）。

\par\textit{原典本文の校正と独自解説は未作成。分類のみの登録項目。}

\subsection{平成26年度 問3-2}
\label{q:h26-reg-basic-3-2}

\noindent\textbf{原典：} sources/H26/basic.pdf、PDF物理9頁。

\textbf{公式正答：} 3（single）。

\par\textit{原典本文の校正と独自解説は未作成。分類のみの登録項目。}

\subsection{平成26年度 問3-3}
\label{q:h26-reg-basic-3-3}

\noindent\textbf{原典：} sources/H26/basic.pdf、PDF物理10頁。

\textbf{公式正答：} 5（single）。

\par\textit{原典本文の校正と独自解説は未作成。分類のみの登録項目。}

\subsection{平成26年度 問3-4}
\label{q:h26-reg-basic-3-4}

\noindent\textbf{原典：} sources/H26/basic.pdf、PDF物理10頁。

\textbf{公式正答：} 1（single）。

\par\textit{原典本文の校正と独自解説は未作成。分類のみの登録項目。}

\subsection{平成26年度 問3-5}
\label{q:h26-reg-basic-3-5}

\noindent\textbf{原典：} sources/H26/basic.pdf、PDF物理11頁。

\textbf{公式正答：} 2（single）。

\par\textit{原典本文の校正と独自解説は未作成。分類のみの登録項目。}

\subsection{平成26年度 問3-6}
\label{q:h26-reg-basic-3-6}

\noindent\textbf{原典：} sources/H26/basic.pdf、PDF物理12頁。

\textbf{公式正答：} 2（single）。

\par\textit{原典本文の校正と独自解説は未作成。分類のみの登録項目。}

\subsection{平成27年度 問3-1}
\label{q:h27-reg-basic-3-1}

\noindent\textbf{原典：} sources/H27/basic.pdf、PDF物理8頁。

\textbf{公式正答：} 2（single）。

\par\textit{原典本文の校正と独自解説は未作成。分類のみの登録項目。}

\subsection{平成27年度 問3-2}
\label{q:h27-reg-basic-3-2}

\noindent\textbf{原典：} sources/H27/basic.pdf、PDF物理8頁。

\textbf{公式正答：} 2（single）。

\par\textit{原典本文の校正と独自解説は未作成。分類のみの登録項目。}

\subsection{平成27年度 問3-3}
\label{q:h27-reg-basic-3-3}

\noindent\textbf{原典：} sources/H27/basic.pdf、PDF物理9頁。

\textbf{公式正答：} 1（single）。

\par\textit{原典本文の校正と独自解説は未作成。分類のみの登録項目。}

\subsection{平成27年度 問3-4}
\label{q:h27-reg-basic-3-4}

\noindent\textbf{原典：} sources/H27/basic.pdf、PDF物理9頁。

\textbf{公式正答：} 3（single）。

\par\textit{原典本文の校正と独自解説は未作成。分類のみの登録項目。}

\subsection{平成27年度 問3-5}
\label{q:h27-reg-basic-3-5}

\noindent\textbf{原典：} sources/H27/basic.pdf、PDF物理10頁。

\textbf{公式正答：} 3（single）。

\par\textit{原典本文の校正と独自解説は未作成。分類のみの登録項目。}

\subsection{平成27年度 問3-6}
\label{q:h27-reg-basic-3-6}

\noindent\textbf{原典：} sources/H27/basic.pdf、PDF物理11頁。

\textbf{公式正答：} 4（single）。

\par\textit{原典本文の校正と独自解説は未作成。分類のみの登録項目。}

\subsection{平成28年度 問3-1}
\label{q:h28-reg-basic-3-1}

\noindent\textbf{原典：} sources/H28/basic.pdf、PDF物理12頁。

\textbf{公式正答：} 3（single）。

\par\textit{原典本文の校正と独自解説は未作成。分類のみの登録項目。}

\subsection{平成28年度 問3-2}
\label{q:h28-reg-basic-3-2}

\noindent\textbf{原典：} sources/H28/basic.pdf、PDF物理12頁。

\textbf{公式正答：} 3（single）。

\par\textit{原典本文の校正と独自解説は未作成。分類のみの登録項目。}

\subsection{平成28年度 問3-3}
\label{q:h28-reg-basic-3-3}

\noindent\textbf{原典：} sources/H28/basic.pdf、PDF物理13頁。

\textbf{公式正答：} 2（single）。

\par\textit{原典本文の校正と独自解説は未作成。分類のみの登録項目。}

\subsection{平成28年度 問3-4}
\label{q:h28-reg-basic-3-4}

\noindent\textbf{原典：} sources/H28/basic.pdf、PDF物理14頁。

\textbf{公式正答：} 5（single）。

\par\textit{原典本文の校正と独自解説は未作成。分類のみの登録項目。}

\subsection{平成28年度 問3-5}
\label{q:h28-reg-basic-3-5}

\noindent\textbf{原典：} sources/H28/basic.pdf、PDF物理15頁。

\textbf{公式正答：} 1（single）。

\par\textit{原典本文の校正と独自解説は未作成。分類のみの登録項目。}

\subsection{平成28年度 問3-6}
\label{q:h28-reg-basic-3-6}

\noindent\textbf{原典：} sources/H28/basic.pdf、PDF物理16頁。

\textbf{公式正答：} 4（single）。

\par\textit{原典本文の校正と独自解説は未作成。分類のみの登録項目。}

\subsection{平成29年度 問3-1}
\label{q:h29-reg-basic-3-1}

\noindent\textbf{原典：} sources/H29/basic.pdf、PDF物理11頁。

\textbf{公式正答：} 5（single）。

\par\textit{原典本文の校正と独自解説は未作成。分類のみの登録項目。}

\subsection{平成29年度 問3-2}
\label{q:h29-reg-basic-3-2}

\noindent\textbf{原典：} sources/H29/basic.pdf、PDF物理11頁。

\textbf{公式正答：} 4（single）。

\par\textit{原典本文の校正と独自解説は未作成。分類のみの登録項目。}

\subsection{平成29年度 問3-3}
\label{q:h29-reg-basic-3-3}

\noindent\textbf{原典：} sources/H29/basic.pdf、PDF物理11頁。

\textbf{公式正答：} 4（single）。

\par\textit{原典本文の校正と独自解説は未作成。分類のみの登録項目。}

\subsection{平成29年度 問3-4}
\label{q:h29-reg-basic-3-4}

\noindent\textbf{原典：} sources/H29/basic.pdf、PDF物理12頁。

\textbf{公式正答：} 2（single）。

\par\textit{原典本文の校正と独自解説は未作成。分類のみの登録項目。}

\subsection{平成29年度 問3-5}
\label{q:h29-reg-basic-3-5}

\noindent\textbf{原典：} sources/H29/basic.pdf、PDF物理13頁。

\textbf{公式正答：} 3（single）。

\par\textit{原典本文の校正と独自解説は未作成。分類のみの登録項目。}

\subsection{平成29年度 問3-6}
\label{q:h29-reg-basic-3-6}

\noindent\textbf{原典：} sources/H29/basic.pdf、PDF物理14頁。

\textbf{公式正答：} 5（single）。

\par\textit{原典本文の校正と独自解説は未作成。分類のみの登録項目。}

\subsection{平成30年度 問3-1}
\label{q:h30-reg-basic-3-1}

\noindent\textbf{原典：} sources/H30/basic.pdf、PDF物理11頁。

\textbf{公式正答：} 1（single）。

\par\textit{原典本文の校正と独自解説は未作成。分類のみの登録項目。}

\subsection{平成30年度 問3-2}
\label{q:h30-reg-basic-3-2}

\noindent\textbf{原典：} sources/H30/basic.pdf、PDF物理11頁。

\textbf{公式正答：} 4（single）。

\par\textit{原典本文の校正と独自解説は未作成。分類のみの登録項目。}

\subsection{平成30年度 問3-3}
\label{q:h30-reg-basic-3-3}

\noindent\textbf{原典：} sources/H30/basic.pdf、PDF物理12頁。

\textbf{公式正答：} 2（single）。

\par\textit{原典本文の校正と独自解説は未作成。分類のみの登録項目。}

\subsection{平成30年度 問3-4}
\label{q:h30-reg-basic-3-4}

\noindent\textbf{原典：} sources/H30/basic.pdf、PDF物理13頁。

\textbf{公式正答：} 2（single）。

\par\textit{原典本文の校正と独自解説は未作成。分類のみの登録項目。}

\subsection{平成30年度 問3-5}
\label{q:h30-reg-basic-3-5}

\noindent\textbf{原典：} sources/H30/basic.pdf、PDF物理14頁。

\textbf{公式正答：} 3（single）。

\par\textit{原典本文の校正と独自解説は未作成。分類のみの登録項目。}

\subsection{平成30年度 問3-6}
\label{q:h30-reg-basic-3-6}

\noindent\textbf{原典：} sources/H30/basic.pdf、PDF物理14頁。

\textbf{公式正答：} 3（single）。

\par\textit{原典本文の校正と独自解説は未作成。分類のみの登録項目。}

\subsection{令和元年度 問3-1}
\label{q:r01-reg-basic-3-1}

\noindent\textbf{原典：} sources/R01/basic.pdf、PDF物理12頁。

\textbf{公式正答：} 5（single）。

\par\textit{原典本文の校正と独自解説は未作成。分類のみの登録項目。}

\subsection{令和元年度 問3-2}
\label{q:r01-reg-basic-3-2}

\noindent\textbf{原典：} sources/R01/basic.pdf、PDF物理12頁。

\textbf{公式正答：} 5（single）。

\par\textit{原典本文の校正と独自解説は未作成。分類のみの登録項目。}

\subsection{令和元年度 問3-3}
\label{q:r01-reg-basic-3-3}

\noindent\textbf{原典：} sources/R01/basic.pdf、PDF物理13頁。

\textbf{公式正答：} 1（single）。

\par\textit{原典本文の校正と独自解説は未作成。分類のみの登録項目。}

\subsection{令和元年度 問3-4}
\label{q:r01-reg-basic-3-4}

\noindent\textbf{原典：} sources/R01/basic.pdf、PDF物理13頁。

\textbf{公式正答：} 4（single）。

\par\textit{原典本文の校正と独自解説は未作成。分類のみの登録項目。}

\subsection{令和元年度 問3-5}
\label{q:r01-reg-basic-3-5}

\noindent\textbf{原典：} sources/R01/basic.pdf、PDF物理14頁。

\textbf{公式正答：} 4（single）。

\par\textit{原典本文の校正と独自解説は未作成。分類のみの登録項目。}

\subsection{令和元年度 問3-6}
\label{q:r01-reg-basic-3-6}

\noindent\textbf{原典：} sources/R01/basic.pdf、PDF物理15頁。

\textbf{公式正答：} 3（single）。

\par\textit{原典本文の校正と独自解説は未作成。分類のみの登録項目。}

\subsection{令和元年度再試験 問3-1}
\label{q:r01-ret-basic-3-1}

\noindent\textbf{原典：} sources/R01\_retest/basic.pdf、PDF物理11頁。

\textbf{公式正答：} 5（single）。

\par\textit{原典本文の校正と独自解説は未作成。分類のみの登録項目。}

\subsection{令和元年度再試験 問3-2}
\label{q:r01-ret-basic-3-2}

\noindent\textbf{原典：} sources/R01\_retest/basic.pdf、PDF物理11頁。

\textbf{公式正答：} 2（single）。

\par\textit{原典本文の校正と独自解説は未作成。分類のみの登録項目。}

\subsection{令和元年度再試験 問3-3}
\label{q:r01-ret-basic-3-3}

\noindent\textbf{原典：} sources/R01\_retest/basic.pdf、PDF物理11頁。

\textbf{公式正答：} 5（single）。

\par\textit{原典本文の校正と独自解説は未作成。分類のみの登録項目。}

\subsection{令和元年度再試験 問3-4}
\label{q:r01-ret-basic-3-4}

\noindent\textbf{原典：} sources/R01\_retest/basic.pdf、PDF物理12頁。

\textbf{公式正答：} 5（single）。

\par\textit{原典本文の校正と独自解説は未作成。分類のみの登録項目。}

\subsection{令和元年度再試験 問3-5}
\label{q:r01-ret-basic-3-5}

\noindent\textbf{原典：} sources/R01\_retest/basic.pdf、PDF物理12頁。

\textbf{公式正答：} 4（single）。

\par\textit{原典本文の校正と独自解説は未作成。分類のみの登録項目。}

\subsection{令和元年度再試験 問3-6}
\label{q:r01-ret-basic-3-6}

\noindent\textbf{原典：} sources/R01\_retest/basic.pdf、PDF物理13頁。

\textbf{公式正答：} 2（single）。

\par\textit{原典本文の校正と独自解説は未作成。分類のみの登録項目。}

\subsection{令和2年度 問3-1}
\label{q:r02-reg-basic-3-1}

\noindent\textbf{原典：} sources/R02/basic.pdf、PDF物理11頁。

\textbf{公式正答：} 2（single）。

\par\textit{原典本文の校正と独自解説は未作成。分類のみの登録項目。}

\subsection{令和2年度 問3-2}
\label{q:r02-reg-basic-3-2}

\noindent\textbf{原典：} sources/R02/basic.pdf、PDF物理11頁。

\textbf{公式正答：} 4（single）。

\par\textit{原典本文の校正と独自解説は未作成。分類のみの登録項目。}

\subsection{令和2年度 問3-3}
\label{q:r02-reg-basic-3-3}

\noindent\textbf{原典：} sources/R02/basic.pdf、PDF物理11頁。

\textbf{公式正答：} 3（single）。

\par\textit{原典本文の校正と独自解説は未作成。分類のみの登録項目。}

\subsection{令和2年度 問3-4}
\label{q:r02-reg-basic-3-4}

\noindent\textbf{原典：} sources/R02/basic.pdf、PDF物理12頁。

\textbf{公式正答：} 4（single）。

\par\textit{原典本文の校正と独自解説は未作成。分類のみの登録項目。}

\subsection{令和2年度 問3-5}
\label{q:r02-reg-basic-3-5}

\noindent\textbf{原典：} sources/R02/basic.pdf、PDF物理13頁。

\textbf{公式正答：} 1（single）。

\par\textit{原典本文の校正と独自解説は未作成。分類のみの登録項目。}

\subsection{令和2年度 問3-6}
\label{q:r02-reg-basic-3-6}

\noindent\textbf{原典：} sources/R02/basic.pdf、PDF物理14頁。

\textbf{公式正答：} 5（single）。

\par\textit{原典本文の校正と独自解説は未作成。分類のみの登録項目。}

\subsection{令和3年度 問3-1}
\label{q:r03-reg-basic-3-1}

\noindent\textbf{原典：} sources/R03/basic.pdf、PDF物理11頁。

\textbf{公式正答：} 4（single）。

\par\textit{原典本文の校正と独自解説は未作成。分類のみの登録項目。}

\subsection{令和3年度 問3-2}
\label{q:r03-reg-basic-3-2}

\noindent\textbf{原典：} sources/R03/basic.pdf、PDF物理11頁。

\textbf{公式正答：} 2（single）。

\par\textit{原典本文の校正と独自解説は未作成。分類のみの登録項目。}

\subsection{令和3年度 問3-3}
\label{q:r03-reg-basic-3-3}

\noindent\textbf{原典：} sources/R03/basic.pdf、PDF物理12頁。

\textbf{公式正答：} 1（single）。

\par\textit{原典本文の校正と独自解説は未作成。分類のみの登録項目。}

\subsection{令和3年度 問3-4}
\label{q:r03-reg-basic-3-4}

\noindent\textbf{原典：} sources/R03/basic.pdf、PDF物理13頁。

\textbf{公式正答：} 3（single）。

\par\textit{原典本文の校正と独自解説は未作成。分類のみの登録項目。}

\subsection{令和3年度 問3-5}
\label{q:r03-reg-basic-3-5}

\noindent\textbf{原典：} sources/R03/basic.pdf、PDF物理14頁。

\textbf{公式正答：} 4（single）。

\par\textit{原典本文の校正と独自解説は未作成。分類のみの登録項目。}

\subsection{令和3年度 問3-6}
\label{q:r03-reg-basic-3-6}

\noindent\textbf{原典：} sources/R03/basic.pdf、PDF物理15頁。

\textbf{公式正答：} 5（single）。

\par\textit{原典本文の校正と独自解説は未作成。分類のみの登録項目。}

\subsection{令和4年度 問3-1}
\label{q:r04-reg-basic-3-1}

\noindent\textbf{原典：} sources/R04/basic.pdf、PDF物理10頁。

\textbf{公式正答：} 4（single）。

\par\textit{原典本文の校正と独自解説は未作成。分類のみの登録項目。}

\subsection{令和4年度 問3-2}
\label{q:r04-reg-basic-3-2}

\noindent\textbf{原典：} sources/R04/basic.pdf、PDF物理10頁。

\textbf{公式正答：} 2（single）。

\par\textit{原典本文の校正と独自解説は未作成。分類のみの登録項目。}

\subsection{令和4年度 問3-3}
\label{q:r04-reg-basic-3-3}

\noindent\textbf{原典：} sources/R04/basic.pdf、PDF物理11頁。

\textbf{公式正答：} 4（single）。

\par\textit{原典本文の校正と独自解説は未作成。分類のみの登録項目。}

\subsection{令和4年度 問3-4}
\label{q:r04-reg-basic-3-4}

\noindent\textbf{原典：} sources/R04/basic.pdf、PDF物理11頁。

\textbf{公式正答：} 2（single）。

\par\textit{原典本文の校正と独自解説は未作成。分類のみの登録項目。}

\subsection{令和4年度 問3-5}
\label{q:r04-reg-basic-3-5}

\noindent\textbf{原典：} sources/R04/basic.pdf、PDF物理12頁。

\textbf{公式正答：} 3（single）。

\par\textit{原典本文の校正と独自解説は未作成。分類のみの登録項目。}

\subsection{令和4年度 問3-6}
\label{q:r04-reg-basic-3-6}

\noindent\textbf{原典：} sources/R04/basic.pdf、PDF物理13頁。

\textbf{公式正答：} 5（single）。

\par\textit{原典本文の校正と独自解説は未作成。分類のみの登録項目。}

\subsection{令和5年度 問3-1}
\label{q:r05-reg-basic-3-1}

\noindent\textbf{原典：} sources/R05/basic.pdf、PDF物理9頁。

\textbf{公式正答：} 4（single）。

\par\textit{原典本文の校正と独自解説は未作成。分類のみの登録項目。}

\subsection{令和5年度 問3-2}
\label{q:r05-reg-basic-3-2}

\noindent\textbf{原典：} sources/R05/basic.pdf、PDF物理9頁。

\textbf{公式正答：} 2（single）。

\par\textit{原典本文の校正と独自解説は未作成。分類のみの登録項目。}

\subsection{令和5年度 問3-3}
\label{q:r05-reg-basic-3-3}

\noindent\textbf{原典：} sources/R05/basic.pdf、PDF物理10頁。

\textbf{公式正答：} 1（single）。

\par\textit{原典本文の校正と独自解説は未作成。分類のみの登録項目。}

\subsection{令和5年度 問3-4}
\label{q:r05-reg-basic-3-4}

\noindent\textbf{原典：} sources/R05/basic.pdf、PDF物理10頁。

\textbf{公式正答：} 5（single）。

\par\textit{原典本文の校正と独自解説は未作成。分類のみの登録項目。}

\subsection{令和5年度 問3-5}
\label{q:r05-reg-basic-3-5}

\noindent\textbf{原典：} sources/R05/basic.pdf、PDF物理11頁。

\textbf{公式正答：} 2（single）。

\par\textit{原典本文の校正と独自解説は未作成。分類のみの登録項目。}

\subsection{令和5年度 問3-6}
\label{q:r05-reg-basic-3-6}

\noindent\textbf{原典：} sources/R05/basic.pdf、PDF物理12頁。

\textbf{公式正答：} 3（single）。

\par\textit{原典本文の校正と独自解説は未作成。分類のみの登録項目。}

\subsection{令和6年度 問3-1}
\label{q:r06-reg-basic-3-1}

\noindent\textbf{原典：} sources/R06/basic.pdf、PDF物理10頁。

\textbf{公式正答：} 4（single）。

\par\textit{原典本文の校正と独自解説は未作成。分類のみの登録項目。}

\subsection{令和6年度 問3-2}
\label{q:r06-reg-basic-3-2}

\noindent\textbf{原典：} sources/R06/basic.pdf、PDF物理10頁。

\textbf{公式正答：} 2（single）。

\par\textit{原典本文の校正と独自解説は未作成。分類のみの登録項目。}

\subsection{令和6年度 問3-3}
\label{q:r06-reg-basic-3-3}

\noindent\textbf{原典：} sources/R06/basic.pdf、PDF物理11頁。

\textbf{公式正答：} 5（single）。

\par\textit{原典本文の校正と独自解説は未作成。分類のみの登録項目。}

\subsection{令和6年度 問3-4}
\label{q:r06-reg-basic-3-4}

\noindent\textbf{原典：} sources/R06/basic.pdf、PDF物理12頁。

\textbf{公式正答：} 3（single）。

\par\textit{原典本文の校正と独自解説は未作成。分類のみの登録項目。}

\subsection{令和6年度 問3-5}
\label{q:r06-reg-basic-3-5}

\noindent\textbf{原典：} sources/R06/basic.pdf、PDF物理13頁。

\textbf{公式正答：} 3（single）。

\par\textit{原典本文の校正と独自解説は未作成。分類のみの登録項目。}

\subsection{令和6年度 問3-6}
\label{q:r06-reg-basic-3-6}

\noindent\textbf{原典：} sources/R06/basic.pdf、PDF物理14頁。

\textbf{公式正答：} 3（single）。

\par\textit{原典本文の校正と独自解説は未作成。分類のみの登録項目。}

\subsection{令和7年度 問3-1}
\label{q:r07-reg-basic-3-1}

\noindent\textbf{原典：} sources/R07/basic.pdf、PDF物理8頁。

\textbf{公式正答：} 2（single）。

\par\textit{原典本文の校正と独自解説は未作成。分類のみの登録項目。}

\subsection{令和7年度 問3-2}
\label{q:r07-reg-basic-3-2}

\noindent\textbf{原典：} sources/R07/basic.pdf、PDF物理8頁。

\textbf{公式正答：} 1（single）。

\par\textit{原典本文の校正と独自解説は未作成。分類のみの登録項目。}

\subsection{令和7年度 問3-3}
\label{q:r07-reg-basic-3-3}

\noindent\textbf{原典：} sources/R07/basic.pdf、PDF物理9頁。

\textbf{公式正答：} 5（single）。

\par\textit{原典本文の校正と独自解説は未作成。分類のみの登録項目。}

\subsection{令和7年度 問3-4}
\label{q:r07-reg-basic-3-4}

\noindent\textbf{原典：} sources/R07/basic.pdf、PDF物理10頁。

\textbf{公式正答：} 2（single）。

\par\textit{原典本文の校正と独自解説は未作成。分類のみの登録項目。}

\subsection{令和7年度 問3-5}
\label{q:r07-reg-basic-3-5}

\noindent\textbf{原典：} sources/R07/basic.pdf、PDF物理11頁。

\textbf{公式正答：} 5（single）。

\par\textit{原典本文の校正と独自解説は未作成。分類のみの登録項目。}

\subsection{令和7年度 問3-6}
\label{q:r07-reg-basic-3-6}

\noindent\textbf{原典：} sources/R07/basic.pdf、PDF物理12頁。

\textbf{公式正答：} 1（single）。

\par\textit{原典本文の校正と独自解説は未作成。分類のみの登録項目。}


\chapter{材料・化学・バイオ}
\section{扱う対象と学習の出発点}
材料は組成だけでなく結合・組織・加工履歴で機能が変わる。化学は物質量・電荷・エネルギーの保存から変化を記述し、生命現象では分子が水中の反応系を構成する。まず何を数え、何に対する割合を取るかを固定すると、燃焼・濃度・核種・合金組成へ共通の見通しが得られる。以下は原子・化学量・機能性セラミックスと、分子から細胞の機能へつなぐ生命科学の初稿である。金属組織・物性・物質の性質も共通理論と全問の解説初稿へ接続する。

\section{原子・核種を数える}
\label{theory:basicchem:isotope}
原子番号\(Z\)は陽子数で元素を決め、質量数\(A\)は陽子と中性子の合計である。\({}^{A}_{Z}X\)の中性子数は\(N=A-Z\)、中性原子の電子数は\(Z\)、電荷\(+ke\)のイオンなら\(Z-k\)となる。質量数は核子数であり、実際の原子質量そのものではない。
\begin{equation}N=A-Z,\qquad n_{\mathrm e}=Z\quad\text{（中性原子）}.\end{equation}
同位体は同じ\(Z\)で異なる\(A\)を持つ核種である\autocite{iupac_isotopes}。Cl-35とCl-37は電子17個、中性子18個と20個である。Ca-40とAr-40は同じ\(A\)でも\(Z=20,18\)が違う同重体で、中性子20個と22個である。同素体は単体の結合・構造の違いで、核種の違いではない。
放射壊変で核の構成が変わると、例えば\(\beta^-\)壊変では中性子が陽子へ変わり\(Z\)が1増える。一方、励起核の\(\gamma\)遷移は\(Z\)を変えない。年代測定・トレーサー・診断等の利用を知るとともに、全壊変で別元素になるとは一般化しない。

\section{同じ電子配置と同位体効果}
\label{theory:basicchem:isotope_effect}
同位体の中性原子は核電荷と電子配置が同じなので、主要な結合様式・元素としての性質はほぼ共通である。しかし質量は分子振動へ入り、二原子振動子では
\begin{equation}\omega=\sqrt{k/\mu},\quad\mu=\frac{m_1m_2}{m_1+m_2},\quad E_0=\frac{\hbar\omega}{2}.\end{equation}
\(k\)は結合の力定数、\(\mu\)は換算質量である。質量変化は振動数・零点エネルギーを変え、速度や平衡定数の同位体効果を生む\autocite{iupac_isotope_effect}。試験の「化学的性質はほぼ同じ」という基礎的近似を、同位体効果が厳密に0という主張へ拡張しない。

\section{物質量は何を数える量か}
\label{theory:basicchem:amount}
反応係数は粒子個数の比を与え、実験では質量・気体体積を測る。物質量は両者をつなぐ。指定粒子の個数\(N_{\mathrm{ent}}\)、質量\(w\)、モル質量\(M\)について
\begin{equation}n=\frac{N_{\mathrm{ent}}}{N_{\mathrm A}}=\frac{w}{M},\quad N_{\mathrm A}=6.02214076\times10^{23}\ \mathrm{mol^{-1}}.\end{equation}
水なら分子、NaClなら式単位、合金なら各元素の原子と対象を明記する。\(N\)を核種の中性子数と混同しない。理想気体では\(n=pV/(RT)\)。0℃、約1気圧なら約22.4 L/molだが、別の温度・圧力へそのモル体積を持ち込まない。現行の定数定義と歴史的な測定精度を分け、指定された原子量・丸め値を計算へ用いる。

\section{反応係数と保存則から収率を求める}
\label{theory:basicchem:stoichiometry}
化学反応では通常の原子核は変わらず、元素ごとの原子数が保存される。\(\Delta n_i=\nu_i\xi\)では\(\nu_i\)が生成側を正、反応側を負に取る反応係数、\(\xi\)が反応進行度である。分子の炭素数が\(c\)、完全燃焼で全炭素がCO2になるなら
\begin{equation}n_{\mathrm{CO_2}}=c\frac{w}{M}.\end{equation}
同じ質量\(w\)の比較は\(c/M\)の比較となる。炭素数だけでなく水素・酸素の質量も分母へ入る。部分燃焼、すす生成、酸素不足は前提の外である。
酸素消費量は生成物に依存する。\(2\mathrm{Cu}+\mathrm{O_2}\to2\mathrm{CuO}\)と\(4\mathrm{Cu}+\mathrm{O_2}\to2\mathrm{Cu_2O}\)では必要量が2倍違う。生成物が未指定なら仮定を示し、可能な生成物でも順位が変わるか確認する。

\section{濃度の分母を先に決める}
\label{theory:basicchem:concentration}
濃度には何で割るかを含む。溶質の物質量\(n_B\)、最終溶液体積\(V\)、溶媒質量\(w_{\mathrm{solvent}}\)について
\begin{equation}c_B=\frac{n_B}{V}\quad[\mathrm{mol/L}],\qquad b_B=\frac{n_B}{w_{\mathrm{solvent}}}\quad[\mathrm{mol/kg}].\end{equation}
前者はモル濃度、後者は質量モル濃度である\autocite{iupac_amount_concentration,iupac_molality}。水1 Lへ溶質1 molを加えても最終溶液が1 Lとは限らない。1 mol/Lの調製は溶解後に全体を1 Lへ定容する。体積は温度・混合で変わるため、溶媒kgを分母とする量とは異なる。
質量10\%は溶質100 gと溶媒900 gから1000 gの溶液を作る場合である。質量ppmは質量分率の\(10^6\)倍。1 mg/Lは質量濃度で、その定義がppmではない。密度約1 kg/Lの希薄水溶液なら質量約1 ppmに一致する。近似と密度・基準の省略を無条件の等号へ変えない。

\section{質量分率とモル分率の変換}
\label{theory:basicchem:fraction}
モル分率は全成分の物質量に対する成分\(i\)の割合である\autocite{iupac_amount_fraction}。
\begin{equation}x_i=\frac{n_i}{\sum_jn_j}=\frac{w_i/M_i}{\sum_j w_j/M_j},\qquad\sum_i x_i=1.\end{equation}
合計100 gを仮に取ると、質量百分率をそのままgとして計算できる。仮の全体質量は分子・分母で消える。原子量が小さい成分は同じ質量でより多くの原子を持つので、質量分率からモル分率へ変換したとき増える方向が期待できる。合金では原子、溶液では分子・式単位と、対象を指定する。

\section{酸化数を電荷収支から求める}
\label{theory:basicchem:oxidation}
酸化数は異種原子間の結合をイオン的に近似した電子配分の指標であり、実際の電荷密度そのものではない\autocite{iupac_oxidation_state}。今回の酸・酸化物はH=+1、O=−2とし、総和を全体電荷へ等置する。
\begin{equation}\sum_i a_i s_i=q.\end{equation}
\(a_i\)は原子数、\(s_i\)は原子1個の酸化数、\(q\)は電気素量単位の電荷である。MnO4\({}^-\)なら\(x-8=-1\)から+7、Cr2O7\({}^{2-}\)なら\(2x-14=-2\)からCrは+6である。単体Mnは0、NH3のNは−3となる。2個のCrの合計+12と1原子の+6を分ける。
過酸化物のO=−1、金属水素化物のH=−1等の例外を今回の化合物が含まないことを確認して使う。酸化数増加が酸化、減少が還元であり、酸素が式にあるかだけでは判定しない。

\section{多原子イオンと式単位}
\label{theory:basicchem:ion}
イオンの価数と個数は違う。NH4\({}^+\)、OH\({}^-\)、NO3\({}^-\)は多原子1価、SO4\({}^{2-}\)は多原子2価、PO4\({}^{3-}\)は多原子3価である。S\({}^{2-}\)、Ca\({}^{2+}\)、Pb\({}^{2+}\)は2価でも単原子である。
Pb(NO3)2の括弧外2はNO3\({}^-\)が2個あることを表し、硝酸イオン1個を−2へ変えない。(NH4)2SO4は2個の+1と−2、Ca3(PO4)2は3個の+2と2個の−3が釣り合う。全体の中性から確かめると名称暗記に頼らず判断できる。

\section{結合エネルギーと反応エンタルピー}
\label{theory:basicchem:enthalpy}
結合切断にはエネルギーを与え、生成では放出する。反応物を気相原子へ分け、生成物の結合を作る仮想経路をヘスの法則で合計する\autocite{mit_bond_enthalpy}。
\begin{equation}\Delta H\simeq\sum D_{\mathrm{broken}}-\sum D_{\mathrm{formed}}.\end{equation}
\(D\)は正の切断エンタルピーである。平均結合エンタルピーは化学環境の違いによる近似を含み、液体・固体の反応では相変化等も必要となる。H2+Cl2\(\to\)2HClで\(D_{HH}=436\)、\(D_{ClCl}=243\)、\(D_{HCl}=432\) kJ/molなら\(\Delta H=679-864=-185\) kJ。HClを2 mol生成する基準なので1 mol当たり−92.5 kJである。
発熱は系の\(\Delta H<0\)だが「外へ発生する熱の大きさ」は正で表すこともある。原典の符号指定と候補を読み、平成23年度の放出熱量と平成26年度の発熱負の指定を区別する。

\section{ハロゲンを一つの順位で表さない}
\label{theory:basicchem:halogen}
電気陰性度は結合中の電子を引く傾向、酸化力は電子を受け取る反応の駆動力、酸性度はプロトン移動の平衡、沸点は凝集相から気相へ移る条件である。同じ元素列でも別現象である。電気陰性度と分子の酸化力は概ねF、Cl、Br、Iの順に減少する\autocite{purdue_halogens,rsc_halogen}。
しかし水中のHFは弱酸、HCl・HBr・HIは強酸である。結合切断の費用、陰イオンの安定性、溶媒和を含むので電気陰性度だけではHXの酸性度を決められない\autocite{purdue_acid_strength,yale_hydrohalic}。同じ分析濃度の希薄水溶液でHFのH\({}^+\)が最大という肢は誤り。強酸3種はほぼ完全解離し、水による水平化で単純な濃度差としては区別しにくい。
通常沸点はHF\(>\)HI\(>\)HBr\(>\)HClである\autocite{nist_halogen_properties,msu_hydrogen_bonding}。HFは水素結合のため特に高く、他3種は大きい分子ほど分散力が強くなる。X2の沸点とHXの沸点も別々に読む。

\section{電子セラミックスを入出力から理解する}
\label{theory:basicchem:ceramics}
機能材料では、何を入力すると何が変わるかを対応させる。BaTiO3系誘電体は高い比誘電率により容量を大きくする\autocite{murata_batio3}。\(C=\varepsilon_0\varepsilon_r A/d\)は幾何学と分極応答を分ける。導電率を上げて電荷を通すことは電極間を絶縁して蓄える機能と異なる。
サーミスタは温度で抵抗が変わる素子である\autocite{murata_thermistor}。圧電体は応力・ひずみと電気応答を結び、変形に伴う電荷を電圧へ変換してセンサーとする\autocite{murata_piezo}。熱膨張とは異なる結合項を使う。
バリスタは非線形な電圧–電流特性を使う\autocite{tdk_varistor}。\(I=KV^\alpha\)の近似なら\(V/I=K^{-1}V^{1-\alpha}\)で、\(\alpha>1\)では電圧上昇で実効抵抗が低下する。\(\alpha=1\)の純抵抗と区別し、サージ時に電流を逃がす機能へ利用する。温度抵抗・圧電・誘電・非線形抵抗を別々の入力出力として理解する。

\section{生命現象を組成・区画・情報から見る}
\label{theory:biology:composition}
生命科学では、同じ化学結合を使って反応を選択し、情報を保存・複製し、細胞の境界を維持する仕組みを扱う。細胞は水中の反応系である。糖・脂肪酸・アミノ酸・ヌクレオチドという小分子の群を、エネルギー源、膜材料、高分子の単位として区別すると用語に役割が付く。
典型的な細胞で水が総重量の約70\%という概数を用いる。一方、H・O・C・Nの主要元素やHの約60\%以上という説明は原子の個数を分母とする。水だけなら3個の原子のうち2個がHであるが、質量ではHは約2/18にすぎない。軽い元素の個数割合と質量割合は一致しない。
\begin{equation}
 x_i^{\rm atom}=\frac{N_i}{\sum_jN_j},\qquad
 w_i=\frac{N_i M_i}{\sum_jN_jM_j}.
\end{equation}
ここで\(N_i\)は元素\(i\)の原子数、\(M_i\)はそのモル質量である。同じ比例定数で原子数を物質量へ換算できるため、式の比では定数が消える。
通常の細胞の有機成分・巨大分子の重量ではタンパク質の寄与が大きい。実験研究が比較する細菌・酵母・哺乳類細胞の組成も同一ではなく、培養状態で変わる\autocite{bio_cell_mass}。典型例の数字は全種・全組織の固定比率ではない。平成26年度問4-6の「すべての生物で同様」はこの範囲に疑義があり、公式の主要元素の誤りと別に追跡する。

\section{膜脂質：水を避ける鎖が境界を作る}
\label{theory:biology:membrane}
リン脂質は親水的な頭部と疎水的な炭化水素鎖を併せ持つ。水中で頭部を水へ向け、鎖を内側に集めることで二重層を作る。通常のグリセロリン脂質では脂肪酸がエステル結合でグリセロールへつながる。脂肪酸の全結合形式がエステルという定義ではなく、スフィンゴ脂質のアミド結合、古細菌のエーテル脂質を区別する。
膜は完全に固定した壁ではなく、脂質等が面内を移動する境界である。温度、鎖長、不飽和度が充填と運動性を変える。同系統の飽和鎖では長い鎖ほど相互作用が増し、融点が高くなる。cis二重結合は鎖を折り曲げ、密な充填を妨げ、同炭素数の飽和鎖より流動性を増す方向に働く\autocite{bio_membrane}。trans型や混合物の全挙動をこの一つの近似へ押し込めない。
温度上昇そのものは流動化する。環境温度が変わる細菌は長鎖化・飽和化等の組成変更でこれを相殺する場合がある。設問の「増える場合がある」を、全細菌で同じ割合になるという法則へ変えない。膜の適切な流動性は輸送・酵素機能を維持する条件であり、単に流動性を最大にすればよいわけではない。

\section{アミノ酸：共通骨格と側鎖を分ける}
\label{theory:biology:amino}
タンパク質を形作る標準20種のアミノ酸は、α炭素にアミノ基、カルボキシ基、H、種類ごとのR側鎖を持つ。水中の電荷状態はpHに依存し、中性付近ではアミノ基の正電荷とカルボキシ基の負電荷を併せ持つ双性イオンとなる場合が多い\autocite{mit_amino}。化学式に未解離のNH2・COOHを書くことと、溶液内の主な状態を区別する。
ロイシン・イソロイシン等の炭化水素側鎖、フェニルアラニンの芳香族側鎖は非極性で疎水的である。システインのRはチオールを持ち、メチオニンのRはチオエーテルを持つ。両者に硫黄があることと、両者が同じS-S架橋を作ることは別である。
グリシンはR=Hでα炭素が不斉でない。他の標準タンパク質アミノ酸は主にLである。D/Lは立体配置の系列を表し、旋光の正負を表す記号ではない。「特殊なものを除く」という出題条件を確認すれば、標準20種と修飾・特殊アミノ酸を混ぜず判断できる。

\section{タンパク質：配列から構造と機能へ}
\label{theory:biology:protein}
隣接するアミノ酸のカルボキシ基とアミノ基に由来する共有のペプチド結合が主鎖を作る。配列を一次構造、局所的なαらせん・βシート等を二次構造、全体の折畳みを三次構造、複数鎖の集合を四次構造という\autocite{mit_protein}。階層の名前と、その維持に働く相互作用の名前を分ける。
配列は翻訳する核酸配列に規定される。一方、機能する立体構造には溶液条件、シャペロン、結合相手等も関わるので、配列の指定を全環境で同じ活性が保証されるという意味にしない。主鎖のペプチドとCys間のジスルフィドは共有結合である。静電相互作用、水素結合、分散力等の非共有相互作用も折畳みへ寄与する。
疎水効果は水中で非極性面を露出する自由エネルギー費用に関係する。通常の水溶性球状タンパク質では非極性側鎖が内部へ偏るが、膜の脂質に接する面などは違う。疎水性相互作用を新しい主鎖の共有結合とはしない。誤答肢の「ジスルフィド等の非共有結合」は、構造への寄与が正しくても結合分類が誤りである。

\section{等電点：正味0と各基の無電荷は違う}
\label{theory:biology:charge}
タンパク質は末端と側鎖に複数の解離基を持ち、それぞれのプロトン化率がpHで変わる。正味電荷0となるpHが等電点pIである。通常は\(pH<pI\)で正、\(pH>pI\)で負となる。酸塩基解離でpH上昇に伴いプロトンを失う方向を追うと符号を理解できる。電気泳動や分離では溶液のpHとpIを比較する\autocite{mit_protein}。
正味0は正負の電荷が相殺することであり、全ての官能基が無電荷という条件ではない。多くの解離基を持つタンパク質のpIを、二つのpKaの平均だけで常に求められるとはしない。

\section{DNA：鎖の内部と鎖同士の会合}
\label{theory:biology:dna}
ヌクレオチドは塩基・糖・リン酸からなる。DNAの糖はデオキシリボース、標準塩基はA・T・G・Cである。RNAはリボースと通常Uを用いる。各鎖でヌクレオチドをつなぐ糖リン酸主鎖はホスホジエステル共有結合であり、二重らせんで塩基は内側、主鎖は外側となる。
A-TとG-Cが相補的に会合し、2鎖は逆平行である\autocite{nhgri_dna}。水素結合は選択性へ、塩基スタッキング等も安定化へ寄与する。会合が解けることと主鎖が分解することを別に扱う。
通常のDNAポリメラーゼは鋳型とプライマーの自由な3端を使ってヌクレオチドを付ける。したがって新生鎖は\(5'\to3'\)、鋳型の読みは逆方向となる。「DNAが二重鎖」の典型説明を、一本鎖DNAを持つウイルス等が存在しないという主張にはしない。

\section{塩基組成：一方の鎖に全体の対称性を課さない}
\label{theory:biology:base_fraction}
一方の鎖の割合を\(a,t,g,c\)、相補鎖を星付きとすれば
\begin{equation}
 a+t+g+c=1,\qquad
 a^{*}=t,\ t^{*}=a,\ g^{*}=c,\ c^{*}=g.
\end{equation}
相補鎖は各塩基の行先を入れ替える写像である。一方の鎖だけに\(a=t\)や\(g=c\)を課す根拠はない。二重鎖全体では対を合算するためその等式が成り立つ。
C20\%、T35\%なら同じ鎖のA+G45\%、相補鎖A+G55\%である。G25\%、A15\%なら相補鎖C25\%、T15\%、G+A60\%である。個々の未指定値は確定しないが、非負と総和1の制約で不可能な肢は排除できる。例えば後者の同じ鎖G+C90\%はC65\%を要求し、C+T60\%に反する。

\section{DNA融解と吸光度：配列を切らずに会合を解く}
\label{theory:biology:melting}
熱又は適切な化学条件で二重鎖は一本鎖へ解離する。通常の変性で主鎖のホスホジエステル結合は残る。変性に伴い約260 nmの紫外吸収が増す増色効果を用いて融解曲線を観測する\autocite{bu_dna_melting}。二重鎖の塩基スタッキングの変化が吸光度へ影響する。
転移の中点を融解温度\(T_m\)という。二状態近似であれば半分が解離した点に対応する。長い不均一なDNAは領域ごとに転移し得るので、分子全体の個数が常に厳密に半数という説明へ固定しない。他条件を揃えるとGC含量が高いほど一般に\(T_m\)は高いが、塩濃度、鎖長、配列、ミスマッチ等も影響する。
相補鎖が存在し、濃度・温度・時間が適切なら冷却で再会合する。これがアニーリングである。物質の共有結合分解を逆に作り直す操作とは異なる。この区別がPCRの熱変性とプライマー会合の理解へ直接つながる。

\section{遺伝暗号：多対一の対応と三文字の読み枠}
\label{theory:biology:code}
転写はDNAの情報からRNAを作り、翻訳ではmRNAを3塩基ずつ読んでアミノ酸鎖を作る。4文字の3文字列は\(4^3=64\)通りであるが、入力の個数と出力種数は一対一である必要はない。標準暗号では61が20種のアミノ酸、3が終止に対応する\autocite{nhgri_codon}。複数のコドンが同じアミノ酸を指定することを縮重という。
残る44文字列が消失した、44の修飾アミノ酸を指定する、転写で20へ書き換わるという説明は違う。縮重があることから第1塩基が無意味とは導けない。第3位置の置換で同義となる場合も全てではない。開始位置が読み枠を定め、同じ文字列でも区切りが変われば指定配列が変わる。
この標準表は特殊な再符号化やミトコンドリア等の暗号差を別にして用いる。DNAのコドンという設問表記は、転写後のmRNAでTに対応してUが使われることと区別する。

\section{突然変異：文字を替えることと区切りをずらすこと}
\label{theory:biology:mutation}
塩基置換は一文字を替え、挿入・欠失は文字数を変える。通常の翻訳領域で3の倍数でない文字数の変化は下流の読み枠をずらすフレームシフトとなる\autocite{nhgri_frameshift}。非翻訳領域の全挿入を、同じ読み枠の式だけで扱わない。
置換は同義となる場合、アミノ酸が変わるミスセンス、途中の終止となるナンセンス等を生む。代表的なHbSのGlu→Valは置換であって1塩基欠失ではない\autocite{nhgri_sickle}。中立性は機能や選択への効果を論じる軸で、ナンセンスと同義の名前ではない。
常染色体の通常の二倍体・完全潜性のモデルでは、正常な対立遺伝子が一方にあると形質が隠れる。両方の状態を見る必要がある。半接合、不完全な優性、組織モザイク等は前提が違う。X線・UVや特定の化学物質はDNA損傷・変異の誘発源となり得るが、物質群の全てを一律に変異原としない。

\section{PCR：温度で分け、目印から増やす}
\label{theory:biology:pcr}
目的領域の両端へ対応する2本のプライマーを用い、熱変性、アニーリング、伸長を繰り返す。DNA主鎖を切るのではなく鎖間会合を解き、プライマーの3端から合成する。プライマー自身の配列は産物の一部となる。
反復加熱に耐えるDNAポリメラーゼを用いる。アニーリング温度が低すぎると非特異的会合が増え、高すぎると目的会合も減り得る。実際の温度・時間はプライマー配列、塩条件、酵素製品等に応じて設定する\autocite{neb_pcr}。伸長速度を\(u\)塩基/秒、標的長を\(L\)塩基とすれば、最小限の時間は概ね\(L/u\)であり、長い配列ほど短くするという判断は逆である。
完全な複製を各回仮定する理想モデルでは
\begin{equation}
 N_{n+1}=2N_n,\qquad N_n=N_0\,2^n.
\end{equation}
\(N_n\)は\(n\)サイクル後の二重鎖相当の分子数で、30回は約\(1.07\times10^9\)倍となる。正確な両端長の産物が最初のサイクルから全分子を占めるわけではなく、実反応は効率・基質枯渇・阻害等でこの理想から外れる。回数をそのまま増幅倍率へ読み替えない。

\section{DNAライブラリと制限酵素：保存された情報と発現した情報}
\label{theory:biology:libraries}
ゲノムライブラリはゲノムDNAの断片群を、cDNAライブラリは採取したRNAから逆転写で作るDNA群を保持する。cDNAは発現したRNAに由来する\autocite{nhgri_cdna}。ゲノムにある遺伝子でも、その組織・時点で未発現なら対応するcDNAはない場合がある。成熟mRNA由来のcDNAでは通常イントロンが除かれるが、UTRは含み得るので「翻訳領域のみ」とは定義しない。
通常の緩衝液でDNAのリン酸主鎖は負に帯電し、ゲル電気泳動では陽極へ進む。分離は長さ等に依存し、陰極へという肢は電荷と移動方向を逆にしている。
6塩基の特定認識配列について、各位置が独立で4塩基が等確率という仮想モデルでは出現確率は\(4^{-6}=1/4096\)である。これは認識開始位置の平均間隔のモデルで、制限酵素切断後の全断片が正確に4096塩基対という主張ではない。ゲノムの組成偏り、配列相関、末端、切断可能性も別に考える。EcoRIの製造者文書も個別の酵素条件を扱う\autocite{neb_ecori}。
組換えインスリンの製法開発は1978年、米国での製品承認は1982年である\autocite{fda_insulin}。実験の成功・製品承認・後年の類似薬を同じ時点にせず、1990年代後半という誤答を識別する。

\section{酵素：平衡を変えず反応経路を変える}
\label{theory:biology:enzymes}
酵素は基質に選択的に作用する生体触媒である。多くはタンパク質だが触媒RNAもある。触媒は低い障壁の反応経路を与え、速度を変える。反応物と生成物の自由エネルギー差や平衡定数を変えることとは異なる。単純な活性化過程の近似\(k\propto\exp[-\Delta G^\ddagger/(RT)]\)では、障壁低下が速度増加へつながる。
リパーゼはトリアシルグリセロールのエステル結合を加水分解する酵素、αアミラーゼはデンプン等のα1→4結合へ作用する酵素である\autocite{iubmb_lipase,iubmb_amylase}。セルロースのβ1→4結合は同じグルコースからなるという理由だけでアミラーゼの標的とはならない。単位分子だけでなく結合の型が酵素の選択性を決める。
温度を上げると反応が速くなる効果と、酵素が変性・失活する効果が競合するため、至適温度と耐熱性は別概念である。Pyrococcus woeseiのαアミラーゼには100℃至適の一次研究があり、20℃以下又は100℃以上の酵素は一切ないという肢を反証する\autocite{hyper_amylase}。一つの存在例で不存在の全称命題を否定できる。

\section{クローン：分化状態を変え、核情報を受け渡す}
\label{theory:biology:cloning}
植物の組織培養ではオーキシン・サイトカイニン等のホルモン条件を用いてカルス形成や器官の再分化を制御する。ポプラの実験研究もホルモン処理とカルス・シュートの段階を比較している\autocite{plant_regeneration}。カルスは設問では不定形の増殖細胞塊、プロトプラストは細胞壁を取り除いた細胞であり別の語である。全カルスの全細胞が完全に同一の未分化状態という一般定義へ拡張しない。
哺乳類の体細胞核移植では、核を除いた卵へ体細胞核を移し、適切な操作で発生を進める\autocite{nhgri_cloning}。核を胚と読み違えたり受容側を卵原細胞とすると手順が成立しない。卵細胞の環境が核の発生プログラムを変えることが必要である。
体細胞クローンの核DNAの由来と、卵由来のミトコンドリアDNA、エピゲノム、発生環境を区別する。核情報が共通であっても外見・機能・健康状態が全て同一とは限らない。

\section{メチル化：配列を替えず使用状態を伝える}
\label{theory:biology:epigenetics}
同じゲノムを持つ細胞でも発現する遺伝子の組合せが違い、異なる役割を担う。DNAメチル化やヒストン修飾は配列を変更せず、その使用状態へ関わる。メチル化状態の維持は細胞分裂後の発現パターンの記憶へ寄与する\autocite{nhgri_epigenomics}。
これは一塩基置換のような配列変異と異なる。全ての修飾が必ず維持される、全遺伝子のメチル化が同じ効果を持つ、個体の全特徴が環境から子孫へ継承されるという主張ではない。設問の「子孫細胞」は細胞分裂での継承として読み、親子個体の世代を無条件に追加しない。

\section{必須アミノ酸：重要性でなく供給条件}
\label{theory:biology:nutrition}
必須アミノ酸は、生体が必要量を体内合成で満たせず外部供給を要するものをいう。非必須もタンパク質の機能へ重要であり、重要でないという分類ではない。通常のヒトの必須9種はHis・Ile・Leu・Lys・Met・Phe・Thr・Trp・Valである。Proは通常非必須に分類され、病態・成長等の条件では条件付き必須となり得る\autocite{nih_amino}。
令和7年度のIle・Val・Leu・Pro・Metでは通常の分類によりProを選ぶ。設問の広い「哺乳類」という語を、全種・全発育段階で栄養条件が同じとする保証へ変えない。前駆体から作れることと、その状況で必要量を十分作れることも区別する。

\section{出題形式と確認手順}
今回の18問では粒子数、同位体の正誤表、酸化数の最大・最小、燃焼収率、濃度調製、組成変換、結合エネルギー、材料語句の組合せを確認した。全年度の頻度は残る分類確認後に更新する。対象と分母を決め、保存則で計算し、候補順と否定方向を確認する。正答番号が合っても密度近似・生成物未指定の条件不足は残す。

\section{過去問}
\section{金属材料：結合・結晶・組織から試験と選択へ}
\label{theory:material:overview}
材料選択では、単に「強い金属」を探すだけでは足りない。小さな荷重での変形、永久変形の始まり、き裂の進展、繰返し負荷、電気抵抗、加工のしやすさは異なる特性である。組成と加工条件が結晶・転位・粒界・欠陥を変え、それが使用中の応答へ結び付く。本節は基礎科目の材料問題から、その因果関係を整理する。

\subsection{金属結合と塑性：電子と転位の役割}
\label{theory:material:bonding}
金属の外側の電子は原子一つに局在せず、正の原子核・内殻電子からなる芯の間で共有される。自由電子の海という近似は金属結合と導電性を理解する入口である\cite{mit_material_atomistic2001}。同位体は中性子数の違う同じ元素の原子であり、自由電子の代わりの結合機構ではない。

結合が特定方向だけに限られないことと、原子面が容易に全て同時に滑ることは同じではない。実際の塑性では格子の線欠陥である転位が移動し、少しずつすべりを進める。転位の移動を妨げる粒界・溶質・析出物は強化に寄与するが、延性との交換条件も生む\cite{mit_material_dislocations2001}。結合だけで全金属の室温延性を保証しない。

\subsection{BCC・FCC・HCP：胞と配位数を数え分ける}
\label{theory:material:crystals}
単位胞は結晶の並進で繰り返す領域である。格子点と、その各点に付く原子の基底を区別する。胞の形を決めずに「一胞に何原子」と尋ねると、原始胞と大きな慣用胞の数が混ざる\cite{cambridge_unitcell,princeton_hcp}。

通常の立方胞では、角の原子は隣り合う8胞で共有し、面の原子は2胞で共有する。この分率から独自に数えると、
\[
 N_{\rm BCC}=8\times\frac18+1=2,\qquad
 N_{\rm FCC}=8\times\frac18+6\times\frac12=4.
\]
BCCの最近接配位数は8、FCCは12である。配位数は一つの原子の周囲にある最近接原子の数であり、胞の中の原子数ではない\cite{openstax_crystals}。

硬球近似で半径を\(R\)、立方胞辺を\(a\)とする。BCCは体対角線、FCCは面対角線で接触するため、幾何から
\[
 \sqrt3a=4R\quad({\rm BCC}),\qquad
 \sqrt2a=4R\quad({\rm FCC})
\]
となる。FCCでは\(a=2\sqrt2R\)である。式の暗記だけでなく、接触する原子列を指定する。

HCPは六方Bravais格子に二原子基底を置いた構造であり、原始並進胞に2原子を含む。よく図示される大きな六角柱では、角12原子の共有分\(12/6\)、上下の面中心2原子の共有分\(2/2\)、内部3原子を合計し、6原子となる。この六角柱は原始並進胞の3倍の体積を持つ。平成29年度と令和6年度の「HCPは2原子」は前者の規約で読み、六角柱の図へ無条件に転記しない。

常温常圧の元素ではFe・NaはBCC、Al・CuはFCC、Ti・ZnはHCPである\cite{princeton_iron,princeton_sodium,princeton_aluminum,princeton_copper,princeton_titanium,princeton_zinc}。相の安定性は温度・圧力・組成に依存する。別資料の温度を指定しない元素例示を、室温の分類にそのまま使わない。

\subsection{荷重と変位から応力・ひずみへ}
\label{theory:material:tension}
初期断面積を\(A_0\)、初期標点距離を\(L_0\)、荷重を\(F\)、標点間の長さを\(L\)とする。公称値は
\[
 \sigma_{\rm e}=\frac{F}{A_0},\qquad
 \varepsilon_{\rm e}=\frac{L-L_0}{L_0}
\]
である。試験機のクロスヘッド変位は機械・つかみの変形を含むことがあるため、試験片の標点間の伸びと対応させる。

現在の断面積\(A\)で割った\(\sigma_{\rm t}=F/A\)が真応力である。真ひずみは現在長さによる微小増分を積分するので、
\[
 \varepsilon_{\rm t}=\int_{L_0}^{L}\frac{d\ell}{\ell}
 =\ln\frac{L}{L_0}=\ln(1+\varepsilon_{\rm e}).
\]
全伸びを変形後の長さで割るだけでは、この積分値にならない\cite{mit_material_stressstrain2001}。微小ひずみでは両ひずみは近い。大変形・くびれ後には、標点間平均と局所値を区別する必要がある。

\subsection{フック則：剛性と強度を分ける}
\label{theory:material:elasticity}
小応力の単軸線形弾性では\(\sigma=E\varepsilon\)となる。比例定数\(E\)がヤング率であり、ひずみが無次元なので応力と同じPaを単位に持つ。棒の剛性は\(k=EA_0/L_0\)であり、材料の\(E\)に加えて形状にも依存する\cite{mit_material_elastic2000}。

降伏応力は塑性が始まる応力、引張強度は通常、公称応力の最大値である。ヤング率はこの二つの強度ではない。ポアソン比は横と軸のひずみの比であり無次元である。ヘンリー則は溶解、ブラッグ則は回折の関係である。

出題のMg45 GPa、W407 GPaは大学教材にもある代表値である\cite{princeton_magnesium,princeton_tungsten}。同じ応力ならMgの方が大きく弾性変形する。通常の金属では昇温に伴う弾性率低下を考えるが、方位・温度区間・相変態を指定し、全材料の単調減少へ一般化しない。弾性は除荷による回復を指し、線形性と常に同義ではない。

\subsection{粒径・加工硬化・再結晶}
\label{theory:material:strengthening}
一般的な多結晶鋼などでは、結晶粒界が転位移動の障害になる。通常の粒径範囲の経験的関係は
\[
 \sigma_y=\sigma_0+\frac{k_y}{\sqrt d}
\]
である。\(d\)が小さいほど第二項が大きくなる。粗大化で強度が上がるという出題肢は、この通常の比較では逆である。

塑性加工で転位が増えると、転位同士も移動を妨げる。これが加工硬化を理解する基本である。加熱による回復・再結晶は蓄積した欠陥を減らし、通常は加工硬化を緩和して延性を回復する\cite{mit_material_dislocations2001}。

Hall--Petch式を原子数が極端に少ないナノ粒まで外挿し、強度が際限なく増すとはしない。粒界すべり・回転・拡散や加工時の欠陥など、支配機構が変わる。逆Hall--Petchの報告にも加工・硬度から降伏応力への換算・試験中粒成長の問題があり、全てのナノ材料が必ず弱化するという保証にしない\cite{hall_petch_limits2020}。

\subsection{欠陥・脆性破壊・疲労・衝撃靭性}
\label{theory:material:fracture}
無欠陥の原子結合を一斉に切る理論強度と、孔・介在物・き裂を含む実材料の強度は異なる。き裂先端には応力が集中し、き裂寸法・形状・負荷と破壊抵抗を一緒に評価する。線形破壊力学の典型的形は\(K_I=Y\sigma\sqrt{\pi a}\)であり、\(a\)の定義と形状係数\(Y\)を指定する。

脆性破壊は顕著な塑性変形が小さい破壊、疲労は繰返し負荷による損傷蓄積である。疲労では一回なら耐えられる荷重でも、き裂が成長して最後に急速破断することがある\cite{mit_material_fatigue2001}。腐食疲労はあるが、疲労を経年腐食だけで定義しない。

リバティ船とコメットは破壊力学発展を促した代表例である\cite{nist_fracture_history2010}。出題は船の脆性と航空機の疲労を対照する。コメットに関する大学教材は窓周辺のリベット孔からのき裂も説明する\cite{mit_material_fracture2001}。原典の「窓の角」という簡略表現を全事故の起点の保証へ変えず、溶接部・開口形状・低温・繰返し圧力などの複合条件を考える。

Charpyは切欠き試験片を衝撃破断し、吸収エネルギーから衝撃靭性を評価する。温度・寸法・切欠き・負荷速度に依存する測定であり、平面ひずみ破壊靭性\(K_{IC}\)の直接測定ではない\cite{nist_charpy}。

\subsection{純鉄の結晶相と磁気転移}
\label{theory:material:ironphases}
常温の純鉄\(\alpha\)-FeはBCCで強磁性、Coも常温の強磁性金属である\cite{princeton_iron,openstax_magnetism}。加熱時には約770°Cで強磁性を失い、BCCのまま常磁性となる。Curie温度1043 Kを摂氏へ換算すると約770°Cである\cite{iron_curie1954}。

その後、約912°Cで\(\gamma\)-FeのFCCへ変態する。加熱時の\(\alpha\to\gamma\)は収縮し、冷却時の逆変態では膨張する\cite{iron_transform2023}。磁気転移と結晶の同素変態は別であり、強磁性が結晶変態温度まで続くとしない。一般の熱膨張と変態による急な変化も別に考える。原子数と実際の格子定数が体積を決めるので、充填率だけで全材料の変態体積を決めない。

\subsection{被膜質量：寸法と百分率の分母}
\label{theory:material:coating}
面積\(S\)、被膜厚\(t_c\)、基材厚\(t_s\)に対し、片面被膜の質量は\(m_c=St_c\rho_c\)、基材は\(m_s=St_s\rho_s\)である。被膜の質量百分率は\(100m_c/(m_c+m_s)\)となる。厚さを同じ長さ単位にそろえ、片面・両面を図示条件に合わせる。

令和4年度基礎4-3の原典値を独自計算すると、
\[
 m_{\rm Ni}=0.50\times1.0\times10\times10^{-6}\times8900
 =0.0445\ {\rm kg},\qquad
 m_s=0.50\times1.0\times0.60\times10^{-3}\times7900
 =2.37\ {\rm kg}.
\]
したがって\(100\times0.0445/(2.37+0.0445)=1.843\ldots\%\)である。候補に対応する丸め値は\(4.5\times10^{-2}\)kg、1.8\%となる。百分率は鋼だけを分母にしない。

Niのレアメタル分類は日本の政策文脈であり、JOGMECは国際的にはベースメタルと呼ばれることも明記する\cite{jogmec_nickel2014}。分類語を普遍的な自然科学上の属性と思わず、計算条件とは別に確認する。

\subsection{材料選択：軽さ・剛性・加工性・使用環境}
\label{theory:material:selection}
軽量化でも、破断・永久変形を避ける選択と、弾性変形を抑える選択は異なる。比強度は\(\sigma_{\rm strength}/\rho\)、比弾性率は\(E/\rho\)である。対象の強度が引張・降伏のどちらかも示す。令和6年度基礎4-4はこの二つの取り違えを問う。

室温代表値では電気抵抗率はCuよりAu、AuよりFeが大きい\cite{openstax_resistivity}。純度と温度をそろえる。ステンレス鋼の一般的な定義はCr最低10.5質量\%、C最大1.2質量\%である\cite{stainless_definition2022}。したがってCr10\%未満は該当しないが、「10\%以上なら全て該当」という逆命題は導けない。

ダイヤモンドの高硬度を切削工具へ利用する例がある\cite{nims_diamond2007}。用途に応じて熱条件・被削材との化学反応も評価する。水素脆化は水素の侵入に伴う延性・破壊抵抗等の機械特性の劣化である。「物性が変化する」というだけでは範囲が広すぎ、改善や無関係な変化まで脆化とは呼ばない\cite{nims_hydrogen2024}。

硫黄快削鋼ではMnS等の介在物を利用して被削性を改善する。ただし工具寿命と仕上面粗さは同じ方向に良くなるとは限らず、脱酸・介在物の形・組織・切削速度も影響する\cite{sulfur_machinability2004}。S量だけから全鋼の被削性を保証せず、延性・溶接性など他の性能との交換条件を考える。

\section{水中の酸性度を共役塩基から考える}
\label{theory:basicchem:acid}
酸の強さは水中でプロトンを渡す平衡であり、分析濃度そのものではない。強酸の希薄溶液ではほぼ完全解離、弱酸では部分解離となる。単プロトン弱酸の分析濃度を\(c\)、解離定数を\(K_a\)とし、解離が十分小さければ\([\mathrm H^+]\simeq\sqrt{K_ac}\)である。水の自己解離が支配するほどの希薄条件や活量の影響を無視できない高濃度にはそのまま使わない。
共役塩基の安定化から構造の影響を読む。カルボキシラートの共鳴と電子求引基は解離を助け、CF3COOHは酢酸より強い。フェノキシドは共鳴安定化するためフェノールはエタノールより強い。HFとHClの比較では結合切断と溶媒和も入り、電気陰性度だけでは順位が決まらない\autocite{purdue_acid_strength,yale_hydrohalic}。炭酸の第一解離と溶存CO2を含む見かけの定義を分ける。

\section{生成エンタルピーの差として反応熱を計算する}
\label{theory:basicchem:hess}
エンタルピーは状態量なので、実際の反応経路によらず始状態と終状態の差で決まる。生成エンタルピーを用いると
\begin{equation}\Delta_r H^\circ=\sum_{\mathrm{products}}\nu_i\Delta_fH_i^\circ-\sum_{\mathrm{reactants}}\nu_i\Delta_fH_i^\circ.\end{equation}
係数と相を指定して足す。原典が生成時の放出熱を正の大きさで示す場合は、対応する生成エンタルピーを負に読み替える。液体水と水蒸気では値が違い、反応熱が小さい差で決まる設問では発熱・吸熱まで変わり得る\autocite{openstax_hess}。単位は指定反応式を一回進める基準か、生成物1 mol当たりかを明記する。

\section{沈殿した後に溶ける両性水酸化物}
\label{theory:basicchem:amphoteric}
Al3+にOH−を加えるとAl(OH)3が沈殿し、さらに過剰OH−でヒドロキソ錯イオンを作って溶ける。
\begin{equation}\mathrm{Al^{3+}+3OH^-\to Al(OH)_3},\qquad
\mathrm{Al(OH)_3+OH^-\rightleftharpoons[Al(OH)_4]^-}.\end{equation}
酸にも塩基にも反応する両性と、任意の沈殿が過剰試薬で溶けることは同じでない。添加剤がNaOHかアンモニアかを確認し、錯形成の条件を取り違えない\autocite{chem55_ampho}。

\section{熱伝導を担い手・散乱・空隙へ分ける}
\label{theory:basicchem:thermal}
通常の金属の熱輸送は自由電子の寄与が大きく、絶縁性結晶では格子振動が主要な担い手となる。合金化や構造の乱れは散乱を増す一般的傾向を持つ。低い熱伝導率の空隙を増せば通常は有効熱伝導率が低下する。ただし気孔中の対流・放射が重要な高温条件や配向の違いは別に検討する。
高分子の輸送も結晶化率だけでなく鎖の配向・界面・温度に依存する。結晶質と非晶質の通常比較を、任意の異組成間の順位へ広げない\autocite{aist_metal_thermal,openstax_thermal_conduction,chem55_polymer}。

\section{平衡収率と反応速度を分ける}
\label{theory:basicchem:equilibrium}
発熱のアンモニア合成\(\mathrm{N_2+3H_2\rightleftharpoons2NH_3}\)では、低温が生成側の平衡に有利である。体積を減らして圧力を上げると、気体4 molから2 molとなる生成側へ移る。温度は平衡定数を変え、圧縮は分圧と反応商を変える\autocite{chem55_equilibrium}。
触媒は平衡へ達する速度を変えるが平衡定数を変えない。低温による収率の増大と速度低下は別の効果なので、工業条件は両方と設備条件を考えて選ぶ。

\section{中和後に残るものからpHを判断する}
\label{theory:basicchem:neutralization}
体積と濃度から物質量を求め、放出・受容できるプロトンの当量を比較する。酸と塩基が当量でも、できた塩の溶液が必ず中性になるわけではない。弱塩基の共役酸NH4+は水へプロトンを渡して酸性を示す\autocite{chem55_hydrolysis}。
硫酸1 molとNaOH2 molは硫酸ナトリウムとなり、通常の中性近似又は硫酸イオンの弱い塩基性で扱う。硫酸1 molとNH3 2 molでは硫酸アンモニウムとなりNH4+の酸性が支配する。強塩基が過剰ならそのOH−を先に扱い、NH3とNH4+が共存する場合は緩衝系を考える。体積加算の近似は精密な濃度計算と区別する。

\section{腐食を環境と表面状態から判断する}
\label{theory:basicchem:corrosion}
腐食は材料と環境の反応で進む。イオン化傾向が大きくても、Alの強固な酸化膜のように反応を抑える不動態があれば通常の環境で耐食性を示す。ステンレスのCrによる不動態も同じ考え方だが、塩化物は孔食・すきま腐食を引き起こし得る\autocite{chem55_corrosion,stainless_definition2022}。海水用途は鋼種・温度・流動・形状等を合わせて選ぶ。
応力腐食割れは特定の材料・環境と引張応力の組合せで起こり、通常の引張強さより低い応力でも破壊し得る。水素脆化は水素侵入・拡散・欠陥との相互作用に関係する\autocite{nims_hydrogen2024}。これらを単なる均一減肉やすべての材料の性質へまとめない。

\section{金属資源の存在量と製錬反応}
\label{theory:basicchem:metallurgy}
地殻の存在量と工業生産量は違う。Alは地殻中で最も多い金属だが、金属生産量の順位をそこから直接決めない。反応性の高いAlは溶融塩中の酸化物を電解して得る\autocite{chem55_crust}。Tiの代表的なクロール法はTiCl4をMgで還元する\autocite{chem55_kroll}。
高炉では鉱石をコークス由来のCO等で還元し、石灰石等で不要成分を分離する\autocite{chem55_iron}。理想化した純粋Fe2O3から純粋Feを得る質量収支は、式量160のうちFeが112なので\(w_{\mathrm{ore}}=(160/112)w_{\mathrm{Fe}}\)。実操業では鉱石品位・歩留まり・銑鉄中の炭素を別に扱う。埋蔵量を生産量で割った可採年数は特定時点の指標であり、不変の枯渇期限ではない。

\section{密度・抵抗率・融点を独立に並べる}
\label{theory:basicchem:properties}
密度は単位体積当たり質量、抵抗率は形状を除いた電気輸送の指標、融点は指定圧力での相転移温度である。純金属の通常値では密度はCu\(>\)Fe\(>\)Al、20℃付近の抵抗率はFe\(>\)Al\(>\)Cu、融点はFe\(>\)Cu\(>\)Alとなる\autocite{openstax_resistivity,chem55_properties,chem55_properties_al,chem55_properties_fe}。三つの順位を一つの「金属の強さ」へ置き換えない。記号表の割当は年度ごとに戻して照合する。

\section{沸点上昇は溶質粒子数を数える}
\label{theory:basicchem:colligative}
不揮発性溶質を溶かすと蒸気圧が下がり、同じ外圧へ達するには温度を上げる必要がある。希薄な理想溶液の近似では\(\Delta T_b=K_b i b\)。\(b\)は溶媒1 kg当たりの溶質式単位の物質量、\(i\)は溶解後の有効粒子数の倍率である\autocite{chem55_colligative}。
同じ質量モル濃度なら、理想的な完全解離モデルの尿素等は1、NaClは2、CaCl2は3の粒子に対応する。実溶液の活量やイオン間相互作用により厳密な倍率からずれる。100℃を基準とする水の沸点は1気圧の条件付きであり、外圧が違えば純水の基準温度も変わる。

\section{材料名から代表的構成元素を読む}
\label{theory:basicchem:elements}
乾電池の代表的負極Zn、シリカ光ファイバーのSi、Al-Cu系ジュラルミンのCu、代表的な磁石のFeを用途と対応させる\autocite{chem55_fiber,chem55_duralumin,openstax_magnetism}。Liイオン電池ではCoを含む正極材料が代表例となるが、Coを含まない正極もあり、全電池共通の必要元素ではない\autocite{openstax_batteries}。材料の例と全称命題を区別し、正極と負極を読み分ける。

\section{ギブズエネルギーの単位をそろえる}
\label{theory:basicchem:gibbs}
\begin{equation}G=H-TS,\qquad\Delta_rG^\circ=\Delta_rH^\circ-T\Delta_rS^\circ.\end{equation}
\(T\)は絶対温度で、エントロピーの単位はエネルギー/温度である。したがって\(TS\)と\(H\)は同じ単位になる\autocite{chem55_gibbs}。一定温度・圧力での反応の駆動力を考えるとき、標準反応量と実際の組成での反応量を分け、後者には反応商の項が入る。

\section{結合の前後から有機反応を分類する}
\label{theory:basicchem:organic}
付加は多重結合に試薬が加わり、脱離は小分子が抜けて多重結合等ができ、置換は一つの原子団が別の原子団へ変わる\autocite{chem55_organic}。アルコールのOHからBrへの変換は置換、アルコールから水が抜ける反応は脱離、アルケンへのHBrは付加となる。カルボン酸からエステルへの変換はアシル基の置換として扱え、同時に縮合という観点でも呼べる。触媒名だけでなく、失われた結合と形成した結合を追う。

\section{反応の分類と反応が進む条件を分ける}
\label{theory:basicchem:redox}
酸化数の増加は酸化、減少は還元で、同一反応では両者を伴う。単体は0、通常の酸化物でOは−2、通常の酸や水でHは+1として電荷収支を取り、例外を確認する\autocite{iupac_oxidation_state}。
反応式の全元素で酸化数が不変ならその式は酸化還元ではない。ただし式が原子数と電荷を保存することは、その反応が指定条件で進むことの十分条件ではない。形式的な塩の交換式では溶解平衡や沈殿の駆動力を別に調べる。

\section{放射線の種類と材料への作用}
\label{theory:basicchem:radiation}
\(\alpha\)線はHe原子核、\(\beta^-\)線は電子、\(\beta^+\)線は陽電子、\(\gamma\)線は光子である。通常\(\alpha\)は強い電離作用と短い飛程を持つ。光子遮蔽の減衰は\(I=I_0e^{-\mu x}\)という単純モデルで表せ、\(\mu\)は材料と光子エネルギーに依存する\autocite{finish56_radiation}。同じ厚さと同じ面密度は別の比較である。照射は結晶欠陥や電子状態を変えて力学・電気特性を変え得る。

\section{構造と用途を対応させる化学物質}
\label{theory:basicchem:substances}
同素体は元素が同じでも結合・構造が違う。黒鉛は層状で軟らかく電気を伝え、ダイヤモンドの硬さ・通常の絶縁性と異なる。Haber--Boschはアンモニア合成、Solvayは炭酸ナトリウムの代表的製法である。
TiO2は光照射による表面反応を利用した抗菌・防汚等の光触媒材料として用いられる\autocite{finish56_titania}。性能は結晶相・複合化・光の波長等に依存する。鉛蓄電池の放電時は負極Pb、正極PbO2で、PbOという別酸化物と混同しない\autocite{openstax_batteries}。

\section{結合極性と分子極性を分ける}
\label{theory:basicchem:polarity}
極性結合を持つ分子でも、各結合の双極子が対称に配置されれば分子の双極子は0となる\autocite{finish56_polarity}。直線CO2、正四面体CH4・CCl4、平面三角形BF3は通常の形で相殺する。エーテルのO付近は折れ曲がり、分子の双極子が残る。結合を作る元素だけでなく、孤立電子対を含む分子形を確かめる。

\section{コロイドの粒子・界面・状態}
\label{theory:basicchem:colloids}
コロイドは分子溶液より大きい粒子等が分散する系で、粒子の界面と電荷が安定性を左右する。疎水コロイドに電解質を加えて反発を弱め、集合・沈殿させるのが凝析。小分子・イオンを膜で通して粒子と分けるのが透析である\autocite{openstax_colloids}。
光散乱で光路が見えるTyndall現象と、電場中で荷電粒子が移動する電気泳動は別の作用である。流動する分散系のゾルと、網目等ができ流動しにくいゲルを区別する。全コロイドが同じ量の電解質で凝析するわけではなく、表面・溶媒・イオン種等が関わる。

\subsection{平成23年度 問4-1}
\label{q:h23-reg-basic-4-1}

\noindent\textbf{原典：} sources/H23/basic.pdf、PDF物理12頁。

\textbf{公式正答：} 2（single）。

\par\textit{原典本文の校正と独自解説は未作成。分類のみの登録項目。}

\subsection{平成23年度 問4-2}
\label{q:h23-reg-basic-4-2}

\noindent\textbf{原典：} sources/H23/basic.pdf、PDF物理12頁。

\textbf{公式正答：} 1（single）。

\par\textit{原典本文の校正と独自解説は未作成。分類のみの登録項目。}

\subsection{平成23年度 問4-3}
\label{q:h23-reg-basic-4-3}

\noindent\textbf{原典：} sources/H23/basic.pdf、PDF物理12頁。

\textbf{公式正答：} 4（single）。

\par\textit{原典本文の校正と独自解説は未作成。分類のみの登録項目。}

\subsection{平成23年度 問4-4}
\label{q:h23-reg-basic-4-4}

\noindent\textbf{原典：} sources/H23/basic.pdf、PDF物理13頁。

\textbf{公式正答：} 4（single）。

\par\textit{原典本文の校正と独自解説は未作成。分類のみの登録項目。}

\subsection{平成23年度 問4-5}
\label{q:h23-reg-basic-4-5}

\noindent\textbf{原典：} sources/H23/basic.pdf、PDF物理14頁。

\textbf{公式正答：} 3（single）。

\par\textit{原典本文の校正と独自解説は未作成。分類のみの登録項目。}

\subsection{平成24年度 問4-1}
\label{q:h24-reg-basic-4-1}

\noindent\textbf{原典：} sources/H24/basic.pdf、PDF物理12頁。

\textbf{公式正答：} 1（single）。

\par\textit{原典本文の校正と独自解説は未作成。分類のみの登録項目。}

\subsection{平成24年度 問4-2}
\label{q:h24-reg-basic-4-2}

\noindent\textbf{原典：} sources/H24/basic.pdf、PDF物理12頁。

\textbf{公式正答：} 3（single）。

\par\textit{原典本文の校正と独自解説は未作成。分類のみの登録項目。}

\subsection{平成24年度 問4-3}
\label{q:h24-reg-basic-4-3}

\noindent\textbf{原典：} sources/H24/basic.pdf、PDF物理13頁。

\textbf{公式正答：} 2（single）。

\par\textit{原典本文の校正と独自解説は未作成。分類のみの登録項目。}

\subsection{平成24年度 問4-4}
\label{q:h24-reg-basic-4-4}

\noindent\textbf{原典：} sources/H24/basic.pdf、PDF物理13頁。

\textbf{公式正答：} 4（single）。

\par\textit{原典本文の校正と独自解説は未作成。分類のみの登録項目。}

\subsection{平成24年度 問4-5}
\label{q:h24-reg-basic-4-5}

\noindent\textbf{原典：} sources/H24/basic.pdf、PDF物理14頁。

\textbf{公式正答：} 3（single）。

\par\textit{原典本文の校正と独自解説は未作成。分類のみの登録項目。}

\subsection{平成25年度 問4-1}
\label{q:h25-reg-basic-4-1}

\noindent\textbf{原典：} sources/H25/basic.pdf、PDF物理15頁。

\textbf{公式正答：} 2（single）。

\par\textit{原典本文の校正と独自解説は未作成。分類のみの登録項目。}

\subsection{平成25年度 問4-2}
\label{q:h25-reg-basic-4-2}

\noindent\textbf{原典：} sources/H25/basic.pdf、PDF物理15頁。

\textbf{公式正答：} 2（single）。

\par\textit{原典本文の校正と独自解説は未作成。分類のみの登録項目。}

\subsection{平成25年度 問4-3}
\label{q:h25-reg-basic-4-3}

\noindent\textbf{原典：} sources/H25/basic.pdf、PDF物理16頁。

\textbf{公式正答：} 4（single）。

\par\textit{原典本文の校正と独自解説は未作成。分類のみの登録項目。}

\subsection{平成25年度 問4-4}
\label{q:h25-reg-basic-4-4}

\noindent\textbf{原典：} sources/H25/basic.pdf、PDF物理17頁。

\textbf{公式正答：} 3（single）。

\par\textit{原典本文の校正と独自解説は未作成。分類のみの登録項目。}

\subsection{平成25年度 問4-5}
\label{q:h25-reg-basic-4-5}

\noindent\textbf{原典：} sources/H25/basic.pdf、PDF物理17頁。

\textbf{公式正答：} 1（single）。

\par\textit{原典本文の校正と独自解説は未作成。分類のみの登録項目。}

\subsection{平成25年度 問4-6}
\label{q:h25-reg-basic-4-6}

\noindent\textbf{原典：} sources/H25/basic.pdf、PDF物理18頁。

\textbf{公式正答：} 3（single）。

\par\textit{原典本文の校正と独自解説は未作成。分類のみの登録項目。}

\subsection{平成26年度 問4-1}
\label{q:h26-reg-basic-4-1}

\noindent\textbf{原典：} sources/H26/basic.pdf、PDF物理13頁。

\textbf{公式正答：} 1（single）。

\par\textit{原典本文の校正と独自解説は未作成。分類のみの登録項目。}

\subsection{平成26年度 問4-2}
\label{q:h26-reg-basic-4-2}

\noindent\textbf{原典：} sources/H26/basic.pdf、PDF物理13頁。

\textbf{公式正答：} 2（single）。

\par\textit{原典本文の校正と独自解説は未作成。分類のみの登録項目。}

\subsection{平成26年度 問4-3}
\label{q:h26-reg-basic-4-3}

\noindent\textbf{原典：} sources/H26/basic.pdf、PDF物理14頁。

\textbf{公式正答：} 5（single）。

\par\textit{原典本文の校正と独自解説は未作成。分類のみの登録項目。}

\subsection{平成26年度 問4-4}
\label{q:h26-reg-basic-4-4}

\noindent\textbf{原典：} sources/H26/basic.pdf、PDF物理14頁。

\textbf{公式正答：} 4（single）。

\par\textit{原典本文の校正と独自解説は未作成。分類のみの登録項目。}

\subsection{平成26年度 問4-5}
\label{q:h26-reg-basic-4-5}

\noindent\textbf{原典：} sources/H26/basic.pdf、PDF物理15頁。

\textbf{公式正答：} 3（single）。

\par\textit{原典本文の校正と独自解説は未作成。分類のみの登録項目。}

\subsection{平成26年度 問4-6}
\label{q:h26-reg-basic-4-6}

\noindent\textbf{原典：} sources/H26/basic.pdf、PDF物理15頁。

\textbf{公式正答：} 1（single）。

\par\textit{原典本文の校正と独自解説は未作成。分類のみの登録項目。}

\subsection{平成27年度 問4-1}
\label{q:h27-reg-basic-4-1}

\noindent\textbf{原典：} sources/H27/basic.pdf、PDF物理12頁。

\textbf{公式正答：} 4（single）。

\par\textit{原典本文の校正と独自解説は未作成。分類のみの登録項目。}

\subsection{平成27年度 問4-2}
\label{q:h27-reg-basic-4-2}

\noindent\textbf{原典：} sources/H27/basic.pdf、PDF物理12頁。

\textbf{公式正答：} 3（single）。

\par\textit{原典本文の校正と独自解説は未作成。分類のみの登録項目。}

\subsection{平成27年度 問4-3}
\label{q:h27-reg-basic-4-3}

\noindent\textbf{原典：} sources/H27/basic.pdf、PDF物理12頁。

\textbf{公式正答：} 5（single）。

\par\textit{原典本文の校正と独自解説は未作成。分類のみの登録項目。}

\subsection{平成27年度 問4-4}
\label{q:h27-reg-basic-4-4}

\noindent\textbf{原典：} sources/H27/basic.pdf、PDF物理13頁。

\textbf{公式正答：} 1（single）。

\par\textit{原典本文の校正と独自解説は未作成。分類のみの登録項目。}

\subsection{平成27年度 問4-5}
\label{q:h27-reg-basic-4-5}

\noindent\textbf{原典：} sources/H27/basic.pdf、PDF物理14頁。

\textbf{公式正答：} 2（single）。

\par\textit{原典本文の校正と独自解説は未作成。分類のみの登録項目。}

\subsection{平成27年度 問4-6}
\label{q:h27-reg-basic-4-6}

\noindent\textbf{原典：} sources/H27/basic.pdf、PDF物理14頁。

\textbf{公式正答：} 3（single）。

\par\textit{原典本文の校正と独自解説は未作成。分類のみの登録項目。}

\subsection{平成28年度 問4-1}
\label{q:h28-reg-basic-4-1}

\noindent\textbf{原典：} sources/H28/basic.pdf、PDF物理17頁。

\textbf{公式正答：} 4（single）。

\par\textit{原典本文の校正と独自解説は未作成。分類のみの登録項目。}

\subsection{平成28年度 問4-2}
\label{q:h28-reg-basic-4-2}

\noindent\textbf{原典：} sources/H28/basic.pdf、PDF物理17頁。

\textbf{公式正答：} 4（single）。

\par\textit{原典本文の校正と独自解説は未作成。分類のみの登録項目。}

\subsection{平成28年度 問4-3}
\label{q:h28-reg-basic-4-3}

\noindent\textbf{原典：} sources/H28/basic.pdf、PDF物理18頁。

\textbf{公式正答：} 3（single）。

\par\textit{原典本文の校正と独自解説は未作成。分類のみの登録項目。}

\subsection{平成28年度 問4-4}
\label{q:h28-reg-basic-4-4}

\noindent\textbf{原典：} sources/H28/basic.pdf、PDF物理18頁。

\textbf{公式正答：} 2（single）。

\par\textit{原典本文の校正と独自解説は未作成。分類のみの登録項目。}

\subsection{平成28年度 問4-5}
\label{q:h28-reg-basic-4-5}

\noindent\textbf{原典：} sources/H28/basic.pdf、PDF物理19頁。

\textbf{公式正答：} 1（single）。

\par\textit{原典本文の校正と独自解説は未作成。分類のみの登録項目。}

\subsection{平成28年度 問4-6}
\label{q:h28-reg-basic-4-6}

\noindent\textbf{原典：} sources/H28/basic.pdf、PDF物理20頁。

\textbf{公式正答：} 3（single）。

\par\textit{原典本文の校正と独自解説は未作成。分類のみの登録項目。}

\subsection{平成29年度 問4-1}
\label{q:h29-reg-basic-4-1}

\noindent\textbf{原典：} sources/H29/basic.pdf、PDF物理15頁。

\textbf{公式正答：} 4（single）。

\par\textit{原典本文の校正と独自解説は未作成。分類のみの登録項目。}

\subsection{平成29年度 問4-2}
\label{q:h29-reg-basic-4-2}

\noindent\textbf{原典：} sources/H29/basic.pdf、PDF物理15頁。

\textbf{公式正答：} 5（single）。

\par\textit{原典本文の校正と独自解説は未作成。分類のみの登録項目。}

\subsection{平成29年度 問4-3}
\label{q:h29-reg-basic-4-3}

\noindent\textbf{原典：} sources/H29/basic.pdf、PDF物理16頁。

\textbf{公式正答：} 4（single）。

\par\textit{原典本文の校正と独自解説は未作成。分類のみの登録項目。}

\subsection{平成29年度 問4-4}
\label{q:h29-reg-basic-4-4}

\noindent\textbf{原典：} sources/H29/basic.pdf、PDF物理16頁。

\textbf{公式正答：} 1（single）。

\par\textit{原典本文の校正と独自解説は未作成。分類のみの登録項目。}

\subsection{平成29年度 問4-5}
\label{q:h29-reg-basic-4-5}

\noindent\textbf{原典：} sources/H29/basic.pdf、PDF物理17頁。

\textbf{公式正答：} 1（single）。

\par\textit{原典本文の校正と独自解説は未作成。分類のみの登録項目。}

\subsection{平成29年度 問4-6}
\label{q:h29-reg-basic-4-6}

\noindent\textbf{原典：} sources/H29/basic.pdf、PDF物理17頁。

\textbf{公式正答：} 1（single）。

\par\textit{原典本文の校正と独自解説は未作成。分類のみの登録項目。}

\subsection{平成30年度 問4-1}
\label{q:h30-reg-basic-4-1}

\noindent\textbf{原典：} sources/H30/basic.pdf、PDF物理15頁。

\textbf{公式正答：} 5（single）。

\par\textit{原典本文の校正と独自解説は未作成。分類のみの登録項目。}

\subsection{平成30年度 問4-2}
\label{q:h30-reg-basic-4-2}

\noindent\textbf{原典：} sources/H30/basic.pdf、PDF物理15頁。

\textbf{公式正答：} 5（single）。

\par\textit{原典本文の校正と独自解説は未作成。分類のみの登録項目。}

\subsection{平成30年度 問4-3}
\label{q:h30-reg-basic-4-3}

\noindent\textbf{原典：} sources/H30/basic.pdf、PDF物理15頁。

\textbf{公式正答：} 3（single）。

\par\textit{原典本文の校正と独自解説は未作成。分類のみの登録項目。}

\subsection{平成30年度 問4-4}
\label{q:h30-reg-basic-4-4}

\noindent\textbf{原典：} sources/H30/basic.pdf、PDF物理16頁。

\textbf{公式正答：} 2（single）。

\par\textit{原典本文の校正と独自解説は未作成。分類のみの登録項目。}

\subsection{平成30年度 問4-5}
\label{q:h30-reg-basic-4-5}

\noindent\textbf{原典：} sources/H30/basic.pdf、PDF物理16頁。

\textbf{公式正答：} 5（single）。

\par\textit{原典本文の校正と独自解説は未作成。分類のみの登録項目。}

\subsection{平成30年度 問4-6}
\label{q:h30-reg-basic-4-6}

\noindent\textbf{原典：} sources/H30/basic.pdf、PDF物理17頁。

\textbf{公式正答：} 3（single）。

\par\textit{原典本文の校正と独自解説は未作成。分類のみの登録項目。}

\subsection{令和元年度 問4-1}
\label{q:r01-reg-basic-4-1}

\noindent\textbf{原典：} sources/R01/basic.pdf、PDF物理16頁。

\textbf{公式正答：} 4（single）。

\par\textit{原典本文の校正と独自解説は未作成。分類のみの登録項目。}

\subsection{令和元年度 問4-2}
\label{q:r01-reg-basic-4-2}

\noindent\textbf{原典：} sources/R01/basic.pdf、PDF物理16頁。

\textbf{公式正答：} 5（single）。

\par\textit{原典本文の校正と独自解説は未作成。分類のみの登録項目。}

\subsection{令和元年度 問4-3}
\label{q:r01-reg-basic-4-3}

\noindent\textbf{原典：} sources/R01/basic.pdf、PDF物理17頁。

\textbf{公式正答：} 4（single）。

\par\textit{原典本文の校正と独自解説は未作成。分類のみの登録項目。}

\subsection{令和元年度 問4-4}
\label{q:r01-reg-basic-4-4}

\noindent\textbf{原典：} sources/R01/basic.pdf、PDF物理17頁。

\textbf{公式正答：} 3（single）。

\par\textit{原典本文の校正と独自解説は未作成。分類のみの登録項目。}

\subsection{令和元年度 問4-5}
\label{q:r01-reg-basic-4-5}

\noindent\textbf{原典：} sources/R01/basic.pdf、PDF物理18頁。

\textbf{公式正答：} 4（single）。

\par\textit{原典本文の校正と独自解説は未作成。分類のみの登録項目。}

\subsection{令和元年度 問4-6}
\label{q:r01-reg-basic-4-6}

\noindent\textbf{原典：} sources/R01/basic.pdf、PDF物理18頁。

\textbf{公式正答：} 3（single）。

\par\textit{原典本文の校正と独自解説は未作成。分類のみの登録項目。}

\subsection{令和元年度再試験 問4-1}
\label{q:r01-ret-basic-4-1}

\noindent\textbf{原典：} sources/R01\_retest/basic.pdf、PDF物理14頁。

\textbf{公式正答：} 2（single）。

\par\textit{原典本文の校正と独自解説は未作成。分類のみの登録項目。}

\subsection{令和元年度再試験 問4-2}
\label{q:r01-ret-basic-4-2}

\noindent\textbf{原典：} sources/R01\_retest/basic.pdf、PDF物理14頁。

\textbf{公式正答：} 3（single）。

\par\textit{原典本文の校正と独自解説は未作成。分類のみの登録項目。}

\subsection{令和元年度再試験 問4-3}
\label{q:r01-ret-basic-4-3}

\noindent\textbf{原典：} sources/R01\_retest/basic.pdf、PDF物理14頁。

\textbf{公式正答：} 2（single）。

\par\textit{原典本文の校正と独自解説は未作成。分類のみの登録項目。}

\subsection{令和元年度再試験 問4-4}
\label{q:r01-ret-basic-4-4}

\noindent\textbf{原典：} sources/R01\_retest/basic.pdf、PDF物理15頁。

\textbf{公式正答：} 4（single）。

\par\textit{原典本文の校正と独自解説は未作成。分類のみの登録項目。}

\subsection{令和元年度再試験 問4-5}
\label{q:r01-ret-basic-4-5}

\noindent\textbf{原典：} sources/R01\_retest/basic.pdf、PDF物理15頁。

\textbf{公式正答：} 3（single）。

\par\textit{原典本文の校正と独自解説は未作成。分類のみの登録項目。}

\subsection{令和元年度再試験 問4-6}
\label{q:r01-ret-basic-4-6}

\noindent\textbf{原典：} sources/R01\_retest/basic.pdf、PDF物理16頁。

\textbf{公式正答：} 3（single）。

\par\textit{原典本文の校正と独自解説は未作成。分類のみの登録項目。}

\subsection{令和2年度 問4-1}
\label{q:r02-reg-basic-4-1}

\noindent\textbf{原典：} sources/R02/basic.pdf、PDF物理15頁。

\textbf{公式正答：} 2（single）。

\par\textit{原典本文の校正と独自解説は未作成。分類のみの登録項目。}

\subsection{令和2年度 問4-2}
\label{q:r02-reg-basic-4-2}

\noindent\textbf{原典：} sources/R02/basic.pdf、PDF物理15頁。

\textbf{公式正答：} 3（single）。

\par\textit{原典本文の校正と独自解説は未作成。分類のみの登録項目。}

\subsection{令和2年度 問4-3}
\label{q:r02-reg-basic-4-3}

\noindent\textbf{原典：} sources/R02/basic.pdf、PDF物理16頁。

\textbf{公式正答：} 4（single）。

\par\textit{原典本文の校正と独自解説は未作成。分類のみの登録項目。}

\subsection{令和2年度 問4-4}
\label{q:r02-reg-basic-4-4}

\noindent\textbf{原典：} sources/R02/basic.pdf、PDF物理17頁。

\textbf{公式正答：} 5（single）。

\par\textit{原典本文の校正と独自解説は未作成。分類のみの登録項目。}

\subsection{令和2年度 問4-5}
\label{q:r02-reg-basic-4-5}

\noindent\textbf{原典：} sources/R02/basic.pdf、PDF物理18頁。

\textbf{公式正答：} 3（single）。

\par\textit{原典本文の校正と独自解説は未作成。分類のみの登録項目。}

\subsection{令和2年度 問4-6}
\label{q:r02-reg-basic-4-6}

\noindent\textbf{原典：} sources/R02/basic.pdf、PDF物理18頁。

\textbf{公式正答：} 4（single）。

\par\textit{原典本文の校正と独自解説は未作成。分類のみの登録項目。}

\subsection{令和3年度 問4-1}
\label{q:r03-reg-basic-4-1}

\noindent\textbf{原典：} sources/R03/basic.pdf、PDF物理16頁。

\textbf{公式正答：} 3（single）。

\par\textit{原典本文の校正と独自解説は未作成。分類のみの登録項目。}

\subsection{令和3年度 問4-2}
\label{q:r03-reg-basic-4-2}

\noindent\textbf{原典：} sources/R03/basic.pdf、PDF物理16頁。

\textbf{公式正答：} 4（single）。

\par\textit{原典本文の校正と独自解説は未作成。分類のみの登録項目。}

\subsection{令和3年度 問4-3}
\label{q:r03-reg-basic-4-3}

\noindent\textbf{原典：} sources/R03/basic.pdf、PDF物理17頁。

\textbf{公式正答：} 2（single）。

\par\textit{原典本文の校正と独自解説は未作成。分類のみの登録項目。}

\subsection{令和3年度 問4-4}
\label{q:r03-reg-basic-4-4}

\noindent\textbf{原典：} sources/R03/basic.pdf、PDF物理18頁。

\textbf{公式正答：} 1（single）。

\par\textit{原典本文の校正と独自解説は未作成。分類のみの登録項目。}

\subsection{令和3年度 問4-5}
\label{q:r03-reg-basic-4-5}

\noindent\textbf{原典：} sources/R03/basic.pdf、PDF物理19頁。

\textbf{公式正答：} 1（single）。

\par\textit{原典本文の校正と独自解説は未作成。分類のみの登録項目。}

\subsection{令和3年度 問4-6}
\label{q:r03-reg-basic-4-6}

\noindent\textbf{原典：} sources/R03/basic.pdf、PDF物理20頁。

\textbf{公式正答：} 2（single）。

\par\textit{原典本文の校正と独自解説は未作成。分類のみの登録項目。}

\subsection{令和4年度 問4-1}
\label{q:r04-reg-basic-4-1}

\noindent\textbf{原典：} sources/R04/basic.pdf、PDF物理14頁。

\textbf{公式正答：} 2（single）。

\par\textit{原典本文の校正と独自解説は未作成。分類のみの登録項目。}

\subsection{令和4年度 問4-2}
\label{q:r04-reg-basic-4-2}

\noindent\textbf{原典：} sources/R04/basic.pdf、PDF物理14頁。

\textbf{公式正答：} 4（single）。

\par\textit{原典本文の校正と独自解説は未作成。分類のみの登録項目。}

\subsection{令和4年度 問4-3}
\label{q:r04-reg-basic-4-3}

\noindent\textbf{原典：} sources/R04/basic.pdf、PDF物理14頁。

\textbf{公式正答：} 1（single）。

\par\textit{原典本文の校正と独自解説は未作成。分類のみの登録項目。}

\subsection{令和4年度 問4-4}
\label{q:r04-reg-basic-4-4}

\noindent\textbf{原典：} sources/R04/basic.pdf、PDF物理15頁。

\textbf{公式正答：} 1（single）。

\par\textit{原典本文の校正と独自解説は未作成。分類のみの登録項目。}

\subsection{令和4年度 問4-5}
\label{q:r04-reg-basic-4-5}

\noindent\textbf{原典：} sources/R04/basic.pdf、PDF物理15頁。

\textbf{公式正答：} 5（single）。

\par\textit{原典本文の校正と独自解説は未作成。分類のみの登録項目。}

\subsection{令和4年度 問4-6}
\label{q:r04-reg-basic-4-6}

\noindent\textbf{原典：} sources/R04/basic.pdf、PDF物理16頁。

\textbf{公式正答：} 5（single）。

\par\textit{原典本文の校正と独自解説は未作成。分類のみの登録項目。}

\subsection{令和5年度 問4-1}
\label{q:r05-reg-basic-4-1}

\noindent\textbf{原典：} sources/R05/basic.pdf、PDF物理13頁。

\textbf{公式正答：} 3（single）。

\par\textit{原典本文の校正と独自解説は未作成。分類のみの登録項目。}

\subsection{令和5年度 問4-2}
\label{q:r05-reg-basic-4-2}

\noindent\textbf{原典：} sources/R05/basic.pdf、PDF物理13頁。

\textbf{公式正答：} 5（single）。

\par\textit{原典本文の校正と独自解説は未作成。分類のみの登録項目。}

\subsection{令和5年度 問4-3}
\label{q:r05-reg-basic-4-3}

\noindent\textbf{原典：} sources/R05/basic.pdf、PDF物理13頁。

\textbf{公式正答：} 3（single）。

\par\textit{原典本文の校正と独自解説は未作成。分類のみの登録項目。}

\subsection{令和5年度 問4-4}
\label{q:r05-reg-basic-4-4}

\noindent\textbf{原典：} sources/R05/basic.pdf、PDF物理14頁。

\textbf{公式正答：} 1（single）。

\par\textit{原典本文の校正と独自解説は未作成。分類のみの登録項目。}

\subsection{令和5年度 問4-5}
\label{q:r05-reg-basic-4-5}

\noindent\textbf{原典：} sources/R05/basic.pdf、PDF物理14頁。

\textbf{公式正答：} 3（single）。

\par\textit{原典本文の校正と独自解説は未作成。分類のみの登録項目。}

\subsection{令和5年度 問4-6}
\label{q:r05-reg-basic-4-6}

\noindent\textbf{原典：} sources/R05/basic.pdf、PDF物理15頁。

\textbf{公式正答：} 5（single）。

\par\textit{原典本文の校正と独自解説は未作成。分類のみの登録項目。}

\subsection{令和6年度 問4-1}
\label{q:r06-reg-basic-4-1}

\noindent\textbf{原典：} sources/R06/basic.pdf、PDF物理15頁。

\textbf{公式正答：} 1（single）。

\par\textit{原典本文の校正と独自解説は未作成。分類のみの登録項目。}

\subsection{令和6年度 問4-2}
\label{q:r06-reg-basic-4-2}

\noindent\textbf{原典：} sources/R06/basic.pdf、PDF物理15頁。

\textbf{公式正答：} 3（single）。

\par\textit{原典本文の校正と独自解説は未作成。分類のみの登録項目。}

\subsection{令和6年度 問4-3}
\label{q:r06-reg-basic-4-3}

\noindent\textbf{原典：} sources/R06/basic.pdf、PDF物理15頁。

\textbf{公式正答：} 3（single）。

\par\textit{原典本文の校正と独自解説は未作成。分類のみの登録項目。}

\subsection{令和6年度 問4-4}
\label{q:r06-reg-basic-4-4}

\noindent\textbf{原典：} sources/R06/basic.pdf、PDF物理16頁。

\textbf{公式正答：} 2（single）。

\par\textit{原典本文の校正と独自解説は未作成。分類のみの登録項目。}

\subsection{令和6年度 問4-5}
\label{q:r06-reg-basic-4-5}

\noindent\textbf{原典：} sources/R06/basic.pdf、PDF物理16頁。

\textbf{公式正答：} 1（single）。

\par\textit{原典本文の校正と独自解説は未作成。分類のみの登録項目。}

\subsection{令和6年度 問4-6}
\label{q:r06-reg-basic-4-6}

\noindent\textbf{原典：} sources/R06/basic.pdf、PDF物理16頁。

\textbf{公式正答：} 4（single）。

\par\textit{原典本文の校正と独自解説は未作成。分類のみの登録項目。}

\subsection{令和7年度 問4-1}
\label{q:r07-reg-basic-4-1}

\noindent\textbf{原典：} sources/R07/basic.pdf、PDF物理13頁。

\textbf{公式正答：} 2（single）。

\par\textit{原典本文の校正と独自解説は未作成。分類のみの登録項目。}

\subsection{令和7年度 問4-2}
\label{q:r07-reg-basic-4-2}

\noindent\textbf{原典：} sources/R07/basic.pdf、PDF物理13頁。

\textbf{公式正答：} 3（single）。

\par\textit{原典本文の校正と独自解説は未作成。分類のみの登録項目。}

\subsection{令和7年度 問4-3}
\label{q:r07-reg-basic-4-3}

\noindent\textbf{原典：} sources/R07/basic.pdf、PDF物理13頁。

\textbf{公式正答：} 3（single）。

\par\textit{原典本文の校正と独自解説は未作成。分類のみの登録項目。}

\subsection{令和7年度 問4-4}
\label{q:r07-reg-basic-4-4}

\noindent\textbf{原典：} sources/R07/basic.pdf、PDF物理14頁。

\textbf{公式正答：} 5（single）。

\par\textit{原典本文の校正と独自解説は未作成。分類のみの登録項目。}

\subsection{令和7年度 問4-5}
\label{q:r07-reg-basic-4-5}

\noindent\textbf{原典：} sources/R07/basic.pdf、PDF物理14頁。

\textbf{公式正答：} 4（single）。

\par\textit{原典本文の校正と独自解説は未作成。分類のみの登録項目。}

\subsection{令和7年度 問4-6}
\label{q:r07-reg-basic-4-6}

\noindent\textbf{原典：} sources/R07/basic.pdf、PDF物理14頁。

\textbf{公式正答：} 4（single）。

\par\textit{原典本文の校正と独自解説は未作成。分類のみの登録項目。}


\chapter{環境・エネルギー・技術}
\noindent 分野解説は執筆中。以下は登録問題と校正済み例の作業用一覧。
\section{環境会計と評価・報告・監査を区別}

\subsection{平成26年度 問5-2}
\label{q:h26-reg-basic-5-2}

\noindent\textbf{原典：} sources/H26/basic.pdf、PDF物理16頁。

\textbf{公式正答：} 5（single）。

\paragraph{問題の要旨} 環境会計と評価・報告・監査を区別。

\Needspace{4\baselineskip}\paragraph{確認した条件・記述（独自要約）}\mbox{}\par

\noindent\textbf{時点：} 平成26年度に指定された資料・制度。\par

\noindent\textbf{照合：} 原典21画像で否定・五肢・組合せ表を確認。後年値で置換しない。\par

\paragraph{問われる核心} 環境会計と評価・報告・監査を区別。

\paragraph{必要知識} 対象、評価目的、会計単位、環境管理手段。

\paragraph{考え方・正答を選ぶ理由（査読待ち初稿）} 環境会計はコストと効果を認識して、貨幣単位又は物量単位で定量化・伝達する仕組みで5。LCAは設備の更新周期だけでなく製品等の全過程。環境報告は社会への情報開示、グリーン購入は必要性と環境負荷、環境監査は環境管理の点検であり、税務監査だけではない。

\Needspace{4\baselineskip}\paragraph{選択肢の検討（査読待ち初稿）}\mbox{}\par

\noindent\textbf{1：不適切} LCAは設備更新時のみ。式・定義・対応の独立判定と一致しない。\par

\noindent\textbf{2：不適切} 環境報告書は設備の自治体届出。式・定義・対応の独立判定と一致しない。\par

\noindent\textbf{3：不適切} グリーン購入は生物資源優先のみ。式・定義・対応の独立判定と一致しない。\par

\noindent\textbf{4：不適切} 環境監査は税務処理のみ。式・定義・対応の独立判定と一致しない。\par

\noindent\textbf{5：適切・選択対象} 環境会計は費用効果を金額物量で。判定と一致する。\par

\paragraph{再利用する考え方} 対象、評価目的、会計単位、環境管理手段。 指定年度・対象範囲・単位・否定を先に固定し、他年度の番号や現行値を転用しない。

\paragraph{共通理論} \hyperref[theory:env:accounting]{第\ref*{theory:env:accounting}節}

\par 参考：\cite{env58_accounting}。

\par 難易度 2（校準初稿）、優先度未評価。

\subsection{平成30年度 問5-2}
\label{q:h30-reg-basic-5-2}

\noindent\textbf{原典：} sources/H30/basic.pdf、PDF物理18頁。

\textbf{公式正答：} 3（single）。

\paragraph{問題の要旨} 環境会計の定義と他手法の誤対応。

\Needspace{4\baselineskip}\paragraph{確認した条件・記述（独自要約）}\mbox{}\par

\noindent\textbf{時点：} 平成30年度に指定された資料・制度。\par

\noindent\textbf{照合：} 原典18画像で否定・五肢・組合せ表を確認。後年値で置換しない。\par

\paragraph{問われる核心} 環境会計の定義と他手法の誤対応。

\paragraph{必要知識} 目的・対象・単位による管理手法の区別。

\paragraph{考え方・正答を選ぶ理由（査読待ち初稿）} 3の費用と効果を貨幣又は物量単位で表す仕組みが環境会計。グリーン購入は必要性と環境負荷全般、環境報告は社会への情報開示、環境監査は環境管理の評価、LCAは全過程の影響評価で、残る肢の限定が違う。

\Needspace{4\baselineskip}\paragraph{選択肢の検討（査読待ち初稿）}\mbox{}\par

\noindent\textbf{1：不適切} グリーン購入はバイオマスのみ。式・定義・対応の独立判定と一致しない。\par

\noindent\textbf{2：不適切} 環境報告は設備の自治体届出。式・定義・対応の独立判定と一致しない。\par

\noindent\textbf{3：適切・選択対象} 会計は費用効果を金額物量。判定と一致する。\par

\noindent\textbf{4：不適切} 環境監査は税務監査。式・定義・対応の独立判定と一致しない。\par

\noindent\textbf{5：不適切} LCAは設備の更新周期。式・定義・対応の独立判定と一致しない。\par

\paragraph{再利用する考え方} 目的・対象・単位による管理手法の区別。 集計対象・時点・法の主体・選択方向を固定し、同型でも元の五肢へ戻す。

\paragraph{共通理論} \hyperref[theory:env:accounting]{第\ref*{theory:env:accounting}節}

\par 参考：\cite{env58_accounting}。

\par 難易度 2（校準初稿）、優先度未評価。

\section{粒径範囲と大気汚染の生成経路}

\subsection{平成26年度 問5-1}
\label{q:h26-reg-basic-5-1}

\noindent\textbf{原典：} sources/H26/basic.pdf、PDF物理16頁。

\textbf{公式正答：} 1（single）。

\par\textbf{検討未確定：} COによるメタン寿命の化学過程と各健康影響は追加一次資料照合。肺沈着と粒径の関係を2.5μm一値へ固定しない。

\paragraph{問題の要旨} 粒径範囲と大気汚染の生成経路。

\Needspace{4\baselineskip}\paragraph{確認した条件・記述（独自要約）}\mbox{}\par

\noindent\textbf{時点：} 平成26年度に指定された資料・制度。\par

\noindent\textbf{照合：} 原典21画像で否定・五肢・組合せ表を確認。後年値で置換しない。\par

\paragraph{問われる核心} 粒径範囲と大気汚染の生成経路。

\paragraph{必要知識} 大きさの範囲と測定、一次・二次汚染物質、生成と除去。

\paragraph{考え方・正答を選ぶ理由（査読待ち初稿）} PM2.5は概ね粒径2.5μm以下の微小粒子で、2.5付近だけに限定する1が誤り。規格上は2.5μmを50\%分離する装置で採取するという測定定義もある。SO2は硫黄、NOxとVOCは二次生成、COは不完全燃焼とヘモグロビンの結合。COによるOH消費はメタン除去に影響する。

\Needspace{4\baselineskip}\paragraph{選択肢の検討（査読待ち初稿）}\mbox{}\par

\noindent\textbf{1：不適切・選択対象} PM2.5は2.5μm付近だけ。判定と一致する。\par

\noindent\textbf{2：条件内で適切} SO2は燃料中S由来。原典の条件内では適切な記述であり、不適切なものを選ぶ問の選択対象ではない。\par

\noindent\textbf{3：条件内で適切} NO2は燃焼と光化学前駆物質。原典の条件内では適切な記述であり、不適切なものを選ぶ問の選択対象ではない。\par

\noindent\textbf{4：条件内で適切} 光化学オキシダントはNOxとVOCの光反応。原典の条件内では適切な記述であり、不適切なものを選ぶ問の選択対象ではない。\par

\noindent\textbf{5：条件内で適切} COは不完全燃焼でメタン寿命にも影響。原典の条件内では適切な記述であり、不適切なものを選ぶ問の選択対象ではない。\par

\paragraph{再利用する考え方} 大きさの範囲と測定、一次・二次汚染物質、生成と除去。 指定年度・対象範囲・単位・否定を先に固定し、他年度の番号や現行値を転用しない。

\paragraph{共通理論} \hyperref[theory:env:air]{第\ref*{theory:env:air}節}

\par 参考：\cite{env58_pmstandard,env58_pm25}。

\par 難易度 3（校準初稿）、優先度未評価。

\subsection{令和5年度 問5-2}
\label{q:r05-reg-basic-5-2}

\noindent\textbf{原典：} sources/R05/basic.pdf、PDF物理16頁。

\textbf{公式正答：} 5（single）。

\paragraph{問題の要旨} PM2.5は2.5μm付近だけではない。

\Needspace{4\baselineskip}\paragraph{確認した条件・記述（独自要約）}\mbox{}\par

\noindent\textbf{原典：} 全五肢と否定方向・組合せ表を画像確認。\par

\noindent\textbf{資料時点：} 令和5年度の指定資料・統計年を保持し、後年値に置換しない。\par

\paragraph{問われる核心} PM2.5は2.5μm付近だけではない。

\paragraph{必要知識} PM分粒、一次排出と二次生成、SO2・NO2・CO・Ox。

\paragraph{考え方・正答を選ぶ理由（査読待ち初稿）} 5が不適切。2.5という値は代表粒径一種類を指すものではなく粒径区分である。各汚染物質は燃料由来・不完全燃焼・大気中の二次生成を分ける。

\Needspace{4\baselineskip}\paragraph{選択肢の検討（査読待ち初稿）}\mbox{}\par

\noindent\textbf{1：条件内で適切} SO2は硫黄燃料燃焼と酸性雨。発生源の硫黄と反応生成物を対応。\par

\noindent\textbf{2：条件内で適切} NOが酸化しNO2となる。燃焼で生じたNOの大気酸化はNO2の生成経路。\par

\noindent\textbf{3：条件内で適切} COは不完全燃焼・酸素運搬阻害。COがヘモグロビンと結合して酸素運搬を妨げる。\par

\noindent\textbf{4：条件内で適切} 光化学オキシダントはNOx・VOC等から二次生成。日射下の化学反応で生じる酸化性物質の総称。\par

\noindent\textbf{5：不適切・選択対象} PM2.5は粒径2.5μm付近のみ。微小粒子の区分は2.5μm以下側の粒径範囲であり、その近傍だけではない。実測は分粒特性で定義される。\par

\paragraph{再利用する考え方} PM分粒、一次排出と二次生成、SO2・NO2・CO・Ox。 否定方向と資料年、分母と単位を確認して五肢へ戻る。

\paragraph{共通理論} \hyperref[theory:env:air]{第\ref*{theory:env:air}節}

\par 参考：\cite{env58_pm25}。

\par 難易度 2（校準初稿）、優先度未評価。

\section{産業廃棄物フローを物質収支で検算}

\subsection{平成28年度 問5-1}
\label{q:h28-reg-basic-5-1}

\noindent\textbf{原典：} sources/H28/basic.pdf、PDF物理21頁。

\textbf{公式正答：} 5（single）。

\paragraph{問題の要旨} 産業廃棄物フローを物質収支で検算。

\Needspace{4\baselineskip}\paragraph{確認した条件・記述（独自要約）}\mbox{}\par

\noindent\textbf{排出：} 379百万t。\par

\noindent\textbf{中間処理：} 290。\par

\noindent\textbf{直接再生：} 83。\par

\noindent\textbf{処理後再生：} 124。\par

\noindent\textbf{処理後最終処分：} 7。\par

\noindent\textbf{全再生：} 207。\par

\noindent\textbf{全最終処分：} 13。\par

\paragraph{問われる核心} 産業廃棄物フローを物質収支で検算。

\paragraph{必要知識} 分岐の収支、直接と処理後、百分率と比率。

\paragraph{考え方・正答を選ぶ理由（査読待ち初稿）} 単位百万t。中間処理290は再利用124、処理後処分7、減量159へ分かれる。総処分13から直接処分は13−7=\allowbreak{}6で、60ではないので5。総排出も83+290+6=\allowbreak{}379、最終の207+159+13=\allowbreak{}379で一致。83/207≈.40、207/379≈.55、207/13≈16、124/83≈1.5、159/379≈.42。

\Needspace{4\baselineskip}\paragraph{選択肢の検討（査読待ち初稿）}\mbox{}\par

\noindent\textbf{1：条件内で適切} 直接再利用83で約40\%。原典の条件内では適切な記述であり、不適切なものを選ぶ問の選択対象ではない。\par

\noindent\textbf{2：条件内で適切} 再利用207で55\%・処分の16倍。原典の条件内では適切な記述であり、不適切なものを選ぶ問の選択対象ではない。\par

\noindent\textbf{3：条件内で適切} 処理後124で直接の1.5倍。原典の条件内では適切な記述であり、不適切なものを選ぶ問の選択対象ではない。\par

\noindent\textbf{4：条件内で適切} 減量159で42\%。原典の条件内では適切な記述であり、不適切なものを選ぶ問の選択対象ではない。\par

\noindent\textbf{5：不適切・選択対象} 直接処分60で16\%。判定と一致する。\par

\paragraph{再利用する考え方} 分岐の収支、直接と処理後、百分率と比率。 指定年度・対象範囲・単位・否定を先に固定し、他年度の番号や現行値を転用しない。

\paragraph{共通理論} \hyperref[theory:env:balance]{第\ref*{theory:env:balance}節}

\paragraph{関連問題} \hyperref[q:r06-reg-basic-5-2]{R06-REG-BASIC-5-2}

\par 難易度 2（校準初稿）、優先度未評価。

\subsection{令和6年度 問5-2}
\label{q:r06-reg-basic-5-2}

\noindent\textbf{原典：} sources/R06/basic.pdf、PDF物理18頁。

\textbf{公式正答：} 5（single）。

\paragraph{問題の要旨} 2020産業廃棄物フローの再生比は約53\%。

\Needspace{4\baselineskip}\paragraph{確認した条件・記述（独自要約）}\mbox{}\par

\noindent\textbf{単位：} 百万トン。\par

\noindent\textbf{原図：} 2020年度の概算値。減量を再生や処分へ加えない。\par

\paragraph{問われる核心} 2020産業廃棄物フローの再生比は約53\%。

\paragraph{必要知識} フロー保存、枝の量と総量、比の分母、複合命題。

\paragraph{考え方・正答を選ぶ理由（査読待ち初稿）} 5が不適切。図は排出374、中間293、直接再生77、中間後再生122、中間後処分5、総再生199、総処分9（百万t）。減量166、直接処分4を保存則から求める。374=\allowbreak{}199+9+166。5の二命題を分け、比の後半だけで適切としない。

\Needspace{4\baselineskip}\paragraph{選択肢の検討（査読待ち初稿）}\mbox{}\par

\noindent\textbf{1：条件内で適切} 直接再生77は再生199の約40\%。77/199≈38.7\%。約40\%に整合。\par

\noindent\textbf{2：条件内で適切} 中間後再生122は処理293の約40\%。122/293≈41.6\%。概算40\%に整合。\par

\noindent\textbf{3：条件内で適切} 減量166は排出374の約45\%。293−122−5=\allowbreak{}166、166/374≈44.4\%。\par

\noindent\textbf{4：条件内で適切} 直接最終4は排出374の約1\%。9−5=\allowbreak{}4、4/374≈1.07\%。\par

\noindent\textbf{5：不適切・選択対象} 再生199は排出374の約20\%・最終9の約22倍。199/374≈53.2\%であり20\%ではない。199/9≈22.1倍の後半が正しくても一肢全体は不適切。\par

\paragraph{再利用する考え方} フロー保存、枝の量と総量、比の分母、複合命題。 否定方向と資料年、分母と単位を確認して五肢へ戻る。

\paragraph{共通理論} \hyperref[theory:env:balance]{第\ref*{theory:env:balance}節}

\paragraph{関連問題} \hyperref[q:h28-reg-basic-5-1]{H28-REG-BASIC-5-1}

\par 難易度 2（校準初稿）、優先度未評価。

\section{生物多様性条約と外来種防除}

\subsection{平成28年度 問5-2}
\label{q:h28-reg-basic-5-2}

\noindent\textbf{原典：} sources/H28/basic.pdf、PDF物理22頁。

\textbf{公式正答：} 5（single）。

\par\textbf{検討未確定：} 1992・1993と議定書、当時の外来生物法の防除条項の一次本文を追加照合。動物愛護上の配慮が不要という意味にはしない。

\paragraph{問題の要旨} 生物多様性条約と外来種防除。

\Needspace{4\baselineskip}\paragraph{確認した条件・記述（独自要約）}\mbox{}\par

\noindent\textbf{時点：} 平成28年度に指定された資料・制度。\par

\noindent\textbf{照合：} 原典21画像で否定・五肢・組合せ表を確認。後年値で置換しない。\par

\paragraph{問われる核心} 生物多様性条約と外来種防除。

\paragraph{必要知識} 条約目的、署名と発効、外来種と防除。

\paragraph{考え方・正答を選ぶ理由（査読待ち初稿）} 外来種による生態系被害の防止には防除制度があり、一律に駆除を禁止する5が誤り。CBDの三目的は保全、持続可能な利用、遺伝資源の利益の公正衡平配分。組換え生物の扱いはCartagena議定書であり、遺伝資源の利益配分のNagoya議定書と混同しない。

\Needspace{4\baselineskip}\paragraph{選択肢の検討（査読待ち初稿）}\mbox{}\par

\noindent\textbf{1：条件内で適切} 組換え生物の議定書。原典の条件内では適切な記述であり、不適切なものを選ぶ問の選択対象ではない。\par

\noindent\textbf{2：条件内で適切} 1992署名開放・1993発効。原典の条件内では適切な記述であり、不適切なものを選ぶ問の選択対象ではない。\par

\noindent\textbf{3：条件内で適切} 保全・持続利用・利益配分。原典の条件内では適切な記述であり、不適切なものを選ぶ問の選択対象ではない。\par

\noindent\textbf{4：条件内で適切} 外来種に重大影響の例。原典の条件内では適切な記述であり、不適切なものを選ぶ問の選択対象ではない。\par

\noindent\textbf{5：不適切・選択対象} 日本は動物愛護で駆除禁止。判定と一致する。\par

\paragraph{再利用する考え方} 条約目的、署名と発効、外来種と防除。 指定年度・対象範囲・単位・否定を先に固定し、他年度の番号や現行値を転用しない。

\paragraph{共通理論} \hyperref[theory:env:biodiversity]{第\ref*{theory:env:biodiversity}節}

\par 難易度 3（校準初稿）、優先度未評価。

\subsection{令和2年度 問5-2}
\label{q:r02-reg-basic-5-2}

\noindent\textbf{原典：} sources/R02/basic.pdf、PDF物理20頁。

\textbf{公式正答：} 3（single）。

\par\textbf{検討未確定：} 出題時の防除条項とCBD全三目的の原条文照合は追加確認。

\paragraph{問題の要旨} 生物多様性保全と外来種防除。

\Needspace{4\baselineskip}\paragraph{確認した条件・記述（独自要約）}\mbox{}\par

\noindent\textbf{時点：} 令和2年度に指定された資料・制度。\par

\noindent\textbf{照合：} 原典18画像で否定・五肢・組合せ表を確認。後年値で置換しない。\par

\paragraph{問われる核心} 生物多様性保全と外来種防除。

\paragraph{必要知識} 条約と議定書、外来種、防除、倫理配慮。

\paragraph{考え方・正答を選ぶ理由（査読待ち初稿）} 3の一律駆除禁止が不適切。防除制度と扱いへの配慮を併せて考える。CBDの三目的とCartagenaの組換え生物の扱いは対象が異なり、Nagoyaの利益配分とも区別する。

\Needspace{4\baselineskip}\paragraph{選択肢の検討（査読待ち初稿）}\mbox{}\par

\noindent\textbf{1：条件内で適切} 組換え生物の国際枠組。条件内で適切な記述であり、不適切な記述を選ぶ対象ではない。\par

\noindent\textbf{2：条件内で適切} 外来種が在来種に影響。条件内で適切な記述であり、不適切な記述を選ぶ対象ではない。\par

\noindent\textbf{3：不適切・選択対象} 駆除は動物愛護で禁止。判定と一致する。\par

\noindent\textbf{4：条件内で適切} 1992署名開放1993発効。条件内で適切な記述であり、不適切な記述を選ぶ対象ではない。\par

\noindent\textbf{5：条件内で適切} 保全持続利用利益配分。条件内で適切な記述であり、不適切な記述を選ぶ対象ではない。\par

\paragraph{再利用する考え方} 条約と議定書、外来種、防除、倫理配慮。 集計対象・時点・法の主体・選択方向を固定し、同型でも元の五肢へ戻す。

\paragraph{共通理論} \hyperref[theory:env:biodiversity]{第\ref*{theory:env:biodiversity}節}

\par 難易度 3（校準初稿）、優先度未評価。

\section{2008年度の温室効果ガスと国際枠組}

\subsection{平成23年度 問5-1}
\label{q:h23-reg-basic-5-1}

\noindent\textbf{原典：} sources/H23/basic.pdf、PDF物理15頁。

\textbf{公式正答：} 1（single）。

\par\textbf{検討未確定：} 当時白書の全ガス内訳・部門図と2009合意の条文の追加照合が必要。CO2九割以上は当時確定値公表表の検索本文でも確認。

\paragraph{問題の要旨} 2008年度の温室効果ガスと国際枠組。

\Needspace{4\baselineskip}\paragraph{確認した条件・記述（独自要約）}\mbox{}\par

\noindent\textbf{時点：} 平成23年度に指定された資料・制度。\par

\noindent\textbf{照合：} 原典21画像で否定・五肢・組合せ表を確認。後年値で置換しない。\par

\paragraph{問われる核心} 2008年度の温室効果ガスと国際枠組。

\paragraph{必要知識} CO2換算、部門の定義、条約と議定書、採択年。

\paragraph{考え方・正答を選ぶ理由（査読待ち初稿）} CO2換算の2008年度総量の九割以上がCO2で、1。部門別では家庭が第二位という組合せが違う。UNFCCCの究極目的に2050という一律期限はなく、1997採択は京都議定書。2009コペンハーゲンの合意を全途上国に法的数量削減義務を課す議定書と扱わない。

\Needspace{4\baselineskip}\paragraph{選択肢の検討（査読待ち初稿）}\mbox{}\par

\noindent\textbf{1：適切・選択対象} CO2が九割以上。判定と一致する。\par

\noindent\textbf{2：不適切} 産業の次は家庭。式・定義・対応の独立判定と一致しない。\par

\noindent\textbf{3：不適切} 条約の濃度安定は2050期限。式・定義・対応の独立判定と一致しない。\par

\noindent\textbf{4：不適切} 1997パリ議定書。式・定義・対応の独立判定と一致しない。\par

\noindent\textbf{5：不適切} 2009途上国に基準年削減数値義務。式・定義・対応の独立判定と一致しない。\par

\paragraph{再利用する考え方} CO2換算、部門の定義、条約と議定書、採択年。 指定年度・対象範囲・単位・否定を先に固定し、他年度の番号や現行値を転用しない。

\paragraph{共通理論} \hyperref[theory:env:climate]{第\ref*{theory:env:climate}節}

\par 参考：\cite{env58_ghg2008}。

\par 難易度 3（校準初稿）、優先度未評価。

\section{暖房と自動車のCO2排出量を比較}

\subsection{平成25年度 問5-2}
\label{q:h25-reg-basic-5-2}

\noindent\textbf{原典：} sources/H25/basic.pdf、PDF物理19頁。

\textbf{公式正答：} 2（single）。

\paragraph{問題の要旨} 暖房と自動車のCO2排出量を比較。

\Needspace{4\baselineskip}\paragraph{確認した条件・記述（独自要約）}\mbox{}\par

\noindent\textbf{電力：} 500W、3h/日×50日、.42kg-CO2/kWh。\par

\noindent\textbf{車：} 10km/L、200km、2.32kg-CO2/L。\par

\paragraph{問われる核心} 暖房と自動車のCO2排出量を比較。

\paragraph{必要知識} エネルギー=電力×時間、燃料使用量、排出係数。

\paragraph{考え方・正答を選ぶ理由（査読待ち初稿）} 暖房は.5kW×3h×50=\allowbreak{}75kWh、排出31.5kg。自動車は200/10=\allowbreak{}20L、排出46.4kg。比31.5/46.4≈.679で2:3に最も近い。WをkWへ直してから時間を掛ける。

\Needspace{4\baselineskip}\paragraph{選択肢の検討（査読待ち初稿）}\mbox{}\par

\noindent\textbf{1：不適切} 1:3。式・定義・対応の独立判定と一致しない。\par

\noindent\textbf{2：適切・選択対象} 2:3。判定と一致する。\par

\noindent\textbf{3：不適切} 1:1。式・定義・対応の独立判定と一致しない。\par

\noindent\textbf{4：不適切} 3:2。式・定義・対応の独立判定と一致しない。\par

\noindent\textbf{5：不適切} 3:1。式・定義・対応の独立判定と一致しない。\par

\paragraph{再利用する考え方} エネルギー=電力×時間、燃料使用量、排出係数。 指定年度・対象範囲・単位・否定を先に固定し、他年度の番号や現行値を転用しない。

\paragraph{共通理論} \hyperref[theory:env:co2]{第\ref*{theory:env:co2}節}

\par 難易度 2（校準初稿）、優先度未評価。

\subsection{平成27年度 問5-3}
\label{q:h27-reg-basic-5-3}

\noindent\textbf{原典：} sources/H27/basic.pdf、PDF物理16頁。

\textbf{公式正答：} 5（single）。

\paragraph{問題の要旨} 利用人数と時間を含めた一人当たり排出。

\Needspace{4\baselineskip}\paragraph{確認した条件・記述（独自要約）}\mbox{}\par

\noindent\textbf{a：} 10km/L、100km、3人、2.32kg/L。\par

\noindent\textbf{b：} 8000kW、200km/h、500km、1000人、.47kg/kWh。\par

\noindent\textbf{c：} 500W、4h/日、50日、4人、.47kg/kWh。\par

\paragraph{問われる核心} 利用人数と時間を含めた一人当たり排出。

\paragraph{必要知識} 排出係数、乗車人数、走行時間、総量と原単位。

\paragraph{考え方・正答を選ぶ理由（査読待ち初稿）} aは100/10×2.32/3=\allowbreak{}7.73kg/人。bは500/200=\allowbreak{}2.5h、8000×2.5×.47/1000=\allowbreak{}9.4kg/人。cは.5×4×50×.47/4=\allowbreak{}11.75kg/人。c>b>aの5。総量の大小ではなく人数で割った量を比べる。

\Needspace{4\baselineskip}\paragraph{選択肢の検討（査読待ち初稿）}\mbox{}\par

\noindent\textbf{1：不適切} a>b>c。式・定義・対応の独立判定と一致しない。\par

\noindent\textbf{2：不適切} a>c>b。式・定義・対応の独立判定と一致しない。\par

\noindent\textbf{3：不適切} b>a>c。式・定義・対応の独立判定と一致しない。\par

\noindent\textbf{4：不適切} c>a>b。式・定義・対応の独立判定と一致しない。\par

\noindent\textbf{5：適切・選択対象} c>b>a。判定と一致する。\par

\paragraph{再利用する考え方} 排出係数、乗車人数、走行時間、総量と原単位。 指定年度・対象範囲・単位・否定を先に固定し、他年度の番号や現行値を転用しない。

\paragraph{共通理論} \hyperref[theory:env:co2]{第\ref*{theory:env:co2}節}

\par 難易度 2（校準初稿）、優先度未評価。

\section{科学と社会をつなぐ双方向の対話}

\subsection{平成25年度 問5-5}
\label{q:h25-reg-basic-5-5}

\noindent\textbf{原典：} sources/H25/basic.pdf、PDF物理22頁。

\textbf{公式正答：} 2（single）。

\paragraph{問題の要旨} 科学と社会をつなぐ双方向の対話。

\Needspace{4\baselineskip}\paragraph{確認した条件・記述（独自要約）}\mbox{}\par

\noindent\textbf{四文の順序：} ア=社会理解、イ=科学技術間だけ、ウ=双方向、エ=媒体へ技術使用\par

\noindent\textbf{問：} 適切な二文。\par

\paragraph{問われる核心} 科学と社会をつなぐ双方向の対話。

\paragraph{必要知識} 専門家と社会、理解と対話、情報発信と相互学習。

\paragraph{考え方・正答を選ぶ理由（査読待ち初稿）} 社会・国民の期待や考えを専門家自身が理解することと、専門家・非専門家間の双方向性が適切。基礎科学と応用技術の内部連携だけをこの領域全体の定義にはできず、媒体制作へ技術を使うことや科学者へ任せることも、社会と対話するという目的の代替ではない。原典の列順で組合せを選ぶ。

\Needspace{4\baselineskip}\paragraph{選択肢の検討（査読待ち初稿）}\mbox{}\par

\noindent\textbf{1：不適切} ア イ。式・定義・対応の独立判定と一致しない。\par

\noindent\textbf{2：適切・選択対象} ア ウ。判定と一致する。\par

\noindent\textbf{3：不適切} イ ウ。式・定義・対応の独立判定と一致しない。\par

\noindent\textbf{4：不適切} イ エ。式・定義・対応の独立判定と一致しない。\par

\noindent\textbf{5：不適切} ウ エ。式・定義・対応の独立判定と一致しない。\par

\paragraph{再利用する考え方} 専門家と社会、理解と対話、情報発信と相互学習。 指定年度・対象範囲・単位・否定を先に固定し、他年度の番号や現行値を転用しない。

\paragraph{共通理論} \hyperref[theory:env:communication]{第\ref*{theory:env:communication}節}

\par 参考：\cite{env58_communication}。

\par 難易度 2（校準初稿）、優先度未評価。

\subsection{平成28年度 問5-5}
\label{q:h28-reg-basic-5-5}

\noindent\textbf{原典：} sources/H28/basic.pdf、PDF物理24頁。

\textbf{公式正答：} 5（single）。

\paragraph{問題の要旨} 科学と社会をつなぐ双方向の対話。

\Needspace{4\baselineskip}\paragraph{確認した条件・記述（独自要約）}\mbox{}\par

\noindent\textbf{四文の順序：} ア=科学技術間だけ、イ=科学者に制作委任、ウ=社会理解、エ=双方向\par

\noindent\textbf{問：} 適切な二文。\par

\paragraph{問われる核心} 科学と社会をつなぐ双方向の対話。

\paragraph{必要知識} 専門家と社会、理解と対話、情報発信と相互学習。

\paragraph{考え方・正答を選ぶ理由（査読待ち初稿）} 社会・国民の期待や考えを専門家自身が理解することと、専門家・非専門家間の双方向性が適切。基礎科学と応用技術の内部連携だけをこの領域全体の定義にはできず、媒体制作へ技術を使うことや科学者へ任せることも、社会と対話するという目的の代替ではない。原典の列順で組合せを選ぶ。

\Needspace{4\baselineskip}\paragraph{選択肢の検討（査読待ち初稿）}\mbox{}\par

\noindent\textbf{1：不適切} ア イ。式・定義・対応の独立判定と一致しない。\par

\noindent\textbf{2：不適切} ア ウ。式・定義・対応の独立判定と一致しない。\par

\noindent\textbf{3：不適切} イ ウ。式・定義・対応の独立判定と一致しない。\par

\noindent\textbf{4：不適切} イ エ。式・定義・対応の独立判定と一致しない。\par

\noindent\textbf{5：適切・選択対象} ウ エ。判定と一致する。\par

\paragraph{再利用する考え方} 専門家と社会、理解と対話、情報発信と相互学習。 指定年度・対象範囲・単位・否定を先に固定し、他年度の番号や現行値を転用しない。

\paragraph{共通理論} \hyperref[theory:env:communication]{第\ref*{theory:env:communication}節}

\par 参考：\cite{env58_communication}。

\par 難易度 2（校準初稿）、優先度未評価。

\section{2010年の電源比とガス技術を分ける}

\subsection{平成24年度 問5-3}
\label{q:h24-reg-basic-5-3}

\noindent\textbf{原典：} sources/H24/basic.pdf、PDF物理16頁。

\textbf{公式正答：} 2（single）。

\par\textbf{検討未確定：} 一般電気事業用という分母の2010年度原表と2005以降の米国生産統計は追加取得。現行の全発電量比率と混同しない。

\paragraph{問題の要旨} 2010年の電源比とガス技術を分ける。

\Needspace{4\baselineskip}\paragraph{確認した条件・記述（独自要約）}\mbox{}\par

\noindent\textbf{時点：} 平成24年度に指定された資料・制度。\par

\noindent\textbf{照合：} 原典21画像で否定・五肢・組合せ表を確認。後年値で置換しない。\par

\paragraph{問われる核心} 2010年の電源比とガス技術を分ける。

\paragraph{必要知識} 歴史的電源構成、複合サイクル、非在来型天然ガス。

\paragraph{考え方・正答を選ぶ理由（査読待ち初稿）} 2010年度の原子力比率は約30\%。天然ガスのコンバインドサイクルはガスタービンと排熱による蒸気タービンを組み合わせる。米国で増産が注目された非在来ガスはシェールガス。三語の組合せ2。

\Needspace{4\baselineskip}\paragraph{選択肢の検討（査読待ち初稿）}\mbox{}\par

\noindent\textbf{1：不適切} 30/再熱再生/タイトサンド。式・定義・対応の独立判定と一致しない。\par

\noindent\textbf{2：適切・選択対象} 30/コンバインド/シェール。判定と一致する。\par

\noindent\textbf{3：不適切} 30/コンバインド/タイトサンド。式・定義・対応の独立判定と一致しない。\par

\noindent\textbf{4：不適切} 20/再熱再生/シェール。式・定義・対応の独立判定と一致しない。\par

\noindent\textbf{5：不適切} 20/コンバインド/シェール。式・定義・対応の独立判定と一致しない。\par

\paragraph{再利用する考え方} 歴史的電源構成、複合サイクル、非在来型天然ガス。 指定年度・対象範囲・単位・否定を先に固定し、他年度の番号や現行値を転用しない。

\paragraph{共通理論} \hyperref[theory:env:energy2010]{第\ref*{theory:env:energy2010}節}

\par 難易度 3（校準初稿）、優先度未評価。

\subsection{令和2年度 問5-4}
\label{q:r02-reg-basic-5-4}

\noindent\textbf{原典：} sources/R02/basic.pdf、PDF物理21頁。

\textbf{公式正答：} 1（single）。

\par\textbf{検討未確定：} 2010の指定分母の原表と米国由来輸入の開始資料は未取得。H24と同じ番号でないことを確認済み。

\paragraph{問題の要旨} 2010原子力比と複合サイクル・シェール。

\Needspace{4\baselineskip}\paragraph{確認した条件・記述（独自要約）}\mbox{}\par

\noindent\textbf{時点：} 令和2年度に指定された資料・制度。\par

\noindent\textbf{照合：} 原典18画像で否定・五肢・組合せ表を確認。後年値で置換しない。\par

\paragraph{問われる核心} 2010原子力比と複合サイクル・シェール。

\paragraph{必要知識} 歴史的電源比、複合サイクル、ガスの供給。

\paragraph{考え方・正答を選ぶ理由（査読待ち初稿）} 2010一般電気事業用原子力約30\%、ガスタービンと排熱蒸気のコンバインド、米国非在来型シェールで1。原子力の指定時点と、輸出への期待から実際のLNG輸入へ変わった本文の時点を別々に読む。

\Needspace{4\baselineskip}\paragraph{選択肢の検討（査読待ち初稿）}\mbox{}\par

\noindent\textbf{1：適切・選択対象} 30/コンバインド/シェール。判定と一致する。\par

\noindent\textbf{2：不適切} 20/コンバインド/シェール。式・定義・対応の独立判定と一致しない。\par

\noindent\textbf{3：不適切} 20/再熱再生/シェール。式・定義・対応の独立判定と一致しない。\par

\noindent\textbf{4：不適切} 30/コンバインド/タイトサンド。式・定義・対応の独立判定と一致しない。\par

\noindent\textbf{5：不適切} 30/再熱再生/タイトサンド。式・定義・対応の独立判定と一致しない。\par

\paragraph{再利用する考え方} 歴史的電源比、複合サイクル、ガスの供給。 集計対象・時点・法の主体・選択方向を固定し、同型でも元の五肢へ戻す。

\paragraph{共通理論} \hyperref[theory:env:energy2010]{第\ref*{theory:env:energy2010}節}

\par 難易度 3（校準初稿）、優先度未評価。

\section{2011年一人当たり一次エネルギー}

\subsection{平成27年度 問5-4}
\label{q:h27-reg-basic-5-4}

\noindent\textbf{原典：} sources/H27/basic.pdf、PDF物理17頁。

\textbf{公式正答：} 5（single）。

\par\textbf{検討未確定：} IEAの指定2011年国別原表は未取得。現行データページを当時値の照合証跡にはせず、Astraの追加確認対象に残す。

\paragraph{問題の要旨} 2011年一人当たり一次エネルギー。

\Needspace{4\baselineskip}\paragraph{確認した条件・記述（独自要約）}\mbox{}\par

\noindent\textbf{時点：} 平成27年度に指定された資料・制度。\par

\noindent\textbf{照合：} 原典21画像で否定・五肢・組合せ表を確認。後年値で置換しない。\par

\paragraph{問われる核心} 2011年一人当たり一次エネルギー。

\paragraph{必要知識} 一次エネルギー、石油換算、人口当たり、年次比較。

\paragraph{考え方・正答を選ぶ理由（査読待ち初稿）} 原典の2011年指標では米加が7toe以上、韓露約5toe、独日約4toeで5。1toe≈42GJを使い、国全体の総量と人口当たりを混同しない。世界1.9、中国2.0という同じ年・同じ指標で比較する。

\Needspace{4\baselineskip}\paragraph{選択肢の検討（査読待ち初稿）}\mbox{}\par

\noindent\textbf{1：不適切} 独日/米加/韓露。式・定義・対応の独立判定と一致しない。\par

\noindent\textbf{2：不適切} 韓露/独日/米加。式・定義・対応の独立判定と一致しない。\par

\noindent\textbf{3：不適切} 韓露/米加/独日。式・定義・対応の独立判定と一致しない。\par

\noindent\textbf{4：不適切} 米加/独日/韓露。式・定義・対応の独立判定と一致しない。\par

\noindent\textbf{5：適切・選択対象} 米加/韓露/独日。判定と一致する。\par

\paragraph{再利用する考え方} 一次エネルギー、石油換算、人口当たり、年次比較。 指定年度・対象範囲・単位・否定を先に固定し、他年度の番号や現行値を転用しない。

\paragraph{共通理論} \hyperref[theory:env:energycapita]{第\ref*{theory:env:energycapita}節}

\par 難易度 3（校準初稿）、優先度未評価。

\subsection{令和3年度 問5-4}
\label{q:r03-reg-basic-5-4}

\noindent\textbf{原典：} sources/R03/basic.pdf、PDF物理24頁。

\textbf{公式正答：} 2（single）。

\par\textbf{検討未確定：} IEA2018国別値の元統計表は追加確認対象。既知の大国総量から値を逆算しない。

\paragraph{問題の要旨} 2018一人当たり一次エネルギーの国群。

\Needspace{4\baselineskip}\paragraph{確認した条件・記述（独自要約）}\mbox{}\par

\noindent\textbf{原典：} 全五肢と否定方向・組合せ表を画像確認。\par

\noindent\textbf{資料時点：} 令和3年度の指定資料・統計年を保持し、後年値に置換しない。\par

\paragraph{問われる核心} 2018一人当たり一次エネルギーの国群。

\paragraph{必要知識} 石油換算トン約42GJ、一人当たりと総量、同年比較。

\paragraph{考え方・正答を選ぶ理由（査読待ち初稿）} 指定されたIEA2018の一人当たり一次供給では、6石油換算トン以上が米国・カナダ、5〜6が韓国・ロシア、3〜4がドイツ・日本。世界1.9、中国2.3と同じ定義・年で比較する。国全体の総量順位や電力だけの順位ではなく、人口で割った一次供給であり2。

\Needspace{4\baselineskip}\paragraph{選択肢の検討（査読待ち初稿）}\mbox{}\par

\noindent\textbf{1：不適切} 米加/独日/韓露。独立に求めた組合せ 米加/韓露/独日 と異なる。\par

\noindent\textbf{2：適切・選択対象} 米加/韓露/独日。独立に求めた組合せ 米加/韓露/独日 と一致する。\par

\noindent\textbf{3：不適切} 独日/米加/韓露。独立に求めた組合せ 米加/韓露/独日 と異なる。\par

\noindent\textbf{4：不適切} 韓露/独日/米加。独立に求めた組合せ 米加/韓露/独日 と異なる。\par

\noindent\textbf{5：不適切} 韓露/米加/独日。独立に求めた組合せ 米加/韓露/独日 と異なる。\par

\paragraph{再利用する考え方} 石油換算トン約42GJ、一人当たりと総量、同年比較。 否定方向と資料年、分母と単位を確認して五肢へ戻る。

\paragraph{共通理論} \hyperref[theory:env:energycapita]{第\ref*{theory:env:energycapita}節}

\par 難易度 3（校準初稿）、優先度未評価。

\section{2015年策定の2030電源見通しを読む}

\subsection{平成28年度 問5-4}
\label{q:h28-reg-basic-5-4}

\noindent\textbf{原典：} sources/H28/basic.pdf、PDF物理23頁。

\textbf{公式正答：} 1（single）。

\par\textbf{検討未確定：} 2015公開関連資料の検索と白書骨子を参照。策定7月の最終PDFの全図表照合は未了で後工程に残す。

\paragraph{問題の要旨} 2015年策定の2030電源見通しを読む。

\Needspace{4\baselineskip}\paragraph{確認した条件・記述（独自要約）}\mbox{}\par

\noindent\textbf{時点：} 平成28年度に指定された資料・制度。\par

\noindent\textbf{照合：} 原典21画像で否定・五肢・組合せ表を確認。後年値で置換しない。\par

\paragraph{問われる核心} 2015年策定の2030電源見通しを読む。

\paragraph{必要知識} 策定年と目標年、電源構成と自給率、見通し。

\paragraph{考え方・正答を選ぶ理由（査読待ち初稿）} 2015の見通しの2030構成は、25–27\%がLNGで石油は約3\%。従って1が不適切。再エネ22–24\%、原子力20–22\%と、需給の省エネ・自給率目標は当時見通しの範囲。2030を対象とする後年改訂計画の比率をこの問題へ遡及しない。

\Needspace{4\baselineskip}\paragraph{選択肢の検討（査読待ち初稿）}\mbox{}\par

\noindent\textbf{1：不適切・選択対象} 石油25–27\%。判定と一致する。\par

\noindent\textbf{2：条件内で適切} 再エネ22–24\%。原典の条件内では適切な記述であり、不適切なものを選ぶ問の選択対象ではない。\par

\noindent\textbf{3：条件内で適切} 原子力20–22\%。原典の条件内では適切な記述であり、不適切なものを選ぶ問の選択対象ではない。\par

\noindent\textbf{4：条件内で適切} 自給率約25\%と準国産。原典の条件内では適切な記述であり、不適切なものを選ぶ問の選択対象ではない。\par

\noindent\textbf{5：条件内で適切} 省エネで最終消費低下。原典の条件内では適切な記述であり、不適切なものを選ぶ問の選択対象ではない。\par

\paragraph{再利用する考え方} 策定年と目標年、電源構成と自給率、見通し。 指定年度・対象範囲・単位・否定を先に固定し、他年度の番号や現行値を転用しない。

\paragraph{共通理論} \hyperref[theory:env:energymix2015]{第\ref*{theory:env:energymix2015}節}

\par 参考：\cite{env58_mix2015}。

\par 難易度 3（校準初稿）、優先度未評価。

\subsection{令和元年度 問5-3}
\label{q:r01-reg-basic-5-3}

\noindent\textbf{原典：} sources/R01/basic.pdf、PDF物理20頁。

\textbf{公式正答：} 3（single）。

\par\textbf{検討未確定：} 2015の7月最終PDF全図表をまだ取得していない。原典と公開骨子に基づく初稿として追加照合。

\paragraph{問題の要旨} 2015策定の2030比率で燃料名を点検。

\Needspace{4\baselineskip}\paragraph{確認した条件・記述（独自要約）}\mbox{}\par

\noindent\textbf{時点：} 令和元年度に指定された資料・制度。\par

\noindent\textbf{照合：} 原典18画像で否定・五肢・組合せ表を確認。後年値で置換しない。\par

\paragraph{問われる核心} 2015策定の2030比率で燃料名を点検。

\paragraph{必要知識} 策定年・目標年・実績と見通し・電源分類。

\paragraph{考え方・正答を選ぶ理由（査読待ち初稿）} 25–27\%はLNG火力の目標であり、石油火力は約3\%。燃料名の誤りの3を選ぶ。策定時点の他の比率・省エネ・自給率を保持し、後年の目標改訂値へ置換しない。

\Needspace{4\baselineskip}\paragraph{選択肢の検討（査読待ち初稿）}\mbox{}\par

\noindent\textbf{1：条件内で適切} 原子力20–22。条件内で適切な記述であり、不適切な記述を選ぶ対象ではない。\par

\noindent\textbf{2：条件内で適切} 再エネ22–24。条件内で適切な記述であり、不適切な記述を選ぶ対象ではない。\par

\noindent\textbf{3：不適切・選択対象} 石油25–27。判定と一致する。\par

\noindent\textbf{4：条件内で適切} 2013より最終消費減。条件内で適切な記述であり、不適切な記述を選ぶ対象ではない。\par

\noindent\textbf{5：条件内で適切} 自給率25程度と準国産。条件内で適切な記述であり、不適切な記述を選ぶ対象ではない。\par

\paragraph{再利用する考え方} 策定年・目標年・実績と見通し・電源分類。 集計対象・時点・法の主体・選択方向を固定し、同型でも元の五肢へ戻す。

\paragraph{共通理論} \hyperref[theory:env:energymix2015]{第\ref*{theory:env:energymix2015}節}

\par 参考：\cite{env58_mix2015}。

\par 難易度 3（校準初稿）、優先度未評価。

\section{専門職責任は職務規定だけでは尽くせない}

\subsection{平成26年度 問5-5}
\label{q:h26-reg-basic-5-5}

\noindent\textbf{原典：} sources/H26/basic.pdf、PDF物理17頁。

\textbf{公式正答：} 3（single）。

\paragraph{問題の要旨} 専門職責任は職務規定だけでは尽くせない。

\Needspace{4\baselineskip}\paragraph{確認した条件・記述（独自要約）}\mbox{}\par

\noindent\textbf{時点：} 平成26年度に指定された資料・制度。\par

\noindent\textbf{照合：} 原典21画像で否定・五肢・組合せ表を確認。後年値で置換しない。\par

\paragraph{問われる核心} 専門職責任は職務規定だけでは尽くせない。

\paragraph{必要知識} 力量、公益、説明責任、法と倫理の範囲。

\paragraph{考え方・正答を選ぶ理由（査読待ち初稿）} 規定に明示がないことを責任免除の根拠とする3が不適切。専門職は力量の範囲、安全・公益、説明と結果への責任を考える。依頼業務が力量を超える場合、能力ある者との協働や適切な引継ぎを含め、無能力のまま責任を引き受けることを避ける。

\Needspace{4\baselineskip}\paragraph{選択肢の検討（査読待ち初稿）}\mbox{}\par

\noindent\textbf{1：条件内で適切} 職業で規範が異なる場合。原典の条件内では適切な記述であり、不適切なものを選ぶ問の選択対象ではない。\par

\noindent\textbf{2：条件内で適切} 力量外を引受けない。原典の条件内では適切な記述であり、不適切なものを選ぶ問の選択対象ではない。\par

\noindent\textbf{3：不適切・選択対象} 明文規定なしなら責任なし。判定と一致する。\par

\noindent\textbf{4：条件内で適切} 行動規範は変化。原典の条件内では適切な記述であり、不適切なものを選ぶ問の選択対象ではない。\par

\noindent\textbf{5：条件内で適切} 専門知識は高い基準。原典の条件内では適切な記述であり、不適切なものを選ぶ問の選択対象ではない。\par

\paragraph{再利用する考え方} 力量、公益、説明責任、法と倫理の範囲。 指定年度・対象範囲・単位・否定を先に固定し、他年度の番号や現行値を転用しない。

\paragraph{共通理論} \hyperref[theory:env:ethics]{第\ref*{theory:env:ethics}節}

\par 参考：\cite{ipej_ethics2011}。

\par 難易度 2（校準初稿）、優先度未評価。

\subsection{平成30年度 問5-6}
\label{q:h30-reg-basic-5-6}

\noindent\textbf{原典：} sources/H30/basic.pdf、PDF物理20頁。

\textbf{公式正答：} 5（single）。

\paragraph{問題の要旨} 倫理責任は明文化の有無だけでは決まらない。

\Needspace{4\baselineskip}\paragraph{確認した条件・記述（独自要約）}\mbox{}\par

\noindent\textbf{時点：} 平成30年度に指定された資料・制度。\par

\noindent\textbf{照合：} 原典18画像で否定・五肢・組合せ表を確認。後年値で置換しない。\par

\paragraph{問われる核心} 倫理責任は明文化の有無だけでは決まらない。

\paragraph{必要知識} 公益、力量、責任、明文と倫理。

\paragraph{考え方・正答を選ぶ理由（査読待ち初稿）} 明文規定がなくても専門職の公益・安全・説明に関する責任は残るので5。力量の範囲を見極め、能力ある者との協働や引継ぎで適切に業務を処理する。

\Needspace{4\baselineskip}\paragraph{選択肢の検討（査読待ち初稿）}\mbox{}\par

\noindent\textbf{1：条件内で適切} 専門家は高い基準。条件内で適切な記述であり、不適切な記述を選ぶ対象ではない。\par

\noindent\textbf{2：条件内で適切} 職業で規範相違。条件内で適切な記述であり、不適切な記述を選ぶ対象ではない。\par

\noindent\textbf{3：条件内で適切} 力量超え仕事を引受けない。条件内で適切な記述であり、不適切な記述を選ぶ対象ではない。\par

\noindent\textbf{4：条件内で適切} 規範は変わる。条件内で適切な記述であり、不適切な記述を選ぶ対象ではない。\par

\noindent\textbf{5：不適切・選択対象} 規定なしなら責任なし。判定と一致する。\par

\paragraph{再利用する考え方} 公益、力量、責任、明文と倫理。 集計対象・時点・法の主体・選択方向を固定し、同型でも元の五肢へ戻す。

\paragraph{共通理論} \hyperref[theory:env:ethics]{第\ref*{theory:env:ethics}節}

\par 参考：\cite{ipej_ethics2011}。

\par 難易度 2（校準初稿）、優先度未評価。

\section{2012開始のFITの原則と拒否事由}

\subsection{平成26年度 問5-3}
\label{q:h26-reg-basic-5-3}

\noindent\textbf{原典：} sources/H26/basic.pdf、PDF物理17頁。

\textbf{公式正答：} 5（single）。

\par\textbf{検討未確定：} 2013METI資料の検索本文で正当拒否事由を確認したがPDF全文取得は失敗。出題時価格決定と対象水力条件の詳細、法令履歴は追加照合。

\paragraph{問題の要旨} 2012開始のFITの原則と拒否事由。

\Needspace{4\baselineskip}\paragraph{確認した条件・記述（独自要約）}\mbox{}\par

\noindent\textbf{時点：} 平成26年度に指定された資料・制度。\par

\noindent\textbf{照合：} 原典21画像で否定・五肢・組合せ表を確認。後年値で置換しない。\par

\paragraph{問われる核心} 2012開始のFITの原則と拒否事由。

\paragraph{必要知識} 固定価格、買取と接続、賦課金、例外。

\paragraph{考え方・正答を選ぶ理由（査読待ち初稿）} 正当な理由がある場合まで接続拒否を一切できないとする5は誤り。2012開始の対象は太陽光・風力・水力・地熱・バイオマスで、調達価格の設定、使用量に応じた賦課、費用調整という仕組みを区別する。例外を削って無条件義務にしない。

\Needspace{4\baselineskip}\paragraph{選択肢の検討（査読待ち初稿）}\mbox{}\par

\noindent\textbf{1：条件内で適切} 五種類の再エネ。原典の条件内では適切な記述であり、不適切なものを選ぶ問の選択対象ではない。\par

\noindent\textbf{2：条件内で適切} 大臣が年度価格設定。原典の条件内では適切な記述であり、不適切なものを選ぶ問の選択対象ではない。\par

\noindent\textbf{3：条件内で適切} 電力量に応じ賦課。原典の条件内では適切な記述であり、不適切なものを選ぶ問の選択対象ではない。\par

\noindent\textbf{4：条件内で適切} 地域負担差を調整。原典の条件内では適切な記述であり、不適切なものを選ぶ問の選択対象ではない。\par

\noindent\textbf{5：不適切・選択対象} 接続は一切拒否不可。判定と一致する。\par

\paragraph{再利用する考え方} 固定価格、買取と接続、賦課金、例外。 指定年度・対象範囲・単位・否定を先に固定し、他年度の番号や現行値を転用しない。

\paragraph{共通理論} \hyperref[theory:env:fit]{第\ref*{theory:env:fit}節}

\par 参考：\cite{env58_fit2013}。

\par 難易度 3（校準初稿）、優先度未評価。

\section{複合発電の高温側から低温側への順}

\subsection{平成28年度 問5-3}
\label{q:h28-reg-basic-5-3}

\noindent\textbf{原典：} sources/H28/basic.pdf、PDF物理22頁。

\textbf{公式正答：} 3（single）。

\paragraph{問題の要旨} 複合発電の高温側から低温側への順。

\Needspace{4\baselineskip}\paragraph{確認した条件・記述（独自要約）}\mbox{}\par

\noindent\textbf{時点：} 平成28年度に指定された資料・制度。\par

\noindent\textbf{照合：} 原典21画像で否定・五肢・組合せ表を確認。後年値で置換しない。\par

\paragraph{問われる核心} 複合発電の高温側から低温側への順。

\paragraph{必要知識} 温度段階、電気化学、静電蓄積、エネルギー変換。

\paragraph{考え方・正答を選ぶ理由（査読待ち初稿）} コンバインドは燃焼ガスでガスタービンを回し、その高温排熱から蒸気を作って蒸気タービンへ進む。逆順の3が不適切。二次電池は再充電、EDLCは界面で蓄える。揚水は発電を消費側から切り離して電気を水の位置エネルギーへ変換する貯蔵で、発電時だけの効率を往復効率としない。

\Needspace{4\baselineskip}\paragraph{選択肢の検討（査読待ち初稿）}\mbox{}\par

\noindent\textbf{1：条件内で適切} 一次使切・二次充電。原典の条件内では適切な記述であり、不適切なものを選ぶ問の選択対象ではない。\par

\noindent\textbf{2：条件内で適切} EDLC急速・長サイクル。原典の条件内では適切な記述であり、不適切なものを選ぶ問の選択対象ではない。\par

\noindent\textbf{3：不適切・選択対象} 先に蒸気タービン次に廃熱ガスタービン。判定と一致する。\par

\noindent\textbf{4：条件内で適切} 燃料電池は水電解の逆。原典の条件内では適切な記述であり、不適切なものを選ぶ問の選択対象ではない。\par

\noindent\textbf{5：条件内で適切} 揚水は位置エネルギーに貯蔵。原典の条件内では適切な記述であり、不適切なものを選ぶ問の選択対象ではない。\par

\paragraph{再利用する考え方} 温度段階、電気化学、静電蓄積、エネルギー変換。 指定年度・対象範囲・単位・否定を先に固定し、他年度の番号や現行値を転用しない。

\paragraph{共通理論} \hyperref[theory:env:generation]{第\ref*{theory:env:generation}節}

\par 難易度 2（校準初稿）、優先度未評価。

\section{発熱量は指定表と単位質量で比較}

\subsection{平成24年度 問5-2}
\label{q:h24-reg-basic-5-2}

\noindent\textbf{原典：} sources/H24/basic.pdf、PDF物理15頁。

\textbf{公式正答：} 4（single）。

\par\textbf{検討未確定：} 2005年度指定表の全四値を同じ定義で取得する作業は未完了。2023改訂資料の歴史表を補助参照し、後年値へ置換しない。

\paragraph{問題の要旨} 発熱量は指定表と単位質量で比較。

\Needspace{4\baselineskip}\paragraph{確認した条件・記述（独自要約）}\mbox{}\par

\noindent\textbf{時点：} 平成24年度に指定された資料・制度。\par

\noindent\textbf{照合：} 原典21画像で否定・五肢・組合せ表を確認。後年値で置換しない。\par

\paragraph{問われる核心} 発熱量は指定表と単位質量で比較。

\paragraph{必要知識} 高位・低位発熱量、単位、燃料品種、密度換算。

\paragraph{考え方・正答を選ぶ理由（査読待ち初稿）} 指定された燃料種の質量当たり発熱量は天然ガスCが最大、原油A、輸入原料炭B、廃材Dの順なので4。天然ガスの表が体積単位なら密度でMJ/kgへ換算する。液体のMJ/Lを質量当たり値へ直接並べない。

\Needspace{4\baselineskip}\paragraph{選択肢の検討（査読待ち初稿）}\mbox{}\par

\noindent\textbf{1：不適切} A>B>C>D。式・定義・対応の独立判定と一致しない。\par

\noindent\textbf{2：不適切} A>B>D>C。式・定義・対応の独立判定と一致しない。\par

\noindent\textbf{3：不適切} B>C>D>A。式・定義・対応の独立判定と一致しない。\par

\noindent\textbf{4：適切・選択対象} C>A>B>D。判定と一致する。\par

\noindent\textbf{5：不適切} C>A>D>B。式・定義・対応の独立判定と一致しない。\par

\paragraph{再利用する考え方} 高位・低位発熱量、単位、燃料品種、密度換算。 指定年度・対象範囲・単位・否定を先に固定し、他年度の番号や現行値を転用しない。

\paragraph{共通理論} \hyperref[theory:env:heating]{第\ref*{theory:env:heating}節}

\par 参考：\cite{env58_heating}。

\par 難易度 3（校準初稿）、優先度未評価。

\subsection{令和元年度再試験 問5-3}
\label{q:r01-ret-basic-5-3}

\noindent\textbf{原典：} sources/R01\_retest/basic.pdf、PDF物理18頁。

\textbf{公式正答：} 3（single）。

\par\textbf{検討未確定：} 指定時版の標準発熱量の全四値は未取得。後年改訂表を指定時表の代替確定にはしない。

\paragraph{問題の要旨} LNG・原油・一般炭・絶乾材の質量発熱量。

\Needspace{4\baselineskip}\paragraph{確認した条件・記述（独自要約）}\mbox{}\par

\noindent\textbf{時点：} 令和元年度再試験に指定された資料・制度。\par

\noindent\textbf{照合：} 原典18画像で否定・五肢・組合せ表を確認。後年値で置換しない。\par

\paragraph{問われる核心} LNG・原油・一般炭・絶乾材の質量発熱量。

\paragraph{必要知識} 燃料指定、乾燥状態、熱量単位と密度換算。

\paragraph{考え方・正答を選ぶ理由（査読待ち初稿）} 指定表の単位質量ではLNG、原油、輸入一般炭、絶乾廃材の順で3。平成24年度は気化ガスと原料炭という異なる指定なので同じ数値を転用しない。ここでも体積単位は密度換算が必要。

\Needspace{4\baselineskip}\paragraph{選択肢の検討（査読待ち初稿）}\mbox{}\par

\noindent\textbf{1：不適切} A>B>C>D。式・定義・対応の独立判定と一致しない。\par

\noindent\textbf{2：不適切} B>A>D>C。式・定義・対応の独立判定と一致しない。\par

\noindent\textbf{3：適切・選択対象} C>A>B>D。判定と一致する。\par

\noindent\textbf{4：不適切} C>B>D>A。式・定義・対応の独立判定と一致しない。\par

\noindent\textbf{5：不適切} D>C>B>A。式・定義・対応の独立判定と一致しない。\par

\paragraph{再利用する考え方} 燃料指定、乾燥状態、熱量単位と密度換算。 集計対象・時点・法の主体・選択方向を固定し、同型でも元の五肢へ戻す。

\paragraph{共通理論} \hyperref[theory:env:heating]{第\ref*{theory:env:heating}節}

\par 参考：\cite{env58_heating}。

\par 難易度 3（校準初稿）、優先度未評価。

\section{発表・発明を古い順へ並べる}

\subsection{平成25年度 問5-6}
\label{q:h25-reg-basic-5-6}

\noindent\textbf{原典：} sources/H25/basic.pdf、PDF物理22頁。

\textbf{公式正答：} 1（single）。

\par\textbf{検討未確定：} 全五業績の初出一次史料は未取得。特にハレー彗星の「発見」は周期性の同定を指す出題表現として留保。

\paragraph{問題の要旨} 発表・発明を古い順へ並べる。

\Needspace{4\baselineskip}\paragraph{確認した条件・記述（独自要約）}\mbox{}\par

\noindent\textbf{時点：} 平成25年度に指定された資料・制度。\par

\noindent\textbf{照合：} 原典21画像で否定・五肢・組合せ表を確認。後年値で置換しない。\par

\paragraph{問われる核心} 発表・発明を古い順へ並べる。

\paragraph{必要知識} 科学史の世紀区分、業績の段階、順序制約。

\paragraph{考え方・正答を選ぶ理由（査読待ち初稿）} ハレーの周期彗星の予測・同定が18世紀初頭、Wattの蒸気機関改良が18世紀後半、Darwin–Wallaceの自然選択発表が1858、一般相対論が1915、Fukuiのフロンティア理論が1952。エ→ア→イ→オ→ウで1。発見と理論発表、特許と実用化のどの段階をいうかも確認する。

\Needspace{4\baselineskip}\paragraph{選択肢の検討（査読待ち初稿）}\mbox{}\par

\noindent\textbf{1：適切・選択対象} エア イオウ。判定と一致する。\par

\noindent\textbf{2：不適切} アイエウオ。式・定義・対応の独立判定と一致しない。\par

\noindent\textbf{3：不適切} エアイウオ。式・定義・対応の独立判定と一致しない。\par

\noindent\textbf{4：不適切} アエオイウ。式・定義・対応の独立判定と一致しない。\par

\noindent\textbf{5：不適切} アイエオウ。式・定義・対応の独立判定と一致しない。\par

\paragraph{再利用する考え方} 科学史の世紀区分、業績の段階、順序制約。 指定年度・対象範囲・単位・否定を先に固定し、他年度の番号や現行値を転用しない。

\paragraph{共通理論} \hyperref[theory:env:historyorder]{第\ref*{theory:env:historyorder}節}

\par 参考：\cite{env58_steam}。

\par 難易度 3（校準初稿）、優先度未評価。

\subsection{平成26年度 問5-6}
\label{q:h26-reg-basic-5-6}

\noindent\textbf{原典：} sources/H26/basic.pdf、PDF物理18頁。

\textbf{公式正答：} 4（single）。

\par\textbf{検討未確定：} Galileo・Wilsonの初出史料は追加照合。ACSの1935生産選定と1939商業化を分ける。

\paragraph{問題の要旨} 観測・動力・検出・高分子の技術史。

\Needspace{4\baselineskip}\paragraph{確認した条件・記述（独自要約）}\mbox{}\par

\noindent\textbf{時点：} 平成26年度に指定された資料・制度。\par

\noindent\textbf{照合：} 原典21画像で否定・五肢・組合せ表を確認。後年値で置換しない。\par

\paragraph{問われる核心} 観測・動力・検出・高分子の技術史。

\paragraph{必要知識} 改良と新機構、検出技術、合成と商業化。

\paragraph{考え方・正答を選ぶ理由（査読待ち初稿）} Galileoの望遠鏡観測が1609年頃、Newcomen機関1712、Wattの改良特許1769、Wilson霧箱は20世紀初頭、Carothersのナイロンは1930年代。ウ→ア→イ→オ→エの4。ナイロンの合成・生産選定と1939の商業化を同一年にせず、霧箱より後という比較に必要な範囲で扱う。

\Needspace{4\baselineskip}\paragraph{選択肢の検討（査読待ち初稿）}\mbox{}\par

\noindent\textbf{1：不適切} イアウオエ。式・定義・対応の独立判定と一致しない。\par

\noindent\textbf{2：不適切} イウアオエ。式・定義・対応の独立判定と一致しない。\par

\noindent\textbf{3：不適切} イアウエオ。式・定義・対応の独立判定と一致しない。\par

\noindent\textbf{4：適切・選択対象} ウアイオエ。判定と一致する。\par

\noindent\textbf{5：不適切} ウイアエオ。式・定義・対応の独立判定と一致しない。\par

\paragraph{再利用する考え方} 改良と新機構、検出技術、合成と商業化。 指定年度・対象範囲・単位・否定を先に固定し、他年度の番号や現行値を転用しない。

\paragraph{共通理論} \hyperref[theory:env:historyorder]{第\ref*{theory:env:historyorder}節}

\par 参考：\cite{env58_steam,env58_nylon}。

\par 難易度 3（校準初稿）、優先度未評価。

\subsection{平成28年度 問5-6}
\label{q:h28-reg-basic-5-6}

\noindent\textbf{原典：} sources/H28/basic.pdf、PDF物理24頁。

\textbf{公式正答：} 4（single）。

\par\textbf{検討未確定：} 雷の実験・三極管・トランジスタ各原報と特許の一次史料は追加照合。三人の貢献を同一作業に単純化しない。

\paragraph{問題の要旨} 電気・化学・電子の代表的業績の順序。

\Needspace{4\baselineskip}\paragraph{確認した条件・記述（独自要約）}\mbox{}\par

\noindent\textbf{時点：} 平成28年度に指定された資料・制度。\par

\noindent\textbf{照合：} 原典21画像で否定・五肢・組合せ表を確認。後年値で置換しない。\par

\paragraph{問われる核心} 電気・化学・電子の代表的業績の順序。

\paragraph{必要知識} 周期律、放射性元素、電子管と半導体、年代順。

\paragraph{考え方・正答を選ぶ理由（査読待ち初稿）} 雷の電気的性質のFranklinが18世紀、Mendeleevの周期律1869、Curie夫妻の新元素1898、De Forestの三極管1906頃、トランジスタ1947。ウ→ア→イ→オ→エの4。電子管から半導体素子への世紀をまたぐ発展と、発見・発明・商用普及の段階を分ける。

\Needspace{4\baselineskip}\paragraph{選択肢の検討（査読待ち初稿）}\mbox{}\par

\noindent\textbf{1：不適切} アイウエオ。式・定義・対応の独立判定と一致しない。\par

\noindent\textbf{2：不適切} アイウオエ。式・定義・対応の独立判定と一致しない。\par

\noindent\textbf{3：不適切} ウイアオエ。式・定義・対応の独立判定と一致しない。\par

\noindent\textbf{4：適切・選択対象} ウアイオエ。判定と一致する。\par

\noindent\textbf{5：不適切} ウイアエオ。式・定義・対応の独立判定と一致しない。\par

\paragraph{再利用する考え方} 周期律、放射性元素、電子管と半導体、年代順。 指定年度・対象範囲・単位・否定を先に固定し、他年度の番号や現行値を転用しない。

\paragraph{共通理論} \hyperref[theory:env:historyorder]{第\ref*{theory:env:historyorder}節}

\par 参考：\cite{env58_curie}。

\par 難易度 3（校準初稿）、優先度未評価。

\subsection{平成30年度 問5-5}
\label{q:h30-reg-basic-5-5}

\noindent\textbf{原典：} sources/H30/basic.pdf、PDF物理20頁。

\textbf{公式正答：} 3（single）。

\par\textbf{検討未確定：} 五項目の初出全原報は未取得。HaberとHahnの単独人物表記は共同の貢献の全体を示すとはしない。

\paragraph{問題の要旨} 動力・通信・化学・核の技術成果の順。

\Needspace{4\baselineskip}\paragraph{確認した条件・記述（独自要約）}\mbox{}\par

\noindent\textbf{時点：} 平成30年度に指定された資料・制度。\par

\noindent\textbf{照合：} 原典18画像で否定・五肢・組合せ表を確認。後年値で置換しない。\par

\paragraph{問われる核心} 動力・通信・化学・核の技術成果の順。

\paragraph{必要知識} 科学的確認、技術発明、工業化の段階。

\paragraph{考え方・正答を選ぶ理由（査読待ち初稿）} Wattの改良が18世紀後半、Bellの電話1876、Hertzの電磁波実験1880年代、Haberのアンモニア合成基礎1909頃、Hahnらの核分裂1938。オ→ウ→エ→ア→イの3。Haberの基礎とBoschによる大規模工業化、核分裂発見と発電開始を分ける。

\Needspace{4\baselineskip}\paragraph{選択肢の検討（査読待ち初稿）}\mbox{}\par

\noindent\textbf{1：不適切} ウエオイア。式・定義・対応の独立判定と一致しない。\par

\noindent\textbf{2：不適切} ウオアエイ。式・定義・対応の独立判定と一致しない。\par

\noindent\textbf{3：適切・選択対象} オウエアイ。判定と一致する。\par

\noindent\textbf{4：不適切} オエウイア。式・定義・対応の独立判定と一致しない。\par

\noindent\textbf{5：不適切} アオウエイ。式・定義・対応の独立判定と一致しない。\par

\paragraph{再利用する考え方} 科学的確認、技術発明、工業化の段階。 集計対象・時点・法の主体・選択方向を固定し、同型でも元の五肢へ戻す。

\paragraph{共通理論} \hyperref[theory:env:historyorder]{第\ref*{theory:env:historyorder}節}

\par 参考：\cite{env58_steam}。

\paragraph{関連問題} \hyperref[q:r03-reg-basic-5-5]{R03-REG-BASIC-5-5}

\par 難易度 3（校準初稿）、優先度未評価。

\subsection{令和元年度再試験 問5-5}
\label{q:r01-ret-basic-5-5}

\noindent\textbf{原典：} sources/R01\_retest/basic.pdf、PDF物理19頁。

\textbf{公式正答：} 4（single）。

\par\textbf{検討未確定：} 初出史料の全五項目照合は未完了。原典の彗星「発見」は周期性の同定として留保。

\paragraph{問題の要旨} 同じ年表でも元の選択順へ戻す。

\Needspace{4\baselineskip}\paragraph{確認した条件・記述（独自要約）}\mbox{}\par

\noindent\textbf{時点：} 令和元年度再試験に指定された資料・制度。\par

\noindent\textbf{照合：} 原典18画像で否定・五肢・組合せ表を確認。後年値で置換しない。\par

\paragraph{問われる核心} 同じ年表でも元の選択順へ戻す。

\paragraph{必要知識} 初出の段階、時代順、候補の並べ替え。

\paragraph{考え方・正答を選ぶ理由（査読待ち初稿）} Halleyの周期性同定が18世紀初頭、Watt18世紀後半、Darwin–Wallace1858、Einstein1915、Fukui1952。エ→ア→イ→オ→ウの4。平成25年度の同型問題の正答1を番号だけ転用しない。

\Needspace{4\baselineskip}\paragraph{選択肢の検討（査読待ち初稿）}\mbox{}\par

\noindent\textbf{1：不適切} アイエウオ。式・定義・対応の独立判定と一致しない。\par

\noindent\textbf{2：不適切} エアイウオ。式・定義・対応の独立判定と一致しない。\par

\noindent\textbf{3：不適切} アエオイウ。式・定義・対応の独立判定と一致しない。\par

\noindent\textbf{4：適切・選択対象} エアイオウ。判定と一致する。\par

\noindent\textbf{5：不適切} アイエオウ。式・定義・対応の独立判定と一致しない。\par

\paragraph{再利用する考え方} 初出の段階、時代順、候補の並べ替え。 集計対象・時点・法の主体・選択方向を固定し、同型でも元の五肢へ戻す。

\paragraph{共通理論} \hyperref[theory:env:historyorder]{第\ref*{theory:env:historyorder}節}

\par 参考：\cite{env58_steam}。

\par 難易度 3（校準初稿）、優先度未評価。

\subsection{令和2年度 問5-6}
\label{q:r02-reg-basic-5-6}

\noindent\textbf{原典：} sources/R02/basic.pdf、PDF物理22頁。

\textbf{公式正答：} 1（single）。

\par\textbf{検討未確定：} 全初出史料の取得は未了。電子管・トランジスタの共同貢献と特許の段階は追加照合。

\paragraph{問題の要旨} 種痘・周期律・新元素・電子素子の順。

\Needspace{4\baselineskip}\paragraph{確認した条件・記述（独自要約）}\mbox{}\par

\noindent\textbf{時点：} 令和2年度に指定された資料・制度。\par

\noindent\textbf{照合：} 原典18画像で否定・五肢・組合せ表を確認。後年値で置換しない。\par

\paragraph{問われる核心} 種痘・周期律・新元素・電子素子の順。

\paragraph{必要知識} 医療経験、元素周期、電子管から半導体。

\paragraph{考え方・正答を選ぶ理由（査読待ち初稿）} Jenner1796、Mendeleev1869、Curie夫妻1898、De Forest三極管1906頃、トランジスタ1947。イ→エ→ア→オ→ウの1。観測された事実、法則化、装置の発明を同じ授賞年へまとめない。

\Needspace{4\baselineskip}\paragraph{選択肢の検討（査読待ち初稿）}\mbox{}\par

\noindent\textbf{1：適切・選択対象} イエアオウ。判定と一致する。\par

\noindent\textbf{2：不適切} イエオウア。式・定義・対応の独立判定と一致しない。\par

\noindent\textbf{3：不適切} イオエアウ。式・定義・対応の独立判定と一致しない。\par

\noindent\textbf{4：不適切} エイオアウ。式・定義・対応の独立判定と一致しない。\par

\noindent\textbf{5：不適切} エオイアウ。式・定義・対応の独立判定と一致しない。\par

\paragraph{再利用する考え方} 医療経験、元素周期、電子管から半導体。 集計対象・時点・法の主体・選択方向を固定し、同型でも元の五肢へ戻す。

\paragraph{共通理論} \hyperref[theory:env:historyorder]{第\ref*{theory:env:historyorder}節}

\par 参考：\cite{env59_jenner,env58_curie}。

\par 難易度 3（校準初稿）、優先度未評価。

\subsection{令和3年度 問5-5}
\label{q:r03-reg-basic-5-5}

\noindent\textbf{原典：} sources/R03/basic.pdf、PDF物理25頁。

\textbf{公式正答：} 4（single）。

\par\textbf{検討未確定：} 業績の発表・特許・工業化の区分と各史料の年の追加照合が必要。

\paragraph{問題の要旨} 蒸気機関から核分裂までの順序。

\Needspace{4\baselineskip}\paragraph{確認した条件・記述（独自要約）}\mbox{}\par

\noindent\textbf{原典：} 全五肢と否定方向・組合せ表を画像確認。\par

\noindent\textbf{資料時点：} 令和3年度の指定資料・統計年を保持し、後年値に置換しない。\par

\paragraph{問われる核心} 蒸気機関から核分裂までの順序。

\paragraph{必要知識} 人物・対象業績・年代の三点照合。

\paragraph{考え方・正答を選ぶ理由（査読待ち初稿）} オWatt蒸気機関改良は18世紀後半、ウBell電話1876、エHertz電磁波確認1880年代後半、アHaberアンモニア合成20世紀初め、イHahn核分裂1938の順。オウエアイの4。発明・実験確認・工業化・受賞の年を一つの年として扱わない。

\Needspace{4\baselineskip}\paragraph{選択肢の検討（査読待ち初稿）}\mbox{}\par

\noindent\textbf{1：不適切} アオウエイ。独立に求めた組合せ オウエアイ と異なる。\par

\noindent\textbf{2：不適切} ウエオイア。独立に求めた組合せ オウエアイ と異なる。\par

\noindent\textbf{3：不適切} ウオアエイ。独立に求めた組合せ オウエアイ と異なる。\par

\noindent\textbf{4：適切・選択対象} オウエアイ。独立に求めた組合せ オウエアイ と一致する。\par

\noindent\textbf{5：不適切} オエウイア。独立に求めた組合せ オウエアイ と異なる。\par

\paragraph{再利用する考え方} 人物・対象業績・年代の三点照合。 否定方向と資料年、分母と単位を確認して五肢へ戻る。

\paragraph{共通理論} \hyperref[theory:env:historyorder]{第\ref*{theory:env:historyorder}節}

\paragraph{関連問題} \hyperref[q:h30-reg-basic-5-5]{H30-REG-BASIC-5-5}

\par 難易度 3（校準初稿）、優先度未評価。

\subsection{令和4年度 問5-6}
\label{q:r04-reg-basic-5-6}

\noindent\textbf{原典：} sources/R04/basic.pdf、PDF物理20頁。

\textbf{公式正答：} 3（single）。

\par\textbf{検討未確定：} 各業績の公刊・特許・開発年の史料は追加照合。ナイロンの開発と販売開始を同一視しない。

\paragraph{問題の要旨} 尿素合成・転炉・赤痢菌・KS鋼・ナイロン。

\Needspace{4\baselineskip}\paragraph{確認した条件・記述（独自要約）}\mbox{}\par

\noindent\textbf{原典：} 全五肢と否定方向・組合せ表を画像確認。\par

\noindent\textbf{資料時点：} 令和4年度の指定資料・統計年を保持し、後年値に置換しない。\par

\paragraph{問われる核心} 尿素合成・転炉・赤痢菌・KS鋼・ナイロン。

\paragraph{必要知識} 人物と対象物、発明・発見・商品化の年代。

\paragraph{考え方・正答を選ぶ理由（査読待ち初稿）} エWöhler尿素合成1828、アBessemer転炉1856、オ志賀赤痢菌1897、イ本多KS鋼1917、ウCarothersナイロン1935の順で3。材料の発明年と量産・販売年は分ける。化学合成から製鋼、細菌、磁性材料、高分子へ時代が進む。

\Needspace{4\baselineskip}\paragraph{選択肢の検討（査読待ち初稿）}\mbox{}\par

\noindent\textbf{1：不適切} アエイオウ。独立に求めた組合せ エアオイウ と異なる。\par

\noindent\textbf{2：不適切} アエオイウ。独立に求めた組合せ エアオイウ と異なる。\par

\noindent\textbf{3：適切・選択対象} エアオイウ。独立に求めた組合せ エアオイウ と一致する。\par

\noindent\textbf{4：不適切} エオアウイ。独立に求めた組合せ エアオイウ と異なる。\par

\noindent\textbf{5：不適切} オエアウイ。独立に求めた組合せ エアオイウ と異なる。\par

\paragraph{再利用する考え方} 人物と対象物、発明・発見・商品化の年代。 否定方向と資料年、分母と単位を確認して五肢へ戻る。

\paragraph{共通理論} \hyperref[theory:env:historyorder]{第\ref*{theory:env:historyorder}節}

\par 参考：\cite{env58_nylon}。

\par 難易度 3（校準初稿）、優先度未評価。

\subsection{令和6年度 問5-6}
\label{q:r06-reg-basic-5-6}

\noindent\textbf{原典：} sources/R06/basic.pdf、PDF物理20頁。

\textbf{公式正答：} 2（single）。

\par\textbf{検討未確定：} 各業績の指定史料と発表・出版年の照合を追加。

\paragraph{問題の要旨} BernoulliからWegenerまでの順序。

\Needspace{4\baselineskip}\paragraph{確認した条件・記述（独自要約）}\mbox{}\par

\noindent\textbf{原典：} 全五肢と否定方向・組合せ表を画像確認。\par

\noindent\textbf{資料時点：} 令和6年度の指定資料・統計年を保持し、後年値に置換しない。\par

\paragraph{問われる核心} BernoulliからWegenerまでの順序。

\paragraph{必要知識} 業績の年代区分と科学・技術の並行発展。

\paragraph{考え方・正答を選ぶ理由（査読待ち初稿）} アBernoulli1738の流体力学、イLavoisier18世紀後半の酸素燃焼理論、オRocket1829、ウDarwin種の起源1859、エWegener大陸移動説1912の順で2。実験、出版、車両製作、学説発表の基準を分ける。

\Needspace{4\baselineskip}\paragraph{選択肢の検討（査読待ち初稿）}\mbox{}\par

\noindent\textbf{1：不適切} アイウエオ。独立に求めた組合せ アイオウエ と異なる。\par

\noindent\textbf{2：適切・選択対象} アイオウエ。独立に求めた組合せ アイオウエ と一致する。\par

\noindent\textbf{3：不適切} イアウオエ。独立に求めた組合せ アイオウエ と異なる。\par

\noindent\textbf{4：不適切} エアイオウ。独立に求めた組合せ アイオウエ と異なる。\par

\noindent\textbf{5：不適切} エイアウオ。独立に求めた組合せ アイオウエ と異なる。\par

\paragraph{再利用する考え方} 業績の年代区分と科学・技術の並行発展。 否定方向と資料年、分母と単位を確認して五肢へ戻る。

\paragraph{共通理論} \hyperref[theory:env:historyorder]{第\ref*{theory:env:historyorder}節}

\par 難易度 3（校準初稿）、優先度未評価。

\section{人物と業績：放射能とラジウムを区別}

\subsection{平成23年度 問5-5}
\label{q:h23-reg-basic-5-5}

\noindent\textbf{原典：} sources/H23/basic.pdf、PDF物理17頁。

\textbf{公式正答：} 4（single）。

\par\textbf{検討未確定：} 原典は「アントワーヌ・ベクレル」と記載し人名が曖昧。放射能発見者Henri Becquerelとの対応を区別。Eastmanの「フィルム乾板」も乾板とロールフィルムをまとめた表現であり、原典を厳密な発明年表に一般化しない。

\paragraph{問題の要旨} 人物と業績：放射能とラジウムを区別。

\Needspace{4\baselineskip}\paragraph{確認した条件・記述（独自要約）}\mbox{}\par

\noindent\textbf{時点：} 平成23年度に指定された資料・制度。\par

\noindent\textbf{照合：} 原典21画像で否定・五肢・組合せ表を確認。後年値で置換しない。\par

\paragraph{問われる核心} 人物と業績：放射能とラジウムを区別。

\paragraph{必要知識} 放射能現象と元素発見、発明と改良・普及の区別。

\paragraph{考え方・正答を選ぶ理由（査読待ち初稿）} ラジウムとポロニウムの発見は1898年のキュリー夫妻の仕事。Henri Becquerelは1896年ウランの自然放射を見いだしたので4が誤り。Galileoは望遠鏡を改良して天体観測へ用いた。Huygensの振子時計、Arkwrightの水力紡績、Eastmanの写真材料・乾板の工業化という他の対応は、出題の代表的業績の範囲で扱う。

\Needspace{4\baselineskip}\paragraph{選択肢の検討（査読待ち初稿）}\mbox{}\par

\noindent\textbf{1：条件内で適切} Galileo/天体望遠鏡で観測。望遠鏡の発明者そのものとする記述ではなく観測への利用。\par

\noindent\textbf{2：条件内で適切} Huygens/振子時計。振子を時間標準にする時計。\par

\noindent\textbf{3：条件内で適切} Arkwright/水力紡績機。産業革命の水力紡績機。\par

\noindent\textbf{4：不適切・選択対象} Becquerel/ラジウム発見。放射能の発見とラジウム発見を取り違える。\par

\noindent\textbf{5：条件内で適切} Eastman/写真用フィルム乾板。乾板と柔軟な写真フィルムの開発・商業化という原典の代表表現。\par

\paragraph{再利用する考え方} 放射能現象と元素発見、発明と改良・普及の区別。 指定年度・対象範囲・単位・否定を先に固定し、他年度の番号や現行値を転用しない。

\paragraph{共通理論} \hyperref[theory:env:historypairs]{第\ref*{theory:env:historypairs}節}

\par 参考：\cite{env58_curie}。

\par 難易度 3（校準初稿）、優先度未評価。

\subsection{平成27年度 問5-6}
\label{q:h27-reg-basic-5-6}

\noindent\textbf{原典：} sources/H27/basic.pdf、PDF物理18頁。

\textbf{公式正答：} 4（single）。

\par\textbf{検討未確定：} 原典は「アントワーヌ・ベクレル」と記載し人名が曖昧。放射能発見者Henri Becquerelとの対応を区別。Eastmanの「フィルム乾板」も乾板とロールフィルムをまとめた表現であり、原典を厳密な発明年表に一般化しない。

\paragraph{問題の要旨} 人物と業績：放射能とラジウムを区別。

\Needspace{4\baselineskip}\paragraph{確認した条件・記述（独自要約）}\mbox{}\par

\noindent\textbf{時点：} 平成27年度に指定された資料・制度。\par

\noindent\textbf{照合：} 原典21画像で否定・五肢・組合せ表を確認。後年値で置換しない。\par

\paragraph{問われる核心} 人物と業績：放射能とラジウムを区別。

\paragraph{必要知識} 放射能現象と元素発見、発明と改良・普及の区別。

\paragraph{考え方・正答を選ぶ理由（査読待ち初稿）} ラジウムとポロニウムの発見は1898年のキュリー夫妻の仕事。Henri Becquerelは1896年ウランの自然放射を見いだしたので4が誤り。Galileoは望遠鏡を改良して天体観測へ用いた。Huygensの振子時計、Arkwrightの水力紡績、Eastmanの写真材料・乾板の工業化という他の対応は、出題の代表的業績の範囲で扱う。

\Needspace{4\baselineskip}\paragraph{選択肢の検討（査読待ち初稿）}\mbox{}\par

\noindent\textbf{1：条件内で適切} Galileo/天体望遠鏡で観測。望遠鏡の発明者そのものとする記述ではなく観測への利用。\par

\noindent\textbf{2：条件内で適切} Huygens/振子時計。振子を時間標準にする時計。\par

\noindent\textbf{3：条件内で適切} Arkwright/水力紡績機。産業革命の水力紡績機。\par

\noindent\textbf{4：不適切・選択対象} Becquerel/ラジウム発見。放射能の発見とラジウム発見を取り違える。\par

\noindent\textbf{5：条件内で適切} Eastman/写真用フィルム乾板。乾板と柔軟な写真フィルムの開発・商業化という原典の代表表現。\par

\paragraph{再利用する考え方} 放射能現象と元素発見、発明と改良・普及の区別。 指定年度・対象範囲・単位・否定を先に固定し、他年度の番号や現行値を転用しない。

\paragraph{共通理論} \hyperref[theory:env:historypairs]{第\ref*{theory:env:historypairs}節}

\par 参考：\cite{env58_curie}。

\par 難易度 3（校準初稿）、優先度未評価。

\section{産業革命の改良・工場化と担い手}

\subsection{平成24年度 問5-5}
\label{q:h24-reg-basic-5-5}

\noindent\textbf{原典：} sources/H24/basic.pdf、PDF物理17頁。

\textbf{公式正答：} 4（single）。

\par\textbf{検討未確定：} からくりからの転用という肢3の個別史料は追加照合。一般論を特定発明の確定年表に拡大しない。

\paragraph{問題の要旨} 産業革命の改良・工場化と担い手。

\Needspace{4\baselineskip}\paragraph{確認した条件・記述（独自要約）}\mbox{}\par

\noindent\textbf{時点：} 平成24年度に指定された資料・制度。\par

\noindent\textbf{照合：} 原典21画像で否定・五肢・組合せ表を確認。後年値で置換しない。\par

\paragraph{問われる核心} 産業革命の改良・工場化と担い手。

\paragraph{必要知識} 発明、改良、普及、研究組織と生産制度。

\paragraph{考え方・正答を選ぶ理由（査読待ち初稿）} 産業革命期の代表的機械は職人・発明家・事業者の改良と事業化で広まり、当時の欧州大学研究者が開発を主導した例が多いという4は不適切。新規性だけでなく既存機構の組合せ、工場制度、生産と社会への影響も見る。Wattと大学との接点のような個別事例は、全体を大学主導とする証明ではない。

\Needspace{4\baselineskip}\paragraph{選択肢の検討（査読待ち初稿）}\mbox{}\par

\noindent\textbf{1：条件内で適切} 既存技術の改良の積重ね。原典の条件内では適切な記述であり、不適切なものを選ぶ問の選択対象ではない。\par

\noindent\textbf{2：条件内で適切} 機械化で大工場へ。原典の条件内では適切な記述であり、不適切なものを選ぶ問の選択対象ではない。\par

\noindent\textbf{3：条件内で適切} 娯楽自動機構から転用もある。原典の条件内では適切な記述であり、不適切なものを選ぶ問の選択対象ではない。\par

\noindent\textbf{4：不適切・選択対象} 大学研究者主導の産学協同例が多い。判定と一致する。\par

\noindent\textbf{5：条件内で適切} 労働と児童労働が社会問題。原典の条件内では適切な記述であり、不適切なものを選ぶ問の選択対象ではない。\par

\paragraph{再利用する考え方} 発明、改良、普及、研究組織と生産制度。 指定年度・対象範囲・単位・否定を先に固定し、他年度の番号や現行値を転用しない。

\paragraph{共通理論} \hyperref[theory:env:industrial]{第\ref*{theory:env:industrial}節}

\par 参考：\cite{env58_steam}。

\par 難易度 3（校準初稿）、優先度未評価。

\subsection{平成29年度 問5-5}
\label{q:h29-reg-basic-5-5}

\noindent\textbf{原典：} sources/H29/basic.pdf、PDF物理19頁。

\textbf{公式正答：} 2（single）。

\par\textbf{検討未確定：} 個別の娯楽機構転用史料は未取得。大学との個別接点も無かったという全称にはしない。

\paragraph{問題の要旨} 産業革命の担い手を大学主導へ一般化しない。

\Needspace{4\baselineskip}\paragraph{確認した条件・記述（独自要約）}\mbox{}\par

\noindent\textbf{時点：} 平成29年度に指定された資料・制度。\par

\noindent\textbf{照合：} 原典18画像で否定・五肢・組合せ表を確認。後年値で置換しない。\par

\paragraph{問われる核心} 産業革命の担い手を大学主導へ一般化しない。

\paragraph{必要知識} 産業革命、改良と普及、制度と労働。

\paragraph{考え方・正答を選ぶ理由（査読待ち初稿）} 2が不適切。職人・発明家・企業家の改良と事業化を大学研究者主導へ一括置換できない。機械技術の利用、生産組織、労働制度という別の側面を合わせて整理する。

\Needspace{4\baselineskip}\paragraph{選択肢の検討（査読待ち初稿）}\mbox{}\par

\noindent\textbf{1：条件内で適切} 既存改良の積重ね。条件内で適切な記述であり、不適切な記述を選ぶ対象ではない。\par

\noindent\textbf{2：不適切・選択対象} 欧州大学研究者主導が多い。判定と一致する。\par

\noindent\textbf{3：条件内で適切} 大工場への機械化。条件内で適切な記述であり、不適切な記述を選ぶ対象ではない。\par

\noindent\textbf{4：条件内で適切} 娯楽機構から転用例。条件内で適切な記述であり、不適切な記述を選ぶ対象ではない。\par

\noindent\textbf{5：条件内で適切} 児童労働等の社会問題。条件内で適切な記述であり、不適切な記述を選ぶ対象ではない。\par

\paragraph{再利用する考え方} 産業革命、改良と普及、制度と労働。 集計対象・時点・法の主体・選択方向を固定し、同型でも元の五肢へ戻す。

\paragraph{共通理論} \hyperref[theory:env:industrial]{第\ref*{theory:env:industrial}節}

\par 参考：\cite{env58_steam}。

\par 難易度 3（校準初稿）、優先度未評価。

\subsection{令和6年度 問5-5}
\label{q:r06-reg-basic-5-5}

\noindent\textbf{原典：} sources/R06/basic.pdf、PDF物理20頁。

\textbf{公式正答：} 3（single）。

\par\textbf{検討未確定：} 娯楽機構の具体的転用例と地域ごとの担い手について原典史料追加照合。

\paragraph{問題の要旨} 産業革命は大学主導の開発がほとんどではない。

\Needspace{4\baselineskip}\paragraph{確認した条件・記述（独自要約）}\mbox{}\par

\noindent\textbf{原典：} 全五肢と否定方向・組合せ表を画像確認。\par

\noindent\textbf{資料時点：} 令和6年度の指定資料・統計年を保持し、後年値に置換しない。\par

\paragraph{問われる核心} 産業革命は大学主導の開発がほとんどではない。

\paragraph{必要知識} 技術改良の担い手、工場制、生産性と労働問題。

\paragraph{考え方・正答を選ぶ理由（査読待ち初稿）} 3が不適切。18〜19世紀の産業技術は職人の経験、発明家の工夫、資本・市場の需要等から発達し、ほぼ全てを現代的産学共同研究と定義できない。科学・技術・経済・労働の関係を区別する。

\Needspace{4\baselineskip}\paragraph{選択肢の検討（査読待ち初稿）}\mbox{}\par

\noindent\textbf{1：条件内で適切} 多くは既存技術の改良の積み重ね。技術革新は段階的な改良を含む。\par

\noindent\textbf{2：条件内で適切} 娯楽用機構のアイデアが転用。機構や技能の用途間移転がある。\par

\noindent\textbf{3：不適切・選択対象} ほとんどが大学研究者主導の産学協同。職人・発明家・事業家等の活動を落とす過度な一般化で不適切。\par

\noindent\textbf{4：条件内で適切} 機械化で大工場へ。動力と機械の集積が工場制へつながる。\par

\noindent\textbf{5：条件内で適切} 効率化と長時間・児童労働等の問題。生産性向上と労働者の環境改善は自動的には同時に実現しない。\par

\paragraph{再利用する考え方} 技術改良の担い手、工場制、生産性と労働問題。 否定方向と資料年、分母と単位を確認して五肢へ戻る。

\paragraph{共通理論} \hyperref[theory:env:industrial]{第\ref*{theory:env:industrial}節}

\par 難易度 3（校準初稿）、優先度未評価。

\section{知的財産の範囲は特許発明より広い}

\subsection{平成27年度 問5-5}
\label{q:h27-reg-basic-5-5}

\noindent\textbf{原典：} sources/H27/basic.pdf、PDF物理17頁。

\textbf{公式正答：} 3（single）。

\paragraph{問題の要旨} 知的財産の範囲は特許発明より広い。

\Needspace{4\baselineskip}\paragraph{確認した条件・記述（独自要約）}\mbox{}\par

\noindent\textbf{時点：} 平成27年度に指定された資料・制度。\par

\noindent\textbf{照合：} 原典21画像で否定・五肢・組合せ表を確認。後年値で置換しない。\par

\paragraph{問われる核心} 知的財産の範囲は特許発明より広い。

\paragraph{必要知識} 発明、知的財産、産業財産権、基本法の責務。

\paragraph{考え方・正答を選ぶ理由（査読待ち初稿）} 特許法の発明は自然法則を利用した高度な技術的思想。知的財産基本法の知財は著作物・商標・商号・営業秘密なども含むので、自然法則利用へ限定する3が不適切。法律名とその定義の範囲を別々に確認する。

\Needspace{4\baselineskip}\paragraph{選択肢の検討（査読待ち初稿）}\mbox{}\par

\noindent\textbf{1：条件内で適切} 特許は発明保護利用と産業発達。原典の条件内では適切な記述であり、不適切なものを選ぶ問の選択対象ではない。\par

\noindent\textbf{2：条件内で適切} 発明は自然法則利用の高度技術思想。原典の条件内では適切な記述であり、不適切なものを選ぶ問の選択対象ではない。\par

\noindent\textbf{3：不適切・選択対象} 知財は自然法則利用のものだけ。判定と一致する。\par

\noindent\textbf{4：条件内で適切} 基本法は創造保護活用の基本事項。原典の条件内では適切な記述であり、不適切なものを選ぶ問の選択対象ではない。\par

\noindent\textbf{5：条件内で適切} 国に施策策定実施責務。原典の条件内では適切な記述であり、不適切なものを選ぶ問の選択対象ではない。\par

\paragraph{再利用する考え方} 発明、知的財産、産業財産権、基本法の責務。 指定年度・対象範囲・単位・否定を先に固定し、他年度の番号や現行値を転用しない。

\paragraph{共通理論} \hyperref[theory:env:ip]{第\ref*{theory:env:ip}節}

\par 参考：\cite{ip_patent_statute,ip_broad_scope}。

\par 難易度 2（校準初稿）、優先度未評価。

\subsection{令和元年度 問5-6}
\label{q:r01-reg-basic-5-6}

\noindent\textbf{原典：} sources/R01/basic.pdf、PDF物理22頁。

\textbf{公式正答：} 5（single）。

\paragraph{問題の要旨} 特許と基本法の定義、国の責務。

\Needspace{4\baselineskip}\paragraph{確認した条件・記述（独自要約）}\mbox{}\par

\noindent\textbf{時点：} 令和元年度に指定された資料・制度。\par

\noindent\textbf{照合：} 原典18画像で否定・五肢・組合せ表を確認。後年値で置換しない。\par

\paragraph{問われる核心} 特許と基本法の定義、国の責務。

\paragraph{必要知識} 特許法1・2条、基本法2・5条、対象の広さ。

\paragraph{考え方・正答を選ぶ理由（査読待ち初稿）} 知的財産基本法は国の施策策定・実施責務を定めるので、無しとする5が誤り。発明の定義と知財全体の定義を区別し、商標や商号は特許発明の要件で排除しない。

\Needspace{4\baselineskip}\paragraph{選択肢の検討（査読待ち初稿）}\mbox{}\par

\noindent\textbf{1：条件内で適切} 発明は自然法則利用高度思想。条件内で適切な記述であり、不適切な記述を選ぶ対象ではない。\par

\noindent\textbf{2：条件内で適切} 特許保護利用と産業。条件内で適切な記述であり、不適切な記述を選ぶ対象ではない。\par

\noindent\textbf{3：条件内で適切} 知財には商標商号も。条件内で適切な記述であり、不適切な記述を選ぶ対象ではない。\par

\noindent\textbf{4：条件内で適切} 基本法は創造保護活用。条件内で適切な記述であり、不適切な記述を選ぶ対象ではない。\par

\noindent\textbf{5：不適切・選択対象} 国は施策責務なし。判定と一致する。\par

\paragraph{再利用する考え方} 特許法1・2条、基本法2・5条、対象の広さ。 集計対象・時点・法の主体・選択方向を固定し、同型でも元の五肢へ戻す。

\paragraph{共通理論} \hyperref[theory:env:ip]{第\ref*{theory:env:ip}節}

\par 参考：\cite{ip_patent_statute,ip_broad_scope}。

\par 難易度 2（校準初稿）、優先度未評価。

\section{気候緩和と適応、LCAの境界}

\subsection{平成25年度 問5-3}
\label{q:h25-reg-basic-5-3}

\noindent\textbf{原典：} sources/H25/basic.pdf、PDF物理20頁。

\textbf{公式正答：} 4（single）。

\par\textbf{検討未確定：} CDMとCFPの出題当時規則全文は未取得。定義の独立判断と条文・規格の確定は分ける。

\paragraph{問題の要旨} 気候緩和と適応、LCAの境界。

\Needspace{4\baselineskip}\paragraph{確認した条件・記述（独自要約）}\mbox{}\par

\noindent\textbf{時点：} 平成25年度に指定された資料・制度。\par

\noindent\textbf{照合：} 原典21画像で否定・五肢・組合せ表を確認。後年値で置換しない。\par

\paragraph{問われる核心} 気候緩和と適応、LCAの境界。

\paragraph{必要知識} 緩和・適応、機能単位とシステム境界、環境基準。

\paragraph{考え方・正答を選ぶ理由（査読待ち初稿）} 排出削減や吸収を通じて原因を減らすのが緩和。影響に対して被害を減らす対策が適応なので4が取り違え。CFPは温室効果ガスのCO2換算に着目し、LCAはより広い環境影響の評価。環境基準は環境の望ましい状態で、個別工場の排出許容値と同じではない。

\Needspace{4\baselineskip}\paragraph{選択肢の検討（査読待ち初稿）}\mbox{}\par

\noindent\textbf{1：条件内で適切} 環境基準は維持が望ましい。原典の条件内では適切な記述であり、不適切なものを選ぶ問の選択対象ではない。\par

\noindent\textbf{2：条件内で適切} CDMは京都の認証削減。原典の条件内では適切な記述であり、不適切なものを選ぶ問の選択対象ではない。\par

\noindent\textbf{3：条件内で適切} CFPは全過程CO2換算。原典の条件内では適切な記述であり、不適切なものを選ぶ問の選択対象ではない。\par

\noindent\textbf{4：不適切・選択対象} 適応は排出削減で原因抑制。判定と一致する。\par

\noindent\textbf{5：条件内で適切} LCAは資源から廃棄まで評価。原典の条件内では適切な記述であり、不適切なものを選ぶ問の選択対象ではない。\par

\paragraph{再利用する考え方} 緩和・適応、機能単位とシステム境界、環境基準。 指定年度・対象範囲・単位・否定を先に固定し、他年度の番号や現行値を転用しない。

\paragraph{共通理論} \hyperref[theory:env:lca]{第\ref*{theory:env:lca}節}

\par 参考：\cite{env58_pmstandard}。

\par 難易度 3（校準初稿）、優先度未評価。

\subsection{令和元年度 問5-2}
\label{q:r01-reg-basic-5-2}

\noindent\textbf{原典：} sources/R01/basic.pdf、PDF物理20頁。

\textbf{公式正答：} 2（single）。

\par\textbf{検討未確定：} JCMの2019運用ルール全文とCFP規格は追加照合。

\paragraph{問題の要旨} 排出原因を減らす緩和と影響への適応。

\Needspace{4\baselineskip}\paragraph{確認した条件・記述（独自要約）}\mbox{}\par

\noindent\textbf{時点：} 令和元年度に指定された資料・制度。\par

\noindent\textbf{照合：} 原典18画像で否定・五肢・組合せ表を確認。後年値で置換しない。\par

\paragraph{問われる核心} 排出原因を減らす緩和と影響への適応。

\paragraph{必要知識} 緩和適応、LCA、CFP、JCM、環境基準。

\paragraph{考え方・正答を選ぶ理由（査読待ち初稿）} 原因の削減を適応と呼ぶ2が誤り。排出削減・吸収は緩和で、影響の被害回避が適応。CFPとLCA、環境基準の定義は対象量と境界を固定して判定する。JCMは削減・吸収への貢献を定量化する枠組みで、効果の無条件な認定ではない。

\Needspace{4\baselineskip}\paragraph{選択肢の検討（査読待ち初稿）}\mbox{}\par

\noindent\textbf{1：条件内で適切} JCMは貢献定量化と目標活用。条件内で適切な記述であり、不適切な記述を選ぶ対象ではない。\par

\noindent\textbf{2：不適切・選択対象} 適応は原因のガス排出削減。判定と一致する。\par

\noindent\textbf{3：条件内で適切} CFP全過程CO2換算。条件内で適切な記述であり、不適切な記述を選ぶ対象ではない。\par

\noindent\textbf{4：条件内で適切} LCA全段階影響評価。条件内で適切な記述であり、不適切な記述を選ぶ対象ではない。\par

\noindent\textbf{5：条件内で適切} 環境基準は望ましい状態。条件内で適切な記述であり、不適切な記述を選ぶ対象ではない。\par

\paragraph{再利用する考え方} 緩和適応、LCA、CFP、JCM、環境基準。 集計対象・時点・法の主体・選択方向を固定し、同型でも元の五肢へ戻す。

\paragraph{共通理論} \hyperref[theory:env:lca]{第\ref*{theory:env:lca}節}

\par 参考：\cite{env59_ndc,env58_pmstandard}。

\par 難易度 3（校準初稿）、優先度未評価。

\section{気液の同質量比較からLNGの体積比}

\subsection{平成23年度 問5-3}
\label{q:h23-reg-basic-5-3}

\noindent\textbf{原典：} sources/H23/basic.pdf、PDF物理16頁。

\textbf{公式正答：} 2（single）。

\paragraph{問題の要旨} 気液の同質量比較からLNGの体積比。

\Needspace{4\baselineskip}\paragraph{確認した条件・記述（独自要約）}\mbox{}\par

\noindent\textbf{気体：} 0℃1気圧、純CH4の理想気体。\par

\noindent\textbf{液体：} 密度425kg/m\textsuperscript{3}一定とする原典モデル。\par

\paragraph{問われる核心} 気液の同質量比較からLNGの体積比。

\paragraph{必要知識} 理想気体モル体積、分子量、密度と体積、単位。

\paragraph{考え方・正答を選ぶ理由（査読待ち初稿）} 同じ物質量を比べる。0℃1気圧で1molの理想気体は約22.4L、メタン質量16g。液体体積は0.016kg/(425kg/m\textsuperscript{3})=\allowbreak{}3.76×10\textasciicircum{}-5m\textsuperscript{3}、気体は.0224m\textsuperscript{3}。比は.00168≈1/595なので最も近い1/600。密度比ρgas/ρliquidでも同じ。冷却の温度条件を気体標準状態と混同しない。

\Needspace{4\baselineskip}\paragraph{選択肢の検討（査読待ち初稿）}\mbox{}\par

\noindent\textbf{1：不適切} 1/400。式・定義・対応の独立判定と一致しない。\par

\noindent\textbf{2：適切・選択対象} 1/600。判定と一致する。\par

\noindent\textbf{3：不適切} 1/800。式・定義・対応の独立判定と一致しない。\par

\noindent\textbf{4：不適切} 1/1000。式・定義・対応の独立判定と一致しない。\par

\noindent\textbf{5：不適切} 1/1200。式・定義・対応の独立判定と一致しない。\par

\paragraph{再利用する考え方} 理想気体モル体積、分子量、密度と体積、単位。 指定年度・対象範囲・単位・否定を先に固定し、他年度の番号や現行値を転用しない。

\paragraph{共通理論} \hyperref[theory:env:lng]{第\ref*{theory:env:lng}節}

\paragraph{関連問題} \hyperref[q:r05-reg-basic-5-4]{R05-REG-BASIC-5-4}

\par 難易度 2（校準初稿）、優先度未評価。

\subsection{平成29年度 問5-3}
\label{q:h29-reg-basic-5-3}

\noindent\textbf{原典：} sources/H29/basic.pdf、PDF物理19頁。

\textbf{公式正答：} 4（single）。

\paragraph{問題の要旨} 標準気体と液体の密度比からLNG体積比。

\Needspace{4\baselineskip}\paragraph{確認した条件・記述（独自要約）}\mbox{}\par

\noindent\textbf{標準状態：} 0℃1気圧、純CH4理想気体。\par

\noindent\textbf{液体：} 425kg/m\textsuperscript{3}一定。\par

\paragraph{問われる核心} 標準気体と液体の密度比からLNG体積比。

\paragraph{必要知識} 理想気体、密度、同質量の体積比較。

\paragraph{考え方・正答を選ぶ理由（査読待ち初稿）} 1molのCH4は約16gで、0℃1気圧の体積22.4L。液体では.016/425=\allowbreak{}3.76×10\textasciicircum{}-5m\textsuperscript{3}、気体.0224m\textsuperscript{3}で比≈1/595。1/600の4が近い。平成23年度の候補順をそのまま使わない。

\Needspace{4\baselineskip}\paragraph{選択肢の検討（査読待ち初稿）}\mbox{}\par

\noindent\textbf{1：不適切} 1/1200。式・定義・対応の独立判定と一致しない。\par

\noindent\textbf{2：不適切} 1/1000。式・定義・対応の独立判定と一致しない。\par

\noindent\textbf{3：不適切} 1/800。式・定義・対応の独立判定と一致しない。\par

\noindent\textbf{4：適切・選択対象} 1/600。判定と一致する。\par

\noindent\textbf{5：不適切} 1/400。式・定義・対応の独立判定と一致しない。\par

\paragraph{再利用する考え方} 理想気体、密度、同質量の体積比較。 集計対象・時点・法の主体・選択方向を固定し、同型でも元の五肢へ戻す。

\paragraph{共通理論} \hyperref[theory:env:lng]{第\ref*{theory:env:lng}節}

\paragraph{関連問題} \hyperref[q:r05-reg-basic-5-4]{R05-REG-BASIC-5-4}

\par 難易度 2（校準初稿）、優先度未評価。

\subsection{令和5年度 問5-4}
\label{q:r05-reg-basic-5-4}

\noindent\textbf{原典：} sources/R05/basic.pdf、PDF物理17頁。

\textbf{公式正答：} 2（single）。

\paragraph{問題の要旨} 0℃1気圧のメタンと液体密度から体積比。

\Needspace{4\baselineskip}\paragraph{確認した条件・記述（独自要約）}\mbox{}\par

\noindent\textbf{気体：} 純CH4理想気体、0℃・1気圧。\par

\noindent\textbf{液体：} 密度425kg/m\textsuperscript{3}一定。\par

\paragraph{問われる核心} 0℃1気圧のメタンと液体密度から体積比。

\paragraph{必要知識} 同質量の体積比、理想気体の標準体積、密度。

\paragraph{考え方・正答を選ぶ理由（査読待ち初稿）} メタン1molの質量は約0.016kg。指定液体密度425kg/m\textsuperscript{3}より液体体積0.016/425=\allowbreak{}3.76×10\textasciicircum{}-5m\textsuperscript{3}、0℃1気圧の理想気体体積は約0.0224m\textsuperscript{3}。比は0.00168≈1/595で1/600の2。液体側に気体の状態方程式を適用しない。

\Needspace{4\baselineskip}\paragraph{選択肢の検討（査読待ち初稿）}\mbox{}\par

\noindent\textbf{1：不適切} 1/400。独立に求めた組合せ 1/600 と異なる。\par

\noindent\textbf{2：適切・選択対象} 1/600。独立に求めた組合せ 1/600 と一致する。\par

\noindent\textbf{3：不適切} 1/800。独立に求めた組合せ 1/600 と異なる。\par

\noindent\textbf{4：不適切} 1/1000。独立に求めた組合せ 1/600 と異なる。\par

\noindent\textbf{5：不適切} 1/1200。独立に求めた組合せ 1/600 と異なる。\par

\paragraph{再利用する考え方} 同質量の体積比、理想気体の標準体積、密度。 否定方向と資料年、分母と単位を確認して五肢へ戻る。

\paragraph{共通理論} \hyperref[theory:env:lng]{第\ref*{theory:env:lng}節}

\paragraph{関連問題} \hyperref[q:h23-reg-basic-5-3]{H23-REG-BASIC-5-3}, \hyperref[q:h29-reg-basic-5-3]{H29-REG-BASIC-5-3}

\par 難易度 2（校準初稿）、優先度未評価。

\section{2012年石油統計の対象範囲を固定}

\subsection{平成25年度 問5-1}
\label{q:h25-reg-basic-5-1}

\noindent\textbf{原典：} sources/H25/basic.pdf、PDF物理19頁。

\textbf{公式正答：} 5（single）。

\par\textbf{検討未確定：} 2012暦年の財務省四項目原表を全て取得していない。取得白書は2013年度83.6\%とサウジ首位を支持する補助で、指定年の確定証跡とは区別する。

\paragraph{問題の要旨} 2012年石油統計の対象範囲を固定。

\Needspace{4\baselineskip}\paragraph{確認した条件・記述（独自要約）}\mbox{}\par

\noindent\textbf{時点：} 平成25年度に指定された資料・制度。\par

\noindent\textbf{照合：} 原典21画像で否定・五肢・組合せ表を確認。後年値で置換しない。\par

\paragraph{問われる核心} 2012年石油統計の対象範囲を固定。

\paragraph{必要知識} 暦年と年度、数量と金額、統計品目、輸送経路。

\paragraph{考え方・正答を選ぶ理由（査読待ち初稿）} 2012年の指定統計では中東依存度約83\%、主要経路ホルムズ海峡、最大輸入先サウジアラビア、原油と石油製品の総輸入額比約21\%で5。数量の依存度と金額比は別であり、石炭・LNGを含む鉱物性燃料全体の比と原油・製品だけの比も違う。

\Needspace{4\baselineskip}\paragraph{選択肢の検討（査読待ち初稿）}\mbox{}\par

\noindent\textbf{1：不適切} 93/マラッカ/クウェート/21。式・定義・対応の独立判定と一致しない。\par

\noindent\textbf{2：不適切} 93/ホルムズ/サウジ/42。式・定義・対応の独立判定と一致しない。\par

\noindent\textbf{3：不適切} 83/マラッカ/サウジ/42。式・定義・対応の独立判定と一致しない。\par

\noindent\textbf{4：不適切} 83/ホルムズ/クウェート/42。式・定義・対応の独立判定と一致しない。\par

\noindent\textbf{5：適切・選択対象} 83/ホルムズ/サウジ/21。判定と一致する。\par

\paragraph{再利用する考え方} 暦年と年度、数量と金額、統計品目、輸送経路。 指定年度・対象範囲・単位・否定を先に固定し、他年度の番号や現行値を転用しない。

\paragraph{共通理論} \hyperref[theory:env:oil2012]{第\ref*{theory:env:oil2012}節}

\par 参考：\cite{env58_oil}。

\par 難易度 3（校準初稿）、優先度未評価。

\section{リサイクル制度の対象と責任主体}

\subsection{平成24年度 問5-1}
\label{q:h24-reg-basic-5-1}

\noindent\textbf{原典：} sources/H24/basic.pdf、PDF物理15頁。

\textbf{公式正答：} 3（single）。

\par\textbf{検討未確定：} 出題時点の建設規模・バーゼル条約対象と例外は追加の法令履歴照合。現行FAQだけで全過去条文の確定とはしない。

\paragraph{問題の要旨} リサイクル制度の対象と責任主体。

\Needspace{4\baselineskip}\paragraph{確認した条件・記述（独自要約）}\mbox{}\par

\noindent\textbf{時点：} 平成24年度に指定された資料・制度。\par

\noindent\textbf{照合：} 原典21画像で否定・五肢・組合せ表を確認。後年値で置換しない。\par

\paragraph{問われる核心} リサイクル制度の対象と責任主体。

\paragraph{必要知識} 法の目的、対象条件、排出者・自治体・小売・製造者の役割。

\paragraph{考え方・正答を選ぶ理由（査読待ち初稿）} 3が家電リサイクル法の小売業者の役割に対応。基本法はまず発生抑制。容器包装は三品目のみではなく、分別収集と再商品化の役割・費用も異なる。建設は規模等の対象条件があり全工事・発注者への実施義務という全称が違う。バーゼルは再利用目的の有害廃棄物でも規制を受け、条約成立の背景の輸送方向も逆。

\Needspace{4\baselineskip}\paragraph{選択肢の検討（査読待ち初稿）}\mbox{}\par

\noindent\textbf{1：不適切} 基本法は再生利用最優先。式・定義・対応の独立判定と一致しない。\par

\noindent\textbf{2：不適切} 容器三品目の全費用を販売者。式・定義・対応の独立判定と一致しない。\par

\noindent\textbf{3：適切・選択対象} 家電小売の引取・製造者への引渡し。判定と一致する。\par

\noindent\textbf{4：不適切} 全建設工事の発注者が分別再資源化。式・定義・対応の独立判定と一致しない。\par

\noindent\textbf{5：不適切} 再利用なら有害物輸出規制なし。式・定義・対応の独立判定と一致しない。\par

\paragraph{再利用する考え方} 法の目的、対象条件、排出者・自治体・小売・製造者の役割。 指定年度・対象範囲・単位・否定を先に固定し、他年度の番号や現行値を転用しない。

\paragraph{共通理論} \hyperref[theory:env:recyclelaw]{第\ref*{theory:env:recyclelaw}節}

\par 参考：\cite{env_circular_law,env_appliance}。

\par 難易度 3（校準初稿）、優先度未評価。

\subsection{令和元年度再試験 問5-2}
\label{q:r01-ret-basic-5-2}

\noindent\textbf{原典：} sources/R01\_retest/basic.pdf、PDF物理18頁。

\textbf{公式正答：} 1（single）。

\par\textbf{検討未確定：} 出題時バーゼル・建設の法令履歴と規模条件全文は追加照合。

\paragraph{問題の要旨} 家電小売の役割と循環制度の限定。

\Needspace{4\baselineskip}\paragraph{確認した条件・記述（独自要約）}\mbox{}\par

\noindent\textbf{時点：} 令和元年度再試験に指定された資料・制度。\par

\noindent\textbf{照合：} 原典18画像で否定・五肢・組合せ表を確認。後年値で置換しない。\par

\paragraph{問われる核心} 家電小売の役割と循環制度の限定。

\paragraph{必要知識} 役割、対象と例外、優先順位、国境移動。

\paragraph{考え方・正答を選ぶ理由（査読待ち初稿）} 1が家電小売の役割。再利用の有害廃棄物取引も規制対象となり得る。容器包装は三品目・販売者全費用ではなく、建設には対象条件と実施主体、循環の優先順には発生抑制がある。

\Needspace{4\baselineskip}\paragraph{選択肢の検討（査読待ち初稿）}\mbox{}\par

\noindent\textbf{1：適切・選択対象} 家電小売の引取と製造者への引渡し。判定と一致する。\par

\noindent\textbf{2：不適切} 再利用ならバーゼル規制なし。式・定義・対応の独立判定と一致しない。\par

\noindent\textbf{3：不適切} 容器三品目を販売者全費用。式・定義・対応の独立判定と一致しない。\par

\noindent\textbf{4：不適切} 全工事の発注者に分別実施。式・定義・対応の独立判定と一致しない。\par

\noindent\textbf{5：不適切} 基本法は再生最優先。式・定義・対応の独立判定と一致しない。\par

\paragraph{再利用する考え方} 役割、対象と例外、優先順位、国境移動。 集計対象・時点・法の主体・選択方向を固定し、同型でも元の五肢へ戻す。

\paragraph{共通理論} \hyperref[theory:env:recyclelaw]{第\ref*{theory:env:recyclelaw}節}

\par 参考：\cite{env_circular_law,env_appliance}。

\par 難易度 3（校準初稿）、優先度未評価。

\section{資源量と太陽光面積の桁を比較}

\subsection{平成26年度 問5-4}
\label{q:h26-reg-basic-5-4}

\noindent\textbf{原典：} sources/H26/basic.pdf、PDF物理17頁。

\textbf{公式正答：} 4（single）。

\par\textbf{検討未確定：} 概算の日射・需要は桁比較であり厳密な上限証明ではない。石炭石油の指定時点埋蔵量、海水U総量、全球純生産と人類エネルギーの各原表は追加照合。初回記録の昼夜連続1kW仮定を解説では日照条件を示す概算へ改めた。

\paragraph{問題の要旨} 資源量と太陽光面積の桁を比較。

\Needspace{4\baselineskip}\paragraph{確認した条件・記述（独自要約）}\mbox{}\par

\noindent\textbf{時点：} 平成26年度に指定された資料・制度。\par

\noindent\textbf{照合：} 原典21画像で否定・五肢・組合せ表を確認。後年値で置換しない。\par

\paragraph{問われる核心} 資源量と太陽光面積の桁を比較。

\paragraph{必要知識} 面積換算、電力とエネルギー、存在量と可採埋蔵量。

\paragraph{考え方・正答を選ぶ理由（査読待ち初稿）} 100km\textsuperscript{2}=10\textasciicircum{}8m\textsuperscript{2}。例えば1kW/m\textsuperscript{2}を昼間平均12h/日、効率を非現実的に100\%としても年約438TWhで、当時の日本の年間電力需要約1000TWhに届かない。現実には夜間・天候・入射角・変換損失でさらに小さいので4。資源の存在量と経済的に採れる確認埋蔵量も同じではなく、後者は探鉱・価格・技術で増え得る。

\Needspace{4\baselineskip}\paragraph{選択肢の検討（査読待ち初稿）}\mbox{}\par

\noindent\textbf{1：条件内で適切} 石炭埋蔵エネルギーが石油超。原典の条件内では適切な記述であり、不適切なものを選ぶ問の選択対象ではない。\par

\noindent\textbf{2：条件内で適切} 石油確認埋蔵量は危機時より増。原典の条件内では適切な記述であり、不適切なものを選ぶ問の選択対象ではない。\par

\noindent\textbf{3：条件内で適切} 海水Uは1億t超。原典の条件内では適切な記述であり、不適切なものを選ぶ問の選択対象ではない。\par

\noindent\textbf{4：不適切・選択対象} 100km\textsuperscript{2}太陽光が日本需要超。判定と一致する。\par

\noindent\textbf{5：条件内で適切} 全球の光合成固定が人類需要超。原典の条件内では適切な記述であり、不適切なものを選ぶ問の選択対象ではない。\par

\paragraph{再利用する考え方} 面積換算、電力とエネルギー、存在量と可採埋蔵量。 指定年度・対象範囲・単位・否定を先に固定し、他年度の番号や現行値を転用しない。

\paragraph{共通理論} \hyperref[theory:env:resource]{第\ref*{theory:env:resource}節}

\par 難易度 3（校準初稿）、優先度未評価。

\section{下水処理の段階と環境対策}

\subsection{平成25年度 問5-4}
\label{q:h25-reg-basic-5-4}

\noindent\textbf{原典：} sources/H25/basic.pdf、PDF物理21頁。

\textbf{公式正答：} 2（single）。

\par\textbf{検討未確定：} 当時土壌対策の掘削除去割合と施設の詳細法定仕様は追加照合。緑化の効果は気象・水・建物条件に依存。

\paragraph{問題の要旨} 下水処理の段階と環境対策。

\Needspace{4\baselineskip}\paragraph{確認した条件・記述（独自要約）}\mbox{}\par

\noindent\textbf{時点：} 平成25年度に指定された資料・制度。\par

\noindent\textbf{照合：} 原典21画像で否定・五肢・組合せ表を確認。後年値で置換しない。\par

\paragraph{問われる核心} 下水処理の段階と環境対策。

\paragraph{必要知識} 一次・二次・高度処理、物理と生物反応、汚染経路。

\paragraph{考え方・正答を選ぶ理由（査読待ち初稿）} 活性汚泥による生物処理は一般に二次処理。一次は沈殿等の物理処理であり、2が不適切。管理型処分場は浸出水等による汚染を防ぐ。緑化は遮蔽・蒸発散、ダイオキシン対策は燃焼と再合成防止の工程を分けて説明する。

\Needspace{4\baselineskip}\paragraph{選択肢の検討（査読待ち初稿）}\mbox{}\par

\noindent\textbf{1：条件内で適切} 管理型の遮水・浸出水処理。原典の条件内では適切な記述であり、不適切なものを選ぶ問の選択対象ではない。\par

\noindent\textbf{2：不適切・選択対象} 活性汚泥は一次処理。判定と一致する。\par

\noindent\textbf{3：条件内で適切} 緑化の表面温度と空調低減。原典の条件内では適切な記述であり、不適切なものを選ぶ問の選択対象ではない。\par

\noindent\textbf{4：条件内で適切} 土壌対策に掘削除去が多い。原典の条件内では適切な記述であり、不適切なものを選ぶ問の選択対象ではない。\par

\noindent\textbf{5：条件内で適切} 完全燃焼と排ガス急冷。原典の条件内では適切な記述であり、不適切なものを選ぶ問の選択対象ではない。\par

\paragraph{再利用する考え方} 一次・二次・高度処理、物理と生物反応、汚染経路。 指定年度・対象範囲・単位・否定を先に固定し、他年度の番号や現行値を転用しない。

\paragraph{共通理論} \hyperref[theory:env:treatment]{第\ref*{theory:env:treatment}節}

\par 難易度 3（校準初稿）、優先度未評価。

\subsection{令和3年度 問5-2}
\label{q:r03-reg-basic-5-2}

\noindent\textbf{原典：} sources/R03/basic.pdf、PDF物理22頁。

\textbf{公式正答：} 5（single）。

\par\textbf{検討未確定：} 土壌汚染対策実績の時点別割合、管理型処分場設備の当時法令条文は追加照合。

\paragraph{問題の要旨} 活性汚泥による生物処理は二次処理。

\Needspace{4\baselineskip}\paragraph{確認した条件・記述（独自要約）}\mbox{}\par

\noindent\textbf{原典：} 全五肢と否定方向・組合せ表を画像確認。\par

\noindent\textbf{資料時点：} 令和3年度の指定資料・統計年を保持し、後年値に置換しない。\par

\paragraph{問われる核心} 活性汚泥による生物処理は二次処理。

\paragraph{必要知識} 一次・二次・高度処理、排ガス急冷、遮水と浸出水。

\paragraph{考え方・正答を選ぶ理由（査読待ち初稿）} 5が不適切。沈砂・沈殿で浮遊固形物等を除く段階と、微生物が有機物を分解する段階を分ける。高度な窒素・リン除去等と生物処理一般の二次処理を混同しない。

\Needspace{4\baselineskip}\paragraph{選択肢の検討（査読待ち初稿）}\mbox{}\par

\noindent\textbf{1：条件内で適切} 焼却排ガスを200℃以下へ急冷。完全燃焼と急冷により生成・再合成を抑える対策。\par

\noindent\textbf{2：条件内で適切} 緑化により熱侵入を抑制。日射遮蔽や蒸散により建物への熱負荷を下げる。\par

\noindent\textbf{3：条件内で適切} 管理型処分場に遮水・浸出水処理。浸出水の環境放出を管理する設備が必要。\par

\noindent\textbf{4：条件内で適切} 原位置浄化もあるが掘削除去実績が多い。原位置は掘削せず処理する意味。実績の多寡は当時資料に依存。\par

\noindent\textbf{5：不適切・選択対象} 活性汚泥は一次処理。一次処理は沈殿等の物理的処理。活性汚泥の生物処理は一般に二次処理である。\par

\paragraph{再利用する考え方} 一次・二次・高度処理、排ガス急冷、遮水と浸出水。 否定方向と資料年、分母と単位を確認して五肢へ戻る。

\paragraph{共通理論} \hyperref[theory:env:treatment]{第\ref*{theory:env:treatment}節}

\par 難易度 3（校準初稿）、優先度未評価。

\section{温対法の排出抑制と吸収作用}

\subsection{平成27年度 問5-2}
\label{q:h27-reg-basic-5-2}

\noindent\textbf{原典：} sources/H27/basic.pdf、PDF物理15頁。

\textbf{公式正答：} 1（single）。

\par\textbf{検討未確定：} 2015施行時の全五条項を履歴条文で一括確認する作業は未完了。現行の2050目標を当時条文へ追加しない。

\paragraph{問題の要旨} 温対法の排出抑制と吸収作用。

\Needspace{4\baselineskip}\paragraph{確認した条件・記述（独自要約）}\mbox{}\par

\noindent\textbf{時点：} 平成27年度に指定された資料・制度。\par

\noindent\textbf{照合：} 原典21画像で否定・五肢・組合せ表を確認。後年値で置換しない。\par

\paragraph{問われる核心} 温対法の排出抑制と吸収作用。

\paragraph{必要知識} 目的規定、吸収と排出、政令対象、計画。

\paragraph{考え方・正答を選ぶ理由（査読待ち初稿）} 目的は排出抑制だけでなく吸収作用の保全・強化も含む。言及しないという1が誤り。対象ガス、算定報告の主体・規模、地方推進体制、計画見直しは別の制度事項なので、吸収と排出の二側面を軸に確認する。

\Needspace{4\baselineskip}\paragraph{選択肢の検討（査読待ち初稿）}\mbox{}\par

\noindent\textbf{1：不適切・選択対象} 森林吸収の保全に言及なし。判定と一致する。\par

\noindent\textbf{2：条件内で適切} 対象七種の温室効果ガス。原典の条件内では適切な記述であり、不適切なものを選ぶ問の選択対象ではない。\par

\noindent\textbf{3：条件内で適切} 一定規模事業所別の算定報告。原典の条件内では適切な記述であり、不適切なものを選ぶ問の選択対象ではない。\par

\noindent\textbf{4：条件内で適切} 知事が推進員委嘱とセンター指定。原典の条件内では適切な記述であり、不適切なものを選ぶ問の選択対象ではない。\par

\noindent\textbf{5：条件内で適切} 少なくとも三年ごと計画検討。原典の条件内では適切な記述であり、不適切なものを選ぶ問の選択対象ではない。\par

\paragraph{再利用する考え方} 目的規定、吸収と排出、政令対象、計画。 指定年度・対象範囲・単位・否定を先に固定し、他年度の番号や現行値を転用しない。

\paragraph{共通理論} \hyperref[theory:env:warminglaw]{第\ref*{theory:env:warminglaw}節}

\par 難易度 3（校準初稿）、優先度未評価。

\section{排出抑制から適正処分までの優先順}

\subsection{平成23年度 問5-2}
\label{q:h23-reg-basic-5-2}

\noindent\textbf{原典：} sources/H23/basic.pdf、PDF物理16頁。

\textbf{公式正答：} 2（single）。

\par\textbf{検討未確定：} 取得した2001資料は基本方針案と施策実績であり、2010白書の引用箇所全文の一致は追加照合。

\paragraph{問題の要旨} 排出抑制から適正処分までの優先順。

\Needspace{4\baselineskip}\paragraph{確認した条件・記述（独自要約）}\mbox{}\par

\noindent\textbf{時点：} 平成23年度に指定された資料・制度。\par

\noindent\textbf{照合：} 原典21画像で否定・五肢・組合せ表を確認。後年値で置換しない。\par

\paragraph{問われる核心} 排出抑制から適正処分までの優先順。

\paragraph{必要知識} 3R、循環利用と処分、基準年度と目標年度。

\paragraph{考え方・正答を選ぶ理由（査読待ち初稿）} 旧来の廃棄物行政の公衆衛生と公害対策に、循環型社会形成が加わる。優先は排出抑制→再使用→再生利用→熱回収→残るものの適正処分。2001基本方針の最終処分目標は1997基準の概ね半分なので2。処分を埋立の一方式だけへ狭めない。

\Needspace{4\baselineskip}\paragraph{選択肢の検討（査読待ち初稿）}\mbox{}\par

\noindent\textbf{1：不適切} 生活環境/再資源化/適正処理/半分。式・定義・対応の独立判定と一致しない。\par

\noindent\textbf{2：適切・選択対象} 公衆衛生/再使用/適正処分/半分。判定と一致する。\par

\noindent\textbf{3：不適切} 生活環境/再資源化/適正処分/1/3。式・定義・対応の独立判定と一致しない。\par

\noindent\textbf{4：不適切} 公衆衛生/再資源化/適正処理/1/3。式・定義・対応の独立判定と一致しない。\par

\noindent\textbf{5：不適切} 生活環境/再使用/適正埋立/1/3。式・定義・対応の独立判定と一致しない。\par

\paragraph{再利用する考え方} 3R、循環利用と処分、基準年度と目標年度。 指定年度・対象範囲・単位・否定を先に固定し、他年度の番号や現行値を転用しない。

\paragraph{共通理論} \hyperref[theory:env:wasteorder]{第\ref*{theory:env:wasteorder}節}

\par 参考：\cite{env_circular_law,env58_wastegoal}。

\par 難易度 3（校準初稿）、優先度未評価。

\section{廃棄物用語の略称と対象}

\subsection{平成27年度 問5-1}
\label{q:h27-reg-basic-5-1}

\noindent\textbf{原典：} sources/H27/basic.pdf、PDF物理15頁。

\textbf{公式正答：} 2（single）。

\par\textbf{検討未確定：} 当時のRDF添加法・産廃業種条件・硫酸ピッチのガス生成に関する技術資料を追加照合。全RDFが同一添加処方という一般化をしない。

\paragraph{問題の要旨} 廃棄物用語の略称と対象。

\Needspace{4\baselineskip}\paragraph{確認した条件・記述（独自要約）}\mbox{}\par

\noindent\textbf{時点：} 平成27年度に指定された資料・制度。\par

\noindent\textbf{照合：} 原典21画像で否定・五肢・組合せ表を確認。後年値で置換しない。\par

\paragraph{問われる核心} 廃棄物用語の略称と対象。

\paragraph{必要知識} 略称、法定分類、処理特性と有害性。

\paragraph{考え方・正答を選ぶ理由（査読待ち初稿）} E-wasteはelectrical/electronic waste、廃電気電子機器。環境に良い又はほぼ全量再資源化可能という定義ではなく2が誤り。RDFは可燃廃物の加工燃料、バイオマスは化石資源を除く。産業廃棄物は事業活動由来なら全てという分類でなく法定種類と業種条件を確認する。

\Needspace{4\baselineskip}\paragraph{選択肢の検討（査読待ち初稿）}\mbox{}\par

\noindent\textbf{1：条件内で適切} RDFはごみ固形化燃料。原典の条件内では適切な記述であり、不適切なものを選ぶ問の選択対象ではない。\par

\noindent\textbf{2：不適切・選択対象} E-wasteはほぼ全て再資源化できるエコ廃物。判定と一致する。\par

\noindent\textbf{3：条件内で適切} バイオマスは化石除く生物資源。原典の条件内では適切な記述であり、不適切なものを選ぶ問の選択対象ではない。\par

\noindent\textbf{4：条件内で適切} 産廃は事業由来の法定20種類。原典の条件内では適切な記述であり、不適切なものを選ぶ問の選択対象ではない。\par

\noindent\textbf{5：条件内で適切} 硫酸ピッチは強酸性油泥。原典の条件内では適切な記述であり、不適切なものを選ぶ問の選択対象ではない。\par

\paragraph{再利用する考え方} 略称、法定分類、処理特性と有害性。 指定年度・対象範囲・単位・否定を先に固定し、他年度の番号や現行値を転用しない。

\paragraph{共通理論} \hyperref[theory:env:wasteterm]{第\ref*{theory:env:wasteterm}節}

\par 難易度 3（校準初稿）、優先度未評価。

\subsection{令和4年度 問5-2}
\label{q:r04-reg-basic-5-2}

\noindent\textbf{原典：} sources/R04/basic.pdf、PDF物理17頁。

\textbf{公式正答：} 1（single）。

\par\textbf{検討未確定：} 当時の全国両区分の総量とメダル事業の回収・製造記録全文を追加確認。

\paragraph{問題の要旨} 一般廃棄物と産業廃棄物の総量を逆転しない。

\Needspace{4\baselineskip}\paragraph{確認した条件・記述（独自要約）}\mbox{}\par

\noindent\textbf{原典：} 全五肢と否定方向・組合せ表を画像確認。\par

\noindent\textbf{資料時点：} 令和4年度の指定資料・統計年を保持し、後年値に置換しない。\par

\paragraph{問われる核心} 一般廃棄物と産業廃棄物の総量を逆転しない。

\paragraph{必要知識} 廃棄物区分・量・性状、再生燃料、資源回収。

\paragraph{考え方・正答を選ぶ理由（査読待ち初稿）} 1が不適切。量は区分別の全国年間排出で比較する。物質の有害性を区分する特別管理、資源由来のバイオマス、再生燃料RPF、回収材利用はそれぞれ別の分類軸である。

\Needspace{4\baselineskip}\paragraph{選択肢の検討（査読待ち初稿）}\mbox{}\par

\noindent\textbf{1：不適切・選択対象} 一般廃棄物の総量が産業廃棄物より多い。同じ近年の年間排出総量では産業廃棄物が一般廃棄物より大きく、逆の記述が不適切。\par

\noindent\textbf{2：条件内で適切} 特別管理産業廃棄物は爆発性・毒性・感染性等。危険性に応じ通常とは別の管理が必要となる。\par

\noindent\textbf{3：条件内で適切} バイオマスは化石資源を除く生物由来有機資源。生物起源でも長期地質過程で生じた化石燃料はここに含めない。\par

\noindent\textbf{4：条件内で適切} RPFは主に廃紙・廃プラスチック固形燃料。略称と対象材料を照合。RDF等の別燃料を混同しない。\par

\noindent\textbf{5：条件内で適切} 東京大会メダルは使用済小型家電の金属を利用。回収金属を資源として使う事業であり、新規採掘だけで製造したとの説明ではない。\par

\paragraph{再利用する考え方} 廃棄物区分・量・性状、再生燃料、資源回収。 否定方向と資料年、分母と単位を確認して五肢へ戻る。

\paragraph{共通理論} \hyperref[theory:env:wasteterm]{第\ref*{theory:env:wasteterm}節}

\par 難易度 3（校準初稿）、優先度未評価。

\section{工業型から都市型の汚染と二次生成}

\subsection{令和元年度 問5-1}
\label{q:r01-reg-basic-5-1}

\noindent\textbf{原典：} sources/R01/basic.pdf、PDF物理19頁。

\textbf{公式正答：} 3（single）。

\paragraph{問題の要旨} 工業型から都市型の汚染と二次生成。

\Needspace{4\baselineskip}\paragraph{確認した条件・記述（独自要約）}\mbox{}\par

\noindent\textbf{時点：} 令和元年度に指定された資料・制度。\par

\noindent\textbf{照合：} 原典18画像で否定・五肢・組合せ表を確認。後年値で置換しない。\par

\paragraph{問われる核心} 工業型から都市型の汚染と二次生成。

\paragraph{必要知識} 発生源、硫黄酸化物、粒子、光化学反応。

\paragraph{考え方・正答を選ぶ理由（査読待ち初稿）} 工業地域の燃料硫黄によるSOx、自動車のNOxとSPM、NOxと炭化水素の光反応から生じる光化学オキシダント。SOx/NOx/SPM/Oxの3。最後の二次生成物をCO2とすると前駆物質と日射の条件に対応しない。

\Needspace{4\baselineskip}\paragraph{選択肢の検討（査読待ち初稿）}\mbox{}\par

\noindent\textbf{1：不適切} SOx/Ox/SPM/CO2。式・定義・対応の独立判定と一致しない。\par

\noindent\textbf{2：不適切} NOx/Ox/CO2/SPM。式・定義・対応の独立判定と一致しない。\par

\noindent\textbf{3：適切・選択対象} SOx/NOx/SPM/Ox。判定と一致する。\par

\noindent\textbf{4：不適切} NOx/SOx/CO2/Ox。式・定義・対応の独立判定と一致しない。\par

\noindent\textbf{5：不適切} SOx/NOx/SPM/CO2。式・定義・対応の独立判定と一致しない。\par

\paragraph{再利用する考え方} 発生源、硫黄酸化物、粒子、光化学反応。 集計対象・時点・法の主体・選択方向を固定し、同型でも元の五肢へ戻す。

\paragraph{共通理論} \hyperref[theory:env:airhistory]{第\ref*{theory:env:airhistory}節}

\par 参考：\cite{env58_pm25}。

\par 難易度 2（校準初稿）、優先度未評価。

\section{緩和・適応とEco-DRRの共便益}

\subsection{令和元年度再試験 問5-1}
\label{q:r01-ret-basic-5-1}

\noindent\textbf{原典：} sources/R01\_retest/basic.pdf、PDF物理17頁。

\textbf{公式正答：} 1（single）。

\paragraph{問題の要旨} 緩和・適応とEco-DRRの共便益。

\Needspace{4\baselineskip}\paragraph{確認した条件・記述（独自要約）}\mbox{}\par

\noindent\textbf{時点：} 令和元年度再試験に指定された資料・制度。\par

\noindent\textbf{照合：} 原典18画像で否定・五肢・組合せ表を確認。後年値で置換しない。\par

\paragraph{問われる核心} 緩和・適応とEco-DRRの共便益。

\paragraph{必要知識} 緩和適応、地域特性、生態系の機能、共便益。

\paragraph{考え方・正答を選ぶ理由（査読待ち初稿）} 排出原因を減らす緩和、避けられない影響に備える適応。湿原の遊水や森林の保全機能を使うEco-DRRは生態系を活用する減災で、生物多様性保全も共便益となる。四語を合わせる1。人材が不要という意味ではなく、略称が指す手法の核心を読む。

\Needspace{4\baselineskip}\paragraph{選択肢の検討（査読待ち初稿）}\mbox{}\par

\noindent\textbf{1：適切・選択対象} 緩和/適応/生態系/生物多様性。判定と一致する。\par

\noindent\textbf{2：不適切} 削減/対応/生態系/地域資源。式・定義・対応の独立判定と一致しない。\par

\noindent\textbf{3：不適切} 緩和/適応/地域人材/地域資源。式・定義・対応の独立判定と一致しない。\par

\noindent\textbf{4：不適切} 緩和/対応/生態系/生物多様性。式・定義・対応の独立判定と一致しない。\par

\noindent\textbf{5：不適切} 削減/対応/地域人材/地域資源。式・定義・対応の独立判定と一致しない。\par

\paragraph{再利用する考え方} 緩和適応、地域特性、生態系の機能、共便益。 集計対象・時点・法の主体・選択方向を固定し、同型でも元の五肢へ戻す。

\paragraph{共通理論} \hyperref[theory:env:ecodrr]{第\ref*{theory:env:ecodrr}節}

\par 参考：\cite{env59_ecodrr}。

\par 難易度 2（校準初稿）、優先度未評価。

\section{部門別エネルギーの歴史的推移}

\subsection{令和2年度 問5-3}
\label{q:r02-reg-basic-5-3}

\noindent\textbf{原典：} sources/R02/basic.pdf、PDF物理20頁。

\textbf{公式正答：} 4（single）。

\par\textbf{検討未確定：} 原典が参照する版のピーク年2005・2002と全五グラフの原表は未取得。後年の統計遡及改訂と区別して追加照合。

\paragraph{問題の要旨} 部門別エネルギーの歴史的推移。

\Needspace{4\baselineskip}\paragraph{確認した条件・記述（独自要約）}\mbox{}\par

\noindent\textbf{時点：} 令和2年度に指定された資料・制度。\par

\noindent\textbf{照合：} 原典18画像で否定・五肢・組合せ表を確認。後年値で置換しない。\par

\paragraph{問われる核心} 部門別エネルギーの歴史的推移。

\paragraph{必要知識} 時系列、部門、GDP原単位、複数要因。

\paragraph{考え方・正答を選ぶ理由（査読待ち初稿）} 家庭は省エネ機器・節電・世帯人員・気温等の影響で、震災後に継続増加したという4が誤り。2019統計の要因分解にも年ごとの増減があり、一様な増加ではない。総量、原単位、世帯数、産業部門の個別と合計を分ける。

\Needspace{4\baselineskip}\paragraph{選択肢の検討（査読待ち初稿）}\mbox{}\par

\noindent\textbf{1：条件内で適切} 最終消費2005ピーク後減。条件内で適切な記述であり、不適切な記述を選ぶ対象ではない。\par

\noindent\textbf{2：条件内で適切} 産業業務全体は省エネで同程度。条件内で適切な記述であり、不適切な記述を選ぶ対象ではない。\par

\noindent\textbf{3：条件内で適切} GDP当たり日本は世界平均未満。条件内で適切な記述であり、不適切な記述を選ぶ対象ではない。\par

\noindent\textbf{4：不適切・選択対象} 家庭は震災後も継続増加。判定と一致する。\par

\noindent\textbf{5：条件内で適切} 旅客2002ピーク後燃費等で減。条件内で適切な記述であり、不適切な記述を選ぶ対象ではない。\par

\paragraph{再利用する考え方} 時系列、部門、GDP原単位、複数要因。 集計対象・時点・法の主体・選択方向を固定し、同型でも元の五肢へ戻す。

\paragraph{共通理論} \hyperref[theory:env:energytrends]{第\ref*{theory:env:energytrends}節}

\par 参考：\cite{env59_energy2019}。

\par 難易度 3（校準初稿）、優先度未評価。

\section{人物の業績とサイクロトロンの発明者}

\subsection{平成29年度 問5-6}
\label{q:h29-reg-basic-5-6}

\noindent\textbf{原典：} sources/H29/basic.pdf、PDF物理20頁。

\textbf{公式正答：} 5（single）。

\par\textbf{検討未確定：} Lawrenceの公式授賞説明を確認。他の四人の全原報は未取得で業績対応は一次史料追加照合を要する。

\paragraph{問題の要旨} 人物の業績とサイクロトロンの発明者。

\Needspace{4\baselineskip}\paragraph{確認した条件・記述（独自要約）}\mbox{}\par

\noindent\textbf{時点：} 平成29年度に指定された資料・制度。\par

\noindent\textbf{照合：} 原典18画像で否定・五肢・組合せ表を確認。後年値で置換しない。\par

\paragraph{問われる核心} 人物の業績とサイクロトロンの発明者。

\paragraph{必要知識} 電池、抗菌物質、放射線、遺伝、加速器。

\paragraph{考え方・正答を選ぶ理由（査読待ち初稿）} サイクロトロンはLawrenceらが開発した加速器で、Edisonという5が不適切。Voltaは定常電流を得る電堆、Flemingはリゾチームとペニシリン、RöntgenはX線、Mendelは形質の遺伝の規則性という対応。発明者名の誤りを、他の正しい加速原理の記述と混ぜて見落とさない。

\Needspace{4\baselineskip}\paragraph{選択肢の検討（査読待ち初稿）}\mbox{}\par

\noindent\textbf{1：条件内で適切} Volta電堆。条件内で適切な記述であり、不適切な記述を選ぶ対象ではない。\par

\noindent\textbf{2：条件内で適切} Flemingリゾチームとペニシリン。条件内で適切な記述であり、不適切な記述を選ぶ対象ではない。\par

\noindent\textbf{3：条件内で適切} RöntgenX線。条件内で適切な記述であり、不適切な記述を選ぶ対象ではない。\par

\noindent\textbf{4：条件内で適切} Mendel遺伝法則。条件内で適切な記述であり、不適切な記述を選ぶ対象ではない。\par

\noindent\textbf{5：不適切・選択対象} Edisonサイクロトロン。判定と一致する。\par

\paragraph{再利用する考え方} 電池、抗菌物質、放射線、遺伝、加速器。 集計対象・時点・法の主体・選択方向を固定し、同型でも元の五肢へ戻す。

\paragraph{共通理論} \hyperref[theory:env:historydiscover]{第\ref*{theory:env:historydiscover}節}

\par 参考：\cite{env59_lawrence}。

\par 難易度 3（校準初稿）、優先度未評価。

\section{家庭エネルギーの用途と消費範囲}

\subsection{平成29年度 問5-4}
\label{q:h29-reg-basic-5-4}

\noindent\textbf{原典：} sources/H29/basic.pdf、PDF物理19頁。

\textbf{公式正答：} 2（single）。

\par\textbf{検討未確定：} 原典「近年」の指定統計年が曖昧で、給湯三割・電気五割・待機五\%の全原表は未取得。比率を現行値の普遍定数にはしない。

\paragraph{問題の要旨} 家庭エネルギーの用途と消費範囲。

\Needspace{4\baselineskip}\paragraph{確認した条件・記述（独自要約）}\mbox{}\par

\noindent\textbf{時点：} 平成29年度に指定された資料・制度。\par

\noindent\textbf{照合：} 原典18画像で否定・五肢・組合せ表を確認。後年値で置換しない。\par

\paragraph{問われる核心} 家庭エネルギーの用途と消費範囲。

\paragraph{必要知識} 用途別消費、エネルギー種、機器効率、待機。

\paragraph{考え方・正答を選ぶ理由（査読待ち初稿）} 全国総和では冷房の年間消費は暖房より小さく、十倍とする2が不適切。電力だけと全燃料を含む家庭エネルギー、機器の効率と使用時間、総量と待機電力の比率を分ける。トップランナーは将来の効率目標の設定方式で、使用する一台ごとの達成義務と同じではない。

\Needspace{4\baselineskip}\paragraph{選択肢の検討（査読待ち初稿）}\mbox{}\par

\noindent\textbf{1：条件内で適切} 給湯が約三割。条件内で適切な記述であり、不適切な記述を選ぶ対象ではない。\par

\noindent\textbf{2：不適切・選択対象} 冷房が暖房の約十倍。判定と一致する。\par

\noindent\textbf{3：条件内で適切} 電気が約五割。条件内で適切な記述であり、不適切な記述を選ぶ対象ではない。\par

\noindent\textbf{4：条件内で適切} 冷蔵庫TV空調はトップランナー。条件内で適切な記述であり、不適切な記述を選ぶ対象ではない。\par

\noindent\textbf{5：条件内で適切} 電力約5\%が待機。条件内で適切な記述であり、不適切な記述を選ぶ対象ではない。\par

\paragraph{再利用する考え方} 用途別消費、エネルギー種、機器効率、待機。 集計対象・時点・法の主体・選択方向を固定し、同型でも元の五肢へ戻す。

\paragraph{共通理論} \hyperref[theory:env:home]{第\ref*{theory:env:home}節}

\par 参考：\cite{env59_energy2019}。

\par 難易度 3（校準初稿）、優先度未評価。

\section{戦前技術史の曖昧な比較と公式複数正答}

\subsection{令和2年度 問5-5}
\label{q:r02-reg-basic-5-5}

\noindent\textbf{原典：} sources/R02/basic.pdf、PDF物理21頁。

\textbf{公式正答：} 4|3（multiple\_accepted）。

\par\textbf{検討未確定：} 公式の複数受理4|3を保持。独立解4は受理対象だが全受理集合を独立に列挙したものではない。3の欧米との差という一般化、戦前統計利用の個別例を追加史料で査読。

\paragraph{問題の要旨} 戦前技術史の曖昧な比較と公式複数正答。

\Needspace{4\baselineskip}\paragraph{確認した条件・記述（独自要約）}\mbox{}\par

\noindent\textbf{時点：} 令和2年度に指定された資料・制度。\par

\noindent\textbf{照合：} 原典18画像で否定・五肢・組合せ表を確認。後年値で置換しない。\par

\paragraph{問われる核心} 戦前技術史の曖昧な比較と公式複数正答。

\paragraph{必要知識} 技術移転、学協会史、科学的管理と品質管理、公式訂正。

\paragraph{考え方・正答を選ぶ理由（査読待ち初稿）} 独立の主な判断は4。日本で統計的品質管理が組織的に普及した節目は日科技連による1950のDeming講義等の戦後であり、戦前のTaylor導入と同時に普及したというまとめは不適切。ただし戦前の個別統計手法の利用まで無かったとはしない。公式は誤解を招く記述があったとして当初4に加え3も正答とした。3の「欧米とは異なり」という一般化は学協会の時期・構成の違いを落とす。独立解4、受理正答3と4をそれぞれ残す。

\Needspace{4\baselineskip}\paragraph{選択肢の検討（査読待ち初稿）}\mbox{}\par

\noindent\textbf{1：条件内で適切} 江戸手工業が西洋技術と発展。条件内で適切な記述であり、不適切な記述を選ぶ対象ではない。\par

\noindent\textbf{2：条件内で適切} 標準化技術を輸入し低賃金競争。条件内で適切な記述であり、不適切な記述を選ぶ対象ではない。\par

\noindent\textbf{3：不適切・公式追加受理} 学協会は欧米と異なり高学歴出身者。欧米との比較表現に誤解を招く記述があるとして公式が3も受理。4との両受理を維持。\par

\noindent\textbf{4：不適切・選択対象} Taylorと同時に統計的品質管理導入。判定と一致する。\par

\noindent\textbf{5：条件内で適切} 技術官僚の地位向上運動。条件内で適切な記述であり、不適切な記述を選ぶ対象ではない。\par

\paragraph{再利用する考え方} 技術移転、学協会史、科学的管理と品質管理、公式訂正。 集計対象・時点・法の主体・選択方向を固定し、同型でも元の五肢へ戻す。

\paragraph{共通理論} \hyperref[theory:env:japanindustrial]{第\ref*{theory:env:japanindustrial}節}

\par 参考：\cite{env59_juse}。

\par 難易度 3（校準初稿）、優先度未評価。

\section{非エネルギー利用を除いた炭素量からCO2へ}

\subsection{令和元年度 問5-4}
\label{q:r01-reg-basic-5-4}

\noindent\textbf{原典：} sources/R01/basic.pdf、PDF物理21頁。

\textbf{公式正答：} 4（single）。

\paragraph{問題の要旨} 非エネルギー利用を除いた炭素量からCO2へ。

\Needspace{4\baselineskip}\paragraph{確認した条件・記述（独自要約）}\mbox{}\par

\noindent\textbf{燃料PJ：} 石炭5044、石油7831、天然ガス都市ガス4696。\par

\noindent\textbf{非燃焼石油：} 約1600PJ。\par

\noindent\textbf{係数t-C/TJ：} 石炭24、石油19、ガス14。\par

\noindent\textbf{換算：} C12、CO2 44、PJ=\allowbreak{}1000TJ。\par

\paragraph{問われる核心} 非エネルギー利用を除いた炭素量からCO2へ。

\paragraph{必要知識} システム境界、非エネルギー利用、PJ/TJ、分子量換算。

\paragraph{考え方・正答を選ぶ理由（査読待ち初稿）} 石油は燃焼に使わない1600PJを引いて6231PJ。炭素排出は[5044×24+6231×19+4696×14]×1000=\allowbreak{}305189000t-C。PJ→TJの1000倍を忘れない。CO2は44/12倍で約1119Mt-CO2なので1100の4。炭素換算値305Mtをそのまま選ばない。

\Needspace{4\baselineskip}\paragraph{選択肢の検討（査読待ち初稿）}\mbox{}\par

\noindent\textbf{1：不適切} 100Mt。式・定義・対応の独立判定と一致しない。\par

\noindent\textbf{2：不適切} 300Mt。式・定義・対応の独立判定と一致しない。\par

\noindent\textbf{3：不適切} 500Mt。式・定義・対応の独立判定と一致しない。\par

\noindent\textbf{4：適切・選択対象} 1100Mt。判定と一致する。\par

\noindent\textbf{5：不適切} 1600Mt。式・定義・対応の独立判定と一致しない。\par

\paragraph{再利用する考え方} システム境界、非エネルギー利用、PJ/TJ、分子量換算。 集計対象・時点・法の主体・選択方向を固定し、同型でも元の五肢へ戻す。

\paragraph{共通理論} \hyperref[theory:env:nationalco2]{第\ref*{theory:env:nationalco2}節}

\par 難易度 2（校準初稿）、優先度未評価。

\section{2017年の原油数量と輸入額比}

\subsection{平成30年度 問5-3}
\label{q:h30-reg-basic-5-3}

\noindent\textbf{原典：} sources/H30/basic.pdf、PDF物理19頁。

\textbf{公式正答：} 3（single）。

\par\textbf{検討未確定：} 2019白書の公式検索本文で中東約87\%を確認。2017財務省暦年の最大輸入先と金額12\%の全原表を取得していない。

\paragraph{問題の要旨} 2017年の原油数量と輸入額比。

\Needspace{4\baselineskip}\paragraph{確認した条件・記述（独自要約）}\mbox{}\par

\noindent\textbf{時点：} 平成30年度に指定された資料・制度。\par

\noindent\textbf{照合：} 原典18画像で否定・五肢・組合せ表を確認。後年値で置換しない。\par

\paragraph{問われる核心} 2017年の原油数量と輸入額比。

\paragraph{必要知識} 統計年度・品目・数量と金額・経路。

\paragraph{考え方・正答を選ぶ理由（査読待ち初稿）} 指定年の中東依存は約87\%、主要経路ホルムズ、最大相手サウジアラビア、石油・石油製品の総輸入額比約12\%で3。依存度は数量、最後は金額であり、他の化石燃料を含む鉱物性燃料の比率へ拡張しない。

\Needspace{4\baselineskip}\paragraph{選択肢の検討（査読待ち初稿）}\mbox{}\par

\noindent\textbf{1：不適切} 67/マラッカ/クウェート/12。式・定義・対応の独立判定と一致しない。\par

\noindent\textbf{2：不適切} 67/ホルムズ/サウジ/32。式・定義・対応の独立判定と一致しない。\par

\noindent\textbf{3：適切・選択対象} 87/ホルムズ/サウジ/12。判定と一致する。\par

\noindent\textbf{4：不適切} 87/マラッカ/クウェート/32。式・定義・対応の独立判定と一致しない。\par

\noindent\textbf{5：不適切} 87/ホルムズ/クウェート/12。式・定義・対応の独立判定と一致しない。\par

\paragraph{再利用する考え方} 統計年度・品目・数量と金額・経路。 集計対象・時点・法の主体・選択方向を固定し、同型でも元の五肢へ戻す。

\paragraph{共通理論} \hyperref[theory:env:oil2017]{第\ref*{theory:env:oil2017}節}

\par 参考：\cite{env59_energy2019}。

\par 難易度 3（校準初稿）、優先度未評価。

\section{パリ協定の削減目標は緩和に対応}

\subsection{平成29年度 問5-2}
\label{q:h29-reg-basic-5-2}

\noindent\textbf{原典：} sources/H29/basic.pdf、PDF物理18頁。

\textbf{公式正答：} 1（single）。

\par\textbf{検討未確定：} 取得UNFCCCはNDC五年と緩和の公式説明。協定2・5・6・13条の本文全体を通して照合する作業は未完了。

\paragraph{問題の要旨} パリ協定の削減目標は緩和に対応。

\Needspace{4\baselineskip}\paragraph{確認した条件・記述（独自要約）}\mbox{}\par

\noindent\textbf{時点：} 平成29年度に指定された資料・制度。\par

\noindent\textbf{照合：} 原典18画像で否定・五肢・組合せ表を確認。後年値で置換しない。\par

\paragraph{問われる核心} パリ協定の削減目標は緩和に対応。

\paragraph{必要知識} NDC、緩和と適応、温度目標、透明性、協力。

\paragraph{考え方・正答を選ぶ理由（査読待ち初稿）} 削減目標NDCの五年提出は緩和の取組であり、適応と呼ぶ1が不適切。適応は影響への備え。温度目標、透明性の情報・レビュー、協力的アプローチと市場、森林の吸収・排出抑制は条文の異なる要素を読む。パリ協定本文がJCMという日本の制度名を直接列挙したとの意味にはしない。

\Needspace{4\baselineskip}\paragraph{選択肢の検討（査読待ち初稿）}\mbox{}\par

\noindent\textbf{1：不適切・選択対象} 五年ごと削減目標で適応を推進。判定と一致する。\par

\noindent\textbf{2：条件内で適切} 2℃十分下方と1.5℃努力。条件内で適切な記述であり、不適切な記述を選ぶ対象ではない。\par

\noindent\textbf{3：条件内で適切} 情報の専門家レビュー。条件内で適切な記述であり、不適切な記述を選ぶ対象ではない。\par

\noindent\textbf{4：条件内で適切} 市場メカニズムの利用。条件内で適切な記述であり、不適切な記述を選ぶ対象ではない。\par

\noindent\textbf{5：条件内で適切} 森林減少・劣化抑制の支援。条件内で適切な記述であり、不適切な記述を選ぶ対象ではない。\par

\paragraph{再利用する考え方} NDC、緩和と適応、温度目標、透明性、協力。 集計対象・時点・法の主体・選択方向を固定し、同型でも元の五肢へ戻す。

\paragraph{共通理論} \hyperref[theory:env:paris]{第\ref*{theory:env:paris}節}

\par 参考：\cite{env59_ndc}。

\par 難易度 3（校準初稿）、優先度未評価。

\section{海洋プラスチックの定義・流出・循環}

\subsection{令和2年度 問5-1}
\label{q:r02-reg-basic-5-1}

\noindent\textbf{原典：} sources/R02/basic.pdf、PDF物理19頁。

\textbf{公式正答：} 2（single）。

\par\textbf{検討未確定：} 原典の2019年6月は政府発表5月31日と相違。厳密に日付も判定すると想定組合せが崩れるためAstra重点査読。2017国内排出・海外再生の全原表も追加照合。

\paragraph{問題の要旨} 海洋プラスチックの定義・流出・循環。

\Needspace{4\baselineskip}\paragraph{確認した条件・記述（独自要約）}\mbox{}\par

\noindent\textbf{時点：} 令和2年度に指定された資料・制度。\par

\noindent\textbf{照合：} 原典18画像で否定・五肢・組合せ表を確認。後年値で置換しない。\par

\paragraph{問われる核心} 海洋プラスチックの定義・流出・循環。

\paragraph{必要知識} 粒径、陸域流出、資源循環、代替資源、日付。

\paragraph{考え方・正答を選ぶ理由（査読待ち初稿）} 原典の主題判定はア正（概ね5mm以下）、イ誤（先進国だけ又は特に先進国由来へ一般化）、ウ誤（海外再生が半数以下という2017説明が違う）、エ正（3R+Renewable）、オ正（陸域の投棄や河川流出への対策）でOXXOOの2。ただし戦略の策定日は政府資料では2019年5月31日で、原典エの「6月」と一致しない。公式2を保持し、内容の判定と日付の疑義を分ける。

\Needspace{4\baselineskip}\paragraph{選択肢の検討（査読待ち初稿）}\mbox{}\par

\noindent\textbf{1：不適切} OOXOX。式・定義・対応の独立判定と一致しない。\par

\noindent\textbf{2：適切・選択対象} OXXOO。判定と一致する。\par

\noindent\textbf{3：不適切} OOOXX。式・定義・対応の独立判定と一致しない。\par

\noindent\textbf{4：不適切} XXOOO。式・定義・対応の独立判定と一致しない。\par

\noindent\textbf{5：不適切} XOXXO。式・定義・対応の独立判定と一致しない。\par

\paragraph{再利用する考え方} 粒径、陸域流出、資源循環、代替資源、日付。 集計対象・時点・法の主体・選択方向を固定し、同型でも元の五肢へ戻す。

\paragraph{共通理論} \hyperref[theory:env:plastics]{第\ref*{theory:env:plastics}節}

\par 参考：\cite{env59_plastic}。

\par 難易度 3（校準初稿）、優先度未評価。

\subsection{令和7年度 問5-2}
\label{q:r07-reg-basic-5-2}

\noindent\textbf{原典：} sources/R07/basic.pdf、PDF物理16頁。

\textbf{公式正答：} 3（single）。

\par\textbf{検討未確定：} 原典オは2019年6月と記すが政府戦略は2019年5月31日。公式選択肢3を保持し日付差を査読事項へ残す。2010推計と2017海外量の元表は追加照合。

\paragraph{問題の要旨} 海洋プラスチック五文は誤誤正誤正。

\Needspace{4\baselineskip}\paragraph{確認した条件・記述（独自要約）}\mbox{}\par

\noindent\textbf{原典：} 全五肢と否定方向・組合せ表を画像確認。\par

\noindent\textbf{資料時点：} 令和7年度の指定資料・統計年を保持し、後年値に置換しない。\par

\paragraph{問われる核心} 海洋プラスチック五文は誤誤正誤正。

\paragraph{必要知識} 微小プラスチック、流出と発生、国内循環と輸出、3R+Renewable。

\paragraph{考え方・正答を選ぶ理由（査読待ち初稿）} アは通常5mm未満を10mmとするため誤。イは2010流出推計を先進国主体とする一般化で誤。ウは陸のごみが川を通じ海へ流れるので陸域対策が有効で正。エは2017海外リサイクルを10万t未満とする桁が誤。オの3R+Renewableという内容は正として誤誤正誤正の3。ただし原典の策定月2019年6月は政府公表の5月31日と区別して留保する。

\Needspace{4\baselineskip}\paragraph{選択肢の検討（査読待ち初稿）}\mbox{}\par

\noindent\textbf{1：不適切} 正正誤誤正。独立に求めた組合せ 誤誤正誤正 と異なる。\par

\noindent\textbf{2：不適切} 正正誤正誤。独立に求めた組合せ 誤誤正誤正 と異なる。\par

\noindent\textbf{3：適切・選択対象} 誤誤正誤正。独立に求めた組合せ 誤誤正誤正 と一致する。\par

\noindent\textbf{4：不適切} 誤正正誤誤。独立に求めた組合せ 誤誤正誤正 と異なる。\par

\noindent\textbf{5：不適切} 誤誤正正正。独立に求めた組合せ 誤誤正誤正 と異なる。\par

\paragraph{再利用する考え方} 微小プラスチック、流出と発生、国内循環と輸出、3R+Renewable。 否定方向と資料年、分母と単位を確認して五肢へ戻る。

\paragraph{共通理論} \hyperref[theory:env:plastics]{第\ref*{theory:env:plastics}節}

\par 参考：\cite{env59_plastic}。

\par 難易度 3（校準初稿）、優先度未評価。

\section{一次供給の再エネ比率を境界に合わせる}

\subsection{令和元年度再試験 問5-4}
\label{q:r01-ret-basic-5-4}

\noindent\textbf{原典：} sources/R01\_retest/basic.pdf、PDF物理19頁。

\textbf{公式正答：} 3（single）。

\par\textbf{検討未確定：} 再試験の本文には内訳がないため通常試験の2017内訳を補助として使った。総合エネルギー統計指定版の原表は追加取得。

\paragraph{問題の要旨} 一次供給の再エネ比率を境界に合わせる。

\Needspace{4\baselineskip}\paragraph{確認した条件・記述（独自要約）}\mbox{}\par

\noindent\textbf{時点：} 令和元年度再試験に指定された資料・制度。\par

\noindent\textbf{照合：} 原典18画像で否定・五肢・組合せ表を確認。後年値で置換しない。\par

\paragraph{問われる核心} 一次供給の再エネ比率を境界に合わせる。

\paragraph{必要知識} 一次供給、発電比との違い、境界、合計と百分率。

\paragraph{考え方・正答を選ぶ理由（査読待ち初稿）} 2017指定統計と通常試験で与えた内訳を使うと、水力710+水力除く再エネ938+未活用596=\allowbreak{}2244PJ。全20095PJで割ると11.17\%で3。これは一次エネルギー供給比であり発電量構成ではない。問題が指定する未活用を含む境界を保持する。

\Needspace{4\baselineskip}\paragraph{選択肢の検討（査読待ち初稿）}\mbox{}\par

\noindent\textbf{1：不適切} 44\%。式・定義・対応の独立判定と一致しない。\par

\noindent\textbf{2：不適切} 22\%。式・定義・対応の独立判定と一致しない。\par

\noindent\textbf{3：適切・選択対象} 11\%。判定と一致する。\par

\noindent\textbf{4：不適切} 2\%。式・定義・対応の独立判定と一致しない。\par

\noindent\textbf{5：不適切} 0.5\%。式・定義・対応の独立判定と一致しない。\par

\paragraph{再利用する考え方} 一次供給、発電比との違い、境界、合計と百分率。 集計対象・時点・法の主体・選択方向を固定し、同型でも元の五肢へ戻す。

\paragraph{共通理論} \hyperref[theory:env:primaryshare]{第\ref*{theory:env:primaryshare}節}

\par 難易度 3（校準初稿）、優先度未評価。

\subsection{令和6年度 問5-4}
\label{q:r06-reg-basic-5-4}

\noindent\textbf{原典：} sources/R06/basic.pdf、PDF物理19頁。

\textbf{公式正答：} 4（single）。

\par\textbf{検討未確定：} 2022確報の数値表から指定三区分を実際に合算する追加再計算が必要。原表索引と再エネ説明を確認した範囲。

\paragraph{問題の要旨} 2022一次供給の再エネ・水力・未活用比。

\Needspace{4\baselineskip}\paragraph{確認した条件・記述（独自要約）}\mbox{}\par

\noindent\textbf{原典：} 全五肢と否定方向・組合せ表を画像確認。\par

\noindent\textbf{資料時点：} 令和6年度の指定資料・統計年を保持し、後年値に置換しない。\par

\paragraph{問われる核心} 2022一次供給の再エネ・水力・未活用比。

\paragraph{必要知識} 一次供給と発電量、揚水除外、区分集計。

\paragraph{考え方・正答を選ぶ理由（査読待ち初稿）} 一次エネルギー供給を分母とし、水力を除く再エネ、水力（揚水除外）、未活用エネルギーを合算すると、指定2022統計では約14\%として4を選ぶ。発電量における再エネ割合や自給率を代入しない。未活用には廃棄・放散される部分の有効利用等が含まれ、区分の重複計上を避ける。

\Needspace{4\baselineskip}\paragraph{選択肢の検討（査読待ち初稿）}\mbox{}\par

\noindent\textbf{1：不適切} 0.5\%。独立に求めた組合せ 14\% と異なる。\par

\noindent\textbf{2：不適切} 2\%。独立に求めた組合せ 14\% と異なる。\par

\noindent\textbf{3：不適切} 7\%。独立に求めた組合せ 14\% と異なる。\par

\noindent\textbf{4：適切・選択対象} 14\%。独立に求めた組合せ 14\% と一致する。\par

\noindent\textbf{5：不適切} 28\%。独立に求めた組合せ 14\% と異なる。\par

\paragraph{再利用する考え方} 一次供給と発電量、揚水除外、区分集計。 否定方向と資料年、分母と単位を確認して五肢へ戻る。

\paragraph{共通理論} \hyperref[theory:env:primaryshare]{第\ref*{theory:env:primaryshare}節}

\par 参考：\cite{env60_energy2022}。

\par 難易度 3（校準初稿）、優先度未評価。

\section{LCA・汚染者負担・拡大生産者責任}

\subsection{平成29年度 問5-1}
\label{q:h29-reg-basic-5-1}

\noindent\textbf{原典：} sources/H29/basic.pdf、PDF物理18頁。

\textbf{公式正答：} 1（single）。

\paragraph{問題の要旨} LCA・汚染者負担・拡大生産者責任。

\Needspace{4\baselineskip}\paragraph{確認した条件・記述（独自要約）}\mbox{}\par

\noindent\textbf{列順：} A=\allowbreak{}LCA、B=汚染者負担、C=拡大生産者、D=税務監査。\par

\noindent\textbf{O/X：} 正/誤。\par

\paragraph{問われる核心} LCA・汚染者負担・拡大生産者責任。

\paragraph{必要知識} ライフサイクル、汚染者負担原則、拡大生産者責任、監査。

\paragraph{考え方・正答を選ぶ理由（査読待ち初稿）} Aは全過程の影響評価でLCA。Bは汚染の防止・回復費用を汚染者が負担するPPP。Cは廃棄・再利用まで生産者の責任を広げるEPR。Dは税務の処理を会計士が見る説明で、環境監査の定義ではない。OOOXの1。

\Needspace{4\baselineskip}\paragraph{選択肢の検討（査読待ち初稿）}\mbox{}\par

\noindent\textbf{1：適切・選択対象} OOOX。判定と一致する。\par

\noindent\textbf{2：不適切} XXXO。式・定義・対応の独立判定と一致しない。\par

\noindent\textbf{3：不適切} XOOX。式・定義・対応の独立判定と一致しない。\par

\noindent\textbf{4：不適切} OOXO。式・定義・対応の独立判定と一致しない。\par

\noindent\textbf{5：不適切} OXXX。式・定義・対応の独立判定と一致しない。\par

\paragraph{再利用する考え方} ライフサイクル、汚染者負担原則、拡大生産者責任、監査。 集計対象・時点・法の主体・選択方向を固定し、同型でも元の五肢へ戻す。

\paragraph{共通理論} \hyperref[theory:env:principles]{第\ref*{theory:env:principles}節}

\par 参考：\cite{env59_epr}。

\par 難易度 2（校準初稿）、優先度未評価。

\section{評価リスクと知覚リスクの隔たり}

\subsection{令和元年度再試験 問5-6}
\label{q:r01-ret-basic-5-6}

\noindent\textbf{原典：} sources/R01\_retest/basic.pdf、PDF物理20頁。

\textbf{公式正答：} 3（single）。

\paragraph{問題の要旨} 評価リスクと知覚リスクの隔たり。

\Needspace{4\baselineskip}\paragraph{確認した条件・記述（独自要約）}\mbox{}\par

\noindent\textbf{時点：} 令和元年度再試験に指定された資料・制度。\par

\noindent\textbf{照合：} 原典18画像で否定・五肢・組合せ表を確認。後年値で置換しない。\par

\paragraph{問われる核心} 評価リスクと知覚リスクの隔たり。

\paragraph{必要知識} リスク評価・管理・対話、知覚、根拠と社会判断。

\paragraph{考え方・正答を選ぶ理由（査読待ち初稿）} 科学的評価と人の知覚は、経験・信頼・自発性・公平性等により隔たり得る。隔たりがないという3が誤り。相互の情報・意見交換によって理解し、評価の不確実性と管理の価値判断も説明する。レギュラトリーサイエンスは科学的根拠を政策等の判断へつなぐが、科学だけで価値判断を消すものではない。

\Needspace{4\baselineskip}\paragraph{選択肢の検討（査読待ち初稿）}\mbox{}\par

\noindent\textbf{1：条件内で適切} 科学評価と費用技術を含む管理。条件内で適切な記述であり、不適切な記述を選ぶ対象ではない。\par

\noindent\textbf{2：条件内で適切} 情報と意見の相互交換。条件内で適切な記述であり、不適切な記述を選ぶ対象ではない。\par

\noindent\textbf{3：不適切・選択対象} 評価と知覚に隔たりなし。判定と一致する。\par

\noindent\textbf{4：条件内で適切} 規制科学が信頼性と波及評価。条件内で適切な記述であり、不適切な記述を選ぶ対象ではない。\par

\noindent\textbf{5：条件内で適切} 規制と合意形成を支援。条件内で適切な記述であり、不適切な記述を選ぶ対象ではない。\par

\paragraph{再利用する考え方} リスク評価・管理・対話、知覚、根拠と社会判断。 集計対象・時点・法の主体・選択方向を固定し、同型でも元の五肢へ戻す。

\paragraph{共通理論} \hyperref[theory:env:risk]{第\ref*{theory:env:risk}節}

\par 参考：\cite{eth54_riskcomm}。

\par 難易度 2（校準初稿）、優先度未評価。

\subsection{令和4年度 問5-5}
\label{q:r04-reg-basic-5-5}

\noindent\textbf{原典：} sources/R04/basic.pdf、PDF物理20頁。

\textbf{公式正答：} 5（single）。

\paragraph{問題の要旨} リスクコミュニケーションでは評価過程を開示。

\Needspace{4\baselineskip}\paragraph{確認した条件・記述（独自要約）}\mbox{}\par

\noindent\textbf{原典：} 全五肢と否定方向・組合せ表を画像確認。\par

\noindent\textbf{資料時点：} 令和4年度の指定資料・統計年を保持し、後年値に置換しない。\par

\paragraph{問われる核心} リスクコミュニケーションでは評価過程を開示。

\paragraph{必要知識} 評価と管理、知覚と科学的評価、双方向と透明性。

\paragraph{考え方・正答を選ぶ理由（査読待ち初稿）} 5が不適切。評価値だけでなく仮定・不確実性・手順を示し、受け手との意見交換を通じて理解を形成する。管理上の費用や実行可能性を評価そのものと混同しない。

\Needspace{4\baselineskip}\paragraph{選択肢の検討（査読待ち初稿）}\mbox{}\par

\noindent\textbf{1：条件内で適切} 評価結果と費用・技術を考え管理。科学的評価と管理上の選択は異なる段階。\par

\noindent\textbf{2：条件内で適切} レギュラトリーサイエンスが合意形成を支援。評価を社会的判断へつなげる科学的基盤。\par

\noindent\textbf{3：条件内で適切} 個人・機関・集団の相互交換。一方向の通知だけとは異なる。\par

\noindent\textbf{4：条件内で適切} 評価リスクと認知リスクに隔たり。知識・価値観・信頼等により受け止めが変わる。\par

\noindent\textbf{5：不適切・選択対象} 混乱回避のため評価過程を開示しない。根拠・不確実性・評価過程を説明して対話する。混乱を理由に非開示を原則とはしない。\par

\paragraph{再利用する考え方} 評価と管理、知覚と科学的評価、双方向と透明性。 否定方向と資料年、分母と単位を確認して五肢へ戻る。

\paragraph{共通理論} \hyperref[theory:env:risk]{第\ref*{theory:env:risk}節}

\par 難易度 2（校準初稿）、優先度未評価。

\section{SDGsの対象は先進国も含む}

\subsection{平成30年度 問5-1}
\label{q:h30-reg-basic-5-1}

\noindent\textbf{原典：} sources/H30/basic.pdf、PDF物理18頁。

\textbf{公式正答：} 2（single）。

\paragraph{問題の要旨} SDGsの対象は先進国も含む。

\Needspace{4\baselineskip}\paragraph{確認した条件・記述（独自要約）}\mbox{}\par

\noindent\textbf{時点：} 平成30年度に指定された資料・制度。\par

\noindent\textbf{照合：} 原典18画像で否定・五肢・組合せ表を確認。後年値で置換しない。\par

\paragraph{問われる核心} SDGsの対象は先進国も含む。

\paragraph{必要知識} 普遍性、統合、ゴールとターゲット、実施体制。

\paragraph{考え方・正答を選ぶ理由（査読待ち初稿）} SDGsは普遍的で全ての国を対象とするので、途上国だけへ限定した2が誤り。17と169、環境・経済・社会の統合、各主体のパートナーシップを一組として読む。MDGsと同じ対象範囲へ戻さない。

\Needspace{4\baselineskip}\paragraph{選択肢の検討（査読待ち初稿）}\mbox{}\par

\noindent\textbf{1：条件内で適切} 2015の2030アジェンダ中核。条件内で適切な記述であり、不適切な記述を選ぶ対象ではない。\par

\noindent\textbf{2：不適切・選択対象} 途上国だけの目標。判定と一致する。\par

\noindent\textbf{3：条件内で適切} 17ゴール169ターゲット。条件内で適切な記述であり、不適切な記述を選ぶ対象ではない。\par

\noindent\textbf{4：条件内で適切} 多様主体の連携。条件内で適切な記述であり、不適切な記述を選ぶ対象ではない。\par

\noindent\textbf{5：条件内で適切} 日本2016実施指針。条件内で適切な記述であり、不適切な記述を選ぶ対象ではない。\par

\paragraph{再利用する考え方} 普遍性、統合、ゴールとターゲット、実施体制。 集計対象・時点・法の主体・選択方向を固定し、同型でも元の五肢へ戻す。

\paragraph{共通理論} \hyperref[theory:env:sdgs]{第\ref*{theory:env:sdgs}節}

\par 参考：\cite{env_sdgs_un,env_sdgs2016}。

\par 難易度 2（校準初稿）、優先度未評価。

\section{スマートメーターとHEMSの機能を分ける}

\subsection{平成30年度 問5-4}
\label{q:h30-reg-basic-5-4}

\noindent\textbf{原典：} sources/H30/basic.pdf、PDF物理19頁。

\textbf{公式正答：} 4（single）。

\paragraph{問題の要旨} スマートメーターとHEMSの機能を分ける。

\Needspace{4\baselineskip}\paragraph{確認した条件・記述（独自要約）}\mbox{}\par

\noindent\textbf{時点：} 平成30年度に指定された資料・制度。\par

\noindent\textbf{照合：} 原典18画像で否定・五肢・組合せ表を確認。後年値で置換しない。\par

\paragraph{問われる核心} スマートメーターとHEMSの機能を分ける。

\paragraph{必要知識} 計測、通信、家庭管理、地域と系統の範囲。

\paragraph{考え方・正答を選ぶ理由（査読待ち初稿）} 家庭内の機器を管理・制御する説明はHEMSに対応。スマートメーターは通信機能を持つ電力量計であり、同じ装置として家庭全体の最適制御を定義した4が不適切。計測情報は管理システムの入力に使えるが、計測と制御の主体を分ける。

\Needspace{4\baselineskip}\paragraph{選択肢の検討（査読待ち初稿）}\mbox{}\par

\noindent\textbf{1：条件内で適切} 系統需給のsmartgrid。条件内で適切な記述であり、不適切な記述を選ぶ対象ではない。\par

\noindent\textbf{2：条件内で適切} 地域の総合管理community。条件内で適切な記述であり、不適切な記述を選ぶ対象ではない。\par

\noindent\textbf{3：条件内で適切} 住宅機器を組み合わせるhouse。条件内で適切な記述であり、不適切な記述を選ぶ対象ではない。\par

\noindent\textbf{4：不適切・選択対象} メーターが家庭需給を最適制御。判定と一致する。\par

\noindent\textbf{5：条件内で適切} smartはかしこい。条件内で適切な記述であり、不適切な記述を選ぶ対象ではない。\par

\paragraph{再利用する考え方} 計測、通信、家庭管理、地域と系統の範囲。 集計対象・時点・法の主体・選択方向を固定し、同型でも元の五肢へ戻す。

\paragraph{共通理論} \hyperref[theory:env:smart]{第\ref*{theory:env:smart}節}

\par 参考：\cite{env59_smart}。

\par 難易度 2（校準初稿）、優先度未評価。

\section{応用が科学的発見より先行する例}

\subsection{令和元年度 問5-5}
\label{q:r01-reg-basic-5-5}

\noindent\textbf{原典：} sources/R01/basic.pdf、PDF物理21頁。

\textbf{公式正答：} 2（single）。

\par\textbf{検討未確定：} 他の四組の初出原報・開発年を全て取得したわけではない。土星の環は観測と環構造の解釈を分ける。

\paragraph{問題の要旨} 応用が科学的発見より先行する例。

\Needspace{4\baselineskip}\paragraph{確認した条件・記述（独自要約）}\mbox{}\par

\noindent\textbf{時点：} 令和元年度に指定された資料・制度。\par

\noindent\textbf{照合：} 原典18画像で否定・五肢・組合せ表を確認。後年値で置換しない。\par

\paragraph{問われる核心} 応用が科学的発見より先行する例。

\paragraph{必要知識} 経験的技術、機序解明、観測装置、時間順。

\paragraph{考え方・正答を選ぶ理由（査読待ち初稿）} Jennerの種痘は1796頃で、ウイルスの発見・概念形成は19世紀末。従って発見を種痘の前とした2が逆。実用的な方法が機序の科学的解明より先行する例であり、科学→技術の一方向の順序だけでは歴史を説明できない。

\Needspace{4\baselineskip}\paragraph{選択肢の検討（査読待ち初稿）}\mbox{}\par

\noindent\textbf{1：条件内で適切} 核分裂後に原子力発電。条件内で適切な記述であり、不適切な記述を選ぶ対象ではない。\par

\noindent\textbf{2：不適切・選択対象} ウイルス発見後に種痘。判定と一致する。\par

\noindent\textbf{3：条件内で適切} 望遠鏡後に土星環確認。条件内で適切な記述であり、不適切な記述を選ぶ対象ではない。\par

\noindent\textbf{4：条件内で適切} 量子力学後にトランジスタ。条件内で適切な記述であり、不適切な記述を選ぶ対象ではない。\par

\noindent\textbf{5：条件内で適切} 電磁波確認後にレーダー。条件内で適切な記述であり、不適切な記述を選ぶ対象ではない。\par

\paragraph{再利用する考え方} 経験的技術、機序解明、観測装置、時間順。 集計対象・時点・法の主体・選択方向を固定し、同型でも元の五肢へ戻す。

\paragraph{共通理論} \hyperref[theory:env:technologysequence]{第\ref*{theory:env:technologysequence}節}

\par 参考：\cite{env59_jenner}。

\par 難易度 3（校準初稿）、優先度未評価。

\subsection{令和5年度 問5-6}
\label{q:r05-reg-basic-5-6}

\noindent\textbf{原典：} sources/R05/basic.pdf、PDF物理18頁。

\textbf{公式正答：} 5（single）。

\par\textbf{検討未確定：} 土星環の解釈・レーダー・発電開始の基準年は各一次史料へ追加照合。

\paragraph{問題の要旨} ワクチンの実践はウイルス発見に先行。

\Needspace{4\baselineskip}\paragraph{確認した条件・記述（独自要約）}\mbox{}\par

\noindent\textbf{原典：} 全五肢と否定方向・組合せ表を画像確認。\par

\noindent\textbf{資料時点：} 令和5年度の指定資料・統計年を保持し、後年値に置換しない。\par

\paragraph{問われる核心} ワクチンの実践はウイルス発見に先行。

\paragraph{必要知識} 観測・理論・応用の双方向性、発見と実践の順序。

\paragraph{考え方・正答を選ぶ理由（査読待ち初稿）} 5が不適切。科学が技術へ一方向に進むとは限らず、種痘の経験的実践はウイルス概念の成立より早い。望遠鏡の発明と環の正しい解釈の時期も区別する。

\Needspace{4\baselineskip}\paragraph{選択肢の検討（査読待ち初稿）}\mbox{}\par

\noindent\textbf{1：条件内で適切} 望遠鏡後に土星の環確認。観測手段が天体の新しい知識を生む例。\par

\noindent\textbf{2：条件内で適切} 量子力学後にトランジスター。20世紀前半の理論と1947の素子発明の順。\par

\noindent\textbf{3：条件内で適切} 電磁波確認後にレーダー。Hertzの確認とその後の電波応用の順。\par

\noindent\textbf{4：条件内で適切} 核分裂発見後に原子力発電。1938核分裂の後に発電利用へ進む。\par

\noindent\textbf{5：不適切・選択対象} ウイルス発見後にワクチン接種開始。Jennerの種痘1796はウイルス研究の19世紀末より先。病原体の詳細理解前にも技術が発達した。\par

\paragraph{再利用する考え方} 観測・理論・応用の双方向性、発見と実践の順序。 否定方向と資料年、分母と単位を確認して五肢へ戻る。

\paragraph{共通理論} \hyperref[theory:env:technologysequence]{第\ref*{theory:env:technologysequence}節}

\par 参考：\cite{env59_jenner}。

\par 難易度 3（校準初稿）、優先度未評価。

\section{AR6観測温暖化は約1.1℃であって3℃ではない}

\subsection{令和4年度 問5-1}
\label{q:r04-reg-basic-5-1}

\noindent\textbf{原典：} sources/R04/basic.pdf、PDF物理17頁。

\textbf{公式正答：} 2（single）。

\par\textbf{検討未確定：} WGI温度箇所を確認。WGII脆弱性とWGIII排出量の各原SPM全段落照合は未完。

\paragraph{問題の要旨} AR6観測温暖化は約1.1℃であって3℃ではない。

\Needspace{4\baselineskip}\paragraph{確認した条件・記述（独自要約）}\mbox{}\par

\noindent\textbf{原典：} 全五肢と否定方向・組合せ表を画像確認。\par

\noindent\textbf{資料時点：} 令和4年度の指定資料・統計年を保持し、後年値に置換しない。\par

\paragraph{問われる核心} AR6観測温暖化は約1.1℃であって3℃ではない。

\paragraph{必要知識} 期間平均の観測差と将来予測、気温と排出量。

\paragraph{考え方・正答を選ぶ理由（査読待ち初稿）} 2が不適切。WGIの2011〜2020対1850〜1900は1.09℃（評価幅0.95〜1.20℃）であり、将来シナリオの温度上昇や特定地域の上昇と区別する。人間影響、地域別脆弱性、排出量は各作業部会の対象を分ける。

\Needspace{4\baselineskip}\paragraph{選択肢の検討（査読待ち初稿）}\mbox{}\par

\noindent\textbf{1：条件内で適切} 人間影響による温暖化は疑う余地がない。WGIの結論は人間影響について明確な評価を示す。\par

\noindent\textbf{2：不適切・選択対象} 2011〜2020は1850〜1900より約3℃高い。指定二期間の世界平均地表温度差は1.09℃、概数1.1℃。約3℃は観測済みの差ではない。\par

\noindent\textbf{3：条件内で適切} 極端現象や酸性化・海面上昇等の影響。物理・生態系影響を一種類に限定しない。\par

\noindent\textbf{4：条件内で適切} 脆弱性は地域内外で異なる。発展形態や曝露、適応能力が脆弱性を変える。\par

\noindent\textbf{5：条件内で適切} 2010〜2019の平均排出量は過去の十年より高い。WGIIIの評価対象は年間平均の正味人為GHG排出量。濃度や年ごとの増加率とは違う。\par

\paragraph{再利用する考え方} 期間平均の観測差と将来予測、気温と排出量。 否定方向と資料年、分母と単位を確認して五肢へ戻る。

\paragraph{共通理論} \hyperref[theory:env:ar6]{第\ref*{theory:env:ar6}節}

\par 参考：\cite{env60_ar6}。

\par 難易度 3（校準初稿）、優先度未評価。

\section{里地里山の危機は利用縮小・管理低下}

\subsection{令和5年度 問5-1}
\label{q:r05-reg-basic-5-1}

\noindent\textbf{原典：} sources/R05/basic.pdf、PDF物理16頁。

\textbf{公式正答：} 3（single）。

\par\textbf{検討未確定：} 指定国家戦略の固有種比率・開発圧力・分布事例の全原文は追加確認対象。

\paragraph{問題の要旨} 里地里山の危機は利用縮小・管理低下。

\Needspace{4\baselineskip}\paragraph{確認した条件・記述（独自要約）}\mbox{}\par

\noindent\textbf{原典：} 全五肢と否定方向・組合せ表を画像確認。\par

\noindent\textbf{資料時点：} 令和5年度の指定資料・統計年を保持し、後年値に置換しない。\par

\paragraph{問われる核心} 里地里山の危機は利用縮小・管理低下。

\paragraph{必要知識} 生物多様性の危機、固有種、二次的自然と人の管理。

\paragraph{考え方・正答を選ぶ理由（査読待ち初稿）} 3が不適切。里山は人間の利用で維持されてきた生態系も含むため、手入れや利用の停止が二次的自然を変える。開発による第一の危機と利用縮小による第二の危機を分け、外来種・化学物質・気候変動の要因も別に読む。

\Needspace{4\baselineskip}\paragraph{選択肢の検討（査読待ち初稿）}\mbox{}\par

\noindent\textbf{1：条件内で適切} 爬虫類約6割・両生類約8割が固有種。固有種割合を全動植物へ一般化しない。\par

\noindent\textbf{2：条件内で適切} 近年大規模開発の圧力は低下。過去の減少と近年の開発圧力の変化を別に記す。\par

\noindent\textbf{3：不適切・選択対象} 里山の利用拡大で生息地減少。里地里山は利用・管理の縮小による環境変化が主要な危機。拡大という原因説明が逆。\par

\noindent\textbf{4：条件内で適切} 他地域からの導入種が在来へ影響。国外だけでなく国内他地域からの移動でも影響し得る。\par

\noindent\textbf{5：条件内で適切} タケや南方系チョウの北上。気候変化と分布の変化の事例。全種が一様に北上するという主張ではない。\par

\paragraph{再利用する考え方} 生物多様性の危機、固有種、二次的自然と人の管理。 否定方向と資料年、分母と単位を確認して五肢へ戻る。

\paragraph{共通理論} \hyperref[theory:env:biodiversity2023]{第\ref*{theory:env:biodiversity2023}節}

\par 参考：\cite{env60_biodiversity}。

\par 難易度 3（校準初稿）、優先度未評価。

\section{2024には藻場吸収量を国連へ報告済み}

\subsection{令和7年度 問5-1}
\label{q:r07-reg-basic-5-1}

\noindent\textbf{原典：} sources/R07/basic.pdf、PDF物理15頁。

\textbf{公式正答：} 5（single）。

\par\textbf{検討未確定：} 報告済みの藻場範囲は確認。他4肢の当時法令要件とZEH・ZEB目標全文の照合は追加対象。

\paragraph{問題の要旨} 2024には藻場吸収量を国連へ報告済み。

\Needspace{4\baselineskip}\paragraph{確認した条件・記述（独自要約）}\mbox{}\par

\noindent\textbf{原典：} 全五肢と否定方向・組合せ表を画像確認。\par

\noindent\textbf{資料時点：} 令和7年度の指定資料・統計年を保持し、後年値に置換しない。\par

\paragraph{問われる核心} 2024には藻場吸収量を国連へ報告済み。

\paragraph{必要知識} ブルーカーボン、評価年と提出年、濃度と排出吸収。

\paragraph{考え方・正答を選ぶ理由（査読待ち初稿）} 5が不適切。令和6年白書の時点では、2022年度排出・吸収量の2024年提出インベントリに海草・海藻藻場の評価が含まれた。海の全炭素吸収を漏れなく算定できたという意味ではない。濃度観測・事業者排出報告・吸収インベントリを別に整理する。

\Needspace{4\baselineskip}\paragraph{選択肢の検討（査読待ち初稿）}\mbox{}\par

\noindent\textbf{1：条件内で適切} 衛星で全球GHG濃度を継続観測。衛星は広域の濃度観測を行う。排出量を各地点で直接計量するという説明ではない。\par

\noindent\textbf{2：条件内で適切} 一定量以上排出する事業者の年度報告。算定・報告・公表制度は排出量の把握と公表を扱う。\par

\noindent\textbf{3：条件内で適切} HFCにも強い温室効果があり回収処理。オゾン層を破壊しないことと温室効果が小さいことは同義でない。\par

\noindent\textbf{4：条件内で適切} 2030以降新築にZEH・ZEB水準を目指す。省エネルギー性能の水準と、全建物が毎時ゼロであることは違う。\par

\noindent\textbf{5：不適切・選択対象} ブルーカーボン吸収を国連に未報告。2024年4月に海草・海藻藻場の吸収量を合わせて国連へ報告したため未報告は不適切。\par

\paragraph{再利用する考え方} ブルーカーボン、評価年と提出年、濃度と排出吸収。 否定方向と資料年、分母と単位を確認して五肢へ戻る。

\paragraph{共通理論} \hyperref[theory:env:blue]{第\ref*{theory:env:blue}節}

\par 参考：\cite{env60_blue}。

\par 難易度 3（校準初稿）、優先度未評価。

\section{代替フロンのオゾン保護と温暖化対策を区別}

\subsection{令和3年度 問5-1}
\label{q:r03-reg-basic-5-1}

\noindent\textbf{原典：} sources/R03/basic.pdf、PDF物理21頁。

\textbf{公式正答：} 2（single）。

\par\textbf{検討未確定：} TCFDの2020賛同200以上の数、FSB構成、各制度の当時の定義全文は追加照合対象。

\paragraph{問題の要旨} 代替フロンのオゾン保護と温暖化対策を区別。

\Needspace{4\baselineskip}\paragraph{確認した条件・記述（独自要約）}\mbox{}\par

\noindent\textbf{原典：} 全五肢と否定方向・組合せ表を画像確認。\par

\noindent\textbf{資料時点：} 令和3年度の指定資料・統計年を保持し、後年値に置換しない。\par

\paragraph{問われる核心} 代替フロンのオゾン保護と温暖化対策を区別。

\paragraph{必要知識} オゾン層破壊係数と温暖化係数、電力と一次エネルギー、目標と開示。

\paragraph{考え方・正答を選ぶ理由（査読待ち初稿）} 2が不適切。CFC・HCFCのオゾン層破壊とHFCの地球温暖化係数は別の問題である。RE100、TCFD、自治体目標、ZEHはそれぞれ電力調達・情報開示・地域目標・建築性能という異なる対象に働く。

\Needspace{4\baselineskip}\paragraph{選択肢の検討（査読待ち初稿）}\mbox{}\par

\noindent\textbf{1：条件内で適切} RE100は使用電力100\%再エネを目指す。目標は企業活動に使用する電力の再エネ化。全エネルギーを必ず電力だけで賄うという定義ではない。\par

\noindent\textbf{2：不適切・選択対象} 代替フロン転換で温暖化対策も十分。HFC等にはオゾン層を破壊しなくても強い温室効果を持つものがあり、転換だけで十分とはいえない。\par

\noindent\textbf{3：条件内で適切} TCFDは投資判断に資する開示を促す。気候関連のリスク・機会を財務情報として示す取組。賛同機関数は2020時点の値として扱う。\par

\noindent\textbf{4：条件内で適切} 自治体の2050実質ゼロ表明。ゼロカーボンシティは自治体の実質ゼロ目標表明を示す。\par

\noindent\textbf{5：条件内で適切} ZEH・ZEH-Mは年間一次エネルギー収支ゼロを目指す。断熱・高効率設備で需要を下げ再エネを導入する。設備導入だけ、又は瞬間ごとのゼロとは限らない。\par

\paragraph{再利用する考え方} オゾン層破壊係数と温暖化係数、電力と一次エネルギー、目標と開示。 否定方向と資料年、分母と単位を確認して五肢へ戻る。

\paragraph{共通理論} \hyperref[theory:env:climateaction]{第\ref*{theory:env:climateaction}節}

\par 参考：\cite{env60_blue}。

\par 難易度 3（校準初稿）、優先度未評価。

\section{太陽電池の国内生産品比率は毎年拡大ではない}

\subsection{令和5年度 問5-3}
\label{q:r05-reg-basic-5-3}

\noindent\textbf{原典：} sources/R05/basic.pdf、PDF物理17頁。

\textbf{公式正答：} 1（single）。

\par\textbf{検討未確定：} 2009〜2020国内生産比率の原表、当時の原油・LNG・風力・地熱数値を追加照合。

\paragraph{問題の要旨} 太陽電池の国内生産品比率は毎年拡大ではない。

\Needspace{4\baselineskip}\paragraph{確認した条件・記述（独自要約）}\mbox{}\par

\noindent\textbf{原典：} 全五肢と否定方向・組合せ表を画像確認。\par

\noindent\textbf{資料時点：} 令和5年度の指定資料・統計年を保持し、後年値に置換しない。\par

\paragraph{問われる核心} 太陽電池の国内生産品比率は毎年拡大ではない。

\paragraph{必要知識} 設備導入と製造比率、資源量と導入、原油とLNG。

\paragraph{考え方・正答を選ぶ理由（査読待ち初稿）} 1が不適切。太陽光発電の国内導入が増えても、国内出荷に占める国内生産品比率が同時に増えるとは限らない。原油・LNG・風力・地熱は輸入構成、設置容量、資源量を別々に評価する。

\Needspace{4\baselineskip}\paragraph{選択肢の検討（査読待ち初稿）}\mbox{}\par

\noindent\textbf{1：不適切・選択対象} 導入量も国内生産品比率も2009〜2020毎年拡大。導入設備の増加と国内製品比率は独立。輸入増加に伴う国内比率低下を落とし、毎年拡大とするのが不適切。\par

\noindent\textbf{2：条件内で適切} 2020原油中東90\%以上でサウジ・UAE上位。原油の数量ベースの依存を読む。\par

\noindent\textbf{3：条件内で適切} LNGの80\%以上が中東以外で豪州最大。原油とは異なる調達構成。LNGを原油と同じ中東依存としない。\par

\noindent\textbf{4：条件内で適切} 2020末風力400万kW超と系統拡大。導入容量の単位kWと発電量kWhは違う。\par

\noindent\textbf{5：条件内で適切} 地熱資源量は日本世界3位。資源量順位と実際の導入容量順位は別。\par

\paragraph{再利用する考え方} 設備導入と製造比率、資源量と導入、原油とLNG。 否定方向と資料年、分母と単位を確認して五肢へ戻る。

\paragraph{共通理論} \hyperref[theory:env:energy2020]{第\ref*{theory:env:energy2020}節}

\par 難易度 3（校準初稿）、優先度未評価。

\section{第7次計画は技術の選択肢を限定しない}

\subsection{令和7年度 問5-3}
\label{q:r07-reg-basic-5-3}

\noindent\textbf{原典：} sources/R07/basic.pdf、PDF物理17頁。

\textbf{公式正答：} 4（single）。

\par\textbf{検討未確定：} 2のS省略は不完全な要約として注意。計画本文と目標の各原段落全体を査読し、単に安全性を不要とする文へ読み替えない。

\paragraph{問題の要旨} 第7次計画は技術の選択肢を限定しない。

\Needspace{4\baselineskip}\paragraph{確認した条件・記述（独自要約）}\mbox{}\par

\noindent\textbf{原典：} 全五肢と否定方向・組合せ表を画像確認。\par

\noindent\textbf{資料時点：} 令和7年度の指定資料・統計年を保持し、後年値に置換しない。\par

\paragraph{問われる核心} 第7次計画は技術の選択肢を限定しない。

\paragraph{必要知識} S+3E、技術選択肢、需要予測、基準年と削減目標。

\paragraph{考え方・正答を選ぶ理由（査読待ち初稿）} 主判断4。脱炭素へ向け使える技術を幅広く活用する方針を、あらかじめ選択肢を限定する方針へ反転している。2は安全性Sが大前提であることを補って読む。電力需要の予測、削減目標、政策の価値基準は別の論点である。

\Needspace{4\baselineskip}\paragraph{選択肢の検討（査読待ち初稿）}\mbox{}\par

\noindent\textbf{1：条件内で適切} データセンター等で電力需要増の可能性。効率改善を見込んでも需要を押し上げる要因がある。\par

\noindent\textbf{2：条件内で適切} 三つのEの最適バランスが基本的視点。安全性Sを大前提に3Eのバランスを追求する。原典の三つのEだけの短縮を完全定義としない。\par

\noindent\textbf{3：条件内で適切} 2035/2040に2013比60\%/73\%削減。目標の基準年は2013年度。前年度比ではない。\par

\noindent\textbf{4：不適切・選択対象} 利用技術を選別し選択肢を限定。不確実性に対応し使える技術の選択肢を幅広く追求する方向と逆で不適切。\par

\noindent\textbf{5：条件内で適切} 脱炭素投資を経済成長へつなげる。供給構造転換と産業投資・成長を一体で扱う。\par

\paragraph{再利用する考え方} S+3E、技術選択肢、需要予測、基準年と削減目標。 否定方向と資料年、分母と単位を確認して五肢へ戻る。

\paragraph{共通理論} \hyperref[theory:env:energyplan7]{第\ref*{theory:env:energyplan7}節}

\par 参考：\cite{env60_plan7}。

\par 難易度 3（校準初稿）、優先度未評価。

\section{2022家庭CO2の積み上げ棒は電気・都市ガス・LP・灯油}

\subsection{令和6年度 問5-3}
\label{q:r06-reg-basic-5-3}

\noindent\textbf{原典：} sources/R06/basic.pdf、PDF物理19頁。

\textbf{公式正答：} 2（single）。

\paragraph{問題の要旨} 2022家庭CO2の積み上げ棒は電気・都市ガス・LP・灯油。

\Needspace{4\baselineskip}\paragraph{確認した条件・記述（独自要約）}\mbox{}\par

\noindent\textbf{指定：} 令和4年度家庭CO2確報、令和6年3月。\par

\noindent\textbf{図：} 2017〜2022年度、A〜Dは下からの順。\par

\paragraph{問われる核心} 2022家庭CO2の積み上げ棒は電気・都市ガス・LP・灯油。

\paragraph{必要知識} 積み上げ棒の凡例、排出係数、世帯平均と燃料量。

\paragraph{考え方・正答を選ぶ理由（査読待ち初稿）} 原典R4棒のA=\allowbreak{}1.74、B=.38、C=.15、D=.32（t-CO2/世帯・年）。確報の区分へ対応させA電気、B都市ガス、C LPガス、D灯油で2。合計2.59。2017等の棒の値を2022へ置換しない。電気は発電段階の排出を係数で世帯へ帰属させ、各家庭が煙突から直接排出した量だけではない。

\Needspace{4\baselineskip}\paragraph{選択肢の検討（査読待ち初稿）}\mbox{}\par

\noindent\textbf{1：不適切} 電気/都市ガス/灯油/LP。独立に求めた組合せ 電気/都市ガス/LP/灯油 と異なる。\par

\noindent\textbf{2：適切・選択対象} 電気/都市ガス/LP/灯油。独立に求めた組合せ 電気/都市ガス/LP/灯油 と一致する。\par

\noindent\textbf{3：不適切} 都市ガス/電気/LP/灯油。独立に求めた組合せ 電気/都市ガス/LP/灯油 と異なる。\par

\noindent\textbf{4：不適切} 都市ガス/電気/灯油/LP。独立に求めた組合せ 電気/都市ガス/LP/灯油 と異なる。\par

\noindent\textbf{5：不適切} 灯油/電気/都市ガス/LP。独立に求めた組合せ 電気/都市ガス/LP/灯油 と異なる。\par

\paragraph{再利用する考え方} 積み上げ棒の凡例、排出係数、世帯平均と燃料量。 否定方向と資料年、分母と単位を確認して五肢へ戻る。

\paragraph{共通理論} \hyperref[theory:env:house2022]{第\ref*{theory:env:house2022}節}

\par 参考：\cite{env60_house2022}。

\par 難易度 2（校準初稿）、優先度未評価。

\section{液体水素の温度・体積・質量発熱と還元}

\subsection{令和4年度 問5-4}
\label{q:r04-reg-basic-5-4}

\noindent\textbf{原典：} sources/R04/basic.pdf、PDF物理19頁。

\textbf{公式正答：} 5（single）。

\par\textbf{検討未確定：} 鉄鉱石全般について温度・相を省いた吸熱断定を避ける。液化値と発熱量の一次物性表は追加照合。

\paragraph{問題の要旨} 液体水素の温度・体積・質量発熱と還元。

\Needspace{4\baselineskip}\paragraph{確認した条件・記述（独自要約）}\mbox{}\par

\noindent\textbf{原典：} 全五肢と否定方向・組合せ表を画像確認。\par

\noindent\textbf{資料時点：} 令和4年度の指定資料・統計年を保持し、後年値に置換しない。\par

\paragraph{問われる核心} 液体水素の温度・体積・質量発熱と還元。

\paragraph{必要知識} 沸点、気液密度、質量と体積の発熱量、反応熱の成立条件。

\paragraph{考え方・正答を選ぶ理由（査読待ち初稿）} 1気圧で水素の液化温度は約−253℃、常温気体と比べ液体体積は約1/800。ガソリンより大きいのは単位質量当たりの発熱量で、単位体積ではない。水素を用いた鉄鉱石還元の工程は熱の供給を要する吸熱側の技術として扱い5。具体的反応エンタルピーは酸化物の種類・温度・生成水の相を指定して求める。水素製造まで含む排出が必ずゼロとの意味ではない。

\Needspace{4\baselineskip}\paragraph{選択肢の検討（査読待ち初稿）}\mbox{}\par

\noindent\textbf{1：不適切} -162/1:600/重量/吸熱。独立に求めた組合せ -253/1:800/重量/吸熱 と異なる。\par

\noindent\textbf{2：不適切} -162/1:800/重量/発熱。独立に求めた組合せ -253/1:800/重量/吸熱 と異なる。\par

\noindent\textbf{3：不適切} -253/1:600/体積/発熱。独立に求めた組合せ -253/1:800/重量/吸熱 と異なる。\par

\noindent\textbf{4：不適切} -253/1:800/体積/発熱。独立に求めた組合せ -253/1:800/重量/吸熱 と異なる。\par

\noindent\textbf{5：適切・選択対象} -253/1:800/重量/吸熱。独立に求めた組合せ -253/1:800/重量/吸熱 と一致する。\par

\paragraph{再利用する考え方} 沸点、気液密度、質量と体積の発熱量、反応熱の成立条件。 否定方向と資料年、分母と単位を確認して五肢へ戻る。

\paragraph{共通理論} \hyperref[theory:env:hydrogen]{第\ref*{theory:env:hydrogen}節}

\par 参考：\cite{env60_hydrogen}。

\par 難易度 3（校準初稿）、優先度未評価。

\section{1949年の最初の日本人ノーベル賞は湯川}

\subsection{令和7年度 問5-5}
\label{q:r07-reg-basic-5-5}

\noindent\textbf{原典：} sources/R07/basic.pdf、PDF物理18頁。

\textbf{公式正答：} 2（single）。

\par\textbf{検討未確定：} 特許庁の出願日を索引確認。Nobelの各受賞一次ページ、ウィーン公式参加、Ford工程資料は追加照合。

\paragraph{問題の要旨} 1949年の最初の日本人ノーベル賞は湯川。

\Needspace{4\baselineskip}\paragraph{確認した条件・記述（独自要約）}\mbox{}\par

\noindent\textbf{原典：} 全五肢と否定方向・組合せ表を画像確認。\par

\noindent\textbf{資料時点：} 令和7年度の指定資料・統計年を保持し、後年値に置換しない。\par

\paragraph{問われる核心} 1949年の最初の日本人ノーベル賞は湯川。

\paragraph{必要知識} 受賞者・年・業績、政府公式参加、出願と登録。

\paragraph{考え方・正答を選ぶ理由（査読待ち初稿）} 2が不適切。人物・受賞年・対象業績が一組でなければならず、1949/中間子は湯川、1965/量子電気力学は朝永。4は出願日の記述であり特許の登録日を問うものではない。

\Needspace{4\baselineskip}\paragraph{選択肢の検討（査読待ち初稿）}\mbox{}\par

\noindent\textbf{1：条件内で適切} Haberはアンモニア合成で化学賞。空中窒素からのアンモニア合成に対する1918化学賞。\par

\noindent\textbf{2：不適切・選択対象} 朝永が1949に中間子予言で最初の物理賞。1949の中間子予言による賞は湯川秀樹。朝永振一郎は1965に量子電気力学で受賞。\par

\noindent\textbf{3：条件内で適切} 明治政府初の公式参加は1873ウィーン。政府公式参加と民間・幕府の先行参加を区別。\par

\noindent\textbf{4：条件内で適切} 堀田の錆止塗料は1885年7月1日出願。特許庁の制度史は出願日7月1日を示す。登録日と入れ替えない。\par

\noindent\textbf{5：条件内で適切} Fordの流れ作業で大量低コスト生産。工程を連続化し作業を分担する生産方式。\par

\paragraph{再利用する考え方} 受賞者・年・業績、政府公式参加、出願と登録。 否定方向と資料年、分母と単位を確認して五肢へ戻る。

\paragraph{共通理論} \hyperref[theory:env:nobeljapan]{第\ref*{theory:env:nobeljapan}節}

\par 参考：\cite{env60_patent,env60_yukawa}。

\par 難易度 3（校準初稿）、優先度未評価。

\section{2019年度原油依存と輸入金額の分母}

\subsection{令和4年度 問5-3}
\label{q:r04-reg-basic-5-3}

\noindent\textbf{原典：} sources/R04/basic.pdf、PDF物理18頁。

\textbf{公式正答：} 1（single）。

\par\textbf{検討未確定：} エネルギー白書2021の2019年度元表・輸入額比の原図は追加確認対象。

\paragraph{問題の要旨} 2019年度原油依存と輸入金額の分母。

\Needspace{4\baselineskip}\paragraph{確認した条件・記述（独自要約）}\mbox{}\par

\noindent\textbf{原典：} 全五肢と否定方向・組合せ表を画像確認。\par

\noindent\textbf{資料時点：} 令和4年度の指定資料・統計年を保持し、後年値に置換しない。\par

\paragraph{問われる核心} 2019年度原油依存と輸入金額の分母。

\paragraph{必要知識} 年度と白書年、数量依存度と金額比、輸送経路。

\paragraph{考え方・正答を選ぶ理由（査読待ち初稿）} 指定白書2021では2019年度の原油中東依存度を概数90\%、経由する海峡をホルムズ、最大相手国をサウジアラビア、石油・石油製品輸入額の総輸入額比を概数10\%として1。数量に占める地域割合と金額に占める品目割合は分母も単位も違う。

\Needspace{4\baselineskip}\paragraph{選択肢の検討（査読待ち初稿）}\mbox{}\par

\noindent\textbf{1：適切・選択対象} 90/ホルムズ/サウジ/10。独立に求めた組合せ 90/ホルムズ/サウジ/10 と一致する。\par

\noindent\textbf{2：不適切} 90/マラッカ/クウェート/32。独立に求めた組合せ 90/ホルムズ/サウジ/10 と異なる。\par

\noindent\textbf{3：不適切} 90/ホルムズ/クウェート/10。独立に求めた組合せ 90/ホルムズ/サウジ/10 と異なる。\par

\noindent\textbf{4：不適切} 67/マラッカ/クウェート/10。独立に求めた組合せ 90/ホルムズ/サウジ/10 と異なる。\par

\noindent\textbf{5：不適切} 67/ホルムズ/サウジ/32。独立に求めた組合せ 90/ホルムズ/サウジ/10 と異なる。\par

\paragraph{再利用する考え方} 年度と白書年、数量依存度と金額比、輸送経路。 否定方向と資料年、分母と単位を確認して五肢へ戻る。

\paragraph{共通理論} \hyperref[theory:env:oil2019]{第\ref*{theory:env:oil2019}節}

\par 難易度 3（校準初稿）、優先度未評価。

\section{2018発電比・2020太陽光費用と風速三乗}

\subsection{令和3年度 問5-3}
\label{q:r03-reg-basic-5-3}

\noindent\textbf{原典：} sources/R03/basic.pdf、PDF物理23頁。

\textbf{公式正答：} 1（single）。

\par\textbf{検討未確定：} 当時の発電統計と調達価格委員会の2020設置単純平均の原表を追加照合。全規模毎年低下という表現は個別年の反例を排除しない。

\paragraph{問題の要旨} 2018発電比・2020太陽光費用と風速三乗。

\Needspace{4\baselineskip}\paragraph{確認した条件・記述（独自要約）}\mbox{}\par

\noindent\textbf{原典：} 全五肢と否定方向・組合せ表を画像確認。\par

\noindent\textbf{資料時点：} 令和3年度の指定資料・統計年を保持し、後年値に置換しない。\par

\paragraph{問われる核心} 2018発電比・2020太陽光費用と風速三乗。

\paragraph{必要知識} 発電割合の水力境界、設備費の単純平均、風のエネルギー流束。

\paragraph{考え方・正答を選ぶ理由（査読待ち初稿）} 水力を除く2018年度再エネ発電割合は約9\%、2020設置の10kW以上太陽光システム費用の単純平均は約25万円/kWとして読む。風の運動エネルギー流束は、通過質量流量ρAvに単位質量エネルギーv\textsuperscript{2}/2を掛けるのでρAv\textsuperscript{3}/2。9/25/3の1。静止した一定質量の運動エネルギーv\textsuperscript{2}と混同しない。

\Needspace{4\baselineskip}\paragraph{選択肢の検討（査読待ち初稿）}\mbox{}\par

\noindent\textbf{1：適切・選択対象} 9/25/3。独立に求めた組合せ 9/25/3 と一致する。\par

\noindent\textbf{2：不適切} 14/25/3。独立に求めた組合せ 9/25/3 と異なる。\par

\noindent\textbf{3：不適切} 14/15/3。独立に求めた組合せ 9/25/3 と異なる。\par

\noindent\textbf{4：不適切} 9/25/2。独立に求めた組合せ 9/25/3 と異なる。\par

\noindent\textbf{5：不適切} 14/15/2。独立に求めた組合せ 9/25/3 と異なる。\par

\paragraph{再利用する考え方} 発電割合の水力境界、設備費の単純平均、風のエネルギー流束。 否定方向と資料年、分母と単位を確認して五肢へ戻る。

\paragraph{共通理論} \hyperref[theory:env:renewables]{第\ref*{theory:env:renewables}節}

\par 難易度 3（校準初稿）、優先度未評価。

\subsection{令和7年度 問5-4}
\label{q:r07-reg-basic-5-4}

\noindent\textbf{原典：} sources/R07/basic.pdf、PDF物理18頁。

\textbf{公式正答：} 2（single）。

\par\textbf{検討未確定：} 22.9\%は官庁発表で確認。2024設置10kW以上単純平均の原表は追加照合。価格算定用想定値や住宅用平均を代入しない。全規模毎年低下の一般化にも留保。

\paragraph{問題の要旨} 2023再エネ比・2024設置費と風速三乗。

\Needspace{4\baselineskip}\paragraph{確認した条件・記述（独自要約）}\mbox{}\par

\noindent\textbf{原典：} 全五肢と否定方向・組合せ表を画像確認。\par

\noindent\textbf{資料時点：} 令和7年度の指定資料・統計年を保持し、後年値に置換しない。\par

\paragraph{問われる核心} 2023再エネ比・2024設置費と風速三乗。

\paragraph{必要知識} 水力含有、対象年、単純平均と想定費用、流束の三乗。

\paragraph{考え方・正答を選ぶ理由（査読待ち初稿）} 水力を含む2023年度の再エネ発電割合22.9\%を約23\%、指定2024設置10kW以上の単純平均設備費を約23万円/kW、風のエネルギー流束をv\textsuperscript{3}とし2。R03の水力除外・2018年度約9\%と同じ値を使わない。風速二乗は一定質量の運動エネルギーであり、単位時間の通過量がさらにvに比例する。

\Needspace{4\baselineskip}\paragraph{選択肢の検討（査読待ち初稿）}\mbox{}\par

\noindent\textbf{1：不適切} 23/23/2。独立に求めた組合せ 23/23/3 と異なる。\par

\noindent\textbf{2：適切・選択対象} 23/23/3。独立に求めた組合せ 23/23/3 と一致する。\par

\noindent\textbf{3：不適切} 23/30/3。独立に求めた組合せ 23/23/3 と異なる。\par

\noindent\textbf{4：不適切} 15/30/2。独立に求めた組合せ 23/23/3 と異なる。\par

\noindent\textbf{5：不適切} 15/23/3。独立に求めた組合せ 23/23/3 と異なる。\par

\paragraph{再利用する考え方} 水力含有、対象年、単純平均と想定費用、流束の三乗。 否定方向と資料年、分母と単位を確認して五肢へ戻る。

\paragraph{共通理論} \hyperref[theory:env:renewables]{第\ref*{theory:env:renewables}節}

\par 参考：\cite{env60_re2023,env60_pv2024}。

\par 難易度 3（校準初稿）、優先度未評価。

\section{安全・標準化・製品責任の制定順}

\subsection{令和5年度 問5-5}
\label{q:r05-reg-basic-5-5}

\noindent\textbf{原典：} sources/R05/basic.pdf、PDF物理18頁。

\textbf{公式正答：} 1（single）。

\par\textbf{検討未確定：} 安全週間1928と工場法1911は官庁史料索引で確認。他3法の当初公布本文との照合は未完。

\paragraph{問題の要旨} 安全・標準化・製品責任の制定順。

\Needspace{4\baselineskip}\paragraph{確認した条件・記述（独自要約）}\mbox{}\par

\noindent\textbf{原典：} 全五肢と否定方向・組合せ表を画像確認。\par

\noindent\textbf{資料時点：} 令和5年度の指定資料・統計年を保持し、後年値に置換しない。\par

\paragraph{問われる核心} 安全・標準化・製品責任の制定順。

\paragraph{必要知識} 制定・施行・初回実施、改称と制定を区別。

\paragraph{考え方・正答を選ぶ理由（査読待ち初稿）} ウ工場法制定1911、エ第1回全国安全週間1928、オ工業標準化法制定1949、ア労働安全衛生法制定1972、イ製造物責任法制定1994の順で1。工場法施行1916、PL法施行1995と制定年を入れ替えない。改称した産業標準化法の現在名称も当時の制定時名称と区別する。

\Needspace{4\baselineskip}\paragraph{選択肢の検討（査読待ち初稿）}\mbox{}\par

\noindent\textbf{1：適切・選択対象} ウエオアイ。独立に求めた組合せ ウエオアイ と一致する。\par

\noindent\textbf{2：不適切} ウオエアイ。独立に求めた組合せ ウエオアイ と異なる。\par

\noindent\textbf{3：不適切} エウオイア。独立に求めた組合せ ウエオアイ と異なる。\par

\noindent\textbf{4：不適切} エオウアイ。独立に求めた組合せ ウエオアイ と異なる。\par

\noindent\textbf{5：不適切} オウアイエ。独立に求めた組合せ ウエオアイ と異なる。\par

\paragraph{再利用する考え方} 制定・施行・初回実施、改称と制定を区別。 否定方向と資料年、分母と単位を確認して五肢へ戻る。

\paragraph{共通理論} \hyperref[theory:env:safetyhistory]{第\ref*{theory:env:safetyhistory}節}

\par 参考：\cite{env60_safety}。

\par 難易度 3（校準初稿）、優先度未評価。

\section{第1〜5期科学技術基本計画の特徴}

\subsection{令和3年度 問5-6}
\label{q:r03-reg-basic-5-6}

\noindent\textbf{原典：} sources/R03/basic.pdf、PDF物理26頁。

\textbf{公式正答：} 4（single）。

\par\textbf{検討未確定：} 2・3期の主要句は一次資料索引で確認。1・4・5期本文の当該段落全体は追加照合。

\paragraph{問題の要旨} 第1〜5期科学技術基本計画の特徴。

\Needspace{4\baselineskip}\paragraph{確認した条件・記述（独自要約）}\mbox{}\par

\noindent\textbf{原典：} 全五肢と否定方向・組合せ表を画像確認。\par

\noindent\textbf{資料時点：} 令和3年度の指定資料・統計年を保持し、後年値に置換しない。\par

\paragraph{問われる核心} 第1〜5期科学技術基本計画の特徴。

\paragraph{必要知識} 計画の期・開始年・特徴句をセットで記憶。

\paragraph{考え方・正答を選ぶ理由（査読待ち初稿）} 第1期1996はウポストドクター等1万人支援、第2期2001はイ社会のため・社会の中の科学技術、第3期2006はア倫理的・法的・社会的課題、第4期2011はオ震災からの復興・再生、第5期2016はエSociety5.0。ウイアオエの4。ELSIの概念が以前になかったという意味ではなく、問の特徴句を指定計画へ対応させる。

\Needspace{4\baselineskip}\paragraph{選択肢の検討（査読待ち初稿）}\mbox{}\par

\noindent\textbf{1：不適切} イアウエオ。独立に求めた組合せ ウイアオエ と異なる。\par

\noindent\textbf{2：不適切} イウアオエ。独立に求めた組合せ ウイアオエ と異なる。\par

\noindent\textbf{3：不適切} ウアイエオ。独立に求めた組合せ ウイアオエ と異なる。\par

\noindent\textbf{4：適切・選択対象} ウイアオエ。独立に求めた組合せ ウイアオエ と一致する。\par

\noindent\textbf{5：不適切} ウイエアオ。独立に求めた組合せ ウイアオエ と異なる。\par

\paragraph{再利用する考え方} 計画の期・開始年・特徴句をセットで記憶。 否定方向と資料年、分母と単位を確認して五肢へ戻る。

\paragraph{共通理論} \hyperref[theory:env:scienceplans]{第\ref*{theory:env:scienceplans}節}

\par 参考：\cite{env60_plan2,env60_plan3}。

\paragraph{関連問題} \hyperref[q:r07-reg-basic-5-6]{R07-REG-BASIC-5-6}

\par 難易度 3（校準初稿）、優先度未評価。

\subsection{令和7年度 問5-6}
\label{q:r07-reg-basic-5-6}

\noindent\textbf{原典：} sources/R07/basic.pdf、PDF物理19頁。

\textbf{公式正答：} 4（single）。

\par\textbf{検討未確定：} 2・3期主要句の一次索引を確認。4〜6期原本文の該当段落と2020改正・施行日を追加査読。

\paragraph{問題の要旨} 第2〜6期の社会・ELSI・復興・Society・総合知。

\Needspace{4\baselineskip}\paragraph{確認した条件・記述（独自要約）}\mbox{}\par

\noindent\textbf{原典：} 全五肢と否定方向・組合せ表を画像確認。\par

\noindent\textbf{資料時点：} 令和7年度の指定資料・統計年を保持し、後年値に置換しない。\par

\paragraph{問われる核心} 第2〜6期の社会・ELSI・復興・Society・総合知。

\paragraph{必要知識} 科学技術政策の期と特徴、総合知、改正と計画開始。

\paragraph{考え方・正答を選ぶ理由（査読待ち初稿）} イ第2期2001は社会のため・社会の中、ア第3期2006はELSI、オ第4期2011は震災復興、エ第5期2016はSociety5.0、ウ第6期2021は人文社会科学と自然科学の融合による総合知。イアオエウの4。2020改正法の名称変更と2021開始の第6期計画を区別する。

\Needspace{4\baselineskip}\paragraph{選択肢の検討（査読待ち初稿）}\mbox{}\par

\noindent\textbf{1：不適切} イアウオエ。独立に求めた組合せ イアオエウ と異なる。\par

\noindent\textbf{2：不適切} アウエオイ。独立に求めた組合せ イアオエウ と異なる。\par

\noindent\textbf{3：不適切} アイオウエ。独立に求めた組合せ イアオエウ と異なる。\par

\noindent\textbf{4：適切・選択対象} イアオエウ。独立に求めた組合せ イアオエウ と一致する。\par

\noindent\textbf{5：不適切} アイオエウ。独立に求めた組合せ イアオエウ と異なる。\par

\paragraph{再利用する考え方} 科学技術政策の期と特徴、総合知、改正と計画開始。 否定方向と資料年、分母と単位を確認して五肢へ戻る。

\paragraph{共通理論} \hyperref[theory:env:scienceplans]{第\ref*{theory:env:scienceplans}節}

\par 参考：\cite{env60_plan2,env60_plan3}。

\paragraph{関連問題} \hyperref[q:r03-reg-basic-5-6]{R03-REG-BASIC-5-6}

\par 難易度 3（校準初稿）、優先度未評価。

\section{ラムサール条約の対象は湿地}

\subsection{令和6年度 問5-1}
\label{q:r06-reg-basic-5-1}

\noindent\textbf{原典：} sources/R06/basic.pdf、PDF物理17頁。

\textbf{公式正答：} 1（single）。

\par\textbf{検討未確定：} 各条約・議定書本文の定義と適用範囲の逐条確認は追加査読対象。

\paragraph{問題の要旨} ラムサール条約の対象は湿地。

\Needspace{4\baselineskip}\paragraph{確認した条件・記述（独自要約）}\mbox{}\par

\noindent\textbf{原典：} 全五肢と否定方向・組合せ表を画像確認。\par

\noindent\textbf{資料時点：} 令和6年度の指定資料・統計年を保持し、後年値に置換しない。\par

\paragraph{問われる核心} ラムサール条約の対象は湿地。

\paragraph{必要知識} 国際環境条約の対象と規制行為、ABS。

\paragraph{考え方・正答を選ぶ理由（査読待ち初稿）} 1が不適切。条約名だけではなく、対象（湿地・種・物質・廃棄物・遺伝資源）、行為（利用・取引・生産・移動・取得）を対応させる。森林に湿地が含まれる事例があっても条約全体を森林保全と定義しない。

\Needspace{4\baselineskip}\paragraph{選択肢の検討（査読待ち初稿）}\mbox{}\par

\noindent\textbf{1：不適切・選択対象} ラムサールは重要森林の保全。重要な湿地の保全と賢明な利用を目的とし、森林一般に置換するのは不適切。\par

\noindent\textbf{2：条件内で適切} ワシントンは希少野生種の国際取引規制。CITESは種の国際取引を管理する。\par

\noindent\textbf{3：条件内で適切} モントリオールはODSの生産消費規制。オゾン層破壊物質を対象とする枠組み。後の改正で対象は広がる。\par

\noindent\textbf{4：条件内で適切} バーゼルは有害廃棄物の越境移動・処分。国境を越える有害廃棄物等の管理。\par

\noindent\textbf{5：条件内で適切} 名古屋は遺伝資源取得と利益配分。ABSは取得機会と利用利益の公正衡平な配分。\par

\paragraph{再利用する考え方} 国際環境条約の対象と規制行為、ABS。 否定方向と資料年、分母と単位を確認して五肢へ戻る。

\paragraph{共通理論} \hyperref[theory:env:treaties]{第\ref*{theory:env:treaties}節}

\par 参考：\cite{env60_ramsar}。

\par 難易度 3（校準初稿）、優先度未評価。

\section{独立正規変数の和差と標本平均}

\subsection{平成23年度 問5-4}
\label{q:h23-reg-basic-5-4}

\noindent\textbf{原典：} sources/H23/basic.pdf、PDF物理17頁。

\textbf{公式正答：} 3（single）。

\paragraph{問題の要旨} 独立な製品質量の標本平均と総質量について、平均と標準偏差を求める。

\Needspace{4\baselineskip}\paragraph{確認した条件・記述（独自要約）}\mbox{}\par

\noindent\textbf{判断方向：} 最も適切な記述又は組合せを選ぶ\par

\noindent\textbf{原典照合：} 14原典画像で16問の条件、数値、単位、表の列と全5候補を確認。\par

\noindent\textbf{モデル：} 独立性、正規性、一定発生率などの前提を個別に保持する。\par

\paragraph{問われる核心} 独立な製品質量の標本平均と総質量について、平均と標準偏差を求める。

\paragraph{必要知識} 独立変数の線形結合、平均の線形性、分散への係数の二乗、正規分布の閉性。

\paragraph{考え方・正答を選ぶ理由（査読待ち初稿）} 各質量を独立正規変数とする。Aの標本平均は \(\bar A=(A_1+\cdots+A_4)/4\) で、期待値80g、分散 \(4(4^2)/4^2=4\,\mathrm{g^2}\) である。総質量は \(T=\sum_{i=1}^4A_i+\sum_{i=1}^4B_i\) なので期待値 \(4(80)+4(120)=800\,\mathrm g\)、分散 \(4(4^2)+4(3^2)=100\,\mathrm{g^2}\) となる。標準偏差はそれぞれ2gと10gで、3を選ぶ。平均を取る操作と、四つの値を合計する操作を混同しない。

\Needspace{4\baselineskip}\paragraph{選択肢の検討（査読待ち初稿）}\mbox{}\par

\noindent\textbf{1：不適切} 平均20・偏差4・合計400・偏差10。Aの標本平均の期待値は80gであり、標本数で平均を割らない。総質量の期待値も800gである。\par

\noindent\textbf{2：不適切} 平均80・偏差4・合計800・偏差7。平均の偏差は4gではなく2gである。合計の偏差も標準偏差をそのまま足した7gではなく10gである。\par

\noindent\textbf{3：適切} 平均80・偏差2・合計800・偏差10。Aの4個平均の偏差は4/2=\allowbreak{}2g、合計の分散は4×16+4×9=\allowbreak{}100g\textsuperscript{2}なので10gである。\par

\noindent\textbf{4：不適切} 平均20・偏差2・合計400・偏差7。平均と総質量の期待値がともに違い、合計の偏差も違う。\par

\noindent\textbf{5：不適切} 平均80・偏差2・合計800・偏差7。平均部分は合うが、合計の標準偏差は7gではない。独立成分の分散を足して最後に平方根を取る。\par

\paragraph{再利用する考え方} 変数の意味と単位を先に定め、平均・分散・確率・判定を順に分ける。品質の要求、工程の監視、ロットの受入れを区別し、成立条件のない一般化を避ける。

\paragraph{共通理論} \hyperref[theory:quality:variance]{第\ref*{theory:quality:variance}節}, \hyperref[theory:quality:normal]{第\ref*{theory:quality:normal}節}, \hyperref[theory:quality:sample_mean]{第\ref*{theory:quality:sample_mean}節}

\par 参考：\cite{stat_nist_normal,stat_psu_linear}。

\paragraph{関連問題} \hyperref[q:r06-reg-basic-1-6]{R06-REG-BASIC-1-6}

\par 難易度 3（校準初稿）、優先度未評価。

\section{工程の品質保証と標準化}

\subsection{平成24年度 問5-4}
\label{q:h24-reg-basic-5-4}

\noindent\textbf{原典：} sources/H24/basic.pdf、PDF物理17頁。

\textbf{公式正答：} 2（single）。

\paragraph{問題の要旨} 最終検査と上流の品質保証を分け、プロセスの標準化と改善の関係を読む。

\Needspace{4\baselineskip}\paragraph{確認した条件・記述（独自要約）}\mbox{}\par

\noindent\textbf{判断方向：} 最も適切な記述又は組合せを選ぶ\par

\noindent\textbf{原典照合：} 14原典画像で16問の条件、数値、単位、表の列と全5候補を確認。\par

\noindent\textbf{モデル：} 独立性、正規性、一定発生率などの前提を個別に保持する。\par

\paragraph{問われる核心} 最終検査と上流の品質保証を分け、プロセスの標準化と改善の関係を読む。

\paragraph{必要知識} 検査と予防、上流工程、標準化、PDCA、プロセスと成果。

\paragraph{考え方・正答を選ぶ理由（査読待ち初稿）} 完成品を検査して不適合品を隔離することと、不適合を生む工程を上流から改善することを分ける。設問の文脈ではアが検査、イが源流管理、ウが標準化、エがプロセスであり2となる。標準を設定して実施し、結果を確認して更新する循環が品質保証を支える。標準化は永久に方法を固定することではなく、改善結果を次の標準へ反映することを含む。

\Needspace{4\baselineskip}\paragraph{選択肢の検討（査読待ち初稿）}\mbox{}\par

\noindent\textbf{1：不適切} 検査・三現管理・基準化・プロセス。現場・現物・現実という三現の確認と、上流で不適合原因を抑える源流管理は同じ説明ではない。\par

\noindent\textbf{2：適切} 検査・源流管理・標準化・プロセス。要求適合を判定する検査、上流からの源流管理、方法を定める標準化、プロセスを管理する流れでつながる。\par

\noindent\textbf{3：不適切} 検査・三現管理・基準化・アウトカム。三現管理・基準化の対応に加え、ここで管理する対象は一連のプロセスである。\par

\noindent\textbf{4：不適切} 評価・源流管理・標準化・アウトカム。評価は広い概念だが本問の要求適合判定は検査である。末尾も成果だけを意味するアウトカムではなくプロセスである。\par

\noindent\textbf{5：不適切} 評価・三現管理・標準化・プロセス。要求適合判定と上流の改善に対応する語が検査・源流管理になっていない。\par

\paragraph{再利用する考え方} 変数の意味と単位を先に定め、平均・分散・確率・判定を順に分ける。品質の要求、工程の監視、ロットの受入れを区別し、成立条件のない一般化を避ける。

\paragraph{共通理論} \hyperref[theory:quality:assurance]{第\ref*{theory:quality:assurance}節}, \hyperref[theory:quality:pdca]{第\ref*{theory:quality:pdca}節}

\par 参考：\cite{quality_process_jsqc,quality_iso_process}。

\paragraph{関連問題} \hyperref[q:r03-reg-basic-1-3]{R03-REG-BASIC-1-3}

\par 難易度 2（校準初稿）、優先度未評価。


\chapter{技術士法・義務責務}
技術士法の第4章は、専門家が高度な知識を持つことに対応した行動上の責任を定める。数式の正しさだけで技術サービスの適否は決まらない。技術の危険を十分に評価し、関係者へ事実を説明し、資格と能力の範囲を明らかにすることが信頼の前提である。本章では条文の語句、具体的な業務判断、資格取得後の継続研鑚を結び付ける。

\section{義務・責務を主語と条件から読む}
\label{theory:engineer_act:duties}
過去問で示された第4章の条文と、2026年10月8日にe-Gov公式APIから取得した2025年6月1日施行版を対照した。以下は本章の設問で扱う範囲の整理であり、全改正履歴の照合を意味しない\cite{engineer_act_current}。
\begin{longtable}{p{.13\linewidth}p{.29\linewidth}p{.47\linewidth}}
条文 & 要点 & 適用対象と条件\\\hline
第44条 & 信用失墜行為の禁止 & 技術士・技術士補。専門家集団の信用を害する行為を禁止する。\\
第45条 & 秘密保持義務 & 業務上知った秘密の漏えい・盗用を正当な理由なしに行わない。資格を失った後も続く。\\
第45条の2 & 公益確保の責務 & 技術士・技術士補。公共の安全、環境その他公益を害さないよう努める。\\
第46条 & 名称表示の場合の義務 & 技術士の名称を業務に関して表示する場合、登録技術部門を明示し、未登録部門を表示しない。\\
第47条 & 技術士補の業務の制限等 & 技術士補の名称での業務は技術士を補助する場合に限る。名称表示の規定を準用する。\\
第47条の2 & 資質向上の責務 & 技術士。その業務の知識技能と資質を継続的に高めるよう努める。\\
\end{longtable}
第4章という見出し全体には「義務」と書かれるが、個々の条文には義務と責務の両方がある。公益確保・資質向上は、努力を求める責務という語で整理する。「責務だからしなくてもよい」という結論は出ない。名称表示をする場合という条件を、すべての業務実施時に無条件に拡張することも避ける。

法令で区分された登録先は「技術部門」である。「専門部門」という日常語へ置き換えない。第47条2項の「準用」は、別の対象へ必要な読み替えを伴って規定を働かせることであり、その規定が元から直接対象とする場合の適用とは区別する。平成26・29・30年度、令和元・2・6年度の空欄問題を、この表と照合すると選択肢の違いを説明できる。

登録部門の力量は取得時に固定されるものではない。技術の変化、法令、失敗事例を学び、自分の業務能力を更新する必要がある。別部門を学ぶことは有益だが、現在の登録部門の学習を不要にするものではない。また顧客の同意だけでは技術士補の名称による独立業務や未登録技術部門の表示が認められることにはならない。

\section{守秘・公益・説明責任を同時に満たす}
\label{theory:engineer_act:public}
守秘義務は、情報を託す相手が安心して技術者に相談できるための基盤である。所属組織の取引先だけでなく、業務上知った秘密一般を対象とし、退職・資格喪失後も存続する。一方、公衆に危険がある場合にも顧客の利益だけを優先することはできない。技術士倫理綱領は、公衆の利益、有能性、真実性、秘密の保持などを相互に関連する行動原則として挙げる\cite{cpd_v3}。

例えば試験データに危険な不具合が表れているのに顧客が報告から除外するよう指示した場合、まず事実と不確かさを整理し、危険と必要な対策を責任ある相手へ説明する。組織内で解決しないときも、公衆の安全を確保する適切な手段を検討する。「顧客命令なら何でも従う」「守秘を理由に危険を放置する」「すべての秘密を直ちに公開する」という一律の処理では解けない。公益確保と守秘の対象・理由・手続を具体的に検討する必要がある。

利益相反は、報酬・所属・人的関係などが専門的判断の公正さに影響し得る状態である。関係を説明し、判断の根拠を客観的な事実へ戻す。関係者の話をよく聞くことと、事実に反した結論を受け入れることは別である。

依頼者の情報から独自に調査したとしても、元の秘密性や契約上の利用制限が自動的に消えるわけではない。さらに情報の所有・利用権と守秘義務は別の問題である。平成23年度問1では権利を当事者の協議へ委ねる表現、平成27年度問1・平成30年度問2・令和5年度問1では独自調査だから制限がなくなるという表現の違いを読む。

\section{適性科目の規範：法の形式と倫理の内容}
\label{theory:engineer_act:norms}
社会規範は、人々の行動を社会の中で評価する基準である。そのうち法規範は法律として制定され、倫理規範は何が責任ある行動かを示す。技術士法第4章は形式上の法律であり、同時に専門家の職業倫理を表している。この二つの性格は両立する。令和元年度再試験問1は、社会規範を上位概念とし、第4章を法規範かつ倫理規範として整理する問題である。

\section{CPD：資格取得後に能力を更新する}
\label{theory:engineer_act:cpd}
CPDは Continuing Professional Development の略である。学習・研修・調査研究などを通して専門家としての能力を継続して高める。指定された第2版は「自己研鑚」、第3版は「継続研鑚」という語を用いる。両版の2.2節では、その対象が本人の能力と立場によって異なり、実施記録と証拠により説明できることを求めている\cite{cpd_v2,cpd_v3}。平成27年度問2・平成30年度問3の空欄は、出題指定資料の語句を確かめる。

現在の知識を使って日常業務を繰り返すことと、その業務に関連して新しい方法を調べ能力を向上させることは別である。研修で何を学んだか、業務でどの能力が増えたかを記録する。実施の事実があっても、その内容を証明できなければ登録上の実績として認められない場合がある。

2021年の大臣通知・省令改正に基づく制度では、希望する技術士のCPD実績を登録簿の資質向上の取組状況へ記載し、条件に応じた認定や公表を行う仕組みが整えられた\cite{cpd_new2021}。令和6年度問15はこの制度も問う。実績の記載・認定制度を、すべての技術士に一律の登録更新義務があるという制度へ読み替えない。

ガイドラインと算定基準は実施年度によって異なる。第1版の対象は2011〜2013年度、第2版は2014〜2016年度、第3版は2017〜2021年度であり、その後の新制度と現在の版は別である\cite{cpd_versions}。過去問を現在の時間算定基準だけで解かない。

平成23年度問4では、資質向上の不履行を「行政罰」に相当するとした表現が複数正答の検討点になる。第36条2項の登録取消し等という行政上の措置、罰則、努力責務の履行評価を区別し、原典の記述と公式の複数正答を保持する。どのようなCPD不足でも自動的に登録が取り消されるという説明にはしない。

\section{コンピテンシー：知識を仕事の能力へ結び付ける}
\label{theory:engineer_act:competency}
平成26年3月の技術士分科会資料は、専門的学識、問題解決、マネジメント、評価、コミュニケーション、リーダーシップ、技術者倫理の7つを示す\cite{mext_competency2014}。令和元年度問2では、列挙された7つすべてがその資料に含まれる。

物理の専門知識を正確に適用するだけでは、複数の制約がある業務を完遂できない。問題解決では矛盾する要求を整理し、マネジメントでは人・設備・資金・情報を配分し、評価では成果を次の改善へつなげる。コミュニケーションとリーダーシップは関係者の理解と調整、技術者倫理は公益と責任を支える。CPDはこれらの能力を業務の変化に合わせて更新する活動である。ここでは出題指定の平成26年資料の項目を扱い、その後の全改訂を同一と見なさない。

\section{教育認定から国際的な業務資格へ}
\label{theory:engineer_act:international}
技術者が国境を越えて活動するには、教育の質、実務能力、資格登録、受入国の法制度を区別する必要がある。教育課程の認定は、その課程で得られる学習成果の質を評価する。JABEEの認定だけで全ての受入国の業務免許が与えられるわけではない。

\begin{description}
 \item[教育と第一次試験免除] 文部科学大臣が指定する認定課程の修了者は、技術士補となる資格を持つ。第一次試験を免除されることと、技術士補として登録を受けることは別の段階である\parencite{engineer_act_current,ipej_associate_registration}。
 \item[国際称号] IPEAやAPECの登録では、教育・実務経験・独立した業務遂行能力・継続学習等の要件を審査する。国内の技術士資格だけで自動的に国際登録となるのではない。IPEAの和文名称は2015年4月に旧EMF国際エンジニアから変更された\parencite{ipej_ipea2014}。
 \item[受入国の業務資格] 共通称号は能力の同等性を示す枠組みである。相互承認には別の二国間枠組みや固有事項審査等が関わり、受入国の法・登録・入国手続まで全て不要になるとはいえない。
\end{description}

\paragraph{平成27年当時と現在の建築教育}
建築系の認定についてJABEEは2012年のUNESCO-UIAシステム認証、2014年のキャンベラ協定暫定加盟、2019年の正式加盟を説明する\parencite{jabee_canberra}。平成27年の肢にあるUNESCO-UIAの要約を現在の正式加盟先に置き換えず、機関の認証と協定への加盟という手続の違いも確認する。

\paragraph{二国間の審査合理化}
文科省は日豪の2003年枠組み、建築技術教育普及センターは2008年の日豪・2009年の日NZ枠組みを説明する\parencite{mext_international,jaeic_mutual2010}。一般的な能力の審査と、受入国固有の法的・技術的事項の審査を分ける。現在の手続や対象分野の表を、2015年の全条件と同一だと扱わない。

\section*{各問題の個別検討}
以下は査読待ち初稿である。平成25年度問1では独立した条文読解と公式の個数が一致しないため、相違を表示する。平成23年度問4の複数正答も一つに確定しない。
\section{継続研鑽：活動・実績・証明をつなぐ}
\label{theory:profession:cpd}
CPDは専門職として能力を継続的に維持・向上させる活動である。仕事に従事している事実だけから、必要な学習と実績記録をすべて満たすとはしない。実務の中で何を学び、能力をどう改善したかを説明できる記録を残す\parencite{cpd_v3,cpd_versions,cpd_new2021}。

令和4年度問題には参加型、発信型、実務型、自己学習型という区分がある。研修の受講、技術指導や執筆、業務成果、自己学習は活動の性格が違う。換算条件と証明は制度・版ごとに確認する。学協会によって資格継続に確認条件があることと、国内の技術士資格が一律に自動失効することも区別する。

\section{2023コンピテンシーの八つのキーワード}
\label{theory:profession:competency2023}
2023年の資料では、専門的学識、問題解決、マネジメント、評価、コミュニケーション、リーダーシップ、技術者倫理、継続研さんの八つを挙げる。2014年の七つのキーワードから継続研さんを追加した時点を確認する\parencite{eth54_competency2023}。

知識を持つことだけでなく、複合的な問題を定義し、制約と利害を調べ、資源を配分し、結果を評価し、関係者と協力して責任ある成果を得る能力を考える。35歳程度での活躍への期待は、法定の受験年齢上限や資格の自動取消条件ではない。文書の版と発表主体を読み、2023版で独立しているコミュニケーションとリーダーシップを一つの項目にまとめない。

\subsection{平成23年度 問1}
\label{q:h23-reg-apt-01}

\noindent\textbf{原典：} sources/H23/aptitude.pdf、PDF物理2頁。

\textbf{公式正答：} 4（single）。

\par\textit{原典本文の校正と独自解説は未作成。分類のみの登録項目。}

\subsection{平成23年度 問4}
\label{q:h23-reg-apt-04}

\noindent\textbf{原典：} sources/H23/aptitude.pdf、PDF物理5頁。

\textbf{公式正答：} 1|5（multiple\_accepted）。

\par\textit{原典本文の校正と独自解説は未作成。分類のみの登録項目。}

\subsection{平成24年度 問1}
\label{q:h24-reg-apt-01}

\noindent\textbf{原典：} sources/H24/aptitude.pdf、PDF物理1頁。

\textbf{公式正答：} 5（single）。

\par\textit{原典本文の校正と独自解説は未作成。分類のみの登録項目。}

\subsection{平成24年度 問14}
\label{q:h24-reg-apt-14}

\noindent\textbf{原典：} sources/H24/aptitude.pdf、PDF物理12頁。

\textbf{公式正答：} 4（single）。

\par\textit{原典本文の校正と独自解説は未作成。分類のみの登録項目。}

\subsection{平成25年度 問1}
\label{q:h25-reg-apt-01}

\noindent\textbf{原典：} sources/H25/aptitude.pdf、PDF物理1頁。

\textbf{公式正答：} 4（single）。

\par\textit{原典本文の校正と独自解説は未作成。分類のみの登録項目。}

\subsection{平成26年度 問1}
\label{q:h26-reg-apt-01}

\noindent\textbf{原典：} sources/H26/aptitude.pdf、PDF物理1頁。

\textbf{公式正答：} 3（single）。

\par\textit{原典本文の校正と独自解説は未作成。分類のみの登録項目。}

\subsection{平成27年度 問1}
\label{q:h27-reg-apt-01}

\noindent\textbf{原典：} sources/H27/aptitude.pdf、PDF物理1頁。

\textbf{公式正答：} 5（single）。

\par\textit{原典本文の校正と独自解説は未作成。分類のみの登録項目。}

\subsection{平成27年度 問2}
\label{q:h27-reg-apt-02}

\noindent\textbf{原典：} sources/H27/aptitude.pdf、PDF物理2頁。

\textbf{公式正答：} 5（single）。

\par\textit{原典本文の校正と独自解説は未作成。分類のみの登録項目。}

\subsection{平成28年度 問1}
\label{q:h28-reg-apt-01}

\noindent\textbf{原典：} sources/H28/aptitude.pdf、PDF物理1頁。

\textbf{公式正答：} 2（single）。

\par\textit{原典本文の校正と独自解説は未作成。分類のみの登録項目。}

\subsection{平成28年度 問2}
\label{q:h28-reg-apt-02}

\noindent\textbf{原典：} sources/H28/aptitude.pdf、PDF物理2頁。

\textbf{公式正答：} 3（single）。

\par\textit{原典本文の校正と独自解説は未作成。分類のみの登録項目。}

\subsection{平成29年度 問1}
\label{q:h29-reg-apt-01}

\noindent\textbf{原典：} sources/H29/aptitude.pdf、PDF物理1頁。

\textbf{公式正答：} 3（single）。

\par\textit{原典本文の校正と独自解説は未作成。分類のみの登録項目。}

\subsection{平成29年度 問2}
\label{q:h29-reg-apt-02}

\noindent\textbf{原典：} sources/H29/aptitude.pdf、PDF物理2頁。

\textbf{公式正答：} 5（single）。

\par\textit{原典本文の校正と独自解説は未作成。分類のみの登録項目。}

\subsection{平成30年度 問1}
\label{q:h30-reg-apt-01}

\noindent\textbf{原典：} sources/H30/aptitude.pdf、PDF物理1頁。

\textbf{公式正答：} 2（single）。

\par\textit{原典本文の校正と独自解説は未作成。分類のみの登録項目。}

\subsection{平成30年度 問2}
\label{q:h30-reg-apt-02}

\noindent\textbf{原典：} sources/H30/aptitude.pdf、PDF物理2頁。

\textbf{公式正答：} 5（single）。

\par\textit{原典本文の校正と独自解説は未作成。分類のみの登録項目。}

\subsection{平成30年度 問3}
\label{q:h30-reg-apt-03}

\noindent\textbf{原典：} sources/H30/aptitude.pdf、PDF物理3頁。

\textbf{公式正答：} 3（single）。

\par\textit{原典本文の校正と独自解説は未作成。分類のみの登録項目。}

\subsection{令和元年度 問1}
\label{q:r01-reg-apt-01}

\noindent\textbf{原典：} sources/R01/aptitude.pdf、PDF物理1頁。

\textbf{公式正答：} 5（single）。

\par\textit{原典本文の校正と独自解説は未作成。分類のみの登録項目。}

\subsection{令和元年度 問2}
\label{q:r01-reg-apt-02}

\noindent\textbf{原典：} sources/R01/aptitude.pdf、PDF物理3頁。

\textbf{公式正答：} 5（single）。

\par\textit{原典本文の校正と独自解説は未作成。分類のみの登録項目。}

\subsection{令和元年度再試験 問1}
\label{q:r01-ret-apt-01}

\noindent\textbf{原典：} sources/R01\_retest/aptitude.pdf、PDF物理1頁。

\textbf{公式正答：} 3（single）。

\par\textit{原典本文の校正と独自解説は未作成。分類のみの登録項目。}

\subsection{令和元年度再試験 問2}
\label{q:r01-ret-apt-02}

\noindent\textbf{原典：} sources/R01\_retest/aptitude.pdf、PDF物理2頁。

\textbf{公式正答：} 3（single）。

\par\textit{原典本文の校正と独自解説は未作成。分類のみの登録項目。}

\subsection{令和2年度 問1}
\label{q:r02-reg-apt-01}

\noindent\textbf{原典：} sources/R02/aptitude.pdf、PDF物理1頁。

\textbf{公式正答：} 5（single）。

\par\textit{原典本文の校正と独自解説は未作成。分類のみの登録項目。}

\subsection{令和3年度 問1}
\label{q:r03-reg-apt-01}

\noindent\textbf{原典：} sources/R03/aptitude.pdf、PDF物理1頁。

\textbf{公式正答：} 3（single）。

\par\textit{原典本文の校正と独自解説は未作成。分類のみの登録項目。}

\subsection{令和4年度 問1}
\label{q:r04-reg-apt-01}

\noindent\textbf{原典：} sources/R04/aptitude.pdf、PDF物理1頁。

\textbf{公式正答：} 4（single）。

\par\textit{原典本文の校正と独自解説は未作成。分類のみの登録項目。}

\subsection{令和5年度 問1}
\label{q:r05-reg-apt-01}

\noindent\textbf{原典：} sources/R05/aptitude.pdf、PDF物理1頁。

\textbf{公式正答：} 3（single）。

\par\textit{原典本文の校正と独自解説は未作成。分類のみの登録項目。}

\subsection{令和6年度 問1}
\label{q:r06-reg-apt-01}

\noindent\textbf{原典：} sources/R06/aptitude.pdf、PDF物理1頁。

\textbf{公式正答：} 4（single）。

\par\textit{原典本文の校正と独自解説は未作成。分類のみの登録項目。}

\subsection{令和6年度 問15}
\label{q:r06-reg-apt-15}

\noindent\textbf{原典：} sources/R06/aptitude.pdf、PDF物理15頁。

\textbf{公式正答：} 2（single）。

\par\textit{原典本文の校正と独自解説は未作成。分類のみの登録項目。}

\subsection{令和7年度 問1}
\label{q:r07-reg-apt-01}

\noindent\textbf{原典：} sources/R07/aptitude.pdf、PDF物理1頁。

\textbf{公式正答：} 3（single）。

\par\textit{原典本文の校正と独自解説は未作成。分類のみの登録項目。}


\chapter{技術者倫理・意思決定}

\section{背景：技術的に可能なことを、実施すべきことへ結び付ける}
技術は人の目的を実現する手段である。しかし、実現できること、依頼者が望むこと、社会が受け入れられることは必ずしも一致しない。材料の変更で納期と原価を改善できても、強度の低下によって使用者へ危険を移してよいとはいえない。専門家の判断は、その技術を十分に理解せず、個別に同意する機会もない人へ影響する。この知識と力の非対称性が技術者の社会的責任の背景にある。

物理モデルは「この条件なら何が起こるか」を説明する。倫理的判断は、その結果を誰が受け、どの負担を許容し、どの手続で決めるかも扱う。事実の把握が不十分なら判断は誤るが、事実だけから価値の優先順位が自動的に出るわけでもない。技術的な計算と価値判断の双方を明示する必要がある。

\section{公衆の優先と専門職への信頼}
\label{theory:ethics:public}
公衆には、顧客だけでなく、施設の周辺住民、製品の使用者、将来世代なども含まれる。依頼者が了承しても第三者の安全を損なう行為が許されるとは限らない。専門職への信頼は、力量、事実に基づく説明、公正な履行を通じて成立する。

\paragraph{出題当時と現在の資料}
2011年改訂の技術士倫理綱領は、公衆の利益、真実性、誠実な履行、正当理由のない秘密の開示の禁止などを示す\parencite{ipej_ethics2011}。2023年版では基本事項に具体的な指針が付され、リスクの評価、代替案、業務内容の説明等を明示する\parencite{ipej_ethics2023}。過去問の語句を現在の番号へ無断で置換しない。

安全を最優先にするとは、「ゼロリスク」を無根拠に断言することでも、全ての新技術を一律に禁止することでもない。危険を調査し、必要な仕様と対策を満たす案を見いだし、残る不確実さを関係者へ説明することである。

\section{倫理原理：結果・義務・徳という異なる問い}
\label{theory:ethics:principles}
\subsection{結果を検討する視点}
功利主義は、行為によって生じる幸福や利益、害の結果を重視する。技術の便益だけでなく、事故、環境負荷、負担の移転まで評価対象を広げる。ただし、全体の利益を合計すると少数者の重い負担が隠れる可能性がある。結果の集計と分配の公平性は別に検討する必要がある。

\subsection{義務と人格を尊重する視点}
義務を重視する立場は、有益な結果が予想されても守るべき約束や権利があると考える。虚偽の報告によって仕事を守ることは、望む結果を得るために相手を単なる手段として扱う問題を含む。人の同意、尊厳、真実へのアクセスを結果の総量だけで処分しない。

\subsection{徳と専門職の人格を問う視点}
徳を重視する立場は、その選択を繰り返すとどのような技術者になるかを問う。誠実さ、慎重さ、勇気、他者への配慮は、規程の空欄を埋める日常的な判断にも働く。能力が高いことと、責任ある判断をする人格は同一ではない。

各視点は問いが異なるので、同じ事例について異なる懸念を発見する。大学の倫理教育資料も複数の視点を用いた検討を提示する\parencite{scu_ethics_framework}。一つの原理が必ず唯一の正解を与えると仮定せず、衝突する理由を明らかにする。

\section{意思決定を手続として組み立てる}
\label{theory:ethics:procedure}
個人の直観だけで結論を出すと、知らない条件や他者への影響が抜けやすい。問題を明確にし、確かめた事実と不明事項を分け、影響を受ける人々を洗い出し、複数の行動案を比較し、相談・実行・振り返りへ進める。このような検討の枠組みが倫理的判断を支える\parencite{davis_seven_steps}。

工程変更なら、まず安全と要求品質が満たされるかを検討する。その制約を満たす案について納期と原価を比較する。「安い」「早い」を最初に固定して、残る要求を都合よく変更しない。ただし、個別設問で選んだ納期と原価の順位を全業務の不変な法則にしない。

予算を扱うときは、成果だけでなく承認と記録も重要である。超過や遅延の見通しを共有し、変更を承認する権限を持つ人へ説明する。私費を投入しても所有・管理・会計の手続を省けるわけではない。使い切るために不要な支出を作らず、未消化額を適切に扱う。

数量化できない評価項目をゼロにしてよいともいえない。安全、組織の信用、社会的な受容、復旧の重要性などは、尺度の定義と評価理由を示して比較する。計数の容易さは重要性の大小ではない。

\section{エシックス・テスト：判断に別の視点を入れる}
\label{theory:ethics:tests}
テストは正答を自動的に返す算法ではなく、見落としを発見する問いである。
\begin{description}
 \item[普遍化可能性] 同じ行為を誰もが行うと、その行為が前提とする約束や社会制度は成り立つかを考える。
 \item[可逆性] 影響を受ける相手の立場に自分を置いても同じ判断を支持するかを考える。
 \item[美徳] この選択を繰り返すことが自分の専門職としての人格をどう形作るかを考える。
 \item[世評] 判断とその理由が公に知られた場合に説明できるかを考える。
 \item[専門家] 所属する専門職の規範や他の専門家の吟味に照らして評価する。
\end{description}
立場交換、人格、公開、専門職の評価という視点はDavisの意思決定教材でも区別される\parencite{davis_seven_steps}。世評を単なる人気投票にせず、美徳を一回の印象だけで判断しない。一つのテストを通過しても、他の責任が免除されるわけではない。

\section{判断を阻害する要因と道徳的自律}
\label{theory:ethics:obstacles}
無知は判断に必要な情報を欠くことであり、知らないことを自覚して調べる姿勢とは異なる。自己中心的志向は自分の立場だけで問題を切り取ることである。微視的視野は細部にとらわれ、周囲や長期の影響を見失うことをいう。細かい調査そのものは悪いことではない。

自己欺瞞は、都合の悪い事実を直視せず、他人のため、今回だけ、皆もしているなどの理由で自分の行為を正当化することである。責任を果たす結果を恐れると、問題を報告せずに済ませたい誘因が生じる。利他性や勇気を一時的な感情で終わらせず、理由と証拠に基づく行動へ結び付ける。

道徳的自律とは、判断の理由を自分で吟味し、その判断への責任を持つことである。専門家への相談や組織での協働と矛盾しない。問題になるのは、指示者の権威だけを理由に判断を放棄すること、集団で合意したからという理由で批判的な検討を省くことである。異論を聞き、事実と評価を点検できる環境が重要となる。

\section{情報開示：警告と秘密保持を具体的に調整する}
\label{theory:ethics:disclosure}
危険の可能性を知った場合、何を確認できているか、被害がどれほど迫るか、誰へどの情報を伝えれば危険を減らせるかを検討する。必要な警告や緊急措置を遅らせる理由として秘密保持を使ってはならない。一方で、「危険の可能性があるなら即時・無条件に全情報を公開する」とするのも、誤情報、個人情報、関係のない秘密を広める危険がある。

報告先と開示範囲、確認状況、緊急性を記録して判断する。内部の検討が有効に機能するか、外部への通報が必要かも条件による。過去問で手続不足の告発が不適切とされたことから、告発全般を否定する結論を導かない。法的な保護要件と具体的な通報手続は法令・説明責任の章で別に扱う。

\section{個人から組織へ：倫理規程と学習の仕組み}
\label{theory:ethics:organization}
不正防止を個人の心掛けだけに任せると、納期や売上の圧力、相談の難しさ、記録の不備が残る。記録と報告の経路、専門能力を持つ人の関与、異論を伝えられる環境、類似事例の共有を整える必要がある。対策会議を開催しても、実際の検討を現場へ押し付けるなら組織としての役割を果たさない。

倫理規程は社会への約束、目指す理想、行動の指針、教育、組織自体を点検する対話の基準になる。本文が存在するだけでは実効性を保証せず、具体的な状況に適用し、運用を見直すことが必要である。JSTの教育資料でも倫理綱領と継続的な組織の取組を関連付ける\parencite{fudano_ethics2021}。ここでの社会との約束は、個別の法律上の契約と同一ではない。

データの真正性を守ることは組織の信用の前提である。誤記の訂正や科学的に妥当な処理には根拠と記録が必要で、依頼者が望む結果に合わせた変更とは違う。顧客の要求や上司の指示は、事実を改変して報告する免許ではない。研究固有の記録と不正行為は研究倫理の章も参照する。

\section{予防倫理と志向倫理}
\label{theory:ethics:aspiration}
予防倫理は有害な行為や不正を防ぎ、守るべき境界を示す。志向倫理は、よりよい意思決定や社会への積極的な貢献へ向かう側面である。違反していないことだけでは、技術者として果たし得る役割を全て果たしたとは限らない。両方を認識して行動することが重要である\parencite{fudano_ethics2021}。

\paragraph{出題当時の公務員倫理}
2017年の人事院教材は、倫理法・規程の範囲を狭義、公務員としての在り方を含む範囲を広義として区別する\parencite{jinji_ethics2017}。平成30年度はこの範囲と技術者倫理の二つの側面を対応させて問う。利害関係者は職務との利害関係に基づく対象であり、単に知人という分類ではない。

\paragraph{現在の扱い}
人事院の説明は、利害関係者との贈与・接待やゴルフ・旅行等の規制を解説し、最新の事例集への参照も案内する\parencite{jinji_ethics_qa}。この教材の語句要約を、例外も含む全規制の説明に置き換えない。出題後のJST資料は予防・志向の理解を補う背景資料として使い、平成30年時点の資料だとは扱わない。

\section{社会的責任と共有価値}
\label{theory:ethics:csr}
企業は経済的な成果を生むとともに、従業員、利用者、地域、環境へ影響する。社会的責任の取組を、利益が全く生じない慈善活動だけに限定できない。有償の回収でも廃棄物削減へ貢献し得るし、社会貢献に伴う信用の向上が売上へ寄与することもあり得る。利益を得た事実だけで社会的効果が消えるわけではない。

共有価値の創出は、社会課題の解決と企業の経済的価値を事業の仕組みとして結び付ける構想である。提唱者所属の研究機関はCSRや寄付との概念差も明示する\parencite{hbs_csv}。CSRとCSVを同一の用語にせず、全ての寄付がCSVになるともいわない。過去問の候補では、両者の定義の細部より、利益と社会的責任が両立しないとの誤った限定が問われる場合がある。

\subsection{社会的責任のガイダンスと認証}
要求事項への適合を認証する仕組みと、組織の行動について指針を与える文書を区別する。
ISO26000は社会的責任のガイダンスであり、認証のための規格ではない\parencite{iso26000}。
H23の例では3を公式正答として保持するが、他肢の対象組織や原則列挙の表現には疑義が残る。当時の文言、公開問題、現在確認できる一次資料の範囲を分ける。一つの正答理由を確認しただけで全肢の検証を完了したとは扱わない。

\section{行動規範：専門能力の限界と未来世代への責任}
\label{theory:ethics:code}
電気学会の2007年行動規範は、技術者の能力には限界があるという認識から、多様な見識を持つ人々との協同を重視する。目標は技術進歩それ自体だけでなく、持続可能な社会と未来世代の生存、安全・健康・福祉への貢献である\parencite{ieej_conduct2007}。同じ専門分野・同じ価値観の人だけで判断すると、影響を受ける他者の事情や、技術外の制度・文化の条件が抜ける。

行動規範は具体的な行動の指針、倫理綱領は基本的な約束として読む。平成30年の前文問題を2007年規範へ対応させ、2021年改訂の綱領と同じ時点の文書とは扱わない。可逆性の点検でも、自分が負担を許容できたことを他者全員の同意へ置き換えず、異なる見識との対話を続ける。

\section{ISO26000：組織の影響を社会的責任へ結び付ける}
\label{theory:ethics:iso}
社会的責任は、組織の活動が社会や環境へ及ぼす影響を判断と行動に反映することである。組織は契約相手だけでなく、雇用者、利用者、地域などの利害関係者と関係する。利害関係者は株主と同義ではなく、非営利組織や行政機関にも存在する。ISO26000の対象は広く組織であり、著名人や公人を独立の適用対象として並列に加える定義ではない。個人が組織の活動に担う役割とは分ける。

2010年発行時のISO資料も、ガイダンスであって認証用の要求事項ではないこと、持続可能な発展への貢献を目的とすることを示す\parencite{iso26000_project2010}。対策を実施した結果を検討するために使用する文書と、適合証明のための要求事項を混同しない。

\subsection{行動の原則と扱う主題を分ける}
七原則は、説明責任、透明性、倫理的行動、利害関係者の利害の尊重、法の支配の尊重、国際行動規範の尊重、人権の尊重である。七つの中核主題は、組織統治、人権、労働慣行、環境、公正な事業慣行、消費者課題、コミュニティへの参画と発展であり、同じリストではない。公式の対応資料はこの二層を分けて示す\parencite{iso_gri2014}。

原則は判断の仕方、主題は検討対象を整理する。たとえば労働慣行について従業員の利害を把握し、判断を透明にし、結果への説明責任を持つというように組み合わせる。「法の支配」を権力者の好みに従うことへ置換すると、法によって権力を制約する意味が失われる。政府の意向を独立の尊重原則として付加する肢は、この区別で吟味できる。透明性も、秘密や個人情報を一律に無条件公開することとは区別する。

\section{NPO：非営利と無収入・無報酬は違う}
\label{theory:ethics:npo}
企業の収益事業だけが技術者の社会参加ではない。技術を用いた災害支援や地域での協働などには、非営利組織としての活動もある。特定非営利活動法人は、法定の活動分野と、不特定多数の利益の増進という目的の条件を持つ。2012年4月施行の改正後は20分野となり、平成24年度問題はこの時点の制度を問う\parencite{cabinet_npo_overview}。

非営利とは利益を構成員へ分配する目的を持たないことである。活動を続ける収入や剰余、職員への給与がない意味ではない。主目的の特定非営利活動に支障がない範囲のその他の事業も可能である。利益分配、公益の事業目的、労働への対価は異なる判断軸である。役員として報酬を受ける人数は役員総数の3分の1以下という条件があり、職員の労働への給与はこの役員報酬とは分ける\parencite{cabinet_npo_officers}。

法定分野の存在から、営利企業を競争から保護するためにNPOを制限する制度だとは導けない。会社員であること自体も参加禁止の理由にはならず、就業規則、利益相反、勤務時間、秘密保持などの具体的条件を点検する。

\section{仕事と生活の調和：選択できる環境を整える}
\label{theory:ethics:worklife}
職場の責任と、育児・介護・地域生活などの責任は、人生の段階によって重なり方が変わる。仕事と生活の調和は、その多様性を前提に、働き続けながら選択を実現できる条件を整える考え方である。2007年の憲章と2010年の改定を背景に、平成27年度の問題を読む\parencite{cabinet_worklife}。

柔軟性には勤務時間、場所、休暇、仕事の分担など複数の軸がある。雇用契約を短期にして不安定さを増すこととは同一ではない。短期非正規への置換を必要条件とする記述は、生活を安定させて選択を支える目的との整合を欠く。採用や定着、就業継続の改善は期待効果であり、企業ごとの実現を無条件に保証しない。国・自治体、労使、個人が条件を改善する役割を分担する。

\section{ユニバーサルデザイン：利用者を設計の前提へ入れる}
\label{theory:ethics:universal}
利用方法を一種類の身体能力や言語、経験に合わせると、それ以外の人は利用しづらくなる。設計段階から多様性を扱うには、利用の公平さ、方法の柔軟性、直観的な操作、必要情報の認知、エラーの害の低減、身体的負担、接近と使用の空間という七つの視点を点検する。ノースカロライナ州立大学の1997年版が七原則の原資料であり、現在のウェブやPDFの掲載年と混同しない\parencite{ncsu_ud1997}。

たとえば操作の結果を視覚と音で知らせる案なら、情報を受け取る条件の多様性を検討できる。エラーへの寛容は、誤操作を起こしてよいという指示ではなく、起こり得る誤操作が重い害へ直結しないよう設計することである。全ての人へ常に完全な適合を保証する標語として扱わず、利用者の条件と安全・費用等を評価する。既存の障壁を除くバリアフリーも必要であり、その対策が常に浪費であるとの一般則を設問の前文から導かない。

\section{出題で確かめる条件}
「即時・無条件」「要求があれば」「全て」「だけ」などの限定を確認する。前半が適切でも、後半に不正を認める例外があれば記述全体は適切でない。数を問う場合は各記述を先に判断してから集計し、組合せでは表の列と選択肢番号を区別する。

ここで扱った各過去問は、原典の条件、独立した判断、公式正答との比較、全5肢の理由へ追跡できる。次の問題ブロックはAstra査読前の初稿である。章内の未執筆問題を代表例で置き換えず、全問の処理を継続する。
\section{告発の倫理：段階と証拠、改善への見込み}
\label{theory:ethics:whistle_conditions}
深刻な危害が見込まれるとき、内部で懸念を伝え有効な是正を求める。
上司が応答しないという一事だけで、他の内部経路も調べず直ちに外部公表へ進むとはしない。
客観的な証拠、危害を伝える相手、公表による改善の見込みと個人に及ぶ危険を考える。
これは倫理的な意思決定を整理する枠組みで、法定の公益通報保護要件そのものではない。\cite{muroran_ethics_cases}

令和元年度再試験の比較では、証拠がなくても外部へ知らせる義務が必ずあるという肢を否定する。
一方、危害が急迫している実際の事態に、全ての内部段階を必ず順番に待つという手順を命じるものでもない。
緊急性と有効な通報先、必要な開示範囲、現在の法定経路を調べる。
原著全章を確認した解説ではないため、倫理モデルの解釈は査読事項として保持する。

\section{宣言の三つの責任}
\label{theory:profession:declaration}
2007年のプロフェッション宣言は、自らの技術、顧客、社会に対する責任を示す。能力を高め、正当な守秘を保ちつつ顧客へ誠実に業務を行い、社会と環境への影響を考える\parencite{eth54_profession2007}。顧客だけ、所属企業だけへ責任を狭めない。

倫理綱領は専門職の行動を支える規範である。団体の会員上の適用と、法律・契約・組織規程の適用は別に考える。非会員だから別の義務がすべてなくなるわけではない。反対に、特定団体の内部規範をすべての他団体へそのまま強制できるとも仮定しない。

\section{多様性とDEI：能力を発揮できる条件を整える}
\label{theory:ethics:diversity}
多様性には性別や国籍などの属性だけでなく、知識、経験、専門性、価値観、働き方の違いもある。人の違いを受け入れ、参加の障壁を調べる。違いの尊重は、他者へ危害を与える行動まで無条件に認めることではない。

日本技術士会の2025年宣言は、多様性、公平性、包摂性の三つを結びつける\parencite{eth54_dei2025}。公平性では、一人ひとりの条件に応じて環境や道具を用意する。同じ道具を配るだけでは利用の障壁が残る場合がある。包摂性は、違いを理由に排除せず、能力を発揮して組織の価値を生む状態を目指す。2021年の多様性問題に用いる場合は、後年の補足資料であることを明示する。

\subsection{平成23年度 問11}
\label{q:h23-reg-apt-11}

\noindent\textbf{原典：} sources/H23/aptitude.pdf、PDF物理12頁。

\textbf{公式正答：} 3（single）。

\paragraph{問題の要旨} 社会的責任の手引きと認証規格の区別を問う。公式正答③を保持し、他の記述の解釈上の疑義を別に残す。

\paragraph{必要知識} 社会的責任、要求事項とガイダンス、認証の適用範囲

\paragraph{考え方（査読待ち初稿）} ISO26000は社会的責任に関するガイダンスであり、認証のための要求事項を規定した規格ではない。③がこの性格を誤っているため公式正答は③。ISOの一次資料でも認証用途ではないことを確認した。

\paragraph{他の選択肢（査読待ち初稿）} 他肢をすべて無条件に正しいと説明しない。①の政府除外の表現、④の原則の列挙範囲には一次資料との照合を要する疑義が残る。②⑤の歴史的な参加・連携の詳細も当時資料の確認を後工程に残す。

\par 参考：\cite{iso26000}。

\par 難易度 3（校準初稿）、優先度未評価。

\subsection{平成24年度 問2}
\label{q:h24-reg-apt-02}

\noindent\textbf{原典：} sources/H24/aptitude.pdf、PDF物理1頁。

\textbf{公式正答：} 3（single）。

\par\textit{原典本文の校正と独自解説は未作成。分類のみの登録項目。}

\subsection{平成24年度 問12}
\label{q:h24-reg-apt-12}

\noindent\textbf{原典：} sources/H24/aptitude.pdf、PDF物理10頁。

\textbf{公式正答：} 2（single）。

\par\textit{原典本文の校正と独自解説は未作成。分類のみの登録項目。}

\subsection{平成24年度 問13}
\label{q:h24-reg-apt-13}

\noindent\textbf{原典：} sources/H24/aptitude.pdf、PDF物理11頁。

\textbf{公式正答：} 3（single）。

\par\textit{原典本文の校正と独自解説は未作成。分類のみの登録項目。}

\subsection{平成25年度 問6}
\label{q:h25-reg-apt-06}

\noindent\textbf{原典：} sources/H25/aptitude.pdf、PDF物理7頁。

\textbf{公式正答：} 5（single）。

\par\textit{原典本文の校正と独自解説は未作成。分類のみの登録項目。}

\subsection{平成25年度 問12}
\label{q:h25-reg-apt-12}

\noindent\textbf{原典：} sources/H25/aptitude.pdf、PDF物理13頁。

\textbf{公式正答：} 5（single）。

\par\textit{原典本文の校正と独自解説は未作成。分類のみの登録項目。}

\subsection{平成26年度 問6}
\label{q:h26-reg-apt-06}

\noindent\textbf{原典：} sources/H26/aptitude.pdf、PDF物理6頁。

\textbf{公式正答：} 5（single）。

\par\textit{原典本文の校正と独自解説は未作成。分類のみの登録項目。}

\subsection{平成27年度 問3}
\label{q:h27-reg-apt-03}

\noindent\textbf{原典：} sources/H27/aptitude.pdf、PDF物理3頁。

\textbf{公式正答：} 2（single）。

\par\textit{原典本文の校正と独自解説は未作成。分類のみの登録項目。}

\subsection{平成27年度 問6}
\label{q:h27-reg-apt-06}

\noindent\textbf{原典：} sources/H27/aptitude.pdf、PDF物理6頁。

\textbf{公式正答：} 3（single）。

\par\textit{原典本文の校正と独自解説は未作成。分類のみの登録項目。}

\subsection{平成27年度 問15}
\label{q:h27-reg-apt-15}

\noindent\textbf{原典：} sources/H27/aptitude.pdf、PDF物理16頁。

\textbf{公式正答：} 4（single）。

\par\textit{原典本文の校正と独自解説は未作成。分類のみの登録項目。}

\subsection{平成29年度 問3}
\label{q:h29-reg-apt-03}

\noindent\textbf{原典：} sources/H29/aptitude.pdf、PDF物理3頁。

\textbf{公式正答：} 3（single）。

\par\textit{原典本文の校正と独自解説は未作成。分類のみの登録項目。}

\subsection{平成29年度 問7}
\label{q:h29-reg-apt-07}

\noindent\textbf{原典：} sources/H29/aptitude.pdf、PDF物理7頁。

\textbf{公式正答：} 4（single）。

\par\textit{原典本文の校正と独自解説は未作成。分類のみの登録項目。}

\subsection{平成29年度 問15}
\label{q:h29-reg-apt-15}

\noindent\textbf{原典：} sources/H29/aptitude.pdf、PDF物理14頁。

\textbf{公式正答：} 2（single）。

\par\textit{原典本文の校正と独自解説は未作成。分類のみの登録項目。}

\subsection{平成30年度 問4}
\label{q:h30-reg-apt-04}

\noindent\textbf{原典：} sources/H30/aptitude.pdf、PDF物理4頁。

\textbf{公式正答：} 1（single）。

\par\textit{原典本文の校正と独自解説は未作成。分類のみの登録項目。}

\subsection{平成30年度 問14}
\label{q:h30-reg-apt-14}

\noindent\textbf{原典：} sources/H30/aptitude.pdf、PDF物理14頁。

\textbf{公式正答：} ALL（all\_awarded）。

\par\textit{原典本文の校正と独自解説は未作成。分類のみの登録項目。}

\subsection{平成30年度 問15}
\label{q:h30-reg-apt-15}

\noindent\textbf{原典：} sources/H30/aptitude.pdf、PDF物理15頁。

\textbf{公式正答：} 4（single）。

\par\textit{原典本文の校正と独自解説は未作成。分類のみの登録項目。}

\subsection{令和元年度 問9}
\label{q:r01-reg-apt-09}

\noindent\textbf{原典：} sources/R01/aptitude.pdf、PDF物理10頁。

\textbf{公式正答：} 2（single）。

\par\textit{原典本文の校正と独自解説は未作成。分類のみの登録項目。}

\subsection{令和元年度 問14}
\label{q:r01-reg-apt-14}

\noindent\textbf{原典：} sources/R01/aptitude.pdf、PDF物理14頁。

\textbf{公式正答：} 4（single）。

\par\textit{原典本文の校正と独自解説は未作成。分類のみの登録項目。}

\subsection{令和元年度再試験 問3}
\label{q:r01-ret-apt-03}

\noindent\textbf{原典：} sources/R01\_retest/aptitude.pdf、PDF物理3頁。

\textbf{公式正答：} 1（single）。

\par\textit{原典本文の校正と独自解説は未作成。分類のみの登録項目。}

\subsection{令和2年度 問2}
\label{q:r02-reg-apt-02}

\noindent\textbf{原典：} sources/R02/aptitude.pdf、PDF物理2頁。

\textbf{公式正答：} 2（single）。

\par\textit{原典本文の校正と独自解説は未作成。分類のみの登録項目。}

\subsection{令和2年度 問11}
\label{q:r02-reg-apt-11}

\noindent\textbf{原典：} sources/R02/aptitude.pdf、PDF物理11頁。

\textbf{公式正答：} 1（single）。

\par\textit{原典本文の校正と独自解説は未作成。分類のみの登録項目。}

\subsection{令和2年度 問15}
\label{q:r02-reg-apt-15}

\noindent\textbf{原典：} sources/R02/aptitude.pdf、PDF物理15頁。

\textbf{公式正答：} 1（single）。

\par\textit{原典本文の校正と独自解説は未作成。分類のみの登録項目。}

\subsection{令和5年度 問2}
\label{q:r05-reg-apt-02}

\noindent\textbf{原典：} sources/R05/aptitude.pdf、PDF物理2頁。

\textbf{公式正答：} 3（single）。

\par\textit{原典本文の校正と独自解説は未作成。分類のみの登録項目。}

\subsection{令和5年度 問9}
\label{q:r05-reg-apt-09}

\noindent\textbf{原典：} sources/R05/aptitude.pdf、PDF物理9頁。

\textbf{公式正答：} 4（single）。

\par\textit{原典本文の校正と独自解説は未作成。分類のみの登録項目。}

\subsection{令和5年度 問11}
\label{q:r05-reg-apt-11}

\noindent\textbf{原典：} sources/R05/aptitude.pdf、PDF物理11頁。

\textbf{公式正答：} 2（single）。

\par\textit{原典本文の校正と独自解説は未作成。分類のみの登録項目。}

\subsection{令和7年度 問3}
\label{q:r07-reg-apt-03}

\noindent\textbf{原典：} sources/R07/aptitude.pdf、PDF物理3頁。

\textbf{公式正答：} 5（single）。

\par\textit{原典本文の校正と独自解説は未作成。分類のみの登録項目。}


\chapter{コンプライアンス・説明責任}

\section{背景：専門家の判断を社会の信頼へつなぐ}
技術者は安全、仕様、費用について依頼者より多くの情報を持つ場合がある。製品の使用者や施設周辺の住民は、その判断の結果を受ける一方で、判断の過程を直接見ることが難しい。適切な計算だけでなく、根拠を記録し、他者の権利を尊重し、判断を説明する仕組みが必要になる。

法令遵守は法律に触れない範囲を知るだけではない。誰のどの利益を保護する制度かを理解し、現実の行為を要件に照らすことである。説明責任、競争の公正、知的創作への報酬、安全上の問題の通報、規格への適合は、異なる仕組みを通じて信頼を支える。本章は原典確認した問題の背景初稿であり、未処理の全問題への執筆を継続する。

\section{説明責任：情報を出すことと理解できること}
\label{theory:compliance:accountability}
説明責任とは、自分の判断の根拠、結果、制約について、関係者へ説明し吟味に応じる責任である。報告書を公開した事実だけで果たせるとは限らない。何を測り、何を仮定し、どこに不確実さがあり、どの代替案と比較したかを、相手が検討できる形で伝える。

専門外の公衆に専門家と同じ知識を要求するのでなく、用語・図・尺度を調整して理解を支援する。説明への質問や異なる価値判断を、知識不足という理由だけで退けない。技術士倫理綱領は真実性と秘密の適切な扱いを示し、電気学会行動規範も相手の専門性に応じた誠実な議論を求める\parencite{ipej_ethics2011,ieej_conduct2007}。

説明責任と秘密保持はともに必要である。第三者の個人情報や無関係な技術秘密を全公開するのでなく、必要な情報・範囲・相手を検討する。安全上の重要情報を秘密として隠すことも許されない。危険の程度、緊急性、通報先、開示の必要性を記録して判断する。

\section{合意形成：事実の共有と公正な手続}
\label{theory:compliance:consensus}
合意形成では、事業者、利用者、住民などの利益と価値観が異なる。騒音を何dBと測るかは技術的問いであるが、誰がどの負担を引き受けるかは測定だけで決まらない。事実の確認と利害の調整を分けることが出発点である。

関係者分析により、参加すべき人、懸念、守りたい利益を把握する。進行役は発言機会と議論の手順を整え、技術助言者は専門的な検討を支える。担当技術者は専門能力を持っていても、事業の当事者として利害を持ち得る。中立性と進行能力は別に確保する必要がある。国土交通省の2006年研究と住民参加資料は、第三者の中立性、関係者分析、専門助言と手続支援を扱う\parencite{mlit_mediation2006,mlit_pi_thirdparty}。第三者の参加だけで必ず合意が成立するとは限らず、不同意の理由も記録する。

\section{知的財産：技術、表現、標識、秘密の保護}
\label{theory:compliance:ip}
研究開発や創作には費用がかかるが、成果を得た後の模倣は比較的容易なことがある。制度は創作への誘因と成果の利用を調整する。保護対象が異なるため、全てを特許と呼ぶことはできない。

\begin{center}
\begin{tabular}{p{.19\linewidth}p{.69\linewidth}}
\hline
権利・利益 & 保護対象の基本的な違い\\ \hline
特許 & 自然法則を利用した技術的な発明。新規性・進歩性等は別途審査する。\\
実用新案 & 物品の形状・構造・組合せに関する考案。方法一般の保護ではない。\\
意匠 & 物品等の外観に関するデザイン。具体的な保護対象は法の時点に従う。\\
商標 & 商品・サービスの出所を識別する標識。技術的仕組みとは異なる。\\
著作権 & 思想・感情の創作的な表現。アイデアだけの独占ではない。\\
営業秘密 & 要件を満たして秘密管理する有用な非公知の技術・営業情報。\\ \hline
\end{tabular}
\end{center}
知的財産基本法の範囲は広く、特許・実用新案・意匠・商標の4つが産業財産権である。著作権も知的財産権であるが、この4種類には含まれない\parencite{ip_basic_act,jpo_ip_overview}。平成28年・令和4年・令和5年の設問で、この範囲の基本的な区別を使用する。後の法改正を過去の登録要件へ遡及して当てはめない。

\subsection{自然法則の利用とソフトウェア}
\label{theory:compliance:invention}
自然法則自体の発見、抽象的な数学式、ゲームの約束、情報の単なる提示と、その法則を利用して問題を解く技術的な創作を区別する。特許と実用新案の定義はともに自然法則の利用を含む\parencite{jpo_patent_utility}。

ソフトウェアは記号を処理するから全て対象外、あるいはプログラムだから全て発明、とはいえない。請求項全体が自然法則を利用した技術的な創作であるかを確認する。抽象的な暗号・算法と、それを具体的な情報処理装置で実現する発明も区別する。特許庁審査基準はソフトウェアを含む全体の判断を示す\parencite{jpo_invention_criteria}。発明該当性と、新規性・進歩性等の充足は別の段階である。

\subsection{営業秘密の3要件}
\label{theory:compliance:secret}
秘密管理性、有用性、非公知性を満たす情報が営業秘密である\parencite{meti_trade_secret}。社内情報という名称だけで足りず、アクセス制御や表示等により秘密として管理する。特許のように公開と登録による独占を得る制度とも異なる。知的財産の創作物に含むかを問う問題と、産業財産権4種類に含むかを問う問題で、対象範囲を取り違えない。

\section{著作権の制限：許諾なしの利用にも要件がある}
\label{theory:compliance:citation}
引用は、自分の検討・説明に他者の表現を必要な範囲で取り込む利用である。公表された著作物について、公正な慣行に合致し、引用の目的上正当な範囲を満たすことを確認し、出所を示す。自分の記述と引用部分を区別し、利用目的・量・方法・影響を検討する。出典を付けた事実だけで長大な転載が適法になるわけではない\parencite{bunka_citation_faq}。

平成27年の設問でも引用を学術研究だけへ限定する説明は誤りである。適法な引用の範囲では翻訳利用も可能であるが、翻訳と自由な翻案・改変を同一視しない\parencite{kyushu_copyright2015}。公表は権利放棄を意味しない。

官公資料には、32条2項の別の仕組みもある。公的機関が周知のため自らの名義で公表する広報・調査統計・報告等を説明材料として刊行物へ転載する場合で、禁止表示などの条件を確認する\parencite{cabinet_copyright2008}。官公庁サイトにある全ての著作物が無条件に対象とはしない。

\subsection{非享受目的の利用とAI}
\label{theory:compliance:nonenjoyment}
通常の鑑賞・読書は表現された思想や感情を享受する。一方、複製機器の性能試験や機械的な情報解析は、表現の鑑賞とは異なる目的で利用する場合がある。2018年改正・2019年施行の30条の4は、この非享受目的について必要な限度で利用を認め、不当な利益侵害を除外する\parencite{bunka_amend2018}。

AI開発の学習データ収集や、その開発を行う第三者への提供が該当し得る。ただし既存作品の表現を出力させる目的が併存する利用や、解析用データベースの販売市場に衝突する利用などは要件を別に検討する。学習段階と生成物の利用段階の権利問題も別である\parencite{bunka_ai2024,bunka_text2026}。令和6年設問の例示を、AI利用一般の包括的な免許へ拡張しない。

\section{公正取引：独立した競争と市場情報の公正}
\label{theory:compliance:competition}
価格・数量を本来各事業者が独立に決める競争で、競争相手同士が共同決定すると、利用者は競争の便益を失う。価格・数量等のカルテル、入札の受注者・金額を事前に決める入札談合を、行為から識別する\parencite{jftc_rules}。入札談合は法的にはカルテルの典型例でもあり、用語の区別は概念が互いに完全に排他的という意味ではない\parencite{jftc_bid_cartel2021}。

\subsection{公共契約の例外と官製談合}
\label{theory:compliance:procurement}
国の契約は一般競争を原則とする一方、競争を許さない性質・目的、緊急等の法定条件では随意契約を認める。例外を適用する理由を具体的に説明し、事務の遅延だけを緊急性へ置き換えない。2006年の財務省通知もこの点を明確にする\parencite{mof_procurement2006}。

談合への行政措置、刑事責任、発注職員の関与に対する制度を区別する。官製談合防止法の職員の入札等妨害への刑事罰は2007年に施行され、平成27年出題時にも存在した\parencite{jftc_kansei}。各責任の構成要件は別に確認する。民間が発注するから競争法の対象外になることもなく、当時の公正取引委員会資料にも民間受注調整の処理実績がある\parencite{jftc_private2018}。

\subsection{金融市場の2種類の問題}
\label{theory:compliance:market}
インサイダー取引は、会社関係者等が法定の重要事実を職務等から知り、公表前に所定の取引をするという情報の公正の問題である。日常語の重要情報全般と法定の重要事実は同一ではなく、対象者・証券・公表・適用除外にも要件がある\parencite{fsa_insider_rules}。

相場操縦は、人為的に作った価格等を自然な需給の結果のように見せて他者の取引を誘引するなど、市場の価格形成を損なう問題である。大きな売買により価格が変わった事実だけで相場操縦と決まらない\parencite{fsa_market2004}。令和4・7年の問題では、競争事業者の共同決定、入札調整、非公開情報、偽った価格形成をそれぞれの名称へ対応させる。

\section{公益通報：問題を知らせる人を守る制度}
\label{theory:compliance:whistleblower}
事業者内部の違反を知る人が報復を恐れて沈黙すれば、事故や消費者被害は拡大し得る。公益通報制度は、対象事実を知らせる人の保護と法令遵守を結び付ける。単なる不満の表明や不正目的の告発、全ての機密の公開を無条件に保護する制度ではない。

\paragraph{出題当時：2018年}
公務員も原則として労働者に含む。内部通報、権限ある行政機関、その他外部の通報先によって要件が異なる。その他外部では相当の理由、被害防止に必要な通報先、報復や証拠隠滅のおそれ・内部調査が行われない・生命身体への急迫危険等の追加事情を確認する。追加事情は全ての同時成立ではなく法定の選択的条件である。改正前の消費者庁ハンドブックにより当時の範囲を確認する\parencite{caa_handbook_pre2022}。内部へ通報しない限り常に保護されないと一般化しない。

\paragraph{現在：2026年10月10日}
2020年改正は2022年6月1日から施行され、退職後1年以内の元労働者や役員等、対象・要件・体制を更新した。通報者の立場による要件の違いも確認する\parencite{caa_amend2020,caa_current_faq,caa_report_current}。300人以下の事業者への体制整備は努力義務であり、法全体が小規模事業者へ適用されないという意味ではない。

\paragraph{施行予定の改正}
2025年改正の主要部分の施行日は2026年12月1日である\parencite{caa_schedule2026}。この教材の確認日では将来の制度なので、施行中の制度や平成30年の正誤へ遡及して使わない。通報先と必要な情報を調べ、手続・確認状況を記録することは、倫理の情報開示の議論と接続する。

\paragraph{省略された条件を全リストにしない}
平成24年度の外部通報の4記述は、旧制度の報復、証拠隠滅、正当な理由のない通報禁止要求、
書面通報後20日以内の調査通知の欠如に対応する。
対象事実や相当の理由等も前提となり、4文だけで全ての要件を尽くすとはしない。\cite{saf53_whistle_oldfaq}
平成27年度の別表の略記は政令の指定を含めて読み、外部通報条件の代表列挙も排他的な全条件とはしない。
公益通報でも、他人の正当な利益や公共の利益を害さないよう配慮する。
第三者の個人情報を無制限に拡散する権利が生じるわけではない。\cite{caa_handbook_pre2022}

\paragraph{原典の改正年表現と対象者}
令和5・7年度の原典は2022年又は令和4年改正と記すが、2020年改正の2022年施行と区別する。
公務員も対象に含まれる。公務員法が適用されるから全て除外という肢を否定する。\cite{saf53_whistle_oldpublic,caa_amend2020}
300人以下の体制整備の努力義務と、通報者の法定保護は別である。
改正の成立、公布、施行、資料の確認日を一つの日付へまとめない。

\section{標準化と認証：規格を決めることと適合を示すこと}
\label{theory:compliance:standards}
規格は寸法、性能、試験方法等の共通の約束を与え、互換性と比較可能性を支える。ISOは産業分野、IECは電気・電子技術分野を主に扱う。JISCの基本説明ではISOの対象から電気・通信・電子技術を除くと整理するが、情報技術等には共同の標準化もある\parencite{jisc_isoiec}。

規格の制定と、具体的な製品等がそれに合うことの認証は別である。JISマークは国登録の第三者機関による審査・認証を受けて表示する。定期審査は少なくとも3年ごと、必要に応じ臨時審査や国の立入検査がある\parencite{jisc_certification}。マークを描くだけで適合性を証明することはできない。

\paragraph{変更の年を分ける}
新JISマーク制度は2005年10月に始まり、旧表示の経過措置は2008年9月までである\parencite{mhlw_jis_history}。2019年7月の改正では規格名が日本工業規格から日本産業規格、法名が工業標準化法から産業標準化法へ変わり、データ・サービス等へ対象が拡大した\parencite{meti_jis2019}。2019年をマークの新旧交替年としない。

令和6年設問の左右の説明は実図と逆である。現行の横並びJIS図はJISCの公式画像で確認できる\parencite{jisc_marks}。図の意味は文字の説明だけで決めず、原典そのものと独立資料を照合する。

\section{出題の読み方と本章の範囲}
正誤表は、各文の判断を先に書いてから行へ対応させる。数の問題は正しい文の数と選択肢番号を区別する。引用と複製、知的財産権と産業財産権、契約方式と競争制限、ISO等の規格と製品認証、改正の成立と施行を切り分ける。制度の専門的な適用には個別条件があり、ここでの試験要旨を例外を含む全規定の説明へ置き換えない。

以下はAstra査読前の独自解説初稿である。未執筆問題の一覧も残し、全問の確認を継続する。
\section{公衆と秘密保持：情報を伝える目的と範囲}
\label{theory:compliance:public_information}
憲法12条は権利の保持と濫用を避ける責任を、民法1条は私権と公共の福祉の適合を定める\parencite{jp_constitution,jp_civil}。技術者倫理では、技術の結果を受けても、その専門的内容を知って同意し又は回避する機会が乏しい人々への責任が重要になる。知識の非対称性は、個人の能力一般を低く評価する根拠ではない。

説明責任は情報を正確に伝え、根拠や不確実さへ応答する責任である。必ず相手が納得するという結果を強制できるわけではない。国家公務員法100条の職務上の秘密保持は退職後も続き、法に定める発表の許可等も区別する\parencite{jp_national_public}。説明責任と秘密保持をどちらか一語だけで決めず、情報の性質、法的根拠、開示先と目的を検討する。

\subsection{通報を是正へ結び付ける}
不正の根拠を検証し、相談・通報手段を活用し、適切な場合は上司へ異議を伝え、是正案を示す。上司の反応だけで終わらず他の理解者や経路を探すことも必要になる。専門職の支持を求めるのは事実を検討するためで、結論に合わせて口裏を合わせることではない。

根拠確認の努力を違法な証拠収集の許可へ広げない。内部の経路を活用するという倫理的助言を、外部通報の保護には必ず内部通報を先行しなければならないという法的条件へ変換しない。公益通報の保護は対象・目的・通報先等の要件で判断する。出題当時から、まさに生じようとしている対象事実を含む点は、後年の対象者拡大とは別である\parencite{caa_handbook_pre2022,caa_current_faq}。
\section{知的財産を分類するための二つの問い}
\label{theory:ip:scope}
知的財産の問題では、まず保護される対象を問い、次にどの制度で保護するかを問う。技術的な解決方法、文章や画像の表現、商品の出所を示す標識、秘密の製造情報は、同じ情報でも保護の根拠が異なる。情報は多人数が同時に利用できるため、物の占有と同じ方法で利用を制御できない。制度は創作・研究開発への誘因と成果の利用を調整する。

特許庁が産業財産権と総称するのは、特許権、実用新案権、意匠権、商標権の4種類である\parencite{ip_industrial_four}。著作権、回路配置利用権、育成者権は知的財産に関係する権利であっても、この4種類に含まれない。営業秘密、商号、商品等表示の保護も広い知的財産の体系に位置付く\parencite{ip_basic_act,ip_broad_scope}。広い体系へ含めることは、全てが登録により絶対的な独占権を生むという意味ではない。

\subsection{集合の補集合と選択肢番号}
\label{theory:ip:count}
列挙された対象の集合を \(L\)、4種類の集合を \(I\) とすると、産業財産権に含まれない対象は \(L\setminus I\) である。平成30年問6は広い知的財産の範囲外の数、令和2年問5は産業財産権の範囲外の数を問う。同じ名称が並んでも、引く集合が違えば数が変わる。対象数を求めた後で、その数を示す選択肢番号に対応させる。

\section{特許：定義と権利取得の要件を分ける}
\label{theory:ip:patent}
自然法則を利用した技術的思想の創作のうち高度なものが、特許法上の発明である。発明に該当しても、直ちに特許を得るわけではない。産業上利用できること、出願前の公知・公用・刊行物等に照らした新規性、進歩性などを別に審査する。公の秩序、善良の風俗又は公衆の衛生を害するおそれがある発明は保護対象から除かれる\parencite{ip_patent_statute}。

\subsection{出願前公開と登録後公開の区別}
\label{theory:ip:novelty}
制度が技術の公開と利用を促すことと、出願前に既に公表されていることを取得要件にすることは別である。通常、出願前公開は新規性に不利に働く。ただし公開された発明が例外なく永久に特許を得られない、とも断定しない。新規性喪失の例外には対象行為、期限、手続の条件がある。

平成30年改正では例外期間が6か月から1年へ延長され、2018年6月9日以後の出願への適用に公開日の経過措置も置かれた\parencite{ip_novelty_grace2018}。令和元年再試験の判断は、この当時すでに例外が存在したことと整合させる。例外は「刊行物で発表済みでなければならない」という要件を作らない。

\section{著作権：表現、創作性、権利の発生}
\label{theory:ip:copyright}
著作物は思想又は感情を創作的に表現したものである。学術的な事実や抽象的なアイデアそのものと、その説明文・図の具体的な表現を区別する。創作性は芸術的な傑作という評価を要求するものではないが、法の定義から削除されたわけでもない。

著作者人格権と著作権の享有に申請や審査を要求しない仕組みを無方式主義という\parencite{ip_copyright_statute}。特許庁へ登録しなければ著作権が発生しないという説明は、産業財産権の手続との混同である。権利の発生と、一定事項の登録制度が存在することも別である。

\subsection{人格的利益と財産的利用}
\label{theory:ip:personality}
公表権、氏名表示権、同一性保持権は著作者の人格的利益に関わる。複製、公衆送信などの利用を管理する財産的な権利と区別する。歌手が歌唱する立場にあることだけで、他者が作った歌詞を自由に変更する権限は生じない。利用許諾があることも、全ての改変への同意があることを当然には意味しない。

\subsection{個人の趣味と公開送信}
\label{theory:ip:public_web}
個人的な趣味であっても、公開Webページへの画像掲載は私的な複製と同じ範囲ではない。写真集の画像について、許諾や引用等の適法な根拠を欠く掲載を、趣味という目的だけで正当化しない。写真の権利と被写体に関わる別の権利も分けて考える。

\section{引用：出典表示だけでは条件を満たさない}
\label{theory:ip:quotation}
32条1項の引用は、公表された著作物を、公正な慣行と引用目的上正当な範囲に従って利用する仕組みである\parencite{ip_copyright_statute}。自分の論述に必要な部分を区別し、利用量と目的の関係を説明できることが重要になる。48条の出所明示も確認するが、出典を書けばどの量でも転載できるという制度ではない。

権利制限を満たす引用は、保護期間内であっても許諾なしにできる。保護期間満了後の利用、権利者の許諾に基づく利用、法が認める引用は異なる根拠である。「無断なら例外なく違法」という断定を避ける。

\subsection{官公資料の範囲と禁止表示}
\label{theory:ip:government}
法令など13条に掲げる対象が著作権の対象から除かれることと、官公庁が作成した全資料を自由に使えることは別である。32条2項は、国等が一般への周知を目的に自らの名義で公表した広報・調査統計・報告等を、説明材料として新聞紙、雑誌その他の刊行物へ転載する仕組みを定めるが、転載禁止の表示がある場合を除く。対象、目的、媒体、禁止表示を確認する\parencite{ip_copyright_statute}。

\subsection{引用のための翻訳と自由な翻案}
\label{theory:ip:translation}
適法な引用に対応する範囲では翻訳利用が認められ、外国語の原文をそのまま使うことだけが唯一の方法ではない。これは他者の文章を自由に翻案・改変してよいという意味ではない。取得した現在の条文では47条の6が翻訳等を扱う。平成25・26年度の問題について当時の条番号を現在番号に置き換えた証跡とはしない。二問の当時版全文の照合は未完として残す。

\section{営業秘密：情報の種類と法的成立を分ける}
\label{theory:ip:secret}
不正競争防止法2条6項の営業秘密は、秘密管理性、有用性、非公知性を同時に満たす情報である\parencite{ip_unfair_statute}。顧客名簿や新規事業計画はその典型になり得るが、資料名や社内で作成された事情だけで営業秘密になるわけではない。秘密として管理され、事業に有用で、公然と知られていないことが必要である。

\subsection{有用性と現実の使用}
\label{theory:ip:usefulness}
情報の有用性は、現に製品へ使われていることだけでは決まらない。将来の事業計画、失敗した実験の記録も、不要な研究費を避けるなどの価値を持ち得る。一方、脱税や有害物質の垂れ流し等の反社会的な活動の情報は、法が保護する正当な事業活動の有用性とは扱われない\parencite{ip_secret_handbook2017,ip_secret_commentary2018}。違法行為の情報を隠すことと、適法な技術情報を守ることを分ける。

\subsection{公開、秘密管理、漏洩経路}
\label{theory:ip:secrecy}
刊行物や公開された特許に記載された内容は、そのまま非公知の秘密とはできない。そこに記載されていない別のノウハウまで一緒に公開されたとは限らない。管理表示やアクセス制御は、対象者が秘密と認識できる管理の具体化である。漏洩経路には従業員、退職者、取引先、共同研究先、サイバー攻撃がある。経路の存在の確認と、特定年に急増したという量的主張の確認は別である。

\subsection{令和2年と令和3年の同じ省略表現}
\label{theory:ip:secret_ambiguity}
両年の設問は、顧客名簿等について、企業活動の過程で生まれたので営業秘密であると述べる。公式表は典型的な該当例として真とする組合せである。一方、三要件を省略しない一般命題として読むと、作成された由来だけでは成立を保証できない。令和3年の独立初回判断はこの読みから4、公式正答は5であった。

本教材は公式値を保持するとともに、初回判断と省略条件を残す。公式と一致するよう法の要件を変更しない。令和2年も同じ条件省略があるため、公式値と一致した事実だけで曖昧さを解消済みにしない。

\section{非享受目的：利用の目的から権利制限を判断する}
\label{theory:ip:nonenjoyment}
30条の4は、表現された思想又は感情を自ら享受し又は他人に享受させることを目的としない場合に、必要な限度で利用を認める。ただし著作権者の利益を不当に害する場合を除く。技術試験、情報解析、人の知覚によらない計算処理は例示であり、「AIだから許される」という機器名の判断ではない\parencite{ip_flexible2019,ip_copyright_statute}。

\subsection{六つの例とそれぞれの境界}
\label{theory:ip:six_cases}
2019年文化庁資料は、AI学習用の収集・第三者提供、プログラムの解析、複製機器や和紙の性能試験、送り仮名等の調査、場所の3DCG化、全文検索用の複製を具体例に挙げる\parencite{ip_flexible2019}。それぞれ表現の鑑賞ではなく、解析・性能・事実情報等を目的とする例である。

AI学習では、既存作品の表現を享受させる目的の併存や、解析用に販売されるデータベースの市場への影響を別に確認する。性能試験の複製物を鑑賞用に頒布する行為まで認めることにはならない。写真から場所を再現する場合も、元写真の創作的な表現を3DCGで承継して享受させる場合は区別する。

検索のための内部複製と検索結果の表示も別である。書名・著者名等の表示を目的とする例と、本文の一部を表示する利用についての47条の5の条件を混同しない。プログラム解析が著作権法上許される場合でも、契約など別の問題が全て解決するわけではない。

\section{法の時点を保った学習の手順}
\label{theory:ip:historical}
4種類への分類、引用の条件、三要件の論理は共通に参照し、保護対象の拡大や条番号、例外の改正は出題時点を確認する。非享受目的の2019年施行の規定を、それ以前の問題にそのまま遡及しない。現在の条文の説明、当時の行政資料で確認した範囲、未取得の過去版全文を区別する。

正誤表では各文の判断を先に定め、原典の列順で5候補と照合する。否定、数、組合せの違いを保ち、同じ公式番号が続くことを同じ答えの内容とみなさない。

\section{雇用と職場：専門能力を発揮できる条件を守る}
\label{theory:employment:background}
技術者は安全な製品を作る責任だけでなく、設計・製造・点検を担う人の尊厳と健康にも関わる。威圧によって異常を報告できない職場では、技術上の問題を発見しても対処へ結び付かない。長時間労働や特定個人への業務集中も判断・継承を損なう。雇用法令と職場の仕組みを、専門能力が継続して発揮されるための条件として理解する。

\subsection{労働法：主体と義務を先に読む}
\label{theory:employment:law}
伝統的な労働三法は労働基準法・労働組合法・労働関係調整法である\parencite{employment_labor_three}。最低条件の保護、集団的な交渉、紛争の調整という役割を区別する。労働安全衛生法は労働者の安全と健康を扱い、規模等に応じた産業医の選任を求める\parencite{employment_industrial_physician}。技術士資格があるだけで、労働者を使う際の義務が免除されるわけではない。

36協定は法定の時間外・休日労働に関する労使協定である。締結と届出を区別し、届出だけで上限を自由に撤廃できるものと読まない\parencite{employment_overtime}。以下では出題当時の記述を保ち、2026年10月12日時点で確認した変更は別に述べる。全ての告示の改正履歴を照合し終えたという意味ではない。

\subsection{パワハラ：職位より関係、叱責より態様}
\label{theory:employment:power}
平成24年WG報告は、職場内の優位性を背景に、適正な業務範囲を超えて苦痛又は環境悪化を生む行為を整理した。上司から部下だけでなく、知識や協力への依存による同僚・部下の優位性も含む。身体攻撃・精神攻撃・隔離・過大要求・過小要求・私生活への過度な介入は典型例であり、網羅的な六項目の暗記だけでは足りない\parencite{employment_power2012}。

検討では、目的に必要な指示か、方法が相当か、相手と周囲の環境へどう影響するかを分ける。手順の誤りを具体的に指摘することと、人前で人格を否定することは同じでない。主観的不満だけで該当性は決まらないが、相手の苦痛を調べずに「指導だから問題なし」としてもいけない。現在の指針も三要素と客観的な適正指導の区別を定める\parencite{employment_power_current}。

\subsection{セクハラ：不利益と環境悪化}
\label{theory:employment:sexual}
性的言動への対応で処遇上の不利益を受ける対価型と、言動によって仕事の環境が害される環境型を区別する。同性間でも生じ、労働者の雇用形態や行為者の所属だけで保護を除外できない。派遣労働者には派遣元・派遣先の関係を確認する\parencite{employment_sexual_current}。

拒絶を口にできない関係や、問題化すると仕事を失うと感じる状況があるため、沈黙だけから同意を推定しない。相談を受けた組織は事実を確認し、情報の扱いと相談者への不利益を防ぐ体制を整える。2018年研修資料も措置義務を説明しており、単なる努力目標という過去問の肢は誤る\parencite{employment_sexual2018}。

\subsection{各類型の範囲を同一にしない}
\label{theory:employment:scope}
日常語のハラスメントは広い。セクハラ、パワハラ、妊娠等に関するハラスメントは法的な条件が同じではない。顧客からの性的言動もセクハラの行為者範囲に入るが、その事実だけで顧客との全ての紛争を同じ職場内のパワハラ定義へ入れない。古いWG報告の顧客例は、職場内の定義とは別の脚注である。

性的指向や性自認を侮辱したり、本人の同意なく周囲へ伝えたりすることは、人格と機微情報に関わる。技術指導に必要な行為かを確認する。過去問の「本人はもとより職場全体のハラスメントとなり得る」は、周囲にも就業環境の悪化が及び得るという判断であり、全員が必ず同一の法的被害者となるという意味ではない。

\subsection{産休と妊婦への配慮}
\label{theory:employment:maternity}
産前産後休業は母体保護を目的とし、育児休業とは根拠と条件が異なる。産前は原則出産予定前6週間、多胎14週間で本人の請求が条件である。産後は原則8週間の就業禁止で、6週間経過後に本人が請求し医師が支障なしと認める業務には例外がある。正規・有期・派遣といった雇用形態で一律に除外しない\parencite{employment_maternity2014,employment_maternity_current}。

妊娠を理由に仕事から一方的に排除することと、健康や業務状況を考えて本人へ調整を相談することは異なる。配慮という名称だけで不利益を正当化せず、目的・本人の意向・法定措置を確認する。2018年の公的研修例も、残業の負担を相談して減らす配慮と、権利利用を妨げる強要を分けている。

\subsection{労働時間：場所と自己管理は免除理由でない}
\label{theory:employment:worktime}
使用者の指揮命令下の時間を考える。命令された業務上の学習は、学習という名称でも労働時間に含まれる。使用者は時間を適切に把握し、自己申告を用いる場合にも記録との食違いを確認する\parencite{employment_hours2017}。

管理監督者は肩書だけで決まらず、職務・権限・勤務の裁量・待遇等の実態から判断する。該当する場合も深夜割増賃金や年次有給休暇まで一律に除外されるものではない\parencite{employment_manager2008}。勤務終了後に必要な報告や研修を行わせながら記録から除外することは、健康管理と制度の趣旨を損なう。

\subsection{裁量労働：資格名と業務の列挙を分ける}
\label{theory:employment:discretion}
専門業務型裁量労働制は、法令で指定する業務の中から労使協定で対象を定め、定めた時間働いたものとみなす仕組みである。平成25年当時は19業務であり、保存した平成26年3月資料にも同じ19業務が並ぶ。この一覧に「技術士の業務」はない\parencite{employment_discretion2014}。

技術士が研究開発など対象業務に従事する可能性と、技術士資格を持てば自動的に対象になる命題は異なる。この混同で平成25年度問4の初回解答を誤ったため、初回3と資料確認後の訂正5を別々に保存した。適用を確かめるには肩書から推測せず、実際の従事業務と指定範囲を照合する。

\paragraph{現在の変更}2024年4月から対象は20業務で、本人同意・同意の撤回等の要件が加わった。現行一覧にも技術士業務という列挙はない\parencite{employment_discretion_current}。みなし時間は健康管理のための勤務状況把握を不要にするものではない。

\subsection{フレックスとテレワーク}
\label{theory:employment:flex}
フレックスは一定の清算期間の総労働時間を定め、その枠内で始業・終業を労働者が選択する仕組みである。期間の上限は2019年4月の改正で1か月から3か月へ延長された\parencite{employment_reform2018}。平成29年度の説明へ3か月を遡及しない。柔軟な時間配分と、総労働時間の把握を両立する。

\label{theory:employment:telework}
テレワークは場所の選択を広げるが、雇用型では労働法令が適用される。勤務場所を示し、時間を把握し、中抜けで自由利用が保障される時間と、移動中も命令された仕事をする時間を分ける。平成30年指針の概要を載せた令和2年白書でもこの区別を確認できる\parencite{employment_telework2020}。時間単位年休には別の制度要件があり、中抜けしたことだけで勝手に年休へ振り替えない。

現在の指針は記録・申告の方法、休止と移動、長時間化対策を整理する\parencite{employment_telework_current}。通信できる環境は常時応答の義務と同じでない。自己管理という説明によって、組織側の時間把握や負荷調整を放棄しない。

\subsection{働き方改革：上限・割増・休暇を分ける}
\label{theory:employment:reform}
法定の時間外労働上限は使用者の指示だけで緩和できない。年10日以上の年休が付与される労働者について、年5日の取得を確保する。本人の請求による取得や計画付与の取得分を含むため、全員へ別枠の5日を追加指定する説明は不正確である\parencite{employment_paidleave}。

月60時間超の時間外労働は割増賃金の問題であり、10時間ごとに代休1日という制度ではない。代替休暇には労使協定と本人意思を要し、引き上げた割増分に対応して換算する仕組みである\parencite{employment_alternative2008}。法定休日に働いた後の代休、事前に休日を振り替える操作、割増分の代替休暇は別の概念である。

\subsection{仕事と生活の調和：属人化と不要作業を減らす}
\label{theory:employment:worklife}
業務を標準化して引き継げる形にし、資料の過剰作成や探索を減らし、予定と負荷を共有する。これらは同じ仕事を少人数へ集中させるためでなく、時間と相互支援を確保する組織改善である。内閣府の十の実践は会議、資料、整理、標準化、時間、分担、相互理解、予定、集中、改善共有を扱う\parencite{employment_worklife2010}。

会議の到達点を明らかにする場合、技術上の証拠が不十分なら追加調査と責任者・期限を決めることも到達点になる。「必ず結論」という標語を、判断の根拠を省略して安全性を断定する命令へ変えない。効率化の取組後も、業務の質、負荷の偏り、実際の労働時間を観察し、改善の効果を確認する。

\section{安全保障貿易管理：仕様・相手・用途を確認する}
\label{theory:export:screening}
同じ測定器が民生研究にも軍事目的にも役立つ場合、買い手の肩書だけで用途は分からない。安全保障貿易管理は、軍事転用可能な貨物や技術が懸念活動へ渡ることを防ぐため、外為法等に基づいて行う\parencite{export_overview}。

まず貨物や技術の仕様を法令のリストへ照合する該非判定を行う。次に、相手・需要者・用途・仕向地等から取引を審査する。リスト非該当でも、補完的なキャッチオール規制の条件に該当する場合がある。最後に、審査した物と発送する物の同一性、許可の取得と内容を出荷時に確認する。少量サンプル、国内で買える製品、自作装置、手荷物での持出しというだけで、この確認を省けない。

メーカーの判定書は重要な資料だが、古い判定書が現在の型式・仕様・法令に合うかを輸出者も確かめる。信頼する取引先からの紹介も、新たな需要者の用途確認を代替しない。営業・技術・管理・出荷の各部門で判断根拠を共有することで、適切な審査を実際の発送へ結び付ける。

\section{技術提供と例外：場所より情報の内容を調べる}
\label{theory:export:exceptions}
技術には設計・製造・使用に必要な情報を、データや助言として提供する場合がある。国内の非居住者への提供も対象となり得るので、日本国内であることだけで規制外と判断しない\parencite{export_scope2022}。大学の非公開な共同研究打合せで技術を交換することと、不特定多数が入手できる公知技術の発表は違う。

装置の部分品特例も、部品が組み込まれているという一事だけで使えるわけではない。公的事例集は、ロボットに内蔵されたセンサーが主要要素でないことを学内の担当部署が確認する手順を示す\parencite{export_cases}。特例の対象と仕様・価格等の条件を検討し、本人だけの思い込みで適用しない。特許出願等に必要な最小限の技術提供についての例外も、手続目的を越えるノウハウ提供まで含めない。当時条文の直接照合は初稿の留保事項である。

\begin{tcolorbox}[title=出題時点と現在の規則]
ここでの過去問は2021--2023年の事例である。経済産業省の現在の案内には2025年10月9日施行の補完的規制見直しが含まれる\parencite{export_current_updates}。現在の品目・対象国・用途・通知条件を過去問へ補って判定しない。取引の実施では、実施時点の規則を確認する。
\end{tcolorbox}

\section{専門職・継続管理・要求性能を別の軸に置く}\label{theory:compliance:management}
高度な専門能力を持つ人材の養成を目的とする専門職大学院の課程は、
博士後期課程そのものではない。専門職の学位と研究の博士課程を区別する。
法務博士という学位名の「博士」だけから博士後期課程と同一視しない。\cite{eng52_professional}

BCMは障害に備える計画を作り、訓練・点検・見直しを続けて事業継続能力を管理する活動、
BCPはその事業継続計画である。計画書を置くだけで管理の全過程が完了するわけではない。
当時の理解は平成21年版等の資料から確認し、後年の制度改訂を遡って加えない。\cite{eng52_bcp}

性能の要求は、果たすべき機能・安全・使用性等を定める。
製作仕様は材料・寸法・方法等を指定する手段であり、両者を同じ定義にしない。
性能規定で要求の満足を検証することと、仕様を守ればよいという評価を区別する。
現実の要求水準には性能と仕様が併存する場合もあるので、一方だけが常に許されるとはしない。
信頼性は必要な機能を条件・時間の中で遂行する度合いを確率的に扱う別の軸である。\cite{eng52_performance}

\subsection{VEの価値と品質の最低条件}
VEは必要な機能を確実に果たし、ライフサイクルを含む資源・コストと機能を比較して価値を改善する。
説明用の比\(V=F/C\)では、分母を下げるだけで機能の必要水準を失う変更を価値向上とはしない。
安い価格と、設計・保守・更新を含む全体コストを分ける。\cite{eng52_ve}

\subsection{内部統制の四目的と時点}
平成19年の基準では、業務の有効性・効率性、財務報告の信頼性、
事業に関わる法令等の遵守、資産の保全という四目的が示される。
組織の全員が業務の中で遂行するプロセスで、互いに関連する目的である。
経営者の整備責任と、その他の者が運用に関与することは両立する。
内部統制全体の目的と、金融商品取引法の財務報告に係る評価・報告・監査の対象を分ける。\cite{eng52_control}
平成23年度の表現を、後年の報告範囲の変更後の用語に無断置換しない。

\section{食品表示：どの情報が選択と安全を支えるか}
\label{theory:compliance:foodlabel}
期限、原産国・原料原産地、成分・アレルギー、事業者情報は、それぞれ別の判断を支える。
賞味期限は定めた保存条件のもとで期待される品質が保たれる期間、
消費期限は安全性を欠くおそれがないと認められる期間である。
いずれも保存条件が前提で、開封後や不適切な温度に置いた食品の無条件保証ではない。
賞味期限を過ぎたら直ちに全て食べられない、という安全期限との混同を避ける。\cite{saf53_food_expiry,saf53_food_loss2013}

輸入加工食品の最終加工国と、国内で加工した食品に使う原料の産地を分ける。
平成25年度の品目・原料の条件を、後年の広い対象へ変えない。
食品表示法の公布2013年と施行2015年も区別する。
アレルギー表示の指定は見直されるため、当時の指定品目を現在の数に書き換えない。\cite{saf53_food_notice}
表示は選択だけでなく、事故時の対象製品と責任を追跡する情報でもある。

\section{不実証広告：表示を裏付ける二つの条件}
\label{theory:compliance:adevidence}
効果を示す資料があるという事実だけでは、広告の主張は確かめられない。
資料自体の客観性と、表示された対象・条件・効果への対応の両方が必要である。
別の材料の試験や特殊条件の結果を、製品の一般的な利用へそのまま拡張しない。
行政が資料を求めても自由に断れるという制度ではなく、所定の不提出には制度上の効果がある。\cite{saf53_ad_overview}

自社試験でも妥当な方法で実証できるが、実施者が専門家であるだけでは足りない。
専門分野で一般に認められない特異な見解や、新しい技術を称することによって
試験・調査の実証を免れるものではない。
経験則を根拠とする場合も、存在を裏付ける専門的な見解等を確認する。\cite{saf53_ad2016}
2016年の説明では資料提出は原則15日で、新しい試験が必要という理由だけで延長できるとはしない。
2014年の設問へ現行条番号や後年の処分制度を補わず、当時の枠組みを保存する。

\section{PDCA：点検した結果を改善へ戻す}
\label{theory:management:pdca}
計画し、実施し、結果を点検し、是正と改善へ反映する順序がPDCAである。目標と評価方法を定め、実行の記録を取り、目標との差と原因を調べて次の計画へつなぐ\parencite{eth54_pdca2017}。

Checkは単に点検作業をしたという事実ではなく、計画に照らして結果と過程を評価する段階である。Actも行動一般ではなく、その結果を踏まえる是正・改善を表す。研究開発の評価は次の研究や資源配分へ役立てる。事前・中間・事後評価の時点と、PDCAの語順は別々に読む。

\subsection{平成23年度 問13}
\label{q:h23-reg-apt-13}

\noindent\textbf{原典：} sources/H23/aptitude.pdf、PDF物理14頁。

\textbf{公式正答：} 3（single）。

\par\textit{原典本文の校正と独自解説は未作成。分類のみの登録項目。}

\subsection{平成23年度 問15}
\label{q:h23-reg-apt-15}

\noindent\textbf{原典：} sources/H23/aptitude.pdf、PDF物理16頁。

\textbf{公式正答：} 1（single）。

\par\textit{原典本文の校正と独自解説は未作成。分類のみの登録項目。}

\subsection{平成27年度 問10}
\label{q:h27-reg-apt-10}

\noindent\textbf{原典：} sources/H27/aptitude.pdf、PDF物理10頁。

\textbf{公式正答：} 5（single）。

\par\textit{原典本文の校正と独自解説は未作成。分類のみの登録項目。}

\subsection{平成27年度 問12}
\label{q:h27-reg-apt-12}

\noindent\textbf{原典：} sources/H27/aptitude.pdf、PDF物理12頁。

\textbf{公式正答：} 2（single）。

\par\textit{原典本文の校正と独自解説は未作成。分類のみの登録項目。}

\subsection{平成28年度 問6}
\label{q:h28-reg-apt-06}

\noindent\textbf{原典：} sources/H28/aptitude.pdf、PDF物理6頁。

\textbf{公式正答：} 1（single）。

\par\textit{原典本文の校正と独自解説は未作成。分類のみの登録項目。}

\subsection{平成28年度 問11}
\label{q:h28-reg-apt-11}

\noindent\textbf{原典：} sources/H28/aptitude.pdf、PDF物理11頁。

\textbf{公式正答：} 4（single）。

\par\textit{原典本文の校正と独自解説は未作成。分類のみの登録項目。}

\subsection{平成29年度 問13}
\label{q:h29-reg-apt-13}

\noindent\textbf{原典：} sources/H29/aptitude.pdf、PDF物理12頁。

\textbf{公式正答：} 3（single）。

\par\textit{原典本文の校正と独自解説は未作成。分類のみの登録項目。}

\subsection{平成30年度 問5}
\label{q:h30-reg-apt-05}

\noindent\textbf{原典：} sources/H30/aptitude.pdf、PDF物理5頁。

\textbf{公式正答：} 4（single）。

\par\textit{原典本文の校正と独自解説は未作成。分類のみの登録項目。}

\subsection{平成30年度 問8}
\label{q:h30-reg-apt-08}

\noindent\textbf{原典：} sources/H30/aptitude.pdf、PDF物理8頁。

\textbf{公式正答：} 3（single）。

\par\textit{原典本文の校正と独自解説は未作成。分類のみの登録項目。}

\subsection{令和4年度 問9}
\label{q:r04-reg-apt-09}

\noindent\textbf{原典：} sources/R04/aptitude.pdf、PDF物理8頁。

\textbf{公式正答：} 5（single）。

\par\textit{原典本文の校正と独自解説は未作成。分類のみの登録項目。}

\subsection{令和4年度 問12}
\label{q:r04-reg-apt-12}

\noindent\textbf{原典：} sources/R04/aptitude.pdf、PDF物理11頁。

\textbf{公式正答：} 5（single）。

\par\textit{原典本文の校正と独自解説は未作成。分類のみの登録項目。}

\subsection{令和5年度 問4}
\label{q:r05-reg-apt-04}

\noindent\textbf{原典：} sources/R05/aptitude.pdf、PDF物理4頁。

\textbf{公式正答：} 2（single）。

\par\textit{原典本文の校正と独自解説は未作成。分類のみの登録項目。}

\subsection{令和6年度 問3}
\label{q:r06-reg-apt-03}

\noindent\textbf{原典：} sources/R06/aptitude.pdf、PDF物理3頁。

\textbf{公式正答：} 1（single）。

\par\textit{原典本文の校正と独自解説は未作成。分類のみの登録項目。}

\subsection{令和6年度 問12}
\label{q:r06-reg-apt-12}

\noindent\textbf{原典：} sources/R06/aptitude.pdf、PDF物理12頁。

\textbf{公式正答：} 3（single）。

\par\textit{原典本文の校正と独自解説は未作成。分類のみの登録項目。}

\subsection{令和7年度 問11}
\label{q:r07-reg-apt-11}

\noindent\textbf{原典：} sources/R07/aptitude.pdf、PDF物理11頁。

\textbf{公式正答：} 4（single）。

\par\textit{原典本文の校正と独自解説は未作成。分類のみの登録項目。}


\chapter{安全・リスク}
\section{背景：危害を生む条件を工程の全体で扱う}
事故を個人の不注意という一語で終えると、同じ事象を起こす設計や運用条件が残る。何が動き、何に接触し、どの防護が機能しなかったのかを先に再現する。その条件を設計、施工、使用、保守、情報の引継ぎまで広げて検討することが、別の業務へ教訓を移す出発点である。

\section{危険源と防護：検知後に止めるだけでよいか}
\label{theory:safety:barriers}
危険源とは危害を生じ得る性質や状況である。重い回転体は大きな運動エネルギーを持ち、人が入る位置で挟まれが起こり得る。リスクを考えるには、その危険へどのように接近し、どの程度の被害が起こり得るかを調べる。

本質的な設計で危害を生むエネルギーや接触を減らす方策、隔離や停止等の防護、残る危険についての情報と使用・保守手順を順に検討する。センサは人を検知した瞬間に機械のエネルギーを消せるわけではない。単純化した回転体で慣性モーメントを\(I\)、角速度を\(\omega\)とすれば運動エネルギーは\(I\omega^2/2\)である。一定の制動トルクの大きさを\(\tau\)とし、制御遅れ等を無視すると、停止までの角変位は
\[
 \Delta\theta=\frac{I\omega^2}{2\tau}
\]
となる。重量・形状や速度を変更したら、検知範囲と停止距離を含めて安全の前提を再評価する必要がある。これは事故機の数値再現ではなく、設計変更が停止性能へ与える影響を示すモデルである。

\section{五つの事例を物理的事象と管理条件へ分ける}
\label{theory:safety:cases}
\begin{description}
 \item[低脂肪乳の食中毒] 停電時の菌増殖と毒素産生が、汚染した脱脂粉乳を通じて最終製品へ及んだ。厚労省の原因説明を用い、菌を殺す工程と既に産生した毒素の問題も分ける\parencite{mhlw_milk2002}。原料の異常から出荷停止・回収までの情報経路は、物理的な原因とともに検討する。
 \item[自動回転ドア] 軽量な設計から機能や外観の変更で重量が増すと、停止できるという前提を再評価する必要がある。大学の事例分析は検知・非常停止への依存と慣性の増大を関連付ける\parencite{muroran_ethics_cases}。顧客の要求に応えたことだけで危険の説明責任が消えるわけではない。
 \item[福知山線] 事故報告の直接原因は速度超過で曲線へ進入したことにあり、ブレーキの遅れと運転士管理方法の関与は、記録に基づく推定や可能性として段階を区別する。曲線の速度超過を防ぐATS機能が使用開始されていれば回避できたと推定されるが、列車に全くATSがなかったわけではない\parencite{jtsb_fukuchiyama2007}。
 \item[戸開走行のエレベーター] 扉が開いたままかごが動くと乗降中の人が挟まれる。2006年の事例はかごの上昇による挟まれであり、ケーブル破断による落下とは機構が異なる。保守の情報と独立した防護の必要性を確認する\parencite{caa_elevator2006}。
 \item[笹子トンネル天井板] 接着系アンカーの脱落等を、材料の劣化という一原因だけへ還元しない。設計の荷重・耐久性、施工と品質記録、点検の方法・頻度、情報の継承を結び付けて扱う\parencite{mlit_sasago2013}。
\end{description}

\section{原因調査と責任の判断を分ける}
\label{theory:safety:responsibility}
再発防止の原因調査は、事象を説明し防護を改善するために行う。刑事責任は、対象被告人の注意義務、予見可能性、結果回避等を個別の事実関係の下で判断する。民事、行政、倫理・組織上の責任も同じ結論になるとは限らない。

歴代3社長を対象とした福知山線の最高裁2017年決定は、個別曲線の危険性の認識等を検討し、無罪の第一審を是認した原判断を相当とした\parencite{court_jr2017}。この結論から、会社の安全上の改善が不要だったとは導けない。事故調査の是正提言と刑事上の注意義務を同一の基準へまとめない。

不起訴処分は、刑事裁判での無罪判決とは異なる。平成30年度問14の公式注記は、笹子の不起訴を判決内容の選択肢へ入れたこと等を理由に不適切出題とする。全員正解という採点措置を、全ての事例の記述が正しいという認定へ置換しない。同問の一部の個別判決資料は未確認であり、問題ブロックにその範囲を残す。

\section{教訓を別の設計へ移す}
事故名と正答番号だけを覚えると、条件が変わった問題へ対応できない。故障機構、危険への接近、防護の作動条件、設計変更、保守で継承すべき情報を取り出す。自分の設備で同じ条件が成立し得るかを点検し、対策の効果と残る危険を記録する。以下では危険源から許容性の判断、対策、記録へ至る共通の枠組みを整理する。
\section{水害・土砂災害：情報を事前の行動計画へ結び付ける}
\label{theory:safety:preparedness}
被害は危険な現象だけで決まらず、曝露する人や設備、避難の可能性にも依存する。
ハザードマップで対象災害、想定の条件、被害範囲、避難先と経路を調べ、
当日の情報をどう行動に変えるか事前に考える。
地図に色がないことを、想定していない現象も含め全て安全だという証明にしない。
避難先と経路は家族の移動能力、地域の危険、所要時間を含めて考える。
備蓄・持出品、連絡方法・集合、保険等は、危険時の移動と被害後の生活の異なる課題を支える。

令和元年度再試験問13の5記述は、当時の一般的な備えとして適切である。
前文のAR5は、評価した排出シナリオのもとで21世紀の全球平均気温上昇を予測する
\cite{env_ar5}。
日本の1901〜2017年の大雨日数増加と降水日数減少も、該当観測資料で扱われる\cite{env_rain2017}。
期間を最新の年へ変えず、全国の集計傾向を全地点・毎年の単調増加へ拡張しない。
温暖化の緩和は将来の危険を減らす課題で、適応・避難の備えを不要にするものではない。

\begin{tcolorbox}[title=出題2019年度と制度改正]
原典は避難準備・高齢者等避難開始、避難勧告、避難指示（緊急）という旧名称を使う。
2021年5月20日から避難勧告は廃止し、従来の勧告の段階から避難指示を行う体系へ一本化した
\cite{env_evac2021}。
旧名称を現在の行動案内としてそのまま使わず、当時の正誤判定を保存する。
\end{tcolorbox}

\begin{tcolorbox}[title=2026年の防災気象情報]
気象庁は2026年5月29日から、情報名称に警戒レベルの数字を付す新たな防災気象情報の運用を開始すると公表した
\cite{env_warning2026}。
市町村等が発令する避難情報と、気象庁等が発表する危険度の情報は主体が異なる。
情報と行動との対応を分かりやすくする改正であり、発令を待つだけで必ず安全になるという保証ではない。
\end{tcolorbox}
注意報・警報等の危険度に応じた発表は、全ての段階が常に順番に来るという約束ではない。
急変で発令や受信が間に合わないこともあるため、危険を認識したときは自ら判断して避難行動を取るという原典の趣旨を理解する。
実際の災害時の個別行動は、現在の地域の情報と自分の条件に基づき決める。
\section{目的と不確かさからリスクを定義する}
\label{theory:safety:definition}
リスクは何が起きるかだけでなく、何を達成したいかに依存する。生産量を増やす設備変更は、
納期を守る機会と、作業者の危険や品質不良の可能性を同時に変え得る。
組織リスクの広い定義は目的への不確かさの影響を扱い、期待からの正負の逸脱を含む。
一方、安全評価では人の危害等の負の側面に焦点を当てる。
この範囲の違いを、互いに矛盾する定義だとしない。
ISOの公開説明も、組織の目的、機会・脅威、資源配分とリスク対応を結び付ける\cite{saf_iso31000}。

\section{危険源・曝露・結果を分ける}
回転する歯車という危険源があっても、閉鎖された領域に人が接近しなければ、
想定する挟まれ事故の経路は成立しない。清掃のためガードを外すと、その経路が開く。
危険源の性質、接触や曝露の機会、防護の状態、起こり得る結果を一つずつ取り出すと、
設計変更、隔離、手順、保護具のどれに作用する対策かが明確になる。
正常運転だけでなく、起動停止、清掃、修理、故障、廃棄でも同じ分解を行う。

\section{確率と損失の積は何を表すか}
\label{theory:safety:quantification}
一つの評価期間に、危害が生じる確率を\(p\)、生じたときの損失を\(L\)とし、
生じないときの損失をゼロとする二状態モデルでは、期待損失は
\[
 \mathbb E[\text{損失}]=pL
\]
である。結果が複数あり、互いに排他的で網羅的な場合には
\[
 \mathbb E[\text{損失}]=\sum_i p_iL_i,\qquad \sum_i p_i=1
\]
と書ける。事象を重複して足さず、期間と曝露集団、損失の単位をそろえる。
これは不確かな結果を平均で評価するモデルであり、最大損失、被害の分布、不公平、
回復不能な危害を全て一つの値だけで表現するものではない。
同じ期待損失でも、稀な大事故と頻繁な軽傷に同じ判断を下すとは限らない。

\section{順序尺度の掛け算と確率を混同しない}
重篤度を1〜5、発生しやすさを1〜5と分類して積を取る手法では、
数値は通常、厳密な損失額や確率ではない。
重篤度4を2のちょうど二倍の被害と解釈したり、分類2を分類1の二倍の確率と
解釈したりするには、別の尺度の根拠が必要である。
行列の評価、加算スコア、段階分岐も可能であり、全てのリスクを確率と損失の積で
定量表現しなければならないという要求にはならない\cite{saf_ra2006}。
手法は対象と利用できる情報に合わせ、重大な危害が平均値に隠れることも点検する。

\section{知識不足とばらつきは異なる}
寸法や故障率をまだ調べていないため分からない状態と、
同じ条件でも故障時点にばらつきがある状態は区別する。
図面、試験、保守記録や外部の事故情報は前者を減らし得る。
ただしデータを増やせば全ての自然変動や未知の事故経路がなくなるわけではない。
既知であることだけでリスクがゼロになるともいえない。
既知の危険な確率事象は依然として安全上のリスクである。

\section{安全と許容性は条件を含む判断である}
\label{theory:safety:acceptability}
安全を許容できないリスクがない状態と捉えると、ゼロリスクであるという証明とは違う。
誰が、どの期間に、どんな危害に曝露されるか、利用の目的と条件、社会の価値、
新たな知見と利用可能な技術を明示する必要がある。
評価基準を一度設定したことは、永久に改訂してはならない理由にならない。
また、限られた設計外力を超える事象があり得ると認識することは、
その被害を社会が無条件に受け入れるべきだという結論ではない。

\section{見積り・分析・評価・対応を分ける}
\label{theory:safety:cycle}
特定は何が起こり得るかを見つける段階である。
見積り・分析は原因、結果、起こりやすさや防護の効き方を調べる段階である。
評価はその情報を基準と比較し、対応の必要性や優先度を判断する段階である。
対応ではリスクを変える方策を選び実施する。
組織リスクの図でアセスメントを特定・分析・評価の三段階として囲んでいても、
職場の実務資料が措置の検討・記録まで含む一連の手法を説明することと直ちに矛盾しない。
設問が採用する用語と囲みの範囲を先に読む。

\section{対策後に戻る理由：残留と新たな危険}
対策の採用後も同じ危険の大きさを再評価する。
遮蔽や囲いを付けたことで温度が上がる、非常停止が通常操作の近くにある、
保守が難しく防護を外したくなる等、新たな経路を生むこともある。
残留リスクとは低減措置後に残るリスクである。
低減したはずという期待だけで完了にせず、許容性と措置の妥当性を確認し、
条件、実施者、判断根拠、残る危険を記録する。
評価書を書くだけでは、接触の可能性や危険源のエネルギーは下がらない。

\section{低減策には優先順位がある}
\label{theory:safety:hierarchy}
製品設計の三段階は、本質的安全設計、防護、使用上の情報という順序で理解する。
職場の2006年指針は、設計・計画での除去低減、工学的対策、管理的対策、
個人用保護具の順を示す\cite{saf_ra2006}。
二つの列挙は用途と粒度が違うため、そのまま同じ四項へ押し込まない。
先に危険な作業をなくす、より危険の小さい材料に替える、動力を減らす等を検討する。
ガード、隔離、排気等で補完し、作業規則と教育、必要な保護具を整える。
人間の注意に依存する方策だけで、上位の検討を省略しない。

\section{ALARPの三つの水準}
\label{theory:safety:alarp}
ALARPは合理的に実行可能な限り低くする原理である。
次の表は、英国議会の報告による許容性の枠組みを独自に整理したものである\cite{saf_parliament2006}。
図形の幅を確率や損失の値と誤読しないよう、数値目盛は付けない。
\begin{center}
\begin{tabular}{p{0.22\linewidth}p{0.65\linewidth}}
\hline
水準 & 判断の意味\\
\hline
許容できない & 低減又は活動の中止を必要とする高い水準である。\\
条件付きで許容 & 合理的に可能な追加低減を行った根拠を要する。\\
広く受容される & 十分低い水準であることと、その状態の維持を検討する。\\
\hline
\end{tabular}
\end{center}
中間へ入っただけで無条件に許容されるわけではない。
下の水準だけを「安全」の唯一の意味とすると、条件付き許容との整合性を失う。
図の説明、規格の適用水準、各国の法的な低減義務を同じ命題へまとめない。

\section{著しい不均衡を予算超過へ置換しない}
追加措置で避ける危害と、その措置に必要な費用、時間、困難等の負担を比較する。
事故の賠償金だけでは回避の便益を表さず、当初事業予算はこの比較の基準とは別である。
例えば一万円でも予算超過なら省略してよいという規則はALARPから導けない。
著しい不均衡を\(C>B\)という単純な大小判定へ置換しない。
ここで\(C\)を費用、\(B\)を同じ尺度で評価できた低減便益としても、
不確かさ、危害の重大性、適用制度の判断が必要である。
あらゆる危険へ使える一定の倍率を、本章で根拠なく作らない。
2006年指針も重大危害と暫定措置の条件を含むため、費用だけで放置を正当化しない。

\section{ALARAは防護の最適化である}
\label{theory:safety:alara}
放射線防護のALARAは合理的に達成可能な限り低くする原理である。
ICRP103第203項の最適化は、線量だけでなく、被ばくの可能性と人数を含み、
経済的・社会的条件を考慮する\cite{saf_icrpaedia}。
防護の改善は目的であるが、線量の数学的な絶対最小値だけを、
生活や社会の他の条件を無視して求める原理ではない。
個人に負担が偏らないような制約も必要である。
食品汚染物質でも国内の分布・摂取等を調べ、国際基準と健康への影響を検討する際に
ALARAが使われてきた\cite{saf_food2023}。
放射線と食品で、同じ定量指標をそのまま流用するわけではない。

\section{低線量モデルと影響の実証を区別する}
ICRPの2007年勧告は確率的影響について低線量の線形応答を防護の基礎に置く
\cite{saf_icrp103}。
モデルの簡略形を\(r(D)=kD\)と書くと、\(D\)は対象とする線量、
\(r\)はその線量に応じた増分リスク、\(k\)は適用する集団・モデルの係数である。
これは生物学的な全影響にしきい値がないという定理ではない。
線量、線量率、影響の種類、集団差と推定の不確かさを切り分ける。
試験の「しきい値がない」を読むときも、確率的影響の防護モデルという範囲を保持する。

\section{リスク対応と負担の移転}
\label{theory:safety:response}
回避は危険な活動を行わない選択、低減は危害の可能性や程度を変える選択、
移転は契約や保険等で負担を他者へ分担させる選択、保有は残る負担を認識して引き受ける選択である。
保険加入で事故の確率や物理的危害そのものが自動的に消えるわけではない。
複数の対応を組み合わせ、契約相手へ危険を知らせず責任を押し付けることと区別する。
対策を実施したという記録だけでなく、望んだ状態への変化を確認する。

\section{リスク認識には同じ式以上のものがある}
\label{theory:safety:perception}
客観的な推定と本人がどれほど危険だと感じるかは別の問題である。
同じ確率でも、本人が選んだ活動か、回避する権限があるか、近い将来か、
誰に利益と損失が分配されるかで受け止め方が違い得る。
同じ「確率と損失の積」という形式を使っても、その二つの主観的入力や価値が違えば
結論は一致しない。
直近の記憶や報道で想起しやすくなること、同じ情報でも利益・損失の提示枠で
判断が変わり得ることも区別する。
これらを全ての人に常に成立する一つの数値法則にせず、設問の一般化を吟味する。

\section{情報交換と市民の参加}
\label{theory:safety:communication}
専門的に細部を説明できない人も、何が不安か、どんな負担を受けるかを表明できる。
情報の不足と価値観の相違を分け、事実、推定、限界、選択肢を伝え、
質問と意見を受け止める。
食品安全委員会の資料は相互作用と参加を重視する\cite{saf_fsc_psych,saf_fsc_glossary}。
専門家の結論を受け入れさせるだけで過程を成功としない。
分からないことは範囲を示し、未知だから既知の被害も説明不要だとしない。

\section{統括管理と全員参加}
\label{theory:safety:governance}
現場の実態は担当者の知識なしには捉えにくい。
その参加を確保しても、必要な人員や資源、方針、措置の実施を統括する責任は残る。
職長の情報収集と、事業の統括管理者の責務を入れ替えない。
新設、変更、非定常作業、新たな事故情報、経年劣化は再調査の契機となる。
調査と記録は説明責任を支えるが、作成だけで法律上の責任が自動免除されるものではない。

\section{参加の効果と残留リスクの規則}
\label{theory:safety:participation}
職場の調査で何が危険かを共有し、合理的に優先度を決める。
残る危険について守る規則と、その必要な理由が分かるようにする。
全員参加の感受性向上という効果と、専門家が分析して危険度を推定する活動は両方あり得る。
公的資料の五つの効果の表現を問う問題では、その列挙との一致を先に確認する
\cite{saf_ra_effects}。
リストにない表現だから一般的にも必ず偽である、という説明を作らない。

\section{経験比は重災害ゼロの保証ではない}
\label{theory:safety:incident}
小さな事故や危険な経験を集めると、重大事故になる前に経路を修正できる。
1:29:300という伝統的な経験比は、その注意を促す説明に用いられる
\cite{saf_welfare}。
集計の期間、業種、危害の定義や報告率が違えば件数は変わる。
29の軽傷要因を除去すれば重大事故が必ずなくなる、各個人が330件で一度重災害を起こす、
という確定的な因果法則へ変えない。
重災害固有の経路も直接調べる。
報告者を責めると情報が隠れるため、早期の報告・保護・改善と共有を結び付ける。

\section{4Sは状態を維持する活動である}
\label{theory:safety:housekeeping}
整理は不要なものを分け除くこと、整頓は必要なものを安全に取り出せる配置へ置くこと、
清掃は汚れを除き異常を見つけること、清潔は良い状態を維持することである
\cite{saf_4s}。
清掃を洗浄へ、清潔をしつけへ置換する列挙は標準的4Sと違う。
しつけを加える列挙は5Sである。
頭文字だけでなく、通路の障害物、漏油、危険物の保管、故障兆候とどう結び付くかを理解する。

\section{警告・使用情報は利用者へ届いて初めて働く}
\label{theory:safety:warnings}
警告は危険源、起こる危害、従わない結果、避ける行動が理解できる必要がある。
表示の可読性、消えにくさ、見つける位置、利用者の言語と能力、使用期間を検討する。
製造国の言語と英語を書いた事実だけで、全ての使用国の利用者が理解できると保証しない。
合理的に予見できる誤使用、保守、運搬、廃棄にも情報を用意する。
ただしあらゆる想像上の誤使用を同じ範囲に含めず、予見可能性の根拠を明示する。
家電の国内指針の適用範囲と、輸出先の法令や表示要求も別である\cite{saf_aeha}。

\section{法令・規格の時点と対象を保存する}
\label{theory:safety:law}
\label{theory:safety:standards}
\begin{tcolorbox}[title=出題当時の制度と法源]
2006年の一般リスク調査の努力義務と、2016年6月1日施行の対象化学物質の調査義務を分ける
\cite{saf_ra2006,saf_chem2016}。
「化学物質」とあるだけで、全製品・全濃度へ無条件の義務を付加しない。
法律本則、政省令、指針、実務解説では規定の場所と性質が違う。
実施時期の詳細が指針にあることを、全部が法律本則の列挙だと言い換えない。
\end{tcolorbox}
SDSは物性・危険有害性・取扱い等の情報を渡す文書で、調査や教育の材料となる。
法的交付・調査義務には指定対象、濃度や用途等の条件がある\cite{saf_chemqa}。
健康診断の実施と結果報告も別であり、一般の定期等の報告の規模条件と
特殊健診の条件を分ける\cite{saf_healthreport}。

\begin{tcolorbox}[title=規格本文の確認範囲]
ISOの公式ページでGuide51の2014版、ISO31000の2018版の情報を確認した。
Guide51の公開全文の取得は失敗しており、本章を全条項の照合済みとはしない
\cite{saf_guide51,saf_iso31000}。
JSAのQ31010:2012の公開表紙はIEC/ISO31010:2009のリスクアセスメント技法を示す
\cite{saf_jis31010}。
平成28年度前文の規格番号の疑義は原典を保存し、正答番号と分離して追加確認へ残す。
平成24年度の1996年規格史と平成25年度の旧家電指針の版も未確認部分を残す。
\end{tcolorbox}

\section{試験での解き方と統合して学ぶ事項}
まず安全の定義、定量モデル、判断基準、措置、記録のどれを問うか決める。
個数問題では各記述の真偽を記録してから数え、候補番号へ変換する。
組合せ問題では表の列順と○×を画像で確かめる。
特に「変わらない」「必ず積」「優先順位なし」「予算に収まらないだけ」という
全称・省略条件を点検する。
公式正答と独立初回が一致しても、他の肢の曖昧さや前文の誤記は解消しない。
頻度は確認済みの概念統計から集計し、異なる回答形式を同一番号の暗記へ縮約しない。
\section{製造物責任を四つの判定に分ける}
\label{theory:pl:framework}
製造物責任法は、製品の欠陥による生命・身体・財産の被害について、製造業者等の民事上の賠償責任を定める。事故が起きたという一事だけで責任が成立するわけではない。まず対象となる物と責任主体を定め、次に欠陥、損害、因果関係を確認し、最後に免責と期間制限を検討する。この順序なら、対象外という判断を安全対策の不要という判断へ飛躍させずに済む。\autocite{pl_statute,pl_qa}

\subsection{過失の証明と欠陥の証明}
\label{theory:pl:proof}
PL法では製造業者の故意又は過失の立証を責任成立の要件としない。ただし、欠陥、損害、両者の因果関係は別の要件である。「過失を証明する必要がない」から「欠陥さえあれば全損害を賠償する」と推論してはならない。H26問7の5とH29問8の2は、救済の変化を短く述べた文として読む一方、独立した法律説明では省略された要件を補って理解する。

\subsection{その製品だけの損害}
\label{theory:pl:selfdamage}
3条ただし書は、損害が当該製造物についてのみ生じた場合をPL法の賠償対象から除く。テレビが発火してテレビだけが壊れた場合と、住宅へ延焼した場合を対比する。前者でも契約等による請求の可能性は別にあり、保証期間が過ぎたという事情だけで全ての救済が消えるとはいえない。自動二輪車の故障も、他の人又は物に被害があるかを先に確かめる。\autocite{pl_statute,pl_qa}

\section{有体物・動産・製造加工の三条件}
\label{theory:pl:scope}
日本法の対象は製造又は加工された動産である。電気、音響、単体ソフトウェア、サービスは有体物という条件を満たさず、土地や家屋自体は動産ではない。ガスやガソリンは有体物であるため、電気と同じ分類に置かない。水道についても、供給サービスと製造加工された物を区別する必要がある。\autocite{pl_definitions2018,pl_qa}

\subsection{単体ソフトウェアと組込み製品}
\label{theory:pl:software}
ソフトウェア自体が対象外でも、その不具合がロボットという製造物の安全性の欠如として現れ、負傷との因果関係が認められるなら、ロボットの製造業者等の責任が問題となる。H24問7のプリインストールされたOSはソフトウェアそのものを、同問のロボットは製品全体を指す。情報の欠陥を有体製品の欠陥へ結び付ける段階が解法の要点である。

\subsection{不動産に取り付けた動産}
\label{theory:pl:attachment}
製造物の引渡し後に建物へ取り付けたという理由だけで、製造時の欠陥の責任が消えるわけではない。エスカレータ、シャッター、サッシ等について、引渡時の動産性と欠陥、損害との因果関係を検討する。建物本体の施工不良とは異なる。逐条解説はエスカレータからの転落事件の裁判例も紹介している。\autocite{pl_definitions2018}

\subsection{冷凍・冷蔵と加工の区別}
\label{theory:pl:processing}
加工は、物の本質を保ちながら新しい属性を付け加える操作として理解する。単なる切断、冷凍、冷蔵、乾燥は基本的にこの加工へ含めない。一方、加熱、調味、塩漬け、燻製等により新しい属性を加えた食品は対象になり得る。冷凍肉とハムを比べる際には、保存状態の変化と新属性の付与を分ける。植物工場のトマトや養殖魚も、生産技術が高度であるという事情だけでは製造加工された物とはいえない。\autocite{pl_definitions2018}

\paragraph{初回解答を訂正する根拠}
H25問7の初回独立解答は、冷凍冷蔵を加工と扱って不適切6件・候補5とした。逐条解説の物理3頁（印刷4頁）を画像で確認し、この分類を訂正した。オを適切とすると、不適切はア・ウ・エ・カ・キの5件となり候補4に対応する。初回記録は保存し、公式解答との一致を訂正理由の代わりにしない。同問イの「品質上の欠陥」は安全性に関係する範囲で読む必要がある。

\section{欠陥は安全性の不足である}
\label{theory:pl:defect}
欠陥の判断には、製品の特性、通常予見される使用形態、引渡時期等を総合して考慮する。単なる性能の不足や外観の不良と、安全性の不足を区別する。設計、製造、指示警告の各段階に安全性の不足が生じ得る。メーカーが意図した操作だけを検討すればよいわけではなく、通常予見できる誤使用も安全設計の検討に含める。\autocite{pl_qa,pl_statute}

\section{製造業者等と免責の条件}
\label{theory:pl:producer}
業として製造・加工・輸入した者に加え、製造業者としての表示、製造業者と誤認させる表示、実質的製造業者と認められる表示をした者が責任主体となり得る。OEMの販売者は「常に除外」でも「常に該当」でもなく、表示と実質を確認する。単なる販売店と輸入者を一括して対象外としてはならない。\autocite{pl_statute,pl_qa}

\subsection{開発危険の抗弁}
\label{theory:pl:development}
引渡時の科学技術の知見では欠陥を認識できなかったことを製造業者等が証明した場合には、4条の免責が問題となる。担当者個人が知らなかったという説明では足りない。「認識できなくても必ず責任を負う」は免責規定を落とした文である。部品又は原材料には、欠陥が専ら完成品側の設計指示によって生じ、かつ部品側に過失がない場合の別の免責もある。

\section{修理された製品と修理行為を分ける}
\label{theory:pl:repair}
性質の回復・維持を行う修理、修繕、整備というサービス自体は、基本的に製造加工ではない。しかし、修理を経た元の製品が製造物でなくなるわけではない。元の引渡時にあった欠陥か、後の修理で生じた不具合かを区別する。H24問7の「使用者が修理した製品」を対象物の定義として読むと対象性は残り、修理行為や修理起因の不具合について読む場合とは数え方が異なる。この問の初回候補3と公式候補4の差は、解釈疑義として保存する。\autocite{pl_qa,pl_definitions2018}

\subsection{再生品と元の製造業者}
\label{theory:pl:recycle}
使用困難になった物の一部から製造加工した再生品は対象となる。再生した者が責任主体となり得るほか、原材料となった元の製品の引渡時の欠陥が再生品の損害と因果関係を持つなら、元の製造業者の責任も検討する。「最後に再生した者が全ての責任を負う」と一律にはいえない。\autocite{pl_definitions2018}

\section{期間制限は出題時点を確認する}
\label{theory:pl:time}
現行法の短期期間は損害及び賠償義務者を知った時から原則3年、生命又は身体侵害では5年である。長期期間は製造業者等の引渡しから10年であり、消費者の購入日から数えるとは限らない。蓄積性物質や潜伏期間のある損害には起算点の特則もある。H30問9の通常のテレビ火災について、欠陥の証明だけで期間経過を無視できるという文は成立しない。\autocite{pl_statute,pl_qa}

\subsection{2020年4月1日の改正と経過措置}
\label{theory:pl:transition}
改正前は生命身体侵害にも短期3年を用いた。2020年4月1日施行の改正ではこれを5年へ延長し、施行時に旧短期期間が完成していない場合の適用を定めた。長期10年についても時効期間であることを明記した。消費者庁の2018年逐条解説は原則として改正後の規定を説明する資料であるため、刊行年をそのまま当時の施行法の年とみなさない。\autocite{pl_transition2020,pl_time2018}

\section{民法・特別法・海外制度との関係}
\label{theory:pl:otherlaw}
PL法に特別の定めがない責任については民法の規定を用いる。対象外となる事例に全く責任がないとは結論できない。原子炉運転等による原子力損害については原子力損害賠償法の規定を確認し、同法4条3項によるPL法不適用を扱う。また、日本法だけを守れば海外の製造物責任制度も満たすという国際共通化の主張は成立しない。逐条解説にも電気の対象性の違いが説明されている。\autocite{pl_nuclear_statute,pl_statute,pl_definitions2018}

\section{賠償責任の回避を超える製品安全}
\label{theory:pl:csr}
R02問12は事故後の賠償だけを問う問題ではない。法令順守後のリスク低減、消費者の期待を踏まえた基準、事故予防、対話、弱者と将来への配慮、緊急時の被害拡大防止、苦情解決の手段を順に検討する。問題に提示された社会的責任の定義に照らすと、7項目はいずれも安全・安心へ向けた行動であり、不適切0件に対応する。ISO26000原規格全文との照合は未完として残す。\autocite{pl_product_csr2016}

\section{維持管理を点検から次の判断へつなぐ}\label{theory:safety:maintenance}
社会資本の機能を持続させるには、状態を点検し、診断し、必要な措置を選び、
その根拠と結果を記録して次の周期へ反映する。
予防的・計画的な修繕は、劣化が大きくなってから一斉更新する負担を避け、
長期の費用の縮減・平準化と利用者に必要な機能の維持へつなげる。
計画を作れば必ず費用が一定になる、という保証ではない。\cite{eng52_infra2014}

平成26年の行動計画は施設分野別の基準・取組を含む。
施設の重要度、状態、材料、環境と分野ごとの最低基準を踏まえて実施するのであり、
全インフラの頻度と内容を一つへ固定するという意味ではない。
分野内での統一した健全度の尺度と、全施設で同じ点検をするという全称を分ける。
構造諸元、点検履歴、修繕記録、損傷位置・変化を蓄積し、
データが不足している状態を把握することも診断の仕事である。

\subsection{新技術と人材・責任を結び付ける}
ロボット、非破壊検査、ICTで効率・観測範囲を改善できるが、
測定の限界・見落とし・適用可能性を評価する能力が必要になる。
施設分野・業務ごとの必要能力を明確にし、担い手を確保・育成する。
資格の保有だけで全分野の経験と判断を代替しない。\cite{eng52_infra2014}

\subsection{行政責任と市民参加は両立する}
管理者の安全確保の責任を、地域住民・民間の協力を一律禁止する根拠にはしない。
国土交通省の平成12年白書にも、行政との合意に基づく道路・河川の清掃等の協働が記録されている。
誰が何を行い、どこから専門点検・管理者の判断が必要になるかを明確にする。
清掃への参加の事実から危険作業や構造診断を無条件に住民へ委任できるとは推論しない。\cite{eng52_adopt}

\section{化学品のGHS：分類と伝達を合わせる}
\label{theory:safety:ghs}
同じ化学品でも国ごとに警告の方法が異なると、運搬・販売・使用の途中で危険を見落とす。
GHSは危険有害性を共通の基準で分類し、その結果をラベルやSDSで伝える枠組みである。
分類の種類と程度を表すクラス・区分、絵表示、注意喚起語、危険有害性情報と注意書きを対応させる。
国内法の表示事項を追加することと国際的な調和は両立する。
当時のGHS案内と平成24年の指針を参照する。\cite{saf53_ghs_old,saf53_ghs_notice}

分類が伝えるのは危険有害性であり、どの作業でも同じ被害確率という意味ではない。
例えば同じ毒性の液体でも、密閉移送と開放した噴霧では曝露の条件が違う。
容器名と供給者が分かることは、追加情報を問い合わせ、事故時に追跡するためにも必要である。
複数のラベル要素の優先規則を、任意に不都合な危険だけ省略できる規則へ変えない。
新しい知見を得たら情報を更新する。当時のJISの細目と現行の表示義務は、版を分けて確認する。

\section{機械安全：本質的な設計と残る危険の情報}
\label{theory:safety:machine}
信頼性は、指定の条件と期間に要求された機能を果たす能力である。
高い信頼性で意図どおり動く装置でも、目的とする動きが人を傷つける設計なら安全とはいえない。
安全性を調べる際は、使用範囲と寿命段階、危険源、危険へ接近する状況を明確にする。
リスク分析の結果を評価し、必要な低減を繰り返す。\cite{saf53_machine2007}

まず危険を生む形状・運動・エネルギーを設計から減らす。
次にガードや保護装置等で接近と危害を防ぐ。それでも残るリスク、正しい使用・点検・訓練等の情報を提供する。
使用上の情報を落として「二つだけの方策」としない。
警告を増やすことだけで本質的な安全設計の検討を代替しない。
機械の製造側と使用側は各段階で役割を持ち、残留リスクがあるから全て使用者の責任という結論にはならない。

システム安全はハードウェアだけでなく、環境・人的条件を含める。
機能安全は必要な安全機能が正しく働くことによる安全で、通常の制御機能全てを意味しない。
安全工学は使用者と作業者双方の危害を扱い、レジリエンスは予防だけでなく
変化への対応や回復できる能力へ目を向ける。原典の対比表現を、全規格の排他的な定義へ広げない。

\section{リコール：危害防止の是正と賠償の成立を分ける}
\label{theory:safety:recall}
安全を損なう製品を把握したら、対象範囲と危害を調べ、回収、修理、交換、使用停止の連絡等を検討する。
劣化や表示の誤りでも是正が必要になり得る。
医薬品に別の品質管理法があることから、回収がないとは言えない。
PMDAは医薬品の回収情報を公表している。\cite{saf53_medicine_recall}

消安法の対象範囲、薬事等の別制度、行政への報告、購入者への連絡、社告を別の問いへ分ける。
対象が全て追跡されて回収できる場合の社告の扱いは、当時の手引きの細条件を要確認とする。
2010手引きの直接全文取得は成功しておらず、平成25年の条件を現行手引きだけから確定しない。\cite{saf53_recall2010}
回収した事実はPL責任の成立を直ちに示さない。欠陥・損害・因果関係等は別に検討する。
一方、賠償が成立しない可能性を、危害防止へ取り組まない根拠にしない。\cite{pl_statute,pl_qa}

\section{重大製品事故：報告と公表の起点を区別する}
\label{theory:safety:productaccident}
重大事故の代表類型は、死亡、治療に要する期間30日以上の重傷病、後遺障害、CO中毒、火災である。
製造・輸入側は事故の発生を知った段階から所定の期限内に報告する。
原因分析が未完という理由で期限を超えてよいとはしない。\cite{saf53_products2012}
事故の日、事業者が知った日、行政が報告を受理した日は一致するとは限らない。
後年取得した事業者用手引きは、知った日を含む10日を示すが、2018施行府令全文の照合と同一視しない。\cite{saf53_products2023}

行政の公表は事業者の報告と別の段階である。
2013年のフローでは、ガス・石油機器や製品起因が疑われる事故の迅速な公表と、
他の段階の報告受理後原則1週間という扱いを区別する。
「原則」を全案件に一律の法定絶対期限へ変えない。\cite{saf53_productflow2013}
通常有すべき安全性と、あらゆる条件での絶対安全は異なる。
合理的に予見できる使用を検討し、残るリスクを説明して減らすことが課題である。

長期使用の点検制度は経年劣化への対応を支える。
平成29年度の9品目は当時の指定として保持し、2021年等の対象見直しを遡及しない。
指定から外れたことも、劣化や保守の問題が消えるという意味ではない。\cite{saf53_oldnine}

\section{食品リスク：評価・管理・対話をつなぐ}
\label{theory:safety:foodroles}
科学的な健康影響を評価する機能と、許容する条件や施策を決め実施する管理の機能を分ける。
食品安全委員会の旧説明は、科学的な評価と管理行政の独立性を示す。\cite{saf53_food_roles2007}
評価結果だけから全ての政策が一意に決まるとは限らず、実行可能性、費用、便益や他の影響も検討する。
管理者の都合で評価結果を改変することも許されない。

リスクコミュニケーションでは、行政が情報を整理し対話の場を設け、
事業者は生産・流通・販売の情報を共有し、消費者は必要な情報を求め意見を述べ、
科学者は根拠と限界を分かりやすく説明する。
全過程における意見と情報の交換であって、結論を一方向に伝えて納得させるだけではない。\cite{saf_fsc_psych,saf_fsc_glossary}
リスク評価と管理の担当を入れ替えた肢は、この機能分担から判断できる。

\section{品質不正：記録の真正性と構造安全を別に検証する}
\label{theory:safety:quality}
測定記録を流用すれば、設計条件どおり施工したかを第三者が検証する根拠を失う。
記録不正だけで全ての構造物が倒壊すると断定できないが、実際の安全性の迅速な確認と
不正行為・管理過程の是正はともに必要である。
2015年の基礎ぐい報告は、責任体制と支持層到達等の技術判断を扱う。\cite{saf53_pile2015}

元請は施工全体を統括し、専門工事の担当は現場で得る情報から技術的に判断する。
設計図があるから現地の地盤確認を省略するとはしない。
機器が不調になったら、作業の停止や再取得、異常の報告と判断の経路を用意する。
ハードウェアの高度化だけでは、記録を偽る組織の圧力や曖昧な責任を取り除けない。
品質不正は他社からも教訓を得られる。トップの優先順位、部門間の役割、検査の独立性まで調べ、
個人の行為と組織条件の両方へ対応する。\cite{saf53_quality2020}

\section{作業の変更と軽微な危害の例外}
\label{theory:safety:workplace}
材料、設備、作業手順や人員の条件を変更したら、以前の評価の前提を確認する。
2006年の指針は、建設物や設備・材料・手順の変更時に調査し、疲労等の付加的な影響も考慮する。
対象の選定では、明らかに軽微な負傷・疾病しか予想されない平坦な歩行等を除外できる。
解説は軽微を医師の治療を要しない程度とし、過去に偶然軽い結果しか出なかったという条件とは区別する。\cite{saf53_workplace2006}

対策の優先順位と実行可能性を検討し、低減を続ける。
ALARPは、追加の低減による便益に対して費用等が著しく不均衡となるかを検討する原則で、
単なる安い対策の選択ではない。法令の最低基準を無視する理由にもならない。\cite{saf_parliament2006}

\section{ヒューマンエラー：結果の手前にある条件}
\label{theory:safety:human}
不適切な操作を認めても、「不注意」で説明を終えない。
知識・経験、注意、伝達、組織、手順の省略、状況への反射、混乱、錯覚、
心身の能力、疲労、単調さなどから、なぜその操作が起きたかを調べる。
教育資料の12分類はこれらの発生要因を含む。\cite{saf53_human12}
疲労や年齢による機能変化は、誤操作という行為自体と同じ概念ではない。
令和2年の問いは発生要因の分類に対応するため0項目を選ぶが、前文の9割以上という数値は
母集団・集計法が未確認である。全分野の普遍的な事故率として使わない。

対策は作業者への注意だけに限らない。認知しやすい表示、誤操作を防ぐ形状、
異常時の停止、負荷と休息、明確な連絡経路、責任分担と記録を組み合わせる。
原因の分類名を覚えるだけで、具体的な工程や条件を調べる課題を代替しない。

\section{2021行動計画：章と七つの取組を読む}
\label{theory:safety:maintenance2021}
第2次インフラ長寿命化計画は2021年6月18日に公表され、令和3〜7年度を計画期間とする。
必要施策の方向性にある七項目は、個別施設計画、点検診断・修繕更新、予算管理、体制、
新技術、情報基盤、基準類である。現状課題の別章と見出し数を混同しない。\cite{saf53_infra2021}
技術継承が独立した七項目の見出しにないことは、人材の確保や能力の継承が不要という意味ではない。
計画の階層を読み、見出しへの所属と実際の施策内容を区別する。

\section{事業継続：原因より重要事業への影響を調べる}
\label{theory:safety:bcp}
BCPは事業継続計画、BCMはその計画を支える経営・管理の活動である。重要な事業を継続し、又は目標とする期間・水準で復旧させるために、事業への影響、必要な資源、依存関係、代替手段を整理する。自然災害に限らず、感染症、設備事故、サイバー攻撃、取引先の中断等も事業へ影響する\parencite{eth54_bcp2013,eth54_bcp2021}。

復旧が容易な仕事を先に再開すればよいとは限らない。供給停止が顧客、収入、雇用、法令・契約、地域機能へ与える影響と許容時間を評価し、重要業務を選ぶ。バックアップは本体と同じ原因で失われないよう、場所、電源、通信、供給者等の共通の依存を点検する。

\paragraph{計画から実行・改善へ}
経営者が責任を持ち、訓練、点検、資源の確保、事業環境の変化への見直しを進める。計画書を作ったことだけでは対応能力を保証しない。重要業務や依存関係の可視化は平時の改善にもつながる。地域と取引先との協力を含め、組織単独で復旧できると仮定しない。

非常時だから法令や契約がすべて無効になるわけではない。緊急措置には固有の条件があり、権限や責任を確認する。「全て遵守しなくてよい」という包括的免除とは区別する。

\paragraph{普及率は調査年と分母を読む}
平成29年度調査では策定済みの割合が大企業64.0\%、中堅企業31.8\%とされる\parencite{eth54_bcp_survey2018}。これはすべての会社を同じ分母で測った数字ではない。策定中も含む率、企業規模、回答数、調査時点を分ける。令和元年度問題を後に公表された調査で説明しない。

\section{予防的取組：不確実性の中で行動を選ぶ}
\label{theory:safety:precaution}
知識が不完全なことと危険がないことは別である。深刻又は不可逆的な影響が懸念されるときは、科学的確実性がないことだけを理由に合理的な低減策を延期しない。利益、費用、代替案、行動しない場合の損失を比較し、対策の程度と見直しの条件を定める\parencite{eth54_rio,eth54_eea}。

平成25年度問題の十二項目は、無知と不確実性、監視、科学の盲点、学際的学習、実世界条件、利益とリスク、代替策、地域の知識、異なる価値、規制の独立性、制度の障壁、分析だけによる停滞を点検する。これらを考慮することは、科学を捨てることやあらゆる技術を禁止することではない。原典前文の宣言の「批准」という表現も、条約の法的手続と同一にしない。

\section{対話の視点と検証可能性}
\label{theory:safety:riskdialogue}
管理する側の集団的な確率評価と、危険を受ける当事者の判断は異なる。自分で選択できるか、利益と危険の配分は公平か、組織を信頼できるか等もリスク認知に関わる。ハザードとアウトレージを合わせる表現は認知モデルであり、単位の異なる物理量を足す計算式ではない\parencite{eth54_riskcomm}。

専門家の知識の伝達だけで対話が完成するわけではない。中立的な媒介者の支援、権限と情報の非対称性への配慮、異なる根拠の提示、検証可能性が必要になる。分かりやすい説明を作るときにも、情報源や不確かさを隠してはならない。

\subsection{平成23年度 問5}
\label{q:h23-reg-apt-05}

\noindent\textbf{原典：} sources/H23/aptitude.pdf、PDF物理6頁。

\textbf{公式正答：} 4（single）。

\par\textit{原典本文の校正と独自解説は未作成。分類のみの登録項目。}

\subsection{平成23年度 問6}
\label{q:h23-reg-apt-06}

\noindent\textbf{原典：} sources/H23/aptitude.pdf、PDF物理7頁。

\textbf{公式正答：} 4（single）。

\par\textit{原典本文の校正と独自解説は未作成。分類のみの登録項目。}

\subsection{平成23年度 問8}
\label{q:h23-reg-apt-08}

\noindent\textbf{原典：} sources/H23/aptitude.pdf、PDF物理9頁。

\textbf{公式正答：} 5（single）。

\par\textit{原典本文の校正と独自解説は未作成。分類のみの登録項目。}

\subsection{平成24年度 問6}
\label{q:h24-reg-apt-06}

\noindent\textbf{原典：} sources/H24/aptitude.pdf、PDF物理5頁。

\textbf{公式正答：} 3（single）。

\par\textit{原典本文の校正と独自解説は未作成。分類のみの登録項目。}

\subsection{平成24年度 問7}
\label{q:h24-reg-apt-07}

\noindent\textbf{原典：} sources/H24/aptitude.pdf、PDF物理6頁。

\textbf{公式正答：} 4（single）。

\par\textit{原典本文の校正と独自解説は未作成。分類のみの登録項目。}

\subsection{平成24年度 問10}
\label{q:h24-reg-apt-10}

\noindent\textbf{原典：} sources/H24/aptitude.pdf、PDF物理9頁。

\textbf{公式正答：} 1（single）。

\par\textit{原典本文の校正と独自解説は未作成。分類のみの登録項目。}

\subsection{平成24年度 問11}
\label{q:h24-reg-apt-11}

\noindent\textbf{原典：} sources/H24/aptitude.pdf、PDF物理10頁。

\textbf{公式正答：} 3（single）。

\par\textit{原典本文の校正と独自解説は未作成。分類のみの登録項目。}

\subsection{平成25年度 問4}
\label{q:h25-reg-apt-04}

\noindent\textbf{原典：} sources/H25/aptitude.pdf、PDF物理5頁。

\textbf{公式正答：} 5（single）。

\par\textit{原典本文の校正と独自解説は未作成。分類のみの登録項目。}

\subsection{平成25年度 問5}
\label{q:h25-reg-apt-05}

\noindent\textbf{原典：} sources/H25/aptitude.pdf、PDF物理6頁。

\textbf{公式正答：} 3（single）。

\par\textit{原典本文の校正と独自解説は未作成。分類のみの登録項目。}

\subsection{平成25年度 問7}
\label{q:h25-reg-apt-07}

\noindent\textbf{原典：} sources/H25/aptitude.pdf、PDF物理8頁。

\textbf{公式正答：} 4（single）。

\par\textit{原典本文の校正と独自解説は未作成。分類のみの登録項目。}

\subsection{平成25年度 問8}
\label{q:h25-reg-apt-08}

\noindent\textbf{原典：} sources/H25/aptitude.pdf、PDF物理9頁。

\textbf{公式正答：} 2（single）。

\par\textit{原典本文の校正と独自解説は未作成。分類のみの登録項目。}

\subsection{平成25年度 問9}
\label{q:h25-reg-apt-09}

\noindent\textbf{原典：} sources/H25/aptitude.pdf、PDF物理10頁。

\textbf{公式正答：} 5（single）。

\par\textit{原典本文の校正と独自解説は未作成。分類のみの登録項目。}

\subsection{平成26年度 問7}
\label{q:h26-reg-apt-07}

\noindent\textbf{原典：} sources/H26/aptitude.pdf、PDF物理7頁。

\textbf{公式正答：} 4（single）。

\par\textit{原典本文の校正と独自解説は未作成。分類のみの登録項目。}

\subsection{平成26年度 問8}
\label{q:h26-reg-apt-08}

\noindent\textbf{原典：} sources/H26/aptitude.pdf、PDF物理8頁。

\textbf{公式正答：} 3（single）。

\par\textit{原典本文の校正と独自解説は未作成。分類のみの登録項目。}

\subsection{平成26年度 問9}
\label{q:h26-reg-apt-09}

\noindent\textbf{原典：} sources/H26/aptitude.pdf、PDF物理9頁。

\textbf{公式正答：} 4（single）。

\par\textit{原典本文の校正と独自解説は未作成。分類のみの登録項目。}

\subsection{平成26年度 問11}
\label{q:h26-reg-apt-11}

\noindent\textbf{原典：} sources/H26/aptitude.pdf、PDF物理11頁。

\textbf{公式正答：} 1（single）。

\par\textit{原典本文の校正と独自解説は未作成。分類のみの登録項目。}

\subsection{平成26年度 問13}
\label{q:h26-reg-apt-13}

\noindent\textbf{原典：} sources/H26/aptitude.pdf、PDF物理13頁。

\textbf{公式正答：} 5（single）。

\par\textit{原典本文の校正と独自解説は未作成。分類のみの登録項目。}

\subsection{平成26年度 問15}
\label{q:h26-reg-apt-15}

\noindent\textbf{原典：} sources/H26/aptitude.pdf、PDF物理15頁。

\textbf{公式正答：} 3（single）。

\par\textit{原典本文の校正と独自解説は未作成。分類のみの登録項目。}

\subsection{平成27年度 問4}
\label{q:h27-reg-apt-04}

\noindent\textbf{原典：} sources/H27/aptitude.pdf、PDF物理4頁。

\textbf{公式正答：} 3（single）。

\par\textit{原典本文の校正と独自解説は未作成。分類のみの登録項目。}

\subsection{平成27年度 問5}
\label{q:h27-reg-apt-05}

\noindent\textbf{原典：} sources/H27/aptitude.pdf、PDF物理5頁。

\textbf{公式正答：} 5（single）。

\par\textit{原典本文の校正と独自解説は未作成。分類のみの登録項目。}

\subsection{平成27年度 問9}
\label{q:h27-reg-apt-09}

\noindent\textbf{原典：} sources/H27/aptitude.pdf、PDF物理9頁。

\textbf{公式正答：} 3（single）。

\par\textit{原典本文の校正と独自解説は未作成。分類のみの登録項目。}

\subsection{平成27年度 問13}
\label{q:h27-reg-apt-13}

\noindent\textbf{原典：} sources/H27/aptitude.pdf、PDF物理13頁。

\textbf{公式正答：} 2（single）。

\par\textit{原典本文の校正と独自解説は未作成。分類のみの登録項目。}

\subsection{平成27年度 問14}
\label{q:h27-reg-apt-14}

\noindent\textbf{原典：} sources/H27/aptitude.pdf、PDF物理14頁。

\textbf{公式正答：} 5（single）。

\par\textit{原典本文の校正と独自解説は未作成。分類のみの登録項目。}

\subsection{平成28年度 問3}
\label{q:h28-reg-apt-03}

\noindent\textbf{原典：} sources/H28/aptitude.pdf、PDF物理3頁。

\textbf{公式正答：} 1（single）。

\par\textit{原典本文の校正と独自解説は未作成。分類のみの登録項目。}

\subsection{平成28年度 問5}
\label{q:h28-reg-apt-05}

\noindent\textbf{原典：} sources/H28/aptitude.pdf、PDF物理5頁。

\textbf{公式正答：} 4（single）。

\par\textit{原典本文の校正と独自解説は未作成。分類のみの登録項目。}

\subsection{平成28年度 問7}
\label{q:h28-reg-apt-07}

\noindent\textbf{原典：} sources/H28/aptitude.pdf、PDF物理7頁。

\textbf{公式正答：} 5（single）。

\par\textit{原典本文の校正と独自解説は未作成。分類のみの登録項目。}

\subsection{平成28年度 問9}
\label{q:h28-reg-apt-09}

\noindent\textbf{原典：} sources/H28/aptitude.pdf、PDF物理9頁。

\textbf{公式正答：} 5（single）。

\par\textit{原典本文の校正と独自解説は未作成。分類のみの登録項目。}

\subsection{平成28年度 問15}
\label{q:h28-reg-apt-15}

\noindent\textbf{原典：} sources/H28/aptitude.pdf、PDF物理15頁。

\textbf{公式正答：} 3（single）。

\par\textit{原典本文の校正と独自解説は未作成。分類のみの登録項目。}

\subsection{平成29年度 問5}
\label{q:h29-reg-apt-05}

\noindent\textbf{原典：} sources/H29/aptitude.pdf、PDF物理5頁。

\textbf{公式正答：} 3（single）。

\par\textit{原典本文の校正と独自解説は未作成。分類のみの登録項目。}

\subsection{平成29年度 問8}
\label{q:h29-reg-apt-08}

\noindent\textbf{原典：} sources/H29/aptitude.pdf、PDF物理8頁。

\textbf{公式正答：} 5（single）。

\par\textit{原典本文の校正と独自解説は未作成。分類のみの登録項目。}

\subsection{平成29年度 問9}
\label{q:h29-reg-apt-09}

\noindent\textbf{原典：} sources/H29/aptitude.pdf、PDF物理9頁。

\textbf{公式正答：} 3（single）。

\par\textit{原典本文の校正と独自解説は未作成。分類のみの登録項目。}

\subsection{平成29年度 問12}
\label{q:h29-reg-apt-12}

\noindent\textbf{原典：} sources/H29/aptitude.pdf、PDF物理11頁。

\textbf{公式正答：} 3（single）。

\par\textit{原典本文の校正と独自解説は未作成。分類のみの登録項目。}

\subsection{平成30年度 問9}
\label{q:h30-reg-apt-09}

\noindent\textbf{原典：} sources/H30/aptitude.pdf、PDF物理9頁。

\textbf{公式正答：} 3（single）。

\par\textit{原典本文の校正と独自解説は未作成。分類のみの登録項目。}

\subsection{平成30年度 問10}
\label{q:h30-reg-apt-10}

\noindent\textbf{原典：} sources/H30/aptitude.pdf、PDF物理10頁。

\textbf{公式正答：} 5（single）。

\par\textit{原典本文の校正と独自解説は未作成。分類のみの登録項目。}

\subsection{平成30年度 問11}
\label{q:h30-reg-apt-11}

\noindent\textbf{原典：} sources/H30/aptitude.pdf、PDF物理11頁。

\textbf{公式正答：} 3（single）。

\par\textit{原典本文の校正と独自解説は未作成。分類のみの登録項目。}

\subsection{令和元年度 問3}
\label{q:r01-reg-apt-03}

\noindent\textbf{原典：} sources/R01/aptitude.pdf、PDF物理4頁。

\textbf{公式正答：} 4（single）。

\par\textit{原典本文の校正と独自解説は未作成。分類のみの登録項目。}

\subsection{令和元年度 問7}
\label{q:r01-reg-apt-07}

\noindent\textbf{原典：} sources/R01/aptitude.pdf、PDF物理8頁。

\textbf{公式正答：} 5（single）。

\par\textit{原典本文の校正と独自解説は未作成。分類のみの登録項目。}

\subsection{令和元年度 問8}
\label{q:r01-reg-apt-08}

\noindent\textbf{原典：} sources/R01/aptitude.pdf、PDF物理9頁。

\textbf{公式正答：} 2（single）。

\par\textit{原典本文の校正と独自解説は未作成。分類のみの登録項目。}

\subsection{令和元年度 問11}
\label{q:r01-reg-apt-11}

\noindent\textbf{原典：} sources/R01/aptitude.pdf、PDF物理12頁。

\textbf{公式正答：} 2（single）。

\par\textit{原典本文の校正と独自解説は未作成。分類のみの登録項目。}

\subsection{令和元年度 問12}
\label{q:r01-reg-apt-12}

\noindent\textbf{原典：} sources/R01/aptitude.pdf、PDF物理13頁。

\textbf{公式正答：} 1（single）。

\par\textit{原典本文の校正と独自解説は未作成。分類のみの登録項目。}

\subsection{令和元年度 問13}
\label{q:r01-reg-apt-13}

\noindent\textbf{原典：} sources/R01/aptitude.pdf、PDF物理13頁。

\textbf{公式正答：} 4（single）。

\par\textit{原典本文の校正と独自解説は未作成。分類のみの登録項目。}

\subsection{令和元年度再試験 問5}
\label{q:r01-ret-apt-05}

\noindent\textbf{原典：} sources/R01\_retest/aptitude.pdf、PDF物理5頁。

\textbf{公式正答：} 3（single）。

\par\textit{原典本文の校正と独自解説は未作成。分類のみの登録項目。}

\subsection{令和元年度再試験 問7}
\label{q:r01-ret-apt-07}

\noindent\textbf{原典：} sources/R01\_retest/aptitude.pdf、PDF物理7頁。

\textbf{公式正答：} 3（single）。

\par\textit{原典本文の校正と独自解説は未作成。分類のみの登録項目。}

\subsection{令和元年度再試験 問12}
\label{q:r01-ret-apt-12}

\noindent\textbf{原典：} sources/R01\_retest/aptitude.pdf、PDF物理13頁。

\textbf{公式正答：} 2（single）。

\par\textit{原典本文の校正と独自解説は未作成。分類のみの登録項目。}

\subsection{令和2年度 問6}
\label{q:r02-reg-apt-06}

\noindent\textbf{原典：} sources/R02/aptitude.pdf、PDF物理6頁。

\textbf{公式正答：} 5（single）。

\par\textit{原典本文の校正と独自解説は未作成。分類のみの登録項目。}

\subsection{令和2年度 問7}
\label{q:r02-reg-apt-07}

\noindent\textbf{原典：} sources/R02/aptitude.pdf、PDF物理7頁。

\textbf{公式正答：} 4（single）。

\par\textit{原典本文の校正と独自解説は未作成。分類のみの登録項目。}

\subsection{令和2年度 問9}
\label{q:r02-reg-apt-09}

\noindent\textbf{原典：} sources/R02/aptitude.pdf、PDF物理9頁。

\textbf{公式正答：} 3（single）。

\par\textit{原典本文の校正と独自解説は未作成。分類のみの登録項目。}

\subsection{令和2年度 問13}
\label{q:r02-reg-apt-13}

\noindent\textbf{原典：} sources/R02/aptitude.pdf、PDF物理13頁。

\textbf{公式正答：} 4（single）。

\par\textit{原典本文の校正と独自解説は未作成。分類のみの登録項目。}

\subsection{令和3年度 問2}
\label{q:r03-reg-apt-02}

\noindent\textbf{原典：} sources/R03/aptitude.pdf、PDF物理2頁。

\textbf{公式正答：} 4（single）。

\par\textit{原典本文の校正と独自解説は未作成。分類のみの登録項目。}

\subsection{令和3年度 問8}
\label{q:r03-reg-apt-08}

\noindent\textbf{原典：} sources/R03/aptitude.pdf、PDF物理8頁。

\textbf{公式正答：} 3（single）。

\par\textit{原典本文の校正と独自解説は未作成。分類のみの登録項目。}

\subsection{令和3年度 問10}
\label{q:r03-reg-apt-10}

\noindent\textbf{原典：} sources/R03/aptitude.pdf、PDF物理10頁。

\textbf{公式正答：} 3（single）。

\par\textit{原典本文の校正と独自解説は未作成。分類のみの登録項目。}

\subsection{令和3年度 問15}
\label{q:r03-reg-apt-15}

\noindent\textbf{原典：} sources/R03/aptitude.pdf、PDF物理15頁。

\textbf{公式正答：} 3（single）。

\par\textit{原典本文の校正と独自解説は未作成。分類のみの登録項目。}

\subsection{令和4年度 問5}
\label{q:r04-reg-apt-05}

\noindent\textbf{原典：} sources/R04/aptitude.pdf、PDF物理4頁。

\textbf{公式正答：} 4（single）。

\par\textit{原典本文の校正と独自解説は未作成。分類のみの登録項目。}

\subsection{令和4年度 問6}
\label{q:r04-reg-apt-06}

\noindent\textbf{原典：} sources/R04/aptitude.pdf、PDF物理5頁。

\textbf{公式正答：} 5（single）。

\par\textit{原典本文の校正と独自解説は未作成。分類のみの登録項目。}

\subsection{令和4年度 問7}
\label{q:r04-reg-apt-07}

\noindent\textbf{原典：} sources/R04/aptitude.pdf、PDF物理6頁。

\textbf{公式正答：} 3（single）。

\par\textit{原典本文の校正と独自解説は未作成。分類のみの登録項目。}

\subsection{令和4年度 問8}
\label{q:r04-reg-apt-08}

\noindent\textbf{原典：} sources/R04/aptitude.pdf、PDF物理7頁。

\textbf{公式正答：} 2（single）。

\par\textit{原典本文の校正と独自解説は未作成。分類のみの登録項目。}

\subsection{令和4年度 問11}
\label{q:r04-reg-apt-11}

\noindent\textbf{原典：} sources/R04/aptitude.pdf、PDF物理10頁。

\textbf{公式正答：} 4（single）。

\par\textit{原典本文の校正と独自解説は未作成。分類のみの登録項目。}

\subsection{令和5年度 問3}
\label{q:r05-reg-apt-03}

\noindent\textbf{原典：} sources/R05/aptitude.pdf、PDF物理3頁。

\textbf{公式正答：} 2（single）。

\par\textit{原典本文の校正と独自解説は未作成。分類のみの登録項目。}

\subsection{令和5年度 問6}
\label{q:r05-reg-apt-06}

\noindent\textbf{原典：} sources/R05/aptitude.pdf、PDF物理6頁。

\textbf{公式正答：} 3（single）。

\par\textit{原典本文の校正と独自解説は未作成。分類のみの登録項目。}

\subsection{令和5年度 問8}
\label{q:r05-reg-apt-08}

\noindent\textbf{原典：} sources/R05/aptitude.pdf、PDF物理8頁。

\textbf{公式正答：} 2（single）。

\par\textit{原典本文の校正と独自解説は未作成。分類のみの登録項目。}

\subsection{令和5年度 問12}
\label{q:r05-reg-apt-12}

\noindent\textbf{原典：} sources/R05/aptitude.pdf、PDF物理12頁。

\textbf{公式正答：} 4（single）。

\par\textit{原典本文の校正と独自解説は未作成。分類のみの登録項目。}

\subsection{令和5年度 問14}
\label{q:r05-reg-apt-14}

\noindent\textbf{原典：} sources/R05/aptitude.pdf、PDF物理14頁。

\textbf{公式正答：} 3（single）。

\par\textit{原典本文の校正と独自解説は未作成。分類のみの登録項目。}

\subsection{令和6年度 問7}
\label{q:r06-reg-apt-07}

\noindent\textbf{原典：} sources/R06/aptitude.pdf、PDF物理7頁。

\textbf{公式正答：} 4（single）。

\par\textit{原典本文の校正と独自解説は未作成。分類のみの登録項目。}

\subsection{令和6年度 問13}
\label{q:r06-reg-apt-13}

\noindent\textbf{原典：} sources/R06/aptitude.pdf、PDF物理13頁。

\textbf{公式正答：} 5（single）。

\par\textit{原典本文の校正と独自解説は未作成。分類のみの登録項目。}

\subsection{令和6年度 問14}
\label{q:r06-reg-apt-14}

\noindent\textbf{原典：} sources/R06/aptitude.pdf、PDF物理14頁。

\textbf{公式正答：} 3（single）。

\par\textit{原典本文の校正と独自解説は未作成。分類のみの登録項目。}

\subsection{令和7年度 問12}
\label{q:r07-reg-apt-12}

\noindent\textbf{原典：} sources/R07/aptitude.pdf、PDF物理12頁。

\textbf{公式正答：} 4（single）。

\par\textit{原典本文の校正と独自解説は未作成。分類のみの登録項目。}

\subsection{令和7年度 問13}
\label{q:r07-reg-apt-13}

\noindent\textbf{原典：} sources/R07/aptitude.pdf、PDF物理13頁。

\textbf{公式正答：} 5（single）。

\par\textit{原典本文の校正と独自解説は未作成。分類のみの登録項目。}

\subsection{令和7年度 問14}
\label{q:r07-reg-apt-14}

\noindent\textbf{原典：} sources/R07/aptitude.pdf、PDF物理14頁。

\textbf{公式正答：} 2（single）。

\par\textit{原典本文の校正と独自解説は未作成。分類のみの登録項目。}

\subsection{令和7年度 問15}
\label{q:r07-reg-apt-15}

\noindent\textbf{原典：} sources/R07/aptitude.pdf、PDF物理15頁。

\textbf{公式正答：} 2（single）。

\par\textit{原典本文の校正と独自解説は未作成。分類のみの登録項目。}


\chapter{環境・持続可能性}
\noindent 分野解説は執筆中。以下は登録問題と校正済み例の作業用一覧。
\section{組換え生物の拡散防止・第一種第二種使用}

\subsection{平成25年度 問11}
\label{q:h25-reg-apt-11}

\noindent\textbf{原典：} sources/H25/aptitude.pdf、PDF物理12頁。

\textbf{公式正答：} 3（single）。

\paragraph{問題の要旨} 組換え生物の第二種使用で、衣服・放流・教育目的・廃棄を判定する。

\Needspace{4\baselineskip}\paragraph{確認した条件・記述（独自要約）}\mbox{}\par

\noindent\textbf{方向：} 個々の記述の○/×に適合する組合せ。\par

\noindent\textbf{並び：} ア・イ・ウ・エ\par

\noindent\textbf{ア：} 実験衣を着用し退出時には脱ぐ。 判定○：衣服を介した付着・移出を防ぐ適切な措置である。文科省Q3-7は必要な付着/感染防止措置を説明しており、全形態で一律の実験衣義務と読み替えない。\par

\noindent\textbf{イ：} 組換えオタマジャクシ・メダカを川へ放流する。 判定×：拡散防止された第二種使用の取扱いではない。動物はP1A等の措置を取り、野外利用は第一種の承認規程という別の手続を検討する。\par

\noindent\textbf{ウ：} 教育目的なら拡散防止措置は不要。 判定×：教育にも同じ法の適用がある。実験の目的だけで環境への拡散防止を除外できない。\par

\noindent\textbf{エ：} 不活化後の廃棄物を廃棄物関係法令等に従い処分。 判定○：不活化で生物学的危険を減らしても廃棄物の化学的・物理的性質に応じた処理の義務は消えない。\par

\paragraph{問われる核心} 拡散防止の使用区分と生物学的処置・廃棄物処理の両立

\paragraph{必要知識} 拡散防止の使用区分と生物学的処置・廃棄物処理の両立。アは○。衣服を介した付着・移出を防ぐ適切な措置である。文科省Q3-7は必要な付着/感染防止措置を説明しており、全形態で一律の実験衣義務と読み替えない。 イは×。拡散防止された第二種使用の取扱いではない。動物はP1A等の措置を取り、野外利用は第一種の承認規程という別の手続を検討する。 ウは×。教育にも同じ法の適用がある。実験の目的だけで環境への拡散防止を除外できない。 エは○。不活化で生物学的危険を減らしても廃棄物の化学的・物理的性質に応じた処理の義務は消えない。 よって○・×・×・○の候補3を選ぶ。

\paragraph{考え方・正答を選ぶ理由（査読待ち初稿）} アは○。衣服を介した付着・移出を防ぐ適切な措置である。文科省Q3-7は必要な付着/感染防止措置を説明しており、全形態で一律の実験衣義務と読み替えない。 イは×。拡散防止された第二種使用の取扱いではない。動物はP1A等の措置を取り、野外利用は第一種の承認規程という別の手続を検討する。 ウは×。教育にも同じ法の適用がある。実験の目的だけで環境への拡散防止を除外できない。 エは○。不活化で生物学的危険を減らしても廃棄物の化学的・物理的性質に応じた処理の義務は消えない。 よって○・×・×・○の候補3を選ぶ。

\Needspace{4\baselineskip}\paragraph{選択肢の検討（査読待ち初稿）}\mbox{}\par

\noindent\textbf{1：整合しない} ×・×・×・×。個別判断と異なるのはア・エである。各記述の理由は確認条件欄に示す。\par

\noindent\textbf{2：整合しない} ○・×・○・×。個別判断と異なるのはウ・エである。各記述の理由は確認条件欄に示す。\par

\noindent\textbf{3：整合する} ○・×・×・○。個別判断の○・×・×・○と全項目で一致する。\par

\noindent\textbf{4：整合しない} ○・×・×・×。個別判断と異なるのはエである。各記述の理由は確認条件欄に示す。\par

\noindent\textbf{5：整合しない} ×・○・○・○。個別判断と異なるのはア・イ・ウである。各記述の理由は確認条件欄に示す。\par

\paragraph{再利用する考え方} 活動・物質・制度の対象と評価の境界を固定し、用語の入替、全部/一部、法的義務/目標、出題時/現在を区別する。拡散防止の使用区分と生物学的処置・廃棄物処理の両立を個別の技術活動に結び付ける。

\paragraph{共通理論} \hyperref[theory:environment:biosafety]{第\ref*{theory:environment:biosafety}節}

\par 参考：\cite{env_biosafety_qa,env_biosafety_school,env_biosafety_law}。

\par 難易度 2（校準初稿）、優先度未評価。

\section{排出削減・オフセットとニュートラルの境界}

\subsection{令和元年度再試験 問11}
\label{q:r01-ret-apt-11}

\noindent\textbf{原典：} sources/R01\_retest/aptitude.pdf、PDF物理12頁。

\textbf{公式正答：} 1（single）。

\paragraph{問題の要旨} 温室効果ガス、低炭素社会、オフセットとニュートラルの定義を判定する。

\Needspace{4\baselineskip}\paragraph{確認した条件・記述（独自要約）}\mbox{}\par

\noindent\textbf{方向：} 個々の記述の○/×に適合する組合せ。\par

\noindent\textbf{並び：} ア・イ・ウ・エ\par

\noindent\textbf{ア：} 温室効果ガスはCO2・CH4・N2Oの3種類だけ。 判定×：HFC/PFC/SF6等もある。当初の京都の6種と自然の水蒸気を含む放射作用という別の集合を区別する。\par

\noindent\textbf{イ：} 暮らしの豊かさと環境保全が両立する低炭素社会を目指す。 判定○：排出を減らす社会構想であり、快適性の放棄又は無制限な排出の肯定ではない。\par

\noindent\textbf{ウ：} オフセットを残余排出全部の埋め合わせ済み状態とする。 判定×：この全量・状態の定義はニュートラルに対応する。オフセットは全部又は一部を埋め合わせる取組。\par

\noindent\textbf{エ：} ニュートラルを残余排出全部又は一部の埋め合わせ取組とする。 判定×：これはオフセットの定義であり、一部だけの相殺では対象全体が中立とは言えない。\par

\paragraph{問われる核心} 排出量会計の対象と全量/一部の相殺

\paragraph{必要知識} 排出量会計の対象と全量/一部の相殺。アは×。HFC/PFC/SF6等もある。当初の京都の6種と自然の水蒸気を含む放射作用という別の集合を区別する。 イは○。排出を減らす社会構想であり、快適性の放棄又は無制限な排出の肯定ではない。 ウは×。この全量・状態の定義はニュートラルに対応する。オフセットは全部又は一部を埋め合わせる取組。 エは×。これはオフセットの定義であり、一部だけの相殺では対象全体が中立とは言えない。 よって×・○・×・×の候補1を選ぶ。

\paragraph{考え方・正答を選ぶ理由（査読待ち初稿）} アは×。HFC/PFC/SF6等もある。当初の京都の6種と自然の水蒸気を含む放射作用という別の集合を区別する。 イは○。排出を減らす社会構想であり、快適性の放棄又は無制限な排出の肯定ではない。 ウは×。この全量・状態の定義はニュートラルに対応する。オフセットは全部又は一部を埋め合わせる取組。 エは×。これはオフセットの定義であり、一部だけの相殺では対象全体が中立とは言えない。 よって×・○・×・×の候補1を選ぶ。

\Needspace{4\baselineskip}\paragraph{選択肢の検討（査読待ち初稿）}\mbox{}\par

\noindent\textbf{1：整合する} ×・○・×・×。個別判断の×・○・×・×と全項目で一致する。\par

\noindent\textbf{2：整合しない} ×・×・○・○。個別判断と異なるのはイ・ウ・エである。各記述の理由は確認条件欄に示す。\par

\noindent\textbf{3：整合しない} ×・○・○・○。個別判断と異なるのはウ・エである。各記述の理由は確認条件欄に示す。\par

\noindent\textbf{4：整合しない} ○・○・×・×。個別判断と異なるのはアである。各記述の理由は確認条件欄に示す。\par

\noindent\textbf{5：整合しない} ○・○・○・○。個別判断と異なるのはア・ウ・エである。各記述の理由は確認条件欄に示す。\par

\paragraph{再利用する考え方} 活動・物質・制度の対象と評価の境界を固定し、用語の入替、全部/一部、法的義務/目標、出題時/現在を区別する。排出量会計の対象と全量/一部の相殺を個別の技術活動に結び付ける。

\paragraph{共通理論} \hyperref[theory:environment:carbon]{第\ref*{theory:environment:carbon}節}

\par 参考：\cite{env_kyoto2005,env_offset_guideline,env_neutral2011}。

\paragraph{関連問題} \hyperref[q:h30-reg-apt-13]{H30-REG-APT-13}, \hyperref[q:r06-reg-apt-05]{R06-REG-APT-05}

\par 難易度 2（校準初稿）、優先度未評価。

\subsection{令和6年度 問5}
\label{q:r06-reg-apt-05}

\noindent\textbf{原典：} sources/R06/aptitude.pdf、PDF物理5頁。

\textbf{公式正答：} 5（single）。

\par\textbf{検討未確定：} 公式5・初回5を維持。イの境界超過で回復不能変化が引き起こされるという強い表現を、原研究のリスク増大と区別。Oとする出題意図と科学的限定をAstra査読対象にする。

\paragraph{問題の要旨} 相殺の2用語、地球の限界、トップランナー基準を比較する。

\Needspace{4\baselineskip}\paragraph{確認した条件・記述（独自要約）}\mbox{}\par

\noindent\textbf{方向：} 個々の記述の○/×に適合する組合せ。\par

\noindent\textbf{並び：} ア・イ・ウ・エ\par

\noindent\textbf{ア：} ニュートラルを自ら削減後の全部又は一部のクレジット埋め合わせ取組とする。 判定×：一部相殺まで含むこの定義はオフセットである。全量の相殺済み状態を区別。\par

\noindent\textbf{イ：} 地球の限界の安全活動範囲を越すと自然資源の回復不能な変化が引き起こされる。 判定○：安全範囲と境界の概念は適切だが、一次研究は急激・不可逆変化のリスク増大と述べる。超過すれば必ず直ちに回復不能となる決定論には疑義。\par

\noindent\textbf{ウ：} 省エネ基準を商品化されている最良機器の性能以上に設定する。 判定○：最良性能を基礎とする目標である。2024原本の非化石エネルギー転換を含む法名を保持。すべての機器へ同じ数値を定める方式ではない。\par

\noindent\textbf{エ：} オフセットを残余排出の全部を埋め合わせた状態とする。 判定×：全量・状態のニュートラル定義との入替である。\par

\paragraph{問われる核心} 排出会計の全量条件と地球の限界・機器効率

\paragraph{必要知識} 排出会計の全量条件と地球の限界・機器効率。アは×。一部相殺まで含むこの定義はオフセットである。全量の相殺済み状態を区別。 イは○。安全範囲と境界の概念は適切だが、一次研究は急激・不可逆変化のリスク増大と述べる。超過すれば必ず直ちに回復不能となる決定論には疑義。 ウは○。最良性能を基礎とする目標である。2024原本の非化石エネルギー転換を含む法名を保持。すべての機器へ同じ数値を定める方式ではない。 エは×。全量・状態のニュートラル定義との入替である。 よって×・○・○・×の候補5を選ぶ。

\paragraph{考え方・正答を選ぶ理由（査読待ち初稿）} アは×。一部相殺まで含むこの定義はオフセットである。全量の相殺済み状態を区別。 イは○。安全範囲と境界の概念は適切だが、一次研究は急激・不可逆変化のリスク増大と述べる。超過すれば必ず直ちに回復不能となる決定論には疑義。 ウは○。最良性能を基礎とする目標である。2024原本の非化石エネルギー転換を含む法名を保持。すべての機器へ同じ数値を定める方式ではない。 エは×。全量・状態のニュートラル定義との入替である。 よって×・○・○・×の候補5を選ぶ。

\Needspace{4\baselineskip}\paragraph{選択肢の検討（査読待ち初稿）}\mbox{}\par

\noindent\textbf{1：整合しない} ○・○・○・○。個別判断と異なるのはア・エである。各記述の理由は確認条件欄に示す。\par

\noindent\textbf{2：整合しない} ×・○・○・○。個別判断と異なるのはエである。各記述の理由は確認条件欄に示す。\par

\noindent\textbf{3：整合しない} ○・×・×・○。個別判断と異なるのはア・イ・ウ・エである。各記述の理由は確認条件欄に示す。\par

\noindent\textbf{4：整合しない} ○・○・×・×。個別判断と異なるのはア・ウである。各記述の理由は確認条件欄に示す。\par

\noindent\textbf{5：整合する} ×・○・○・×。個別判断の×・○・○・×と全項目で一致する。\par

\paragraph{再利用する考え方} 活動・物質・制度の対象と評価の境界を固定し、用語の入替、全部/一部、法的義務/目標、出題時/現在を区別する。排出会計の全量条件と地球の限界・機器効率を個別の技術活動に結び付ける。

\paragraph{共通理論} \hyperref[theory:environment:carbon]{第\ref*{theory:environment:carbon}節}, \hyperref[theory:environment:planetary]{第\ref*{theory:environment:planetary}節}, \hyperref[theory:environment:toprunner]{第\ref*{theory:environment:toprunner}節}

\par 参考：\cite{env_neutral2011,env_offset_guideline,env_planetary,env_toprunner2019}。

\paragraph{関連問題} \hyperref[q:r01-ret-apt-11]{R01-RET-APT-11}

\par 難易度 2（校準初稿）、優先度未評価。

\section{循環資源・再使用再生利用と処理優先}

\subsection{平成24年度 問5}
\label{q:h24-reg-apt-05}

\noindent\textbf{原典：} sources/H24/aptitude.pdf、PDF物理4頁。

\textbf{公式正答：} 3（single）。

\paragraph{問題の要旨} 循環資源の五段階の処理優先順位を比較する。

\Needspace{4\baselineskip}\paragraph{確認した条件・記述（独自要約）}\mbox{}\par

\noindent\textbf{方向：} 最も適切な順\par

\paragraph{問われる核心} 資源価値を維持する利用段階と処理の一般優先

\paragraph{必要知識} 資源価値を維持する利用段階と処理の一般優先。第5〜7条の一般原則は発生抑制、再使用、再生利用、熱回収、適正処分であり3。第7条は技術的経済的可能性と環境負荷の低減を考慮する例外を設け、全物品に機械的に同じ処理を強制するものではない。排出者責任と拡大生産者責任は主体を別にした責任である。

\paragraph{考え方・正答を選ぶ理由（査読待ち初稿）} 第5〜7条の一般原則は発生抑制、再使用、再生利用、熱回収、適正処分であり3。第7条は技術的経済的可能性と環境負荷の低減を考慮する例外を設け、全物品に機械的に同じ処理を強制するものではない。排出者責任と拡大生産者責任は主体を別にした責任である。

\Needspace{4\baselineskip}\paragraph{選択肢の検討（査読待ち初稿）}\mbox{}\par

\noindent\textbf{1：整合しない} 再使用→発生抑制→再生利用→熱回収→適正処分。発生を避ける段階が再使用より後になっている。\par

\noindent\textbf{2：整合しない} 発生抑制→再使用→再生利用→適正処分→熱回収。熱回収可能なものを先に適正処分とする順が逆。\par

\noindent\textbf{3：整合する} 発生抑制→再使用→再生利用→熱回収→適正処分。発生を避け、形を保ち、材料を回収し、熱利用し、残りを処分する一般順位に一致。\par

\noindent\textbf{4：整合しない} 再使用→再生利用→発生抑制→熱回収→適正処分。発生抑制を既に発生した資源の利用より後に置く。\par

\noindent\textbf{5：整合しない} 再生利用→再使用→発生抑制→適正処分→熱回収。発生抑制・再使用・再生利用の順と熱回収/処分の順の双方が逆。\par

\paragraph{再利用する考え方} 活動・物質・制度の対象と評価の境界を固定し、用語の入替、全部/一部、法的義務/目標、出題時/現在を区別する。資源価値を維持する利用段階と処理の一般優先を個別の技術活動に結び付ける。

\paragraph{共通理論} \hyperref[theory:environment:circular]{第\ref*{theory:environment:circular}節}, \hyperref[theory:environment:responsibility]{第\ref*{theory:environment:responsibility}節}

\par 参考：\cite{env_circular_law,env_circular_outline}。

\paragraph{関連問題} \hyperref[q:r04-reg-apt-10]{R04-REG-APT-10}

\par 難易度 2（校準初稿）、優先度未評価。

\subsection{平成28年度 問14}
\label{q:h28-reg-apt-14}

\noindent\textbf{原典：} sources/H28/aptitude.pdf、PDF物理14頁。

\textbf{公式正答：} 5（single）。

\par\textbf{検討未確定：} 公式5・初回5を保持。エの小売引取義務は過去販売/同種買替条件が省略され、無条件義務と読むと不正確。2016年のテレビ対象と現在の有機EL追加を分ける。

\paragraph{問題の要旨} 基本法・廃棄物処理・管理票・家電の主体別責任を比較する。

\Needspace{4\baselineskip}\paragraph{確認した条件・記述（独自要約）}\mbox{}\par

\noindent\textbf{方向：} 個々の記述の○/×に適合する組合せ。\par

\noindent\textbf{並び：} ア・イ・ウ・エ\par

\noindent\textbf{ア：} 基本法と容器包装・家電・食品・建設資材等の個別法を整備。 判定○：基本原則と、製品ごとの回収処理制度という階層が異なる。\par

\noindent\textbf{イ：} 産業廃棄物は排出事業者の責任で適正に処理。 判定○：委託しても排出者の責任を無条件に移転できない。一般廃棄物/産業廃棄物は法の分類で、危険性の大小だけによらない。\par

\noindent\textbf{ウ：} 第三者委託時も最終処分まで必要な措置に努め管理票を使う。 判定○：管理票は運搬・中間処理・最終処分を確認する連絡記録である。費用を払っただけで管理が完了するわけではない。\par

\noindent\textbf{エ：} 家電4群で排出者が費用を負担し小売が引取り指定場所へ渡す。 判定○：制度概要は整合するが、小売の引取義務は過去販売したもの又は同種買替などの条件を持つ。原典は引取りを求められた対象機器とだけ記し条件を省略する。\par

\paragraph{問われる核心} 排出者責任と製品回収における主体・条件

\paragraph{必要知識} 排出者責任と製品回収における主体・条件。アは○。基本原則と、製品ごとの回収処理制度という階層が異なる。 イは○。委託しても排出者の責任を無条件に移転できない。一般廃棄物/産業廃棄物は法の分類で、危険性の大小だけによらない。 ウは○。管理票は運搬・中間処理・最終処分を確認する連絡記録である。費用を払っただけで管理が完了するわけではない。 エは○。制度概要は整合するが、小売の引取義務は過去販売したもの又は同種買替などの条件を持つ。原典は引取りを求められた対象機器とだけ記し条件を省略する。 よって○・○・○・○の候補5を選ぶ。

\paragraph{考え方・正答を選ぶ理由（査読待ち初稿）} アは○。基本原則と、製品ごとの回収処理制度という階層が異なる。 イは○。委託しても排出者の責任を無条件に移転できない。一般廃棄物/産業廃棄物は法の分類で、危険性の大小だけによらない。 ウは○。管理票は運搬・中間処理・最終処分を確認する連絡記録である。費用を払っただけで管理が完了するわけではない。 エは○。制度概要は整合するが、小売の引取義務は過去販売したもの又は同種買替などの条件を持つ。原典は引取りを求められた対象機器とだけ記し条件を省略する。 よって○・○・○・○の候補5を選ぶ。

\Needspace{4\baselineskip}\paragraph{選択肢の検討（査読待ち初稿）}\mbox{}\par

\noindent\textbf{1：整合しない} ○・○・○・×。個別判断と異なるのはエである。各記述の理由は確認条件欄に示す。\par

\noindent\textbf{2：整合しない} ○・○・×・○。個別判断と異なるのはウである。各記述の理由は確認条件欄に示す。\par

\noindent\textbf{3：整合しない} ○・×・○・○。個別判断と異なるのはイである。各記述の理由は確認条件欄に示す。\par

\noindent\textbf{4：整合しない} ×・○・○・○。個別判断と異なるのはアである。各記述の理由は確認条件欄に示す。\par

\noindent\textbf{5：整合する} ○・○・○・○。個別判断の○・○・○・○と全項目で一致する。\par

\paragraph{再利用する考え方} 活動・物質・制度の対象と評価の境界を固定し、用語の入替、全部/一部、法的義務/目標、出題時/現在を区別する。排出者責任と製品回収における主体・条件を個別の技術活動に結び付ける。

\paragraph{共通理論} \hyperref[theory:environment:responsibility]{第\ref*{theory:environment:responsibility}節}, \hyperref[theory:environment:appliance]{第\ref*{theory:environment:appliance}節}

\par 参考：\cite{env_circular_law,env_waste2003,env_appliance,env_appliance2024}。

\par 難易度 2（校準初稿）、優先度未評価。

\subsection{令和4年度 問10}
\label{q:r04-reg-apt-10}

\noindent\textbf{原典：} sources/R04/aptitude.pdf、PDF物理9頁。

\textbf{公式正答：} 4（single）。

\paragraph{問題の要旨} 循環型社会、循環的利用、再生利用、処理順位の定義を照合する。

\Needspace{4\baselineskip}\paragraph{確認した条件・記述（独自要約）}\mbox{}\par

\noindent\textbf{方向：} 個々の記述の○/×に適合する組合せ。\par

\noindent\textbf{並び：} ア・イ・ウ・エ\par

\noindent\textbf{ア：} 発生抑制・循環利用・適正処分で天然資源消費と環境負荷を低減。 判定○：第2条の循環型社会の目的構造に一致。\par

\noindent\textbf{イ：} 循環的利用は再使用・再生利用・熱回収。 判定○：利用の3方式と適正処分は別である。\par

\noindent\textbf{ウ：} 再生利用はそのまま製品又は部品として使用。 判定×：これは再使用の定義である。再生利用は原材料として利用する。\par

\noindent\textbf{エ：} 優先は抑制・再生利用・再使用・熱回収・処分。 判定×：一般順位の2位再使用と3位再生利用が入れ替わっている。\par

\paragraph{問われる核心} 製品・原材料・熱という資源利用形態と優先順位

\paragraph{必要知識} 製品・原材料・熱という資源利用形態と優先順位。アは○。第2条の循環型社会の目的構造に一致。 イは○。利用の3方式と適正処分は別である。 ウは×。これは再使用の定義である。再生利用は原材料として利用する。 エは×。一般順位の2位再使用と3位再生利用が入れ替わっている。 よって○・○・×・×の候補4を選ぶ。

\paragraph{考え方・正答を選ぶ理由（査読待ち初稿）} アは○。第2条の循環型社会の目的構造に一致。 イは○。利用の3方式と適正処分は別である。 ウは×。これは再使用の定義である。再生利用は原材料として利用する。 エは×。一般順位の2位再使用と3位再生利用が入れ替わっている。 よって○・○・×・×の候補4を選ぶ。

\Needspace{4\baselineskip}\paragraph{選択肢の検討（査読待ち初稿）}\mbox{}\par

\noindent\textbf{1：整合しない} ○・○・○・○。個別判断と異なるのはウ・エである。各記述の理由は確認条件欄に示す。\par

\noindent\textbf{2：整合しない} ×・○・×・○。個別判断と異なるのはア・エである。各記述の理由は確認条件欄に示す。\par

\noindent\textbf{3：整合しない} ○・×・○・×。個別判断と異なるのはイ・ウである。各記述の理由は確認条件欄に示す。\par

\noindent\textbf{4：整合する} ○・○・×・×。個別判断の○・○・×・×と全項目で一致する。\par

\noindent\textbf{5：整合しない} ○・×・○・○。個別判断と異なるのはイ・ウ・エである。各記述の理由は確認条件欄に示す。\par

\paragraph{再利用する考え方} 活動・物質・制度の対象と評価の境界を固定し、用語の入替、全部/一部、法的義務/目標、出題時/現在を区別する。製品・原材料・熱という資源利用形態と優先順位を個別の技術活動に結び付ける。

\paragraph{共通理論} \hyperref[theory:environment:circular]{第\ref*{theory:environment:circular}節}

\par 参考：\cite{env_circular_law}。

\paragraph{関連問題} \hyperref[q:h24-reg-apt-05]{H24-REG-APT-05}, \hyperref[q:r07-reg-apt-10]{R07-REG-APT-10}

\par 難易度 2（校準初稿）、優先度未評価。

\subsection{令和7年度 問10}
\label{q:r07-reg-apt-10}

\noindent\textbf{原典：} sources/R07/aptitude.pdf、PDF物理10頁。

\textbf{公式正答：} 3（single）。

\paragraph{問題の要旨} 循環資源、再使用、再生利用、熱回収、優先順位の5記述の適切数を求める。

\Needspace{4\baselineskip}\paragraph{確認した条件・記述（独自要約）}\mbox{}\par

\noindent\textbf{方向：} 適切な記述の数。\par

\noindent\textbf{ア：} 有価無価を問わない廃棄物等のうち有用なものを循環資源とする。○：法の対象を有価品だけに限定しない。\par

\noindent\textbf{イ：} 製品/部品として使うことを再生利用とする。×：再使用の定義。修理して利用することも再使用。\par

\noindent\textbf{ウ：} 原材料として利用することを再使用とする。×：再生利用の定義。\par

\noindent\textbf{エ：} 燃焼可能な資源を熱を得ることへ使う熱回収。○：材料回収とは利用形態が異なる。\par

\noindent\textbf{オ：} 抑制→再生利用→再使用→熱回収→処分。×：一般順位の再使用と再生利用が逆。\par

\paragraph{問われる核心} 循環資源の利用形態と法定一般順位

\paragraph{必要知識} 循環資源の利用形態と法定一般順位。ア○・イ×・ウ×・エ○・オ×で適切2件、候補3。再使用は製品や部品として形・機能を保つ利用、再生利用は原材料としての利用である。熱回収を加えた3利用と、発生抑制・適正処分を含む一般5段階を区別する。

\paragraph{考え方・正答を選ぶ理由（査読待ち初稿）} ア○・イ×・ウ×・エ○・オ×で適切2件、候補3。再使用は製品や部品として形・機能を保つ利用、再生利用は原材料としての利用である。熱回収を加えた3利用と、発生抑制・適正処分を含む一般5段階を区別する。

\Needspace{4\baselineskip}\paragraph{選択肢の検討（査読待ち初稿）}\mbox{}\par

\noindent\textbf{1：整合しない} 4。3誤りがあるため4件ではない。\par

\noindent\textbf{2：整合しない} 3。イウオを誤りとすると適切2件であり3件ではない。\par

\noindent\textbf{3：整合する} 2。アとエの2件で一致。\par

\noindent\textbf{4：整合しない} 1。正しいアエのどちらかを取り落とす。\par

\noindent\textbf{5：整合しない} 0。アエは法の定義と一致し0件ではない。\par

\paragraph{再利用する考え方} 活動・物質・制度の対象と評価の境界を固定し、用語の入替、全部/一部、法的義務/目標、出題時/現在を区別する。循環資源の利用形態と法定一般順位を個別の技術活動に結び付ける。

\paragraph{共通理論} \hyperref[theory:environment:circular]{第\ref*{theory:environment:circular}節}

\par 参考：\cite{env_circular_law}。

\paragraph{関連問題} \hyperref[q:r04-reg-apt-10]{R04-REG-APT-10}

\par 難易度 2（校準初稿）、優先度未評価。

\section{環境法・循環・持続可能性}

\subsection{平成23年度 問14}
\label{q:h23-reg-apt-14}

\noindent\textbf{原典：} sources/H23/aptitude.pdf、PDF物理15頁。

\textbf{公式正答：} 5（single）。

\par\textbf{検討未確定：} 肢1の2011年までの取組企業数増加とISO普及・費用上昇との因果は一次時系列未確認。定義の検証と分ける。

\paragraph{問題の要旨} 最終処分ゼロ、生物濃縮、多様性、持続可能性、排出取引の5概念を区別する。

\Needspace{4\baselineskip}\paragraph{確認した条件・記述（独自要約）}\mbox{}\par

\noindent\textbf{方向：} 不適切\par

\paragraph{問われる核心} 資源循環・生態系・世代間公平と排出の埋め合わせの境界

\paragraph{必要知識} 資源循環・生態系・世代間公平と排出の埋め合わせの境界。削減証明や排出枠は埋め合わせに使える手段であるが、手段・制度の名称とオフセットという活動は同一でないため5が不適切。1の狭義の定義と企業数の時系列主張を分離し、後者を未確認として残す。

\paragraph{考え方・正答を選ぶ理由（査読待ち初稿）} 削減証明や排出枠は埋め合わせに使える手段であるが、手段・制度の名称とオフセットという活動は同一でないため5が不適切。1の狭義の定義と企業数の時系列主張を分離し、後者を未確認として残す。

\Needspace{4\baselineskip}\paragraph{選択肢の検討（査読待ち初稿）}\mbox{}\par

\noindent\textbf{1：定義適切・時系列要確認} 狭義のゼロエミッションは生産廃棄物の最終処分ゼロで、国内の取組企業が増加。。最終処分量を対象にする狭い定義は企業事例と整合。ISO14000は規格群である。2011時点の取組企業数増加とその因果の一次統計は未確認。\par

\noindent\textbf{2：適切（蓄積性物質）} 蓄積性物質が食物連鎖上位の生物で高濃度となる。。難分解・蓄積性の物質では捕食による取込と排出の差で上位濃度が上がり得る。すべての物質で必ず上昇するという主張ではない。\par

\noindent\textbf{3：適切} 多様性は種・遺伝子・生態系の3階層と時間変化を含む。。同じ種の中の遺伝的差、種の間、相互作用する生態系の差を別階層とする。静止した種数だけに限定しない。\par

\noindent\textbf{4：適切} 将来のニーズを満たす能力を損なわず現在のニーズを満たす開発。。WCEDの定義は現在の開発を認めつつ将来能力を保持する世代間条件である。\par

\noindent\textbf{5：不適切} オフセットは先進国間で取引できる排出削減証明や排出枠の取引制度。。これは京都の排出量取引・クレジットの説明である。オフセットは自らの削減後の残余を他所の削減吸収で埋め合わせる取組で、政府間取引に限定されない。\par

\paragraph{再利用する考え方} 活動・物質・制度の対象と評価の境界を固定し、用語の入替、全部/一部、法的義務/目標、出題時/現在を区別する。資源循環・生態系・世代間公平と排出の埋め合わせの境界を個別の技術活動に結び付ける。

\paragraph{共通理論} \hyperref[theory:environment:boundary]{第\ref*{theory:environment:boundary}節}, \hyperref[theory:environment:zero]{第\ref*{theory:environment:zero}節}, \hyperref[theory:environment:accumulation]{第\ref*{theory:environment:accumulation}節}, \hyperref[theory:environment:biodiversity]{第\ref*{theory:environment:biodiversity}節}, \hyperref[theory:environment:development]{第\ref*{theory:environment:development}節}, \hyperref[theory:environment:carbon]{第\ref*{theory:environment:carbon}節}

\par 参考：\cite{env_zero2008,env_foodchain,env_biodiversity,env_development,env_offset2008}。

\par 難易度 2（校準初稿）、優先度未評価。

\subsection{平成30年度 問13}
\label{q:h30-reg-apt-13}

\noindent\textbf{原典：} sources/H30/aptitude.pdf、PDF物理13頁。

\textbf{公式正答：} 5（single）。

\par\textbf{検討未確定：} 公式5・初回5を保持。エ後半の生活に不要な物質の濃縮を一律に異常・環境問題とする表現は、毒性・曝露・自然濃縮を省略し一般化に疑義。

\paragraph{問題の要旨} 排出の埋め合わせ、持続可能性、廃棄物ゼロ、生物濃縮を判定する。

\Needspace{4\baselineskip}\paragraph{確認した条件・記述（独自要約）}\mbox{}\par

\noindent\textbf{方向：} 個々の記述の○/×に適合する組合せ。\par

\noindent\textbf{並び：} ア・イ・ウ・エ\par

\noindent\textbf{ア：} オフセットは削減後の残余を削減活動への投資等で埋め合わせる。 判定○：自らの削減と他所の削減吸収という二段の区別に一致。\par

\noindent\textbf{イ：} 1987年WCED報告で将来と現在のニーズを共に考える。 判定○：将来世代がニーズを満たす能力を損なわないという世代間条件を読む。後代の需要を完全予測する要求ではない。\par

\noindent\textbf{ウ：} 産業間で廃棄物・副産物を資源化し社会の排出ゼロを目指す。 判定○：社会システムの構想であり、現実の全施設で既に完全な物質・エネルギー排出ゼロを達成したとの主張ではない。\par

\noindent\textbf{エ：} 外界より高濃度の蓄積、特に不要元素物質の濃縮は異常・環境問題。 判定○：前半は濃縮の説明。後半を一律に読むと強過ぎる。健康影響には物質の毒性・濃度・曝露期間が必要で、濃縮や不要という性質だけで被害は決まらない。\par

\paragraph{問われる核心} 資源と排出の境界、濃縮と有害性の区別

\paragraph{必要知識} 資源と排出の境界、濃縮と有害性の区別。アは○。自らの削減と他所の削減吸収という二段の区別に一致。 イは○。将来世代がニーズを満たす能力を損なわないという世代間条件を読む。後代の需要を完全予測する要求ではない。 ウは○。社会システムの構想であり、現実の全施設で既に完全な物質・エネルギー排出ゼロを達成したとの主張ではない。 エは○。前半は濃縮の説明。後半を一律に読むと強過ぎる。健康影響には物質の毒性・濃度・曝露期間が必要で、濃縮や不要という性質だけで被害は決まらない。 よって○・○・○・○の候補5を選ぶ。

\paragraph{考え方・正答を選ぶ理由（査読待ち初稿）} アは○。自らの削減と他所の削減吸収という二段の区別に一致。 イは○。将来世代がニーズを満たす能力を損なわないという世代間条件を読む。後代の需要を完全予測する要求ではない。 ウは○。社会システムの構想であり、現実の全施設で既に完全な物質・エネルギー排出ゼロを達成したとの主張ではない。 エは○。前半は濃縮の説明。後半を一律に読むと強過ぎる。健康影響には物質の毒性・濃度・曝露期間が必要で、濃縮や不要という性質だけで被害は決まらない。 よって○・○・○・○の候補5を選ぶ。

\Needspace{4\baselineskip}\paragraph{選択肢の検討（査読待ち初稿）}\mbox{}\par

\noindent\textbf{1：整合しない} ×・○・○・○。個別判断と異なるのはアである。各記述の理由は確認条件欄に示す。\par

\noindent\textbf{2：整合しない} ○・×・○・○。個別判断と異なるのはイである。各記述の理由は確認条件欄に示す。\par

\noindent\textbf{3：整合しない} ○・○・×・○。個別判断と異なるのはウである。各記述の理由は確認条件欄に示す。\par

\noindent\textbf{4：整合しない} ○・○・○・×。個別判断と異なるのはエである。各記述の理由は確認条件欄に示す。\par

\noindent\textbf{5：整合する} ○・○・○・○。個別判断の○・○・○・○と全項目で一致する。\par

\paragraph{再利用する考え方} 活動・物質・制度の対象と評価の境界を固定し、用語の入替、全部/一部、法的義務/目標、出題時/現在を区別する。資源と排出の境界、濃縮と有害性の区別を個別の技術活動に結び付ける。

\paragraph{共通理論} \hyperref[theory:environment:carbon]{第\ref*{theory:environment:carbon}節}, \hyperref[theory:environment:development]{第\ref*{theory:environment:development}節}, \hyperref[theory:environment:zero]{第\ref*{theory:environment:zero}節}, \hyperref[theory:environment:accumulation]{第\ref*{theory:environment:accumulation}節}

\par 参考：\cite{env_offset2008,env_development,env_zero2008,env_accumulation}。

\paragraph{関連問題} \hyperref[q:r01-ret-apt-11]{R01-RET-APT-11}

\par 難易度 2（校準初稿）、優先度未評価。

\subsection{令和5年度 問15}
\label{q:r05-reg-apt-15}

\noindent\textbf{原典：} sources/R05/aptitude.pdf、PDF物理15頁。

\textbf{公式正答：} 3（single）。

\paragraph{問題の要旨} 環境基本法の典型7公害として列挙されないものを選ぶ。

\Needspace{4\baselineskip}\paragraph{確認した条件・記述（独自要約）}\mbox{}\par

\noindent\textbf{方向：} 列挙されないもの\par

\paragraph{問われる核心} 環境法の公害類型と行為・被害の区別

\paragraph{必要知識} 環境法の公害類型と行為・被害の区別。7項目は大気汚染・水質汚濁・土壌汚染・騒音・振動・地盤沈下・悪臭である。3は原因となる行為の名前であり列挙された作用の種類ではない。法の限定列挙と日常語の広い環境問題を区別する。

\paragraph{考え方・正答を選ぶ理由（査読待ち初稿）} 7項目は大気汚染・水質汚濁・土壌汚染・騒音・振動・地盤沈下・悪臭である。3は原因となる行為の名前であり列挙された作用の種類ではない。法の限定列挙と日常語の広い環境問題を区別する。

\Needspace{4\baselineskip}\paragraph{選択肢の検討（査読待ち初稿）}\mbox{}\par

\noindent\textbf{1：列挙される} 騒音。第2条3項に列挙される。\par

\noindent\textbf{2：条件付きで列挙される} 地盤の沈下。法は鉱物掘採のための土地掘削によるものを除く条件で列挙する。\par

\noindent\textbf{3：列挙されない} 廃棄物投棄。それ自体は7項目の列挙名称ではない。投棄が水質・土壌等の公害を生む場合や別法違反となることを否定しない。\par

\noindent\textbf{4：列挙される} 悪臭。第2条3項に列挙される。\par

\noindent\textbf{5：列挙される} 振動。第2条3項に列挙される。\par

\paragraph{再利用する考え方} 活動・物質・制度の対象と評価の境界を固定し、用語の入替、全部/一部、法的義務/目標、出題時/現在を区別する。環境法の公害類型と行為・被害の区別を個別の技術活動に結び付ける。

\paragraph{共通理論} \hyperref[theory:environment:pollution]{第\ref*{theory:environment:pollution}節}

\par 参考：\cite{env_pollution}。

\par 難易度 2（校準初稿）、優先度未評価。

\section{地域エネルギーと熱電併給・ZEH}

\subsection{令和2年度 問10}
\label{q:r02-reg-apt-10}

\noindent\textbf{原典：} sources/R02/aptitude.pdf、PDF物理10頁。

\textbf{公式正答：} 1（single）。

\par\textbf{検討未確定：} 2020原典のZEHの総称表現を保持。2018定義の対象用途と評価境界を限定し、CEMSを含む当時の全公的定義の逐語対照は追加照合。

\paragraph{問題の要旨} 地域のエネルギー管理・熱電併給・ZEH。

\Needspace{4\baselineskip}\paragraph{確認した条件・記述（独自要約）}\mbox{}\par

\noindent\textbf{原典：} 問題文・否定・全五候補と組合せ表を保存原典36画像で確認。\par

\noindent\textbf{時点：} 令和2年度の本文表現を保持。後の制度改正を遡及しない。\par

\paragraph{問われる核心} 地域のエネルギー管理・熱電併給・ZEH。

\paragraph{必要知識} 再生可能資源、Community Energy Management System、熱電併給、年間一次エネルギー収支。

\paragraph{考え方・正答を選ぶ理由（査読待ち初稿）} アの再生可能エネルギーの例、イの地域全体のCEMS、ウの電力と熱を同時に利用するコージェネレーション、エの年間の正味エネルギー収支をゼロ以下へ近づけるZEHは適切で1。ZEHは各瞬間の購入電力ゼロや、家電を含む全エネルギーを無条件にゼロとすることではない。

\Needspace{4\baselineskip}\paragraph{選択肢の検討（査読待ち初稿）}\mbox{}\par

\noindent\textbf{1：条件内で適切} OOOO。各文又は定義の判定を先に行う。全項目が一致する。\par

\noindent\textbf{2：不適切} XOOO。各文又は定義の判定を先に行う。並びが判定 OOOO と異なる。\par

\noindent\textbf{3：不適切} OXOO。各文又は定義の判定を先に行う。並びが判定 OOOO と異なる。\par

\noindent\textbf{4：不適切} OOXO。各文又は定義の判定を先に行う。並びが判定 OOOO と異なる。\par

\noindent\textbf{5：不適切} OOOX。各文又は定義の判定を先に行う。並びが判定 OOOO と異なる。\par

\paragraph{再利用する考え方} 再生可能資源、Community Energy Management System、熱電併給、年間一次エネルギー収支。 設問の否定・全称・基準年を確かめ、根拠と評価範囲を分けて判断する。

\paragraph{共通理論} \hyperref[theory:sustainability:energyterms]{第\ref*{theory:sustainability:energyterms}節}

\par 参考：\cite{eth54_cogen2018,eth54_zeh2018}。

\par 難易度 3（校準初稿）、優先度未評価。

\section{再生可能資源・熱移送と二次エネルギー}

\subsection{令和3年度 問11}
\label{q:r03-reg-apt-11}

\noindent\textbf{原典：} sources/R03/aptitude.pdf、PDF物理11頁。

\textbf{公式正答：} 3（single）。

\par\textbf{検討未確定：} 公式3・初回3を維持。問題幹の再生可能エネルギーは温室効果ガスを排出しないとの全称は、バイオガス燃焼やライフサイクルを含めれば不正確。パリ協定の採択2015/発効2016を区別。個別5記述の資源分類と幹の一般化を分離。

\paragraph{問題の要旨} 石炭、空気熱、水素、潮汐、バイオガスを資源の再生性と媒体で分類する。

\Needspace{4\baselineskip}\paragraph{確認した条件・記述（独自要約）}\mbox{}\par

\noindent\textbf{方向：} 個々の記述の○/×に適合する組合せ。\par

\noindent\textbf{並び：} ア・イ・ウ・エ・オ\par

\noindent\textbf{ア：} 石炭も植物から再び作られるため再生可能。 判定×：地質時間の形成は人間の利用時間で補充される再生可能性に当たらない。\par

\noindent\textbf{イ：} 空気熱をヒートポンプで温熱・冷熱利用。 判定○：自然界の熱を仕事を加えて移す技術である。温熱ではQ高=Q低+Wで、自然熱利用は入力仕事ゼロを意味しない。\par

\noindent\textbf{ウ：} 水素燃料は各種原料から製造するのでそれ自体は再生可能資源でない。 判定○：水素は二次エネルギー媒体で、製造元により再生可能性・排出が異なる。天然水素が絶対存在しないとはしない。\par

\noindent\textbf{エ：} 周期的潮汐のエネルギーを発電に利用。 判定○：人間の利用時間で周期的に補充される自然の流れを使う再生可能な供給である。\par

\noindent\textbf{オ：} 微生物分解で得るバイオガスは再生可能でない。 判定×：持続的に補充される生物資源由来のメタンはバイオマスエネルギーに含まれる。\par

\paragraph{問われる核心} 資源の再生時間、熱移送とエネルギー媒体

\paragraph{必要知識} 資源の再生時間、熱移送とエネルギー媒体。アは×。地質時間の形成は人間の利用時間で補充される再生可能性に当たらない。 イは○。自然界の熱を仕事を加えて移す技術である。温熱ではQ高=Q低+Wで、自然熱利用は入力仕事ゼロを意味しない。 ウは○。水素は二次エネルギー媒体で、製造元により再生可能性・排出が異なる。天然水素が絶対存在しないとはしない。 エは○。人間の利用時間で周期的に補充される自然の流れを使う再生可能な供給である。 オは×。持続的に補充される生物資源由来のメタンはバイオマスエネルギーに含まれる。 よって×・○・○・○・×の候補3を選ぶ。

\paragraph{考え方・正答を選ぶ理由（査読待ち初稿）} アは×。地質時間の形成は人間の利用時間で補充される再生可能性に当たらない。 イは○。自然界の熱を仕事を加えて移す技術である。温熱ではQ高=Q低+Wで、自然熱利用は入力仕事ゼロを意味しない。 ウは○。水素は二次エネルギー媒体で、製造元により再生可能性・排出が異なる。天然水素が絶対存在しないとはしない。 エは○。人間の利用時間で周期的に補充される自然の流れを使う再生可能な供給である。 オは×。持続的に補充される生物資源由来のメタンはバイオマスエネルギーに含まれる。 よって×・○・○・○・×の候補3を選ぶ。

\Needspace{4\baselineskip}\paragraph{選択肢の検討（査読待ち初稿）}\mbox{}\par

\noindent\textbf{1：整合しない} ○・○・○・○・○。個別判断と異なるのはア・オである。各記述の理由は確認条件欄に示す。\par

\noindent\textbf{2：整合しない} ○・×・○・×・○。個別判断と異なるのはア・イ・エ・オである。各記述の理由は確認条件欄に示す。\par

\noindent\textbf{3：整合する} ×・○・○・○・×。個別判断の×・○・○・○・×と全項目で一致する。\par

\noindent\textbf{4：整合しない} ×・○・×・○・×。個別判断と異なるのはウである。各記述の理由は確認条件欄に示す。\par

\noindent\textbf{5：整合しない} ×・×・×・×・○。個別判断と異なるのはイ・ウ・エ・オである。各記述の理由は確認条件欄に示す。\par

\paragraph{再利用する考え方} 活動・物質・制度の対象と評価の境界を固定し、用語の入替、全部/一部、法的義務/目標、出題時/現在を区別する。資源の再生時間、熱移送とエネルギー媒体を個別の技術活動に結び付ける。

\paragraph{共通理論} \hyperref[theory:environment:energy]{第\ref*{theory:environment:energy}節}

\par 参考：\cite{env_renewable2022,env_hydrogen2015,env_biomass,env_paris}。

\par 難易度 2（校準初稿）、優先度未評価。

\section{生物多様性と食品安全の別審査}

\subsection{令和2年度 問14}
\label{q:r02-reg-apt-14}

\noindent\textbf{原典：} sources/R02/aptitude.pdf、PDF物理14頁。

\textbf{公式正答：} 1（single）。

\par\textbf{検討未確定：} 2020当時の食品添加物の全例外と認定手続を網羅していない。未審査の全てを全技術・全品目へ広げず、規格基準の対象と手続の適用範囲を査読対象とする。2024後の行政所管変更を遡及しない。

\paragraph{問題の要旨} GMOの生物多様性審査と食品安全審査を分ける。

\Needspace{4\baselineskip}\paragraph{確認した条件・記述（独自要約）}\mbox{}\par

\noindent\textbf{原典：} 問題文・否定・全五候補と組合せ表を保存原典27画像で確認。\par

\noindent\textbf{時点：} 令和2年度の本文表現を保持。後の制度改正を遡及しない。\par

\paragraph{問われる核心} GMOの生物多様性審査と食品安全審査を分ける。

\paragraph{必要知識} 組換え形質、環境への拡散と生物多様性、食品衛生法上の審査、食品安全委員会の健康影響評価。

\paragraph{考え方・正答を選ぶ理由（査読待ち初稿）} 新しい形質の便益と生態系への影響を検討するア、カルタヘナ法の規制を述べるイ、食品の審査を述べるウ・エはいずれも正、1。環境での生物多様性影響と、食べた時の健康影響は別の評価対象であり、一方の承認だけで他方を不要としない。

\Needspace{4\baselineskip}\paragraph{選択肢の検討（査読待ち初稿）}\mbox{}\par

\noindent\textbf{1：条件内で適切} OOOO。各文又は定義の判定を先に行う。全項目が一致する。\par

\noindent\textbf{2：不適切} OOOX。各文又は定義の判定を先に行う。並びが判定 OOOO と異なる。\par

\noindent\textbf{3：不適切} OOXO。各文又は定義の判定を先に行う。並びが判定 OOOO と異なる。\par

\noindent\textbf{4：不適切} OXOO。各文又は定義の判定を先に行う。並びが判定 OOOO と異なる。\par

\noindent\textbf{5：不適切} XOOO。各文又は定義の判定を先に行う。並びが判定 OOOO と異なる。\par

\paragraph{再利用する考え方} 対象・責任主体・時点・要件を分け、文の否定と「全て」「不要」の範囲を点検する。知識を当てはめる前に何を証明する資料かを確かめる。

\paragraph{共通理論} \hyperref[theory:sustainability:gmo]{第\ref*{theory:sustainability:gmo}節}

\par 参考：\cite{saf53_cartagena,saf53_gmo_food2004,saf53_gmo_standards2000}。

\par 難易度 3（校準初稿）、優先度未評価。

\section{SDGsの普遍性・統合・人間の安全保障}

\subsection{令和元年度 問15}
\label{q:r01-reg-apt-15}

\noindent\textbf{原典：} sources/R01/aptitude.pdf、PDF物理14頁。

\textbf{公式正答：} 5（single）。

\paragraph{問題の要旨} SDGsの対象、原則、年限、目標数、参加者、範囲、評価の7記述を数える。

\Needspace{4\baselineskip}\paragraph{確認した条件・記述（独自要約）}\mbox{}\par

\noindent\textbf{方向：} 正しい記述の数。\par

\noindent\textbf{並び：} ア・イ・ウ・エ・オ・カ・キ\par

\noindent\textbf{ア：} 開発途上国だけの目標。 ×：先進国を含むすべての国へ適用する普遍性を持つ。\par

\noindent\textbf{イ：} 普遍性・包摂性・参画型・統合性・透明性。 ○：日本の実施指針が掲げる実施原則に一致し、最後は透明性と説明責任である。\par

\noindent\textbf{ウ：} 2030年を年限とする。 ○：2030アジェンダの年限である。個別ターゲットには2020年など異なる期限もある。\par

\noindent\textbf{エ：} 17の国際目標。 ○：目標は17、ターゲット169、指標はそれらと別の測定体系。\par

\noindent\textbf{オ：} 日本の取組は大企業・業界団体だけ。 ×：自治体・市民・個人等も参画し企業規模で排除しない。\par

\noindent\textbf{カ：} 環境問題に特化。 ×：経済・社会・環境を統合し、貧困・教育・平等なども含む。\par

\noindent\textbf{キ：} 指標と定期フォローアップ・評価公表。 ○：目標の達成を測る説明責任とレビューを必要とする。\par

\paragraph{問われる核心} SDGsの普遍性・統合性と説明責任

\paragraph{必要知識} SDGsの普遍性・統合性と説明責任。ア×・イ○・ウ○・エ○・オ×・カ×・キ○で、正しい4件、候補5。個々の理由は確認条件欄へ記録。2030期限と個別ターゲットの期限、目標/ターゲット/指標の階層を混同しない。

\paragraph{考え方・正答を選ぶ理由（査読待ち初稿）} ア×・イ○・ウ○・エ○・オ×・カ×・キ○で、正しい4件、候補5。個々の理由は確認条件欄へ記録。2030期限と個別ターゲットの期限、目標/ターゲット/指標の階層を混同しない。

\Needspace{4\baselineskip}\paragraph{選択肢の検討（査読待ち初稿）}\mbox{}\par

\noindent\textbf{1：整合しない} 0。イウエキの4件を全て除外している。\par

\noindent\textbf{2：整合しない} 1。4件中3件を取り落としている。\par

\noindent\textbf{3：整合しない} 2。4件中2件を取り落としている。\par

\noindent\textbf{4：整合しない} 3。4件中1件を取り落としている。\par

\noindent\textbf{5：整合する} 4。正しいイウエキの4件と一致。\par

\paragraph{再利用する考え方} 活動・物質・制度の対象と評価の境界を固定し、用語の入替、全部/一部、法的義務/目標、出題時/現在を区別する。SDGsの普遍性・統合性と説明責任を個別の技術活動に結び付ける。

\paragraph{共通理論} \hyperref[theory:environment:sdgs]{第\ref*{theory:environment:sdgs}節}, \hyperref[theory:environment:implementation]{第\ref*{theory:environment:implementation}節}

\par 参考：\cite{env_sdgs_agenda,env_sdgs2016,env_sdgs_un}。

\paragraph{関連問題} \hyperref[q:r04-reg-apt-14]{R04-REG-APT-14}

\par 難易度 2（校準初稿）、優先度未評価。

\subsection{令和元年度再試験 問14}
\label{q:r01-ret-apt-14}

\noindent\textbf{原典：} sources/R01\_retest/aptitude.pdf、PDF物理15頁。

\textbf{公式正答：} ALL（all\_awarded）。

\par\textbf{独立解答と公式正答に相違あり：} 独立解答 3。原典候補3の感梁症という誤植で的確な選択肢なし、公式全員得点。初回3は語を見落とした判断として不変保存。学習内容の指針照合と採点措置を分離。

\paragraph{問題の要旨} 2016年SDGs実施指針の4空欄と、原典選択肢の誤植による全員得点を記録する。

\Needspace{4\baselineskip}\paragraph{確認した条件・記述（独自要約）}\mbox{}\par

\noindent\textbf{空欄ア：} 地球規模の連鎖課題：感染症。\par

\noindent\textbf{空欄イ：} 従来の開発課題：開発途上国。\par

\noindent\textbf{空欄ウ：} 経済・社会と並ぶ環境。\par

\noindent\textbf{空欄エ：} 17目標と169ターゲット。\par

\noindent\textbf{原典：} 候補3の感梁症を感染症へ無断で校正しない。\par

\noindent\textbf{採点：} 公式正答ALL、全員得点。初回独立解答3。\par

\paragraph{問われる核心} SDGs実施指針の統合・普遍性と誤植の採点措置

\paragraph{必要知識} SDGs実施指針の統合・普遍性と誤植の採点措置。指針原文は感染症・開発途上国・環境・17となる。初回独立解答は3を保存したが、原本候補3には感梁症と印刷されている。公式注記は正答としていた候補の一部に誤植があり的確な選択肢がなく、受験者全員に得点を与える措置である。公式ALLを維持し、学習用に正しい語を示すことと原典の改変を区別する。

\paragraph{考え方・正答を選ぶ理由（査読待ち初稿）} 指針原文は感染症・開発途上国・環境・17となる。初回独立解答は3を保存したが、原本候補3には感梁症と印刷されている。公式注記は正答としていた候補の一部に誤植があり的確な選択肢がなく、受験者全員に得点を与える措置である。公式ALLを維持し、学習用に正しい語を示すことと原典の改変を区別する。

\Needspace{4\baselineskip}\paragraph{選択肢の検討（査読待ち初稿）}\mbox{}\par

\noindent\textbf{1：意図と不整合・全員得点} 国際紛争・先進国・環境・15。指針の感染症・開発途上国・環境・17とア・イ・エが異なる。\par

\noindent\textbf{2：意図と不整合・全員得点} 感染症・先進国・教育・15。イ・ウ・エが異なる。\par

\noindent\textbf{3：原典誤植・全員得点} 感梁症・開発途上国・環境・17（原典の字）。意図上近いが原典は感梁症という誤植であり、的確な候補でない。初回は感染症として3を選んだが、この見落としを後続の原本再照合で発見した。\par

\noindent\textbf{4：意図と不整合・全員得点} 国際紛争・開発途上国・教育・17。ア・ウが異なる。\par

\noindent\textbf{5：意図と不整合・全員得点} 感染症・開発途上国・教育・17。ウの教育が三側面の環境と異なる。\par

\paragraph{再利用する考え方} 活動・物質・制度の対象と評価の境界を固定し、用語の入替、全部/一部、法的義務/目標、出題時/現在を区別する。SDGs実施指針の統合・普遍性と誤植の採点措置を個別の技術活動に結び付ける。

\paragraph{共通理論} \hyperref[theory:environment:sdgs]{第\ref*{theory:environment:sdgs}節}, \hyperref[theory:environment:implementation]{第\ref*{theory:environment:implementation}節}

\par 参考：\cite{env_sdgs2016}。

\paragraph{関連問題} \hyperref[q:r06-reg-apt-06]{R06-REG-APT-06}

\par 難易度 2（校準初稿）、優先度未評価。

\subsection{令和3年度 問5}
\label{q:r03-reg-apt-05}

\noindent\textbf{原典：} sources/R03/aptitude.pdf、PDF物理5頁。

\textbf{公式正答：} 4（single）。

\paragraph{問題の要旨} SDGsの参加者、法的拘束力、背景、未来からの計画を判定する。

\Needspace{4\baselineskip}\paragraph{確認した条件・記述（独自要約）}\mbox{}\par

\noindent\textbf{方向：} 個々の記述の○/×に適合する組合せ。\par

\noindent\textbf{並び：} ア・イ・ウ・エ\par

\noindent\textbf{ア：} すべての参加者による誰一人取り残さない枠組み。 判定○：全ての国とstakeholderを対象とする。\par

\noindent\textbf{イ：} SDGsそのものに法的拘束力があり罰則対象。 判定×：SDGsの政治的目標を、それ自体罰則付きの条約義務としない。関連する国内法の遵守義務は別。\par

\noindent\textbf{ウ：} 環境・格差・紛争など将来への危機感を背景とする。 判定○：経済社会環境の相互作用を扱う課題設定である。\par

\noindent\textbf{エ：} バックキャストでなく現在からのフォーキャストが重要。 判定×：目標の将来像から現在へ必要条件を逆に求める方法が目標実現を設計する。実現可能性の分析を不要とするのでなく、将来像を現状延長だけで置き換えない。\par

\paragraph{問われる核心} SDGsの目標と法的義務、将来像からの計画

\paragraph{必要知識} SDGsの目標と法的義務、将来像からの計画。アは○。全ての国とstakeholderを対象とする。 イは×。SDGsの政治的目標を、それ自体罰則付きの条約義務としない。関連する国内法の遵守義務は別。 ウは○。経済社会環境の相互作用を扱う課題設定である。 エは×。目標の将来像から現在へ必要条件を逆に求める方法が目標実現を設計する。実現可能性の分析を不要とするのでなく、将来像を現状延長だけで置き換えない。 よって○・×・○・×の候補4を選ぶ。

\paragraph{考え方・正答を選ぶ理由（査読待ち初稿）} アは○。全ての国とstakeholderを対象とする。 イは×。SDGsの政治的目標を、それ自体罰則付きの条約義務としない。関連する国内法の遵守義務は別。 ウは○。経済社会環境の相互作用を扱う課題設定である。 エは×。目標の将来像から現在へ必要条件を逆に求める方法が目標実現を設計する。実現可能性の分析を不要とするのでなく、将来像を現状延長だけで置き換えない。 よって○・×・○・×の候補4を選ぶ。

\Needspace{4\baselineskip}\paragraph{選択肢の検討（査読待ち初稿）}\mbox{}\par

\noindent\textbf{1：整合しない} ○・×・○・○。個別判断と異なるのはエである。各記述の理由は確認条件欄に示す。\par

\noindent\textbf{2：整合しない} ○・○・○・×。個別判断と異なるのはイである。各記述の理由は確認条件欄に示す。\par

\noindent\textbf{3：整合しない} ×・○・×・○。個別判断と異なるのはア・イ・ウ・エである。各記述の理由は確認条件欄に示す。\par

\noindent\textbf{4：整合する} ○・×・○・×。個別判断の○・×・○・×と全項目で一致する。\par

\noindent\textbf{5：整合しない} ×・×・○・○。個別判断と異なるのはア・エである。各記述の理由は確認条件欄に示す。\par

\paragraph{再利用する考え方} 活動・物質・制度の対象と評価の境界を固定し、用語の入替、全部/一部、法的義務/目標、出題時/現在を区別する。SDGsの目標と法的義務、将来像からの計画を個別の技術活動に結び付ける。

\paragraph{共通理論} \hyperref[theory:environment:sdgs]{第\ref*{theory:environment:sdgs}節}, \hyperref[theory:environment:planning]{第\ref*{theory:environment:planning}節}

\par 参考：\cite{env_sdgs_agenda,env_sdgs_binding,env_sdgs2016}。

\par 難易度 2（校準初稿）、優先度未評価。

\subsection{令和4年度 問14}
\label{q:r04-reg-apt-14}

\noindent\textbf{原典：} sources/R04/aptitude.pdf、PDF物理13頁。

\textbf{公式正答：} 3（single）。

\paragraph{問題の要旨} SDGsの対象国、MDGsとの継承、3側面、人権と参加の7記述を判定する。

\Needspace{4\baselineskip}\paragraph{確認した条件・記述（独自要約）}\mbox{}\par

\noindent\textbf{方向：} 個々の記述の○/×に適合する組合せ。\par

\noindent\textbf{並び：} ア・イ・ウ・エ・オ・カ・キ\par

\noindent\textbf{ア：} 先進国だけが実行する。 判定×：先進国と途上国の双方を含む普遍性に反する。\par

\noindent\textbf{イ：} MDGsの未達成課題を引き継いで全うする。 判定○：前文のMDGsを基に残課題を完成する趣旨。\par

\noindent\textbf{ウ：} 経済・社会・環境を調和させる。 判定○：持続可能な開発の3側面を統合する。\par

\noindent\textbf{エ：} 誰一人取り残さない。 判定○：脆弱な人々と不平等へ焦点を当てる包摂性。\par

\noindent\textbf{オ：} 人権とジェンダー平等・女性女児の能力強化。 判定○：前文と目標5、段落20に対応。\par

\noindent\textbf{カ：} 全stakeholderが協同的パートナーシップで実行。 判定○：政府だけに限定しない。\par

\noindent\textbf{キ：} 環境問題に特化する。 判定×：教育・貧困・平等・経済等も含むため限定が誤り。\par

\paragraph{問われる核心} SDGsの普遍・包摂・統合と協働

\paragraph{必要知識} SDGsの普遍・包摂・統合と協働。アは×。先進国と途上国の双方を含む普遍性に反する。 イは○。前文のMDGsを基に残課題を完成する趣旨。 ウは○。持続可能な開発の3側面を統合する。 エは○。脆弱な人々と不平等へ焦点を当てる包摂性。 オは○。前文と目標5、段落20に対応。 カは○。政府だけに限定しない。 キは×。教育・貧困・平等・経済等も含むため限定が誤り。 よって×・○・○・○・○・○・×の候補3を選ぶ。

\paragraph{考え方・正答を選ぶ理由（査読待ち初稿）} アは×。先進国と途上国の双方を含む普遍性に反する。 イは○。前文のMDGsを基に残課題を完成する趣旨。 ウは○。持続可能な開発の3側面を統合する。 エは○。脆弱な人々と不平等へ焦点を当てる包摂性。 オは○。前文と目標5、段落20に対応。 カは○。政府だけに限定しない。 キは×。教育・貧困・平等・経済等も含むため限定が誤り。 よって×・○・○・○・○・○・×の候補3を選ぶ。

\Needspace{4\baselineskip}\paragraph{選択肢の検討（査読待ち初稿）}\mbox{}\par

\noindent\textbf{1：整合しない} ×・×・○・○・○・○・○。個別判断と異なるのはイ・キである。各記述の理由は確認条件欄に示す。\par

\noindent\textbf{2：整合しない} ×・○・×・○・×・○・×。個別判断と異なるのはウ・オである。各記述の理由は確認条件欄に示す。\par

\noindent\textbf{3：整合する} ×・○・○・○・○・○・×。個別判断の×・○・○・○・○・○・×と全項目で一致する。\par

\noindent\textbf{4：整合しない} ○・×・○・×・○・×・○。個別判断と異なるのはア・イ・エ・カ・キである。各記述の理由は確認条件欄に示す。\par

\noindent\textbf{5：整合しない} ×・○・○・○・○・×・×。個別判断と異なるのはカである。各記述の理由は確認条件欄に示す。\par

\paragraph{再利用する考え方} 活動・物質・制度の対象と評価の境界を固定し、用語の入替、全部/一部、法的義務/目標、出題時/現在を区別する。SDGsの普遍・包摂・統合と協働を個別の技術活動に結び付ける。

\paragraph{共通理論} \hyperref[theory:environment:sdgs]{第\ref*{theory:environment:sdgs}節}

\par 参考：\cite{env_sdgs_agenda,env_sdgs_un}。

\paragraph{関連問題} \hyperref[q:r01-reg-apt-15]{R01-REG-APT-15}

\par 難易度 2（校準初稿）、優先度未評価。

\subsection{令和6年度 問2}
\label{q:r06-reg-apt-02}

\noindent\textbf{原典：} sources/R06/aptitude.pdf、PDF物理2頁。

\textbf{公式正答：} 5（single）。

\paragraph{問題の要旨} 2030アジェンダのジェンダー段落の4空欄に入らない語を選ぶ。

\Needspace{4\baselineskip}\paragraph{確認した条件・記述（独自要約）}\mbox{}\par

\noindent\textbf{並び：} ア能力強化（エンパワーメント）・イ雇用・ウリーダーシップ・エ意思決定。\par

\noindent\textbf{方向：} 4空欄のいずれにも入らない語。\par

\paragraph{問われる核心} SDGsのジェンダー平等と指定文脈の能力・参画

\paragraph{必要知識} SDGsのジェンダー平等と指定文脈の能力・参画。段落20の並びは能力強化、雇用、リーダーシップ、意思決定である。よってケア5がどの空欄にも入らない。文脈の範囲を指定段落に固定し、他のターゲットに語があることと区別する。

\paragraph{考え方・正答を選ぶ理由（査読待ち初稿）} 段落20の並びは能力強化、雇用、リーダーシップ、意思決定である。よってケア5がどの空欄にも入らない。文脈の範囲を指定段落に固定し、他のターゲットに語があることと区別する。

\Needspace{4\baselineskip}\paragraph{選択肢の検討（査読待ち初稿）}\mbox{}\par

\noindent\textbf{1：イに対応} 雇用。イに入り、働く機会の平等を示す。\par

\noindent\textbf{2：ウに対応} リーダーシップ。ウに入り、指導的役割へ参画する機会。\par

\noindent\textbf{3：エに対応} 意思決定。エに入り、決定過程で同等の機会。\par

\noindent\textbf{4：アに対応} エンパワーメント。アに入り、女性女児の能力強化。外務省仮訳は能力強化と表す。\par

\noindent\textbf{5：どれにも入らない} ケア。この段落の4空欄には入らない。ただし目標5.4は無償のケア/家事労働を扱い、SDGs全体でケアが無関係という結論ではない。\par

\paragraph{再利用する考え方} 活動・物質・制度の対象と評価の境界を固定し、用語の入替、全部/一部、法的義務/目標、出題時/現在を区別する。SDGsのジェンダー平等と指定文脈の能力・参画を個別の技術活動に結び付ける。

\paragraph{共通理論} \hyperref[theory:environment:gender]{第\ref*{theory:environment:gender}節}

\par 参考：\cite{env_sdgs_agenda}。

\paragraph{関連問題} \hyperref[q:r06-reg-apt-06]{R06-REG-APT-06}

\par 難易度 2（校準初稿）、優先度未評価。

\subsection{令和6年度 問6}
\label{q:r06-reg-apt-06}

\noindent\textbf{原典：} sources/R06/aptitude.pdf、PDF物理6頁。

\textbf{公式正答：} 2（single）。

\paragraph{問題の要旨} SDGsの3側面、統合性、普遍性、人間の安全保障の4空欄に入らない語を選ぶ。

\Needspace{4\baselineskip}\paragraph{確認した条件・記述（独自要約）}\mbox{}\par

\noindent\textbf{方向：} 指定4空欄のどれにも入らない。\par

\noindent\textbf{並び：} ア経済・イ統合・ウ普遍・エ安全。\par

\paragraph{問われる核心} SDGsの統合・普遍性と人間の安全保障

\paragraph{必要知識} SDGsの統合・普遍性と人間の安全保障。ア経済、イ統合、ウ普遍、エ安全と対応し、候補2の平和が空欄にない。安全保障は国家間の軍事だけでなく、脆弱な一人ひとりが尊厳ある生活を送れるという人間中心の概念である。

\paragraph{考え方・正答を選ぶ理由（査読待ち初稿）} ア経済、イ統合、ウ普遍、エ安全と対応し、候補2の平和が空欄にない。安全保障は国家間の軍事だけでなく、脆弱な一人ひとりが尊厳ある生活を送れるという人間中心の概念である。

\Needspace{4\baselineskip}\paragraph{選択肢の検討（査読待ち初稿）}\mbox{}\par

\noindent\textbf{1：エに対応} 安全。エの人間の安全保障を構成する。\par

\noindent\textbf{2：空欄にない} 平和。この4空欄には入らない。平和はSDG16や前文の重要事項であり、SDGsの対象外という意味ではない。\par

\noindent\textbf{3：アに対応} 経済。アに入り社会・経済・環境の3側面を成す。\par

\noindent\textbf{4：ウに対応} 普遍。ウに入りすべての国の目標という普遍性。\par

\noindent\textbf{5：イに対応} 統合。イに入り相互関係のある目標を統合的に解決する。\par

\paragraph{再利用する考え方} 活動・物質・制度の対象と評価の境界を固定し、用語の入替、全部/一部、法的義務/目標、出題時/現在を区別する。SDGsの統合・普遍性と人間の安全保障を個別の技術活動に結び付ける。

\paragraph{共通理論} \hyperref[theory:environment:sdgs]{第\ref*{theory:environment:sdgs}節}, \hyperref[theory:environment:implementation]{第\ref*{theory:environment:implementation}節}

\par 参考：\cite{env_sdgs_pamphlet,env_humansecurity}。

\paragraph{関連問題} \hyperref[q:r01-ret-apt-14]{R01-RET-APT-14}, \hyperref[q:r06-reg-apt-02]{R06-REG-APT-02}

\par 難易度 2（校準初稿）、優先度未評価。


\chapter{研究・情報倫理}
\section{背景：他者が検証できる知識を蓄積する}
研究では観察、仮説、検証、批判を通じて知識を更新する。結論が後で否定される可能性を含むため、結論が常に正しいことと、研究過程が誠実であることは同じではない。例えば計測器の応答を合理的に評価した上で報告した結果に、新しい測定から修正が必要と分かる場合がある。この更新を不正として罰すれば、新しい仮説を検証する営みを妨げる。一方、都合の悪い測定点を隠して結論を作れば、第三者の判断材料を壊してしまう。

本章では、知識の\emph{根拠}、功績と責任の\emph{帰属}、判断に影響する\emph{利害}という3つを追跡する。技術者が研究成果を設計や安全判断へ利用するときにも、値だけでなく、それがどの過程で得られ、誰が説明し、どの利害に影響され得るかを確認する必要がある。

\section{科学者の自律性と社会への説明}
\label{theory:research:norms}
日本学術会議の2013年改訂規範は、人文・社会科学を含む研究者と科学的知識を利用する専門職を対象とする。独立した専門判断は社会との対話と両立し、科学的助言では不確実性や意見の相違も説明する\parencite{scj_code2013}。

例えば地盤の安全評価で2つのモデルが異なる予測を与える場合、モデルの仮定、利用可能な観測、結論が変わる範囲を示す。その上で、許容できる危険、費用、生活への影響などを含めた政策判断を決定者が行う。多数決でモデル間の差を消して1つの値だけを提示すれば、どの条件で結論が変わるかという情報を失う。科学者は助言の根拠へ責任を持つが、科学的結果だけから社会の選好が一意に定まるとは考えない。

研究の便益と悪用の可能性がともにあることを利用の両義性という。便益が期待されるという判断だけでは、公表する内容や方法についての検討を尽くしたことにならない。平成26年度問2は独立性と対話、令和元年度再試験問6は両義性、令和5年度問7は相違の説明と政策判断をそれぞれ問う。

\section{FFP、特定不正行為、誠実な誤り}
\label{theory:research:ffp}
捏造は架空のデータ・結果の作成、改ざんは取得や処理を操作して結果を不真正にすること、盗用は他者の成果を適切な了解・表示なく流用することである。2014年文部科学省ガイドラインの特定不正行為は、故意又は基本的注意義務を著しく怠った、公表成果中のこの3種を対象とする。正当に得られた結果が後に誤りと判明した場合は除外される\parencite{mext_misconduct2014}。

\begin{center}
\begin{tabular}{@{}p{0.25\linewidth}p{0.70\linewidth}@{}}
\toprule
問い & 判断の手掛かり\\
\midrule
データが存在したか & 実施しない測定の値を作るなら捏造である。\\
実データをどう扱ったか & 採否や補正の操作によって、報告結果が本来の記録と違うものになっていないかを確認する。\\
知的な寄与を誰へ帰属させたか & 他者の考案・解析・文章などを自分の成果として提示していないかを確認する。\\
結論が後に変わった理由は何か & 新しい証拠による更新と、元の根拠の不正な作成・操作を区別する。\\
\bottomrule
\end{tabular}
\end{center}

特定不正行為に該当しないことは、倫理的に許されることの十分条件ではない。例えば同じ結果を独立した新成果に見せて複数誌へ投稿すれば、知識の新規性と業績評価をゆがめる。平成28年度問4では「特定」が付いた限定と、不正一般を混同する記述が誤答を誘う。研究費の不正使用も、研究結果のFFPと対象が違う。実際には両者が同時に起こり得るため、用語だけで相互排他と考えない。

\section{研究記録：再現の経路と変更の履歴}
\label{theory:research:records}
研究ノートの役割は、過去の自分だけが分かる備忘録にとどまらない。例えば波速を求めた測定では、試料番号、寸法、温度、装置、波形ファイル、到達時刻の決め方、解析版を対応付ければ、別の者が報告値へ至る経路を追える。観測事実と解釈を分けることで、結論を更新しても元データの意味を失わずに済む。

誤りを訂正する場合は、元記録を消す代わりに、何をいつなぜ訂正したかを追加する。紙なら差替えや消去を防ぐ形式、電子なら編集履歴とアクセス管理などにより、この目的を実現する。ノートの帰属・保管・異動時の扱いは、個人が書いたという事実だけで決まらず、契約や機関規程を確認する\parencite{jsps_integrity2025}。平成27年度問8の「原則として綴じた紙と消せない筆記具」という記述を、電子記録の一律禁止へ拡張してはいけない。

知的財産をめぐる説明でも研究の経緯を示す記録は役立つ。ただしノートがあるだけで、発明者、契約上の帰属、特許の各要件がすべて満たされるわけではない。記録の証拠としての役割と、権利成立の要件は別に検討する。

\section{著者資格：功績を受けることと内容を説明すること}
\label{theory:research:authorship}
論文の著者表示は研究への寄与を認めると同時に、内容へ責任を負う者を示す。資金獲得だけでは資格を満たさず、資格のない者を追加する贈与著者、資格ある者を省く幽霊著者は不適切である。貢献と責任の基準を確認し、著者資格に至らない援助は謝辞などで適切に表示する\parencite{jsps_integrity2025}。

共同実験で測定を担当した者は、少なくとも測定方法、記録、結果を説明し、原稿でどのように使われたかを確認する必要がある。他の担当者の分析に疑問が生じた場合も、単に自分の部分ではないと応答して済ませるものではない。ただし役割分担を明示した共同研究の責任と、全著者が全実験を独力で再実施することは同義ではない。平成29年度問14は、説明できない実験貢献者の著者資格と、責任範囲の一律限定を問う。

発表の単位にも透明性が必要である。同じ成果を別言語へ翻訳しただけで新しい独立成果にはならない。共同成果の再利用では共著者との合意に加え、出典表示、出版側の権利、投稿・出版規程を確認する。研究を分割した複数論文についても、それぞれに独立した問いと新規な寄与があるか、重複部分を読者が把握できるかを検討する。多くの論文を書くという数だけでは、適否を決められない。

\section{利益相反：判断への利害を開示して管理する}
\label{theory:research:coi}
利益相反とは、公的・専門的な判断に対して私的な利害が影響し得る状態である。実際に判断をゆがめたことの証明だけが問題になるのではない。自社株を持つ評価者が「私は公平である」と考えても、評価結果と自己の利益が関係しているという構造は残る。家族の経済的利害や査読での競争関係も検討対象となり、開示と機関による管理が必要である\parencite{jsps_integrity2025}。

判断の手順は、\emph{役割を列挙する}、\emph{利害の経路を示す}、\emph{第三者へ相談する}、\emph{条件や回避措置を定める}という順になる。大学教員の兼業、設備の利用、企業の株式、査読では経路が異なる。適切な事前相談は、私企業への無償設備利用が必ず許可されることを意味しない。許可する者が機関の規程と資金の条件を検討できるようにする行為である。令和2年度問3ではこの違いを読む。

\section{予防の仕組みと制度の時点}
\label{theory:research:prevention}
研究不正を予防するには、個人の誠実さを、検証可能な記録、明確な役割、相談できる仕組みに結び付ける。研究者が異動してもデータの所在が分かること、若手が疑問を相談できること、公開前に方法と結果を互いに確認することは、その具体例である。

\begin{tcolorbox}[title=出題資料と現在の参照資料]
この初稿で扱う日本学術会議の問題は2013年改訂声明、平成28年度・令和元年度再試験の問題は2014年文部科学省ガイドラインに基づく。2025年公開の日本学術振興会教材第2版は現在の背景確認に用い、2015〜2024年の問題が第2版を引用したとは記さない。資金に係る措置では対象となる研究活動・関与者・資金と事案の条件を確認し、倫理責任が全研究にあることから全研究費へ一律の返還義務を推論しない。
\end{tcolorbox}

\section{問題で確認する順序}
まず適切・不適切のどちらを数えるかを確定する。組合せ問題では各記述を独立に判定してから候補を選ぶ。「だけ」「必要がない」「必ず」「個人判断」という限定は、どの範囲を変えているかを確認する。以下は原典を確認した解説初稿と、未処理問題の作業一覧である。初稿はAstra査読前の段階である。
\section{秘密情報：接近・持出し・気付きの機会を設計する}
\label{theory:research:secret_management}
情報の価値を保つには、必要な人が利用できることと、権限のない利用を防ぐことを両立させる。秘密を扱う人数をただ最小にすればよいわけではない。相互の確認ができない作業は、不正だけでなく誤送信や紛失も見つかりにくくなる。

まず情報を分類し、業務上の必要性に応じてアクセス権を決める。次に媒体の表示と取扱ルールで秘密の認識を支え、コピー・印刷・外部送信などの持出し経路を制御する。さらに作業の確認や配置を工夫し、異常へ気付く機会を確保する。これらは互いに補い合う対策である。

経産省のハンドブックは、可能な限り複数人での作業体制を整え、単独が必要なら責任者が事前の必要性と事後の作業内容を確認する案を示す\parencite{meti_secret_handbook2024}。複数人による確認と、誰でも秘密へアクセスさせることは違う。配置の見通しをよくする際も、無権限者へ秘密の画面を見せることを避ける。

暗号化とID・端末等による閲覧制限を適切に組み合わせれば、持ち出されたデータの閲覧を制限できる。ただし暗号だけで誤送信、権限者の不正、画面撮影まで全て消えるわけではない。ペーパーレス化も、アクセス管理と組み合わせて紙の放置等の機会を減らす対策であり、電子化だけの完全防御ではない。法的に営業秘密として保護される条件は第\ref{theory:compliance:secret}節で別に扱う。

\paragraph{資料の版}
この背景資料は2024年の第6版である。令和元年・令和5年の過去問の文言はそれぞれの原典を保持し、現在版の対策例を出題当時の義務へ置換しない。対策を一律の法律上の強制要件ではなく、情報と業務の条件に応じた低減策として読む。
\section{情報を扱う責任：公正な条件と説明の相手}
\label{theory:information:fairness}
情報通信技術は、場所や時間を越えて情報へ接近する能力を増す。試験中の検索や外部への質問が可能でも、その使用が試験制度の目的に適うとは限らない。試験は一定の条件の下で能力を比較するための測定である。ある受験者だけが禁止された助力を受ければ、点数が示す意味が変わり、他の受験者との比較を損なう。

発覚による自分の不利益を引き受けるという判断だけでは、他人の選考機会への影響を取り除けない。技術倫理では、行為者に生じる便益と危険だけでなく、情報にアクセスできない人、評価される人、結果に基づいて判断する組織へ及ぶ影響を考える。刑法の保護する利益も身体や物だけではない。平成23年度問3の注記は当時の刑法233条を挙げるが、逮捕という記載から個別事件の有罪を推論しない。

\section{暗黙知と形式知：技能の共有と知識の再利用}
\label{theory:information:knowledge}
装置の異音を聞いて故障の前兆へ気付く技能は、手順を読んだだけではすぐに身に付かない。一方、測定の設定・判定の条件を文章や図で残せば、他の担当者が比較や改善を行いやすい。前者の経験・技能として保持される知を暗黙知、後者の明示的な表現を持つ知を形式知という。二つは相互に補完する\parencite{jaist_knowledge}。

熟練者と共に作業して勘所を学ぶこと、経験を条件付きの手順へ表すこと、複数の手順や記録を組み合わせること、手順を実際に試して技能へ取り込むことをそれぞれ共同化・表出化・連結化・内面化と整理できる。ただし技能を完全に文章化できると仮定するモデルではない。手順書に温度や振動の条件が書かれていても、異常時の判断には実地の訓練を要することがある。共有可能性と、知識が成立した文脈の保持を両立させる。

\section{機密性・完全性・可用性：何を失う危険か}
\label{theory:information:cia}
情報セキュリティは秘密を守ることだけではない。患者の記録が漏れなくても数値が改ざんされれば診療判断を誤る。正しい記録が保管されていても必要時に閲覧できなければ業務が止まる。機密性は権限のない開示を防ぐこと、完全性は情報や処理の正確さと完全な状態を保つこと、可用性は認可された人が必要時に利用できることを扱う。

\begin{center}
\begin{tabular}{@{}p{0.18\linewidth}p{0.35\linewidth}p{0.40\linewidth}@{}}
\toprule
要素 & 失敗の例 & 対策の目的\\
\midrule
機密性 & 顧客へ他の顧客の連絡先を送る & 役割に応じた権限と送信先確認\\
完全性 & 測定値や履歴が変更される & 正確な更新、変更記録、整合検査\\
可用性 & 必要時に記録へ接近できない & 運用継続、復旧、適切な冗長化\\
\bottomrule
\end{tabular}
\end{center}
即時性や保存性という日常語を、この三要素の標準的名称と取り違えない。三要素は同じ対策で同時に改善する場合も、緊急時のアクセス制限と利用可能性のように調整を要する場合もある。資産と業務の条件に応じて要求水準を定める。

\section{方針と管理体系：文書から継続する運用へ}
\label{theory:information:management}
情報セキュリティポリシーは組織の方針・行動指針を示す。情報セキュリティマネジメントシステム（ISMS）は、対象資産とリスクを把握し、方針を計画と実施へ結び付け、結果を評価して改善する体系である。方針を作っただけで、その実施が保証されるわけではない。幹部を含む教育、担当と権限の設定、監視、事故時の初動と復旧、原因分析を継続して行う。

平成23年度問10の原典は2005年のISO27000と記すが、公開書誌のISO/IEC27000初版は2009年である\parencite{iso27000_2009}。原典の名称を黙って別規格へ書き換えず、空欄のCIA定義の判断と書誌上の疑義を分ける。規格全条項はこの初稿で照合しておらず、一般的な改善策を各版の逐語的要求事項とは断定しない。

\section{個人情報の範囲：公開・評価・映像でも識別性を見る}
\label{theory:privacy:scope}
個人情報保護は、情報の利用による便益と、本人の権利利益への影響を調整する。氏名等の事実だけでなく、個人についての評価や、識別可能な音声・映像も対象となり得る。死者だけに関する情報は生存する個人の定義から外れるが、それが生存する家族の情報でもある場合は、その家族についての識別性を調べる。新聞掲載や公開サイトに載ったというだけでも個人情報性は消えない\parencite{privacy_general}。

組織内の従業員番号等は、番号から本人へ容易に照合できるかを確認する。意味のない乱数でも照合表を管理していれば、単なる数字だから除外するという判断はできない。一方、全ての識別不能な番号へ同じ結論を無条件に当てはめない。個人情報、検索可能な集合としてのデータベース、その構成要素である個人データという段階も区別する。

\subsection{旧小規模事業者の基準と改正}
\label{theory:privacy:history}
平成24年度問4の5000人基準は当時の事業者適用範囲についての問題である。個々のデータベースを5000人以下へ分割すれば事業者の規律を免れる、という判断は成立しない。2017年5月30日施行の改正で5000人以下を適用除外する規定は廃止された\parencite{privacy_small2017}。現在の制度の説明へ旧閾値を残さない。旧政令の計数の全条件と当時版の直接照合は査読へ残す。

\section{本人同意：利用目的と第三者提供を分ける}
\label{theory:privacy:consent}
情報を取得する際の目的、内部で使う目的、他の主体へ提供することは、同じ行為ではない。ポイント管理のために得た顧客電話番号を求人へ転用する場合には、取得目的に必要な範囲かを検討する。名簿を内部で作成できることから、その名簿を他の家庭へ配布できるとは推論しない。

本人同意は、示された取扱方法について承諾する意思表示である。口頭や書面等の明示だけに限られず、黙示の同意が成立するかは個別具体的な事実で判断する\parencite{privacy_consent2017,privacy_consent_current}。したがって、問い合わせを受けたこと、説明の便宜があること、本人に役立つと事業者が考えることだけで、第三者へ連絡先を渡す承諾を補わない。

令和3年度問14では、カード請求の調査依頼から加盟店への連絡先提供まで黙示同意があると初回に解したが、公式正答は全記述不適切である。提供方法への承諾が確認できる契約・発言等が設問にない点を記録し、初回解答を保持して査読対象とする。黙示同意が法的に一切存在しない、と単純化して公式へ合わせない。CC一斉送信も、宛先同士へアドレスが見えるという実際の情報流れから判断する。

\subsection{配送：宛名情報と封入物の中身}
\label{theory:privacy:delivery}
発送のために氏名・住所を業者へ伝える場合、その宛名情報の取扱いは委託となり得る。他方、封入された媒体の中身を通常の配送業者は取り扱わない。中身の取扱合意等がなければ、その中身の個人データの取扱いを委託したとは一般に考えない\parencite{privacy_delivery}。この区別は、発送者の安全管理義務や配送契約上の責任を全て消すものではない。安全な配送方法、適切な業者、暗号化等を情報の条件に応じて検討する。

\subsection{災害時：必要性と提供範囲を限定する}
\label{theory:privacy:evacuation}
2013年の災害対策基本法改正は避難行動要支援者名簿の作成を市町村へ義務付けた。平時の支援者への名簿提供と、発災又は発生のおそれのある場合の提供は条件が違う\parencite{privacy_evacuation2013,privacy_evacuation_scope}。災害時には生命身体の保護へ特に必要な場合、必要な限度で同意なし提供が認められる。個人情報だから一切提供できない、又は災害対策という目的なら全情報を自由に共有できる、のどちらにも一般化しない。

\section{個人識別符号：生体特徴の列挙と変換条件}
\label{theory:privacy:biometrics}
政令はDNA、顔、虹彩、声、歩行、静脈、指紋・掌紋の七つの生体特徴を列挙する。電子計算機向けに変換され、特定個人の識別に足りる基準に合う符号という条件がある\parencite{privacy_general}。生体の特徴全てがそのまま個人識別符号になるわけではない。令和元年度問4は七分類が全て含まれること、令和6年度問9は骨振動の特性という非列挙項目を見分けることを問う。認証技術として実現可能かという一般論を、法令が何を列挙しているかという設問へ置換しない。

\section{説明責任と守秘：何を誰へ伝えるか}
\label{theory:information:accountability}
説明責任では、判断に影響した条件、根拠、限界、利益相反を示す。研究の実施場所や領域が違っても、その価値は共有される\parencite{singapore2010}。相手が質問し、理由を吟味する双方向の対話を支えることと、相手へ無条件の納得を要求することは違う。守秘を要する場合には、情報を抽象化する、説明範囲を確認する、権限のある監督者へ相談するなどにより、説明と秘密保持の双方を満たす方法を考える。

\section{情報のCIAと不審な依頼への初動}
\label{theory:research:security}
情報の機密性は許可された者に利用を限ること、完全性は情報の正確さ・完全さ、
可用性は必要なときに使えることである。
例えば設計値が書き換えられれば完全性、正しい値を必要時に取り出せなければ可用性の問題となる。
秘密を守るアクセス制限だけでは、他の二つを保証しない。\cite{saf53_ipa_cia}

方針を定め、対策を実施し、監査や監視で有効性を評価して見直す。
幹部を含む教育、事故時の初動、分析、復旧と再発防止を組織の役割へつなぐ。
見知った送信者名でも、時期や依頼内容に異常があれば、メール内の窓口に確認を求めるだけでは独立確認にならない。
既知の担当部署へ速やかに報告し、担当者が真偽と他の受信者への影響を調べる。
2017年のIPA資料は、開封率だけでなく報告から対応する組織の体制も訓練対象とする。\cite{saf53_ipa_mail}
一つの違和感は攻撃の確定証拠ではないが、被害が出るまで待つ理由にもならない。

\section{期待と異なるデータを検証する}
\label{theory:research:validation}
測定結果が仮説や依頼者の期待に反したときには、期待からの距離をデータの品質の尺度にしない。装置、校正、測定範囲、試料、時刻、記録の処理を点検し、除外には科学的な理由を記録する。予測値と実測値は区別する。測定できなかった部分を予測で補う場合にも、観測として扱わず方法と限界を示す\parencite{scj_code2013,mext_misconduct2014}。

報告期限と装置故障が重なる場合は、修理や代替測定の見通しを調べ、指導者・依頼者と期限、暫定報告、再測定を協議する。すべての未確認値を無説明で提出することも、都合の悪い値を無根拠に捨てることも妥当でない。生データと判断の根拠を残すことで、結論を後から検証できる。

\paragraph{説明と記録を分けて管理する}
依頼者が詳しくないことは、重要な仕様変更や危険を説明しなくてよい根拠にはならない。リスクを相手の判断に使える形で示す。技術情報を他者へ伝える場合は公開性、契約、守秘範囲を確かめる。すべての技術的会話が秘密漏洩ではない一方、業務上の非公開情報を友人へ話してよいともしない。

測定記録の保存・廃棄には責任、期間、アクセスの管理が必要である。記憶に基づく説明がすべて不正というわけではないが、記録があるのに確認せず数値報告を作れば正確性と検証可能性を失う。パスワードや施錠は対策の一部であり、更新・権限・保存期間も点検する。

\section{研究不正への個人責任と機関責任}
\label{theory:research:guideline2014}
平成26年の研究不正ガイドラインは、個人の不正の認定と措置に加え、研究機関の教育・予防・調査の責任を扱う。機関への管理条件、間接経費の削減、配分停止等の段階があり、「行政措置の根拠にはならない」とする説明は適切でない\parencite{mext_misconduct2014}。

教育対象を研究者本人だけに狭めない。学生や研究支援者など、実際に研究へ関わる人々に対する説明と組織体制を考える。研究活動の不正、研究費の不正使用、日常の不適切行為には異なる規程や手続がある。出題年のガイドラインの範囲・条件を確認し、現行の改定や別の資金制度の例外を一括して遡及しない。

\section{公益・説明責任}

\subsection{令和3年度 問3}
\label{q:r03-reg-apt-03}

\noindent\textbf{原典：} sources/R03/aptitude.pdf、PDF物理3頁。

\textbf{公式正答：} 1（single）。

\par\textbf{検討未確定：} アの受入努力は無条件の同意ではない。雇用者へCOI開示する技術者規範の当時版との最終照合は査読対象に残す。

\paragraph{問題の要旨} 説明責任・利益相反・研究公正・守秘義務の両立を判定する。

\Needspace{4\baselineskip}\paragraph{確認した条件・記述（独自要約）}\mbox{}\par

\noindent\textbf{原典確認：} 問題幹と各記述の主体・否定・全称・表の5候補を保存原典画像で確認。\par

\noindent\textbf{制度時点：} R03年度の原典表現を保持。現行資料の条番号や仕様を過去へ補わない。\par

\noindent\textbf{各記述の判断：} ア：○：受け手も理解のため努力する双方向の対話である。；イ：○：利益相反の可能性を雇用者等へ開示する。；ウ：○：研究実施の説明責任は共通価値である。；エ：○：守秘と説明の双方を満たす方法を考える。\par

\noindent\textbf{前提：} 説明責任は理由・根拠を相手が理解できる形で示すことである。研究領域を越える説明責任の価値と、利益相反情報の開示を確認する。守秘情報を無制限に公開する義務ではない。\par

\paragraph{問われる核心} 説明責任・利益相反・研究公正・守秘義務の両立を判定する。

\paragraph{必要知識} 説明責任は理由・根拠を相手が理解できる形で示すことである。研究領域を越える説明責任の価値と、利益相反情報の開示を確認する。守秘情報を無制限に公開する義務ではない。

\paragraph{考え方・正答を選ぶ理由（査読待ち初稿）} 四記述とも○、1である。アは受け手へ理由を吟味する姿勢を求める趣旨で、説明者の責務を受け手へ転嫁したり結論への無条件の賛成を要求したりする意味ではない。機密部分を抽象化し、権限のある相手や監督者へ説明することで守秘との両立を図る。

\Needspace{4\baselineskip}\paragraph{選択肢の検討（査読待ち初稿）}\mbox{}\par

\noindent\textbf{1：適切} ○・○・○・○。各記述の判断と全て一致する。\par

\noindent\textbf{2：不適切} ×・○・○・○。記述アの真偽が逆である。\par

\noindent\textbf{3：不適切} ○・×・○・○。記述イの真偽が逆である。\par

\noindent\textbf{4：不適切} ○・○・×・○。記述ウの真偽が逆である。\par

\noindent\textbf{5：不適切} ○・○・○・×。記述エの真偽が逆である。\par

\paragraph{再利用する考え方} 相手・権限・範囲を指定して説明し、情報公開と説明責任を同一視しない。

\paragraph{共通理論} \hyperref[theory:research:coi]{第\ref*{theory:research:coi}節}, \hyperref[theory:information:accountability]{第\ref*{theory:information:accountability}節}

\par 参考：\cite{singapore2010,scj_code2013}。

\par 難易度 3（校準初稿）、優先度未評価。

\section{研究・情報倫理}

\subsection{平成23年度 問2}
\label{q:h23-reg-apt-02}

\noindent\textbf{原典：} sources/H23/aptitude.pdf、PDF物理3頁。

\textbf{公式正答：} 3（single）。

\paragraph{問題の要旨} 2006年の科学者行動規範から、社会への説明を理由としたデータ改変の不適切さを判断する。

\Needspace{4\baselineskip}\paragraph{確認した条件・記述（独自要約）}\mbox{}\par

\noindent\textbf{原典確認：} 問題幹と各記述の主体・否定・全称・表の5候補を保存原典画像で確認。\par

\noindent\textbf{制度時点：} H23年度の原典表現を保持。現行資料の条番号や仕様を過去へ補わない。\par

\noindent\textbf{各記述の判断：} 規範：2006年声明の2・4・5・9・11項に照合。\par

\noindent\textbf{前提：} 科学者の自律性は社会の信頼に支えられ、正直な判断と厳格な記録を必要とする。分かりやすい表現と結果自体の改変は違う。\par

\paragraph{問われる核心} 2006年の科学者行動規範から、社会への説明を理由としたデータ改変の不適切さを判断する。

\paragraph{必要知識} 科学者の自律性は社会の信頼に支えられ、正直な判断と厳格な記録を必要とする。分かりやすい表現と結果自体の改変は違う。

\paragraph{考え方・正答を選ぶ理由（査読待ち初稿）} 3の前半の厳格な記録は適切だが、説明のためにデータを修正してよいという後半が改ざん禁止に反する。他の四候補は正直・中立・相互批判・利益相反への対応であり、2006年規範に対応する。

\Needspace{4\baselineskip}\paragraph{選択肢の検討（査読待ち初稿）}\mbox{}\par

\noindent\textbf{1：適切} 正直な行動と科学者相互の評価。規範2の正直・誠実と相互評価に対応する。\par

\noindent\textbf{2：適切} 影響を中立・客観的に公開する。規範4の説明と公開に対応する。\par

\noindent\textbf{3：不適切} 理解を得るため多少のデータ修正を認める。規範5の改ざん禁止と矛盾する。説明用の図の整理や誤記訂正を、結論を変える改ざんへ拡張しない。\par

\noindent\textbf{4：適切} 批判に耳を傾け他者の成果を尊重する。規範9の他者との関係に対応する。\par

\noindent\textbf{5：適切} 利益の衝突に注意して公共性へ配慮する。規範11の利益相反への対応に合う。\par

\paragraph{再利用する考え方} 説明を受けやすくするための表現改善と、根拠を不真正にする操作を分ける。

\paragraph{共通理論} \hyperref[theory:research:norms]{第\ref*{theory:research:norms}節}, \hyperref[theory:research:ffp]{第\ref*{theory:research:ffp}節}

\par 参考：\cite{scj_code2006}。

\par 難易度 3（校準初稿）、優先度未評価。

\subsection{平成23年度 問3}
\label{q:h23-reg-apt-03}

\noindent\textbf{原典：} sources/H23/aptitude.pdf、PDF物理3頁。

\textbf{公式正答：} 5（single）。

\par\textbf{検討未確定：} 原典注記の刑法233条は2011年の懲役表現を保持する。具体的な事件の判決と現行刑罰条文全体は未照合であり、逮捕から有罪を補わない。

\paragraph{問題の要旨} 通信技術を用いた入試不正を、公正な比較条件と倫理・法の関係から評価する。

\Needspace{4\baselineskip}\paragraph{確認した条件・記述（独自要約）}\mbox{}\par

\noindent\textbf{原典確認：} 問題幹と各記述の主体・否定・全称・表の5候補を保存原典画像で確認。\par

\noindent\textbf{制度時点：} H23年度の原典表現を保持。現行資料の条番号や仕様を過去へ補わない。\par

\noindent\textbf{各記述の判断：} ア：○：共通条件の公正を維持する必要がある。；イ：×：業務や信用に対する犯罪もあり得る。；ウ：×：本人がリスクを負うだけでは他者への不公正を解消しない。；エ：○：新技術への法整備と倫理観の育成が必要である。\par

\noindent\textbf{前提：} 試験の成績は共通条件で比較する指標である。条件を破った得点は比較可能性を壊す。法が保護する利益は身体や物に限られない。\par

\paragraph{問われる核心} 通信技術を用いた入試不正を、公正な比較条件と倫理・法の関係から評価する。

\paragraph{必要知識} 試験の成績は共通条件で比較する指標である。条件を破った得点は比較可能性を壊す。法が保護する利益は身体や物に限られない。

\paragraph{考え方・正答を選ぶ理由（査読待ち初稿）} ア・エが○、イ・ウが×である。身体侵害や窃盗がないというだけでは犯罪成立の可能性は除外されず、自分が発覚リスクを引き受けても他の受験者の公正な選考は損なわれる。○××○、5である。逮捕を有罪判決と同一視しない。

\Needspace{4\baselineskip}\paragraph{選択肢の検討（査読待ち初稿）}\mbox{}\par

\noindent\textbf{1：不適切} ×・○・○・×。記述ア・イ・ウ・エの真偽が逆である。\par

\noindent\textbf{2：不適切} ○・○・×・○。記述イの真偽が逆である。\par

\noindent\textbf{3：不適切} ○・○・○・○。記述イ・ウの真偽が逆である。\par

\noindent\textbf{4：不適切} ○・×・○・×。記述ウ・エの真偽が逆である。\par

\noindent\textbf{5：適切} ○・×・×・○。各記述の判断と全て一致する。\par

\paragraph{再利用する考え方} 新技術が行為を可能にしたことから、その行為の許容性を推論しない。

\paragraph{共通理論} \hyperref[theory:information:fairness]{第\ref*{theory:information:fairness}節}

\par 参考：\cite{scj_code2006}。

\par 難易度 3（校準初稿）、優先度未評価。

\subsection{平成26年度 問2}
\label{q:h26-reg-apt-02}

\noindent\textbf{原典：} sources/H26/aptitude.pdf、PDF物理2頁。

\textbf{公式正答：} 5（single）。

\paragraph{問題の要旨} 2013年改訂の科学者の行動規範から、最も不適切な記述を選ぶ。

\Needspace{4\baselineskip}\paragraph{確認した条件・記述（独自要約）}\mbox{}\par

\noindent\textbf{時点：} 2006年制定、2013年改訂の日本学術会議声明。\par

\noindent\textbf{方向：} 不適切な記述を1つ選ぶ。\par

\paragraph{問われる核心} 科学の自律性と社会への責任

\paragraph{必要知識} 科学者の範囲、研究の功績と責任、説明責任、社会との対話。

\paragraph{考え方・正答を選ぶ理由（査読待ち初稿）} 各肢の主語と責任を確認する。1〜4は科学者の範囲、品質責任、説明、研究公正という規範の要素に対応する。5は科学の自律性を社会からの孤立と読み替え、対話と助言を否定しているため選ぶ。

\Needspace{4\baselineskip}\paragraph{選択肢の検討（査読待ち初稿）}\mbox{}\par

\noindent\textbf{1：問の条件に一致しない候補} 科学者の範囲は所属に依存せず人文・社会・自然科学の研究者と知識を活用する専門職を含む。前文での科学者の定義に合う。自然科学者や大学所属者に限定しない。\par

\noindent\textbf{2：問の条件に一致しない候補} 知識・技術の質を保証し、健康・福祉・安全・地球環境へ貢献する責任を持つ。基本的責任の規定に合う。専門知識の質と社会への貢献の双方を求める。\par

\noindent\textbf{3：問の条件に一致しない候補} 研究の意義と社会・環境への影響を中立・客観的に説明・公表する。説明と公開の責務に合う。専門性を理由に説明責任を免れるものではない。\par

\noindent\textbf{4：問の条件に一致しない候補} 役割に応じた功績と責任を負い、データを保存してFFPを行わず加担しない。研究活動の公正さに関する責務に合う。功績のみを受け取る立場ではない。\par

\noindent\textbf{5：問の条件に一致する候補} 独立性のため市民対話に否定的であり、政策への科学的助言も避ける。独立した科学的判断は対話や助言を避ける理由にならない。規範は市民との交流と政策への助言を促すので不適切である。\par

\paragraph{再利用する考え方} 制度の対象・定義と倫理上の責任を区別し、科学の自律性と社会への責任を具体的行為に適用する。

\paragraph{共通理論} \hyperref[theory:research:norms]{第\ref*{theory:research:norms}節}

\par 参考：\cite{scj_code2013}。

\par 難易度 2（校準初稿）、優先度未評価。

\subsection{平成26年度 問3}
\label{q:h26-reg-apt-03}

\noindent\textbf{原典：} sources/H26/aptitude.pdf、PDF物理3頁。

\textbf{公式正答：} 1（single）。

\paragraph{問題の要旨} 捏造・改ざん・盗用という3種の研究不正の定義を判定する。

\Needspace{4\baselineskip}\paragraph{確認した条件・記述（独自要約）}\mbox{}\par

\noindent\textbf{判断順：} 記述アから順に○=適切、×=不適切。\par

\noindent\textbf{ア：} 捏造は存在しないデータや研究結果を作ること。\par

\noindent\textbf{イ：} 改ざんは資料・機器・過程などを操作し、データや結果を真正でないものに加工すること。\par

\noindent\textbf{ウ：} 盗用は他者のアイデア・分析法・データ・結果・論文・用語を、了解若しくは適切な表示なしに流用すること。\par

\paragraph{問われる核心} 研究不正FFPの定義

\paragraph{必要知識} 捏造・改ざん・盗用の対象、データ処理と不真正化の区別。

\paragraph{考え方・正答を選ぶ理由（査読待ち初稿）} アは架空のデータの作成、イは取得・処理過程を操作した不真正化、ウは他者の成果の不適切な流用という各定義に合う。3記述とも○なので第1候補である。引用を示す責任と、著作物利用の許諾の問題は別に検討し、単に許諾があれば他者の成果を自分のものと表示してよいとは考えない。

\Needspace{4\baselineskip}\paragraph{選択肢の検討（査読待ち初稿）}\mbox{}\par

\noindent\textbf{1：問の条件に一致する候補} ○○○。全記述の判定○○○に一致する。\par

\noindent\textbf{2：問の条件に一致しない候補} ×○○。記述アの判定が逆である。独立判定は○○○となる。\par

\noindent\textbf{3：問の条件に一致しない候補} ○×○。記述イの判定が逆である。独立判定は○○○となる。\par

\noindent\textbf{4：問の条件に一致しない候補} ××○。記述ア・イの判定が逆である。独立判定は○○○となる。\par

\noindent\textbf{5：問の条件に一致しない候補} ×××。記述ア・イ・ウの判定が逆である。独立判定は○○○となる。\par

\paragraph{再利用する考え方} 制度の対象・定義と倫理上の責任を区別し、研究不正FFPの定義を具体的行為に適用する。

\paragraph{共通理論} \hyperref[theory:research:ffp]{第\ref*{theory:research:ffp}節}

\par 参考：\cite{mext_misconduct2014}。

\par 難易度 2（校準初稿）、優先度未評価。

\subsection{平成27年度 問8}
\label{q:h27-reg-apt-08}

\noindent\textbf{原典：} sources/H27/aptitude.pdf、PDF物理8頁。

\textbf{公式正答：} 1（single）。

\paragraph{問題の要旨} 実験ノートの記録・保存・訂正・異動時の持出しを判定する。

\Needspace{4\baselineskip}\paragraph{確認した条件・記述（独自要約）}\mbox{}\par

\noindent\textbf{判断順：} 記述アから順に○=適切、×=不適切。\par

\noindent\textbf{ア：} 原則として綴じたノートを消せない筆記具で記載する。\par

\noindent\textbf{イ：} 研究ノートは特許紛争回避や権利関係の明確化にも役立ち、企業研究者にも有用である。\par

\noindent\textbf{ウ：} 研究機関は一定期間の研究データ保存と必要時の開示を義務付ける規程を整える。\par

\noindent\textbf{エ：} 記述の誤りがあるページは廃棄し、見やすく修正するのが望ましい。\par

\noindent\textbf{オ：} 研究者個人に帰属するので異動時には個人判断で持ち出せる。\par

\paragraph{問われる核心} 研究記録の追跡可能性と機関管理

\paragraph{必要知識} 元記録の保存、訂正の履歴、データ帰属と異動時の管理。

\paragraph{考え方・正答を選ぶ理由（査読待ち初稿）} アは紙ノートの追記・差替えを追跡できるようにする原則である。イは研究の経緯を立証する記録として有用という記述であり、ノートだけで特許権が自動的に決まるとの意味ではない。ウは機関の保存規程に対応する。エのページ廃棄は元の経過を失わせる。オの個人所有・自由持出しの一律断定も不適切であり、資金・契約・機関規程による管理が必要である。○○○××の第1候補を選ぶ。

\Needspace{4\baselineskip}\paragraph{選択肢の検討（査読待ち初稿）}\mbox{}\par

\noindent\textbf{1：問の条件に一致する候補} ○○○××。全記述の判定○○○××に一致する。\par

\noindent\textbf{2：問の条件に一致しない候補} ×○○××。記述アの判定が逆である。独立判定は○○○××となる。\par

\noindent\textbf{3：問の条件に一致しない候補} ○×○××。記述イの判定が逆である。独立判定は○○○××となる。\par

\noindent\textbf{4：問の条件に一致しない候補} ×○×○×。記述ア・ウ・エの判定が逆である。独立判定は○○○××となる。\par

\noindent\textbf{5：問の条件に一致しない候補} ○○○×○。記述オの判定が逆である。独立判定は○○○××となる。\par

\paragraph{再利用する考え方} 制度の対象・定義と倫理上の責任を区別し、研究記録の追跡可能性と機関管理を具体的行為に適用する。

\paragraph{共通理論} \hyperref[theory:research:records]{第\ref*{theory:research:records}節}

\par 参考：\cite{mext_misconduct2014,jsps_integrity2025}。

\par 難易度 2（校準初稿）、優先度未評価。

\subsection{平成28年度 問4}
\label{q:h28-reg-apt-04}

\noindent\textbf{原典：} sources/H28/aptitude.pdf、PDF物理4頁。

\textbf{公式正答：} 2（single）。

\paragraph{問題の要旨} 2014年ガイドラインの特定不正行為と研究倫理全般を区別して判定する。

\Needspace{4\baselineskip}\paragraph{確認した条件・記述（独自要約）}\mbox{}\par

\noindent\textbf{判断順：} 記述アから順に○=適切、×=不適切。\par

\noindent\textbf{ア：} 故意又は基本的注意義務を著しく怠った、発表成果中の捏造・改ざん・盗用が特定不正行為である。\par

\noindent\textbf{イ：} 確認された不正に係る競争的資金について配分機関は事案に応じ交付取消しや一部・全部の返還を求められる。\par

\noindent\textbf{ウ：} 二重投稿や不適切な著者表示は研究者倫理に反するが不正行為ではない。\par

\noindent\textbf{エ：} 不正の対応は研究者・コミュニティ・機関の規律と自浄作用に基づく。\par

\noindent\textbf{オ：} 研究倫理教育は倫理を向上させるが、不正予防の環境整備にはならない。\par

\paragraph{問われる核心} 特定不正行為と広い研究倫理違反

\paragraph{必要知識} FFPの限定、一般的な不正、競争的資金の措置、研究倫理教育。

\paragraph{考え方・正答を選ぶ理由（査読待ち初稿）} アは特定不正行為の要件である。イは対象の競争的資金と事案に応じた措置に限る記述で妥当であり、全資金の自動返還ではない。ウは特定不正行為の限定を不正一般へ広げている。二重投稿等も倫理に反する不正として認識されるので×である。エは自己規律・自律の位置付けに合い、オは教育の予防的役割を否定するため×である。○○×○×の第2候補となる。

\Needspace{4\baselineskip}\paragraph{選択肢の検討（査読待ち初稿）}\mbox{}\par

\noindent\textbf{1：問の条件に一致しない候補} ○○○○×。記述ウの判定が逆である。独立判定は○○×○×となる。\par

\noindent\textbf{2：問の条件に一致する候補} ○○×○×。全記述の判定○○×○×に一致する。\par

\noindent\textbf{3：問の条件に一致しない候補} ○×○×○。記述イ・ウ・エ・オの判定が逆である。独立判定は○○×○×となる。\par

\noindent\textbf{4：問の条件に一致しない候補} ×○×○×。記述アの判定が逆である。独立判定は○○×○×となる。\par

\noindent\textbf{5：問の条件に一致しない候補} ○×○××。記述イ・ウ・エの判定が逆である。独立判定は○○×○×となる。\par

\paragraph{再利用する考え方} 制度の対象・定義と倫理上の責任を区別し、特定不正行為と広い研究倫理違反を具体的行為に適用する。

\paragraph{共通理論} \hyperref[theory:research:ffp]{第\ref*{theory:research:ffp}節}, \hyperref[theory:research:prevention]{第\ref*{theory:research:prevention}節}

\par 参考：\cite{mext_misconduct2014}。

\par 難易度 2（校準初稿）、優先度未評価。

\subsection{平成29年度 問14}
\label{q:h29-reg-apt-14}

\noindent\textbf{原典：} sources/H29/aptitude.pdf、PDF物理13頁。

\textbf{公式正答：} 3（single）。

\paragraph{問題の要旨} 研究の自由、記録、資金、著者責任、引用に関する7記述中の不適切数を求める。

\Needspace{4\baselineskip}\paragraph{確認した条件・記述（独自要約）}\mbox{}\par

\noindent\textbf{方向：} 不適切な記述の個数。候補1〜5はそれぞれ0〜4個。\par

\noindent\textbf{ア：} 研究の自由は真理又は公益、独創性と正確さ、結果への責任を伴う。\par

\noindent\textbf{イ：} 客観的検証のため記録を保存し、FFPを禁じる。\par

\noindent\textbf{ウ：} 研究費は計画・申請に基づき公正に使い、目的外利用・不正操作をしない。\par

\noindent\textbf{エ：} 貢献者を著者とし論文の正確さなどへ責任を負うが、共著者は自分に関係した部分だけへ責任を持てばよい。\par

\noindent\textbf{オ：} 実験の大きな貢献があれば、論文の内容と正確さを説明できなくても共著者になれる。\par

\noindent\textbf{カ：} 先行研究を明示し、自らの独創的寄与を区別する。\par

\noindent\textbf{キ：} 著作物の引用部分と原著者を明示し、引用の範囲を超えるコピーには原則として許諾を得る。\par

\paragraph{問われる核心} 著者の貢献と説明責任

\paragraph{必要知識} 功績の帰属、共同研究での責任、引用と著作物利用の区別。

\paragraph{考え方・正答を選ぶ理由（査読待ち初稿）} ア・イ・ウは研究の目的、検証、資金責任、カ・キは先行成果と利用条件の明示で適切である。エは一律に自分の部分だけへ責任を限定しており不適切である。役割分担を明示する場合も、共同研究の信頼性の確認を放棄することにはならない。オは実験への貢献だけで内容を説明できない人を著者にできるとする点が不適切である。誤りは2件、第3候補となる。

\Needspace{4\baselineskip}\paragraph{選択肢の検討（査読待ち初稿）}\mbox{}\par

\noindent\textbf{1：問の条件に一致しない候補} 0個。エ・オの2つの誤りを数えていない。\par

\noindent\textbf{2：問の条件に一致しない候補} 1個。エの責任限定とオの説明不能の2つが不適切である。\par

\noindent\textbf{3：問の条件に一致する候補} 2個。エ・オの2記述を不適切と数えた結果に一致する。\par

\noindent\textbf{4：問の条件に一致しない候補} 3個。ア・イ・ウ・カ・キには、記載された倫理原則に反する内容はない。\par

\noindent\textbf{5：問の条件に一致しない候補} 4個。不適切はエ・オだけであり、4件ではない。\par

\paragraph{再利用する考え方} 制度の対象・定義と倫理上の責任を区別し、著者の貢献と説明責任を具体的行為に適用する。

\paragraph{共通理論} \hyperref[theory:research:authorship]{第\ref*{theory:research:authorship}節}

\par 参考：\cite{jsps_integrity2025}。

\par 難易度 2（校準初稿）、優先度未評価。

\subsection{令和元年度再試験 問10}
\label{q:r01-ret-apt-10}

\noindent\textbf{原典：} sources/R01\_retest/aptitude.pdf、PDF物理10頁。

\textbf{公式正答：} 2（single）。

\paragraph{問題の要旨} 誠実な研究の誤り、私費研究、言語の異なる重複発表、共著論文の翻訳利用を判定する。

\Needspace{4\baselineskip}\paragraph{確認した条件・記述（独自要約）}\mbox{}\par

\noindent\textbf{判断順：} 記述アから順に○=適切、×=不適切。\par

\noindent\textbf{ア：} 適切な方法で正当に得た研究成果が結果的に誤りでも、ガイドライン公表後は不正行為となった。\par

\noindent\textbf{イ：} 個人の資金と労力で研究・公表するならガイドラインの内容を考慮する必要がない。\par

\noindent\textbf{ウ：} 同一成果を日本語と英語で別の学会に発表すれば二重投稿にはならない。\par

\noindent\textbf{エ：} AとBの共著英語論文の一部を、Aが単独日本語書籍に翻訳利用する前にBへ相談して了解を得た。\par

\paragraph{問われる核心} 研究の誠実さと公表時の透明性

\paragraph{必要知識} 誠実な誤りの除外、資金源に依存しない倫理、重複発表と共同成果利用。

\paragraph{考え方・正答を選ぶ理由（査読待ち初稿）} アは誠実な誤りを不正とする点、イは資金源だけで倫理責任を免れる点、ウは言語だけで重複発表を正当化する点がいずれも誤りである。エは共同成果の再利用を独断で行わず共著者へ相談した行為として適切である。出版側の許諾・引用表示などが不要との記述ではなく、実利用時には別途確認する。×××○の第2候補である。

\Needspace{4\baselineskip}\paragraph{選択肢の検討（査読待ち初稿）}\mbox{}\par

\noindent\textbf{1：問の条件に一致しない候補} ×○×○。記述イの判定が逆である。独立判定は×××○となる。\par

\noindent\textbf{2：問の条件に一致する候補} ×××○。全記述の判定×××○に一致する。\par

\noindent\textbf{3：問の条件に一致しない候補} ○××○。記述アの判定が逆である。独立判定は×××○となる。\par

\noindent\textbf{4：問の条件に一致しない候補} ○○○×。記述ア・イ・ウ・エの判定が逆である。独立判定は×××○となる。\par

\noindent\textbf{5：問の条件に一致しない候補} ××○○。記述ウの判定が逆である。独立判定は×××○となる。\par

\paragraph{再利用する考え方} 制度の対象・定義と倫理上の責任を区別し、研究の誠実さと公表時の透明性を具体的行為に適用する。

\paragraph{共通理論} \hyperref[theory:research:ffp]{第\ref*{theory:research:ffp}節}, \hyperref[theory:research:authorship]{第\ref*{theory:research:authorship}節}

\par 参考：\cite{mext_misconduct2014,jsps_integrity2025}。

\par 難易度 2（校準初稿）、優先度未評価。

\subsection{令和元年度再試験 問6}
\label{q:r01-ret-apt-06}

\noindent\textbf{原典：} sources/R01\_retest/aptitude.pdf、PDF物理6頁。

\textbf{公式正答：} 4（single）。

\paragraph{問題の要旨} 科学者の範囲・誠実さ・両義的利用・研究環境の責任を判定する。

\Needspace{4\baselineskip}\paragraph{確認した条件・記述（独自要約）}\mbox{}\par

\noindent\textbf{判断順：} 記述アから順に○=適切、×=不適切。\par

\noindent\textbf{ア：} 科学者は所属にかかわらず全学術分野の研究者や科学知識の利活用を行う専門職を含む。\par

\noindent\textbf{イ：} 正直・誠実に行動し、専門能力の維持向上と知識の正確さの提示に努める。\par

\noindent\textbf{ウ：} 意図に反し悪用され得ても、社会へ貢献すると判断すれば速やかに研究・公表を積極的に行う。\par

\noindent\textbf{エ：} 公正な研究環境の維持と不正抑止の教育啓発に継続して取り組む責任がある。\par

\paragraph{問われる核心} 科学研究の利用の両義性

\paragraph{必要知識} 研究の便益と悪用リスク、社会に許容される研究・公表方法。

\paragraph{考え方・正答を選ぶ理由（査読待ち初稿）} アは前文の定義、イは誠実な姿勢、エは環境整備と教育の責任に対応する。ウは善用の可能性だけで悪用への配慮を省いている。規範は研究と公表の際に社会へ許容される手段・方法を選ぶことを求めるので×である。○○×○の第4候補である。

\Needspace{4\baselineskip}\paragraph{選択肢の検討（査読待ち初稿）}\mbox{}\par

\noindent\textbf{1：問の条件に一致しない候補} ○○○○。記述ウの判定が逆である。独立判定は○○×○となる。\par

\noindent\textbf{2：問の条件に一致しない候補} ×○○○。記述ア・ウの判定が逆である。独立判定は○○×○となる。\par

\noindent\textbf{3：問の条件に一致しない候補} ○×○○。記述イ・ウの判定が逆である。独立判定は○○×○となる。\par

\noindent\textbf{4：問の条件に一致する候補} ○○×○。全記述の判定○○×○に一致する。\par

\noindent\textbf{5：問の条件に一致しない候補} ○○○×。記述ウ・エの判定が逆である。独立判定は○○×○となる。\par

\paragraph{再利用する考え方} 制度の対象・定義と倫理上の責任を区別し、科学研究の利用の両義性を具体的行為に適用する。

\paragraph{共通理論} \hyperref[theory:research:norms]{第\ref*{theory:research:norms}節}, \hyperref[theory:research:prevention]{第\ref*{theory:research:prevention}節}

\par 参考：\cite{scj_code2013}。

\par 難易度 2（校準初稿）、優先度未評価。

\subsection{令和2年度 問3}
\label{q:r02-reg-apt-03}

\noindent\textbf{原典：} sources/R02/aptitude.pdf、PDF物理3頁。

\textbf{公式正答：} 3（single）。

\paragraph{問題の要旨} 兼業変更、競争相手の査読、配偶者名義の株式、公的設備の企業利用の利益相反管理を判定する。

\Needspace{4\baselineskip}\paragraph{確認した条件・記述（独自要約）}\mbox{}\par

\noindent\textbf{判断順：} 記述アから順に○=適切、×=不適切。\par

\noindent\textbf{ア：} 水曜日に許可された有償兼業について土曜の技術会議参加が必要となり、大学COI委員会へ相談した。\par

\noindent\textbf{イ：} 自己の研究と競争関係にある論文の公正な査読に不安があるため辞退した。\par

\noindent\textbf{ウ：} 公的資金で評価した企業の未公開株を少額で配偶者名義にて買い、COI委員会への相談を省いた。\par

\noindent\textbf{エ：} 大学の兼業許可を得て設立した企業が公的資金の研究設備を無償利用する必要が生じ、事前に大学COI委員会へ相談した。\par

\paragraph{問われる核心} 利益相反の開示と第三者管理

\paragraph{必要知識} 経済的利益と責務の衝突、家族の利害、査読の競争関係、機関資源。

\paragraph{考え方・正答を選ぶ理由（査読待ち初稿）} アは許可条件の変更を自己判断せず相談する対応、イは査読の公正さを確保する辞退で適切である。ウは家族名義や少額を理由に利害関係を隠しており不適切である。エは所属機関の資源と私企業の利益を分ける事前相談として適切である。相談すれば無償利用が必ず認められるとの意味ではない。○○×○の第3候補を選ぶ。

\Needspace{4\baselineskip}\paragraph{選択肢の検討（査読待ち初稿）}\mbox{}\par

\noindent\textbf{1：問の条件に一致しない候補} ○○○○。記述ウの判定が逆である。独立判定は○○×○となる。\par

\noindent\textbf{2：問の条件に一致しない候補} ○○○×。記述ウ・エの判定が逆である。独立判定は○○×○となる。\par

\noindent\textbf{3：問の条件に一致する候補} ○○×○。全記述の判定○○×○に一致する。\par

\noindent\textbf{4：問の条件に一致しない候補} ○×○○。記述イ・ウの判定が逆である。独立判定は○○×○となる。\par

\noindent\textbf{5：問の条件に一致しない候補} ×○○○。記述ア・ウの判定が逆である。独立判定は○○×○となる。\par

\paragraph{再利用する考え方} 制度の対象・定義と倫理上の責任を区別し、利益相反の開示と第三者管理を具体的行為に適用する。

\paragraph{共通理論} \hyperref[theory:research:coi]{第\ref*{theory:research:coi}節}

\par 参考：\cite{jsps_integrity2025}。

\par 難易度 2（校準初稿）、優先度未評価。

\subsection{令和5年度 問7}
\label{q:r05-reg-apt-07}

\noindent\textbf{原典：} sources/R05/aptitude.pdf、PDF物理7頁。

\textbf{公式正答：} 2（single）。

\paragraph{問題の要旨} 科学者の研究公正、対話、科学的助言と政策判断の関係を判定する。

\Needspace{4\baselineskip}\paragraph{確認した条件・記述（独自要約）}\mbox{}\par

\noindent\textbf{判断順：} 記述アから順に○=適切、×=不適切。\par

\noindent\textbf{ア：} 公表は役割に応じた功績と責任を伴い、データを保存してFFPを行わず加担しない。\par

\noindent\textbf{イ：} 市民対話と政策助言に努め、科学者に意見の相違がある場合は多数決で統一見解を決めて助言する。\par

\noindent\textbf{ウ：} 客観的根拠に基づき権威を濫用せず、科学的不確実性と見解の多様性も説明する。\par

\noindent\textbf{エ：} 科学的知見は政策の唯一の根拠ではなく、助言と異なる決定には必要に応じ社会への説明を求める。\par

\paragraph{問われる核心} 科学的助言の不確実性と政策責任

\paragraph{必要知識} 合意と意見相違の説明、科学的根拠と政策決定の他の判断材料。

\paragraph{考え方・正答を選ぶ理由（査読待ち初稿）} ア・ウ・エはそれぞれ研究公正、助言の質、科学と政策の役割に合う。イの前半は適切だが、多数決で相違を消すという後半が誤りである。規範は意見が異なる場合に相違を説明することを求めている。○×○○の第2候補である。

\Needspace{4\baselineskip}\paragraph{選択肢の検討（査読待ち初稿）}\mbox{}\par

\noindent\textbf{1：問の条件に一致しない候補} ×○○○。記述ア・イの判定が逆である。独立判定は○×○○となる。\par

\noindent\textbf{2：問の条件に一致する候補} ○×○○。全記述の判定○×○○に一致する。\par

\noindent\textbf{3：問の条件に一致しない候補} ○○×○。記述イ・ウの判定が逆である。独立判定は○×○○となる。\par

\noindent\textbf{4：問の条件に一致しない候補} ○○○×。記述イ・エの判定が逆である。独立判定は○×○○となる。\par

\noindent\textbf{5：問の条件に一致しない候補} ○○○○。記述イの判定が逆である。独立判定は○×○○となる。\par

\paragraph{再利用する考え方} 制度の対象・定義と倫理上の責任を区別し、科学的助言の不確実性と政策責任を具体的行為に適用する。

\paragraph{共通理論} \hyperref[theory:research:norms]{第\ref*{theory:research:norms}節}

\par 参考：\cite{scj_code2013}。

\par 難易度 2（校準初稿）、優先度未評価。

\subsection{令和6年度 問8}
\label{q:r06-reg-apt-08}

\noindent\textbf{原典：} sources/R06/aptitude.pdf、PDF物理8頁。

\textbf{公式正答：} 3（single）。

\paragraph{問題の要旨} 贈与・幽霊著者、資金獲得のみの著者資格、細切れ出版を判定する。

\Needspace{4\baselineskip}\paragraph{確認した条件・記述（独自要約）}\mbox{}\par

\noindent\textbf{判断順：} 記述アから順に○=適切、×=不適切。\par

\noindent\textbf{ア：} 資格を満たさない者を好意や強制により著者へ入れる行為はギフト・オーサーシップという倫理違反である。\par

\noindent\textbf{イ：} 研究費を獲得した者は研究遂行へ寄与したので論文著者の資格を満たす。\par

\noindent\textbf{ウ：} 1つの研究を複数の小研究に分け細切れに公表することをサラミ又はボローニャ出版という。\par

\noindent\textbf{エ：} 実質的貢献と著者要件を満たす者を著者として表示しない行為はゴースト・オーサーシップという倫理違反である。\par

\paragraph{問われる核心} 著者資格と研究成果の適切な単位

\paragraph{必要知識} 貢献と責任、謝辞、ギフト・ゴースト著者、細切れ出版。

\paragraph{考え方・正答を選ぶ理由（査読待ち初稿）} アは不相応な著者の追加、エは資格ある著者の省略という定義に合う。イは資金獲得のみで著者資格を充足するとする誤りである。ウは細切れ公表の用語を説明しており○である。独立した科学的問いを持つ複数論文すべてが倫理違反という意味に一般化しない。○×○○の第3候補を選ぶ。

\Needspace{4\baselineskip}\paragraph{選択肢の検討（査読待ち初稿）}\mbox{}\par

\noindent\textbf{1：問の条件に一致しない候補} ○○○○。記述イの判定が逆である。独立判定は○×○○となる。\par

\noindent\textbf{2：問の条件に一致しない候補} ×○○○。記述ア・イの判定が逆である。独立判定は○×○○となる。\par

\noindent\textbf{3：問の条件に一致する候補} ○×○○。全記述の判定○×○○に一致する。\par

\noindent\textbf{4：問の条件に一致しない候補} ○○×○。記述イ・ウの判定が逆である。独立判定は○×○○となる。\par

\noindent\textbf{5：問の条件に一致しない候補} ○○○×。記述イ・エの判定が逆である。独立判定は○×○○となる。\par

\paragraph{再利用する考え方} 制度の対象・定義と倫理上の責任を区別し、著者資格と研究成果の適切な単位を具体的行為に適用する。

\paragraph{共通理論} \hyperref[theory:research:authorship]{第\ref*{theory:research:authorship}節}

\par 参考：\cite{jsps_integrity2025}。

\par 難易度 2（校準初稿）、優先度未評価。

\section{研究データの検証と期限}

\subsection{平成24年度 問15}
\label{q:h24-reg-apt-15}

\noindent\textbf{原典：} sources/H24/aptitude.pdf、PDF物理13頁。

\textbf{公式正答：} 5（single）。

\paragraph{問題の要旨} 期待と異なる測定結果の検証と相談。

\Needspace{4\baselineskip}\paragraph{確認した条件・記述（独自要約）}\mbox{}\par

\noindent\textbf{状況：} 企業Cの期待と違う夜間自動測定、装置故障、再測定不能、報告期限。\par

\noindent\textbf{判断：} データの妥当性が未確定であり、期待との差だけでは除外理由にならない。\par

\paragraph{問われる核心} 期待と異なる測定結果の検証と相談。

\paragraph{必要知識} 生データの保存、測定の妥当性、予測と実測の区別、期限と説明責任。

\paragraph{考え方・正答を選ぶ理由（査読待ち初稿）} 翌朝の装置故障により直ちに再測定できない。新しい自動記録が予測と違うだけで捏造や無効と決めず、装置と測定条件の検証、修理・再測定の見通しを調べて教授へ相談する5。締切を理由に都合のよい値を作らない。

\Needspace{4\baselineskip}\paragraph{選択肢の検討（査読待ち初稿）}\mbox{}\par

\noindent\textbf{1：不適切} 未確認の新データを期限に合わせ提出。検証できていないデータの条件と不確かさを確認せず、期限だけで結論を急ぐ。\par

\noindent\textbf{2：不適切} 新データを除外して従来の予測どおり報告。期待と違うこと自体は除外の科学的理由にならない。\par

\noindent\textbf{3：不適切} 予測から作った値を実測として報告。予測値と観測値を偽って同一視する。\par

\noindent\textbf{4：不適切} 将来の契約を優先して新データを伏せる。事業上の都合を研究結果の取捨選択の根拠にしている。\par

\noindent\textbf{5：条件内で適切} 装置の修理見通しを調べ教授と協議。再測定の可能性、納期、データ品質を整理して判断し、結論の限界を伝える。\par

\paragraph{再利用する考え方} 生データの保存、測定の妥当性、予測と実測の区別、期限と説明責任。 設問の否定・全称・基準年を確かめ、根拠と評価範囲を分けて判断する。

\paragraph{共通理論} \hyperref[theory:research:validation]{第\ref*{theory:research:validation}節}

\par 参考：\cite{scj_code2013,mext_misconduct2014}。

\par 難易度 2（校準初稿）、優先度未評価。

\section{2014研究不正ガイドラインの措置}

\subsection{平成27年度 問7}
\label{q:h27-reg-apt-07}

\noindent\textbf{原典：} sources/H27/aptitude.pdf、PDF物理7頁。

\textbf{公式正答：} 1（single）。

\par\textbf{検討未確定：} 検索本文で措置を確認。当時全文の各選択肢への逐条対応、適用範囲のすべての例外は追加照合対象。

\paragraph{問題の要旨} 2014研究不正ガイドラインの組織責任と行政措置。

\Needspace{4\baselineskip}\paragraph{確認した条件・記述（独自要約）}\mbox{}\par

\noindent\textbf{原典：} 問題文・否定・全五候補と組合せ表を保存原典36画像で確認。\par

\noindent\textbf{時点：} 平成27年度の本文表現を保持。後の制度改正を遡及しない。\par

\paragraph{問われる核心} 2014研究不正ガイドラインの組織責任と行政措置。

\paragraph{必要知識} 研究倫理教育、適用範囲、機関への措置、管理条件と配分機関。

\paragraph{考え方・正答を選ぶ理由（査読待ち初稿）} 不適切1。平成26年ガイドラインは研究機関に管理条件を付し、履行されない場合の間接経費削減・配分停止を段階的に示す。個人の不正と機関の予防・調査体制は別の責任である。

\Needspace{4\baselineskip}\paragraph{選択肢の検討（査読待ち初稿）}\mbox{}\par

\noindent\textbf{1：不適切} ガイドラインは行政措置の根拠ではない。管理条件、間接経費削減、配分停止等の措置を規定している。\par

\noindent\textbf{2：条件内で適切} 管理条件未履行に間接経費等の措置。段階と条件を満たして研究機関への措置が行われる。自動的な全資金停止とはしない。\par

\noindent\textbf{3：条件内で適切} 学生や研究支援者も倫理教育の対象。研究者本人だけでなく研究に関わる者への教育を組織的に整える。\par

\noindent\textbf{4：条件内で適切} 競争的資金以外を含む省の予算による研究。当該ガイドラインの適用範囲は競争的資金という名称だけで区切らない。\par

\noindent\textbf{5：条件内で適切} 機関として責任ある不正防止体制。個人への処分に加えて組織の対応・管理責任を問う。\par

\paragraph{再利用する考え方} 研究倫理教育、適用範囲、機関への措置、管理条件と配分機関。 設問の否定・全称・基準年を確かめ、根拠と評価範囲を分けて判断する。

\paragraph{共通理論} \hyperref[theory:research:guideline2014]{第\ref*{theory:research:guideline2014}節}

\par 参考：\cite{mext_misconduct2014}。

\par 難易度 3（校準初稿）、優先度未評価。

\section{技術情報と説明・記録の行動}

\subsection{令和元年度 問10}
\label{q:r01-reg-apt-10}

\noindent\textbf{原典：} sources/R01/aptitude.pdf、PDF物理11頁。

\textbf{公式正答：} 3（single）。

\par\textbf{検討未確定：} アの技術情報の非公開性・守秘対象は前提を補い過ぎない。公開情報の会話まで禁止しない。カも記憶を使うこと自体を捏造とせず、情報が十分でないのに不確かさを確認せず回答する条件で判断する。

\paragraph{問題の要旨} 技術情報と説明・記録の六つの行動。

\Needspace{4\baselineskip}\paragraph{確認した条件・記述（独自要約）}\mbox{}\par

\noindent\textbf{ア：} 業務中の技術情報を友人へ話す。公開情報なら評価が変わり得る。\par

\noindent\textbf{イ：} 相手が専門知識を持たないとして仕様変更を説明しない。\par

\noindent\textbf{ウ：} 依頼者が求めなくても重要なリスクを説明する。\par

\noindent\textbf{エ：} 不要な専有情報の記憶媒体を家庭のゴミ箱へ捨てる。\par

\noindent\textbf{オ：} パスワードと施錠によるアクセス管理。これだけで全対策が完成するとはしない。\par

\noindent\textbf{カ：} 顧客の使用製品に関する質問へ十分なデータを持たず記憶で答える。\par

\paragraph{問われる核心} 技術情報と説明・記録の六つの行動。

\paragraph{必要知識} 守秘の範囲、仕様変更の説明、重要リスク、データ保存、アクセス制御。

\paragraph{考え方・正答を選ぶ理由（査読待ち初稿）} アは現在の業務の技術的内容を友人へ話す、イは仕様変更を専門知識のない顧客へ説明しない、エは専有情報入りの記憶媒体を家庭のゴミ箱へ捨てる、カは十分なデータを持たずに記憶で顧客の問い合わせに答えるため、設問では不適切4件の3。ウの重要リスクの説明とオの個人情報のパスワード・施錠管理は適切。

\Needspace{4\baselineskip}\paragraph{選択肢の検討（査読待ち初稿）}\mbox{}\par

\noindent\textbf{1：不適切} 2件。不適切はア・イ・エ・カの4件。 判定を先に行った集計と一致しない。\par

\noindent\textbf{2：不適切} 3件。不適切はア・イ・エ・カの4件。 判定を先に行った集計と一致しない。\par

\noindent\textbf{3：条件内で適切} 4件。不適切はア・イ・エ・カの4件。 集計と一致する。\par

\noindent\textbf{4：不適切} 5件。不適切はア・イ・エ・カの4件。 判定を先に行った集計と一致しない。\par

\noindent\textbf{5：不適切} 6件。不適切はア・イ・エ・カの4件。 判定を先に行った集計と一致しない。\par

\paragraph{再利用する考え方} 守秘の範囲、仕様変更の説明、重要リスク、データ保存、アクセス制御。 設問の否定・全称・基準年を確かめ、根拠と評価範囲を分けて判断する。

\paragraph{共通理論} \hyperref[theory:research:validation]{第\ref*{theory:research:validation}節}

\par 参考：\cite{ipej_ethics2011,scj_code2013}。

\par 難易度 3（校準初稿）、優先度未評価。

\section{暗黙知・形式知の変換}

\subsection{平成23年度 問10}
\label{q:h23-reg-apt-10}

\noindent\textbf{原典：} sources/H23/aptitude.pdf、PDF物理11頁。

\textbf{公式正答：} 3（single）。

\par\textbf{検討未確定：} 原典のISO27000と2005年の組合せは、27000:2009公開書誌と不一致である。原典を27001へ無断修正しない。規格全文の条項は未取得。

\paragraph{問題の要旨} 暗黙知と形式知、および情報セキュリティの三要素を語句と定義から対応させる。

\Needspace{4\baselineskip}\paragraph{確認した条件・記述（独自要約）}\mbox{}\par

\noindent\textbf{原典確認：} 問題幹と各記述の主体・否定・全称・表の5候補を保存原典画像で確認。\par

\noindent\textbf{制度時点：} H23年度の原典表現を保持。現行資料の条番号や仕様を過去へ補わない。\par

\noindent\textbf{各記述の判断：} 対応：ア暗黙知・イ形式知・ウ機密性・エ完全性・オ可用性。；規格疑義：原典は2005年ISO27000と記すが、公開書誌の27000初版は2009年である。\par

\noindent\textbf{前提：} 暗黙知は経験や技能として身に付く知、形式知は文章・図表等で表現された知である。機密性は権限、完全性は正確さ、可用性は必要時の利用を扱う。\par

\paragraph{問われる核心} 暗黙知と形式知、および情報セキュリティの三要素を語句と定義から対応させる。

\paragraph{必要知識} 暗黙知は経験や技能として身に付く知、形式知は文章・図表等で表現された知である。機密性は権限、完全性は正確さ、可用性は必要時の利用を扱う。

\paragraph{考え方・正答を選ぶ理由（査読待ち初稿）} 技能の勘所を暗黙知、マニュアルを形式知へ対応させる。認可された者だけという制限は機密性、正確で完全な状態は完全性、必要時に利用できることは可用性である。全て合う3を選ぶ。即時性は可用性の名称ではない。

\Needspace{4\baselineskip}\paragraph{選択肢の検討（査読待ち初稿）}\mbox{}\par

\noindent\textbf{1：不適切} 暗黙知・形式知・機密性・完全性・即時性。最後は可用性であり、即時性ではない。\par

\noindent\textbf{2：不適切} 経験知・客観知・特権性・完全性・即時性。知識管理の対語とCIAの名称が合わない。\par

\noindent\textbf{3：適切} 暗黙知・形式知・機密性・完全性・可用性。五つの定義に全て対応する。\par

\noindent\textbf{4：不適切} 暗黙知・客観知・機密性・保存性・可用性。形式知と完全性の位置を別語にしている。\par

\noindent\textbf{5：不適切} 経験知・形式知・特権性・保存性・可用性。暗黙知・機密性・完全性の標準名称に合わない。\par

\paragraph{再利用する考え方} 技能共有と情報保護を別の体系として整理し、用語を定義から判断する。

\paragraph{共通理論} \hyperref[theory:information:knowledge]{第\ref*{theory:information:knowledge}節}, \hyperref[theory:information:cia]{第\ref*{theory:information:cia}節}

\par 参考：\cite{jaist_knowledge,iso27000_2009}。

\paragraph{関連問題} \hyperref[q:r04-reg-apt-13]{R04-REG-APT-13}

\par 難易度 3（校準初稿）、優先度未評価。

\section{身体特徴の個人識別符号}

\subsection{令和元年度 問4}
\label{q:r01-reg-apt-04}

\noindent\textbf{原典：} sources/R01/aptitude.pdf、PDF物理5頁。

\textbf{公式正答：} 1（single）。

\par\textbf{検討未確定：} 現在の公開通則の列挙を確認。令和元年の施行令版の逐条差分までは未照合である。

\paragraph{問題の要旨} 身体特徴を変換した個人識別符号の七分類を判定する。

\Needspace{4\baselineskip}\paragraph{確認した条件・記述（独自要約）}\mbox{}\par

\noindent\textbf{原典確認：} 問題幹と各記述の主体・否定・全称・表の5候補を保存原典画像で確認。\par

\noindent\textbf{制度時点：} R01年度の原典表現を保持。現行資料の条番号や仕様を過去へ補わない。\par

\noindent\textbf{各記述の判断：} ア：DNA；イ：顔；ウ：虹彩；エ：声；オ：歩行；カ：静脈；キ：指紋・掌紋\par

\noindent\textbf{前提：} 政令のDNA・顔・虹彩・声・歩行・静脈・指紋掌紋に対応する。生の特徴全てではなく、電子計算機向けに変換され識別に足りる基準を満たす符号という前提である。\par

\paragraph{問われる核心} 身体特徴を変換した個人識別符号の七分類を判定する。

\paragraph{必要知識} 政令のDNA・顔・虹彩・声・歩行・静脈・指紋掌紋に対応する。生の特徴全てではなく、電子計算機向けに変換され識別に足りる基準を満たす符号という前提である。

\paragraph{考え方・正答を選ぶ理由（査読待ち初稿）} 七つの記述は政令第1条の七生体分類と全て対応する。含まれないもの0個、1である。列挙された身体そのものをそのまま符号と考えず、問題幹の変換・識別の前提を各記述に適用する。

\Needspace{4\baselineskip}\paragraph{選択肢の検討（査読待ち初稿）}\mbox{}\par

\noindent\textbf{1：適切} 0。含まれないなものは0個であり、一致する。\par

\noindent\textbf{2：不適切} 1。含まれないなものは0個であり、この個数とは一致しない。\par

\noindent\textbf{3：不適切} 2。含まれないなものは0個であり、この個数とは一致しない。\par

\noindent\textbf{4：不適切} 3。含まれないなものは0個であり、この個数とは一致しない。\par

\noindent\textbf{5：不適切} 4。含まれないなものは0個であり、この個数とは一致しない。\par

\paragraph{再利用する考え方} 生体の特徴と、認証用に処理された識別符号を区別する。

\paragraph{共通理論} \hyperref[theory:privacy:biometrics]{第\ref*{theory:privacy:biometrics}節}

\par 参考：\cite{privacy_general}。

\paragraph{関連問題} \hyperref[q:r06-reg-apt-09]{R06-REG-APT-09}

\par 難易度 3（校準初稿）、優先度未評価。

\subsection{令和6年度 問9}
\label{q:r06-reg-apt-09}

\noindent\textbf{原典：} sources/R06/aptitude.pdf、PDF物理9頁。

\textbf{公式正答：} 5（single）。

\par\textbf{検討未確定：} 出題当時の施行令版全文との差分は未照合。5の原典生理表現を肯定する説明ではなく、法令列挙の設問として扱う。

\paragraph{問題の要旨} 施行令に列挙される生体特徴から、骨の振動特性という非列挙項目を区別する。

\Needspace{4\baselineskip}\paragraph{確認した条件・記述（独自要約）}\mbox{}\par

\noindent\textbf{原典確認：} 問題幹と各記述の主体・否定・全称・表の5候補を保存原典画像で確認。\par

\noindent\textbf{制度時点：} R06年度の原典表現を保持。現行資料の条番号や仕様を過去へ補わない。\par

\noindent\textbf{各記述の判断：} 列挙：1指紋掌紋・2顔・3声・4歩行は該当、5の振動特性は列挙にない。\par

\noindent\textbf{前提：} 個人識別符号の生体分類は政令の列挙を確認する。生体認証に使えそうな特徴という一般論と、政令に定める分類は同じではない。\par

\paragraph{問われる核心} 施行令に列挙される生体特徴から、骨の振動特性という非列挙項目を区別する。

\paragraph{必要知識} 個人識別符号の生体分類は政令の列挙を確認する。生体認証に使えそうな特徴という一般論と、政令に定める分類は同じではない。

\paragraph{考え方・正答を選ぶ理由（査読待ち初稿）} 1から4は七生体分類の一部である。5の骨を振動させたときの特性はこの列挙にないので5である。原典の内耳を介さずという医学的表現を正しい生理記述へ校正し直して答えを作らない。

\Needspace{4\baselineskip}\paragraph{選択肢の検討（査読待ち初稿）}\mbox{}\par

\noindent\textbf{1：記載あり} 指紋又は掌紋。施行令の指紋・掌紋分類に対応する。\par

\noindent\textbf{2：記載あり} 顔の骨格・皮膚・部位配置の容貌。顔の特徴分類に対応する。\par

\noindent\textbf{3：記載あり} 発声時の声帯・声門・声道の変化。声の特徴分類に対応する。\par

\noindent\textbf{4：記載あり} 歩行の姿勢・腕の動作・歩幅等。歩行の特徴分類に対応する。\par

\noindent\textbf{5：記載なし} 骨振動による蝸牛への振動特性。政令の七生体分類にこの項目は列挙されていない。\par

\paragraph{再利用する考え方} 技術的な識別能力と法令上の列挙を区別する。

\paragraph{共通理論} \hyperref[theory:privacy:biometrics]{第\ref*{theory:privacy:biometrics}節}

\par 参考：\cite{privacy_general}。

\paragraph{関連問題} \hyperref[q:r01-reg-apt-04]{R01-REG-APT-04}

\par 難易度 3（校準初稿）、優先度未評価。

\section{本人同意・第三者提供・利用目的}

\subsection{平成26年度 問4}
\label{q:h26-reg-apt-04}

\noindent\textbf{原典：} sources/H26/aptitude.pdf、PDF物理4頁。

\textbf{公式正答：} 4（single）。

\par\textbf{検討未確定：} 学校名簿の配布は同意原則の問題であり、法定例外を全て否定するものではない。2013年指針の当時版全文は未照合。

\paragraph{問題の要旨} 公開情報の保護と、災害時の名簿共有・学校名簿の作成と配布の違いを判定する。

\Needspace{4\baselineskip}\paragraph{確認した条件・記述（独自要約）}\mbox{}\par

\noindent\textbf{原典確認：} 問題幹と各記述の主体・否定・全称・表の5候補を保存原典画像で確認。\par

\noindent\textbf{制度時点：} H26年度の原典表現を保持。現行資料の条番号や仕様を過去へ補わない。\par

\noindent\textbf{各記述の判断：} ア：○：識別できる評価・映像・音声も対象である。；イ：×：公知だから除外という規定ではない。；ウ：○：法定災害条件と必要範囲内の名簿共有である。；エ：○：通常の学校名簿配布で同意を必要とする。\par

\noindent\textbf{前提：} 適正取得・目的通知、利用、第三者提供を分ける。災害発生又は発生のおそれで生命身体保護へ特に必要な場合は、避難支援に必要な限度で同意なし提供が認められる。\par

\paragraph{問われる核心} 公開情報の保護と、災害時の名簿共有・学校名簿の作成と配布の違いを判定する。

\paragraph{必要知識} 適正取得・目的通知、利用、第三者提供を分ける。災害発生又は発生のおそれで生命身体保護へ特に必要な場合は、避難支援に必要な限度で同意なし提供が認められる。

\paragraph{考え方・正答を選ぶ理由（査読待ち初稿）} 公開済みという理由で保護を除くイが×、他は通常の条件下で○となり○×○○、4である。災害時の例外を平時の名簿の自由な提供へ拡張しない。名簿を内部で作る行為と家庭等へ配布する行為では規律が異なる。

\Needspace{4\baselineskip}\paragraph{選択肢の検討（査読待ち初稿）}\mbox{}\par

\noindent\textbf{1：不適切} ○・×・×・○。記述ウの真偽が逆である。\par

\noindent\textbf{2：不適切} ×・○・○・○。記述ア・イの真偽が逆である。\par

\noindent\textbf{3：不適切} ×・○・×・○。記述ア・イ・ウの真偽が逆である。\par

\noindent\textbf{4：適切} ○・×・○・○。各記述の判断と全て一致する。\par

\noindent\textbf{5：不適切} ○・○・○・×。記述イ・エの真偽が逆である。\par

\paragraph{再利用する考え方} 取得・利用・提供を時点と目的に沿って分け、例外の必要範囲を確認する。

\paragraph{共通理論} \hyperref[theory:privacy:scope]{第\ref*{theory:privacy:scope}節}, \hyperref[theory:privacy:consent]{第\ref*{theory:privacy:consent}節}, \hyperref[theory:privacy:evacuation]{第\ref*{theory:privacy:evacuation}節}

\par 参考：\cite{privacy_general,privacy_evacuation2013,privacy_evacuation_scope}。

\paragraph{関連問題} \hyperref[q:h24-reg-apt-04]{H24-REG-APT-04}

\par 難易度 3（校準初稿）、優先度未評価。

\subsection{令和3年度 問14}
\label{q:r03-reg-apt-14}

\noindent\textbf{原典：} sources/R03/aptitude.pdf、PDF物理14頁。

\textbf{公式正答：} 5（single）。

\par\textbf{独立解答と公式正答に相違あり：} 独立解答 2。初回独立2と公式5は一致しない。黙示同意自体を一律禁止する説明は不正確である。加盟店への連絡を承諾した具体的事実、契約条件、当時の直接Q\&Aは未確認であり、公式へ合わせて条件を補わない。

\paragraph{問題の要旨} 謝罪目的の連絡先提供・カード照会・目的外の求人勧誘・CC送信を本人同意と漏えいから評価する。

\Needspace{4\baselineskip}\paragraph{確認した条件・記述（独自要約）}\mbox{}\par

\noindent\textbf{原典確認：} 問題幹と各記述の主体・否定・全称・表の5候補を保存原典画像で確認。\par

\noindent\textbf{制度時点：} R03年度の原典表現を保持。現行資料の条番号や仕様を過去へ補わない。\par

\noindent\textbf{各記述の判断：} ア：×：被害側の同意なしに加害側の謝罪目的で連絡先を渡さない。；イ：公式では×：調査依頼と加盟店への第三者提供同意を同一視しない。；ウ：×：ポイント目的から求人勧誘への変更には目的範囲と同意の検討が必要である。；エ：×：CCは他の宛先へアドレスを開示する。\par

\noindent\textbf{前提：} 本人同意は示された取扱方法への承諾の意思表示である。明示に限られないが黙示同意は具体的事実で判断する。本人の調査依頼だけで加盟店への連絡先提供を当然に認めない。\par

\paragraph{問われる核心} 謝罪目的の連絡先提供・カード照会・目的外の求人勧誘・CC送信を本人同意と漏えいから評価する。

\paragraph{必要知識} 本人同意は示された取扱方法への承諾の意思表示である。明示に限られないが黙示同意は具体的事実で判断する。本人の調査依頼だけで加盟店への連絡先提供を当然に認めない。

\paragraph{考え方・正答を選ぶ理由（査読待ち初稿）} 公式は四つとも×、5である。アは謝罪する側の事情で被害側の同意を代替している。ウは取得時の利用目的を求人へ広げ、エは顧客間でアドレスを漏らす。イの初回は調査依頼から黙示の承諾があると解して2としたが、PPCの定義では提供方法への承諾を認識できる具体的条件が必要であり、問題にその条件は示されない。公式5を保持し、初回2との相違と同意条件の疑義を査読へ残す。

\Needspace{4\baselineskip}\paragraph{選択肢の検討（査読待ち初稿）}\mbox{}\par

\noindent\textbf{1：不適切} ○・×・×・×。記述アの真偽が逆である。\par

\noindent\textbf{2：不適切} ×・○・×・×。記述イの真偽が逆である。\par

\noindent\textbf{3：不適切} ×・×・○・×。記述ウの真偽が逆である。\par

\noindent\textbf{4：不適切} ×・×・×・○。記述エの真偽が逆である。\par

\noindent\textbf{5：適切} ×・×・×・×。各記述の判断と全て一致する。\par

\paragraph{再利用する考え方} 便宜・善意・契約履行の必要性と、本人が承諾した提供範囲を別々に確認する。

\paragraph{共通理論} \hyperref[theory:privacy:consent]{第\ref*{theory:privacy:consent}節}

\par 参考：\cite{privacy_general,privacy_consent2017,privacy_consent_current}。

\par 難易度 3（校準初稿）、優先度未評価。

\section{個人情報の定義・公開と識別}

\subsection{平成23年度 問7}
\label{q:h23-reg-apt-07}

\noindent\textbf{原典：} sources/H23/aptitude.pdf、PDF物理8頁。

\textbf{公式正答：} 5（single）。

\par\textbf{検討未確定：} 配送物の中身を取り扱わない通常条件に依存する。当時ガイドラインの該当段落と責任主体の直接照合は未完了。現行Q7-35は中身の合意・発送者の安全管理を条件としている。

\paragraph{問題の要旨} 死者の家族・公開情報・監視画像・封入媒体の運送を、個人情報と取扱委託の条件から判定する。

\Needspace{4\baselineskip}\paragraph{確認した条件・記述（独自要約）}\mbox{}\par

\noindent\textbf{原典確認：} 問題幹と各記述の主体・否定・全称・表の5候補を保存原典画像で確認。\par

\noindent\textbf{制度時点：} H23年度の原典表現を保持。現行資料の条番号や仕様を過去へ補わない。\par

\noindent\textbf{各記述の判断：} ア：×：生存家族について識別できる情報は対象となる。；イ：×：公開は個人情報性を消さない。；ウ：×：個人を識別できる画像は対象となる。；エ：×：中身を取り扱わない通常の配送だけで中身の個人データ取扱委託としない。\par

\noindent\textbf{前提：} 死者自身の情報と生存する家族の情報を分ける。公開済みでも個人情報である。識別可能な映像は対象となる。宛名情報と配送物内部のCDデータは別の取扱いである。\par

\paragraph{問われる核心} 死者の家族・公開情報・監視画像・封入媒体の運送を、個人情報と取扱委託の条件から判定する。

\paragraph{必要知識} 死者自身の情報と生存する家族の情報を分ける。公開済みでも個人情報である。識別可能な映像は対象となる。宛名情報と配送物内部のCDデータは別の取扱いである。

\paragraph{考え方・正答を選ぶ理由（査読待ち初稿）} ア・イ・ウの一律除外は定義に反する。エは通常の配送で中身の取扱合意がないという条件で、中身のデータの委託先とはみなさない。四つとも×、5である。誤配に契約上等の責任が一切ないとはいえず、発送者の安全管理義務も残る。

\Needspace{4\baselineskip}\paragraph{選択肢の検討（査読待ち初稿）}\mbox{}\par

\noindent\textbf{1：不適切} ○・×・×・×。記述アの真偽が逆である。\par

\noindent\textbf{2：不適切} ×・○・×・×。記述イの真偽が逆である。\par

\noindent\textbf{3：不適切} ×・×・○・×。記述ウの真偽が逆である。\par

\noindent\textbf{4：不適切} ×・×・×・○。記述エの真偽が逆である。\par

\noindent\textbf{5：適切} ×・×・×・×。各記述の判断と全て一致する。\par

\paragraph{再利用する考え方} 情報の対象者と、事業者が実際に扱う情報の範囲を個別に決める。

\paragraph{共通理論} \hyperref[theory:privacy:scope]{第\ref*{theory:privacy:scope}節}, \hyperref[theory:privacy:delivery]{第\ref*{theory:privacy:delivery}節}

\par 参考：\cite{privacy_general,privacy_delivery}。

\paragraph{関連問題} \hyperref[q:h24-reg-apt-04]{H24-REG-APT-04}

\par 難易度 3（校準初稿）、優先度未評価。

\subsection{平成24年度 問4}
\label{q:h24-reg-apt-04}

\noindent\textbf{原典：} sources/H24/aptitude.pdf、PDF物理3頁。

\textbf{公式正答：} 2（single）。

\par\textbf{検討未確定：} 5000人基準の旧政令全文と重複人数の計数は追加照合が必要である。番号だけで一律に個人情報とするのでなく容易な照合の条件を明示する。

\paragraph{問題の要旨} 旧5000人基準と、評価情報・識別番号・防犯映像の個人情報性を判定する。

\Needspace{4\baselineskip}\paragraph{確認した条件・記述（独自要約）}\mbox{}\par

\noindent\textbf{原典確認：} 問題幹と各記述の主体・否定・全称・表の5候補を保存原典画像で確認。\par

\noindent\textbf{制度時点：} H24年度の原典表現を保持。現行資料の条番号や仕様を過去へ補わない。\par

\noindent\textbf{各記述の判断：} ア：○：判断・評価も個人に関する情報である。；イ：×：個々のDBだけを数えて判断しない。；ウ：○：本人へ容易に照合できる従業員・学籍番号等を想定する。；エ：×：防犯目的は個人情報性の除外理由ではない。\par

\noindent\textbf{前提：} 当時の事業者適用除外の人数は事業者全体で判断し、DBを分割しただけでは適用を免れない。番号は容易な照合による識別性を確認する。2017年に小規模事業者の除外は廃止された。\par

\paragraph{問われる核心} 旧5000人基準と、評価情報・識別番号・防犯映像の個人情報性を判定する。

\paragraph{必要知識} 当時の事業者適用除外の人数は事業者全体で判断し、DBを分割しただけでは適用を免れない。番号は容易な照合による識別性を確認する。2017年に小規模事業者の除外は廃止された。

\paragraph{考え方・正答を選ぶ理由（査読待ち初稿）} アとウは○である。ウは番号から本人へ容易に照合できる組織内の番号という通常条件に沿う。イは事業者全体の人数をDB単位へ分割して判断している。エは識別できる映像を目的だけで除外する。○×○×、2である。

\Needspace{4\baselineskip}\paragraph{選択肢の検討（査読待ち初稿）}\mbox{}\par

\noindent\textbf{1：不適切} ×・○・×・○。記述ア・イ・ウ・エの真偽が逆である。\par

\noindent\textbf{2：適切} ○・×・○・×。各記述の判断と全て一致する。\par

\noindent\textbf{3：不適切} ○・×・×・○。記述ウ・エの真偽が逆である。\par

\noindent\textbf{4：不適切} ○・○・×・×。記述イ・ウの真偽が逆である。\par

\noindent\textbf{5：不適切} ×・○・○・×。記述ア・イの真偽が逆である。\par

\paragraph{再利用する考え方} 適用対象の旧人数基準と、情報自体の識別性を混同しない。

\paragraph{共通理論} \hyperref[theory:privacy:scope]{第\ref*{theory:privacy:scope}節}, \hyperref[theory:privacy:history]{第\ref*{theory:privacy:history}節}

\par 参考：\cite{privacy_general,privacy_small2017}。

\paragraph{関連問題} \hyperref[q:h23-reg-apt-07]{H23-REG-APT-07}, \hyperref[q:h26-reg-apt-04]{H26-REG-APT-04}

\par 難易度 3（校準初稿）、優先度未評価。

\section{CIAと不審メールの初動}

\subsection{平成28年度 問13}
\label{q:h28-reg-apt-13}

\noindent\textbf{原典：} sources/H28/aptitude.pdf、PDF物理13頁。

\textbf{公式正答：} 4（single）。

\paragraph{問題の要旨} 完全性は正確さ、可用性は必要な時に使える性質。

\Needspace{4\baselineskip}\paragraph{確認した条件・記述（独自要約）}\mbox{}\par

\noindent\textbf{原典：} 問題文・否定・全五候補と組合せ表を保存原典27画像で確認。\par

\noindent\textbf{時点：} 平成28年度の本文表現を保持。後の制度改正を遡及しない。\par

\paragraph{問われる核心} 完全性は正確さ、可用性は必要な時に使える性質。

\paragraph{必要知識} 機密性・完全性・可用性、PDCA、事故の初動・分析・復旧・予防。

\paragraph{考え方・正答を選ぶ理由（査読待ち初稿）} アのCIA、ウの方針と改善、エの全員の教育、オの事故対応と見直しは正。イが可用性を正確さ・改ざん防止という完全性の意味へ取り違えているので誤、4。秘密を守るだけでは、正しい情報が必要時に使えることを保証しない。

\Needspace{4\baselineskip}\paragraph{選択肢の検討（査読待ち初稿）}\mbox{}\par

\noindent\textbf{1：不適切} XOOXO。各文又は定義の判定を先に行う。並びが判定 OXOOO と異なる。\par

\noindent\textbf{2：不適切} XXOOX。各文又は定義の判定を先に行う。並びが判定 OXOOO と異なる。\par

\noindent\textbf{3：不適切} OOXOX。各文又は定義の判定を先に行う。並びが判定 OXOOO と異なる。\par

\noindent\textbf{4：条件内で適切} OXOOO。各文又は定義の判定を先に行う。全項目が一致する。\par

\noindent\textbf{5：不適切} OOXOO。各文又は定義の判定を先に行う。並びが判定 OXOOO と異なる。\par

\paragraph{再利用する考え方} 対象・責任主体・時点・要件を分け、文の否定と「全て」「不要」の範囲を点検する。知識を当てはめる前に何を証明する資料かを確かめる。

\paragraph{共通理論} \hyperref[theory:research:security]{第\ref*{theory:research:security}節}

\par 参考：\cite{saf53_ipa_cia}。

\paragraph{関連問題} \hyperref[q:h29-reg-apt-06]{H29-REG-APT-06}

\par 難易度 2（校準初稿）、優先度未評価。

\subsection{平成29年度 問6}
\label{q:h29-reg-apt-06}

\noindent\textbf{原典：} sources/H29/aptitude.pdf、PDF物理6頁。

\textbf{公式正答：} 1（single）。

\paragraph{問題の要旨} 不審な依頼は既知の担当窓口へ速やかに報告する。

\Needspace{4\baselineskip}\paragraph{確認した条件・記述（独自要約）}\mbox{}\par

\noindent\textbf{原典：} 問題文・否定・全五候補と組合せ表を保存原典27画像で確認。\par

\noindent\textbf{時点：} 平成29年度の本文表現を保持。後の制度改正を遡及しない。\par

\paragraph{問われる核心} 不審な依頼は既知の担当窓口へ速やかに報告する。

\paragraph{必要知識} 送信者の表示と認証の区別、独立した連絡経路、早期報告、情報共有と被害拡大防止。

\paragraph{考え方・正答を選ぶ理由（査読待ち初稿）} 適切1。毎年の時期と異なる依頼は攻撃の確定証拠ではないが、確認へ進む合理的な兆候である。メールの送信名や記載窓口ではなく、平常から把握している社内担当へ報告し、専門部署が組織の対応を判断する。

\Needspace{4\baselineskip}\paragraph{選択肢の検討（査読待ち初稿）}\mbox{}\par

\noindent\textbf{1：条件内で適切} 社内担当部署へ報告。既知の社内経路で疑いを共有すれば、担当者が他の受信者への影響と真偽を確認できる。メール内の連絡先を信用の根拠にしない。\par

\noindent\textbf{2：不適切} メール記載の担当者へ連絡。不審なメッセージが示した窓口も攻撃者の可能性があり、独立確認にならない。\par

\noindent\textbf{3：不適切} 再度違和感が出るまで待つ。既に異常を認めているのに情報共有を遅らせる。担当部署で確認する機会を失う。\par

\noindent\textbf{4：不適切} そのまま回答。依頼の真偽が未確認で、添付操作や情報の提供が危害を広げ得る。\par

\noindent\textbf{5：不適切} 自部署だけへ伝えて待つ。専門担当部署へ届かず、組織全体の対応や情報共有を遅らせる。\par

\paragraph{再利用する考え方} 対象・責任主体・時点・要件を分け、文の否定と「全て」「不要」の範囲を点検する。知識を当てはめる前に何を証明する資料かを確かめる。

\paragraph{共通理論} \hyperref[theory:research:security]{第\ref*{theory:research:security}節}

\par 参考：\cite{saf53_ipa_mail}。

\paragraph{関連問題} \hyperref[q:h28-reg-apt-13]{H28-REG-APT-13}

\par 難易度 2（校準初稿）、優先度未評価。

\section{秘密情報の接近制御・相互確認と持出し防止}

\subsection{令和元年度 問9}
\label{q:r01-reg-apt-09}

\noindent\textbf{原典：} sources/R01/aptitude.pdf、PDF物理10頁。

\textbf{公式正答：} 2（single）。

\paragraph{問題の要旨} 秘密情報のアクセス、相互確認、表示、配置、暗号・閲覧制限、紙削減の6記述で不適切な数を求める。

\Needspace{4\baselineskip}\paragraph{確認した条件・記述（独自要約）}\mbox{}\par

\noindent\textbf{方向：} 不適切な記述の数。候補0/1/2/3/4。\par

\noindent\textbf{構成：} ア〜カの6項目。\par

\noindent\textbf{限定：} オは外部では閲覧できないように設定する対策。暗号自体の万能性や外部送信の許可ではない。\par

\paragraph{問われる核心} 秘密情報の接近制御と相互確認

\paragraph{必要知識} 秘密情報の接近制御と相互確認。ア○：情報の分類に応じたアクセス権で、必要のない人の接近を制限する。イ×：相互の確認や責任者のチェックを減らす単独作業の一律推奨は逆方向である。ウ○：秘密の表示で取扱者の認識を支える。エ○：資料・作業が死角になることを減らすが、権限のない者へ画面を見せてよい意味ではない。オ○：暗号と登録ID・端末等の閲覧権制限を適切に設定する対策である。暗号化だけで漏えいが絶対不可能になるという主張ではなく、無断送信を許すものでもない。カ○：電子的アクセス管理と組み合わせ、無権限者が紙に接する機会を減らす。不適切はイの1項目で2。

\paragraph{考え方・正答を選ぶ理由（査読待ち初稿）} ア○：情報の分類に応じたアクセス権で、必要のない人の接近を制限する。イ×：相互の確認や責任者のチェックを減らす単独作業の一律推奨は逆方向である。ウ○：秘密の表示で取扱者の認識を支える。エ○：資料・作業が死角になることを減らすが、権限のない者へ画面を見せてよい意味ではない。オ○：暗号と登録ID・端末等の閲覧権制限を適切に設定する対策である。暗号化だけで漏えいが絶対不可能になるという主張ではなく、無断送信を許すものでもない。カ○：電子的アクセス管理と組み合わせ、無権限者が紙に接する機会を減らす。不適切はイの1項目で2。

\Needspace{4\baselineskip}\paragraph{選択肢の検討（査読待ち初稿）}\mbox{}\par

\noindent\textbf{1：選択対象ではない} 0項目。単独作業を可能な限り推奨するイが不適切なので0ではない。\par

\noindent\textbf{2：選択対象} 1項目。イだけが不適切であり1項目となる。\par

\noindent\textbf{3：選択対象ではない} 2項目。イに加えて他の適切な対策も否定すると数え過ぎとなる。\par

\noindent\textbf{4：選択対象ではない} 3項目。アクセス権の設定・表示・配置・適切な閲覧制限・紙削減の全てを危険低減策として区別すれば3ではない。\par

\noindent\textbf{5：選択対象ではない} 4項目。適切な対策まで不適切に数えている。秘密保護と一人作業の推奨を混同しない。\par

\paragraph{再利用する考え方} 対象、条件、時点を先に確認する。秘密情報の接近制御と相互確認を具体的な事例へ適用し、技術的原因・法的要件・倫理的判断のどれを問うか区別する。

\paragraph{共通理論} \hyperref[theory:research:secret_management]{第\ref*{theory:research:secret_management}節}, \hyperref[theory:compliance:secret]{第\ref*{theory:compliance:secret}節}

\par 参考：\cite{meti_secret_handbook2024}。

\paragraph{関連問題} \hyperref[q:r05-reg-apt-02]{R05-REG-APT-02}

\par 難易度 2（校準初稿）、優先度未評価。

\subsection{令和5年度 問2}
\label{q:r05-reg-apt-02}

\noindent\textbf{原典：} sources/R05/aptitude.pdf、PDF物理2頁。

\textbf{公式正答：} 3（single）。

\paragraph{問題の要旨} アクセス権、秘密表示、単独作業、電子的コピー制限、座席配置という漏えい対策を判定する。

\Needspace{4\baselineskip}\paragraph{確認した条件・記述（独自要約）}\mbox{}\par

\noindent\textbf{方向：} 最も不適切。\par

\noindent\textbf{時点：} 令和5年。2024年版ハンドブックは現在の背景資料として区別する。\par

\noindent\textbf{範囲：} 5つの対策案。電子的制限の完全防御を主張していない。\par

\paragraph{問われる核心} 秘密保護におけるアクセス権と相互確認

\paragraph{必要知識} 秘密保護におけるアクセス権と相互確認。情報の接近制御、秘密の認識、持出しの困難化、相互チェックを組み合わせる。3だけはチェックを減らす単独作業を一律に推奨しているので不適切。複数人にアクセス権を無制限に与えることとは区別する。

\paragraph{考え方・正答を選ぶ理由（査読待ち初稿）} 情報の接近制御、秘密の認識、持出しの困難化、相互チェックを組み合わせる。3だけはチェックを減らす単独作業を一律に推奨しているので不適切。複数人にアクセス権を無制限に与えることとは区別する。

\Needspace{4\baselineskip}\paragraph{選択肢の検討（査読待ち初稿）}\mbox{}\par

\noindent\textbf{1：適切} 情報分類ごとにアクセス権の設定ルールと適切な範囲を定める。。必要性に基づくアクセス範囲の制御である。権限を無制限に広げる対策ではない。\par

\noindent\textbf{2：適切} 媒体へ秘密情報であることが分かる表示をする。。取扱者が秘密の分類と注意を認識する仕組みを支える。\par

\noindent\textbf{3：不適切} 複数人を避け、可能な限り単独で秘密情報を扱う。。必要に応じ複数人で相互確認できる体制を整え、単独が必要なら責任者が事前・事後に確認する。単独を最大化する方針は逆である。\par

\noindent\textbf{4：適切} 印刷・複製・外部メディアへのコピー等を制限する設定をする。。持出しの経路を減らす技術的対策となる。印刷禁止やコピーガードが撮影・再入力まで完全に防ぐとの保証ではない。\par

\noindent\textbf{5：適切} 作業・資料が死角になりにくい配置を検討する。。相互に作業へ注意が向く環境を作る。秘密画面をアクセス権のない人にさらすことを避けるという別条件も守る。\par

\paragraph{再利用する考え方} 対象、条件、時点を先に確認する。秘密保護におけるアクセス権と相互確認を具体的な事例へ適用し、技術的原因・法的要件・倫理的判断のどれを問うか区別する。

\paragraph{共通理論} \hyperref[theory:research:secret_management]{第\ref*{theory:research:secret_management}節}, \hyperref[theory:compliance:secret]{第\ref*{theory:compliance:secret}節}

\par 参考：\cite{meti_secret_handbook2024}。

\paragraph{関連問題} \hyperref[q:r01-reg-apt-09]{R01-REG-APT-09}

\par 難易度 2（校準初稿）、優先度未評価。

\section{CIA三要素と方針・管理体系}

\subsection{令和4年度 問13}
\label{q:r04-reg-apt-13}

\noindent\textbf{原典：} sources/R04/aptitude.pdf、PDF物理12頁。

\textbf{公式正答：} 3（single）。

\par\textbf{検討未確定：} 規格全文は未取得であり各措置を特定条項の逐語義務とは断定しない。

\paragraph{問題の要旨} CIA三要素と、方針・教育・監視・事故対応による継続改善を判定する。

\Needspace{4\baselineskip}\paragraph{確認した条件・記述（独自要約）}\mbox{}\par

\noindent\textbf{原典確認：} 問題幹と各記述の主体・否定・全称・表の5候補を保存原典画像で確認。\par

\noindent\textbf{制度時点：} R04年度の原典表現を保持。現行資料の条番号や仕様を過去へ補わない。\par

\noindent\textbf{各記述の判断：} ア：○：CIAを維持する。；イ：×：これは可用性でなく完全性である。；ウ：○：方針の有効性を継続的に確認する。；エ：○：幹部も含む教育を実施する。；オ：○：初動・分析・復旧・再発防止を方針へ反映する。\par

\noindent\textbf{前提：} 可用性は必要時の利用、完全性は情報や処理の正確さと完全な状態である。方針を実施し監視・事故分析から改善するのが管理体系である。\par

\paragraph{問われる核心} CIA三要素と、方針・教育・監視・事故対応による継続改善を判定する。

\paragraph{必要知識} 可用性は必要時の利用、完全性は情報や処理の正確さと完全な状態である。方針を実施し監視・事故分析から改善するのが管理体系である。

\paragraph{考え方・正答を選ぶ理由（査読待ち初稿）} イだけ×であり、○×○○○、3である。正しい内容が書かれていても、対応する名称が可用性に入れ替わっている。ウ・エ・オは方針を文書で作って終わらせず、運用結果を点検して改善する枠組みに対応する。

\Needspace{4\baselineskip}\paragraph{選択肢の検討（査読待ち初稿）}\mbox{}\par

\noindent\textbf{1：不適切} ×・○・○・×・○。記述ア・イ・エの真偽が逆である。\par

\noindent\textbf{2：不適切} ×・×・○・○・○。記述アの真偽が逆である。\par

\noindent\textbf{3：適切} ○・×・○・○・○。各記述の判断と全て一致する。\par

\noindent\textbf{4：不適切} ○・○・×・○・×。記述イ・ウ・オの真偽が逆である。\par

\noindent\textbf{5：不適切} ○・○・×・○・○。記述イ・ウの真偽が逆である。\par

\paragraph{再利用する考え方} 名称の定義と、運用を継続改善する仕組みを別々に確かめる。

\paragraph{共通理論} \hyperref[theory:information:cia]{第\ref*{theory:information:cia}節}, \hyperref[theory:information:management]{第\ref*{theory:information:management}節}

\par 参考：\cite{iso27000_2009}。

\paragraph{関連問題} \hyperref[q:h23-reg-apt-10]{H23-REG-APT-10}, \hyperref[q:r07-reg-apt-05]{R07-REG-APT-05}

\par 難易度 3（校準初稿）、優先度未評価。

\subsection{令和7年度 問5}
\label{q:r07-reg-apt-05}

\noindent\textbf{原典：} sources/R07/aptitude.pdf、PDF物理5頁。

\textbf{公式正答：} 1（single）。

\par\textbf{検討未確定：} 27001の該当版の全条項は未取得。CIA名称と方針・体系の概念差を初稿として説明する。

\paragraph{問題の要旨} セキュリティ方針と管理体系の名称の入替えを、CIA三要素と比較する。

\Needspace{4\baselineskip}\paragraph{確認した条件・記述（独自要約）}\mbox{}\par

\noindent\textbf{原典確認：} 問題幹と各記述の主体・否定・全称・表の5候補を保存原典画像で確認。\par

\noindent\textbf{制度時点：} R07年度の原典表現を保持。現行資料の条番号や仕様を過去へ補わない。\par

\noindent\textbf{各記述の判断：} ア：×：ISO/IEC27001の管理体系を方針という名称にしている。；イ：○：必要時の認可利用は可用性である。；ウ：○：正確さと完全な状態は完全性である。；エ：○：認可者だけのアクセスは機密性である。；オ：×：方針の定義を管理体系という名称にしている。\par

\noindent\textbf{前提：} ポリシーは組織の方針と行動指針、ISMSはその方針を計画・実施・評価・改善する管理体系である。CIAの定義と独立に名称の対応を確認する。\par

\paragraph{問われる核心} セキュリティ方針と管理体系の名称の入替えを、CIA三要素と比較する。

\paragraph{必要知識} ポリシーは組織の方針と行動指針、ISMSはその方針を計画・実施・評価・改善する管理体系である。CIAの定義と独立に名称の対応を確認する。

\paragraph{考え方・正答を選ぶ理由（査読待ち初稿）} イ・ウ・エは各三要素の定義に合う。アとオは管理体系と方針の名称を入れ替えた説明であるため×。×○○○×、1である。方針は管理体系に含まれるが、方針文書だけで管理体系の全てを代替しない。

\Needspace{4\baselineskip}\paragraph{選択肢の検討（査読待ち初稿）}\mbox{}\par

\noindent\textbf{1：適切} ×・○・○・○・×。各記述の判断と全て一致する。\par

\noindent\textbf{2：不適切} ○・×・○・○・○。記述ア・イ・オの真偽が逆である。\par

\noindent\textbf{3：不適切} ○・○・×・×・○。記述ア・ウ・エ・オの真偽が逆である。\par

\noindent\textbf{4：不適切} ×・○・○・×・○。記述エ・オの真偽が逆である。\par

\noindent\textbf{5：不適切} ○・○・○・○・×。記述アの真偽が逆である。\par

\paragraph{再利用する考え方} 同じ情報管理に関する語でも方針・実施・評価を分けて読む。

\paragraph{共通理論} \hyperref[theory:information:cia]{第\ref*{theory:information:cia}節}, \hyperref[theory:information:management]{第\ref*{theory:information:management}節}

\par 参考：\cite{iso27000_2009}。

\paragraph{関連問題} \hyperref[q:r04-reg-apt-13]{R04-REG-APT-13}

\par 難易度 3（校準初稿）、優先度未評価。


\chapter{AIと社会変革}
\section{技術能力と社会の選択}
\label{theory:ai:society}
AIが予測や生成を行えることと、その利用が社会に望ましいことは別の判断である。2017年内閣府報告書は倫理・法・経済・教育・社会・研究開発の六つの観点を置き、実例から課題を整理した\parencite{ai_society2017}。人の感情や判断への気付かない誘導、創作の価値の評価、データ利用と権利保護、利用能力の差による新たな格差は、その異なる側面である。

例えば推薦サービスは利用者に役立つ情報を選ぶ一方、見える情報の範囲を狭める場合がある。精度という指標だけでは、誰が目的を決め、何を推薦対象から除き、利用者へどのような選択の自由を与えるかは決まらない。AIを使わない選択にも配慮し、人とシステムの役割分担を具体的に設計する。

\section{AI利活用の十原則と主体ごとの対策}
\label{theory:ai:principles}
2019年総務省ガイドラインは、適正利用・適正学習・連携・安全・セキュリティ・プライバシー・尊厳自律・公平性・透明性・説明責任の十原則を整理した\parencite{ai_use2019}。適正利用は目的と役割、適正学習は入力データの質、連携はネットワーク化によるリスク、最後の四つは人の判断や待遇と説明の責任へ目を向ける。

開発者だけでなくサービス提供者、業務利用者、消費者、データ提供者で必要な対策は異なる。用途と文脈に応じて原則と措置を選ぶソフトローであり、全利用へ全措置を同じ形で要求する法律ではない。透明性も、全てのソースコードや個人データを公開する義務と同じではない。検証と説明ができる範囲を、秘密とプライバシーの保護とともに検討する。

\section{知識集約社会と対話型生成AI}
\label{theory:ai:knowledge_society}
2024年白書は、知識を価値創造の中心とし、スマート化によって製品やサービスの高付加価値化が進む社会への変化を述べる\parencite{mext_ai_whitepaper2024}。生成AIの対話型インターフェースは専門家以外にも接近しやすい入口となる。ただし入口の容易さは出力の正確さを保証しない。利用者は目的と根拠を確認し、誤りの影響が大きい判断では適切な検証を行う。

\paragraph{出題傾向の読み方}
本章の五問は、AIと社会変革について異なる資料と時点を扱う。これだけでAIが毎年必ず出題されるとはいえない。政策文書の語句と、技術が社会へ及ぼす影響を考える枠組みを分けて学ぶ。
\section{Society5.0：技術融合と人間中心の目的}
\label{theory:society:five}
狩猟社会、農耕社会、工業社会、情報社会に続く第五の社会という整理を読む。サイバー空間と現実空間の融合によって、経済的発展と社会的課題の解決の両立を目指す。中心に置くのは人間であり、技術の導入量だけを目的にしない\parencite{eth54_society2021}。

2016年の提案と2021年の計画を分ける。2021年には持続可能性、強靱性、安全・安心、多様な幸福への説明が具体化される。人文・社会科学も含む知識をつなぎ、倫理・法・社会への影響を検討する。過去問の穴埋めでは、社会段階、社会的課題、人間中心、持続可能性を別の空欄として確かめる。後年の政策用語を過去の指定語へ無条件に置き換えない。

\section{AIの社会的論点と知識集約社会}

\subsection{令和元年度再試験 問15}
\label{q:r01-ret-apt-15}

\noindent\textbf{原典：} sources/R01\_retest/aptitude.pdf、PDF物理16頁。

\textbf{公式正答：} 1（single）。

\paragraph{問題の要旨} AIの感情誘導・価値評価・データ保護・技術格差を2017年の六論点から判定する。

\Needspace{4\baselineskip}\paragraph{確認した条件・記述（独自要約）}\mbox{}\par

\noindent\textbf{原典確認：} 問題幹と各記述の主体・否定・全称・表の5候補を保存原典画像で確認。\par

\noindent\textbf{制度時点：} R01年度の原典表現を保持。現行資料の条番号や仕様を過去へ補わない。\par

\noindent\textbf{各記述の判断：} ア：○：本人が気付かない感情等の誘導は倫理課題である。；イ：○：生産性と創作価値の多様性を考える必要がある。；ウ：○：プライバシー保護と利活用の過度な萎縮防止を両立させる。；エ：○：AIリテラシーの差によるデバイドが起こり得る。\par

\noindent\textbf{前提：} 2017年報告書は倫理・法・経済・教育・社会・研究開発の六観点で便益と課題を検討する。社会の価値と機械の精度は別であり、利用できる能力の格差も課題となる。\par

\paragraph{問われる核心} AIの感情誘導・価値評価・データ保護・技術格差を2017年の六論点から判定する。

\paragraph{必要知識} 2017年報告書は倫理・法・経済・教育・社会・研究開発の六観点で便益と課題を検討する。社会の価値と機械の精度は別であり、利用できる能力の格差も課題となる。

\paragraph{考え方・正答を選ぶ理由（査読待ち初稿）} 四記述は報告書の倫理・法・社会の検討に対応する。不適切0個、1である。生産性向上は健全な利活用への期待であり、どの導入でも必ず達成されるという定理ではない。格差の可能性を、利用しない人への強制と取り違えない。

\Needspace{4\baselineskip}\paragraph{選択肢の検討（査読待ち初稿）}\mbox{}\par

\noindent\textbf{1：適切} 0。不適切なものは0個であり、一致する。\par

\noindent\textbf{2：不適切} 1。不適切なものは0個であり、この個数とは一致しない。\par

\noindent\textbf{3：不適切} 2。不適切なものは0個であり、この個数とは一致しない。\par

\noindent\textbf{4：不適切} 3。不適切なものは0個であり、この個数とは一致しない。\par

\noindent\textbf{5：不適切} 4。不適切なものは0個であり、この個数とは一致しない。\par

\paragraph{再利用する考え方} AIの技術能力、制度、価値判断、利用の機会を別軸で考える。

\paragraph{共通理論} \hyperref[theory:ai:society]{第\ref*{theory:ai:society}節}

\par 参考：\cite{ai_society2017}。

\paragraph{関連問題} \hyperref[q:r03-reg-apt-06]{R03-REG-APT-06}

\par 難易度 3（校準初稿）、優先度未評価。

\subsection{令和7年度 問7}
\label{q:r07-reg-apt-07}

\noindent\textbf{原典：} sources/R07/aptitude.pdf、PDF物理7頁。

\textbf{公式正答：} 1（single）。

\paragraph{問題の要旨} 2024年白書の知識集約社会と生成AIの普及を四つの語句へ対応させる。

\Needspace{4\baselineskip}\paragraph{確認した条件・記述（独自要約）}\mbox{}\par

\noindent\textbf{原典確認：} 問題幹と各記述の主体・否定・全称・表の5候補を保存原典画像で確認。\par

\noindent\textbf{制度時点：} R07年度の原典表現を保持。現行資料の条番号や仕様を過去へ補わない。\par

\noindent\textbf{各記述の判断：} ア：知識；イ：高付加価値化；ウ：医療；エ：対話\par

\noindent\textbf{前提：} 白書は知識を価値創造の中心とする社会、製品やサービスの高付加価値化、医療への影響、対話型インターフェースの普及を説明する。\par

\paragraph{問われる核心} 2024年白書の知識集約社会と生成AIの普及を四つの語句へ対応させる。

\paragraph{必要知識} 白書は知識を価値創造の中心とする社会、製品やサービスの高付加価値化、医療への影響、対話型インターフェースの普及を説明する。

\paragraph{考え方・正答を選ぶ理由（査読待ち初稿）} 2024年白書の対象本文は知識集約型、高付加価値化、医療、対話型生成AIと記している。四語が全て合う1である。AIの利活用の一般的な倫理課題が存在することから、本文の医療という産業分野を倫理へ置き換えない。

\Needspace{4\baselineskip}\paragraph{選択肢の検討（査読待ち初稿）}\mbox{}\par

\noindent\textbf{1：適切} 知識・高付加価値化・医療・対話。四空欄が白書の対象記述に対応する。\par

\noindent\textbf{2：不適切} 情報・標準規定化・倫理・対話。知識集約・高付加価値化・医療の三語が違う。\par

\noindent\textbf{3：不適切} 知識・標準規定化・医療・学習。高付加価値化と対話型の二語が違う。\par

\noindent\textbf{4：不適切} 情報・標準規定化・医療・対話。知識集約と高付加価値化の二語が違う。\par

\noindent\textbf{5：不適切} 知識・高付加価値化・倫理・学習。医療と対話型の二語が違う。\par

\paragraph{再利用する考え方} 政策資料の時点と対象箇所を固定し、一般的な連想で空欄を補わない。

\paragraph{共通理論} \hyperref[theory:ai:knowledge_society]{第\ref*{theory:ai:knowledge_society}節}

\par 参考：\cite{mext_ai_whitepaper2024}。

\par 難易度 3（校準初稿）、優先度未評価。

\section{AI利活用の十原則}

\subsection{令和3年度 問6}
\label{q:r03-reg-apt-06}

\noindent\textbf{原典：} sources/R03/aptitude.pdf、PDF物理6頁。

\textbf{公式正答：} 1（single）。

\paragraph{問題の要旨} 2019年AI利活用ガイドラインの十原則に含まれないものを数える。

\Needspace{4\baselineskip}\paragraph{確認した条件・記述（独自要約）}\mbox{}\par

\noindent\textbf{原典確認：} 問題幹と各記述の主体・否定・全称・表の5候補を保存原典画像で確認。\par

\noindent\textbf{制度時点：} R03年度の原典表現を保持。現行資料の条番号や仕様を過去へ補わない。\par

\noindent\textbf{各記述の判断：} ア：適正利用；イ：適正学習；ウ：連携；エ：安全；オ：セキュリティ；カ：プライバシー；キ：尊厳・自律；ク：公平性；ケ：透明性；コ：アカウンタビリティ\par

\noindent\textbf{前提：} 利用者にはサービス提供者・業務利用者・消費者が含まれる。適正利用、適正学習、連携、安全、セキュリティ、プライバシー、尊厳自律、公平性、透明性、説明責任の十原則を整理する。\par

\paragraph{問われる核心} 2019年AI利活用ガイドラインの十原則に含まれないものを数える。

\paragraph{必要知識} 利用者にはサービス提供者・業務利用者・消費者が含まれる。適正利用、適正学習、連携、安全、セキュリティ、プライバシー、尊厳自律、公平性、透明性、説明責任の十原則を整理する。

\paragraph{考え方・正答を選ぶ理由（査読待ち初稿）} 十記述は2019年ガイドライン9--10頁の十原則に全て対応する。含まれないもの0個、1である。ただしガイドライン5--6頁は利用の目的・文脈に応じた原則の選定を説明しており、全ての利用者へ全措置を同じ形で要求する法律ではない。

\Needspace{4\baselineskip}\paragraph{選択肢の検討（査読待ち初稿）}\mbox{}\par

\noindent\textbf{1：適切} 0。含まれないなものは0個であり、一致する。\par

\noindent\textbf{2：不適切} 1。含まれないなものは0個であり、この個数とは一致しない。\par

\noindent\textbf{3：不適切} 2。含まれないなものは0個であり、この個数とは一致しない。\par

\noindent\textbf{4：不適切} 3。含まれないなものは0個であり、この個数とは一致しない。\par

\noindent\textbf{5：不適切} 4。含まれないなものは0個であり、この個数とは一致しない。\par

\paragraph{再利用する考え方} 原則名の列挙と、各主体へ必要な具体策の選定を区別する。

\paragraph{共通理論} \hyperref[theory:ai:principles]{第\ref*{theory:ai:principles}節}

\par 参考：\cite{ai_use2019}。

\paragraph{関連問題} \hyperref[q:r01-ret-apt-15]{R01-RET-APT-15}

\par 難易度 3（校準初稿）、優先度未評価。

\section{Society5.0と人間中心の目的}

\subsection{令和4年度 問4}
\label{q:r04-reg-apt-04}

\noindent\textbf{原典：} sources/R04/aptitude.pdf、PDF物理3頁。

\textbf{公式正答：} 1（single）。

\paragraph{問題の要旨} Society5.0の歴史的位置と人間中心の目的。

\Needspace{4\baselineskip}\paragraph{確認した条件・記述（独自要約）}\mbox{}\par

\noindent\textbf{原典：} 問題文・否定・全五候補と組合せ表を保存原典36画像で確認。\par

\noindent\textbf{時点：} 令和4年度の本文表現を保持。後の制度改正を遡及しない。\par

\paragraph{問われる核心} Society5.0の歴史的位置と人間中心の目的。

\paragraph{必要知識} 社会の五段階、サイバー空間と現実空間、人間中心、経済と社会的課題、持続可能性。

\paragraph{考え方・正答を選ぶ理由（査読待ち初稿）} 狩猟社会1.0、農耕社会2.0、工業社会3.0、情報社会4.0に続く5.0である。技術を融合する目的は経済発展と社会的課題の解決の両立で、人間中心を核に持続可能な社会を目指す。五語の組合せを選び、2021計画の背景を後年の新計画へ置き換えない。

\Needspace{4\baselineskip}\paragraph{選択肢の検討（査読待ち初稿）}\mbox{}\par

\noindent\textbf{1：条件内で適切} 狩猟・農耕・社会・人間・持続可能。各文又は定義の判定を先に行う。全項目が一致する。\par

\noindent\textbf{2：不適切} 農耕・狩猟・社会・人間・持続可能。各文又は定義の判定を先に行う。並びが判定 狩猟・農耕・社会・人間・持続可能 と異なる。\par

\noindent\textbf{3：不適切} 狩猟・農耕・社会・人間・即応。各文又は定義の判定を先に行う。並びが判定 狩猟・農耕・社会・人間・持続可能 と異なる。\par

\noindent\textbf{4：不適切} 農耕・狩猟・技術・自然・即応。各文又は定義の判定を先に行う。並びが判定 狩猟・農耕・社会・人間・持続可能 と異なる。\par

\noindent\textbf{5：不適切} 狩猟・農耕・技術・自然・即応。各文又は定義の判定を先に行う。並びが判定 狩猟・農耕・社会・人間・持続可能 と異なる。\par

\paragraph{再利用する考え方} 社会の五段階、サイバー空間と現実空間、人間中心、経済と社会的課題、持続可能性。 設問の否定・全称・基準年を確かめ、根拠と評価範囲を分けて判断する。

\paragraph{共通理論} \hyperref[theory:society:five]{第\ref*{theory:society:five}節}

\par 参考：\cite{eth54_society2021}。

\paragraph{関連問題} \hyperref[q:r07-reg-apt-09]{R07-REG-APT-09}

\par 難易度 2（校準初稿）、優先度未評価。

\subsection{令和7年度 問9}
\label{q:r07-reg-apt-09}

\noindent\textbf{原典：} sources/R07/aptitude.pdf、PDF物理9頁。

\textbf{公式正答：} 4（single）。

\paragraph{問題の要旨} Society5.0の歴史的位置と人間中心の目的。

\Needspace{4\baselineskip}\paragraph{確認した条件・記述（独自要約）}\mbox{}\par

\noindent\textbf{原典：} 問題文・否定・全五候補と組合せ表を保存原典36画像で確認。\par

\noindent\textbf{時点：} 令和7年度の本文表現を保持。後の制度改正を遡及しない。\par

\paragraph{問われる核心} Society5.0の歴史的位置と人間中心の目的。

\paragraph{必要知識} 社会の五段階、サイバー空間と現実空間、人間中心、経済と社会的課題、持続可能性。

\paragraph{考え方・正答を選ぶ理由（査読待ち初稿）} 狩猟社会1.0、農耕社会2.0、工業社会3.0、情報社会4.0に続く5.0である。技術を融合する目的は経済発展と社会的課題の解決の両立で、人間中心を核に持続可能な社会を目指す。五語の組合せを選び、2021計画の背景を後年の新計画へ置き換えない。

\Needspace{4\baselineskip}\paragraph{選択肢の検討（査読待ち初稿）}\mbox{}\par

\noindent\textbf{1：不適切} 狩猟・農耕・社会・人間・即応。各文又は定義の判定を先に行う。並びが判定 狩猟・農耕・社会・人間・持続可能 と異なる。\par

\noindent\textbf{2：不適切} 農耕・狩猟・技術・自然・持続可能。各文又は定義の判定を先に行う。並びが判定 狩猟・農耕・社会・人間・持続可能 と異なる。\par

\noindent\textbf{3：不適切} 農耕・狩猟・社会・自然・即応。各文又は定義の判定を先に行う。並びが判定 狩猟・農耕・社会・人間・持続可能 と異なる。\par

\noindent\textbf{4：条件内で適切} 狩猟・農耕・社会・人間・持続可能。各文又は定義の判定を先に行う。全項目が一致する。\par

\noindent\textbf{5：不適切} 狩猟・農耕・技術・人間・持続可能。各文又は定義の判定を先に行う。並びが判定 狩猟・農耕・社会・人間・持続可能 と異なる。\par

\paragraph{再利用する考え方} 社会の五段階、サイバー空間と現実空間、人間中心、経済と社会的課題、持続可能性。 設問の否定・全称・基準年を確かめ、根拠と評価範囲を分けて判断する。

\paragraph{共通理論} \hyperref[theory:society:five]{第\ref*{theory:society:five}節}

\par 参考：\cite{eth54_society2021}。

\paragraph{関連問題} \hyperref[q:r04-reg-apt-04]{R04-REG-APT-04}

\par 難易度 2（校準初稿）、優先度未評価。


\chapter{力学}
本章では運動を決める力、保存量、物質の変形を結び付ける。物理学で既習の運動方程式を再導入するより、試験で落としやすい条件の区別を明確にする。力は加速度を決め、エネルギーは始状態と終状態を結び、応力は物体内部に力がどう分配されるかを表す。速度と加速度、静止摩擦と動摩擦、浮力と合力、半径と高度を区別すれば、同じ式を異なる年度にも適用できる。

\section{力と加速度：斜面・摩擦・見掛けの重さ}
\label{theory:mechanics:newton}
慣性系では $m\ddot{\boldsymbol r}=\sum\boldsymbol F$ である。合力がゼロなら速度は一定であり、必ずしも静止していない。鉛直投射の最高点では速度が一瞬ゼロになるが、重力による加速度はゼロにならない。この区別が\hyperref[q:r04-reg-applied-04]{令和4年度問4}の核心である。

斜面の上向きを正とし、引く力 $F$ が斜面に平行なら $N=mg\cos\theta$ である。上向きに滑り始める限界では摩擦が下向きに働くため
\[
 F=mg\sin\theta+\mu_s mg\cos\theta.
\]
力を加えずに下向きに滑り始める角度 $\theta_c$ が分かれば、$mg\sin\theta_c=\mu_s mg\cos\theta_c$ より $\mu_s=\tan\theta_c$ と決まる。静止摩擦は常に $\mu_sN$ なのではなく、その値以下でつり合いに必要な大きさを取る。滑り続ける物体では動摩擦係数 $\mu_k$ を使う。平成23年度問1は後者、平成28・令和3・令和7年度問1は前者である。

体重計は質量を直接測らず、足から受ける力を標準の重力加速度で換算している。上向き加速度 $a$ のエレベーターなら $N=m(g+a)$、一定速度なら $a=0$ なので読みは変わらない。自由落下で読みがゼロになっても質量は消失しない。質量は力と加速度の関係や運動量からも測定できる。

\section{仕事とエネルギー：途中の経路を消去する}
\label{theory:mechanics:energy}
運動方程式と変位の内積を取ると $d(mv^2/2)=\boldsymbol F\cdot d\boldsymbol r$ となる。重力と理想ばねだけが仕事をする場合、
\[
 E=\frac12 mv^2+mgh+\frac12 kx^2
\]
が保存される。抵抗・摩擦による仕事がある場合には、この三項だけの保存を仮定しない。

高さ $h$ から初速の大きさ $v_0$ で投射すると、着地時の速さは $v^2=v_0^2+2gh$ である。投射方向は飛行時間や着地点を変えるが、同じ始終高度なら速さを変えない。令和元年度再試験問2はこの不変量を問う。平成24年度問1は $v_0=0$、令和4年度問1は水平ばねのエネルギー $kx^2/2$ を運動エネルギーへ変換する場合である。数値を代入する前に質量が消えるかを確かめれば、余計な計算を避けられる。

\section{弾性：材料定数と形状による剛性を分ける}
\label{theory:mechanics:elasticity}
応力 $\sigma=F/A$ は単位面積当たりの内力、ひずみ $\epsilon=\Delta L/L$ は相対変形である。小変形・線形弾性の一軸引張では $\sigma=E\epsilon$ が成り立つ。これがフックの法則の材料表示であり、ヤング率 $E$ の単位は Pa である。一様棒のばね定数は $k=EA/L$ なので、形状で変わる剛性 $k$ と材料の $E$ を取り違えない。

横ひずみ $\epsilon_\perp$ と軸ひずみ $\epsilon_\parallel$ からポアソン比を $\nu=-\epsilon_\perp/\epsilon_\parallel$ と定義する。通常の正の $\nu$ を持つ材料は引張方向に伸び、横方向に縮む。一次の体積ひずみは三方向のひずみの和なので
\[
 \frac{\Delta V}{V}\simeq\epsilon_x+\epsilon_y+\epsilon_z=(1-2\nu)\epsilon_\parallel.
\]
したがって令和5年度問1で与えられた $0.1\leq\nu\leq0.45$ の材料は引張で体積が増加する。$\nu=1/2$ は非圧縮の極限である。すべての材料がこの数値範囲に入ると一般化しない。

せん断では形状の角度変化 $\gamma$ に対し $\tau=G\gamma$ であり、$G$ がせん断弾性率である。等方線形弾性では $E=2G(1+\nu)$ となる。名称の空欄を埋める平成23年度問4・平成25年度問6・平成30年度問2も、軸変形と角度変形を区別する問題である。これらの構成関係はMITの材料力学講義にも整理されている\cite{mit_torsion}。

\section{丸棒のねじり：最大せん断応力は表面に生じる}
\label{theory:mechanics:torsion}
一様な丸棒の線形弾性ねじりでは、中心から距離 $r$ のせん断ひずみは $r\,d\phi/dz$、応力は $\tau=Gr\,d\phi/dz$ である。断面全体のトルクを積分すると
\[
 T=\int_A\tau r\,dA=G\frac{d\phi}{dz}J,
 \qquad J=\int_A r^2\,dA=\frac{\pi d^4}{32}.
\]
よって $\tau=Tr/J$、表面 $r=d/2$ では $\tau_{\max}=16T/(\pi d^3)$ となる。$J$ は断面二次極モーメントで、単位は $\mathrm{m^4}$ である。令和元年度問2では mm を m に直してから三乗する。直径を半径として代入すると8倍の誤差になる。この導出と円形断面・線形弾性という適用条件はMIT講義のねじり節を参照されたい\cite{mit_torsion}。非円形断面や塑性化した棒へ同じ分布を無条件に適用しない。

\section{静水圧と浮力：置き換えた流体の重さ}
\label{theory:mechanics:hydrostatics}
静止流体の鉛直上向き座標を $z$ とすると、薄い流体要素の力のつり合いから $dp/dz=-\rho_f g$ となる。物体表面の圧力を積分した上向き合力は、同じ体積を占める流体の重さと等しい。密度一定なら
\[
 F_B=\rho_f gV_{\mathrm{displaced}}.
\]
完全に水没した物体では排除体積が物体全体の体積であり、浮力に物体密度は入らない。物体密度 $\rho_b$ が入るのは重力との合力 $(\rho_f-\rho_b)Vg$ や、浮いているときの水没割合 $V_{\mathrm{displaced}}/V=\rho_b/\rho_f$ である。平成26年度問3と令和元年度問1は同じ式を異なる番号の選択肢で問うため、番号ではなく式を記憶する。

\section{流れのエネルギーとピトー管}
\label{theory:mechanics:bernoulli}
定常・非粘性・非圧縮の流れで、同じ流線に沿ってオイラー方程式を積分すると
\[
 p+\frac12\rho v^2+\rho gz=\mathrm{const.}
\]
となる。各項は単位体積当たりのエネルギーであり、圧力の単位を持つ。水平流では静圧 $p$ と動圧 $q=\rho v^2/2$ の和が全圧 $p_0$ である。流れを正面の孔で停止させるピトー管は全圧を、流れに沿う側孔は静圧を測り、差から $v=\sqrt{2(p_0-p)/\rho}$ を求める。この測定配置はMITの実験講義資料で確認できる\cite{mit_pitot}。

平成24年度問9・令和2年度問2では「全圧と静圧を測定し、その差が動圧」という関係を問う。粘性抵抗を扱うストークスの法則と混同しない。大きな圧縮性、圧力損失、異なる流線間では、この簡単な式の条件を再点検する必要がある。

\section{円軌道：第一宇宙速度と静止衛星}
\label{theory:mechanics:orbit}
万有引力が円運動の向心力になるので $GMm/r^2=mv^2/r$、すなわち $v=\sqrt{GM/r}$ である。地表近くで $g=GM/R^2$ とおけば第一宇宙速度は $\sqrt{gR}$ になる。衛星質量が消える点が平成26年度問1・平成28年度問2の識別点である。

角速度 $2\pi/T$ を使うと軌道半径は
\[
 r=\left(\frac{GMT^2}{4\pi^2}\right)^{1/3}.
\]
静止衛星は地球の自転と同じ角速度で赤道上を順行する。概算では $T\simeq24$ 時間を秒に換算する。地心距離 $r$ は約 $4.2\times10^7$ m、地表からの高度は $r-R$ であり別の量である。平成30年度問1は m、令和6年度問1は km で選択肢を示す。厳密には恒星日を使うが、提示された桁の選択には概算との差は影響しない。

\section{衝突と反発：保存則の適用区間を選ぶ}
\label{theory:mechanics:impact}
短い衝突時間では外力の力積を無視できる条件の下で、衝突する二物体の運動量が保存される。反発係数は離れる相対速度と近づく相対速度の比であり、一直線上なら $v_B-v_A=e(u_A-u_B)$ である。$e=1$ で運動エネルギーも保存されるが、$e<1$ では減少する。

質量 $m$、速度 $u$ の物体が静止する質量 $M$ の物体へ完全弾性衝突すると
\[
 v_A=\frac{m-M}{m+M}u,\qquad v_B=\frac{2m}{m+M}u.
\]
令和3年度問2では求める物体がBであることを確認する。地面で跳ね返る球なら、落下区間では $v^2=2gh$、衝突で速さが $e$ 倍、上昇区間で高さが $e^2$ 倍となる。2回目の反発高さは $e^4h_0$ であり、令和2年度問1の「2回」を1回の式へ代入しない。

\section{剛体のつり合い：支点を選んで未知力を減らす}
\label{theory:mechanics:equilibrium}
剛体の静止条件は合力と合モーメントがともにゼロである。床に立てた一様棒の長さを $L$、床となす角を $\theta$ とする。壁が滑らかなら壁の反力 $N_w$ は水平、床では鉛直反力 $N_f$ と摩擦力 $F$ が働く。床との接点を支点に選ぶと $N_f,F$ のモーメントはゼロになり、
\[
 N_wL\sin\theta=Mg\frac L2\cos\theta,
 \qquad F=N_w=\frac{Mg}{2\tan\theta},\qquad N_f=Mg
\]
と決まる。令和6年度問2では $L$ が消えることを次元解析でも確認できる。ただしこの静止状態の実現には $F\leq\mu_sN_f$ が必要である。必要な摩擦力を求める式と、摩擦限界で倒れる条件は別である。

\section{振動：復元力の傾きが周期を決める}
\label{theory:mechanics:oscillation}
平衡点からの変位 $x$ に対して復元力が $-kx$ なら $m\ddot x+kx=0$、$\omega=\sqrt{k/m}$、$T=2\pi\sqrt{m/k}$ である。鉛直ばねは静的伸び $\delta=mg/k$ を用いると $T=2\pi\sqrt{\delta/g}$ となる。重力は平衡点を移すが、線形ばねの復元力の傾きを変えない。平成27年度問1・令和7年度問2はこの形で整理できる。

単振り子は $mL^2\ddot\theta=-mgL\sin\theta$ であり、小振幅で $\sin\theta\simeq\theta$ と近似すれば $T=2\pi\sqrt{L/g}$ となる。質量は消える。平成29年度問1・令和5年度問2の短い数値計算はこの近似を使う。大振幅では周期が振幅に依存するため、任意振幅への一般化はしない。基礎科目の\hyperref[q:h24-reg-basic-3-2]{平成24年度III-2}および\hyperref[q:h26-reg-basic-3-5]{平成26年度III-5}とも同じ線形振動の構造を共有する。

\section*{原典確認・独立解法を保存した問題}
以下の解説は査読待ちの初稿である。章の共通理論へ戻り、年度ごとの条件と各選択肢を検討する。正答番号の一致のみを最終検証とはしない。
\subsection{平成23年度 問1}
\label{q:h23-reg-applied-01}

\noindent\textbf{原典：} sources/H23/17\_applied\_science.pdf、PDF物理1頁。

\textbf{公式正答：} 3（single）。

\par\textit{原典本文の校正と独自解説は未作成。分類のみの登録項目。}

\subsection{平成23年度 問4}
\label{q:h23-reg-applied-04}

\noindent\textbf{原典：} sources/H23/17\_applied\_science.pdf、PDF物理2頁。

\textbf{公式正答：} 1（single）。

\par\textit{原典本文の校正と独自解説は未作成。分類のみの登録項目。}

\subsection{平成24年度 問1}
\label{q:h24-reg-applied-01}

\noindent\textbf{原典：} sources/H24/17\_applied\_science.pdf、PDF物理1頁。

\textbf{公式正答：} 3（single）。

\par\textit{原典本文の校正と独自解説は未作成。分類のみの登録項目。}

\subsection{平成24年度 問9}
\label{q:h24-reg-applied-09}

\noindent\textbf{原典：} sources/H24/17\_applied\_science.pdf、PDF物理3頁。

\textbf{公式正答：} 1（single）。

\par\textit{原典本文の校正と独自解説は未作成。分類のみの登録項目。}

\subsection{平成25年度 問1}
\label{q:h25-reg-applied-01}

\noindent\textbf{原典：} sources/H25/17\_applied\_science.pdf、PDF物理1頁。

\textbf{公式正答：} 2（single）。

\par\textit{原典本文の校正と独自解説は未作成。分類のみの登録項目。}

\subsection{平成25年度 問6}
\label{q:h25-reg-applied-06}

\noindent\textbf{原典：} sources/H25/17\_applied\_science.pdf、PDF物理3頁。

\textbf{公式正答：} 1（single）。

\par\textit{原典本文の校正と独自解説は未作成。分類のみの登録項目。}

\subsection{平成26年度 問1}
\label{q:h26-reg-applied-01}

\noindent\textbf{原典：} sources/H26/17\_applied\_science.pdf、PDF物理1頁。

\textbf{公式正答：} 1（single）。

\par\textit{原典本文の校正と独自解説は未作成。分類のみの登録項目。}

\subsection{平成26年度 問3}
\label{q:h26-reg-applied-03}

\noindent\textbf{原典：} sources/H26/17\_applied\_science.pdf、PDF物理1頁。

\textbf{公式正答：} 2（single）。

\par\textit{原典本文の校正と独自解説は未作成。分類のみの登録項目。}

\subsection{平成28年度 問1}
\label{q:h28-reg-applied-01}

\noindent\textbf{原典：} sources/H28/17\_applied\_science.pdf、PDF物理1頁。

\textbf{公式正答：} 5（single）。

\par\textit{原典本文の校正と独自解説は未作成。分類のみの登録項目。}

\subsection{平成28年度 問2}
\label{q:h28-reg-applied-02}

\noindent\textbf{原典：} sources/H28/17\_applied\_science.pdf、PDF物理1頁。

\textbf{公式正答：} 1（single）。

\par\textit{原典本文の校正と独自解説は未作成。分類のみの登録項目。}

\subsection{平成30年度 問1}
\label{q:h30-reg-applied-01}

\noindent\textbf{原典：} sources/H30/17\_applied\_science.pdf、PDF物理1頁。

\textbf{公式正答：} 4（single）。

\par\textit{原典本文の校正と独自解説は未作成。分類のみの登録項目。}

\subsection{平成30年度 問2}
\label{q:h30-reg-applied-02}

\noindent\textbf{原典：} sources/H30/17\_applied\_science.pdf、PDF物理1頁。

\textbf{公式正答：} 2（single）。

\par\textit{原典本文の校正と独自解説は未作成。分類のみの登録項目。}

\subsection{令和元年度 問1}
\label{q:r01-reg-applied-01}

\noindent\textbf{原典：} sources/R01/17\_applied\_science.pdf、PDF物理1頁。

\textbf{公式正答：} 5（single）。

\par\textit{原典本文の校正と独自解説は未作成。分類のみの登録項目。}

\subsection{令和元年度 問2}
\label{q:r01-reg-applied-02}

\noindent\textbf{原典：} sources/R01/17\_applied\_science.pdf、PDF物理1頁。

\textbf{公式正答：} 2（single）。

\par\textit{原典本文の校正と独自解説は未作成。分類のみの登録項目。}

\subsection{令和元年度再試験 問2}
\label{q:r01-ret-applied-02}

\noindent\textbf{原典：} sources/R01\_retest/17\_applied\_science.pdf、PDF物理1頁。

\textbf{公式正答：} 4（single）。

\par\textit{原典本文の校正と独自解説は未作成。分類のみの登録項目。}

\subsection{令和2年度 問1}
\label{q:r02-reg-applied-01}

\noindent\textbf{原典：} sources/R02/17\_applied\_science.pdf、PDF物理1頁。

\textbf{公式正答：} 2（single）。

\par\textit{原典本文の校正と独自解説は未作成。分類のみの登録項目。}

\subsection{令和2年度 問2}
\label{q:r02-reg-applied-02}

\noindent\textbf{原典：} sources/R02/17\_applied\_science.pdf、PDF物理1頁。

\textbf{公式正答：} 1（single）。

\par\textit{原典本文の校正と独自解説は未作成。分類のみの登録項目。}

\subsection{令和3年度 問1}
\label{q:r03-reg-applied-01}

\noindent\textbf{原典：} sources/R03/17\_applied\_science.pdf、PDF物理1頁。

\textbf{公式正答：} 4（single）。

\par\textit{原典本文の校正と独自解説は未作成。分類のみの登録項目。}

\subsection{令和3年度 問2}
\label{q:r03-reg-applied-02}

\noindent\textbf{原典：} sources/R03/17\_applied\_science.pdf、PDF物理1頁。

\textbf{公式正答：} 5（single）。

\par\textit{原典本文の校正と独自解説は未作成。分類のみの登録項目。}

\subsection{令和4年度 問1}
\label{q:r04-reg-applied-01}

\noindent\textbf{原典：} sources/R04/17\_applied\_science.pdf、PDF物理1頁。

\textbf{公式正答：} 3（single）。

\par\textit{原典本文の校正と独自解説は未作成。分類のみの登録項目。}

\subsection{令和4年度 問2}
\label{q:r04-reg-applied-02}

\noindent\textbf{原典：} sources/R04/17\_applied\_science.pdf、PDF物理1頁。

\textbf{公式正答：} 5（single）。

\par\textit{原典本文の校正と独自解説は未作成。分類のみの登録項目。}

\subsection{令和5年度 問1}
\label{q:r05-reg-applied-01}

\noindent\textbf{原典：} sources/R05/17\_applied\_science.pdf、PDF物理1頁。

\textbf{公式正答：} 5（single）。

\par\textit{原典本文の校正と独自解説は未作成。分類のみの登録項目。}

\subsection{令和5年度 問2}
\label{q:r05-reg-applied-02}

\noindent\textbf{原典：} sources/R05/17\_applied\_science.pdf、PDF物理1頁。

\textbf{公式正答：} 1（single）。

\par\textit{原典本文の校正と独自解説は未作成。分類のみの登録項目。}

\subsection{令和6年度 問1}
\label{q:r06-reg-applied-01}

\noindent\textbf{原典：} sources/R06/17\_applied\_science.pdf、PDF物理1頁。

\textbf{公式正答：} 3（single）。

\par\textit{原典本文の校正と独自解説は未作成。分類のみの登録項目。}

\subsection{令和6年度 問2}
\label{q:r06-reg-applied-02}

\noindent\textbf{原典：} sources/R06/17\_applied\_science.pdf、PDF物理1頁。

\textbf{公式正答：} 5（single）。

\par\textit{原典本文の校正と独自解説は未作成。分類のみの登録項目。}

\subsection{令和7年度 問1}
\label{q:r07-reg-applied-01}

\noindent\textbf{原典：} sources/R07/17\_applied\_science.pdf、PDF物理1頁。

\textbf{公式正答：} 3（single）。

\par\textit{原典本文の校正と独自解説は未作成。分類のみの登録項目。}


\chapter{振動・波動}
\section{局所的な振動から伝播へ}
振動は自由度が時間的に変わる運動であり、波動は隣接する自由度の結合を通して擾乱が空間へ伝わる現象である。
ばねの単振動と小角の振り子は第\ref{theory:mechanics:oscillation}節で扱った。
ここでは、変位の運動と位相の伝播を区別する。波が右へ進んでも媒質の各粒子が右へ運ばれ続けるとは限らない。
地震探査、音響測定、信号伝送へ共通するのは、媒質の復元力が速さを決め、幾何と境界が位相差を決めることである。

\section{位相を空間と時間で読む}\label{theory:wave:phase}
正弦波を \(y=A\sin(\omega t-kx+\phi_0)\) とする。
\(\omega\) は時間に対する位相の変化率、\(k\) は空間に対する位相の変化率であり、単位はそれぞれ rad/s、rad/m である。
時間又は空間で位相が \(2\pi\) 変化する間隔を定めれば
\[
 T=\frac{2\pi}{\omega},\qquad \lambda=\frac{2\pi}{k},\qquad f=\frac1T
\]
となる。同じ山を追うには \(\omega\,dt-k\,dx=0\) を満たすので、位相速度は
\[
 v_{\mathrm{ph}}=\frac{dx}{dt}=\frac{\omega}{k}=\frac{\lambda}{T}
\]
である。\(kx+\omega t\) と書けば進行方向は逆になる。
この結果は位相の運動を表し、分散性の媒質では波束の群速度 \(d\omega/dk\) と同じとは限らない。
試験の単一正弦波へ、波束を明示しないまま群速度を追加する必要はない。
周期で表した式の時間係数が \(2\pi/T\) であることと、位相速度に余分な \(2\pi\) が入らないことを一緒に確認する。

\section{復元力・密度と縦波・横波}\label{theory:wave:elastic}
等方線形弾性体の変位 \(\boldsymbol u\) に対する運動方程式を、密度 \(\rho\)、体積弾性率 \(K\)、剛性率 \(\mu\) で書くと
\[
 \rho\,\partial_t^2\boldsymbol u
 =(K+\mu/3)\nabla(\nabla\cdot\boldsymbol u)+\mu\nabla^2\boldsymbol u
\]
である。波数方向に平行な変位と垂直な変位へ分解することで
\[
 c_P=\sqrt{\frac{K+4\mu/3}{\rho}},\qquad
 c_S=\sqrt{\frac{\mu}{\rho}}
\]
を得る。縦波には体積変形とせん断変形の双方の復元力が入り、横波にはせん断の復元力が入る。
安定な等方モデルでは \(K,\mu>0\) であり、\(c_P>c_S\) となる。
これは異方性結晶の任意の方向へ無条件に適用する式ではない。
通常の流体は静的せん断に抗する \(\mu\) を持たず、体積内の音は圧縮の縦波である。
真空では音を伝える媒質がない。液体表面の重力波などを、流体の体積音波と同じ分類にしない。

気体では音波の小振幅の圧縮・膨張をほぼ断熱とみなして \(K_s=\gamma p\) とする。
理想気体式 \(p=\rho RT/M\) より
\[
 c=\sqrt{\frac{\gamma p}{\rho}}=\sqrt{\frac{\gamma RT}{M}}
\]
となる。\(T\) は絶対温度、\(R\) は気体定数、\(M\) はモル質量、\(\gamma=C_p/C_v\) である。
気体組成などを同じとして温度を増すと速くなる。
水は気体より密度も大きいが、体積弾性率の増大が支配的で、海水の音速は通常空気より大きい。
代表値は空気約340 m/s、海水約1500 m/sであり、精密測定では温度・圧力・塩分を指定する\cite{os_sound}。

\section{音源の移動と観測者の移動}\label{theory:wave:doppler}
音速 \(c\) を一定とし、速度は媒質に対して測る。
音源周期 \(T_s\) の間に波面は \(cT_s\) 進む。
音源が観測者から速さ \(u_s\) で遠ざかると、観測者側の波面間隔は \((c+u_s)T_s\) へ広がる。
静止観測者に対し
\[
 f_{\mathrm{obs}}=\frac{c}{(c+u_s)T_s}=f_s\frac{c}{c+u_s}
\]
となる。音源が近付けば分母は \(c-u_s\) となる。
一方、静止音源へ観測者が速さ \(u_o\) で近付けば、変わるのは波長ではなく波面の通過速度であり
\(f_{\mathrm{obs}}=f_s(c+u_o)/c\) となる。
両方が動く場合は、この二つの効果を別々に組み合わせる\cite{os_doppler}。
音源と観測者の相対速度だけを一つの分母へ入れる操作は一般に正しくない。

\section{動く船と反射音の距離}\label{theory:wave:echo}
音の反射を使う距離測定では、送受信位置が同じか移動したかが重要である。
船が岸壁へ速さ \(u\) で近付くとする。発音時距離を \(D_0\)、戻るまでの時間を \(t\) とすれば、
音は行きに \(D_0\)、帰りに受信時の距離 \(D_1=D_0-ut\) を進む。
\[
 ct=D_0+D_1=2D_0-ut,\qquad
 D_0=\frac{(c+u)t}{2},\qquad D_1=\frac{(c-u)t}{2}
\]
となる。岸壁での反射時刻を途中で明示しても同じ結果となる。
この式の二つの距離を区別すれば、「汽笛を鳴らしたとき」と「反射音を聞いたとき」の取り違えを防げる。
単に \(ct/2\) とするのは \(u=0\) の場合である。流れや風による音速の方向差は、問題が与えたときに別途扱う。

\section{二つの音源の経路差と位相差}\label{theory:wave:interference}
同じ角振動数の線形波を重ねると、受信点での合成振幅の二乗は
\[
 A^2=A_1^2+A_2^2+2A_1A_2\cos\delta,\qquad
 \delta=k(r_2-r_1)+\delta_0
\]
となる。\(r_i\) は経路長、\(\delta_0\) は放射時の相対位相である。
極小は \(\delta=(2n+1)\pi\) に現れるが、\(A_1\ne A_2\) なら残る振幅は \(|A_1-A_2|\) で、必ずしも無音ではない\cite{os_superposition}。
スピーカが間隔 \(a\) で並び、正面の聴取線まで \(L\) とする。
中央で強め合い、片方の正面で弱め合う設定なら、経路差は
\[
 \Delta=\sqrt{L^2+a^2}-L=\frac{a^2}{\sqrt{L^2+a^2}+L}
\]
である。後式は桁落ちを避けるにも有用である。
同位相の放射に対して \(\Delta=(2n+1)\lambda/2\) となり、
\[
 f_n=\frac{(2n+1)c}{2\Delta},\qquad n=0,1,2,\ldots
\]
である。最初の極小の指定がない場合、選択肢の周波数範囲も照合して次数を決める。
距離 \(L\) とスピーカ間隔 \(a\) の単純な差は経路差ではない。

\section{振動の合成・衝撃波・信号変調}\label{theory:wave:applications}
垂直な単振動 \(x=A\sin(\omega_1t+\alpha)\)、\(y=B\sin(\omega_2t+\beta)\) の軌道はリサジュー図形となる。
周波数比が有理なら周期的な閉軌道となり、位相差により直線や楕円などへ変わる。
一方、衝撃波は大振幅・非線形な圧縮で密度・圧力等が急に変わる前面である。
前面速度は未擾乱側の小振幅音速を超える。線形の正弦音波と同じ一定振幅モデルをそのまま当てない。
振幅変調では、例えば \([A_0+a\cos(\Omega t)]\cos(\omega_ct)\) と書き、
搬送波の振幅包絡へ信号を入れる。変調周波数 \(\Omega\) と搬送周波数 \(\omega_c\) を分ける。
これらは「波」という名称が共通でも、軌道・伝播前面・情報の表現という違う対象を問う。

\section{典型的な解法と出題の確認範囲}
正弦波では位相係数と次元、ドップラーでは媒質基準の速度、反射音では距離の時点、
干渉では経路差と次数を先に決める。
本章の振動3問に続き、S7で波動10問を原典確認した。
同じ正弦波や音源運動が再出題されても、選択肢の配列は一致しないことがある。
頻度の全体集計は索引・統計を参照し、この確認範囲だけから将来出題を断定しない。

\section{登録問題と解説初稿}
\section{音源運動・ドップラー・反射音}

\subsection{平成23年度 問3}
\label{q:h23-reg-applied-03}

\noindent\textbf{原典：} sources/H23/17\_applied\_science.pdf、PDF物理1頁。

\textbf{公式正答：} 2（single）。

\paragraph{問題の要旨} 静止した観測者から遠ざかる音源の観測振動数を求める。

\Needspace{4\baselineskip}\paragraph{確認した条件・記述（独自要約）}\mbox{}\par

\noindent\textbf{音源：} 440 Hz、観測者から4 m/sで遠ざかる。\par

\noindent\textbf{音速：} 340 m/s。観測者は静止。\par

\noindent\textbf{方向：} 観測振動数に最も近い値。\par

\paragraph{問われる核心} 媒質基準の音源運動によるドップラー効果

\paragraph{必要知識} 媒質基準の音源運動によるドップラー効果。\(f_{\rm obs}=f_s c/(c+u)=440\times340/344\simeq434.9\) Hzで2。音源が1周期の間に遠ざかる距離を波長へ加える。観測者の移動による波面の通過速度の変化とは分ける。

\paragraph{考え方・正答を選ぶ理由（査読待ち初稿）} \(f_{\rm obs}=f_s c/(c+u)=440\times340/344\simeq434.9\) Hzで2。音源が1周期の間に遠ざかる距離を波長へ加える。観測者の移動による波面の通過速度の変化とは分ける。

\Needspace{4\baselineskip}\paragraph{選択肢の検討（査読待ち初稿）}\mbox{}\par

\noindent\textbf{1：不適切} 430 Hz。後退で下がる方向は正しいが、計算値約434.9 Hzには435の方が近い。\par

\noindent\textbf{2：最適} 435 Hz。波面間隔が広がるため約434.9 Hzとなり、この候補が最も近い。\par

\noindent\textbf{3：不適切} 440 Hz。運動による波長変化を無視した音源そのものの振動数。\par

\noindent\textbf{4：不適切} 445 Hz。遠ざかるときの変化を接近と取り違えた候補。\par

\noindent\textbf{5：不適切} 450 Hz。遠ざかる音源で観測周波数が増えることはない。\par

\paragraph{再利用する考え方} 媒質基準の音源運動によるドップラー効果を暗記だけでなく、位相・幾何・符号と成立条件から導く。選択肢番号を類似年度へそのまま移さない。

\paragraph{共通理論} \hyperref[theory:wave:doppler]{第\ref*{theory:wave:doppler}節}

\par 参考：\cite{os_doppler}。

\paragraph{関連問題} \hyperref[q:h25-reg-applied-05]{H25-REG-APPLIED-05}

\par 難易度 2（校準初稿）、優先度未評価。

\subsection{平成25年度 問5}
\label{q:h25-reg-applied-05}

\noindent\textbf{原典：} sources/H25/17\_applied\_science.pdf、PDF物理3頁。

\textbf{公式正答：} 2（single）。

\paragraph{問題の要旨} 静止した観測者から遠ざかる音源の観測振動数を求める。

\Needspace{4\baselineskip}\paragraph{確認した条件・記述（独自要約）}\mbox{}\par

\noindent\textbf{音源：} 440 Hz、観測者から4 m/sで遠ざかる。\par

\noindent\textbf{音速：} 340 m/s。観測者は静止。\par

\noindent\textbf{方向：} 観測振動数に最も近い値。\par

\paragraph{問われる核心} 媒質基準の音源運動によるドップラー効果

\paragraph{必要知識} 媒質基準の音源運動によるドップラー効果。\(f_{\rm obs}=f_s c/(c+u)=440\times340/344\simeq434.9\) Hzで2。音源が1周期の間に遠ざかる距離を波長へ加える。観測者の移動による波面の通過速度の変化とは分ける。

\paragraph{考え方・正答を選ぶ理由（査読待ち初稿）} \(f_{\rm obs}=f_s c/(c+u)=440\times340/344\simeq434.9\) Hzで2。音源が1周期の間に遠ざかる距離を波長へ加える。観測者の移動による波面の通過速度の変化とは分ける。

\Needspace{4\baselineskip}\paragraph{選択肢の検討（査読待ち初稿）}\mbox{}\par

\noindent\textbf{1：不適切} 430 Hz。後退で下がる方向は正しいが、計算値約434.9 Hzには435の方が近い。\par

\noindent\textbf{2：最適} 435 Hz。波面間隔が広がるため約434.9 Hzとなり、この候補が最も近い。\par

\noindent\textbf{3：不適切} 440 Hz。運動による波長変化を無視した音源そのものの振動数。\par

\noindent\textbf{4：不適切} 445 Hz。遠ざかるときの変化を接近と取り違えた候補。\par

\noindent\textbf{5：不適切} 450 Hz。遠ざかる音源で観測周波数が増えることはない。\par

\paragraph{再利用する考え方} 媒質基準の音源運動によるドップラー効果を暗記だけでなく、位相・幾何・符号と成立条件から導く。選択肢番号を類似年度へそのまま移さない。

\paragraph{共通理論} \hyperref[theory:wave:doppler]{第\ref*{theory:wave:doppler}節}

\par 参考：\cite{os_doppler}。

\paragraph{関連問題} \hyperref[q:h23-reg-applied-03]{H23-REG-APPLIED-03}

\par 難易度 2（校準初稿）、優先度未評価。

\subsection{平成30年度 問5}
\label{q:h30-reg-applied-05}

\noindent\textbf{原典：} sources/H30/17\_applied\_science.pdf、PDF物理2頁。

\textbf{公式正答：} 2（single）。

\paragraph{問題の要旨} 岸壁へ接近する船が汽笛の反射音を聞いた時点の岸壁までの距離を求める。

\Needspace{4\baselineskip}\paragraph{確認した条件・記述（独自要約）}\mbox{}\par

\noindent\textbf{船速：} 岸壁へ20 m/sで等速接近。\par

\noindent\textbf{音速・時間：} 340 m/s、汽笛から反射音到着まで3 s。\par

\noindent\textbf{問う時点：} 反射音を聞いた時点。\par

\paragraph{問われる核心} 移動する観測位置と反射音の往復幾何

\paragraph{必要知識} 移動する観測位置と反射音の往復幾何。発音時の距離をD0、戻るまでt=\allowbreak{}3 sとすると音の総行程は\(ct=D_0+(D_0-ut)\)。\(D_0=(c+u)t/2=540\) m、戻った時点では\(D_1=D_0-ut=480\) mとなり2。

\paragraph{考え方・正答を選ぶ理由（査読待ち初稿）} 発音時の距離をD0、戻るまでt=\allowbreak{}3 sとすると音の総行程は\(ct=D_0+(D_0-ut)\)。\(D_0=(c+u)t/2=540\) m、戻った時点では\(D_1=D_0-ut=480\) mとなり2。

\Needspace{4\baselineskip}\paragraph{選択肢の検討（査読待ち初稿）}\mbox{}\par

\noindent\textbf{1：不適切} 450 m。船の移動60 mを正しい往復の幾何から減じた480 mとは一致しない。\par

\noindent\textbf{2：最適} 480 m。発音時540 mから移動60 mを差し引いた到着時距離。\par

\noindent\textbf{3：不適切} 510 m。静止する船としてct/2を計算した値で、船の移動を無視している。\par

\noindent\textbf{4：不適切} 540 m。発音時の距離であり、反射音を聞いた時点の距離ではない。\par

\noindent\textbf{5：不適切} 570 m。移動方向が接近であることと、音の往復距離の減少を正しく反映していない。\par

\paragraph{再利用する考え方} 移動する観測位置と反射音の往復幾何を暗記だけでなく、位相・幾何・符号と成立条件から導く。選択肢番号を類似年度へそのまま移さない。

\paragraph{共通理論} \hyperref[theory:wave:echo]{第\ref*{theory:wave:echo}節}

\par 難易度 2（校準初稿）、優先度未評価。

\section{音の経路差・位相差・干渉}

\subsection{令和元年度 問5}
\label{q:r01-reg-applied-05}

\noindent\textbf{原典：} sources/R01/17\_applied\_science.pdf、PDF物理2頁。

\textbf{公式正答：} 2（single）。

\paragraph{問題の要旨} 2スピーカの正面での強め合いと、片方の正面での極小から音の振動数を求める。

\Needspace{4\baselineskip}\paragraph{確認した条件・記述（独自要約）}\mbox{}\par

\noindent\textbf{配置：} スピーカ間隔2 m、正面の聴取線まで3.3 m。\par

\noindent\textbf{観測：} 中央は強め合い、片方の正面で最小。音速340 m/s。\par

\noindent\textbf{次数：} 初回極小という明示なし。候補範囲も照合。\par

\paragraph{問われる核心} 幾何学的経路差とコヒーレントな音の位相差

\paragraph{必要知識} 幾何学的経路差とコヒーレントな音の位相差。中央では両経路が等しく強め合う。片方の正面では距離\(L\)と\(\sqrt{L^2+a^2}\)の差\(\Delta=\sqrt{L^2+a^2}-L\)がある。逆位相の条件\(\Delta=(2n+1)\lambda/2\)の最初n=\allowbreak{}0では\(f=c/(2\Delta)\)≈304.2 Hzで2。次のn=\allowbreak{}1は約912.7 Hzで全候補より大きい。原典に最初の極小とは明示されないが、候補の範囲で次数を確認できる。距離差により振幅が異なるなら極小は必ずしも無音ではない。

\paragraph{考え方・正答を選ぶ理由（査読待ち初稿）} 中央では両経路が等しく強め合う。片方の正面では距離\(L\)と\(\sqrt{L^2+a^2}\)の差\(\Delta=\sqrt{L^2+a^2}-L\)がある。逆位相の条件\(\Delta=(2n+1)\lambda/2\)の最初n=\allowbreak{}0では\(f=c/(2\Delta)\)≈304.2 Hzで2。次のn=\allowbreak{}1は約912.7 Hzで全候補より大きい。原典に最初の極小とは明示されないが、候補の範囲で次数を確認できる。距離差により振幅が異なるなら極小は必ずしも無音ではない。

\Needspace{4\baselineskip}\paragraph{選択肢の検討（査読待ち初稿）}\mbox{}\par

\noindent\textbf{1：不適切} 280 Hz。最初の逆位相による極小の振動数約304.2 Hzから、正答候補より離れている。次の奇数次数は候補範囲外。\par

\noindent\textbf{2：最適} 310 Hz。道のり差約0.5588 mが半波長となる振動数約304.2 Hzに最も近い。\par

\noindent\textbf{3：不適切} 370 Hz。最初の逆位相による極小の振動数約304.2 Hzから、正答候補より離れている。次の奇数次数は候補範囲外。\par

\noindent\textbf{4：不適切} 430 Hz。最初の逆位相による極小の振動数約304.2 Hzから、正答候補より離れている。次の奇数次数は候補範囲外。\par

\noindent\textbf{5：不適切} 490 Hz。最初の逆位相による極小の振動数約304.2 Hzから、正答候補より離れている。次の奇数次数は候補範囲外。\par

\paragraph{再利用する考え方} 幾何学的経路差とコヒーレントな音の位相差を暗記だけでなく、位相・幾何・符号と成立条件から導く。選択肢番号を類似年度へそのまま移さない。

\paragraph{共通理論} \hyperref[theory:wave:interference]{第\ref*{theory:wave:interference}節}

\par 参考：\cite{os_superposition}。

\paragraph{関連問題} \hyperref[q:r05-reg-applied-05]{R05-REG-APPLIED-05}

\par 難易度 2（校準初稿）、優先度未評価。

\subsection{令和5年度 問5}
\label{q:r05-reg-applied-05}

\noindent\textbf{原典：} sources/R05/17\_applied\_science.pdf、PDF物理2頁。

\textbf{公式正答：} 3（single）。

\paragraph{問題の要旨} 2スピーカの正面での強め合いと、片方の正面での極小から音の振動数を求める。

\Needspace{4\baselineskip}\paragraph{確認した条件・記述（独自要約）}\mbox{}\par

\noindent\textbf{配置：} スピーカ間隔2.4 m、正面の聴取線まで4.0 m。\par

\noindent\textbf{観測：} 中央は強め合い、片方の正面で最小。音速340 m/s。\par

\noindent\textbf{次数：} 初回極小という明示なし。候補範囲も照合。\par

\paragraph{問われる核心} 幾何学的経路差とコヒーレントな音の位相差

\paragraph{必要知識} 幾何学的経路差とコヒーレントな音の位相差。中央では両経路が等しく強め合う。片方の正面では距離\(L\)と\(\sqrt{L^2+a^2}\)の差\(\Delta=\sqrt{L^2+a^2}-L\)がある。逆位相の条件\(\Delta=(2n+1)\lambda/2\)の最初n=\allowbreak{}0では\(f=c/(2\Delta)\)≈255.7 Hzで3。次のn=\allowbreak{}1は約767.2 Hzで全候補より大きい。原典に最初の極小とは明示されないが、候補の範囲で次数を確認できる。距離差により振幅が異なるなら極小は必ずしも無音ではない。

\paragraph{考え方・正答を選ぶ理由（査読待ち初稿）} 中央では両経路が等しく強め合う。片方の正面では距離\(L\)と\(\sqrt{L^2+a^2}\)の差\(\Delta=\sqrt{L^2+a^2}-L\)がある。逆位相の条件\(\Delta=(2n+1)\lambda/2\)の最初n=\allowbreak{}0では\(f=c/(2\Delta)\)≈255.7 Hzで3。次のn=\allowbreak{}1は約767.2 Hzで全候補より大きい。原典に最初の極小とは明示されないが、候補の範囲で次数を確認できる。距離差により振幅が異なるなら極小は必ずしも無音ではない。

\Needspace{4\baselineskip}\paragraph{選択肢の検討（査読待ち初稿）}\mbox{}\par

\noindent\textbf{1：不適切} 160 Hz。最初の逆位相による極小の振動数約255.7 Hzから、正答候補より離れている。次の奇数次数は候補範囲外。\par

\noindent\textbf{2：不適切} 210 Hz。最初の逆位相による極小の振動数約255.7 Hzから、正答候補より離れている。次の奇数次数は候補範囲外。\par

\noindent\textbf{3：最適} 260 Hz。道のり差約0.6648 mが半波長となる振動数約255.7 Hzに最も近い。\par

\noindent\textbf{4：不適切} 320 Hz。最初の逆位相による極小の振動数約255.7 Hzから、正答候補より離れている。次の奇数次数は候補範囲外。\par

\noindent\textbf{5：不適切} 400 Hz。最初の逆位相による極小の振動数約255.7 Hzから、正答候補より離れている。次の奇数次数は候補範囲外。\par

\paragraph{再利用する考え方} 幾何学的経路差とコヒーレントな音の位相差を暗記だけでなく、位相・幾何・符号と成立条件から導く。選択肢番号を類似年度へそのまま移さない。

\paragraph{共通理論} \hyperref[theory:wave:interference]{第\ref*{theory:wave:interference}節}

\par 参考：\cite{os_superposition}。

\paragraph{関連問題} \hyperref[q:r01-reg-applied-05]{R01-REG-APPLIED-05}

\par 難易度 2（校準初稿）、優先度未評価。

\section{音波・弾性波・媒質の復元力}

\subsection{平成24年度 問2}
\label{q:h24-reg-applied-02}

\noindent\textbf{原典：} sources/H24/17\_applied\_science.pdf、PDF物理1頁。

\textbf{公式正答：} 4（single）。

\paragraph{問題の要旨} 音の伝播、温度依存、流体の振動方向と固体の波速を判定する。

\Needspace{4\baselineskip}\paragraph{確認した条件・記述（独自要約）}\mbox{}\par

\noindent\textbf{方向：} 最も不適切。\par

\noindent\textbf{範囲：} 通常の音波。固体の比較は一般的な等方弾性体。\par

\paragraph{問われる核心} 圧縮波・せん断波の復元力と波速

\paragraph{必要知識} 圧縮波・せん断波の復元力と波速。媒質の体積弾性と密度が音速を決める。海水中の代表的な値は空気中より大きく、4が不適切。他の肢は通常の流体音波又は等方弾性体の一般的な説明として適切である。

\paragraph{考え方・正答を選ぶ理由（査読待ち初稿）} 媒質の体積弾性と密度が音速を決める。海水中の代表的な値は空気中より大きく、4が不適切。他の肢は通常の流体音波又は等方弾性体の一般的な説明として適切である。

\Needspace{4\baselineskip}\paragraph{選択肢の検討（査読待ち初稿）}\mbox{}\par

\noindent\textbf{1：適切} 真空中では音は伝わらない。。音は媒質の圧縮・変位の伝播であり、真空にはその媒質がない。\par

\noindent\textbf{2：適切} 空気の温度が高いほど音速は大きい。。理想気体の\(c=\sqrt{\gamma RT/M}\)では絶対温度の平方根に比例する。組成等を同じとして比較する。\par

\noindent\textbf{3：適切} 流体中を伝わる音は縦波である。。静的せん断弾性を持たない流体では、通常の体積音波は圧縮の縦波となる。表面波とは区別。\par

\noindent\textbf{4：不適切} 海水中の音速は空気中より小さい。。海水では概ね1500 m/s、空気では概ね340 m/sであり比較方向が逆。状態により変わるがこの大小は通常逆転しない。\par

\noindent\textbf{5：適切（一般的モデル）} 一般に固体中の縦波は横波より速い。。安定な等方線形弾性体で\(c_P^2-c_S^2=(K+\mu/3)/\rho>0\)である。異方性媒質の全方向をこの式だけで扱わない。\par

\paragraph{再利用する考え方} 圧縮波・せん断波の復元力と波速を暗記だけでなく、位相・幾何・符号と成立条件から導く。選択肢番号を類似年度へそのまま移さない。

\paragraph{共通理論} \hyperref[theory:wave:elastic]{第\ref*{theory:wave:elastic}節}

\par 参考：\cite{os_sound}。

\paragraph{関連問題} \hyperref[q:r04-reg-applied-05]{R04-REG-APPLIED-05}

\par 難易度 2（校準初稿）、優先度未評価。

\subsection{令和4年度 問5}
\label{q:r04-reg-applied-05}

\noindent\textbf{原典：} sources/R04/17\_applied\_science.pdf、PDF物理2頁。

\textbf{公式正答：} 2（single）。

\paragraph{問題の要旨} リサジュー図形、波の振動方向、衝撃波、縦横波の速さ、AMを判定する。

\Needspace{4\baselineskip}\paragraph{確認した条件・記述（独自要約）}\mbox{}\par

\noindent\textbf{方向：} 最も不適切。\par

\noindent\textbf{範囲：} 振動・波動の一般的説明。\par

\paragraph{問われる核心} 波の偏り、非線形な圧縮波、変調と単振動合成

\paragraph{必要知識} 波の偏り、非線形な圧縮波、変調と単振動合成。縦波・横波は速さではなく、変位方向と伝播方向の関係で定義する。同方向を横波とした2が不適切。固体の波速の一般的な比較は等方弾性体モデルで確認する。

\paragraph{考え方・正答を選ぶ理由（査読待ち初稿）} 縦波・横波は速さではなく、変位方向と伝播方向の関係で定義する。同方向を横波とした2が不適切。固体の波速の一般的な比較は等方弾性体モデルで確認する。

\Needspace{4\baselineskip}\paragraph{選択肢の検討（査読待ち初稿）}\mbox{}\par

\noindent\textbf{1：適切} 垂直な単振動の合成軌道は一般にリサジュー図形という。。xとyの単振動の周波数比・位相差で軌道が変わる。比が有理なら閉曲線となる。\par

\noindent\textbf{2：不適切} 伝播方向と媒質の振動方向が一致する波を横波という。。一致する波は縦波である。横波は伝播に垂直な振動を持つ。\par

\noindent\textbf{3：適切} 衝撃波は媒質中を超音速で伝播する圧縮波である。。不連続な圧力等の変化を伴う圧縮の前面が、未擾乱側の小振幅音速を超えて進む。小振幅正弦波とは分ける。\par

\noindent\textbf{4：適切（一般的モデル）} 一般に固体中の縦波は横波より速い。。安定な等方線形弾性体では縦波の復元力はK+4µ/3、横波はµなので、密度が同じなら縦波が速い。\par

\noindent\textbf{5：適切} AMは信号に応じて搬送波の振幅を変える方式である。。振幅包絡へ信号を入れる方式で、周波数を変えるFMと異なる。\par

\paragraph{再利用する考え方} 波の偏り、非線形な圧縮波、変調と単振動合成を暗記だけでなく、位相・幾何・符号と成立条件から導く。選択肢番号を類似年度へそのまま移さない。

\paragraph{共通理論} \hyperref[theory:wave:elastic]{第\ref*{theory:wave:elastic}節}, \hyperref[theory:wave:applications]{第\ref*{theory:wave:applications}節}

\par 参考：\cite{os_sound}。

\paragraph{関連問題} \hyperref[q:h24-reg-applied-02]{H24-REG-APPLIED-02}

\par 難易度 2（校準初稿）、優先度未評価。

\section{ばね・振り子}

\subsection{平成27年度 問1}
\label{q:h27-reg-applied-01}

\noindent\textbf{原典：} sources/H27/17\_applied\_science.pdf、PDF物理1頁。

\textbf{公式正答：} 4（single）。

\paragraph{問題の要旨} ばね定数1.0 N/mの軽いばねに質量8.0 kgをつるした鉛直単振動の周期を求める。

\Needspace{4\baselineskip}\paragraph{確認した条件・記述（独自要約）}\mbox{}\par

\noindent\textbf{spring\_N\_m：} 1.0\par

\noindent\textbf{mass\_kg：} 8.0\par

\noindent\textbf{spring\_mass：} negligible\par

\paragraph{問われる核心} 復元力の線形化と固有角振動数

\paragraph{必要知識} \(m\ddot x+kx=0\)なら\(T=2\pi\sqrt{m/k}\)。単振り子は小角近似で\(\ddot\theta+(g/\ell)\theta=0\)。

\paragraph{考え方（査読待ち初稿）} \(T=2\pi\sqrt{m/k}=2\pi\sqrt8\simeq17.8\) sで④。重力は平衡位置をずらすが、平衡位置からの振動周期には入らない。

\paragraph{正答を選ぶ理由} 独立に得た条件適合値・組合せは選択肢4である。導出した周期に最も近い。角周波数と周期を区別して2πを含める。

\Needspace{4\baselineskip}\paragraph{選択肢の検討（査読待ち初稿）}\mbox{}\par

\noindent\textbf{1：不採用} 1.4 s。指定条件を振動周期式へ代入した値と一致しない。\par

\noindent\textbf{2：不採用} 2.8 s。指定条件を振動周期式へ代入した値と一致しない。\par

\noindent\textbf{3：不採用} 8.8 s。指定条件を振動周期式へ代入した値と一致しない。\par

\noindent\textbf{4：採用} 18 s。導出した周期に最も近い。角周波数と周期を区別して2πを含める。\par

\noindent\textbf{5：不採用} 26 s。指定条件を振動周期式へ代入した値と一致しない。\par

\paragraph{再利用する考え方} 復元力の線形化と固有角振動数を、条件と式の対応から説明できること。

\paragraph{共通理論} \hyperref[theory:mechanics:oscillation]{第\ref*{theory:mechanics:oscillation}節}

\par 難易度 1（校準初稿）、優先度未評価。

\subsection{平成29年度 問1}
\label{q:h29-reg-applied-01}

\noindent\textbf{原典：} sources/H29/17\_applied\_science.pdf、PDF物理1頁。

\textbf{公式正答：} 3（single）。

\paragraph{問題の要旨} 長さ0.1 mのひもに質量0.1 kgの小さいおもりをつけた小角単振り子の周期を求める。

\Needspace{4\baselineskip}\paragraph{確認した条件・記述（独自要約）}\mbox{}\par

\noindent\textbf{length\_m：} 0.1\par

\noindent\textbf{mass\_kg：} 0.1\par

\noindent\textbf{g\_m\_s2：} 9.8\par

\noindent\textbf{angle：} small\par

\noindent\textbf{bob\_size：} negligible\par

\paragraph{問われる核心} 復元力の線形化と固有角振動数

\paragraph{必要知識} \(m\ddot x+kx=0\)なら\(T=2\pi\sqrt{m/k}\)。単振り子は小角近似で\(\ddot\theta+(g/\ell)\theta=0\)。

\paragraph{考え方（査読待ち初稿）} 小角近似\(\sin\theta\simeq\theta\)から\(T=2\pi\sqrt{\ell/g}\)。\(T=2\pi\sqrt{0.1/9.8}\simeq0.635\) sで③。質量は消える。

\paragraph{正答を選ぶ理由} 独立に得た条件適合値・組合せは選択肢3である。導出した周期に最も近い。角周波数と周期を区別して2πを含める。

\Needspace{4\baselineskip}\paragraph{選択肢の検討（査読待ち初稿）}\mbox{}\par

\noindent\textbf{1：不採用} 0.1 s。指定条件を振動周期式へ代入した値と一致しない。\par

\noindent\textbf{2：不採用} 0.3 s。指定条件を振動周期式へ代入した値と一致しない。\par

\noindent\textbf{3：採用} 0.6 s。導出した周期に最も近い。角周波数と周期を区別して2πを含める。\par

\noindent\textbf{4：不採用} 1.0 s。指定条件を振動周期式へ代入した値と一致しない。\par

\noindent\textbf{5：不採用} 6.0 s。指定条件を振動周期式へ代入した値と一致しない。\par

\paragraph{再利用する考え方} 復元力の線形化と固有角振動数を、条件と式の対応から説明できること。

\paragraph{共通理論} \hyperref[theory:mechanics:oscillation]{第\ref*{theory:mechanics:oscillation}節}

\par 難易度 1（校準初稿）、優先度未評価。

\subsection{令和7年度 問2}
\label{q:r07-reg-applied-02}

\noindent\textbf{原典：} sources/R07/17\_applied\_science.pdf、PDF物理1頁。

\textbf{公式正答：} 1（single）。

\paragraph{問題の要旨} 質量0.80 kgを鉛直ばねにつるすと4.9 cm伸びた。平衡位置から引いて放す単振動の周期を求める。

\Needspace{4\baselineskip}\paragraph{確認した条件・記述（独自要約）}\mbox{}\par

\noindent\textbf{mass\_kg：} 0.8\par

\noindent\textbf{static\_extension\_m：} 0.049\par

\noindent\textbf{g\_m\_s2：} 9.8\par

\noindent\textbf{spring\_mass：} light\par

\paragraph{問われる核心} 復元力の線形化と固有角振動数

\paragraph{必要知識} \(m\ddot x+kx=0\)なら\(T=2\pi\sqrt{m/k}\)。単振り子は小角近似で\(\ddot\theta+(g/\ell)\theta=0\)。

\paragraph{考え方（査読待ち初稿）} 静的釣合い\(k\Delta\ell=mg\)から\(k=0.80\times9.8/0.049=160\) N/m。\(T=2\pi\sqrt{m/k}=2\pi\sqrt{0.80/160}\simeq0.444\) sで①。cmをmへ変換してから計算する。

\paragraph{正答を選ぶ理由} 独立に得た条件適合値・組合せは選択肢1である。導出した周期に最も近い。角周波数と周期を区別して2πを含める。

\Needspace{4\baselineskip}\paragraph{選択肢の検討（査読待ち初稿）}\mbox{}\par

\noindent\textbf{1：採用} \(4.4\times10^{-1}\) s。導出した周期に最も近い。角周波数と周期を区別して2πを含める。\par

\noindent\textbf{2：不採用} \(8.9\times10^{-1}\) s。指定条件を振動周期式へ代入した値と一致しない。\par

\noindent\textbf{3：不採用} 1.0 s。指定条件を振動周期式へ代入した値と一致しない。\par

\noindent\textbf{4：不採用} 1.8 s。指定条件を振動周期式へ代入した値と一致しない。\par

\noindent\textbf{5：不採用} 2.0 s。指定条件を振動周期式へ代入した値と一致しない。\par

\paragraph{再利用する考え方} 復元力の線形化と固有角振動数を、条件と式の対応から説明できること。

\paragraph{共通理論} \hyperref[theory:mechanics:oscillation]{第\ref*{theory:mechanics:oscillation}節}

\par 難易度 2（校準初稿）、優先度未評価。

\section{進行波の位相・波長・角振動数}

\subsection{平成25年度 問2}
\label{q:h25-reg-applied-02}

\noindent\textbf{原典：} sources/H25/17\_applied\_science.pdf、PDF物理1頁。

\textbf{公式正答：} 3（single）。

\paragraph{問題の要旨} 正弦進行波の位相から波長と位相速度の組合せを選ぶ。

\Needspace{4\baselineskip}\paragraph{確認した条件・記述（独自要約）}\mbox{}\par

\noindent\textbf{波：} \(y=A\sin(\omega t-kx)\)、A・ω・kは正。\par

\noindent\textbf{方向：} 波長λと伝播速度vの適切な組合せ。\par

\paragraph{問われる核心} 進行波の位相・波数・角振動数

\paragraph{必要知識} 進行波の位相・波数・角振動数。空間で同位相が繰り返す条件は\(k\lambda=2\pi\)。一定位相を微分すると\(\omega\,dt-k\,dx=0\)で\(v=dx/dt=\omega/k\)。両式が一致する3を選ぶ。角振動数と周期を取り違えない。

\paragraph{考え方・正答を選ぶ理由（査読待ち初稿）} 空間で同位相が繰り返す条件は\(k\lambda=2\pi\)。一定位相を微分すると\(\omega\,dt-k\,dx=0\)で\(v=dx/dt=\omega/k\)。両式が一致する3を選ぶ。角振動数と周期を取り違えない。

\Needspace{4\baselineskip}\paragraph{選択肢の検討（査読待ち初稿）}\mbox{}\par

\noindent\textbf{1：不適切} \(\lambda=2\pi/\omega,\ v=2\pi\omega/k\)。波長の分母は波数であり、角振動数ではない。速度に余分な2πもある。\par

\noindent\textbf{2：不適切} \(\lambda=2\pi/\omega,\ v=2\pi k/\omega\)。波長は角振動数の逆数ではない。速度は逆数に加えて2πが余分。\par

\noindent\textbf{3：最適} \(\lambda=2\pi/k,\ v=\omega/k\)。空間位相が2π増える距離と位相一定の移動から導かれ、両方正しい。\par

\noindent\textbf{4：不適切} \(\lambda=2\pi/k,\ v=k/\omega\)。k/ωは速度の逆数となり、速度の次元が違う。\par

\noindent\textbf{5：不適切} \(\lambda=2\pi/k,\ v=2\pi\omega/k\)。波長は正しいが位相速度に余分な2πを掛けている。\par

\paragraph{再利用する考え方} 進行波の位相・波数・角振動数を暗記だけでなく、位相・幾何・符号と成立条件から導く。選択肢番号を類似年度へそのまま移さない。

\paragraph{共通理論} \hyperref[theory:wave:phase]{第\ref*{theory:wave:phase}節}

\paragraph{関連問題} \hyperref[q:h29-reg-applied-04]{H29-REG-APPLIED-04}

\par 難易度 2（校準初稿）、優先度未評価。

\subsection{平成29年度 問4}
\label{q:h29-reg-applied-04}

\noindent\textbf{原典：} sources/H29/17\_applied\_science.pdf、PDF物理2頁。

\textbf{公式正答：} 5（single）。

\paragraph{問題の要旨} 正弦進行波の位相から波長と位相速度の組合せを選ぶ。

\Needspace{4\baselineskip}\paragraph{確認した条件・記述（独自要約）}\mbox{}\par

\noindent\textbf{波：} \(y=A\sin(\omega t-kx)\)、A・ω・kは正。\par

\noindent\textbf{方向：} 波長λと伝播速度vの適切な組合せ。\par

\paragraph{問われる核心} 進行波の位相・波数・角振動数

\paragraph{必要知識} 進行波の位相・波数・角振動数。空間で同位相が繰り返す条件は\(k\lambda=2\pi\)。一定位相を微分すると\(\omega\,dt-k\,dx=0\)で\(v=dx/dt=\omega/k\)。両式が一致する5を選ぶ。角振動数と周期を取り違えない。

\paragraph{考え方・正答を選ぶ理由（査読待ち初稿）} 空間で同位相が繰り返す条件は\(k\lambda=2\pi\)。一定位相を微分すると\(\omega\,dt-k\,dx=0\)で\(v=dx/dt=\omega/k\)。両式が一致する5を選ぶ。角振動数と周期を取り違えない。

\Needspace{4\baselineskip}\paragraph{選択肢の検討（査読待ち初稿）}\mbox{}\par

\noindent\textbf{1：不適切} \(\lambda=2\pi/\omega,\ v=2\pi\omega/k\)。波長の分母は波数であり、角振動数ではない。速度に余分な2πもある。\par

\noindent\textbf{2：不適切} \(\lambda=2\pi/\omega,\ v=2\pi k/\omega\)。波長は角振動数の逆数ではない。速度は逆数に加えて2πが余分。\par

\noindent\textbf{3：不適切} \(\lambda=2\pi/k,\ v=2\pi\omega/k\)。波長は正しいが位相速度に余分な2πを掛けている。\par

\noindent\textbf{4：不適切} \(\lambda=2\pi/k,\ v=k/\omega\)。k/ωは速度の逆数となり、速度の次元が違う。\par

\noindent\textbf{5：最適} \(\lambda=2\pi/k,\ v=\omega/k\)。空間位相が2π増える距離と位相一定の移動から導かれ、両方正しい。\par

\paragraph{再利用する考え方} 進行波の位相・波数・角振動数を暗記だけでなく、位相・幾何・符号と成立条件から導く。選択肢番号を類似年度へそのまま移さない。

\paragraph{共通理論} \hyperref[theory:wave:phase]{第\ref*{theory:wave:phase}節}

\paragraph{関連問題} \hyperref[q:h25-reg-applied-02]{H25-REG-APPLIED-02}, \hyperref[q:r07-reg-applied-05]{R07-REG-APPLIED-05}

\par 難易度 2（校準初稿）、優先度未評価。

\subsection{令和7年度 問5}
\label{q:r07-reg-applied-05}

\noindent\textbf{原典：} sources/R07/17\_applied\_science.pdf、PDF物理2頁。

\textbf{公式正答：} 4（single）。

\paragraph{問題の要旨} 周期Tと波長λで書かれた正弦波の位相速度と角振動数を求める。

\Needspace{4\baselineskip}\paragraph{確認した条件・記述（独自要約）}\mbox{}\par

\noindent\textbf{波：} \(y=A\sin\{2\pi(t/T-x/\lambda)\}\)。A,T,λは正。\par

\noindent\textbf{方向：} 伝播速度vと角振動数ωの適切な組合せ。\par

\noindent\textbf{候補5：} 角振動数は2πλT（原画像照合、OCR誤認を採用しない）。\par

\paragraph{問われる核心} 周期・波長による位相速度と角振動数

\paragraph{必要知識} 周期・波長による位相速度と角振動数。\(y=A\sin\{2\pi(t/T-x/\lambda)\}\)の時間係数は\(\omega=2\pi/T\)。一定位相では\(dt/T=dx/\lambda\)なので\(v=\lambda/T\)で4。

\paragraph{考え方・正答を選ぶ理由（査読待ち初稿）} \(y=A\sin\{2\pi(t/T-x/\lambda)\}\)の時間係数は\(\omega=2\pi/T\)。一定位相では\(dt/T=dx/\lambda\)なので\(v=\lambda/T\)で4。

\Needspace{4\baselineskip}\paragraph{選択肢の検討（査読待ち初稿）}\mbox{}\par

\noindent\textbf{1：不適切} \(v=2\pi/\lambda,\ \omega=\lambda/T\)。速度と角振動数の双方の次元が異なる。\par

\noindent\textbf{2：不適切} \(v=\lambda T,\ \omega=2\pi/\lambda\)。速度は周期で割る量であり積ではない。角振動数も空間の波数ではない。\par

\noindent\textbf{3：不適切} \(v=2\pi\lambda/T,\ \omega=2\pi/T\)。角振動数は正しいが、位相速度に余分な2πがある。\par

\noindent\textbf{4：最適} \(v=\lambda/T,\ \omega=2\pi/T\)。位相一定の微分と時間位相係数を両方満たす。\par

\noindent\textbf{5：不適切} \(v=\lambda/T,\ \omega=2\pi\lambda T\)。速度は正しいが角振動数に波長と周期を掛け、次元が異なる。\par

\paragraph{再利用する考え方} 周期・波長による位相速度と角振動数を暗記だけでなく、位相・幾何・符号と成立条件から導く。選択肢番号を類似年度へそのまま移さない。

\paragraph{共通理論} \hyperref[theory:wave:phase]{第\ref*{theory:wave:phase}節}

\paragraph{関連問題} \hyperref[q:h29-reg-applied-04]{H29-REG-APPLIED-04}

\par 難易度 2（校準初稿）、優先度未評価。


\chapter{熱力学}
\section{背景：保存される量と、変換できる量を分ける}
発電機関では燃料から得たエネルギーの全てを仕事へ変換できない。断熱材ではエネルギーが保存されても流れる速さが重要であり、圧力計ではその熱移動を信号として利用する。熱力学は、巨視的な系を少数の状態変数で表して、エネルギーの保存、変化の方向、平衡条件を調べる枠組みである。輸送現象はそこへ時間・空間の勾配を加える。

物理専攻の読者には、公式の再暗記よりも\emph{系の境界、仕事の向き、モデルの条件}を問題ごとに確認することが重要である。温度の数値が°Cであっても、状態方程式の比には絶対温度が必要である。「断熱」と「仕事をしない」、「理想気体」と「実在気体」、「比熱」と「熱伝導率」は異なる概念である。以下の理論から個別問題の条件へ接続する。

\section{第一法則：内部エネルギーと熱・仕事の符号}
\label{theory:thermo:firstlaw}
閉じた系について、系に入る熱量を$Q>0$、系へ外部から行われる仕事を$W_{\mathrm{on}}>0$とすると、全体の運動・位置エネルギー変化を無視した第一法則は
\[
 \Delta U=Q+W_{\mathrm{on}}=Q-W_{\mathrm{out}}
\]
である。$U$は状態量であり、$Q,W$は経路に依存する。体積仕事だけなら
\[
 \delta W_{\mathrm{on}}=-p_{\mathrm{ext}}\,dV,
 \qquad W_{\mathrm{out}}=\int p_{\mathrm{ext}}\,dV.
\]
圧縮では$dV<0$なので系へされる仕事は正である。準静的な力学平衡では外圧と系の圧力がほぼ一致するが、急速な膨張で系内の圧力をそのまま仕事の式へ代入してよいわけではない。符号と経路の扱いは大学講義の第一法則の整理と整合する\parencite{mit_thermo_firstlaw2008}。

断熱は境界からの熱流がない条件であり、電気入力や攪拌仕事を否定しない。断熱シリンダーの物質Aをヒーターで気化させる問題では、電気的な入力が物質Aのエネルギーを増やす。求める仕事が$p\Delta V$なら、ヒーターの総入力や潜熱と取り違えない。

\section{理想気体：状態方程式と熱容量の差}
\label{theory:thermo:ideal}
理想気体は分子間相互作用と占有体積の効果を無視するモデルである。量$n$を固定すると
\[
 pV=nRT,
 \qquad T[\mathrm{K}]=t[{}^\circ\mathrm{C}]+273.15.
\]
等温では$V\propto1/p$、定圧では$V\propto T$、定積では$p\propto T$となる。273へ丸めると本章の27°C、327°C、−123°Cは300、600、150Kとなり、候補の精度には十分である。温度差1°Cと1Kは等しいが、温度の比は°Cの値から作れない。

\subsection{定圧と定積で必要熱量が異なる理由}
\label{theory:thermo:capacity}
熱容量は系全体の温度変化当たりのエネルギー、比熱は質量当たり、モル熱容量は物質量当たりである。単相・一定組成の理想気体では$U=U(T)$であり、
\[
 C_{V,m}=\frac1n\left(\frac{\partial U}{\partial T}\right)_V,
 \qquad H=U+pV=U+nRT,
 \qquad C_{p,m}-C_{V,m}=R.
\]
定積加熱は内部エネルギーを増やし、定圧加熱はそれに加え膨張仕事を行う。このため同じ理想気体の同じ昇温に必要な熱は定圧の方が多い。熱容量がほぼ一定の区間なら$Q_V=nC_{V,m}\Delta T$、$Q_p=nC_{p,m}\Delta T$であり、大きく変わる区間では積分する。

\section{分子自由度と比熱：温度尺度を読む}
\label{theory:thermo:degrees}
分子のエネルギーには並進、回転、振動と、分子間の相互作用が寄与する。古典的な二次形式の自由度は1つ当たり平均$k_BT/2$を持つ。単原子理想気体では並進の3自由度から$C_{V,m}=3R/2$となる。二原子分子では回転の2自由度が励起され、振動が凍結している範囲なら
\[
 C_{V,m}=\frac52R,\qquad C_{p,m}=\frac72R,
 \qquad \gamma=\frac{C_{p,m}}{C_{V,m}}=\frac75.
\]
振動は運動とポテンシャルの双方を持つので、1振動モードが古典的に励起されると$R$の寄与を加える。低温では励起準位の間隔が$k_BT$より大きくなり、その自由度の熱容量への寄与は凍結する。従って「比較的低い温度で7/5」は、振動の励起温度より低く、回転の励起温度より高い中間範囲を読む必要がある。内部エネルギーの絶対零点を分子種間でそろえる問題と、温度によって増える熱エネルギーの自由度比較は区別する。

\section{断熱膨張：外部仕事がある場合と自由膨張}
\label{theory:thermo:adiabatic}
体積仕事だけの断熱過程では$nC_{V,m}dT=-p_{\mathrm{ext}}dV$である。外部へ正の仕事を行えば内部エネルギーが減り、理想気体は冷える。可逆で熱容量一定なら$p_{\mathrm{ext}}\simeq p=nRT/V$より
\[
 \frac{dT}{T}=-\frac{R}{C_{V,m}}\frac{dV}{V},
 \qquad TV^{\gamma-1}=\text{一定},\qquad pV^\gamma=\text{一定}.
\]
可逆断熱と、一定外圧に抗する不可逆断熱は区別して扱う\parencite{mit_thermo_adiabatic2008}。真空への自由膨張では外圧が0で仕事も0となり、断熱理想気体の$\Delta U$と$\Delta T$は0である。膨張が起きることだけから、常に温度低下すると推論してはいけない。

\begin{tcolorbox}[title=原典の限定不足を保持する]
平成25年度問9と平成30年度問6は、断熱膨張の記述に仕事又は可逆性の指定がない。定積・定圧の必要熱量を逆転した肢は明確に誤るが、断熱膨張の肢も自由膨張を含めれば無条件には正しくない。教材側で原典を可逆膨張へ書き換えず、独立判定と解釈疑義を保存する。
\end{tcolorbox}

\section{第二法則とカルノー機関：温度比の意味}
\label{theory:thermo:entropy}
熱が高温物体から低温物体へ受動的に流れるとき、エネルギーは保存されるが全体のエントロピーは増える。高温源$T_H$から熱$Q>0$が低温源$T_L$へ移れば
\[
 \Delta S_{\mathrm{total}}=-\frac{Q}{T_H}+\frac{Q}{T_L}>0.
\]
冷凍機は外部仕事で低温側の熱を運び出し、装置と周囲を含めた全体で第二法則を満たす。「低温から高温へ熱を移す装置がある」ことは受動的熱移動の法則への反例ではない。

\subsection{可逆サイクルの効率の上限}
\label{theory:thermo:carnot}
2熱源の間で作動する可逆サイクルでは、正の大きさを使って吸熱$Q_H$、放熱$Q_L$を表すと、
\[
 \frac{Q_H}{T_H}=\frac{Q_L}{T_L},
 \qquad W=Q_H-Q_L,
 \qquad \eta=\frac{W}{Q_H}=1-\frac{T_L}{T_H}.
\]
これがカルノー効率である。不可逆なサイクルでは、熱源も含めたエントロピー生成が非負であるため$Q_L/T_L-Q_H/T_H\ge0$、従って効率はこの値を超えない。$T_L\to T_H$で効率は0へ向かう。この極限は分母を温度の和にした誤式を見分ける手掛かりとなる。放熱の問題は$Q_L=Q_HT_L/T_H$、仕事の問題はその差を使う。

\section{仕事の経路と断熱混合の熱収支}
\label{theory:thermo:workpaths}
定圧・定積の段階を表へ分けると、各段階で不変な量を取り違えずに済む。固定$n$の定圧段階なら
\[
 W_{\mathrm{on}}=-p\Delta V=-nR\Delta T,
 \qquad nR=\frac{p_0V_0}{T_0}.
\]
定積段階の体積仕事は0である。これは「第3段階」の仕事を聞く問題で、サイクル全体の仕事や、初期圧力を全段階へ使う誤りを避ける。27°C、227°C、27°C、−73°Cの順の令和7年度問6では、第3段階の冷却幅は約100Kであり、初期量から$nR$を求めるだけでも仕事が出る。

気化の定圧仕事では、液体体積$V_l$を無視して理想蒸気の終体積を使えば
\[
 W_{\mathrm{out}}=p(V_g-V_l)\simeq pV_g=nRT.
\]
外圧が打ち消されるのは理想蒸気と体積無視という仮定の結果である。一般の相変化で圧力を無視してよいわけではない。

\subsection{熱容量を重みにする接触熱収支}
\label{theory:thermo:mixheat}
2物体を接触させ、相変化と外界への熱流を無視し比熱を一定とすると、
\[
 m_1c_1(T_f-T_1)+m_2c_2(T_f-T_2)=0,
 \qquad T_f=\frac{m_1c_1T_1+m_2c_2T_2}{m_1c_1+m_2c_2}.
\]
最終温度は熱容量を重みにした平均である。等質量でも比熱が違えば単純平均ではない。平成26年度問6は水の熱容量が金属球の10倍なので、最終温度は水の初期25°Cに近い。容器も吸熱する現実の実験では、その熱容量を収支へ加える。

\section{相平衡、混合、実在気体}
\label{theory:thermo:phase}
相平衡では各相の化学ポテンシャルが等しい。純物質で$d\mu=-s_m\,dT+v_m\,dp$を液相と気相へ適用し、共存曲線上で両者を等置すると
\[
 \frac{dp}{dT}=\frac{s_{g,m}-s_{l,m}}{v_{g,m}-v_{l,m}}
 =\frac{L_m}{T\Delta v_m}.
\]
これがクラペイロンの式である。液体の体積を無視し、気体を理想気体とすると$d\ln p/dT=L_m/(RT^2)$となる。さらに潜熱$L_m$を区間内で一定と近似すれば
\[
 \ln\frac{p_2}{p_1}=-\frac{L_m}{R}\left(\frac1{T_2}-\frac1{T_1}\right).
\]
このクラウジウス＝クラペイロンの関係で外圧から沸点を求める\parencite{mit_clapeyron2008}。外圧が低い山頂では沸点が低い。三重点の温度・圧力だけで共存曲線の全体が一意に決まるわけではなく、潜熱などの情報が必要である。式の名前を選ぶ問題へ、指定されない山頂の沸点の数値を補わない。

一次相転移は通常、$G$の一階微分である$S,V$などに不連続性を持ち、潜熱は$T\Delta S$である。二次・連続転移では一階微分が連続でも二階微分に跳びや発散が現れ得る。従って二次転移なら比熱も必ず連続という主張は誤る。過冷却は核生成が遅れて準安定相が続く速度論的な状態であり、平衡相図と区別する。

\subsection{異種気体の混合と理想溶液}
\label{theory:thermo:mixing}
同温同圧の異種理想気体を混ぜると、組成$x_i$について$\Delta S_{\mathrm{mix}}=-nR\sum_i x_i\ln x_i>0$、$\Delta G_{\mathrm{mix}}=nRT\sum_i x_i\ln x_i<0$となる。理想溶液では液相の活動度がモル分率と等しくなり、蒸気圧は$p_i=x_i p_i^*$に従う。これはラウールの法則であり、現実の分子間相互作用が大きく異なる溶液では活動度の補正が必要になる\parencite{mit_solutions2008}。同じ物質の分割を混ぜることを異種混合と数えることはしない。

\subsection{実在気体の圧縮係数}
\label{theory:thermo:realgas}
実在気体の圧縮係数$Z=pV_m/(RT)$は、同じ$T,p$の理想モル体積に対する実在モル体積の比である。ファンデルワールス式なら
\[
 \left(p+\frac{a}{V_m^2}\right)(V_m-b)=RT,
 \qquad Z=\frac{V_m}{V_m-b}-\frac{a}{RTV_m}.
\]
第1項の排除体積は$Z$を増やし、第2項の引力は減らす。温度・密度によって$Z$は1より小さくも大きくもなり、希薄極限で1へ向かう。この式は定性的にずれを理解する近似であり、実用の高精度状態方程式を代替する保証ではない。

\section{熱伝導・放射・計測：輸送の担体を区別する}
\label{theory:thermo:transport}
線形・等方の熱伝導では熱流密度$\boldsymbol{j}_q=-\kappa\nabla T$である。熱伝導率$\kappa$の単位はW/(m・K)、比熱の単位はJ/(kg・K)であり、前者は輸送のしやすさ、後者は温度を変えるエネルギーを表す。結晶では格子振動が熱を運び、金属では伝導電子も重要である。通常の縮退金属の電子輸送では
\[
 \frac{\kappa_e}{\sigma T}\simeq L_0,
 \qquad L_0=\frac{\pi^2}{3}\left(\frac{k_B}{e}\right)^2
\]
というウィーデマン＝フランツ則が成り立つ近似がある。格子伝導が大きい材料まで、電気伝導の大小から熱伝導を判定できるわけではない。

産総研の室温付近の値では、Alは226〜237、Feは72〜80.4、Cuは386〜402、Agは429W/(m・K)程度である\parencite{aist_metal_thermal}。高品質バルクダイヤモンドは格子による熱伝導が大きく、通常の銅・鉄鋼・板ガラスを上回るが、粒径や欠陥で変わる\parencite{wang_sun_diamond2024}。板ガラスの目安は約1W/(m・K)である\parencite{agc_glass_thermal}。これらは試料の代表値であり、物質名だけで全ての試料の厳密値を決める表ではない。

多孔質の断熱材は固体の伝導経路を減らし、気体の小さい熱伝導を利用する。ただし孔が大きく対流が起きる場合や高温での放射は別に評価する。熱放射は電磁波によるエネルギー移動であり、物質媒体の移動を必要としない。

\subsection{熱輸送と熱電能を利用する計測}
\label{theory:thermo:measurement}
ピラニ真空計は加熱フィラメントから気体へ逃げる熱が圧力で変わることを利用する。温度維持の加熱電力又は抵抗変化を圧力へ換算する原理である\parencite{ulvac_pirani}。吸着で生じる圧力変化から熱伝導率を測る説明とは、入力と出力が違う。実測では気体種や校正範囲も確認する。

熱電対は異種材料の熱電能の差によるゼーベック効果を用いる。基準温度$T_r$と測定温度$T_h$に対する電圧は、向きの定義に応じて$\int_{T_r}^{T_h}(S_A-S_B)\,dT$である。熱起電力を読むには基準側の温度も必要である。電流を流した接合で熱が出入りするペルティエ効果とは役割を区別する。

\section{固体比熱：電子・格子・相転移を分ける}
\label{theory:thermo:solidcapacity}
電子の励起はフェルミ準位の近傍に限られ、滑らかな状態密度$D(E_F)$を持つ通常の金属の低温電子比熱は
\[
 C_e=\gamma T,\qquad\gamma=\frac{\pi^2}{3}k_B^2D(E_F).
\]
$D$を系全体の状態数/エネルギーと定義すれば$C_e$も系全体の熱容量である。単位体積当たりの状態密度なら体積当たりとなる。反比例ではなく比例であり、熱的に励起できる状態が多いほど比熱が大きい\parencite{mit_electronic_heat2012}。

格子ではデバイ温度$\Theta_D$が振動の励起の尺度を与え、1molの原子の比熱はデバイモデルで
\[
 C_{V,m}=9R\left(\frac{T}{\Theta_D}\right)^3
 \int_0^{\Theta_D/T}\frac{x^4 e^x}{(e^x-1)^2}\,dx.
\]
低温で$T^3$、高温で$3R$へ近づく。このデュロン＝プティ則の高温とは、室温という固定値でなく$T\gg\Theta_D$の意味である\parencite{mit_debye2005}。ダイヤモンドのように励起温度の高い物質では室温の格子比熱は$3R$からずれる。格子比熱が小さいことと熱伝導率が大きいことは矛盾しない。後者には熱を運ぶ励起の速度と散乱の少なさも効く。

固体の熱容量差は、体積膨張率$\alpha_V=(1/V)(\partial V/\partial T)_p$、等温圧縮率$\kappa_T=-(1/V)(\partial V/\partial p)_T$を使って
\[
 C_p-C_V=\frac{TV\alpha_V^2}{\kappa_T}
\]
と表せる。一般に低温で熱膨張が小さくなるため差は小さくなるが、「100K以下なら全物質で無視できる」という普遍的境界ではない。転移近傍や特殊な材料の例外を無視しない。

\section{問題の読み方とこの初稿の確認範囲}
計算問題では、対象量を仕事・熱・内部エネルギーに分け、向きと単位を決めてから不変量を使う。知識問題では、理想モデル、温度範囲、相転移の次数、計測の入出力を確認する。複数の近似を含む5肢の問題で、最も明確な誤りの選択と他肢の普遍的な正しさは同じ判断ではない。令和3年度問7の温度範囲の省略は、断熱膨張の2問とともに査読事項として保持する。

このバッチは熱力学22問を原典画像から確認した初稿である。頻度表は確認済み部分集合と全体の暫定集計を区別し、将来出題を保証しない。以下の個別解説は全てAstra査読前であり、他章の未執筆問題も含めて作業を継続する。
\section{熱力学・熱機関}

\subsection{平成23年度 問5}
\label{q:h23-reg-applied-05}

\noindent\textbf{原典：} sources/H23/17\_applied\_science.pdf、PDF物理2頁。

\textbf{公式正答：} 1（single）。

\paragraph{問題の要旨} 気体の状態変化、第一法則、第二法則の5記述から最も不適切なものを選ぶ。

\Needspace{4\baselineskip}\paragraph{確認した条件・記述（独自要約）}\mbox{}\par

\noindent\textbf{方向：} 最も不適切な記述を選ぶ。\par

\noindent\textbf{原典の範囲：} 一定質量の気体、温度又は圧力一定。理想気体という語は明記されていない。\par

\paragraph{問われる核心} 第一法則の仕事の符号と気体の状態方程式

\paragraph{必要知識} 閉じた系、理想気体近似、絶対温度、熱と仕事の符号、受動的な熱移動。

\paragraph{考え方・正答を選ぶ理由（査読待ち初稿）} 等温の気体の体積と圧力の関係を\(pV=nRT\)で判定する。体積は圧力へ反比例するので1が不適切である。2は理想気体近似のシャルルの法則、3と4はエネルギー保存、5は第二法則である。気体法則は全ての実在気体へ厳密に成立するという主張へ拡張しない。

\Needspace{4\baselineskip}\paragraph{選択肢の検討（査読待ち初稿）}\mbox{}\par

\noindent\textbf{1：誤り} 温度一定・一定質量の気体の体積は圧力に比例する。理想気体なら\(pV=nRT\)であり、固定\(n,T\)では\(V=nRT/p\)と反比例する。比例という記述が誤りである。\par

\noindent\textbf{2：理想気体近似で正しい} 圧力一定・一定質量の気体の体積は絶対温度に比例する。理想気体又はその近似が成り立つ範囲では固定\(n,p\)により\(V=(nR/p)T\)である。\par

\noindent\textbf{3：正しい} 外部からの熱と系へされた仕事の和が内部エネルギー増加に等しい。熱を流入正、仕事を系へ入力正とすれば\(\Delta U=Q+W_{\mathrm{on}}\)である。系全体の運動・位置エネルギー変化はこの閉じた熱力学系では考えない。\par

\noindent\textbf{4：正しい} エネルギー変換で関与する全エネルギーの和が保存される。第一法則のエネルギー保存の内容である。対象系から外へ出るエネルギーを数え落とさない。\par

\noindent\textbf{5：正しい} 熱は高温から低温へ移り、自然に低温から高温へは移らない。第二法則の巨視的な受動的熱移動の方向である。冷凍機は外部仕事を使うので反例ではない。\par

\paragraph{再利用する考え方} モデルの前提を明示し、第一法則の仕事の符号と気体の状態方程式から記述や数値を判定する。

\paragraph{共通理論} \hyperref[theory:thermo:firstlaw]{第\ref*{theory:thermo:firstlaw}節}, \hyperref[theory:thermo:ideal]{第\ref*{theory:thermo:ideal}節}, \hyperref[theory:thermo:entropy]{第\ref*{theory:thermo:entropy}節}

\par 参考：\cite{mit_thermo_firstlaw2008}。

\paragraph{関連問題} \hyperref[q:h26-reg-applied-08]{H26-REG-APPLIED-08}

\par 難易度 1（校準初稿）、優先度未評価。

\subsection{平成24年度 問8}
\label{q:h24-reg-applied-08}

\noindent\textbf{原典：} sources/H24/17\_applied\_science.pdf、PDF物理2頁。

\textbf{公式正答：} 3（single）。

\paragraph{問題の要旨} 内部エネルギーの要素と単原子・二原子分子の自由度を、4つの空欄へ対応させる。

\Needspace{4\baselineskip}\paragraph{確認した条件・記述（独自要約）}\mbox{}\par

\noindent\textbf{A：} 分子間の力に由来するエネルギー。\par

\noindent\textbf{B：} 分子の熱運動によるエネルギー。\par

\noindent\textbf{C：} 二原子分子ではBに加わる分子内自由度。\par

\noindent\textbf{D：} 同温での分子1個当たりの二原子分子の平均内部エネルギーを、単原子分子と比較する。\par

\paragraph{問われる核心} 分子自由度と内部エネルギーの構成

\paragraph{必要知識} 並進・回転・振動、分子間相互作用、等分配と量子凍結。

\paragraph{考え方・正答を選ぶ理由（査読待ち初稿）} 分子間の力は位置エネルギー、分子の並進は運動エネルギーに対応する。二原子分子には回転と振動があり、励起される自由度は熱エネルギーに寄与するのでA位置・B運動・C振動回転・D大きいの3となる。量子論では低温に振動・回転が凍結し、常に全自由度が等分配されるわけではない。また分子種間の内部エネルギーの絶対零点を比較する問題とは区別する。

\Needspace{4\baselineskip}\paragraph{選択肢の検討（査読待ち初稿）}\mbox{}\par

\noindent\textbf{1：条件に不一致} A運動／B振動・回転／C位置／D小さい。Aは分子間相互作用の位置エネルギーであり、Bは並進の運動エネルギーである。A〜Cの対応が違う。\par

\noindent\textbf{2：条件に不一致} A運動／B振動・回転／C位置／D大きい。大きいというDだけでは足りず、A〜Cの対応が違う。\par

\noindent\textbf{3：条件に一致} A位置／B運動／C振動・回転／D大きい。A位置、B運動、C振動・回転、D大きいという対応に一致する。\par

\noindent\textbf{4：条件に不一致} A位置／B運動／C振動・回転／D小さい。二原子分子の回転等が熱的に励起される条件では平均の熱エネルギーが加わる。小さいは逆である。\par

\noindent\textbf{5：条件に不一致} A運動／B位置／C振動・回転／D大きい。相互作用の位置エネルギーと並進の運動エネルギーを逆にしている。\par

\paragraph{再利用する考え方} モデルの前提を明示し、分子自由度と内部エネルギーの構成から記述や数値を判定する。

\paragraph{共通理論} \hyperref[theory:thermo:degrees]{第\ref*{theory:thermo:degrees}節}

\par 参考：\cite{mit_debye2005}。

\par 難易度 1（校準初稿）、優先度未評価。

\subsection{平成25年度 問8}
\label{q:h25-reg-applied-08}

\noindent\textbf{原典：} sources/H25/17\_applied\_science.pdf、PDF物理4頁。

\textbf{公式正答：} 5（single）。

\paragraph{問題の要旨} 高温・低温の2熱源で作動する熱機関の最大効率を、機関名と式の組合せで選ぶ。

\Needspace{4\baselineskip}\paragraph{確認した条件・記述（独自要約）}\mbox{}\par

\noindent\textbf{温度：} 絶対温度\(T_H>T_L>0\)。\par

\noindent\textbf{効率：} 外部へ行う仕事／高温熱源から吸収する熱。\par

\noindent\textbf{空欄：} アは最大効率を与える熱機関名、イはその効率。\par

\paragraph{問われる核心} 第二法則が決める熱機関効率の上限

\paragraph{必要知識} 2熱源、可逆性、熱収支、絶対温度、サイクルのエントロピー収支。

\paragraph{考え方・正答を選ぶ理由（査読待ち初稿）} 熱収支は\(W=Q_H-Q_L\)、可逆サイクルでは\(Q_H/T_H=Q_L/T_L\)である。従って\(\eta=W/Q_H=1-T_L/T_H=(T_H-T_L)/T_H\)となり、カルノー機関とこの式を組み合わせた5である。不可逆損失がある実機の効率はこの上限以下になる。

\Needspace{4\baselineskip}\paragraph{選択肢の検討（査読待ち初稿）}\mbox{}\par

\noindent\textbf{1：条件に不一致} カルノー／\((T_H-T_L)/(T_H+T_L)\)。機関名は合うが効率の分母は吸熱源の絶対温度であり、温度の和ではない。\par

\noindent\textbf{2：条件に不一致} フェルマー／\((T_H-T_L)/(T_H+T_L)\)。機関名と分母の両方が違う。\par

\noindent\textbf{3：条件に不一致} カルノー／\(T_H/(T_H+T_L)\)。低温源が高温源に近づくと効率は0へ向かうべきであり、この式はその条件を満たさない。\par

\noindent\textbf{4：条件に不一致} フェルマー／\((T_H-T_L)/T_H\)。効率式は合うが機関名が違う。\par

\noindent\textbf{5：条件に一致} カルノー／\((T_H-T_L)/T_H\)。可逆なカルノー機関とその効率に一致する。\par

\paragraph{再利用する考え方} モデルの前提を明示し、第二法則が決める熱機関効率の上限から記述や数値を判定する。

\paragraph{共通理論} \hyperref[theory:thermo:carnot]{第\ref*{theory:thermo:carnot}節}, \hyperref[theory:thermo:entropy]{第\ref*{theory:thermo:entropy}節}

\paragraph{関連問題} \hyperref[q:h28-reg-applied-08]{H28-REG-APPLIED-08}

\par 難易度 1（校準初稿）、優先度未評価。

\subsection{平成26年度 問6}
\label{q:h26-reg-applied-06}

\noindent\textbf{原典：} sources/H26/17\_applied\_science.pdf、PDF物理3頁。

\textbf{公式正答：} 1（single）。

\paragraph{問題の要旨} 水と高温金属球を外界との熱交換なしで接触させ、共通の最終温度を求める。

\Needspace{4\baselineskip}\paragraph{確認した条件・記述（独自要約）}\mbox{}\par

\noindent\textbf{水：} 100g、25°C、比熱4.2J/(g・K)\par

\noindent\textbf{金属：} 100g、80°C、比熱0.42J/(g・K)\par

\noindent\textbf{無視するもの：} ビーカーへの伝熱と外界への熱損失。\par

\paragraph{問われる核心} 断熱混合のエネルギー保存と熱容量の重み

\paragraph{必要知識} 熱容量mc、比熱と熱容量の違い、相変化なし、熱収支。

\paragraph{考え方・正答を選ぶ理由（査読待ち初稿）} 水が受け取る熱と金属が放出する熱を等置し、\(100\times4.2(T-25)=100\times0.42(80-T)\)とする。整理すると\(420(T-25)=42(80-T)\)、\(462T=13860\)より\(T=30\,{}^\circ\mathrm{C}\)となる。第1候補である。相変化はなく温度差だけなので°Cの差をそのまま使える。

\Needspace{4\baselineskip}\paragraph{選択肢の検討（査読待ち初稿）}\mbox{}\par

\noindent\textbf{1：条件に一致} 約30°C。水420J/Kと金属42J/Kの熱容量で重みを付けると30°Cになる。\par

\noindent\textbf{2：条件に不一致} 約40°C。金属の熱容量は水の1/10であり、共通温度は水の初期温度にもっと近い。\par

\noindent\textbf{3：条件に不一致} 約50°C。等質量でも比熱が異なるため単純平均ではない。\par

\noindent\textbf{4：条件に不一致} 約60°C。水の大きな熱容量を無視している。\par

\noindent\textbf{5：条件に不一致} 約70°C。水の大きな熱容量を無視している。\par

\paragraph{再利用する考え方} モデルの前提を明示し、断熱混合のエネルギー保存と熱容量の重みから記述や数値を判定する。

\paragraph{共通理論} \hyperref[theory:thermo:mixheat]{第\ref*{theory:thermo:mixheat}節}

\par 難易度 2（校準初稿）、優先度未評価。

\subsection{平成26年度 問8}
\label{q:h26-reg-applied-08}

\noindent\textbf{原典：} sources/H26/17\_applied\_science.pdf、PDF物理3頁。

\textbf{公式正答：} 1（single）。

\paragraph{問題の要旨} 気体の状態変化、第一法則、第二法則の5記述から最も不適切なものを選ぶ。

\Needspace{4\baselineskip}\paragraph{確認した条件・記述（独自要約）}\mbox{}\par

\noindent\textbf{方向：} 最も不適切な記述を選ぶ。\par

\noindent\textbf{原典の範囲：} 一定質量の気体、温度又は圧力一定。理想気体という語は明記されていない。\par

\paragraph{問われる核心} 第一法則の仕事の符号と気体の状態方程式

\paragraph{必要知識} 閉じた系、理想気体近似、絶対温度、熱と仕事の符号、受動的な熱移動。

\paragraph{考え方・正答を選ぶ理由（査読待ち初稿）} 等温の気体の体積と圧力の関係を\(pV=nRT\)で判定する。体積は圧力へ反比例するので1が不適切である。2は理想気体近似のシャルルの法則、3と4はエネルギー保存、5は第二法則である。気体法則は全ての実在気体へ厳密に成立するという主張へ拡張しない。

\Needspace{4\baselineskip}\paragraph{選択肢の検討（査読待ち初稿）}\mbox{}\par

\noindent\textbf{1：誤り} 温度一定・一定質量の気体の体積は圧力に比例する。理想気体なら\(pV=nRT\)であり、固定\(n,T\)では\(V=nRT/p\)と反比例する。比例という記述が誤りである。\par

\noindent\textbf{2：理想気体近似で正しい} 圧力一定・一定質量の気体の体積は絶対温度に比例する。理想気体又はその近似が成り立つ範囲では固定\(n,p\)により\(V=(nR/p)T\)である。\par

\noindent\textbf{3：正しい} 外部からの熱と系へされた仕事の和が内部エネルギー増加に等しい。熱を流入正、仕事を系へ入力正とすれば\(\Delta U=Q+W_{\mathrm{on}}\)である。系全体の運動・位置エネルギー変化はこの閉じた熱力学系では考えない。\par

\noindent\textbf{4：正しい} エネルギー変換で関与する全エネルギーの和が保存される。第一法則のエネルギー保存の内容である。対象系から外へ出るエネルギーを数え落とさない。\par

\noindent\textbf{5：正しい} 熱は高温から低温へ移り、自然に低温から高温へは移らない。第二法則の巨視的な受動的熱移動の方向である。冷凍機は外部仕事を使うので反例ではない。\par

\paragraph{再利用する考え方} モデルの前提を明示し、第一法則の仕事の符号と気体の状態方程式から記述や数値を判定する。

\paragraph{共通理論} \hyperref[theory:thermo:firstlaw]{第\ref*{theory:thermo:firstlaw}節}, \hyperref[theory:thermo:ideal]{第\ref*{theory:thermo:ideal}節}, \hyperref[theory:thermo:entropy]{第\ref*{theory:thermo:entropy}節}

\par 参考：\cite{mit_thermo_firstlaw2008}。

\paragraph{関連問題} \hyperref[q:h23-reg-applied-05]{H23-REG-APPLIED-05}, \hyperref[q:r01-ret-applied-04]{R01-RET-APPLIED-04}

\par 難易度 1（校準初稿）、優先度未評価。

\subsection{平成27年度 問8}
\label{q:h27-reg-applied-08}

\noindent\textbf{原典：} sources/H27/17\_applied\_science.pdf、PDF物理3頁。

\textbf{公式正答：} 3（single）。

\paragraph{問題の要旨} カルノーサイクルの吸熱と2熱源の温度から、1サイクルの低温熱源への放熱を求める。

\Needspace{4\baselineskip}\paragraph{確認した条件・記述（独自要約）}\mbox{}\par

\noindent\textbf{高温熱源：} 1500K\par

\noindent\textbf{低温熱源：} 1200K\par

\noindent\textbf{高温側吸熱：} 750J\par

\noindent\textbf{求める量：} 低温源へ放出する熱量の正の大きさ。\par

\paragraph{問われる核心} カルノーサイクルの可逆エントロピー収支

\paragraph{必要知識} 絶対温度、吸熱と放熱の正の大きさ、仕事と熱収支。

\paragraph{考え方・正答を選ぶ理由（査読待ち初稿）} 可逆サイクルのエントロピー変化は0なので\(Q_L=Q_HT_L/T_H\)である。数値を代入すると\(750\times1200/1500=600\,\mathrm{J}\)となり、3を選ぶ。取り出した仕事は吸熱と放熱の差であり、求める放熱と混同しない。

\Needspace{4\baselineskip}\paragraph{選択肢の検討（査読待ち初稿）}\mbox{}\par

\noindent\textbf{1：条件に不一致} 480J。可逆性から得る低温側放熱600Jと一致しない。\par

\noindent\textbf{2：条件に不一致} 540J。可逆性から得る低温側放熱600Jと一致しない。\par

\noindent\textbf{3：条件に一致} 600J。吸熱に低温／高温の温度比を掛けた放熱に一致する。\par

\noindent\textbf{4：条件に不一致} 660J。可逆性から得る低温側放熱600Jと一致しない。\par

\noindent\textbf{5：条件に不一致} 720J。可逆性から得る低温側放熱600Jと一致しない。\par

\paragraph{再利用する考え方} モデルの前提を明示し、カルノーサイクルの可逆エントロピー収支から記述や数値を判定する。

\paragraph{共通理論} \hyperref[theory:thermo:carnot]{第\ref*{theory:thermo:carnot}節}

\paragraph{関連問題} \hyperref[q:r04-reg-applied-06]{R04-REG-APPLIED-06}

\par 難易度 2（校準初稿）、優先度未評価。

\subsection{平成28年度 問8}
\label{q:h28-reg-applied-08}

\noindent\textbf{原典：} sources/H28/17\_applied\_science.pdf、PDF物理3頁。

\textbf{公式正答：} 5（single）。

\paragraph{問題の要旨} 高温・低温の2熱源で作動する熱機関の最大効率を、機関名と式の組合せで選ぶ。

\Needspace{4\baselineskip}\paragraph{確認した条件・記述（独自要約）}\mbox{}\par

\noindent\textbf{温度：} 絶対温度\(T_H>T_L>0\)。\par

\noindent\textbf{効率：} 外部へ行う仕事／高温熱源から吸収する熱。\par

\noindent\textbf{空欄：} アは最大効率を与える熱機関名、イはその効率。\par

\paragraph{問われる核心} 第二法則が決める熱機関効率の上限

\paragraph{必要知識} 2熱源、可逆性、熱収支、絶対温度、サイクルのエントロピー収支。

\paragraph{考え方・正答を選ぶ理由（査読待ち初稿）} 熱収支は\(W=Q_H-Q_L\)、可逆サイクルでは\(Q_H/T_H=Q_L/T_L\)である。従って\(\eta=W/Q_H=1-T_L/T_H=(T_H-T_L)/T_H\)となり、カルノー機関とこの式を組み合わせた5である。不可逆損失がある実機の効率はこの上限以下になる。

\Needspace{4\baselineskip}\paragraph{選択肢の検討（査読待ち初稿）}\mbox{}\par

\noindent\textbf{1：条件に不一致} カルノー／\((T_H-T_L)/(T_H+T_L)\)。機関名は合うが効率の分母は吸熱源の絶対温度であり、温度の和ではない。\par

\noindent\textbf{2：条件に不一致} フェルマー／\((T_H-T_L)/(T_H+T_L)\)。機関名と分母の両方が違う。\par

\noindent\textbf{3：条件に不一致} カルノー／\(T_H/(T_H+T_L)\)。低温源が高温源に近づくと効率は0へ向かうべきであり、この式はその条件を満たさない。\par

\noindent\textbf{4：条件に不一致} フェルマー／\((T_H-T_L)/T_H\)。効率式は合うが機関名が違う。\par

\noindent\textbf{5：条件に一致} カルノー／\((T_H-T_L)/T_H\)。可逆なカルノー機関とその効率に一致する。\par

\paragraph{再利用する考え方} モデルの前提を明示し、第二法則が決める熱機関効率の上限から記述や数値を判定する。

\paragraph{共通理論} \hyperref[theory:thermo:carnot]{第\ref*{theory:thermo:carnot}節}, \hyperref[theory:thermo:entropy]{第\ref*{theory:thermo:entropy}節}

\paragraph{関連問題} \hyperref[q:h25-reg-applied-08]{H25-REG-APPLIED-08}

\par 難易度 1（校準初稿）、優先度未評価。

\subsection{平成29年度 問6}
\label{q:h29-reg-applied-06}

\noindent\textbf{原典：} sources/H29/17\_applied\_science.pdf、PDF物理2頁。

\textbf{公式正答：} 4（single）。

\paragraph{問題の要旨} ヒーターで1molの物質Aを一定外圧で全て気体にしたとき、物質Aが外部へ行う仕事を求める。

\Needspace{4\baselineskip}\paragraph{確認した条件・記述（独自要約）}\mbox{}\par

\noindent\textbf{系：} 断熱シリンダー内の1molの液体物質Aと、加熱用ヒーター。\par

\noindent\textbf{温度：} 物質Aの沸点373K\par

\noindent\textbf{外圧：} 0.1MPa一定\par

\noindent\textbf{仮定：} 液体体積を無視、気体状態の物質Aは理想気体。\par

\noindent\textbf{気体定数：} 8.3J/(mol・K)\par

\noindent\textbf{求める量：} 物質Aが外部へ行う仕事の正の大きさ。\par

\paragraph{問われる核心} 気化で増えた体積による定圧仕事

\paragraph{必要知識} 一定外圧の仕事、理想気体、液体体積無視、壁からの熱流と電気入力の区別。

\paragraph{考え方・正答を選ぶ理由（査読待ち初稿）} 外部への体積仕事は\(W_{\mathrm{out}}=p_{\mathrm{ext}}(V_g-V_l)\)である。液体体積を無視し、終状態の理想気体に\(p_{\mathrm{ext}}V_g=nRT\)を適用すると\(W_{\mathrm{out}}=1\times8.3\times373=3095.9\,\mathrm{J}\)である。最も近い3100Jの4を選ぶ。加熱の総エネルギーや潜熱を問う問題ではない。断熱壁でもヒーターから仕事としてエネルギーが入る。

\Needspace{4\baselineskip}\paragraph{選択肢の検討（査読待ち初稿）}\mbox{}\par

\noindent\textbf{1：条件に不一致} 0.023J。外圧×体積増加=nRTに対して桁又は値が合わない。\par

\noindent\textbf{2：条件に不一致} 830J。外圧×体積増加=nRTに対して桁又は値が合わない。\par

\noindent\textbf{3：条件に不一致} 2300J。外圧×体積増加=nRTに対して桁又は値が合わない。\par

\noindent\textbf{4：条件に一致} 3100J。理想蒸気のpV=\allowbreak{}nRTに一致する。\par

\noindent\textbf{5：条件に不一致} 8300J。外圧×体積増加=nRTに対して桁又は値が合わない。\par

\paragraph{再利用する考え方} モデルの前提を明示し、気化で増えた体積による定圧仕事から記述や数値を判定する。

\paragraph{共通理論} \hyperref[theory:thermo:firstlaw]{第\ref*{theory:thermo:firstlaw}節}, \hyperref[theory:thermo:workpaths]{第\ref*{theory:thermo:workpaths}節}

\par 参考：\cite{mit_thermo_firstlaw2008}。

\paragraph{関連問題} \hyperref[q:r06-reg-applied-06]{R06-REG-APPLIED-06}

\par 難易度 2（校準初稿）、優先度未評価。

\subsection{平成30年度 問7}
\label{q:h30-reg-applied-07}

\noindent\textbf{原典：} sources/H30/17\_applied\_science.pdf、PDF物理3頁。

\textbf{公式正答：} 1（single）。

\paragraph{問題の要旨} 室温での熱伝導率を4材料について大きい順に並べる。

\Needspace{4\baselineskip}\paragraph{確認した条件・記述（独自要約）}\mbox{}\par

\noindent\textbf{条件：} 室温で通常のバルク材料の代表的な熱伝導率を比較する。\par

\paragraph{問われる核心} 熱を運ぶ電子と格子振動・材料依存

\paragraph{必要知識} 熱伝導率と電気伝導率の区別、室温、バルク材料、純度と欠陥。

\paragraph{考え方・正答を選ぶ理由（査読待ち初稿）} ダイヤモンドの結晶格子は室温で熱をよく運び、高品質バルクでは銅より熱伝導率が大きい。通常の銅は鉄鋼より大きく、板ガラスは小さい。従ってダイヤモンド→銅→鉄鋼→板ガラスの1である。電気絶縁体なら熱伝導率も低い、という推論はダイヤモンドでは成立しない。結晶欠陥・粒径・純度・合金種で値は変わるので、全ての試料に同じ序列を断定しない。

\Needspace{4\baselineskip}\paragraph{選択肢の検討（査読待ち初稿）}\mbox{}\par

\noindent\textbf{1：条件に一致} ダイヤモンド→銅→鉄鋼→板ガラス。代表的な室温の大小関係に一致する。\par

\noindent\textbf{2：条件に不一致} ダイヤモンド→鉄鋼→銅→板ガラス。鉄鋼と銅が逆である。\par

\noindent\textbf{3：条件に不一致} 銅→鉄鋼→ダイヤモンド→板ガラス。高品質のバルクダイヤモンドを金属より下に置いている。\par

\noindent\textbf{4：条件に不一致} 銅→鉄鋼→板ガラス→ダイヤモンド。バルクダイヤモンドを最小にしている。\par

\noindent\textbf{5：条件に不一致} 銅→ダイヤモンド→鉄鋼→板ガラス。バルクダイヤモンドと銅が逆である。\par

\paragraph{再利用する考え方} モデルの前提を明示し、熱を運ぶ電子と格子振動・材料依存から記述や数値を判定する。

\paragraph{共通理論} \hyperref[theory:thermo:transport]{第\ref*{theory:thermo:transport}節}

\par 参考：\cite{aist_metal_thermal,wang_sun_diamond2024,agc_glass_thermal}。

\par 難易度 1（校準初稿）、優先度未評価。

\subsection{令和元年度 問6}
\label{q:r01-reg-applied-06}

\noindent\textbf{原典：} sources/R01/17\_applied\_science.pdf、PDF物理2頁。

\textbf{公式正答：} 1（single）。

\paragraph{問題の要旨} 熱伝導、計測、熱起電力の記述から最も不適切なものを選ぶ。

\Needspace{4\baselineskip}\paragraph{確認した条件・記述（独自要約）}\mbox{}\par

\noindent\textbf{方向：} 最も不適切。\par

\paragraph{問われる核心} 熱輸送を利用した圧力計測と熱電効果

\paragraph{必要知識} ピラニ計のフィラメント、フォノン、電子、熱電能、フーリエの法則。

\paragraph{考え方・正答を選ぶ理由（査読待ち初稿）} 1は計測の対象と手段を取り違えている。ピラニ計は気体の熱移動による加熱線の熱損失を利用して真空の圧力を測る。2〜5は各材料・線形応答の条件で妥当であり、吸着による圧力変化から熱伝導率を求めるという1を選ぶ。

\Needspace{4\baselineskip}\paragraph{選択肢の検討（査読待ち初稿）}\mbox{}\par

\noindent\textbf{1：誤り} ピラニ真空計は吸着による気体圧力の変化で熱伝導率を求める。ピラニ計は加熱フィラメントから気体へ逃げる熱の圧力依存性を、抵抗や加熱電力で読み圧力へ換算する。吸着を測って熱伝導率を求める説明ではない。\par

\noindent\textbf{2：正しい} 固体の熱伝導はフォノンによるエネルギー移動を伴う。結晶格子の振動励起がエネルギーを運ぶ。金属では伝導電子の寄与もあるので、唯一の担体とは考えない。\par

\noindent\textbf{3：通常の室温材料で正しい} Alの熱伝導率はFeより大きい。産総研の室温付近のデータではAlは226〜237、Feは72〜80.4W/(m・K)程度であり大小が合う。\par

\noindent\textbf{4：正しい} 異種金属の接合部の温度差は導線方向の電場を生む熱起電力に関係する。ゼーベック効果である。異種材料の温度依存する熱電能の差を利用する。\par

\noindent\textbf{5：線形等方近似で正しい} フーリエの法則では熱流密度が温度勾配の大きさに比例する。熱は温度が下がる方向へ流れるのでベクトルとしては温度勾配と逆向きである。\par

\paragraph{再利用する考え方} モデルの前提を明示し、熱輸送を利用した圧力計測と熱電効果から記述や数値を判定する。

\paragraph{共通理論} \hyperref[theory:thermo:transport]{第\ref*{theory:thermo:transport}節}, \hyperref[theory:thermo:measurement]{第\ref*{theory:thermo:measurement}節}

\par 参考：\cite{ulvac_pirani,aist_metal_thermal}。

\par 難易度 1（校準初稿）、優先度未評価。

\subsection{令和元年度再試験 問4}
\label{q:r01-ret-applied-04}

\noindent\textbf{原典：} sources/R01\_retest/17\_applied\_science.pdf、PDF物理2頁。

\textbf{公式正答：} 4（single）。

\paragraph{問題の要旨} 第一法則・第二法則と気体の体積変化を判定し、最も不適切な記述を選ぶ。

\Needspace{4\baselineskip}\paragraph{確認した条件・記述（独自要約）}\mbox{}\par

\noindent\textbf{方向：} 最も不適切。\par

\noindent\textbf{原典の範囲：} 一定質量の気体。理想気体という語は明記されていない。\par

\paragraph{問われる核心} 気体法則とエネルギー保存の対応

\paragraph{必要知識} 第一法則、第二法則、絶対温度と理想気体近似。

\paragraph{考え方・正答を選ぶ理由（査読待ち初稿）} 1と5は第一法則、2は第二法則、3は理想気体近似の定圧関係である。4の等温体積が圧力に比例するという主張は\(V=nRT/p\)に反するので選ぶ。類似問題と肢の位置が異なるため、以前の番号を転用しない。

\Needspace{4\baselineskip}\paragraph{選択肢の検討（査読待ち初稿）}\mbox{}\par

\noindent\textbf{1：正しい} 外部からの熱と系へされた仕事の和が内部エネルギー増加に等しい。熱を流入正、仕事を系へ入力正とすれば\(\Delta U=Q+W_{\mathrm{on}}\)である。系全体の運動・位置エネルギー変化はこの閉じた熱力学系では考えない。\par

\noindent\textbf{2：正しい} 熱は高温から低温へ移り、自然に低温から高温へは移らない。第二法則の巨視的な受動的熱移動の方向である。冷凍機は外部仕事を使うので反例ではない。\par

\noindent\textbf{3：理想気体近似で正しい} 圧力一定・一定質量の気体の体積は絶対温度に比例する。理想気体又はその近似が成り立つ範囲では固定\(n,p\)により\(V=(nR/p)T\)である。\par

\noindent\textbf{4：誤り} 温度一定・一定質量の気体の体積は圧力に比例する。理想気体なら\(pV=nRT\)であり、固定\(n,T\)では\(V=nRT/p\)と反比例する。比例という記述が誤りである。\par

\noindent\textbf{5：正しい} エネルギー変換で関与する全エネルギーの和が保存される。第一法則のエネルギー保存の内容である。対象系から外へ出るエネルギーを数え落とさない。\par

\paragraph{再利用する考え方} モデルの前提を明示し、気体法則とエネルギー保存の対応から記述や数値を判定する。

\paragraph{共通理論} \hyperref[theory:thermo:firstlaw]{第\ref*{theory:thermo:firstlaw}節}, \hyperref[theory:thermo:ideal]{第\ref*{theory:thermo:ideal}節}, \hyperref[theory:thermo:entropy]{第\ref*{theory:thermo:entropy}節}

\par 参考：\cite{mit_thermo_firstlaw2008}。

\paragraph{関連問題} \hyperref[q:h26-reg-applied-08]{H26-REG-APPLIED-08}

\par 難易度 1（校準初稿）、優先度未評価。

\subsection{令和2年度 問6}
\label{q:r02-reg-applied-06}

\noindent\textbf{原典：} sources/R02/17\_applied\_science.pdf、PDF物理2頁。

\textbf{公式正答：} 4（single）。

\paragraph{問題の要旨} 理想気体の定圧・定積4段階変化のうち、第3段階で外部からされる仕事を求める。

\Needspace{4\baselineskip}\paragraph{確認した条件・記述（独自要約）}\mbox{}\par

\noindent\textbf{初期：} \(p_0=2.0\times10^5\,\mathrm{Pa},V_0=100\,\mathrm{m}^3,T_0=27\,{}^\circ\mathrm{C}\)。\par

\noindent\textbf{第1段階：} 定圧で327°Cへ加熱。\par

\noindent\textbf{第2段階：} 定積で27°Cへ冷却。\par

\noindent\textbf{第3段階：} 定圧で−123°Cへ冷却。\par

\noindent\textbf{第4段階：} 定積で27°Cへ加熱。\par

\noindent\textbf{求める量：} 第3段階で気体へされる仕事。\par

\paragraph{問われる核心} 定圧仕事の経路依存と絶対温度の換算

\paragraph{必要知識} 一定気体量のpV/T保存、定積p∝T、定圧V∝T、圧縮仕事の正符号。

\paragraph{考え方・正答を選ぶ理由（査読待ち初稿）} 温度を約300,600,300,150Kとする。第1段階で体積は200m\textsuperscript{3}、第2段階で圧力は\(1.0\times10^5\,\mathrm{Pa}\)になる。第3段階は150Kへの定圧冷却なので体積は100m\textsuperscript{3}となる。気体へされる仕事は\(W_{\mathrm{on}}=p(V_{\mathrm{start}}-V_{\mathrm{end}})=10^5(200-100)=10^7\,\mathrm{J}\)、第4候補である。体積は100m\textsuperscript{3}であり100mLではない。

\Needspace{4\baselineskip}\paragraph{選択肢の検討（査読待ち初稿）}\mbox{}\par

\noindent\textbf{1：条件に不一致} \(1\times10^6\,\mathrm{J}\)。圧縮仕事の計算値より1桁小さい。\par

\noindent\textbf{2：条件に不一致} \(2\times10^6\,\mathrm{J}\)。圧縮の温度幅と気体量の両方を反映していない。\par

\noindent\textbf{3：条件に不一致} \(5\times10^6\,\mathrm{J}\)。第3段階の体積差100m\textsuperscript{3}と圧力10\textasciicircum{}5Paの積の半分である。\par

\noindent\textbf{4：条件に一致} \(1\times10^7\,\mathrm{J}\)。圧力10\textasciicircum{}5Pa×減少体積100m\textsuperscript{3}に一致する。\par

\noindent\textbf{5：条件に不一致} \(2\times10^7\,\mathrm{J}\)。第3段階へ初期圧力2×10\textasciicircum{}5Paをそのまま使った値である。\par

\paragraph{再利用する考え方} モデルの前提を明示し、定圧仕事の経路依存と絶対温度の換算から記述や数値を判定する。

\paragraph{共通理論} \hyperref[theory:thermo:workpaths]{第\ref*{theory:thermo:workpaths}節}, \hyperref[theory:thermo:ideal]{第\ref*{theory:thermo:ideal}節}

\paragraph{関連問題} \hyperref[q:r07-reg-applied-06]{R07-REG-APPLIED-06}

\par 難易度 2（校準初稿）、優先度未評価。

\subsection{令和3年度 問7}
\label{q:r03-reg-applied-07}

\noindent\textbf{原典：} sources/R03/17\_applied\_science.pdf、PDF物理3頁。

\textbf{公式正答：} 1（single）。

\par\textbf{検討未確定：} 1の反比例は明確な誤りで独立選択1。2の室温以上で3Rという一律範囲は高デバイ温度物質で成立しない。3の低温は回転励起・振動凍結の中間範囲、5の100Kも材料依存である。原典の限定不足をAstraの解釈査読対象とし、部分校正を保持する。

\paragraph{問題の要旨} 電子・格子・気体の比熱と相転移の5記述から最も不適切なものを選ぶ。

\Needspace{4\baselineskip}\paragraph{確認した条件・記述（独自要約）}\mbox{}\par

\noindent\textbf{方向：} 最も不適切。\par

\noindent\textbf{記号：} Rは気体定数。\par

\noindent\textbf{原典の省略：} 格子比熱は「室温又はそれ以上」、二原子分子は「比較的低い温度」、固体の熱容量差は「100K以下」と記す。物質や固有温度は指定しない。\par

\paragraph{問われる核心} 量子励起の温度尺度と比熱

\paragraph{必要知識} 電子比熱γT、デバイ温度、二原子分子の自由度、相転移、熱膨張と圧縮率。

\paragraph{考え方・正答を選ぶ理由（査読待ち初稿）} 状態密度\(D(E_F)\)と温度Tへの比例が電子比熱の特徴であり、反比例とする1を最も明確な誤りとして選ぶ。2は高温の古典極限、3は回転のみ励起、4は転移の比熱異常、5は小さな熱膨張に対応する。原典の温度範囲は物質固有の励起温度を記さないため、2・3・5を全材料・全低温に成り立つ法則として説明しない。

\Needspace{4\baselineskip}\paragraph{選択肢の検討（査読待ち初稿）}\mbox{}\par

\noindent\textbf{1：誤り} 電子比熱はフェルミ準位の状態密度と温度に反比例する。通常の縮退電子では\(C_e=(\pi^2/3)k_B^2D(E_F)T\)で、状態密度と温度へ比例し、反比例ではない。\par

\noindent\textbf{2：デバイ温度に比べ高温なら近似的に正しい／原典の範囲に疑義} 格子モル比熱は室温又はそれ以上の高温で3Rとなる。デュロン＝プティ則は\(T\gg\Theta_D\)の古典極限である。室温というだけでは高いデバイ温度を持つ物質には適用できない。\par

\noindent\textbf{3：回転励起・振動凍結の範囲で正しい} 二原子分子の定圧モル比熱は比較的低温で定積の7/5倍である。並進3・回転2自由度なら\(C_{V,m}=5R/2\)、\(C_{p,m}=7R/2\)で比は7/5である。回転まで凍結する極低温や振動の励起される高温では違う。\par

\noindent\textbf{4：二次相転移の比熱異常として妥当} 二次相転移では不連続な変化を伴う異常比熱が現れる。ギブズエネルギーの二階微分に異常が出る。有限跳びや発散があり、全ての連続相転移が同じ形の跳びを持つとは限らない。\par

\noindent\textbf{5：通常の固体での低温近似／物質依存} 固体の定積・定圧モル比熱の差は100K以下で無視できるほど小さい。\(C_p-C_V=TV\alpha_V^2/\kappa_T\)であり低温では一般に小さくなる。ただし100Kという数値だけから全物質で無視可能と断定はできない。\par

\paragraph{再利用する考え方} モデルの前提を明示し、量子励起の温度尺度と比熱から記述や数値を判定する。

\paragraph{共通理論} \hyperref[theory:thermo:degrees]{第\ref*{theory:thermo:degrees}節}, \hyperref[theory:thermo:solidcapacity]{第\ref*{theory:thermo:solidcapacity}節}, \hyperref[theory:thermo:phase]{第\ref*{theory:thermo:phase}節}

\par 参考：\cite{mit_electronic_heat2012,mit_debye2005}。

\par 難易度 2（校準初稿）、優先度未評価。

\subsection{令和4年度 問17}
\label{q:r04-reg-applied-17}

\noindent\textbf{原典：} sources/R04/17\_applied\_science.pdf、PDF物理6頁。

\textbf{公式正答：} 5（single）。

\paragraph{問題の要旨} 実在気体の圧縮係数と分子間力の記述から最も不適切なものを選ぶ。

\Needspace{4\baselineskip}\paragraph{確認した条件・記述（独自要約）}\mbox{}\par

\noindent\textbf{方向：} 最も不適切。\par

\noindent\textbf{定義：} 同温・同圧で実在気体のモル体積を理想気体のモル体積で割った比を圧縮係数とする。\par

\paragraph{問われる核心} 圧縮係数による理想気体からのずれの表現

\paragraph{必要知識} モル体積、理想気体極限、引力と排除体積、ファンデルワールス式。

\paragraph{考え方・正答を選ぶ理由（査読待ち初稿）} 圧縮係数を\(Z=pV_m/(RT)\)と定義する。引力は同温同密度で圧力を下げ、分子の大きさは圧力を上げる方向へ効く。従ってZは1の上下いずれにもなり得て、温度・圧力にかかわらず1より大きいという5は誤りである。

\Needspace{4\baselineskip}\paragraph{選択肢の検討（査読待ち初稿）}\mbox{}\par

\noindent\textbf{1：一般に正しい} 高圧では理想気体の法則から外れる。分子の占有体積や相互作用が無視できず、低密度の理想近似から外れやすい。\par

\noindent\textbf{2：低密度極限で正しい} 低圧では分子間力の影響が小さい。密度が下がると分子間相互作用による圧力への補正は小さくなる。\par

\noindent\textbf{3：正しい} 圧縮係数は同温同圧の実在／理想のモル体積比である。理想モル体積は\(RT/p\)なので比は\(Z=pV_m/(RT)\)である。\par

\noindent\textbf{4：正しい} ファンデルワールス式は引力と分子の大きさを考慮する。\((p+a/V_m^2)(V_m-b)=RT\)でaが引力、bが排除体積の補正を表す。\par

\noindent\textbf{5：誤り} 圧縮係数は温度や圧力に関係なく1より大きい。引力が優勢なら\(Z<1\)、排除体積が優勢なら\(Z>1\)になり得る。低密度では1へ近づく。\par

\paragraph{再利用する考え方} モデルの前提を明示し、圧縮係数による理想気体からのずれの表現から記述や数値を判定する。

\paragraph{共通理論} \hyperref[theory:thermo:realgas]{第\ref*{theory:thermo:realgas}節}

\par 難易度 1（校準初稿）、優先度未評価。

\subsection{令和4年度 問6}
\label{q:r04-reg-applied-06}

\noindent\textbf{原典：} sources/R04/17\_applied\_science.pdf、PDF物理2頁。

\textbf{公式正答：} 1（single）。

\paragraph{問題の要旨} カルノーサイクルの吸熱と2熱源の温度から、1サイクルの低温熱源への放熱を求める。

\Needspace{4\baselineskip}\paragraph{確認した条件・記述（独自要約）}\mbox{}\par

\noindent\textbf{高温熱源：} 1000K\par

\noindent\textbf{低温熱源：} 800K\par

\noindent\textbf{高温側吸熱：} 500J\par

\noindent\textbf{求める量：} 低温源へ放出する熱量の正の大きさ。\par

\paragraph{問われる核心} カルノーサイクルの可逆エントロピー収支

\paragraph{必要知識} 絶対温度、吸熱と放熱の正の大きさ、仕事と熱収支。

\paragraph{考え方・正答を選ぶ理由（査読待ち初稿）} 可逆サイクルのエントロピー変化は0なので\(Q_L=Q_HT_L/T_H\)である。数値を代入すると\(500\times800/1000=400\,\mathrm{J}\)となり、1を選ぶ。取り出した仕事は吸熱と放熱の差であり、求める放熱と混同しない。

\Needspace{4\baselineskip}\paragraph{選択肢の検討（査読待ち初稿）}\mbox{}\par

\noindent\textbf{1：条件に一致} 400J。吸熱に低温／高温の温度比を掛けた放熱に一致する。\par

\noindent\textbf{2：条件に不一致} 450J。可逆性から得る低温側放熱400Jと一致しない。\par

\noindent\textbf{3：条件に不一致} 500J。可逆性から得る低温側放熱400Jと一致しない。\par

\noindent\textbf{4：条件に不一致} 550J。可逆性から得る低温側放熱400Jと一致しない。\par

\noindent\textbf{5：条件に不一致} 600J。可逆性から得る低温側放熱400Jと一致しない。\par

\paragraph{再利用する考え方} モデルの前提を明示し、カルノーサイクルの可逆エントロピー収支から記述や数値を判定する。

\paragraph{共通理論} \hyperref[theory:thermo:carnot]{第\ref*{theory:thermo:carnot}節}

\paragraph{関連問題} \hyperref[q:h27-reg-applied-08]{H27-REG-APPLIED-08}

\par 難易度 2（校準初稿）、優先度未評価。

\subsection{令和4年度 問7}
\label{q:r04-reg-applied-07}

\noindent\textbf{原典：} sources/R04/17\_applied\_science.pdf、PDF物理2頁。

\textbf{公式正答：} 3（single）。

\paragraph{問題の要旨} 熱電対、金属の伝導、多孔質、放射について最も不適切な記述を選ぶ。

\Needspace{4\baselineskip}\paragraph{確認した条件・記述（独自要約）}\mbox{}\par

\noindent\textbf{方向：} 最も不適切。\par

\paragraph{問われる核心} 熱伝導と電気伝導の関係の適用範囲

\paragraph{必要知識} 熱電対、ウィーデマン＝フランツ則、フォノン、多孔質の有効伝導、電磁放射。

\paragraph{考え方・正答を選ぶ理由（査読待ち初稿）} 典型的なバルクダイヤモンドは室温で銀より熱伝導率が大きいので3が不適切である。他の肢は熱電対、金属の電子輸送の近似、多孔質断熱、熱放射に対応する。熱伝導の担体は電子だけではなく、ダイヤモンドでは格子振動が重要である。

\Needspace{4\baselineskip}\paragraph{選択肢の検討（査読待ち初稿）}\mbox{}\par

\noindent\textbf{1：正しい} 異種金属又は合金2種で温度計測できる。熱電対は2種の材料の熱電能の差から温度差を電圧へ変換する。\par

\noindent\textbf{2：通常の金属の電子輸送近似で正しい} 金属の熱伝導率／電気伝導率は極低温を除き絶対温度へ比例する。ウィーデマン＝フランツ則の電子輸送の関係である。散乱の機構や格子熱伝導の大きさによって厳密には変わる。\par

\noindent\textbf{3：通常のバルク比較で誤り} ダイヤモンドの熱伝導率は銀より小さい。高品質バルクダイヤモンドは銀より大きい。金属の電気伝導の序列から熱伝導の序列を全材料へ拡張できない。\par

\noindent\textbf{4：一般的な断熱材料で正しい} 空孔にある気体は熱を運びにくく、多孔質材料は熱伝導率が低い。固体の連続した伝導経路を減らし気体の低い伝導率を利用する。ただし大きな孔の対流や高温の放射は別に評価する。\par

\noindent\textbf{5：正しい} 熱放射は電磁波でのエネルギー移動である。媒体の巨視的移動が必要な対流とは区別する。\par

\paragraph{再利用する考え方} モデルの前提を明示し、熱伝導と電気伝導の関係の適用範囲から記述や数値を判定する。

\paragraph{共通理論} \hyperref[theory:thermo:transport]{第\ref*{theory:thermo:transport}節}, \hyperref[theory:thermo:measurement]{第\ref*{theory:thermo:measurement}節}

\par 参考：\cite{aist_metal_thermal,wang_sun_diamond2024}。

\par 難易度 1（校準初稿）、優先度未評価。

\subsection{令和5年度 問6}
\label{q:r05-reg-applied-06}

\noindent\textbf{原典：} sources/R05/17\_applied\_science.pdf、PDF物理3頁。

\textbf{公式正答：} 3（single）。

\paragraph{問題の要旨} 相転移の次数、熱力学量、強磁性、過冷却の記述から最も不適切なものを選ぶ。

\Needspace{4\baselineskip}\paragraph{確認した条件・記述（独自要約）}\mbox{}\par

\noindent\textbf{方向：} 最も不適切。\par

\noindent\textbf{対象：} 融解・気化と連続相転移の通常の分類。\par

\paragraph{問われる核心} 自由エネルギーの微分から見る相転移と準安定状態

\paragraph{必要知識} ギブズエネルギー、潜熱、エントロピー、熱容量、秩序変数、核生成。

\paragraph{考え方・正答を選ぶ理由（査読待ち初稿）} 相転移の次数は自由エネルギーの微分の性質で整理する。一次ではエントロピーなど一階微分が跳び、二次・連続転移では比熱など二階微分が異常を示す。熱容量は連続であると一律に述べる3が不適切である。過冷却は平衡相図と変化の速度を区別する例である。

\Needspace{4\baselineskip}\paragraph{選択肢の検討（査読待ち初稿）}\mbox{}\par

\noindent\textbf{1：臨界点から離れた相共存で正しい} 融解と気化は一次相転移である。融解・気化は潜熱を伴う通常の一次転移である。液気臨界点で相の区別がなくなる場合は別である。\par

\noindent\textbf{2：正しい} 一次相転移ではエントロピーが不連続に変わる。潜熱Lと転移温度Tにより\(\Delta S=L/T\)でエントロピーの跳びがある。\par

\noindent\textbf{3：誤り} 二次相転移では熱容量は連続に変わる。二次転移ではギブズエネルギーの一階微分は連続でも、比熱のような二階微分に跳びや発散が出る。\par

\noindent\textbf{4：通常の零磁場のキュリー転移として正しい} 鉄の強磁性から常磁性への転移は二次相転移である。磁化という秩序変数が零磁場で連続に消える転移であり、比熱には異常が現れる。\par

\noindent\textbf{5：正しい} 液体や気体を平衡転移温度より冷やしても相変化しない過冷却がある。新相の核生成の障壁などで準安定状態が続く。平衡の共存温度そのものが変わったとの意味ではない。\par

\paragraph{再利用する考え方} モデルの前提を明示し、自由エネルギーの微分から見る相転移と準安定状態から記述や数値を判定する。

\paragraph{共通理論} \hyperref[theory:thermo:phase]{第\ref*{theory:thermo:phase}節}

\par 難易度 1（校準初稿）、優先度未評価。

\subsection{令和6年度 問6}
\label{q:r06-reg-applied-06}

\noindent\textbf{原典：} sources/R06/17\_applied\_science.pdf、PDF物理2頁。

\textbf{公式正答：} 3（single）。

\paragraph{問題の要旨} ヒーターで1molの物質Aを一定外圧で全て気体にしたとき、物質Aが外部へ行う仕事を求める。

\Needspace{4\baselineskip}\paragraph{確認した条件・記述（独自要約）}\mbox{}\par

\noindent\textbf{系：} 断熱シリンダー内の1molの液体物質Aと、加熱用ヒーター。\par

\noindent\textbf{温度：} 物質Aの沸点506K\par

\noindent\textbf{外圧：} 0.1MPa一定\par

\noindent\textbf{仮定：} 液体体積を無視、気体状態の物質Aは理想気体。\par

\noindent\textbf{気体定数：} 8.3J/(mol・K)\par

\noindent\textbf{求める量：} 物質Aが外部へ行う仕事の正の大きさ。\par

\paragraph{問われる核心} 気化で増えた体積による定圧仕事

\paragraph{必要知識} 一定外圧の仕事、理想気体、液体体積無視、壁からの熱流と電気入力の区別。

\paragraph{考え方・正答を選ぶ理由（査読待ち初稿）} 外部への体積仕事は\(W_{\mathrm{out}}=p_{\mathrm{ext}}(V_g-V_l)\)である。液体体積を無視し、終状態の理想気体に\(p_{\mathrm{ext}}V_g=nRT\)を適用すると\(W_{\mathrm{out}}=1\times8.3\times506=4199.8\,\mathrm{J}\)である。最も近い4200Jの3を選ぶ。加熱の総エネルギーや潜熱を問う問題ではない。断熱壁でもヒーターから仕事としてエネルギーが入る。

\Needspace{4\baselineskip}\paragraph{選択肢の検討（査読待ち初稿）}\mbox{}\par

\noindent\textbf{1：条件に不一致} 830J。外圧×体積増加=nRTに対して桁又は値が合わない。\par

\noindent\textbf{2：条件に不一致} 2240J。外圧×体積増加=nRTに対して桁又は値が合わない。\par

\noindent\textbf{3：条件に一致} 4200J。理想蒸気のpV=\allowbreak{}nRTに一致する。\par

\noindent\textbf{4：条件に不一致} 8300J。外圧×体積増加=nRTに対して桁又は値が合わない。\par

\noindent\textbf{5：条件に不一致} 22400J。外圧×体積増加=nRTに対して桁又は値が合わない。\par

\paragraph{再利用する考え方} モデルの前提を明示し、気化で増えた体積による定圧仕事から記述や数値を判定する。

\paragraph{共通理論} \hyperref[theory:thermo:firstlaw]{第\ref*{theory:thermo:firstlaw}節}, \hyperref[theory:thermo:workpaths]{第\ref*{theory:thermo:workpaths}節}

\par 参考：\cite{mit_thermo_firstlaw2008}。

\paragraph{関連問題} \hyperref[q:h29-reg-applied-06]{H29-REG-APPLIED-06}

\par 難易度 2（校準初稿）、優先度未評価。

\subsection{令和7年度 問6}
\label{q:r07-reg-applied-06}

\noindent\textbf{原典：} sources/R07/17\_applied\_science.pdf、PDF物理3頁。

\textbf{公式正答：} 1（single）。

\paragraph{問題の要旨} 理想気体を定圧・定積で4段階に変化させ、第3段階の外部からの仕事を求める。

\Needspace{4\baselineskip}\paragraph{確認した条件・記述（独自要約）}\mbox{}\par

\noindent\textbf{初期：} \(p_0=5.0\times10^3\,\mathrm{Pa},V_0=120\,\mathrm{m}^3,T_0=27\,{}^\circ\mathrm{C}\)。\par

\noindent\textbf{第1段階：} 定圧で227°Cへ加熱。\par

\noindent\textbf{第2段階：} 定積で27°Cへ冷却。\par

\noindent\textbf{第3段階：} 定圧で−73°Cへ冷却。\par

\noindent\textbf{第4段階：} 定積で27°Cへ加熱。\par

\noindent\textbf{求める量：} 第3段階で気体へされる仕事。\par

\paragraph{問われる核心} 状態方程式から気体量を消去する定圧仕事

\paragraph{必要知識} pV=\allowbreak{}nRT、絶対温度、定圧仕事−nRΔT、仕事の向き。

\paragraph{考え方・正答を選ぶ理由（査読待ち初稿）} 絶対温度を約300,500,300,200Kとする。気体量を\(nR=p_0V_0/T_0=(5.0\times10^3)120/300=2000\,\mathrm{J/K}\)と置く。第3段階の定圧圧縮仕事は\(W_{\mathrm{on}}=-nR\Delta T=2000(300-200)=2.0\times10^5\,\mathrm{J}\)、第1候補である。途中状態はV=\allowbreak{}200m\textsuperscript{3}、p=\allowbreak{}3000Pa、終体積400/3m\textsuperscript{3}であり、pΔVからも同じ値になる。

\Needspace{4\baselineskip}\paragraph{選択肢の検討（査読待ち初稿）}\mbox{}\par

\noindent\textbf{1：条件に一致} \(2\times10^5\,\mathrm{J}\)。nR=初期pV/Tと圧縮温度幅100Kの積に一致する。\par

\noindent\textbf{2：条件に不一致} \(5\times10^5\,\mathrm{J}\)。原典の圧力の指数と圧縮温度幅を用いた結果と合わない。\par

\noindent\textbf{3：条件に不一致} \(1\times10^6\,\mathrm{J}\)。計算値の5倍であり、段階の温度と圧力を区別していない。\par

\noindent\textbf{4：条件に不一致} \(2\times10^6\,\mathrm{J}\)。計算値より1桁大きい。初期圧力の指数は3である。\par

\noindent\textbf{5：条件に不一致} \(5\times10^6\,\mathrm{J}\)。圧力の指数又は体積変化の計算が一致しない。\par

\paragraph{再利用する考え方} モデルの前提を明示し、状態方程式から気体量を消去する定圧仕事から記述や数値を判定する。

\paragraph{共通理論} \hyperref[theory:thermo:workpaths]{第\ref*{theory:thermo:workpaths}節}, \hyperref[theory:thermo:ideal]{第\ref*{theory:thermo:ideal}節}

\paragraph{関連問題} \hyperref[q:r02-reg-applied-06]{R02-REG-APPLIED-06}

\par 難易度 2（校準初稿）、優先度未評価。

\section{相平衡・溶液}

\subsection{平成25年度 問9}
\label{q:h25-reg-applied-09}

\noindent\textbf{原典：} sources/H25/17\_applied\_science.pdf、PDF物理5頁。

\textbf{公式正答：} 4（single）。

\par\textbf{検討未確定：} 原典の断熱膨張の肢は外部仕事又は可逆性を明記しない。理想気体の断熱自由膨張は温度不変であり、この肢も無条件には正しくない。熱容量の逆転肢を最も明確な誤りとして独立選択したが、原典の省略を追加条件で書き換えず部分校正としてAstraへ送る。

\paragraph{問題の要旨} 理想気体と溶液・相平衡の記述から、最も不適切なものを選ぶ。

\Needspace{4\baselineskip}\paragraph{確認した条件・記述（独自要約）}\mbox{}\par

\noindent\textbf{方向：} 最も不適切。\par

\noindent\textbf{前提の確認：} 断熱膨張の肢に外部仕事又は可逆性の指定はない。\par

\paragraph{問われる核心} 熱容量と膨張仕事・相平衡の成立条件

\paragraph{必要知識} 理想気体の熱容量差、断熱と自由膨張、液気平衡、混合エントロピー、ラウールの法則。

\paragraph{考え方・正答を選ぶ理由（査読待ち初稿）} 同じ理想気体を同じだけ昇温させると定積では\(Q_V=nC_{V,m}\Delta T\)、定圧では\(Q_p=nC_{p,m}\Delta T\)を要し、\(C_{p,m}-C_{V,m}=R>0\)である。定積の方が熱量が大きいという肢4を選ぶ。断熱膨張の冷却は正の外部仕事を行う場合に限り、自由膨張という反例があるので解釈疑義を保持する。

\Needspace{4\baselineskip}\paragraph{選択肢の検討（査読待ち初稿）}\mbox{}\par

\noindent\textbf{1：通常の液気平衡で正しい} 高圧ほど純液体の沸点が上がる。液気平衡で\(dp/dT=L/(T\Delta V)>0\)なので外圧上昇で沸点が上がる。臨界点から離れた通常の沸騰を考える。\par

\noindent\textbf{2：仕事を外へ行う膨張なら正しい／前提省略} 理想気体を断熱膨張させると温度が下がる。断熱では\(\Delta U=-W_{\mathrm{out}}\)。正の仕事を外へ行う理想気体は冷えるが、真空への自由膨張なら\(W=0\)で温度は変わらない。膨張だけで冷却は保証されない。\par

\noindent\textbf{3：理想的な異種気体で正しい} 同温・同圧で異種気体の混合が自発的に起こる。非反応の異種理想気体の混合では\(\Delta G=nRT\sum_i x_i\ln x_i<0\)となる。\par

\noindent\textbf{4：誤り} 理想気体の同じ温度上昇には定積の方が定圧より多い熱が必要。定圧では膨張仕事も必要であり、理想気体のモル熱容量は定圧が定積よりR大きい。定積の方が多いという記述は逆である。\par

\noindent\textbf{5：理想溶液で正しい} 理想溶液の各成分の蒸気圧は純成分の蒸気圧×液相モル分率である。ラウールの法則\(p_i=x_i p_i^*\)である。実在溶液の全組成範囲へ無条件に使う法則ではない。\par

\paragraph{再利用する考え方} モデルの前提を明示し、熱容量と膨張仕事・相平衡の成立条件から記述や数値を判定する。

\paragraph{共通理論} \hyperref[theory:thermo:capacity]{第\ref*{theory:thermo:capacity}節}, \hyperref[theory:thermo:adiabatic]{第\ref*{theory:thermo:adiabatic}節}, \hyperref[theory:thermo:phase]{第\ref*{theory:thermo:phase}節}, \hyperref[theory:thermo:mixing]{第\ref*{theory:thermo:mixing}節}

\par 参考：\cite{mit_thermo_adiabatic2008,mit_clapeyron2008,mit_solutions2008}。

\paragraph{関連問題} \hyperref[q:h30-reg-applied-06]{H30-REG-APPLIED-06}

\par 難易度 2（校準初稿）、優先度未評価。

\subsection{平成30年度 問6}
\label{q:h30-reg-applied-06}

\noindent\textbf{原典：} sources/H30/17\_applied\_science.pdf、PDF物理2頁。

\textbf{公式正答：} 5（single）。

\par\textbf{検討未確定：} 原典の断熱膨張の肢は外部仕事又は可逆性を明記しない。理想気体の断熱自由膨張は温度不変であり、この肢も無条件には正しくない。熱容量の逆転肢を最も明確な誤りとして独立選択したが、原典の省略を追加条件で書き換えず部分校正としてAstraへ送る。

\paragraph{問題の要旨} 理想気体と溶液・相平衡の記述から、最も不適切なものを選ぶ。

\Needspace{4\baselineskip}\paragraph{確認した条件・記述（独自要約）}\mbox{}\par

\noindent\textbf{方向：} 最も不適切。\par

\noindent\textbf{前提の確認：} 断熱膨張の肢に外部仕事又は可逆性の指定はない。\par

\paragraph{問われる核心} 熱容量と膨張仕事・相平衡の成立条件

\paragraph{必要知識} 理想気体の熱容量差、断熱と自由膨張、液気平衡、混合エントロピー、ラウールの法則。

\paragraph{考え方・正答を選ぶ理由（査読待ち初稿）} 同じ理想気体を同じだけ昇温させると定積では\(Q_V=nC_{V,m}\Delta T\)、定圧では\(Q_p=nC_{p,m}\Delta T\)を要し、\(C_{p,m}-C_{V,m}=R>0\)である。定積の方が熱量が大きいという肢5を選ぶ。断熱膨張の冷却は正の外部仕事を行う場合に限り、自由膨張という反例があるので解釈疑義を保持する。

\Needspace{4\baselineskip}\paragraph{選択肢の検討（査読待ち初稿）}\mbox{}\par

\noindent\textbf{1：仕事を外へ行う膨張なら正しい／前提省略} 理想気体を断熱膨張させると温度が下がる。断熱では\(\Delta U=-W_{\mathrm{out}}\)。正の仕事を外へ行う理想気体は冷えるが、真空への自由膨張なら\(W=0\)で温度は変わらない。膨張だけで冷却は保証されない。\par

\noindent\textbf{2：理想的な異種気体で正しい} 同温・同圧で異種気体の混合が自発的に起こる。非反応の異種理想気体の混合では\(\Delta G=nRT\sum_i x_i\ln x_i<0\)となる。\par

\noindent\textbf{3：理想溶液で正しい} 理想溶液の各成分の蒸気圧は純成分の蒸気圧×液相モル分率である。ラウールの法則\(p_i=x_i p_i^*\)である。実在溶液の全組成範囲へ無条件に使う法則ではない。\par

\noindent\textbf{4：通常の液気平衡で正しい} 高圧ほど純液体の沸点が上がる。液気平衡で\(dp/dT=L/(T\Delta V)>0\)なので外圧上昇で沸点が上がる。臨界点から離れた通常の沸騰を考える。\par

\noindent\textbf{5：誤り} 理想気体の同じ温度上昇には定積の方が定圧より多い熱が必要。定圧では膨張仕事も必要であり、理想気体のモル熱容量は定圧が定積よりR大きい。定積の方が多いという記述は逆である。\par

\paragraph{再利用する考え方} モデルの前提を明示し、熱容量と膨張仕事・相平衡の成立条件から記述や数値を判定する。

\paragraph{共通理論} \hyperref[theory:thermo:capacity]{第\ref*{theory:thermo:capacity}節}, \hyperref[theory:thermo:adiabatic]{第\ref*{theory:thermo:adiabatic}節}, \hyperref[theory:thermo:phase]{第\ref*{theory:thermo:phase}節}, \hyperref[theory:thermo:mixing]{第\ref*{theory:thermo:mixing}節}

\par 参考：\cite{mit_thermo_adiabatic2008,mit_clapeyron2008,mit_solutions2008}。

\paragraph{関連問題} \hyperref[q:h25-reg-applied-09]{H25-REG-APPLIED-09}

\par 難易度 2（校準初稿）、優先度未評価。

\subsection{令和4年度 問11}
\label{q:r04-reg-applied-11}

\noindent\textbf{原典：} sources/R04/17\_applied\_science.pdf、PDF物理4頁。

\textbf{公式正答：} 2（single）。

\paragraph{問題の要旨} 水の三重点を既知として、富士山頂の大気圧下の沸点を求めるための理論式を選ぶ。

\Needspace{4\baselineskip}\paragraph{確認した条件・記述（独自要約）}\mbox{}\par

\noindent\textbf{既知のもの：} 水の三重点の温度と蒸気圧、山頂の大気圧。\par

\noindent\textbf{求める対象：} 純水の液気共存温度。数値の沸点を答える問題ではない。\par

\paragraph{問われる核心} 相共存曲線と圧力による沸点の変化

\paragraph{必要知識} 相平衡の化学ポテンシャル一致、潜熱、蒸気圧、理想蒸気近似。

\paragraph{考え方・正答を選ぶ理由（査読待ち初稿）} 沸点では蒸気圧が外圧に等しく、液気共存曲線をたどればよい。関係式は\(dp/dT=L/[T(V_g-V_l)]\)で、蒸気を理想気体とし液体体積を無視すると\(d\ln p/dT=L_m/(RT^2)\)となる。クラウジウス＝クラペイロンの式の2を選ぶ。有限区間の数値を求めるには潜熱やその温度依存なども必要であり、指定されない山頂の沸点を補わない。

\Needspace{4\baselineskip}\paragraph{選択肢の検討（査読待ち初稿）}\mbox{}\par

\noindent\textbf{1：条件に不一致} ボルンの式。イオンの静電的な溶媒和などを扱う式で、直接の液気共存曲線ではない。\par

\noindent\textbf{2：条件に一致} クラウジウス＝クラペイロンの式。液気共存の圧力と温度の関係を与える。\par

\noindent\textbf{3：条件に不一致} ファントホッフの式。平衡定数の温度依存を扱うが、この液気平衡を直接記す名称としてはクラウジウス＝クラペイロンが適切である。\par

\noindent\textbf{4：条件に不一致} プランクの式。黒体放射のスペクトルを扱う。\par

\noindent\textbf{5：条件に不一致} レイリー＝ジーンズの式。黒体放射の古典的な低周波近似を扱う。\par

\paragraph{再利用する考え方} モデルの前提を明示し、相共存曲線と圧力による沸点の変化から記述や数値を判定する。

\paragraph{共通理論} \hyperref[theory:thermo:phase]{第\ref*{theory:thermo:phase}節}

\par 参考：\cite{mit_clapeyron2008}。

\par 難易度 1（校準初稿）、優先度未評価。


\chapter{電磁気・物性}
\section{電荷・場・回路をつなぐ}
電磁気の問題では、電荷と電流が場を作る過程、場が電荷へ力を及ぼす過程、
場に蓄えられるエネルギーを分けて整理する。
回路は、導線と素子へ分布する場を端子電圧と電流で近似したモデルである。
電荷一定か電圧一定か、静電平衡か時間変化する場か、直列か並列かを最初に確定する。
この共通本文は静電場・容量・直流・磁場・誘導・RLCの27問に対応する初稿である。
磁性・物性等の未執筆問題も一覧へ保持し、章全体の完成とはしない。

\section{点電荷と電場}\label{theory:em:coulomb}
真空中の点電荷間の力の大きさは
\[
 F=\frac{|q_1q_2|}{4\pi\varepsilon_0r^2},\qquad \mathbf F=q\mathbf E
\]
である。第一式は二つの電荷が作る相互作用、第二式の$\mathbf E$は対象電荷以外が作る場である。
正電荷の力は電場と同じ方向、負電荷では逆方向になる。
数値問題では電荷をC、距離をmへそろえ、最後にNになることを確かめる。
同符号か異符号かは力の向きを決め、大きさを求める絶対値とは別に扱う。\autocite{mit_em02}

\section{電位差と外力の仕事}\label{theory:em:potential}
静電場では$\mathbf E=-\nabla V$であり、PからQへの移動で
\[
 V_Q-V_P=-\int_P^Q\mathbf E\cdot d\boldsymbol\ell,
 \quad \Delta U=q(V_Q-V_P),\quad W_{\rm field}=-\Delta U
\]
となる。ゆっくり動かし、運動エネルギー変化を無視できるとき、外力の仕事は$W_{\rm ext}=\Delta U$である。
電場方向へ負電荷を動かす場合、電位は下がるが$q<0$なので位置エネルギーは増える。
したがって正の外力仕事と負の電位変化は矛盾しない。
候補が正の電位差だけを示す場合には、大小差$|V_Q-V_P|$と、向きを定めた差を本文で明示する。
時間変化する磁束のある回路では誘導電場が生じるため、静電位だけで全起電力を表せるとは限らない。\autocite{mit_em03}

\section{ガウス法則と導体}\label{theory:em:gauss}
真空中では$\int_S\mathbf E\cdot\mathbf n\,dS=Q/\varepsilon_0$である。
一様帯電平面をまたぐ薄い柱を取ると、側面の寄与は0、両端面の寄与は各$EA$となるため
$2EA=|\sigma|A/\varepsilon_0$、すなわち$E=|\sigma|/(2\varepsilon_0)$を得る。
一枚の平面と、逆符号二枚の間で場を足した場合の係数を区別する。
静電平衡の導体内部は$E=0$である。導体内部に場があれば自由電荷の移動が続くためである。
孤立した球の外側にある逆二乗場を、導体内部へ延長しない。\autocite{mit_em04}

\section{容量と接続}\label{theory:em:capacitor}
容量は$C=Q/V$で、理想平行平板では$C=\varepsilon A/d$となる。
容量の単位FはC/Vである。$Q$は片方の極板の電荷の大きさを表す。
未充電のコンデンサーを直列接続すると、孤立した中間導体の正味電荷が0のまま保存されるため、
各素子の電気量の大きさは等しくなる。その電圧を足すと$1/C_{\rm eq}=\sum 1/C_i$である。
並列では電圧が共通なので電荷を足し、$C_{\rm eq}=\sum C_i$となる。
抵抗の直列・並列とは逆の合成関係に見えるが、共通量と加算量から導けば暗記を必要としない。\autocite{mit_em05}

\section{電場エネルギー}\label{theory:em:energy}
固定形状の線形コンデンサーに微小電荷$dq$を運ぶ仕事は$dW=Vdq=(q/C)dq$である。
0からQまで積分して
\[
 U=\int_0^Q\frac qC\,dq=\frac{Q^2}{2C}=\frac12CV^2
\]
を得る。$QV$では充電中の電圧変化を無視しており、2倍になる。
エネルギーの数値計算で$\mu$を$10^{-6}$へ変換し、電圧の二乗に注意する。
LC回路で用いる$LI^2/2$も場へ蓄えるエネルギーであり、抵抗で熱になる消費エネルギーとは違う。

\section{分極・強誘電性・損失}\label{theory:em:dielectric}
分極$\mathbf P$は単位体積あたりの電気双極子モーメントで、
$\mathbf D=\varepsilon_0\mathbf E+\mathbf P$と表す。
通常の誘電体で外場に応じて分極することと、強誘電体の自発分極を区別する。
強誘電性は外部電場なしでも生じる分極と、電場によるその反転を特徴とする。
交流で分極の応答に遅れや散逸があると、供給エネルギーの一部が熱となる誘電損失が生じる。
電解コンデンサーの酸化アルミニウム・酸化タンタル膜は薄い絶縁層として使われる。
誘電率が大きいことだけで、損失や耐電圧もよいとは判断しない。\autocite{mit_em05,mit_ferroelectric}

\section{固定電荷での極板間引力}\label{theory:em:force}
端効果のない一様な線形誘電体を仮定する。Q一定で間隔$d$を変えると
\[
 U(Q,d)=\frac{Q^2d}{2\varepsilon A},\qquad
 F_{\rm attraction}=\left(\frac{\partial U}{\partial d}\right)_Q
 =\frac{Q^2}{2\varepsilon A}
\]
である。これは引力の大きさであり、間隔を増やす方向の力成分は負になる。
Qが一定なら間隔の変更だけでは引力は変わらないが、誘電率を5倍にすれば引力は1/5となる。
容量が変更前と同じでも、間隔微分が同じとは限らない。
電圧一定の場合は電源とのエネルギーの授受があるため、$CV^2/2$を固定Qの式と同じように微分しない。
固定Vでの大きさは$\varepsilon AV^2/(2d^2)$である。条件を先に選ぶ必要がある。

\section{直流と抵抗}\label{theory:em:dc}
電流は断面を通る符号付き電荷の時間変化率$I=dQ/dt$で、1 A=1 C/sである。
負の担い手の移動方向と電流方向は逆になる。電子に限らず、イオンも電流を運べる。
オーム性導体では$V=IR$、電力は$P=VI=I^2R=V^2/R$となる。
抵抗を並列につなぐと電圧が共通で$I=\sum V/R_i$より$1/R_{\rm eq}=\sum1/R_i$を得る。
合成抵抗が最小の構成抵抗より小さいことも検算になる。\autocite{mit_em06,mit_em07}

\section{抵抗率と担い手}\label{theory:em:resistivity}
均一な棒の抵抗は$R=\rho L/A$で、材料の抵抗率と形状の寄与を分離する。
金属では格子振動、不純物、欠陥等による散乱が抵抗へ寄与する。
室温近傍の変化は$\rho(T)\simeq\rho(T_0)[1+\alpha(T-T_0)]$で近似することが多いが、全温度域の厳密な法則ではない。
真性半導体では温度上昇で電子と正孔が熱励起され、担い手の増加が抵抗減少の主因となる。
低温の金属での特定のべき則を、多くの導体についての一般則へ広げない。\autocite{mit_em06,mit_semiconductor}

\section{電流による場と荷電粒子の運動}\label{theory:em:biot}\label{theory:em:lorentz}
定常電流素片が作る磁束密度はビオ・サバールの法則
$d\mathbf B=\mu_0I(d\boldsymbol\ell\times\hat{\mathbf r})/(4\pi r^2)$で与えられる。
直線電流の周囲の場は円周方向へ向き、局所的に電流方向へ垂直になる。
一方、場の中で動く粒子への磁気力は$q\mathbf v\times\mathbf B$である。
その仕事率は$\mathbf F\cdot\mathbf v=0$なので速さを変えず、方向を変える。
一様Bへ垂直な速度で入り、他の力がなければ円運動となり、半径は$mv/(|q|B)$である。
場を作る法則と、その場から力を受ける法則を混同しない。\autocite{mit_em09,mit_lorentz}

\section{ソレノイドと透磁率}\label{theory:em:solenoid}
長いコイルの端部から離れた内部では、巻数密度$n=N/\ell$を使い$B=\mu nI$と近似する。
全巻数Nをそのまま掛けると長さの次元が足りない。
磁性体があるときは、出題で指定された線形範囲の$\mu=\mu_0\mu_r$を用いる。
有限長の端部の場、材料の飽和や履歴はこの簡単な式の外にある。
TからmTへの換算では$10^3$を掛ける。\autocite{mit_em09}

\section{誘導とマクスウェル方程式}\label{theory:em:induction}\label{theory:em:maxwell}
一巻きの磁束を$\Phi_B$とすると$\mathcal E=-N\,d\Phi_B/dt$である。
負号は変化を妨げる向きであり、大きさの問題では$N|d\Phi_B/dt|$を使う。
真電荷を自由電荷として扱う巨視的な表記では
\begin{align*}
 \nabla\cdot\mathbf D&=\rho_{\rm free},&\nabla\cdot\mathbf B&=0,\\
 \nabla\times\mathbf E&=-\frac{\partial\mathbf B}{\partial t},&
 \nabla\times\mathbf H&=\mathbf i+\frac{\partial\mathbf D}{\partial t}.
\end{align*}
磁束変化から電場を得る法則と、電束密度変化から変位電流を得る項を区別する。
$\nabla\cdot\mathbf B=0$は、通常の電磁気学で磁気単極子を源としないことを表す。
電荷保存は$\partial_t\rho+\nabla\cdot\mathbf i=0$で、流出を正とした総電荷変化には負号が付く。
「出入り」を流入と流出の単なる和として扱わない。\autocite{mit_em10,mit_em13}

\section{LC振動と直列共振}\label{theory:em:lc}\label{theory:em:rlc}
無抵抗LC回路では$Q^2/(2C)+LI^2/2$が保存する。
最大電流時はQ=0となり、初期にコンデンサーへ蓄えたエネルギーが全てコイルへ移る。
交流電源で駆動する理想直列RLC回路では
\[
 Z=R+j\left(\omega L-\frac1{\omega C}\right),\quad
 |Z|^2=R^2+\left(\omega L-\frac1{\omega C}\right)^2
\]
となる。括弧が0となる$f_0=1/(2\pi\sqrt{LC})$で$|Z|$は最小Rになる。
抵抗で減衰する自由振動の周波数と、交流定常状態の電流共振を混同しない。
L,CをSIへ直してから計算し、最後にHzをkHzへ換算する。\autocite{mit_em11,mit_em12}

\section{場のエネルギー流}\label{theory:em:poynting}
電磁場のエネルギー流密度は$\mathbf S=\mathbf E\times\mathbf H$、単位はW/m$^2$である。
外積なのでE,H双方へ垂直であり、流れる方向と単位面積あたりの仕事率を同時に表す。
力の向きや電荷の移動方向と同一視しない。
電磁波だけでなく、回路へエネルギーが運ばれる記述にも使われる。\autocite{mit_em13}

\section{磁化を符号・秩序・履歴に分ける}\label{theory:em:magnetism}
単位体積あたりの磁気モーメントを磁化\(\mathbf M\)とする。SIでは
\(\mathbf B=\mu_0(\mathbf H+\mathbf M)\)。等方的な弱い線形応答では
\(\mathbf M=\chi\mathbf H\)で、反磁性は負の\(\chi\)、常磁性は小さな正の\(\chi\)を持つ。
強磁性は外場なしの自発磁化を持つ秩序であり、多磁区試料の正味磁化が小さいことと両立する。
反強磁性は典型的には反平行の副格子磁化が打ち消し合う秩序で、Néel温度を超えると秩序を失う。
通常の強磁性の秩序を失う温度はCurie温度と呼ぶ。残留磁化は外場を取り去った後の履歴による磁化で、
磁区配置と保磁力を含む材料の状態に依存する。\cite{openstax_magnetism}

室温の代表的強磁性元素はFe・Co・Niである。過去問のMnをCoに読み替えない。
磁性は温度・結晶相・組成にも依存する。中性子回折は核による構造散乱だけでなく磁気散乱にも使える。
Slater--Pauling曲線は合金の電子数と磁気モーメント、Ising模型はスピン相互作用と秩序を扱う。
一方、Wiedemann--Franz則の\(\kappa_e/(\sigma T)\)は電子の熱電輸送係数の比である。
磁性に関する相対比較で区別し、磁場やスピンの影響を全く受けないと断定しない。\cite{em51_mit_magnetic_readings,em51_wf}

\section{局在モーメントと伝導電子}\label{theory:em:pauli}
独立した局在モーメント、弱い外場、高温側のモデルでは、磁化の方向が熱揺らぎで乱され、
磁化率はCurie則\(\chi=C/T\)に従う。Tは絶対温度である。
相互作用を含むCurie--Weiss近似の\(C/(T-\theta)\)も、転移点近傍や全温度の厳密則ではない。
「常磁性なら必ず同じ温度依存」という暗記を避ける。

Pauli常磁性は金属の伝導電子のスピンへの応答で、占有のFermi統計が重要になる。
低温では全電子が独立に方向を変える局在モーメントのモデルを使えず、Fermi準位近傍の状態が応答する。
原子核まわりの軌道運動を起源とするという過去問の説明はこの機構と違う。
自由電子の軌道応答であるLandau反磁性、束縛電子の応答を分ける。\cite{em51_pauli}

超伝導のMeissner状態では磁束が排除されるが、温度・磁場の条件がある。
II型の混合状態では量子化磁束が侵入するため、「超伝導の全状態で完全排除」とはしない。
電気抵抗ゼロという輸送性質だけから、磁化の全状態を決めない。

\section{形状で変わる反磁場と磁場勾配の力}\label{theory:em:demag}
試料の磁化は周囲と内部の場を変える。均一磁化した楕円体では、主軸方向ごとに
\[
 H_{\mathrm d,i}=-N_iM_i,\qquad N_x+N_y+N_z=1
\]
と書ける。NはSIでの無次元の反磁場係数で、形状に依存する。
一般の不均一磁化・任意形状を一つのスカラーNで表すとは限らない。\cite{em51_demag}
反磁場は磁化から生じる場であり、反磁性という材料応答の名前とは別である。

黒鉛にも反磁性があるため、磁場中で「全く力を受けない」という全称は誤りである。
弱い等方的な誘起磁化の小試料なら、背景場を変えない近似で
\(\mathbf F\simeq\chi V\nabla(B^2)/(2\mu_0)\)と考えられる。
反磁性では強い場から弱い場へ向かう。実際の黒鉛の異方性や試料形状は別に考える。
理想的な一様場の並進力ゼロと、非一様場で働き得る力を区別する。

\section{巻数比を磁束変化から導く}\label{theory:em:transformer}
理想変圧器の一次・二次が共通の磁束\(\Phi(t)\)を共有するなら、
誘導起電力の大きさはそれぞれ\(N_1|d\Phi/dt|\)、\(N_2|d\Phi/dt|\)である。
したがって同じ位相規約での電圧の大きさ、実効値又は最大値同士では
\[
 \frac{E_2}{E_1}=\frac{N_2}{N_1}.
\]
巻数が多い側ほど電圧が大きいことが検算となる。損失を無視した電力保存から
\(E_1I_1=E_2I_2\)なので電流比は逆になる。
磁束漏れ、巻線抵抗、飽和、周波数応答を含む実機の限界は、この理想比の条件の外である。\cite{em51_transformer_ref}

\section{誘導加熱と誘電加熱を損失で分ける}\label{theory:em:heating}
誘導加熱は時間変化する磁場で導電体に渦電流を作り、そのJoule損を利用する。
磁化履歴損は強磁性体等に寄与するが、すべての導電体に必須の損失ではない。
誘電加熱は交流電場に対する分極応答の損失を利用し、代表的な対象は誘電体である。
有限の導電率を持つ材料では導電損と分極損が共存する場合がある。
マイクロ波加熱も誘電損失を利用できるが、高周波加熱の全方式がGHz帯という意味ではない。\cite{em51_heating}

良導体の線形な半空間に入る調和場では、表面からの場・電流の振幅が
\(\exp(-z/\delta)\)の形で減衰する。代表的な表皮深さは
\[
 \delta=\sqrt{\frac{2}{\omega\mu\sigma}}
       =\sqrt{\frac{\rho}{\pi f\mu}} .
\]
深さ\(\delta\)で振幅が\(1/e\)になるため、内部への侵入が厳密にゼロではない。
抵抗率・透磁率・周波数の条件を読み、有限試料の全温度分布をこの式だけで決めない。
加熱された表面から内部への熱伝導は電磁場の侵入とは別の過程である。\cite{em51_heating,em51_skin}
プラズマも高周波場から加熱できる。融合プラズマのRF加熱、無電極放電トーチ、
固体の渦電流加熱の回路と粒子の吸収機構は同一ではない。\cite{em51_plasma}

\section{副殻の容量と半導体の四面体結合}\label{theory:em:electrons}
一個の空間軌道にはスピンが異なる最大2電子を収容する。
s・p・d・f副殻の空間軌道数は1・3・5・7で、総容量は2・6・10・14である。
原典の「軌道」の語を、一個の軌道と副殻群で区別する。
Ge原子の外殻配置\(4s^24p^2\)と、通常の結晶の4近接・四面体のsp3結合モデルを分ける。
sp2の三角平面構造をダイヤモンド型へ当てはめない。\cite{atomic_quantum,em51_ge,em51_ge_structure}

SiへBをドープするとアクセプターを介して正孔が主要な担い手になる。
金属での負電子の平均ドリフト、半導体の正孔、孤立原子のスピンと
バルク金属のPauli応答を区別する。\cite{mit_semiconductor,em51_pauli}

\section{圧電性の逆命題を使わない}\label{theory:em:ferroelectric}
圧電性は応力による電気応答と電場によるひずみの結合である。
強誘電性は外場なしの自発分極と、その向きの電場による反転を特徴とする。
石英は圧電性を持つが、通常の強誘電体のような反転する自発分極を持たない。
圧電性から強誘電性を保証する逆命題は成立しない。\cite{mit_ferroelectric,em51_quartz}
多ドメイン試料では局所の自発分極が存在しても、正味の分極は相殺され得る。

典型的な強誘電体ではCurie温度を越えて常誘電相へ移る。
一次相転移は秩序変数が不連続に変わる場合で、強誘電相転移では自発分極を用いて考える。
有限の測定電場での飽和値・履歴と、電場ゼロでの熱力学的な秩序変数を無条件に同一視しない。
原子の変位、配向の秩序、結晶対称性と分極をつなぎ、微視的な機構を全物質で一種類に固定しない。\cite{em51_ferro_transition}

\section{特性線・連続放射と法令上の電波}\label{theory:em:spectral}
真空中の電磁波では\(c=f\lambda\)、光子エネルギーは\(E=hf\)である。
赤外・可視・紫外の順に波長が短くなる。
X線管の放射は、電子の減速による連続的な制動放射に、内殻準位差に対応する特性線が重なる。
管全体のスペクトルが連続成分を含むことと、特性X線自体を連続スペクトルと呼ぶことは違う。\cite{em51_xray}
取得できた電波法第二条の周波数上限は300万MHz以下で、\(3\times10^{12}\)Hz以下へ換算される。
令和5年度の施行版の直接取得には失敗したため、その時点の条文照合は個別問題で未確定として残す。
用途に基づくマイクロ波の波長区分と、法令の電波の範囲を別々に読む。\cite{em51_radio_law}

\section{過去問での確認}
\section{静電場・容量}

\subsection{平成23年度 問8}
\label{q:h23-reg-applied-08}

\noindent\textbf{原典：} sources/H23/17\_applied\_science.pdf、PDF物理3頁。

\textbf{公式正答：} 3（single）。

\paragraph{問題の要旨} 充電したコンデンサーに蓄えられるエネルギーを計算する。

\Needspace{4\baselineskip}\paragraph{確認した条件・記述（独自要約）}\mbox{}\par

\noindent\textbf{容量：} 2.0 μF\par

\noindent\textbf{電圧：} 300 V\par

\paragraph{問われる核心} 容量に蓄える電場エネルギー

\paragraph{必要知識} 静電容量、μ=10\textsuperscript{-6}、電圧上昇の積分。

\paragraph{考え方・正答を選ぶ理由（査読待ち初稿）} \(U=CV^2/2=(2.0\times10^{-6})300^2/2=0.09\) Jで3を選ぶ。充電中の微小仕事Vdqを積分するため、最終電圧を全電荷へ掛けたQVそのものではない。

\Needspace{4\baselineskip}\paragraph{選択肢の検討（査読待ち初稿）}\mbox{}\par

\noindent\textbf{1：条件に不一致} \(9.0\times10^{-3}\) J。正しい値0.09 Jの1/10である。\par

\noindent\textbf{2：条件に不一致} \(1.8\times10^{-2}\) J。正しい値0.09 Jの1/5である。\par

\noindent\textbf{3：条件に一致} \(9.0\times10^{-2}\) J。容量を2.0×10\textsuperscript{-6} Fへ換算したCV\textsuperscript{2}/2に一致する。\par

\noindent\textbf{4：条件に不一致} \(1.8\times10^{-1}\) J。CV\textsuperscript{2}とするとこの値になるが、充電中の電圧上昇を積分する係数1/2がない。\par

\noindent\textbf{5：条件に不一致} \(9.0\times10^{-1}\) J。正しい値の10倍で、接頭語又は指数の扱いが違う。\par

\paragraph{再利用する考え方} 容量に蓄える電場エネルギーを、成立条件・単位・符号から再現する。電場と電束密度、磁場と磁束密度、電荷一定と電圧一定を読み分ける。

\paragraph{共通理論} \hyperref[theory:em:capacitor]{第\ref*{theory:em:capacitor}節}, \hyperref[theory:em:energy]{第\ref*{theory:em:energy}節}

\par 参考：\cite{mit_em05}。

\par 難易度 2（校準初稿）、優先度未評価。

\subsection{平成24年度 問6}
\label{q:h24-reg-applied-06}

\noindent\textbf{原典：} sources/H24/17\_applied\_science.pdf、PDF物理2頁。

\textbf{公式正答：} 3（single）。

\paragraph{問題の要旨} 真空中の2点電荷間の静電気力の大きさを求める。

\Needspace{4\baselineskip}\paragraph{確認した条件・記述（独自要約）}\mbox{}\par

\noindent\textbf{電荷：} +2.0×10\textsuperscript{-6} C、+1.5×10\textsuperscript{-6} C\par

\noindent\textbf{距離：} 0.30 m\par

\noindent\textbf{係数：} 9.0×10\textsuperscript{9} N·m\textsuperscript{2}/C\textsuperscript{2}\par

\paragraph{問われる核心} 点電荷のクーロン力

\paragraph{必要知識} 電荷の符号、逆二乗則、単位と接頭語。

\paragraph{考え方・正答を選ぶ理由（査読待ち初稿）} \(F=k|q_1q_2|/r^2=9.0\times10^9(2.0\times10^{-6})(1.5\times10^{-6})/0.30^2=0.30\) N、3である。同符号なので互いに離れる方向の反発力であり、ここでは大きさを問う。

\Needspace{4\baselineskip}\paragraph{選択肢の検討（査読待ち初稿）}\mbox{}\par

\noindent\textbf{1：条件に不一致} 0.10 N。距離の逆二乗で評価した0.30 Nに一致しない。\par

\noindent\textbf{2：条件に不一致} 0.15 N。0.30 Nの半分であり、クーロン力へ1/2を掛ける理由がない。\par

\noindent\textbf{3：条件に一致} 0.30 N。電荷の積と距離の逆二乗から0.30 Nとなる。\par

\noindent\textbf{4：条件に不一致} 1.0 N。与えられた電荷と距離では0.30 Nである。\par

\noindent\textbf{5：条件に不一致} 3.0 N。正しい値の10倍である。\par

\paragraph{再利用する考え方} 点電荷のクーロン力を、成立条件・単位・符号から再現する。電場と電束密度、磁場と磁束密度、電荷一定と電圧一定を読み分ける。

\paragraph{共通理論} \hyperref[theory:em:coulomb]{第\ref*{theory:em:coulomb}節}

\par 参考：\cite{mit_em02}。

\par 難易度 2（校準初稿）、優先度未評価。

\subsection{平成25年度 問3}
\label{q:h25-reg-applied-03}

\noindent\textbf{原典：} sources/H25/17\_applied\_science.pdf、PDF物理2頁。

\textbf{公式正答：} 2（single）。

\paragraph{問題の要旨} 静電場の閉曲面積分の法則名と、無限一様帯電平面の電場の大きさを選ぶ。

\Needspace{4\baselineskip}\paragraph{確認した条件・記述（独自要約）}\mbox{}\par

\noindent\textbf{式：} \(\int_S E_n\,dS=Q/\varepsilon_0\)\par

\noindent\textbf{平面：} 面電荷密度σ。両側に等しい大きさで逆向きの垂直電場。\par

\noindent\textbf{組合せ：} 法則名／電場の大きさ。\par

\paragraph{問われる核心} ガウス法則と対称性

\paragraph{必要知識} 電束、外向き法線、無限平面、二面の寄与。

\paragraph{考え方・正答を選ぶ理由（査読待ち初稿）} 平面をまたぐ薄い柱を閉曲面に選ぶ。側面の電場法線成分は0、二つの端面の寄与は各EAである。内部電荷の大きさは|σ|Aなので\(2EA=|\sigma|A/\varepsilon_0\)、2を選ぶ。一枚の帯電平面と、逆符号二枚の間の電場を混同しない。

\Needspace{4\baselineskip}\paragraph{選択肢の検討（査読待ち初稿）}\mbox{}\par

\noindent\textbf{1：条件に不一致} ガウス／\(|\sigma|/\varepsilon_0\)。法則名は合うが、両面から出る電束の係数2が抜ける。\par

\noindent\textbf{2：条件に一致} ガウス／\(|\sigma|/(2\varepsilon_0)\)。ガウス法則を薄い柱へ適用した2EA=|σ|A/ε0に一致する。\par

\noindent\textbf{3：条件に不一致} ガウス／\(2|\sigma|/\varepsilon_0\)。両面の電束を足すと電場は半分になり、2倍ではない。\par

\noindent\textbf{4：条件に不一致} クーロン／\(|\sigma|/\varepsilon_0\)。原典の閉曲面積分はガウス法則で、電場の値も2倍である。\par

\noindent\textbf{5：条件に不一致} クーロン／\(|\sigma|/(2\varepsilon_0)\)。電場値は合うが提示された法則名と違う。\par

\paragraph{再利用する考え方} ガウス法則と対称性を、成立条件・単位・符号から再現する。電場と電束密度、磁場と磁束密度、電荷一定と電圧一定を読み分ける。

\paragraph{共通理論} \hyperref[theory:em:gauss]{第\ref*{theory:em:gauss}節}

\par 参考：\cite{mit_em04}。

\par 難易度 2（校準初稿）、優先度未評価。

\subsection{平成26年度 問4}
\label{q:h26-reg-applied-04}

\noindent\textbf{原典：} sources/H26/17\_applied\_science.pdf、PDF物理2頁。

\textbf{公式正答：} 4（single）。

\paragraph{問題の要旨} 負電荷を電場方向へゆっくり動かす外力の仕事から、電位差と電場の大きさを選ぶ。

\Needspace{4\baselineskip}\paragraph{確認した条件・記述（独自要約）}\mbox{}\par

\noindent\textbf{電荷：} −2 C\par

\noindent\textbf{移動：} PからQ。電場の方向と移動方向は同じ。\par

\noindent\textbf{外力の仕事：} 80 J。ゆっくり移動し運動エネルギー変化を無視。\par

\noindent\textbf{距離：} 0.1 m\par

\noindent\textbf{候補の順：} 電位差／電場。電位差は候補に合わせ大きさを表す。\par

\paragraph{問われる核心} 電位・外力仕事と負電荷の符号

\paragraph{必要知識} 準静的移動、位置エネルギーqV、E=−grad V、V/m=\allowbreak{}N/C。

\paragraph{考え方・正答を選ぶ理由（査読待ち初稿）} 外力の仕事は\(W_{ext}=q(V_Q-V_P)\)で、80=\allowbreak{}(−2)ΔVよりΔV=−40 Vである。電場方向へ進むと電位が下がるが、負電荷の位置エネルギーは増える。候補の電位差の大きさは40 V、\(|E|=|\Delta V|/d\)=400 V/mなので4である。電場の仕事は−80 Jと逆符号になる。

\Needspace{4\baselineskip}\paragraph{選択肢の検討（査読待ち初稿）}\mbox{}\par

\noindent\textbf{1：条件に不一致} 160 V／320 V·C。外力80 Jを電荷の絶対値2 Cで割ると40 Vであり、160 Vではない。\par

\noindent\textbf{2：条件に不一致} 160 V／1600 V/m。電位差の大きさは40 Vで、電場も400 V/mである。\par

\noindent\textbf{3：条件に不一致} 40 V／80 V·C。40 Vは合うがV·Cはエネルギーの単位で、電場ではない。\par

\noindent\textbf{4：条件に一致} 40 V／400 V/m。電位差の大きさ40 V、電場の大きさ400 V/mで両単位も合う。\par

\noindent\textbf{5：条件に不一致} 8 V／16 V·C。電位差は40 Vであり8 Vではなく、電場の単位も違う。\par

\paragraph{再利用する考え方} 電位・外力仕事と負電荷の符号を、成立条件・単位・符号から再現する。電場と電束密度、磁場と磁束密度、電荷一定と電圧一定を読み分ける。

\paragraph{共通理論} \hyperref[theory:em:potential]{第\ref*{theory:em:potential}節}

\par 参考：\cite{mit_em03}。

\paragraph{関連問題} \hyperref[q:r01-ret-applied-08]{R01-RET-APPLIED-08}

\par 難易度 2（校準初稿）、優先度未評価。

\subsection{平成28年度 問4}
\label{q:h28-reg-applied-04}

\noindent\textbf{原典：} sources/H28/17\_applied\_science.pdf、PDF物理2頁。

\textbf{公式正答：} 2（single）。

\paragraph{問題の要旨} 誘電体とコンデンサーについて、最も不適切な説明を選ぶ。

\Needspace{4\baselineskip}\paragraph{確認した条件・記述（独自要約）}\mbox{}\par

\noindent\textbf{方向：} 最も不適切。\par

\noindent\textbf{モデル：} 容量の比例関係は一様誘電体と平行平板の近似。\par

\paragraph{問われる核心} 誘電分極と強誘電性の区別

\paragraph{必要知識} 分極、容量、交流損失、電解コンデンサーの酸化膜。

\paragraph{考え方・正答を選ぶ理由（査読待ち初稿）} 電場による電気双極子の誘起・配向は分極であり、全ての誘電分極を強誘電性と呼ぶことはできない。外場0でも自発分極を持ち、それを電場で反転できることが強誘電体の特徴である。2を選ぶ。損失、材料、幾何学による容量変化は別の性質である。

\Needspace{4\baselineskip}\paragraph{選択肢の検討（査読待ち初稿）}\mbox{}\par

\noindent\textbf{1：適切} 静電容量の単位はファラドFである。。C=\allowbreak{}Q/Vなので単位はC/V=\allowbreak{}Fである。\par

\noindent\textbf{2：不適切} 絶縁体を電場へ置いて電気双極子モーメントを持つことを強誘電性という。。これは一般的な誘電分極である。強誘電性は外部電場なしの自発分極とその反転を特徴とする。\par

\noindent\textbf{3：適切} 交流電場で電気エネルギーの一部が熱となることを誘電損失という。。分極の応答に損失があると電場からのエネルギーが熱へ散逸する。\par

\noindent\textbf{4：適切} 可動電極の対向面積を変えると容量が変化する。。端効果を無視するとC=εA/dであり対向面積Aに依存する。\par

\noindent\textbf{5：適切} 誘電体として酸化アルミニウムや酸化タンタルが使われる。。金属酸化物の薄い絶縁層を使う電解コンデンサーの例である。\par

\paragraph{再利用する考え方} 誘電分極と強誘電性の区別を、成立条件・単位・符号から再現する。電場と電束密度、磁場と磁束密度、電荷一定と電圧一定を読み分ける。

\paragraph{共通理論} \hyperref[theory:em:dielectric]{第\ref*{theory:em:dielectric}節}, \hyperref[theory:em:capacitor]{第\ref*{theory:em:capacitor}節}

\par 参考：\cite{mit_em05,mit_ferroelectric}。

\paragraph{関連問題} \hyperref[q:r02-reg-applied-03]{R02-REG-APPLIED-03}

\par 難易度 2（校準初稿）、優先度未評価。

\subsection{平成29年度 問2}
\label{q:h29-reg-applied-02}

\noindent\textbf{原典：} sources/H29/17\_applied\_science.pdf、PDF物理1頁。

\textbf{公式正答：} 1（single）。

\paragraph{問題の要旨} 平行平板コンデンサーの容量と直列接続について、最も不適切なものを選ぶ。

\Needspace{4\baselineskip}\paragraph{確認した条件・記述（独自要約）}\mbox{}\par

\noindent\textbf{方向：} 最も不適切。\par

\noindent\textbf{モデル：} 平板の端効果を無視。容量比較では形状を固定。直列の2個は初め未充電。\par

\paragraph{問われる核心} 容量の幾何学と固定電荷

\paragraph{必要知識} 平行平板近似、比誘電率、V=\allowbreak{}Q/C、直列の電荷保存。

\paragraph{考え方・正答を選ぶ理由（査読待ち初稿）} \(C=\varepsilon A/d\)より誘電率、面積の増加と距離の減少は容量を大きくする。誘電率増加で容量が減るとする1が誤りである。Q一定の電圧はQ/C、未充電直列の同じ電気量は中間導体の電荷保存から確認する。

\Needspace{4\baselineskip}\paragraph{選択肢の検討（査読待ち初稿）}\mbox{}\par

\noindent\textbf{1：不適切} 比誘電率の大きな物質を極板間へ入れると容量は減少する。。C=εA/dなので同じ形状では誘電率の増加で容量は増え、減少しない。\par

\noindent\textbf{2：適切} 極板間距離を小さくすると容量は増える。。分母dが小さくなるとCが増える。\par

\noindent\textbf{3：適切} 極板面積に比例して容量は増える。。面積Aが増えるとCも同じ割合で増える。\par

\noindent\textbf{4：適切} 電気量が同じなら容量の小さい方が大きな電位差を持つ。。V=\allowbreak{}Q/CなのでQ一定ではCが小さいほどVが大きい。\par

\noindent\textbf{5：適切} 未充電の2個を直列につないで電圧を加えると等しい電気量が蓄えられる。。漏れのない直列回路で中間導体が初めに中性なら、電荷保存で対向極板の電荷の大きさは等しい。\par

\paragraph{再利用する考え方} 容量の幾何学と固定電荷を、成立条件・単位・符号から再現する。電場と電束密度、磁場と磁束密度、電荷一定と電圧一定を読み分ける。

\paragraph{共通理論} \hyperref[theory:em:capacitor]{第\ref*{theory:em:capacitor}節}, \hyperref[theory:em:dielectric]{第\ref*{theory:em:dielectric}節}

\par 参考：\cite{mit_em05}。

\paragraph{関連問題} \hyperref[q:r07-reg-applied-03]{R07-REG-APPLIED-03}

\par 難易度 2（校準初稿）、優先度未評価。

\subsection{令和元年度再試験 問8}
\label{q:r01-ret-applied-08}

\noindent\textbf{原典：} sources/R01\_retest/17\_applied\_science.pdf、PDF物理3頁。

\textbf{公式正答：} 3（single）。

\paragraph{問題の要旨} 負電荷を電場方向へゆっくり動かす外力の仕事から、電位差と電場の大きさを選ぶ。

\Needspace{4\baselineskip}\paragraph{確認した条件・記述（独自要約）}\mbox{}\par

\noindent\textbf{電荷：} −2 C\par

\noindent\textbf{移動：} PからQ。電場の方向と移動方向は同じ。\par

\noindent\textbf{外力の仕事：} 80 J。ゆっくり移動し運動エネルギー変化を無視。\par

\noindent\textbf{距離：} 0.2 m\par

\noindent\textbf{候補の順：} 電位差／電場。電位差は候補に合わせ大きさを表す。\par

\paragraph{問われる核心} 電位・外力仕事と負電荷の符号

\paragraph{必要知識} 準静的移動、位置エネルギーqV、E=−grad V、V/m=\allowbreak{}N/C。

\paragraph{考え方・正答を選ぶ理由（査読待ち初稿）} 外力の仕事は\(W_{ext}=q(V_Q-V_P)\)で、80=\allowbreak{}(−2)ΔVよりΔV=−40 Vである。電場方向へ進むと電位が下がるが、負電荷の位置エネルギーは増える。候補の電位差の大きさは40 V、\(|E|=|\Delta V|/d\)=200 V/mなので3である。電場の仕事は−80 Jと逆符号になる。

\Needspace{4\baselineskip}\paragraph{選択肢の検討（査読待ち初稿）}\mbox{}\par

\noindent\textbf{1：条件に不一致} 160 V／32 V·C。外力80 Jを電荷の絶対値2 Cで割ると40 Vであり、160 Vではない。\par

\noindent\textbf{2：条件に不一致} 160 V／800 V/m。電位差の大きさは40 Vで、電場も200 V/mである。\par

\noindent\textbf{3：条件に一致} 40 V／200 V/m。電位差の大きさ40 V、電場の大きさ200 V/mで両単位も合う。\par

\noindent\textbf{4：条件に不一致} 40 V／400 V/m。40 Vは合うが距離は0.2 mなので電場は200 V/mである。\par

\noindent\textbf{5：条件に不一致} 40 V／8 V·C。V·CはJに等しく、電場の単位V/mではない。\par

\paragraph{再利用する考え方} 電位・外力仕事と負電荷の符号を、成立条件・単位・符号から再現する。電場と電束密度、磁場と磁束密度、電荷一定と電圧一定を読み分ける。

\paragraph{共通理論} \hyperref[theory:em:potential]{第\ref*{theory:em:potential}節}

\par 参考：\cite{mit_em03}。

\paragraph{関連問題} \hyperref[q:h26-reg-applied-04]{H26-REG-APPLIED-04}

\par 難易度 2（校準初稿）、優先度未評価。

\subsection{令和2年度 問3}
\label{q:r02-reg-applied-03}

\noindent\textbf{原典：} sources/R02/17\_applied\_science.pdf、PDF物理2頁。

\textbf{公式正答：} 1（single）。

\paragraph{問題の要旨} 誘電体とコンデンサーについて、最も不適切な説明を選ぶ。

\Needspace{4\baselineskip}\paragraph{確認した条件・記述（独自要約）}\mbox{}\par

\noindent\textbf{方向：} 最も不適切。\par

\noindent\textbf{モデル：} 容量の比例関係は一様誘電体と平行平板の近似。\par

\paragraph{問われる核心} 誘電分極と強誘電性の区別

\paragraph{必要知識} 分極、容量、交流損失、電解コンデンサーの酸化膜。

\paragraph{考え方・正答を選ぶ理由（査読待ち初稿）} 電場による電気双極子の誘起・配向は分極であり、全ての誘電分極を強誘電性と呼ぶことはできない。外場0でも自発分極を持ち、それを電場で反転できることが強誘電体の特徴である。1を選ぶ。損失、材料、幾何学による容量変化は別の性質である。

\Needspace{4\baselineskip}\paragraph{選択肢の検討（査読待ち初稿）}\mbox{}\par

\noindent\textbf{1：不適切} 絶縁体を電場へ置いて電気双極子モーメントを持つことを強誘電性という。。これは一般的な誘電分極である。強誘電性は外部電場なしの自発分極とその反転を特徴とする。\par

\noindent\textbf{2：適切} 交流電場で電気エネルギーの一部が熱となることを誘電損失という。。分極の応答に損失があると電場からのエネルギーが熱へ散逸する。\par

\noindent\textbf{3：適切} 可動電極の対向面積を変えると容量が変化する。。端効果を無視するとC=εA/dであり対向面積Aに依存する。\par

\noindent\textbf{4：適切} 誘電体として酸化アルミニウムや酸化タンタルが使われる。。金属酸化物の薄い絶縁層を使う電解コンデンサーの例である。\par

\noindent\textbf{5：適切} 平行平板容量は誘電率に比例し、誘電体の厚さへ反比例する。。誘電体が一様に隙間を満たす平行平板近似ではC=εA/dである。\par

\paragraph{再利用する考え方} 誘電分極と強誘電性の区別を、成立条件・単位・符号から再現する。電場と電束密度、磁場と磁束密度、電荷一定と電圧一定を読み分ける。

\paragraph{共通理論} \hyperref[theory:em:dielectric]{第\ref*{theory:em:dielectric}節}, \hyperref[theory:em:capacitor]{第\ref*{theory:em:capacitor}節}

\par 参考：\cite{mit_em05,mit_ferroelectric}。

\paragraph{関連問題} \hyperref[q:h28-reg-applied-04]{H28-REG-APPLIED-04}

\par 難易度 2（校準初稿）、優先度未評価。

\subsection{令和5年度 問3}
\label{q:r05-reg-applied-03}

\noindent\textbf{原典：} sources/R05/17\_applied\_science.pdf、PDF物理2頁。

\textbf{公式正答：} 2（single）。

\paragraph{問題の要旨} 電荷を固定した平行平板の間隔と誘電率を変え、極板間引力の比を選ぶ。

\Needspace{4\baselineskip}\paragraph{確認した条件・記述（独自要約）}\mbox{}\par

\noindent\textbf{初期：} 真空、間隔d、帯電量Q、引力F0。\par

\noindent\textbf{変更：} 間隔5d、比誘電率5の誘電体で全間隙を満たす。\par

\noindent\textbf{固定量：} Qは変わらない。\par

\noindent\textbf{モデル：} 同じ面積A、端効果なし、一様な線形誘電体。\par

\paragraph{問われる核心} 固定電荷のエネルギー微分と静電力

\paragraph{必要知識} 容量、固定する量、平板近似、力とエネルギーの微分。

\paragraph{考え方・正答を選ぶ理由（査読待ち初稿）} 固定Qのエネルギーは\(U=Q^2/(2C)=Q^2d/(2\varepsilon A)\)であり、間隔方向の引力の大きさは\(F=\partial U/\partial d=Q^2/(2\varepsilon A)\)となる。εが5ε0へ変わるのでF=\allowbreak{}F0/5、2である。間隔5倍で容量は初期と同じになるが、力はエネルギーの間隔微分であり同じとは限らない。

\Needspace{4\baselineskip}\paragraph{選択肢の検討（査読待ち初稿）}\mbox{}\par

\noindent\textbf{1：条件に不一致} \(F=F_0/25\)。固定Qで引力は間隔に依存しないため、距離の逆二乗だけを使うのは違う。\par

\noindent\textbf{2：条件に一致} \(F=F_0/5\)。F=\allowbreak{}Q\textsuperscript{2}/(2εA)でεが5倍となり1/5である。\par

\noindent\textbf{3：条件に不一致} \(F=F_0\)。εの増加による引力減少を無視している。\par

\noindent\textbf{4：条件に不一致} \(F=5F_0\)。固定Qでは引力はεに反比例し、比例ではない。\par

\noindent\textbf{5：条件に不一致} \(F=25F_0\)。電荷固定の平板モデルでは25倍となる要因はない。\par

\paragraph{再利用する考え方} 固定電荷のエネルギー微分と静電力を、成立条件・単位・符号から再現する。電場と電束密度、磁場と磁束密度、電荷一定と電圧一定を読み分ける。

\paragraph{共通理論} \hyperref[theory:em:force]{第\ref*{theory:em:force}節}, \hyperref[theory:em:capacitor]{第\ref*{theory:em:capacitor}節}

\par 参考：\cite{mit_em05}。

\par 難易度 2（校準初稿）、優先度未評価。

\subsection{令和7年度 問3}
\label{q:r07-reg-applied-03}

\noindent\textbf{原典：} sources/R07/17\_applied\_science.pdf、PDF物理1頁。

\textbf{公式正答：} 3（single）。

\paragraph{問題の要旨} 平行平板コンデンサーの容量と直列接続について、最も不適切なものを選ぶ。

\Needspace{4\baselineskip}\paragraph{確認した条件・記述（独自要約）}\mbox{}\par

\noindent\textbf{方向：} 最も不適切。\par

\noindent\textbf{モデル：} 平板の端効果を無視。容量比較では形状を固定。直列の2個は初め未充電。\par

\paragraph{問われる核心} 容量の幾何学と固定電荷

\paragraph{必要知識} 平行平板近似、比誘電率、V=\allowbreak{}Q/C、直列の電荷保存。

\paragraph{考え方・正答を選ぶ理由（査読待ち初稿）} \(C=\varepsilon A/d\)より誘電率、面積の増加と距離の減少は容量を大きくする。誘電率増加で容量が減るとする3が誤りである。Q一定の電圧はQ/C、未充電直列の同じ電気量は中間導体の電荷保存から確認する。

\Needspace{4\baselineskip}\paragraph{選択肢の検討（査読待ち初稿）}\mbox{}\par

\noindent\textbf{1：適切} 極板面積に比例して容量は増える。。面積Aが増えるとCも同じ割合で増える。\par

\noindent\textbf{2：適切} 極板間距離を小さくすると容量は増える。。分母dが小さくなるとCが増える。\par

\noindent\textbf{3：不適切} 比誘電率の大きな物質を極板間へ入れると容量は減少する。。C=εA/dなので同じ形状では誘電率の増加で容量は増え、減少しない。\par

\noindent\textbf{4：適切} 未充電の2個を直列につないで電圧を加えると等しい電気量が蓄えられる。。漏れのない直列回路で中間導体が初めに中性なら、電荷保存で対向極板の電荷の大きさは等しい。\par

\noindent\textbf{5：適切} 電気量が同じなら容量の小さい方が大きな電位差を持つ。。V=\allowbreak{}Q/CなのでQ一定ではCが小さいほどVが大きい。\par

\paragraph{再利用する考え方} 容量の幾何学と固定電荷を、成立条件・単位・符号から再現する。電場と電束密度、磁場と磁束密度、電荷一定と電圧一定を読み分ける。

\paragraph{共通理論} \hyperref[theory:em:capacitor]{第\ref*{theory:em:capacitor}節}, \hyperref[theory:em:dielectric]{第\ref*{theory:em:dielectric}節}

\par 参考：\cite{mit_em05}。

\paragraph{関連問題} \hyperref[q:h29-reg-applied-02]{H29-REG-APPLIED-02}

\par 難易度 2（校準初稿）、優先度未評価。

\section{電磁誘導・磁性}

\subsection{平成27年度 問2}
\label{q:h27-reg-applied-02}

\noindent\textbf{原典：} sources/H27/17\_applied\_science.pdf、PDF物理1頁。

\textbf{公式正答：} 3（single）。

\paragraph{問題の要旨} コイルを貫く磁束が一定速度で変わるときの誘導起電力の大きさを求める。

\Needspace{4\baselineskip}\paragraph{確認した条件・記述（独自要約）}\mbox{}\par

\noindent\textbf{巻数：} 2000\par

\noindent\textbf{磁束変化：} 各巻きを貫く磁束が20秒で0.01 Wb変化。\par

\noindent\textbf{速度：} 一定の割合。1 Wb=\allowbreak{}1 V·s。\par

\paragraph{問われる核心} 磁束鎖交数の変化と誘導

\paragraph{必要知識} ファラデーの法則、巻数、Wb/s=\allowbreak{}V。

\paragraph{考え方・正答を選ぶ理由（査読待ち初稿）} \(|\mathcal E|=N|\Delta\Phi|/\Delta t=2000(0.01)/20=1.0\) Vで3を選ぶ。一定速度なので平均変化率が任意の時点の変化率に等しい。向きは磁束変化を妨げるレンツの法則で決まり、ここでは大きさのみを求める。

\Needspace{4\baselineskip}\paragraph{選択肢の検討（査読待ち初稿）}\mbox{}\par

\noindent\textbf{1：条件に不一致} 0.10 V。2000×0.01/20=\allowbreak{}1 Vと違う。\par

\noindent\textbf{2：条件に不一致} 0.50 V。正しい値の1/2であり余分な係数がある。\par

\noindent\textbf{3：条件に一致} 1.0 V。巻数×各巻きの磁束変化率に一致する。\par

\noindent\textbf{4：条件に不一致} 2.5 V。磁束変化の時間20 sで割った1 Vと違う。\par

\noindent\textbf{5：条件に不一致} 5.0 V。与えられた変化率では1 Vである。\par

\paragraph{再利用する考え方} 磁束鎖交数の変化と誘導を、成立条件・単位・符号から再現する。電場と電束密度、磁場と磁束密度、電荷一定と電圧一定を読み分ける。

\paragraph{共通理論} \hyperref[theory:em:induction]{第\ref*{theory:em:induction}節}

\par 参考：\cite{mit_em10}。

\par 難易度 2（校準初稿）、優先度未評価。

\subsection{平成28年度 問3}
\label{q:h28-reg-applied-03}

\noindent\textbf{原典：} sources/H28/17\_applied\_science.pdf、PDF物理1頁。

\textbf{公式正答：} 3（single）。

\paragraph{問題の要旨} 電束密度・磁束密度を用いた電磁気の法則で、最も不適切なものを選ぶ。

\Needspace{4\baselineskip}\paragraph{確認した条件・記述（独自要約）}\mbox{}\par

\noindent\textbf{記号：} E=電場、D=電束密度、H=磁場、B=磁束密度、ρ=真電荷密度、i=伝導電流密度。\par

\noindent\textbf{方向：} 最も不適切な記述。\par

\paragraph{問われる核心} マクスウェル方程式と場の発散

\paragraph{必要知識} 自由電荷と束縛電荷、回転と発散、変位電流、磁気単極子。

\paragraph{考え方・正答を選ぶ理由（査読待ち初稿）} 磁場に関するガウス法則\(\operatorname{div}B=0\)は、閉曲面を通る正味磁束が0であること、すなわち通常の理論で磁気単極子がないことを表す。「存在」を反映するとする3が誤りである。他の式は真電荷、磁束の変化、変位電流をそれぞれ扱う。

\Needspace{4\baselineskip}\paragraph{選択肢の検討（査読待ち初稿）}\mbox{}\par

\noindent\textbf{1：適切} 閉曲面を貫く電束は、内部にある真電荷の総量に等しい。。Dの閉曲面電束は自由電荷、原典の真電荷の総量に等しい。束縛電荷まで含むEの式と区別する。\par

\noindent\textbf{2：適切} \(\operatorname{rot}E=-\partial B/\partial t\)は時間変化する磁束密度が電場を誘導する法則である。。ファラデーの法則の局所形で、負号は向きを表す。\par

\noindent\textbf{3：不適切} \(\operatorname{div}B=0\)は磁気単極子の存在を反映する。。発散0は磁気単極子による湧き出し・吸込みがないという通常の電磁気学の式であり、存在と逆である。\par

\noindent\textbf{4：適切} \(\operatorname{rot}H=i+\partial D/\partial t\)で\(\partial D/\partial t\)は変位電流密度である。。伝導電流に加える変位電流の項を正しく述べる。\par

\noindent\textbf{5：適切} 磁場のガウス法則は時間変化する磁場にも成立する。。div B=\allowbreak{}0は静磁場だけの制限ではなく、マクスウェル方程式の一つである。\par

\paragraph{再利用する考え方} マクスウェル方程式と場の発散を、成立条件・単位・符号から再現する。電場と電束密度、磁場と磁束密度、電荷一定と電圧一定を読み分ける。

\paragraph{共通理論} \hyperref[theory:em:maxwell]{第\ref*{theory:em:maxwell}節}, \hyperref[theory:em:gauss]{第\ref*{theory:em:gauss}節}

\par 参考：\cite{mit_em13}。

\paragraph{関連問題} \hyperref[q:r03-reg-applied-04]{R03-REG-APPLIED-04}

\par 難易度 2（校準初稿）、優先度未評価。

\subsection{平成30年度 問4}
\label{q:h30-reg-applied-04}

\noindent\textbf{原典：} sources/H30/17\_applied\_science.pdf、PDF物理2頁。

\textbf{公式正答：} 3（single）。

\paragraph{問題の要旨} 小さな電流素片が作る磁場を与える法則名を選ぶ。

\Needspace{4\baselineskip}\paragraph{確認した条件・記述（独自要約）}\mbox{}\par

\noindent\textbf{対象：} 定常電流の微小素片が作る磁場。\par

\paragraph{問われる核心} 電流素片から磁場を作る法則

\paragraph{必要知識} 定常電流、外積、重ね合わせ、磁束の時間変化との違い。

\paragraph{考え方・正答を選ぶ理由（査読待ち初稿）} 微小素片の寄与は\(dB=\mu_0 I(d\ell\times\hat r)/(4\pi r^2)\)というビオ・サバールの法則で与えられ、3である。素片を積分して導線全体の場を求める。誘導起電力の法則と、電流を源とする場の法則を区別する。

\Needspace{4\baselineskip}\paragraph{選択肢の検討（査読待ち初稿）}\mbox{}\par

\noindent\textbf{1：条件に不一致} クーロンの法則。点電荷の静電気力を与え、電流素片の磁場ではない。\par

\noindent\textbf{2：条件に不一致} プランクの法則。黒体の熱放射スペクトルの法則である。\par

\noindent\textbf{3：条件に一致} ビオ・サバールの法則。定常電流素片Idℓから磁束密度dBを与える法則である。\par

\noindent\textbf{4：条件に不一致} ファラデーの誘導法則。磁束の時間変化による誘導起電力を与える。\par

\noindent\textbf{5：条件に不一致} キュリー・ワイスの法則。磁性体の磁化率の温度依存を扱う。\par

\paragraph{再利用する考え方} 電流素片から磁場を作る法則を、成立条件・単位・符号から再現する。電場と電束密度、磁場と磁束密度、電荷一定と電圧一定を読み分ける。

\paragraph{共通理論} \hyperref[theory:em:biot]{第\ref*{theory:em:biot}節}

\par 参考：\cite{mit_em09}。

\paragraph{関連問題} \hyperref[q:r04-reg-applied-03]{R04-REG-APPLIED-03}

\par 難易度 2（校準初稿）、優先度未評価。

\subsection{令和元年度 問4}
\label{q:r01-reg-applied-04}

\noindent\textbf{原典：} sources/R01/17\_applied\_science.pdf、PDF物理2頁。

\textbf{公式正答：} 3（single）。

\paragraph{問題の要旨} 電磁気の基本法則の説明で、最も不適切なものを選ぶ。

\Needspace{4\baselineskip}\paragraph{確認した条件・記述（独自要約）}\mbox{}\par

\noindent\textbf{注意：} ③の原典は二箇所とも電束と記す。磁束へ書き換えて読まない。\par

\noindent\textbf{方向：} 最も不適切。\par

\paragraph{問われる核心} 電束の変化と磁束の変化

\paragraph{必要知識} ガウス、ビオ・サバール、ファラデー、連続の式、変位電流。

\paragraph{考え方・正答を選ぶ理由（査読待ち初稿）} 3の決め手は原典の「電束」である。ファラデーの法則は\(\mathcal E=-d\Phi_B/dt\)と磁束を用いる。電束変化による変位電流と混同しているため3を選ぶ。保存則の「出入り」は流入を正にした正味量として区別する。

\Needspace{4\baselineskip}\paragraph{選択肢の検討（査読待ち初稿）}\mbox{}\par

\noindent\textbf{1：適切} 閉曲面を貫く電束は内部の真電荷の量に等しく、電場のガウス法則である。。Dの閉曲面電束を内部真電荷へ対応させる説明である。\par

\noindent\textbf{2：適切} ビオ・サバールの法則は定常電流が作る磁場を決める。。電流素片から場を積分する定常磁場の法則である。\par

\noindent\textbf{3：不適切} 閉回路を貫く電束が変化すると、電束変化を妨げる起電力が生じることをファラデーの法則という。。誘導起電力は磁束の変化に対応する。電束変化はアンペール・マクスウェルの変位電流側の量である。\par

\noindent\textbf{4：適切} 領域の総電荷の時間変化は閉曲面を通じた電荷の正味の出入りと対応し、電気量保存を表す。。外向き電流を正に取ればdQ/dt=−∫i·n dSである。原典の出入りは符号付きの正味量として読む。\par

\noindent\textbf{5：適切} マクスウェル方程式は電磁場の基本方程式で19世紀に提出された。。マクスウェルが19世紀に電磁気の法則をまとめたという説明に合う。\par

\paragraph{再利用する考え方} 電束の変化と磁束の変化を、成立条件・単位・符号から再現する。電場と電束密度、磁場と磁束密度、電荷一定と電圧一定を読み分ける。

\paragraph{共通理論} \hyperref[theory:em:induction]{第\ref*{theory:em:induction}節}, \hyperref[theory:em:maxwell]{第\ref*{theory:em:maxwell}節}

\par 参考：\cite{mit_em10,mit_em13}。

\par 難易度 2（校準初稿）、優先度未評価。

\subsection{令和3年度 問4}
\label{q:r03-reg-applied-04}

\noindent\textbf{原典：} sources/R03/17\_applied\_science.pdf、PDF物理2頁。

\textbf{公式正答：} 2（single）。

\paragraph{問題の要旨} 電束密度・磁束密度を用いた電磁気の法則で、最も不適切なものを選ぶ。

\Needspace{4\baselineskip}\paragraph{確認した条件・記述（独自要約）}\mbox{}\par

\noindent\textbf{記号：} E=電場、D=電束密度、H=磁場、B=磁束密度、ρ=真電荷密度、i=伝導電流密度。\par

\noindent\textbf{方向：} 最も不適切な記述。\par

\paragraph{問われる核心} マクスウェル方程式と場の発散

\paragraph{必要知識} 自由電荷と束縛電荷、回転と発散、変位電流、磁気単極子。

\paragraph{考え方・正答を選ぶ理由（査読待ち初稿）} 磁場に関するガウス法則\(\operatorname{div}B=0\)は、閉曲面を通る正味磁束が0であること、すなわち通常の理論で磁気単極子がないことを表す。「存在」を反映するとする2が誤りである。他の式は真電荷、磁束の変化、変位電流をそれぞれ扱う。

\Needspace{4\baselineskip}\paragraph{選択肢の検討（査読待ち初稿）}\mbox{}\par

\noindent\textbf{1：適切} \(\operatorname{rot}E=-\partial B/\partial t\)は時間変化する磁束密度が電場を誘導する法則である。。ファラデーの法則の局所形で、負号は向きを表す。\par

\noindent\textbf{2：不適切} \(\operatorname{div}B=0\)は磁気単極子の存在を反映する。。発散0は磁気単極子による湧き出し・吸込みがないという通常の電磁気学の式であり、存在と逆である。\par

\noindent\textbf{3：適切} 磁場のガウス法則は時間変化する磁場にも成立する。。div B=\allowbreak{}0は静磁場だけの制限ではなく、マクスウェル方程式の一つである。\par

\noindent\textbf{4：適切} 閉曲面を貫く電束は、内部にある真電荷の総量に等しい。。Dの閉曲面電束は自由電荷、原典の真電荷の総量に等しい。束縛電荷まで含むEの式と区別する。\par

\noindent\textbf{5：適切} \(\operatorname{rot}H=i+\partial D/\partial t\)で\(\partial D/\partial t\)は変位電流密度である。。伝導電流に加える変位電流の項を正しく述べる。\par

\paragraph{再利用する考え方} マクスウェル方程式と場の発散を、成立条件・単位・符号から再現する。電場と電束密度、磁場と磁束密度、電荷一定と電圧一定を読み分ける。

\paragraph{共通理論} \hyperref[theory:em:maxwell]{第\ref*{theory:em:maxwell}節}, \hyperref[theory:em:gauss]{第\ref*{theory:em:gauss}節}

\par 参考：\cite{mit_em13}。

\paragraph{関連問題} \hyperref[q:h28-reg-applied-03]{H28-REG-APPLIED-03}

\par 難易度 2（校準初稿）、優先度未評価。

\subsection{令和4年度 問3}
\label{q:r04-reg-applied-03}

\noindent\textbf{原典：} sources/R04/17\_applied\_science.pdf、PDF物理1頁。

\textbf{公式正答：} 4（single）。

\paragraph{問題の要旨} 小さな電流素片が作る磁場を与える法則名を選ぶ。

\Needspace{4\baselineskip}\paragraph{確認した条件・記述（独自要約）}\mbox{}\par

\noindent\textbf{対象：} 定常電流の微小素片が作る磁場。\par

\paragraph{問われる核心} 電流素片から磁場を作る法則

\paragraph{必要知識} 定常電流、外積、重ね合わせ、磁束の時間変化との違い。

\paragraph{考え方・正答を選ぶ理由（査読待ち初稿）} 微小素片の寄与は\(dB=\mu_0 I(d\ell\times\hat r)/(4\pi r^2)\)というビオ・サバールの法則で与えられ、4である。素片を積分して導線全体の場を求める。誘導起電力の法則と、電流を源とする場の法則を区別する。

\Needspace{4\baselineskip}\paragraph{選択肢の検討（査読待ち初稿）}\mbox{}\par

\noindent\textbf{1：条件に不一致} プランクの法則。黒体の熱放射スペクトルの法則である。\par

\noindent\textbf{2：条件に不一致} クーロンの法則。点電荷の静電気力を与え、電流素片の磁場ではない。\par

\noindent\textbf{3：条件に不一致} ファラデーの誘導法則。磁束の時間変化による誘導起電力を与える。\par

\noindent\textbf{4：条件に一致} ビオ・サバールの法則。定常電流素片Idℓから磁束密度dBを与える法則である。\par

\noindent\textbf{5：条件に不一致} キュリー・ワイスの法則。磁性体の磁化率の温度依存を扱う。\par

\paragraph{再利用する考え方} 電流素片から磁場を作る法則を、成立条件・単位・符号から再現する。電場と電束密度、磁場と磁束密度、電荷一定と電圧一定を読み分ける。

\paragraph{共通理論} \hyperref[theory:em:biot]{第\ref*{theory:em:biot}節}

\par 参考：\cite{mit_em09}。

\paragraph{関連問題} \hyperref[q:h30-reg-applied-04]{H30-REG-APPLIED-04}

\par 難易度 2（校準初稿）、優先度未評価。

\subsection{令和6年度 問4}
\label{q:r06-reg-applied-04}

\noindent\textbf{原典：} sources/R06/17\_applied\_science.pdf、PDF物理2頁。

\textbf{公式正答：} 2（single）。

\paragraph{問題の要旨} 静電場、磁場中の運動、誘導、エネルギー流の説明で最も不適切なものを選ぶ。

\Needspace{4\baselineskip}\paragraph{確認した条件・記述（独自要約）}\mbox{}\par

\noindent\textbf{方向：} 最も不適切。\par

\noindent\textbf{導体：} 孤立導体球を帯電させた静電平衡。\par

\noindent\textbf{運動：} 磁場に垂直な初速度、磁場以外の力なし。\par

\paragraph{問われる核心} 静電平衡導体の内外と電磁場

\paragraph{必要知識} 導体遮蔽、クーロン則、ローレンツ力、磁束、外積と流密度。

\paragraph{考え方・正答を選ぶ理由（査読待ち初稿）} 導体内に非零の電場が残ると自由電荷が移動して平衡ではなくなる。静電平衡では内部電場は0で、表面の外側には点電荷に相当する逆二乗場がある。内外を入れ替えた2を選ぶ。磁気力は仕事をせず、エネルギー輸送はポインティングベクトルで記述する。

\Needspace{4\baselineskip}\paragraph{選択肢の検討（査読待ち初稿）}\mbox{}\par

\noindent\textbf{1：適切} 点電荷間の力は距離の二乗に反比例する。。クーロン力の逆二乗則に一致する。\par

\noindent\textbf{2：不適切} 帯電導体球の内部電場は中心からの距離の二乗に反比例する。。静電平衡の導体内部はE=\allowbreak{}0である。逆二乗場は孤立球の外側の説明である。\par

\noindent\textbf{3：適切} 磁場に垂直な速度を持つ荷電粒子は、速度に垂直な力で円運動する。。qv×Bは速度へ垂直で速さを変えず、向きを変える。\par

\noindent\textbf{4：適切} 起電力の大きさは回路を貫く磁束の時間変化率に比例する。。ファラデーの法則の大きさに一致する。\par

\noindent\textbf{5：適切} ポインティングベクトルはエネルギー流の密度で、電場・磁場に垂直である。。S=\allowbreak{}E×Hでエネルギー流密度W/m\textsuperscript{2}を表し、外積なので双方へ垂直である。\par

\paragraph{再利用する考え方} 静電平衡導体の内外と電磁場を、成立条件・単位・符号から再現する。電場と電束密度、磁場と磁束密度、電荷一定と電圧一定を読み分ける。

\paragraph{共通理論} \hyperref[theory:em:gauss]{第\ref*{theory:em:gauss}節}, \hyperref[theory:em:lorentz]{第\ref*{theory:em:lorentz}節}, \hyperref[theory:em:poynting]{第\ref*{theory:em:poynting}節}

\par 参考：\cite{mit_em04,mit_em13,mit_lorentz}。

\par 難易度 2（校準初稿）、優先度未評価。

\subsection{令和7年度 問4}
\label{q:r07-reg-applied-04}

\noindent\textbf{原典：} sources/R07/17\_applied\_science.pdf、PDF物理2頁。

\textbf{公式正答：} 4（single）。

\paragraph{問題の要旨} 磁性体で満たしたソレノイド内部の磁束密度を計算する。

\Needspace{4\baselineskip}\paragraph{確認した条件・記述（独自要約）}\mbox{}\par

\noindent\textbf{長さ：} 40 cm=\allowbreak{}0.40 m\par

\noindent\textbf{全巻数：} 150\par

\noindent\textbf{比透磁率：} 100\par

\noindent\textbf{電流：} 1.5 A\par

\noindent\textbf{磁気定数：} \(4\pi\times10^{-7}\) T·m/A\par

\noindent\textbf{モデル：} 長いソレノイドの内部近似。比透磁率は与えられた値を使用。\par

\paragraph{問われる核心} ソレノイドの巻数密度と磁束密度

\paragraph{必要知識} アンペール則、透磁率、全巻数と巻数密度、cmとm、TとmT。

\paragraph{考え方・正答を選ぶ理由（査読待ち初稿）} 単位長さの巻数は150/0.40=\allowbreak{}375 m\textsuperscript{-1}である。\(B=\mu_0\mu_r(N/\ell)I=4\pi\times10^{-7}\times100\times375\times1.5\)≃0.0707 T=\allowbreak{}70.7 mTなので4を選ぶ。端部や磁性体の非線形性はこの出題の近似へ含めず、与えられた比透磁率を用いる。

\Needspace{4\baselineskip}\paragraph{選択肢の検討（査読待ち初稿）}\mbox{}\par

\noindent\textbf{1：条件に不一致} \(2.8\times10^{-3}\) mT。与えられた透磁率・巻数密度では約70.7 mTであり一致しない。\par

\noindent\textbf{2：条件に不一致} \(1.4\times10^{-1}\) mT。透磁率と巻数密度を含めた70.7 mTと一致しない。\par

\noindent\textbf{3：条件に不一致} \(1.1\times10^1\) mT。約70.7 mTより小さく、全巻数を単位長さ巻数へ換算する必要がある。\par

\noindent\textbf{4：条件に一致} \(7.1\times10^1\) mT。B=μ0μr(N/ℓ)Iから約70.7 mTで最も近い。\par

\noindent\textbf{5：条件に不一致} \(1.1\times10^2\) mT。同じ計算では約70.7 mTで110 mTではない。\par

\paragraph{再利用する考え方} ソレノイドの巻数密度と磁束密度を、成立条件・単位・符号から再現する。電場と電束密度、磁場と磁束密度、電荷一定と電圧一定を読み分ける。

\paragraph{共通理論} \hyperref[theory:em:solenoid]{第\ref*{theory:em:solenoid}節}

\par 参考：\cite{mit_em09}。

\par 難易度 2（校準初稿）、優先度未評価。

\section{回路・抵抗}

\subsection{平成24年度 問4}
\label{q:h24-reg-applied-04}

\noindent\textbf{原典：} sources/H24/17\_applied\_science.pdf、PDF物理1頁。

\textbf{公式正答：} 3（single）。

\paragraph{問題の要旨} 10 Ωと5 Ωを並列接続した合成抵抗を選ぶ。

\Needspace{4\baselineskip}\paragraph{確認した条件・記述（独自要約）}\mbox{}\par

\noindent\textbf{接続：} 同じ二端子間へ並列に接続。\par

\paragraph{問われる核心} 並列抵抗の電流加算

\paragraph{必要知識} オーム則、コンダクタンス、直列と並列。

\paragraph{考え方・正答を選ぶ理由（査読待ち初稿）} \(1/R=1/10+1/5=3/10\)よりR=\allowbreak{}10/3 Ω≃3.3 Ωで3を選ぶ。並列では電圧が共通、電流が加算されるため、抵抗そのものではなくコンダクタンスを足す。

\Needspace{4\baselineskip}\paragraph{選択肢の検討（査読待ち初稿）}\mbox{}\par

\noindent\textbf{1：条件に不一致} 約0.2 Ω。0.2は5 Ωの逆数だが単位はSで、抵抗ではない。\par

\noindent\textbf{2：条件に不一致} 約1.5 Ω。逆数の和を逆に戻した10/3 Ωと違う。\par

\noindent\textbf{3：条件に一致} 約3.3 Ω。共通電圧で各電流を足した等価抵抗10/3 Ωに最も近い。\par

\noindent\textbf{4：条件に不一致} 約7.5 Ω。合成抵抗が小さい方の5 Ωを超えており並列の性質に反する。\par

\noindent\textbf{5：条件に不一致} 約15 Ω。10+5は直列接続の場合である。\par

\paragraph{再利用する考え方} 並列抵抗の電流加算を、成立条件・単位・符号から再現する。電場と電束密度、磁場と磁束密度、電荷一定と電圧一定を読み分ける。

\paragraph{共通理論} \hyperref[theory:em:dc]{第\ref*{theory:em:dc}節}

\par 参考：\cite{mit_em07}。

\paragraph{関連問題} \hyperref[q:h26-reg-applied-02]{H26-REG-APPLIED-02}

\par 難易度 2（校準初稿）、優先度未評価。

\subsection{平成26年度 問2}
\label{q:h26-reg-applied-02}

\noindent\textbf{原典：} sources/H26/17\_applied\_science.pdf、PDF物理1頁。

\textbf{公式正答：} 3（single）。

\paragraph{問題の要旨} 10 Ωと5 Ωを並列接続した合成抵抗を選ぶ。

\Needspace{4\baselineskip}\paragraph{確認した条件・記述（独自要約）}\mbox{}\par

\noindent\textbf{接続：} 同じ二端子間へ並列に接続。\par

\paragraph{問われる核心} 並列抵抗の電流加算

\paragraph{必要知識} オーム則、コンダクタンス、直列と並列。

\paragraph{考え方・正答を選ぶ理由（査読待ち初稿）} \(1/R=1/10+1/5=3/10\)よりR=\allowbreak{}10/3 Ω≃3.3 Ωで3を選ぶ。並列では電圧が共通、電流が加算されるため、抵抗そのものではなくコンダクタンスを足す。

\Needspace{4\baselineskip}\paragraph{選択肢の検討（査読待ち初稿）}\mbox{}\par

\noindent\textbf{1：条件に不一致} 約0.2 Ω。0.2は5 Ωの逆数だが単位はSで、抵抗ではない。\par

\noindent\textbf{2：条件に不一致} 約1.5 Ω。逆数の和を逆に戻した10/3 Ωと違う。\par

\noindent\textbf{3：条件に一致} 約3.3 Ω。共通電圧で各電流を足した等価抵抗10/3 Ωに最も近い。\par

\noindent\textbf{4：条件に不一致} 約7.5 Ω。合成抵抗が小さい方の5 Ωを超えており並列の性質に反する。\par

\noindent\textbf{5：条件に不一致} 約15 Ω。10+5は直列接続の場合である。\par

\paragraph{再利用する考え方} 並列抵抗の電流加算を、成立条件・単位・符号から再現する。電場と電束密度、磁場と磁束密度、電荷一定と電圧一定を読み分ける。

\paragraph{共通理論} \hyperref[theory:em:dc]{第\ref*{theory:em:dc}節}

\par 参考：\cite{mit_em07}。

\paragraph{関連問題} \hyperref[q:h24-reg-applied-04]{H24-REG-APPLIED-04}

\par 難易度 2（校準初稿）、優先度未評価。

\subsection{平成29年度 問3}
\label{q:h29-reg-applied-03}

\noindent\textbf{原典：} sources/H29/17\_applied\_science.pdf、PDF物理1頁。

\textbf{公式正答：} 3（single）。

\paragraph{問題の要旨} 直列RLC回路のインピーダンスの大きさが最小となる交流周波数を求める。

\Needspace{4\baselineskip}\paragraph{確認した条件・記述（独自要約）}\mbox{}\par

\noindent\textbf{抵抗：} 100 Ω\par

\noindent\textbf{インダクタンス：} 1 mH\par

\noindent\textbf{容量：} 0.1 μF\par

\noindent\textbf{接続：} 抵抗・コイル・コンデンサーを直列。理想的な周波数一定素子の定常交流。\par

\paragraph{問われる核心} 直列RLCのリアクタンス相殺

\paragraph{必要知識} 複素インピーダンス、角周波数と周波数、m=\allowbreak{}10\textsuperscript{-3}、μ=10\textsuperscript{-6}。

\paragraph{考え方・正答を選ぶ理由（査読待ち初稿）} \(Z=R+j(\omega L-1/(\omega C))\)なので大きさの二乗は\(R^2+(\omega L-1/(\omega C))^2\)となる。非負の第2項が0となる\(f_0=1/(2\pi\sqrt{LC})\)で最小である。mHとμFをHとFへ換算すると約15.9155 kHz、3を選ぶ。Rは最小値を決めるが、この理想モデルの共振周波数は変えない。

\Needspace{4\baselineskip}\paragraph{選択肢の検討（査読待ち初稿）}\mbox{}\par

\noindent\textbf{1：条件に不一致} 1 kHz。SIへ換算した共振周波数は約15.9155 kHzであり、この候補より離れる。\par

\noindent\textbf{2：条件に不一致} 5.1 kHz。SIへ換算した共振周波数は約15.9155 kHzであり、この候補より離れる。\par

\noindent\textbf{3：条件に一致} 16 kHz。リアクタンスが相殺する約15.9155 kHzに最も近い。\par

\noindent\textbf{4：条件に不一致} 100 kHz。SIへ換算した共振周波数は約15.9155 kHzであり、この候補より離れる。\par

\noindent\textbf{5：条件に不一致} 630 kHz。SIへ換算した共振周波数は約15.9155 kHzであり、この候補より離れる。\par

\paragraph{再利用する考え方} 直列RLCのリアクタンス相殺を、成立条件・単位・符号から再現する。電場と電束密度、磁場と磁束密度、電荷一定と電圧一定を読み分ける。

\paragraph{共通理論} \hyperref[theory:em:rlc]{第\ref*{theory:em:rlc}節}

\par 参考：\cite{mit_em12}。

\paragraph{関連問題} \hyperref[q:r02-reg-applied-04]{R02-REG-APPLIED-04}

\par 難易度 2（校準初稿）、優先度未評価。

\subsection{平成30年度 問3}
\label{q:h30-reg-applied-03}

\noindent\textbf{原典：} sources/H30/17\_applied\_science.pdf、PDF物理2頁。

\textbf{公式正答：} 3（single）。

\paragraph{問題の要旨} 電流と電気抵抗について、最も不適切な記述を選ぶ。

\Needspace{4\baselineskip}\paragraph{確認した条件・記述（独自要約）}\mbox{}\par

\noindent\textbf{対象：} 一定の抵抗率を持つ一様な導体。オーム則を使う際は温度等を固定。\par

\paragraph{問われる核心} 抵抗率と導体の幾何学

\paragraph{必要知識} 電流の定義、電子の電荷、オーム則、ジュール電力。

\paragraph{考え方・正答を選ぶ理由（査読待ち初稿）} 一様導体では長いほど衝突する経路が増え、太いほど電流の経路が増える。\(R=\rho L/A\)なので3の比例関係が逆である。電力と電流の関係はオーム性・温度一定という条件で読む。

\Needspace{4\baselineskip}\paragraph{選択肢の検討（査読待ち初稿）}\mbox{}\par

\noindent\textbf{1：適切} 電流は断面を単位時間に通過する電気量である。。I=\allowbreak{}dQ/dtという定義に一致する。\par

\noindent\textbf{2：適切} 自由電子は負電荷なので電場と逆向きの力を受けて流れる。。電気力qEでq<0なので力はEと逆向きである。\par

\noindent\textbf{3：不適切} 抵抗は断面積に比例し長さに反比例する。。R=ρL/Aであり、長さと断面積の比例・反比例が逆である。\par

\noindent\textbf{4：適切} 電力は単位時間の仕事で、電圧の二乗に比例し抵抗へ反比例する。。抵抗性負荷ではP=\allowbreak{}VI=\allowbreak{}V\textsuperscript{2}/Rとなる。\par

\noindent\textbf{5：適切} 電流は電圧に比例し、抵抗が大きいほど流れにくい。。オーム性導体のI=\allowbreak{}V/Rに一致する。\par

\paragraph{再利用する考え方} 抵抗率と導体の幾何学を、成立条件・単位・符号から再現する。電場と電束密度、磁場と磁束密度、電荷一定と電圧一定を読み分ける。

\paragraph{共通理論} \hyperref[theory:em:dc]{第\ref*{theory:em:dc}節}, \hyperref[theory:em:resistivity]{第\ref*{theory:em:resistivity}節}

\par 参考：\cite{mit_em06}。

\par 難易度 2（校準初稿）、優先度未評価。

\subsection{令和元年度 問3}
\label{q:r01-reg-applied-03}

\noindent\textbf{原典：} sources/R01/17\_applied\_science.pdf、PDF物理1頁。

\textbf{公式正答：} 5（single）。

\paragraph{問題の要旨} 電流の定義、向き、磁場、単位の4空欄の組合せを選ぶ。

\Needspace{4\baselineskip}\paragraph{確認した条件・記述（独自要約）}\mbox{}\par

\noindent\textbf{A：} 電流が表す何の流れか。\par

\noindent\textbf{B：} 電流方向と逆に移動するAの符号。\par

\noindent\textbf{C：} 電流方向に対する、直線電流の周囲の磁場方向。\par

\noindent\textbf{D：} アンペアの表現。\par

\noindent\textbf{並び：} A／B／C／D。\par

\paragraph{問われる核心} 電流の符号規約と単位

\paragraph{必要知識} 電荷輸送、正負の担い手、右ねじ、アンペア。

\paragraph{考え方・正答を選ぶ理由（査読待ち初稿）} 電流は電荷の移動を数える。正電荷の移動方向を正とするため、負電荷の移動と電流は逆向きになる。直線電流の周囲の磁場は電流を囲む円周方向、1 A=\allowbreak{}1 C/sで5を選ぶ。電流の担い手を電子だけへ限定しない。

\Needspace{4\baselineskip}\paragraph{選択肢の検討（査読待ち初稿）}\mbox{}\par

\noindent\textbf{1：条件に不一致} 原子／負／垂直／C/s。電流は一般に電荷の流れで、原子に限定されない。\par

\noindent\textbf{2：条件に不一致} 電子／正／垂直／C/s。電子は負であり、正の電子が逆方向へ進むという表現は合わない。\par

\noindent\textbf{3：条件に不一致} イオン／正／平行／C/m\textsuperscript{2}。正電荷の移動方向は電流方向と同じで、磁場と単位も違う。\par

\noindent\textbf{4：条件に不一致} 電荷／負／平行／C/m\textsuperscript{2}。最初の二語は合うが、直線電流の磁場は周方向で電流へ垂直、単位はC/sである。\par

\noindent\textbf{5：条件に一致} 電荷／負／垂直／C/s。電荷の流れ、負の移動と逆向き、垂直、C/sの4条件が全て合う。\par

\paragraph{再利用する考え方} 電流の符号規約と単位を、成立条件・単位・符号から再現する。電場と電束密度、磁場と磁束密度、電荷一定と電圧一定を読み分ける。

\paragraph{共通理論} \hyperref[theory:em:dc]{第\ref*{theory:em:dc}節}, \hyperref[theory:em:biot]{第\ref*{theory:em:biot}節}

\par 参考：\cite{mit_em06,mit_em09}。

\par 難易度 2（校準初稿）、優先度未評価。

\subsection{令和元年度再試験 問6}
\label{q:r01-ret-applied-06}

\noindent\textbf{原典：} sources/R01\_retest/17\_applied\_science.pdf、PDF物理3頁。

\textbf{公式正答：} 4（single）。

\paragraph{問題の要旨} 直流回路と抵抗の説明で、最も不適切なものを選ぶ。

\Needspace{4\baselineskip}\paragraph{確認した条件・記述（独自要約）}\mbox{}\par

\noindent\textbf{対象：} 定常直流、真性半導体、純金属の抵抗。\par

\noindent\textbf{方向：} 最も不適切。\par

\paragraph{問われる核心} 抵抗の温度依存と材料の区別

\paragraph{必要知識} 金属の散乱、残留抵抗、真性キャリア、近似の温度範囲。

\paragraph{考え方・正答を選ぶ理由（査読待ち初稿）} 4は材料と温度域を限定せず二乗比例を一般則とする点が誤りである。金属の室温近傍では\(\rho(T)\simeq\rho(T_0)[1+\alpha(T-T_0)]\)と近似し、真性半導体では担い手の熱励起が支配する。低温でT\textsuperscript{2}項が現れる場合まで否定するのではない。

\Needspace{4\baselineskip}\paragraph{選択肢の検討（査読待ち初稿）}\mbox{}\par

\noindent\textbf{1：適切} 二点間の電位差はどの経路で測っても同じである。。静電的・定常直流の近似では電位は一価な関数であり経路に依存しない。\par

\noindent\textbf{2：適切} 真性半導体では温度上昇で電気抵抗が減る。。熱励起により担い手が増え、真性領域では抵抗が減る。\par

\noindent\textbf{3：適切} 純金属に少量の不純物を混ぜると抵抗は増える。。不純物による電子散乱が増え、残留抵抗等が大きくなる。\par

\noindent\textbf{4：不適切} 多くの導体の抵抗率は温度の二乗に比例する。。室温近傍の金属では基準温度からの変化を一次近似することが多い。特定低温域のT\textsuperscript{2}項を多くの導体の一般則にできない。\par

\noindent\textbf{5：適切} 消費電力は単位時間あたり消費するエネルギーである。。P=\allowbreak{}dW/dtであり、単位J/s=\allowbreak{}Wである。\par

\paragraph{再利用する考え方} 抵抗の温度依存と材料の区別を、成立条件・単位・符号から再現する。電場と電束密度、磁場と磁束密度、電荷一定と電圧一定を読み分ける。

\paragraph{共通理論} \hyperref[theory:em:resistivity]{第\ref*{theory:em:resistivity}節}, \hyperref[theory:em:potential]{第\ref*{theory:em:potential}節}

\par 参考：\cite{mit_em06,mit_semiconductor}。

\par 難易度 2（校準初稿）、優先度未評価。

\subsection{令和2年度 問4}
\label{q:r02-reg-applied-04}

\noindent\textbf{原典：} sources/R02/17\_applied\_science.pdf、PDF物理2頁。

\textbf{公式正答：} 1（single）。

\paragraph{問題の要旨} 直列RLC回路のインピーダンスの大きさが最小となる交流周波数を求める。

\Needspace{4\baselineskip}\paragraph{確認した条件・記述（独自要約）}\mbox{}\par

\noindent\textbf{抵抗：} 2 kΩ\par

\noindent\textbf{インダクタンス：} 5 mH\par

\noindent\textbf{容量：} 2 μF\par

\noindent\textbf{接続：} 抵抗・コイル・コンデンサーを直列。理想的な周波数一定素子の定常交流。\par

\paragraph{問われる核心} 直列RLCのリアクタンス相殺

\paragraph{必要知識} 複素インピーダンス、角周波数と周波数、m=\allowbreak{}10\textsuperscript{-3}、μ=10\textsuperscript{-6}。

\paragraph{考え方・正答を選ぶ理由（査読待ち初稿）} \(Z=R+j(\omega L-1/(\omega C))\)なので大きさの二乗は\(R^2+(\omega L-1/(\omega C))^2\)となる。非負の第2項が0となる\(f_0=1/(2\pi\sqrt{LC})\)で最小である。mHとμFをHとFへ換算すると約1.5915 kHz、1を選ぶ。Rは最小値を決めるが、この理想モデルの共振周波数は変えない。

\Needspace{4\baselineskip}\paragraph{選択肢の検討（査読待ち初稿）}\mbox{}\par

\noindent\textbf{1：条件に一致} 1.6 kHz。リアクタンスが相殺する約1.5915 kHzに最も近い。\par

\noindent\textbf{2：条件に不一致} 5.1 kHz。SIへ換算した共振周波数は約1.5915 kHzであり、この候補より離れる。\par

\noindent\textbf{3：条件に不一致} 10 kHz。SIへ換算した共振周波数は約1.5915 kHzであり、この候補より離れる。\par

\noindent\textbf{4：条件に不一致} 100 kHz。SIへ換算した共振周波数は約1.5915 kHzであり、この候補より離れる。\par

\noindent\textbf{5：条件に不一致} 630 kHz。SIへ換算した共振周波数は約1.5915 kHzであり、この候補より離れる。\par

\paragraph{再利用する考え方} 直列RLCのリアクタンス相殺を、成立条件・単位・符号から再現する。電場と電束密度、磁場と磁束密度、電荷一定と電圧一定を読み分ける。

\paragraph{共通理論} \hyperref[theory:em:rlc]{第\ref*{theory:em:rlc}節}

\par 参考：\cite{mit_em12}。

\paragraph{関連問題} \hyperref[q:h29-reg-applied-03]{H29-REG-APPLIED-03}, \hyperref[q:r06-reg-applied-03]{R06-REG-APPLIED-03}

\par 難易度 2（校準初稿）、優先度未評価。

\subsection{令和4年度 問4}
\label{q:r04-reg-applied-04}

\noindent\textbf{原典：} sources/R04/17\_applied\_science.pdf、PDF物理2頁。

\textbf{公式正答：} 3（single）。

\paragraph{問題の要旨} 充電したコンデンサーを抵抗0のコイルへ接続し、最大電流時のコイルエネルギーを求める。

\Needspace{4\baselineskip}\paragraph{確認した条件・記述（独自要約）}\mbox{}\par

\noindent\textbf{容量：} 2.0×10\textsuperscript{-10} F\par

\noindent\textbf{初期電圧：} 5.0 V\par

\noindent\textbf{インダクタンス：} 5.0×10\textsuperscript{-3} H\par

\noindent\textbf{抵抗：} 0。理想LC振動。\par

\paragraph{問われる核心} LC振動の電気・磁気エネルギー交換

\paragraph{必要知識} 無散逸振動、コンデンサーとコイルのエネルギー、最大値の時点。

\paragraph{考え方・正答を選ぶ理由（査読待ち初稿）} 抵抗0なので全エネルギー\(Q^2/(2C)+LI^2/2\)は一定である。電流最大時はQ=\allowbreak{}0、初期の\(CV_0^2/2=(2.0\times10^{-10})5.0^2/2=2.5\times10^{-9}\) Jが全てコイルにある。3を選ぶ。Lは最大電流を決めるが、問われる全エネルギーは初期条件から決まる。

\Needspace{4\baselineskip}\paragraph{選択肢の検討（査読待ち初稿）}\mbox{}\par

\noindent\textbf{1：条件に不一致} 0 J。最大電流時はコイルの磁場エネルギーが最大であり0でない。\par

\noindent\textbf{2：条件に不一致} \(1.3\times10^{-9}\) J。初期エネルギーの約半分だが、最大電流では電気エネルギーは0となり全量が磁気側へ移る。\par

\noindent\textbf{3：条件に一致} \(2.5\times10^{-9}\) J。初期のCV\textsuperscript{2}/2=\allowbreak{}2.5×10\textsuperscript{-9} Jに一致する。\par

\noindent\textbf{4：条件に不一致} \(5.0\times10^{-9}\) J。係数1/2を落としたCV\textsuperscript{2}の値である。\par

\noindent\textbf{5：条件に不一致} \(7.5\times10^{-9}\) J。保存される初期エネルギーの3倍である。\par

\paragraph{再利用する考え方} LC振動の電気・磁気エネルギー交換を、成立条件・単位・符号から再現する。電場と電束密度、磁場と磁束密度、電荷一定と電圧一定を読み分ける。

\paragraph{共通理論} \hyperref[theory:em:lc]{第\ref*{theory:em:lc}節}, \hyperref[theory:em:energy]{第\ref*{theory:em:energy}節}

\par 参考：\cite{mit_em11}。

\par 難易度 2（校準初稿）、優先度未評価。

\subsection{令和6年度 問3}
\label{q:r06-reg-applied-03}

\noindent\textbf{原典：} sources/R06/17\_applied\_science.pdf、PDF物理2頁。

\textbf{公式正答：} 1（single）。

\paragraph{問題の要旨} 直列RLC回路のインピーダンスの大きさが最小となる交流周波数を求める。

\Needspace{4\baselineskip}\paragraph{確認した条件・記述（独自要約）}\mbox{}\par

\noindent\textbf{抵抗：} 200 Ω\par

\noindent\textbf{インダクタンス：} 5.0 mH\par

\noindent\textbf{容量：} 0.50 μF\par

\noindent\textbf{接続：} 抵抗・コイル・コンデンサーを直列。理想的な周波数一定素子の定常交流。\par

\paragraph{問われる核心} 直列RLCのリアクタンス相殺

\paragraph{必要知識} 複素インピーダンス、角周波数と周波数、m=\allowbreak{}10\textsuperscript{-3}、μ=10\textsuperscript{-6}。

\paragraph{考え方・正答を選ぶ理由（査読待ち初稿）} \(Z=R+j(\omega L-1/(\omega C))\)なので大きさの二乗は\(R^2+(\omega L-1/(\omega C))^2\)となる。非負の第2項が0となる\(f_0=1/(2\pi\sqrt{LC})\)で最小である。mHとμFをHとFへ換算すると約3.1831 kHz、1を選ぶ。Rは最小値を決めるが、この理想モデルの共振周波数は変えない。

\Needspace{4\baselineskip}\paragraph{選択肢の検討（査読待ち初稿）}\mbox{}\par

\noindent\textbf{1：条件に一致} 3.2 kHz。リアクタンスが相殺する約3.1831 kHzに最も近い。\par

\noindent\textbf{2：条件に不一致} 6.3 kHz。SIへ換算した共振周波数は約3.1831 kHzであり、この候補より離れる。\par

\noindent\textbf{3：条件に不一致} 20 kHz。SIへ換算した共振周波数は約3.1831 kHzであり、この候補より離れる。\par

\noindent\textbf{4：条件に不一致} 64 kHz。SIへ換算した共振周波数は約3.1831 kHzであり、この候補より離れる。\par

\noindent\textbf{5：条件に不一致} 130 kHz。SIへ換算した共振周波数は約3.1831 kHzであり、この候補より離れる。\par

\paragraph{再利用する考え方} 直列RLCのリアクタンス相殺を、成立条件・単位・符号から再現する。電場と電束密度、磁場と磁束密度、電荷一定と電圧一定を読み分ける。

\paragraph{共通理論} \hyperref[theory:em:rlc]{第\ref*{theory:em:rlc}節}

\par 参考：\cite{mit_em12}。

\paragraph{関連問題} \hyperref[q:r02-reg-applied-04]{R02-REG-APPLIED-04}

\par 難易度 2（校準初稿）、優先度未評価。

\section{自発分極・圧電と強誘電相転移}

\subsection{令和7年度 問7}
\label{q:r07-reg-applied-07}

\noindent\textbf{原典：} sources/R07/17\_applied\_science.pdf、PDF物理3頁。

\textbf{公式正答：} 3（single）。

\par\textbf{検討未確定：} 多ドメイン試料の正味分極、有限外場での飽和分極の測定条件を留保。一次転移の理論的秩序変数の不連続と無条件に同一視しない。

\paragraph{問題の要旨} 自発分極・圧電性と強誘電相転移の包含関係を読む。

\Needspace{4\baselineskip}\paragraph{確認した条件・記述（独自要約）}\mbox{}\par

\noindent\textbf{原典確認：} S7\_051の11原典画像で13問の全五肢・指数・単位・否定を照合。\par

\noindent\textbf{対象：} 自発分極・圧電性と強誘電相転移の包含関係を読む。\par

\paragraph{問われる核心} 自発分極・圧電性と強誘電相転移の包含関係を読む。

\paragraph{必要知識} 3は包含関係を逆にした全称命題。圧電性は応力と分極・電場とひずみの結合、強誘電性は外場なしの分極とその電場反転を特徴とする。磁性のCurie温度と、分極のCurie温度も異なる秩序変数である。

\paragraph{考え方・正答を選ぶ理由（査読待ち初稿）} 3は包含関係を逆にした全称命題。圧電性は応力と分極・電場とひずみの結合、強誘電性は外場なしの分極とその電場反転を特徴とする。磁性のCurie温度と、分極のCurie温度も異なる秩序変数である。

\Needspace{4\baselineskip}\paragraph{選択肢の検討（査読待ち初稿）}\mbox{}\par

\noindent\textbf{1：代表的条件で適切} 外場なしで自発的な双極子。強誘電相の各分極ドメインに自発分極がある。多ドメイン試料で正味の分極が相殺されても局所の自発分極は否定されない。\par

\noindent\textbf{2：代表的条件で適切} 一般にCurie温度以上は常誘電相。典型的な強誘電体の相転移の説明。全物質の全温度経路の普遍的な単一転移へは広げない。\par

\noindent\textbf{3：不適切} 圧電結晶はすべて強誘電体。石英のように圧電性を持ちながら電場で反転する自発分極を持たない結晶がある。強誘電体は圧電性を持つという通常の関係の逆を保証しない。\par

\noindent\textbf{4：代表的条件で適切} 一次転移点で飽和分極が不連続。典型的な一次の強誘電相転移では秩序変数が不連続に変わる。有限電場での測定量や履歴の条件と熱力学の相境界を区別する。\par

\noindent\textbf{5：代表的条件で適切} 強誘電転移は構造相転移。通常の強誘電性ではイオン位置・配向秩序等の構造と分極が関係する。変位型と秩序無秩序型等の微視的機構を区別する。\par

\paragraph{再利用する考え方} 場・応答・担い手・材料・温度を区別し、数値の指数と単位、代表モデルの適用条件を確認する。

\paragraph{共通理論} \hyperref[theory:em:dielectric]{第\ref*{theory:em:dielectric}節}, \hyperref[theory:em:ferroelectric]{第\ref*{theory:em:ferroelectric}節}

\par 参考：\cite{mit_ferroelectric,em51_quartz,em51_ferro_transition}。

\par 難易度 3（校準初稿）、優先度未評価。

\section{磁気模型・測定と熱電輸送則}

\subsection{平成23年度 問2}
\label{q:h23-reg-applied-02}

\noindent\textbf{原典：} sources/H23/17\_applied\_science.pdf、PDF物理1頁。

\textbf{公式正答：} 1（single）。

\paragraph{問題の要旨} 磁性の測定・模型と熱電輸送則の役割を区別する。

\Needspace{4\baselineskip}\paragraph{確認した条件・記述（独自要約）}\mbox{}\par

\noindent\textbf{原典確認：} S7\_051の11原典画像で13問の全五肢・指数・単位・否定を照合。\par

\noindent\textbf{対象：} 磁性の測定・模型と熱電輸送則の役割を区別する。\par

\paragraph{問われる核心} 磁性の測定・模型と熱電輸送則の役割を区別する。

\paragraph{必要知識} 輸送係数の関係である1と、磁気モーメントの模型・温度依存・磁気構造の測定を分ける。Wiedemann–Franz則はκ/(σT)=\allowbreak{}Lで、低温の弾性散乱などの条件でLorenz数を比較する。物質の磁性と全く無関係という排他命題へは変えない。

\paragraph{考え方・正答を選ぶ理由（査読待ち初稿）} 輸送係数の関係である1と、磁気モーメントの模型・温度依存・磁気構造の測定を分ける。Wiedemann–Franz則はκ/(σT)=\allowbreak{}Lで、低温の弾性散乱などの条件でLorenz数を比較する。物質の磁性と全く無関係という排他命題へは変えない。

\Needspace{4\baselineskip}\paragraph{選択肢の検討（査読待ち初稿）}\mbox{}\par

\noindent\textbf{1：不適切} Wiedemann–Franz則。金属の熱伝導率と電気伝導率の関係が中心。磁性との相対比較では最も関係が薄い。磁場の影響を受けないという意味ではない。\par

\noindent\textbf{2：代表的条件で適切} Slater–Pauling曲線。合金の電子数・組成と磁気モーメントの関係を表す。\par

\noindent\textbf{3：代表的条件で適切} Curie則。局在磁気モーメントの常磁性磁化率と温度の関係。\par

\noindent\textbf{4：代表的条件で適切} Ising模型。離散的なスピン変数と相互作用で磁気秩序・相転移を扱う模型。\par

\noindent\textbf{5：代表的条件で適切} 中性子回折。中性子の磁気モーメントによる散乱で磁気秩序を調べられる。結晶構造の回折にも使われる。\par

\paragraph{再利用する考え方} 場・応答・担い手・材料・温度を区別し、数値の指数と単位、代表モデルの適用条件を確認する。

\paragraph{共通理論} \hyperref[theory:em:magnetism]{第\ref*{theory:em:magnetism}節}

\par 参考：\cite{em51_mit_magnetic_readings,em51_wf}。

\par 難易度 2（校準初稿）、優先度未評価。

\section{磁気秩序・Curie/Néelと反磁場}

\subsection{平成27年度 問7}
\label{q:h27-reg-applied-07}

\noindent\textbf{原典：} sources/H27/17\_applied\_science.pdf、PDF物理3頁。

\textbf{公式正答：} 5（single）。

\par\textbf{検討未確定：} 原典の常磁性全体への表現は広すぎる。PauliやVan Vleck等を含む全常磁性に単一の温度法則を適用せず、局在モーメントのCurie領域として説明する。

\paragraph{問題の要旨} 磁気秩序・残留磁化とCurie則の温度依存を区別する。

\Needspace{4\baselineskip}\paragraph{確認した条件・記述（独自要約）}\mbox{}\par

\noindent\textbf{原典確認：} S7\_051の11原典画像で13問の全五肢・指数・単位・否定を照合。\par

\noindent\textbf{対象：} 磁気秩序・残留磁化とCurie則の温度依存を区別する。\par

\paragraph{問われる核心} 磁気秩序・残留磁化とCurie則の温度依存を区別する。

\paragraph{必要知識} 局在モーメントの通常のCurie常磁性でχ=C/Tとなるので5。自発磁化の消失、反磁性の符号、Néel温度、残留磁化の定義を分ける。

\paragraph{考え方・正答を選ぶ理由（査読待ち初稿）} 局在モーメントの通常のCurie常磁性でχ=C/Tとなるので5。自発磁化の消失、反磁性の符号、Néel温度、残留磁化の定義を分ける。

\Needspace{4\baselineskip}\paragraph{選択肢の検討（査読待ち初稿）}\mbox{}\par

\noindent\textbf{1：代表的条件で適切} 昇温で自発磁化が消えるCurie温度。通常の強磁性相は温度上昇で秩序が弱まり、Curie温度で自発磁化を失う。\par

\noindent\textbf{2：代表的条件で適切} 外場に反対向きの反磁性。線形応答の磁化率が負で、誘起磁化が外場と反対向きになる。\par

\noindent\textbf{3：代表的条件で適切} 反強磁性から常磁性へのNéel温度。反平行の長距離秩序が失われる温度。強磁性のCurie温度とは区別する。\par

\noindent\textbf{4：代表的条件で適切} 磁場ゼロで残る残留磁化。磁化履歴のある試料に、外場を取り去っても残る磁化。全強磁性試料が必ず大きな値を持つとは限らない。\par

\noindent\textbf{5：不適切} 常磁性磁化率が温度とともに増加。局在モーメントのCurie領域ではχ=C/Tで、絶対温度の上昇に対して減少する。伝導電子のPauli常磁性の弱い温度依存と混同しない。\par

\paragraph{再利用する考え方} 場・応答・担い手・材料・温度を区別し、数値の指数と単位、代表モデルの適用条件を確認する。

\paragraph{共通理論} \hyperref[theory:em:magnetism]{第\ref*{theory:em:magnetism}節}, \hyperref[theory:em:pauli]{第\ref*{theory:em:pauli}節}

\par 参考：\cite{openstax_magnetism,em51_pauli}。

\par 難易度 3（校準初稿）、優先度未評価。

\subsection{平成29年度 問9}
\label{q:h29-reg-applied-09}

\noindent\textbf{原典：} sources/H29/17\_applied\_science.pdf、PDF物理3頁。

\textbf{公式正答：} 2（single）。

\par\textbf{検討未確定：} 元素の分類は通常の室温・安定相を想定する。常磁性の全モデルへの逆温度則の一般化、黒鉛の一様場と勾配場の混同を避ける。

\paragraph{問題の要旨} Néel温度・反磁場の形状依存・元素の磁性を照合する。

\Needspace{4\baselineskip}\paragraph{確認した条件・記述（独自要約）}\mbox{}\par

\noindent\textbf{原典確認：} S7\_051の11原典画像で13問の全五肢・指数・単位・否定を照合。\par

\noindent\textbf{対象：} Néel温度・反磁場の形状依存・元素の磁性を照合する。\par

\paragraph{問われる核心} Néel温度・反磁場の形状依存・元素の磁性を照合する。

\paragraph{必要知識} 適切な肢2は反強磁性秩序を失うNéel温度。ほかはCurie領域の逆温度依存、MnとCoの違い、形状による反磁場、黒鉛に働き得る磁場勾配力を確認する。

\paragraph{考え方・正答を選ぶ理由（査読待ち初稿）} 適切な肢2は反強磁性秩序を失うNéel温度。ほかはCurie領域の逆温度依存、MnとCoの違い、形状による反磁場、黒鉛に働き得る磁場勾配力を確認する。

\Needspace{4\baselineskip}\paragraph{選択肢の検討（査読待ち初稿）}\mbox{}\par

\noindent\textbf{1：不適切} 常磁性磁化率が温度とともに増加。局在モーメントのCurie領域ではχ=C/Tで、絶対温度の上昇に対して減少する。伝導電子のPauli常磁性の弱い温度依存と混同しない。\par

\noindent\textbf{2：代表的条件で適切} 反強磁性から常磁性へのNéel温度。反平行の長距離秩序が失われる温度。強磁性のCurie温度とは区別する。\par

\noindent\textbf{3：不適切} Fe・Mn・Niが単体強磁性。通常の室温の代表的強磁性元素はFe・Co・Ni。原典のMnをCoと読み替えない。Mnの相・温度依存は別に扱う。\par

\noindent\textbf{4：不適切} 反磁場が形状に依存しない。磁化で生じる反磁場は試料形状と磁化分布に依存する。均一磁化した楕円体でHd=−NMと表せる。\par

\noindent\textbf{5：不適切} 黒鉛は磁場中で全く力を受けない。黒鉛は反磁性を持ち、不均一な磁場で力を受け得る。一様な理想場では正味の並進力がゼロでも、全条件で力がないという記述は誤り。\par

\paragraph{再利用する考え方} 場・応答・担い手・材料・温度を区別し、数値の指数と単位、代表モデルの適用条件を確認する。

\paragraph{共通理論} \hyperref[theory:em:magnetism]{第\ref*{theory:em:magnetism}節}, \hyperref[theory:em:demag]{第\ref*{theory:em:demag}節}

\par 参考：\cite{openstax_magnetism,em51_pauli,em51_demag}。

\paragraph{関連問題} \hyperref[q:r06-reg-applied-07]{R06-REG-APPLIED-07}

\par 難易度 3（校準初稿）、優先度未評価。

\subsection{令和6年度 問7}
\label{q:r06-reg-applied-07}

\noindent\textbf{原典：} sources/R06/17\_applied\_science.pdf、PDF物理3頁。

\textbf{公式正答：} 3（single）。

\par\textbf{検討未確定：} 元素の分類は通常の室温・安定相を想定する。常磁性の全モデルへの逆温度則の一般化、黒鉛の一様場と勾配場の混同を避ける。

\paragraph{問題の要旨} Néel温度・反磁場の形状依存・元素の磁性を照合する。

\Needspace{4\baselineskip}\paragraph{確認した条件・記述（独自要約）}\mbox{}\par

\noindent\textbf{原典確認：} S7\_051の11原典画像で13問の全五肢・指数・単位・否定を照合。\par

\noindent\textbf{対象：} Néel温度・反磁場の形状依存・元素の磁性を照合する。\par

\paragraph{問われる核心} Néel温度・反磁場の形状依存・元素の磁性を照合する。

\paragraph{必要知識} 適切な肢3は反強磁性秩序を失うNéel温度。ほかはCurie領域の逆温度依存、MnとCoの違い、形状による反磁場、黒鉛に働き得る磁場勾配力を確認する。

\paragraph{考え方・正答を選ぶ理由（査読待ち初稿）} 適切な肢3は反強磁性秩序を失うNéel温度。ほかはCurie領域の逆温度依存、MnとCoの違い、形状による反磁場、黒鉛に働き得る磁場勾配力を確認する。

\Needspace{4\baselineskip}\paragraph{選択肢の検討（査読待ち初稿）}\mbox{}\par

\noindent\textbf{1：不適切} 常磁性磁化率が温度とともに増加。局在モーメントのCurie領域ではχ=C/Tで、絶対温度の上昇に対して減少する。伝導電子のPauli常磁性の弱い温度依存と混同しない。\par

\noindent\textbf{2：不適切} 黒鉛は磁場中で全く力を受けない。黒鉛は反磁性を持ち、不均一な磁場で力を受け得る。一様な理想場では正味の並進力がゼロでも、全条件で力がないという記述は誤り。\par

\noindent\textbf{3：代表的条件で適切} 反強磁性から常磁性へのNéel温度。反平行の長距離秩序が失われる温度。強磁性のCurie温度とは区別する。\par

\noindent\textbf{4：不適切} Fe・Mn・Niが単体強磁性。通常の室温の代表的強磁性元素はFe・Co・Ni。原典のMnをCoと読み替えない。Mnの相・温度依存は別に扱う。\par

\noindent\textbf{5：不適切} 反磁場が形状に依存しない。磁化で生じる反磁場は試料形状と磁化分布に依存する。均一磁化した楕円体でHd=−NMと表せる。\par

\paragraph{再利用する考え方} 場・応答・担い手・材料・温度を区別し、数値の指数と単位、代表モデルの適用条件を確認する。

\paragraph{共通理論} \hyperref[theory:em:magnetism]{第\ref*{theory:em:magnetism}節}, \hyperref[theory:em:demag]{第\ref*{theory:em:demag}節}

\par 参考：\cite{openstax_magnetism,em51_pauli,em51_demag}。

\paragraph{関連問題} \hyperref[q:h29-reg-applied-09]{H29-REG-APPLIED-09}

\par 難易度 3（校準初稿）、優先度未評価。

\section{局在スピン・Fermi統計とPauli応答}

\subsection{令和元年度 問8}
\label{q:r01-reg-applied-08}

\noindent\textbf{原典：} sources/R01/17\_applied\_science.pdf、PDF物理3頁。

\textbf{公式正答：} 5（single）。

\par\textbf{検討未確定：} 肢2の完全反磁性はMeissner状態の条件付き。II型超伝導の混合状態を含めて全状態へ断定しない。

\paragraph{問題の要旨} Pauliのスピン応答と軌道反磁性を区別する。

\Needspace{4\baselineskip}\paragraph{確認した条件・記述（独自要約）}\mbox{}\par

\noindent\textbf{原典確認：} S7\_051の11原典画像で13問の全五肢・指数・単位・否定を照合。\par

\noindent\textbf{対象：} Pauliのスピン応答と軌道反磁性を区別する。\par

\paragraph{問われる核心} Pauliのスピン応答と軌道反磁性を区別する。

\paragraph{必要知識} 誤りは5の起源の取り違え。局在スピンのCurie則と、Fermi準位近傍の伝導電子によるPauli常磁性を分ける。Meissner効果も超伝導の全磁場領域での完全排除へは広げない。

\paragraph{考え方・正答を選ぶ理由（査読待ち初稿）} 誤りは5の起源の取り違え。局在スピンのCurie則と、Fermi準位近傍の伝導電子によるPauli常磁性を分ける。Meissner効果も超伝導の全磁場領域での完全排除へは広げない。

\Needspace{4\baselineskip}\paragraph{選択肢の検討（査読待ち初稿）}\mbox{}\par

\noindent\textbf{1：代表的条件で適切} 室温の鉄は強磁性。通常圧・通常の純Feは室温で強磁性。結晶構造転移と磁気秩序の転移を区別する。\par

\noindent\textbf{2：条件・語義を留保} Meissner効果による完全反磁性。Meissner状態では磁束を排除する。臨界磁場・温度以下の条件が必要で、II型の混合状態では量子化磁束が侵入する。\par

\noindent\textbf{3：代表的条件で適切} 反強磁性から常磁性へのNéel温度。反平行の長距離秩序が失われる温度。強磁性のCurie温度とは区別する。\par

\noindent\textbf{4：代表的条件で適切} Curie点以上で常磁性。通常の強磁性物質ではCurie温度を超えると自発磁化秩序を失う。\par

\noindent\textbf{5：不適切} Pauli常磁性が原子核まわりの軌道運動。伝導電子のスピンに外場が作用し、Fermi統計を反映した応答となる。自由電子の軌道応答はLandau反磁性として区別する。\par

\paragraph{再利用する考え方} 場・応答・担い手・材料・温度を区別し、数値の指数と単位、代表モデルの適用条件を確認する。

\paragraph{共通理論} \hyperref[theory:em:pauli]{第\ref*{theory:em:pauli}節}

\par 参考：\cite{em51_pauli,openstax_magnetism}。

\par 難易度 3（校準初稿）、優先度未評価。

\section{金属抵抗率・温度と単位桁}

\subsection{平成26年度 問5}
\label{q:h26-reg-applied-05}

\noindent\textbf{原典：} sources/H26/17\_applied\_science.pdf、PDF物理2頁。

\textbf{公式正答：} 5（single）。

\paragraph{問題の要旨} 金属抵抗の形状・温度依存と抵抗率の桁を照合する。

\Needspace{4\baselineskip}\paragraph{確認した条件・記述（独自要約）}\mbox{}\par

\noindent\textbf{原典確認：} S7\_051の11原典画像で13問の全五肢・指数・単位・否定を照合。\par

\noindent\textbf{対象：} 金属抵抗の形状・温度依存と抵抗率の桁を照合する。\par

\paragraph{問われる核心} 金属抵抗の形状・温度依存と抵抗率の桁を照合する。

\paragraph{必要知識} 各回の誤りは肢5の抵抗率の桁である。R=ρL/AとI=\allowbreak{}V/Rを、材料・断面・温度一定という条件から使う。残留抵抗と温度依存を分け、H26のΩcm、R01とR05のΩmを独立に校正した。通常金属の温度傾向を特殊な転移や全合金へ広げない。

\paragraph{考え方・正答を選ぶ理由（査読待ち初稿）} 各回の誤りは肢5の抵抗率の桁である。R=ρL/AとI=\allowbreak{}V/Rを、材料・断面・温度一定という条件から使う。残留抵抗と温度依存を分け、H26のΩcm、R01とR05のΩmを独立に校正した。通常金属の温度傾向を特殊な転移や全合金へ広げない。

\Needspace{4\baselineskip}\paragraph{選択肢の検討（査読待ち初稿）}\mbox{}\par

\noindent\textbf{1：代表的条件で適切} 300Kと100Kの抵抗。通常の金属では格子振動による散乱が高温ほど増え、室温側の抵抗が大きい。相転移や異常な材料を含む全温度の保証ではない。\par

\noindent\textbf{2：代表的条件で適切} 不純物・格子欠陥と残留抵抗。低温で格子振動の寄与が減っても不純物・欠陥による散乱は残る。組成・構造による寄与を考える。\par

\noindent\textbf{3：代表的条件で適切} 一定断面でRが長さに比例。均一な導線はR=ρL/A。同じ材料・温度・断面積を保つ条件が必要。\par

\noindent\textbf{4：代表的条件で適切} 電流が電位差に比例。オーム性・温度一定ならI=\allowbreak{}V/R。自己加熱でRが変わる高電圧域まで無条件には広げない。\par

\noindent\textbf{5：不適切} 室温の金属が10\textsuperscript{2}〜10\textsuperscript{4}Ωcm。原典は正の指数でΩcm。Cuの代表値1.68×10\textsuperscript{-8}Ωmは1.68×10\textsuperscript{-6}Ωcm。元の範囲は1〜100Ωmに相当し、金属の典型抵抗率から大きく外れる。\par

\paragraph{再利用する考え方} 場・応答・担い手・材料・温度を区別し、数値の指数と単位、代表モデルの適用条件を確認する。

\paragraph{共通理論} \hyperref[theory:em:resistivity]{第\ref*{theory:em:resistivity}節}

\par 参考：\cite{openstax_resistivity,em51_resistance}。

\paragraph{関連問題} \hyperref[q:r01-reg-applied-07]{R01-REG-APPLIED-07}

\par 難易度 2（校準初稿）、優先度未評価。

\subsection{令和元年度 問7}
\label{q:r01-reg-applied-07}

\noindent\textbf{原典：} sources/R01/17\_applied\_science.pdf、PDF物理3頁。

\textbf{公式正答：} 3（single）。

\paragraph{問題の要旨} 金属抵抗の形状・温度依存と抵抗率の桁を照合する。

\Needspace{4\baselineskip}\paragraph{確認した条件・記述（独自要約）}\mbox{}\par

\noindent\textbf{原典確認：} S7\_051の11原典画像で13問の全五肢・指数・単位・否定を照合。\par

\noindent\textbf{対象：} 金属抵抗の形状・温度依存と抵抗率の桁を照合する。\par

\paragraph{問われる核心} 金属抵抗の形状・温度依存と抵抗率の桁を照合する。

\paragraph{必要知識} 各回の誤りは肢3の抵抗率の桁である。R=ρL/AとI=\allowbreak{}V/Rを、材料・断面・温度一定という条件から使う。残留抵抗と温度依存を分け、H26のΩcm、R01とR05のΩmを独立に校正した。通常金属の温度傾向を特殊な転移や全合金へ広げない。

\paragraph{考え方・正答を選ぶ理由（査読待ち初稿）} 各回の誤りは肢3の抵抗率の桁である。R=ρL/AとI=\allowbreak{}V/Rを、材料・断面・温度一定という条件から使う。残留抵抗と温度依存を分け、H26のΩcm、R01とR05のΩmを独立に校正した。通常金属の温度傾向を特殊な転移や全合金へ広げない。

\Needspace{4\baselineskip}\paragraph{選択肢の検討（査読待ち初稿）}\mbox{}\par

\noindent\textbf{1：代表的条件で適切} 一定断面でRが長さに比例。均一な導線はR=ρL/A。同じ材料・温度・断面積を保つ条件が必要。\par

\noindent\textbf{2：代表的条件で適切} 電流が電位差に比例。オーム性・温度一定ならI=\allowbreak{}V/R。自己加熱でRが変わる高電圧域まで無条件には広げない。\par

\noindent\textbf{3：不適切} 室温の金属が10\textsuperscript{2}〜10\textsuperscript{4}Ωm。原典はΩmで、H26のΩcmと同じ数値扱いにはしない。一般的金属の代表値10\textsuperscript{-8}〜10\textsuperscript{-6}Ωm程度と桁が違う。\par

\noindent\textbf{4：代表的条件で適切} 300Kと100Kの抵抗。通常の金属では格子振動による散乱が高温ほど増え、室温側の抵抗が大きい。相転移や異常な材料を含む全温度の保証ではない。\par

\noindent\textbf{5：代表的条件で適切} 不純物・格子欠陥と残留抵抗。低温で格子振動の寄与が減っても不純物・欠陥による散乱は残る。組成・構造による寄与を考える。\par

\paragraph{再利用する考え方} 場・応答・担い手・材料・温度を区別し、数値の指数と単位、代表モデルの適用条件を確認する。

\paragraph{共通理論} \hyperref[theory:em:resistivity]{第\ref*{theory:em:resistivity}節}

\par 参考：\cite{openstax_resistivity,em51_resistance}。

\paragraph{関連問題} \hyperref[q:h26-reg-applied-05]{H26-REG-APPLIED-05}, \hyperref[q:r05-reg-applied-08]{R05-REG-APPLIED-08}

\par 難易度 2（校準初稿）、優先度未評価。

\subsection{令和5年度 問8}
\label{q:r05-reg-applied-08}

\noindent\textbf{原典：} sources/R05/17\_applied\_science.pdf、PDF物理3頁。

\textbf{公式正答：} 4（single）。

\paragraph{問題の要旨} 金属抵抗の形状・温度依存と抵抗率の桁を照合する。

\Needspace{4\baselineskip}\paragraph{確認した条件・記述（独自要約）}\mbox{}\par

\noindent\textbf{原典確認：} S7\_051の11原典画像で13問の全五肢・指数・単位・否定を照合。\par

\noindent\textbf{対象：} 金属抵抗の形状・温度依存と抵抗率の桁を照合する。\par

\paragraph{問われる核心} 金属抵抗の形状・温度依存と抵抗率の桁を照合する。

\paragraph{必要知識} 各回の誤りは肢4の抵抗率の桁である。R=ρL/AとI=\allowbreak{}V/Rを、材料・断面・温度一定という条件から使う。残留抵抗と温度依存を分け、H26のΩcm、R01とR05のΩmを独立に校正した。通常金属の温度傾向を特殊な転移や全合金へ広げない。

\paragraph{考え方・正答を選ぶ理由（査読待ち初稿）} 各回の誤りは肢4の抵抗率の桁である。R=ρL/AとI=\allowbreak{}V/Rを、材料・断面・温度一定という条件から使う。残留抵抗と温度依存を分け、H26のΩcm、R01とR05のΩmを独立に校正した。通常金属の温度傾向を特殊な転移や全合金へ広げない。

\Needspace{4\baselineskip}\paragraph{選択肢の検討（査読待ち初稿）}\mbox{}\par

\noindent\textbf{1：代表的条件で適切} 不純物・格子欠陥と残留抵抗。低温で格子振動の寄与が減っても不純物・欠陥による散乱は残る。組成・構造による寄与を考える。\par

\noindent\textbf{2：代表的条件で適切} 一定断面でRが長さに比例。均一な導線はR=ρL/A。同じ材料・温度・断面積を保つ条件が必要。\par

\noindent\textbf{3：代表的条件で適切} −195℃と0℃の抵抗。それぞれ約78Kと273K。通常の金属の比較では0℃側が大きい。摂氏温度を比で使わない。\par

\noindent\textbf{4：不適切} 室温の金属が1〜10\textsuperscript{2}Ωm。原典は1〜100Ωm。一般的な金属の抵抗率の典型値より何桁も大きい。特定合金の値を全金属へ固定しない。\par

\noindent\textbf{5：代表的条件で適切} 電流が電位差に比例。オーム性・温度一定ならI=\allowbreak{}V/R。自己加熱でRが変わる高電圧域まで無条件には広げない。\par

\paragraph{再利用する考え方} 場・応答・担い手・材料・温度を区別し、数値の指数と単位、代表モデルの適用条件を確認する。

\paragraph{共通理論} \hyperref[theory:em:resistivity]{第\ref*{theory:em:resistivity}節}

\par 参考：\cite{openstax_resistivity,em51_resistance}。

\paragraph{関連問題} \hyperref[q:r01-reg-applied-07]{R01-REG-APPLIED-07}

\par 難易度 2（校準初稿）、優先度未評価。

\section{誘電損失・渦電流と表皮深さ}

\subsection{令和3年度 問3}
\label{q:r03-reg-applied-03}

\noindent\textbf{原典：} sources/R03/17\_applied\_science.pdf、PDF物理1頁。

\textbf{公式正答：} 1（single）。

\par\textbf{検討未確定：} 公式1を保持するが肢3の全導電体への履歴損、肢4の厳密な侵入ゼロには疑義。無電極トーチの原回路も未照合でAstraへ残す。

\paragraph{問題の要旨} 誘電損失・渦電流と表皮深さを分けて高周波加熱を読む。

\Needspace{4\baselineskip}\paragraph{確認した条件・記述（独自要約）}\mbox{}\par

\noindent\textbf{原典確認：} S7\_051の11原典画像で13問の全五肢・指数・単位・否定を照合。\par

\noindent\textbf{対象：} 誘電損失・渦電流と表皮深さを分けて高周波加熱を読む。\par

\paragraph{問われる核心} 誘電損失・渦電流と表皮深さを分けて高周波加熱を読む。

\paragraph{必要知識} 公式1は誘電体と導電体の入れ替えに対応する。誘電損失、導電電流のJoule損、磁化履歴損を別々に考える。表皮深さδが有限であるため、肢4の内部へ全く侵入しないという字義を正確な物理定義として教えない。

\paragraph{考え方・正答を選ぶ理由（査読待ち初稿）} 公式1は誘電体と導電体の入れ替えに対応する。誘電損失、導電電流のJoule損、磁化履歴損を別々に考える。表皮深さδが有限であるため、肢4の内部へ全く侵入しないという字義を正確な物理定義として教えない。

\Needspace{4\baselineskip}\paragraph{選択肢の検討（査読待ち初稿）}\mbox{}\par

\noindent\textbf{1：不適切} 誘電加熱の対象を導電体とする。代表的な誘電加熱は誘電体の交流電場に対する損失を利用する。導電体を主対象とする誘導加熱と入れ替えている。有限導電率の試料で両方の損失が共存することはある。\par

\noindent\textbf{2：代表的条件で適切} GHz帯の誘電加熱をマイクロ波加熱。マイクロ波電場で誘電損失を生じさせる方式。高周波誘電加熱の全周波数がGHzという意味ではない。\par

\noindent\textbf{3：条件・語義を留保} 誘導加熱で渦電流損・履歴損。導電体に誘起された渦電流のJoule損が基本。磁化履歴による損失は強磁性体等で寄与し、すべての導電体が履歴損を持つわけではない。\par

\noindent\textbf{4：条件・語義を留保} 内部に侵入しない表皮効果。電流・場は表面近傍へ集中するが、良導体半空間モデルでも深さに対し指数的に減衰し、数学的にゼロの侵入ではない。原典の近似表現を留保する。\par

\noindent\textbf{5：代表的条件で適切} 無極放電・融合プラズマの高周波加熱。プラズマにも高周波電磁場からエネルギーを供給できる。QST資料で融合プラズマのRF加熱を確認。トーチの具体的な無電極回路は別途照合する。\par

\paragraph{再利用する考え方} 場・応答・担い手・材料・温度を区別し、数値の指数と単位、代表モデルの適用条件を確認する。

\paragraph{共通理論} \hyperref[theory:em:heating]{第\ref*{theory:em:heating}節}

\par 参考：\cite{em51_heating,em51_skin,em51_plasma}。

\par 難易度 3（校準初稿）、優先度未評価。

\section{理想変圧器の電圧巻数比}

\subsection{令和5年度 問4}
\label{q:r05-reg-applied-04}

\noindent\textbf{原典：} sources/R05/17\_applied\_science.pdf、PDF物理2頁。

\textbf{公式正答：} 3（single）。

\paragraph{問題の要旨} RLC共振・ソレノイド・理想変圧器の巻数比を照合する。

\Needspace{4\baselineskip}\paragraph{確認した条件・記述（独自要約）}\mbox{}\par

\noindent\textbf{原典確認：} S7\_051の11原典画像で13問の全五肢・指数・単位・否定を照合。\par

\noindent\textbf{対象：} RLC共振・ソレノイド・理想変圧器の巻数比を照合する。\par

\paragraph{問われる核心} RLC共振・ソレノイド・理想変圧器の巻数比を照合する。

\paragraph{必要知識} 誤りは3の逆比。各巻の起電力が共通の磁束変化に比例するので、二次の巻数が多ければ二次電圧も大きい。共振の駆動条件と長いソレノイドの近似条件も別々に確認する。

\paragraph{考え方・正答を選ぶ理由（査読待ち初稿）} 誤りは3の逆比。各巻の起電力が共通の磁束変化に比例するので、二次の巻数が多ければ二次電圧も大きい。共振の駆動条件と長いソレノイドの近似条件も別々に確認する。

\Needspace{4\baselineskip}\paragraph{選択肢の検討（査読待ち初稿）}\mbox{}\par

\noindent\textbf{1：代表的条件で適切} 直列RLCのf=\allowbreak{}1/(2π√LC)。交流電源で駆動する直列回路の電流共振ではωL=\allowbreak{}1/(ωC)。抵抗で減衰する自由振動周波数とは別である。\par

\noindent\textbf{2：代表的条件で適切} ソレノイドBが巻数密度nに比例。長く端効果が小さい、電流・透磁率一定のモデルでB=μnI。飽和域でμを一定としない。\par

\noindent\textbf{3：不適切} E2=\allowbreak{}(N1/N2)E1。理想変圧器は同じ時間変化磁束を共有し、E2/E1=\allowbreak{}N2/N1。原典は分子分母が逆。電圧の最大値同士でも同じ巻数比になる。\par

\noindent\textbf{4：代表的条件で適切} 弱い磁化が外場と同じ向きなら常磁性。通常の線形応答で小さな正の磁化率を持つ。外場なしの自発磁化を持つ強磁性と分ける。\par

\noindent\textbf{5：代表的条件で適切} Fe・Ni・Co等の強磁性体をコアに利用。代表的な強磁性元素で磁束を集中できる。実際の交流コアの選定では透磁率だけでなく保磁力・渦電流損・周波数を考える。\par

\paragraph{再利用する考え方} 場・応答・担い手・材料・温度を区別し、数値の指数と単位、代表モデルの適用条件を確認する。

\paragraph{共通理論} \hyperref[theory:em:rlc]{第\ref*{theory:em:rlc}節}, \hyperref[theory:em:solenoid]{第\ref*{theory:em:solenoid}節}, \hyperref[theory:em:transformer]{第\ref*{theory:em:transformer}節}

\par 参考：\cite{mit_em12,mit_em09,em51_transformer_ref,openstax_magnetism}。

\par 難易度 2（校準初稿）、優先度未評価。

\section{熱電効果・可逆熱と散逸}

\subsection{平成29年度 問7}
\label{q:h29-reg-applied-07}

\noindent\textbf{原典：} sources/H29/17\_applied\_science.pdf、PDF物理3頁。

\textbf{公式正答：} 5（single）。

\par\textbf{検討未確定：} 候補4は材料・端点温度・開回路を省略している。起電力と負荷時の端子電圧の区別を付し、一般に形状が無関係とはしない。

\paragraph{問題の要旨} ゼーベック・ペルティエ・トムソン効果を、温度依存と電流依存で区別する。

\Needspace{4\baselineskip}\paragraph{確認した条件・記述（独自要約）}\mbox{}\par

\noindent\textbf{確認範囲：} 問題幹、数値、式、否定、全5候補を原典画像で照合した。\par

\noindent\textbf{成立条件：} 熱起電力は開回路で同じ材料・同じ両端温度を比較する。端点温度と材料を固定した熱起電力と、負荷を接続した端子電圧を区別する。\par

\paragraph{問われる核心} ゼーベック・ペルティエ・トムソン効果を、温度依存と電流依存で区別する。

\paragraph{必要知識} 熱起電力は開回路で同じ材料・同じ両端温度を比較する。端点温度と材料を固定した熱起電力と、負荷を接続した端子電圧を区別する。

\paragraph{考え方・正答を選ぶ理由（査読待ち初稿）} 温度勾配を持つ導体のトムソン熱は電流に一次であり、電流反転で吸熱と発熱が入れ替わる。\(I^2R\) に従い電流反転で符号が変わらないのはジュール熱である。したがって5が不適切である。開回路の熱起電力は材料と端点温度が同じなら熱電能の温度積分で決まり、長さ・断面積を含まない。

\Needspace{4\baselineskip}\paragraph{選択肢の検討（査読待ち初稿）}\mbox{}\par

\noindent\textbf{1：適切} ゼーベック効果は温度差測定に利用。異種材料の接点温度差から生じる起電力を熱電対の測温に用いる。\par

\noindent\textbf{2：適切} ペルティエ効果は熱電冷却に応用。異種材料の接合部で電流方向に応じて吸熱又は発熱が起こる。\par

\noindent\textbf{3：適切} 熱電能は温度依存。ゼーベック係数S(T)は一般に材料と温度に依存し、一定値とするのは限定した温度範囲の近似である。\par

\noindent\textbf{4：条件付きで適切} 熱起電力は長さや太さに無関係。材料と端点温度を固定した開回路の起電力には幾何形状が直接入らない。電気抵抗や負荷時の端子電圧、温度分布まで形状に依存しないという意味ではない。\par

\noindent\textbf{5：不適切} トムソン熱は電流の2乗に比例。トムソン熱は電流と温度勾配に一次で可逆である。二乗依存はジュール熱の特徴である。\par

\paragraph{再利用する考え方} 接合部のペルティエ熱、単一導体内のトムソン熱、散逸によるジュール熱の場所と電流反転時の符号を比べる。

\paragraph{共通理論} \hyperref[theory:electrochem:seebeck]{第\ref*{theory:electrochem:seebeck}節}, \hyperref[theory:electrochem:peltier]{第\ref*{theory:electrochem:peltier}節}, \hyperref[theory:electrochem:thomson]{第\ref*{theory:electrochem:thomson}節}, \hyperref[theory:electrochem:joule]{第\ref*{theory:electrochem:joule}節}

\par 参考：\cite{mit_thermoelectric2024}。

\paragraph{関連問題} \hyperref[q:r03-reg-applied-06]{R03-REG-APPLIED-06}

\par 難易度 2（校準初稿）、優先度未評価。

\subsection{令和3年度 問6}
\label{q:r03-reg-applied-06}

\noindent\textbf{原典：} sources/R03/17\_applied\_science.pdf、PDF物理3頁。

\textbf{公式正答：} 2（single）。

\paragraph{問題の要旨} 熱電現象の名称から光学現象を区別する。

\Needspace{4\baselineskip}\paragraph{確認した条件・記述（独自要約）}\mbox{}\par

\noindent\textbf{確認範囲：} 問題幹、数値、式、否定、全5候補を原典画像で照合した。\par

\noindent\textbf{成立条件：} 最も不適切な用語を選ぶ。カー効果には電気光学と磁気光学の用法があるが、いずれもこの熱電効果の名称ではない。\par

\paragraph{問われる核心} 熱電現象の名称から光学現象を区別する。

\paragraph{必要知識} 最も不適切な用語を選ぶ。カー効果には電気光学と磁気光学の用法があるが、いずれもこの熱電効果の名称ではない。

\paragraph{考え方・正答を選ぶ理由（査読待ち初稿）} ゼーベック、ペルティエ、トムソンと熱起電力はいずれも熱と電気の結合に関わる。カー効果は電場による複屈折又は磁気光学的な偏光変化などの光学現象の名称であり、列挙された熱電現象と区別して2を選ぶ。

\Needspace{4\baselineskip}\paragraph{選択肢の検討（査読待ち初稿）}\mbox{}\par

\noindent\textbf{1：適切} ゼーベック効果。温度差により起電力が生じる熱電現象である。\par

\noindent\textbf{2：不適切} カー効果。光の屈折・偏光に関わる名称であり、温度差と電流の相互作用を表す熱電効果ではない。\par

\noindent\textbf{3：適切} ペルティエ効果。異種材料接合部を流れる電流に伴う吸放熱である。\par

\noindent\textbf{4：適切} トムソン効果。温度勾配を持つ材料内の電流に伴う可逆的な吸放熱である。\par

\noindent\textbf{5：適切} 熱起電力。温度差による起電力そのものを表す用語である。\par

\paragraph{再利用する考え方} 似た人名の効果を、外部の原因と観測量で分類する。熱電は温度・電流・熱流、カー効果は電磁場と光学応答である。

\paragraph{共通理論} \hyperref[theory:electrochem:seebeck]{第\ref*{theory:electrochem:seebeck}節}, \hyperref[theory:electrochem:peltier]{第\ref*{theory:electrochem:peltier}節}, \hyperref[theory:electrochem:thomson]{第\ref*{theory:electrochem:thomson}節}, \hyperref[theory:electrochem:kerr]{第\ref*{theory:electrochem:kerr}節}

\par 参考：\cite{mit_thermoelectric2024,mit_kerr_calendar2005}。

\paragraph{関連問題} \hyperref[q:h29-reg-applied-07]{H29-REG-APPLIED-07}

\par 難易度 2（校準初稿）、優先度未評価。


\chapter{光学}
\section{光線の幾何と場の位相}
幾何光学は光の波長より十分大きい構造で光線と像を扱い、波動光学は重ね合わせと位相を扱う。
レンズの実像・虚像は光線の交点の問題であり、回折・干渉の明暗は場の位相差の問題である。
同じ屈折率が両方へ関係するが、条件を取り違えないことが重要である。
光電効果と電子線回折の比較は現代物理の第\ref{theory:modern:duality}節へ分類する。
幾何から波動へという構成は大学の光学講義にも共通する\cite{mit_optics2014}。

\section{薄レンズの符号で実像と虚像を分ける}\label{theory:optics:lens}
近軸・薄レンズ近似において、凸レンズの焦点距離を \(f>0\)、入射側の実物体までを \(a>0\) とする。
出射側の像を \(b>0\)、入射側の虚像を \(b<0\) と定めれば
\[
 \frac1f=\frac1a+\frac1b,\qquad b=\frac{fa}{a-f},\qquad M=-\frac ba
\]
である。\(M\) は像高／物体高であり、正なら正立、負なら倒立である。
焦点外 \(a>f\) は倒立実像、焦点内 \(0<a<f\) は正立虚像となる。
\(a=f\) なら出射光線は平行となり、有限の像面はない\cite{os_lens}。
「虚像だから小さい」という推論は正しくない。焦点内の凸レンズでは \(|M|>1\) となる。
実像の移動方向は
\[
 \frac{db}{da}=-\frac{f^2}{(a-f)^2}<0
\]
から直接判定する。焦点外の物体がレンズへ近付くと像は遠ざかる。
焦点を越えて物体を動かす場合、有限の実像がそのままレンズを通過して虚像になるわけではなく、無限遠を経由する。

\section{二枚のレンズは中間像の位置を保存する}\label{theory:optics:cascade}
第一レンズを \(x=0\)、第二を \(x=\ell\) とし、第一の像位置を \(x=b_1\) とする。
第二へ来る光は、第一の像位置から発散又はその位置へ収束する束である。
第二からの物体距離は
\[
 a_2=\ell-b_1
\]
であり、\(b_1<0\) の第一の虚像なら \(a_2>\ell>0\) となる。
第一の「虚像」が第二に対しても必ず「虚物体」というわけではない。
逆に、第一の実像予定位置が第二のさらに後方にあれば \(a_2<0\) の虚物体となる。
各レンズの式を順に解き、最後に同じ座標軸で位置を照合する。
レンズ間隔を各物体距離へそのまま置く方法では中間像の位置を失ってしまう。

\section{屈折率・全反射・色収差}\label{theory:optics:refraction}
屈折率は光の位相速度を真空の光速 \(c\) で割った逆数 \(n=c/v_{\mathrm{ph}}\) である。
境界に平行な波数成分の保存から \(n_1\sin\theta_1=n_2\sin\theta_2\) を得る。
高屈折率側から低屈折率側へ入射し、\(\sin\theta_1>n_2/n_1\) となると、
伝播する屈折波を作れず全反射となる。
したがって光ファイバのコアはクラッドより高い屈折率を持つ\cite{os_fiber}。
境界外にはエバネッセントな場があり、「外側の電磁場が必ず厳密に零」とはしない。
通常の透明域の正常分散では紫の屈折率が赤より大きい。
同形の薄い凸レンズを空気中に置くと、曲率に依存する定数 \(C>0\) に対して \(1/f=(n-1)C\) となるので、
紫は焦点距離が短く、赤は遠くへ集光する。これは正常分散の範囲であり、共鳴近傍の異常分散まで無条件に含めない。

\section{反射位相を経路差へ足す}\label{theory:optics:reflection}
損失のない誘電体へほぼ垂直に入射する電場の反射係数は \(r=(n_1-n_2)/(n_1+n_2)\) である。
低屈折率側から高屈折率側なら \(r<0\) で位相が \(\pi\) 反転し、逆なら反転しない。
膜の両面からの反射光を比較するときは、光学的経路差と反射位相を別に加える\cite{os_thinfilm}。
角度・偏光・吸収まで一般化すると反射係数はこれより複雑になる。
問題に角度がない場合の通常の説明を、全入射角に同じ位相反転が成立する普遍則として教えない。

\section{ニュートンリングの幾何と次数}\label{theory:optics:newton}
曲率半径 \(R\) の凸面を平面へ置くと、接点から半径 \(x\) における空気層の厚さは
\[
 t(x)=R-\sqrt{R^2-x^2}\simeq\frac{x^2}{2R}\quad(x\ll R)
\]
となる。垂直入射の反射光を観察し、一方だけ位相反転すると、暗環は \(2t=n\lambda\)、明環は
\[
 2t=(n+1/2)\lambda,\qquad x_n^2=(n+1/2)R\lambda,\qquad n=0,1,2,\ldots
\]
である。接点は暗く、最初の明環は \(x^2=R\lambda/2\) となる。
ここで \(n\) は零始まりの干渉次数である。
「第 \(m\) 番目」を \(m=1,2,\ldots\) とする序数なら \(n=m-1\) であり、\(x_m^2=(m-1/2)R\lambda\) と書く。
開始値を指定せず両方を同じ記号で使わない。
令和3年度の原典は「\(m\) 番目」とする一方、公式肢は零始まりの形であるため解釈疑義を残している。
同じ曲率・同じ次数で波長が長いと半径が増すことは、この開始値の違いによらない。

\section{回折格子：位置・幅・強度を分ける}\label{theory:optics:grating}
垂直入射、間隔 \(d\) の等間隔スリットから遠方へ向かう光の隣接経路差は \(d\sin\theta\) である。
全てのスリットが同位相となる主極大は
\[
 d\sin\theta=m\lambda,\qquad m=0,\pm1,\pm2,\ldots
\]
であり、スクリーン距離 \(L\) と小角近似で主極大間隔は \(\Delta x\simeq L\lambda/d\) となる\cite{os_grating}。
入射強度は線形領域で全体の明るさを変え、この位置を変えない。
スリット数 \(N\) の効果は、位相差 \(\delta=2\pi d\sin\theta/\lambda\) を用いて和を取ると分かる。
\[
 \left|\sum_{j=0}^{N-1}e^{ij\delta}\right|^2
 =\left(\frac{\sin(N\delta/2)}{\sin(\delta/2)}\right)^2
\]
であり、\(\delta=2\pi m\) は極限値 \(N^2\) の主極大である。
主極大間には \(N-1\) 個の零点と \(N-2\) 個の副極大がある。
本数を増すと主極大は狭くなり、同じ主極大間の細かい暗線は増える。
有限幅スリットの単スリット包絡も別に掛かる\cite{os_double}。
したがって「縞間隔」を主極大間隔と読むか、全ての暗線の間隔と読むかを区別する。
平成27年度の「明線（もしくは暗線）」という表現には、この適用範囲の疑義を保存した。
格子本数は \(1/d\) である。波長を nm、密度を本/mmで求める場合は、まず \(d\) を mmへ換算する。
第一次数の650 nmと30度なら \(d=1.3\,\mu\mathrm m\) で、約769本/mmとなる。
これは主極大の方位を測る問題であり、縞の幅や副極大を数える問題ではない。

\section{電磁波領域と光子エネルギー}\label{theory:optics:spectrum}
真空中では \(\lambda f=c\) である。
波長の長い順は、ここで問う4領域では赤外・可視・紫外・γである\cite{os_spectrum}。
光子エネルギー \(E=hf=hc/\lambda\) の大きい順なら逆になる。
ビーム全体のエネルギーは光子数にも依存するので、波長の短い光が任意の照射条件で必ず大きな総エネルギーを持つとは言えない。
この区別は分光、放射線、光電効果の問題へつながる。

\section{出題の型と原典確認の範囲}
像の問題は物体と像の位置の符号、干渉は反射位相・経路差・次数、格子は間隔と本数の逆数、
スペクトルは比較する量を先に確定する。
S7で光学の11登録問を確認し、電子線回折の1問を現代物理へ移した。
対象14問の解説初稿を本章へ置く。表の空欄や次数の疑義は解説と未解決表で追跡する。

\section{登録問題と解説初稿}
\subsection{平成23年度 問7}
\label{q:h23-reg-applied-07}

\noindent\textbf{原典：} sources/H23/17\_applied\_science.pdf、PDF物理3頁。

\textbf{公式正答：} 2（single）。

\par\textit{原典本文の校正と独自解説は未作成。分類のみの登録項目。}

\subsection{平成24年度 問3}
\label{q:h24-reg-applied-03}

\noindent\textbf{原典：} sources/H24/17\_applied\_science.pdf、PDF物理1頁。

\textbf{公式正答：} 1（single）。

\par\textit{原典本文の校正と独自解説は未作成。分類のみの登録項目。}

\subsection{平成24年度 問7}
\label{q:h24-reg-applied-07}

\noindent\textbf{原典：} sources/H24/17\_applied\_science.pdf、PDF物理2頁。

\textbf{公式正答：} 2（single）。

\par\textit{原典本文の校正と独自解説は未作成。分類のみの登録項目。}

\subsection{平成26年度 問7}
\label{q:h26-reg-applied-07}

\noindent\textbf{原典：} sources/H26/17\_applied\_science.pdf、PDF物理3頁。

\textbf{公式正答：} 2（single）。

\par\textit{原典本文の校正と独自解説は未作成。分類のみの登録項目。}

\subsection{平成27年度 問3}
\label{q:h27-reg-applied-03}

\noindent\textbf{原典：} sources/H27/17\_applied\_science.pdf、PDF物理1頁。

\textbf{公式正答：} 1（single）。

\par\textit{原典本文の校正と独自解説は未作成。分類のみの登録項目。}

\subsection{平成28年度 問5}
\label{q:h28-reg-applied-05}

\noindent\textbf{原典：} sources/H28/17\_applied\_science.pdf、PDF物理2頁。

\textbf{公式正答：} 3（single）。

\par\textit{原典本文の校正と独自解説は未作成。分類のみの登録項目。}

\subsection{令和元年度再試験 問3}
\label{q:r01-ret-applied-03}

\noindent\textbf{原典：} sources/R01\_retest/17\_applied\_science.pdf、PDF物理2頁。

\textbf{公式正答：} 5（single）。

\par\textit{原典本文の校正と独自解説は未作成。分類のみの登録項目。}

\subsection{令和2年度 問5}
\label{q:r02-reg-applied-05}

\noindent\textbf{原典：} sources/R02/17\_applied\_science.pdf、PDF物理2頁。

\textbf{公式正答：} 5（single）。

\par\textit{原典本文の校正と独自解説は未作成。分類のみの登録項目。}

\subsection{令和2年度 問7}
\label{q:r02-reg-applied-07}

\noindent\textbf{原典：} sources/R02/17\_applied\_science.pdf、PDF物理3頁。

\textbf{公式正答：} 4（single）。

\par\textit{原典本文の校正と独自解説は未作成。分類のみの登録項目。}

\subsection{令和3年度 問5}
\label{q:r03-reg-applied-05}

\noindent\textbf{原典：} sources/R03/17\_applied\_science.pdf、PDF物理2頁。

\textbf{公式正答：} 3（single）。

\par\textit{原典本文の校正と独自解説は未作成。分類のみの登録項目。}

\subsection{令和6年度 問5}
\label{q:r06-reg-applied-05}

\noindent\textbf{原典：} sources/R06/17\_applied\_science.pdf、PDF物理2頁。

\textbf{公式正答：} 5（single）。

\par\textit{原典本文の校正と独自解説は未作成。分類のみの登録項目。}


\chapter{現代物理・原子核}
\section{背景：観測する量と反応で保存する量を分ける}
量子論では波動的な伝播・干渉と離散的なエネルギー授受が両立する。光と電子を「波か粒子か」の一語へ分類するのでなく、どの実験で何を観測したかを問う。
原子核では粒子の放出によって質量数Aや原子番号Zが変わる場合と、核種を保ったまま準位が変わる場合がある。放射線の透過性、壊変率、吸収エネルギー、生物影響を考慮した防護量も別の観測量である。
本章には先行の光量子比較1問と今回の現代物理8問を収める。原典と全候補を確認した解説初稿であり、平成23年度のβ安定への全称は科学的疑義を保持する。

\section{波と粒子のどちらを観測しているか}\label{theory:modern:duality}
量子論では光のエネルギー・運動量は \(E=hf\)、\(p=h/\lambda\) である。
光電効果では入射光子のエネルギーから仕事関数 \(W\) を差し引き、
\(K_{\max}=hf-W\) として光電子の最大運動エネルギーを説明する\cite{os_photoelectric}。
強度を上げることと一光子のエネルギーを上げることを分ける。
コンプトン散乱は光子と電子のエネルギー・運動量保存から波長変化を説明する例である。
一方、電子線回折は電子が運動量 \(p\) に対する物質波の波長 \(\lambda=h/p\) を持つことを示す。
結晶の周期構造による干渉は電子の波動性の観測であり、光の粒子性そのものの直接観測ではない\cite{os_matter}。
粒子名だけで「粒子性」と判断せず、現象が示す対象と性質を識別する。
平成24年度の比較問題をこの節へ移している。光電効果で変わる一光子エネルギーと、電子線の運動量に対応する干渉波長を混同しない。

\section{物質波：同速・同エネルギーは別の条件}\label{theory:modern:matter}
粒子の物質波長は \(\lambda=h/p\) である。同速vの自由粒子を比較すると \(p=\gamma mv\) の \(\gamma=(1-v^2/c^2)^{-1/2}\) が共通なので \(\lambda\propto1/m\) となる。電子、陽子、α核の静止質量はこの順に大きく、波長はこの順に長い。
非相対論的な同運動エネルギーKなら \(p=\sqrt{2mK}\) なので \(\lambda\propto m^{-1/2}\) となる。同速と同Kで質量の依存指数は異なる。同じ運動量で比較するなら波長は等しい。
電子線の回折から運動量と結晶構造を読む用途へつながる。電子の電荷の大きさだけでは波長順を求められない\cite{os_matter}。

\section{連続X線の最短端：一光子へのエネルギー上限}\label{theory:modern:xray}
初期運動エネルギーが無視できる電子を電圧Vで加速すれば、電場の仕事は \(eV\) である。金属標的での減速に伴う制動放射は連続スペクトルを持つが、一光子のエネルギーは電子が持つエネルギーを超えない。
\[
 hf_{\max}=eV,\qquad
 \lambda_{\min}=\frac{hc}{eV}.
\]
この上限はエネルギー保存と光子 \(E=hc/\lambda\) による。電子速度を \(mv^2/2\) から求める非相対論近似は不要である。金属の特性X線の離散線と連続線の端を区別する。
平成27年度は指定の \(h=6.6\times10^{-34}\)、\(e=1.6\times10^{-19}\)、\(c=3.0\times10^8\)、35kVを使う。現在の定義定数に勝手に置き換えず、約 \(3.54\times10^{-11}\,\mathrm m\) を得る。波長の長い光子が存在し得ることは、それを最短端の正答にする理由にならない\cite{os_photoelectric}。

\section{光素子：吸収・再結合・誘導放出}\label{theory:modern:devices}
半導体や分子の電子にはエネルギー状態がある。光子を吸収して高い状態へ移る過程と、励起された状態から低い状態へ移って光子を出す過程を使い分ける。
太陽電池は吸収した光によるキャリア生成と分離を使い、電流を取り出す。LEDは注入した電子・正孔の放射再結合を使う。有機ELは有機半導体の電子励起・励起子等の放射緩和を使う。全ての生成キャリアが光へ戻るわけではなく非放射過程も競合する\cite{modern_pv,modern_led,modern_oled}。
半導体レーザーでは、励起された活性層で誘導放出を増幅し共振器の帰還と組み合わせる。LEDの自発放出と、光場に相関した増幅を行うレーザーの機構を区別する。励起と光が関係するという同じ枠組みでも、素子の役割は同じではない\cite{modern_laser}。

\section{超伝導を光素子との相対比較で読む}\label{theory:modern:superconduct}
超伝導は臨界条件下で直流の無抵抗と特徴的な磁場応答を示す巨視的量子状態である。通常のBCS理論では電子対の凝縮を説明し、全ての超伝導体の微視的媒介機構を同じフォノン機構へ固定するわけではない。
電子が光を吸収・放出することを中心原理とする前節の4素子に対し、超伝導という現象の中心は凝縮状態の電気磁気応答である。令和5年度の比較はこの違いで判断する。超伝導体が電磁場と相互作用しない、光で状態を変えられないという断言へ広げない\cite{modern_superconduct,modern_pairs}。

\section{核種と放射線：A・Zの帳簿}\label{theory:modern:nuclear}
核種は陽子数Zと中性子数Nで指定し、A=Z+Nである。同位体は同Z・異Aであり、同質量数核とは異なる。
α放出では \((A,Z)\mapsto(A-4,Z-2)\)、βマイナスでは \((A,Z)\mapsto(A,Z+1)\)、βプラス又は電子捕獲では \((A,Z)\mapsto(A,Z-1)\) となる。γ脱励起はA・Zを保ち励起エネルギーを放射へ渡す。ニュートリノを含むエネルギー・運動量等の保存も必要である。
U・Thの主要天然壊変系列はα・βの連なりを経て安定Pb同位体へ至る。系列中でZが変わることと、個々のγ放出がZを変えないことは矛盾しない\cite{modern_mit_nuclear,modern_series}。

\section{β安定：質量最小という省略の限界}\label{theory:modern:stability}
同Aで質量をZの関数にすると、液滴模型のクーロン項と非対称項からおおよそ放物線を得る。ただし核の偶奇に伴う対エネルギーがあり、偶数Aでは偶偶核と奇奇核の曲線が異なる。
β壊変は隣接Zの終状態へQ値が正となる場合にエネルギー的に可能になる。原子質量を用いたβマイナスは、電子束縛等を省く近似で \(Q_{\beta^-}=[M(A,Z)-M(A,Z+1)]c^2\) である。βプラスの原子質量式には \(2m_ec^2\) のしきい値がある。可能性と寿命は核状態の選択則等も含み同じではない。
\begin{tcolorbox}[title=平成23年度の最小質量への全称を保持しない]
偶数Aでは単一βに対して安定な核が複数あり得る。中間の奇奇核が高い質量にあるため、一つのβ遷移で最小質量へ進めない場合がある。二重β壊変の観測はあるが、その経路・極めて長い時間尺度は単一β連鎖と同一ではない。「全ての同A核が最小質量核へβ壊変する」という原典の省略を、無条件の定理として教材へ移さない。
\end{tcolorbox}
公式2という熱中性子の取り違えは独立に説明できるが、3の一般傾向と全称の科学的一意性を分けて査読へ送る\cite{modern_mit_mass,modern_doublebeta,modern_mit_nuclear}。

\section{核分裂性と中性子のエネルギー}\label{theory:modern:fission}
重い核が中性子を吸収し、複合核の励起が分裂障壁等に対して有効となると核分裂が起こり得る。核種によって必要な中性子エネルギーが違う。
235U等は熱中性子で核分裂しやすい核分裂性核種である。238Uは高速中性子による核分裂が可能で、熱中性子の捕獲からβ壊変を経てPuへつながる親物質でもある。
熱中性子炉の連鎖反応で、自然存在量の多い238Uをそのまま熱核分裂の主役として説明しない。核分裂可能と熱核分裂性、単発反応と持続連鎖を分ける\cite{modern_fastfission}。

\section{壊変率と半減期：何が減った割合か}\label{theory:modern:decay}
同じ核種の単位時間の壊変確率を一定の \(\lambda\) とし、生成・流出入がなければ
\[
 \frac{dN}{dt}=-\lambda N,\qquad
 N(t)=N_0e^{-\lambda t},\qquad
 T_{1/2}=\frac{\ln2}{\lambda},\qquad
 \frac{N(t)}{N_0}=2^{-t/T_{1/2}}.
\]
半減期は一つの核が必ずその時刻に壊変する時間ではなく、集団の平均残存が半分になる時間である。放射能 \(\mathcal A=\lambda N\) も同じ核種なら同じ比で減る。
90年/30年は3半減期で1/8、30日/8日は約0.0743、30日/9日は約0.0992である。整数の半減期数に丸めず指数で評価する。元の親核種の残存比と壊変済み割合 \(1-N/N_0\)、娘核種量を区別する。娘も壊変する系列では、娘量はこの単純な差だけではない\cite{modern_mit_nuclear}。

\section{吸収線量・等価線量・実効線量}\label{theory:modern:dose}\label{theory:modern:activity}
吸収線量Dは単位質量当たりに吸収された放射線エネルギーで、Gy=J/kgである。Bqは壊変率であり、核種の活動度だけからそのままGyへ換算できない。
組織Tに放射線Rが与える平均吸収線量 \(D_{T,R}\) に放射線の種類等の無次元重み \(w_R\) を付け、
\[
 H_T=\sum_Rw_RD_{T,R},\qquad E=\sum_Tw_TH_T
\]
を等価線量と実効線量とする。どちらもSvである。実効線量の \(w_T\) は組織の寄与を重み付ける。原典の「荷重係数」という表記を校正で別の語へ置換せず、重みとして意味を説明する。
GyとSvは同じJ/kgの次元であるが量の定義が違う。等価線量・実効線量は防護上の量であり、高線量の確定的影響や個人への治療判断をそれだけで表す量へ転用しない。
α線は外部透過が弱くても、内部取込みでは高い密度でエネルギーを与え得る。実際の影響には量・体内分布・経路等が必要である。「紙で止まるから常に影響なし」「内部なら常に同じ被害」のどちらも単位の定義からは導けない\cite{modern_dose,modern_alpha}。

\section{出題形式と理解の接続}
今回の8問は、同速の物質波1、制動X線端1、半減期3、核反応比較1、線量対応1、光電子素子比較1である。先行の光量子比較を加え9問を収録する。これらは確認済みの当該問題集合の集計であり将来出題の保証ではない。
エネルギー上限、運動量と波長、保存するA・Z、指数減衰、測定量の意味の順に整理すると、式を暗記するだけでなく条件の変更へ転用できる。地球の同位体年代では親娘閉鎖と試料反応がさらに必要であり、本章の壊変式だけで地質年代を確定しない。

\section{過去問と全候補の検討}
\subsection{平成23年度 問10}
\label{q:h23-reg-applied-10}

\noindent\textbf{原典：} sources/H23/17\_applied\_science.pdf、PDF物理4頁。

\textbf{公式正答：} 2（single）。

\par\textit{原典本文の校正と独自解説は未作成。分類のみの登録項目。}

\subsection{平成25年度 問7}
\label{q:h25-reg-applied-07}

\noindent\textbf{原典：} sources/H25/17\_applied\_science.pdf、PDF物理4頁。

\textbf{公式正答：} 5（single）。

\par\textit{原典本文の校正と独自解説は未作成。分類のみの登録項目。}

\subsection{平成26年度 問9}
\label{q:h26-reg-applied-09}

\noindent\textbf{原典：} sources/H26/17\_applied\_science.pdf、PDF物理4頁。

\textbf{公式正答：} 2（single）。

\par\textit{原典本文の校正と独自解説は未作成。分類のみの登録項目。}

\subsection{平成27年度 問5}
\label{q:h27-reg-applied-05}

\noindent\textbf{原典：} sources/H27/17\_applied\_science.pdf、PDF物理2頁。

\textbf{公式正答：} 2（single）。

\par\textit{原典本文の校正と独自解説は未作成。分類のみの登録項目。}

\subsection{平成27年度 問6}
\label{q:h27-reg-applied-06}

\noindent\textbf{原典：} sources/H27/17\_applied\_science.pdf、PDF物理3頁。

\textbf{公式正答：} 5（single）。

\par\textit{原典本文の校正と独自解説は未作成。分類のみの登録項目。}

\subsection{平成28年度 問9}
\label{q:h28-reg-applied-09}

\noindent\textbf{原典：} sources/H28/17\_applied\_science.pdf、PDF物理4頁。

\textbf{公式正答：} 3（single）。

\par\textit{原典本文の校正と独自解説は未作成。分類のみの登録項目。}

\subsection{令和4年度 問9}
\label{q:r04-reg-applied-09}

\noindent\textbf{原典：} sources/R04/17\_applied\_science.pdf、PDF物理3頁。

\textbf{公式正答：} 2（single）。

\par\textit{原典本文の校正と独自解説は未作成。分類のみの登録項目。}

\subsection{令和5年度 問7}
\label{q:r05-reg-applied-07}

\noindent\textbf{原典：} sources/R05/17\_applied\_science.pdf、PDF物理3頁。

\textbf{公式正答：} 4（single）。

\par\textit{原典本文の校正と独自解説は未作成。分類のみの登録項目。}


\chapter{原子・分子・結晶}
\section{背景：状態・構造・物性を三段階でつなぐ}
原子の電子状態を記述し、複数原子がどのように結合するかを考え、並んだ結晶の物性へ進む。電子配置から導電性までを一つの暗記表にすると、共有結合なら導電性が高い、胞の原子数なら配位数と同じ、といった取り違えが生じる。物理の量子状態と幾何を使い、それぞれの定義へ戻って判断する。

\section{空間軌道と電子状態}\label{theory:atomic:orbital}
非相対論的な中心力近似で、空間軌道は主量子数\(n\)、方位量子数\(l\)、磁気量子数\(m_l\)で表す。
\[
 n=1,2,\ldots,\qquad l=0,\ldots,n-1,\qquad m_l=-l,\ldots,l.
\]
\(l=0,1,2,3\)をs,p,d,fと呼ぶ。電子状態にはさらに\(m_s=\pm1/2\)が必要であり、水素原子の電子にもスピンはある。空間軌道と、そこに置く電子の状態は同じ対象ではない\cite{atomic_quantum}。

Pauli原理により同じ空間軌道の電子は最大2個で逆スピンである。従ってs副殻は2、pは6、dは10、fは14電子まで入る。中性酸素は8電子で\(1s^2 2s^2 2p^4\)。総電子数が8でも\(2s^3\)や\(2s^4\)は容量を超える。原子番号と電子数が等しいのは中性原子で、イオンでは電子増減を数える。

通常の基底状態の構築順では3pの後に4sが入るが、全原子・励起状態・イオンで一定の軌道エネルギー順があるとはしない。Cuは\([\mathrm{Ar}]3d^{10}4s^1\)、Znは\([\mathrm{Ar}]3d^{10}4s^2\)である\cite{atomic_nist2019}。

\begin{tcolorbox}[title=公式正答と独立判断の相違を残す]
令和元年度問14の公式は1、独立初回は4。原典の4は水素の空間軌道と多電子原子のスピンを対比し、電子状態の意味で読むと水素のスピンも必要になる。一方、1に含まれる\(\phi\)を「存在しない」と断言すると、2014年のthorocene研究で報告された\(\phi\)軌道相互作用と整合しない\cite{atomic_phi2014}。多中心の軌道相互作用と純粋な二原子\(\phi\)結合を同一視もしない。公式の採点を保持し、この不一致を未解決の査読対象とする。
\end{tcolorbox}

\section{結合次数と電子対配置}\label{theory:atomic:bonding}
分子軌道の結合次数は\((N_b-N_a)/2\)で、結合性だけでなく反結合性の電子も数える。代表的な\(\mathrm{N_2,O_2,F_2}\)は3,2,1であり、順に共有電子が増え結合が強くなるという説明は逆である。個別分子の結合長・エネルギーは分子軌道占有と合わせて読む。

VSEPRは中心原子まわりの電子領域を数える近似である。水は結合対2、孤立対2のAB2E2。4電子領域の配置と分子形を区別し、CH4は四面体、NH3は三角錐、H2Oは折れ線となる。sp3という共通モデルは、三分子の全原子が同じ四面体形という意味ではない。d関与の形式的混成表記も、全超原子価分子の実在結合をd軌道占有だけで説明する保証ではない。

\section{分子間力：ポテンシャルを微分して力を得る}\label{theory:atomic:interactions}
無極性分子でも瞬間的な電子密度の偏りが隣の分子に双極子を誘起し、分散力が生じる。典型的な水素結合では、電気陰性度の高い原子に結合したHと別の電子対を持つ原子が供与体・受容体として相互作用する。同種分子に限らない\cite{atomic_nonbond}。対称なCH4は永久双極子を持たず、強い分子極性を前提とする水素結合の説明にはならない。

12--6 Lennard--Jonesモデルは
\[
 U(r)=4\varepsilon\left[(\sigma/r)^{12}-(\sigma/r)^6\right],\qquad
 F_r=-\frac{dU}{dr}=24\varepsilon\left[\frac{2\sigma^{12}}{r^{13}}-\frac{\sigma^6}{r^7}\right].
\]
ポテンシャルの反発は正の12乗項、引力は負の6乗項で、力の指数は13と7へ変わる\cite{atomic_lj}。平成29年度問11では反発・引力の逆転を説明できるが、Wの凝集エネルギーを配位数だけへ帰する別候補にも疑義が残る。金属の多電子状態を一つの原子対のモデルへ無条件還元しない。

\section{格子・基底・最近接距離}\label{theory:atomic:geometry}
格子は等価な点を並進で繰り返す数学的集合、基底は各格子点へ付く原子群である。3次元Bravais格子は14種類である\cite{atomic_bravais_ref}。慣用立方胞のFCCは原子数4、最近接配位数12、BCCは原子数2、最近接配位数8。同じ数ではない\cite{openstax_crystals}。

FCCの隅から面中心へのベクトルは\((a/2,a/2,0)\)、BCCの隅から体中心へは\((a/2,a/2,a/2)\)だから
\[
 d_{\mathrm{FCC}}=a/\sqrt2,\qquad d_{\mathrm{BCC}}=\sqrt3a/2.
\]
過去問のdが最近接距離なのか格子定数なのかを各回で確認する。FCCでdからaなら\(a=\sqrt2d\)。BCCでaからdなら\(d=\sqrt3a/2\)。原子半径をrとし接触硬球を置く場合はd=2rである。

HCPの慣用胞規約・理想軸比・すべり条件は第\ref{theory:material:crystals}節と合わせて読む。元素例の構造は温度・圧力・組成に依存する。純Feの通常圧の912℃構造変態と、770℃付近の磁気転移を同一視しない。アモルファスは長距離並進秩序を持たないが短距離の局所秩序はあり得る。

逆格子の原点まわりでWigner--Seitz胞を作った領域が第一Brillouin域である。実格子の胞と逆格子の胞を区別する。diamondはFCCに二原子基底を付け、第一近接4、第二近接12を持つ。

\section{単位胞の密度：分子量と長さの単位をそろえる}\label{theory:atomic:density}
慣用立方胞の原子数z、モル質量M、格子定数aなら
\[
 \rho=\frac{zM}{N_Aa^3}.
\]
胞質量と胞体積の比であり、配位数をzへ代入しない。平成26年度はBCCのz=2、M=52g/mol、a=0.29nm=\(2.9\times10^{-8}\)cmから7.11g/cm3。令和4年度はFCCのz=4、M=0.197kg/mol、a=\(4.1\times10^{-10}\)mから\(1.9056\times10^4\)kg/m3。指定\(N_A=6.0\times10^{23}\)を使用する。OCRの指数33は原典画像の23へ校正し、未校正原文も保存する。

\section{NaClとCsCl：異種の中心はBCC格子点ではない}\label{theory:atomic:ionic}
NaCl型は二つのFCC副格子がずれた構造で異種の最近接配位数6。CsCl型は単純立方格子に異種イオンの基底が付いて配位数8となる。胞の中心と隅が違う元素なので、全点が同等なBCC Bravais格子ではない\cite{openstax_crystals}。Madelung定数はイオンの並びを反映する格子和で、同じ1:1組成というだけで同値とはならない。

NaClに接触硬球のイオン半径\(r_+,r_-\)を置くなら\(a/2=r_++r_-\)。同種の最近接距離\(a/\sqrt2\)は\(\sqrt2(r_++r_-)\)。前半は格子幾何、後半の半径和はモデル仮定である。電子雲に固い表面があり全隣接が厳密に触れるという意味ではない。

\section{炭素の構造と電子状態}\label{theory:atomic:carbon}
diamondは四面体sp3の強い結合網を持ち、通常純結晶は大きなバンドギャップのため良導体ではない\cite{atomic_diamond}。高硬度と高導電性は同じ物性ではない。黒鉛は層内の六角網面と比較的弱い層間結合を持つ。CNTは網面を巻いた単層又は多層の管で、径と巻き方が電子構造へ影響する\cite{atomic_cnt}。

ホウ素ドープしたdiamondやCNTの導電性の例がある\cite{atomic_cnt}。通常純結晶の比較と、ドープ・欠陥・接触条件で変えた材料を分ける。Mohsの最高位10も全人工材料・全方向・全荷重の絶対的順位ではない。

\section{格子振動と磁区応答}\label{theory:atomic:phonon}
格子の正規振動を量子化した励起がphononで、Einsteinモデルは共通の振動数の独立振動子、Debyeモデルは振動数分布を考える。高温の格子比熱はモル当たり3Rへ近づき、低温では励起が減り、3次元Debye近似で格子比熱は\(T^3\)に比例する\cite{atomic_phonon}。金属全体の比熱には電子の寄与もある。

Debye--Waller因子は原子位置の熱変位等による散乱強度の減衰を表す。一方、Barkhausen効果は磁区壁のピン止めと不連続な磁化応答が中心である\cite{atomic_barkhausen}。格子欠陥が磁区壁へ作用し得ることを認めながら、過去問の相対比較では磁区応答を選ぶ。

\section{分子光吸収・光反応と線幅}\label{theory:atomic:optical}
IR振動吸収には双極子の変化、Raman散乱には分極率の変化が関係する\cite{atomic_ir_raman}。CO2の温室効果の主な振動吸収は熱赤外域であり、可視光の振動吸収という説明は適切でない\cite{atomic_ir_co2}。光子のエネルギーとどの自由度を励起するかを先に確認する。

視物質の11cisレチナールの光異性化は視覚の起点になる\cite{atomic_retinal}。光合成ではクロロフィルが光を吸収し電子励起を生じ、NADP+は電子を受け取りNADPHへ還元される役割を持つ\cite{atomic_photosynthesis}。二つを同じ一次光捕集体としない。光素子と超伝導の比較は第\ref{theory:modern:devices}節を参照する。

許容遷移の選択則は線幅が不変という意味ではない。速度分布によるDoppler幅と衝突による幅が存在し、NISTの分子吸収解析でも温度・分子質量・圧力による幅を扱う\cite{atomic_linewidth}。

\section{溶解度と同位体：定義の分母と基準}\label{theory:atomic:solubility}\label{theory:atomic:isotopes}
平成27年度問11の溶解度は100gの飽和溶液中の溶質量である。60℃の水は47.8g、20℃の溶質/水は24/76だから、溶存量15.0947g、析出量37.1053gとなる。100g水当たりの定義へ読み替えず、水の蒸発と結晶への水取込みを無視する条件を明示する。

同位体は同じZでNが違う核種である。単体の構造の違いは同素体、同じAは同重体、周期表の族は別分類。自然試料のδは第\ref{theory:isotope:delta}節の基準比へ戻る。海水硫酸由来NBS-127の\(\delta^{34}\mathrm S\)は約+20.3‰であり、旧CDT脚注もある\cite{atomic_sulfate}。全自然海水を厳密に同値とせず、基準尺度と代表値を区別する。

\section{学習の順序と典型的な誤り}
まず空間軌道の容量と電子状態を区別し、次に電子対配置と結合次数、分子間力へ進む。結晶では格子点と基底を分けて距離を幾何から導き、その後胞質量と体積から密度を求める。導電性・硬度・比熱は電子又は格子のどの自由度で説明するかを確認する。疑義のある全称・不一致を公式正答の番号だけで消さない。

\section{過去問と全候補の検討}
\section{結合次数と電子対配置のモデル}

\subsection{令和元年度 問17}
\label{q:r01-reg-applied-17}

\noindent\textbf{原典：} sources/R01/17\_applied\_science.pdf、PDF物理5頁。

\textbf{公式正答：} 1（single）。

\par\textbf{検討未確定：} Arを分子性結晶と呼ぶ広義分類と単原子性を区別。水の常温液体は圧力等を要する。

\paragraph{問題の要旨} 結合次数・混成モデル・分子間相互作用を比較する。

\Needspace{4\baselineskip}\paragraph{確認した条件・記述（独自要約）}\mbox{}\par

\noindent\textbf{判断対象：} 結合次数・混成モデル・分子間相互作用を比較する。\par

\noindent\textbf{原典確認：} 全5候補・指数・否定・記号を保存原典22頁画像と照合。\par

\paragraph{問われる核心} 結合次数・混成モデル・分子間相互作用を比較する。

\paragraph{必要知識} N2,O2,F2の結合次数は3,2,1で、共有する結合電子対が増えるという1は逆。

\paragraph{考え方・正答を選ぶ理由（査読待ち初稿）} N2,O2,F2の結合次数は3,2,1で、共有する結合電子対が増えるという1は逆。sp3の電子対配置と実際の分子形は違い、NH3やH2Oを正四面体分子と呼ばない。Arの単原子固体を分子性の分類へ含める用語と、通常常圧で水が液体という状態条件を分ける。

\Needspace{4\baselineskip}\paragraph{選択肢の検討（査読待ち初稿）}\mbox{}\par

\noindent\textbf{1：選ぶ候補} N2/O2/F2順に共有電子増で強い。結合次数は3/2/1で逆。反結合性電子の占有も数える。\par

\noindent\textbf{2：選ばない候補} CH4/NH3/H2Oはsp3を反映。4電子領域のモデル。分子形は四面体/三角錐/折れ線と異なる。\par

\noindent\textbf{3：選ばない候補} NaClは静電引力と立方胞。NaCl型の慣用立方胞。\par

\noindent\textbf{4：選ばない候補} 水はHの正電荷と水素結合で常温液体。通常の常圧条件で分子間水素結合が液体を安定化。\par

\noindent\textbf{5：選ばない候補} Ar/ナフタレンは分子性結晶。Arは単原子の分散力固体を広い分子性分類で含める用語。Ar分子や常温固体を意味しない。\par

\paragraph{再利用する考え方} 軌道と電子状態、格子と基底、単位胞原子数と配位数、定義の分母を区別して候補を全て評価する。

\paragraph{共通理論} \hyperref[theory:atomic:bonding]{第\ref*{theory:atomic:bonding}節}

\par 参考：\cite{atomic_nonbond}。

\par 難易度 2（校準初稿）、優先度未評価。

\section{Bravais格子・基底と逆格子}

\subsection{令和7年度 問8}
\label{q:r07-reg-applied-08}

\noindent\textbf{原典：} sources/R07/17\_applied\_science.pdf、PDF物理3頁。

\textbf{公式正答：} 2（single）。

\paragraph{問題の要旨} 格子の並進対称性・Bravais種類と基底を区別する。

\Needspace{4\baselineskip}\paragraph{確認した条件・記述（独自要約）}\mbox{}\par

\noindent\textbf{判断対象：} 格子の並進対称性・Bravais種類と基底を区別する。\par

\noindent\textbf{原典確認：} 全5候補・指数・否定・記号を保存原典22頁画像と照合。\par

\paragraph{問われる核心} 格子の並進対称性・Bravais種類と基底を区別する。

\paragraph{必要知識} 3次元Bravais格子は14種類なので8とする2が不適切。

\paragraph{考え方・正答を選ぶ理由（査読待ち初稿）} 3次元Bravais格子は14種類なので8とする2が不適切。結晶は格子と基底で構成され、NaClの異種イオンやdiamondの二原子基底を格子点の種類として別カウントしない。逆格子で原点まわりのWigner–Seitz胞を第一Brillouin域と呼ぶ。胞原子数と第一/第二配位数は別。

\Needspace{4\baselineskip}\paragraph{選択肢の検討（査読待ち初稿）}\mbox{}\par

\noindent\textbf{1：選ばない候補} 3次元格子は並進対称。格子点は整数並進で互いに等価。\par

\noindent\textbf{2：選ぶ候補} 3次元Bravaisは8種。3次元Bravais格子は14種類。\par

\noindent\textbf{3：選ばない候補} 逆格子Wigner–Seitz胞はBrillouin域。逆格子点まわりのWigner–Seitz胞は第一Brillouin域。\par

\noindent\textbf{4：選ばない候補} diamondの隣接は4と12。四面体第一近接4と同じ副格子の第二近接12。\par

\noindent\textbf{5：選ばない候補} NaClはFCCに属す。FCC Bravais格子にNa/Clの基底が付く。\par

\paragraph{再利用する考え方} 軌道と電子状態、格子と基底、単位胞原子数と配位数、定義の分母を区別して候補を全て評価する。

\paragraph{共通理論} \hyperref[theory:atomic:geometry]{第\ref*{theory:atomic:geometry}節}

\par 参考：\cite{atomic_bravais_ref,openstax_crystals}。

\par 難易度 2（校準初稿）、優先度未評価。

\section{単位胞質量・体積と密度}

\subsection{平成26年度 問14}
\label{q:h26-reg-applied-14}

\noindent\textbf{原典：} sources/H26/17\_applied\_science.pdf、PDF物理6頁。

\textbf{公式正答：} 2（single）。

\paragraph{問題の要旨} BCCの胞原子数と体積から密度を求める。

\Needspace{4\baselineskip}\paragraph{確認した条件・記述（独自要約）}\mbox{}\par

\noindent\textbf{構造：} BCC\par

\noindent\textbf{格子定数：} 0.29nm\par

\noindent\textbf{原子量：} 52\par

\noindent\textbf{指定NA：} 6.0×10\textasciicircum{}23mol\textasciicircum{}-1\par

\noindent\textbf{求める単位：} g/cm3\par

\paragraph{問われる核心} BCCの胞原子数と体積から密度を求める。

\paragraph{必要知識} BCC慣用立方胞では隅8×1/8と体中心1を合わせ2原子。

\paragraph{考え方・正答を選ぶ理由（査読待ち初稿）} BCC慣用立方胞では隅8×1/8と体中心1を合わせ2原子。0.29nm=\allowbreak{}2.9×10\textasciicircum{}-8cmなので、ρ=(2×52g/mol)/(6.0×10\textasciicircum{}23mol\textasciicircum{}-1×(2.9×10\textasciicircum{}-8cm)\textasciicircum{}3)=\allowbreak{}7.11g/cm3、2。長さを立方する前に単位を変える。現在の定義定数へ置き換えず原典指定値を使う。

\Needspace{4\baselineskip}\paragraph{選択肢の検討（査読待ち初稿）}\mbox{}\par

\noindent\textbf{1：選ばない候補} 3.6g/cm3。胞原子数を1とした約半分の値。\par

\noindent\textbf{2：選ぶ候補} 7.1g/cm3。2×52/(6.0×10\textasciicircum{}23×(2.9×10\textasciicircum{}-8)\textasciicircum{}3)=約7.11。\par

\noindent\textbf{3：選ばない候補} 7.9g/cm3。指定量からの計算値より大きい。\par

\noindent\textbf{4：選ばない候補} 10.1g/cm3。指定量からの計算値より大きい。\par

\noindent\textbf{5：選ばない候補} 12.3g/cm3。指定量からの計算値より大きい。\par

\paragraph{再利用する考え方} 軌道と電子状態、格子と基底、単位胞原子数と配位数、定義の分母を区別して候補を全て評価する。

\paragraph{共通理論} \hyperref[theory:atomic:density]{第\ref*{theory:atomic:density}節}

\par 参考：\cite{openstax_crystals}。

\paragraph{関連問題} \hyperref[q:r04-reg-applied-12]{R04-REG-APPLIED-12}

\par 難易度 3（校準初稿）、優先度未評価。

\subsection{令和4年度 問12}
\label{q:r04-reg-applied-12}

\noindent\textbf{原典：} sources/R04/17\_applied\_science.pdf、PDF物理4頁。

\textbf{公式正答：} 4（single）。

\paragraph{問題の要旨} FCCの胞質量をSI体積で割り密度を求める。

\Needspace{4\baselineskip}\paragraph{確認した条件・記述（独自要約）}\mbox{}\par

\noindent\textbf{構造：} FCC\par

\noindent\textbf{格子定数：} 0.41nm\par

\noindent\textbf{原子量：} 197\par

\noindent\textbf{指定NA：} 6.0×10\textasciicircum{}23mol\textasciicircum{}-1\par

\noindent\textbf{求める単位：} kg/m3\par

\paragraph{問われる核心} FCCの胞質量をSI体積で割り密度を求める。

\paragraph{必要知識} FCC慣用胞の原子数は8×1/8+6×1/2=\allowbreak{}4。

\paragraph{考え方・正答を選ぶ理由（査読待ち初稿）} FCC慣用胞の原子数は8×1/8+6×1/2=\allowbreak{}4。胞質量は4×0.197/(6.0×10\textasciicircum{}23)kg、体積は(4.1×10\textasciicircum{}-10)\textasciicircum{}3m3。ρ=1.9056×10\textasciicircum{}4kg/m3なので4。197g/molをkg/molへ直す。H26と同じ方法だが構造・単位・数値が異なる。

\Needspace{4\baselineskip}\paragraph{選択肢の検討（査読待ち初稿）}\mbox{}\par

\noindent\textbf{1：選ばない候補} 7.9×10\textasciicircum{}3kg/m3。計算密度に達しない。\par

\noindent\textbf{2：選ばない候補} 1.2×10\textasciicircum{}4kg/m3。計算密度より小さい。\par

\noindent\textbf{3：選ばない候補} 1.5×10\textasciicircum{}4kg/m3。計算密度より小さい。\par

\noindent\textbf{4：選ぶ候補} 1.9×10\textasciicircum{}4kg/m3。胞原子数4とa=\allowbreak{}4.1×10\textasciicircum{}-10mから約1.91×10\textasciicircum{}4。\par

\noindent\textbf{5：選ばない候補} 2.3×10\textasciicircum{}4kg/m3。計算密度より大きい。\par

\paragraph{再利用する考え方} 軌道と電子状態、格子と基底、単位胞原子数と配位数、定義の分母を区別して候補を全て評価する。

\paragraph{共通理論} \hyperref[theory:atomic:density]{第\ref*{theory:atomic:density}節}

\par 参考：\cite{openstax_crystals}。

\paragraph{関連問題} \hyperref[q:h26-reg-applied-14]{H26-REG-APPLIED-14}

\par 難易度 3（校準初稿）、優先度未評価。

\section{格子・配位数と最近接距離}

\subsection{平成24年度 問5}
\label{q:h24-reg-applied-05}

\noindent\textbf{原典：} sources/H24/17\_applied\_science.pdf、PDF物理2頁。

\textbf{公式正答：} 2（single）。

\paragraph{問題の要旨} FCCの最近接距離から格子定数を求める。

\Needspace{4\baselineskip}\paragraph{確認した条件・記述（独自要約）}\mbox{}\par

\noindent\textbf{判断対象：} FCCの最近接距離から格子定数を求める。\par

\noindent\textbf{原典確認：} 全5候補・指数・否定・記号を保存原典22頁画像と照合。\par

\paragraph{問われる核心} FCCの最近接距離から格子定数を求める。

\paragraph{必要知識} FCCでは面対角線上の隅と面中心が最近接で、d=√2a/2。

\paragraph{考え方・正答を選ぶ理由（査読待ち初稿）} FCCでは面対角線上の隅と面中心が最近接で、d=√2a/2。従ってa=√2dで2。面中心と隅のベクトル(a/2,a/2,0)の長さからも導ける。単位胞の原子数4と配位数12を、この距離式の係数として代入しない。

\Needspace{4\baselineskip}\paragraph{選択肢の検討（査読待ち初稿）}\mbox{}\par

\noindent\textbf{1：選ばない候補} d。FCCの最近接距離は胞辺より短い。\par

\noindent\textbf{2：選ぶ候補} √2d。面対角線を半分にした距離d=\allowbreak{}a/√2からa=√2d。\par

\noindent\textbf{3：選ばない候補} d/√2。距離と格子定数を逆にした値。\par

\noindent\textbf{4：選ばない候補} 2d/√3。BCCなら現れる関係でFCCではない。\par

\noindent\textbf{5：選ばない候補} √3d。FCCの面対角線の幾何に合わない。\par

\paragraph{再利用する考え方} 軌道と電子状態、格子と基底、単位胞原子数と配位数、定義の分母を区別して候補を全て評価する。

\paragraph{共通理論} \hyperref[theory:atomic:geometry]{第\ref*{theory:atomic:geometry}節}

\par 参考：\cite{openstax_crystals}。

\paragraph{関連問題} \hyperref[q:h28-reg-applied-07]{H28-REG-APPLIED-07}, \hyperref[q:r01-ret-applied-07]{R01-RET-APPLIED-07}

\par 難易度 3（校準初稿）、優先度未評価。

\subsection{平成28年度 問7}
\label{q:h28-reg-applied-07}

\noindent\textbf{原典：} sources/H28/17\_applied\_science.pdf、PDF物理3頁。

\textbf{公式正答：} 2（single）。

\paragraph{問題の要旨} FCCの最近接距離から格子定数を求める。

\Needspace{4\baselineskip}\paragraph{確認した条件・記述（独自要約）}\mbox{}\par

\noindent\textbf{判断対象：} FCCの最近接距離から格子定数を求める。\par

\noindent\textbf{原典確認：} 全5候補・指数・否定・記号を保存原典22頁画像と照合。\par

\paragraph{問われる核心} FCCの最近接距離から格子定数を求める。

\paragraph{必要知識} FCCでは面対角線上の隅と面中心が最近接で、d=√2a/2。

\paragraph{考え方・正答を選ぶ理由（査読待ち初稿）} FCCでは面対角線上の隅と面中心が最近接で、d=√2a/2。従ってa=√2dで2。面中心と隅のベクトル(a/2,a/2,0)の長さからも導ける。単位胞の原子数4と配位数12を、この距離式の係数として代入しない。

\Needspace{4\baselineskip}\paragraph{選択肢の検討（査読待ち初稿）}\mbox{}\par

\noindent\textbf{1：選ばない候補} d。FCCの最近接距離は胞辺より短い。\par

\noindent\textbf{2：選ぶ候補} √2d。面対角線を半分にした距離d=\allowbreak{}a/√2からa=√2d。\par

\noindent\textbf{3：選ばない候補} d/√2。距離と格子定数を逆にした値。\par

\noindent\textbf{4：選ばない候補} 2d/√3。BCCなら現れる関係でFCCではない。\par

\noindent\textbf{5：選ばない候補} √3d。FCCの面対角線の幾何に合わない。\par

\paragraph{再利用する考え方} 軌道と電子状態、格子と基底、単位胞原子数と配位数、定義の分母を区別して候補を全て評価する。

\paragraph{共通理論} \hyperref[theory:atomic:geometry]{第\ref*{theory:atomic:geometry}節}

\par 参考：\cite{openstax_crystals}。

\paragraph{関連問題} \hyperref[q:h24-reg-applied-05]{H24-REG-APPLIED-05}

\par 難易度 3（校準初稿）、優先度未評価。

\subsection{令和元年度再試験 問7}
\label{q:r01-ret-applied-07}

\noindent\textbf{原典：} sources/R01\_retest/17\_applied\_science.pdf、PDF物理3頁。

\textbf{公式正答：} 1（single）。

\paragraph{問題の要旨} BCCの格子定数から最近接距離を求める。

\Needspace{4\baselineskip}\paragraph{確認した条件・記述（独自要約）}\mbox{}\par

\noindent\textbf{判断対象：} BCCの格子定数から最近接距離を求める。\par

\noindent\textbf{原典確認：} 全5候補・指数・否定・記号を保存原典22頁画像と照合。\par

\paragraph{問われる核心} BCCの格子定数から最近接距離を求める。

\paragraph{必要知識} BCCの体中心(d/2,d/2,d/2)と隅(0,0,0)の距離は√3d/2で1。

\paragraph{考え方・正答を選ぶ理由（査読待ち初稿）} BCCの体中心(d/2,d/2,d/2)と隅(0,0,0)の距離は√3d/2で1。FCCの面対角線と混同しない。ここで原典のdは格子定数であり、H24等のd=最近接距離と記号の役割が違う。

\Needspace{4\baselineskip}\paragraph{選択肢の検討（査読待ち初稿）}\mbox{}\par

\noindent\textbf{1：選ぶ候補} √3d/2。体対角線の隅から体中心までの半分。\par

\noindent\textbf{2：選ばない候補} d/√2。FCCの関係。\par

\noindent\textbf{3：選ばない候補} √3d。体対角線全長で最近接の2倍。\par

\noindent\textbf{4：選ばない候補} √2d。面対角線全長。\par

\noindent\textbf{5：選ばない候補} d。隅どうしの胞辺で体中心はさらに近い。\par

\paragraph{再利用する考え方} 軌道と電子状態、格子と基底、単位胞原子数と配位数、定義の分母を区別して候補を全て評価する。

\paragraph{共通理論} \hyperref[theory:atomic:geometry]{第\ref*{theory:atomic:geometry}節}

\par 参考：\cite{openstax_crystals}。

\paragraph{関連問題} \hyperref[q:h24-reg-applied-05]{H24-REG-APPLIED-05}

\par 難易度 3（校準初稿）、優先度未評価。

\subsection{令和3年度 問8}
\label{q:r03-reg-applied-08}

\noindent\textbf{原典：} sources/R03/17\_applied\_science.pdf、PDF物理3頁。

\textbf{公式正答：} 2（single）。

\par\textbf{検討未確定：} 元素例は通常相を想定。HCPを全元素で理想最密軸比とする読みは避け、出題時全相条件は未指定。

\paragraph{問題の要旨} 結晶の配位数と塑性すべり条件を区別する。

\Needspace{4\baselineskip}\paragraph{確認した条件・記述（独自要約）}\mbox{}\par

\noindent\textbf{判断対象：} 結晶の配位数と塑性すべり条件を区別する。\par

\noindent\textbf{原典確認：} 全5候補・指数・否定・記号を保存原典22頁画像と照合。\par

\paragraph{問われる核心} 結晶の配位数と塑性すべり条件を区別する。

\paragraph{必要知識} FCCの配位数12を8とした2が不適切。

\paragraph{考え方・正答を選ぶ理由（査読待ち初稿）} FCCの配位数12を8とした2が不適切。BCCは8で充填68\%。胞原子数BCC2/FCC4と最近接配位数8/12を混同しない。すべりは結晶面と方向へ分解した剪断応力で判断する。HCP元素例を全温度圧力や理想軸比の保証へ広げない。

\Needspace{4\baselineskip}\paragraph{選択肢の検討（査読待ち初稿）}\mbox{}\par

\noindent\textbf{1：選ばない候補} αFe/Cr/MoはBCCで充填約68\%。BCC硬球充填率π√3/8≈0.680。\par

\noindent\textbf{2：選ぶ候補} Al/Ni/CuはFCCで配位数8。FCC最近接は12で8はBCC。\par

\noindent\textbf{3：選ばない候補} Mg/Zn/Cd/αTiはHCP。通常相の六方配列。Zn等の軸比を理想最密値へ固定しない。\par

\noindent\textbf{4：選ばない候補} 分解剪断応力が臨界に達しすべる。τ=σcosφcosλのSchmid則に対応。\par

\noindent\textbf{5：選ばない候補} 温度等で構造変化は同素変態。同元素の結晶構造が相条件で変わる。\par

\paragraph{再利用する考え方} 軌道と電子状態、格子と基底、単位胞原子数と配位数、定義の分母を区別して候補を全て評価する。

\paragraph{共通理論} \hyperref[theory:atomic:geometry]{第\ref*{theory:atomic:geometry}節}, \hyperref[theory:material:crystals]{第\ref*{theory:material:crystals}節}

\par 参考：\cite{openstax_crystals}。

\paragraph{関連問題} \hyperref[q:r04-reg-applied-08]{R04-REG-APPLIED-08}

\par 難易度 2（校準初稿）、優先度未評価。

\subsection{令和4年度 問8}
\label{q:r04-reg-applied-08}

\noindent\textbf{原典：} sources/R04/17\_applied\_science.pdf、PDF物理3頁。

\textbf{公式正答：} 5（single）。

\par\textbf{検討未確定：} 元素例の全相境界・出題時測定条件は未指定。912℃は通常圧純Feの代表構造変態値。

\paragraph{問題の要旨} 金属の通常相とアモルファスの長距離秩序を比較する。

\Needspace{4\baselineskip}\paragraph{確認した条件・記述（独自要約）}\mbox{}\par

\noindent\textbf{判断対象：} 金属の通常相とアモルファスの長距離秩序を比較する。\par

\noindent\textbf{原典確認：} 全5候補・指数・否定・記号を保存原典22頁画像と照合。\par

\paragraph{問われる核心} 金属の通常相とアモルファスの長距離秩序を比較する。

\paragraph{必要知識} アモルファスは長距離の周期的結晶格子を持たないので5が不適切。

\paragraph{考え方・正答を選ぶ理由（査読待ち初稿）} アモルファスは長距離の周期的結晶格子を持たないので5が不適切。完全に全原子が無相関という意味ではない。Feの912℃の構造変態と770℃付近の磁気転移は別である。元素例は通常相を想定し全温度圧力に固定しない。

\Needspace{4\baselineskip}\paragraph{選択肢の検討（査読待ち初稿）}\mbox{}\par

\noindent\textbf{1：選ばない候補} Al/Cu/NiはFCC。通常の常温常圧相。\par

\noindent\textbf{2：選ばない候補} W/Mo/NbはBCC。通常の常温常圧相。\par

\noindent\textbf{3：選ばない候補} Mg/Zn/YはHCP。通常の六方相。\par

\noindent\textbf{4：選ばない候補} Feは912℃で格子型変態。純FeのαBCC→γFCCの相変態。\par

\noindent\textbf{5：選ぶ候補} アモルファスは規則的結晶格子。長距離の周期的並進秩序を持たない。短距離秩序はあり得る。\par

\paragraph{再利用する考え方} 軌道と電子状態、格子と基底、単位胞原子数と配位数、定義の分母を区別して候補を全て評価する。

\paragraph{共通理論} \hyperref[theory:atomic:geometry]{第\ref*{theory:atomic:geometry}節}

\par 参考：\cite{princeton_iron,princeton_tungsten,openstax_crystals}。

\paragraph{関連問題} \hyperref[q:r03-reg-applied-08]{R03-REG-APPLIED-08}

\par 難易度 2（校準初稿）、優先度未評価。

\section{電子配置・容量と原子データ}

\subsection{令和元年度再試験 問14}
\label{q:r01-ret-applied-14}

\noindent\textbf{原典：} sources/R01\_retest/17\_applied\_science.pdf、PDF物理5頁。

\textbf{公式正答：} 2（single）。

\paragraph{問題の要旨} 酸素基底状態の電子数とs軌道容量を確認する。

\Needspace{4\baselineskip}\paragraph{確認した条件・記述（独自要約）}\mbox{}\par

\noindent\textbf{判断対象：} 酸素基底状態の電子数とs軌道容量を確認する。\par

\noindent\textbf{原典確認：} 全5候補・指数・否定・記号を保存原典22頁画像と照合。\par

\paragraph{問われる核心} 酸素基底状態の電子数とs軌道容量を確認する。

\paragraph{必要知識} 酸素はZ=\allowbreak{}8で中性なら8電子。

\paragraph{考え方・正答を選ぶ理由（査読待ち初稿）} 酸素はZ=\allowbreak{}8で中性なら8電子。s副殻は1空間軌道に逆スピン2電子が最大で、基底配置1s2 2s2 2p4、2。候補は全て合計8でも、総数だけで判定できず容量とエネルギー順を確認する。

\Needspace{4\baselineskip}\paragraph{選択肢の検討（査読待ち初稿）}\mbox{}\par

\noindent\textbf{1：選ばない候補} 1s1 2s3 2p4。2sに3電子は不可能。\par

\noindent\textbf{2：選ぶ候補} 1s2 2s2 2p4。低い1s/2sを各2、2pを4で計8。\par

\noindent\textbf{3：選ばない候補} 1s1 2s4 2p3。2sに4電子は不可能。\par

\noindent\textbf{4：選ばない候補} 1s2 2s3 2p3。2sに3電子は不可能。\par

\noindent\textbf{5：選ばない候補} 1s2 2s4 2p2。2sに4電子は不可能。\par

\paragraph{再利用する考え方} 軌道と電子状態、格子と基底、単位胞原子数と配位数、定義の分母を区別して候補を全て評価する。

\paragraph{共通理論} \hyperref[theory:atomic:orbital]{第\ref*{theory:atomic:orbital}節}

\par 参考：\cite{atomic_quantum,atomic_nist2019}。

\paragraph{関連問題} \hyperref[q:r02-reg-applied-10]{R02-REG-APPLIED-10}

\par 難易度 2（校準初稿）、優先度未評価。

\subsection{令和2年度 問10}
\label{q:r02-reg-applied-10}

\noindent\textbf{原典：} sources/R02/17\_applied\_science.pdf、PDF物理4頁。

\textbf{公式正答：} ALL（all\_awarded）。

\par\textbf{独立解答と公式正答に相違あり：} 独立解答 3|5。公式ALLを維持。初回3|5は公式前の独立判断として保持。原典の銅・原子番号・電子数を別元素へ変更しない。公式措置理由の全文は未取得。

\paragraph{問題の要旨} 原典の複数の誤りと全員正解の公式採点を保持する。

\Needspace{4\baselineskip}\paragraph{確認した条件・記述（独自要約）}\mbox{}\par

\noindent\textbf{判断対象：} 原典の複数の誤りと全員正解の公式採点を保持する。\par

\noindent\textbf{原典確認：} 全5候補・指数・否定・記号を保存原典22頁画像と照合。\par

\paragraph{問われる核心} 原典の複数の誤りと全員正解の公式採点を保持する。

\paragraph{必要知識} 独立初回は3|5、公式はALL（全員正解）。

\paragraph{考え方・正答を選ぶ理由（査読待ち初稿）} 独立初回は3|5、公式はALL（全員正解）。原典には実際に銅とあり、Gaへ黙って修正しない。Cuは29、Znは30で4s電子も1と2。さらにKrの第一IEはArより小さいため5も不適切。単一候補へ絞らず公式採点と二つの科学的誤りを併記する。公式の全員正解措置の理由文は未取得なので、理由を公式発表として捏造しない。

\Needspace{4\baselineskip}\paragraph{選択肢の検討（査読待ち初稿）}\mbox{}\par

\noindent\textbf{1：公式採点は全員正解} Neは1s2 2s2 2p6。中性Neの10電子配置。\par

\noindent\textbf{2：公式採点は全員正解} Krは[Ar]4s2 3d10 4p6。中性Krの36電子配置。記載順と容量を区別。\par

\noindent\textbf{3：公式採点は全員正解} ZnよりZが1大きいCuは4sが1多い。原典のCuを保持。Cu29はZn30より1小さく4sは1対2で少ない。\par

\noindent\textbf{4：公式採点は全員正解} MgのIEはNaより大きい。NIST値Mg7.6462eV>Na5.1391eV。\par

\noindent\textbf{5：公式採点は全員正解} KrのIEはArより大きい。NIST値Kr13.9996eV<Ar15.7596eVで逆。\par

\paragraph{再利用する考え方} 軌道と電子状態、格子と基底、単位胞原子数と配位数、定義の分母を区別して候補を全て評価する。

\paragraph{共通理論} \hyperref[theory:atomic:orbital]{第\ref*{theory:atomic:orbital}節}

\par 参考：\cite{atomic_nist2019}。

\paragraph{関連問題} \hyperref[q:r01-ret-applied-14]{R01-RET-APPLIED-14}

\par 難易度 2（校準初稿）、優先度未評価。

\subsection{令和6年度 問8}
\label{q:r06-reg-applied-08}

\noindent\textbf{原典：} sources/R06/17\_applied\_science.pdf、PDF物理3頁。

\textbf{公式正答：} 4（single）。

\par\textbf{検討未確定：} 肢3の軌道という原典語は副殻群の総容量として条件を補う。肢5も孤立原子とバルク金属の常磁性機構が同じとはしない。

\paragraph{問題の要旨} 副殻容量・半導体の担い手とGeの四面体結合を区別する。

\Needspace{4\baselineskip}\paragraph{確認した条件・記述（独自要約）}\mbox{}\par

\noindent\textbf{原典確認：} S7\_051の11原典画像で13問の全五肢・指数・単位・否定を照合。\par

\noindent\textbf{対象：} 副殻容量・半導体の担い手とGeの四面体結合を区別する。\par

\paragraph{問われる核心} 副殻容量・半導体の担い手とGeの四面体結合を区別する。

\paragraph{必要知識} 4は四面体構造をsp\textsuperscript{2}へ置き換えたことが誤り。原子の4s\textsuperscript{2}4p\textsuperscript{2}配置と結晶の共有結合モデルを分ける。副殻の総容量と一個の空間軌道の容量、孤立原子と金属の磁気応答も分離する。

\paragraph{考え方・正答を選ぶ理由（査読待ち初稿）} 4は四面体構造をsp\textsuperscript{2}へ置き換えたことが誤り。原子の4s\textsuperscript{2}4p\textsuperscript{2}配置と結晶の共有結合モデルを分ける。副殻の総容量と一個の空間軌道の容量、孤立原子と金属の磁気応答も分離する。

\Needspace{4\baselineskip}\paragraph{選択肢の検討（査読待ち初稿）}\mbox{}\par

\noindent\textbf{1：代表的条件で適切} 電池で自由電子が正極側へ動く。電子の負電荷による平均ドリフトは通常の電流方向と反対。各電子の熱運動を含む瞬間的な全軌道が正極へ直線的に進むわけではない。\par

\noindent\textbf{2：代表的条件で適切} SiへBを添加するとp型。Bはアクセプターとして働き正孔が主要な担い手となる。結晶全体が単に正に帯電する意味ではない。\par

\noindent\textbf{3：条件・語義を留保} s・p・d・fの収容数2・6・10・14。同じ副殻に属する空間軌道の数1・3・5・7へ各2電子を入れた総容量。一個の空間軌道はどの種類も2電子なので原典の軌道という語を副殻群として説明する。\par

\noindent\textbf{4：不適切} Geの4s\textsuperscript{2}4p\textsuperscript{2}がsp\textsuperscript{2}でダイヤモンド構造。原典の電子配置は4s\textsuperscript{2}4p\textsuperscript{2}でNISTにも一致。通常のGeは四面体の4近接結合を持つダイヤモンド型で、教科書的混成モデルはsp\textsuperscript{3}。sp\textsuperscript{2}の三角平面と取り違え。\par

\noindent\textbf{5：条件・語義を留保} Li・Naの不対スピンと常磁性。孤立原子のs電子の軌道角運動量ゼロとスピンを区別する。ただしバルク金属のPauli応答はFermi統計を含み、孤立原子のCurie則をそのまま使わない。\par

\paragraph{再利用する考え方} 場・応答・担い手・材料・温度を区別し、数値の指数と単位、代表モデルの適用条件を確認する。

\paragraph{共通理論} \hyperref[theory:atomic:orbital]{第\ref*{theory:atomic:orbital}節}, \hyperref[theory:atomic:geometry]{第\ref*{theory:atomic:geometry}節}, \hyperref[theory:em:pauli]{第\ref*{theory:em:pauli}節}

\par 参考：\cite{atomic_quantum,em51_ge,em51_ge_structure,mit_semiconductor,em51_pauli}。

\par 難易度 3（校準初稿）、優先度未評価。

\section{分子間力・分散とLJのエネルギー}

\subsection{平成26年度 問11}
\label{q:h26-reg-applied-11}

\noindent\textbf{原典：} sources/H26/17\_applied\_science.pdf、PDF物理5頁。

\textbf{公式正答：} 2（single）。

\paragraph{問題の要旨} 分子極性と水素結合の供与・受容を判断する。

\Needspace{4\baselineskip}\paragraph{確認した条件・記述（独自要約）}\mbox{}\par

\noindent\textbf{判断対象：} 分子極性と水素結合の供与・受容を判断する。\par

\noindent\textbf{原典確認：} 全5候補・指数・否定・記号を保存原典22頁画像と照合。\par

\paragraph{問われる核心} 分子極性と水素結合の供与・受容を判断する。

\paragraph{必要知識} CH4は対称な正四面体で永久双極子がなく、極性が強いため水素結合を生じるという2が不適切。

\paragraph{考え方・正答を選ぶ理由（査読待ち初稿）} CH4は対称な正四面体で永久双極子がなく、極性が強いため水素結合を生じるという2が不適切。水・HFの典型的水素結合や異種分子間の供与体/受容体を区別する。水素を含めば全て同等の水素結合が生じるという推論はできない。

\Needspace{4\baselineskip}\paragraph{選択肢の検討（査読待ち初稿）}\mbox{}\par

\noindent\textbf{1：選ばない候補} 極性の強い水素化合物に多い。典型的なN/O/Fに結合したH等を入口とする。\par

\noindent\textbf{2：選ぶ候補} CH4は強極性で水素結合。正四面体CH4の結合双極子は相殺し、強極性という前提が誤り。\par

\noindent\textbf{3：選ばない候補} 水の沸点はH2Sより高い。水の分子間水素結合が高い沸点を説明する。\par

\noindent\textbf{4：選ばない候補} HFは高い電気陰性度に関係。HFの分極した結合と電子対による相互作用。\par

\noindent\textbf{5：選ばない候補} 異種分子間にも水素結合。供与体と受容体が違う分子でも成立する。\par

\paragraph{再利用する考え方} 軌道と電子状態、格子と基底、単位胞原子数と配位数、定義の分母を区別して候補を全て評価する。

\paragraph{共通理論} \hyperref[theory:atomic:interactions]{第\ref*{theory:atomic:interactions}節}

\par 参考：\cite{atomic_nonbond}。

\par 難易度 2（校準初稿）、優先度未評価。

\subsection{平成29年度 問11}
\label{q:h29-reg-applied-11}

\noindent\textbf{原典：} sources/H29/17\_applied\_science.pdf、PDF物理4頁。

\textbf{公式正答：} 5（single）。

\par\textbf{検討未確定：} 候補2は凝集エネルギーを配位数のみへ帰する排他的主張で疑義。候補5はポテンシャルと力の語も混在。公式5と科学的一意性を分けて査読。

\paragraph{問題の要旨} 結合多重度・分散力とLennard–Jonesの符号を確認する。

\Needspace{4\baselineskip}\paragraph{確認した条件・記述（独自要約）}\mbox{}\par

\noindent\textbf{判断対象：} 結合多重度・分散力とLennard–Jonesの符号を確認する。\par

\noindent\textbf{原典確認：} 全5候補・指数・否定・記号を保存原典22頁画像と照合。\par

\paragraph{問われる核心} 結合多重度・分散力とLennard–Jonesの符号を確認する。

\paragraph{必要知識} 通常U=\allowbreak{}4ε[(σ/r)\textasciicircum{}12−(σ/r)\textasciicircum{}6]で短距離反発が正のr\textasciicircum{}-12、引力が負のr\textasciicircum{}-6。

\paragraph{考え方・正答を選ぶ理由（査読待ち初稿）} 通常U=\allowbreak{}4ε[(σ/r)\textasciicircum{}12−(σ/r)\textasciicircum{}6]で短距離反発が正のr\textasciicircum{}-12、引力が負のr\textasciicircum{}-6。5は両者が逆なので独立初回も5。ただし力F=−dU/drでは指数13と7となり、原典の力という語も区別する。候補2のW凝集を配位数だけで説明する断定には別の疑義を残し、全他肢が科学的に無条件正しいとはしない。

\Needspace{4\baselineskip}\paragraph{選択肢の検討（査読待ち初稿）}\mbox{}\par

\noindent\textbf{1：選ばない候補} 炭素結合の多重度増で強く短くなる。代表的C-C/C=\allowbreak{}C/C≡Cの傾向。\par

\noindent\textbf{2：選ばない候補} Wの大凝集エネルギーは結合エネルギーでなく配位数だけ。配位数だけとの排他的説明に疑義。BCCの配位数8だけで凝集を保証できない。\par

\noindent\textbf{3：選ばない候補} CO2は電子対反発を避け直線。VSEPRでCまわり2電子領域は180度。\par

\noindent\textbf{4：選ばない候補} 分散は瞬間双極子による。瞬間双極子が他分子へ双極子を誘起。\par

\noindent\textbf{5：選ぶ候補} LJはr\textasciicircum{}-6反発とr\textasciicircum{}-12引力。通常の12–6ポテンシャルの反発と引力が逆。力の指数もポテンシャルとは異なる。\par

\paragraph{再利用する考え方} 軌道と電子状態、格子と基底、単位胞原子数と配位数、定義の分母を区別して候補を全て評価する。

\paragraph{共通理論} \hyperref[theory:atomic:interactions]{第\ref*{theory:atomic:interactions}節}

\par 参考：\cite{atomic_lj,atomic_nonbond}。

\par 難易度 2（校準初稿）、優先度未評価。

\section{NaCl/CsClと格子・基底}

\subsection{平成27年度 問9}
\label{q:h27-reg-applied-09}

\noindent\textbf{原典：} sources/H27/17\_applied\_science.pdf、PDF物理4頁。

\textbf{公式正答：} 3（single）。

\par\textbf{検討未確定：} 全隣接イオンが接するという記述はイオン半径の硬球近似。半径の定義・実結晶条件は原典で未指定。

\paragraph{問題の要旨} NaClの配位数と硬球近似の距離を区別する。

\Needspace{4\baselineskip}\paragraph{確認した条件・記述（独自要約）}\mbox{}\par

\noindent\textbf{判断対象：} NaClの配位数と硬球近似の距離を区別する。\par

\noindent\textbf{原典確認：} 全5候補・指数・否定・記号を保存原典22頁画像と照合。\par

\paragraph{問われる核心} NaClの配位数と硬球近似の距離を区別する。

\paragraph{必要知識} NaCl型は最近接異種イオンの配位数6で、4とする3が不適切。

\paragraph{考え方・正答を選ぶ理由（査読待ち初稿）} NaCl型は最近接異種イオンの配位数6で、4とする3が不適切。陽陰それぞれのFCC副格子がずれて重なる。硬球接触を置くとa/2=\allowbreak{}rNa+rCl、同種最近接a/√2=√2(rNa+rCl)。格子の幾何と、半径の和を実距離とするモデル仮定を分ける。

\Needspace{4\baselineskip}\paragraph{選択肢の検討（査読待ち初稿）}\mbox{}\par

\noindent\textbf{1：選ばない候補} 単位格子は立方体。慣用胞は立方。\par

\noindent\textbf{2：選ばない候補} 静電引力による結晶。陽陰イオンの静電相互作用。\par

\noindent\textbf{3：選ぶ候補} 最近接Clは4個。Naの周囲のClは上下左右前後の6個。\par

\noindent\textbf{4：選ばない候補} 隣接Na/Clは全て接触。イオン半径による硬球近似の記述で、電子雲の厳密境界ではない。\par

\noindent\textbf{5：選ばない候補} 最近接Na-Naは√2(rNa+rCl)。Na-Cl距離a/2=\allowbreak{}rNa+rClとNa-Na=\allowbreak{}a/√2から導く近似式。\par

\paragraph{再利用する考え方} 軌道と電子状態、格子と基底、単位胞原子数と配位数、定義の分母を区別して候補を全て評価する。

\paragraph{共通理論} \hyperref[theory:atomic:ionic]{第\ref*{theory:atomic:ionic}節}

\par 参考：\cite{openstax_crystals}。

\paragraph{関連問題} \hyperref[q:r03-reg-applied-14]{R03-REG-APPLIED-14}

\par 難易度 2（校準初稿）、優先度未評価。

\subsection{令和3年度 問14}
\label{q:r03-reg-applied-14}

\noindent\textbf{原典：} sources/R03/17\_applied\_science.pdf、PDF物理5頁。

\textbf{公式正答：} 3（single）。

\par\textbf{検討未確定：} 他候補は混成・典型結合強度・分子固体というモデル/一般傾向。全条件への無限定一般化は保持しない。

\paragraph{問題の要旨} NaClとCsClの格子・配位を比較する。

\Needspace{4\baselineskip}\paragraph{確認した条件・記述（独自要約）}\mbox{}\par

\noindent\textbf{判断対象：} NaClとCsClの格子・配位を比較する。\par

\noindent\textbf{原典確認：} 全5候補・指数・否定・記号を保存原典22頁画像と照合。\par

\paragraph{問われる核心} NaClとCsClの格子・配位を比較する。

\paragraph{必要知識} NaClは陽陰のFCC副格子を重ね配位6、CsClは単純立方の隅と異種の中心で配位8。

\paragraph{考え方・正答を選ぶ理由（査読待ち初稿）} NaClは陽陰のFCC副格子を重ね配位6、CsClは単純立方の隅と異種の中心で配位8。異なる幾何なので同構造/同Madelung定数という3は不適切。CsClをBCC Bravais格子と呼ばない。異種イオンの中心は隅と並進で同等ではない。Madelung定数の数値暗記を使わず構造の相違から判断する。

\Needspace{4\baselineskip}\paragraph{選択肢の検討（査読待ち初稿）}\mbox{}\par

\noindent\textbf{1：選ばない候補} O2/F2はσ分裂がπより大きい。指定2p軌道の重なりの向きの比較。\par

\noindent\textbf{2：選ばない候補} H2Oはsp3電子対で折れ線。結合対2と孤立対2のモデル。\par

\noindent\textbf{3：選ぶ候補} NaClとCsClは同構造・同Madelung定数。NaClはFCC基底で配位6、CsClは単純立方基底で配位8で異なる。\par

\noindent\textbf{4：選ばない候補} 水素結合は多くの場合共有結合より弱い。典型的分子間水素結合の比較で、全化学種の普遍的上限とはしない。\par

\noindent\textbf{5：選ばない候補} 分子性固体は一般に軟らかく低融点。相互作用の代表傾向で全分子固体の数値保証ではない。\par

\paragraph{再利用する考え方} 軌道と電子状態、格子と基底、単位胞原子数と配位数、定義の分母を区別して候補を全て評価する。

\paragraph{共通理論} \hyperref[theory:atomic:ionic]{第\ref*{theory:atomic:ionic}節}

\par 参考：\cite{openstax_crystals}。

\paragraph{関連問題} \hyperref[q:h27-reg-applied-09]{H27-REG-APPLIED-09}

\par 難易度 2（校準初稿）、優先度未評価。

\section{格子比熱・散乱と磁区応答}

\subsection{平成30年度 問9}
\label{q:h30-reg-applied-09}

\noindent\textbf{原典：} sources/H30/17\_applied\_science.pdf、PDF物理3頁。

\textbf{公式正答：} 5（single）。

\paragraph{問題の要旨} 格子振動の比熱・散乱と磁区の応答を区別する。

\Needspace{4\baselineskip}\paragraph{確認した条件・記述（独自要約）}\mbox{}\par

\noindent\textbf{判断対象：} 格子振動の比熱・散乱と磁区の応答を区別する。\par

\noindent\textbf{原典確認：} 全5候補・指数・否定・記号を保存原典22頁画像と照合。\par

\paragraph{問われる核心} 格子振動の比熱・散乱と磁区の応答を区別する。

\paragraph{必要知識} 1から4は格子の振動自由度又は熱変位を直接扱う。

\paragraph{考え方・正答を選ぶ理由（査読待ち初稿）} 1から4は格子の振動自由度又は熱変位を直接扱う。Barkhausen効果は磁区壁がピン止めを越えて不連続に動く磁化応答で5。磁性と格子が一切相互作用しないという断言へ広げない。高温3Rと低温で減る格子比熱は振動子の量子化の同じ枠組みで説明できる。

\Needspace{4\baselineskip}\paragraph{選択肢の検討（査読待ち初稿）}\mbox{}\par

\noindent\textbf{1：選ばない候補} Dulong–Petit則。高温の格子比熱3Rに対応。\par

\noindent\textbf{2：選ばない候補} Debye–Waller因子。原子の変位による散乱強度の減衰。\par

\noindent\textbf{3：選ばない候補} Debye温度。格子振動スペクトルの特徴温度。\par

\noindent\textbf{4：選ばない候補} Einstein比熱。量子調和振動子による格子比熱モデル。\par

\noindent\textbf{5：選ぶ候補} Barkhausen効果。磁区壁の不連続運動による磁化変化で、比較中は格子振動との関係が薄い。\par

\paragraph{再利用する考え方} 軌道と電子状態、格子と基底、単位胞原子数と配位数、定義の分母を区別して候補を全て評価する。

\paragraph{共通理論} \hyperref[theory:atomic:phonon]{第\ref*{theory:atomic:phonon}節}

\par 参考：\cite{atomic_phonon,atomic_barkhausen}。

\par 難易度 2（校準初稿）、優先度未評価。

\section{空間軌道・スピンと結合対称性}

\subsection{令和元年度 問14}
\label{q:r01-reg-applied-14}

\noindent\textbf{原典：} sources/R01/17\_applied\_science.pdf、PDF物理4頁。

\textbf{公式正答：} 1（single）。

\par\textbf{独立解答と公式正答に相違あり：} 独立解答 4。初回4/公式1の不一致。φ相互作用と純φ結合の分類、水素の空間軌道と電子状態、d混成近似の解釈を要査読。公式に合わせて不存在の説明を作らない。

\paragraph{問題の要旨} 空間軌道・スピンと結合対称性を区別する。

\Needspace{4\baselineskip}\paragraph{確認した条件・記述（独自要約）}\mbox{}\par

\noindent\textbf{判断対象：} 空間軌道・スピンと結合対称性を区別する。\par

\noindent\textbf{原典確認：} 全5候補・指数・否定・記号を保存原典22頁画像と照合。\par

\paragraph{問われる核心} 空間軌道・スピンと結合対称性を区別する。

\paragraph{必要知識} 独立初回は4、公式は1で不一致を保持。

\paragraph{考え方・正答を選ぶ理由（査読待ち初稿）} 独立初回は4、公式は1で不一致を保持。n,l,mは空間軌道、スピンは電子状態の属性で水素にも存在する。候補4を多電子にのみスピンがあると読むと問題がある。一方、σ/π/δ/φは軸まわりの軌道対称性に対応し、2014年のthorocene研究にはφ型軌道相互作用の実験的根拠がある。これを単純な二原子純φ結合の確立と同一視もしない。公式1の採点を表示しつつ、1の一意の不適切性は未解決としてAstraへ送る。

\Needspace{4\baselineskip}\paragraph{選択肢の検討（査読待ち初稿）}\mbox{}\par

\noindent\textbf{1：選ぶ候補} σ/π/δ/φ等の結合。公式はこの候補1。ただし2014年研究にはφ軌道相互作用があり、φの概念の不存在では説明できない。\par

\noindent\textbf{2：選ばない候補} H2OはAB2E2。O周囲は結合対2と孤立対2。\par

\noindent\textbf{3：選ばない候補} dを含む混成もある。d関与の形式的混成モデルはある。超原子価をd混成だけで実在説明するのは別問題。\par

\noindent\textbf{4：選ばない候補} 水素はn,l,mだが多電子では加えてスピン。空間軌道は3量子数、電子状態は水素にもスピンがある。多電子にのみ追加する読みには疑義。\par

\noindent\textbf{5：選ばない候補} 3p充填後は4s。中性原子の通常の構築順でAr後Kに4sが入る。全励起状態・イオンの順序とはしない。\par

\paragraph{再利用する考え方} 軌道と電子状態、格子と基底、単位胞原子数と配位数、定義の分母を区別して候補を全て評価する。

\paragraph{共通理論} \hyperref[theory:atomic:orbital]{第\ref*{theory:atomic:orbital}節}

\par 参考：\cite{atomic_quantum,atomic_phi2014}。

\par 難易度 2（校準初稿）、優先度未評価。

\section{炭素同素体の構造と物性}

\subsection{平成28年度 問14}
\label{q:h28-reg-applied-14}

\noindent\textbf{原典：} sources/H28/17\_applied\_science.pdf、PDF物理5頁。

\textbf{公式正答：} 2（single）。

\par\textbf{検討未確定：} ダイヤモンドの硬度をMohs尺度・通常純結晶へ限定。ドーピング等の導電性例外を保持。

\paragraph{問題の要旨} 炭素同素体の結合構造と導電性を比較する。

\Needspace{4\baselineskip}\paragraph{確認した条件・記述（独自要約）}\mbox{}\par

\noindent\textbf{判断対象：} 炭素同素体の結合構造と導電性を比較する。\par

\noindent\textbf{原典確認：} 全5候補・指数・否定・記号を保存原典22頁画像と照合。\par

\paragraph{問われる核心} 炭素同素体の結合構造と導電性を比較する。

\paragraph{必要知識} 通常の純粋ダイヤモンドはsp3結合網で硬いが良導体ではない。

\paragraph{考え方・正答を選ぶ理由（査読待ち初稿）} 通常の純粋ダイヤモンドはsp3結合網で硬いが良導体ではない。共有結合なら導電性が高いとする候補2が不適切。黒鉛の層内/層間と、CNTの巻き方依存の電子構造を区別する。ホウ素ドープによる導電・超伝導ダイヤモンドの例外は通常純結晶の説明と別に扱う。

\Needspace{4\baselineskip}\paragraph{選択肢の検討（査読待ち初稿）}\mbox{}\par

\noindent\textbf{1：選ばない候補} ダイヤモンドは四面体結合で全鉱物中最高硬度。通常の鉱物のMohs尺度の最高位10。全人工材料・全荷重の絶対比較とはしない。\par

\noindent\textbf{2：選ぶ候補} ダイヤモンドは共有結合結晶で電気をよく通す。通常の純粋ダイヤモンドは大きなバンドギャップを持ち良導体ではない。ドープ材の例外と分ける。\par

\noindent\textbf{3：選ばない候補} 黒鉛は六角網面で層間が弱い。層内の強い共有結合と層間相互作用を分ける。\par

\noindent\textbf{4：選ばない候補} CNTは単層又は多層同軸のナノ管。六角網面の管状構造。原典の同軸形状を保持。\par

\noindent\textbf{5：選ばない候補} CNTの直径・配列で導電性が変わる。巻き方のカイラリティと径が電子構造へ関係する。\par

\paragraph{再利用する考え方} 軌道と電子状態、格子と基底、単位胞原子数と配位数、定義の分母を区別して候補を全て評価する。

\paragraph{共通理論} \hyperref[theory:atomic:carbon]{第\ref*{theory:atomic:carbon}節}

\par 参考：\cite{atomic_cnt,atomic_diamond}。

\paragraph{関連問題} \hyperref[q:r05-reg-applied-10]{R05-REG-APPLIED-10}

\par 難易度 2（校準初稿）、優先度未評価。

\subsection{令和5年度 問10}
\label{q:r05-reg-applied-10}

\noindent\textbf{原典：} sources/R05/17\_applied\_science.pdf、PDF物理4頁。

\textbf{公式正答：} 5（single）。

\par\textbf{検討未確定：} ダイヤモンドの硬度をMohs尺度・通常純結晶へ限定。ドーピング等の導電性例外を保持。

\paragraph{問題の要旨} 炭素同素体の結合構造と導電性を比較する。

\Needspace{4\baselineskip}\paragraph{確認した条件・記述（独自要約）}\mbox{}\par

\noindent\textbf{判断対象：} 炭素同素体の結合構造と導電性を比較する。\par

\noindent\textbf{原典確認：} 全5候補・指数・否定・記号を保存原典22頁画像と照合。\par

\paragraph{問われる核心} 炭素同素体の結合構造と導電性を比較する。

\paragraph{必要知識} 通常の純粋ダイヤモンドはsp3結合網で硬いが良導体ではない。

\paragraph{考え方・正答を選ぶ理由（査読待ち初稿）} 通常の純粋ダイヤモンドはsp3結合網で硬いが良導体ではない。共有結合なら導電性が高いとする候補5が不適切。黒鉛の層内/層間と、CNTの巻き方依存の電子構造を区別する。ホウ素ドープによる導電・超伝導ダイヤモンドの例外は通常純結晶の説明と別に扱う。

\Needspace{4\baselineskip}\paragraph{選択肢の検討（査読待ち初稿）}\mbox{}\par

\noindent\textbf{1：選ばない候補} CNTは単層又は多層同軸のナノ管。六角網面の管状構造。原典の同軸形状を保持。\par

\noindent\textbf{2：選ばない候補} CNTの直径・配列で導電性が変わる。巻き方のカイラリティと径が電子構造へ関係する。\par

\noindent\textbf{3：選ばない候補} 黒鉛は六角網面で層間が弱い。層内の強い共有結合と層間相互作用を分ける。\par

\noindent\textbf{4：選ばない候補} ダイヤモンドは四面体結合で全鉱物中最高硬度。通常の鉱物のMohs尺度の最高位10。全人工材料・全荷重の絶対比較とはしない。\par

\noindent\textbf{5：選ぶ候補} ダイヤモンドは共有結合結晶で電気をよく通す。通常の純粋ダイヤモンドは大きなバンドギャップを持ち良導体ではない。ドープ材の例外と分ける。\par

\paragraph{再利用する考え方} 軌道と電子状態、格子と基底、単位胞原子数と配位数、定義の分母を区別して候補を全て評価する。

\paragraph{共通理論} \hyperref[theory:atomic:carbon]{第\ref*{theory:atomic:carbon}節}

\par 参考：\cite{atomic_cnt,atomic_diamond}。

\paragraph{関連問題} \hyperref[q:h28-reg-applied-14]{H28-REG-APPLIED-14}

\par 難易度 2（校準初稿）、優先度未評価。

\section{同位体の基礎}

\subsection{平成28年度 問10}
\label{q:h28-reg-applied-10}

\noindent\textbf{原典：} sources/H28/17\_applied\_science.pdf、PDF物理4頁。

\textbf{公式正答：} 2（single）。

\paragraph{問題の要旨} 原子番号と質量数から同位体を定義する。

\Needspace{4\baselineskip}\paragraph{確認した条件・記述（独自要約）}\mbox{}\par

\noindent\textbf{判断対象：} 原子番号と質量数から同位体を定義する。\par

\noindent\textbf{原典確認：} 全5候補・指数・否定・記号を保存原典22頁画像と照合。\par

\paragraph{問われる核心} 原子番号と質量数から同位体を定義する。

\paragraph{必要知識} 同位体はZが等しくNが違う原子の組で、A=\allowbreak{}Z+Nは異なる。

\paragraph{考え方・正答を選ぶ理由（査読待ち初稿）} 同位体はZが等しくNが違う原子の組で、A=\allowbreak{}Z+Nは異なる。2。1は同素体、3は同重体、4は周期表の族、5は有機同族列。単体の結晶構造が違うことと核の中性子数が違うことを混同しない。

\Needspace{4\baselineskip}\paragraph{選択肢の検討（査読待ち初稿）}\mbox{}\par

\noindent\textbf{1：選ばない候補} 硫黄の異なる単体。同素体の説明。\par

\noindent\textbf{2：選ぶ候補} 原子番号同一・質量数不同。同じZ、違う中性子数N、A=\allowbreak{}Z+Nが異なる。\par

\noindent\textbf{3：選ばない候補} 陽子と中性子の和が同じ。同じAの同重体でZ同一とは限らない。\par

\noindent\textbf{4：選ばない候補} 周期表の同族。電子配置による化学的族の分類。\par

\noindent\textbf{5：選ばない候補} アルカン等の同族列。有機化合物の同族列で核種ではない。\par

\paragraph{再利用する考え方} 軌道と電子状態、格子と基底、単位胞原子数と配位数、定義の分母を区別して候補を全て評価する。

\paragraph{共通理論} \hyperref[theory:atomic:isotopes]{第\ref*{theory:atomic:isotopes}節}

\par 参考：\cite{modern_mit_nuclear}。

\paragraph{関連問題} \hyperref[q:h29-reg-applied-18]{H29-REG-APPLIED-18}

\par 難易度 2（校準初稿）、優先度未評価。


\chapter{物理化学}
\noindent 分野解説は執筆中。以下は登録問題と校正済み例の作業用一覧。
\section{酸塩基の定義と溶媒依存性}

\subsection{平成30年度 問11}
\label{q:h30-reg-applied-11}

\noindent\textbf{原典：} sources/H30/17\_applied\_science.pdf、PDF物理4頁。

\textbf{公式正答：} 4（single）。

\par\textbf{検討未確定：} 4は濃度と温度を指定しない。初回記録のルクス・フラッドに関する短文は電子対と酸化物イオンを混同しており、原記録を保持して本文では定義を訂正する。

\paragraph{問題の要旨} 溶媒系・酸化物イオン・電子対の酸塩基定義と、硫酸の酸化性を区別する。

\Needspace{4\baselineskip}\paragraph{確認した条件・記述（独自要約）}\mbox{}\par

\noindent\textbf{確認範囲：} 問題幹、否定、全5肢、数式、数値、単位を原典画像で確認した。\par

\noindent\textbf{成立条件：} 最も不適切な記述。酸の強さ、酸化還元、溶媒の相対的なプロトン授受を分離する。\par

\paragraph{問われる核心} 溶媒系・酸化物イオン・電子対の酸塩基定義と、硫酸の酸化性を区別する。

\paragraph{必要知識} 最も不適切な記述。酸の強さ、酸化還元、溶媒の相対的なプロトン授受を分離する。

\paragraph{考え方・正答を選ぶ理由（査読待ち初稿）} 溶媒系では自己解離で生じる陽イオンを増やす物質が酸、ルクス・フラッドでは酸化物イオン受容体、ルイスでは電子対受容体が酸である。酢酸もより強い酸である硫酸溶媒中ではプロトンを受ける塩基として働く。硫酸の酸化性は濃度・温度に依存し、一般に非常に強い酸化剤とする4は適切でない。

\Needspace{4\baselineskip}\paragraph{選択肢の検討（査読待ち初稿）}\mbox{}\par

\noindent\textbf{1：適切} 溶媒の自己解離陽イオンを生じさせる物質は酸。溶媒系の定義であり、水に限らない。\par

\noindent\textbf{2：適切} 酸化物イオンを受け取る物質は酸。ルクス・フラッドの定義である。電子対受容の定義とは区別する。\par

\noindent\textbf{3：適切} 電子対を受け取る物質は酸。ルイスの定義である。\par

\noindent\textbf{4：不適切} 硫酸は非常に強い酸化剤。希硫酸と熱濃硫酸を同じ酸化剤として一括できない。酸性の強さから酸化性を導けない。\par

\noindent\textbf{5：適切} 酢酸は水中で酸、硫酸中で塩基。強酸溶媒中のプロトン受容という相対的役割である。\par

\paragraph{再利用する考え方} 反応速度と平衡位置、全圧と分圧、濃度と活量を分ける。式の適用条件を確認し、原典の不正確な記述を科学的に正しい文へ置き換えない。

\paragraph{共通理論} \hyperref[theory:chem_eq:definitions]{第\ref*{theory:chem_eq:definitions}節}

\par 参考：\cite{iupac_greenbook2012}。

\paragraph{関連問題} \hyperref[q:r06-reg-applied-12]{R06-REG-APPLIED-12}

\par 難易度 2（校準初稿）、優先度未評価。

\subsection{令和6年度 問12}
\label{q:r06-reg-applied-12}

\noindent\textbf{原典：} sources/R06/17\_applied\_science.pdf、PDF物理4頁。

\textbf{公式正答：} 5（single）。

\par\textbf{検討未確定：} 3の原典表現は酸解離定数pKaである。pKaとKaを同一視する文を科学的に修正した原文として保存せず、用語の不正確さを注記する。

\paragraph{問題の要旨} ルイス・ブレンステッド・アレニウスの酸、pKa、塩素のオキソ酸の強さを比較する。

\Needspace{4\baselineskip}\paragraph{確認した条件・記述（独自要約）}\mbox{}\par

\noindent\textbf{確認範囲：} 問題幹、否定、全5肢、数式、数値、単位を原典画像で確認した。\par

\noindent\textbf{成立条件：} 酸解離定数は活量に基づく無次元量でありpKaはその負常用対数。HClOとHClO2は同じ中心元素で比較。\par

\paragraph{問われる核心} ルイス・ブレンステッド・アレニウスの酸、pKa、塩素のオキソ酸の強さを比較する。

\paragraph{必要知識} 酸解離定数は活量に基づく無次元量でありpKaはその負常用対数。HClOとHClO2は同じ中心元素で比較。

\paragraph{考え方・正答を選ぶ理由（査読待ち初稿）} ルイス酸は電子対受容体、ブレンステッド酸はプロトン供与体、アレニウス酸は水中で水素イオンを増やす物質である。\(\mathrm{p}K_a=-\log_{10}K_a\) なので小さいほど強い酸となる。塩素中心のオキソ酸では亜塩素酸HClO2が次亜塩素酸HClOより強く、逆を述べる5が誤りである。

\Needspace{4\baselineskip}\paragraph{選択肢の検討（査読待ち初稿）}\mbox{}\par

\noindent\textbf{1：適切} ルイス酸は電子対受容体。共有結合形成に電子対を受け取る。\par

\noindent\textbf{2：適切} ブレンステッド酸はプロトン供与体。相手の塩基にプロトンを渡す。\par

\noindent\textbf{3：条件付きで適切} 酸解離定数pKaが小さいほど強酸。数値の大小関係はpKaについて正しい。pKaそのものを酸解離定数と呼ぶ点は不正確である。\par

\noindent\textbf{4：適切} アレニウス酸は水中で水素イオンを生じる。水溶液に限定した定義である。\par

\noindent\textbf{5：不適切} 次亜塩素酸は亜塩素酸より強い酸。Purdueの表ではHOClのKaが2.9の10のマイナス8乗、HOClOが1.1の10のマイナス2乗であり、逆である。\par

\paragraph{再利用する考え方} 反応速度と平衡位置、全圧と分圧、濃度と活量を分ける。式の適用条件を確認し、原典の不正確な記述を科学的に正しい文へ置き換えない。

\paragraph{共通理論} \hyperref[theory:chem_eq:definitions]{第\ref*{theory:chem_eq:definitions}節}, \hyperref[theory:chem_eq:acid_strength]{第\ref*{theory:chem_eq:acid_strength}節}

\par 参考：\cite{iupac_greenbook2012,purdue_acid_strength}。

\paragraph{関連問題} \hyperref[q:h30-reg-applied-11]{H30-REG-APPLIED-11}

\par 難易度 2（校準初稿）、優先度未評価。

\section{分離操作の原理と共沸による蒸留限界を判断する。}

\subsection{平成30年度 問18}
\label{q:h30-reg-applied-18}

\noindent\textbf{原典：} sources/H30/17\_applied\_science.pdf、PDF物理6頁。

\textbf{公式正答：} 1（single）。

\par\textbf{検討未確定：} 逆浸透膜の全分子・イオンの完全阻止という記述は理想化。公式は蒸留の共沸例外を選ぶが、実膜の阻止率と対象化学種の条件を留保。

\paragraph{問題の要旨} 分離操作の原理と共沸による蒸留限界を判断する。

\Needspace{4\baselineskip}\paragraph{確認した条件・記述（独自要約）}\mbox{}\par

\noindent\textbf{原典確認：} S7\_049の26原典画像で全五肢・否定・添字を確認。\par

\noindent\textbf{対象：} 分離操作の原理と共沸による蒸留限界を判断する。\par

\paragraph{問われる核心} 分離操作の原理と共沸による蒸留限界を判断する。

\paragraph{必要知識} 蒸留は気液平衡の組成差を使う。共沸点でy=\allowbreak{}xとなる例があるため、必ず異なるという全称が誤りである。吸収は液体内部への溶解、吸着は表面や細孔への保持、抽出は非混和液相間の分配として区別する。逆浸透は圧力で水を透過させる代表的操作だが、理想的な完全溶質阻止をすべての膜に保証しない。

\paragraph{考え方・正答を選ぶ理由（査読待ち初稿）} 蒸留は気液平衡の組成差を使う。共沸点でy=\allowbreak{}xとなる例があるため、必ず異なるという全称が誤りである。吸収は液体内部への溶解、吸着は表面や細孔への保持、抽出は非混和液相間の分配として区別する。逆浸透は圧力で水を透過させる代表的操作だが、理想的な完全溶質阻止をすべての膜に保証しない。

\Needspace{4\baselineskip}\paragraph{選択肢の検討（査読待ち初稿）}\mbox{}\par

\noindent\textbf{1：不適切} 蒸留で気液組成は必ず違う。共沸組成では気液の組成が一致し、通常の単純な蒸留を反復してもその組成を越えて分離できない。\par

\noindent\textbf{2：通常の条件で適切} 気体成分の液体への吸収。気体を液体と接触させ、溶けやすい成分を液相へ移す。表面への吸着と区別する。\par

\noindent\textbf{3：通常の条件で適切} 非混和液相への抽出。二相間の分配平衡を利用し、溶質を片方の相へ移す。\par

\noindent\textbf{4：通常の条件で適切} 活性炭・シリカゲル・ゼオライト。表面や細孔への相互作用により成分を保持する吸着剤として用いられる。\par

\noindent\textbf{5：条件付き・解釈留保} 逆浸透膜の水の選択透過。溶媒を選択的に透過させ、浸透圧差に逆らう圧力で濃縮・脱塩する。実膜は全溶質を100\%阻止する理想膜とは限らない。\par

\paragraph{再利用する考え方} 測定対象と分離対象、定義の分母、絶対量と相対校正、全称と成立条件を分けて判断する。

\paragraph{共通理論} \hyperref[theory:analysis:operations]{第\ref*{theory:analysis:operations}節}

\par 参考：\cite{nptel_azeotrope,inorganic_zeolite}。

\paragraph{関連問題} \hyperref[q:r02-reg-applied-17]{R02-REG-APPLIED-17}

\par 難易度 4（校準初稿）、優先度未評価。

\subsection{令和2年度 問17}
\label{q:r02-reg-applied-17}

\noindent\textbf{原典：} sources/R02/17\_applied\_science.pdf、PDF物理6頁。

\textbf{公式正答：} 4（single）。

\par\textbf{検討未確定：} 逆浸透膜の全分子・イオンの完全阻止という記述は理想化。公式は蒸留の共沸例外を選ぶが、実膜の阻止率と対象化学種の条件を留保。

\paragraph{問題の要旨} 分離操作の原理と共沸による蒸留限界を判断する。

\Needspace{4\baselineskip}\paragraph{確認した条件・記述（独自要約）}\mbox{}\par

\noindent\textbf{原典確認：} S7\_049の26原典画像で全五肢・否定・添字を確認。\par

\noindent\textbf{対象：} 分離操作の原理と共沸による蒸留限界を判断する。\par

\paragraph{問われる核心} 分離操作の原理と共沸による蒸留限界を判断する。

\paragraph{必要知識} 蒸留は気液平衡の組成差を使う。共沸点でy=\allowbreak{}xとなる例があるため、必ず異なるという全称が誤りである。吸収は液体内部への溶解、吸着は表面や細孔への保持、抽出は非混和液相間の分配として区別する。逆浸透は圧力で水を透過させる代表的操作だが、理想的な完全溶質阻止をすべての膜に保証しない。

\paragraph{考え方・正答を選ぶ理由（査読待ち初稿）} 蒸留は気液平衡の組成差を使う。共沸点でy=\allowbreak{}xとなる例があるため、必ず異なるという全称が誤りである。吸収は液体内部への溶解、吸着は表面や細孔への保持、抽出は非混和液相間の分配として区別する。逆浸透は圧力で水を透過させる代表的操作だが、理想的な完全溶質阻止をすべての膜に保証しない。

\Needspace{4\baselineskip}\paragraph{選択肢の検討（査読待ち初稿）}\mbox{}\par

\noindent\textbf{1：通常の条件で適切} 活性炭・シリカゲル・ゼオライト。表面や細孔への相互作用により成分を保持する吸着剤として用いられる。\par

\noindent\textbf{2：通常の条件で適切} 非混和液相への抽出。二相間の分配平衡を利用し、溶質を片方の相へ移す。\par

\noindent\textbf{3：通常の条件で適切} 気体成分の液体への吸収。気体を液体と接触させ、溶けやすい成分を液相へ移す。表面への吸着と区別する。\par

\noindent\textbf{4：不適切} 蒸留で気液組成は必ず違う。共沸組成では気液の組成が一致し、通常の単純な蒸留を反復してもその組成を越えて分離できない。\par

\noindent\textbf{5：条件付き・解釈留保} 逆浸透膜の水の選択透過。溶媒を選択的に透過させ、浸透圧差に逆らう圧力で濃縮・脱塩する。実膜は全溶質を100\%阻止する理想膜とは限らない。\par

\paragraph{再利用する考え方} 測定対象と分離対象、定義の分母、絶対量と相対校正、全称と成立条件を分けて判断する。

\paragraph{共通理論} \hyperref[theory:analysis:operations]{第\ref*{theory:analysis:operations}節}

\par 参考：\cite{nptel_azeotrope,inorganic_zeolite}。

\paragraph{関連問題} \hyperref[q:h30-reg-applied-18]{H30-REG-APPLIED-18}

\par 難易度 4（校準初稿）、優先度未評価。

\section{溶液分母の溶解度と析出収支}

\subsection{平成27年度 問11}
\label{q:h27-reg-applied-11}

\noindent\textbf{原典：} sources/H27/17\_applied\_science.pdf、PDF物理4頁。

\textbf{公式正答：} 4（single）。

\par\textbf{検討未確定：} 水の蒸発なし・析出結晶への水の取込みなしを置く物質収支。原典の分母は飽和溶液。

\paragraph{問題の要旨} 飽和溶液質量を分母にした溶解度から析出量を求める。

\Needspace{4\baselineskip}\paragraph{確認した条件・記述（独自要約）}\mbox{}\par

\noindent\textbf{60℃溶質：} 飽和溶液100g中52.2g\par

\noindent\textbf{20℃溶質：} 飽和溶液100g中24.0g\par

\noindent\textbf{水の初期量：} 47.8g\par

\paragraph{問われる核心} 飽和溶液質量を分母にした溶解度から析出量を求める。

\paragraph{必要知識} 原典の定義は100gの飽和水溶液中の溶質量。

\paragraph{考え方・正答を選ぶ理由（査読待ち初稿）} 原典の定義は100gの飽和水溶液中の溶質量。60℃の100gには溶質52.2g、水47.8gがある。蒸発を無視し20℃で溶質/水=24/76とすると残る溶質は47.8×24/76=\allowbreak{}15.0947g。析出52.2−15.0947=\allowbreak{}37.1053gで4。100g水当たりという通常の別定義へ読み替えない。

\Needspace{4\baselineskip}\paragraph{選択肢の検討（査読待ち初稿）}\mbox{}\par

\noindent\textbf{1：選ばない候補} 15.1g。冷却後の溶存量で析出量ではない。\par

\noindent\textbf{2：選ばない候補} 24.0g。冷却前の水量を100gと誤認。\par

\noindent\textbf{3：選ばない候補} 28.2g。52.2−24.0は溶液量が変わることを無視。\par

\noindent\textbf{4：選ぶ候補} 37.1g。水47.8gを保持し溶存15.0947gを差し引く。\par

\noindent\textbf{5：選ばない候補} 59.0g。初期溶質52.2gより大きく質量保存に反する。\par

\paragraph{再利用する考え方} 軌道と電子状態、格子と基底、単位胞原子数と配位数、定義の分母を区別して候補を全て評価する。

\paragraph{共通理論} \hyperref[theory:atomic:solubility]{第\ref*{theory:atomic:solubility}節}

\par 難易度 3（校準初稿）、優先度未評価。

\section{ボルンハーバーサイクル}

\subsection{平成26年度 問15}
\label{q:h26-reg-applied-15}

\noindent\textbf{原典：} sources/H26/17\_applied\_science.pdf、PDF物理6頁。

\textbf{公式正答：} 2（single）。

\par\textbf{検討未確定：} 原典は「生成熱」を負の候補で表す一方、他の熱量は正の大きさで与える。現代的な生成エンタルピーと熱化学式の右辺の発熱量の符号を区別して説明する。

\paragraph{問題の要旨} ボルン・ハーバーサイクルの各過程に符号を付けてNaClの生成エンタルピーを求める。

\Needspace{4\baselineskip}\paragraph{確認した条件・記述（独自要約）}\mbox{}\par

\noindent\textbf{確認範囲：} 問題幹、数値、式、否定、全5候補を原典画像で照合した。\par

\noindent\textbf{成立条件：} 原典の電子親和力354と格子エネルギー788は正の大きさとして与えられている。電子獲得と結晶生成はエンタルピー低下であり負符号、Cl原子1 molにはCl2を半mol解離する。\par

\paragraph{問われる核心} ボルン・ハーバーサイクルの各過程に符号を付けてNaClの生成エンタルピーを求める。

\paragraph{必要知識} 原典の電子親和力354と格子エネルギー788は正の大きさとして与えられている。電子獲得と結晶生成はエンタルピー低下であり負符号、Cl原子1 molにはCl2を半mol解離する。

\paragraph{考え方・正答を選ぶ理由（査読待ち初稿）} Na(s)の昇華+108、Na(g)のイオン化+502、Cl2(g)半molの解離+243/2、Cl(g)の電子獲得−354、気体イオンからNaCl(s)の生成−788を加える。\(108+502+121.5-354-788=-410.5\,\mathrm{kJ\,mol^{-1}}\)。最も近い−411 kJ/molの2を選ぶ。

\Needspace{4\baselineskip}\paragraph{選択肢の検討（査読待ち初稿）}\mbox{}\par

\noindent\textbf{1：不適切} −289 kJ/mol。Cl2の解離熱243を半分にせず全量使うと−289になる。NaCl 1 molに必要なCl原子は1 molである。\par

\noindent\textbf{2：適切} −411 kJ/mol。半molの塩素分子解離と吸熱・発熱の方向をすべて反映する。\par

\noindent\textbf{3：不適切} −471 kJ/mol。指定された5過程のエンタルピー和−410.5と合わない。\par

\noindent\textbf{4：不適切} −597 kJ/mol。原典の昇華、イオン化、半解離、電子獲得、格子形成の和と合わない。\par

\noindent\textbf{5：不適切} −1166 kJ/mol。吸熱の昇華・イオン化・解離を含めた生成エンタルピーより過剰な負値である。\par

\paragraph{再利用する考え方} 電子親和力と格子エネルギーは文献で符号慣習が異なる。名称だけに符号を割り当てず、反応方向を書いてから足す。

\paragraph{共通理論} \hyperref[theory:electrochem:born_haber]{第\ref*{theory:electrochem:born_haber}節}, \hyperref[theory:electrochem:heat_sign]{第\ref*{theory:electrochem:heat_sign}節}

\par 参考：\cite{openstax_born_haber}。

\paragraph{関連問題} \hyperref[q:h24-reg-applied-14]{H24-REG-APPLIED-14}

\par 難易度 2（校準初稿）、優先度未評価。

\section{共役酸塩基対・緩衝容量・pH}

\subsection{令和3年度 問17}
\label{q:r03-reg-applied-17}

\noindent\textbf{原典：} sources/R03/17\_applied\_science.pdf、PDF物理6頁。

\textbf{公式正答：} 2（single）。

\par\textbf{検討未確定：} 5は混合比を指定しない。部分中和という成立条件を注記し、一律に全混合比で適切とは扱わない。

\paragraph{問題の要旨} 緩衝溶液の定義、血液、式、酸塩基対の調製条件を判定する。

\Needspace{4\baselineskip}\paragraph{確認した条件・記述（独自要約）}\mbox{}\par

\noindent\textbf{確認範囲：} 問題幹、否定、全5肢、数式、数値、単位を原典画像で確認した。\par

\noindent\textbf{成立条件：} 最も不適切な文。緩衝は有限量の酸塩基添加に対する変化の小ささであり、pH完全不変ではない。弱塩基の強酸による部分中和。\par

\paragraph{問われる核心} 緩衝溶液の定義、血液、式、酸塩基対の調製条件を判定する。

\paragraph{必要知識} 最も不適切な文。緩衝は有限量の酸塩基添加に対する変化の小ささであり、pH完全不変ではない。弱塩基の強酸による部分中和。

\paragraph{考え方・正答を選ぶ理由（査読待ち初稿）} 緩衝溶液には弱酸と共役塩基、又は弱塩基と共役酸の両方が必要である。血液のpH維持にも緩衝作用が関与するので、それを否定する2が誤りである。\(K_a=a_{\mathrm{H^+}}a_{\mathrm{A^-}}/a_{\mathrm{HA}}\) を対数化すると \(\mathrm{pH=p}K_a+\log_{10}(a_{\mathrm{A^-}}/a_{\mathrm{HA}})\) となる。弱塩基と強酸の混合は部分中和で両成分を残す条件が必要である。

\Needspace{4\baselineskip}\paragraph{選択肢の検討（査読待ち初稿）}\mbox{}\par

\noindent\textbf{1：適切} 少量の強酸強塩基添加でpHが大きく変わらない。緩衝容量内での性質である。\par

\noindent\textbf{2：不適切} 血液のpH7.4維持に緩衝は関与しない。重炭酸系やタンパク等の酸塩基対が関与する。\par

\noindent\textbf{3：条件付きで適切} ヘンダーソン・ハッセルバルヒ式でpHを求める。活量比の式は平衡から導かれる。濃度比や混合時の物質量比に置換するには近似条件が必要である。\par

\noindent\textbf{4：適切} 弱酸とその共役塩基の塩で調製。両成分が有効量存在するように調製する。\par

\noindent\textbf{5：条件付きで適切} 弱塩基と強酸で緩衝溶液を作る。部分中和で弱塩基と共役酸の両方を残す。強酸過剰では一般に緩衝溶液にならない。\par

\paragraph{再利用する考え方} 反応速度と平衡位置、全圧と分圧、濃度と活量を分ける。式の適用条件を確認し、原典の不正確な記述を科学的に正しい文へ置き換えない。

\paragraph{共通理論} \hyperref[theory:chem_eq:buffer]{第\ref*{theory:chem_eq:buffer}節}

\par 参考：\cite{iupac_greenbook2012,mit_buffer2008}。

\par 難易度 2（校準初稿）、優先度未評価。

\section{触媒・反応経路・平衡不変}

\subsection{平成23年度 問12}
\label{q:h23-reg-applied-12}

\noindent\textbf{原典：} sources/H23/17\_applied\_science.pdf、PDF物理5頁。

\textbf{公式正答：} 3（single）。

\par\textbf{検討未確定：} 負触媒という出題時の用語と現行IUPAC定義の不一致、及び自ら変化せずの正味再生という省略を保持する。

\paragraph{問題の要旨} 触媒が変える反応速度と、変えない平衡定数・反応熱を区別する。

\Needspace{4\baselineskip}\paragraph{確認した条件・記述（独自要約）}\mbox{}\par

\noindent\textbf{確認範囲：} 問題幹、否定、全5肢、数式、数値、単位を原典画像で確認した。\par

\noindent\textbf{成立条件：} 触媒の正味の再生、同一温度での平衡不変、反応熱の状態関数性。\par

\paragraph{問われる核心} 触媒が変える反応速度と、変えない平衡定数・反応熱を区別する。

\paragraph{必要知識} 触媒の正味の再生、同一温度での平衡不変、反応熱の状態関数性。

\paragraph{考え方・正答を選ぶ理由（査読待ち初稿）} 触媒は別の反応経路を与えるが、反応前後の状態が同じなら反応熱を変えない。したがって反応熱の減少を述べる3が不適切である。4の負触媒は歴史的用語として読み、現行IUPACでは速度を下げる作用を阻害として区別する。

\Needspace{4\baselineskip}\paragraph{選択肢の検討（査読待ち初稿）}\mbox{}\par

\noindent\textbf{1：条件付きで適切} 自ら変化せず速度を変え又は反応を開始する。反応過程では中間体として変化し得るが、正味で再生されるという意味である。開始は実用上観測できる速度への増大として読む。\par

\noindent\textbf{2：適切} 平衡に影響しない。同じ温度で反応の標準ギブズエネルギー差を変えない。\par

\noindent\textbf{3：不適切} 反応熱を減少させる。活性化障壁の低下と反応前後のエンタルピー差は別である。\par

\noindent\textbf{4：条件付きで適切} 負触媒は速度を低下させる。歴史的呼称ではその意味であるが、IUPACは負触媒という用語を採用せず阻害を区別する。\par

\noindent\textbf{5：適切} 選択性を持ち特定生成物を与える。競合する反応経路の相対速度を変えて生成物選択性を持ち得る。\par

\paragraph{再利用する考え方} 反応速度と平衡位置、全圧と分圧、濃度と活量を分ける。式の適用条件を確認し、原典の不正確な記述を科学的に正しい文へ置き換えない。

\paragraph{共通理論} \hyperref[theory:chem_eq:catalyst]{第\ref*{theory:chem_eq:catalyst}節}

\par 参考：\cite{iupac_greenbook2012,iupac_catalyst}。

\paragraph{関連問題} \hyperref[q:r01-ret-applied-15]{R01-RET-APPLIED-15}

\par 難易度 2（校準初稿）、優先度未評価。

\subsection{令和元年度再試験 問15}
\label{q:r01-ret-applied-15}

\noindent\textbf{原典：} sources/R01\_retest/17\_applied\_science.pdf、PDF物理6頁。

\textbf{公式正答：} 5（single）。

\paragraph{問題の要旨} 活性化障壁・動的平衡・触媒選択性と、速度と平衡位置の独立性を問う。

\Needspace{4\baselineskip}\paragraph{確認した条件・記述（独自要約）}\mbox{}\par

\noindent\textbf{確認範囲：} 問題幹、否定、全5肢、数式、数値、単位を原典画像で確認した。\par

\noindent\textbf{成立条件：} 温度・圧力一定。同じ反応物と生成物について平衡と速度を比較する。活性化エネルギーは問題の反応経路上の障壁模型として読む。\par

\paragraph{問われる核心} 活性化障壁・動的平衡・触媒選択性と、速度と平衡位置の独立性を問う。

\paragraph{必要知識} 温度・圧力一定。同じ反応物と生成物について平衡と速度を比較する。活性化エネルギーは問題の反応経路上の障壁模型として読む。

\paragraph{考え方・正答を選ぶ理由（査読待ち初稿）} 触媒は障壁の低い別経路を与えるが、温度一定で平衡定数を変えない。平衡では正逆の反応が停止するのではなく同じ速度で継続する。定温定圧で反応速度が増すと平衡位置も変わるとする5が誤りである。

\Needspace{4\baselineskip}\paragraph{選択肢の検討（査読待ち初稿）}\mbox{}\par

\noindent\textbf{1：適切} 反応物から遷移状態への差が障壁。指定した反応経路のエネルギー障壁を説明する。経験的アレニウス量と無条件に同一視しない。\par

\noindent\textbf{2：適切} 平衡では正逆反応速度が等しい。濃度の正味変化が0でも各方向の反応は進む。\par

\noindent\textbf{3：適切} 触媒は低い障壁の経路を与える。同じ反応の別経路を通じて速度を上げる。\par

\noindent\textbf{4：適切} 触媒は選択性を持つ。競合経路の相対速度を変え得る。\par

\noindent\textbf{5：不適切} 速度増大で定温定圧の平衡位置が変わる。速度は平衡への到達時間に関わり、平衡位置は自由エネルギーで決まる。\par

\paragraph{再利用する考え方} 反応速度と平衡位置、全圧と分圧、濃度と活量を分ける。式の適用条件を確認し、原典の不正確な記述を科学的に正しい文へ置き換えない。

\paragraph{共通理論} \hyperref[theory:chem_eq:catalyst]{第\ref*{theory:chem_eq:catalyst}節}, \hyperref[theory:chem_eq:activation]{第\ref*{theory:chem_eq:activation}節}, \hyperref[theory:chem_eq:rates]{第\ref*{theory:chem_eq:rates}節}

\par 参考：\cite{iupac_greenbook2012}。

\paragraph{関連問題} \hyperref[q:h23-reg-applied-12]{H23-REG-APPLIED-12}

\par 難易度 2（校準初稿）、優先度未評価。

\section{束一的性質・凝固点降下}

\subsection{平成24年度 問11}
\label{q:h24-reg-applied-11}

\noindent\textbf{原典：} sources/H24/17\_applied\_science.pdf、PDF物理3頁。

\textbf{公式正答：} 4（single）。

\paragraph{問題の要旨} 完全解離する食塩の溶液について、粒子数を数えて凝固点降下を計算する。

\Needspace{4\baselineskip}\paragraph{確認した条件・記述（独自要約）}\mbox{}\par

\noindent\textbf{確認範囲：} 問題幹、条件、否定、数値、単位、全5肢を原典画像で照合した。\par

\noindent\textbf{成立条件：} 水1.0 kg、NaCl 0.20 mol、\(K_f=1.85\,\mathrm{K\,kg\,mol^{-1}}\)、完全イオン解離の指定。希薄溶液の束一的性質の式を用いる。\par

\paragraph{問われる核心} 完全解離する食塩の溶液について、粒子数を数えて凝固点降下を計算する。

\paragraph{必要知識} 水1.0 kg、NaCl 0.20 mol、\(K_f=1.85\,\mathrm{K\,kg\,mol^{-1}}\)、完全イオン解離の指定。希薄溶液の束一的性質の式を用いる。

\paragraph{考え方・正答を選ぶ理由（査読待ち初稿）} 溶媒の水は1.0 kg、NaClは0.20 molである。完全解離によりNaとClのイオンが各0.20 molとなり、粒子基準の質量モル濃度は \(2\times0.20/1.0=0.40\,\mathrm{mol\,kg^{-1}}\) である。正の凝固点降下量は \(\Delta T_f=1.85\times0.40=0.74\,\mathrm K\)。純水の凝固点を0°Cとして溶液は−0.74°Cとなり4を選ぶ。

\Needspace{4\baselineskip}\paragraph{選択肢の検討（査読待ち初稿）}\mbox{}\par

\noindent\textbf{1：不適切} −0.19°C。完全解離後の粒子濃度0.40 mol/kgを使った降下量0.74 Kと合わない。\par

\noindent\textbf{2：不適切} −0.37°C。NaClの式量単位だけを数え、2種類のイオンへの解離を無視した値である。\par

\noindent\textbf{3：不適切} −0.56°C。粒子濃度と係数の積0.74 Kに合わない。\par

\noindent\textbf{4：適切} −0.74°C。解離による粒子数2倍と凝固点の低下方向を共に反映する。\par

\noindent\textbf{5：不適切} −1.9°C。指定濃度は粒子基準でも0.40 mol/kgであり、1 mol/kg相当の降下とは異なる。\par

\paragraph{再利用する考え方} 質量モル濃度の分母は溶液体積でなく溶媒質量である。完全解離の仮定と、濃厚電解質の実際の活量補正を区別する。

\paragraph{共通理論} \hyperref[theory:phase_sep:freezing]{第\ref*{theory:phase_sep:freezing}節}

\par 参考：\cite{mit_colligative2008}。

\par 難易度 2（校準初稿）、優先度未評価。

\section{衝突・反応断面積・拡散律速}

\subsection{令和元年度 問15}
\label{q:r01-reg-applied-15}

\noindent\textbf{原典：} sources/R01/17\_applied\_science.pdf、PDF物理5頁。

\textbf{公式正答：} 3（single）。

\par\textbf{検討未確定：} 衝突頻度の対象、立体因子模型、アレニウス活性化エネルギーを障壁差で代表する省略を原典通り保持し、適用範囲を教材で限定する。

\paragraph{問題の要旨} 衝突模型、拡散律速、平衡の熱力学、活性化エネルギーの意味を比較する。

\Needspace{4\baselineskip}\paragraph{確認した条件・記述（独自要約）}\mbox{}\par

\noindent\textbf{確認範囲：} 問題幹、否定、全5肢、数式、数値、単位を原典画像で確認した。\par

\noindent\textbf{成立条件：} 気体圧力比較は温度・組成一定。衝突頻度の単位と対象を明記する。アレニウス式の経験的活性化エネルギーと電子的障壁高さは一般には別。\par

\paragraph{問われる核心} 衝突模型、拡散律速、平衡の熱力学、活性化エネルギーの意味を比較する。

\paragraph{必要知識} 気体圧力比較は温度・組成一定。衝突頻度の単位と対象を明記する。アレニウス式の経験的活性化エネルギーと電子的障壁高さは一般には別。

\paragraph{考え方・正答を選ぶ理由（査読待ち初稿）} 出会った後のプロトン移動が十分速い中和では、分子やイオンが出会う拡散が速度を制限する。反応そのものが律速とする3が不適切である。1は温度と組成一定の分子1個あたりの衝突頻度、2は単純な立体因子模型、5は障壁模型という適用範囲を持つ。

\Needspace{4\baselineskip}\paragraph{選択肢の検討（査読待ち初稿）}\mbox{}\par

\noindent\textbf{1：条件付きで適切} 衝突頻度は圧力に比例。分子1個あたりの頻度は数密度に比例する。単位体積あたりの衝突数なら固定組成では圧力の2乗に比例する。\par

\noindent\textbf{2：条件付きで適切} 反応断面積は衝突断面積と立体因子の積。配向の割合を掛ける単純模型である。一般の反応断面積の定義を尽くさない。\par

\noindent\textbf{3：不適切} 速い酸塩基中和は反応律速。遭遇後の反応が速い極限では拡散による遭遇が律速する。\par

\noindent\textbf{4：適切} 標準自由エネルギー変化は平衡定数を決める。\(K=\exp(-\Delta_rG^\circ/RT)\) に従う。\par

\noindent\textbf{5：条件付きで適切} アレニウス活性化エネルギーは活性錯合体との差。単純な障壁模型の説明である。厳密な経験的定義は \(E_a=RT^2d\ln k/dT\) であり、常に障壁高さと同じではない。\par

\paragraph{再利用する考え方} 反応速度と平衡位置、全圧と分圧、濃度と活量を分ける。式の適用条件を確認し、原典の不正確な記述を科学的に正しい文へ置き換えない。

\paragraph{共通理論} \hyperref[theory:chem_eq:collision]{第\ref*{theory:chem_eq:collision}節}, \hyperref[theory:chem_eq:diffusion]{第\ref*{theory:chem_eq:diffusion}節}, \hyperref[theory:chem_eq:activation]{第\ref*{theory:chem_eq:activation}節}

\par 参考：\cite{iupac_greenbook2012}。

\paragraph{関連問題} \hyperref[q:r02-reg-applied-13]{R02-REG-APPLIED-13}

\par 難易度 2（校準初稿）、優先度未評価。

\subsection{令和2年度 問13}
\label{q:r02-reg-applied-13}

\noindent\textbf{原典：} sources/R02/17\_applied\_science.pdf、PDF物理5頁。

\textbf{公式正答：} 2（single）。

\par\textbf{独立解答と公式正答に相違あり：} 独立解答 3。初回3、公式2の不一致を保持する。原典3の速度比は一般に誤りであり、2も比較する断面積の定義に依存する。最も不適切を一意に説明する根拠は未解決であり、公式に合わせて3を速度定数に書き換えない。

\paragraph{問題の要旨} 衝突断面積と反応断面積、反応速度と速度定数、拡散律速の定義を区別する。

\Needspace{4\baselineskip}\paragraph{確認した条件・記述（独自要約）}\mbox{}\par

\noindent\textbf{確認範囲：} 問題幹、否定、全5肢、数式、数値、単位を原典画像で確認した。\par

\noindent\textbf{成立条件：} 原典3は逆反応の反応速度と正反応の反応速度の比。速度定数へ置換しない。断面積は入射流束に対する有効面積であり、比較する衝突の定義とエネルギーを指定する。\par

\paragraph{問われる核心} 衝突断面積と反応断面積、反応速度と速度定数、拡散律速の定義を区別する。

\paragraph{必要知識} 原典3は逆反応の反応速度と正反応の反応速度の比。速度定数へ置換しない。断面積は入射流束に対する有効面積であり、比較する衝突の定義とエネルギーを指定する。

\paragraph{考え方・正答を選ぶ理由（査読待ち初稿）} 1は分子1個あたりの \(z=\sigma\bar v_{\mathrm{rel}}n\)、4は活性化模型、5は拡散律速を述べる。3の原文は速度定数ではなく反応速度の比である。平衡では正逆の速度が等しく比は1なので、一般の平衡定数を表せず、初回は3を選んだ。公式正答は2である。2について剛体球の幾何学断面積と実効的反応断面積を区別する必要があるが、断面積の比較条件は原典に明示されない。2の公式判定だけで3の記述を適切に直すことはできない。

\Needspace{4\baselineskip}\paragraph{選択肢の検討（査読待ち初稿）}\mbox{}\par

\noindent\textbf{1：適切} 衝突頻度は断面積・相対速さ・数密度に比例。分子1個あたりの衝突頻度の式である。\par

\noindent\textbf{2：解釈未解決} 反応断面積は衝突断面積を上回らない。同じ遭遇集合のうち反応する割合という模型なら成立する。剛体球幾何学面積との比較を一般則とすることはできない。公式2の根拠となる断面積の意味は追加査読が必要である。\par

\noindent\textbf{3：不適切} 平衡定数は逆反応速度と正反応速度の比。動的平衡では速度比は1である。一次可逆反応の速度定数比でも順序は正方向の定数を逆方向で割る。\par

\noindent\textbf{4：条件付きで適切} 二分子速度定数は障壁を超える衝突割合に依存。活性化障壁を持つ単純衝突模型である。無障壁反応や拡散律速を包括する一般式ではない。\par

\noindent\textbf{5：適切} 遭遇対が分離より速く生成物を作るなら拡散で決まる。出会った後の反応確率が高い拡散律速の極限である。\par

\paragraph{再利用する考え方} 反応速度と平衡位置、全圧と分圧、濃度と活量を分ける。式の適用条件を確認し、原典の不正確な記述を科学的に正しい文へ置き換えない。

\paragraph{共通理論} \hyperref[theory:chem_eq:collision]{第\ref*{theory:chem_eq:collision}節}, \hyperref[theory:chem_eq:diffusion]{第\ref*{theory:chem_eq:diffusion}節}, \hyperref[theory:chem_eq:rates]{第\ref*{theory:chem_eq:rates}節}

\par 参考：\cite{iupac_greenbook2012}。

\paragraph{関連問題} \hyperref[q:r01-reg-applied-15]{R01-REG-APPLIED-15}

\par 難易度 2（校準初稿）、優先度未評価。

\section{コロイド分散・会合・乳化}

\subsection{令和6年度 問11}
\label{q:r06-reg-applied-11}

\noindent\textbf{原典：} sources/R06/17\_applied\_science.pdf、PDF物理4頁。

\textbf{公式正答：} 2（single）。

\par\textbf{検討未確定：} 溶解という語の範囲、石けんの濃度・温度・ミセル形成条件を原典は限定しない。公式2と初回2を保持し、全称の説明はしない。

\paragraph{問題の要旨} コロイドの分散と分子・イオンの溶解を、乳濁液・光散乱・石けんで区別する。

\Needspace{4\baselineskip}\paragraph{確認した条件・記述（独自要約）}\mbox{}\par

\noindent\textbf{確認範囲：} 問題幹、条件、否定、数値、単位、全5肢を原典画像で照合した。\par

\noindent\textbf{成立条件：} 最も関係が薄い用語の選択。石けんの会合は濃度・温度に依存し、通常の分子・イオンの溶液とコロイド分散を区別する。\par

\paragraph{問われる核心} コロイドの分散と分子・イオンの溶解を、乳濁液・光散乱・石けんで区別する。

\paragraph{必要知識} 最も関係が薄い用語の選択。石けんの会合は濃度・温度に依存し、通常の分子・イオンの溶液とコロイド分散を区別する。

\paragraph{考え方・正答を選ぶ理由（査読待ち初稿）} エマルションは液滴の分散、チンダル現象は分散粒子による光散乱、マヨネーズは乳濁液の具体例、石けん水はミセル等の会合体を含み得る。通常の分子・イオンとして均一に溶ける溶解との直接性が最も低い2を選ぶ。ただしコロイド溶液という慣用語や高分子の溶解もあり、溶解とコロイドが一切無関係という意味ではない。

\Needspace{4\baselineskip}\paragraph{選択肢の検討（査読待ち初稿）}\mbox{}\par

\noindent\textbf{1：適切} エマルション。互いに混ざりにくい液体の一方が微小な液滴として他方へ分散した乳濁液である。\par

\noindent\textbf{2：条件付きで適切} 溶解。通常の小分子・イオンの溶解はコロイドの粒子分散とは区別する。高分子や会合体を含む系まで語だけで排除できない。\par

\noindent\textbf{3：適切} チンダル現象。分散粒子が光を散乱し、光の経路が見える現象である。\par

\noindent\textbf{4：適切} マヨネーズ。油滴が水性の連続相に分散し、卵黄成分などが乳化を安定化する例である。\par

\noindent\textbf{5：条件付きで適切} せっけん水。石けんは親水部と疎水部を持ち、適切な温度・濃度でミセル等の会合体を作る。全濃度の水溶液が同じコロイド状態という意味ではない。\par

\paragraph{再利用する考え方} 物質名より、何がどの尺度で分散しているかを読む。乳化した油滴と、石けん分子だけが集まったミセルを区別する。

\paragraph{共通理論} \hyperref[theory:phase_sep:colloid]{第\ref*{theory:phase_sep:colloid}節}, \hyperref[theory:phase_sep:micelle]{第\ref*{theory:phase_sep:micelle}節}

\par 参考：\cite{openstax_colloids}。

\par 難易度 2（校準初稿）、優先度未評価。

\section{真空排気・低温凝縮吸着}

\subsection{平成30年度 問8}
\label{q:h30-reg-applied-08}

\noindent\textbf{原典：} sources/H30/17\_applied\_science.pdf、PDF物理3頁。

\textbf{公式正答：} 1（single）。

\paragraph{問題の要旨} 気体の低温凝縮又は吸着を利用する真空ポンプを識別する。

\Needspace{4\baselineskip}\paragraph{確認した条件・記述（独自要約）}\mbox{}\par

\noindent\textbf{確認範囲：} 問題幹、条件、否定、数値、単位、全5肢を原典画像で照合した。\par

\noindent\textbf{成立条件：} 気体の凝縮又は吸着を利用する装置。拡散ポンプの油蒸気の循環と、排気対象を冷面へ捕捉する原理を区別する。\par

\paragraph{問われる核心} 気体の低温凝縮又は吸着を利用する真空ポンプを識別する。

\paragraph{必要知識} 気体の凝縮又は吸着を利用する装置。拡散ポンプの油蒸気の循環と、排気対象を冷面へ捕捉する原理を区別する。

\paragraph{考え方・正答を選ぶ理由（査読待ち初稿）} クライオポンプは低温面で気体を凝縮させ、凝縮しにくい気体を低温の吸着材へ捕捉する。気体分子を系内で捕捉する原理を問うため1を選ぶ。他の4候補は運動量輸送又は機械的な排気を主体とする。

\Needspace{4\baselineskip}\paragraph{選択肢の検討（査読待ち初稿）}\mbox{}\par

\noindent\textbf{1：適切} クライオポンプ。低温面での凝縮と、活性炭等の吸着材への低温吸着で排気対象を捕捉する。\par

\noindent\textbf{2：不適切} ターボ分子ポンプ。高速回転翼が気体分子へ方向を持つ運動量を与えて排出する。\par

\noindent\textbf{3：不適切} 拡散ポンプ。噴出する作動油等の蒸気が排気対象の気体を運ぶ。作動流体の凝縮があっても、対象気体を凝縮捕捉する主原理ではない。\par

\noindent\textbf{4：不適切} ダイヤフラムポンプ。膜の往復変形と弁で容積を変え、気体を吸入・排出する。\par

\noindent\textbf{5：不適切} ルーツ型ポンプ。互いにかみ合うロータで気体を吸入口から排気側へ移送する。\par

\paragraph{再利用する考え方} 排気対象の捕捉と、作動流体の相変化を分ける。クライオポンプには捕捉容量があり、再生が必要になる。

\paragraph{共通理論} \hyperref[theory:phase_sep:vacuum]{第\ref*{theory:phase_sep:vacuum}節}

\par 参考：\cite{leybold_cryopump}。

\par 難易度 2（校準初稿）、優先度未評価。

\section{蒸留・気液平衡比・共沸}

\subsection{平成27年度 問15}
\label{q:h27-reg-applied-15}

\noindent\textbf{原典：} sources/H27/17\_applied\_science.pdf、PDF物理5頁。

\textbf{公式正答：} 1（single）。

\par\textbf{検討未確定：} 平衡定数の意味を原典は限定しない。気液平衡比にもK値があり、他肢より関係が薄いことの一意な根拠はAstra査読へ残す。

\paragraph{問題の要旨} 蒸留の装置・操作・気液平衡を表す5用語の直接性を比較する。

\Needspace{4\baselineskip}\paragraph{確認した条件・記述（独自要約）}\mbox{}\par

\noindent\textbf{確認範囲：} 問題幹、条件、否定、数値、単位、全5肢を原典画像で照合した。\par

\noindent\textbf{成立条件：} 最も関係が薄い用語を選ぶ。反応平衡定数と相分配の気液平衡比を区別する。\par

\paragraph{問われる核心} 蒸留の装置・操作・気液平衡を表す5用語の直接性を比較する。

\paragraph{必要知識} 最も関係が薄い用語を選ぶ。反応平衡定数と相分配の気液平衡比を区別する。

\paragraph{考え方・正答を選ぶ理由（査読待ち初稿）} 理論段数、相対揮発度、還流、充填塔は精留の設計又は操作に直接現れる。公式は平衡定数の1を選ぶ。化学反応の平衡定数として読めばこの選択になるが、気液平衡のK値も蒸留に使用する。平衡という概念自体が蒸留と無関係という説明は成立しない。

\Needspace{4\baselineskip}\paragraph{選択肢の検討（査読待ち初稿）}\mbox{}\par

\noindent\textbf{1：解釈未解決} 平衡定数。化学反応の平衡定数という意味なら他候補より直接性が低いが、蒸留の気液平衡比 \(K_i=y_i/x_i\) もK値と呼ばれるため、関係がないとは断定できない。\par

\noindent\textbf{2：適切} 理論段数。各段で気液が平衡に達する理想化に必要な接触段数であり、精留塔の分離能力を表す。\par

\noindent\textbf{3：適切} 相対揮発度。2成分の気液平衡比の比であり、蒸留の分離の容易さを表す。\par

\noindent\textbf{4：適切} 還流。凝縮した塔頂液の一部を塔へ戻し、上昇蒸気と接触させる。\par

\noindent\textbf{5：適切} 充填塔。充填物の表面で液体と蒸気を接触させる装置であり、蒸留に使用する。\par

\paragraph{再利用する考え方} 蒸留は平衡を利用する分離操作である。\(K_i=y_i/x_i\) と反応の平衡定数の区別を、用語だけでなく定義から行う。

\paragraph{共通理論} \hyperref[theory:phase_sep:distillation]{第\ref*{theory:phase_sep:distillation}節}, \hyperref[theory:phase_sep:stages]{第\ref*{theory:phase_sep:stages}節}

\par 参考：\cite{nptel_distillation1}。

\paragraph{関連問題} \hyperref[q:r07-reg-applied-18]{R07-REG-APPLIED-18}

\par 難易度 2（校準初稿）、優先度未評価。

\subsection{令和元年度再試験 問10}
\label{q:r01-ret-applied-10}

\noindent\textbf{原典：} sources/R01\_retest/17\_applied\_science.pdf、PDF物理4頁。

\textbf{公式正答：} 4（single）。

\par\textbf{検討未確定：} 96\%の濃度基準を原典は指定しない。常圧の代表的な共沸値は質量基準約96\%と整合するが、厳密数値としての同定は追加査読事項に残す。

\paragraph{問題の要旨} 精留・水蒸気蒸留・減圧蒸留の目的と、共沸による通常精留の限界を比較する。

\Needspace{4\baselineskip}\paragraph{確認した条件・記述（独自要約）}\mbox{}\par

\noindent\textbf{確認範囲：} 問題幹、条件、否定、数値、単位、全5肢を原典画像で照合した。\par

\noindent\textbf{成立条件：} 外圧 \(1.013\times10^5\,\mathrm{Pa}\)。原典のエタノール96\%には質量・体積・モルの基準が明示されない。共沸の数値の解釈と、通常精留不能の原理を分ける。\par

\paragraph{問われる核心} 精留・水蒸気蒸留・減圧蒸留の目的と、共沸による通常精留の限界を比較する。

\paragraph{必要知識} 外圧 \(1.013\times10^5\,\mathrm{Pa}\)。原典のエタノール96\%には質量・体積・モルの基準が明示されない。共沸の数値の解釈と、通常精留不能の原理を分ける。

\paragraph{考え方・正答を選ぶ理由（査読待ち初稿）} 共沸組成では気相組成と液相組成が一致し、段数や還流を増してもその圧力での通常精留による分離は進まない。完全に水と分離できるとする4が不適切である。減圧は沸点を下げる操作、水蒸気蒸留はほぼ不混和な物質を水蒸気と留出させる操作であり、役割が異なる。

\Needspace{4\baselineskip}\paragraph{選択肢の検討（査読待ち初稿）}\mbox{}\par

\noindent\textbf{1：適切} より揮発性が高い成分は蒸気、低い成分は液に濃縮。同じ状態で気相への移りやすさの異なる成分を分ける。共沸点ではその差が失われるため一般操作の原理として読む。\par

\noindent\textbf{2：適切} 充填塔や棚段塔で還流・向流接触・蒸発と分縮を反復。上昇する蒸気と下降する液を接触させ、多段の分離に相当する効果を得る。\par

\noindent\textbf{3：適切} ほぼ混和しない物質へ水蒸気を吹き込み揮発成分を留出。水と対象の分圧が全圧へ寄与し、両者が共存する条件で水蒸気と共に成分を運ぶ。\par

\noindent\textbf{4：不適切} 約96\%のエタノール水溶液の共沸も通常の精留で完全分離できる。指定圧力の共沸では気液の組成が同じになり、通常の精留による濃縮が進まない。\par

\noindent\textbf{5：適切} 減圧蒸留は沸点を下げ熱分解を避けるのに適する。低い外圧で沸騰でき、熱に弱い成分の分離に適する。ただし減圧だけで全共沸が必ず消えるわけではない。\par

\paragraph{再利用する考え方} 分離装置の能力を増すことと、平衡曲線そのものを変えることは異なる。濃度パーセントは基準を確認し、96\%をモル分率と決め付けない。

\paragraph{共通理論} \hyperref[theory:phase_sep:distillation]{第\ref*{theory:phase_sep:distillation}節}, \hyperref[theory:phase_sep:azeotrope]{第\ref*{theory:phase_sep:azeotrope}節}, \hyperref[theory:phase_sep:steam]{第\ref*{theory:phase_sep:steam}節}

\par 参考：\cite{nptel_distillation1,nptel_azeotrope}。

\paragraph{関連問題} \hyperref[q:r06-reg-applied-15]{R06-REG-APPLIED-15}

\par 難易度 2（校準初稿）、優先度未評価。

\subsection{令和6年度 問15}
\label{q:r06-reg-applied-15}

\noindent\textbf{原典：} sources/R06/17\_applied\_science.pdf、PDF物理5頁。

\textbf{公式正答：} 1（single）。

\par\textbf{検討未確定：} 96\%の濃度基準を原典は指定しない。常圧の代表的な共沸値は質量基準約96\%と整合するが、厳密数値としての同定は追加査読事項に残す。

\paragraph{問題の要旨} 精留・水蒸気蒸留・減圧蒸留の目的と、共沸による通常精留の限界を比較する。

\Needspace{4\baselineskip}\paragraph{確認した条件・記述（独自要約）}\mbox{}\par

\noindent\textbf{確認範囲：} 問題幹、条件、否定、数値、単位、全5肢を原典画像で照合した。\par

\noindent\textbf{成立条件：} 外圧1気圧。原典のエタノール96\%には質量・体積・モルの基準が明示されない。共沸の数値の解釈と、通常精留不能の原理を分ける。\par

\paragraph{問われる核心} 精留・水蒸気蒸留・減圧蒸留の目的と、共沸による通常精留の限界を比較する。

\paragraph{必要知識} 外圧1気圧。原典のエタノール96\%には質量・体積・モルの基準が明示されない。共沸の数値の解釈と、通常精留不能の原理を分ける。

\paragraph{考え方・正答を選ぶ理由（査読待ち初稿）} 共沸組成では気相組成と液相組成が一致し、段数や還流を増してもその圧力での通常精留による分離は進まない。完全に水と分離できるとする1が不適切である。減圧は沸点を下げる操作、水蒸気蒸留はほぼ不混和な物質を水蒸気と留出させる操作であり、役割が異なる。

\Needspace{4\baselineskip}\paragraph{選択肢の検討（査読待ち初稿）}\mbox{}\par

\noindent\textbf{1：不適切} 約96\%のエタノール水溶液の共沸も通常の精留で完全分離できる。指定圧力の共沸では気液の組成が同じになり、通常の精留による濃縮が進まない。\par

\noindent\textbf{2：適切} 減圧蒸留は沸点を下げ熱分解を避けるのに適する。低い外圧で沸騰でき、熱に弱い成分の分離に適する。ただし減圧だけで全共沸が必ず消えるわけではない。\par

\noindent\textbf{3：適切} より揮発性が高い成分は蒸気、低い成分は液に濃縮。同じ状態で気相への移りやすさの異なる成分を分ける。共沸点ではその差が失われるため一般操作の原理として読む。\par

\noindent\textbf{4：適切} ほぼ混和しない物質へ水蒸気を吹き込み揮発成分を留出。水と対象の分圧が全圧へ寄与し、両者が共存する条件で水蒸気と共に成分を運ぶ。\par

\noindent\textbf{5：適切} 充填塔や棚段塔で還流・向流接触・蒸発と分縮を反復。上昇する蒸気と下降する液を接触させ、多段の分離に相当する効果を得る。\par

\paragraph{再利用する考え方} 分離装置の能力を増すことと、平衡曲線そのものを変えることは異なる。濃度パーセントは基準を確認し、96\%をモル分率と決め付けない。

\paragraph{共通理論} \hyperref[theory:phase_sep:distillation]{第\ref*{theory:phase_sep:distillation}節}, \hyperref[theory:phase_sep:azeotrope]{第\ref*{theory:phase_sep:azeotrope}節}, \hyperref[theory:phase_sep:steam]{第\ref*{theory:phase_sep:steam}節}

\par 参考：\cite{nptel_distillation1,nptel_azeotrope}。

\paragraph{関連問題} \hyperref[q:r01-ret-applied-10]{R01-RET-APPLIED-10}

\par 難易度 2（校準初稿）、優先度未評価。

\subsection{令和7年度 問18}
\label{q:r07-reg-applied-18}

\noindent\textbf{原典：} sources/R07/17\_applied\_science.pdf、PDF物理6頁。

\textbf{公式正答：} 4（single）。

\par\textbf{検討未確定：} 平衡定数の意味を原典は限定しない。気液平衡比にもK値があり、他肢より関係が薄いことの一意な根拠はAstra査読へ残す。

\paragraph{問題の要旨} 蒸留の装置・操作・気液平衡を表す5用語の直接性を比較する。

\Needspace{4\baselineskip}\paragraph{確認した条件・記述（独自要約）}\mbox{}\par

\noindent\textbf{確認範囲：} 問題幹、条件、否定、数値、単位、全5肢を原典画像で照合した。\par

\noindent\textbf{成立条件：} 最も関係が薄い用語を選ぶ。反応平衡定数と相分配の気液平衡比を区別する。\par

\paragraph{問われる核心} 蒸留の装置・操作・気液平衡を表す5用語の直接性を比較する。

\paragraph{必要知識} 最も関係が薄い用語を選ぶ。反応平衡定数と相分配の気液平衡比を区別する。

\paragraph{考え方・正答を選ぶ理由（査読待ち初稿）} 理論段数、相対揮発度、還流、充填塔は精留の設計又は操作に直接現れる。公式は平衡定数の4を選ぶ。化学反応の平衡定数として読めばこの選択になるが、気液平衡のK値も蒸留に使用する。平衡という概念自体が蒸留と無関係という説明は成立しない。

\Needspace{4\baselineskip}\paragraph{選択肢の検討（査読待ち初稿）}\mbox{}\par

\noindent\textbf{1：適切} 還流。凝縮した塔頂液の一部を塔へ戻し、上昇蒸気と接触させる。\par

\noindent\textbf{2：適切} 充填塔。充填物の表面で液体と蒸気を接触させる装置であり、蒸留に使用する。\par

\noindent\textbf{3：適切} 相対揮発度。2成分の気液平衡比の比であり、蒸留の分離の容易さを表す。\par

\noindent\textbf{4：解釈未解決} 平衡定数。化学反応の平衡定数という意味なら他候補より直接性が低いが、蒸留の気液平衡比 \(K_i=y_i/x_i\) もK値と呼ばれるため、関係がないとは断定できない。\par

\noindent\textbf{5：適切} 理論段数。各段で気液が平衡に達する理想化に必要な接触段数であり、精留塔の分離能力を表す。\par

\paragraph{再利用する考え方} 蒸留は平衡を利用する分離操作である。\(K_i=y_i/x_i\) と反応の平衡定数の区別を、用語だけでなく定義から行う。

\paragraph{共通理論} \hyperref[theory:phase_sep:distillation]{第\ref*{theory:phase_sep:distillation}節}, \hyperref[theory:phase_sep:stages]{第\ref*{theory:phase_sep:stages}節}

\par 参考：\cite{nptel_distillation1}。

\paragraph{関連問題} \hyperref[q:h27-reg-applied-15]{H27-REG-APPLIED-15}

\par 難易度 2（校準初稿）、優先度未評価。

\section{電池電位・反応自由エネルギー}

\subsection{令和4年度 問15}
\label{q:r04-reg-applied-15}

\noindent\textbf{原典：} sources/R04/17\_applied\_science.pdf、PDF物理5頁。

\textbf{公式正答：} 3（single）。

\paragraph{問題の要旨} 電池の可逆起電力と反応ギブズエネルギー、平衡、標準電極電位の関係を判定する。

\Needspace{4\baselineskip}\paragraph{確認した条件・記述（独自要約）}\mbox{}\par

\noindent\textbf{確認範囲：} 問題幹、数値、式、否定、全5候補を原典画像で照合した。\par

\noindent\textbf{成立条件：} 可逆・開回路の起電力を扱う。標準電極電位は共通の還元電位表で比較し、全反応の平衡と個別界面の平衡を分ける。\par

\paragraph{問われる核心} 電池の可逆起電力と反応ギブズエネルギー、平衡、標準電極電位の関係を判定する。

\paragraph{必要知識} 可逆・開回路の起電力を扱う。標準電極電位は共通の還元電位表で比較し、全反応の平衡と個別界面の平衡を分ける。

\paragraph{考え方・正答を選ぶ理由（査読待ち初稿）} 反応1 molの進行で外部へ運ばれる電荷はnFであり、可逆に取り出す電気仕事の最大値は \(-\Delta_rG\) である。したがって \(\Delta_rG=-nFE\) であり、無関係とする3が不適切である。全電池反応の化学平衡ではE=\allowbreak{}0となるが、個々の電極が局所平衡の開回路電池で起電力が0とは限らない。

\Needspace{4\baselineskip}\paragraph{選択肢の検討（査読待ち初稿）}\mbox{}\par

\noindent\textbf{1：適切} 酸化電極をアノード、還元電極をカソードと呼ぶ。反応による定義である。放電電池と電気分解では正負の対応が異なっても、酸化・還元による名称は変わらない。\par

\noindent\textbf{2：適切} 電池反応全体が化学平衡なら電位差0。全反応の \(\Delta_rG=0\) から \(E=0\)。個別電極界面が局所平衡である開回路状態だけとは区別する。\par

\noindent\textbf{3：不適切} 電池起電力は反応ギブズエネルギーと無関係。可逆起電力は \(\Delta_rG=-nFE\) で結ばれる。nは釣合い全反応で移動する電子のmol数である。\par

\noindent\textbf{4：適切} 標準水素電極を0とする基準。標準電極電位はこの参照電極に対する相対値として表す。単独電極の絶対電位を直接測定する意味ではない。\par

\noindent\textbf{5：適切} 標準電池電位は両標準電極電位の差。双方を還元電位として \(E^\circ=E^\circ_{\mathrm{cathode}}-E^\circ_{\mathrm{anode}}\)。反応の係数倍で電位を倍にしない。\par

\paragraph{再利用する考え方} 起電力は示強量である。反応式を倍にすると反応エネルギーと電子数が共に倍になり、Eは変わらない。

\paragraph{共通理論} \hyperref[theory:electrochem:gibbs]{第\ref*{theory:electrochem:gibbs}節}, \hyperref[theory:electrochem:nernst]{第\ref*{theory:electrochem:nernst}節}, \hyperref[theory:electrochem:electrodes]{第\ref*{theory:electrochem:electrodes}節}

\par 参考：\cite{openstax_cell_gibbs,mit_redox2014}。

\paragraph{関連問題} \hyperref[q:r07-reg-applied-10]{R07-REG-APPLIED-10}

\par 難易度 2（校準初稿）、優先度未評価。

\subsection{令和7年度 問10}
\label{q:r07-reg-applied-10}

\noindent\textbf{原典：} sources/R07/17\_applied\_science.pdf、PDF物理4頁。

\textbf{公式正答：} 1（single）。

\paragraph{問題の要旨} 電池の可逆起電力と反応ギブズエネルギー、平衡、標準電極電位の関係を判定する。

\Needspace{4\baselineskip}\paragraph{確認した条件・記述（独自要約）}\mbox{}\par

\noindent\textbf{確認範囲：} 問題幹、数値、式、否定、全5候補を原典画像で照合した。\par

\noindent\textbf{成立条件：} 可逆・開回路の起電力を扱う。標準電極電位は共通の還元電位表で比較し、全反応の平衡と個別界面の平衡を分ける。\par

\paragraph{問われる核心} 電池の可逆起電力と反応ギブズエネルギー、平衡、標準電極電位の関係を判定する。

\paragraph{必要知識} 可逆・開回路の起電力を扱う。標準電極電位は共通の還元電位表で比較し、全反応の平衡と個別界面の平衡を分ける。

\paragraph{考え方・正答を選ぶ理由（査読待ち初稿）} 反応1 molの進行で外部へ運ばれる電荷はnFであり、可逆に取り出す電気仕事の最大値は \(-\Delta_rG\) である。したがって \(\Delta_rG=-nFE\) であり、無関係とする1が不適切である。全電池反応の化学平衡ではE=\allowbreak{}0となるが、個々の電極が局所平衡の開回路電池で起電力が0とは限らない。

\Needspace{4\baselineskip}\paragraph{選択肢の検討（査読待ち初稿）}\mbox{}\par

\noindent\textbf{1：不適切} 電池起電力は反応ギブズエネルギーと無関係。可逆起電力は \(\Delta_rG=-nFE\) で結ばれる。nは釣合い全反応で移動する電子のmol数である。\par

\noindent\textbf{2：適切} 電池反応全体が化学平衡なら電位差0。全反応の \(\Delta_rG=0\) から \(E=0\)。個別電極界面が局所平衡である開回路状態だけとは区別する。\par

\noindent\textbf{3：適切} 酸化電極をアノード、還元電極をカソードと呼ぶ。反応による定義である。放電電池と電気分解では正負の対応が異なっても、酸化・還元による名称は変わらない。\par

\noindent\textbf{4：適切} 標準水素電極を0とする基準。標準電極電位はこの参照電極に対する相対値として表す。単独電極の絶対電位を直接測定する意味ではない。\par

\noindent\textbf{5：適切} 標準電池電位は両標準電極電位の差。双方を還元電位として \(E^\circ=E^\circ_{\mathrm{cathode}}-E^\circ_{\mathrm{anode}}\)。反応の係数倍で電位を倍にしない。\par

\paragraph{再利用する考え方} 起電力は示強量である。反応式を倍にすると反応エネルギーと電子数が共に倍になり、Eは変わらない。

\paragraph{共通理論} \hyperref[theory:electrochem:gibbs]{第\ref*{theory:electrochem:gibbs}節}, \hyperref[theory:electrochem:nernst]{第\ref*{theory:electrochem:nernst}節}, \hyperref[theory:electrochem:electrodes]{第\ref*{theory:electrochem:electrodes}節}

\par 参考：\cite{openstax_cell_gibbs,mit_redox2014}。

\paragraph{関連問題} \hyperref[q:r04-reg-applied-15]{R04-REG-APPLIED-15}

\par 難易度 2（校準初稿）、優先度未評価。

\section{ファラデーの法則・電気分解}

\subsection{平成25年度 問12}
\label{q:h25-reg-applied-12}

\noindent\textbf{原典：} sources/H25/17\_applied\_science.pdf、PDF物理5頁。

\textbf{公式正答：} 4（single）。

\paragraph{問題の要旨} 二価の銅イオンの析出量から必要電荷と通電時間を求める。

\Needspace{4\baselineskip}\paragraph{確認した条件・記述（独自要約）}\mbox{}\par

\noindent\textbf{確認範囲：} 問題幹、数値、式、否定、全5候補を原典画像で照合した。\par

\noindent\textbf{成立条件：} 電流5.36 A、Cu 1.27 g、モル質量63.5 g/mol、F=\allowbreak{}96500 C/mol。全電流が電気分解に使われるという原典指定により電流効率1とする。\par

\paragraph{問われる核心} 二価の銅イオンの析出量から必要電荷と通電時間を求める。

\paragraph{必要知識} 電流5.36 A、Cu 1.27 g、モル質量63.5 g/mol、F=\allowbreak{}96500 C/mol。全電流が電気分解に使われるという原典指定により電流効率1とする。

\paragraph{考え方・正答を選ぶ理由（査読待ち初稿）} 析出した銅は \(1.27/63.5=0.0200\,\mathrm{mol}\)。\(\mathrm{Cu^{2+}}+2e^-\to\mathrm{Cu}\) に必要な電子は0.0400 molである。\(Q=n_eF=0.0400\times96500=3860\,\mathrm C\)、\(t=Q/I=3860/5.36=720.15\,\mathrm s\) である。60で割ると約12.0分となり4を選ぶ。

\Needspace{4\baselineskip}\paragraph{選択肢の検討（査読待ち初稿）}\mbox{}\par

\noindent\textbf{1：不適切} 2分。必要電荷3860 Cに対し5.36×120=\allowbreak{}643.2 Cしか流れない。\par

\noindent\textbf{2：不適切} 6分。1929.6 Cであり、Cu2+に2電子が必要なことを無視するとこの時間に近付く。\par

\noindent\textbf{3：不適切} 10分。3216 Cであり、指定析出量に必要な3860 Cより不足する。\par

\noindent\textbf{4：適切} 12分。3859.2 Cとなり丸めた必要電荷と一致する。\par

\noindent\textbf{5：不適切} 20分。6432 Cとなり、指定析出量に必要な電荷より大きい。\par

\paragraph{再利用する考え方} 物質量→電子量→電荷→秒→分の順に単位をつないで計算する。Fを銅1 molの電荷とせず電子1 molの電荷として使う。

\paragraph{共通理論} \hyperref[theory:electrochem:faraday]{第\ref*{theory:electrochem:faraday}節}, \hyperref[theory:electrochem:copper]{第\ref*{theory:electrochem:copper}節}

\par 参考：\cite{mit_redox2014}。

\paragraph{関連問題} \hyperref[q:h26-reg-applied-13]{H26-REG-APPLIED-13}

\par 難易度 2（校準初稿）、優先度未評価。

\section{反応熱・ヘスの法則}

\subsection{平成24年度 問14}
\label{q:h24-reg-applied-14}

\noindent\textbf{原典：} sources/H24/17\_applied\_science.pdf、PDF物理5頁。

\textbf{公式正答：} 4（single）。

\paragraph{問題の要旨} 熱化学方程式を組み合わせ、指定量の液体ベンゼンの完全燃焼熱を求める。

\Needspace{4\baselineskip}\paragraph{確認した条件・記述（独自要約）}\mbox{}\par

\noindent\textbf{確認範囲：} 問題幹、数値、式、否定、全5候補を原典画像で照合した。\par

\noindent\textbf{成立条件：} 原典の水は液体、炭素は黒鉛、ベンゼンは液体。式に付く熱量の符号を反応エンタルピーへ変換し、指定物質量を最後に掛ける。\par

\paragraph{問われる核心} 熱化学方程式を組み合わせ、指定量の液体ベンゼンの完全燃焼熱を求める。

\paragraph{必要知識} 原典の水は液体、炭素は黒鉛、ベンゼンは液体。式に付く熱量の符号を反応エンタルピーへ変換し、指定物質量を最後に掛ける。

\paragraph{考え方・正答を選ぶ理由（査読待ち初稿）} ベンゼンの生成式の右辺にある−49 kJは吸熱を表し、生成エンタルピーは+49 kJ/molである。燃焼式は \(\mathrm{C_6H_6(l)}+(15/2)\mathrm{O_2(g)}\to6\mathrm{CO_2(g)}+3\mathrm{H_2O(l)}\)。ヘスの法則で \(\Delta H=6(-394)+3(-286)-49=-3271\,\mathrm{kJ\,mol^{-1}}\) となる。1 molの放出熱量は3271 kJであり、最も近い3300 kJの4を選ぶ。

\Needspace{4\baselineskip}\paragraph{選択肢の検討（査読待ち初稿）}\mbox{}\par

\noindent\textbf{1：不適切} 900 kJ。生成物のCO2六つと水三つの寄与を合計した値より大きく不足する。\par

\noindent\textbf{2：不適切} 1700 kJ。炭素と水素の両方の酸化による発熱を正しく合計した値に達しない。\par

\noindent\textbf{3：不適切} 2500 kJ。生成物と反応物のエンタルピー差から求めた値より小さい。\par

\noindent\textbf{4：適切} 3300 kJ。原典指定の量の放出熱量に最も近い。\par

\noindent\textbf{5：不適切} 4100 kJ。原典の反応係数と指定物質量から求めた値より大きい。\par

\paragraph{再利用する考え方} 反応を逆にするとエンタルピーの符号が逆転し、係数倍すると熱量も同じ倍率になる。放出熱量の正値と負の反応エンタルピーを区別する。

\paragraph{共通理論} \hyperref[theory:electrochem:hess]{第\ref*{theory:electrochem:hess}節}, \hyperref[theory:electrochem:heat_sign]{第\ref*{theory:electrochem:heat_sign}節}, \hyperref[theory:electrochem:combustion]{第\ref*{theory:electrochem:combustion}節}

\par 参考：\cite{openstax_hess}。

\paragraph{関連問題} \hyperref[q:h26-reg-applied-15]{H26-REG-APPLIED-15}, \hyperref[q:h28-reg-applied-12]{H28-REG-APPLIED-12}, \hyperref[q:r01-ret-applied-12]{R01-RET-APPLIED-12}

\par 難易度 2（校準初稿）、優先度未評価。

\subsection{平成28年度 問12}
\label{q:h28-reg-applied-12}

\noindent\textbf{原典：} sources/H28/17\_applied\_science.pdf、PDF物理5頁。

\textbf{公式正答：} 4（single）。

\paragraph{問題の要旨} 熱化学方程式を組み合わせ、指定量の液体ベンゼンの完全燃焼熱を求める。

\Needspace{4\baselineskip}\paragraph{確認した条件・記述（独自要約）}\mbox{}\par

\noindent\textbf{確認範囲：} 問題幹、数値、式、否定、全5候補を原典画像で照合した。\par

\noindent\textbf{成立条件：} 原典の水は液体、炭素は黒鉛、ベンゼンは液体。式に付く熱量の符号を反応エンタルピーへ変換し、指定物質量を最後に掛ける。\par

\paragraph{問われる核心} 熱化学方程式を組み合わせ、指定量の液体ベンゼンの完全燃焼熱を求める。

\paragraph{必要知識} 原典の水は液体、炭素は黒鉛、ベンゼンは液体。式に付く熱量の符号を反応エンタルピーへ変換し、指定物質量を最後に掛ける。

\paragraph{考え方・正答を選ぶ理由（査読待ち初稿）} ベンゼンの生成式の右辺にある−49 kJは吸熱を表し、生成エンタルピーは+49 kJ/molである。燃焼式は \(\mathrm{C_6H_6(l)}+(15/2)\mathrm{O_2(g)}\to6\mathrm{CO_2(g)}+3\mathrm{H_2O(l)}\)。ヘスの法則で \(\Delta H=6(-394)+3(-286)-49=-3271\,\mathrm{kJ\,mol^{-1}}\) となる。1 molの放出熱量は3271 kJであり、最も近い3300 kJの4を選ぶ。

\Needspace{4\baselineskip}\paragraph{選択肢の検討（査読待ち初稿）}\mbox{}\par

\noindent\textbf{1：不適切} 900 kJ。生成物のCO2六つと水三つの寄与を合計した値より大きく不足する。\par

\noindent\textbf{2：不適切} 1700 kJ。炭素と水素の両方の酸化による発熱を正しく合計した値に達しない。\par

\noindent\textbf{3：不適切} 2500 kJ。生成物と反応物のエンタルピー差から求めた値より小さい。\par

\noindent\textbf{4：適切} 3300 kJ。原典指定の量の放出熱量に最も近い。\par

\noindent\textbf{5：不適切} 4100 kJ。原典の反応係数と指定物質量から求めた値より大きい。\par

\paragraph{再利用する考え方} 反応を逆にするとエンタルピーの符号が逆転し、係数倍すると熱量も同じ倍率になる。放出熱量の正値と負の反応エンタルピーを区別する。

\paragraph{共通理論} \hyperref[theory:electrochem:hess]{第\ref*{theory:electrochem:hess}節}, \hyperref[theory:electrochem:heat_sign]{第\ref*{theory:electrochem:heat_sign}節}, \hyperref[theory:electrochem:combustion]{第\ref*{theory:electrochem:combustion}節}

\par 参考：\cite{openstax_hess}。

\paragraph{関連問題} \hyperref[q:h24-reg-applied-14]{H24-REG-APPLIED-14}, \hyperref[q:r01-ret-applied-12]{R01-RET-APPLIED-12}

\par 難易度 2（校準初稿）、優先度未評価。

\subsection{令和元年度再試験 問12}
\label{q:r01-ret-applied-12}

\noindent\textbf{原典：} sources/R01\_retest/17\_applied\_science.pdf、PDF物理5頁。

\textbf{公式正答：} 4（single）。

\paragraph{問題の要旨} 熱化学方程式を組み合わせ、指定量の液体ベンゼンの完全燃焼熱を求める。

\Needspace{4\baselineskip}\paragraph{確認した条件・記述（独自要約）}\mbox{}\par

\noindent\textbf{確認範囲：} 問題幹、数値、式、否定、全5候補を原典画像で照合した。\par

\noindent\textbf{成立条件：} 原典の水は液体、炭素は黒鉛、ベンゼンは液体。式に付く熱量の符号を反応エンタルピーへ変換し、指定物質量を最後に掛ける。\par

\paragraph{問われる核心} 熱化学方程式を組み合わせ、指定量の液体ベンゼンの完全燃焼熱を求める。

\paragraph{必要知識} 原典の水は液体、炭素は黒鉛、ベンゼンは液体。式に付く熱量の符号を反応エンタルピーへ変換し、指定物質量を最後に掛ける。

\paragraph{考え方・正答を選ぶ理由（査読待ち初稿）} ベンゼンの生成式の右辺にある−49 kJは吸熱を表し、生成エンタルピーは+49 kJ/molである。燃焼式は \(\mathrm{C_6H_6(l)}+(15/2)\mathrm{O_2(g)}\to6\mathrm{CO_2(g)}+3\mathrm{H_2O(l)}\)。ヘスの法則で \(\Delta H=6(-394)+3(-286)-49=-3271\,\mathrm{kJ\,mol^{-1}}\) となる。0.5 molでは放出熱量が \(0.5\times3271=1635.5\,\mathrm{kJ}\) なので1650 kJの4を選ぶ。

\Needspace{4\baselineskip}\paragraph{選択肢の検討（査読待ち初稿）}\mbox{}\par

\noindent\textbf{1：不適切} 450 kJ。生成物のCO2六つと水三つの寄与を合計した値より大きく不足する。\par

\noindent\textbf{2：不適切} 850 kJ。炭素と水素の両方の酸化による発熱を正しく合計した値に達しない。\par

\noindent\textbf{3：不適切} 1250 kJ。生成物と反応物のエンタルピー差から求めた値より小さい。\par

\noindent\textbf{4：適切} 1650 kJ。原典指定の量の放出熱量に最も近い。\par

\noindent\textbf{5：不適切} 2050 kJ。原典の反応係数と指定物質量から求めた値より大きい。\par

\paragraph{再利用する考え方} 反応を逆にするとエンタルピーの符号が逆転し、係数倍すると熱量も同じ倍率になる。放出熱量の正値と負の反応エンタルピーを区別する。

\paragraph{共通理論} \hyperref[theory:electrochem:hess]{第\ref*{theory:electrochem:hess}節}, \hyperref[theory:electrochem:heat_sign]{第\ref*{theory:electrochem:heat_sign}節}, \hyperref[theory:electrochem:combustion]{第\ref*{theory:electrochem:combustion}節}

\par 参考：\cite{openstax_hess}。

\paragraph{関連問題} \hyperref[q:h24-reg-applied-14]{H24-REG-APPLIED-14}, \hyperref[q:h28-reg-applied-12]{H28-REG-APPLIED-12}

\par 難易度 2（校準初稿）、優先度未評価。

\section{鉛蓄電池・電解液質量収支}

\subsection{平成26年度 問13}
\label{q:h26-reg-applied-13}

\noindent\textbf{原典：} sources/H26/17\_applied\_science.pdf、PDF物理5頁。

\textbf{公式正答：} 4（single）。

\par\textbf{検討未確定：} 電流効率と外部の水損失を原典は明示しない。「使用後」を放電と解釈した主反応のみのモデルで公式4と一致する。充電や副反応を含む一般条件へ広げない。

\paragraph{問題の要旨} 鉛蓄電池の放電反応から、硫酸量と電解液量の両方を更新して濃度を求める。

\Needspace{4\baselineskip}\paragraph{確認した条件・記述（独自要約）}\mbox{}\par

\noindent\textbf{確認範囲：} 問題幹、数値、式、否定、全5候補を原典画像で照合した。\par

\noindent\textbf{成立条件：} 30質量\%の電解液1000 g、100 A、965 s、F=\allowbreak{}96500 C/mol。使用後を放電と読み、全電荷が主反応へ使われ、外部との液交換や蒸発がないモデルを用いる。\par

\paragraph{問われる核心} 鉛蓄電池の放電反応から、硫酸量と電解液量の両方を更新して濃度を求める。

\paragraph{必要知識} 30質量\%の電解液1000 g、100 A、965 s、F=\allowbreak{}96500 C/mol。使用後を放電と読み、全電荷が主反応へ使われ、外部との液交換や蒸発がないモデルを用いる。

\paragraph{考え方・正答を選ぶ理由（査読待ち初稿）} 放電の全反応は \(\mathrm{Pb+PbO_2+2H_2SO_4\to2PbSO_4+2H_2O}\) で、2 molの電子に対し硫酸2 molを消費し水2 molを生成する。通過電子は \(100\times965/96500=1.00\,\mathrm{mol}\)。硫酸98 gが減り水18 gが増える。硫酸は \(300-98=202\,\mathrm g\)、電解液全体は \(1000-98+18=920\,\mathrm g\) となる。質量百分率 \(100\times202/920=21.96\%\) に最も近い22\%の4を選ぶ。

\Needspace{4\baselineskip}\paragraph{選択肢の検討（査読待ち初稿）}\mbox{}\par

\noindent\textbf{1：不適切} 11\%。電子1 molに対する硫酸消費量1 molと水生成量1 molを使った質量収支に合わない。\par

\noindent\textbf{2：不適切} 16\%。硫酸の減少と水の生成を合わせた最終質量920 gに対する202 gの比ではない。\par

\noindent\textbf{3：不適切} 20\%。202/1000を約20\%とするのは電解液全質量を初期値のまま固定する誤りである。\par

\noindent\textbf{4：適切} 22\%。硫酸202 gを最終電解液920 gで割った約21.96\%に最も近い。\par

\noindent\textbf{5：不適切} 25\%。反応による硫酸消費98 gを反映した最終比率より大きい。\par

\paragraph{再利用する考え方} 電極上の固体PbSO4へ移る質量は液相の分母から除く。硫酸が水になるという言い換えだけで分母を1000 gのまま保たない。

\paragraph{共通理論} \hyperref[theory:electrochem:lead]{第\ref*{theory:electrochem:lead}節}, \hyperref[theory:electrochem:solution_mass]{第\ref*{theory:electrochem:solution_mass}節}, \hyperref[theory:electrochem:faraday]{第\ref*{theory:electrochem:faraday}節}

\par 参考：\cite{openstax_batteries,mit_redox2014}。

\paragraph{関連問題} \hyperref[q:h25-reg-applied-12]{H25-REG-APPLIED-12}

\par 難易度 3（校準初稿）、優先度未評価。

\section{リチウム電池・イオン挿入}

\subsection{平成24年度 問17}
\label{q:h24-reg-applied-17}

\noindent\textbf{原典：} sources/H24/17\_applied\_science.pdf、PDF物理6頁。

\textbf{公式正答：} 4（single）。

\par\textbf{検討未確定：} 「最も卑」は電極反応・溶媒・標準状態を原典が明記しない。通常の水溶液標準還元電位の比較として扱い、全溶媒系への全称とはしない。

\paragraph{問題の要旨} リチウムイオン電池の材料と、金属リチウム負極の課題を5下線部で判定する。

\Needspace{4\baselineskip}\paragraph{確認した条件・記述（独自要約）}\mbox{}\par

\noindent\textbf{確認範囲：} 問題幹、数値、式、否定、全5候補を原典画像で照合した。\par

\noindent\textbf{成立条件：} 原典の下線4はOCRの「自金」でなく「白金」である。標準電極電位は通常の水溶液系還元電位の基準、負極は放電時の役割として読む。\par

\paragraph{問われる核心} リチウムイオン電池の材料と、金属リチウム負極の課題を5下線部で判定する。

\paragraph{必要知識} 原典の下線4はOCRの「自金」でなく「白金」である。標準電極電位は通常の水溶液系還元電位の基準、負極は放電時の役割として読む。

\paragraph{考え方・正答を選ぶ理由（査読待ち初稿）} 金属リチウムの析出によるデンドライトは内部短絡を起こし得る。実用化で負極にはリチウムイオンを取り込む炭素系材料が採用され、白金という下線4が不適切である。初期の石油コークスと後の黒鉛を歴史上同一材料としない。正極のコバルト酸リチウムはLiCoO2である。

\Needspace{4\baselineskip}\paragraph{選択肢の検討（査読待ち初稿）}\mbox{}\par

\noindent\textbf{1：適切} 卑な。通常の標準還元電位でLi+/Liは非常に負である。金属リチウムを酸化させる高電圧電池の負極としての性質につながる。\par

\noindent\textbf{2：適切} 小さい。リチウムは金属元素で最も小さいモル質量を持つ。1 molあたりの移動電子数と合わせて質量当たり容量を考える。\par

\noindent\textbf{3：適切} 負極。放電時にリチウムを酸化させ、電子を外部へ供給する側である。\par

\noindent\textbf{4：不適切} 白金。この実用化経路の負極は炭素系材料であり、白金ではない。白金の不活性電極としての別の用途と混同しない。\par

\noindent\textbf{5：適切} コバルト酸リチウム。LiCoO2はリチウムイオンの出入りを利用する実用正極材料である。\par

\paragraph{再利用する考え方} 高い標準電位差だけで実用電池は決まらない。可逆にイオンを取り込むホスト、電解液、界面、短絡の抑制を合わせて読む。

\paragraph{共通理論} \hyperref[theory:electrochem:lithium]{第\ref*{theory:electrochem:lithium}節}, \hyperref[theory:electrochem:electrodes]{第\ref*{theory:electrochem:electrodes}節}

\par 参考：\cite{nobel_lithium2019,openstax_batteries}。

\par 難易度 2（校準初稿）、優先度未評価。

\section{相状態・臨界・実在気体}

\subsection{平成26年度 問10}
\label{q:h26-reg-applied-10}

\noindent\textbf{原典：} sources/H26/17\_applied\_science.pdf、PDF物理4頁。

\textbf{公式正答：} 3（single）。

\par\textbf{検討未確定：} 熱容量の差を仕事だけで説明する原典肢には理想気体の指定がない。純水の相平衡温度は丸めた常用値である。

\paragraph{問題の要旨} 相平衡の化学ポテンシャル、臨界点、実在気体と定圧熱容量の条件を比較する。

\Needspace{4\baselineskip}\paragraph{確認した条件・記述（独自要約）}\mbox{}\par

\noindent\textbf{確認範囲：} 問題幹、条件、否定、数値、単位、全5肢を原典画像で照合した。\par

\noindent\textbf{成立条件：} 標準状態でのCO2と窒素の比較。純物質の相平衡、圧縮による液化、昇華、同じ昇温の熱容量比較。\par

\paragraph{問われる核心} 相平衡の化学ポテンシャル、臨界点、実在気体と定圧熱容量の条件を比較する。

\paragraph{必要知識} 標準状態でのCO2と窒素の比較。純物質の相平衡、圧縮による液化、昇華、同じ昇温の熱容量比較。

\paragraph{考え方・正答を選ぶ理由（査読待ち初稿）} 相平衡は相ごとの化学ポテンシャルの一致であり、臨界点では気液共存が終わる。CO2の方が窒素より理想性のずれが小さいという比較が逆であるため、3を選ぶ。熱容量差を外部仕事だけで説明する肢は理想気体を想定する必要がある。

\Needspace{4\baselineskip}\paragraph{選択肢の検討（査読待ち初稿）}\mbox{}\par

\noindent\textbf{1：条件付きで適切} 純水のモルGは1気圧100°Cで気液、0°Cで固液が等しい。純物質の相平衡条件である。原典の0°C・100°Cは常用の丸め値であり、精密な共存温度の定義ではない。\par

\noindent\textbf{2：適切} 臨界点では気液の区別が消え、臨界温度以上で圧縮だけでは通常の液化ができない。臨界点は気液共存曲線の終端である。超臨界の流体を高密度にすることと、気液界面を持つ液化は異なる。\par

\noindent\textbf{3：不適切} 標準状態で二酸化炭素は窒素より理想気体からのずれが小さい。低圧の標準的条件ではCO2の相互作用による非理想性が窒素より大きい。理想性の比較は同じ温度・圧力で行い、気体名だけで全状態へ一般化しない。\par

\noindent\textbf{4：条件付きで適切} 定圧加熱が定積より多く熱を要する差は外部仕事による。理想気体の同じ初期状態と同じ昇温なら \(U=U(T)\) であり差は膨張仕事となる。実在気体では最終体積の違いに伴う内部エネルギー差もあり、仕事だけという説明は一般化できない。\par

\noindent\textbf{5：適切} 固体が液相を経ず気体に昇華することがある。固相と気相が直接関わる相変化であり、氷も温度・圧力条件に応じて昇華する。\par

\paragraph{再利用する考え方} 相平衡の等式と、実在気体への近似を区別する。定圧・定積で同じ最終温度でも最終体積は異なり、実在気体の内部エネルギーは一般に同じにならない。

\paragraph{共通理論} \hyperref[theory:phase_sep:chemical_potential]{第\ref*{theory:phase_sep:chemical_potential}節}, \hyperref[theory:phase_sep:critical]{第\ref*{theory:phase_sep:critical}節}, \hyperref[theory:phase_sep:realgas]{第\ref*{theory:phase_sep:realgas}節}, \hyperref[theory:phase_sep:capacity]{第\ref*{theory:phase_sep:capacity}節}

\par 参考：\cite{iupac_greenbook2012,mit_thermo_firstlaw2008}。

\paragraph{関連問題} \hyperref[q:r01-reg-applied-16]{R01-REG-APPLIED-16}

\par 難易度 2（校準初稿）、優先度未評価。

\subsection{令和元年度 問16}
\label{q:r01-reg-applied-16}

\noindent\textbf{原典：} sources/R01/17\_applied\_science.pdf、PDF物理5頁。

\textbf{公式正答：} 4（single）。

\par\textbf{検討未確定：} 熱容量の差を仕事だけで説明する原典肢には理想気体の指定がない。

\paragraph{問題の要旨} 相平衡の化学ポテンシャル、臨界点、実在気体と定圧熱容量の条件を比較する。

\Needspace{4\baselineskip}\paragraph{確認した条件・記述（独自要約）}\mbox{}\par

\noindent\textbf{確認範囲：} 問題幹、条件、否定、数値、単位、全5肢を原典画像で照合した。\par

\noindent\textbf{成立条件：} 標準状態でのCO2と窒素の比較。純物質の相平衡、圧縮による液化、昇華、同じ昇温の熱容量比較。\par

\paragraph{問われる核心} 相平衡の化学ポテンシャル、臨界点、実在気体と定圧熱容量の条件を比較する。

\paragraph{必要知識} 標準状態でのCO2と窒素の比較。純物質の相平衡、圧縮による液化、昇華、同じ昇温の熱容量比較。

\paragraph{考え方・正答を選ぶ理由（査読待ち初稿）} 相平衡は相ごとの化学ポテンシャルの一致であり、臨界点では気液共存が終わる。CO2の方が窒素より理想性のずれが小さいという比較が逆であるため、4を選ぶ。熱容量差を外部仕事だけで説明する肢は理想気体を想定する必要がある。

\Needspace{4\baselineskip}\paragraph{選択肢の検討（査読待ち初稿）}\mbox{}\par

\noindent\textbf{1：適切} 相平衡にある固体と液体の化学ポテンシャルは等しい。相間の物質移動でギブズエネルギーが低下しない条件である。純物質の化学ポテンシャルはモルギブズエネルギーである。\par

\noindent\textbf{2：適切} 臨界点では気液の区別が消え、臨界温度以上で圧縮だけでは通常の液化ができない。臨界点は気液共存曲線の終端である。超臨界の流体を高密度にすることと、気液界面を持つ液化は異なる。\par

\noindent\textbf{3：適切} 固体が液相を経ず気体に昇華することがある。固相と気相が直接関わる相変化であり、氷も温度・圧力条件に応じて昇華する。\par

\noindent\textbf{4：不適切} 標準状態で二酸化炭素は窒素より理想気体からのずれが小さい。低圧の標準的条件ではCO2の相互作用による非理想性が窒素より大きい。理想性の比較は同じ温度・圧力で行い、気体名だけで全状態へ一般化しない。\par

\noindent\textbf{5：条件付きで適切} 定圧加熱が定積より多く熱を要する差は外部仕事による。理想気体の同じ初期状態と同じ昇温なら \(U=U(T)\) であり差は膨張仕事となる。実在気体では最終体積の違いに伴う内部エネルギー差もあり、仕事だけという説明は一般化できない。\par

\paragraph{再利用する考え方} 相平衡の等式と、実在気体への近似を区別する。定圧・定積で同じ最終温度でも最終体積は異なり、実在気体の内部エネルギーは一般に同じにならない。

\paragraph{共通理論} \hyperref[theory:phase_sep:chemical_potential]{第\ref*{theory:phase_sep:chemical_potential}節}, \hyperref[theory:phase_sep:critical]{第\ref*{theory:phase_sep:critical}節}, \hyperref[theory:phase_sep:realgas]{第\ref*{theory:phase_sep:realgas}節}, \hyperref[theory:phase_sep:capacity]{第\ref*{theory:phase_sep:capacity}節}

\par 参考：\cite{iupac_greenbook2012,mit_thermo_firstlaw2008}。

\paragraph{関連問題} \hyperref[q:h26-reg-applied-10]{H26-REG-APPLIED-10}

\par 難易度 2（校準初稿）、優先度未評価。

\subsection{令和3年度 問15}
\label{q:r03-reg-applied-15}

\noindent\textbf{原典：} sources/R03/17\_applied\_science.pdf、PDF物理6頁。

\textbf{公式正答：} 3（single）。

\par\textbf{検討未確定：} 気体の種類を原典は指定しない。冷却の全称には理想気体の反例があり、ギブズエネルギーの判定条件には温度一定等が省略される。

\paragraph{問題の要旨} 外圧0の断熱膨張での保存量を確認し、気体が必ず冷えるかを検討する。

\Needspace{4\baselineskip}\paragraph{確認した条件・記述（独自要約）}\mbox{}\par

\noindent\textbf{確認範囲：} 問題幹、条件、否定、数値、単位、全5肢を原典画像で照合した。\par

\noindent\textbf{成立条件：} 一定量の気体、外圧0、断熱、前後は平衡。外部仕事を行う断熱膨張と自由膨張を区別する。天秤の浮力は空気の密度と物体体積に依存する。\par

\paragraph{問われる核心} 外圧0の断熱膨張での保存量を確認し、気体が必ず冷えるかを検討する。

\paragraph{必要知識} 一定量の気体、外圧0、断熱、前後は平衡。外部仕事を行う断熱膨張と自由膨張を区別する。天秤の浮力は空気の密度と物体体積に依存する。

\paragraph{考え方・正答を選ぶ理由（査読待ち初稿）} 断熱で \(Q=0\)、外圧0で体積仕事 \(W=0\) なので \(\Delta U=0\) である。理想気体なら \(U=U(T)\) より温度は変わらない。実在気体では \((\partial T/\partial V)_U=[p-T(\partial p/\partial T)_V]/C_V\) の符号が状態方程式に依存する。必ず冷却するとする3は一般に成立しない。

\Needspace{4\baselineskip}\paragraph{選択肢の検討（査読待ち初稿）}\mbox{}\par

\noindent\textbf{1：適切} 三相が同時に平衡共存する点は三重点。一成分系の三相平衡を表す。相図上の一点であり臨界点とは異なる。\par

\noindent\textbf{2：条件付きで適切} ギブズエネルギーは定圧条件の自由エネルギー。圧力を自然変数とする \(G=U+pV-TS\) である。自発変化や極小の判定には温度も一定、閉じた系で非体積仕事がない等の条件が必要である。\par

\noindent\textbf{3：不適切} 真空への断熱膨張後は必ず冷却される。理想気体では温度不変が反例となる。実在気体でも内部エネルギー一定に沿う温度変化は状態方程式に依存する。\par

\noindent\textbf{4：適切} ファンデルワールス式は実在気体の体積と分子間相互作用を考慮。原典の実在気体の体積を、分子の有限サイズによる排除体積bの補正として説明する。引力はaの項で表す。\par

\noindent\textbf{5：適切} 天秤の測定では温湿度による空気密度の変化に留意。排除する空気質量による浮力が試料と分銅の見掛けの重量を変える。温湿度は空気密度を変える。\par

\paragraph{再利用する考え方} 断熱という条件だけでなく外圧と仕事を読む。自由膨張の内部エネルギー一定を、絞り流れのエンタルピー一定と混同しない。

\paragraph{共通理論} \hyperref[theory:phase_sep:free]{第\ref*{theory:phase_sep:free}節}, \hyperref[theory:phase_sep:gibbs]{第\ref*{theory:phase_sep:gibbs}節}, \hyperref[theory:phase_sep:realgas]{第\ref*{theory:phase_sep:realgas}節}, \hyperref[theory:phase_sep:buoyancy]{第\ref*{theory:phase_sep:buoyancy}節}

\par 参考：\cite{mit_thermo_firstlaw2008,iupac_greenbook2012}。

\par 難易度 2（校準初稿）、優先度未評価。

\subsection{令和5年度 問16}
\label{q:r05-reg-applied-16}

\noindent\textbf{原典：} sources/R05/17\_applied\_science.pdf、PDF物理6頁。

\textbf{公式正答：} 1（single）。

\par\textbf{検討未確定：} 平均密度・初期充填率が原典で指定されない。気相密度の単調上昇や任意容器の臨界到達を断定せず、臨界密度を持つ反例として公式1を説明する。

\paragraph{問題の要旨} 剛体密閉容器の平均密度を固定した加熱と、臨界点への到達条件を検討する。

\Needspace{4\baselineskip}\paragraph{確認した条件・記述（独自要約）}\mbox{}\par

\noindent\textbf{確認範囲：} 問題幹、条件、否定、数値、単位、全5肢を原典画像で照合した。\par

\noindent\textbf{成立条件：} 変形しない閉じた容器、加熱、平均密度一定。臨界密度322 kg/m3、臨界温度647.096 KはIAPWS-95の参照定数である。\par

\paragraph{問われる核心} 剛体密閉容器の平均密度を固定した加熱と、臨界点への到達条件を検討する。

\paragraph{必要知識} 変形しない閉じた容器、加熱、平均密度一定。臨界密度322 kg/m3、臨界温度647.096 KはIAPWS-95の参照定数である。

\paragraph{考え方・正答を選ぶ理由（査読待ち初稿）} 臨界点へ二相共存状態から近づくと液体と蒸気の密度差は0となる。剛体閉容器では平均密度 \(\bar\rho=M/V\) が一定であり、\(\bar\rho=\rho_c\) に選べば臨界点へ達する経路が可能である。したがって密度が一致することはないという1の全称は誤る。ただし任意の充填量で必ず臨界へ達するわけではなく、液相又は気相が先に消える場合もある。

\Needspace{4\baselineskip}\paragraph{選択肢の検討（査読待ち初稿）}\mbox{}\par

\noindent\textbf{1：不適切} 加熱した蒸気の密度は液体の密度に一致することがない。平均密度を臨界密度に設定した加熱が反例となる。任意密閉容器での一致を主張する必要はない。\par

\noindent\textbf{2：適切} ファンデルワールス式は分子間相互作用を考慮。分子間引力のa項を持つ。排除体積のb項も理想気体との差となる。\par

\noindent\textbf{3：条件付きで適切} ギブズエネルギーは定圧条件の自由エネルギー。圧力を自然変数とする \(G=U+pV-TS\) である。自発変化や極小の判定には温度も一定、閉じた系で非体積仕事がない等の条件が必要である。\par

\noindent\textbf{4：適切} 気液固が同時に平衡共存する点は三重点。一成分三相共存の定義である。\par

\noindent\textbf{5：適切} 気液平衡では化学ポテンシャルが等しい。相間の物質移動が全ギブズエネルギーを低下させない条件である。\par

\paragraph{再利用する考え方} 常に起こらないという全称は一つの可能な経路で反証できる。臨界点への加熱には温度だけでなく平均密度の条件が必要である。

\paragraph{共通理論} \hyperref[theory:phase_sep:critical]{第\ref*{theory:phase_sep:critical}節}, \hyperref[theory:phase_sep:isochore]{第\ref*{theory:phase_sep:isochore}節}, \hyperref[theory:phase_sep:gibbs]{第\ref*{theory:phase_sep:gibbs}節}

\par 参考：\cite{iapws95_2018}。

\par 難易度 2（校準初稿）、優先度未評価。

\section{反応商による化学平衡移動}

\subsection{平成28年度 問18}
\label{q:h28-reg-applied-18}

\noindent\textbf{原典：} sources/H28/17\_applied\_science.pdf、PDF物理7頁。

\textbf{公式正答：} 3（single）。

\paragraph{問題の要旨} 発熱するアンモニア合成の温度・圧力・窒素追加・触媒の効果を判定する。

\Needspace{4\baselineskip}\paragraph{確認した条件・記述（独自要約）}\mbox{}\par

\noindent\textbf{確認範囲：} 問題幹、否定、全5肢、数式、数値、単位を原典画像で確認した。\par

\noindent\textbf{成立条件：} \(\mathrm{N_2+3H_2\rightleftharpoons2NH_3}+92\,\mathrm{kJ}\)。窒素追加は温度・体積一定。気体の分圧を用いる理想気体近似。\par

\paragraph{問われる核心} 発熱するアンモニア合成の温度・圧力・窒素追加・触媒の効果を判定する。

\paragraph{必要知識} \(\mathrm{N_2+3H_2\rightleftharpoons2NH_3}+92\,\mathrm{kJ}\)。窒素追加は温度・体積一定。気体の分圧を用いる理想気体近似。

\paragraph{考え方・正答を選ぶ理由（査読待ち初稿）} 反応式は窒素1、水素3、アンモニア2であり、生成側に92 kJを持つ。窒素追加は定温定積で反応商を低下させる。低温は発熱方向、高圧は気体量の少ない生成方向を平衡上有利にする。触媒による平衡移動だけが誤りであり、3を選ぶ。

\Needspace{4\baselineskip}\paragraph{選択肢の検討（査読待ち初稿）}\mbox{}\par

\noindent\textbf{1：適切} 発熱反応である。式の生成側の92 kJは反応式1回分の放熱を表す。\par

\noindent\textbf{2：適切} 定温定積で窒素を加えると生成側へ進む。直後の \(Q=p_{\mathrm{NH_3}}^2/(p_{\mathrm{N_2}}p_{\mathrm{H_2}}^3)\) が小さくなる。\par

\noindent\textbf{3：不適切} 平衡後の触媒追加で生成側へ移る。触媒は平衡定数を変えず、正逆の反応を促進する。\par

\noindent\textbf{4：適切} 鉄触媒のハーバー・ボッシュ法で製造する。鉄系触媒による工業合成である。実用温度は反応速度も考慮して選ぶ。\par

\noindent\textbf{5：適切} 低温・高圧ほど生成に有利である。発熱方向は低温で平衡上有利となり、気体の化学量論数が4から2へ減る方向は圧縮で有利となる。\par

\paragraph{再利用する考え方} 反応速度と平衡位置、全圧と分圧、濃度と活量を分ける。式の適用条件を確認し、原典の不正確な記述を科学的に正しい文へ置き換えない。

\paragraph{共通理論} \hyperref[theory:chem_eq:quotient]{第\ref*{theory:chem_eq:quotient}節}, \hyperref[theory:chem_eq:ammonia]{第\ref*{theory:chem_eq:ammonia}節}, \hyperref[theory:chem_eq:catalyst]{第\ref*{theory:chem_eq:catalyst}節}

\par 参考：\cite{iupac_greenbook2012}。

\paragraph{関連問題} \hyperref[q:h30-reg-applied-16]{H30-REG-APPLIED-16}, \hyperref[q:r04-reg-applied-10]{R04-REG-APPLIED-10}, \hyperref[q:r07-reg-applied-13]{R07-REG-APPLIED-13}

\par 難易度 2（校準初稿）、優先度未評価。

\subsection{平成30年度 問16}
\label{q:h30-reg-applied-16}

\noindent\textbf{原典：} sources/H30/17\_applied\_science.pdf、PDF物理5頁。

\textbf{公式正答：} 4（single）。

\paragraph{問題の要旨} 発熱するアンモニア合成の温度・圧力・窒素追加・触媒の効果を判定する。

\Needspace{4\baselineskip}\paragraph{確認した条件・記述（独自要約）}\mbox{}\par

\noindent\textbf{確認範囲：} 問題幹、否定、全5肢、数式、数値、単位を原典画像で確認した。\par

\noindent\textbf{成立条件：} \(\mathrm{N_2+3H_2\rightleftharpoons2NH_3}+92\,\mathrm{kJ}\)。窒素追加は温度・体積一定。気体の分圧を用いる理想気体近似。\par

\paragraph{問われる核心} 発熱するアンモニア合成の温度・圧力・窒素追加・触媒の効果を判定する。

\paragraph{必要知識} \(\mathrm{N_2+3H_2\rightleftharpoons2NH_3}+92\,\mathrm{kJ}\)。窒素追加は温度・体積一定。気体の分圧を用いる理想気体近似。

\paragraph{考え方・正答を選ぶ理由（査読待ち初稿）} 反応式は窒素1、水素3、アンモニア2であり、生成側に92 kJを持つ。窒素追加は定温定積で反応商を低下させる。低温は発熱方向、高圧は気体量の少ない生成方向を平衡上有利にする。触媒による平衡移動だけが誤りであり、4を選ぶ。

\Needspace{4\baselineskip}\paragraph{選択肢の検討（査読待ち初稿）}\mbox{}\par

\noindent\textbf{1：適切} 発熱反応である。式の生成側の92 kJは反応式1回分の放熱を表す。\par

\noindent\textbf{2：適切} 定温定積で窒素を加えると生成側へ進む。直後の \(Q=p_{\mathrm{NH_3}}^2/(p_{\mathrm{N_2}}p_{\mathrm{H_2}}^3)\) が小さくなる。\par

\noindent\textbf{3：適切} 低温・高圧ほど生成に有利である。発熱方向は低温で平衡上有利となり、気体の化学量論数が4から2へ減る方向は圧縮で有利となる。\par

\noindent\textbf{4：不適切} 平衡後の触媒追加で生成側へ移る。触媒は平衡定数を変えず、正逆の反応を促進する。\par

\noindent\textbf{5：適切} 鉄触媒のハーバー・ボッシュ法で製造する。鉄系触媒による工業合成である。実用温度は反応速度も考慮して選ぶ。\par

\paragraph{再利用する考え方} 反応速度と平衡位置、全圧と分圧、濃度と活量を分ける。式の適用条件を確認し、原典の不正確な記述を科学的に正しい文へ置き換えない。

\paragraph{共通理論} \hyperref[theory:chem_eq:quotient]{第\ref*{theory:chem_eq:quotient}節}, \hyperref[theory:chem_eq:ammonia]{第\ref*{theory:chem_eq:ammonia}節}, \hyperref[theory:chem_eq:catalyst]{第\ref*{theory:chem_eq:catalyst}節}

\par 参考：\cite{iupac_greenbook2012}。

\paragraph{関連問題} \hyperref[q:h28-reg-applied-18]{H28-REG-APPLIED-18}

\par 難易度 2（校準初稿）、優先度未評価。

\subsection{令和4年度 問10}
\label{q:r04-reg-applied-10}

\noindent\textbf{原典：} sources/R04/17\_applied\_science.pdf、PDF物理3頁。

\textbf{公式正答：} 3（single）。

\paragraph{問題の要旨} 発熱するアンモニア合成の温度・圧力・窒素追加・触媒の効果を判定する。

\Needspace{4\baselineskip}\paragraph{確認した条件・記述（独自要約）}\mbox{}\par

\noindent\textbf{確認範囲：} 問題幹、否定、全5肢、数式、数値、単位を原典画像で確認した。\par

\noindent\textbf{成立条件：} \(\mathrm{N_2+3H_2\rightleftharpoons2NH_3}+92\,\mathrm{kJ}\)。窒素追加は温度・体積一定。気体の分圧を用いる理想気体近似。\par

\paragraph{問われる核心} 発熱するアンモニア合成の温度・圧力・窒素追加・触媒の効果を判定する。

\paragraph{必要知識} \(\mathrm{N_2+3H_2\rightleftharpoons2NH_3}+92\,\mathrm{kJ}\)。窒素追加は温度・体積一定。気体の分圧を用いる理想気体近似。

\paragraph{考え方・正答を選ぶ理由（査読待ち初稿）} 反応式は窒素1、水素3、アンモニア2であり、生成側に92 kJを持つ。窒素追加は定温定積で反応商を低下させる。低温は発熱方向、高圧は気体量の少ない生成方向を平衡上有利にする。触媒による平衡移動だけが誤りであり、3を選ぶ。

\Needspace{4\baselineskip}\paragraph{選択肢の検討（査読待ち初稿）}\mbox{}\par

\noindent\textbf{1：適切} 発熱反応である。式の生成側の92 kJは反応式1回分の放熱を表す。\par

\noindent\textbf{2：適切} 定温定積で窒素を加えると生成側へ進む。直後の \(Q=p_{\mathrm{NH_3}}^2/(p_{\mathrm{N_2}}p_{\mathrm{H_2}}^3)\) が小さくなる。\par

\noindent\textbf{3：不適切} 平衡後の触媒追加で生成側へ移る。触媒は平衡定数を変えず、正逆の反応を促進する。\par

\noindent\textbf{4：適切} 鉄触媒のハーバー・ボッシュ法で製造する。鉄系触媒による工業合成である。実用温度は反応速度も考慮して選ぶ。\par

\noindent\textbf{5：適切} 低温・高圧ほど生成に有利である。発熱方向は低温で平衡上有利となり、気体の化学量論数が4から2へ減る方向は圧縮で有利となる。\par

\paragraph{再利用する考え方} 反応速度と平衡位置、全圧と分圧、濃度と活量を分ける。式の適用条件を確認し、原典の不正確な記述を科学的に正しい文へ置き換えない。

\paragraph{共通理論} \hyperref[theory:chem_eq:quotient]{第\ref*{theory:chem_eq:quotient}節}, \hyperref[theory:chem_eq:ammonia]{第\ref*{theory:chem_eq:ammonia}節}, \hyperref[theory:chem_eq:catalyst]{第\ref*{theory:chem_eq:catalyst}節}

\par 参考：\cite{iupac_greenbook2012}。

\paragraph{関連問題} \hyperref[q:h28-reg-applied-18]{H28-REG-APPLIED-18}

\par 難易度 2（校準初稿）、優先度未評価。

\subsection{令和7年度 問13}
\label{q:r07-reg-applied-13}

\noindent\textbf{原典：} sources/R07/17\_applied\_science.pdf、PDF物理5頁。

\textbf{公式正答：} 3（single）。

\paragraph{問題の要旨} 四酸化二窒素の解離平衡で、反応種の分圧と不活性気体による全圧変化を区別する。

\Needspace{4\baselineskip}\paragraph{確認した条件・記述（独自要約）}\mbox{}\par

\noindent\textbf{確認範囲：} 問題幹、否定、全5肢、数式、数値、単位を原典画像で確認した。\par

\noindent\textbf{成立条件：} 1、2、3、5は温度・体積一定、4は温度一定で体積増加。理想気体近似。不活性な窒素と固体触媒は反応商に含めない。\par

\paragraph{問われる核心} 四酸化二窒素の解離平衡で、反応種の分圧と不活性気体による全圧変化を区別する。

\paragraph{必要知識} 1、2、3、5は温度・体積一定、4は温度一定で体積増加。理想気体近似。不活性な窒素と固体触媒は反応商に含めない。

\paragraph{考え方・正答を選ぶ理由（査読待ち初稿）} \(\mathrm{N_2O_4\rightleftharpoons2NO_2}\) の \(Q=p_{\mathrm{NO_2}}^2/p_{\mathrm{N_2O_4}}\) を各操作直後に比較する。一定温度・体積で窒素を加えても反応種の \(p_i=n_iRT/V\) は変わらず、Qも変わらない。全圧が上がるだけで左へ進むとする3が誤りである。体積を増やす場合は全反応種の分圧が下がり、Qは体積に反比例して下がるため右へ進む。

\Needspace{4\baselineskip}\paragraph{選択肢の検討（査読待ち初稿）}\mbox{}\par

\noindent\textbf{1：適切} 定温定積で四酸化二窒素を加えると右へ。反応商の分母が増えてQが下がる。\par

\noindent\textbf{2：適切} 定温定積で二酸化窒素を除くと右へ。反応商の分子が減ってQが下がる。\par

\noindent\textbf{3：不適切} 窒素を加えて圧力を上げると左へ。一定体積では反応種の分圧が変わらず、圧縮による圧力上昇とは異なる。\par

\noindent\textbf{4：適切} 一定温度で体積を増すと右へ。反応商が下がり気体の化学量論数の多い側へ進む。\par

\noindent\textbf{5：適切} 定温定積で固体触媒を加えても平衡は変わらない。触媒はKを変えず、反応商の因子にもならない。\par

\paragraph{再利用する考え方} 反応速度と平衡位置、全圧と分圧、濃度と活量を分ける。式の適用条件を確認し、原典の不正確な記述を科学的に正しい文へ置き換えない。

\paragraph{共通理論} \hyperref[theory:chem_eq:quotient]{第\ref*{theory:chem_eq:quotient}節}, \hyperref[theory:chem_eq:inert]{第\ref*{theory:chem_eq:inert}節}

\par 参考：\cite{iupac_greenbook2012}。

\paragraph{関連問題} \hyperref[q:h28-reg-applied-18]{H28-REG-APPLIED-18}

\par 難易度 2（校準初稿）、優先度未評価。

\section{反応ギブズエネルギー・活量・平衡定数}

\subsection{令和3年度 問10}
\label{q:r03-reg-applied-10}

\noindent\textbf{原典：} sources/R03/17\_applied\_science.pdf、PDF物理4頁。

\textbf{公式正答：} 4（single）。

\paragraph{問題の要旨} 一定温度・圧力のAとBの可逆反応で、実際の反応ギブズエネルギーと標準量を区別する。

\Needspace{4\baselineskip}\paragraph{確認した条件・記述（独自要約）}\mbox{}\par

\noindent\textbf{確認範囲：} 問題幹、否定、全5肢、数式、数値、単位を原典画像で確認した。\par

\noindent\textbf{成立条件：} 定温定圧の \(\mathrm{A\rightleftharpoons B}\)。標準生成ギブズエネルギーの差と実際の組成による反応量を区別する。濃度比で活量比を近似する条件を明記する。\par

\paragraph{問われる核心} 一定温度・圧力のAとBの可逆反応で、実際の反応ギブズエネルギーと標準量を区別する。

\paragraph{必要知識} 定温定圧の \(\mathrm{A\rightleftharpoons B}\)。標準生成ギブズエネルギーの差と実際の組成による反応量を区別する。濃度比で活量比を近似する条件を明記する。

\paragraph{考え方・正答を選ぶ理由（査読待ち初稿）} \(\Delta_rG=\Delta_rG^\circ+RT\ln Q\) であり、AとBの1対1反応では \(Q=a_B/a_A\) である。一定温度でも組成が変われば実際の反応ギブズエネルギーは変わる。平衡の \(\Delta_rG=0\) から \(K=\exp(-\Delta_rG^\circ/RT)\) が得られる。濃度比変化を無視する4が不適切である。

\Needspace{4\baselineskip}\paragraph{選択肢の検討（査読待ち初稿）}\mbox{}\par

\noindent\textbf{1：適切} 反応ギブズエネルギー負なら正方向へ進む。定温定圧で正方向の反応進行が系のギブズエネルギーを低下させる。\par

\noindent\textbf{2：適切} 平衡では反応ギブズエネルギーが0。反応進行度に関するギブズエネルギーの微分が0となる。\par

\noindent\textbf{3：適切} 標準生成量から反応標準量を求める。生成物の標準生成ギブズエネルギーの和から反応物の和を化学量論係数付きで引く。\par

\noindent\textbf{4：不適切} 一定温度なら濃度比が変わっても反応量は一定。標準反応量と異なり、実際の反応量は反応商に依存する。\par

\noindent\textbf{5：適切} 平衡定数は標準生成ギブズエネルギーから求まる。標準反応量を求めて指数式に代入する。\par

\paragraph{再利用する考え方} 反応速度と平衡位置、全圧と分圧、濃度と活量を分ける。式の適用条件を確認し、原典の不正確な記述を科学的に正しい文へ置き換えない。

\paragraph{共通理論} \hyperref[theory:chem_eq:gibbs]{第\ref*{theory:chem_eq:gibbs}節}, \hyperref[theory:chem_eq:quotient]{第\ref*{theory:chem_eq:quotient}節}

\par 参考：\cite{iupac_greenbook2012}。

\paragraph{関連問題} \hyperref[q:r05-reg-applied-17]{R05-REG-APPLIED-17}

\par 難易度 2（校準初稿）、優先度未評価。

\subsection{令和5年度 問17}
\label{q:r05-reg-applied-17}

\noindent\textbf{原典：} sources/R05/17\_applied\_science.pdf、PDF物理6頁。

\textbf{公式正答：} 1（single）。

\paragraph{問題の要旨} 一定温度・圧力のAとBの可逆反応で、実際の反応ギブズエネルギーと標準量を区別する。

\Needspace{4\baselineskip}\paragraph{確認した条件・記述（独自要約）}\mbox{}\par

\noindent\textbf{確認範囲：} 問題幹、否定、全5肢、数式、数値、単位を原典画像で確認した。\par

\noindent\textbf{成立条件：} 定温定圧の \(\mathrm{A\rightleftharpoons B}\)。標準生成ギブズエネルギーの差と実際の組成による反応量を区別する。濃度比で活量比を近似する条件を明記する。\par

\paragraph{問われる核心} 一定温度・圧力のAとBの可逆反応で、実際の反応ギブズエネルギーと標準量を区別する。

\paragraph{必要知識} 定温定圧の \(\mathrm{A\rightleftharpoons B}\)。標準生成ギブズエネルギーの差と実際の組成による反応量を区別する。濃度比で活量比を近似する条件を明記する。

\paragraph{考え方・正答を選ぶ理由（査読待ち初稿）} \(\Delta_rG=\Delta_rG^\circ+RT\ln Q\) であり、AとBの1対1反応では \(Q=a_B/a_A\) である。一定温度でも組成が変われば実際の反応ギブズエネルギーは変わる。平衡の \(\Delta_rG=0\) から \(K=\exp(-\Delta_rG^\circ/RT)\) が得られる。濃度比変化を無視する1が不適切である。

\Needspace{4\baselineskip}\paragraph{選択肢の検討（査読待ち初稿）}\mbox{}\par

\noindent\textbf{1：不適切} 一定温度なら濃度比が変わっても反応量は一定。標準反応量と異なり、実際の反応量は反応商に依存する。\par

\noindent\textbf{2：適切} 反応ギブズエネルギー負なら正方向へ進む。定温定圧で正方向の反応進行が系のギブズエネルギーを低下させる。\par

\noindent\textbf{3：適切} 平衡定数は標準生成ギブズエネルギーから求まる。標準反応量を求めて指数式に代入する。\par

\noindent\textbf{4：適切} 標準生成量から反応標準量を求める。生成物の標準生成ギブズエネルギーの和から反応物の和を化学量論係数付きで引く。\par

\noindent\textbf{5：適切} 平衡では反応ギブズエネルギーが0。反応進行度に関するギブズエネルギーの微分が0となる。\par

\paragraph{再利用する考え方} 反応速度と平衡位置、全圧と分圧、濃度と活量を分ける。式の適用条件を確認し、原典の不正確な記述を科学的に正しい文へ置き換えない。

\paragraph{共通理論} \hyperref[theory:chem_eq:gibbs]{第\ref*{theory:chem_eq:gibbs}節}, \hyperref[theory:chem_eq:quotient]{第\ref*{theory:chem_eq:quotient}節}

\par 参考：\cite{iupac_greenbook2012}。

\paragraph{関連問題} \hyperref[q:r03-reg-applied-10]{R03-REG-APPLIED-10}

\par 難易度 2（校準初稿）、優先度未評価。

\section{酸化数・電子収支}

\subsection{平成25年度 問11}
\label{q:h25-reg-applied-11}

\noindent\textbf{原典：} sources/H25/17\_applied\_science.pdf、PDF物理5頁。

\textbf{公式正答：} 1（single）。

\paragraph{問題の要旨} 反応前後の酸化数を求め、還元に当たる変化を選ぶ。

\Needspace{4\baselineskip}\paragraph{確認した条件・記述（独自要約）}\mbox{}\par

\noindent\textbf{確認範囲：} 問題幹、数値、式、否定、全5候補を原典画像で照合した。\par

\noindent\textbf{成立条件：} 単体の酸化数は0、単原子イオンは電荷に等しい。この問題の通常の酸化物でOは−2、H2SでHは+1である。式は変化の概略であって全反応の釣合い式ではない。\par

\paragraph{問われる核心} 反応前後の酸化数を求め、還元に当たる変化を選ぶ。

\paragraph{必要知識} 単体の酸化数は0、単原子イオンは電荷に等しい。この問題の通常の酸化物でOは−2、H2SでHは+1である。式は変化の概略であって全反応の釣合い式ではない。

\paragraph{考え方・正答を選ぶ理由（査読待ち初稿）} \(\mathrm{MnO_4^-}\) ではMnの酸化数が \(x+4(-2)=-1\) より+7である。\(\mathrm{Mn^{2+}}\) では+2なので5電子を受け取る還元であり1を選ぶ。他の変化は酸化数不変又は増加である。

\Needspace{4\baselineskip}\paragraph{選択肢の検討（査読待ち初稿）}\mbox{}\par

\noindent\textbf{1：適切} \(\mathrm{MnO_4^-}\to\mathrm{Mn^{2+}}\)。Mnは+7から+2へ低下し還元である。\par

\noindent\textbf{2：不適切} NaCl→NaOH。Naは両側とも+1であり、Naの還元ではない。ClはCl−として別に移るなど全反応の収支が省略された表示である。\par

\noindent\textbf{3：不適切} \(2\mathrm{I^-}\to\mathrm{I_2}\)。Iは−1から0へ増加し、2電子を放出する酸化である。\par

\noindent\textbf{4：不適切} Mg→MgCl2。Mgは0から+2へ増加し酸化である。\par

\noindent\textbf{5：不適切} H2S→S。Sは−2から0へ増加し酸化である。\par

\paragraph{再利用する考え方} 酸素の付加だけで判定せず、対象元素の酸化数と電子の受け渡しを追う。酸化数は形式的な電荷であり、共有結合内の実電荷そのものではない。

\paragraph{共通理論} \hyperref[theory:electrochem:oxidation_number]{第\ref*{theory:electrochem:oxidation_number}節}, \hyperref[theory:electrochem:half_reaction]{第\ref*{theory:electrochem:half_reaction}節}

\par 参考：\cite{mit_redox2014}。

\par 難易度 2（校準初稿）、優先度未評価。

\section{中和後の塩の加水分解とpH}

\subsection{平成24年度 問10}
\label{q:h24-reg-applied-10}

\noindent\textbf{原典：} sources/H24/17\_applied\_science.pdf、PDF物理3頁。

\textbf{公式正答：} 5（single）。

\paragraph{問題の要旨} 同濃度・同体積で中和した3組の酸塩基について、生成塩の加水分解からpHの順序を求める。

\Needspace{4\baselineskip}\paragraph{確認した条件・記述（独自要約）}\mbox{}\par

\noindent\textbf{確認範囲：} 問題幹、否定、全5肢、数式、数値、単位を原典画像で確認した。\par

\noindent\textbf{成立条件：} 各水溶液は同じモル濃度、同じ体積。1対1の中和を先に処理する。強酸強塩基塩は水の自己解離に対して中性、弱酸塩は塩基性、弱塩基塩は酸性。\par

\paragraph{問われる核心} 同濃度・同体積で中和した3組の酸塩基について、生成塩の加水分解からpHの順序を求める。

\paragraph{必要知識} 各水溶液は同じモル濃度、同じ体積。1対1の中和を先に処理する。強酸強塩基塩は水の自己解離に対して中性、弱酸塩は塩基性、弱塩基塩は酸性。

\paragraph{考え方・正答を選ぶ理由（査読待ち初稿）} Aの酢酸と水酸化ナトリウムから生じる酢酸イオンは加水分解で水酸化物イオンを生じ、塩基性となる。Bの塩酸と水酸化ナトリウムは中性の塩、Cの塩酸とアンモニアはアンモニウムイオンを含む酸性の塩である。低いpHからC、B、Aとなり5を選ぶ。

\Needspace{4\baselineskip}\paragraph{選択肢の検討（査読待ち初稿）}\mbox{}\par

\noindent\textbf{1：不適切} A・B・C。酸性C、中性B、塩基性Aの順に一致しない。\par

\noindent\textbf{2：不適切} B・C・A。酸性C、中性B、塩基性Aの順に一致しない。\par

\noindent\textbf{3：不適切} B・A・C。酸性C、中性B、塩基性Aの順に一致しない。\par

\noindent\textbf{4：不適切} C・A・B。酸性C、中性B、塩基性Aの順に一致しない。\par

\noindent\textbf{5：適切} C・B・A。酸性C、中性B、塩基性Aの順に一致する。\par

\paragraph{再利用する考え方} 反応速度と平衡位置、全圧と分圧、濃度と活量を分ける。式の適用条件を確認し、原典の不正確な記述を科学的に正しい文へ置き換えない。

\paragraph{共通理論} \hyperref[theory:chem_eq:hydrolysis]{第\ref*{theory:chem_eq:hydrolysis}節}

\par 参考：\cite{iupac_greenbook2012,mit_buffer2008}。

\paragraph{関連問題} \hyperref[q:h30-reg-applied-17]{H30-REG-APPLIED-17}

\par 難易度 2（校準初稿）、優先度未評価。

\subsection{平成30年度 問17}
\label{q:h30-reg-applied-17}

\noindent\textbf{原典：} sources/H30/17\_applied\_science.pdf、PDF物理6頁。

\textbf{公式正答：} 3（single）。

\par\textbf{検討未確定：} 形式濃度1 mol/dm3の塩溶液に活量係数1を用いる計算上の近似を明示する。原典は温度・活量係数・実在溶液の調製可能性を指定しないため、実測pHの断定とは区別する。

\paragraph{問題の要旨} 安息香酸と等モルの強塩基を混合し、生成塩の加水分解でpHを計算する。

\Needspace{4\baselineskip}\paragraph{確認した条件・記述（独自要約）}\mbox{}\par

\noindent\textbf{確認範囲：} 問題幹、否定、全5肢、数式、数値、単位を原典画像で確認した。\par

\noindent\textbf{成立条件：} \(K_a=10^{-4}\,\mathrm{mol\,dm^{-3}}\)、\(K_w=10^{-14}\,\mathrm{mol^2\,dm^{-6}}\)。体積加成、濃度で平衡式を扱う理想溶液近似、加水分解量が1 mol/dm3に比べ十分小さいこと。\par

\paragraph{問われる核心} 安息香酸と等モルの強塩基を混合し、生成塩の加水分解でpHを計算する。

\paragraph{必要知識} \(K_a=10^{-4}\,\mathrm{mol\,dm^{-3}}\)、\(K_w=10^{-14}\,\mathrm{mol^2\,dm^{-6}}\)。体積加成、濃度で平衡式を扱う理想溶液近似、加水分解量が1 mol/dm3に比べ十分小さいこと。

\paragraph{考え方・正答を選ぶ理由（査読待ち初稿）} 両溶液は \(2\,\mathrm{mol\,dm^{-3}}\)、各 \(10\,\mathrm{cm^3}=0.010\,\mathrm{dm^3}\) である。各物質量は0.020 molなので過剰の強塩基は残らない。混合体積0.020 dm3より安息香酸塩の形式濃度は1 mol/dm3である。\(K_b=K_w/K_a=10^{-10}\,\mathrm{mol\,dm^{-3}}\) とし、加水分解量 \(x\) は \(x^2/(1-x)=10^{-10}\) から約 \(10^{-5}\,\mathrm{mol\,dm^{-3}}\) となる。与えた \(K_w=10^{-14}\) に対応してpOH=\allowbreak{}5、pH=\allowbreak{}9であり3を選ぶ。

\Needspace{4\baselineskip}\paragraph{選択肢の検討（査読待ち初稿）}\mbox{}\par

\noindent\textbf{1：不適切} 7。等モル中和は過剰酸塩基を消すが、弱酸の塩の加水分解を消さない。\par

\noindent\textbf{2：不適切} 8。加水分解の平方根から得たpOH=\allowbreak{}5と対応しない。\par

\noindent\textbf{3：適切} 9。水酸化物イオン濃度が約10のマイナス5乗なので対応する。\par

\noindent\textbf{4：不適切} 10。形式濃度1とKbの積の平方根を取った結果と対応しない。\par

\noindent\textbf{5：不適切} 14。未反応の1 mol/dm3強塩基が残る場合ではない。\par

\paragraph{再利用する考え方} 反応速度と平衡位置、全圧と分圧、濃度と活量を分ける。式の適用条件を確認し、原典の不正確な記述を科学的に正しい文へ置き換えない。

\paragraph{共通理論} \hyperref[theory:chem_eq:hydrolysis]{第\ref*{theory:chem_eq:hydrolysis}節}

\par 参考：\cite{iupac_greenbook2012,mit_buffer2008}。

\paragraph{関連問題} \hyperref[q:h24-reg-applied-10]{H24-REG-APPLIED-10}

\par 難易度 3（校準初稿）、優先度未評価。

\section{溶解度積・pH・共通イオン}

\subsection{平成29年度 問13}
\label{q:h29-reg-applied-13}

\noindent\textbf{原典：} sources/H29/17\_applied\_science.pdf、PDF物理4頁。

\textbf{公式正答：} 3（single）。

\paragraph{問題の要旨} 一次可逆反応、反応性衝突、溶解度積、連鎖反応、熱暴走を比較する。

\Needspace{4\baselineskip}\paragraph{確認した条件・記述（独自要約）}\mbox{}\par

\noindent\textbf{確認範囲：} 問題幹、否定、全5肢、数式、数値、単位を原典画像で確認した。\par

\noindent\textbf{成立条件：} 1のAとBの可逆反応は双方一次。2はClとNOClの衝突。3は \(\mathrm{Mg(OH)_2}\)。4は連鎖開始・伝播・停止。5は発熱と温度上昇の正帰還。\par

\paragraph{問われる核心} 一次可逆反応、反応性衝突、溶解度積、連鎖反応、熱暴走を比較する。

\paragraph{必要知識} 1のAとBの可逆反応は双方一次。2はClとNOClの衝突。3は \(\mathrm{Mg(OH)_2}\)。4は連鎖開始・伝播・停止。5は発熱と温度上昇の正帰還。

\paragraph{考え方・正答を選ぶ理由（査読待ち初稿）} 水酸化マグネシウムの飽和溶液では \(K_{\mathrm{sp}}=a_{\mathrm{Mg^{2+}}}a_{\mathrm{OH^-}}^2\) である。一定温度でもpHにより水酸化物イオンの活量が変わるため、溶解度は一般に変化する。pH依存性を否定する3が不適切である。

\Needspace{4\baselineskip}\paragraph{選択肢の検討（査読待ち初稿）}\mbox{}\par

\noindent\textbf{1：適切} 一次正逆反応の速度定数比が平衡定数。平衡で \(k_f[A]=k_r[B]\) なので濃度基準の \(K=[B]/[A]=k_f/k_r\) となる。\par

\noindent\textbf{2：適切} 十分なエネルギーの衝突でも必ず反応するとは限らない。分子の向きや反応経路にも依存する。\par

\noindent\textbf{3：不適切} 水酸化マグネシウムの溶解度はpHに依存しない。共通イオン効果を通じて水酸化物イオンの活量が溶解量を変える。\par

\noindent\textbf{4：適切} 連鎖反応は開始・伝播・停止を含む。反応性中間体の生成、再生、消滅を区別する。\par

\noindent\textbf{5：適切} 発熱が速度を上げて熱爆発を起こし得る。発熱と除熱の釣合いが失われると温度上昇と反応加速が循環する。\par

\paragraph{再利用する考え方} 反応速度と平衡位置、全圧と分圧、濃度と活量を分ける。式の適用条件を確認し、原典の不正確な記述を科学的に正しい文へ置き換えない。

\paragraph{共通理論} \hyperref[theory:chem_eq:solubility]{第\ref*{theory:chem_eq:solubility}節}, \hyperref[theory:chem_eq:rates]{第\ref*{theory:chem_eq:rates}節}

\par 参考：\cite{iupac_greenbook2012}。

\par 難易度 2（校準初稿）、優先度未評価。

\section{飽和蒸気圧・沸騰・密封条件}

\subsection{平成24年度 問13}
\label{q:h24-reg-applied-13}

\noindent\textbf{原典：} sources/H24/17\_applied\_science.pdf、PDF物理4頁。

\textbf{公式正答：} 3（single）。

\par\textbf{検討未確定：} 原典は純物質という限定、膨張後も液体が残る条件、密封加熱時の空気・気相部を明示しない。標準的な想定での公式判定と、任意容器への全称を区別する。H24の肢2には平衡状態の明記もなく、過渡状態を排除できない。

\paragraph{問題の要旨} 飽和蒸気圧、動的平衡、定圧沸騰と密封加熱を条件ごとに判定する。

\Needspace{4\baselineskip}\paragraph{確認した条件・記述（独自要約）}\mbox{}\par

\noindent\textbf{確認範囲：} 問題幹、条件、否定、数値、単位、全5肢を原典画像で照合した。\par

\noindent\textbf{成立条件：} 定温、気液共存と平衡、定圧沸騰、室温大気圧での密封、不揮発性溶質という条件。蒸気の分圧と容器の全圧を区別する。\par

\paragraph{問われる核心} 飽和蒸気圧、動的平衡、定圧沸騰と密封加熱を条件ごとに判定する。

\paragraph{必要知識} 定温、気液共存と平衡、定圧沸騰、室温大気圧での密封、不揮発性溶質という条件。蒸気の分圧と容器の全圧を区別する。

\paragraph{考え方・正答を選ぶ理由（査読待ち初稿）} 純液体の定圧沸騰では、二相が残る間の温度は一定である。上昇し続けるとする3を選ぶ。密封容器の全圧と水の分圧を分け、定温膨張は液体が残る場合に限って飽和蒸気圧へ戻ると読む。

\Needspace{4\baselineskip}\paragraph{選択肢の検討（査読待ち初稿）}\mbox{}\par

\noindent\textbf{1：条件付きで適切} 定温で二相共存の密封容器を広げても蒸気圧は変わらない。膨張後も液体が残り平衡に達すると、追加蒸発により同じ温度の飽和蒸気圧に戻る。全液体が蒸発した後は成立しない。\par

\noindent\textbf{2：条件付きで適切} 密封容器に液体と気体が共存すると蒸発速度と凝縮速度が等しい。動的平衡に達した場合に等しい。原典はこの肢に平衡状態の語がなく、共存だけでは正味蒸発が起こる過渡状態も含まれる。\par

\noindent\textbf{3：不適切} 大気圧の沸騰中も温度は上昇し続ける。純液体の定圧気液共存では加熱が蒸発潜熱へ使われ、温度は沸点に保たれる。\par

\noindent\textbf{4：条件付きで適切} 室温・大気圧で水を入れ密封し加熱すると100°Cで沸騰しない。気相部に空気が残る定積容器なら、100°Cで水蒸気分圧が約1気圧でも全圧には空気分圧が加わる。沸騰の判定は全外圧による。空気を除いた容器などへ無条件に一般化しない。\par

\noindent\textbf{5：適切} 不揮発性溶質を溶かした希薄溶液の蒸気圧は下がる。同じ温度で溶媒の活量が低下し、理想希薄近似では \(p_A=x_Ap_A^*\) となる。溶質は蒸気圧に寄与しない。\par

\paragraph{再利用する考え方} 圧力を一定に保つ開放容器と、物質量・体積を固定する密封容器は異なる。液体が残る条件を飽和蒸気圧の議論に付ける。

\paragraph{共通理論} \hyperref[theory:phase_sep:vapor]{第\ref*{theory:phase_sep:vapor}節}, \hyperref[theory:phase_sep:boiling]{第\ref*{theory:phase_sep:boiling}節}, \hyperref[theory:phase_sep:sealed]{第\ref*{theory:phase_sep:sealed}節}, \hyperref[theory:phase_sep:raoult]{第\ref*{theory:phase_sep:raoult}節}

\par 参考：\cite{mit_solutions2008,mit_clapeyron2008}。

\paragraph{関連問題} \hyperref[q:h27-reg-applied-17]{H27-REG-APPLIED-17}

\par 難易度 2（校準初稿）、優先度未評価。

\subsection{平成27年度 問17}
\label{q:h27-reg-applied-17}

\noindent\textbf{原典：} sources/H27/17\_applied\_science.pdf、PDF物理6頁。

\textbf{公式正答：} 3（single）。

\par\textbf{検討未確定：} 原典は純物質という限定、膨張後も液体が残る条件、密封加熱時の空気・気相部を明示しない。標準的な想定での公式判定と、任意容器への全称を区別する。

\paragraph{問題の要旨} 飽和蒸気圧、動的平衡、定圧沸騰と密封加熱を条件ごとに判定する。

\Needspace{4\baselineskip}\paragraph{確認した条件・記述（独自要約）}\mbox{}\par

\noindent\textbf{確認範囲：} 問題幹、条件、否定、数値、単位、全5肢を原典画像で照合した。\par

\noindent\textbf{成立条件：} 定温、気液共存と平衡、定圧沸騰、室温大気圧での密封、不揮発性溶質という条件。蒸気の分圧と容器の全圧を区別する。\par

\paragraph{問われる核心} 飽和蒸気圧、動的平衡、定圧沸騰と密封加熱を条件ごとに判定する。

\paragraph{必要知識} 定温、気液共存と平衡、定圧沸騰、室温大気圧での密封、不揮発性溶質という条件。蒸気の分圧と容器の全圧を区別する。

\paragraph{考え方・正答を選ぶ理由（査読待ち初稿）} 純液体の定圧沸騰では、二相が残る間の温度は一定である。上昇し続けるとする3を選ぶ。密封容器の全圧と水の分圧を分け、定温膨張は液体が残る場合に限って飽和蒸気圧へ戻ると読む。

\Needspace{4\baselineskip}\paragraph{選択肢の検討（査読待ち初稿）}\mbox{}\par

\noindent\textbf{1：条件付きで適切} 定温で二相共存の密封容器を広げても蒸気圧は変わらない。膨張後も液体が残り平衡に達すると、追加蒸発により同じ温度の飽和蒸気圧に戻る。全液体が蒸発した後は成立しない。\par

\noindent\textbf{2：適切} 平衡の液体と気体では蒸発速度と凝縮速度が等しい。界面を往復する移動が続いても、物質量の正味変化は0となる。平衡を移動停止と混同しない。\par

\noindent\textbf{3：不適切} 大気圧の沸騰中も温度は上昇し続ける。純液体の定圧気液共存では加熱が蒸発潜熱へ使われ、温度は沸点に保たれる。\par

\noindent\textbf{4：条件付きで適切} 室温・大気圧で水を入れ密封し加熱すると100°Cで沸騰しない。気相部に空気が残る定積容器なら、100°Cで水蒸気分圧が約1気圧でも全圧には空気分圧が加わる。沸騰の判定は全外圧による。空気を除いた容器などへ無条件に一般化しない。\par

\noindent\textbf{5：適切} 不揮発性溶質を溶かした希薄溶液の蒸気圧は下がる。同じ温度で溶媒の活量が低下し、理想希薄近似では \(p_A=x_Ap_A^*\) となる。溶質は蒸気圧に寄与しない。\par

\paragraph{再利用する考え方} 圧力を一定に保つ開放容器と、物質量・体積を固定する密封容器は異なる。液体が残る条件を飽和蒸気圧の議論に付ける。

\paragraph{共通理論} \hyperref[theory:phase_sep:vapor]{第\ref*{theory:phase_sep:vapor}節}, \hyperref[theory:phase_sep:boiling]{第\ref*{theory:phase_sep:boiling}節}, \hyperref[theory:phase_sep:sealed]{第\ref*{theory:phase_sep:sealed}節}, \hyperref[theory:phase_sep:raoult]{第\ref*{theory:phase_sep:raoult}節}

\par 参考：\cite{mit_solutions2008,mit_clapeyron2008}。

\paragraph{関連問題} \hyperref[q:h24-reg-applied-13]{H24-REG-APPLIED-13}, \hyperref[q:h29-reg-applied-14]{H29-REG-APPLIED-14}, \hyperref[q:r01-reg-applied-10]{R01-REG-APPLIED-10}

\par 難易度 2（校準初稿）、優先度未評価。

\subsection{平成29年度 問14}
\label{q:h29-reg-applied-14}

\noindent\textbf{原典：} sources/H29/17\_applied\_science.pdf、PDF物理5頁。

\textbf{公式正答：} 3（single）。

\par\textbf{検討未確定：} 原典は純物質という限定、膨張後も液体が残る条件、密封加熱時の空気・気相部を明示しない。標準的な想定での公式判定と、任意容器への全称を区別する。

\paragraph{問題の要旨} 飽和蒸気圧、動的平衡、定圧沸騰と密封加熱を条件ごとに判定する。

\Needspace{4\baselineskip}\paragraph{確認した条件・記述（独自要約）}\mbox{}\par

\noindent\textbf{確認範囲：} 問題幹、条件、否定、数値、単位、全5肢を原典画像で照合した。\par

\noindent\textbf{成立条件：} 定温、気液共存と平衡、定圧沸騰、室温大気圧での密封、不揮発性溶質という条件。蒸気の分圧と容器の全圧を区別する。\par

\paragraph{問われる核心} 飽和蒸気圧、動的平衡、定圧沸騰と密封加熱を条件ごとに判定する。

\paragraph{必要知識} 定温、気液共存と平衡、定圧沸騰、室温大気圧での密封、不揮発性溶質という条件。蒸気の分圧と容器の全圧を区別する。

\paragraph{考え方・正答を選ぶ理由（査読待ち初稿）} 純液体の定圧沸騰では、二相が残る間の温度は一定である。上昇し続けるとする3を選ぶ。密封容器の全圧と水の分圧を分け、定温膨張は液体が残る場合に限って飽和蒸気圧へ戻ると読む。

\Needspace{4\baselineskip}\paragraph{選択肢の検討（査読待ち初稿）}\mbox{}\par

\noindent\textbf{1：条件付きで適切} 定温で二相共存の密封容器を広げても蒸気圧は変わらない。膨張後も液体が残り平衡に達すると、追加蒸発により同じ温度の飽和蒸気圧に戻る。全液体が蒸発した後は成立しない。\par

\noindent\textbf{2：適切} 平衡の液体と気体では蒸発速度と凝縮速度が等しい。界面を往復する移動が続いても、物質量の正味変化は0となる。平衡を移動停止と混同しない。\par

\noindent\textbf{3：不適切} 大気圧の沸騰中も温度は上昇し続ける。純液体の定圧気液共存では加熱が蒸発潜熱へ使われ、温度は沸点に保たれる。\par

\noindent\textbf{4：条件付きで適切} 室温・大気圧で水を入れ密封し加熱すると100°Cで沸騰しない。気相部に空気が残る定積容器なら、100°Cで水蒸気分圧が約1気圧でも全圧には空気分圧が加わる。沸騰の判定は全外圧による。空気を除いた容器などへ無条件に一般化しない。\par

\noindent\textbf{5：適切} 不揮発性溶質を溶かした希薄溶液の蒸気圧は下がる。同じ温度で溶媒の活量が低下し、理想希薄近似では \(p_A=x_Ap_A^*\) となる。溶質は蒸気圧に寄与しない。\par

\paragraph{再利用する考え方} 圧力を一定に保つ開放容器と、物質量・体積を固定する密封容器は異なる。液体が残る条件を飽和蒸気圧の議論に付ける。

\paragraph{共通理論} \hyperref[theory:phase_sep:vapor]{第\ref*{theory:phase_sep:vapor}節}, \hyperref[theory:phase_sep:boiling]{第\ref*{theory:phase_sep:boiling}節}, \hyperref[theory:phase_sep:sealed]{第\ref*{theory:phase_sep:sealed}節}, \hyperref[theory:phase_sep:raoult]{第\ref*{theory:phase_sep:raoult}節}

\par 参考：\cite{mit_solutions2008,mit_clapeyron2008}。

\paragraph{関連問題} \hyperref[q:h27-reg-applied-17]{H27-REG-APPLIED-17}

\par 難易度 2（校準初稿）、優先度未評価。

\subsection{令和元年度 問10}
\label{q:r01-reg-applied-10}

\noindent\textbf{原典：} sources/R01/17\_applied\_science.pdf、PDF物理3頁。

\textbf{公式正答：} 4（single）。

\par\textbf{検討未確定：} 原典は純物質という限定、膨張後も液体が残る条件、密封加熱時の空気・気相部を明示しない。標準的な想定での公式判定と、任意容器への全称を区別する。

\paragraph{問題の要旨} 飽和蒸気圧、動的平衡、定圧沸騰と密封加熱を条件ごとに判定する。

\Needspace{4\baselineskip}\paragraph{確認した条件・記述（独自要約）}\mbox{}\par

\noindent\textbf{確認範囲：} 問題幹、条件、否定、数値、単位、全5肢を原典画像で照合した。\par

\noindent\textbf{成立条件：} 定温、気液共存と平衡、定圧沸騰、室温大気圧での密封、不揮発性溶質という条件。蒸気の分圧と容器の全圧を区別する。\par

\paragraph{問われる核心} 飽和蒸気圧、動的平衡、定圧沸騰と密封加熱を条件ごとに判定する。

\paragraph{必要知識} 定温、気液共存と平衡、定圧沸騰、室温大気圧での密封、不揮発性溶質という条件。蒸気の分圧と容器の全圧を区別する。

\paragraph{考え方・正答を選ぶ理由（査読待ち初稿）} 二相共存が続く純物質の飽和蒸気圧は温度で決まる。体積を増すと液体が蒸発して気相の量が増え、平衡圧力を維持する。4を不適切とする公式判定はこの条件下で成立する。液体が尽きるほどの膨張とは区別する。

\Needspace{4\baselineskip}\paragraph{選択肢の検討（査読待ち初稿）}\mbox{}\par

\noindent\textbf{1：条件付きで適切} 室温・大気圧で水を入れ密封し加熱すると100°Cで沸騰しない。気相部に空気が残る定積容器なら、100°Cで水蒸気分圧が約1気圧でも全圧には空気分圧が加わる。沸騰の判定は全外圧による。空気を除いた容器などへ無条件に一般化しない。\par

\noindent\textbf{2：適切} 不揮発性溶質を溶かした希薄溶液の蒸気圧は下がる。同じ温度で溶媒の活量が低下し、理想希薄近似では \(p_A=x_Ap_A^*\) となる。溶質は蒸気圧に寄与しない。\par

\noindent\textbf{3：条件付きで適切} 大気圧で沸騰中の液体の温度は上昇しない。純物質の定圧で気液が共存する間は、加熱が潜熱となる。多成分混合物の組成変化や全液体の蒸発後まで一定とはしない。\par

\noindent\textbf{4：不適切} 定温の二相共存容器を広げると蒸気圧が変化する。膨張後も二相が残る標準的な想定では、同じ温度の飽和蒸気圧に戻るため変化しない。\par

\noindent\textbf{5：適切} 平衡の液体と気体では蒸発速度と凝縮速度が等しい。界面を往復する移動が続いても、物質量の正味変化は0となる。平衡を移動停止と混同しない。\par

\paragraph{再利用する考え方} 圧力を一定に保つ開放容器と、物質量・体積を固定する密封容器は異なる。液体が残る条件を飽和蒸気圧の議論に付ける。

\paragraph{共通理論} \hyperref[theory:phase_sep:vapor]{第\ref*{theory:phase_sep:vapor}節}, \hyperref[theory:phase_sep:boiling]{第\ref*{theory:phase_sep:boiling}節}, \hyperref[theory:phase_sep:sealed]{第\ref*{theory:phase_sep:sealed}節}, \hyperref[theory:phase_sep:raoult]{第\ref*{theory:phase_sep:raoult}節}

\par 参考：\cite{mit_solutions2008,mit_clapeyron2008}。

\paragraph{関連問題} \hyperref[q:h27-reg-applied-17]{H27-REG-APPLIED-17}, \hyperref[q:r04-reg-applied-14]{R04-REG-APPLIED-14}

\par 難易度 2（校準初稿）、優先度未評価。

\subsection{令和3年度 問11}
\label{q:r03-reg-applied-11}

\noindent\textbf{原典：} sources/R03/17\_applied\_science.pdf、PDF物理4頁。

\textbf{公式正答：} 5（single）。

\paragraph{問題の要旨} 純物質の蒸発と沸騰を区別し、平衡と相変化熱の関係を判定する。

\Needspace{4\baselineskip}\paragraph{確認した条件・記述（独自要約）}\mbox{}\par

\noindent\textbf{確認範囲：} 問題幹、条件、否定、数値、単位、全5肢を原典画像で照合した。\par

\noindent\textbf{成立条件：} 純物質の気液相変化。平衡で \(\Delta S_{\mathrm{vap}}=\Delta H_{\mathrm{vap}}/T\)、同じ端点間の逆変化でエンタルピー差の符号が逆になる。クラウジウス＝クラペイロン式は蒸気を理想気体とし液体体積を無視する近似。\par

\paragraph{問われる核心} 純物質の蒸発と沸騰を区別し、平衡と相変化熱の関係を判定する。

\paragraph{必要知識} 純物質の気液相変化。平衡で \(\Delta S_{\mathrm{vap}}=\Delta H_{\mathrm{vap}}/T\)、同じ端点間の逆変化でエンタルピー差の符号が逆になる。クラウジウス＝クラペイロン式は蒸気を理想気体とし液体体積を無視する近似。

\paragraph{考え方・正答を選ぶ理由（査読待ち初稿）} 沸騰は液体内部にも蒸気泡が成長できる条件であるが、表面からの蒸発は沸点未満でも起こる。液体が安定な条件にも飽和蒸気が共存し得る。沸点未満では気化しないとする5を選ぶ。

\Needspace{4\baselineskip}\paragraph{選択肢の検討（査読待ち初稿）}\mbox{}\par

\noindent\textbf{1：条件付きで適切} 蒸気圧の温度依存性はクラウジウス＝クラペイロン式で表す。液相体積を無視し蒸気を理想気体とする近似である。より一般的な関係はクラペイロン式であり、積分時は潜熱一定などの追加仮定が必要となる。\par

\noindent\textbf{2：適切} 蒸発エントロピーは蒸発エンタルピーと気化温度で表せる。平衡での可逆相変化に対し、エントロピー変化は潜熱を絶対温度で割った量である。\par

\noindent\textbf{3：適切} 平衡気液の化学ポテンシャルは等しい。相間の物質移動に対する自由エネルギー変化が0である。\par

\noindent\textbf{4：適切} 蒸発・凝縮エンタルピーは等大逆符号。同じ温度・圧力の端点を逆向きに結ぶ場合の状態関数差である。\par

\noindent\textbf{5：不適切} 沸点に満たない液体は気化しない。蒸発は表面から起こり、沸騰の開始条件を必要としない。\par

\paragraph{再利用する考え方} 蒸発、沸騰、凝縮を場所と圧力条件で区別する。相の安定性が液体優位でも、気相の物質量が必ず0になるわけではない。

\paragraph{共通理論} \hyperref[theory:phase_sep:chemical_potential]{第\ref*{theory:phase_sep:chemical_potential}節}, \hyperref[theory:phase_sep:boiling]{第\ref*{theory:phase_sep:boiling}節}, \hyperref[theory:phase_sep:latent]{第\ref*{theory:phase_sep:latent}節}

\par 参考：\cite{mit_clapeyron2008}。

\par 難易度 2（校準初稿）、優先度未評価。

\subsection{令和4年度 問14}
\label{q:r04-reg-applied-14}

\noindent\textbf{原典：} sources/R04/17\_applied\_science.pdf、PDF物理5頁。

\textbf{公式正答：} 2（single）。

\par\textbf{検討未確定：} 原典は純物質という限定、膨張後も液体が残る条件、密封加熱時の空気・気相部を明示しない。標準的な想定での公式判定と、任意容器への全称を区別する。

\paragraph{問題の要旨} 飽和蒸気圧、動的平衡、定圧沸騰と密封加熱を条件ごとに判定する。

\Needspace{4\baselineskip}\paragraph{確認した条件・記述（独自要約）}\mbox{}\par

\noindent\textbf{確認範囲：} 問題幹、条件、否定、数値、単位、全5肢を原典画像で照合した。\par

\noindent\textbf{成立条件：} 定温、気液共存と平衡、定圧沸騰、室温大気圧での密封、不揮発性溶質という条件。蒸気の分圧と容器の全圧を区別する。\par

\paragraph{問われる核心} 飽和蒸気圧、動的平衡、定圧沸騰と密封加熱を条件ごとに判定する。

\paragraph{必要知識} 定温、気液共存と平衡、定圧沸騰、室温大気圧での密封、不揮発性溶質という条件。蒸気の分圧と容器の全圧を区別する。

\paragraph{考え方・正答を選ぶ理由（査読待ち初稿）} 体積2倍と十分な静置だけでは液体残存の指定にならないが、二相共存が続く通常の読みにおいて飽和蒸気圧は変わらない。2が公式の不適切肢である。3は開放容器と明示され、密封加熱時の全圧上昇と区別できる。

\Needspace{4\baselineskip}\paragraph{選択肢の検討（査読待ち初稿）}\mbox{}\par

\noindent\textbf{1：条件付きで適切} 室温・大気圧で水を入れ密封し加熱すると100°Cで沸騰しない。気相部に空気が残る定積容器なら、100°Cで水蒸気分圧が約1気圧でも全圧には空気分圧が加わる。沸騰の判定は全外圧による。空気を除いた容器などへ無条件に一般化しない。\par

\noindent\textbf{2：不適切} 定温で容器体積を2倍にして十分静置すると蒸気圧が変わる。2倍後も液体が残れば追加蒸発で元の飽和蒸気圧へ戻る。静置時間は平衡到達を保証しても液体の残存量を保証しない。\par

\noindent\textbf{3：条件付きで適切} 大気圧で沸騰中の液体の温度は上昇しない。純物質の定圧で気液が共存する間は、加熱が潜熱となる。多成分混合物の組成変化や全液体の蒸発後まで一定とはしない。\par

\noindent\textbf{4：適切} 平衡の液体と気体では蒸発速度と凝縮速度が等しい。界面を往復する移動が続いても、物質量の正味変化は0となる。平衡を移動停止と混同しない。\par

\noindent\textbf{5：適切} 不揮発性溶質を溶かした希薄溶液の蒸気圧は下がる。同じ温度で溶媒の活量が低下し、理想希薄近似では \(p_A=x_Ap_A^*\) となる。溶質は蒸気圧に寄与しない。\par

\paragraph{再利用する考え方} 圧力を一定に保つ開放容器と、物質量・体積を固定する密封容器は異なる。液体が残る条件を飽和蒸気圧の議論に付ける。

\paragraph{共通理論} \hyperref[theory:phase_sep:vapor]{第\ref*{theory:phase_sep:vapor}節}, \hyperref[theory:phase_sep:boiling]{第\ref*{theory:phase_sep:boiling}節}, \hyperref[theory:phase_sep:sealed]{第\ref*{theory:phase_sep:sealed}節}, \hyperref[theory:phase_sep:raoult]{第\ref*{theory:phase_sep:raoult}節}

\par 参考：\cite{mit_solutions2008,mit_clapeyron2008}。

\paragraph{関連問題} \hyperref[q:r01-reg-applied-10]{R01-REG-APPLIED-10}

\par 難易度 2（校準初稿）、優先度未評価。

\section{水の相変化・過冷却}

\subsection{令和2年度 問16}
\label{q:r02-reg-applied-16}

\noindent\textbf{原典：} sources/R02/17\_applied\_science.pdf、PDF物理6頁。

\textbf{公式正答：} 5（single）。

\par\textbf{検討未確定：} 比誘電率の温度・周波数条件が省略され、三重点の0.01°Cがどの温度尺度か明記されない。過冷却の反例と公式正答は一致するが、全肢の限定を保持する。

\paragraph{問題の要旨} 水の誘電性・三重点・臨界状態と、過冷却という準安定状態を区別する。

\Needspace{4\baselineskip}\paragraph{確認した条件・記述（独自要約）}\mbox{}\par

\noindent\textbf{確認範囲：} 問題幹、条件、否定、数値、単位、全5肢を原典画像で照合した。\par

\noindent\textbf{成立条件：} 純水、大気圧、臨界圧力と温度の両指定。誘電率比較は通常の室温静的値、三重点はITS-90の慣用値として読む。\par

\paragraph{問われる核心} 水の誘電性・三重点・臨界状態と、過冷却という準安定状態を区別する。

\paragraph{必要知識} 純水、大気圧、臨界圧力と温度の両指定。誘電率比較は通常の室温静的値、三重点はITS-90の慣用値として読む。

\paragraph{考え方・正答を選ぶ理由（査読待ち初稿）} 大気圧で0°C未満の水が直ちに全て氷になるとは限らない。核生成には障壁があり、準安定の液体が過冷却として残る。したがって過冷却を否定する5が不適切である。三重点の常用温度とSIでの温度の定義は区別する。

\Needspace{4\baselineskip}\paragraph{選択肢の検討（査読待ち初稿）}\mbox{}\par

\noindent\textbf{1：条件付きで適切} 純水の比誘電率は多くの有機溶媒より大きい。室温付近の静的な比誘電率の典型比較として読む。全有機溶媒・全温度への順位ではない。\par

\noindent\textbf{2：適切} 固体が液相を経ず気体に昇華することがある。固相と気相が直接関わる相変化であり、氷も温度・圧力条件に応じて昇華する。\par

\noindent\textbf{3：適切} 純水が臨界圧力・臨界温度になると気液密度が一致し界面が消える。指定された臨界状態では共存する二相の差が消える。任意の定積容器を加熱すれば必ずその状態になるという命題ではない。\par

\noindent\textbf{4：条件付きで適切} 純水の三重点は0.01°Cである。ITS-90における慣用温度として正しい。2019年以降のSIでは水の三重点を厳密な熱力学温度の定義に用いない。\par

\noindent\textbf{5：不適切} 大気圧で液体の水を氷点下まで冷やせない。氷の核生成が遅れれば、平衡では固体が安定な温度でも準安定な液体が残り、過冷却となる。\par

\paragraph{再利用する考え方} 平衡相図は最も安定な相を示すが、相転移が起こるまでの時間を指定しない。温度尺度の慣用値と、熱力学温度の定義を分離する。

\paragraph{共通理論} \hyperref[theory:phase_sep:water]{第\ref*{theory:phase_sep:water}節}, \hyperref[theory:phase_sep:critical]{第\ref*{theory:phase_sep:critical}節}, \hyperref[theory:phase_sep:metastable]{第\ref*{theory:phase_sep:metastable}節}

\par 参考：\cite{iapws_supercooled2015,iapws95_2018,si_brochure9}。

\paragraph{関連問題} \hyperref[q:r06-reg-applied-10]{R06-REG-APPLIED-10}

\par 難易度 2（校準初稿）、優先度未評価。

\subsection{令和6年度 問10}
\label{q:r06-reg-applied-10}

\noindent\textbf{原典：} sources/R06/17\_applied\_science.pdf、PDF物理3頁。

\textbf{公式正答：} 1（single）。

\par\textbf{検討未確定：} 比誘電率の温度・周波数条件が省略され、三重点の0.01°Cがどの温度尺度か明記されない。過冷却の反例と公式正答は一致するが、全肢の限定を保持する。

\paragraph{問題の要旨} 水の誘電性・三重点・臨界状態と、過冷却という準安定状態を区別する。

\Needspace{4\baselineskip}\paragraph{確認した条件・記述（独自要約）}\mbox{}\par

\noindent\textbf{確認範囲：} 問題幹、条件、否定、数値、単位、全5肢を原典画像で照合した。\par

\noindent\textbf{成立条件：} 純水、大気圧、臨界圧力と温度の両指定。誘電率比較は通常の室温静的値、三重点はITS-90の慣用値として読む。\par

\paragraph{問われる核心} 水の誘電性・三重点・臨界状態と、過冷却という準安定状態を区別する。

\paragraph{必要知識} 純水、大気圧、臨界圧力と温度の両指定。誘電率比較は通常の室温静的値、三重点はITS-90の慣用値として読む。

\paragraph{考え方・正答を選ぶ理由（査読待ち初稿）} 大気圧で0°C未満の水が直ちに全て氷になるとは限らない。核生成には障壁があり、準安定の液体が過冷却として残る。したがって過冷却を否定する1が不適切である。三重点の常用温度とSIでの温度の定義は区別する。

\Needspace{4\baselineskip}\paragraph{選択肢の検討（査読待ち初稿）}\mbox{}\par

\noindent\textbf{1：不適切} 大気圧で液体の水を氷点下まで冷やせない。氷の核生成が遅れれば、平衡では固体が安定な温度でも準安定な液体が残り、過冷却となる。\par

\noindent\textbf{2：条件付きで適切} 純水の三重点は0.01°Cである。ITS-90における慣用温度として正しい。2019年以降のSIでは水の三重点を厳密な熱力学温度の定義に用いない。\par

\noindent\textbf{3：適切} 固体が液相を経ず気体に昇華することがある。固相と気相が直接関わる相変化であり、氷も温度・圧力条件に応じて昇華する。\par

\noindent\textbf{4：条件付きで適切} 純水の比誘電率は多くの有機溶媒より大きい。室温付近の静的な比誘電率の典型比較として読む。全有機溶媒・全温度への順位ではない。\par

\noindent\textbf{5：適切} 純水が臨界圧力・臨界温度になると気液密度が一致し界面が消える。指定された臨界状態では共存する二相の差が消える。任意の定積容器を加熱すれば必ずその状態になるという命題ではない。\par

\paragraph{再利用する考え方} 平衡相図は最も安定な相を示すが、相転移が起こるまでの時間を指定しない。温度尺度の慣用値と、熱力学温度の定義を分離する。

\paragraph{共通理論} \hyperref[theory:phase_sep:water]{第\ref*{theory:phase_sep:water}節}, \hyperref[theory:phase_sep:critical]{第\ref*{theory:phase_sep:critical}節}, \hyperref[theory:phase_sep:metastable]{第\ref*{theory:phase_sep:metastable}節}

\par 参考：\cite{iapws_supercooled2015,iapws95_2018,si_brochure9}。

\paragraph{関連問題} \hyperref[q:r02-reg-applied-16]{R02-REG-APPLIED-16}

\par 難易度 2（校準初稿）、優先度未評価。


\chapter{無機化学}
無機化学の判断では、元素の位置、通常状態、反応条件、材料の用途を別々に照合する。「全て」「必ず」「反応しない」は、代表的条件の知識だけでは証明できない。以下では出題時の五候補を維持し、一般化の疑義を各問に残す。
\section{希ガスと外殻例外}\label{theory:inorganic:noble}
Heの基底配置は\(1s^2\)で外殻2電子、Ne等の典型的外殻は8電子である。通常の単体気体は単原子で、低圧放電では励起・脱励起による特有の発光を示す。化学反応性が低いことは化合物が一切ないことを意味しない。Xeのフッ化物が反例になる。軽い希ガスの低沸点を高反応性と結び付けない。\cite{inorganic_noble_ref,atomic_nist2019}
\section{気体の性状と電子移動}\label{theory:inorganic:gases}
溶解性、酸塩基、酸化還元、有害性は異なる性質である。NH$_3$は水中で塩基として働く。H$_2$S中のSは酸化数\(-2\)で、通常のハロゲンとの反応では還元剤になる。Cl$_2$の水中反応で生じる次亜塩素酸等の酸化作用と区別する。石灰水では
\[\mathrm{CO_2+Ca(OH)_2\to CaCO_3\downarrow+H_2O}\]
となり、過剰CO$_2$では炭酸水素イオンを伴う溶解が進む。白濁の生成と過剰添加後の変化を二段階で読む。\cite{inorganic_halogen_ref,inorganic_hydrogen_ref}
\section{ゼオライトのイオン交換}\label{theory:inorganic:exchange}
アルミノケイ酸塩骨格の負電荷を補償する陽イオンは交換可能である。Na型の交換部位をZで表せば
\[2\mathrm{ZNa}+\mathrm{Ca^{2+}}\rightleftharpoons\mathrm{Z_2Ca}+2\mathrm{Na^+}\]
という電荷収支で硬度イオンを除去する。吸着という広い機能もあってよいが、軟化に問われる具体的機構はイオン交換である。Caの価数を変える酸化還元ではない。\cite{inorganic_zeolite}
\section{周期・族・ブロック}\label{theory:inorganic:periodic}
原子番号は陽子数であり原子質量の順位ではない。周期は横列、族は縦列として読む。副殻s,p,d,fの最大電子数は2,6,10,14である。個別の基底電子配置に例外があっても周期表のブロックを毎回移動させない。Cuは\([\mathrm{Ar}]3d^{10}4s^1\)、Znは\([\mathrm{Ar}]3d^{10}4s^2\)であり、両元素はd系列に置かれる。遷移元素という用語の12族の扱いとブロック分類は区別する。\cite{atomic_nist2019}
LiからCsへ通常の融点は低下する。第一イオン化エネルギー、同電荷のイオン半径、融点は別の測定量で、同じ傾向と決めつけない。ランタノイドとアクチノイドの区別は放射性の有無ではなく系列の位置で行う。\cite{inorganic_alkali}
\section{沈殿と錯形成}\label{theory:inorganic:complex}
沈殿はイオン積と溶解度積、錯形成は配位子の活量と生成定数に左右される。少量のNH$_3$水でZn(OH)$_2$が析出し、過剰NH$_3$では\([\mathrm{Zn(NH_3)_4}]^{2+}\)の形成が溶解を促す。白色沈殿という色だけで金属を一意に決めない。BaSO$_4$は白色、AgClは光分解で暗色化し、日光による溶解とは異なる。ZnSの酸性での析出は酸性の程度と濃度に依存するため、全酸性で不可能という命題に広げない。\cite{inorganic_zinc_ref,qualitative_stanford}
\section{Alの両性と製錬}\label{theory:inorganic:aluminum}
Alおよびその水酸化物は両性を示す。強塩基ではアルミン酸・水酸化錯イオンの形成により反応が進む。不動態は表面保護による条件付きの反応抑制であり、全酸・全温度への不反応ではない。ボーキサイト精製、アルミナ、氷晶石を含む融液中の電解を順に区別する。水溶液でAl金属を容易に得る工程と取り違えない。ミョウバンは複塩、ジュラルミンはAl主成分の合金で分類が異なる。\cite{inorganic_al_ref}
\section{ハロゲンとHF}\label{theory:inorganic:halogen}
通常状態ではF$_2$とCl$_2$は気体、Br$_2$は液体、I$_2$は固体である。典型的な酸化力系列はF$_2$からI$_2$へ低下する。通常の希薄水溶液でHFは弱酸に分類されるが、腐食性・毒性が低いことを意味しない。HF原液、濃厚水溶液、希薄水溶液の発煙条件を区別する。濃度・溶媒が変わる系へ酸の分類を無条件に延長しない。\cite{inorganic_halogen_ref,niosh_hf}
\section{分子形・格子・配位数}\label{theory:inorganic:structure}
P$_4$の四面体、CCl$_4$やSiCl$_4$の中心四配位、代表的ホウ素結晶のB$_{12}$単位は異なる構造階層である。NaClは6配位、CsClは8配位で結晶型が異なる。CsClの見かけの立方体中心は同種点のBCC格子を意味せず、単純立方格子に基底を加えた構造である。ヘムは脱酸素時の五配位に酸素が加わって六配位となる代表モデルを使うが、全結合角が理想形という保証にはしない。\cite{openstax_crystals,inorganic_metalloid_ref,inorganic_heme}
\section{フッ素系の膜除去と樹脂}\label{theory:inorganic:fluorine}
HFの半導体工程での役割はSiO$_2$膜を反応させて除去するエッチングであり、生成と逆にしない。フッ素系樹脂PTFEはTFEの付加重合体である。代表融点約327℃は再結晶材の値として読む。新生材の融解挙動や高溶融粘度を考慮し、通常の熱可塑性樹脂と同じ加工性を仮定しない。\cite{inorganic_hf_etch,inorganic_ptfe}
\section{Siの化合物と耐薬品性}\label{theory:inorganic:silicon}
通常結晶Siはdiamond型の四面体結合を持ち、高純度材は半導体に利用される。地殻中のSiは主にSiO$_2$やケイ酸塩として存在する。SiCとSiB$_6$等があるため、CやBとは結合しないという全称は誤りである。通常の酸への耐性を、酸化剤を含むHF混酸等にも一律適用しない。\cite{inorganic_metalloid_ref,inorganic_sic,inorganic_sib6,inorganic_hf_etch}
\section{元素と機器の対応}\label{theory:inorganic:materials}
代表的な青色LEDはGaN/InGaN系である。InはITO透明電極、Ndは永久磁石、Ptは排ガス浄化触媒という代表用途で整理する。特定元素が全方式に不可欠という全称には別の検証が必要である。海洋鉱床の存在、資源量評価、経済的採掘、商業生産は異なる到達段階である。未照合の当時版や数量は各問の保留事項として残す。\cite{inorganic_blueled,inorganic_rareuses,inorganic_searesources}
\section{水素化物の分類}\label{theory:inorganic:hydride}
NaH等の塩型水素化物ではH$^-$、FやO等に結合した共有分子では通常H側が正に分極する。金属格子の間隙占有による不定比水素化物を別の分類として扱う。B$_2$H$_6$の架橋は三中心二電子結合で、全結合を通常の二中心結合だけで説明しない。BH$_3$の特殊条件や配位付加体の存在と、通常持続する中性ホウ素水素化物の比較を区別する。\cite{inorganic_hydrogen_ref,inorganic_metalloid_ref}
\section{生成気体は物質収支で区別する}\label{theory:inorganic:gas_generation}
酸を加えれば全て同じ気体が出るわけではない。脱水、酸塩基、酸化還元の三種類を分ける。
\[
 \mathrm{HCOOH}\longrightarrow\mathrm{CO}+\mathrm{H_2O}.
\]
ギ酸への濃硫酸の作用による脱水では一酸化炭素を生じ、二酸化炭素へ書き換えない\cite{analysis49_co_ref}。
NaClと濃硫酸の塩化水素、Znと希硫酸の水素、MnO2と濃塩酸の塩素、Cuと濃硝酸の二酸化窒素を区別する。塩素生成は
\[
 \mathrm{MnO_2}+4\mathrm{HCl}\longrightarrow
 \mathrm{MnCl_2}+\mathrm{Cl_2}+2\mathrm{H_2O}
\]
でMnの還元とClの酸化が同時に起こる\cite{inorganic_halogen_ref}。
濃硝酸と希硝酸、濃硫酸と希硫酸では酸化性や代表生成物が違うので、原典の濃度条件を保つ。

\section{カルシウム塩の溶解・水和・潮解}\label{theory:inorganic:calcium}
炭酸カルシウムは通常純水へ難溶だが、CO2を含む水では
\[
 \mathrm{CaCO_3}+\mathrm{CO_2}+\mathrm{H_2O}
 \rightleftharpoons\mathrm{Ca^{2+}}+2\mathrm{HCO_3^-}
\]
という平衡で溶解が促進される。酸を加えた場合は炭酸塩からCO2が放出される。
CO2の溶存、pH、溶解・析出は地下水の水質と石灰岩地形の説明につながる\cite{inorganic_waterchem}。
骨・歯の無機成分にはCaを含む水酸アパタイトがある。半水石膏は加水して二水石膏となり、結晶成長によって硬化する\cite{analysis49_gypsum_ref}。
塩化カルシウムの吸湿・潮解は空気中の水分を取り込む性質であり、石膏の硬化とは別の現象である。

\section{カルシウム塩の溶解・水和・潮解を比較する。}

\subsection{平成28年度 問13}
\label{q:h28-reg-applied-13}

\noindent\textbf{原典：} sources/H28/17\_applied\_science.pdf、PDF物理5頁。

\textbf{公式正答：} 2（single）。

\paragraph{問題の要旨} カルシウム塩の溶解・水和・潮解を比較する。

\Needspace{4\baselineskip}\paragraph{確認した条件・記述（独自要約）}\mbox{}\par

\noindent\textbf{原典確認：} S7\_049の26原典画像で全五肢・否定・添字を確認。\par

\noindent\textbf{対象：} カルシウム塩の溶解・水和・潮解を比較する。\par

\paragraph{問われる核心} カルシウム塩の溶解・水和・潮解を比較する。

\paragraph{必要知識} 溶解度だけでなく化学平衡を確認する。CaCO3+CO2+H2O⇌Ca2++2HCO3−によりCO2水中での溶解が促進されるので2が不適切。半水石膏の硬化は水分が蒸発するだけでなく水和反応を含み、CaCl2の潮解とは機構が異なる。

\paragraph{考え方・正答を選ぶ理由（査読待ち初稿）} 溶解度だけでなく化学平衡を確認する。CaCO3+CO2+H2O⇌Ca2++2HCO3−によりCO2水中での溶解が促進されるので2が不適切。半水石膏の硬化は水分が蒸発するだけでなく水和反応を含み、CaCl2の潮解とは機構が異なる。

\Needspace{4\baselineskip}\paragraph{選択肢の検討（査読待ち初稿）}\mbox{}\par

\noindent\textbf{1：通常の条件で適切} CaCO3と塩酸でCO2。炭酸塩が酸で分解されCa2+・CO2・水となる。\par

\noindent\textbf{2：不適切} CO2水に不溶で純水に易溶。通常は逆。純水に難溶で、CO2の溶解で酸性となる水では炭酸水素イオンの生成を伴い溶解が促進される。\par

\noindent\textbf{3：通常の条件で適切} 骨・歯の水酸アパタイト。Caを含むリン酸塩で骨や歯の無機成分の主成分をなす。\par

\noindent\textbf{4：通常の条件で適切} 半水石膏は加水で硬化。CaSO4·1/2H2Oが二水石膏へ水和し、結晶が成長して固まる。\par

\noindent\textbf{5：通常の条件で適切} CaCl2の潮解性と吸湿。空気中の水分を取り込み、水溶液となり得るので乾燥用途に用いられる。\par

\paragraph{再利用する考え方} 測定対象と分離対象、定義の分母、絶対量と相対校正、全称と成立条件を分けて判断する。

\paragraph{共通理論} \hyperref[theory:inorganic:calcium]{第\ref*{theory:inorganic:calcium}節}

\par 参考：\cite{analysis49_gypsum_ref,inorganic_waterchem}。

\par 難易度 3（校準初稿）、優先度未評価。

\section{無機気体の生成物を反応式で確認する。}

\subsection{平成27年度 問12}
\label{q:h27-reg-applied-12}

\noindent\textbf{原典：} sources/H27/17\_applied\_science.pdf、PDF物理5頁。

\textbf{公式正答：} 1（single）。

\paragraph{問題の要旨} 無機気体の生成物を反応式で確認する。

\Needspace{4\baselineskip}\paragraph{確認した条件・記述（独自要約）}\mbox{}\par

\noindent\textbf{原典確認：} S7\_049の26原典画像で全五肢・否定・添字を確認。\par

\noindent\textbf{対象：} 無機気体の生成物を反応式で確認する。\par

\paragraph{問われる核心} 無機気体の生成物を反応式で確認する。

\paragraph{必要知識} HCOOH→CO+H2Oが脱水の物質収支である。1のCO2は脱水と別の反応を混同している。他の四肢は酸塩基又は酸化還元の標準条件として成立する。操作の実施手順を教材で補わず、濃度・加熱という原典条件と生成物を対応させる。

\paragraph{考え方・正答を選ぶ理由（査読待ち初稿）} HCOOH→CO+H2Oが脱水の物質収支である。1のCO2は脱水と別の反応を混同している。他の四肢は酸塩基又は酸化還元の標準条件として成立する。操作の実施手順を教材で補わず、濃度・加熱という原典条件と生成物を対応させる。

\Needspace{4\baselineskip}\paragraph{選択肢の検討（査読待ち初稿）}\mbox{}\par

\noindent\textbf{1：不適切} ギ酸の濃硫酸脱水でCO2。脱水生成物はCOと水。CO2への変更では酸素・水素の収支が合わない。\par

\noindent\textbf{2：通常の条件で適切} NaClと濃硫酸でHCl。不揮発性の酸により塩化物から塩化水素を発生させる。温度・当量で硫酸塩の組成は変わる。\par

\noindent\textbf{3：通常の条件で適切} Znと希硫酸でH2。Znが酸化されH+が還元される。硫酸の濃度を濃硫酸へ置き換えない。\par

\noindent\textbf{4：通常の条件で適切} MnO2と濃塩酸でCl2。Mn(IV)を還元しCl−を酸化する。MnO2+4HCl→MnCl2+Cl2+2H2O。\par

\noindent\textbf{5：通常の条件で適切} Cuと濃硝酸でNO2。酸化性酸でCuが溶け、濃硝酸では代表的にNO2を生じる。希硝酸の代表生成物NOと区別する。\par

\paragraph{再利用する考え方} 測定対象と分離対象、定義の分母、絶対量と相対校正、全称と成立条件を分けて判断する。

\paragraph{共通理論} \hyperref[theory:inorganic:gas_generation]{第\ref*{theory:inorganic:gas_generation}節}

\par 参考：\cite{analysis49_co_ref,inorganic_halogen_ref}。

\par 難易度 3（校準初稿）、優先度未評価。

\section{フッ化水素の物性と酸解離}

\subsection{令和元年度 問9}
\label{q:r01-reg-applied-09}

\noindent\textbf{原典：} sources/R01/17\_applied\_science.pdf、PDF物理3頁。

\textbf{公式正答：} 5（single）。

\par\textbf{検討未確定：} 純液体HFの誘電率と単独HFの金・白金非反応性について条件付き一次資料の全文確認が残る。発煙も濃度条件を要する。初回5と公式5の一致だけで全肢校正完了としない。

\paragraph{問題の要旨} フッ化水素の物性・ガラスへの反応・危険性と、酸解離の強さを区別する。

\Needspace{4\baselineskip}\paragraph{確認した条件・記述（独自要約）}\mbox{}\par

\noindent\textbf{確認範囲：} 問題幹、否定、全5肢、数式、数値、単位を原典画像で確認した。\par

\noindent\textbf{成立条件：} 他のハロゲン化水素との比較。純液体の誘電性と水溶液の酸解離は別。単独のHFと他の酸化剤を併用する系を区別する。\par

\paragraph{問われる核心} フッ化水素の物性・ガラスへの反応・危険性と、酸解離の強さを区別する。

\paragraph{必要知識} 他のハロゲン化水素との比較。純液体の誘電性と水溶液の酸解離は別。単独のHFと他の酸化剤を併用する系を区別する。

\paragraph{考え方・正答を選ぶ理由（査読待ち初稿）} HFは強い水素結合を形成し、他のハロゲン化水素より高い沸点を持つ。ガラスのケイ酸成分と反応する。高濃度水溶液は発煙し、弱酸であっても深刻な有害性を持つ。5は有害性を強酸であることに結び付けるため誤りである。

\Needspace{4\baselineskip}\paragraph{選択肢の検討（査読待ち初稿）}\mbox{}\par

\noindent\textbf{1：適切} 他のハロゲン化水素より沸点が高い。分子間水素結合による会合が寄与する。\par

\noindent\textbf{2：解釈未解決} 液体HFは大きな誘電率を持つ。定性的な候補として初回は適切としたが、温度条件付きの定量一次資料の照合が残る。\par

\noindent\textbf{3：条件付きで適切} ガラスと反応するが金・白金とは反応しない。ガラス腐食はNIOSHで確認できる。金・白金の非反応性を他の酸化剤との混合条件まで一般化しない。\par

\noindent\textbf{4：条件付きで適切} 水溶液は空気中で発煙する。NIOSHは濃度40パーセント超の発煙を説明する。濃度を指定しない全水溶液への全称とは読まない。\par

\noindent\textbf{5：不適切} 強い毒性を持つ強酸である。毒性・腐食性は強いが、通常の希薄水溶液における解離は弱酸の分類である。\par

\paragraph{再利用する考え方} 反応速度と平衡位置、全圧と分圧、濃度と活量を分ける。式の適用条件を確認し、原典の不正確な記述を科学的に正しい文へ置き換えない。

\paragraph{共通理論} \hyperref[theory:chem_eq:acid_strength]{第\ref*{theory:chem_eq:acid_strength}節}

\par 参考：\cite{purdue_acid_strength,niosh_hf}。

\par 難易度 2（校準初稿）、優先度未評価。

\section{Alの両性・不動態・塩と製錬を比較する。}

\subsection{平成26年度 問12}
\label{q:h26-reg-applied-12}

\noindent\textbf{原典：} sources/H26/17\_applied\_science.pdf、PDF物理5頁。

\textbf{公式正答：} 1（single）。

\par\textbf{検討未確定：} 濃硝酸の温度・濃度・表面条件、ミョウバン結晶化条件は未指定。原典の代表工程を全条件の保証にしない。

\paragraph{問題の要旨} Alの両性・不動態・塩と製錬を比較する。

\Needspace{4\baselineskip}\paragraph{確認した条件・記述（独自要約）}\mbox{}\par

\noindent\textbf{判断対象：} Alの両性・不動態・塩と製錬を比較する。\par

\noindent\textbf{原典確認：} 22原典頁画像で五候補・否定・化学式を照合。\par

\paragraph{問われる核心} Alの両性・不動態・塩と製錬を比較する。

\paragraph{必要知識} Alは両性で酸にも強塩基にも反応し、水溶液中の強塩基ではアルミン酸/水酸化錯イオン形成を伴うため1。

\paragraph{考え方・正答を選ぶ理由（査読待ち初稿）} Alは両性で酸にも強塩基にも反応し、水溶液中の強塩基ではアルミン酸/水酸化錯イオン形成を伴うため1。不動態は酸への全条件無反応を意味しない。ボーキサイト精製→アルミナ→氷晶石を含む融液の電解と、Al主合金・複塩の分類を別に整理する。

\Needspace{4\baselineskip}\paragraph{選択肢の検討（査読待ち初稿）}\mbox{}\par

\noindent\textbf{1：選ぶ候補} 酸と反応するが塩基とは反応しない。Alは強塩基とも反応し1が誤り。\par

\noindent\textbf{2：選ばない候補} 濃硝酸で酸化膜・不動態。通常の濃硝酸で表面保護が起きるという条件付き説明。\par

\noindent\textbf{3：選ばない候補} ジュラルミンはAl主成分。Alを基材とする合金分類。\par

\noindent\textbf{4：選ばない候補} Al硫酸塩/K硫酸塩からミョウバン。KAl(SO4)2・12H2Oという複塩。\par

\noindent\textbf{5：選ばない候補} ボーキサイト/アルミナ/氷晶石電解。アルミナを溶融塩に溶かして還元する。\par

\paragraph{再利用する考え方} 原子番号と周期・族、物性と反応性、酸強度と危険性、生成と除去、定義の分母を別々に確かめる。

\paragraph{共通理論} \hyperref[theory:inorganic:aluminum]{第\ref*{theory:inorganic:aluminum}節}

\par 参考：\cite{inorganic_al_ref}。

\par 難易度 2（校準初稿）、優先度未評価。

\section{レアメタルの元素と機器用途を対応させる。}

\subsection{令和4年度 問13}
\label{q:r04-reg-applied-13}

\noindent\textbf{原典：} sources/R04/17\_applied\_science.pdf、PDF物理4頁。

\textbf{公式正答：} 1（single）。

\par\textbf{検討未確定：} 機器は代表方式・出題時用途を扱う。全車両モータNd不可欠の全称、海洋鉱床の経済性・数量は未確定。

\paragraph{問題の要旨} レアメタルの元素と機器用途を対応させる。

\Needspace{4\baselineskip}\paragraph{確認した条件・記述（独自要約）}\mbox{}\par

\noindent\textbf{判断対象：} レアメタルの元素と機器用途を対応させる。\par

\noindent\textbf{原典確認：} 22原典頁画像で五候補・否定・化学式を照合。\par

\paragraph{問われる核心} レアメタルの元素と機器用途を対応させる。

\paragraph{必要知識} 代表的青色LEDはGaN/InGaNの窒化物半導体で、Asとする1が不適切。

\paragraph{考え方・正答を選ぶ理由（査読待ち初稿）} 代表的青色LEDはGaN/InGaNの窒化物半導体で、Asとする1が不適切。In透明電極、Nd磁石、Pt触媒は元素の電子構造と機器機能の対応。海洋資源の存在確認を商業採掘済みと読み替えず、Ndを全HV/EV方式に不可欠とする全称にも留保を残す。

\Needspace{4\baselineskip}\paragraph{選択肢の検討（査読待ち初稿）}\mbox{}\par

\noindent\textbf{1：選ぶ候補} 青LEDにAs。代表的青LEDはGaN/InGaN系で、Asの対応が不適切。\par

\noindent\textbf{2：選ばない候補} 透明電極にIn。ITO等の透明導電膜のIn。\par

\noindent\textbf{3：選ばない候補} 日本EEZに熱水/Coクラスト/RE泥。存在確認と経済的採掘・埋蔵量評価は別。\par

\noindent\textbf{4：選ばない候補} HV/EVモータにNd。Nd永久磁石モータの代表用途。全モータ方式に不可欠ではない。\par

\noindent\textbf{5：選ばない候補} 排ガス浄化触媒にPt。白金族金属の代表的触媒用途。\par

\paragraph{再利用する考え方} 原子番号と周期・族、物性と反応性、酸強度と危険性、生成と除去、定義の分母を別々に確かめる。

\paragraph{共通理論} \hyperref[theory:inorganic:materials]{第\ref*{theory:inorganic:materials}節}

\par 参考：\cite{inorganic_blueled,inorganic_rareuses,inorganic_searesources}。

\par 難易度 2（校準初稿）、優先度未評価。

\section{HFの酸化膜除去とPTFEの原料・加工を比較する。}

\subsection{令和元年度再試験 問9}
\label{q:r01-ret-applied-09}

\noindent\textbf{原典：} sources/R01\_retest/17\_applied\_science.pdf、PDF物理4頁。

\textbf{公式正答：} 3（single）。

\par\textbf{検討未確定：} PTFEの原料製造・全グレード加工の出題当時本文は未照合。初回と再結晶で融解温度が異なり得る。

\paragraph{問題の要旨} HFの酸化膜除去とPTFEの原料・加工を比較する。

\Needspace{4\baselineskip}\paragraph{確認した条件・記述（独自要約）}\mbox{}\par

\noindent\textbf{判断対象：} HFの酸化膜除去とPTFEの原料・加工を比較する。\par

\noindent\textbf{原典確認：} 22原典頁画像で五候補・否定・化学式を照合。\par

\paragraph{問われる核心} HFの酸化膜除去とPTFEの原料・加工を比較する。

\paragraph{必要知識} HFの半導体工程での用途は酸化膜の除去であり、生成とする3が逆。

\paragraph{考え方・正答を選ぶ理由（査読待ち初稿）} HFの半導体工程での用途は酸化膜の除去であり、生成とする3が逆。SiO2とFの反応による可溶性/揮発性フッ化物を使う。PTFEはTFE重合体で高い溶融粘度を持ち、特殊加工を使うが327℃を全熱履歴の普遍融点へ固定しない。

\Needspace{4\baselineskip}\paragraph{選択肢の検討（査読待ち初稿）}\mbox{}\par

\noindent\textbf{1：選ばない候補} FはPauling EN最大。指定尺度でFが最大。\par

\noindent\textbf{2：選ばない候補} HFはホタル石由来。CaF2からHFを製造。\par

\noindent\textbf{3：選ぶ候補} HFはウェハSiO2膜生成。HFはSiO2をエッチして除去するので生成とは逆。\par

\noindent\textbf{4：選ばない候補} TFE重合でPTFE。CF2=\allowbreak{}CF2の付加重合。\par

\noindent\textbf{5：選ばない候補} PTFE融点327℃・特殊成形。再結晶材の代表融点で、新生材などの熱履歴差を保持。\par

\paragraph{再利用する考え方} 原子番号と周期・族、物性と反応性、酸強度と危険性、生成と除去、定義の分母を別々に確かめる。

\paragraph{共通理論} \hyperref[theory:inorganic:fluorine]{第\ref*{theory:inorganic:fluorine}節}

\par 参考：\cite{inorganic_metalloid_ref,inorganic_hf_etch,inorganic_ptfe}。

\par 難易度 2（校準初稿）、優先度未評価。

\section{気体の水溶性・酸塩基・酸化還元を比較する。}

\subsection{平成23年度 問13}
\label{q:h23-reg-applied-13}

\noindent\textbf{原典：} sources/H23/17\_applied\_science.pdf、PDF物理5頁。

\textbf{公式正答：} 5（single）。

\par\textbf{検討未確定：} 最軽量は通常分子気体の同温同圧比較。NH3工業合成史の初めてという全称・全気体性状の原典当時全文は未照合。

\paragraph{問題の要旨} 気体の水溶性・酸塩基・酸化還元を比較する。

\Needspace{4\baselineskip}\paragraph{確認した条件・記述（独自要約）}\mbox{}\par

\noindent\textbf{判断対象：} 気体の水溶性・酸塩基・酸化還元を比較する。\par

\noindent\textbf{原典確認：} 22原典頁画像で五候補・否定・化学式を照合。\par

\paragraph{問われる核心} 気体の水溶性・酸塩基・酸化還元を比較する。

\paragraph{必要知識} H2SのSは−2で、例えばハロゲンへ電子を渡してSなどへ酸化される還元剤として扱う。

\paragraph{考え方・正答を選ぶ理由（査読待ち初稿）} H2SのSは−2で、例えばハロゲンへ電子を渡してSなどへ酸化される還元剤として扱う。5が誤り。硫化物の沈殿はpHと溶解度積の別機構である。CO2+Ca(OH)2→CaCO3+H2O、さらにCO2と水で炭酸水素イオンを含む溶液となる。臭い・毒性と酸化還元性を同一視しない。

\Needspace{4\baselineskip}\paragraph{選択肢の検討（査読待ち初稿）}\mbox{}\par

\noindent\textbf{1：選ばない候補} H2は最軽量で難溶。通常の中性分子気体の同温同圧比較。\par

\noindent\textbf{2：選ばない候補} Cl2は黄緑・酸化・漂白。水中HOCl等の酸化作用。\par

\noindent\textbf{3：選ばない候補} NH3は易溶アルカリ/Haber–Bosch。NH3の塩基性とN2/H2の合成。歴史の初めての範囲は留保。\par

\noindent\textbf{4：選ばない候補} CO2で石灰水白濁/過剰で溶解。CaCO3生成後に炭酸水素イオンを伴って溶解する。\par

\noindent\textbf{5：選ぶ候補} H2Sは腐卵臭・酸化性・硫化物沈殿。通常の反応比較ではH2Sは還元剤、酸化性という語が不適切。\par

\paragraph{再利用する考え方} 原子番号と周期・族、物性と反応性、酸強度と危険性、生成と除去、定義の分母を別々に確かめる。

\paragraph{共通理論} \hyperref[theory:inorganic:gases]{第\ref*{theory:inorganic:gases}節}

\par 参考：\cite{inorganic_halogen_ref,inorganic_hydrogen_ref}。

\par 難易度 2（校準初稿）、優先度未評価。

\section{ハロゲン単体の状態・酸化力とHFの酸強度を比較する。}

\subsection{平成28年度 問11}
\label{q:h28-reg-applied-11}

\noindent\textbf{原典：} sources/H28/17\_applied\_science.pdf、PDF物理4頁。

\textbf{公式正答：} 5（single）。

\par\textbf{検討未確定：} HFの酸強度は希薄水溶液の通常分類。濃厚系の活量・溶媒依存を全条件へ拡張しない。

\paragraph{問題の要旨} ハロゲン単体の状態・酸化力とHFの酸強度を比較する。

\Needspace{4\baselineskip}\paragraph{確認した条件・記述（独自要約）}\mbox{}\par

\noindent\textbf{判断対象：} ハロゲン単体の状態・酸化力とHFの酸強度を比較する。\par

\noindent\textbf{原典確認：} 22原典頁画像で五候補・否定・化学式を照合。\par

\paragraph{問われる核心} ハロゲン単体の状態・酸化力とHFの酸強度を比較する。

\paragraph{必要知識} 通常希薄水溶液ではHFは弱酸なので全強酸という候補5が不適切。

\paragraph{考え方・正答を選ぶ理由（査読待ち初稿）} 通常希薄水溶液ではHFは弱酸なので全強酸という候補5が不適切。F2単体の酸化力の強さ、HF水溶液の解離、毒性・腐食性は別の性質。酸化力系列と常温の状態を同じ五候補の順で照合する。

\Needspace{4\baselineskip}\paragraph{選択肢の検討（査読待ち初稿）}\mbox{}\par

\noindent\textbf{1：選ばない候補} 四単体は二原子分子で有毒。F2,Cl2,Br2,I2の分子単体と有害性。\par

\noindent\textbf{2：選ばない候補} F2/Cl2気体、Br2液体、I2固体。常温常圧の代表状態。\par

\noindent\textbf{3：選ばない候補} F2は水と激しく反応、I2は難溶。F2の酸化性とI2水への低溶解度。\par

\noindent\textbf{4：選ばない候補} 酸化力F2>Cl2>Br2>I2。標準的な水溶液の酸化還元比較。\par

\noindent\textbf{5：選ぶ候補} HF/HCl/HBr/HIは全強酸。通常希薄水溶液のHFは弱酸で、毒性と酸強度は別。\par

\paragraph{再利用する考え方} 原子番号と周期・族、物性と反応性、酸強度と危険性、生成と除去、定義の分母を別々に確かめる。

\paragraph{共通理論} \hyperref[theory:inorganic:halogen]{第\ref*{theory:inorganic:halogen}節}

\par 参考：\cite{inorganic_halogen_ref,niosh_hf}。

\paragraph{関連問題} \hyperref[q:r03-reg-applied-13]{R03-REG-APPLIED-13}, \hyperref[q:r07-reg-applied-17]{R07-REG-APPLIED-17}

\par 難易度 2（校準初稿）、優先度未評価。

\subsection{令和3年度 問13}
\label{q:r03-reg-applied-13}

\noindent\textbf{原典：} sources/R03/17\_applied\_science.pdf、PDF物理5頁。

\textbf{公式正答：} 1（single）。

\par\textbf{検討未確定：} 誘電率の全温度依存・Au/Ptの全混合物耐性は未照合。原液HFと水溶液濃度の条件を保持。

\paragraph{問題の要旨} HFの弱酸性と水素結合・腐食性を比較する。

\Needspace{4\baselineskip}\paragraph{確認した条件・記述（独自要約）}\mbox{}\par

\noindent\textbf{判断対象：} HFの弱酸性と水素結合・腐食性を比較する。\par

\noindent\textbf{原典確認：} 22原典頁画像で五候補・否定・化学式を照合。\par

\paragraph{問われる核心} HFの弱酸性と水素結合・腐食性を比較する。

\paragraph{必要知識} 通常HF水溶液を強酸とする1が不適切。

\paragraph{考え方・正答を選ぶ理由（査読待ち初稿）} 通常HF水溶液を強酸とする1が不適切。水中の解離が小さいことと腐食・有害性が低いことは同義でない。HF分子間水素結合はHX系列での高い沸点に関係する。NIOSHはHF液体/気体と40\%を越える水溶液の発煙を区別しており、原典の空気中発煙を全希薄水溶液の保証にはしない。

\Needspace{4\baselineskip}\paragraph{選択肢の検討（査読待ち初稿）}\mbox{}\par

\noindent\textbf{1：選ぶ候補} HF水溶液は毒性の強い強酸。毒性は強いが通常希薄水溶液は弱酸で酸強度の語が誤り。\par

\noindent\textbf{2：選ばない候補} 液体HFの高い誘電率。分極と液体の配向応答。値は温度等に依存。\par

\noindent\textbf{3：選ばない候補} 他HXより高沸点。分子間水素結合が異常な高沸点を説明。\par

\noindent\textbf{4：選ばない候補} 空気中で発煙。HF原液・濃厚水溶液で発煙。全希薄溶液とはしない。\par

\noindent\textbf{5：選ばない候補} SiO2/ガラスを腐食、Au/Ptを侵さない。ガラス除去性と貴金属の通常耐性。全混合条件保証ではない。\par

\paragraph{再利用する考え方} 原子番号と周期・族、物性と反応性、酸強度と危険性、生成と除去、定義の分母を別々に確かめる。

\paragraph{共通理論} \hyperref[theory:inorganic:halogen]{第\ref*{theory:inorganic:halogen}節}

\par 参考：\cite{niosh_hf,yale_hydrohalic}。

\paragraph{関連問題} \hyperref[q:h28-reg-applied-11]{H28-REG-APPLIED-11}

\par 難易度 2（校準初稿）、優先度未評価。

\subsection{令和7年度 問17}
\label{q:r07-reg-applied-17}

\noindent\textbf{原典：} sources/R07/17\_applied\_science.pdf、PDF物理6頁。

\textbf{公式正答：} 3（single）。

\par\textbf{検討未確定：} HFの酸強度は希薄水溶液の通常分類。濃厚系の活量・溶媒依存を全条件へ拡張しない。

\paragraph{問題の要旨} ハロゲン単体の状態・酸化力とHFの酸強度を比較する。

\Needspace{4\baselineskip}\paragraph{確認した条件・記述（独自要約）}\mbox{}\par

\noindent\textbf{判断対象：} ハロゲン単体の状態・酸化力とHFの酸強度を比較する。\par

\noindent\textbf{原典確認：} 22原典頁画像で五候補・否定・化学式を照合。\par

\paragraph{問われる核心} ハロゲン単体の状態・酸化力とHFの酸強度を比較する。

\paragraph{必要知識} 通常希薄水溶液ではHFは弱酸なので全強酸という候補3が不適切。

\paragraph{考え方・正答を選ぶ理由（査読待ち初稿）} 通常希薄水溶液ではHFは弱酸なので全強酸という候補3が不適切。F2単体の酸化力の強さ、HF水溶液の解離、毒性・腐食性は別の性質。酸化力系列と常温の状態を同じ五候補の順で照合する。

\Needspace{4\baselineskip}\paragraph{選択肢の検討（査読待ち初稿）}\mbox{}\par

\noindent\textbf{1：選ばない候補} F2/Cl2気体、Br2液体、I2固体。常温常圧の代表状態。\par

\noindent\textbf{2：選ばない候補} 酸化力F2>Cl2>Br2>I2。標準的な水溶液の酸化還元比較。\par

\noindent\textbf{3：選ぶ候補} HF/HCl/HBr/HIは全強酸。通常希薄水溶液のHFは弱酸で、毒性と酸強度は別。\par

\noindent\textbf{4：選ばない候補} F2は水と激しく反応、I2は難溶。F2の酸化性とI2水への低溶解度。\par

\noindent\textbf{5：選ばない候補} 四単体は二原子分子で有毒。F2,Cl2,Br2,I2の分子単体と有害性。\par

\paragraph{再利用する考え方} 原子番号と周期・族、物性と反応性、酸強度と危険性、生成と除去、定義の分母を別々に確かめる。

\paragraph{共通理論} \hyperref[theory:inorganic:halogen]{第\ref*{theory:inorganic:halogen}節}

\par 参考：\cite{inorganic_halogen_ref,niosh_hf}。

\paragraph{関連問題} \hyperref[q:h28-reg-applied-11]{H28-REG-APPLIED-11}

\par 難易度 2（校準初稿）、優先度未評価。

\section{水素化物の電荷・分極と多中心結合を比較する。}

\subsection{令和7年度 問16}
\label{q:r07-reg-applied-16}

\noindent\textbf{原典：} sources/R07/17\_applied\_science.pdf、PDF物理6頁。

\textbf{公式正答：} 5（single）。

\par\textbf{検討未確定：} B2H6最小安定の通常条件とBH3の特殊な単量体・配位付加体を区別。全遷移金属水素化物を不定比にしない。

\paragraph{問題の要旨} 水素化物の電荷・分極と多中心結合を比較する。

\Needspace{4\baselineskip}\paragraph{確認した条件・記述（独自要約）}\mbox{}\par

\noindent\textbf{判断対象：} 水素化物の電荷・分極と多中心結合を比較する。\par

\noindent\textbf{原典確認：} 22原典頁画像で五候補・否定・化学式を照合。\par

\paragraph{問われる核心} 水素化物の電荷・分極と多中心結合を比較する。

\paragraph{必要知識} 非金属のF/O等に結合したHは正に分極するので5が逆。

\paragraph{考え方・正答を選ぶ理由（査読待ち初稿）} 非金属のF/O等に結合したHは正に分極するので5が逆。金属塩型のH−、共有分子のHδ+、金属格子間隙の水素を区別する。B2H6の架橋B-H-Bには三中心二電子結合があり、全結合を普通の二中心結合だけで説明しない。最小安定という語はBH3の特殊条件での単量体存在を排除する普遍定義ではない。

\Needspace{4\baselineskip}\paragraph{選択肢の検討（査読待ち初稿）}\mbox{}\par

\noindent\textbf{1：選ばない候補} Hは外殻1/希ガスに1不足。H1s1はHe1s2へ1不足。Na等との共通点と独自性。\par

\noindent\textbf{2：選ばない候補} NaH/CaH2はイオン性。電気陽性金属とのH−を含む塩型。\par

\noindent\textbf{3：選ばない候補} 遷移金属間隙で不定比水素化物。格子の間隙占有率で組成が変わる系がある。\par

\noindent\textbf{4：選ばない候補} B2H6は最小安定ホウ素水素化物で3中心。通常持続する中性boraneの代表。BH3単量体の特殊条件存在と区別。\par

\noindent\textbf{5：選ぶ候補} ハロゲン/カルコゲンではHが負に分極。F/O等がHより電気陰性度大でHは通常正に分極、符号が逆。\par

\paragraph{再利用する考え方} 原子番号と周期・族、物性と反応性、酸強度と危険性、生成と除去、定義の分母を別々に確かめる。

\paragraph{共通理論} \hyperref[theory:inorganic:hydride]{第\ref*{theory:inorganic:hydride}節}

\par 参考：\cite{inorganic_hydrogen_ref,inorganic_metalloid_ref}。

\par 難易度 2（校準初稿）、優先度未評価。

\section{ゼオライトによる硬度イオン除去の機構を判断する。}

\subsection{平成23年度 問17}
\label{q:h23-reg-applied-17}

\noindent\textbf{原典：} sources/H23/17\_applied\_science.pdf、PDF物理6頁。

\textbf{公式正答：} 1（single）。

\paragraph{問題の要旨} ゼオライトによる硬度イオン除去の機構を判断する。

\Needspace{4\baselineskip}\paragraph{確認した条件・記述（独自要約）}\mbox{}\par

\noindent\textbf{判断対象：} ゼオライトによる硬度イオン除去の機構を判断する。\par

\noindent\textbf{原典確認：} 22原典頁画像で五候補・否定・化学式を照合。\par

\paragraph{問われる核心} ゼオライトによる硬度イオン除去の機構を判断する。

\paragraph{必要知識} ゼオライトの骨格負電荷を補償する交換可能陽イオンと、水中のCa2+/Mg2+が交換して硬度を下げるので1。

\paragraph{考え方・正答を選ぶ理由（査読待ち初稿）} ゼオライトの骨格負電荷を補償する交換可能陽イオンと、水中のCa2+/Mg2+が交換して硬度を下げるので1。2Na型+Ca2+⇌Ca型+2Na+という電荷収支で理解する。吸着機能もあることと、軟化に問われる主機構は別。

\Needspace{4\baselineskip}\paragraph{選択肢の検討（査読待ち初稿）}\mbox{}\par

\noindent\textbf{1：選ぶ候補} イオン交換。交換可能なNa等とCa/Mgを交換。\par

\noindent\textbf{2：選ばない候補} 吸着。表面への保持一般で、軟化の指定機構を十分に特定しない。\par

\noindent\textbf{3：選ばない候補} 吸収。内部への取込み一般で、対イオン交換を特定しない。\par

\noindent\textbf{4：選ばない候補} 還元。Ca/Mgの価数を減らす操作ではない。\par

\noindent\textbf{5：選ばない候補} 酸化。Ca/Mgを酸化する操作ではない。\par

\paragraph{再利用する考え方} 原子番号と周期・族、物性と反応性、酸強度と危険性、生成と除去、定義の分母を別々に確かめる。

\paragraph{共通理論} \hyperref[theory:inorganic:exchange]{第\ref*{theory:inorganic:exchange}節}

\par 参考：\cite{inorganic_zeolite}。

\par 難易度 2（校準初稿）、優先度未評価。

\section{希ガスの通常状態・外殻例外と放電を比較する。}

\subsection{平成23年度 問11}
\label{q:h23-reg-applied-11}

\noindent\textbf{原典：} sources/H23/17\_applied\_science.pdf、PDF物理4頁。

\textbf{公式正答：} 4（single）。

\paragraph{問題の要旨} 希ガスの通常状態・外殻例外と放電を比較する。

\Needspace{4\baselineskip}\paragraph{確認した条件・記述（独自要約）}\mbox{}\par

\noindent\textbf{判断対象：} 希ガスの通常状態・外殻例外と放電を比較する。\par

\noindent\textbf{原典確認：} 22原典頁画像で五候補・否定・化学式を照合。\par

\paragraph{問われる核心} 希ガスの通常状態・外殻例外と放電を比較する。

\paragraph{必要知識} 希ガス原子の励起・脱励起が特有の放電色を与えるので4。

\paragraph{考え方・正答を選ぶ理由（査読待ち初稿）} 希ガス原子の励起・脱励起が特有の放電色を与えるので4。通常の空気中は単体、常温常圧で気体でありHe外殻は2。Xe化合物の存在から希ガスは一切反応しないとする逆の全称にも広げない。沸点の低さと化学的反応性は別の性質。

\Needspace{4\baselineskip}\paragraph{選択肢の検討（査読待ち初稿）}\mbox{}\par

\noindent\textbf{1：選ばない候補} 空気中は化合物。通常の空気では単体原子で存在。\par

\noindent\textbf{2：選ばない候補} 常温常圧で気体と液体。出題範囲のHeからRnは通常気体。\par

\noindent\textbf{3：選ばない候補} 軽いほど低沸点・高反応。沸点傾向はよいが軽いほど反応性に富むは適切でない。\par

\noindent\textbf{4：選ぶ候補} 低圧放電で特有色。励起原子の準位遷移の発光。\par

\noindent\textbf{5：選ばない候補} 全て外殻8電子。Heの1s2は外殻2電子という例外。\par

\paragraph{再利用する考え方} 原子番号と周期・族、物性と反応性、酸強度と危険性、生成と除去、定義の分母を別々に確かめる。

\paragraph{共通理論} \hyperref[theory:inorganic:noble]{第\ref*{theory:inorganic:noble}節}

\par 参考：\cite{inorganic_noble_ref}。

\paragraph{関連問題} \hyperref[q:h27-reg-applied-10]{H27-REG-APPLIED-10}

\par 難易度 2（校準初稿）、優先度未評価。

\subsection{平成27年度 問10}
\label{q:h27-reg-applied-10}

\noindent\textbf{原典：} sources/H27/17\_applied\_science.pdf、PDF物理4頁。

\textbf{公式正答：} 4（single）。

\paragraph{問題の要旨} 希ガスの通常状態・外殻例外と放電を比較する。

\Needspace{4\baselineskip}\paragraph{確認した条件・記述（独自要約）}\mbox{}\par

\noindent\textbf{判断対象：} 希ガスの通常状態・外殻例外と放電を比較する。\par

\noindent\textbf{原典確認：} 22原典頁画像で五候補・否定・化学式を照合。\par

\paragraph{問われる核心} 希ガスの通常状態・外殻例外と放電を比較する。

\paragraph{必要知識} 希ガス原子の励起・脱励起が特有の放電色を与えるので4。

\paragraph{考え方・正答を選ぶ理由（査読待ち初稿）} 希ガス原子の励起・脱励起が特有の放電色を与えるので4。通常の空気中は単体、常温常圧で気体でありHe外殻は2。Xe化合物の存在から希ガスは一切反応しないとする逆の全称にも広げない。沸点の低さと化学的反応性は別の性質。

\Needspace{4\baselineskip}\paragraph{選択肢の検討（査読待ち初稿）}\mbox{}\par

\noindent\textbf{1：選ばない候補} 空気中は化合物。通常の空気では単体原子で存在。\par

\noindent\textbf{2：選ばない候補} 常温常圧で気体と液体。出題範囲のHeからRnは通常気体。\par

\noindent\textbf{3：選ばない候補} 軽いほど低沸点・高反応。沸点傾向はよいが軽いほど反応性に富むは適切でない。\par

\noindent\textbf{4：選ぶ候補} 低圧放電で特有色。励起原子の準位遷移の発光。\par

\noindent\textbf{5：選ばない候補} 全て外殻8電子。Heの1s2は外殻2電子という例外。\par

\paragraph{再利用する考え方} 原子番号と周期・族、物性と反応性、酸強度と危険性、生成と除去、定義の分母を別々に確かめる。

\paragraph{共通理論} \hyperref[theory:inorganic:noble]{第\ref*{theory:inorganic:noble}節}

\par 参考：\cite{inorganic_noble_ref}。

\paragraph{関連問題} \hyperref[q:h23-reg-applied-11]{H23-REG-APPLIED-11}, \hyperref[q:r01-ret-applied-16]{R01-RET-APPLIED-16}

\par 難易度 2（校準初稿）、優先度未評価。

\subsection{令和元年度再試験 問16}
\label{q:r01-ret-applied-16}

\noindent\textbf{原典：} sources/R01\_retest/17\_applied\_science.pdf、PDF物理6頁。

\textbf{公式正答：} 2（single）。

\paragraph{問題の要旨} 希ガスの単原子性とHeの電子例外を確認する。

\Needspace{4\baselineskip}\paragraph{確認した条件・記述（独自要約）}\mbox{}\par

\noindent\textbf{判断対象：} 希ガスの単原子性とHeの電子例外を確認する。\par

\noindent\textbf{原典確認：} 22原典頁画像で五候補・否定・化学式を照合。\par

\paragraph{問われる核心} 希ガスの単原子性とHeの電子例外を確認する。

\paragraph{必要知識} 原典の1原子分子は単原子気体という意味で2を選ぶ。

\paragraph{考え方・正答を選ぶ理由（査読待ち初稿）} 原典の1原子分子は単原子気体という意味で2を選ぶ。共有結合した多原子分子ではない。Heの1s2とXeの化合物の例外を保持し、通常状態と極低温や特殊な分子種を混同しない。

\Needspace{4\baselineskip}\paragraph{選択肢の検討（査読待ち初稿）}\mbox{}\par

\noindent\textbf{1：選ばない候補} 空気中は化合物。通常空気は単体。\par

\noindent\textbf{2：選ぶ候補} 1原子分子からなる。通常の単体気体は単原子で2が適切。\par

\noindent\textbf{3：選ばない候補} 常温常圧で気体と液体。出題範囲の通常単体は気体。\par

\noindent\textbf{4：選ばない候補} 軽いほど低沸点・高反応。軽いほど高反応性とは言えない。\par

\noindent\textbf{5：選ばない候補} 全て外殻8電子。Heは2で例外。\par

\paragraph{再利用する考え方} 原子番号と周期・族、物性と反応性、酸強度と危険性、生成と除去、定義の分母を別々に確かめる。

\paragraph{共通理論} \hyperref[theory:inorganic:noble]{第\ref*{theory:inorganic:noble}節}

\par 参考：\cite{inorganic_noble_ref}。

\paragraph{関連問題} \hyperref[q:h27-reg-applied-10]{H27-REG-APPLIED-10}

\par 難易度 2（校準初稿）、優先度未評価。

\section{アルカリ金属の半径・イオン化と融点を比較する。}

\subsection{平成24年度 問12}
\label{q:h24-reg-applied-12}

\noindent\textbf{原典：} sources/H24/17\_applied\_science.pdf、PDF物理4頁。

\textbf{公式正答：} 4（single）。

\paragraph{問題の要旨} アルカリ金属の半径・イオン化と融点を比較する。

\Needspace{4\baselineskip}\paragraph{確認した条件・記述（独自要約）}\mbox{}\par

\noindent\textbf{判断対象：} アルカリ金属の半径・イオン化と融点を比較する。\par

\noindent\textbf{原典確認：} 22原典頁画像で五候補・否定・化学式を照合。\par

\paragraph{問われる核心} アルカリ金属の半径・イオン化と融点を比較する。

\paragraph{必要知識} Li,Na,K,Rb,Csは外殻ns1の第1族で、殻が増え陽イオン半径が大きく第一IEは低くなる。

\paragraph{考え方・正答を選ぶ理由（査読待ち初稿）} Li,Na,K,Rb,Csは外殻ns1の第1族で、殻が増え陽イオン半径が大きく第一IEは低くなる。通常圧の融点はLiからCsへ下がるので4が不適切。融点は金属結晶の凝集と相の自由エネルギーに関係し、原子のIEの同義ではない。

\Needspace{4\baselineskip}\paragraph{選択肢の検討（査読待ち初稿）}\mbox{}\par

\noindent\textbf{1：選ばない候補} 外殻s1が失われ希ガス配置。電子を1個失う通常の1価陽イオン。\par

\noindent\textbf{2：選ばない候補} Li→Csで陽イオン半径増。同じ電荷で殻が増える系列。\par

\noindent\textbf{3：選ばない候補} Li→CsでIE低下。外殻が遠く遮蔽されて第一IEが下がる。\par

\noindent\textbf{4：選ぶ候補} Li→Csで融点上昇。通常の融点は系列で低下するので逆。\par

\noindent\textbf{5：選ばない候補} 固有の炎色。原子励起の発光スペクトルの違い。\par

\paragraph{再利用する考え方} 原子番号と周期・族、物性と反応性、酸強度と危険性、生成と除去、定義の分母を別々に確かめる。

\paragraph{共通理論} \hyperref[theory:inorganic:periodic]{第\ref*{theory:inorganic:periodic}節}

\par 参考：\cite{inorganic_alkali,atomic_nist2019}。

\par 難易度 2（校準初稿）、優先度未評価。

\subsection{平成29年度 問8}
\label{q:h29-reg-applied-08}

\noindent\textbf{原典：} sources/H29/17\_applied\_science.pdf、PDF物理3頁。

\textbf{公式正答：} 5（single）。

\paragraph{問題の要旨} 周期表の副殻容量と希土類の電子配置を比較する。

\Needspace{4\baselineskip}\paragraph{確認した条件・記述（独自要約）}\mbox{}\par

\noindent\textbf{判断対象：} 周期表の副殻容量と希土類の電子配置を比較する。\par

\noindent\textbf{原典確認：} 22原典頁画像で五候補・否定・化学式を照合。\par

\paragraph{問われる核心} 周期表の副殻容量と希土類の電子配置を比較する。

\paragraph{必要知識} 希土類をd副殻に1個ずつ増える系列とする5は誤り。

\paragraph{考え方・正答を選ぶ理由（査読待ち初稿）} 希土類をd副殻に1個ずつ増える系列とする5は誤り。ランタノイドのf電子数にも例外があり、希土類の定義はSc/Yを含む。閉副殻Beと閉殻希ガスの化学性も同一視しない。原子番号・副殻・族の三つの分類を分ける。

\Needspace{4\baselineskip}\paragraph{選択肢の検討（査読待ち初稿）}\mbox{}\par

\noindent\textbf{1：選ばない候補} 第3周期はn3まで占有。中性基底の第3周期の配置。\par

\noindent\textbf{2：選ばない候補} s/p/dは2/6/10電子。副殻軌道数1/3/5と各2電子。\par

\noindent\textbf{3：選ばない候補} Beの2s満でも不活性ではない。閉じた2s副殻だけで希ガス的全不活性にならない。\par

\noindent\textbf{4：選ばない候補} K19番目は4s。通常K基底[Ar]4s1。\par

\noindent\textbf{5：選ぶ候補} 希土類はd電子が1ずつ増える。主にf配置が変わる系列で、dの単調増加ではない。\par

\paragraph{再利用する考え方} 原子番号と周期・族、物性と反応性、酸強度と危険性、生成と除去、定義の分母を別々に確かめる。

\paragraph{共通理論} \hyperref[theory:inorganic:periodic]{第\ref*{theory:inorganic:periodic}節}

\par 参考：\cite{atomic_nist2019,inorganic_periodic_ref}。

\paragraph{関連問題} \hyperref[q:h30-reg-applied-10]{H30-REG-APPLIED-10}

\par 難易度 2（校準初稿）、優先度未評価。

\subsection{平成30年度 問10}
\label{q:h30-reg-applied-10}

\noindent\textbf{原典：} sources/H30/17\_applied\_science.pdf、PDF物理3頁。

\textbf{公式正答：} 3（single）。

\par\textbf{検討未確定：} 主族/遷移元素の定義には12族の取扱い差。dblockの分類と区別する。

\paragraph{問題の要旨} 電子配置の例外と周期表ブロックの定義を比較する。

\Needspace{4\baselineskip}\paragraph{確認した条件・記述（独自要約）}\mbox{}\par

\noindent\textbf{判断対象：} 電子配置の例外と周期表ブロックの定義を比較する。\par

\noindent\textbf{原典確認：} 22原典頁画像で五候補・否定・化学式を照合。\par

\paragraph{問われる核心} 電子配置の例外と周期表ブロックの定義を比較する。

\paragraph{必要知識} CuとZnの基底配置の4s電子数だけを見てsblockと結論する3は誤り。

\paragraph{考え方・正答を選ぶ理由（査読待ち初稿）} CuとZnの基底配置の4s電子数だけを見てsblockと結論する3は誤り。ブロックは周期表のd系列の位置として分類し、個別基底配置の例外だけで移動させない。遷移元素の語では12族を除外する定義があり、dblockとの同義を強制しない。

\Needspace{4\baselineskip}\paragraph{選択肢の検討（査読待ち初稿）}\mbox{}\par

\noindent\textbf{1：選ばない候補} s/pは主族。通常のs/p主族分類。\par

\noindent\textbf{2：選ばない候補} dは遷移、12族を除く場合あり。遷移元素の定義とdblockの範囲を区別。\par

\noindent\textbf{3：選ぶ候補} Cu4s1d10/Zn4s2d10だからsblock。電子配置は正しくてもCu/Znの周期表位置はdblock。\par

\noindent\textbf{4：選ばない候補} lanthanoid/actinoid内部遷移。内部系列の代表分類。\par

\noindent\textbf{5：選ばない候補} fblockは1原子番号で必ずf1増でない。基底配置の例外を認める。\par

\paragraph{再利用する考え方} 原子番号と周期・族、物性と反応性、酸強度と危険性、生成と除去、定義の分母を別々に確かめる。

\paragraph{共通理論} \hyperref[theory:inorganic:periodic]{第\ref*{theory:inorganic:periodic}節}

\par 参考：\cite{atomic_nist2019}。

\paragraph{関連問題} \hyperref[q:h29-reg-applied-08]{H29-REG-APPLIED-08}

\par 難易度 2（校準初稿）、優先度未評価。

\subsection{令和5年度 問12}
\label{q:r05-reg-applied-12}

\noindent\textbf{原典：} sources/R05/17\_applied\_science.pdf、PDF物理5頁。

\textbf{公式正答：} 5（single）。

\paragraph{問題の要旨} 元素の原子番号を周期表で比較する。

\Needspace{4\baselineskip}\paragraph{確認した条件・記述（独自要約）}\mbox{}\par

\noindent\textbf{判断対象：} 元素の原子番号を周期表で比較する。\par

\noindent\textbf{原典確認：} 22原典頁画像で五候補・否定・化学式を照合。\par

\paragraph{問われる核心} 元素の原子番号を周期表で比較する。

\paragraph{必要知識} 原子番号は陽子数で質量順位ではない。

\paragraph{考え方・正答を選ぶ理由（査読待ち初稿）} 原子番号は陽子数で質量順位ではない。Nb41/Ga31、Pm61/Dy66、B5/P15、I53/Kr36の比較は正しいがRe75はTa73より大きく5が不適切。原子質量や天然存在量の大小から推測しない。

\Needspace{4\baselineskip}\paragraph{選択肢の検討（査読待ち初稿）}\mbox{}\par

\noindent\textbf{1：選ばない候補} Nb>Ga。41>31。\par

\noindent\textbf{2：選ばない候補} Pm<Dy。61<66。\par

\noindent\textbf{3：選ばない候補} B<P。5<15。\par

\noindent\textbf{4：選ばない候補} I>Kr。53>36。\par

\noindent\textbf{5：選ぶ候補} Re<Ta。Re75>Ta73で小さいは逆。\par

\paragraph{再利用する考え方} 原子番号と周期・族、物性と反応性、酸強度と危険性、生成と除去、定義の分母を別々に確かめる。

\paragraph{共通理論} \hyperref[theory:inorganic:periodic]{第\ref*{theory:inorganic:periodic}節}

\par 参考：\cite{atomic_nist2019}。

\paragraph{関連問題} \hyperref[q:r06-reg-applied-13]{R06-REG-APPLIED-13}

\par 難易度 2（校準初稿）、優先度未評価。

\subsection{令和6年度 問13}
\label{q:r06-reg-applied-13}

\noindent\textbf{原典：} sources/R06/17\_applied\_science.pdf、PDF物理4頁。

\textbf{公式正答：} 5（single）。

\paragraph{問題の要旨} 周期と族を分けて元素の位置を比較する。

\Needspace{4\baselineskip}\paragraph{確認した条件・記述（独自要約）}\mbox{}\par

\noindent\textbf{判断対象：} 周期と族を分けて元素の位置を比較する。\par

\noindent\textbf{原典確認：} 22原典頁画像で五候補・否定・化学式を照合。\par

\paragraph{問われる核心} 周期と族を分けて元素の位置を比較する。

\paragraph{必要知識} MoとIrの周期を同じ6とする5が不適切。

\paragraph{考え方・正答を選ぶ理由（査読待ち初稿）} MoとIrの周期を同じ6とする5が不適切。周期は横の列、族は縦の列。Mg/Si第3、In/B13族、Nb/Tc第5、Au/Cu11族の五対応を元素の位置から確認する。

\Needspace{4\baselineskip}\paragraph{選択肢の検討（査読待ち初稿）}\mbox{}\par

\noindent\textbf{1：選ばない候補} Mg/Siは第3周期。Mg12/Si14は第3。\par

\noindent\textbf{2：選ばない候補} In/Bは同族。In49/B5は13族。\par

\noindent\textbf{3：選ばない候補} Nb/Tcは第5周期。Nb41/Tc43は第5。\par

\noindent\textbf{4：選ばない候補} Au/Cuは同族。Au79/Cu29は11族。\par

\noindent\textbf{5：選ぶ候補} Mo/Irは第6周期。Mo42は第5、Ir77は第6。\par

\paragraph{再利用する考え方} 原子番号と周期・族、物性と反応性、酸強度と危険性、生成と除去、定義の分母を別々に確かめる。

\paragraph{共通理論} \hyperref[theory:inorganic:periodic]{第\ref*{theory:inorganic:periodic}節}

\par 参考：\cite{atomic_nist2019}。

\paragraph{関連問題} \hyperref[q:r05-reg-applied-12]{R05-REG-APPLIED-12}

\par 難易度 2（校準初稿）、優先度未評価。

\subsection{令和7年度 問15}
\label{q:r07-reg-applied-15}

\noindent\textbf{原典：} sources/R07/17\_applied\_science.pdf、PDF物理5頁。

\textbf{公式正答：} 2（single）。

\paragraph{問題の要旨} ランタノイドとアクチノイドを区別する。

\Needspace{4\baselineskip}\paragraph{確認した条件・記述（独自要約）}\mbox{}\par

\noindent\textbf{判断対象：} ランタノイドとアクチノイドを区別する。\par

\noindent\textbf{原典確認：} 22原典頁画像で五候補・否定・化学式を照合。\par

\paragraph{問われる核心} ランタノイドとアクチノイドを区別する。

\paragraph{必要知識} Th90はアクチノイドなので2。

\paragraph{考え方・正答を選ぶ理由（査読待ち初稿）} Th90はアクチノイドなので2。Sm62,Nd60,Pm61,Ho67はランタノイド。放射性か安定かという軸は系列分類と別で、放射性Pmを系列外としない。希土類の広い範囲とランタノイド15元素も区別する。

\Needspace{4\baselineskip}\paragraph{選択肢の検討（査読待ち初稿）}\mbox{}\par

\noindent\textbf{1：選ばない候補} Sm。Sm62はlanthanoid。\par

\noindent\textbf{2：選ぶ候補} Th。Th90はactinoidで提示系列に属さない。\par

\noindent\textbf{3：選ばない候補} Nd。Nd60はlanthanoid。\par

\noindent\textbf{4：選ばない候補} Pm。Pm61は放射性だがlanthanoid。\par

\noindent\textbf{5：選ばない候補} Ho。Ho67はlanthanoid。\par

\paragraph{再利用する考え方} 原子番号と周期・族、物性と反応性、酸強度と危険性、生成と除去、定義の分母を別々に確かめる。

\paragraph{共通理論} \hyperref[theory:inorganic:periodic]{第\ref*{theory:inorganic:periodic}節}

\par 参考：\cite{atomic_nist2019}。

\par 難易度 2（校準初稿）、優先度未評価。

\section{金属イオンの沈殿・錯体と反応条件を確認する。}

\subsection{平成25年度 問13}
\label{q:h25-reg-applied-13}

\noindent\textbf{原典：} sources/H25/17\_applied\_science.pdf、PDF物理6頁。

\textbf{公式正答：} 4（single）。

\par\textbf{検討未確定：} 酸性ZnSの析出はpH・濃度に依存し3の一律排除をしない。濃硝酸Feの温度・濃度・表面状態も未指定。

\paragraph{問題の要旨} 金属イオンの沈殿・錯体と反応条件を確認する。

\Needspace{4\baselineskip}\paragraph{確認した条件・記述（独自要約）}\mbox{}\par

\noindent\textbf{判断対象：} 金属イオンの沈殿・錯体と反応条件を確認する。\par

\noindent\textbf{原典確認：} 22原典頁画像で五候補・否定・化学式を照合。\par

\paragraph{問われる核心} 金属イオンの沈殿・錯体と反応条件を確認する。

\paragraph{必要知識} NH3水は塩基としてOH−を供給してZn(OH)2を析出させ、過剰NH3は[Zn(NH3)4]2+として溶解を促すので4。

\paragraph{考え方・正答を選ぶ理由（査読待ち初稿）} NH3水は塩基としてOH−を供給してZn(OH)2を析出させ、過剰NH3は[Zn(NH3)4]2+として溶解を促すので4。色の誤り、光分解と溶解、酸の酸化性と不動態を分ける。3の酸性という広い表現はZnSが全酸性で析出不可という証明にはならず、pH・金属濃度・S2−の活量に留保を残す。

\Needspace{4\baselineskip}\paragraph{選択肢の検討（査読待ち初稿）}\mbox{}\par

\noindent\textbf{1：選ばない候補} BaSO4黒色沈殿。BaSO4の通常沈殿は白色。\par

\noindent\textbf{2：選ばない候補} AgCl白沈殿が日光で溶解。光分解で銀等を生じ暗色化し、溶解とは異なる。\par

\noindent\textbf{3：選ばない候補} 酸性ZnにH2Sで白沈殿。ZnS白色はあり得るが、酸性の程度・濃度で析出が変わる。\par

\noindent\textbf{4：選ぶ候補} Znに少量NH3で白沈殿/過剰で溶解。Zn(OH)2沈殿と過剰NH3のammine錯体で溶解。\par

\noindent\textbf{5：選ばない候補} Fe濃硝酸でH2と溶解。通常濃硝酸はFeを不動態化。酸化性酸からH2発生とも説明できない。\par

\paragraph{再利用する考え方} 原子番号と周期・族、物性と反応性、酸強度と危険性、生成と除去、定義の分母を別々に確かめる。

\paragraph{共通理論} \hyperref[theory:inorganic:complex]{第\ref*{theory:inorganic:complex}節}

\par 参考：\cite{qualitative_stanford,inorganic_zinc_ref}。

\par 難易度 2（校準初稿）、優先度未評価。

\section{Siの通常性状と炭化物・ホウ化物形成を比較する。}

\subsection{令和2年度 問14}
\label{q:r02-reg-applied-14}

\noindent\textbf{原典：} sources/R02/17\_applied\_science.pdf、PDF物理5頁。

\textbf{公式正答：} 4（single）。

\par\textbf{検討未確定：} 水/無機酸不溶の全称には混酸等の例外。純度・酸種・温度・表面条件未指定。公式誤りと他候補の一般化疑義を分ける。

\paragraph{問題の要旨} Siの通常性状と炭化物・ホウ化物形成を比較する。

\Needspace{4\baselineskip}\paragraph{確認した条件・記述（独自要約）}\mbox{}\par

\noindent\textbf{判断対象：} Siの通常性状と炭化物・ホウ化物形成を比較する。\par

\noindent\textbf{原典確認：} 22原典頁画像で五候補・否定・化学式を照合。\par

\paragraph{問われる核心} Siの通常性状と炭化物・ホウ化物形成を比較する。

\paragraph{必要知識} SiCやSiのホウ化物が存在するため、C/Bと結合しないという候補4が不適切。

\paragraph{考え方・正答を選ぶ理由（査読待ち初稿）} SiCやSiのホウ化物が存在するため、C/Bと結合しないという候補4が不適切。半導体用途、地殻での酸化物/ケイ酸塩、diamond型構造を区別する。一方、無機酸に全く溶けないという他候補は酸種・混合・酸化剤で変わり得るので留保を残す。

\Needspace{4\baselineskip}\paragraph{選択肢の検討（査読待ち初稿）}\mbox{}\par

\noindent\textbf{1：選ばない候補} 高純度Siは半導体材料。通常の電子材料用途。\par

\noindent\textbf{2：選ばない候補} SiO2/ケイ酸塩で地殻に多い。元素Siの単体ではなく化合物の分布。\par

\noindent\textbf{3：選ばない候補} 高純度Siは水/無機酸に不溶。通常の非酸化性酸の代表耐性。酸化剤とHFを組合せたエッチングなどに例外。\par

\noindent\textbf{4：選ぶ候補} 高温でもC/Bと直接結合しない。SiCとSiB6等の化合物が存在し無結合の全称は誤り。\par

\noindent\textbf{5：選ばない候補} Siにはdiamond構造。通常結晶Siは四面体結合のdiamond型。\par

\paragraph{再利用する考え方} 原子番号と周期・族、物性と反応性、酸強度と危険性、生成と除去、定義の分母を別々に確かめる。

\paragraph{共通理論} \hyperref[theory:inorganic:silicon]{第\ref*{theory:inorganic:silicon}節}

\par 参考：\cite{inorganic_metalloid_ref,inorganic_sic,inorganic_sib6,inorganic_hf_etch}。

\paragraph{関連問題} \hyperref[q:r06-reg-applied-09]{R06-REG-APPLIED-09}

\par 難易度 2（校準初稿）、優先度未評価。

\subsection{令和6年度 問9}
\label{q:r06-reg-applied-09}

\noindent\textbf{原典：} sources/R06/17\_applied\_science.pdf、PDF物理3頁。

\textbf{公式正答：} 5（single）。

\par\textbf{検討未確定：} 水/無機酸不溶の全称には混酸等の例外。純度・酸種・温度・表面条件未指定。公式誤りと他候補の一般化疑義を分ける。

\paragraph{問題の要旨} Siの通常性状と炭化物・ホウ化物形成を比較する。

\Needspace{4\baselineskip}\paragraph{確認した条件・記述（独自要約）}\mbox{}\par

\noindent\textbf{判断対象：} Siの通常性状と炭化物・ホウ化物形成を比較する。\par

\noindent\textbf{原典確認：} 22原典頁画像で五候補・否定・化学式を照合。\par

\paragraph{問われる核心} Siの通常性状と炭化物・ホウ化物形成を比較する。

\paragraph{必要知識} SiCやSiのホウ化物が存在するため、C/Bと結合しないという候補5が不適切。

\paragraph{考え方・正答を選ぶ理由（査読待ち初稿）} SiCやSiのホウ化物が存在するため、C/Bと結合しないという候補5が不適切。半導体用途、地殻での酸化物/ケイ酸塩、diamond型構造を区別する。一方、無機酸に全く溶けないという他候補は酸種・混合・酸化剤で変わり得るので留保を残す。

\Needspace{4\baselineskip}\paragraph{選択肢の検討（査読待ち初稿）}\mbox{}\par

\noindent\textbf{1：選ばない候補} SiO2/ケイ酸塩で地殻に多い。元素Siの単体ではなく化合物の分布。\par

\noindent\textbf{2：選ばない候補} 高純度Siは半導体材料。通常の電子材料用途。\par

\noindent\textbf{3：選ばない候補} 高純度Siは水/無機酸に不溶。通常の非酸化性酸の代表耐性。酸化剤とHFを組合せたエッチングなどに例外。\par

\noindent\textbf{4：選ばない候補} Siにはdiamond構造。通常結晶Siは四面体結合のdiamond型。\par

\noindent\textbf{5：選ぶ候補} 高温でもC/Bと直接結合しない。SiCとSiB6等の化合物が存在し無結合の全称は誤り。\par

\paragraph{再利用する考え方} 原子番号と周期・族、物性と反応性、酸強度と危険性、生成と除去、定義の分母を別々に確かめる。

\paragraph{共通理論} \hyperref[theory:inorganic:silicon]{第\ref*{theory:inorganic:silicon}節}

\par 参考：\cite{inorganic_metalloid_ref,inorganic_sic,inorganic_sib6,inorganic_hf_etch}。

\paragraph{関連問題} \hyperref[q:r02-reg-applied-14]{R02-REG-APPLIED-14}

\par 難易度 2（校準初稿）、優先度未評価。

\section{分子・結晶・錯体の配位を比較する。}

\subsection{平成29年度 問10}
\label{q:h29-reg-applied-10}

\noindent\textbf{原典：} sources/H29/17\_applied\_science.pdf、PDF物理3頁。

\textbf{公式正答：} 4（single）。

\par\textbf{検討未確定：} ホウ素全同素体の一般化とヘムの厳密理想配位には留保。原典の五配位/六配位と分子全構造を区別。

\paragraph{問題の要旨} 分子・結晶・錯体の配位を比較する。

\Needspace{4\baselineskip}\paragraph{確認した条件・記述（独自要約）}\mbox{}\par

\noindent\textbf{判断対象：} 分子・結晶・錯体の配位を比較する。\par

\noindent\textbf{原典確認：} 22原典頁画像で五候補・否定・化学式を照合。\par

\paragraph{問われる核心} 分子・結晶・錯体の配位を比較する。

\paragraph{必要知識} NaClとCsClは異なる構造で4が不適切。

\paragraph{考え方・正答を選ぶ理由（査読待ち初稿）} NaClとCsClは異なる構造で4が不適切。NaClはFCC基底の6配位、CsClは単純立方基底の8配位。P4とBXのクラスター、四面体分子、ヘムの配位数は別の階層。ヘモグロビンの近似配位を完全な理想八面体全角度の保証にしない。

\Needspace{4\baselineskip}\paragraph{選択肢の検討（査読待ち初稿）}\mbox{}\par

\noindent\textbf{1：選ばない候補} 白リンP4は正四面体。P4の4原子の分子構造。\par

\noindent\textbf{2：選ばない候補} Bは12頂点二十面体単位。代表的な結晶BのB12単位。全同素体の例外は留保。\par

\noindent\textbf{3：選ばない候補} CCl4/SiCl4は四面体。中心原子に4結合の分子形。\par

\noindent\textbf{4：選ぶ候補} NaCl/CsClは同一結晶型。NaCl6/CsCl8の配位とBravaisが違う。\par

\noindent\textbf{5：選ばない候補} deoxyheme5配位、oxyheme6配位。酸素結合で第6配位子が加わる代表近似。平面からの変位等もある。\par

\paragraph{再利用する考え方} 原子番号と周期・族、物性と反応性、酸強度と危険性、生成と除去、定義の分母を別々に確かめる。

\paragraph{共通理論} \hyperref[theory:inorganic:structure]{第\ref*{theory:inorganic:structure}節}

\par 参考：\cite{openstax_crystals,inorganic_metalloid_ref,inorganic_heme}。

\par 難易度 2（校準初稿）、優先度未評価。


\chapter{分析化学}
\section{陽イオンの系統分析}
\label{theory:analytical:separation}
\subsection{背景と基本概念}
混合溶液の成分を、一度に名前だけで判定するのでなく、条件を変えながら分離する。
沈殿ができた段階でその成分はろ液から除かれる。後の手順を解くときは、既に除いたイオンを候補に戻さない。
\subsection{物理専攻者への接続}
沈殿の選択性は平衡条件と溶解度で説明できる。イオン積と溶解度積を比較し、pHが反応種の濃度へ与える影響を考える。
硫化水素の酸解離を二段で扱う近似では、条件を満たす範囲で$[\mathrm{S}^{2-}]\propto[\mathrm{H}^+]^{-2}$となり、酸性と塩基性で硫化物の沈殿対象が変わる。
\subsection{過去問の条件}
H26の例は塩化物、酸性下の硫化物、水酸化物、アンモニア性条件の硫化物を順に分離する。
原典の手順IVを手順Iと読むと対象が変わる。実験条件の詳細は教材の簡略化から推測で補わない\cite{qualitative_stanford}。

\section{観測することと分けること}\label{theory:analysis:separation}
分析では、試料から何を知るかと、どの操作で成分を分けるかを区別する。NMRは磁場中の核スピンの共鳴を測り、通常は試料を成分別に回収しない。遠心分離は沈降、電気泳動は電場中の移動度、クロマトグラフィーは固定相と移動相の間の保持の差を利用する。
MSは作ったイオンを質量電荷比で選別するため装置内部の分離を含むが、元の中性混合物を大量に分取する方法という意味とは異なる\cite{analysis49_ms_ref}。原典の「成分分離」の広さを確認する。

\section{分子量はどの平均か}\label{theory:analysis:molar}
分子量の違う鎖が共存する高分子では、一つの値を述べるだけでは分布が分からない。分子数 \(N_i\)、分子量 \(M_i\) に対し
\[
 M_n=\frac{\sum N_iM_i}{\sum N_i},\qquad
 M_w=\frac{\sum N_iM_i^2}{\sum N_iM_i}.
\]
粒子数に依存する浸透圧は数平均に対応する。溶質質量濃度 \(c\)、絶対温度 \(T\)、気体定数 \(R\) を使い、非電解質高分子の希薄溶液では
\[
 \frac{\Pi}{c}=\frac{RT}{M_n}+O(c).
\]
濃度を変えて零へ外挿し、相互作用項を分けて切片を使う。
GPCは溶液中の大きさに応じた細孔への侵入の差を使う。標準物質による保持の校正から相対的な分子量と分布も測定できる\cite{analysis49_gpc_ref}。
鎖形状や溶媒の違いは校正へ影響する。「絶対測定」を明記しない問いで、GPCを分子量測定でないと排除しない。

\section{略称を励起・分離・検出へ分解する}\label{theory:analysis:acronyms}
ESRはElectron Spin Resonance、EPRはElectron Paramagnetic Resonanceで不対電子を対象とする。核磁気共鳴NMRとは対象が違う\cite{analysis49_esr_ref}。
HPLCの通常名はHigh Performance Liquid ChromatographyでPrecisionではない。
MALDIはMatrix-Assisted Laser Desorption/Ionizationでイオン化、TOFは飛行時間型の質量選別を表す\cite{analysis49_hplc_ref,analysis49_maldi_ref}。
CARSはCoherent Anti-Stokes Raman Spectroscopy、ESCAはElectron Spectroscopy for Chemical Analysisである\cite{analysis49_cars_ref,analysis49_xps_ref}。
FTIRのFTはFourier Transformで、核分裂飛跡と区別する。EPMAは電子線による局所分析、NAAは中性子放射化分析、GC/MSは気体クロマトグラフィーと質量分析の組合せである\cite{analysis49_epma_ref,analysis49_naa_ref}。
装置名に含まれる「何を照射するか」「何で分けるか」「何を観測するか」を別々に読む。

\section{吸光・発光と濃度}\label{theory:analysis:beer}
光強度の透過率 \(T_{\rm opt}=I/I_0\)、吸光度 \(A=-\log_{10}T_{\rm opt}\) を定める。
同じ化学種、波長、光路長 \(l\) を保つ希薄溶液で
\[
 A=\varepsilon cl,\qquad \log_{10}T_{\rm opt}=-\varepsilon cl.
\]
透過率自体は濃度に線形ではないが、その対数は負の係数で比例する\cite{analysis49_beer_ref}。濃度で会合・解離が変わる試料や強い散乱、迷光を無条件に同じ式へ入れない。
炎光分析は励起した原子が低いエネルギー状態へ戻る発光を使う。UV--Visは電子遷移、IRは主に振動遷移という励起対象の違いを押さえる。

\section{温度計と天秤にも測定条件がある}\label{theory:analysis:metrology}
温度計の感温部を入れれば常に真の温度を得るわけではない。全浸没型液体温度計と部分浸没型は校正条件が違い、後者では露出柱の温度も関係する\cite{analysis49_thermo_ref}。「水銀柱全部が同温度」は全型を網羅する規則ではない。
天秤では空気の浮力により概略の有効力が \(mg-\rho_{\rm air}Vg\) になる。同質量でも密度の違う分銅と試料は体積が違う。空気密度を温度・圧力・湿度から求め、必要なら補正する\cite{analysis49_buoy_ref}。
物理専攻者は力の釣合いへ戻って誤差の向きを判断できる。高精度測定での補正と、通常の秤量で無視する近似の大きさを区別する。

\section{内殻を測っても結合状態を調べられる}\label{theory:analysis:xray}
内殻電子を励起・電離すると空孔ができる。別の電子がその空孔へ移ると元素に特徴的な蛍光X線を出し得る。入射光子は吸収端より高エネルギー、すなわち短波長であることが必要になる。XRFの複数ピークから複数元素を分析できる\cite{analysis49_xrf_ref}。
XPSでは放出した光電子の運動エネルギーから結合エネルギーを求める。
\[
 E_{\rm bind}=h\nu-E_{\rm kin}-\phi
\]
ここで \(\phi\) は装置の仕事関数等の基準補正である。化学状態で内殻結合エネルギーに差が生じるので「内殻だから化学結合に鈍感」としない\cite{analysis49_xps_ref}。
X線発光にも高分解能で電子状態の差を扱う方法がある\cite{analysis49_sxes_ref}。微小シフトがあることと、任意の通常装置で必ずそれを分解できることは違う。

\section{分離と検出と校正}\label{theory:analysis:chrom}
保持値は固定相・移動相・温度等に依存する。同条件で標準と一致すれば候補を推定できるが、共溶出があり保持一致だけでは唯一の同定とはならない。MS検出器は質量スペクトルという追加情報を得られる\cite{analysis49_library_ref}。
外部標準の検量線では既知量と信号の関係を作る。内標準法は対象と分離される物質を加え、応答比で注入量等の変動を補正する。標準添加は試料の分取へ既知量を加え、そのマトリックスで応答を測る\cite{analysis49_chrom_ref}。
全検出器が一律に線形応答というわけではない。例えばFPD硫黄モードの非線形性がある。検出器の機構と検量線の適合を確認する。

\section{分離操作と共沸の例外}\label{theory:analysis:operations}
蒸留は気液の組成差を使うが、共沸では \(y=x\) となる。単に沸点の違う二成分だから必ず気液組成が違うとは言えない\cite{nptel_azeotrope}。
吸収は気体成分を液体内部へ溶かす操作、吸着は表面や細孔へ保持する操作、液液抽出は非混和二相間の分配を使う操作である。
逆浸透は膜を挟む浸透圧差に逆らう圧力差を与える。水を優先的に透過させる代表的な脱塩法だが、実膜の全溶質阻止率を常に100\%としない。理想膜での原理と実際の透過性能を分ける。

\section{IRは伸縮だけではない}\label{theory:analysis:ir}
二原子の近似では \(\omega\simeq\sqrt{k/\mu_{\rm red}}\) で、力定数と換算質量が振動数を決める。多原子分子には伸縮、結合角を変える変角、面外変形等の多数の正規振動がある。振動に伴って双極子が変わるものがIR吸収へ現れる\cite{analysis49_ir_ref,atomic_ir_raman}。
官能基が同じなら似た吸収領域を持つことが多いが、共役・水素結合・分子環境で変わり得る。Ramanの分極率変化という条件と、IRの双極子変化を交換しない。

\section{発光・解離・軌道エネルギー}\label{theory:analysis:photophysics}
通常の蛍光は同じスピン多重度の励起状態からの発光、りん光は典型的に三重項から一重項への遷移である。系間交差、無放射緩和、発光を一つの過程へまとめない。核は電子に比べて重く、電子遷移では核配置をほぼ固定した垂直遷移と振動波動関数の重なりを扱う。
前期解離は束縛的な状態と解離性の別状態の結合により生じ得る非放射解離である。自身の解離極限を超えた直接解離の単なる別名ではない\cite{analysis49_pred_ref}。
Hartree--Fockで他の軌道を凍結して一電子を取り去るKoopmans近似では
\[
 I_i\simeq-\varepsilon_i.
\]
軌道緩和と電子相関の変化を省く近似であり、全固体・全分子の観測値への厳密な等式ではない\cite{analysis49_koop_ref}。
発光の種類には遅延蛍光等もあり、過去問の典型的な図式と全過程の一般論を区別する。

\section{学習順序と原典の疑義}
まず対象を原子核・電子・振動・元素・分子量へ分け、次に分離の原理と検出信号を整理する。最後に検量線・測定環境・成立条件を確認する。公式の番号を保存しつつ、GPC、検出器の全称、温度計の校正などの限定不足は問題別の未解決記録へ残す。

\section{原典と全五肢の解説初稿}
\section{ESRを核磁気共鳴と取り違えた略称を判断する。}

\subsection{平成24年度 問27}
\label{q:h24-reg-applied-27}

\noindent\textbf{原典：} sources/H24/17\_applied\_science.pdf、PDF物理8頁。

\textbf{公式正答：} 4（single）。

\paragraph{問題の要旨} ESRを核磁気共鳴と取り違えた略称を判断する。

\Needspace{4\baselineskip}\paragraph{確認した条件・記述（独自要約）}\mbox{}\par

\noindent\textbf{原典確認：} S7\_049の26原典画像で全五肢・否定・添字を確認。\par

\noindent\textbf{対象：} ESRを核磁気共鳴と取り違えた略称を判断する。\par

\paragraph{問われる核心} ESRを核磁気共鳴と取り違えた略称を判断する。

\paragraph{必要知識} 磁気共鳴で対象となるスピンを区別する。ESR/EPRは不対電子、NMRはスピンを持つ原子核の共鳴を扱う。電子と核を交換した4が不適切である。

\paragraph{考え方・正答を選ぶ理由（査読待ち初稿）} 磁気共鳴で対象となるスピンを区別する。ESR/EPRは不対電子、NMRはスピンを持つ原子核の共鳴を扱う。電子と核を交換した4が不適切である。

\Needspace{4\baselineskip}\paragraph{選択肢の検討（査読待ち初稿）}\mbox{}\par

\noindent\textbf{1：通常の条件で適切} XRF／蛍光X線。X-ray Fluorescenceである。\par

\noindent\textbf{2：通常の条件で適切} SEM／走査型電子顕微鏡。Scanning Electron Microscopeである。\par

\noindent\textbf{3：通常の条件で適切} ICP-MS／誘導結合プラズマ質量分析。Inductively Coupled Plasma Mass Spectrometryで、プラズマで作ったイオンを質量分析する。\par

\noindent\textbf{4：不適切} ESR／核磁気共鳴。ESRはElectron Spin Resonance。核磁気共鳴はNMRである。\par

\noindent\textbf{5：通常の条件で適切} FTIR／フーリエ変換赤外分光。Fourier Transform Infraredで、干渉信号から赤外スペクトルを求める。\par

\paragraph{再利用する考え方} 測定対象と分離対象、定義の分母、絶対量と相対校正、全称と成立条件を分けて判断する。

\paragraph{共通理論} \hyperref[theory:analysis:acronyms]{第\ref*{theory:analysis:acronyms}節}

\par 参考：\cite{analysis49_esr_ref}。

\par 難易度 3（校準初稿）、優先度未評価。

\subsection{平成30年度 問13}
\label{q:h30-reg-applied-13}

\noindent\textbf{原典：} sources/H30/17\_applied\_science.pdf、PDF物理4頁。

\textbf{公式正答：} 2（single）。

\paragraph{問題の要旨} HPLCのPerformanceとPrecisionを区別する。

\Needspace{4\baselineskip}\paragraph{確認した条件・記述（独自要約）}\mbox{}\par

\noindent\textbf{原典確認：} S7\_049の26原典画像で全五肢・否定・添字を確認。\par

\noindent\textbf{対象：} HPLCのPerformanceとPrecisionを区別する。\par

\paragraph{問われる核心} HPLCのPerformanceとPrecisionを区別する。

\paragraph{必要知識} HPLCのPを高精度のPrecisionへ交換した肢が不適切。通常はHigh Performance Liquid Chromatographyである。他の略称は観測又はイオン化機構を表す。MALDIはイオン源、TOF等は質量分析器であり、それぞれの機能を分けて覚える。

\paragraph{考え方・正答を選ぶ理由（査読待ち初稿）} HPLCのPを高精度のPrecisionへ交換した肢が不適切。通常はHigh Performance Liquid Chromatographyである。他の略称は観測又はイオン化機構を表す。MALDIはイオン源、TOF等は質量分析器であり、それぞれの機能を分けて覚える。

\Needspace{4\baselineskip}\paragraph{選択肢の検討（査読待ち初稿）}\mbox{}\par

\noindent\textbf{1：通常の条件で適切} EPRの正式名。Electron Paramagnetic Resonance。電子常磁性共鳴でESRと同じ対象を扱う。\par

\noindent\textbf{2：不適切} HPLCをHigh Precisionとする。現在の通常名はHigh Performance Liquid Chromatography。高精度のPrecisionではない。\par

\noindent\textbf{3：通常の条件で適切} ESCAの正式名。Electron Spectroscopy for Chemical Analysis。XPSによる元素・化学状態の分析に用いられる。\par

\noindent\textbf{4：通常の条件で適切} CARSの正式名。Coherent Anti-Stokes Raman Spectroscopy。コヒーレントな反Stokes Raman信号を観測する。\par

\noindent\textbf{5：通常の条件で適切} MALDIの正式名。Matrix-Assisted Laser Desorption/Ionization。マトリックス支援レーザー脱離イオン化を表す。\par

\paragraph{再利用する考え方} 測定対象と分離対象、定義の分母、絶対量と相対校正、全称と成立条件を分けて判断する。

\paragraph{共通理論} \hyperref[theory:analysis:acronyms]{第\ref*{theory:analysis:acronyms}節}

\par 参考：\cite{analysis49_hplc_ref,analysis49_esr_ref,analysis49_maldi_ref,analysis49_cars_ref}。

\paragraph{関連問題} \hyperref[q:r06-reg-applied-14]{R06-REG-APPLIED-14}

\par 難易度 3（校準初稿）、優先度未評価。

\subsection{令和6年度 問14}
\label{q:r06-reg-applied-14}

\noindent\textbf{原典：} sources/R06/17\_applied\_science.pdf、PDF物理4頁。

\textbf{公式正答：} 4（single）。

\paragraph{問題の要旨} HPLCのPerformanceとPrecisionを区別する。

\Needspace{4\baselineskip}\paragraph{確認した条件・記述（独自要約）}\mbox{}\par

\noindent\textbf{原典確認：} S7\_049の26原典画像で全五肢・否定・添字を確認。\par

\noindent\textbf{対象：} HPLCのPerformanceとPrecisionを区別する。\par

\paragraph{問われる核心} HPLCのPerformanceとPrecisionを区別する。

\paragraph{必要知識} HPLCのPを高精度のPrecisionへ交換した肢が不適切。通常はHigh Performance Liquid Chromatographyである。他の略称は観測又はイオン化機構を表す。MALDIはイオン源、TOF等は質量分析器であり、それぞれの機能を分けて覚える。

\paragraph{考え方・正答を選ぶ理由（査読待ち初稿）} HPLCのPを高精度のPrecisionへ交換した肢が不適切。通常はHigh Performance Liquid Chromatographyである。他の略称は観測又はイオン化機構を表す。MALDIはイオン源、TOF等は質量分析器であり、それぞれの機能を分けて覚える。

\Needspace{4\baselineskip}\paragraph{選択肢の検討（査読待ち初稿）}\mbox{}\par

\noindent\textbf{1：通常の条件で適切} CARSの正式名。Coherent Anti-Stokes Raman Spectroscopy。コヒーレントな反Stokes Raman信号を観測する。\par

\noindent\textbf{2：通常の条件で適切} EPRの正式名。Electron Paramagnetic Resonance。電子常磁性共鳴でESRと同じ対象を扱う。\par

\noindent\textbf{3：通常の条件で適切} ESCAの正式名。Electron Spectroscopy for Chemical Analysis。XPSによる元素・化学状態の分析に用いられる。\par

\noindent\textbf{4：不適切} HPLCをHigh Precisionとする。現在の通常名はHigh Performance Liquid Chromatography。高精度のPrecisionではない。\par

\noindent\textbf{5：通常の条件で適切} MALDIの正式名。Matrix-Assisted Laser Desorption/Ionization。マトリックス支援レーザー脱離イオン化を表す。\par

\paragraph{再利用する考え方} 測定対象と分離対象、定義の分母、絶対量と相対校正、全称と成立条件を分けて判断する。

\paragraph{共通理論} \hyperref[theory:analysis:acronyms]{第\ref*{theory:analysis:acronyms}節}

\par 参考：\cite{analysis49_hplc_ref,analysis49_esr_ref,analysis49_maldi_ref,analysis49_cars_ref}。

\paragraph{関連問題} \hyperref[q:h30-reg-applied-13]{H30-REG-APPLIED-13}

\par 難易度 3（校準初稿）、優先度未評価。

\subsection{令和6年度 問19}
\label{q:r06-reg-applied-19}

\noindent\textbf{原典：} sources/R06/17\_applied\_science.pdf、PDF物理6頁。

\textbf{公式正答：} 1（single）。

\paragraph{問題の要旨} FTIRとフィッショントラックを区別する。

\Needspace{4\baselineskip}\paragraph{確認した条件・記述（独自要約）}\mbox{}\par

\noindent\textbf{原典確認：} S7\_049の26原典画像で全五肢・否定・添字を確認。\par

\noindent\textbf{対象：} FTIRとフィッショントラックを区別する。\par

\paragraph{問われる核心} FTIRとフィッショントラックを区別する。

\paragraph{必要知識} フーリエ変換赤外というFTIRと、年代測定等で使う核分裂飛跡を区別する。EPMAは電子線、XRFはX線、NAAは中性子を試料へ与える。GC/MSは分離器と検出分析器の組合せである。

\paragraph{考え方・正答を選ぶ理由（査読待ち初稿）} フーリエ変換赤外というFTIRと、年代測定等で使う核分裂飛跡を区別する。EPMAは電子線、XRFはX線、NAAは中性子を試料へ与える。GC/MSは分離器と検出分析器の組合せである。

\Needspace{4\baselineskip}\paragraph{選択肢の検討（査読待ち初稿）}\mbox{}\par

\noindent\textbf{1：不適切} FTIR／フィッショントラック。FTIRはFourier Transform Infrared。フィッショントラックは核分裂飛跡を用いる別方法。\par

\noindent\textbf{2：通常の条件で適切} EPMA／電子線マイクロアナライザー。Electron Probe Microanalyzerで局所に電子線を照射し特性X線等を測る。\par

\noindent\textbf{3：通常の条件で適切} GC/MS／ガスクロマトグラフ質量分析。GCで分離した成分をMSで測定する複合方法。\par

\noindent\textbf{4：通常の条件で適切} NAA／中性子放射化分析。Neutron Activation Analysis。中性子照射後に生じる放射性核種等の信号から元素を分析する。\par

\noindent\textbf{5：通常の条件で適切} XRF／蛍光X線。X-ray Fluorescenceである。\par

\paragraph{再利用する考え方} 測定対象と分離対象、定義の分母、絶対量と相対校正、全称と成立条件を分けて判断する。

\paragraph{共通理論} \hyperref[theory:analysis:acronyms]{第\ref*{theory:analysis:acronyms}節}

\par 参考：\cite{analysis49_ir_ref,analysis49_epma_ref,analysis49_naa_ref}。

\par 難易度 3（校準初稿）、優先度未評価。

\section{クロマトグラフィーの保持・検出・検量線を区別する。}

\subsection{令和5年度 問22}
\label{q:r05-reg-applied-22}

\noindent\textbf{原典：} sources/R05/17\_applied\_science.pdf、PDF物理8頁。

\textbf{公式正答：} 4（single）。

\par\textbf{検討未確定：} 公式4を保持。肢1はMS等の成分識別情報とFPD硫黄の非線形応答に反例がある。一般濃度検出器という限定が原典になく、全五肢の排他性を査読。

\paragraph{問題の要旨} クロマトグラフィーの保持・検出・検量線を区別する。

\Needspace{4\baselineskip}\paragraph{確認した条件・記述（独自要約）}\mbox{}\par

\noindent\textbf{原典確認：} S7\_049の26原典画像で全五肢・否定・添字を確認。\par

\noindent\textbf{対象：} クロマトグラフィーの保持・検出・検量線を区別する。\par

\paragraph{問われる核心} クロマトグラフィーの保持・検出・検量線を区別する。

\paragraph{必要知識} 内標準と対象成分のピークは分離される必要がある。比を取る二信号を分けて測れない4が公式の不適切肢である。ただし1も検出器を全称化するとMSや非線形検出器に合わない。2の保持一致は同一性の推定であって単独の厳密な証明ではない。

\paragraph{考え方・正答を選ぶ理由（査読待ち初稿）} 内標準と対象成分のピークは分離される必要がある。比を取る二信号を分けて測れない4が公式の不適切肢である。ただし1も検出器を全称化するとMSや非線形検出器に合わない。2の保持一致は同一性の推定であって単独の厳密な証明ではない。

\Needspace{4\baselineskip}\paragraph{選択肢の検討（査読待ち初稿）}\mbox{}\par

\noindent\textbf{1：条件付き・解釈留保} 検出器は常に比例応答で成分同定不可。一般の濃度検出器の略説には対応するがMS検出器は質量スペクトル等を得る。FPD硫黄モードのように応答が量の二乗となる方式もあり、検出器全ての性質ではない。\par

\noindent\textbf{2：条件付き・解釈留保} 保持値一致で同一物質と推定。同条件なら同定の手掛かりになるが、異なる成分が共溶出することもある。推定と確証を区別する。\par

\noindent\textbf{3：通常の条件で適切} 絶対検量線は量とピーク面積の関係。既知濃度の標準試料から応答曲線を作り未知試料へ適用する。名称に反して測定値が校正不要という意味ではない。\par

\noindent\textbf{4：不適切} 内標準に重なるピークを選ぶ。各対象ピークから分離される内標準を選び、面積比等で変動を補正する。重なると双方の信号を独立に測れない。\par

\noindent\textbf{5：通常の条件で適切} 標準添加でマトリックス影響を考慮。同じ試料の分取へ既知量を添加することで試料環境での応答を校正する。直線性・加算性等は確認する。\par

\paragraph{再利用する考え方} 測定対象と分離対象、定義の分母、絶対量と相対校正、全称と成立条件を分けて判断する。

\paragraph{共通理論} \hyperref[theory:analysis:chrom]{第\ref*{theory:analysis:chrom}節}

\par 参考：\cite{analysis49_chrom_ref,analysis49_library_ref,analysis49_hplc_ref}。

\par 難易度 4（校準初稿）、優先度未評価。

\section{IRの伸縮と変角を区別し分光法の対象を対応させる。}

\subsection{令和5年度 問15}
\label{q:r05-reg-applied-15}

\noindent\textbf{原典：} sources/R05/17\_applied\_science.pdf、PDF物理5頁。

\textbf{公式正答：} 2（single）。

\par\textbf{検討未確定：} モル吸光係数を構造だけの無条件定数とはしない。波長・溶媒・化学種等を指定する条件付きの説明として保持。

\paragraph{問題の要旨} IRの伸縮と変角を区別し分光法の対象を対応させる。

\Needspace{4\baselineskip}\paragraph{確認した条件・記述（独自要約）}\mbox{}\par

\noindent\textbf{原典確認：} S7\_049の26原典画像で全五肢・否定・添字を確認。\par

\noindent\textbf{対象：} IRの伸縮と変角を区別し分光法の対象を対応させる。\par

\paragraph{問われる核心} IRの伸縮と変角を区別し分光法の対象を対応させる。

\paragraph{必要知識} IRは伸縮だけでなく変角・曲げ等の振動も観測できる。したがってのみという限定の肢を選ぶ。NMRは核スピン、UV–Visは電子遷移を対象とする。官能基の吸収領域は近似的な構造指標であり、水素結合や共役で変わり得る。

\paragraph{考え方・正答を選ぶ理由（査読待ち初稿）} IRは伸縮だけでなく変角・曲げ等の振動も観測できる。したがってのみという限定の肢を選ぶ。NMRは核スピン、UV–Visは電子遷移を対象とする。官能基の吸収領域は近似的な構造指標であり、水素結合や共役で変わり得る。

\Needspace{4\baselineskip}\paragraph{選択肢の検討（査読待ち初稿）}\mbox{}\par

\noindent\textbf{1：通常の条件で適切} NMRは核スピンとラジオ波の共鳴。静磁場中でスピンを持つ核のエネルギー差に対応する高周波を用いる。\par

\noindent\textbf{2：不適切} IRは伸縮運動のみ。結合角を変える変角・曲げ振動等も双極子を変えれば吸収する。\par

\noindent\textbf{3：通常の条件で適切} UV–Visは電子遷移。光子吸収で分子の電子状態を励起する。\par

\noindent\textbf{4：条件付き・解釈留保} 強度と構造固有のモル吸光係数。εは化合物の構造に関係するが、波長・溶媒・温度・化学種状態等も指定する必要がある。\par

\noindent\textbf{5：通常の条件で適切} 同じ官能基は近い吸収領域。局所的な結合の力定数と換算質量が似るため。共役・水素結合・環境によるずれはある。\par

\paragraph{再利用する考え方} 測定対象と分離対象、定義の分母、絶対量と相対校正、全称と成立条件を分けて判断する。

\paragraph{共通理論} \hyperref[theory:analysis:ir]{第\ref*{theory:analysis:ir}節}

\par 参考：\cite{analysis49_ir_ref,analysis49_beer_ref}。

\paragraph{関連問題} \hyperref[q:r07-reg-applied-14]{R07-REG-APPLIED-14}

\par 難易度 4（校準初稿）、優先度未評価。

\subsection{令和7年度 問14}
\label{q:r07-reg-applied-14}

\noindent\textbf{原典：} sources/R07/17\_applied\_science.pdf、PDF物理5頁。

\textbf{公式正答：} 4（single）。

\par\textbf{検討未確定：} モル吸光係数を構造だけの無条件定数とはしない。波長・溶媒・化学種等を指定する条件付きの説明として保持。

\paragraph{問題の要旨} IRの伸縮と変角を区別し分光法の対象を対応させる。

\Needspace{4\baselineskip}\paragraph{確認した条件・記述（独自要約）}\mbox{}\par

\noindent\textbf{原典確認：} S7\_049の26原典画像で全五肢・否定・添字を確認。\par

\noindent\textbf{対象：} IRの伸縮と変角を区別し分光法の対象を対応させる。\par

\paragraph{問われる核心} IRの伸縮と変角を区別し分光法の対象を対応させる。

\paragraph{必要知識} IRは伸縮だけでなく変角・曲げ等の振動も観測できる。したがってのみという限定の肢を選ぶ。NMRは核スピン、UV–Visは電子遷移を対象とする。官能基の吸収領域は近似的な構造指標であり、水素結合や共役で変わり得る。

\paragraph{考え方・正答を選ぶ理由（査読待ち初稿）} IRは伸縮だけでなく変角・曲げ等の振動も観測できる。したがってのみという限定の肢を選ぶ。NMRは核スピン、UV–Visは電子遷移を対象とする。官能基の吸収領域は近似的な構造指標であり、水素結合や共役で変わり得る。

\Needspace{4\baselineskip}\paragraph{選択肢の検討（査読待ち初稿）}\mbox{}\par

\noindent\textbf{1：通常の条件で適切} NMRは核スピンとラジオ波の共鳴。静磁場中でスピンを持つ核のエネルギー差に対応する高周波を用いる。\par

\noindent\textbf{2：通常の条件で適切} UV–Visは電子遷移。光子吸収で分子の電子状態を励起する。\par

\noindent\textbf{3：条件付き・解釈留保} 強度と構造固有のモル吸光係数。εは化合物の構造に関係するが、波長・溶媒・温度・化学種状態等も指定する必要がある。\par

\noindent\textbf{4：不適切} IRは伸縮運動のみ。結合角を変える変角・曲げ振動等も双極子を変えれば吸収する。\par

\noindent\textbf{5：通常の条件で適切} 同じ官能基は近い吸収領域。局所的な結合の力定数と換算質量が似るため。共役・水素結合・環境によるずれはある。\par

\paragraph{再利用する考え方} 測定対象と分離対象、定義の分母、絶対量と相対校正、全称と成立条件を分けて判断する。

\paragraph{共通理論} \hyperref[theory:analysis:ir]{第\ref*{theory:analysis:ir}節}

\par 参考：\cite{analysis49_ir_ref,analysis49_beer_ref}。

\paragraph{関連問題} \hyperref[q:r05-reg-applied-15]{R05-REG-APPLIED-15}

\par 難易度 4（校準初稿）、優先度未評価。

\section{分析信号と水銀温度計・天秤の環境依存を比較する。}

\subsection{平成29年度 問12}
\label{q:h29-reg-applied-12}

\noindent\textbf{原典：} sources/H29/17\_applied\_science.pdf、PDF物理4頁。

\textbf{公式正答：} 5（single）。

\par\textbf{検討未確定：} 水銀温度計の全浸没型という条件が原典に明記されていない。部分浸没型の校正まで全水銀柱同温と一般化しない。Beer則にも化学種・波長・散乱等の成立条件がある。

\paragraph{問題の要旨} 分析信号と水銀温度計・天秤の環境依存を比較する。

\Needspace{4\baselineskip}\paragraph{確認した条件・記述（独自要約）}\mbox{}\par

\noindent\textbf{原典確認：} S7\_049の26原典画像で全五肢・否定・添字を確認。\par

\noindent\textbf{対象：} 分析信号と水銀温度計・天秤の環境依存を比較する。\par

\paragraph{問われる核心} 分析信号と水銀温度計・天秤の環境依存を比較する。

\paragraph{必要知識} 天秤で下向きの有効力は概略mg−ρair Vg。空気密度は温度・圧力・湿度に依存し、試料と分銅の体積差によって浮力補正が必要となる。全く変わらない5を選ぶ。Beer則A=−log10 T=εclなのでH29のlog透過度も負の比例であり、R02の吸光度も正の比例となる。

\paragraph{考え方・正答を選ぶ理由（査読待ち初稿）} 天秤で下向きの有効力は概略mg−ρair Vg。空気密度は温度・圧力・湿度に依存し、試料と分銅の体積差によって浮力補正が必要となる。全く変わらない5を選ぶ。Beer則A=−log10 T=εclなのでH29のlog透過度も負の比例であり、R02の吸光度も正の比例となる。

\Needspace{4\baselineskip}\paragraph{選択肢の検討（査読待ち初稿）}\mbox{}\par

\noindent\textbf{1：通常の条件で適切} 炎中の原子励起と発光。高い状態から低い状態への遷移による特性発光を測る。\par

\noindent\textbf{2：通常の条件で適切} 透過度の対数の濃度比例。固定波長・光路長でBeer則A=−log10T=εcl。logTも係数の符号が負の比例関係となる。\par

\noindent\textbf{3：条件付き・解釈留保} 水銀柱全部が同温度なら正しい。全浸没用温度計ではその条件に対応する。一方、部分浸没型は指定浸没位置と露出柱の校正条件があり、全型に無条件の規則ではない。\par

\noindent\textbf{4：通常の条件で適切} XRF波長が化学結合状態で変わり得る。内殻の信号にも周辺電子状態に応じた小さいエネルギー差があり、高分解能では化学シフトを扱える。\par

\noindent\textbf{5：不適切} 天秤は温度・圧力・湿度で変化しない。空気密度と浮力が変わる。試料と分銅の密度差があると見かけの比較が変わる。\par

\paragraph{再利用する考え方} 測定対象と分離対象、定義の分母、絶対量と相対校正、全称と成立条件を分けて判断する。

\paragraph{共通理論} \hyperref[theory:analysis:metrology]{第\ref*{theory:analysis:metrology}節}

\par 参考：\cite{analysis49_beer_ref,analysis49_thermo_ref,analysis49_buoy_ref,analysis49_sxes_ref}。

\paragraph{関連問題} \hyperref[q:r02-reg-applied-11]{R02-REG-APPLIED-11}

\par 難易度 4（校準初稿）、優先度未評価。

\subsection{令和2年度 問11}
\label{q:r02-reg-applied-11}

\noindent\textbf{原典：} sources/R02/17\_applied\_science.pdf、PDF物理4頁。

\textbf{公式正答：} 5（single）。

\par\textbf{検討未確定：} 水銀温度計の全浸没型という条件が原典に明記されていない。部分浸没型の校正まで全水銀柱同温と一般化しない。Beer則にも化学種・波長・散乱等の成立条件がある。

\paragraph{問題の要旨} 分析信号と水銀温度計・天秤の環境依存を比較する。

\Needspace{4\baselineskip}\paragraph{確認した条件・記述（独自要約）}\mbox{}\par

\noindent\textbf{原典確認：} S7\_049の26原典画像で全五肢・否定・添字を確認。\par

\noindent\textbf{対象：} 分析信号と水銀温度計・天秤の環境依存を比較する。\par

\paragraph{問われる核心} 分析信号と水銀温度計・天秤の環境依存を比較する。

\paragraph{必要知識} 天秤で下向きの有効力は概略mg−ρair Vg。空気密度は温度・圧力・湿度に依存し、試料と分銅の体積差によって浮力補正が必要となる。全く変わらない5を選ぶ。Beer則A=−log10 T=εclなのでH29のlog透過度も負の比例であり、R02の吸光度も正の比例となる。

\paragraph{考え方・正答を選ぶ理由（査読待ち初稿）} 天秤で下向きの有効力は概略mg−ρair Vg。空気密度は温度・圧力・湿度に依存し、試料と分銅の体積差によって浮力補正が必要となる。全く変わらない5を選ぶ。Beer則A=−log10 T=εclなのでH29のlog透過度も負の比例であり、R02の吸光度も正の比例となる。

\Needspace{4\baselineskip}\paragraph{選択肢の検討（査読待ち初稿）}\mbox{}\par

\noindent\textbf{1：通常の条件で適切} 炎中の原子励起と発光。高い状態から低い状態への遷移による特性発光を測る。\par

\noindent\textbf{2：通常の条件で適切} 吸光度の濃度比例。固定波長・光路長でBeer則A=−log10T=εcl。logTも係数の符号が負の比例関係となる。\par

\noindent\textbf{3：条件付き・解釈留保} 水銀柱全部が同温度なら正しい。全浸没用温度計ではその条件に対応する。一方、部分浸没型は指定浸没位置と露出柱の校正条件があり、全型に無条件の規則ではない。\par

\noindent\textbf{4：通常の条件で適切} XRF波長が化学結合状態で変わり得る。内殻の信号にも周辺電子状態に応じた小さいエネルギー差があり、高分解能では化学シフトを扱える。\par

\noindent\textbf{5：不適切} 天秤は温度・圧力・湿度で変化しない。空気密度と浮力が変わる。試料と分銅の密度差があると見かけの比較が変わる。\par

\paragraph{再利用する考え方} 測定対象と分離対象、定義の分母、絶対量と相対校正、全称と成立条件を分けて判断する。

\paragraph{共通理論} \hyperref[theory:analysis:metrology]{第\ref*{theory:analysis:metrology}節}

\par 参考：\cite{analysis49_beer_ref,analysis49_thermo_ref,analysis49_buoy_ref,analysis49_sxes_ref}。

\paragraph{関連問題} \hyperref[q:h29-reg-applied-12]{H29-REG-APPLIED-12}

\par 難易度 4（校準初稿）、優先度未評価。

\section{高分子分子量の相対GPC測定と浸透圧による数平均を比較する。}

\subsection{平成25年度 問16}
\label{q:h25-reg-applied-16}

\noindent\textbf{原典：} sources/H25/17\_applied\_science.pdf、PDF物理7頁。

\textbf{公式正答：} 3（single）。

\par\textbf{独立解答と公式正答に相違あり：} 独立解答 1|3。公式3を保持。初回独立1|3を保存。1も校正されたGPCによる分子量測定として成立し、単一正答の排他性は未解決。

\paragraph{問題の要旨} 高分子分子量の相対GPC測定と浸透圧による数平均を比較する。

\Needspace{4\baselineskip}\paragraph{確認した条件・記述（独自要約）}\mbox{}\par

\noindent\textbf{原典確認：} S7\_049の26原典画像で全五肢・否定・添字を確認。\par

\noindent\textbf{対象：} 高分子分子量の相対GPC測定と浸透圧による数平均を比較する。\par

\paragraph{問われる核心} 高分子分子量の相対GPC測定と浸透圧による数平均を比較する。

\paragraph{必要知識} 質量濃度c、数平均分子量Mnなら希薄溶液でΠ/c=\allowbreak{}RT/Mn+O(c)となる。濃度を変えてゼロへ外挿すれば切片からMnを得るので公式は3。一方、1のGPCも標準試料による相対校正で分子量を測定する実用法である。絶対測定の優位性で3を説明できるが、原典の一般的な問いから1を排除する十分な限定はない。

\paragraph{考え方・正答を選ぶ理由（査読待ち初稿）} 質量濃度c、数平均分子量Mnなら希薄溶液でΠ/c=\allowbreak{}RT/Mn+O(c)となる。濃度を変えてゼロへ外挿すれば切片からMnを得るので公式は3。一方、1のGPCも標準試料による相対校正で分子量を測定する実用法である。絶対測定の優位性で3を説明できるが、原典の一般的な問いから1を排除する十分な限定はない。

\Needspace{4\baselineskip}\paragraph{選択肢の検討（査読待ち初稿）}\mbox{}\par

\noindent\textbf{1：条件付き・解釈留保} ゲルろ過の溶離時間。サイズ排除クロマトグラフィーとして標準物質で校正すれば分子量と分布を測定できる。原典に絶対測定という限定はなく、一般の分子量測定法としても成立する。\par

\noindent\textbf{2：不適切} 中和滴定から解離度。示されたpH・解離定数の式は電離状態を調べる。末端基の数など追加情報なしに一般の高分子分子量を与えない。\par

\noindent\textbf{3：通常の条件で適切} 半透膜による浸透圧。希薄溶液の粒子数に依存する浸透圧から数平均分子量を求める。濃度ゼロへの外挿で相互作用の寄与を減らす。\par

\noindent\textbf{4：不適切} ガラス転移点の測定。分子量と相関することはあるが、化学構造・可塑剤等にも依存し、記述だけでは一意に分子量が決まらない。\par

\noindent\textbf{5：不適切} 円偏光二色性。左右円偏光に対する吸収差を測り、立体構造や二次構造を調べる。一般の分子量直接測定ではない。\par

\paragraph{再利用する考え方} 測定対象と分離対象、定義の分母、絶対量と相対校正、全称と成立条件を分けて判断する。

\paragraph{共通理論} \hyperref[theory:analysis:molar]{第\ref*{theory:analysis:molar}節}

\par 参考：\cite{analysis49_gpc_ref}。

\par 難易度 4（校準初稿）、優先度未評価。

\section{成分の分離とNMRによる観測を区別する。}

\subsection{平成25年度 問15}
\label{q:h25-reg-applied-15}

\noindent\textbf{原典：} sources/H25/17\_applied\_science.pdf、PDF物理6頁。

\textbf{公式正答：} 1（single）。

\par\textbf{検討未確定：} 質量分析でのイオン選別と、中性試料の調製・分取という分離の意味は同一でない。原典の広い成分分離の比較として公式のNMR選択を説明する。

\paragraph{問題の要旨} 成分の分離とNMRによる観測を区別する。

\Needspace{4\baselineskip}\paragraph{確認した条件・記述（独自要約）}\mbox{}\par

\noindent\textbf{原典確認：} S7\_049の26原典画像で全五肢・否定・添字を確認。\par

\noindent\textbf{対象：} 成分の分離とNMRによる観測を区別する。\par

\paragraph{問われる核心} 成分の分離とNMRによる観測を区別する。

\paragraph{必要知識} NMRは核の磁気共鳴を観測する分析法で、通常は成分を位置・経路別に分ける操作を行わない。遠心分離・電気泳動・HPLCは移動速度や保持の違いを利用する。MSは装置内でイオンをm/zで選別するという意味で分離を含む。候補番号の入替えを確認し、NMRを選ぶ。

\paragraph{考え方・正答を選ぶ理由（査読待ち初稿）} NMRは核の磁気共鳴を観測する分析法で、通常は成分を位置・経路別に分ける操作を行わない。遠心分離・電気泳動・HPLCは移動速度や保持の違いを利用する。MSは装置内でイオンをm/zで選別するという意味で分離を含む。候補番号の入替えを確認し、NMRを選ぶ。

\Needspace{4\baselineskip}\paragraph{選択肢の検討（査読待ち初稿）}\mbox{}\par

\noindent\textbf{1：不適切} 核磁気共鳴。共鳴周波数・化学シフト等を観測して構造を分析する。通常の装置は混合物を成分別に回収する分離器ではない。\par

\noindent\textbf{2：条件付き・解釈留保} 質量分析。イオンを質量電荷比で選別して測定する装置内の分離を含む。一方、元の中性混合物を大量に分取する手法とは異なる。\par

\noindent\textbf{3：通常の条件で適切} 遠心分離。回転場で密度・大きさに応じた沈降速度差を利用する。\par

\noindent\textbf{4：通常の条件で適切} 電気泳動。電荷と抵抗に依存する移動度差を利用して成分を分ける。\par

\noindent\textbf{5：通常の条件で適切} 高速液体クロマトグラフィー。固定相と移動相への分配等により保持時間が異なる成分を分ける。\par

\paragraph{再利用する考え方} 測定対象と分離対象、定義の分母、絶対量と相対校正、全称と成立条件を分けて判断する。

\paragraph{共通理論} \hyperref[theory:analysis:separation]{第\ref*{theory:analysis:separation}節}

\par 参考：\cite{analysis49_ms_ref}。

\paragraph{関連問題} \hyperref[q:h28-reg-applied-16]{H28-REG-APPLIED-16}

\par 難易度 4（校準初稿）、優先度未評価。

\subsection{平成28年度 問16}
\label{q:h28-reg-applied-16}

\noindent\textbf{原典：} sources/H28/17\_applied\_science.pdf、PDF物理6頁。

\textbf{公式正答：} 1（single）。

\par\textbf{検討未確定：} 質量分析でのイオン選別と、中性試料の調製・分取という分離の意味は同一でない。原典の広い成分分離の比較として公式のNMR選択を説明する。

\paragraph{問題の要旨} 成分の分離とNMRによる観測を区別する。

\Needspace{4\baselineskip}\paragraph{確認した条件・記述（独自要約）}\mbox{}\par

\noindent\textbf{原典確認：} S7\_049の26原典画像で全五肢・否定・添字を確認。\par

\noindent\textbf{対象：} 成分の分離とNMRによる観測を区別する。\par

\paragraph{問われる核心} 成分の分離とNMRによる観測を区別する。

\paragraph{必要知識} NMRは核の磁気共鳴を観測する分析法で、通常は成分を位置・経路別に分ける操作を行わない。遠心分離・電気泳動・HPLCは移動速度や保持の違いを利用する。MSは装置内でイオンをm/zで選別するという意味で分離を含む。候補番号の入替えを確認し、NMRを選ぶ。

\paragraph{考え方・正答を選ぶ理由（査読待ち初稿）} NMRは核の磁気共鳴を観測する分析法で、通常は成分を位置・経路別に分ける操作を行わない。遠心分離・電気泳動・HPLCは移動速度や保持の違いを利用する。MSは装置内でイオンをm/zで選別するという意味で分離を含む。候補番号の入替えを確認し、NMRを選ぶ。

\Needspace{4\baselineskip}\paragraph{選択肢の検討（査読待ち初稿）}\mbox{}\par

\noindent\textbf{1：不適切} 核磁気共鳴。共鳴周波数・化学シフト等を観測して構造を分析する。通常の装置は混合物を成分別に回収する分離器ではない。\par

\noindent\textbf{2：条件付き・解釈留保} 質量分析。イオンを質量電荷比で選別して測定する装置内の分離を含む。一方、元の中性混合物を大量に分取する手法とは異なる。\par

\noindent\textbf{3：通常の条件で適切} 遠心分離。回転場で密度・大きさに応じた沈降速度差を利用する。\par

\noindent\textbf{4：通常の条件で適切} 電気泳動。電荷と抵抗に依存する移動度差を利用して成分を分ける。\par

\noindent\textbf{5：通常の条件で適切} 高速液体クロマトグラフィー。固定相と移動相への分配等により保持時間が異なる成分を分ける。\par

\paragraph{再利用する考え方} 測定対象と分離対象、定義の分母、絶対量と相対校正、全称と成立条件を分けて判断する。

\paragraph{共通理論} \hyperref[theory:analysis:separation]{第\ref*{theory:analysis:separation}節}

\par 参考：\cite{analysis49_ms_ref}。

\paragraph{関連問題} \hyperref[q:h25-reg-applied-15]{H25-REG-APPLIED-15}, \hyperref[q:r01-reg-applied-12]{R01-REG-APPLIED-12}

\par 難易度 4（校準初稿）、優先度未評価。

\subsection{令和元年度 問12}
\label{q:r01-reg-applied-12}

\noindent\textbf{原典：} sources/R01/17\_applied\_science.pdf、PDF物理4頁。

\textbf{公式正答：} 4（single）。

\par\textbf{検討未確定：} 質量分析でのイオン選別と、中性試料の調製・分取という分離の意味は同一でない。原典の広い成分分離の比較として公式のNMR選択を説明する。

\paragraph{問題の要旨} 成分の分離とNMRによる観測を区別する。

\Needspace{4\baselineskip}\paragraph{確認した条件・記述（独自要約）}\mbox{}\par

\noindent\textbf{原典確認：} S7\_049の26原典画像で全五肢・否定・添字を確認。\par

\noindent\textbf{対象：} 成分の分離とNMRによる観測を区別する。\par

\paragraph{問われる核心} 成分の分離とNMRによる観測を区別する。

\paragraph{必要知識} NMRは核の磁気共鳴を観測する分析法で、通常は成分を位置・経路別に分ける操作を行わない。遠心分離・電気泳動・HPLCは移動速度や保持の違いを利用する。MSは装置内でイオンをm/zで選別するという意味で分離を含む。候補番号の入替えを確認し、NMRを選ぶ。

\paragraph{考え方・正答を選ぶ理由（査読待ち初稿）} NMRは核の磁気共鳴を観測する分析法で、通常は成分を位置・経路別に分ける操作を行わない。遠心分離・電気泳動・HPLCは移動速度や保持の違いを利用する。MSは装置内でイオンをm/zで選別するという意味で分離を含む。候補番号の入替えを確認し、NMRを選ぶ。

\Needspace{4\baselineskip}\paragraph{選択肢の検討（査読待ち初稿）}\mbox{}\par

\noindent\textbf{1：条件付き・解釈留保} 質量分析。イオンを質量電荷比で選別して測定する装置内の分離を含む。一方、元の中性混合物を大量に分取する手法とは異なる。\par

\noindent\textbf{2：通常の条件で適切} 電気泳動。電荷と抵抗に依存する移動度差を利用して成分を分ける。\par

\noindent\textbf{3：通常の条件で適切} 遠心分離。回転場で密度・大きさに応じた沈降速度差を利用する。\par

\noindent\textbf{4：不適切} 核磁気共鳴。共鳴周波数・化学シフト等を観測して構造を分析する。通常の装置は混合物を成分別に回収する分離器ではない。\par

\noindent\textbf{5：通常の条件で適切} 高速液体クロマトグラフィー。固定相と移動相への分配等により保持時間が異なる成分を分ける。\par

\paragraph{再利用する考え方} 測定対象と分離対象、定義の分母、絶対量と相対校正、全称と成立条件を分けて判断する。

\paragraph{共通理論} \hyperref[theory:analysis:separation]{第\ref*{theory:analysis:separation}節}

\par 参考：\cite{analysis49_ms_ref}。

\paragraph{関連問題} \hyperref[q:h28-reg-applied-16]{H28-REG-APPLIED-16}

\par 難易度 4（校準初稿）、優先度未評価。

\section{多元素同時分析に向く方法を比較する。}

\subsection{平成24年度 問16}
\label{q:h24-reg-applied-16}

\noindent\textbf{原典：} sources/H24/17\_applied\_science.pdf、PDF物理5頁。

\textbf{公式正答：} 5（single）。

\paragraph{問題の要旨} 多元素同時分析に向く方法を比較する。

\Needspace{4\baselineskip}\paragraph{確認した条件・記述（独自要約）}\mbox{}\par

\noindent\textbf{原典確認：} S7\_049の26原典画像で全五肢・否定・添字を確認。\par

\noindent\textbf{対象：} 多元素同時分析に向く方法を比較する。\par

\paragraph{問われる核心} 多元素同時分析に向く方法を比較する。

\paragraph{必要知識} 元素の種類を区別する観測量を探す。特性X線のエネルギー又は波長を検出するXRFなら、複数元素のピークを識別できる。各滴定の総当量だけでは混合元素を一度に分解できない。装置方式による同時・逐次検出の差はあるが、五方法の比較では5を選ぶ。

\paragraph{考え方・正答を選ぶ理由（査読待ち初稿）} 元素の種類を区別する観測量を探す。特性X線のエネルギー又は波長を検出するXRFなら、複数元素のピークを識別できる。各滴定の総当量だけでは混合元素を一度に分解できない。装置方式による同時・逐次検出の差はあるが、五方法の比較では5を選ぶ。

\Needspace{4\baselineskip}\paragraph{選択肢の検討（査読待ち初稿）}\mbox{}\par

\noindent\textbf{1：不適切} 酸化還元滴定。当量関係を使う濃度測定で、複数元素の特性線を一度に識別する方法ではない。\par

\noindent\textbf{2：不適切} 赤外吸収。主に振動と官能基を調べる。通常の多元素同時定量の第一選択ではない。\par

\noindent\textbf{3：不適切} キレート滴定。金属イオンと錯形成する試薬の当量を使い、各元素を同時に独立識別する方式ではない。\par

\noindent\textbf{4：不適切} 吸光光度。特定の吸収帯や呈色反応を用いる。ここでの全元素比較ではXRFが適する。\par

\noindent\textbf{5：通常の条件で適切} 蛍光X線。元素ごとの内殻遷移エネルギーを分け、複数の元素の信号を一分析で得られる。\par

\paragraph{再利用する考え方} 測定対象と分離対象、定義の分母、絶対量と相対校正、全称と成立条件を分けて判断する。

\paragraph{共通理論} \hyperref[theory:analysis:xray]{第\ref*{theory:analysis:xray}節}

\par 参考：\cite{analysis49_xrf_ref}。

\par 難易度 3（校準初稿）、優先度未評価。

\subsection{令和3年度 問16}
\label{q:r03-reg-applied-16}

\noindent\textbf{原典：} sources/R03/17\_applied\_science.pdf、PDF物理6頁。

\textbf{公式正答：} 3（single）。

\paragraph{問題の要旨} 内殻光電子の化学シフトと核・電子スピンを比較する。

\Needspace{4\baselineskip}\paragraph{確認した条件・記述（独自要約）}\mbox{}\par

\noindent\textbf{原典確認：} S7\_049の26原典画像で全五肢・否定・添字を確認。\par

\noindent\textbf{対象：} 内殻光電子の化学シフトと核・電子スピンを比較する。\par

\paragraph{問われる核心} 内殻光電子の化学シフトと核・電子スピンを比較する。

\paragraph{必要知識} 内殻の信号だから結合状態に無関係、という推論を避ける。XPSの結合エネルギーには化学シフトがあり3が不適切。励起X線はE=\allowbreak{}hc/λなので短波長ほど高エネルギーで、吸収端を超える内殻励起が蛍光X線を生む。

\paragraph{考え方・正答を選ぶ理由（査読待ち初稿）} 内殻の信号だから結合状態に無関係、という推論を避ける。XPSの結合エネルギーには化学シフトがあり3が不適切。励起X線はE=\allowbreak{}hc/λなので短波長ほど高エネルギーで、吸収端を超える内殻励起が蛍光X線を生む。

\Needspace{4\baselineskip}\paragraph{選択肢の検討（査読待ち初稿）}\mbox{}\par

\noindent\textbf{1：通常の条件で適切} ESRと不対電子。不対電子の磁気共鳴を測る。\par

\noindent\textbf{2：通常の条件で適切} 13Cはスピン1/2でNMR対象。13Cは核スピン1/2を持つ。豊富な12Cは核スピン0なので同じ測定対象ではない。\par

\noindent\textbf{3：不適切} XPSは結合状態に鈍感で化学分析不向き。周辺結合状態による内殻結合エネルギーの化学シフトを測定でき、元素と化学状態の分析に用いられる。\par

\noindent\textbf{4：通常の条件で適切} Ramanで振動・回転を調べる。非弾性散乱の周波数差が振動・回転遷移に対応する。\par

\noindent\textbf{5：通常の条件で適切} 吸収端より短いX線でXRF。光子エネルギーが内殻電離の閾値を超えると空孔が作られ、その後の遷移から特性X線を放射し得る。\par

\paragraph{再利用する考え方} 測定対象と分離対象、定義の分母、絶対量と相対校正、全称と成立条件を分けて判断する。

\paragraph{共通理論} \hyperref[theory:analysis:xray]{第\ref*{theory:analysis:xray}節}

\par 参考：\cite{analysis49_xps_ref,analysis49_esr_ref,atomic_ir_raman}。

\par 難易度 3（校準初稿）、優先度未評価。

\section{陽イオンの系統分析}

\subsection{平成26年度 問16}
\label{q:h26-reg-applied-16}

\noindent\textbf{原典：} sources/H26/17\_applied\_science.pdf、PDF物理7頁。

\textbf{公式正答：} 4（single）。

\par\textbf{検討未確定：} 酸性度・濃度・錯形成の条件は出題の操作列に依存し一般条件へ無限定に拡張しない。

\paragraph{問題の要旨} 塩化物、酸性下の硫化物、水酸化物を順に除いたろ液から、最後に硫化物として分離する陽イオンを特定する。

\Needspace{4\baselineskip}\paragraph{確認した条件・記述（独自要約）}\mbox{}\par

\noindent\textbf{I：} 塩酸を滴下し沈殿をろ別\par

\noindent\textbf{II：} Iのろ液中でチオアセトアミドCH3CSNH2の加水分解を行い、生成した沈殿をろ別\par

\noindent\textbf{III：} IIのろ液中へアンモニア水を滴下し、生成した沈殿をろ別\par

\noindent\textbf{IV：} IIIのろ液中へ硫化水素を滴下し、生成した沈殿をろ別\par

\paragraph{問われる核心} 分析化学

\paragraph{必要知識} 系統的定性分析、溶解度積、硫化水素の酸解離、pHと沈殿選択性

\paragraph{考え方・正答を選ぶ理由（査読待ち初稿）} 手順IでAg+をAgClとして除く。手順IIの酸性条件でCu2+が難溶性の硫化物となる。手順IIIでFe3+を水酸化物として除く。アンモニア性のろ液に硫化水素を加える手順IVではZn2+をZnSとして分離するので④。Na+はこれらの分離後も水溶液に残る。これは出題された古典的系統分析の条件に基づく説明で、実験条件一般への無条件な適用はしない。

\Needspace{4\baselineskip}\paragraph{選択肢の検討（査読待ち初稿）}\mbox{}\par

\noindent\textbf{1：不適切} Na+。操作I〜IVでの難溶沈殿として分離されず水溶液に残る。\par

\noindent\textbf{2：不適切} Ag+。Iの塩酸でAgClとして先に除く。\par

\noindent\textbf{3：不適切} Cu2+。IIの酸性硫化物の段階で先に除く。\par

\noindent\textbf{4：適切・選択対象} Zn2+。III後のアンモニア性ろ液にH2Sを加えるIVでZnSとして分離。\par

\noindent\textbf{5：不適切} Fe3+。IIIで水酸化物として先に除く。\par

\paragraph{再利用する考え方} 系統的定性分析、溶解度積、硫化水素の酸解離、pHと沈殿選択性

\paragraph{共通理論} \hyperref[theory:analytical:separation]{第\ref*{theory:analytical:separation}節}

\par 参考：\cite{qualitative_stanford}。

\par 難易度 3（校準初稿）、優先度未評価。


\chapter{単位・計測}
\section{背景：量の意味と数値を共有するための単位}
観測量を数値で比較するには、同じ物理量を同じ尺度で表す必要がある。単位の定義は尺度の基準を定め、実現は実験でその尺度へ接続する方法である。定義・物理法則・測定手順を同じものとして扱わない。
SIは長さm、質量kg、時間s、電流A、熱力学温度K、物質量mol、光度cdの7基本単位を持つ。平成23・27年度の原典の導入はcdを欠く6単位の列挙であり、光度を含む全物理量をその6単位だけで組み立てられるという説明は保持できない。誤った選択肢の判定と、この導入の疑義を別に記録する\cite{si_brochure9}。
本章は対象11問の原典指数・条件を確認した初稿である。正答を導けても、2導入と改定後の三重点の厳密さ1問には疑義を残す。Astra本文査読前である。

\section{量方程式と数値方程式}\label{theory:units:quantity}
量は「数値×単位」である。量方程式 \(F=ma\) は単位の選択に依存しないが、数値だけの式は単位を揃える必要がある。
量の次元を \(M,L,T,I\) 等で表せば、力は \(MLT^{-2}\)、エネルギーは \(ML^2T^{-2}\) である。加算する量は同次元でなければならず、指数・対数の引数は無次元である。
次元一致は必要条件であって十分条件ではない。同じ次元を持つ仕事とトルク、周波数と放射能でも物理的意味は異なる。定数・係数・方向や対象を次元だけで確定しない。

\section{力学量：力から圧力・仕事・仕事率へ}\label{theory:units:derived}
運動方程式、面積への力、仕事、単位時間の仕事から順に
\[
 [F]=\mathrm{kg\,m\,s^{-2}}=\mathrm N,\qquad
 [p]=[F]/[S]=\mathrm{kg\,m^{-1}s^{-2}}=\mathrm{Pa},
\]
\[
 [E]=[F][\ell]=\mathrm{kg\,m^2s^{-2}}=\mathrm J,\qquad
 [P]=[E]/[t]=\mathrm{kg\,m^2s^{-3}}=\mathrm W
\]
となる。仕事は一般には \(\int\boldsymbol F\cdot d\boldsymbol r\) であるため、単純積は一定力とその方向の移動距離の場合を表す。次元の導出には同じ長さ因子を使う。
平成27年度のWの長さ指数1は、仕事を作る移動距離を一つ欠く。Paは面積で割るのでmの指数が \(-1\) になる。力の単位をそのまま圧力へ使わない\cite{si_nist811}。

\section{電気量：電荷・電圧・抵抗・磁束密度}\label{theory:units:electrical}
一定電流なら \(Q=It\)、一般には \(Q=\int I\,dt\) でありC=Asである。電位差は電荷当たりの仕事なので
\[
 [V]=[\text{仕事}]/[Q]=\mathrm{kg\,m^2s^{-3}A^{-1}},\qquad
 [R]=[V]/[I]=\mathrm{kg\,m^2s^{-3}A^{-2}}.
\]
原典のVのm指数1やΩのA指数 \(-1\) はこの接続を欠く。電場は力/電荷でN/C、電位差はJ/Cであり、両者は長さ一つ違う。
一様磁場に垂直な直線電流の磁気力 \(F=I\ell B\) から
\[
 [B]=\mathrm{N/(A\,m)}=\mathrm{kg\,s^{-2}A^{-1}}=\mathrm T
\]
である。テスラへ不要なmを残したりAを二回で割ったりしない。単位はオームの法則の適用材料・一定温度等の条件を保証するものではない\cite{si_nist811}。

\section{周波数・角度・放射能の意味}\label{theory:units:angle}\label{theory:units:activity}
周期運動の周波数は \(f=1/T\) でHz=\(\mathrm{s^{-1}}\)、平面角は弧長/半径でradの次元は1である。角の意味を残すためradを用い、角周波数 \(\omega\) と周期周波数 \(f\) を \(\omega=2\pi f\) で区別する。
放射能は核壊変数の時間率でBq=\(\mathrm{s^{-1}}\) である。HzとBqは同次元だが、周期数と壊変数の意味は異なる。線量GyやSvに変換するには壊変ごとのエネルギー、吸収先・質量等の追加情報が必要である。
仕事とトルクにも同じ次元が現れる。数値・次元が一致しても、何を測る量かを読み落とさない\cite{si_brochure9}。

\section{定義の改定：原器から固定した定数へ}\label{theory:units:history}
単位を人工物へ結び付けると、長期変化や各国への伝達に限界がある。再現可能な自然定数と原子の周波数へ結び付けることで、各地の実験から一つの尺度へ接続できる。
\begin{tcolorbox}[title=出題時の定義を固定する]
平成28・29年度は2019年5月20日より前である。kgは国際キログラム原器、Kは規定水の三重点、molは0.012kgの炭素12中の原子数と等しい数の指定要素粒子に基づいた。これらの出題時の適切肢を、現在の定義だけで誤りへ変えない。
\end{tcolorbox}
令和元年度再試験は問題幹が2019年5月の改定後を明記する。原器によるkg定義はそこで不適切となる。改定によって既存の測定値が一斉に別の数値へ跳ぶように設計されたわけではない。連続性を保ちつつ、以前固定されていた物質の量に測定不確かさが移る\cite{si_cgpm2018,si_nist811}。

\section{7基本単位と定義定数}\label{theory:units:definition}
現行SIの定義定数との接続は次の通りである。定数と単位が互いに孤立した一対一の定義ではなく、既に固定した時間・長さ等を組み合わせる点に注意する。
\begin{center}
\begin{tabular}{lll}
量・基本単位 & 主な定義定数 & 固定する数値\\\hline
時間 s & Cs133超微細遷移周波数 & \(9192631770\,\mathrm{Hz}\)\\
長さ m & 真空光速 \(c\) & \(299792458\,\mathrm{m/s}\)\\
質量 kg & プランク定数 \(h\) & \(6.62607015\times10^{-34}\,\mathrm{J\,s}\)\\
電流 A & 素電荷 \(e\) & \(1.602176634\times10^{-19}\,\mathrm C\)\\
熱力学温度 K & ボルツマン定数 \(k\) & \(1.380649\times10^{-23}\,\mathrm{J/K}\)\\
物質量 mol & アボガドロ定数 \(N_{\rm A}\) & \(6.02214076\times10^{23}\,\mathrm{mol^{-1}}\)\\
光度 cd & 視感効果度 \(K_{\rm cd}\) & \(683\,\mathrm{lm/W}\)（540THz）\\
\end{tabular}
\end{center}
たとえば \(h\) の次元は \(\mathrm{kg\,m^2/s}\) であるため、mとsの定義を踏まえてkgの尺度を得る。重力定数 \(G\) をkgの定義に固定しているわけではない。kgを天秤の重さ測定へ実現する際に重力加速度を使うこととは区別する。
電流には \(e\) と秒、温度には \(k\) とエネルギー尺度、物質量には指定粒子数を使う。今回の6定義問題は、改定の時点又はこの対応の取り違えを問う\cite{si_cgpm2018}。

\section{水の三重点・凝固点・Celsius・温度尺度}\label{theory:units:temperature}
三重点は固体・液体・気体の三相が平衡にある温度と圧力の点である。通常圧力での凝固点とは違う。Celsius温度の数値は \(t/{}^\circ\mathrm C=T/\mathrm K-273.15\) である。
旧K定義では規定同位体組成の水の三重点は厳密に273.16K、したがって0.01℃であった。水組成・圧力の条件を省き、すべての水の凝固がその値になるとしない\cite{si_kelvin2007}。
\begin{tcolorbox}[title=改定後の三重点は熱力学温度の厳密定義値ではない]
2019年後のKは \(k\) の固定で定義する。BIPM実現文書は三重点の最良推定273.16K、標準不確かさ0.1mKを示す。実用温度尺度ITS-90の \(T_{90}\) の固定点と、熱力学温度 \(T\) の厳密定義を区別する。令和元年度再試験の0.01℃は提示桁の近似として理解できるが、厳密な定義値という読みに疑義を残す。
\end{tcolorbox}
「値がほぼ同じ」と「定義が同じ」を混同しない\cite{si_kelvin_realization}。

\section{光度：黒体基準と単色放射の尺度}\label{theory:units:photometry}
放射エネルギーの尺度だけでは人の視覚に対する明るさを同じ数値で表せない。周波数による視感度を重み付けした測光量が必要になる。光度cdは方向当たりの光束であり、エネルギー流のWや面積当たりの照度lxとは異なる。
現在は周波数 \(540\times10^{12}\,\mathrm{Hz}\) の単色放射に対する \(K_{\rm cd}=683\,\mathrm{lm/W}\) を固定する。その放射の放射強度が \(1/683\,\mathrm{W/sr}\) の方向の光度が1cdに対応する。
黒体を使った歴史的定義と現在の定義、黒体光源を使う実験上の実現は別である。令和3年度の「黒体で定義」は不適切、令和7年度の「単色放射で定義」はこの固定周波数・視感度の省略として適切である\cite{si_brochure9}。

\section{換算：面積指数・力と質量・動粘度}\label{theory:units:conversion}
単位を換えるときは量を変えず、長さ換算を指数の回数だけ適用する。\(1\,\mathrm{cm^2}=10^{-4}\,\mathrm{m^2}\) であり \(10^{-2}\) ではない。
kgfは標準重力加速度 \(g_{\rm n}=9.80665\,\mathrm{m/s^2}\) により \(1\,\mathrm{kgf}=9.80665\,\mathrm N\) とした力の単位であり、kgは質量である。したがって
\[
 1\,\mathrm{kgf/cm^2}=\frac{9.80665}{10^{-4}}\,\mathrm{N/m^2}
 =9.80665\times10^4\,\mathrm{Pa}.
\]
標準atmは101325Pa、barは \(10^5\) Pa、通常のmmHgは約133.322Pa、psiは約6894.757Paである。液柱の実際の圧力 \(\rho g h\) は流体の温度や重力条件に依存する。通常の換算表の係数と実測条件を区別する。
動粘度 \(\nu=\mu/\rho\) の次元は \(L^2/T\)。Stは \(\mathrm{cm^2/s}\) なので \(10^{-4}\,\mathrm{m^2/s}\) となる。非SI単位であっても、そのSIへの換算が正しい場合がある\cite{si_nist811}。

\section{今回の問われ方と学習順序}
主概念は組立単位・換算5問と定義・改定6問である。まず物理量の量方程式からm・s・Aの指数を導き、次に定義の基準日を置き、最後に非SI係数を照合する。単位を丸暗記すると、原器と三重点の改定や電場と電位の違いを取りこぼす。
この集計は本章11問の確認結果であり将来の出題保証ではない。疑義の3問も含めて全候補を追跡する。

\section{過去問と全候補の検討}
\subsection{平成23年度 問6}
\label{q:h23-reg-applied-06}

\noindent\textbf{原典：} sources/H23/17\_applied\_science.pdf、PDF物理3頁。

\textbf{公式正答：} 3（single）。

\par\textit{原典本文の校正と独自解説は未作成。分類のみの登録項目。}

\subsection{平成25年度 問4}
\label{q:h25-reg-applied-04}

\noindent\textbf{原典：} sources/H25/17\_applied\_science.pdf、PDF物理3頁。

\textbf{公式正答：} 5（single）。

\par\textit{原典本文の校正と独自解説は未作成。分類のみの登録項目。}

\subsection{平成27年度 問4}
\label{q:h27-reg-applied-04}

\noindent\textbf{原典：} sources/H27/17\_applied\_science.pdf、PDF物理2頁。

\textbf{公式正答：} 4（single）。

\par\textit{原典本文の校正と独自解説は未作成。分類のみの登録項目。}

\subsection{平成28年度 問6}
\label{q:h28-reg-applied-06}

\noindent\textbf{原典：} sources/H28/17\_applied\_science.pdf、PDF物理2頁。

\textbf{公式正答：} 4（single）。

\par\textit{原典本文の校正と独自解説は未作成。分類のみの登録項目。}

\subsection{平成29年度 問5}
\label{q:h29-reg-applied-05}

\noindent\textbf{原典：} sources/H29/17\_applied\_science.pdf、PDF物理2頁。

\textbf{公式正答：} 4（single）。

\par\textit{原典本文の校正と独自解説は未作成。分類のみの登録項目。}

\subsection{令和元年度再試験 問1}
\label{q:r01-ret-applied-01}

\noindent\textbf{原典：} sources/R01\_retest/17\_applied\_science.pdf、PDF物理1頁。

\textbf{公式正答：} 1（single）。

\par\textit{原典本文の校正と独自解説は未作成。分類のみの登録項目。}

\subsection{令和2年度 問9}
\label{q:r02-reg-applied-09}

\noindent\textbf{原典：} sources/R02/17\_applied\_science.pdf、PDF物理3頁。

\textbf{公式正答：} 2（single）。

\par\textit{原典本文の校正と独自解説は未作成。分類のみの登録項目。}

\subsection{令和3年度 問9}
\label{q:r03-reg-applied-09}

\noindent\textbf{原典：} sources/R03/17\_applied\_science.pdf、PDF物理4頁。

\textbf{公式正答：} 1（single）。

\par\textit{原典本文の校正と独自解説は未作成。分類のみの登録項目。}

\subsection{令和5年度 問13}
\label{q:r05-reg-applied-13}

\noindent\textbf{原典：} sources/R05/17\_applied\_science.pdf、PDF物理5頁。

\textbf{公式正答：} 5（single）。

\par\textit{原典本文の校正と独自解説は未作成。分類のみの登録項目。}

\subsection{令和6年度 問16}
\label{q:r06-reg-applied-16}

\noindent\textbf{原典：} sources/R06/17\_applied\_science.pdf、PDF物理5頁。

\textbf{公式正答：} 3（single）。

\par\textit{原典本文の校正と独自解説は未作成。分類のみの登録項目。}

\subsection{令和7年度 問9}
\label{q:r07-reg-applied-09}

\noindent\textbf{原典：} sources/R07/17\_applied\_science.pdf、PDF物理3頁。

\textbf{公式正答：} 3（single）。

\par\textit{原典本文の校正と独自解説は未作成。分類のみの登録項目。}


\chapter{有機・高分子化学}
\noindent 分野解説は執筆中。以下は登録問題と校正済み例の作業用一覧。
\subsection{平成23年度 問9}
\label{q:h23-reg-applied-09}

\noindent\textbf{原典：} sources/H23/17\_applied\_science.pdf、PDF物理4頁。

\textbf{公式正答：} 1（single）。

\par\textit{原典本文の校正と独自解説は未作成。分類のみの登録項目。}

\subsection{平成23年度 問14}
\label{q:h23-reg-applied-14}

\noindent\textbf{原典：} sources/H23/17\_applied\_science.pdf、PDF物理5頁。

\textbf{公式正答：} 4（single）。

\par\textit{原典本文の校正と独自解説は未作成。分類のみの登録項目。}

\subsection{平成23年度 問15}
\label{q:h23-reg-applied-15}

\noindent\textbf{原典：} sources/H23/17\_applied\_science.pdf、PDF物理6頁。

\textbf{公式正答：} 2（single）。

\par\textit{原典本文の校正と独自解説は未作成。分類のみの登録項目。}

\subsection{平成23年度 問16}
\label{q:h23-reg-applied-16}

\noindent\textbf{原典：} sources/H23/17\_applied\_science.pdf、PDF物理6頁。

\textbf{公式正答：} 3（single）。

\par\textit{原典本文の校正と独自解説は未作成。分類のみの登録項目。}

\subsection{平成24年度 問15}
\label{q:h24-reg-applied-15}

\noindent\textbf{原典：} sources/H24/17\_applied\_science.pdf、PDF物理5頁。

\textbf{公式正答：} 4（single）。

\par\textit{原典本文の校正と独自解説は未作成。分類のみの登録項目。}

\subsection{平成24年度 問18}
\label{q:h24-reg-applied-18}

\noindent\textbf{原典：} sources/H24/17\_applied\_science.pdf、PDF物理6頁。

\textbf{公式正答：} 3（single）。

\par\textit{原典本文の校正と独自解説は未作成。分類のみの登録項目。}

\subsection{平成25年度 問14}
\label{q:h25-reg-applied-14}

\noindent\textbf{原典：} sources/H25/17\_applied\_science.pdf、PDF物理6頁。

\textbf{公式正答：} 1（single）。

\par\textit{原典本文の校正と独自解説は未作成。分類のみの登録項目。}

\subsection{平成25年度 問17}
\label{q:h25-reg-applied-17}

\noindent\textbf{原典：} sources/H25/17\_applied\_science.pdf、PDF物理7頁。

\textbf{公式正答：} 4（single）。

\par\textit{原典本文の校正と独自解説は未作成。分類のみの登録項目。}

\subsection{平成26年度 問17}
\label{q:h26-reg-applied-17}

\noindent\textbf{原典：} sources/H26/17\_applied\_science.pdf、PDF物理7頁。

\textbf{公式正答：} 2（single）。

\par\textit{原典本文の校正と独自解説は未作成。分類のみの登録項目。}

\subsection{平成27年度 問13}
\label{q:h27-reg-applied-13}

\noindent\textbf{原典：} sources/H27/17\_applied\_science.pdf、PDF物理5頁。

\textbf{公式正答：} 2（single）。

\par\textit{原典本文の校正と独自解説は未作成。分類のみの登録項目。}

\subsection{平成27年度 問14}
\label{q:h27-reg-applied-14}

\noindent\textbf{原典：} sources/H27/17\_applied\_science.pdf、PDF物理5頁。

\textbf{公式正答：} 4（single）。

\par\textit{原典本文の校正と独自解説は未作成。分類のみの登録項目。}

\subsection{平成27年度 問16}
\label{q:h27-reg-applied-16}

\noindent\textbf{原典：} sources/H27/17\_applied\_science.pdf、PDF物理6頁。

\textbf{公式正答：} 2（single）。

\par\textit{原典本文の校正と独自解説は未作成。分類のみの登録項目。}

\subsection{平成28年度 問15}
\label{q:h28-reg-applied-15}

\noindent\textbf{原典：} sources/H28/17\_applied\_science.pdf、PDF物理6頁。

\textbf{公式正答：} 5（single）。

\par\textit{原典本文の校正と独自解説は未作成。分類のみの登録項目。}

\subsection{平成28年度 問17}
\label{q:h28-reg-applied-17}

\noindent\textbf{原典：} sources/H28/17\_applied\_science.pdf、PDF物理6頁。

\textbf{公式正答：} 4（single）。

\par\textit{原典本文の校正と独自解説は未作成。分類のみの登録項目。}

\subsection{平成29年度 問15}
\label{q:h29-reg-applied-15}

\noindent\textbf{原典：} sources/H29/17\_applied\_science.pdf、PDF物理5頁。

\textbf{公式正答：} 2（single）。

\par\textit{原典本文の校正と独自解説は未作成。分類のみの登録項目。}

\subsection{平成29年度 問16}
\label{q:h29-reg-applied-16}

\noindent\textbf{原典：} sources/H29/17\_applied\_science.pdf、PDF物理5頁。

\textbf{公式正答：} 1（single）。

\par\textit{原典本文の校正と独自解説は未作成。分類のみの登録項目。}

\subsection{平成29年度 問17}
\label{q:h29-reg-applied-17}

\noindent\textbf{原典：} sources/H29/17\_applied\_science.pdf、PDF物理5頁。

\textbf{公式正答：} 2（single）。

\par\textit{原典本文の校正と独自解説は未作成。分類のみの登録項目。}

\subsection{平成30年度 問15}
\label{q:h30-reg-applied-15}

\noindent\textbf{原典：} sources/H30/17\_applied\_science.pdf、PDF物理5頁。

\textbf{公式正答：} 5（single）。

\par\textit{原典本文の校正と独自解説は未作成。分類のみの登録項目。}

\subsection{令和元年度 問11}
\label{q:r01-reg-applied-11}

\noindent\textbf{原典：} sources/R01/17\_applied\_science.pdf、PDF物理4頁。

\textbf{公式正答：} 3（single）。

\par\textit{原典本文の校正と独自解説は未作成。分類のみの登録項目。}

\subsection{令和元年度 問13}
\label{q:r01-reg-applied-13}

\noindent\textbf{原典：} sources/R01/17\_applied\_science.pdf、PDF物理4頁。

\textbf{公式正答：} 3（single）。

\par\textit{原典本文の校正と独自解説は未作成。分類のみの登録項目。}

\subsection{令和元年度再試験 問11}
\label{q:r01-ret-applied-11}

\noindent\textbf{原典：} sources/R01\_retest/17\_applied\_science.pdf、PDF物理4頁。

\textbf{公式正答：} 4（single）。

\par\textit{原典本文の校正と独自解説は未作成。分類のみの登録項目。}

\subsection{令和元年度再試験 問13}
\label{q:r01-ret-applied-13}

\noindent\textbf{原典：} sources/R01\_retest/17\_applied\_science.pdf、PDF物理5頁。

\textbf{公式正答：} 1（single）。

\par\textit{原典本文の校正と独自解説は未作成。分類のみの登録項目。}

\subsection{令和元年度再試験 問17}
\label{q:r01-ret-applied-17}

\noindent\textbf{原典：} sources/R01\_retest/17\_applied\_science.pdf、PDF物理6頁。

\textbf{公式正答：} 1（single）。

\par\textit{原典本文の校正と独自解説は未作成。分類のみの登録項目。}

\subsection{令和2年度 問15}
\label{q:r02-reg-applied-15}

\noindent\textbf{原典：} sources/R02/17\_applied\_science.pdf、PDF物理5頁。

\textbf{公式正答：} 2（single）。

\par\textit{原典本文の校正と独自解説は未作成。分類のみの登録項目。}

\subsection{令和2年度 問18}
\label{q:r02-reg-applied-18}

\noindent\textbf{原典：} sources/R02/17\_applied\_science.pdf、PDF物理6頁。

\textbf{公式正答：} 3（single）。

\par\textit{原典本文の校正と独自解説は未作成。分類のみの登録項目。}

\subsection{令和3年度 問12}
\label{q:r03-reg-applied-12}

\noindent\textbf{原典：} sources/R03/17\_applied\_science.pdf、PDF物理5頁。

\textbf{公式正答：} 2（single）。

\par\textit{原典本文の校正と独自解説は未作成。分類のみの登録項目。}

\subsection{令和4年度 問16}
\label{q:r04-reg-applied-16}

\noindent\textbf{原典：} sources/R04/17\_applied\_science.pdf、PDF物理5頁。

\textbf{公式正答：} 4（single）。

\par\textit{原典本文の校正と独自解説は未作成。分類のみの登録項目。}

\subsection{令和4年度 問18}
\label{q:r04-reg-applied-18}

\noindent\textbf{原典：} sources/R04/17\_applied\_science.pdf、PDF物理6頁。

\textbf{公式正答：} 4（single）。

\par\textit{原典本文の校正と独自解説は未作成。分類のみの登録項目。}

\subsection{令和5年度 問11}
\label{q:r05-reg-applied-11}

\noindent\textbf{原典：} sources/R05/17\_applied\_science.pdf、PDF物理4頁。

\textbf{公式正答：} 4（single）。

\par\textit{原典本文の校正と独自解説は未作成。分類のみの登録項目。}

\subsection{令和5年度 問14}
\label{q:r05-reg-applied-14}

\noindent\textbf{原典：} sources/R05/17\_applied\_science.pdf、PDF物理5頁。

\textbf{公式正答：} 5（single）。

\par\textit{原典本文の校正と独自解説は未作成。分類のみの登録項目。}

\subsection{令和6年度 問17}
\label{q:r06-reg-applied-17}

\noindent\textbf{原典：} sources/R06/17\_applied\_science.pdf、PDF物理5頁。

\textbf{公式正答：} 3（single）。

\par\textit{原典本文の校正と独自解説は未作成。分類のみの登録項目。}

\subsection{令和7年度 問11}
\label{q:r07-reg-applied-11}

\noindent\textbf{原典：} sources/R07/17\_applied\_science.pdf、PDF物理4頁。

\textbf{公式正答：} 5（single）。

\par\textit{原典本文の校正と独自解説は未作成。分類のみの登録項目。}

\subsection{令和7年度 問12}
\label{q:r07-reg-applied-12}

\noindent\textbf{原典：} sources/R07/17\_applied\_science.pdf、PDF物理4頁。

\textbf{公式正答：} 2（single）。

\par\textit{原典本文の校正と独自解説は未作成。分類のみの登録項目。}


\chapter{材料物性}
\noindent 分野解説は執筆中。以下は登録問題と校正済み例の作業用一覧。
\section{金属の強化機構を整理する}
\subsection{背景}
塑性変形には転位の運動が関わる。加工、固溶、析出、結晶粒の変化を、単に「金属」という語だけで無機化学へ一括分類しない。
\subsection{原典と正しい説明を分ける}
R02の例の⑤は、原典に意図的な誤答として「バウシンガー変態」と書かれている。
校正時にこの語をマルテンサイトへ置換すると、誤り選択問題の意味を失う。
解説では鋼の焼入れに関わるマルテンサイト変態を示し\cite{martensite_cambridge}、原典の文言と区別する。

\section{金属の強化機構}

\subsection{令和2年度 問8}
\label{q:r02-reg-applied-08}

\noindent\textbf{原典：} sources/R02/17\_applied\_science.pdf、PDF物理3頁。

\textbf{公式正答：} 5（single）。

\paragraph{問題の要旨} 金属の強化機構を識別し、焼入れに関する変態名の誤りを見抜く。

\Needspace{4\baselineskip}\paragraph{確認した条件・記述（独自要約）}\mbox{}\par

\noindent\textbf{原典：} 指定頁を実画像で再確認。\par

\paragraph{問われる核心} 材料物性

\paragraph{必要知識} 転位、加工硬化、固溶・析出・粒界強化、マルテンサイト変態

\paragraph{考え方・正答を選ぶ理由（査読待ち初稿）} 鋼の焼入れ硬化で中心となるのはマルテンサイト変態である。⑤の「バウシンガー変態」は原典の誤答として保存し、解説中で正しい用語と区別する。したがって公式正答は⑤。

\Needspace{4\baselineskip}\paragraph{選択肢の検討（査読待ち初稿）}\mbox{}\par

\noindent\textbf{1：記述と範囲を確認} 加工硬化。塑性変形に伴う転位等により変形抵抗が増す。\par

\noindent\textbf{2：記述と範囲を確認} 固溶強化。溶質原子が格子・転位へ影響し強化する。\par

\noindent\textbf{3：記述と範囲を確認} 析出強化。固溶限の温度依存を利用した析出が転位運動を妨げる。\par

\noindent\textbf{4：記述と範囲を確認} 結晶粒微細化。粒界が転位の移動を妨げ通常降伏強さが増す。極端な粒径や温度まで無限定にしない。\par

\noindent\textbf{5：不適切・独立選択} 焼入れはバウシンガー変態。焼入れ硬化はマルテンサイト変態。Bauschingerは逆方向の降伏挙動に関する効果で焼入れ変態名ではない。\par

\paragraph{再利用する考え方} 転位、加工硬化、固溶・析出・粒界強化、マルテンサイト変態

\paragraph{共通理論} \hyperref[theory:material:strengthening]{第\ref*{theory:material:strengthening}節}

\par 難易度 2（校準初稿）、優先度未評価。


\chapter{地震学}
\noindent 原典条件と共通理論を結び付ける解説初稿。各問の成立条件と参照先を確認する。
\section{内陸地殻内地震}

\subsection{平成26年度 問32}
\label{q:h26-reg-applied-32}

\noindent\textbf{原典：} sources/H26/17\_applied\_science.pdf、PDF物理13頁。

\textbf{公式正答：} 5（single）。

\par\textbf{検討未確定：} 正答5の「最も活発」「約1000年」は母集団と平均期間の定義が不足する。公式5と独立初回5は保持し、厳密な下限の主張は未解決とする。

\paragraph{問題の要旨} 内陸地殻内地震の発生位置・規模・深さと、同じ活断層の再来間隔を区別する。

\Needspace{4\baselineskip}\paragraph{確認した条件・記述（独自要約）}\mbox{}\par

\noindent\textbf{確認範囲：} 問題幹、数値、単位、否定、全5候補を原典画像で照合した。\par

\noindent\textbf{成立条件：} 内陸という語を震央が必ず陸上という意味に限定しない。気象庁マグニチュードとモーメントマグニチュードを混ぜない。再来間隔は全国の地震間隔でなく一つの断層である。\par

\paragraph{問われる核心} 内陸地殻内地震の発生位置・規模・深さと、同じ活断層の再来間隔を区別する。

\paragraph{必要知識} 内陸という語を震央が必ず陸上という意味に限定しない。気象庁マグニチュードとモーメントマグニチュードを混ぜない。再来間隔は全国の地震間隔でなく一つの断層である。

\paragraph{考え方・正答を選ぶ理由（査読待ち初稿）} 地殻内地震は上盤側プレートの内部でも発生し、プレート境界地震と異なる。1995年兵庫県南部地震の気象庁マグニチュードは7.3であり、7を超えないという2には出題以前の反例がある。同じ活断層の活動間隔は通常人の寿命より長く、公式は5である。ただし「最も活発でも約1000年」を厳密な下限とするには断層群の定義と歴史資料の追加確認が必要である。

\Needspace{4\baselineskip}\paragraph{選択肢の検討（査読待ち初稿）}\mbox{}\par

\noindent\textbf{1：不適切} プレート境界で発生。地殻内の活断層活動は上盤側のプレート内部でも起きる。境界面のすべりとの区別が必要である。\par

\noindent\textbf{2：不適切} 最大級でもM7を超えない。出題前の兵庫県南部地震の気象庁M7.3が反例である。\par

\noindent\textbf{3：不適切} 震源はほとんど20 kmより深い。上部地殻の脆性変形を主体とする浅い地震であり、厚い地殻のすべてが地震を起こすわけではない。\par

\noindent\textbf{4：不適切} 震央は陸上。地殻内地震という発生機構の区分には沿岸の海域に震央を持つものも含み得る。地表の海陸区分だけで機構を決めない。\par

\noindent\textbf{5：解釈要確認} 最も活発な活断層でも約1000年に1回。公式が採る代表的な長期尺度である。同一断層の繰返しと各地の地震を区別する意図は妥当だが、全活断層の下限が1000年であるとの定量保証は未確定である。\par

\paragraph{再利用する考え方} 発生機構・震源の深さ・震央の海陸・Mの尺度を分ける。平均間隔が長いことは近日の発生が不可能という意味ではない。

\paragraph{共通理論} \hyperref[theory:interior:crust_quakes]{第\ref*{theory:interior:crust_quakes}節}

\par 参考：\cite{jma_hyogo_history,herp_longterm_intervals}。

\par 難易度 2（校準初稿）、優先度未評価。


\chapter{地球内部・プレート}
\noindent 分野解説は執筆中。以下は登録問題と校正済み例の作業用一覧。
\section{大陸地殻の組成・上下構造}

\subsection{令和2年度 問28}
\label{q:r02-reg-applied-28}

\noindent\textbf{原典：} sources/R02/17\_applied\_science.pdf、PDF物理10頁。

\textbf{公式正答：} 3（single）。

\par\textbf{検討未確定：} 年代・厚さ・密度・面積は代表値と傾向である。二層の均一性や局所的な密度序列、約40\%の厳密な境界定義まで原典は指定しない。

\paragraph{問題の要旨} 大陸地殻の上下の組成と密度、典型厚さ、年代、面積を整理する。

\Needspace{4\baselineskip}\paragraph{確認した条件・記述（独自要約）}\mbox{}\par

\noindent\textbf{確認範囲：} 問題幹、数値、単位、否定、全5候補を原典画像で照合した。\par

\noindent\textbf{成立条件：} 地殻全体の表面積を分母とする。平均と全地点、現存岩石年代と地球誕生、岩石名と平均化学組成を区別する。\par

\paragraph{問われる核心} 大陸地殻の上下の組成と密度、典型厚さ、年代、面積を整理する。

\paragraph{必要知識} 地殻全体の表面積を分母とする。平均と全地点、現存岩石年代と地球誕生、岩石名と平均化学組成を区別する。

\paragraph{考え方・正答を選ぶ理由（査読待ち初稿）} 大陸地殻の上部は平均的に珪長質に富み、下部はそれより苦鉄質に富むというモデルである。下部が主に花崗岩質という3が不適切である。大陸地殻は陸地だけでなく海面下の大陸棚にも分布するため、面積約40\%と陸地約30\%を混同しない。

\Needspace{4\baselineskip}\paragraph{選択肢の検討（査読待ち初稿）}\mbox{}\par

\noindent\textbf{1：条件付きで適切} 岩石の年代は約40億年前から現在。大陸には約40億年前の古い岩石が残り、現在も岩石が形成される。地球の形成年代や最古ジルコンの年代そのものとは区別する。\par

\noindent\textbf{2：条件付きで適切} 平均厚さ30〜40 km。代表的平均値であり、厚い山脈下や薄いリフトを含む全地点の範囲ではない。\par

\noindent\textbf{3：不適切} 下部は主に花崗岩質。一般的な組成モデルで下部は上部より苦鉄質に富む。上部の珪長質組成を下部へ置き換えた記述である。\par

\noindent\textbf{4：条件付きで適切} 密度は上部より下部が大きい傾向。下部の組成・圧力に対応する代表的傾向である。個々の局所岩体を密度だけで機械的に区分しない。\par

\noindent\textbf{5：条件付きで適切} 地殻表面積の約40\%が大陸地殻。大陸棚など海面下の大陸地殻を含む。陸地の面積比約3割とは分母対象が異なる。\par

\paragraph{再利用する考え方} 岩石名・上下・密度を関連付けて読み、40億年、40 km、40\%という似た数字の次元と分母を分ける。

\paragraph{共通理論} \hyperref[theory:interior:continental]{第\ref*{theory:interior:continental}節}, \hyperref[theory:interior:ages]{第\ref*{theory:interior:ages}節}

\par 参考：\cite{umich_crust_teaching,tokyo_crust_gallery}。

\paragraph{関連問題} \hyperref[q:r01-ret-applied-34]{R01-RET-APPLIED-34}, \hyperref[q:r04-reg-applied-26]{R04-REG-APPLIED-26}

\par 難易度 2（校準初稿）、優先度未評価。

\subsection{令和4年度 問26}
\label{q:r04-reg-applied-26}

\noindent\textbf{原典：} sources/R04/17\_applied\_science.pdf、PDF物理9頁。

\textbf{公式正答：} 1（single）。

\par\textbf{検討未確定：} 年代・厚さ・密度・面積は代表値と傾向である。二層の均一性や局所的な密度序列、約40\%の厳密な境界定義まで原典は指定しない。

\paragraph{問題の要旨} 大陸地殻の上下の組成と密度、典型厚さ、年代、面積を整理する。

\Needspace{4\baselineskip}\paragraph{確認した条件・記述（独自要約）}\mbox{}\par

\noindent\textbf{確認範囲：} 問題幹、数値、単位、否定、全5候補を原典画像で照合した。\par

\noindent\textbf{成立条件：} 地殻全体の表面積を分母とする。平均と全地点、現存岩石年代と地球誕生、岩石名と平均化学組成を区別する。\par

\paragraph{問われる核心} 大陸地殻の上下の組成と密度、典型厚さ、年代、面積を整理する。

\paragraph{必要知識} 地殻全体の表面積を分母とする。平均と全地点、現存岩石年代と地球誕生、岩石名と平均化学組成を区別する。

\paragraph{考え方・正答を選ぶ理由（査読待ち初稿）} 大陸地殻の上部は平均的に珪長質に富み、下部はそれより苦鉄質に富むというモデルである。下部が主に花崗岩質という1が不適切である。大陸地殻は陸地だけでなく海面下の大陸棚にも分布するため、面積約40\%と陸地約30\%を混同しない。

\Needspace{4\baselineskip}\paragraph{選択肢の検討（査読待ち初稿）}\mbox{}\par

\noindent\textbf{1：不適切} 下部は主に花崗岩質。一般的な組成モデルで下部は上部より苦鉄質に富む。上部の珪長質組成を下部へ置き換えた記述である。\par

\noindent\textbf{2：条件付きで適切} 岩石の年代は約40億年前から現在。大陸には約40億年前の古い岩石が残り、現在も岩石が形成される。地球の形成年代や最古ジルコンの年代そのものとは区別する。\par

\noindent\textbf{3：条件付きで適切} 平均厚さ30〜40 km。代表的平均値であり、厚い山脈下や薄いリフトを含む全地点の範囲ではない。\par

\noindent\textbf{4：条件付きで適切} 密度は上部より下部が大きい傾向。下部の組成・圧力に対応する代表的傾向である。個々の局所岩体を密度だけで機械的に区分しない。\par

\noindent\textbf{5：条件付きで適切} 地殻表面積の約40\%が大陸地殻。大陸棚など海面下の大陸地殻を含む。陸地の面積比約3割とは分母対象が異なる。\par

\paragraph{再利用する考え方} 岩石名・上下・密度を関連付けて読み、40億年、40 km、40\%という似た数字の次元と分母を分ける。

\paragraph{共通理論} \hyperref[theory:interior:continental]{第\ref*{theory:interior:continental}節}, \hyperref[theory:interior:ages]{第\ref*{theory:interior:ages}節}

\par 参考：\cite{umich_crust_teaching,tokyo_crust_gallery}。

\paragraph{関連問題} \hyperref[q:r02-reg-applied-28]{R02-REG-APPLIED-28}

\par 難易度 2（校準初稿）、優先度未評価。

\section{地殻とモホ}

\subsection{平成26年度 問30}
\label{q:h26-reg-applied-30}

\noindent\textbf{原典：} sources/H26/17\_applied\_science.pdf、PDF物理12頁。

\textbf{公式正答：} 4（single）。

\par\textbf{検討未確定：} 二つの均質な岩石層という記述と30〜40 kmの数値は近似である。原典の一般化を地球全域の厳密な組成・厚さの保証へ拡張しない。

\paragraph{問題の要旨} 地殻の位置・組成・厚さ・生成場所を、マントルとプレートの区分から判定する。

\Needspace{4\baselineskip}\paragraph{確認した条件・記述（独自要約）}\mbox{}\par

\noindent\textbf{確認範囲：} 問題幹、数値、単位、否定、全5候補を原典画像で照合した。\par

\noindent\textbf{成立条件：} 地殻とプレートは別概念である。厚さには海水を含めない。花崗岩層・玄武岩層という二層像は平均的組成の近似として読む。\par

\paragraph{問われる核心} 地殻の位置・組成・厚さ・生成場所を、マントルとプレートの区分から判定する。

\paragraph{必要知識} 地殻とプレートは別概念である。厚さには海水を含めない。花崗岩層・玄武岩層という二層像は平均的組成の近似として読む。

\paragraph{考え方・正答を選ぶ理由（査読待ち初稿）} 海洋地殻は主として発散境界の海嶺でマグマが固まって形成される。海溝は海洋プレートが沈み込む側であり、海溝で形成されるとする4が不適切である。大陸と海洋の組成・厚さの典型的相違はこの形成と再循環の違いに対応する。

\Needspace{4\baselineskip}\paragraph{選択肢の検討（査読待ち初稿）}\mbox{}\par

\noindent\textbf{1：適切} 地殻は大きな層構造の最外層。地殻・マントル・核という区分では最外側に位置する。大気をこの固体地球の区分に含めない。\par

\noindent\textbf{2：条件付きで適切} 大陸は花崗岩層と玄武岩層、海洋は玄武岩層。大陸上部は珪長質寄り、下部はより苦鉄質、海洋地殻は玄武岩・斑れい岩を主体とする。実岩石が全域で二つの均一層を作るという断定ではない。\par

\noindent\textbf{3：条件付きで適切} 大陸30〜40 km、海水を除く深海の地殻は大部分10 km未満。代表的厚さの比較であり、山脈下の厚い地殻や海台などの例外を除く典型像である。\par

\noindent\textbf{4：不適切} 海洋地殻は海溝で形成。新しい海洋地殻の主な生成場所は海嶺である。海溝の付加体形成と海洋地殻の生成を同一視しない。\par

\noindent\textbf{5：適切} 地殻とマントルの境界はモホロビチッチ不連続面。地震波の速度変化で認識される地殻底の境界である。リソスフェア底面とは異なる。\par

\paragraph{再利用する考え方} 岩石組成の区分と力学的な硬さの区分を先に分け、生成する海嶺と沈み込む海溝を対にして覚える。

\paragraph{共通理論} \hyperref[theory:interior:crust]{第\ref*{theory:interior:crust}節}, \hyperref[theory:interior:moho]{第\ref*{theory:interior:moho}節}, \hyperref[theory:interior:boundaries]{第\ref*{theory:interior:boundaries}節}

\par 参考：\cite{usgs_interior_structure,usgs_plate_boundaries,umich_crust_teaching}。

\paragraph{関連問題} \hyperref[q:h28-reg-applied-30]{H28-REG-APPLIED-30}, \hyperref[q:h29-reg-applied-26]{H29-REG-APPLIED-26}, \hyperref[q:r07-reg-applied-24]{R07-REG-APPLIED-24}

\par 難易度 2（校準初稿）、優先度未評価。

\subsection{平成28年度 問30}
\label{q:h28-reg-applied-30}

\noindent\textbf{原典：} sources/H28/17\_applied\_science.pdf、PDF物理12頁。

\textbf{公式正答：} 4（single）。

\par\textbf{検討未確定：} 二つの均質な岩石層という記述と30〜40 kmの数値は近似である。原典の一般化を地球全域の厳密な組成・厚さの保証へ拡張しない。

\paragraph{問題の要旨} 地殻の位置・組成・厚さ・生成場所を、マントルとプレートの区分から判定する。

\Needspace{4\baselineskip}\paragraph{確認した条件・記述（独自要約）}\mbox{}\par

\noindent\textbf{確認範囲：} 問題幹、数値、単位、否定、全5候補を原典画像で照合した。\par

\noindent\textbf{成立条件：} 地殻とプレートは別概念である。厚さには海水を含めない。花崗岩質層・玄武岩質層という二層像は平均的組成の近似として読む。\par

\paragraph{問われる核心} 地殻の位置・組成・厚さ・生成場所を、マントルとプレートの区分から判定する。

\paragraph{必要知識} 地殻とプレートは別概念である。厚さには海水を含めない。花崗岩質層・玄武岩質層という二層像は平均的組成の近似として読む。

\paragraph{考え方・正答を選ぶ理由（査読待ち初稿）} 海洋地殻は主として発散境界の海嶺でマグマが固まって形成される。海溝は海洋プレートが沈み込む側であり、海溝で形成されるとする4が不適切である。大陸と海洋の組成・厚さの典型的相違はこの形成と再循環の違いに対応する。

\Needspace{4\baselineskip}\paragraph{選択肢の検討（査読待ち初稿）}\mbox{}\par

\noindent\textbf{1：適切} 地殻は大きな層構造の最外層。地殻・マントル・核という区分では最外側に位置する。大気をこの固体地球の区分に含めない。\par

\noindent\textbf{2：条件付きで適切} 大陸は花崗岩質層と玄武岩質層、海洋は玄武岩質層。大陸上部は珪長質寄り、下部はより苦鉄質、海洋地殻は玄武岩・斑れい岩を主体とする。実岩石が全域で二つの均一層を作るという断定ではない。\par

\noindent\textbf{3：条件付きで適切} 大陸30〜40 km、海水を除く深海の地殻は大部分10 km未満。代表的厚さの比較であり、山脈下の厚い地殻や海台などの例外を除く典型像である。\par

\noindent\textbf{4：不適切} 海洋地殻は海溝で形成。新しい海洋地殻の主な生成場所は海嶺である。海溝の付加体形成と海洋地殻の生成を同一視しない。\par

\noindent\textbf{5：適切} 地殻とマントルの境界はモホロビチッチ不連続面。地震波の速度変化で認識される地殻底の境界である。リソスフェア底面とは異なる。\par

\paragraph{再利用する考え方} 岩石組成の区分と力学的な硬さの区分を先に分け、生成する海嶺と沈み込む海溝を対にして覚える。

\paragraph{共通理論} \hyperref[theory:interior:crust]{第\ref*{theory:interior:crust}節}, \hyperref[theory:interior:moho]{第\ref*{theory:interior:moho}節}, \hyperref[theory:interior:boundaries]{第\ref*{theory:interior:boundaries}節}

\par 参考：\cite{usgs_interior_structure,usgs_plate_boundaries,umich_crust_teaching}。

\paragraph{関連問題} \hyperref[q:h26-reg-applied-30]{H26-REG-APPLIED-30}, \hyperref[q:h29-reg-applied-26]{H29-REG-APPLIED-26}

\par 難易度 2（校準初稿）、優先度未評価。

\subsection{平成29年度 問26}
\label{q:h29-reg-applied-26}

\noindent\textbf{原典：} sources/H29/17\_applied\_science.pdf、PDF物理9頁。

\textbf{公式正答：} 4（single）。

\par\textbf{検討未確定：} 二つの均質な岩石層という記述と30〜40 kmの数値は近似である。原典の一般化を地球全域の厳密な組成・厚さの保証へ拡張しない。

\paragraph{問題の要旨} 地殻の位置・組成・厚さ・生成場所を、マントルとプレートの区分から判定する。

\Needspace{4\baselineskip}\paragraph{確認した条件・記述（独自要約）}\mbox{}\par

\noindent\textbf{確認範囲：} 問題幹、数値、単位、否定、全5候補を原典画像で照合した。\par

\noindent\textbf{成立条件：} 地殻とプレートは別概念である。厚さには海水を含めない。花崗岩層・玄武岩層という二層像は平均的組成の近似として読む。\par

\paragraph{問われる核心} 地殻の位置・組成・厚さ・生成場所を、マントルとプレートの区分から判定する。

\paragraph{必要知識} 地殻とプレートは別概念である。厚さには海水を含めない。花崗岩層・玄武岩層という二層像は平均的組成の近似として読む。

\paragraph{考え方・正答を選ぶ理由（査読待ち初稿）} 海洋地殻は主として発散境界の海嶺でマグマが固まって形成される。海溝は海洋プレートが沈み込む側であり、海溝で形成されるとする4が不適切である。大陸と海洋の組成・厚さの典型的相違はこの形成と再循環の違いに対応する。

\Needspace{4\baselineskip}\paragraph{選択肢の検討（査読待ち初稿）}\mbox{}\par

\noindent\textbf{1：適切} 地殻は大きな層構造の最外層。地殻・マントル・核という区分では最外側に位置する。大気をこの固体地球の区分に含めない。\par

\noindent\textbf{2：条件付きで適切} 大陸は花崗岩層と玄武岩層、海洋は玄武岩層。大陸上部は珪長質寄り、下部はより苦鉄質、海洋地殻は玄武岩・斑れい岩を主体とする。実岩石が全域で二つの均一層を作るという断定ではない。\par

\noindent\textbf{3：条件付きで適切} 大陸30〜40 km、海水を除く深海の地殻は大部分10 km未満。代表的厚さの比較であり、山脈下の厚い地殻や海台などの例外を除く典型像である。\par

\noindent\textbf{4：不適切} 海洋地殻は海溝で形成。新しい海洋地殻の主な生成場所は海嶺である。海溝の付加体形成と海洋地殻の生成を同一視しない。\par

\noindent\textbf{5：適切} 地殻とマントルの境界はモホロビチッチ不連続面。地震波の速度変化で認識される地殻底の境界である。リソスフェア底面とは異なる。\par

\paragraph{再利用する考え方} 岩石組成の区分と力学的な硬さの区分を先に分け、生成する海嶺と沈み込む海溝を対にして覚える。

\paragraph{共通理論} \hyperref[theory:interior:crust]{第\ref*{theory:interior:crust}節}, \hyperref[theory:interior:moho]{第\ref*{theory:interior:moho}節}, \hyperref[theory:interior:boundaries]{第\ref*{theory:interior:boundaries}節}

\par 参考：\cite{usgs_interior_structure,usgs_plate_boundaries,umich_crust_teaching}。

\paragraph{関連問題} \hyperref[q:h26-reg-applied-30]{H26-REG-APPLIED-30}, \hyperref[q:h28-reg-applied-30]{H28-REG-APPLIED-30}

\par 難易度 2（校準初稿）、優先度未評価。

\section{島弧配置と熱流量}

\subsection{令和元年度再試験 問33}
\label{q:r01-ret-applied-33}

\noindent\textbf{原典：} sources/R01\_retest/17\_applied\_science.pdf、PDF物理11頁。

\textbf{公式正答：} 4（single）。

\par\textbf{検討未確定：} 5島弧の区分と縁海の同義扱いは慣用的である。熱流量の広域傾向と地点ごとの地下水・火山の熱輸送を混同しない。

\paragraph{問題の要旨} 島弧の海溝・火山帯・背弧海盆・熱流量・沈み込む地震面の配置を判定する。

\Needspace{4\baselineskip}\paragraph{確認した条件・記述（独自要約）}\mbox{}\par

\noindent\textbf{確認範囲：} 問題幹、数値、単位、否定、全5候補を原典画像で照合した。\par

\noindent\textbf{成立条件：} 海溝側を前弧、反対の大陸側を背弧とする。日本の5島弧は通常の区分による。地殻熱流量は地表を通る熱流束であり、地下温度そのものではない。\par

\paragraph{問われる核心} 島弧の海溝・火山帯・背弧海盆・熱流量・沈み込む地震面の配置を判定する。

\paragraph{必要知識} 海溝側を前弧、反対の大陸側を背弧とする。日本の5島弧は通常の区分による。地殻熱流量は地表を通る熱流束であり、地下温度そのものではない。

\paragraph{考え方・正答を選ぶ理由（査読待ち初稿）} 東北日本弧の広域的な地殻熱流量は、太平洋側の前弧が低く、日本海側の火山・背弧側が高いという分布である。太平洋側で高く大陸側で低いとする4は、この方向を逆にしている。局所的な地下水の移流や火山近傍の異常を広域傾向と区別する。

\Needspace{4\baselineskip}\paragraph{選択肢の検討（査読待ち初稿）}\mbox{}\par

\noindent\textbf{1：条件付きで適切} 5つの弧状列島から構成。千島・東北日本・西南日本・伊豆小笠原・琉球という通常の区分で読む。島弧をさらに分割する流儀まで固定しない。\par

\noindent\textbf{2：適切} 東北日本弧の太平洋側に6000 m以上の海溝が平行。日本海溝と東北日本弧の沈み込み側の配置に対応する。\par

\noindent\textbf{3：条件付きで適切} 大陸側に縁海・背弧海盆。日本海などの典型配置である。縁海という地理区分と背弧海盆という形成過程の区分は完全な同義ではない。\par

\noindent\textbf{4：不適切} 中部軸に火山帯、東北日本の熱流量は太平洋側で高く大陸側で低い。火山帯の記述と合わせた複合肢のうち、熱流量の東西の大小が逆である。\par

\noindent\textbf{5：適切} 海溝から大陸側へ深くなる深発地震面。沈み込むスラブ内の地震分布が深部へ傾く和達・ベニオフ帯に対応する。\par

\paragraph{再利用する考え方} 前弧・火山フロント・背弧を位置の関係で把握し、地温勾配、熱流束、地震面の深度を別の観測量として読む。

\paragraph{共通理論} \hyperref[theory:interior:arc]{第\ref*{theory:interior:arc}節}, \hyperref[theory:interior:heat]{第\ref*{theory:interior:heat}節}

\par 参考：\cite{gsj_heat_flow1991,usgs_plate_boundaries}。

\par 難易度 2（校準初稿）、優先度未評価。

\section{リソスフェアとアセノスフェア}

\subsection{令和7年度 問24}
\label{q:r07-reg-applied-24}

\noindent\textbf{原典：} sources/R07/17\_applied\_science.pdf、PDF物理9頁。

\textbf{公式正答：} 4（single）。

\par\textbf{検討未確定：} 中心温度5000〜6000℃の推定根拠の精査は未完了である。主要プレート枚数と長期的硬さの区分の限定を保持する。

\paragraph{問題の要旨} 地殻・マントルの組成境界とリソスフェア・アセノスフェアの力学境界を区別する。

\Needspace{4\baselineskip}\paragraph{確認した条件・記述（独自要約）}\mbox{}\par

\noindent\textbf{確認範囲：} 問題幹、数値、単位、否定、全5候補を原典画像で照合した。\par

\noindent\textbf{成立条件：} 地球中心の温度は実験と内部モデルからの推定である。地殻・マントルという層と力学的硬さの層は境界位置が一致しない。主要プレート枚数は細分方法による。\par

\paragraph{問われる核心} 地殻・マントルの組成境界とリソスフェア・アセノスフェアの力学境界を区別する。

\paragraph{必要知識} 地球中心の温度は実験と内部モデルからの推定である。地殻・マントルという層と力学的硬さの層は境界位置が一致しない。主要プレート枚数は細分方法による。

\paragraph{考え方・正答を選ぶ理由（査読待ち初稿）} モホロビチッチ不連続面は地殻とマントルの境界である。リソスフェアには地殻だけでなく最上部マントルが含まれる。したがってリソスフェアとアセノスフェアの境界をモホとする4が不適切である。軟らかく流れるという性質は長い時間尺度の変形であり、マントル全体が液体という意味ではない。

\Needspace{4\baselineskip}\paragraph{選択肢の検討（査読待ち初稿）}\mbox{}\par

\noindent\textbf{1：条件付きで適切} 温度は深部ほど高く中心約5000〜6000℃。中心温度の代表的推定範囲である。全深度の厳密な単調性や直接測定値とはしない。\par

\noindent\textbf{2：適切} 上部マントルに流れやすいアセノスフェア。温度が高く長期的に変形しやすい領域であり、通常は主に固体である。\par

\noindent\textbf{3：適切} その上の硬いリソスフェア。地殻と最上部マントルを含む、長期にも比較的強い層である。\par

\noindent\textbf{4：不適切} 両者の境界はモホ。モホは地殻底であり、リソスフェア底の力学境界ではない。\par

\noindent\textbf{5：条件付きで適切} 10数〜20枚程度のプレートの実体はリソスフェア。主要プレートという典型区分では妥当である。微小プレートを数えると枚数が増える。\par

\paragraph{再利用する考え方} 物質の種類、地震波の速度、変形の時間尺度による区分を分ける。硬い層を地殻だけとみなさない。

\paragraph{共通理論} \hyperref[theory:interior:lithosphere]{第\ref*{theory:interior:lithosphere}節}, \hyperref[theory:interior:moho]{第\ref*{theory:interior:moho}節}

\par 参考：\cite{usgs_lithosphere2002,kagoshima_plate2019}。

\paragraph{関連問題} \hyperref[q:h26-reg-applied-30]{H26-REG-APPLIED-30}

\par 難易度 2（校準初稿）、優先度未評価。

\section{静岩圧・上載荷重}

\subsection{平成27年度 問32}
\label{q:h27-reg-applied-32}

\noindent\textbf{原典：} sources/H27/17\_applied\_science.pdf、PDF物理12頁。

\textbf{公式正答：} 4（single）。

\paragraph{問題の要旨} 上載岩石の密度と厚さからモホ面の静岩圧を求める。

\Needspace{4\baselineskip}\paragraph{確認した条件・記述（独自要約）}\mbox{}\par

\noindent\textbf{確認範囲：} 問題幹、数値、単位、否定、全5候補を原典画像で照合した。\par

\noindent\textbf{成立条件：} 厚さ35 km、平均密度2.8×10\textasciicircum{}3 kg/m\textsuperscript{3}、重力加速度9.8 m/s\textsuperscript{2}。上載荷重のみを考え、間隙流体の影響はないという原典指定を使用する。\par

\paragraph{問われる核心} 上載岩石の密度と厚さからモホ面の静岩圧を求める。

\paragraph{必要知識} 厚さ35 km、平均密度2.8×10\textasciicircum{}3 kg/m\textsuperscript{3}、重力加速度9.8 m/s\textsuperscript{2}。上載荷重のみを考え、間隙流体の影響はないという原典指定を使用する。

\paragraph{考え方・正答を選ぶ理由（査読待ち初稿）} 水平面の単位面積当たりに載る岩石の質量はρhである。上載荷重だけを考えるので \(p=\rho gh=(2.8\times10^3)(9.8)(35\times10^3)=9.604\times10^8\,\mathrm{Pa}=0.9604\,\mathrm{GPa}\)。最も近い1 GPaの4を選ぶ。

\Needspace{4\baselineskip}\paragraph{選択肢の検討（査読待ち初稿）}\mbox{}\par

\noindent\textbf{1：不適切} 1 MPa。9.604×10\textasciicircum{}8 Paより約3桁小さい。\par

\noindent\textbf{2：不適切} 10 MPa。導出圧力より約2桁小さい。\par

\noindent\textbf{3：不適切} 100 MPa。導出圧力0.9604 GPaより約1桁小さい。\par

\noindent\textbf{4：適切} 1 GPa。9.604×10\textasciicircum{}8 Paを10\textasciicircum{}9 Paで割った0.9604 GPaに最も近い。\par

\noindent\textbf{5：不適切} 10 GPa。導出圧力より約1桁大きい。\par

\paragraph{再利用する考え方} kmをmへ直し、密度×加速度×長さの単位がPaになることを確認する。静岩圧と間隙水圧及び有効応力を別にする。

\paragraph{共通理論} \hyperref[theory:interior:pressure]{第\ref*{theory:interior:pressure}節}

\par 参考：\cite{usgs_lithosphere2002}。

\par 難易度 2（校準初稿）、優先度未評価。

\section{海洋底と大陸の比較}

\subsection{令和元年度再試験 問34}
\label{q:r01-ret-applied-34}

\noindent\textbf{原典：} sources/R01\_retest/17\_applied\_science.pdf、PDF物理11頁。

\textbf{公式正答：} 2（single）。

\par\textbf{検討未確定：} 大洋底の範囲・モホ深度の基準・堆積速度の地域比較が不明である。候補3の存在命題には例外の精査が必要であり、公式2だけを根拠に候補3を断定しない。

\paragraph{問題の要旨} 海洋底と大陸を、地殻厚さ・堆積・組成・年代の異なる尺度で比較する。

\Needspace{4\baselineskip}\paragraph{確認した条件・記述（独自要約）}\mbox{}\par

\noindent\textbf{確認範囲：} 問題幹、数値、単位、否定、全5候補を原典画像で照合した。\par

\noindent\textbf{成立条件：} 典型的な外洋底と大陸の比較であり、海底すべてを海洋地殻と同一視しない。モホの深さには海面又は固体表面という基準が必要である。\par

\paragraph{問われる核心} 海洋底と大陸を、地殻厚さ・堆積・組成・年代の異なる尺度で比較する。

\paragraph{必要知識} 典型的な外洋底と大陸の比較であり、海底すべてを海洋地殻と同一視しない。モホの深さには海面又は固体表面という基準が必要である。

\paragraph{考え方・正答を選ぶ理由（査読待ち初稿）} 外洋の遠洋性海底では陸源物質の供給が少なく、堆積が遅いことが典型比較となり公式は2である。海洋地殻は薄く、海嶺生成と沈み込みにより現存海洋底の年代は通常数億年以内であり30億年ではない。花崗岩質岩石の全称否定には海洋斜長花崗岩や大陸棚などの範囲の問題があり、3の不適切性を無条件には確定しない。

\Needspace{4\baselineskip}\paragraph{選択肢の検討（査読待ち初稿）}\mbox{}\par

\noindent\textbf{1：不適切} 大洋底のモホは大陸より深い。典型的海洋地殻は薄いため固体表面からのモホ深度は浅い。海面基準でも典型的大陸地殻より浅いが、基準と地域が未指定である。\par

\noindent\textbf{2：条件付きで適切} 大洋底の堆積速度は大陸より遅い。陸源物質の供給が小さい遠洋域の典型である。大陸縁辺の高い堆積速度や大陸の無堆積・侵食域は例外である。\par

\noindent\textbf{3：解釈要確認} 大洋底に花崗岩質岩石が存在。典型的海洋地殻の主体が玄武岩・斑れい岩という説明だけでは、少量の珪長質岩や沈水した大陸地殻の存在を否定できない。原典の大洋底の範囲を要確認とする。\par

\noindent\textbf{4：不適切} 海洋地殻は大陸地殻より厚い。典型的な厚さの大小が逆である。\par

\noindent\textbf{5：不適切} 大洋底の年代は0〜30億年。プレートの再循環が続く現存外洋底は代表的に数億年以内であり、大陸の古い岩石の年代を当てはめている。\par

\paragraph{再利用する考え方} 一般的な大小比較と存在命題を分ける。「主体でない」ことから「存在しない」を導かない。

\paragraph{共通理論} \hyperref[theory:interior:crust]{第\ref*{theory:interior:crust}節}, \hyperref[theory:interior:ages]{第\ref*{theory:interior:ages}節}

\par 参考：\cite{usgs_plate_boundaries,umich_crust_teaching}。

\paragraph{関連問題} \hyperref[q:r02-reg-applied-28]{R02-REG-APPLIED-28}

\par 難易度 2（校準初稿）、優先度未評価。

\section{プレート境界の相対運動}

\subsection{令和3年度 問31}
\label{q:r03-reg-applied-31}

\noindent\textbf{原典：} sources/R03/17\_applied\_science.pdf、PDF物理11頁。

\textbf{公式正答：} 3（single）。

\par\textbf{検討未確定：} 原典の海嶺高度と低速拡大時の溶融度は典型像であり、超低速域の地域条件は省略される。トランスフォームと非活動断裂帯は本文で区別する。

\paragraph{問題の要旨} 三種のプレート境界とリフト・背弧の応力場を、生成・消費・横ずれで区別する。

\Needspace{4\baselineskip}\paragraph{確認した条件・記述（独自要約）}\mbox{}\par

\noindent\textbf{確認範囲：} 問題幹、数値、単位、否定、全5候補を原典画像で照合した。\par

\noindent\textbf{成立条件：} 海嶺間の活動的なトランスフォーム区間と、その外側の非活動的な断裂帯を区別する。海嶺の高さと部分溶融は典型モデルとして扱う。\par

\paragraph{問われる核心} 三種のプレート境界とリフト・背弧の応力場を、生成・消費・横ずれで区別する。

\paragraph{必要知識} 海嶺間の活動的なトランスフォーム区間と、その外側の非活動的な断裂帯を区別する。海嶺の高さと部分溶融は典型モデルとして扱う。

\paragraph{考え方・正答を選ぶ理由（査読待ち初稿）} 海嶺区間の間にある活動的なトランスフォーム断層は両プレートの相対的な横ずれを担う。正断層は伸張で上盤が相対的に下がる断層であり、3の説明とは異なる。断裂帯全長が現在も活動するプレート境界という意味でもない。

\Needspace{4\baselineskip}\paragraph{選択肢の検討（査読待ち初稿）}\mbox{}\par

\noindent\textbf{1：条件付きで適切} 高温の新しい海洋プレートが低密度で海嶺は2〜3 km高い。熱膨張と浮力による海嶺の隆起を説明する典型像である。特定の地形差の厳密値ではない。\par

\noindent\textbf{2：条件付きで適切} 遅い拡大では溶融度が小さくマグマをほとんど伴わない場合。特に超低速拡大の一部ではマグマ供給が少なくマントル岩が露出し得る。遅い海嶺の全域が無マグマという意味ではない。\par

\noindent\textbf{3：不適切} 海嶺をつなぐ断裂帯で正断層運動しトランスフォームとなる。活動区間での基本運動は横ずれである。周辺に局所的な正断層があることとは別の分類である。\par

\noindent\textbf{4：適切} 大陸発散境界の正断層で伸張・薄化・地溝形成。伸張する地殻の相対沈降とリフト形成を対応させた記述である。\par

\noindent\textbf{5：条件付きで適切} 沈み込み帯は圧縮だが背弧が伸張し拡大する場合。沈み込みの存在は背弧の全域が常に圧縮であることを保証しない。海溝後退等に伴う背弧拡大がある。\par

\paragraph{再利用する考え方} 境界の分類を断層名の連想で決めず、相対速度の方向を見て発散・収束・横ずれを判定する。

\paragraph{共通理論} \hyperref[theory:interior:boundaries]{第\ref*{theory:interior:boundaries}節}, \hyperref[theory:interior:ridge]{第\ref*{theory:interior:ridge}節}

\par 参考：\cite{usgs_plate_boundaries,umich_crust_teaching,noaa_ultraslow2004}。

\par 難易度 2（校準初稿）、優先度未評価。

\section{相対プレート速度}

\subsection{令和元年度 問30}
\label{q:r01-reg-applied-30}

\noindent\textbf{原典：} sources/R01/17\_applied\_science.pdf、PDF物理10頁。

\textbf{公式正答：} 5（single）。

\par\textbf{検討未確定：} 北アメリカ固定は明記されるが比較地点とオイラー極モデルは不明である。大学の解説は日本列島に対する代表速度であり、原典の参照系の厳密な数値検証は未完了である。

\paragraph{問題の要旨} 固定した基準プレートに対する日本付近の二つの海洋プレートの速度の桁を比較する。

\Needspace{4\baselineskip}\paragraph{確認した条件・記述（独自要約）}\mbox{}\par

\noindent\textbf{確認範囲：} 問題幹、数値、単位、否定、全5候補を原典画像で照合した。\par

\noindent\textbf{成立条件：} 原典は北アメリカプレート固定を明記する。剛体プレートの相対角速度と地点の位置から速度が決まるため、地点と運動モデルの省略を残す。\par

\paragraph{問われる核心} 固定した基準プレートに対する日本付近の二つの海洋プレートの速度の桁を比較する。

\paragraph{必要知識} 原典は北アメリカプレート固定を明記する。剛体プレートの相対角速度と地点の位置から速度が決まるため、地点と運動モデルの省略を残す。

\paragraph{考え方・正答を選ぶ理由（査読待ち初稿）} 日本付近の代表的沈み込み速度は太平洋プレートが年8〜9 cm、フィリピン海プレートが年4〜5 cmの桁である。候補では9 cm/年と4.5 cm/年の5が対応する。北アメリカ固定という原典の基準を保持し、一般的な日本列島に対する値と完全に同じモデルと証明したとは扱わない。

\Needspace{4\baselineskip}\paragraph{選択肢の検討（査読待ち初稿）}\mbox{}\par

\noindent\textbf{1：不適切} 太平洋6.5 mm/年、フィリピン海2 mm/年。両方が代表的cm/年の運動より一桁程度小さい。\par

\noindent\textbf{2：不適切} 太平洋8 mm/年、フィリピン海7 mm/年。mm/年とcm/年の違いにより代表的な沈み込み速度に達しない。\par

\noindent\textbf{3：不適切} 太平洋2.5 cm/年、フィリピン海1.5 cm/年。両方とも日本付近の代表的速度より小さい組である。\par

\noindent\textbf{4：不適切} 太平洋3 cm/年、フィリピン海8.5 cm/年。太平洋の方が大きい代表比較が逆になっている。\par

\noindent\textbf{5：条件付きで適切} 太平洋9 cm/年、フィリピン海4.5 cm/年。代表的な大きさと大小関係に対応する。原典にない地点・モデルを補って精密値としない。\par

\paragraph{再利用する考え方} cm/年は絶対値ではなく参照系と場所を伴う速度である。拡大速度、沈み込み速度、陸上GNSS点の速度を同じものとして比較しない。

\paragraph{共通理論} \hyperref[theory:interior:velocity]{第\ref*{theory:interior:velocity}節}

\par 参考：\cite{kagoshima_plate2019}。

\par 難易度 2（校準初稿）、優先度未評価。


\chapter{重力・測地}
\noindent 分野解説は執筆中。以下は登録問題と校正済み例の作業用一覧。
\subsection{平成27年度 問33}
\label{q:h27-reg-applied-33}

\noindent\textbf{原典：} sources/H27/17\_applied\_science.pdf、PDF物理13頁。

\textbf{公式正答：} 3（single）。

\par\textit{原典本文の校正と独自解説は未作成。分類のみの登録項目。}


\chapter{地磁気}
\section{磁場の方向と地理的な北}
\label{theory:geomagnetism:declination}
方位磁針が向く方向は地理的な真北と一般に一致しない。地磁気の水平成分を真北方向\(X\)、東方向\(Y\)に分けると、東偏を正とする規約では偏角\(D\)を\(D=\operatorname{atan2}(Y,X)\)で定める。国土地理院の磁気図で西偏を正の大きさとして表示する場合は符号規約が異なる。角の大きさ、東西方向、符号を別々に確かめる。

地球中心の理想軸双極子だけなら非常に単純な分布となるが、実際の磁場には非双極子成分と局所的な磁気異常がある。同じ緯度又は同じ経度にあるだけでは偏角は同一にならない。また地球内部の流体運動などに伴って永年変化するため、磁気図には基準時刻である元期が付される。

令和元年度問27の「現在」は2019年の試験時点である。試験前に公表された国土地理院の2015.0年磁気図を参照すると、日本周辺は西偏で、沖縄本島付近は約5度、北海道では概ね9〜10度である\cite{gsi_declination2015}。これらは当時の問題を判断するための概略値であり、2026年現在の各地点の測量値を表すものではない。緯度だけの普遍的単調則を主張せず、地図上の分布に基づいて誤答肢を判断する。

\section{誘導磁化・残留磁化と磁気探査}
\label{theory:geomagnetism:magnetization}
磁気探査は岩石や埋設物の磁化が生む磁場の差を測り、地下構造を推定する。岩石の磁化には、現在の地磁場に応答する誘導磁化と、過去の場などを記録する残留磁化がある。線形近似では
\[
 \boldsymbol{M}=\kappa\boldsymbol{H}+\boldsymbol{M}_r
\]
である。\(\kappa\)は磁化率、\(\boldsymbol{H}\)は磁場の強さ、\(\boldsymbol{M}_r\)は自然残留磁化である。両者の方向が一致するとは限らないので、磁気異常の正負を磁性物質の有無だけで決めることはできない。

ケーニヒスベルガー比は両者の大きさを比べる無次元量であり、
\[
 Q=\frac{|\boldsymbol{M}_r|}{|\kappa\boldsymbol{H}|}
\]
と書ける。\(Q\ll1\)なら誘導磁化が支配的、\(Q\gg1\)なら残留磁化を無視できない。地質調査研究報告の花崗岩の物理定数研究にもこの定義が使われている\cite{kanaya_okuma2011}。比を知っても磁化方向まで確定するわけではなく、密度の比、磁化率そのもの、地震波速度の比ではない。

強磁性・フェリ磁性体の自発磁化は熱運動による磁気秩序の変化と結び付いている。キュリー温度は自発磁化が消失する転移温度であり、磁性鉱物や磁気を帯びる層の解釈に関係する。一方、プランク関数は熱平衡放射のスペクトルを表す。岩石の温度を扱うという広い共通点だけで、磁気探査の基本用語と同列にしてはならない。

\section{探査データから製品名まで断定できるか}
磁化が同じでも形状・深度が違えば類似した磁気異常を生みうる。位置の概略と磁性体の存在を探ることと、細かな製品種類まで一意に識別することの間には逆問題の非一意性がある。平成23年度問29はこの通常の磁気探査上の限界を述べる。他の観測や形状モデルを組み合わせた現代の識別技術全般まで否定するものと読まない。共通理論は物理探査章でも観測量の違いから整理する。
\subsection{令和元年度 問27}
\label{q:r01-reg-applied-27}

\noindent\textbf{原典：} sources/R01/17\_applied\_science.pdf、PDF物理9頁。

\textbf{公式正答：} 2（single）。

\par\textit{原典本文の校正と独自解説は未作成。分類のみの登録項目。}


\chapter{火山・マグマ}
\section{火砕密度流を考える背景}
\label{theory:volcanology:pyroflow}
噴火で破砕された粒子と高温のガスが混合し、周囲より重い混合物として斜面を流下するのが火砕密度流である。大気を取り込み、温度と濃度を変えながら進む。濃い流れと希薄なサージでは地形に対する挙動も異なる。堆積物は一般に分級が悪い塊状となりやすいが、全てが一様という意味ではない。

\section{ドーム崩落と噴煙柱からの発生}
\label{theory:volcanology:source}
溶岩ドームの崩落は岩塊と灰からなる流れを生む。メラピ1994年の論文では、谷を埋める流れと、それから分離して地形を横切る希薄サージが記録される。爆発的破砕と重力による崩落を分ける。噴煙柱の崩壊や大規模な噴出は別の発生経路であり、ドーム起源の小規模流とカルデラ形成を伴う大規模な活動を同一視しない。過去問のプレー型・メラピ型という歴史的呼称と大規模流の一般化は、個別解説で確認範囲を残す。\cite{volc_usgs_merapi}

\section{堆積分布は風と地形を分けて読む}
\label{theory:volcanology:distribution}
降下火山灰は風によって火口の風下へ運ばれ、細粒ほど遠くまで届きやすい。地面を流れる火砕流では谷、尾根、斜面変化、流れの濃度が経路を変えるため、舌状・ローブ状・扇状の分布や複数谷の埋積を考える。火口を一端とする単純な長円形を一般的な火砕流の分布として決め付けない。元問題図の転載をせず、移動媒体と支配条件の違いを読み取る。\cite{volc_usgs_msh_flow,volc_usgs_merapi}

\section{固体の内部で局所的に融解する}\label{theory:volcano:melting}
マントルの多くは固体である。地震波を伝える固体であることと、地質学的な時間にゆっくり流動することは両立する。マグマは地球内部を一面に覆う液体ではなく、温度・圧力・組成の経路が融解条件を越える場所で生じる。
多成分岩石では融解開始の固相線と全てが融ける液相線が異なる。間では固相とメルトが共存し、部分溶融となる。初生マグマは原岩の組成をそのまま全量コピーしたものではない\cite{vol50_melt_ref}。
物理での相図の読み方を使うと、融解開始は概念的に \(T_{\rm path}(P)>T_{\rm solidus}(P,X)\) と表せる。上昇するマントルでは圧力が下がり、ほぼ断熱的な温度経路と固相線の交差で減圧融解が始まる。海嶺の成因を「低温の物質が突然加熱される」と説明しない。ホットスポットの熱異常と上昇減圧は別の要因である。
沈み込みでは含水物質の脱水により流体がマントルウェッジへ供給され、融解開始温度を下げる\cite{vol50_subduction_ref}。温度を上げる、圧力を下げる、水等で固相線を下げるという三経路を分ける。
メルトが周囲より低密度なら浮力で上昇し得る。概略の駆動力は体積 \(V\) 当たり \((\rho_{\rm rock}-\rho_{\rm melt})gV\) だが、実際の移動には割れ目・応力・粘性・冷却も関係する。全マグマだまりを中立浮力だけで決定しない。

\section{原岩と別マグマのどちらと比較するか}
初生的なマントル融解液は代表的に玄武岩質である。結晶が分離すれば残る液体の組成は変わり、地殻の融解や同化・混合も別のマグマを作る\cite{vol50_melt_ref}。
玄武岩質は流紋岩質等よりSiO2に乏しい側だが、原岩マントルに比べればNa・K等の不適合元素がメルトへ入り濃集する場合がある。「富む」は比較対象を明記する。
JMAの1999年表では玄武岩質の代表噴出温度が1000--1200℃、流紋岩質が700--900℃で、珪酸分が多い側は一般に粘性が大きい\cite{jma_magma1999}。これは原典時点の分類と代表値であり、全組成・全含水量・全結晶量の固定境界ではない。

\section{海溝・前弧・火山弧・背弧}\label{theory:volcano:arc}
地表で海溝から陸側へ前弧・火山弧・背弧を並べ、地下では沈み込むスラブとその上のマントルウェッジを対応させる。火山前線は火山の分布の海溝側限界で、海溝そのものと違う\cite{usgs_plate_boundaries}。
火山が少ない前弧でもプレート境界やスラブの地震活動があり、「火山がないから地震もない」は誤りである。背弧に拡大海盆ができる場合はあるが、全島弧が同じ拡大状態を持つとはしない。
アルカリ元素と水を一つの変数として扱わない。島弧の背弧方向のアルカリ濃度には一般的な増加傾向を学ぶ一方、GSJの東北日本弧の発表には含水量が背弧へ減少する分布の報告がある\cite{vol50_arcwater_ref}。融解度・供給源・脱ガスの違いを考え、全弧の水濃度へ固定した順位を使わない。
中央海嶺は発散境界で新しい海洋地殻を作る。USGS概説では毎年の地球の噴出物の約4分の3が拡大海嶺起源と推計される\cite{vol50_ridge_ref}。陸上噴火の個数や被害の印象と、全世界の体積率を区別する。

\section{間欠爆発と連続溶岩噴泉}\label{theory:volcano:styles}
低粘性の玄武岩質マグマの連続的な噴泉や流出はハワイ式の代表である。ストロンボリ式はガスの断続的な解放により小爆発を繰り返し、火山弾・スコリアを放出する\cite{jma_magma1999,vol50_styles_ref}。同じ火山が常に一つの様式だけで噴火するわけではない。
高粘性の溶岩は流れにくく火口近くへ蓄積してドームを作り得る。ガスが抜ける時間、供給量、結晶量等も噴火様式へ関係する。岩種名だけで爆発の有無を必ず決めない。
マグマ水蒸気爆発では外来水とマグマの急な熱交換・蒸気膨張・破砕が関与する。水蒸気噴火で既存岩石が破砕される場合と、新しいマグマが外来水へ接触する場合を区別する。

\section{溶岩表面の形と冷却}\label{theory:volcano:lava}
パホイホイは代表的に玄武岩質の流動的な溶岩で、滑らか・縄状等の表面を持つ。アアは粗いクリンカーに覆われる。同じ玄武岩でも冷却・結晶化・変形速度等で一方から他方へ変化し得る\cite{vol50_lava_ref}。滑らかだから流紋岩という対応は誤りである。
ブロック溶岩では大きい岩塊の比較的平滑な破断面が特徴となる。表面を形態で分ける軸と、岩石の化学組成を分ける軸を混同しない。
厚さ／長さ等のアスペクト比は地形を比較する指標である。原典の1/8区分は分類上の目安として保存し、普遍的な相転移や流動停止の定数とはしない。長さの測り方・流量・地形・冷却の影響を考える。

\section{ガスの主要成分と平衡条件}\label{theory:volcano:gases}
通常の火山ガスの最多成分は水蒸気であり、CO2、SO2、H2S、水素ハロゲン化物等も放出される\cite{vol50_gas_ref}。危険性の大きさと濃度順位は別の情報である。
減圧ではメルト中の揮発性分の溶解度が下がり、気泡が分離し膨張する。CO2と水の脱ガス段階は同じではない\cite{vol50_degassing_ref}。閉塞で過圧が生じる場合と、開いた系でガスが抜ける場合を分ける。外圧低下だけからガス絶対圧力の単調増加を推論しない。
高温マグマのガス、冷却された噴気、地下水と反応した熱水ガスは組成が違い得る。CH4を通常高温の主要ガスの比較で選外とすることは、全自然条件で存在しないことを意味しない。令和7年度の100気圧・1200℃・玄武岩質という条件も、酸化還元状態・全元素量なしに任意の正確なモル割合を作れるわけではない。

\section{粒子とガス、水、岩屑を分ける}\label{theory:volcano:hazards}
火砕流は高温のガス・粒子混合、火山泥流は水・粒子混合、岩屑なだれは崩壊岩体の移動を中心に考える。速さの目安と材料を同時に確認する。噴煙柱の崩壊やドーム崩落と、融雪・湖水・降雨による泥流形成を同一視しない。火砕密度流の背景は第\ref{theory:volcanology:pyroflow}節へ戻る。
大量の物質が海へ入り水を急激に動かすと津波が生じる場合がある。噴火・海底地殻変動・崩壊による水塊の変位を考え、全噴火が津波を起こすとはしない。
噴煙が高くても、灰の短期的な遮蔽とSO2由来硫酸エアロゾルの長距離・長期間の気候影響を区別する\cite{vol50_climate_ref}。粒子の沈降時間、化学反応、成層圏での滞留という物理と化学の段階を追う。

\section{火炎状構造の「火」は成因ではない}\label{theory:volcano:terms}
カルデラ、火山岩の安山岩、溶岩片が付着したアグルチネート、マグマ水蒸気爆発は火山の直接用語である。
火炎状構造は未固結の軽い堆積層が上の重い層の間へ押し上げられる等の変形を指し、名称は形の比喩である\cite{vol50_terms_ref}。堆積物の形成地域が火山である可能性を認めても、変形の直接機構を炎や噴火としない。名前を覚えた後に、物質・応力・流れの機構を説明できることを学習目標とする。


\section{過去問と個別の解法}
\section{島弧沈み込み系と中央海嶺の用語を区別する。}

\subsection{平成25年度 問26}
\label{q:h25-reg-applied-26}

\noindent\textbf{原典：} sources/H25/17\_applied\_science.pdf、PDF物理10頁。

\textbf{公式正答：} 5（single）。

\paragraph{問題の要旨} 島弧沈み込み系と中央海嶺の用語を区別する。

\Needspace{4\baselineskip}\paragraph{確認した条件・記述（独自要約）}\mbox{}\par

\noindent\textbf{原典確認：} S7\_050の29原典画像で34問の全五肢・否定・数値を照合。\par

\noindent\textbf{対象：} 島弧沈み込み系と中央海嶺の用語を区別する。\par

\paragraph{問われる核心} 島弧沈み込み系と中央海嶺の用語を区別する。

\paragraph{必要知識} 海溝から陸側へ前弧・火山弧・背弧を並べ、スラブとその上のマントルウェッジを考える。中央海嶺は発散境界であり島弧沈み込み系と直接の用語が違うので5。背弧海盆にも拡大軸があり得るが、地球規模の中央海嶺を島弧そのものへ交換しない。

\paragraph{考え方・正答を選ぶ理由（査読待ち初稿）} 海溝から陸側へ前弧・火山弧・背弧を並べ、スラブとその上のマントルウェッジを考える。中央海嶺は発散境界であり島弧沈み込み系と直接の用語が違うので5。背弧海盆にも拡大軸があり得るが、地球規模の中央海嶺を島弧そのものへ交換しない。

\Needspace{4\baselineskip}\paragraph{選択肢の検討（査読待ち初稿）}\mbox{}\par

\noindent\textbf{1：代表的条件で適切} マントルウェッジ。沈み込むスラブの上のくさび状マントルで、流体供給と島弧マグマ生成に関係する。\par

\noindent\textbf{2：代表的条件で適切} 背弧。火山弧の海溝と反対側で、伸張する背弧盆地等が形成され得る。\par

\noindent\textbf{3：代表的条件で適切} 火山前線。島弧火山の分布の海溝側限界を示す線。海溝そのものと違う。\par

\noindent\textbf{4：代表的条件で適切} 沈み込み帯。収束するプレートが地下へ沈み込む島弧活動の主要な生成場。\par

\noindent\textbf{5：不適切} 中央海嶺。プレートが離れ、新しい海洋地殻を作る発散境界。島弧の直接構成用語ではない。\par

\paragraph{再利用する考え方} 地質環境・生成機構・物質の名称を対応させ、平均値と全称、典型例と例外、当時の分類と現在の知見を分ける。

\paragraph{共通理論} \hyperref[theory:volcano:arc]{第\ref*{theory:volcano:arc}節}

\par 参考：\cite{usgs_plate_boundaries}。

\par 難易度 3（校準初稿）、優先度未評価。

\subsection{令和元年度再試験 問23}
\label{q:r01-ret-applied-23}

\noindent\textbf{原典：} sources/R01\_retest/17\_applied\_science.pdf、PDF物理8頁。

\textbf{公式正答：} 3（single）。

\par\textbf{検討未確定：} 公式3。背弧方向のアルカリと水の濃度は別の変数で、水まで一般に増加と改変しない。全マグマの上部マントル起源という全称と温度の全噴火への適用も留保。アルカリ傾向の原論文による追加照合が必要。

\paragraph{問題の要旨} 火山前線・背弧の組成傾向とマグマ温度を比較する。

\Needspace{4\baselineskip}\paragraph{確認した条件・記述（独自要約）}\mbox{}\par

\noindent\textbf{原典確認：} S7\_050の29原典画像で34問の全五肢・否定・数値を照合。\par

\noindent\textbf{対象：} 火山前線・背弧の組成傾向とマグマ温度を比較する。\par

\paragraph{問われる核心} 火山前線・背弧の組成傾向とマグマ温度を比較する。

\paragraph{必要知識} 火山前線の海溝と反対側へ進むと、玄武岩中のアルカリ濃度が一般に高まる傾向がある。原典3は小さくなるとするため不適切。ただし水の濃度を常にアルカリと同方向としない。GSJの東北日本弧の調査には含水量が背弧へ減る分布の報告があり、アルカリと水を分けて比較する。

\paragraph{考え方・正答を選ぶ理由（査読待ち初稿）} 火山前線の海溝と反対側へ進むと、玄武岩中のアルカリ濃度が一般に高まる傾向がある。原典3は小さくなるとするため不適切。ただし水の濃度を常にアルカリと同方向としない。GSJの東北日本弧の調査には含水量が背弧へ減る分布の報告があり、アルカリと水を分けて比較する。

\Needspace{4\baselineskip}\paragraph{選択肢の検討（査読待ち初稿）}\mbox{}\par

\noindent\textbf{1：条件・解釈を留保} マントル部分融解によるマグマ。初生的な玄武岩質マグマの主要な生成経路。地殻の融解からもマグマが作られるので全種類の唯一の経路とはしない。\par

\noindent\textbf{2：代表的条件で適切} 海嶺・ホットスポットは主に玄武岩質。代表的な初生噴出物は玄武岩質。分化した溶岩等の例外はある。\par

\noindent\textbf{3：不適切} 背弧へアルカリと水が系統的に低下。島弧の代表的アルカリ濃度は背弧へ増す傾向で、アルカリと水を同方向の普遍則へ結び付けた記述は適切でない。水の傾向には別の検証が必要。\par

\noindent\textbf{4：代表的条件で適切} 島弧にはシリカに富む分化マグマもある。安山岩・デイサイト・流紋岩等も噴出し、初生玄武岩だけではない。\par

\noindent\textbf{5：代表的条件で適切} 玄武岩1000–1200℃・流紋岩700–900℃。JMAの1999年表にも記載された代表的な噴出温度範囲で、全噴火の厳密な境界ではない。\par

\paragraph{再利用する考え方} 地質環境・生成機構・物質の名称を対応させ、平均値と全称、典型例と例外、当時の分類と現在の知見を分ける。

\paragraph{共通理論} \hyperref[theory:volcano:arc]{第\ref*{theory:volcano:arc}節}

\par 参考：\cite{vol50_melt_ref,jma_magma1999,vol50_arcwater_ref}。

\par 難易度 4（校準初稿）、優先度未評価。

\subsection{令和6年度 問25}
\label{q:r06-reg-applied-25}

\noindent\textbf{原典：} sources/R06/17\_applied\_science.pdf、PDF物理8頁。

\textbf{公式正答：} 5（single）。

\par\textbf{検討未確定：} 原典4の島弧・陸弧を海の有無だけで定義する単純化と、前弧に火山が全くないという全称は留保。公式5の地震活動の誤りを保持。

\paragraph{問題の要旨} 海溝・前弧・火山弧・背弧と地震活動を整理する。

\Needspace{4\baselineskip}\paragraph{確認した条件・記述（独自要約）}\mbox{}\par

\noindent\textbf{原典確認：} S7\_050の29原典画像で34問の全五肢・否定・数値を照合。\par

\noindent\textbf{対象：} 海溝・前弧・火山弧・背弧と地震活動を整理する。\par

\paragraph{問われる核心} 海溝・前弧・火山弧・背弧と地震活動を整理する。

\paragraph{必要知識} 沈み込むプレートを地下まで追うと、火山前線より海溝側でも地震が起きる。火山分布と地震分布の範囲を同一視した5が不適切。海溝→前弧→火山弧→背弧の配列を地表と断面で対応させる。

\paragraph{考え方・正答を選ぶ理由（査読待ち初稿）} 沈み込むプレートを地下まで追うと、火山前線より海溝側でも地震が起きる。火山分布と地震分布の範囲を同一視した5が不適切。海溝→前弧→火山弧→背弧の配列を地表と断面で対応させる。

\Needspace{4\baselineskip}\paragraph{選択肢の検討（査読待ち初稿）}\mbox{}\par

\noindent\textbf{1：代表的条件で適切} 海洋側に海溝。島弧海溝の基本配列では海溝が沈み込みの海洋側にある。\par

\noindent\textbf{2：代表的条件で適切} 陸側へ前弧・火山弧・背弧。海溝から火山前線までの前弧、火山弧、その背後の背弧を順に読む。\par

\noindent\textbf{3：代表的条件で適切} 背弧で海が形成され得る。伸張して背弧盆地・海が形成される場合がある。必ず拡大海盆を持つわけではない。\par

\noindent\textbf{4：条件・解釈を留保} 背後の海による島弧・陸弧の呼称。過去問の基本図式は大陸から海で隔たる島弧と大陸縁の陸弧。海の有無だけで全遷移形態を排他的に分類する厳密定義には留保する。\par

\noindent\textbf{5：不適切} 前弧は地震活動も静穏。沈み込み境界やスラブに関係する地震があり、火山が少ないことから地震も静穏とは言えない。\par

\paragraph{再利用する考え方} 地質環境・生成機構・物質の名称を対応させ、平均値と全称、典型例と例外、当時の分類と現在の知見を分ける。

\paragraph{共通理論} \hyperref[theory:volcano:arc]{第\ref*{theory:volcano:arc}節}

\par 参考：\cite{usgs_plate_boundaries,vol50_subduction_ref}。

\par 難易度 4（校準初稿）、優先度未評価。

\section{ハワイ式とストロンボリ式・火砕流を区別する。}

\subsection{平成25年度 問25}
\label{q:h25-reg-applied-25}

\noindent\textbf{原典：} sources/H25/17\_applied\_science.pdf、PDF物理9頁。

\textbf{公式正答：} 3（single）。

\paragraph{問題の要旨} ハワイ式とストロンボリ式・火砕流を区別する。

\Needspace{4\baselineskip}\paragraph{確認した条件・記述（独自要約）}\mbox{}\par

\noindent\textbf{原典確認：} S7\_050の29原典画像で34問の全五肢・否定・数値を照合。\par

\noindent\textbf{対象：} ハワイ式とストロンボリ式・火砕流を区別する。\par

\paragraph{問われる核心} ハワイ式とストロンボリ式・火砕流を区別する。

\paragraph{必要知識} マグマの粘性とガス放出の仕方を対応させる。ハワイ式の代表は低粘性の玄武岩質で、流紋岩とした3が誤り。ストロンボリ式の間欠爆発、火砕流の高速なガス粒子混合、外来水との爆発、ドーム形成をそれぞれ別の現象として整理する。

\paragraph{考え方・正答を選ぶ理由（査読待ち初稿）} マグマの粘性とガス放出の仕方を対応させる。ハワイ式の代表は低粘性の玄武岩質で、流紋岩とした3が誤り。ストロンボリ式の間欠爆発、火砕流の高速なガス粒子混合、外来水との爆発、ドーム形成をそれぞれ別の現象として整理する。

\Needspace{4\baselineskip}\paragraph{選択肢の検討（査読待ち初稿）}\mbox{}\par

\noindent\textbf{1：代表的条件で適切} 間欠的な火山弾・スコリア放出。ストロンボリ式は比較的低粘性のマグマからガスが断続的に解放される小爆発を伴う。\par

\noindent\textbf{2：代表的条件で適切} 高温ガスと火砕物の高速流下。火砕流は高温のガス・粒子混合流で、地表を高速に移動する。\par

\noindent\textbf{3：不適切} 低粘性流紋岩のハワイ式。典型的なハワイ式は流動性に富む玄武岩質マグマの溶岩噴泉等。流紋岩へ交換した点が誤り。\par

\noindent\textbf{4：代表的条件で適切} 外来水とマグマの爆発的相互作用。マグマ水蒸気爆発では急な熱交換・水の蒸発・破砕が関与する。マグマ由来の揮発性分だけの発泡と分ける。\par

\noindent\textbf{5：代表的条件で適切} 高粘性溶岩のドーム。安山岩質～デイサイト質等の高粘性溶岩は火口近くに堆積してドームを作り得る。組成だけで全ての形を固定しない。\par

\paragraph{再利用する考え方} 地質環境・生成機構・物質の名称を対応させ、平均値と全称、典型例と例外、当時の分類と現在の知見を分ける。

\paragraph{共通理論} \hyperref[theory:volcano:styles]{第\ref*{theory:volcano:styles}節}

\par 参考：\cite{jma_magma1999,vol50_styles_ref}。

\paragraph{関連問題} \hyperref[q:h30-reg-applied-29]{H30-REG-APPLIED-29}

\par 難易度 3（校準初稿）、優先度未評価。

\subsection{平成30年度 問29}
\label{q:h30-reg-applied-29}

\noindent\textbf{原典：} sources/H30/17\_applied\_science.pdf、PDF物理10頁。

\textbf{公式正答：} 3（single）。

\paragraph{問題の要旨} ハワイ式とストロンボリ式・火砕流を区別する。

\Needspace{4\baselineskip}\paragraph{確認した条件・記述（独自要約）}\mbox{}\par

\noindent\textbf{原典確認：} S7\_050の29原典画像で34問の全五肢・否定・数値を照合。\par

\noindent\textbf{対象：} ハワイ式とストロンボリ式・火砕流を区別する。\par

\paragraph{問われる核心} ハワイ式とストロンボリ式・火砕流を区別する。

\paragraph{必要知識} マグマの粘性とガス放出の仕方を対応させる。ハワイ式の代表は低粘性の玄武岩質で、流紋岩とした3が誤り。ストロンボリ式の間欠爆発、火砕流の高速なガス粒子混合、外来水との爆発、ドーム形成をそれぞれ別の現象として整理する。

\paragraph{考え方・正答を選ぶ理由（査読待ち初稿）} マグマの粘性とガス放出の仕方を対応させる。ハワイ式の代表は低粘性の玄武岩質で、流紋岩とした3が誤り。ストロンボリ式の間欠爆発、火砕流の高速なガス粒子混合、外来水との爆発、ドーム形成をそれぞれ別の現象として整理する。

\Needspace{4\baselineskip}\paragraph{選択肢の検討（査読待ち初稿）}\mbox{}\par

\noindent\textbf{1：代表的条件で適切} 間欠的な火山弾・スコリア放出。ストロンボリ式は比較的低粘性のマグマからガスが断続的に解放される小爆発を伴う。\par

\noindent\textbf{2：代表的条件で適切} 高温ガスと火砕物の高速流下。火砕流は高温のガス・粒子混合流で、地表を高速に移動する。\par

\noindent\textbf{3：不適切} 低粘性流紋岩のハワイ式。典型的なハワイ式は流動性に富む玄武岩質マグマの溶岩噴泉等。流紋岩へ交換した点が誤り。\par

\noindent\textbf{4：代表的条件で適切} 外来水とマグマの爆発的相互作用。マグマ水蒸気爆発では急な熱交換・水の蒸発・破砕が関与する。マグマ由来の揮発性分だけの発泡と分ける。\par

\noindent\textbf{5：代表的条件で適切} 高粘性溶岩のドーム。安山岩質～デイサイト質等の高粘性溶岩は火口近くに堆積してドームを作り得る。組成だけで全ての形を固定しない。\par

\paragraph{再利用する考え方} 地質環境・生成機構・物質の名称を対応させ、平均値と全称、典型例と例外、当時の分類と現在の知見を分ける。

\paragraph{共通理論} \hyperref[theory:volcano:styles]{第\ref*{theory:volcano:styles}節}

\par 参考：\cite{jma_magma1999,vol50_styles_ref}。

\paragraph{関連問題} \hyperref[q:h25-reg-applied-25]{H25-REG-APPLIED-25}

\par 難易度 3（校準初稿）、優先度未評価。

\subsection{令和元年度再試験 問24}
\label{q:r01-ret-applied-24}

\noindent\textbf{原典：} sources/R01\_retest/17\_applied\_science.pdf、PDF物理8頁。

\textbf{公式正答：} 5（single）。

\paragraph{問題の要旨} 噴火現象を気体・水・溶岩・岩屑の移動へ分ける。

\Needspace{4\baselineskip}\paragraph{確認した条件・記述（独自要約）}\mbox{}\par

\noindent\textbf{原典確認：} S7\_050の29原典画像で34問の全五肢・否定・数値を照合。\par

\noindent\textbf{対象：} 噴火現象を気体・水・溶岩・岩屑の移動へ分ける。\par

\paragraph{問われる核心} 噴火現象を気体・水・溶岩・岩屑の移動へ分ける。

\paragraph{必要知識} 流れる物質と速度を確認する。火砕流は高温ガスと火砕粒子の混合が高速に流下する現象であり5の低速度が誤り。外来水を含む爆発・泥流、液状溶岩、崩壊する岩屑の移動はそれぞれ異なる。

\paragraph{考え方・正答を選ぶ理由（査読待ち初稿）} 流れる物質と速度を確認する。火砕流は高温ガスと火砕粒子の混合が高速に流下する現象であり5の低速度が誤り。外来水を含む爆発・泥流、液状溶岩、崩壊する岩屑の移動はそれぞれ異なる。

\Needspace{4\baselineskip}\paragraph{選択肢の検討（査読待ち初稿）}\mbox{}\par

\noindent\textbf{1：代表的条件で適切} 液状マグマの流出。液体として流れ出て移動する溶岩流。気体粒子混合の火砕流と分ける。\par

\noindent\textbf{2：代表的条件で適切} ストロンボリの間欠爆発。比較的低粘性のマグマから気泡が断続的に放出される。\par

\noindent\textbf{3：代表的条件で適切} 外来水の急な水蒸気化。マグマと水の急な熱交換で蒸気の膨張・破砕が生じ得る。\par

\noindent\textbf{4：代表的条件で適切} 崩壊山体の岩屑なだれ。不安定な火山体が崩れ高速に移動する場合がある。\par

\noindent\textbf{5：不適切} 火砕流は時速数km以下。一般的な火砕流は時速数十kmから100km以上等の高速で、ゆっくりした数km以下の流れという説明が不適切。\par

\paragraph{再利用する考え方} 地質環境・生成機構・物質の名称を対応させ、平均値と全称、典型例と例外、当時の分類と現在の知見を分ける。

\paragraph{共通理論} \hyperref[theory:volcano:styles]{第\ref*{theory:volcano:styles}節}

\par 参考：\cite{jma_magma1999,volc_usgs_msh_flow}。

\paragraph{関連問題} \hyperref[q:r03-reg-applied-26]{R03-REG-APPLIED-26}

\par 難易度 3（校準初稿）、優先度未評価。

\subsection{令和3年度 問26}
\label{q:r03-reg-applied-26}

\noindent\textbf{原典：} sources/R03/17\_applied\_science.pdf、PDF物理9頁。

\textbf{公式正答：} 3（single）。

\paragraph{問題の要旨} ストロンボリの間欠爆発と連続溶岩噴泉を区別する。

\Needspace{4\baselineskip}\paragraph{確認した条件・記述（独自要約）}\mbox{}\par

\noindent\textbf{原典確認：} S7\_050の29原典画像で34問の全五肢・否定・数値を照合。\par

\noindent\textbf{対象：} ストロンボリの間欠爆発と連続溶岩噴泉を区別する。\par

\paragraph{問われる核心} ストロンボリの間欠爆発と連続溶岩噴泉を区別する。

\paragraph{必要知識} 間欠か連続かを原典から読む。ストロンボリ式は気泡の解放による間欠的小爆発で、割れ目からの連続的噴泉とした3が誤り。平成・令和の似た五肢でも、誤りが岩種か速度か噴出様式かを確かめる。

\paragraph{考え方・正答を選ぶ理由（査読待ち初稿）} 間欠か連続かを原典から読む。ストロンボリ式は気泡の解放による間欠的小爆発で、割れ目からの連続的噴泉とした3が誤り。平成・令和の似た五肢でも、誤りが岩種か速度か噴出様式かを確かめる。

\Needspace{4\baselineskip}\paragraph{選択肢の検討（査読待ち初稿）}\mbox{}\par

\noindent\textbf{1：代表的条件で適切} 時速数十から100km以上の火砕流。高温のガス粒子混合流の代表的高速移動。速度の全現象一律境界ではない。\par

\noindent\textbf{2：代表的条件で適切} 液状溶岩流。破片の混合ではなく液体としてのマグマが流出する。\par

\noindent\textbf{3：不適切} 連続噴泉をストロンボリ式とする。連続的な溶岩噴泉は代表的にハワイ式。ストロンボリは間欠的な爆発が特徴。\par

\noindent\textbf{4：代表的条件で適切} 山体崩壊の岩屑なだれ。火山体の不安定部が崩れ高速に流下する。\par

\noindent\textbf{5：代表的条件で適切} 外来水の相変化を伴う爆発。マグマ水蒸気爆発で急な熱交換・膨張等が起こる。\par

\paragraph{再利用する考え方} 地質環境・生成機構・物質の名称を対応させ、平均値と全称、典型例と例外、当時の分類と現在の知見を分ける。

\paragraph{共通理論} \hyperref[theory:volcano:styles]{第\ref*{theory:volcano:styles}節}

\par 参考：\cite{jma_magma1999,vol50_styles_ref,volc_usgs_msh_flow}。

\paragraph{関連問題} \hyperref[q:r01-ret-applied-24]{R01-RET-APPLIED-24}

\par 難易度 3（校準初稿）、優先度未評価。

\section{火山ガスの主要成分と温度・酸化還元条件を区別する。}

\subsection{平成28年度 問28}
\label{q:h28-reg-applied-28}

\noindent\textbf{原典：} sources/H28/17\_applied\_science.pdf、PDF物理11頁。

\textbf{公式正答：} 5（single）。

\par\textbf{検討未確定：} 理論的平均と主要の厳密な濃度基準は未提示。特にR07の100気圧・1200℃平衡の定量表と酸化還元条件の照合が残る。通常の火山ガスの代表比較として説明し、CH4の全条件不在やH2の固定割合を主張しない。

\paragraph{問題の要旨} 火山ガスの主要成分と温度・酸化還元条件を区別する。

\Needspace{4\baselineskip}\paragraph{確認した条件・記述（独自要約）}\mbox{}\par

\noindent\textbf{原典確認：} S7\_050の29原典画像で34問の全五肢・否定・数値を照合。\par

\noindent\textbf{対象：} 火山ガスの主要成分と温度・酸化還元条件を区別する。\par

\paragraph{問われる核心} 火山ガスの主要成分と温度・酸化還元条件を区別する。

\paragraph{必要知識} 水蒸気、CO2、硫黄ガス、水素ハロゲン化物等が通常の火山ガスの代表で、候補中ではCH4を選ぶ。高温マグマのガスと低温熱水・堆積物由来のガスを分ける。令和7年度は100気圧・1200℃・玄武岩質の平衡条件が原典にあり、平成28年度・令和元再試験へその条件を遡って追加しない。

\paragraph{考え方・正答を選ぶ理由（査読待ち初稿）} 水蒸気、CO2、硫黄ガス、水素ハロゲン化物等が通常の火山ガスの代表で、候補中ではCH4を選ぶ。高温マグマのガスと低温熱水・堆積物由来のガスを分ける。令和7年度は100気圧・1200℃・玄武岩質の平衡条件が原典にあり、平成28年度・令和元再試験へその条件を遡って追加しない。

\Needspace{4\baselineskip}\paragraph{選択肢の検討（査読待ち初稿）}\mbox{}\par

\noindent\textbf{1：代表的条件で適切} H2O。通常の火山ガスの最多成分は水蒸気であり、酸性有毒ガスが量の大部分という意味ではない。\par

\noindent\textbf{2：代表的条件で適切} CO2。マグマから放出される代表的な炭素ガス。水蒸気との量比は環境に依存する。\par

\noindent\textbf{3：代表的条件で適切} SO2。硫黄ガスの代表で酸化的な高温火山ガスに重要。冷却・熱水系では化学種や吸収が変わる。\par

\noindent\textbf{4：代表的条件で適切} H2S。硫黄の還元的な化学種であり、火山ガス・熱水系に含まれる。\par

\noindent\textbf{5：不適切} CH4。通常の高温マグマに由来する平均主要ガスの比較では適さない。低温・還元環境の存在を含め、厳密に常に零とする意味ではない。\par

\paragraph{再利用する考え方} 地質環境・生成機構・物質の名称を対応させ、平均値と全称、典型例と例外、当時の分類と現在の知見を分ける。

\paragraph{共通理論} \hyperref[theory:volcano:gases]{第\ref*{theory:volcano:gases}節}

\par 参考：\cite{vol50_gas_ref}。

\paragraph{関連問題} \hyperref[q:r01-ret-applied-22]{R01-RET-APPLIED-22}

\par 難易度 4（校準初稿）、優先度未評価。

\subsection{令和元年度再試験 問22}
\label{q:r01-ret-applied-22}

\noindent\textbf{原典：} sources/R01\_retest/17\_applied\_science.pdf、PDF物理7頁。

\textbf{公式正答：} 4（single）。

\par\textbf{検討未確定：} 理論的平均と主要の厳密な濃度基準は未提示。特にR07の100気圧・1200℃平衡の定量表と酸化還元条件の照合が残る。通常の火山ガスの代表比較として説明し、CH4の全条件不在やH2の固定割合を主張しない。

\paragraph{問題の要旨} 火山ガスの主要成分と温度・酸化還元条件を区別する。

\Needspace{4\baselineskip}\paragraph{確認した条件・記述（独自要約）}\mbox{}\par

\noindent\textbf{原典確認：} S7\_050の29原典画像で34問の全五肢・否定・数値を照合。\par

\noindent\textbf{対象：} 火山ガスの主要成分と温度・酸化還元条件を区別する。\par

\paragraph{問われる核心} 火山ガスの主要成分と温度・酸化還元条件を区別する。

\paragraph{必要知識} 水蒸気、CO2、硫黄ガス、水素ハロゲン化物等が通常の火山ガスの代表で、候補中ではCH4を選ぶ。高温マグマのガスと低温熱水・堆積物由来のガスを分ける。令和7年度は100気圧・1200℃・玄武岩質の平衡条件が原典にあり、平成28年度・令和元再試験へその条件を遡って追加しない。

\paragraph{考え方・正答を選ぶ理由（査読待ち初稿）} 水蒸気、CO2、硫黄ガス、水素ハロゲン化物等が通常の火山ガスの代表で、候補中ではCH4を選ぶ。高温マグマのガスと低温熱水・堆積物由来のガスを分ける。令和7年度は100気圧・1200℃・玄武岩質の平衡条件が原典にあり、平成28年度・令和元再試験へその条件を遡って追加しない。

\Needspace{4\baselineskip}\paragraph{選択肢の検討（査読待ち初稿）}\mbox{}\par

\noindent\textbf{1：代表的条件で適切} CO2。マグマから放出される代表的な炭素ガス。水蒸気との量比は環境に依存する。\par

\noindent\textbf{2：代表的条件で適切} H2S。硫黄の還元的な化学種であり、火山ガス・熱水系に含まれる。\par

\noindent\textbf{3：代表的条件で適切} HCl。火山ガス中の代表的な水素ハロゲン化物。地下水への吸収で地表の濃度は変わり得る。\par

\noindent\textbf{4：不適切} CH4。通常の高温マグマに由来する平均主要ガスの比較では適さない。低温・還元環境の存在を含め、厳密に常に零とする意味ではない。\par

\noindent\textbf{5：代表的条件で適切} SO2。硫黄ガスの代表で酸化的な高温火山ガスに重要。冷却・熱水系では化学種や吸収が変わる。\par

\paragraph{再利用する考え方} 地質環境・生成機構・物質の名称を対応させ、平均値と全称、典型例と例外、当時の分類と現在の知見を分ける。

\paragraph{共通理論} \hyperref[theory:volcano:gases]{第\ref*{theory:volcano:gases}節}

\par 参考：\cite{vol50_gas_ref}。

\paragraph{関連問題} \hyperref[q:h28-reg-applied-28]{H28-REG-APPLIED-28}, \hyperref[q:r07-reg-applied-20]{R07-REG-APPLIED-20}

\par 難易度 4（校準初稿）、優先度未評価。

\subsection{令和7年度 問20}
\label{q:r07-reg-applied-20}

\noindent\textbf{原典：} sources/R07/17\_applied\_science.pdf、PDF物理7頁。

\textbf{公式正答：} 2（single）。

\par\textbf{検討未確定：} 理論的平均と主要の厳密な濃度基準は未提示。特にR07の100気圧・1200℃平衡の定量表と酸化還元条件の照合が残る。通常の火山ガスの代表比較として説明し、CH4の全条件不在やH2の固定割合を主張しない。

\paragraph{問題の要旨} 火山ガスの主要成分と温度・酸化還元条件を区別する。

\Needspace{4\baselineskip}\paragraph{確認した条件・記述（独自要約）}\mbox{}\par

\noindent\textbf{原典確認：} S7\_050の29原典画像で34問の全五肢・否定・数値を照合。\par

\noindent\textbf{対象：} 火山ガスの主要成分と温度・酸化還元条件を区別する。\par

\paragraph{問われる核心} 火山ガスの主要成分と温度・酸化還元条件を区別する。

\paragraph{必要知識} 水蒸気、CO2、硫黄ガス、水素ハロゲン化物等が通常の火山ガスの代表で、候補中ではCH4を選ぶ。高温マグマのガスと低温熱水・堆積物由来のガスを分ける。令和7年度は100気圧・1200℃・玄武岩質の平衡条件が原典にあり、平成28年度・令和元再試験へその条件を遡って追加しない。

\paragraph{考え方・正答を選ぶ理由（査読待ち初稿）} 水蒸気、CO2、硫黄ガス、水素ハロゲン化物等が通常の火山ガスの代表で、候補中ではCH4を選ぶ。高温マグマのガスと低温熱水・堆積物由来のガスを分ける。令和7年度は100気圧・1200℃・玄武岩質の平衡条件が原典にあり、平成28年度・令和元再試験へその条件を遡って追加しない。

\Needspace{4\baselineskip}\paragraph{選択肢の検討（査読待ち初稿）}\mbox{}\par

\noindent\textbf{1：代表的条件で適切} CO2。マグマから放出される代表的な炭素ガス。水蒸気との量比は環境に依存する。\par

\noindent\textbf{2：不適切} CH4。通常の高温マグマに由来する平均主要ガスの比較では適さない。低温・還元環境の存在を含め、厳密に常に零とする意味ではない。\par

\noindent\textbf{3：条件・解釈を留保} H2。高温ガスに含まれる還元的な水素種。量は温度・圧力だけでなく酸化還元・全体組成等にも依存する。\par

\noindent\textbf{4：代表的条件で適切} HCl。火山ガス中の代表的な水素ハロゲン化物。地下水への吸収で地表の濃度は変わり得る。\par

\noindent\textbf{5：代表的条件で適切} SO2。硫黄ガスの代表で酸化的な高温火山ガスに重要。冷却・熱水系では化学種や吸収が変わる。\par

\paragraph{再利用する考え方} 地質環境・生成機構・物質の名称を対応させ、平均値と全称、典型例と例外、当時の分類と現在の知見を分ける。

\paragraph{共通理論} \hyperref[theory:volcano:gases]{第\ref*{theory:volcano:gases}節}

\par 参考：\cite{vol50_gas_ref}。

\paragraph{関連問題} \hyperref[q:r01-ret-applied-22]{R01-RET-APPLIED-22}

\par 難易度 4（校準初稿）、優先度未評価。

\section{火山ガス・気候影響・火砕流・泥流・津波を比較する。}

\subsection{平成28年度 問35}
\label{q:h28-reg-applied-35}

\noindent\textbf{原典：} sources/H28/17\_applied\_science.pdf、PDF物理14頁。

\textbf{公式正答：} 2（single）。

\par\textbf{検討未確定：} 公式2を保持。肢1の全球冷却を大量火山灰へ帰す説明には、灰の短期影響と硫酸エアロゾルの長期影響の区別が必要。噴煙20–40kmも代表範囲。

\paragraph{問題の要旨} 火山ガス・気候影響・火砕流・泥流・津波を比較する。

\Needspace{4\baselineskip}\paragraph{確認した条件・記述（独自要約）}\mbox{}\par

\noindent\textbf{原典確認：} S7\_050の29原典画像で34問の全五肢・否定・数値を照合。\par

\noindent\textbf{対象：} 火山ガス・気候影響・火砕流・泥流・津波を比較する。\par

\paragraph{問われる核心} 火山ガス・気候影響・火砕流・泥流・津波を比較する。

\paragraph{必要知識} 火山ガスの量と有害性を分ける。通常最大成分は水蒸気なので2が誤り。火砕流はガス粒子、泥流は水粒子の混合、津波は海水の変位が中心である。1の気候説明は硫黄ガスから作る硫酸エアロゾルの長期的役割を補足し、灰だけで全球冷却を説明しない。

\paragraph{考え方・正答を選ぶ理由（査読待ち初稿）} 火山ガスの量と有害性を分ける。通常最大成分は水蒸気なので2が誤り。火砕流はガス粒子、泥流は水粒子の混合、津波は海水の変位が中心である。1の気候説明は硫黄ガスから作る硫酸エアロゾルの長期的役割を補足し、灰だけで全球冷却を説明しない。

\Needspace{4\baselineskip}\paragraph{選択肢の検討（査読待ち初稿）}\mbox{}\par

\noindent\textbf{1：条件・解釈を留保} 高い噴煙と灰による世界冷却。高い噴煙と日射遮蔽はあり得るが、長距離・長期間の全球冷却は主にSO2由来の硫酸エアロゾル。火山灰だけの説明を一般化しない。\par

\noindent\textbf{2：不適切} 大部分が有害な酸性ガス。量の大部分は通常水蒸気。少量の有害ガスの危険性と量の割合を混同している。\par

\noindent\textbf{3：代表的条件で適切} 高温ガスと粒子の火砕流。ガスと火砕粒子の混合流が高速に移動する。\par

\noindent\textbf{4：代表的条件で適切} 融雪・湖水等を取り込む火山泥流。水と火山砕屑物が混ざり、土石流状に流下する。噴火と同時以外の降雨でも発生し得る。\par

\noindent\textbf{5：代表的条件で適切} 火山島噴火や海への岩屑で津波。水塊の急激な変位により津波を起こす場合がある。必ず発生するという全称ではない。\par

\paragraph{再利用する考え方} 地質環境・生成機構・物質の名称を対応させ、平均値と全称、典型例と例外、当時の分類と現在の知見を分ける。

\paragraph{共通理論} \hyperref[theory:volcano:hazards]{第\ref*{theory:volcano:hazards}節}

\par 参考：\cite{vol50_gas_ref,vol50_climate_ref,volc_usgs_msh_flow}。

\par 難易度 4（校準初稿）、優先度未評価。

\section{パホイホイ・アア・ブロックと溶岩形態を比較する。}

\subsection{令和2年度 問24}
\label{q:r02-reg-applied-24}

\noindent\textbf{原典：} sources/R02/17\_applied\_science.pdf、PDF物理9頁。

\textbf{公式正答：} 3（single）。

\par\textbf{検討未確定：} 原典のアスペクト比1/8の分類出典・測定長さと、アアの低温・組成に関する全例の一般化は追加照合。パホイホイを単に急冷と定義しない。

\paragraph{問題の要旨} パホイホイ・アア・ブロックと溶岩形態を比較する。

\Needspace{4\baselineskip}\paragraph{確認した条件・記述（独自要約）}\mbox{}\par

\noindent\textbf{原典確認：} S7\_050の29原典画像で34問の全五肢・否定・数値を照合。\par

\noindent\textbf{対象：} パホイホイ・アア・ブロックと溶岩形態を比較する。\par

\paragraph{問われる核心} パホイホイ・アア・ブロックと溶岩形態を比較する。

\paragraph{必要知識} パホイホイとアアは同じ玄武岩質でも流動・冷却・変形の条件で異なる形になり得る。パホイホイを流紋岩質へ置換した3が誤り。滑らかな流れ、粗いクリンカー、大きな岩塊の表面を岩種と別の形態軸として整理する。

\paragraph{考え方・正答を選ぶ理由（査読待ち初稿）} パホイホイとアアは同じ玄武岩質でも流動・冷却・変形の条件で異なる形になり得る。パホイホイを流紋岩質へ置換した3が誤り。滑らかな流れ、粗いクリンカー、大きな岩塊の表面を岩種と別の形態軸として整理する。

\Needspace{4\baselineskip}\paragraph{選択肢の検討（査読待ち初稿）}\mbox{}\par

\noindent\textbf{1：条件・解釈を留保} 厚さ／長さ1/8で流れとドームを分ける。原典の形態分類の目安を保持する。普遍的な流体力学上の相転移閾値ではなく、測り方と分類文献の追加確認が必要。\par

\noindent\textbf{2：代表的条件で適切} 粘性・冷却・流速で表面が変化。組成・温度・結晶量等が流動性と表面破砕へ影響する。\par

\noindent\textbf{3：不適切} パホイホイは流紋岩質の急冷。滑らか・縄状等の表面を持つパホイホイは代表的に流動性に富む玄武岩質溶岩。岩種を交換している。\par

\noindent\textbf{4：代表的条件で適切} アアは粗いクリンカー表面。相対的に高粘性・大きい変形速度等で表面が粗い破片となる。温度だけで状態を決めない。\par

\noindent\textbf{5：代表的条件で適切} ブロックは大きな平滑破断面の岩塊。高粘性の中間質～珪長質等で、アアの小さい粗いクリンカーと区別する表面形態。\par

\paragraph{再利用する考え方} 地質環境・生成機構・物質の名称を対応させ、平均値と全称、典型例と例外、当時の分類と現在の知見を分ける。

\paragraph{共通理論} \hyperref[theory:volcano:lava]{第\ref*{theory:volcano:lava}節}

\par 参考：\cite{vol50_lava_ref,jma_magma1999}。

\par 難易度 4（校準初稿）、優先度未評価。

\section{マグマの局所生成・部分溶融・組成・発泡を比較する。}

\subsection{平成29年度 問29}
\label{q:h29-reg-applied-29}

\noindent\textbf{原典：} sources/H29/17\_applied\_science.pdf、PDF物理10頁。

\textbf{公式正答：} 3（single）。

\par\textbf{検討未確定：} 公式H29=\allowbreak{}3/R01=\allowbreak{}2。富むの比較対象が原典に明記されないため、原岩に対するNa・K等の濃集と珪長質マグマに対する組成差を区別する。発泡の絶対圧力と浮力だけによる貯留の全称も留保。

\paragraph{問題の要旨} マグマの局所生成・部分溶融・組成・発泡を比較する。

\Needspace{4\baselineskip}\paragraph{確認した条件・記述（独自要約）}\mbox{}\par

\noindent\textbf{原典確認：} S7\_050の29原典画像で34問の全五肢・否定・数値を照合。\par

\noindent\textbf{対象：} マグマの局所生成・部分溶融・組成・発泡を比較する。\par

\paragraph{問われる核心} マグマの局所生成・部分溶融・組成・発泡を比較する。

\paragraph{必要知識} 地球内部の多くは固体で、固相線を越えた条件で部分溶融する。典型的なマントルの融解液は玄武岩質であり、流紋岩等と比較すればSiO2とアルカリが少ない側なので、原典でSi・Na・Kに富むとする肢を公式は選ぶ。ただし比較を原岩マントルへ変えると不適合元素の濃集があるため、比較対象を明示して説明する。気泡の膨張と圧力の単調増加を同一視せず、マグマ貯留も中立浮力だけで全てを説明しない。

\paragraph{考え方・正答を選ぶ理由（査読待ち初稿）} 地球内部の多くは固体で、固相線を越えた条件で部分溶融する。典型的なマントルの融解液は玄武岩質であり、流紋岩等と比較すればSiO2とアルカリが少ない側なので、原典でSi・Na・Kに富むとする肢を公式は選ぶ。ただし比較を原岩マントルへ変えると不適合元素の濃集があるため、比較対象を明示して説明する。気泡の膨張と圧力の単調増加を同一視せず、マグマ貯留も中立浮力だけで全てを説明しない。

\Needspace{4\baselineskip}\paragraph{選択肢の検討（査読待ち初稿）}\mbox{}\par

\noindent\textbf{1：代表的条件で適切} 地殻・マントルは通常固体。地下全体に液状マグマが遍在するわけではない。条件が固相線を越える局所で融解が始まる。\par

\noindent\textbf{2：代表的条件で適切} 一部が融ける部分溶融。多成分の岩石では溶けやすい成分からメルトが生じ、残留固相と液相が共存する。\par

\noindent\textbf{3：条件・解釈を留保} 玄武岩質はSi・Na・Kに富む。流紋岩等と比較する通常の分類では玄武岩はSiO2・アルカリが少ない側。ただしマントル原岩に比べれば不適合元素がメルトへ濃集するため、比較対象なしの富むは曖昧。\par

\noindent\textbf{4：条件・解釈を留保} 浮力上昇とマグマだまり。通常メルトは周囲より低密度で上昇し得るが、貯留は中立浮力だけでなく応力・層境界・冷却等にも影響される。\par

\noindent\textbf{5：条件・解釈を留保} 減圧で揮発性分が分離し圧力が高まる。減圧で溶解度が下がり気泡が生じ膨張する。閉塞等で過圧を作る場合はあるが、外圧低下からガスの絶対圧力が必ず上昇するとは限らない。\par

\paragraph{再利用する考え方} 地質環境・生成機構・物質の名称を対応させ、平均値と全称、典型例と例外、当時の分類と現在の知見を分ける。

\paragraph{共通理論} \hyperref[theory:volcano:melting]{第\ref*{theory:volcano:melting}節}

\par 参考：\cite{vol50_melt_ref,jma_magma1999,vol50_degassing_ref}。

\paragraph{関連問題} \hyperref[q:r01-reg-applied-23]{R01-REG-APPLIED-23}

\par 難易度 4（校準初稿）、優先度未評価。

\subsection{令和元年度 問23}
\label{q:r01-reg-applied-23}

\noindent\textbf{原典：} sources/R01/17\_applied\_science.pdf、PDF物理8頁。

\textbf{公式正答：} 2（single）。

\par\textbf{検討未確定：} 公式H29=\allowbreak{}3/R01=\allowbreak{}2。富むの比較対象が原典に明記されないため、原岩に対するNa・K等の濃集と珪長質マグマに対する組成差を区別する。発泡の絶対圧力と浮力だけによる貯留の全称も留保。

\paragraph{問題の要旨} マグマの局所生成・部分溶融・組成・発泡を比較する。

\Needspace{4\baselineskip}\paragraph{確認した条件・記述（独自要約）}\mbox{}\par

\noindent\textbf{原典確認：} S7\_050の29原典画像で34問の全五肢・否定・数値を照合。\par

\noindent\textbf{対象：} マグマの局所生成・部分溶融・組成・発泡を比較する。\par

\paragraph{問われる核心} マグマの局所生成・部分溶融・組成・発泡を比較する。

\paragraph{必要知識} 地球内部の多くは固体で、固相線を越えた条件で部分溶融する。典型的なマントルの融解液は玄武岩質であり、流紋岩等と比較すればSiO2とアルカリが少ない側なので、原典でSi・Na・Kに富むとする肢を公式は選ぶ。ただし比較を原岩マントルへ変えると不適合元素の濃集があるため、比較対象を明示して説明する。気泡の膨張と圧力の単調増加を同一視せず、マグマ貯留も中立浮力だけで全てを説明しない。

\paragraph{考え方・正答を選ぶ理由（査読待ち初稿）} 地球内部の多くは固体で、固相線を越えた条件で部分溶融する。典型的なマントルの融解液は玄武岩質であり、流紋岩等と比較すればSiO2とアルカリが少ない側なので、原典でSi・Na・Kに富むとする肢を公式は選ぶ。ただし比較を原岩マントルへ変えると不適合元素の濃集があるため、比較対象を明示して説明する。気泡の膨張と圧力の単調増加を同一視せず、マグマ貯留も中立浮力だけで全てを説明しない。

\Needspace{4\baselineskip}\paragraph{選択肢の検討（査読待ち初稿）}\mbox{}\par

\noindent\textbf{1：条件・解釈を留保} 減圧で揮発性分が分離し圧力が高まる。減圧で溶解度が下がり気泡が生じ膨張する。閉塞等で過圧を作る場合はあるが、外圧低下からガスの絶対圧力が必ず上昇するとは限らない。\par

\noindent\textbf{2：条件・解釈を留保} 玄武岩質はSi・Na・Kに富む。流紋岩等と比較する通常の分類では玄武岩はSiO2・アルカリが少ない側。ただしマントル原岩に比べれば不適合元素がメルトへ濃集するため、比較対象なしの富むは曖昧。\par

\noindent\textbf{3：条件・解釈を留保} 浮力上昇とマグマだまり。通常メルトは周囲より低密度で上昇し得るが、貯留は中立浮力だけでなく応力・層境界・冷却等にも影響される。\par

\noindent\textbf{4：代表的条件で適切} 一部が融ける部分溶融。多成分の岩石では溶けやすい成分からメルトが生じ、残留固相と液相が共存する。\par

\noindent\textbf{5：代表的条件で適切} 地殻・マントルは通常固体。地下全体に液状マグマが遍在するわけではない。条件が固相線を越える局所で融解が始まる。\par

\paragraph{再利用する考え方} 地質環境・生成機構・物質の名称を対応させ、平均値と全称、典型例と例外、当時の分類と現在の知見を分ける。

\paragraph{共通理論} \hyperref[theory:volcano:melting]{第\ref*{theory:volcano:melting}節}

\par 参考：\cite{vol50_melt_ref,jma_magma1999,vol50_degassing_ref}。

\paragraph{関連問題} \hyperref[q:h29-reg-applied-29]{H29-REG-APPLIED-29}

\par 難易度 4（校準初稿）、優先度未評価。

\subsection{令和5年度 問25}
\label{q:r05-reg-applied-25}

\noindent\textbf{原典：} sources/R05/17\_applied\_science.pdf、PDF物理9頁。

\textbf{公式正答：} 3（single）。

\par\textbf{検討未確定：} 密度の釣合いだけで全マグマだまりを説明しない。R07原典1の融点は多成分岩石では固相線・液相線という範囲をもつ。公式R05=\allowbreak{}3/R07=\allowbreak{}2を保持して条件を補足。

\paragraph{問題の要旨} 減圧融解・含水固相線低下・部分融解を区別する。

\Needspace{4\baselineskip}\paragraph{確認した条件・記述（独自要約）}\mbox{}\par

\noindent\textbf{原典確認：} S7\_050の29原典画像で34問の全五肢・否定・数値を照合。\par

\noindent\textbf{対象：} 減圧融解・含水固相線低下・部分融解を区別する。\par

\paragraph{問われる核心} 減圧融解・含水固相線低下・部分融解を区別する。

\paragraph{必要知識} 海嶺等では上昇するマントルが減圧され、温度圧力経路が固相線を越えて融解する。低温物質の急な昇温だけを原因とした肢を選ぶ。沈み込みでは含水物質の脱水流体が固相線を下げる。純物質の一つの融点とは違い、岩石の部分溶融では固相とメルトが共存する。浮力上昇と貯留には周囲の構造や応力も関係する。

\paragraph{考え方・正答を選ぶ理由（査読待ち初稿）} 海嶺等では上昇するマントルが減圧され、温度圧力経路が固相線を越えて融解する。低温物質の急な昇温だけを原因とした肢を選ぶ。沈み込みでは含水物質の脱水流体が固相線を下げる。純物質の一つの融点とは違い、岩石の部分溶融では固相とメルトが共存する。浮力上昇と貯留には周囲の構造や応力も関係する。

\Needspace{4\baselineskip}\paragraph{選択肢の検討（査読待ち初稿）}\mbox{}\par

\noindent\textbf{1：代表的条件で適切} 昇温又は減圧で融解。温度圧力経路が固相線を越えると融解が始まる。全圧力域で同一勾配の融点を仮定しない。\par

\noindent\textbf{2：代表的条件で適切} 残留固相のある部分溶融。一部の成分が融け、他の固相が残る多成分岩石の挙動。\par

\noindent\textbf{3：不適切} 低温マントルの急な昇温が海嶺・ホットスポット成因。代表的には高温マントルの上昇に伴う減圧が融解を促す。ホットスポットの熱異常はあるが、冷たい物質の急加熱だけという説明ではない。\par

\noindent\textbf{4：代表的条件で適切} スラブ由来水による固相線低下。沈み込む含水物質の脱水で流体が供給され、上のマントルの融解が促進される。\par

\noindent\textbf{5：条件・解釈を留保} 低密度メルトの浮力と貯留。浮力は重要だが、貯留の位置は密度差だけでなく応力・構造・結晶化等に依存する。\par

\paragraph{再利用する考え方} 地質環境・生成機構・物質の名称を対応させ、平均値と全称、典型例と例外、当時の分類と現在の知見を分ける。

\paragraph{共通理論} \hyperref[theory:volcano:melting]{第\ref*{theory:volcano:melting}節}

\par 参考：\cite{vol50_melt_ref,vol50_subduction_ref}。

\paragraph{関連問題} \hyperref[q:r07-reg-applied-26]{R07-REG-APPLIED-26}

\par 難易度 4（校準初稿）、優先度未評価。

\subsection{令和7年度 問26}
\label{q:r07-reg-applied-26}

\noindent\textbf{原典：} sources/R07/17\_applied\_science.pdf、PDF物理10頁。

\textbf{公式正答：} 2（single）。

\par\textbf{検討未確定：} 密度の釣合いだけで全マグマだまりを説明しない。R07原典1の融点は多成分岩石では固相線・液相線という範囲をもつ。公式R05=\allowbreak{}3/R07=\allowbreak{}2を保持して条件を補足。

\paragraph{問題の要旨} 減圧融解・含水固相線低下・部分融解を区別する。

\Needspace{4\baselineskip}\paragraph{確認した条件・記述（独自要約）}\mbox{}\par

\noindent\textbf{原典確認：} S7\_050の29原典画像で34問の全五肢・否定・数値を照合。\par

\noindent\textbf{対象：} 減圧融解・含水固相線低下・部分融解を区別する。\par

\paragraph{問われる核心} 減圧融解・含水固相線低下・部分融解を区別する。

\paragraph{必要知識} 海嶺等では上昇するマントルが減圧され、温度圧力経路が固相線を越えて融解する。低温物質の急な昇温だけを原因とした肢を選ぶ。沈み込みでは含水物質の脱水流体が固相線を下げる。純物質の一つの融点とは違い、岩石の部分溶融では固相とメルトが共存する。浮力上昇と貯留には周囲の構造や応力も関係する。

\paragraph{考え方・正答を選ぶ理由（査読待ち初稿）} 海嶺等では上昇するマントルが減圧され、温度圧力経路が固相線を越えて融解する。低温物質の急な昇温だけを原因とした肢を選ぶ。沈み込みでは含水物質の脱水流体が固相線を下げる。純物質の一つの融点とは違い、岩石の部分溶融では固相とメルトが共存する。浮力上昇と貯留には周囲の構造や応力も関係する。

\Needspace{4\baselineskip}\paragraph{選択肢の検討（査読待ち初稿）}\mbox{}\par

\noindent\textbf{1：代表的条件で適切} 昇温又は減圧で融解。温度圧力経路が固相線を越えると融解が始まる。全圧力域で同一勾配の融点を仮定しない。\par

\noindent\textbf{2：不適切} 低温マントルの急な昇温が海嶺・ホットスポット成因。代表的には高温マントルの上昇に伴う減圧が融解を促す。ホットスポットの熱異常はあるが、冷たい物質の急加熱だけという説明ではない。\par

\noindent\textbf{3：代表的条件で適切} スラブ由来水による固相線低下。沈み込む含水物質の脱水で流体が供給され、上のマントルの融解が促進される。\par

\noindent\textbf{4：代表的条件で適切} 残留固相のある部分溶融。一部の成分が融け、他の固相が残る多成分岩石の挙動。\par

\noindent\textbf{5：条件・解釈を留保} 低密度メルトの浮力と貯留。浮力は重要だが、貯留の位置は密度差だけでなく応力・構造・結晶化等に依存する。\par

\paragraph{再利用する考え方} 地質環境・生成機構・物質の名称を対応させ、平均値と全称、典型例と例外、当時の分類と現在の知見を分ける。

\paragraph{共通理論} \hyperref[theory:volcano:melting]{第\ref*{theory:volcano:melting}節}

\par 参考：\cite{vol50_melt_ref,vol50_subduction_ref}。

\paragraph{関連問題} \hyperref[q:r05-reg-applied-25]{R05-REG-APPLIED-25}

\par 難易度 4（校準初稿）、優先度未評価。

\section{火山用語と軟堆積物の火炎状構造を区別する。}

\subsection{平成29年度 問23}
\label{q:h29-reg-applied-23}

\noindent\textbf{原典：} sources/H29/17\_applied\_science.pdf、PDF物理8頁。

\textbf{公式正答：} 4（single）。

\paragraph{問題の要旨} 火山用語と軟堆積物の火炎状構造を区別する。

\Needspace{4\baselineskip}\paragraph{確認した条件・記述（独自要約）}\mbox{}\par

\noindent\textbf{原典確認：} S7\_050の29原典画像で34問の全五肢・否定・数値を照合。\par

\noindent\textbf{対象：} 火山用語と軟堆積物の火炎状構造を区別する。\par

\paragraph{問われる核心} 火山用語と軟堆積物の火炎状構造を区別する。

\paragraph{必要知識} 火炎状構造の火炎は見た目の比喩であり、成因は未固結堆積物の変形である。火山活動の直接用語として最も弱い4を選ぶ。火砕堆積物でも変形構造ができ得るので、火山地域に絶対に存在しないという意味ではない。

\paragraph{考え方・正答を選ぶ理由（査読待ち初稿）} 火炎状構造の火炎は見た目の比喩であり、成因は未固結堆積物の変形である。火山活動の直接用語として最も弱い4を選ぶ。火砕堆積物でも変形構造ができ得るので、火山地域に絶対に存在しないという意味ではない。

\Needspace{4\baselineskip}\paragraph{選択肢の検討（査読待ち初稿）}\mbox{}\par

\noindent\textbf{1：代表的条件で適切} カルデラ。マグマの大量放出等に伴う陥没などで生じる大きな凹地で、火山活動と関係する。\par

\noindent\textbf{2：代表的条件で適切} 安山岩。島弧で代表的に見られる中間質火山岩。\par

\noindent\textbf{3：代表的条件で適切} アグルチネート。高温で放出された溶岩片等が付着・溶結してできる火砕岩。\par

\noindent\textbf{4：不適切} 火炎状構造。未固結の軽い層が上の重い堆積物の間へ押し上げられた荷重変形などで生じ、火炎という形の名は噴火を意味しない。\par

\noindent\textbf{5：代表的条件で適切} マグマ水蒸気爆発。マグマと外来水の相互作用を伴う爆発的噴火。\par

\paragraph{再利用する考え方} 地質環境・生成機構・物質の名称を対応させ、平均値と全称、典型例と例外、当時の分類と現在の知見を分ける。

\paragraph{共通理論} \hyperref[theory:volcano:terms]{第\ref*{theory:volcano:terms}節}

\par 参考：\cite{vol50_terms_ref}。

\par 難易度 3（校準初稿）、優先度未評価。

\section{噴出物の全世界体積率と火山の個数を区別する。}

\subsection{平成28年度 問29}
\label{q:h28-reg-applied-29}

\noindent\textbf{原典：} sources/H28/17\_applied\_science.pdf、PDF物理12頁。

\textbf{公式正答：} 4（single）。

\paragraph{問題の要旨} 噴出物の全世界体積率と火山の個数を区別する。

\Needspace{4\baselineskip}\paragraph{確認した条件・記述（独自要約）}\mbox{}\par

\noindent\textbf{原典確認：} S7\_050の29原典画像で34問の全五肢・否定・数値を照合。\par

\noindent\textbf{対象：} 噴出物の全世界体積率と火山の個数を区別する。\par

\paragraph{問われる核心} 噴出物の全世界体積率と火山の個数を区別する。

\paragraph{必要知識} 単位はkm3／年で、火山の個数・噴火回数・陸上被害の大きさではない。海嶺拡大が作る全世界の新しい海洋地殻に対応する4が最大。USGSの約4分の3は代表的推計で、厳密な年度別観測値や未来の一定比率ではない。

\paragraph{考え方・正答を選ぶ理由（査読待ち初稿）} 単位はkm3／年で、火山の個数・噴火回数・陸上被害の大きさではない。海嶺拡大が作る全世界の新しい海洋地殻に対応する4が最大。USGSの約4分の3は代表的推計で、厳密な年度別観測値や未来の一定比率ではない。

\Needspace{4\baselineskip}\paragraph{選択肢の検討（査読待ち初稿）}\mbox{}\par

\noindent\textbf{1：不適切} 島弧・陸弧。爆発的な陸上噴火が目立つが、全世界の体積噴出率の最大域ではない。\par

\noindent\textbf{2：不適切} 大陸地溝帯。伸張に伴う火山活動はあるが、長大な海嶺系の地殻生産に比べて小さい。\par

\noindent\textbf{3：不適切} 縁海。活動的な背弧海盆等もあるが全縁海が活動中の火山域というわけではない。\par

\noindent\textbf{4：代表的条件で適切} 中央海嶺の軸。全世界の長大な拡大軸で連続的に海洋地殻を作り、USGS概説でも年噴出物の約4分の3と見積もる。\par

\noindent\textbf{5：不適切} 海山・火山島。ハワイ等は大きな火山を作るが、個別の大型火山と全世界海嶺系の総量を比較する。\par

\paragraph{再利用する考え方} 地質環境・生成機構・物質の名称を対応させ、平均値と全称、典型例と例外、当時の分類と現在の知見を分ける。

\paragraph{共通理論} \hyperref[theory:volcano:arc]{第\ref*{theory:volcano:arc}節}

\par 参考：\cite{vol50_ridge_ref}。

\par 難易度 3（校準初稿）、優先度未評価。

\section{火砕密度流の形成と地形に沿う分布}

\subsection{令和元年度 問25}
\label{q:r01-reg-applied-25}

\noindent\textbf{原典：} sources/R01/17\_applied\_science.pdf、PDF物理9頁。

\textbf{公式正答：} 5（single）。

\par\textbf{検討未確定：} 大規模火砕流/プリニー/カルデラの一般化とプレー型の用語範囲は全類型の一次資料照合が未完了である。公式/独立5と確認済みの非単純分布を保持し部分校正に残す。

\paragraph{問題の要旨} 火砕流の発生・混合物・堆積性状と、地形を反映する分布を比較する。

\Needspace{4\baselineskip}\paragraph{確認した条件・記述（独自要約）}\mbox{}\par

\noindent\textbf{判断方向：} 最も不適切・誤りを選ぶ\par

\noindent\textbf{原典照合：} 15原典画像で18問の問題幹、全5肢、否定、専門語、化学式、候補順を確認。\par

\noindent\textbf{適用条件：} 流れと構造の典型対応は一意成因を保証しない。用語、時間間隙、保存、試料処理の条件を個別に示す。\par

\paragraph{問われる核心} 火砕流の発生・混合物・堆積性状と、地形を反映する分布を比較する。

\paragraph{必要知識} 火砕密度流の粒子と気体、ドーム崩落、噴煙柱、地形と分布、降下と流下の区別。

\paragraph{考え方・正答を選ぶ理由（査読待ち初稿）} 火砕流は地面に沿う密度流で、経路と堆積域は谷や尾根、流れの濃度に依存する。USGSには重なった扇状・舌状・ローブ状の分布及びメラピの谷埋積と希薄サージ分離の報告がある。従って単純な長円形が一般的だとする5が最も不適切である。降下火山灰は風に運ばれるため別の分布モデルを使う。大規模噴出・カルデラやプレー型という呼称は、この資料確認だけで全例を普遍化せず原典の典型説明として保持する。

\Needspace{4\baselineskip}\paragraph{選択肢の検討（査読待ち初稿）}\mbox{}\par

\noindent\textbf{1：適切（条件を含めて検討）} 一般に分級が悪く塊状となる。多粒径を混合して急速に堆積する流れの代表的特徴である。希薄サージ等の全堆積物へ塊状を強制しない。\par

\noindent\textbf{2：適切（条件を含めて検討）} 大規模火砕流は高発泡のプリニー式に伴い通常カルデラを作る。高発泡マグマの大規模噴出とカルデラ形成の典型を述べるが、全活動様式の必要十分条件としては確認していない。\par

\noindent\textbf{3：適切（条件を含めて検討）} 高温火砕物・火山ガス・取り込む大気が地面に沿って流れる。高温粒子と気体の混合物が地形に沿って流下する現象に対応する。固体だけ又は溶岩だけではない。\par

\noindent\textbf{4：適切（条件を含めて検討）} ドームの爆発型をプレー型、重力崩落型をメラピ型とする。メラピのドーム崩落起源は一次論文で確認した。プレー型という原典呼称と全類型の整理は追加確認を要する。\par

\noindent\textbf{5：誤り・最も不適切} 分布は火口を長軸の一端に持つ単純な長円が多い。谷沿いの舌状・ローブ状、複数谷の埋積などを取り得る。風で伸びる降下物の模式分布と混同する一般化である。\par

\paragraph{再利用する考え方} 観察できる構造と成因・環境の解釈を区別する。代表モデルに適用条件を添え、単一の形や測定値から唯一の環境を断定しない。

\paragraph{共通理論} \hyperref[theory:volcanology:pyroflow]{第\ref*{theory:volcanology:pyroflow}節}, \hyperref[theory:volcanology:source]{第\ref*{theory:volcanology:source}節}, \hyperref[theory:volcanology:distribution]{第\ref*{theory:volcanology:distribution}節}

\par 参考：\cite{volc_usgs_merapi,volc_usgs_msh_flow}。

\par 難易度 3（校準初稿）、優先度未評価。


\chapter{気象・大気}
\section{地球上の流体を保存則から読む}
大気は回転する球面上の圧縮性流体であり、重力で鉛直に成層し、日射と地球放射の差で加熱・冷却される。
気象の名称を個別に暗記する前に、運動量、水の質量、エネルギーの保存則を分けると、既知の物理学から多くを理解できる。
一方、組成や代表高度は理論だけでは確定しない観測量である。標準モデル、全球平均、地域・季節ごとの値を同じ精度で扱わない。
本章では23問の全5肢を原典で確認した初稿を収録する。正答一致と数値・文言の適用範囲の確認は別であり、部分校正の4問には疑義を残した。

\section{乾燥大気の組成：量と作用は別である}
\label{theory:atmosphere:composition}
理想気体近似では体積比はモル比に等しい。水蒸気を除いた乾燥大気では窒素約78\%、酸素約21\%、アルゴン約0.93\%の順であり、
二酸化炭素はアルゴンより少ない\cite{atm_noaa_composition}。水蒸気を含む全体を分母にすれば、乾燥成分の見かけの割合も変わる。
主要成分が一様という記述は、乱流混合で均される下層の長寿命気体を指す。水蒸気やオゾンの鉛直分布まで一様とはしない。
上方では分子拡散が重要となり、分子質量による分離と紫外線による解離が起こる。優勢種の切替高度は太陽活動等にも依存するため、
特定高度で必ず単一成分だけになると教えない。
酸素の現在の存在量は地史を通じた生物・酸化還元過程の結果である。
二酸化炭素の観測曲線は長期傾向に季節変動が重なり、メタンは湿地・家畜・水田・化石燃料等の多様な源を持つ\cite{atm_noaa_ghg}。
本章は現在の精密な濃度値を過年度の問題へ遡及して代入しない。

\section{温度で区分する層と、密度で測る薄さ}
\label{theory:atmosphere:layers}
下から対流圏、成層圏、中間圏、熱圏という区分は主として気温の鉛直勾配を基準にする。
対流圏は上方へ低温、成層圏は上方へ高温、中間圏では再び低温、熱圏では高エネルギー放射の吸収等により高温となる。
成層圏界面は約50km、中間圏界面は約80〜85kmの代表高度である\cite{atm_noaa_layers}。
対流圏界面は低緯度で高く、極域で低い。大気質量が下層へ集中するため、対流圏は全球大気質量の約3/4という概数を持つが、
界面は緯度・季節で動き、厳密な固定比ではない。
高い高度で温度が高いことは、単位体積の熱容量が大きいことと同じではない。少ない粒子の平均運動エネルギーと、全粒子数を区別する。
気圧、質量密度、混合比も別の量である。

\section{静水圧と状態方程式から密度を見積もる}
\label{theory:atmosphere:hydrostatic}
大きい時間・空間尺度の鉛直平均では、重力と圧力勾配が釣り合う静水圧近似が有効である。高さを\(z\)、重力加速度を\(g\)、
乾燥空気の比気体定数を\(R_d\)とすると
\[
 \frac{dp}{dz}=-\rho g,\qquad p=\rho R_dT.
\]
等温なら\(p(z)=p(0)\exp[-z/H]\)、スケールハイトは\(H=R_dT/g\)となる。
温度が\(T=T_0-Lz\)で減る層では、一定\(g,R_d,L\)のモデルで積分して
\[
 \frac{p(z)}{p_0}=\left(1-\frac{Lz}{T_0}\right)^{g/(R_dL)},\qquad
 \frac{\rho(z)}{\rho_0}=\left(1-\frac{Lz}{T_0}\right)^{g/(R_dL)-1}
\]
を得る。圧力と密度は温度因子の分だけ比が違う。標準の\(T_0=288.15\,{\rm K}\)、\(L=0.0065\,{\rm K/m}\)、
\(R_d\simeq287.05\,{\rm J/(kg\,K)}\)、\(g\simeq9.80665\,{\rm m/s^2}\)を置くと10kmで圧力比約0.26、密度比約0.34である。
標準大気は現実の全地域の精密保証ではなく、比較用モデルである\cite{atm_us1976}。
令和元年度問19と令和7年度問21の原典は「密度が1/4程度」と記す。圧力に読み替える校正は行わず、
上の計算との差を概数精度の査読事項として残す。

\section{環境減率と断熱減率を分ける}
\label{theory:atmosphere:lapse}
実際の周囲の温度勾配を環境気温減率という。標準モデルの対流圏の代表値は6.5K/kmであり、
観測では雲・放射・移流等によって変化し、逆転層もある。
これに対し、断熱減率は移動する気塊の温度変化である。乾燥理想気体の断熱第一法則を単位質量あたりで書けば
\[
 c_p\,dT-\frac{dp}{\rho}=0,\qquad
 \frac{dT}{dz}=-\frac{g}{c_p}
\]
であり、空気の\(c_p\simeq1004\,{\rm J/(kg\,K)}\)なら約9.8K/kmとなる。
凝結を伴う飽和気塊では潜熱放出が冷却を補い、湿潤断熱減率はより小さく、温度・圧力・水蒸気量に依存する。
環境がどの程度冷えているかと、持ち上げた気塊がどの程度冷えるかを比較して浮力の符号を判定する。
「対流圏は3K/km」という一般化は、湿潤過程の一例を平均環境へ置き換える誤りである。

\section{オゾンの紫外吸収と成層圏加熱}
\label{theory:atmosphere:ozone}
大気中オゾンのおよそ9割は成層圏にある。紫外線を吸収して地表への到達を減らし、同時に大気へエネルギーを与える。
この作用が成層圏の上方への温度増加と関係する\cite{atm_jma_ozone}。
赤外線を吸収する温室効果と、オゾンの紫外吸収による鉛直加熱を同じ説明へまとめない。
オゾンの生成・消滅は動的な化学反応の収支であり、CFC由来の塩素やN2O由来の窒素酸化物は触媒反応へ関与する。
微量成分が減るだけでも大きい放射・化学効果を持ち得る。地表付近にもオゾンはあり、成層圏の有用な吸収作用を
地表オゾンの高濃度が無害であるという意味へ拡張しない。

\section{太陽定数と受光断面：円と球を混同しない}
\label{theory:atmosphere:geometry}
太陽定数\(S_0\)は大気上端で太陽光に垂直な単位面積が受ける放射束である。
大気下端では吸収・散乱と入射角が影響する。地球半径\(R\)、光線が平行という近似で受光総量は
\[
 P=\int_{\rm illuminated} S_0\cos\theta\,dA=S_0\pi R^2.
\]
球の前面を光線方向へ投影すると円になるためである。全球球面積に均すと
\[
 F_{\rm in}=\frac{P}{4\pi R^2}=\frac{S_0}{4}.
\]
総量Wを問うのか、平均放射束W/m²を問うのかで計算が違う。6370kmはmへ換算してから二乗する。
問題が与える1368W/m²を用いれば総量約\(1.74\times10^{17}\,{\rm W}\)、全球平均342W/m²となる。
これらは問題の指定値であり、現在の精密観測で常に1368であるとの主張ではない。

\section{短波・長波と全球放射収支}
\label{theory:atmosphere:radiation}
太陽の高温の放射は主に可視・近赤外を含む短波帯であり、より低温の地球からの放射は主に熱赤外の長波帯である。
太陽放射の半分をガンマ線・X線とする説明はスペクトルの主帯域を取り違える。
全球アルベド\(\alpha\)を宇宙へ反射する割合とすれば
\[
 F_{\rm refl}=\alpha S_0/4,\qquad F_{\rm abs}=(1-\alpha)S_0/4.
\]
定常の簡略モデルでは長波流出が吸収量に等しくなる。代表的な収支図は入射342、反射107、吸収235W/m²である
\cite{atm_nasa_budget}。現在の観測や平均期間で値は変わり、反射率は概ね約3割である\cite{atm_nasa_albedo}。
太陽距離の季節変化は日射へ影響するが、四季の主因は地軸傾斜に伴う太陽高度・日長の変化である。
全球平均の式だけでは緯度別の日射は決まらず、南北半球の逆の季節も考える。

\section{温室効果と水蒸気のフィードバック}
\label{theory:atmosphere:greenhouse}
温室気体は地表等が放出する赤外を吸収し、温度に応じて放射する。地表のエネルギー収支には下向き赤外も入り、
単純な大気なしの放射平衡とは違う温度になる。
現在の自然温室効果に水蒸気が大きく寄与することと、近代の人為的な追加放射強制力でCO2が重要であることは両立する
\cite{atm_jma_ghg2012}。大気の水蒸気量は温度と凝結によって速く応答し、気候変化を増幅する側面を持つ。
分子の大気量、1分子あたりの吸収、現在の自然寄与、人為増加の寄与を同じランキングにしない。
窒素と酸素は多量でも、CO2や水蒸気と同じ赤外吸収特性を持つという結論にはならない。

\section{潜熱と緯度間の熱輸送}
\label{theory:atmosphere:transport}
緯度による日射差は地球放射との差を生む。低緯度の正味の加熱と高緯度の冷却を、大気・海洋が熱輸送で補う
\cite{atm_nasa_transport}。この背景を欠くと、全球平衡だからどの緯度でも同じ収支という誤解が生じる。
蒸発は液体から気体への変化に必要な潜熱を地表から受け取り、凝結は大気で潜熱を放出する。
凝結地点が蒸発地点と違えば水蒸気の移流がエネルギー輸送となる。潜熱フラックスは水の蒸発質量率\(\dot m/A\)と
単位質量の蒸発潜熱\(L_v\)の積である。
代表収支の潜熱78W/m²を入射342W/m²で割ると約23\%であり、約1割ではない。
吸収日射、地表正味放射、全球入射など分母を変えれば割合も変わる。気象の割合問題では分母を先に書き出す。

\section{回転系の運動方程式と地衡風}
\label{theory:atmosphere:wind}
水平速度を\(\mathbf u\)、鉛直上向き単位ベクトルを\(\mathbf k\)、コリオリ係数を\(f=2\Omega\sin\varphi\)とすれば、
局所水平面での基本式は
\[
 \frac{D\mathbf u}{Dt}+f\mathbf k\times\mathbf u
       =-\frac{1}{\rho}\nabla_h p+\hbox{摩擦等}.
\]
転向加速度は左辺項を移した\(-f\mathbf k\times\mathbf u\)であり、北半球で運動方向の右へ向く。
加速度と摩擦が小さい近似では
\[
 f\mathbf k\times\mathbf u_g=-\frac{1}{\rho}\nabla_h p,\qquad
 \mathbf u_g=\frac{1}{\rho f}\mathbf k\times\nabla_h p
\]
となる。地衡風は等圧線に沿い、北半球では進行方向の左に低圧側がある\cite{atm_noaa_winds,atm_jma_dynamic}。
圧力傾度力が等圧線に直交することから風まで直交すると結論しない。
近似の目安はロスビー数\(U/(|f|L)\ll1\)である。赤道近く、強い曲率や地表摩擦がある条件では同じ式を無条件に使わない。
地表から約1kmの境界層は摩擦・乱流・熱交換が支配的な代表領域で、厚さは気象条件に依存する。

\section{局地風、大循環、ロスビー波}
\label{theory:atmosphere:local}
海陸風は海と陸の熱応答差に由来する。日中は陸が速く暖まり、地表付近では海から陸へ海風が吹く。
夜間は陸が速く冷え、基本の循環は陸から海への陸風となる。強い一般風等があると単純な日周期が明瞭でない。
全球の帯状平均の模式は低緯度のハドレー、中緯度のフェレル、高緯度の極循環である\cite{atm_noaa_winds}。
瞬時の大気が3つの独立した定常セルだけで構成されるという保証ではなく、中緯度の擾乱は熱・運動量輸送に重要である。
コリオリ係数の緯度変化は渦度保存を介してロスビー波の復元性となる。偏西風の大規模な蛇行を一般の重力波と区別する。
季節的な海陸加熱差の影響を受けるモンスーンを、赤道帯全体のハドレー循環や日周期の海陸風と同じものにしない。

\section{ENSO：大気と海洋が共有する数年規模の変動}
\label{theory:atmosphere:enso}
熱帯太平洋の海面水温と大気循環は結合して変動する。エルニーニョ・ラニーニャは主に海洋側の状態を、
南方振動は東西の地上気圧のシーソー状変化を捉えたものと理解できる\cite{atm_jma_enso}。
両者を総合するのがENSOである。数年ごとに生じる変動を、数十年規模の変動と同一視しない。
数年は厳密に同じ周期の正弦波という意味ではない。指数の符号、基準平均、海域を区別して観測する。
本章では機構の位置付けを扱い、過去問の1肢から特定の将来天候や発生年を断定しない。

\section{水収支、貯留と平均滞留時間}
\label{theory:atmosphere:water}
水がどこにあるかを表す貯留量と、どれだけ移るかを表す流量は別である。大気貯留\(W\)の収支は
\[
 \frac{dW}{dt}=E-P-\nabla_h\cdot\mathbf Q
\]
の形を持つ。ここで面積あたりの蒸発量率\(E\)、降水量率\(P\)、鉛直積分した水平水蒸気輸送\(\mathbf Q\)を同じ質量単位へ合わせる。
全球長期平均では水平発散の積分が消え、蒸発散と降水が概ね釣り合う。海洋は蒸発超過、陸域は降水超過で、
海から陸への大気輸送と陸から海への流出が結ぶ。
定常に近い貯留を出る流量率\(F\)で割った\(W/F\)が平均交換時間の目安である。
大気の代表貯留12900km³と日蒸発散1170km³というUSGS推定なら約11日となり\cite{atm_usgs_water}、
他推定の8〜10日と合わせて一桁〜十日程度という比較を行える。
氷床は大量を保持するが交換に長い時間を要する\cite{atm_usgs_cycle}。大きい貯留が大きい年間輸送を意味しない。
全球の更新可能性は、地域の渇水や地下水の過剰取水が起こらないという保証でもない。

\section{可降水量と雨の形成}
\label{theory:atmosphere:rain}
可降水量はある時点の大気柱の水蒸気をすべて水へ換算した、単位底面積あたりの水量である\cite{atm_noaa_pwat}。
日降水量は24時間の流出の積分であり、風で次々と水蒸気が補給されれば初期の柱貯留を超え得る。
在庫と一日の販売量の違いに相当し、保存則に反しない。
雲滴の衝突併合による雨と、氷晶等が成長して落下・融解する雨を区別する。
後者の「冷たい雨」は氷相を経る過程の名称で、地上の雨の温度を測った分類ではない。
日本の降水の大部分には冷たい雲の過程が関わる\cite{atm_mri_cloud}。
0℃未満でも液体の過冷却水滴が存在するため、雲中の全水分が温度だけで即座に凍ると仮定しない。

\section{水を通す係数と蓄える係数}
\label{theory:atmosphere:aquifer}
帯水層の透水係数\(K\)は水理勾配に応じた流速の尺度で、単位は長さ/時間である。厚さ\(b\)まで積分して
水平方向の通水能力を表す透水量係数は\(T=Kb\)、単位は面積/時間となる\cite{atm_usgs_transmissivity}。
比貯留率\(S_s\)は水頭変化に対する単位体積の放出・吸収を表し、被圧層の貯留係数\(S=S_sb\)は無次元である。
透水と貯留を厚さという共通の乗数だけで取り違えず、次元と物理的役割を確認する。
不飽和帯に低透水層があれば、その上に局所的な地下水面を持つ宙水が生じ得る\cite{atm_usgs_perched}。
水循環の問題は大気から地下へつながるため、この用語は地下水の章にも接続する。

\section{降水pH：対数、炭酸平衡、中和}
\label{theory:atmosphere:acid}
pHの厳密な定義は水素イオン活量の負の常用対数であり、希薄な降水の比較では濃度へ近似して
\[
 {\rm pH}\simeq-\log_{10}[{\rm H}^+],\qquad
 [{\rm H}^+]\simeq10^{-{\rm pH}}\ {\rm mol/L}
\]
と計算する。pH差1はH+濃度の10倍差である。pH3.3は約501μmol/Lで、30μmol/Lではない。
25℃付近のCO2と平衡する希薄水は約pH5.6を目安にするが、これは全降水の不変値ではない。
酸性の硫酸・硝酸等と、土壌粒子やアンモニア等の中和成分を加えると値が変わる\cite{atm_jma_acid}。
海洋上の雨は海水そのものではなく、海水のpHを雨へ直接代入できない。
平成27年度の日本の傾向は、2008〜2012調査の地点年平均4.48〜5.37、全地点5年4.72という当時資料で読む
\cite{atm_jma_acid2014}。その観測値を現在の各地点へ保証しない。
酸性沈着の影響は酸性度だけでなく降水量にも依存し、水・土壌・生態系・建造物という受け手の条件もある。

\section{過去問の判断手順}
数値には対象、分母、単位、平均期間を添える。正誤判定では、温度区分と成分、密度と圧力、全球平均と地域保証、
環境減率と気塊の断熱減率を分ける。同じ理論の年度違いでも正答番号は変わり、正しい肢と誤った肢の入替えがある。
以下の解説はAstra査読前の初稿で、疑義を公式正答に合わせて消していない。
\section{大気循環と地衡風・局地風}

\subsection{平成27年度 問24}
\label{q:h27-reg-applied-24}

\noindent\textbf{原典：} sources/H27/17\_applied\_science.pdf、PDF物理9頁。

\textbf{公式正答：} 4（single）。

\paragraph{問題の要旨} 大循環、鉛直層序、大気組成、海陸風、地衡風の記述から適切なものを選ぶ。

\Needspace{4\baselineskip}\paragraph{確認した条件・記述（独自要約）}\mbox{}\par

\noindent\textbf{方向：} 最も適切。\par

\paragraph{問われる核心} 鉛直構造と水平循環を支配する別の条件

\paragraph{必要知識} 鉛直構造と水平循環を支配する別の条件。大循環は緯度、層序は温度勾配、上層組成は混合と拡散、風は力学の条件で判定する。日中の海→陸を述べる4が適切。海陸風を大循環や地衡風の定常近似と同一視しない。

\paragraph{考え方・正答を選ぶ理由（査読待ち初稿）} 大循環は緯度、層序は温度勾配、上層組成は混合と拡散、風は力学の条件で判定する。日中の海→陸を述べる4が適切。海陸風を大循環や地衡風の定常近似と同一視しない。

\Needspace{4\baselineskip}\paragraph{選択肢の検討（査読待ち初稿）}\mbox{}\par

\noindent\textbf{1：不適切} 赤道にモンスーン、極域にハドレー循環。。ハドレーは低緯度の循環であり、極域への割当が違う。モンスーンを低緯度帯全体の基本循環へ一括置換しない。\par

\noindent\textbf{2：不適切} 中間圏オゾンUV加熱で80km温度極大。。オゾン加熱の代表は成層圏で、約50kmに温度極大。中間圏では上方へ冷却し約80〜85kmで極小となる。\par

\noindent\textbf{3：不適切} 成層圏の主成分はO原子とHe。。成層圏でも主成分はN2とO2。O原子・He優勢はさらに高い高度の拡散分離・解離と関係する。\par

\noindent\textbf{4：最適} 日中は地表付近で海から陸へ海風。。日射で陸が海より速く暖まり地表で海→陸の循環を生じる。強い一般風がある場合は局地循環が明瞭でない。\par

\noindent\textbf{5：不適切} 北半球転向力は左向き、地衡風は等圧線直交。。北半球の転向力は進行方向の右。地衡風は気圧傾度力と釣り合い等圧線にほぼ沿う。\par

\paragraph{再利用する考え方} 全球平均・地域差、貯留量・流量、自然温室効果・追加強制力、代表値・普遍則を区別する。鉛直構造と水平循環を支配する別の条件の対象と成立条件を確認する。

\paragraph{共通理論} \hyperref[theory:atmosphere:layers]{第\ref*{theory:atmosphere:layers}節}, \hyperref[theory:atmosphere:wind]{第\ref*{theory:atmosphere:wind}節}, \hyperref[theory:atmosphere:local]{第\ref*{theory:atmosphere:local}節}

\par 参考：\cite{atm_noaa_layers,atm_noaa_winds,atm_jma_dynamic}。

\paragraph{関連問題} \hyperref[q:r02-reg-applied-21]{R02-REG-APPLIED-21}

\par 難易度 2（校準初稿）、優先度未評価。

\subsection{令和2年度 問21}
\label{q:r02-reg-applied-21}

\noindent\textbf{原典：} sources/R02/17\_applied\_science.pdf、PDF物理8頁。

\textbf{公式正答：} 4（single）。

\par\textbf{検討未確定：} 肢3の170km超O・1000km付近He主要の数値高度は一次数表未確認。原典の数字は保持し部分校正。

\paragraph{問題の要旨} 大循環、成層圏加熱、上層組成、海陸風と地衡風の記述を比較する。

\Needspace{4\baselineskip}\paragraph{確認した条件・記述（独自要約）}\mbox{}\par

\noindent\textbf{方向：} 最も不適切。\par

\paragraph{問われる核心} 鉛直組成と循環モデルの適用条件

\paragraph{必要知識} 鉛直組成と循環モデルの適用条件。日中の陸加熱と夜間の陸冷却を分けると夜間海→陸の4が不適切となる。大循環は緯度帯、地衡風は運動方程式の近似で判定する。3の原子O・Heの数値高度は定性的な分離の理解とは別に一次数表を要するため、全肢確認へ昇格させない。

\paragraph{考え方・正答を選ぶ理由（査読待ち初稿）} 日中の陸加熱と夜間の陸冷却を分けると夜間海→陸の4が不適切となる。大循環は緯度帯、地衡風は運動方程式の近似で判定する。3の原子O・Heの数値高度は定性的な分離の理解とは別に一次数表を要するため、全肢確認へ昇格させない。

\Needspace{4\baselineskip}\paragraph{選択肢の検討（査読待ち初稿）}\mbox{}\par

\noindent\textbf{1：適切（平均模式）} 低緯度ハドレー、中緯度フェレル、高緯度極循環。。緯度帯に対応する帯状平均循環の基本模式である。瞬時の風を3つの閉じた固定セルだけで表さない。\par

\noindent\textbf{2：適切} 成層圏オゾンUV吸収で約50km温度極大。。層序とオゾンの放射作用が整合する。\par

\noindent\textbf{3：数値高度は要確認} 80kmまでほぼ一定、170km超でO原子、1000km付近でHe主要。。均質圏から上方の解離と拡散分離の定性的説明は合う。優勢成分が切り替わる高度は太陽活動等に依存し、この2数値の一次数表検証は未完了。\par

\noindent\textbf{4：不適切} 夜間は地表で海から陸へ海風。。夜間は陸が速く冷え、基本の地表循環は陸→海の陸風。\par

\noindent\textbf{5：適切（近似条件）} 傾度力と転向力の釣合いで等圧線沿いの地衡風。。摩擦・加速度が小さく赤道から離れた条件での理論風である。\par

\paragraph{再利用する考え方} 全球平均・地域差、貯留量・流量、自然温室効果・追加強制力、代表値・普遍則を区別する。鉛直組成と循環モデルの適用条件の対象と成立条件を確認する。

\paragraph{共通理論} \hyperref[theory:atmosphere:composition]{第\ref*{theory:atmosphere:composition}節}, \hyperref[theory:atmosphere:wind]{第\ref*{theory:atmosphere:wind}節}, \hyperref[theory:atmosphere:local]{第\ref*{theory:atmosphere:local}節}

\par 参考：\cite{atm_noaa_winds,atm_noaa_layers,atm_us1976}。

\paragraph{関連問題} \hyperref[q:h27-reg-applied-24]{H27-REG-APPLIED-24}

\par 難易度 2（校準初稿）、優先度未評価。

\section{乾燥大気の組成と微量気体の循環}

\subsection{平成24年度 問34}
\label{q:h24-reg-applied-34}

\noindent\textbf{原典：} sources/H24/17\_applied\_science.pdf、PDF物理10頁。

\textbf{公式正答：} 2（single）。

\paragraph{問題の要旨} 大気微量成分の発生源、温室効果、成層圏オゾンへの作用を比較する。

\Needspace{4\baselineskip}\paragraph{確認した条件・記述（独自要約）}\mbox{}\par

\noindent\textbf{方向：} 最も不適切。\par

\paragraph{問われる核心} 微量気体の源・寿命・放射と化学反応

\paragraph{必要知識} 微量気体の源・寿命・放射と化学反応。大気成分の量と放射・化学上の作用は別である。微量でも温室効果やオゾン反応へ寄与する。メタンの源は化石燃料採掘のほか湿地、水田、家畜等に広がるため2が不適切。N2OをN2やNO2と取り違えず、オゾンのUV作用と温室気体の赤外作用を区別する。

\paragraph{考え方・正答を選ぶ理由（査読待ち初稿）} 大気成分の量と放射・化学上の作用は別である。微量でも温室効果やオゾン反応へ寄与する。メタンの源は化石燃料採掘のほか湿地、水田、家畜等に広がるため2が不適切。N2OをN2やNO2と取り違えず、オゾンのUV作用と温室気体の赤外作用を区別する。

\Needspace{4\baselineskip}\paragraph{選択肢の検討（査読待ち初稿）}\mbox{}\par

\noindent\textbf{1：適切（条件付き概数）} N2Oは温室効果とオゾン破壊に関与。。N2Oから生成されるNO等がオゾンの触媒反応へ関わり、赤外吸収も持つ。\par

\noindent\textbf{2：不適切} メタンは湿地・水田・反すう動物の寄与が小さく採掘由来がほぼ全て。。嫌気的な微生物過程や家畜・水田なども重要な源で、化石燃料採掘だけに限定できない。\par

\noindent\textbf{3：適切（条件付き概数）} CO2は季節変動を伴い長期増加。。原典当時の観測と整合。季節周期の上へ長期増加が重なる。\par

\noindent\textbf{4：適切（条件付き概数）} CFC11/12は安定で輸送後、成層圏でUV分解しClを放出。。対流圏で壊れにくい性質が成層圏への輸送を可能にし、UVで遊離した塩素が反応へ関与する。\par

\noindent\textbf{5：適切（条件付き概数）} 成層圏オゾンは有害UVを吸収し大気を加熱。。オゾンの吸収が成層圏の正の温度勾配を作る要因である。\par

\paragraph{再利用する考え方} 全球平均・地域差、貯留量・流量、自然温室効果・追加強制力、代表値・普遍則を区別する。微量気体の源・寿命・放射と化学反応の対象と成立条件を確認する。

\paragraph{共通理論} \hyperref[theory:atmosphere:composition]{第\ref*{theory:atmosphere:composition}節}, \hyperref[theory:atmosphere:ozone]{第\ref*{theory:atmosphere:ozone}節}

\par 参考：\cite{atm_noaa_ghg,atm_jma_ozone}。

\par 難易度 2（校準初稿）、優先度未評価。

\subsection{平成30年度 問20}
\label{q:h30-reg-applied-20}

\noindent\textbf{原典：} sources/H30/17\_applied\_science.pdf、PDF物理7頁。

\textbf{公式正答：} 4（single）。

\paragraph{問題の要旨} 乾燥大気の主要成分の量順と、微量温室気体の発生・季節変動を判定する。

\Needspace{4\baselineskip}\paragraph{確認した条件・記述（独自要約）}\mbox{}\par

\noindent\textbf{方向：} 最も不適切。\par

\paragraph{問われる核心} 乾燥大気の混合比と変動性の違い

\paragraph{必要知識} 乾燥大気の混合比と変動性の違い。乾燥大気のN2約78\%、O2約21\%、Ar約0.93\%の順を先に確かめる。CO2はArより少ないので順位を入れ替えた肢が不適切。微量でも放射・生物循環の作用は大きく、体積比の多寡から温室作用の大きさを直接決めない。

\paragraph{考え方・正答を選ぶ理由（査読待ち初稿）} 乾燥大気のN2約78\%、O2約21\%、Ar約0.93\%の順を先に確かめる。CO2はArより少ないので順位を入れ替えた肢が不適切。微量でも放射・生物循環の作用は大きく、体積比の多寡から温室作用の大きさを直接決めない。

\Needspace{4\baselineskip}\paragraph{選択肢の検討（査読待ち初稿）}\mbox{}\par

\noindent\textbf{1：適切（条件付き概数）} N2は安定で全球の混合比がほぼ一様。。よく混合された下層大気の乾燥混合比の代表。局所的な反応が全くないという意味ではない。\par

\noindent\textbf{2：適切（条件付き概数）} O2は原始大気に乏しく地史で増加。。酸素発生型光合成等と地球表層の酸化還元過程の長期変化に関係する。現在の組成を原始大気へ遡及しない。\par

\noindent\textbf{3：適切（条件付き概数）} CO2は光合成・呼吸で季節変動。。特に北半球の植生季節活動が観測曲線へ現れる。長期人為増加とは別成分。\par

\noindent\textbf{4：不適切} ArはN2・O2・CO2の次に多い。。乾燥大気の量順はN2、O2、Ar、CO2。Ar約0.93\%はCO2の約0.04\%より大きい。\par

\noindent\textbf{5：適切（条件付き概数）} CH4は嫌気的条件で生成する温室気体で近年増加。。生物起源の生成は酸素の乏しい環境で起こる。源は多様で、ここでいう増加は問題出題の近年に関する。\par

\paragraph{再利用する考え方} 全球平均・地域差、貯留量・流量、自然温室効果・追加強制力、代表値・普遍則を区別する。乾燥大気の混合比と変動性の違いの対象と成立条件を確認する。

\paragraph{共通理論} \hyperref[theory:atmosphere:composition]{第\ref*{theory:atmosphere:composition}節}

\par 参考：\cite{atm_noaa_composition,atm_noaa_ghg}。

\paragraph{関連問題} \hyperref[q:r05-reg-applied-18]{R05-REG-APPLIED-18}

\par 難易度 2（校準初稿）、優先度未評価。

\subsection{令和5年度 問18}
\label{q:r05-reg-applied-18}

\noindent\textbf{原典：} sources/R05/17\_applied\_science.pdf、PDF物理6頁。

\textbf{公式正答：} 1（single）。

\paragraph{問題の要旨} 乾燥大気の主要成分の量順と、微量温室気体の発生・季節変動を判定する。

\Needspace{4\baselineskip}\paragraph{確認した条件・記述（独自要約）}\mbox{}\par

\noindent\textbf{方向：} 最も不適切。\par

\paragraph{問われる核心} 乾燥大気の混合比と変動性の違い

\paragraph{必要知識} 乾燥大気の混合比と変動性の違い。乾燥大気のN2約78\%、O2約21\%、Ar約0.93\%の順を先に確かめる。CO2はArより少ないので順位を入れ替えた肢が不適切。微量でも放射・生物循環の作用は大きく、体積比の多寡から温室作用の大きさを直接決めない。

\paragraph{考え方・正答を選ぶ理由（査読待ち初稿）} 乾燥大気のN2約78\%、O2約21\%、Ar約0.93\%の順を先に確かめる。CO2はArより少ないので順位を入れ替えた肢が不適切。微量でも放射・生物循環の作用は大きく、体積比の多寡から温室作用の大きさを直接決めない。

\Needspace{4\baselineskip}\paragraph{選択肢の検討（査読待ち初稿）}\mbox{}\par

\noindent\textbf{1：不適切} ArはN2・O2・CO2の次に多い。。乾燥大気の量順はN2、O2、Ar、CO2。Ar約0.93\%はCO2の約0.04\%より大きい。\par

\noindent\textbf{2：適切（条件付き概数）} CH4は嫌気的条件で生成する温室気体で近年増加。。生物起源の生成は酸素の乏しい環境で起こる。源は多様で、ここでいう増加は問題出題の近年に関する。\par

\noindent\textbf{3：適切（条件付き概数）} N2は安定で全球の混合比がほぼ一様。。よく混合された下層大気の乾燥混合比の代表。局所的な反応が全くないという意味ではない。\par

\noindent\textbf{4：適切（条件付き概数）} O2は原始大気に乏しく地史で増加。。酸素発生型光合成等と地球表層の酸化還元過程の長期変化に関係する。現在の組成を原始大気へ遡及しない。\par

\noindent\textbf{5：適切（条件付き概数）} CO2は光合成・呼吸で季節変動。。特に北半球の植生季節活動が観測曲線へ現れる。長期人為増加とは別成分。\par

\paragraph{再利用する考え方} 全球平均・地域差、貯留量・流量、自然温室効果・追加強制力、代表値・普遍則を区別する。乾燥大気の混合比と変動性の違いの対象と成立条件を確認する。

\paragraph{共通理論} \hyperref[theory:atmosphere:composition]{第\ref*{theory:atmosphere:composition}節}

\par 参考：\cite{atm_noaa_composition,atm_noaa_ghg}。

\paragraph{関連問題} \hyperref[q:h30-reg-applied-20]{H30-REG-APPLIED-20}

\par 難易度 2（校準初稿）、優先度未評価。

\section{大気の鉛直温度・密度構造}

\subsection{令和元年度 問19}
\label{q:r01-reg-applied-19}

\noindent\textbf{原典：} sources/R01/17\_applied\_science.pdf、PDF物理6頁。

\textbf{公式正答：} 4（single）。

\par\textbf{検討未確定：} 肢3の10km密度1/4と標準モデル約0.34に差。原典密度を圧力へ改変しない。概数の許容精度をAstra査読待ち。

\paragraph{問題の要旨} 大気の層序、組成、密度低下、成層圏加熱と均質圏を比較する。

\Needspace{4\baselineskip}\paragraph{確認した条件・記述（独自要約）}\mbox{}\par

\noindent\textbf{方向：} 最も不適切。\par

\paragraph{問われる核心} 鉛直温度構造と質量密度・混合比を区別する

\paragraph{必要知識} 鉛直温度構造と質量密度・混合比を区別する。50km付近の温度極大をCO2の赤外吸収に帰す4が不適切である。オゾンのUV吸収で成層圏を加熱する。ただし3の密度1/4は標準モデル約0.34との精度差を残し、圧力への無断校正をしない。公式4と独立4の一致は全ての概数の厳密保証ではない。

\paragraph{考え方・正答を選ぶ理由（査読待ち初稿）} 50km付近の温度極大をCO2の赤外吸収に帰す4が不適切である。オゾンのUV吸収で成層圏を加熱する。ただし3の密度1/4は標準モデル約0.34との精度差を残し、圧力への無断校正をしない。公式4と独立4の一致は全ての概数の厳密保証ではない。

\Needspace{4\baselineskip}\paragraph{選択肢の検討（査読待ち初稿）}\mbox{}\par

\noindent\textbf{1：適切} 下から対流圏・成層圏・中間圏・熱圏。。主として気温の鉛直勾配を基準とする区分である。\par

\noindent\textbf{2：適切（概数）} 空気は体積でN2約8割、O2約2割。。端数の微量成分や水蒸気を除いた粗い概数として適切。\par

\noindent\textbf{3：概数精度に疑義} 10kmで質量密度は地表の約1/4。。原典は密度と記載。標準モデルで約0.414kg/m\textsuperscript{3}対地表1.225、比約0.338となり1/4とは差がある。圧力比約0.26との混同があり得るが、原典を圧力へ改変しない。\par

\noindent\textbf{4：不適切} 50km付近の温度極大はCO2赤外吸収による。。代表的な成層圏界面の加熱はオゾンの紫外吸収。温室気体一般の赤外作用へ置換できない。\par

\noindent\textbf{5：適切（主要乾燥成分）} 乾燥主要成分比は80km程度までほぼ一定。。乱流混合が分子拡散より強い均質圏の主要長寿命気体を指す。水蒸気・オゾン等まで一様とはしない。\par

\paragraph{再利用する考え方} 全球平均・地域差、貯留量・流量、自然温室効果・追加強制力、代表値・普遍則を区別する。鉛直温度構造と質量密度・混合比を区別するの対象と成立条件を確認する。

\paragraph{共通理論} \hyperref[theory:atmosphere:layers]{第\ref*{theory:atmosphere:layers}節}, \hyperref[theory:atmosphere:hydrostatic]{第\ref*{theory:atmosphere:hydrostatic}節}, \hyperref[theory:atmosphere:ozone]{第\ref*{theory:atmosphere:ozone}節}

\par 参考：\cite{atm_noaa_layers,atm_us1976,atm_jma_ozone}。

\paragraph{関連問題} \hyperref[q:r07-reg-applied-21]{R07-REG-APPLIED-21}

\par 難易度 2（校準初稿）、優先度未評価。

\subsection{令和元年度再試験 問21}
\label{q:r01-ret-applied-21}

\noindent\textbf{原典：} sources/R01\_retest/17\_applied\_science.pdf、PDF物理7頁。

\textbf{公式正答：} 2（single）。

\paragraph{問題の要旨} 対流圏の質量・気温減率・界面高度と成層圏温度、オゾンの分布を判定する。

\Needspace{4\baselineskip}\paragraph{確認した条件・記述（独自要約）}\mbox{}\par

\noindent\textbf{方向：} 最も不適切。\par

\paragraph{問われる核心} 大気の鉛直構造と環境・断熱気温減率

\paragraph{必要知識} 大気の鉛直構造と環境・断熱気温減率。環境減率の標準値は約6.5K/kmであり、3K/kmという肢が不適切。温度の観測平均と、気塊の断熱移動を表す9.8K/km等を混ぜない。界面は赤道側で高く、オゾンが多い成層圏ではUV加熱による温度上昇が生じる。

\paragraph{考え方・正答を選ぶ理由（査読待ち初稿）} 環境減率の標準値は約6.5K/kmであり、3K/kmという肢が不適切。温度の観測平均と、気塊の断熱移動を表す9.8K/km等を混ぜない。界面は赤道側で高く、オゾンが多い成層圏ではUV加熱による温度上昇が生じる。

\Needspace{4\baselineskip}\paragraph{選択肢の検討（査読待ち初稿）}\mbox{}\par

\noindent\textbf{1：適切（条件付き概数）} 対流圏質量は気圏の約3/4。。大気質量が下層へ集中する代表的な全球概数。界面の緯度・季節差で厳密値は変わる。\par

\noindent\textbf{2：不適切} 対流圏の気温減率は3℃/km。。標準的な対流圏環境減率は約6.5K/km。乾燥断熱9.8K/kmや湿潤断熱・実際の逆転層と区別する。\par

\noindent\textbf{3：適切（条件付き概数）} 対流圏界面は低緯度で高く高緯度で低い。。低緯度の高い温度と深い対流により界面は高くなる。\par

\noindent\textbf{4：適切（条件付き概数）} 成層圏では高度とともに気温上昇。。オゾン吸収による加熱が上方の温度増加を作る。成層圏の下部がほぼ等温の場合を無視した厳密な全高度単調則ではない。\par

\noindent\textbf{5：適切（条件付き概数）} オゾンの約9割は成層圏。。全球オゾン鉛直分布の代表値。地表オゾン汚染の存在を否定しない。\par

\paragraph{再利用する考え方} 全球平均・地域差、貯留量・流量、自然温室効果・追加強制力、代表値・普遍則を区別する。大気の鉛直構造と環境・断熱気温減率の対象と成立条件を確認する。

\paragraph{共通理論} \hyperref[theory:atmosphere:layers]{第\ref*{theory:atmosphere:layers}節}, \hyperref[theory:atmosphere:lapse]{第\ref*{theory:atmosphere:lapse}節}, \hyperref[theory:atmosphere:ozone]{第\ref*{theory:atmosphere:ozone}節}

\par 参考：\cite{atm_noaa_layers,atm_us1976,atm_jma_ozone}。

\paragraph{関連問題} \hyperref[q:r04-reg-applied-22]{R04-REG-APPLIED-22}

\par 難易度 2（校準初稿）、優先度未評価。

\subsection{令和4年度 問22}
\label{q:r04-reg-applied-22}

\noindent\textbf{原典：} sources/R04/17\_applied\_science.pdf、PDF物理7頁。

\textbf{公式正答：} 1（single）。

\paragraph{問題の要旨} 対流圏の質量・気温減率・界面高度と成層圏温度、オゾンの分布を判定する。

\Needspace{4\baselineskip}\paragraph{確認した条件・記述（独自要約）}\mbox{}\par

\noindent\textbf{方向：} 最も不適切。\par

\paragraph{問われる核心} 大気の鉛直構造と環境・断熱気温減率

\paragraph{必要知識} 大気の鉛直構造と環境・断熱気温減率。環境減率の標準値は約6.5K/kmであり、3K/kmという肢が不適切。温度の観測平均と、気塊の断熱移動を表す9.8K/km等を混ぜない。界面は赤道側で高く、オゾンが多い成層圏ではUV加熱による温度上昇が生じる。

\paragraph{考え方・正答を選ぶ理由（査読待ち初稿）} 環境減率の標準値は約6.5K/kmであり、3K/kmという肢が不適切。温度の観測平均と、気塊の断熱移動を表す9.8K/km等を混ぜない。界面は赤道側で高く、オゾンが多い成層圏ではUV加熱による温度上昇が生じる。

\Needspace{4\baselineskip}\paragraph{選択肢の検討（査読待ち初稿）}\mbox{}\par

\noindent\textbf{1：不適切} 対流圏の気温減率は3℃/km。。標準的な対流圏環境減率は約6.5K/km。乾燥断熱9.8K/kmや湿潤断熱・実際の逆転層と区別する。\par

\noindent\textbf{2：適切（条件付き概数）} オゾンの約9割は成層圏。。全球オゾン鉛直分布の代表値。地表オゾン汚染の存在を否定しない。\par

\noindent\textbf{3：適切（条件付き概数）} 成層圏では高度とともに気温上昇。。オゾン吸収による加熱が上方の温度増加を作る。成層圏の下部がほぼ等温の場合を無視した厳密な全高度単調則ではない。\par

\noindent\textbf{4：適切（条件付き概数）} 対流圏界面は低緯度で高く高緯度で低い。。低緯度の高い温度と深い対流により界面は高くなる。\par

\noindent\textbf{5：適切（条件付き概数）} 対流圏質量は気圏の約3/4。。大気質量が下層へ集中する代表的な全球概数。界面の緯度・季節差で厳密値は変わる。\par

\paragraph{再利用する考え方} 全球平均・地域差、貯留量・流量、自然温室効果・追加強制力、代表値・普遍則を区別する。大気の鉛直構造と環境・断熱気温減率の対象と成立条件を確認する。

\paragraph{共通理論} \hyperref[theory:atmosphere:layers]{第\ref*{theory:atmosphere:layers}節}, \hyperref[theory:atmosphere:lapse]{第\ref*{theory:atmosphere:lapse}節}, \hyperref[theory:atmosphere:ozone]{第\ref*{theory:atmosphere:ozone}節}

\par 参考：\cite{atm_noaa_layers,atm_us1976,atm_jma_ozone}。

\paragraph{関連問題} \hyperref[q:r01-ret-applied-21]{R01-RET-APPLIED-21}

\par 難易度 2（校準初稿）、優先度未評価。

\subsection{令和7年度 問21}
\label{q:r07-reg-applied-21}

\noindent\textbf{原典：} sources/R07/17\_applied\_science.pdf、PDF物理7頁。

\textbf{公式正答：} 2（single）。

\par\textbf{検討未確定：} 肢1の10km密度1/4は標準モデル約0.34と差があり概数精度を査読待ち。原典の数値・密度という語は保持。

\paragraph{問題の要旨} 大気密度の概数、界面の緯度差、主要成分の混合、オゾン分布と層序を比較する。

\Needspace{4\baselineskip}\paragraph{確認した条件・記述（独自要約）}\mbox{}\par

\noindent\textbf{方向：} 最も不適切。\par

\paragraph{問われる核心} 温度層序・密度と緯度変化

\paragraph{必要知識} 温度層序・密度と緯度変化。低緯度の対流圏界面は高く、極域は低いので2が不適切。ただし1は標準モデルの10km密度比約0.34と1/4の精度差があり、原典の密度を圧力へ書き換えず部分校正に残す。公式正答2と別肢の概数精度の査読を分ける。

\paragraph{考え方・正答を選ぶ理由（査読待ち初稿）} 低緯度の対流圏界面は高く、極域は低いので2が不適切。ただし1は標準モデルの10km密度比約0.34と1/4の精度差があり、原典の密度を圧力へ書き換えず部分校正に残す。公式正答2と別肢の概数精度の査読を分ける。

\Needspace{4\baselineskip}\paragraph{選択肢の検討（査読待ち初稿）}\mbox{}\par

\noindent\textbf{1：概数精度に疑義} 10kmで質量密度は地表の約1/4。。原典は密度と記載。標準モデルで約0.414kg/m\textsuperscript{3}対地表1.225、比約0.338となり1/4とは差がある。圧力比約0.26との混同があり得るが、原典を圧力へ改変しない。\par

\noindent\textbf{2：不適切} 対流圏界面は赤道より極域で高い。。低緯度では深い対流と高温により界面は高い。逆にしたこの肢が明白な誤り。\par

\noindent\textbf{3：適切（主要成分）} N2/O2/Arの組成は数10kmまでほぼ一定。。均質圏の主要長寿命成分の混合比の概数。微量気体や水蒸気の分布は同じではない。\par

\noindent\textbf{4：適切（概数）} 全オゾンの約10\%が対流圏。。約90\%成層圏という代表分布の補集合であり整合する。\par

\noindent\textbf{5：適切} 下から対流圏・成層圏・中間圏・熱圏。。主として気温の鉛直勾配を基準とする区分である。\par

\paragraph{再利用する考え方} 全球平均・地域差、貯留量・流量、自然温室効果・追加強制力、代表値・普遍則を区別する。温度層序・密度と緯度変化の対象と成立条件を確認する。

\paragraph{共通理論} \hyperref[theory:atmosphere:layers]{第\ref*{theory:atmosphere:layers}節}, \hyperref[theory:atmosphere:hydrostatic]{第\ref*{theory:atmosphere:hydrostatic}節}, \hyperref[theory:atmosphere:ozone]{第\ref*{theory:atmosphere:ozone}節}

\par 参考：\cite{atm_noaa_layers,atm_us1976,atm_jma_ozone}。

\paragraph{関連問題} \hyperref[q:r01-reg-applied-19]{R01-REG-APPLIED-19}

\par 難易度 2（校準初稿）、優先度未評価。

\section{地球の受光幾何と放射収支}

\subsection{平成26年度 問26}
\label{q:h26-reg-applied-26}

\noindent\textbf{原典：} sources/H26/17\_applied\_science.pdf、PDF物理11頁。

\textbf{公式正答：} 1（single）。

\paragraph{問題の要旨} 地球の反射率、太陽放射帯域、季節原因、太陽定数と緯度間の熱収支を判定する。

\Needspace{4\baselineskip}\paragraph{確認した条件・記述（独自要約）}\mbox{}\par

\noindent\textbf{方向：} 最も適切。\par

\paragraph{問われる核心} 短波入射と長波放射の地球規模の収支

\paragraph{必要知識} 短波入射と長波放射の地球規模の収支。反射と吸収、短波日射と長波地球放射を分ける。反射の約31\%が問題採用の適切な概数である。四季は地軸傾斜、太陽定数は大気上端の垂直面、放射不均衡は低緯度から高緯度への輸送につながる。31\%を現在の精密観測で必ず同じとするのではなく、問題の比較精度で扱う。

\paragraph{考え方・正答を選ぶ理由（査読待ち初稿）} 反射と吸収、短波日射と長波地球放射を分ける。反射の約31\%が問題採用の適切な概数である。四季は地軸傾斜、太陽定数は大気上端の垂直面、放射不均衡は低緯度から高緯度への輸送につながる。31\%を現在の精密観測で必ず同じとするのではなく、問題の比較精度で扱う。

\Needspace{4\baselineskip}\paragraph{選択肢の検討（査読待ち初稿）}\mbox{}\par

\noindent\textbf{1：最適} 入射の約31\%が雲・大気・地表で反射散乱され宇宙へ戻る。。問題で採る全球平均反射率約3割の代表値と合う。時間・雲・地表で変化するので正確な不変定数ではない。\par

\noindent\textbf{2：不適切} 太陽放射の約半分はガンマ線・X線。。太陽放射の主な帯域は可視・近赤外であり、高エネルギー放射が半分ではない。\par

\noindent\textbf{3：不適切} 地球太陽距離の変化が四季の原因。。地軸の傾きで太陽高度と日長が変わることが主因であり、南北半球の季節は逆になる。\par

\noindent\textbf{4：不適切} 大気下端の太陽垂直面の受光を太陽定数という。。大気の吸収散乱の影響を除いた大気上端で、太陽光に垂直な単位面積あたりの放射束で定義する。\par

\noindent\textbf{5：不適切} 吸収日射の緯度差は小さく、高緯度の放出が低緯度より大きい。。低緯度の日射吸収は大きく、正味収支の緯度差を大気・海洋の熱輸送が補う。\par

\paragraph{再利用する考え方} 全球平均・地域差、貯留量・流量、自然温室効果・追加強制力、代表値・普遍則を区別する。短波入射と長波放射の地球規模の収支の対象と成立条件を確認する。

\paragraph{共通理論} \hyperref[theory:atmosphere:radiation]{第\ref*{theory:atmosphere:radiation}節}, \hyperref[theory:atmosphere:transport]{第\ref*{theory:atmosphere:transport}節}

\par 参考：\cite{atm_nasa_budget,atm_nasa_albedo,atm_nasa_transport}。

\paragraph{関連問題} \hyperref[q:r01-reg-applied-22]{R01-REG-APPLIED-22}, \hyperref[q:r06-reg-applied-20]{R06-REG-APPLIED-20}

\par 難易度 2（校準初稿）、優先度未評価。

\subsection{平成28年度 問26}
\label{q:h28-reg-applied-26}

\noindent\textbf{原典：} sources/H28/17\_applied\_science.pdf、PDF物理11頁。

\textbf{公式正答：} 1（single）。

\paragraph{問題の要旨} 指定された地球半径と太陽定数から、地球が毎秒受け取る日射の総量を求める。

\Needspace{4\baselineskip}\paragraph{確認した条件・記述（独自要約）}\mbox{}\par

\noindent\textbf{方向：} 最も近い値。\par

\noindent\textbf{半径：} 6370km=\allowbreak{}6.370×10\textasciicircum{}6m。\par

\noindent\textbf{太陽定数：} 1.368×10\textasciicircum{}3W/m\textsuperscript{2}。\par

\noindent\textbf{大気：} H28は反射吸収無視、H30は大気上端で定義。\par

\paragraph{問われる核心} 地球の受光断面と球面平均を区別する

\paragraph{必要知識} 地球の受光断面と球面平均を区別する。平行入射光を遮る投影面積は\(\pi R^2\)で、球表面積ではない。\(P=S_0\pi R^2=1368\pi(6.370\times10^6)^2\simeq1.744\times10^{17}\ {\rm W}\)。Wは毎秒のエネルギーを表す。全球平均面積当たりならこの総量を\(4\pi R^2\)で割り\(S_0/4\)となるが、今回は総量を問う。

\paragraph{考え方・正答を選ぶ理由（査読待ち初稿）} 平行入射光を遮る投影面積は\(\pi R^2\)で、球表面積ではない。\(P=S_0\pi R^2=1368\pi(6.370\times10^6)^2\simeq1.744\times10^{17}\ {\rm W}\)。Wは毎秒のエネルギーを表す。全球平均面積当たりならこの総量を\(4\pi R^2\)で割り\(S_0/4\)となるが、今回は総量を問う。

\Needspace{4\baselineskip}\paragraph{選択肢の検討（査読待ち初稿）}\mbox{}\par

\noindent\textbf{1：最適} \(1.74\times10^{17}\ {\rm W}\)。受光円の計算1.74×10\textasciicircum{}17Wと一致。\par

\noindent\textbf{2：不適切} \(3.49\times10^{17}\ {\rm W}\)。計算値の約2倍。受光断面と球表面を混同せず、単位換算と面積を確認する。\par

\noindent\textbf{3：不適切} \(6.97\times10^{17}\ {\rm W}\)。計算値の約4倍。受光断面と球表面を混同せず、単位換算と面積を確認する。\par

\noindent\textbf{4：不適切} \(1.39\times10^{18}\ {\rm W}\)。計算値の約8倍。受光断面と球表面を混同せず、単位換算と面積を確認する。\par

\noindent\textbf{5：不適切} \(2.79\times10^{18}\ {\rm W}\)。計算値の約16倍。受光断面と球表面を混同せず、単位換算と面積を確認する。\par

\paragraph{再利用する考え方} 全球平均・地域差、貯留量・流量、自然温室効果・追加強制力、代表値・普遍則を区別する。地球の受光断面と球面平均を区別するの対象と成立条件を確認する。

\paragraph{共通理論} \hyperref[theory:atmosphere:geometry]{第\ref*{theory:atmosphere:geometry}節}

\par 参考：\cite{atm_nasa_budget}。

\paragraph{関連問題} \hyperref[q:h30-reg-applied-22]{H30-REG-APPLIED-22}, \hyperref[q:r01-ret-applied-20]{R01-RET-APPLIED-20}

\par 難易度 2（校準初稿）、優先度未評価。

\subsection{平成30年度 問22}
\label{q:h30-reg-applied-22}

\noindent\textbf{原典：} sources/H30/17\_applied\_science.pdf、PDF物理7頁。

\textbf{公式正答：} 2（single）。

\paragraph{問題の要旨} 指定された地球半径と太陽定数から、地球が毎秒受け取る日射の総量を求める。

\Needspace{4\baselineskip}\paragraph{確認した条件・記述（独自要約）}\mbox{}\par

\noindent\textbf{方向：} 最も近い値。\par

\noindent\textbf{半径：} 6370km=\allowbreak{}6.370×10\textasciicircum{}6m。\par

\noindent\textbf{太陽定数：} 1.368×10\textasciicircum{}3W/m\textsuperscript{2}。\par

\noindent\textbf{大気：} H28は反射吸収無視、H30は大気上端で定義。\par

\paragraph{問われる核心} 地球の受光断面と球面平均を区別する

\paragraph{必要知識} 地球の受光断面と球面平均を区別する。平行入射光を遮る投影面積は\(\pi R^2\)で、球表面積ではない。\(P=S_0\pi R^2=1368\pi(6.370\times10^6)^2\simeq1.744\times10^{17}\ {\rm W}\)。Wは毎秒のエネルギーを表す。全球平均面積当たりならこの総量を\(4\pi R^2\)で割り\(S_0/4\)となるが、今回は総量を問う。

\paragraph{考え方・正答を選ぶ理由（査読待ち初稿）} 平行入射光を遮る投影面積は\(\pi R^2\)で、球表面積ではない。\(P=S_0\pi R^2=1368\pi(6.370\times10^6)^2\simeq1.744\times10^{17}\ {\rm W}\)。Wは毎秒のエネルギーを表す。全球平均面積当たりならこの総量を\(4\pi R^2\)で割り\(S_0/4\)となるが、今回は総量を問う。

\Needspace{4\baselineskip}\paragraph{選択肢の検討（査読待ち初稿）}\mbox{}\par

\noindent\textbf{1：不適切} \(8.71\times10^{16}\ {\rm W}\)。計算値の約0.5倍。受光断面と球表面を混同せず、単位換算と面積を確認する。\par

\noindent\textbf{2：最適} \(1.74\times10^{17}\ {\rm W}\)。受光円の計算1.74×10\textasciicircum{}17Wと一致。\par

\noindent\textbf{3：不適切} \(3.49\times10^{17}\ {\rm W}\)。計算値の約2倍。受光断面と球表面を混同せず、単位換算と面積を確認する。\par

\noindent\textbf{4：不適切} \(6.97\times10^{17}\ {\rm W}\)。計算値の約4倍。受光断面と球表面を混同せず、単位換算と面積を確認する。\par

\noindent\textbf{5：不適切} \(1.39\times10^{18}\ {\rm W}\)。計算値の約8倍。受光断面と球表面を混同せず、単位換算と面積を確認する。\par

\paragraph{再利用する考え方} 全球平均・地域差、貯留量・流量、自然温室効果・追加強制力、代表値・普遍則を区別する。地球の受光断面と球面平均を区別するの対象と成立条件を確認する。

\paragraph{共通理論} \hyperref[theory:atmosphere:geometry]{第\ref*{theory:atmosphere:geometry}節}

\par 参考：\cite{atm_nasa_budget}。

\paragraph{関連問題} \hyperref[q:h28-reg-applied-26]{H28-REG-APPLIED-26}

\par 難易度 2（校準初稿）、優先度未評価。

\subsection{令和元年度 問22}
\label{q:r01-reg-applied-22}

\noindent\textbf{原典：} sources/R01/17\_applied\_science.pdf、PDF物理7頁。

\textbf{公式正答：} 3（single）。

\paragraph{問題の要旨} 季節・緯度収支、太陽放射の主帯域、反射率、太陽定数の定義を比較する。

\Needspace{4\baselineskip}\paragraph{確認した条件・記述（独自要約）}\mbox{}\par

\noindent\textbf{方向：} 最も適切。\par

\paragraph{問われる核心} 放射の波長帯・反射率・幾何の区別

\paragraph{必要知識} 放射の波長帯・反射率・幾何の区別。太陽放射の波長分布は可視・近赤外が中心である。3の約半分可視が適切。他肢の距離を季節原因とする説明、緯度収支の逆転、15\%反射、大気下端定義を分けて否定する。

\paragraph{考え方・正答を選ぶ理由（査読待ち初稿）} 太陽放射の波長分布は可視・近赤外が中心である。3の約半分可視が適切。他肢の距離を季節原因とする説明、緯度収支の逆転、15\%反射、大気下端定義を分けて否定する。

\Needspace{4\baselineskip}\paragraph{選択肢の検討（査読待ち初稿）}\mbox{}\par

\noindent\textbf{1：不適切} 地球太陽距離の変化が四季の原因。。地軸の傾きで太陽高度と日長が変わることが主因であり、南北半球の季節は逆になる。\par

\noindent\textbf{2：不適切} 吸収日射の緯度差は小さく、高緯度の放出が低緯度より大きい。。低緯度の日射吸収は大きく、正味収支の緯度差を大気・海洋の熱輸送が補う。\par

\noindent\textbf{3：最適} 太陽放射の約半分は可視光線。。太陽光スペクトルの主な一部が可視であるという概数と合う。残りは主に赤外と一部紫外。\par

\noindent\textbf{4：不適切} 入射の約15\%が反射散乱され宇宙へ戻る。。全球反射は約3割の代表値であり、15\%は小さすぎる。\par

\noindent\textbf{5：不適切} 大気下端の太陽垂直面の受光を太陽定数という。。大気の吸収散乱の影響を除いた大気上端で、太陽光に垂直な単位面積あたりの放射束で定義する。\par

\paragraph{再利用する考え方} 全球平均・地域差、貯留量・流量、自然温室効果・追加強制力、代表値・普遍則を区別する。放射の波長帯・反射率・幾何の区別の対象と成立条件を確認する。

\paragraph{共通理論} \hyperref[theory:atmosphere:radiation]{第\ref*{theory:atmosphere:radiation}節}, \hyperref[theory:atmosphere:geometry]{第\ref*{theory:atmosphere:geometry}節}

\par 参考：\cite{atm_nasa_budget,atm_nasa_transport}。

\paragraph{関連問題} \hyperref[q:h26-reg-applied-26]{H26-REG-APPLIED-26}

\par 難易度 2（校準初稿）、優先度未評価。

\subsection{令和元年度再試験 問20}
\label{q:r01-ret-applied-20}

\noindent\textbf{原典：} sources/R01\_retest/17\_applied\_science.pdf、PDF物理7頁。

\textbf{公式正答：} 2（single）。

\paragraph{問題の要旨} 指定された全球平均入射342W/m\textsuperscript{2}に対し、反射散乱で宇宙へ戻る放射束を選ぶ。

\Needspace{4\baselineskip}\paragraph{確認した条件・記述（独自要約）}\mbox{}\par

\noindent\textbf{方向：} 最も適切。\par

\noindent\textbf{指定：} 大気上端の平均342W/m\textsuperscript{2}。宇宙へ戻る反射散乱の和。\par

\paragraph{問われる核心} 全球アルベドと球面平均放射束

\paragraph{必要知識} 全球アルベドと球面平均放射束。問題の採る全球平均値では反射率\(\alpha\simeq0.31\)なので\(F_{\rm refl}=\alpha F_{\rm in}\simeq0.31\times342\simeq106\ {\rm W/m^2}\)。最近の107を選ぶ。342は球面平均であり、ここで再び4で割らない。NASAの代表収支図にも342/107があるが、現在のすべての年の厳密な観測値とはしない。

\paragraph{考え方・正答を選ぶ理由（査読待ち初稿）} 問題の採る全球平均値では反射率\(\alpha\simeq0.31\)なので\(F_{\rm refl}=\alpha F_{\rm in}\simeq0.31\times342\simeq106\ {\rm W/m^2}\)。最近の107を選ぶ。342は球面平均であり、ここで再び4で割らない。NASAの代表収支図にも342/107があるが、現在のすべての年の厳密な観測値とはしない。

\Needspace{4\baselineskip}\paragraph{選択肢の検討（査読待ち初稿）}\mbox{}\par

\noindent\textbf{1：不適切} 68W/m\textsuperscript{2}。入射の約20\%で代表約31\%より小さい。\par

\noindent\textbf{2：最適} 107W/m\textsuperscript{2}。107/342=約0.313で約31\%の全球反射と合う。\par

\noindent\textbf{3：不適切} 137W/m\textsuperscript{2}。約40\%で問題の全球代表反射率より大きい。\par

\noindent\textbf{4：不適切} 171W/m\textsuperscript{2}。約50\%で問題の全球代表反射率より大きい。\par

\noindent\textbf{5：不適切} 205W/m\textsuperscript{2}。約60\%で反射を過大評価している。\par

\paragraph{再利用する考え方} 全球平均・地域差、貯留量・流量、自然温室効果・追加強制力、代表値・普遍則を区別する。全球アルベドと球面平均放射束の対象と成立条件を確認する。

\paragraph{共通理論} \hyperref[theory:atmosphere:radiation]{第\ref*{theory:atmosphere:radiation}節}, \hyperref[theory:atmosphere:geometry]{第\ref*{theory:atmosphere:geometry}節}

\par 参考：\cite{atm_nasa_budget,atm_nasa_albedo}。

\paragraph{関連問題} \hyperref[q:h28-reg-applied-26]{H28-REG-APPLIED-26}

\par 難易度 2（校準初稿）、優先度未評価。

\subsection{令和6年度 問20}
\label{q:r06-reg-applied-20}

\noindent\textbf{原典：} sources/R06/17\_applied\_science.pdf、PDF物理6頁。

\textbf{公式正答：} 3（single）。

\paragraph{問題の要旨} 地球の反射率、太陽放射帯域、季節原因、太陽定数と緯度間の熱収支を判定する。

\Needspace{4\baselineskip}\paragraph{確認した条件・記述（独自要約）}\mbox{}\par

\noindent\textbf{方向：} 最も適切。\par

\paragraph{問われる核心} 短波入射と長波放射の地球規模の収支

\paragraph{必要知識} 短波入射と長波放射の地球規模の収支。反射と吸収、短波日射と長波地球放射を分ける。反射の約31\%が問題採用の適切な概数である。四季は地軸傾斜、太陽定数は大気上端の垂直面、放射不均衡は低緯度から高緯度への輸送につながる。31\%を現在の精密観測で必ず同じとするのではなく、問題の比較精度で扱う。

\paragraph{考え方・正答を選ぶ理由（査読待ち初稿）} 反射と吸収、短波日射と長波地球放射を分ける。反射の約31\%が問題採用の適切な概数である。四季は地軸傾斜、太陽定数は大気上端の垂直面、放射不均衡は低緯度から高緯度への輸送につながる。31\%を現在の精密観測で必ず同じとするのではなく、問題の比較精度で扱う。

\Needspace{4\baselineskip}\paragraph{選択肢の検討（査読待ち初稿）}\mbox{}\par

\noindent\textbf{1：不適切} 地球太陽距離の変化が四季の原因。。地軸の傾きで太陽高度と日長が変わることが主因であり、南北半球の季節は逆になる。\par

\noindent\textbf{2：不適切} 吸収日射の緯度差は小さく、高緯度の放出が低緯度より大きい。。低緯度の日射吸収は大きく、正味収支の緯度差を大気・海洋の熱輸送が補う。\par

\noindent\textbf{3：最適} 入射の約31\%が雲・大気・地表で反射散乱され宇宙へ戻る。。問題で採る全球平均反射率約3割の代表値と合う。時間・雲・地表で変化するので正確な不変定数ではない。\par

\noindent\textbf{4：不適切} 大気下端の太陽垂直面の受光を太陽定数という。。大気の吸収散乱の影響を除いた大気上端で、太陽光に垂直な単位面積あたりの放射束で定義する。\par

\noindent\textbf{5：不適切} 太陽放射の約半分はガンマ線・X線。。太陽放射の主な帯域は可視・近赤外であり、高エネルギー放射が半分ではない。\par

\paragraph{再利用する考え方} 全球平均・地域差、貯留量・流量、自然温室効果・追加強制力、代表値・普遍則を区別する。短波入射と長波放射の地球規模の収支の対象と成立条件を確認する。

\paragraph{共通理論} \hyperref[theory:atmosphere:radiation]{第\ref*{theory:atmosphere:radiation}節}, \hyperref[theory:atmosphere:transport]{第\ref*{theory:atmosphere:transport}節}

\par 参考：\cite{atm_nasa_budget,atm_nasa_albedo,atm_nasa_transport}。

\paragraph{関連問題} \hyperref[q:h26-reg-applied-26]{H26-REG-APPLIED-26}

\par 難易度 2（校準初稿）、優先度未評価。

\section{大気海洋相互作用と南方振動}

\subsection{平成29年度 問21}
\label{q:h29-reg-applied-21}

\noindent\textbf{原典：} sources/H29/17\_applied\_science.pdf、PDF物理7頁。

\textbf{公式正答：} 4（single）。

\paragraph{問題の要旨} 地衡風、境界層、海陸風、ENSO、ロスビー波を時間・空間規模から比較する。

\Needspace{4\baselineskip}\paragraph{確認した条件・記述（独自要約）}\mbox{}\par

\noindent\textbf{方向：} 最も不適切。\par

\paragraph{問われる核心} 循環現象の力学と時間尺度

\paragraph{必要知識} 循環現象の力学と時間尺度。地衡風の力の釣合い、境界層の地表交換、海陸風の日周期、南方振動の年〜数年を分ける。南方振動は太平洋東西の気圧がシーソー状に変化する大気側の現象で、4の数十年は時間尺度が違う。ロスビー波と同じ「振動」という語だけで同じ周期にまとめない。

\paragraph{考え方・正答を選ぶ理由（査読待ち初稿）} 地衡風の力の釣合い、境界層の地表交換、海陸風の日周期、南方振動の年〜数年を分ける。南方振動は太平洋東西の気圧がシーソー状に変化する大気側の現象で、4の数十年は時間尺度が違う。ロスビー波と同じ「振動」という語だけで同じ周期にまとめない。

\Needspace{4\baselineskip}\paragraph{選択肢の検討（査読待ち初稿）}\mbox{}\par

\noindent\textbf{1：適切（条件付き概数）} 地衡風は傾度力とコリオリ力が釣り合い等圧線沿い。。摩擦が小さく加速度が相対的に小さい回転流の近似である。\par

\noindent\textbf{2：適切（条件付き概数）} 地表から約1kmの乱流支配層は大気境界層。。地表の摩擦・熱・水蒸気の交換が強く影響する層の代表厚さ。固定で1kmに鋭い境界があるとはしない。\par

\noindent\textbf{3：適切（条件付き概数）} 海陸風は温度差による日変化の局地風。。海と陸の加熱・冷却の差が水平気圧勾配を作る。\par

\noindent\textbf{4：不適切} 南方振動は熱帯太平洋の気圧の数十年規模東西振動。。ENSOに伴う南方振動は数年規模で、数十年規模の変動とは区別する。\par

\noindent\textbf{5：適切（条件付き概数）} 中緯度偏西風帯の波動はロスビー波。。コリオリ係数の緯度変化が復元性を与える惑星波。中緯度の擾乱と熱輸送に関わる。\par

\paragraph{再利用する考え方} 全球平均・地域差、貯留量・流量、自然温室効果・追加強制力、代表値・普遍則を区別する。循環現象の力学と時間尺度の対象と成立条件を確認する。

\paragraph{共通理論} \hyperref[theory:atmosphere:wind]{第\ref*{theory:atmosphere:wind}節}, \hyperref[theory:atmosphere:local]{第\ref*{theory:atmosphere:local}節}, \hyperref[theory:atmosphere:enso]{第\ref*{theory:atmosphere:enso}節}

\par 参考：\cite{atm_jma_enso,atm_noaa_winds,atm_jma_dynamic}。

\par 難易度 2（校準初稿）、優先度未評価。

\section{降水の酸塩基平衡と酸性沈着}

\subsection{平成27年度 問26}
\label{q:h27-reg-applied-26}

\noindent\textbf{原典：} sources/H27/17\_applied\_science.pdf、PDF物理10頁。

\textbf{公式正答：} 2（single）。

\paragraph{問題の要旨} 降水のpHと水素イオン濃度、炭酸平衡・中和、当時の日本の酸性度を判定する。

\Needspace{4\baselineskip}\paragraph{確認した条件・記述（独自要約）}\mbox{}\par

\noindent\textbf{方向：} 最も不適切。\par

\paragraph{問われる核心} 対数尺度と降水の酸塩基平衡

\paragraph{必要知識} 対数尺度と降水の酸塩基平衡。pHは厳密にはH+活量の負対数であり、希薄溶液ではモル濃度に近似する。\([{\rm H}^+]\simeq10^{-3.3}=5.01\times10^{-4}\ {\rm mol/L}\)なので約501μmol/Lとなり2が不適切。2008〜2012の観測を出題当時の根拠として扱い、現在すべての地点で同じという説明へ流用しない。

\paragraph{考え方・正答を選ぶ理由（査読待ち初稿）} pHは厳密にはH+活量の負対数であり、希薄溶液ではモル濃度に近似する。\([{\rm H}^+]\simeq10^{-3.3}=5.01\times10^{-4}\ {\rm mol/L}\)なので約501μmol/Lとなり2が不適切。2008〜2012の観測を出題当時の根拠として扱い、現在すべての地点で同じという説明へ流用しない。

\Needspace{4\baselineskip}\paragraph{選択肢の検討（査読待ち初稿）}\mbox{}\par

\noindent\textbf{1：適切（条件付き概数）} CO2のみでも25℃の降水pHは約5.6。。CO2と平衡する希薄水の代表的目安であり、全降水の固定値ではない。\par

\noindent\textbf{2：不適切} pH3.3は約30μmol/LのH+。。希薄近似で10\textasciicircum{}-3.3 mol/L=約501μmol/L。30μmol/LならpH約4.52である。\par

\noindent\textbf{3：適切（条件付き概数）} pH5はpH6の10倍のH+。。対数が1違うと濃度比は10となる。pH値を単純に比で比較しない。\par

\noindent\textbf{4：適切（条件付き概数）} 日本には年平均pH5未満の地点が多い。。気象庁が紹介する2008〜2012調査は地点年平均4.48〜5.37、全地点5年4.72で、原典当時の定性的傾向と整合する。全地点5未満ではない。\par

\noindent\textbf{5：適切（条件付き概数）} 乾燥地ではアルカリ粒子・アンモニア中和でpH7に近い場合がある。。土壌の炭酸塩等やNH3は酸を中和する。地域の組成次第で、全乾燥地で固定7とはしない。\par

\paragraph{再利用する考え方} 全球平均・地域差、貯留量・流量、自然温室効果・追加強制力、代表値・普遍則を区別する。対数尺度と降水の酸塩基平衡の対象と成立条件を確認する。

\paragraph{共通理論} \hyperref[theory:atmosphere:acid]{第\ref*{theory:atmosphere:acid}節}

\par 参考：\cite{atm_jma_acid,atm_jma_acid2014}。

\paragraph{関連問題} \hyperref[q:r05-reg-applied-20]{R05-REG-APPLIED-20}

\par 難易度 2（校準初稿）、優先度未評価。

\subsection{令和5年度 問20}
\label{q:r05-reg-applied-20}

\noindent\textbf{原典：} sources/R05/17\_applied\_science.pdf、PDF物理7頁。

\textbf{公式正答：} 1（single）。

\paragraph{問題の要旨} 降水pHの地域差、対数比、炭酸平衡、中和と酸性沈着の影響を判定する。

\Needspace{4\baselineskip}\paragraph{確認した条件・記述（独自要約）}\mbox{}\par

\noindent\textbf{方向：} 最も不適切。\par

\paragraph{問われる核心} 降水の酸塩基平衡と酸性沈着の対象

\paragraph{必要知識} 降水の酸塩基平衡と酸性沈着の対象。海に接している雨と海水は別の相であり、海水のpHをそのまま降水へ代入できない。CO2平衡・硫酸硝酸・塩基性土壌成分等の収支で判断するため1が不適切。pH5.6を人為汚染がない全降水の固定値にもせず、中和と自然酸性源を考える。

\paragraph{考え方・正答を選ぶ理由（査読待ち初稿）} 海に接している雨と海水は別の相であり、海水のpHをそのまま降水へ代入できない。CO2平衡・硫酸硝酸・塩基性土壌成分等の収支で判断するため1が不適切。pH5.6を人為汚染がない全降水の固定値にもせず、中和と自然酸性源を考える。

\Needspace{4\baselineskip}\paragraph{選択肢の検討（査読待ち初稿）}\mbox{}\par

\noindent\textbf{1：不適切} 海洋上・極域の降水はpH8以上が多い。。海水の弱アルカリ性を降水へ移せない。雨のpHはCO2と酸性/中和成分で決まり、場所名だけで一般に8以上としない。\par

\noindent\textbf{2：適切（条件付き概数）} pH5はpH6の10倍のH+。。10\textasciicircum{}-5/10\textasciicircum{}-6=\allowbreak{}10。対数尺度の比較である。\par

\noindent\textbf{3：適切（条件付き概数）} 大気CO2と平衡する水は約pH5.6。。25℃付近の希薄水の代表的目安。CO2分圧や温度、他の溶質が変われば同じ値ではない。\par

\noindent\textbf{4：適切（条件付き概数）} 乾燥地でNH3等中和によりpH7近くなる場合が多い。。アンモニア・土壌由来アルカリ成分は酸を中和する。全乾燥地で必ず7の保証ではない。\par

\noindent\textbf{5：適切（条件付き概数）} 酸性雨は森林・魚類・歴史建造物へ被害を起こす。。酸性沈着による水・土壌の変化と鉱物・金属への作用が生態系と構造物へ影響する。\par

\paragraph{再利用する考え方} 全球平均・地域差、貯留量・流量、自然温室効果・追加強制力、代表値・普遍則を区別する。降水の酸塩基平衡と酸性沈着の対象の対象と成立条件を確認する。

\paragraph{共通理論} \hyperref[theory:atmosphere:acid]{第\ref*{theory:atmosphere:acid}節}

\par 参考：\cite{atm_jma_acid,atm_jma_acid2014}。

\paragraph{関連問題} \hyperref[q:h27-reg-applied-26]{H27-REG-APPLIED-26}

\par 難易度 2（校準初稿）、優先度未評価。

\section{水収支・貯留・平均滞留時間}

\subsection{平成23年度 問20}
\label{q:h23-reg-applied-20}

\noindent\textbf{原典：} sources/H23/17\_applied\_science.pdf、PDF物理7頁。

\textbf{公式正答：} 4（single）。

\paragraph{問題の要旨} 大気中水蒸気の全球平均滞留時間に最も近い日数を選ぶ。

\Needspace{4\baselineskip}\paragraph{確認した条件・記述（独自要約）}\mbox{}\par

\noindent\textbf{方向：} 最も適切。\par

\paragraph{問われる核心} 平均滞留時間を貯留量と流量から定義する

\paragraph{必要知識} 平均滞留時間を貯留量と流量から定義する。定常に近い水の貯留量Wと流出率Fなら平均交換時間の目安は\(\tau=W/F\)である。USGSの大気12900立方km、日蒸発散1170立方kmという別推定の値でも約11日となる。推定法により8〜10日等の差はあるが、本肢では10日が最近。個々の水分子が必ず10日で降るという周期ではない。

\paragraph{考え方・正答を選ぶ理由（査読待ち初稿）} 定常に近い水の貯留量Wと流出率Fなら平均交換時間の目安は\(\tau=W/F\)である。USGSの大気12900立方km、日蒸発散1170立方kmという別推定の値でも約11日となる。推定法により8〜10日等の差はあるが、本肢では10日が最近。個々の水分子が必ず10日で降るという周期ではない。

\Needspace{4\baselineskip}\paragraph{選択肢の検討（査読待ち初稿）}\mbox{}\par

\noindent\textbf{1：不適切} 200日。代表的な約10日に対し約20倍で長すぎる。\par

\noindent\textbf{2：不適切} 120日。約10日の一桁日数とは整合しない。\par

\noindent\textbf{3：不適切} 50日。大気の小さい貯留量と大きい降水流量の比より長すぎる。\par

\noindent\textbf{4：最適} 10日。全球水量を出入りする水量率で割った約8〜11日の代表範囲に最近。\par

\noindent\textbf{5：不適切} 36時間。約1.5日で代表的な全球交換時間より短い。\par

\paragraph{再利用する考え方} 全球平均・地域差、貯留量・流量、自然温室効果・追加強制力、代表値・普遍則を区別する。平均滞留時間を貯留量と流量から定義するの対象と成立条件を確認する。

\paragraph{共通理論} \hyperref[theory:atmosphere:water]{第\ref*{theory:atmosphere:water}節}

\par 参考：\cite{atm_usgs_water,atm_usgs_cycle}。

\par 難易度 2（校準初稿）、優先度未評価。

\subsection{平成28年度 問25}
\label{q:h28-reg-applied-25}

\noindent\textbf{原典：} sources/H28/17\_applied\_science.pdf、PDF物理10頁。

\textbf{公式正答：} 2（single）。

\paragraph{問題の要旨} 水蒸気柱・降水、帯水層係数、氷晶過程、宙水、海陸水収支を比較する。

\Needspace{4\baselineskip}\paragraph{確認した条件・記述（独自要約）}\mbox{}\par

\noindent\textbf{方向：} 最も不適切。\par

\paragraph{問われる核心} 水の貯留と輸送の測定量を分ける

\paragraph{必要知識} 水の貯留と輸送の測定量を分ける。貯留を表す係数と輸送を表す係数を次元で区別する。透水量係数は\(T=Kb\)で単位は面積/時間、比貯留率×厚さは無次元の貯留係数なので2が不適切。可降水量と日雨量も貯留対流量であり、一日の雨が柱の初期貯留を超えても保存則には反しない。

\paragraph{考え方・正答を選ぶ理由（査読待ち初稿）} 貯留を表す係数と輸送を表す係数を次元で区別する。透水量係数は\(T=Kb\)で単位は面積/時間、比貯留率×厚さは無次元の貯留係数なので2が不適切。可降水量と日雨量も貯留対流量であり、一日の雨が柱の初期貯留を超えても保存則には反しない。

\Needspace{4\baselineskip}\paragraph{選択肢の検討（査読待ち初稿）}\mbox{}\par

\noindent\textbf{1：適切（条件付き概数）} 日降水量は可降水量を超えることがある。。可降水量はある時点の柱の貯留で、日降水は時間積分した流量。横からの水蒸気補給があるので上限とはならない。\par

\noindent\textbf{2：不適切} 透水量係数は比貯留率と厚さの積。。T=\allowbreak{}KbでKは透水係数。Ss bは被圧帯水層の貯留係数で、流れやすさとは別量。\par

\noindent\textbf{3：適切（条件付き概数）} 氷晶過程を経る冷たい雨は中緯度降雨に多い。。氷粒子の成長後の融解が雨となる。地上の雨温度が冷たいという命名ではない。\par

\noindent\textbf{4：適切（条件付き概数）} 不飽和帯の低透水層の上に宙水が形成されることがある。。透水性の小さい層が下方への排水を制限し、局所的な地下水面を作り得る。\par

\noindent\textbf{5：適切（条件付き概数）} 海洋は蒸発超過、陸域は降水超過。。長期・全球平均では海洋から陸へ水蒸気が移り、陸から河川等で戻る。個々の地点の常時の関係ではない。\par

\paragraph{再利用する考え方} 全球平均・地域差、貯留量・流量、自然温室効果・追加強制力、代表値・普遍則を区別する。水の貯留と輸送の測定量を分けるの対象と成立条件を確認する。

\paragraph{共通理論} \hyperref[theory:atmosphere:water]{第\ref*{theory:atmosphere:water}節}, \hyperref[theory:atmosphere:rain]{第\ref*{theory:atmosphere:rain}節}, \hyperref[theory:atmosphere:aquifer]{第\ref*{theory:atmosphere:aquifer}節}

\par 参考：\cite{atm_noaa_pwat,atm_mri_cloud,atm_usgs_transmissivity,atm_usgs_perched}。

\par 難易度 2（校準初稿）、優先度未評価。

\subsection{平成29年度 問19}
\label{q:h29-reg-applied-19}

\noindent\textbf{原典：} sources/H29/17\_applied\_science.pdf、PDF物理6頁。

\textbf{公式正答：} 4（single）。

\paragraph{問題の要旨} 潜熱・温室効果、海陸の水収支、氷床の貯留と輸送、平均滞留時間を判定する。

\Needspace{4\baselineskip}\paragraph{確認した条件・記述（独自要約）}\mbox{}\par

\noindent\textbf{方向：} 最も不適切。\par

\paragraph{問われる核心} 水の貯留・交換とエネルギー輸送

\paragraph{必要知識} 水の貯留・交換とエネルギー輸送。交換時間は貯留量/流量である。氷床は大量の水を保持するが長い時間をかけて交換するので、貯留と輸送がともに大きいという肢が不適切。水蒸気の約10日の滞留、潜熱輸送、自然温室効果を同じ分子の別の役割として整理する。水収支の海陸差は地域ごとの恒常的な降水保証ではない。

\paragraph{考え方・正答を選ぶ理由（査読待ち初稿）} 交換時間は貯留量/流量である。氷床は大量の水を保持するが長い時間をかけて交換するので、貯留と輸送がともに大きいという肢が不適切。水蒸気の約10日の滞留、潜熱輸送、自然温室効果を同じ分子の別の役割として整理する。水収支の海陸差は地域ごとの恒常的な降水保証ではない。

\Needspace{4\baselineskip}\paragraph{選択肢の検討（査読待ち初稿）}\mbox{}\par

\noindent\textbf{1：適切（条件付き概数）} 水蒸気の潜熱は低緯度から中緯度の熱輸送に重要。。蒸発時に地表からエネルギーを受け取り移流後の凝結で放出する。\par

\noindent\textbf{2：適切（条件付き概数）} 水蒸気は現在の自然温室効果へ最大に寄与。。自然温室効果への寄与と、人為的な追加放射強制力を区別する。CO2の温暖化上の重要性を否定しない。\par

\noindent\textbf{3：適切（条件付き概数）} 陸は降水超過、海は蒸発超過。。全球長期平均では海から陸への大気輸送と陸から海への流出が釣り合う。\par

\noindent\textbf{4：不適切} 氷床は貯留量も輸送量も大きい。。氷床は大きい貯留だが交換に長い時間を要する。大きい貯留から大きい年間流量を結論できない。\par

\noindent\textbf{5：適切（条件付き概数）} 海洋・水蒸気・地下水・土壌水では水蒸気の滞留が最短。。大気は貯留が小さく出入りが大きいため約10日。地下水等には広い時間分布があり代表平均として比較する。\par

\paragraph{再利用する考え方} 全球平均・地域差、貯留量・流量、自然温室効果・追加強制力、代表値・普遍則を区別する。水の貯留・交換とエネルギー輸送の対象と成立条件を確認する。

\paragraph{共通理論} \hyperref[theory:atmosphere:water]{第\ref*{theory:atmosphere:water}節}, \hyperref[theory:atmosphere:greenhouse]{第\ref*{theory:atmosphere:greenhouse}節}, \hyperref[theory:atmosphere:transport]{第\ref*{theory:atmosphere:transport}節}

\par 参考：\cite{atm_usgs_water,atm_usgs_cycle,atm_jma_ghg2012}。

\paragraph{関連問題} \hyperref[q:r01-reg-applied-21]{R01-REG-APPLIED-21}

\par 難易度 2（校準初稿）、優先度未評価。

\subsection{令和元年度 問21}
\label{q:r01-reg-applied-21}

\noindent\textbf{原典：} sources/R01/17\_applied\_science.pdf、PDF物理7頁。

\textbf{公式正答：} 5（single）。

\paragraph{問題の要旨} 潜熱・温室効果、海陸の水収支、氷床の貯留と輸送、平均滞留時間を判定する。

\Needspace{4\baselineskip}\paragraph{確認した条件・記述（独自要約）}\mbox{}\par

\noindent\textbf{方向：} 最も不適切。\par

\paragraph{問われる核心} 水の貯留・交換とエネルギー輸送

\paragraph{必要知識} 水の貯留・交換とエネルギー輸送。交換時間は貯留量/流量である。氷床は大量の水を保持するが長い時間をかけて交換するので、貯留と輸送がともに大きいという肢が不適切。水蒸気の約10日の滞留、潜熱輸送、自然温室効果を同じ分子の別の役割として整理する。水収支の海陸差は地域ごとの恒常的な降水保証ではない。

\paragraph{考え方・正答を選ぶ理由（査読待ち初稿）} 交換時間は貯留量/流量である。氷床は大量の水を保持するが長い時間をかけて交換するので、貯留と輸送がともに大きいという肢が不適切。水蒸気の約10日の滞留、潜熱輸送、自然温室効果を同じ分子の別の役割として整理する。水収支の海陸差は地域ごとの恒常的な降水保証ではない。

\Needspace{4\baselineskip}\paragraph{選択肢の検討（査読待ち初稿）}\mbox{}\par

\noindent\textbf{1：適切（条件付き概数）} 水蒸気は現在の自然温室効果へ最大に寄与。。自然温室効果への寄与と、人為的な追加放射強制力を区別する。CO2の温暖化上の重要性を否定しない。\par

\noindent\textbf{2：適切（条件付き概数）} 水蒸気の潜熱は低緯度から中緯度の熱輸送に重要。。蒸発時に地表からエネルギーを受け取り移流後の凝結で放出する。\par

\noindent\textbf{3：適切（条件付き概数）} 陸は降水超過、海は蒸発超過。。全球長期平均では海から陸への大気輸送と陸から海への流出が釣り合う。\par

\noindent\textbf{4：適切（条件付き概数）} 海洋・水蒸気・地下水・土壌水では水蒸気の滞留が最短。。大気は貯留が小さく出入りが大きいため約10日。地下水等には広い時間分布があり代表平均として比較する。\par

\noindent\textbf{5：不適切} 氷床は貯留量も輸送量も大きい。。氷床は大きい貯留だが交換に長い時間を要する。大きい貯留から大きい年間流量を結論できない。\par

\paragraph{再利用する考え方} 全球平均・地域差、貯留量・流量、自然温室効果・追加強制力、代表値・普遍則を区別する。水の貯留・交換とエネルギー輸送の対象と成立条件を確認する。

\paragraph{共通理論} \hyperref[theory:atmosphere:water]{第\ref*{theory:atmosphere:water}節}, \hyperref[theory:atmosphere:greenhouse]{第\ref*{theory:atmosphere:greenhouse}節}, \hyperref[theory:atmosphere:transport]{第\ref*{theory:atmosphere:transport}節}

\par 参考：\cite{atm_usgs_water,atm_usgs_cycle,atm_jma_ghg2012}。

\paragraph{関連問題} \hyperref[q:h29-reg-applied-19]{H29-REG-APPLIED-19}

\par 難易度 2（校準初稿）、優先度未評価。

\subsection{令和2年度 問22}
\label{q:r02-reg-applied-22}

\noindent\textbf{原典：} sources/R02/17\_applied\_science.pdf、PDF物理8頁。

\textbf{公式正答：} 1（single）。

\par\textbf{検討未確定：} 公式1を保持。肢4は全球更新と地域枯渇を区別する適用範囲疑義、肢5は日本の流出/蒸発散の比較分母・期間が未指定で定量的な一次統計未確認。

\paragraph{問題の要旨} 蒸発潜熱のエネルギー割合、温室効果、海陸輸送、淡水更新と日本の水収支を比較する。

\Needspace{4\baselineskip}\paragraph{確認した条件・記述（独自要約）}\mbox{}\par

\noindent\textbf{方向：} 最も不適切。\par

\paragraph{問われる核心} 水の再生流量と熱収支の分母

\paragraph{必要知識} 水の再生流量と熱収支の分母。NASA代表収支の潜熱78W/m\textsuperscript{2}と入射342W/m\textsuperscript{2}の比は約23\%である。入射の約1割とする1が不適切で独立解答・公式とも1。ただし4の淡水が枯渇しないという表現は、全球循環の更新と地域資源の不足を区別する必要がある。5も全球又は他国との比較範囲が明示されないので、定量的に全国一律の低蒸発散とは教えない。

\paragraph{考え方・正答を選ぶ理由（査読待ち初稿）} NASA代表収支の潜熱78W/m\textsuperscript{2}と入射342W/m\textsuperscript{2}の比は約23\%である。入射の約1割とする1が不適切で独立解答・公式とも1。ただし4の淡水が枯渇しないという表現は、全球循環の更新と地域資源の不足を区別する必要がある。5も全球又は他国との比較範囲が明示されないので、定量的に全国一律の低蒸発散とは教えない。

\Needspace{4\baselineskip}\paragraph{選択肢の検討（査読待ち初稿）}\mbox{}\par

\noindent\textbf{1：不適切} 入射太陽エネルギーの約1割が蒸発潜熱になる。。代表収支の78/342≈0.23で約1/4、約1割ではない。地表正味日射を分母にすればさらに比が変わる。\par

\noindent\textbf{2：適切} 水蒸気は赤外吸収で温暖な気候へ寄与。。水蒸気は自然温室効果を支える。人為起源の追加強制力と同一のランキングへしない。\par

\noindent\textbf{3：適切（全球長期）} 陸の降水超過は海洋からの水蒸気輸送による。。長期の海陸区分の全球平均収支と整合する。地域や時期の全てで降水超過を保証しない。\par

\noindent\textbf{4：適用範囲に疑義} 淡水資源が枯渇しないのは大気滞留が短く豊富な降水があるため。。全球的に水が再生されることは説明できるが、地域の渇水・地下水過剰取水等の枯渇がないという普遍的保証にはならない。\par

\noindent\textbf{5：定性傾向・定量範囲要確認} 日本は流出が大きく蒸発散が小さい。。日本の多雨と流出優勢という代表的説明。年・流域による変動があり、原典の比較分母や対象期間は未指定である。\par

\paragraph{再利用する考え方} 全球平均・地域差、貯留量・流量、自然温室効果・追加強制力、代表値・普遍則を区別する。水の再生流量と熱収支の分母の対象と成立条件を確認する。

\paragraph{共通理論} \hyperref[theory:atmosphere:water]{第\ref*{theory:atmosphere:water}節}, \hyperref[theory:atmosphere:radiation]{第\ref*{theory:atmosphere:radiation}節}, \hyperref[theory:atmosphere:transport]{第\ref*{theory:atmosphere:transport}節}

\par 参考：\cite{atm_nasa_budget,atm_usgs_cycle}。

\par 難易度 2（校準初稿）、優先度未評価。


\chapter{海洋}
\section{海洋を水・運動量・物質の輸送系として理解する}
\label{theory:ocean:overview}
海洋は塩水の巨大な貯蔵庫であると同時に、風と密度差により運動し、熱と炭素と栄養塩を運ぶ系である。海水の組成、循環、生物による物質移動、海底での堆積を一つの収支としてつなぐ。津波では同じ海水でも長波の動的応答が中心となる。

\subsection{水の割合は何を分母にするか}
\label{theory:ocean:reservoir}
海水の体積は地球に存在する水の総体積の代表的推計で約96.5\%を占める。地球表面のうち海が占める面積の割合とは異なる。水量を時間不変の厳密値とせず、氷期の海水準変動なども含む貯蔵量と理解する。\cite{usgs_ocean_water_volume}

\subsection{35‰と塩類組成の85\%}
\label{theory:ocean:salinity}
原典の外洋の平均塩分は35‰である。質量比として\(35\,\mathrm g/1000\,\mathrm g=0.035\)なので約3.5\%となり、35\%ではない。海水全体に占める塩類と、その塩類中に占めるCl\(^{-}\)・Na\(^{+}\)の約85\%は異なる分母である。塩分の地域差とこの成分比を混同しない。\cite{noaa_sea_salt}
保存的なイオンの質量\(m_i,m_j\)を変えず水だけを蒸発又は加える理想モデルでは、濃度\(c_i=m_i/M\)の分母\(M\)が共通に変わり、比\(c_i/c_j=m_i/m_j\)は保たれる。これは保存性を前提とした収支である。反応、河川由来の異なる組成、熱水、微量の栄養塩には同じ比を保証しない。実用塩分という観測尺度と質量による絶対塩分の厳密な差まで、原典の概数に読み込まない。

\subsection{表層の風成循環と地衡流}
\label{theory:ocean:gyres}
風応力は海へ運動量を渡し、コリオリ力と圧力勾配が大規模な流れの方向に関わる。密度一定、水平流速\((u,v)\)、鉛直上向き、コリオリ係数\(f=2\Omega\sin\phi\)とする地衡流近似は
\[
 -fv=-\rho^{-1}\partial_xp,\qquad
 fu=-\rho^{-1}\partial_yp
\]
である。海面が高い領域の外向き圧力勾配に釣り合う循環は、\(f>0\)の北半球では時計回り、\(f<0\)の南半球では反時計回りとなる。これを亜熱帯環流の典型像へ対応させる。亜寒帯環流の海面の高低と回転方向は別であり、南半球の全ての渦を同じ向きとしない。\cite{noaa_subtropical_ecco}
赤道付近で\(f\)が小さいとこの単純な地衡流式だけでは応答を表せない。表層流が単に風向きと一致するわけでもない。

\subsection{密度差と深層循環の時間尺度}
\label{theory:ocean:overturning}
冷却と塩分増加は一般に海水の密度を増し、高密度水の沈降に関わる。北大西洋深層水や南極底層水など複数の水塊の形成、混合、湧昇を通じて海洋の深部と表層はつながる。千年程度という模式的な循環時間を、一つの流体粒子の一意の閉軌道や全水塊の絶対的最大年齢としない。「最大2000年」という原典は定義の追加確認が必要である。\cite{noaa_overturning}

\subsection{栄養塩：表層で消費され深部で再生される}
\label{theory:ocean:nutrients}
光の届く表層では植物プランクトンがリンやケイ素などを取り込む。沈降する有機物や殻が分解・溶解して水へ戻すと、深部の栄養塩が増える。典型的な表層低・深層高という向きは、この生物による移動の収支から理解できる。湧昇、季節変動、水塊の年齢、種類ごとの溶解を含む実分布は単調ではない。\cite{noaa_nutrient_observations}
濃度を\(C\)、流速を\(\boldsymbol u\)、拡散係数を\(D\)、再生と消費を\(R,U\)とする概念モデルは
\[
 \partial_t C+\boldsymbol u\cdot\nabla C
 =\nabla\cdot(D\nabla C)+R-U
\]
である。深さの関数だけでなく、移流と生物作用を含む点が保存的な主要塩類との違いである。

\subsection{炭酸塩の飽和・溶解速度・補償深度}
\label{theory:ocean:carbonate}
炭酸塩の飽和度は、所定の温度・圧力・塩分での溶解度積に対するイオン積の比で表せる。熱力学的に溶けやすいことと、一定時間にどれだけ溶けるかは別である。実際の速度は飽和度、粒子、境界層と海域に関わり、「深海ほど必ず速い」とする単調則は採らない。\cite{feely_carbonate2004}
炭酸塩補償深度は、海底への炭酸塩の供給と溶解の釣合いで残存が変わる深度である。したがって「深海」という名称だけでは石灰質軟泥の有無を確定できず、補償深度より上の深海底では残り得る。飽和深度、溶解が顕著になる深度、補償深度を同一の面としない。\cite{noaa_ccd_definition,feely_carbonate2004}

\subsection{団塊とクラストを区別する}
\label{theory:ocean:nodules}
マンガン団塊は深海底の核の周囲にゆっくり成長する鉄・マンガン酸化物の集合である。海山表面のマンガンクラストとは形態と位置が違う。成長が遅いため、速い堆積で埋まらず海底面に保たれる環境が形成・残存に関わるというモデルで理解する。堆積の全域に一律の速度を与えない。\cite{jamstec_deep_resources}

\subsection{CO2の吸収と海水の酸性化}
\label{theory:ocean:carbon}
溶けたCO2は炭酸水素イオンなどへ分配される。簡略化した反応は
\[
 \mathrm{CO_2(aq)+H_2O\rightleftharpoons H^++HCO_3^-},\qquad
 \mathrm{HCO_3^-\rightleftharpoons H^++CO_3^{2-}}
\]
である。CO2の増加でpHが低下する酸性化は、海水が直ちに酸性のpH7未満になることを意味しない。溶存無機炭素、pH、全アルカリ度は異なる変数である。海洋への人為的炭素の蓄積は大気の増加を和らげる一方、海洋化学の状態を変える。全球の正味吸収と局所的な放出域を区別する。\cite{noaa_carbon_chemistry}

\subsection{潮汐は場所ごとの応答を含む}
\label{theory:ocean:tides}
月と太陽の引力の空間差が海洋を周期的に駆動する。月の半日周の主要成分は約12時間25分、太陽の半日周成分は約12時間である。1日に満潮と干潮が二つずつという半日周潮に加え、日周潮と混合潮もある。原典の「通常2回」は代表像であり、すべての海岸の厳密な頻度ではない。地形、海底、水の慣性と回転による局所応答を含めて潮位を測る。\cite{noaa_tidal_patterns}

\subsection{長波として津波速度を導く}
\label{theory:ocean:tsunami}
水深\(h\)、小さい海面変位\(\eta\)、水深平均水平流速\(u\)を考える。非粘性、一定水深、波長が水深より長く、振幅が水深より小さい一方向の線形浅水モデルは
\[
 \partial_t\eta+h\partial_xu=0,\qquad
 \partial_tu+g\partial_x\eta=0.
\]
二式から\(\partial_t^2\eta-gh\partial_x^2\eta=0\)となり、波速は\(c=\sqrt{gh}\)である。水深の二乗に比例しない。例えば\(h=4000\)mなら\(c\simeq198\)m/s、約713 km/hである。これは一様水深のモデル計算であり実経路の予報値ではない。\cite{jma_tsunami_prediction}

\subsection{浅水変形・湾の集束・河川遡上}
\label{theory:ocean:shoaling}
反射・砕波・散逸が小さく、ゆっくり変化する水深、一定波面幅を仮定する。振幅\(a\)の長波のエネルギー密度は\(a^2\)に比例するため、流束保存\(a^2c\simeq\text{一定}\)と\(c\propto h^{1/2}\)から\(a\propto h^{-1/4}\)を得る。これがグリーンの法則の単純な形である。湾の幅、反射と共振、海岸の砕波があれば別の効果が重なる。\cite{jma_tsunami_prediction}
V字湾の奥では波が集まって増幅し得る。川へ入った津波も陸の内部まで進むことがある。海岸からの距離だけで判定せず、河道と地形を考える。\cite{jma_tsunami_generation,jma_tsunami_rivers}

\subsection{過去問に戻るときの確認項目}
塩分は単位と分母、環流は半球と種類、栄養塩は取込みと再生、堆積物は供給と溶解、津波は速度と高さを先に分ける。原典の「最大」「どこも」「深海ほど」を一般傾向から全称命題へ変えない。頻出度と学習優先度の最終集計は全問の確認後に行う。

\section{過去問とその判別点}
\subsection{平成30年度 問19}
\label{q:h30-reg-applied-19}

\noindent\textbf{原典：} sources/H30/17\_applied\_science.pdf、PDF物理6頁。

\textbf{公式正答：} 4（single）。

\par\textit{原典本文の校正と独自解説は未作成。分類のみの登録項目。}

\subsection{令和元年度 問18}
\label{q:r01-reg-applied-18}

\noindent\textbf{原典：} sources/R01/17\_applied\_science.pdf、PDF物理6頁。

\textbf{公式正答：} 1（single）。

\par\textit{原典本文の校正と独自解説は未作成。分類のみの登録項目。}

\subsection{令和元年度再試験 問19}
\label{q:r01-ret-applied-19}

\noindent\textbf{原典：} sources/R01\_retest/17\_applied\_science.pdf、PDF物理7頁。

\textbf{公式正答：} 1（single）。

\par\textit{原典本文の校正と独自解説は未作成。分類のみの登録項目。}

\subsection{令和2年度 問23}
\label{q:r02-reg-applied-23}

\noindent\textbf{原典：} sources/R02/17\_applied\_science.pdf、PDF物理9頁。

\textbf{公式正答：} 4（single）。

\par\textit{原典本文の校正と独自解説は未作成。分類のみの登録項目。}

\subsection{令和4年度 問21}
\label{q:r04-reg-applied-21}

\noindent\textbf{原典：} sources/R04/17\_applied\_science.pdf、PDF物理7頁。

\textbf{公式正答：} 1（single）。

\par\textit{原典本文の校正と独自解説は未作成。分類のみの登録項目。}

\subsection{令和5年度 問19}
\label{q:r05-reg-applied-19}

\noindent\textbf{原典：} sources/R05/17\_applied\_science.pdf、PDF物理7頁。

\textbf{公式正答：} 2（single）。

\par\textit{原典本文の校正と独自解説は未作成。分類のみの登録項目。}


\chapter{地球化学}
\section{背景：水を介して岩石・大気・生物を結ぶ}
\label{theory:geochemistry:boundary}
地球化学は、元素がどこに存在し、どの形態で移り、何が濃度を変えるかを扱う。
物理専攻者には、保存則と輸送方程式が入口になる。ただし元素の総量を保存しても、
溶存イオン、吸着、沈殿、気体、生体の間で配分が変わる。濃度の変化を水量の変化だけで説明できない。
降水が土壌を通り、岩石と反応し、地下水となり、河川から海へ運ぶ流れに沿って、
酸塩基・電子授受・溶解平衡・反応速度を結び付ける。

この初稿は地球化学の登録22問を原典19頁画像で確認した範囲である。
各問の公式正答と独立解答、各肢の成立条件を分け、原典の不十分な名称や強い一般化を改変せず記録する。
同位体の詳細は同位体地球化学章、地下水の水理は地下水章にも接続する。

\section{貯留・流量・濃度：何を比較しているか}
\label{theory:geochemistry:reservoir}
ある区画の水量を \(V\)、溶存成分濃度を \(C\) とすれば質量は \(M=VC\) である。
水の流量 \(Q\) は体積/時間、溶質の移動率 \(QC\) は質量/時間で、貯留量とは異なる。
区画を十分混合と近似すると
\[
 \frac{d(VC)}{dt}=\sum_j Q_jC_j-Q_{\rm out}C+P-L
\]
である。\(P,L\) は反応・供給・除去の質量/時間である。
蒸発は水を除いて塩を残すため、成分を追加しなくても \(C\) が上がる。
沈殿や吸着は溶存量を下げるが、系内の元素全量を必ず減らすわけではない。

USGSの代表貯留表は総水量約 \(1.386\times10^9\,\mathrm{km^3}\)、海約96.54\%、
河川約2120、大気約12900、全地下水約2340万 \(\mathrm{km^3}\) とする\cite{gc_water}。
大きな年間河川流出量を、河川内の瞬時貯留へ代入しない。
地下水は提示された湖・河・湿地・土壌より大きいが、取り出せる淡水資源の量と全地下水量は別である。
\begin{tcolorbox}[title=「約70\%」の分母]
氷河・雪氷は淡水貯留の約68.7\%である。海洋以外の全水に塩性地下水まで含めると、
同じ約70\%にはならない。平成26年度問25と令和元年度再試験問18は「陸水」という語を使うため、
教材では原語を残して淡水を分母とした意図と範囲疑義を併記する。
\end{tcolorbox}
平均滞留時間 \(V/Q\) は定常・同じ区画の出入りを近似する指標であり、
個々の水分子の年齢が全てその値になるという意味ではない。

\section{汚濁指標と起源を追うトレーサー}
\label{theory:geochemistry:indicators}
BODは生物分解に伴う酸素要求量で、分解可能な有機負荷を追う。
SSは浮遊粒子、全窒素は窒素成分の合計、大腸菌群数は出題当時の衛生汚濁の指標である。
これらは同じ「水質」でも対象が違う。BODを全有機炭素量に、SSを溶存量に置き換えない。
測定対象・方法・単位を確かめる必要がある\cite{gc_hem}。
水素同位体比は蒸発、降水、混合等で変わり、水の起源や移動を追う。
比が高いだけで汚濁量が多いとは判断できない。
汚染研究に同位体を使うことと、同位体比が一般的な汚濁負荷の直接指標であることは別である。

\section{河川水は海水を薄めたものだけではない}
\label{theory:geochemistry:river}
雨に取り込まれた海塩、土壌の反応、岩石風化、温泉・鉱山・人為排水、蒸発、海水混合が組成を変える。
世界平均河川の主要イオンにはCaと炭酸水素が重要で、NaClが世界の大半を支配するという一般化は誤る。
沿岸の降水に海塩が多いことを、全河川の溶存成分の起源へ拡張しない。

北野・加藤の1970年論文の比較は、小林1960とLivingstone1963の集計を用いる。
表4では日本の蒸発残渣74.8mg/L、溶存SiO2の割合26.84\%、世界100mg/L、11.67\%である。
割合を濃度へ戻すと
\[
 C_{\mathrm{SiO_2,Japan}}=74.8\times0.2684\simeq20.1\,\mathrm{mg/L},\qquad
 C_{\mathrm{SiO_2,world}}=100\times0.1167=11.67\,\mathrm{mg/L}
\]
となる。同論文は多雨・蒸発濃縮の小ささと火山地質を説明する\cite{gc_river1970}。
古い試料群の一般傾向を、現在の全流域が必ず同じ順位であるという命題にしない。
Si表示とSiO2換算表示も分ける。

\section{炭酸平衡：石灰岩を溶かす条件と析出させる条件}
\label{theory:geochemistry:carbonate}
土壌呼吸で供給されたCO2が水に溶けると炭酸系を作る。
\[
 \mathrm{CO_2(g)\rightleftharpoons CO_2(aq)},\qquad
 \mathrm{CO_2(aq)+H_2O\rightleftharpoons H^++HCO_3^-},\qquad
 \mathrm{HCO_3^-\rightleftharpoons H^++CO_3^{2-}} .
\]
第一式の気液分配と後二式の酸解離は異なる平衡である。
ここで第二式のCO2には水和した炭酸を合わせた慣用的な表現を用いた。
各平衡定数は温度・活量基準に依存する。
石灰岩の代表鉱物である方解石について
\[
 \mathrm{CaCO_3+CO_2+H_2O\rightleftharpoons Ca^{2+}+2HCO_3^-}
\]
が、Caと炭酸水素を増やす反応を示す。
CO2を多く含む水は溶解側へ、洞窟でCO2を失う水は析出側へ動き、
鍾乳石や石筍を形成する\cite{gc_hem,gc_cave}。
CO2添加を普通の洞窟での析出原因とする令和6年度問22の肢2は方向が逆である。
\begin{tcolorbox}[title=原典の水質名を科学的な文へ書き換えない]
平成25年度問31と令和元年度問20は「炭酸カルシウム型」と呼ぶ。
通常、主要溶存種のCaとHCO3に基づく名称は炭酸水素カルシウム型、すなわちCa-HCO3型である。
固体CaCO3と、Ca・HCO3が別々に溶存する水を混同しない\cite{gc_hem}。
原典はそのまま保存し、名称の疑義を問題ごとに保持する。
\end{tcolorbox}

\section{採水後に変わるpH：測った時点も条件である}
\label{theory:geochemistry:sampling}
地下水が大気より高いCO2分圧の土壌と接した場合、採取・攪拌で脱気するとpHが上昇し得る。
したがって現場の非攪拌水と室内放置水は同じ組成の測定とは限らない。
ただし酸化反応、沈殿、空気からのCO2吸収等もあり、全試料で必ず上がるという規則ではない。
温度、曝気、保存時間を記録して現場条件を優先する\cite{gc_hem}。
pHは水素イオン活量の負の常用対数であり、濃度との同一視は希薄な近似である。
炭酸系の平衡は酸化還元を伴わなくてもpHを変える。

\section{活量：濃度だけで平衡を記述できない理由}
\label{theory:geochemistry:activity}
理想溶液では化学ポテンシャルの組成依存を濃度比で表せるが、イオンは互いの電場に影響される。
規定した標準状態に対し
\[
 \mu_i=\mu_i^\circ+RT\ln a_i,\qquad a_i=\gamma_i\frac{c_i}{c^\circ},\qquad
 I=\frac12\sum_i c_i z_i^2
\]
と置く。\(a_i,\gamma_i\) は無次元、\(z_i\) は価数で、イオン強度 \(I\) の濃度基準を揃える。
二価イオンは同じ濃度でも \(z_i^2=4\) の重みを持つ。
Debye--Huckelの極限則は希薄溶液の静電効果から
\[
 \log_{10}\gamma_i\simeq-Az_i^2\sqrt I
\]
という形を与える。\(A\) の値と単位は温度・溶媒・濃度基準に依存する。
濃厚塩水ではこの極限形を外挿せず、より適切な活量モデルを使う。
単一イオン活量の分離には慣例が必要である\cite{gc_debye}。
これは鉱物の風化抵抗順位そのもののモデルではない。

\section{平衡でどこまで溶けるかと、どれほど速く溶けるか}
\label{theory:geochemistry:equilibrium}
固定した温度・圧力・化学条件で、残存する固相と水の化学ポテンシャルが釣り合うと溶解平衡となる。
方解石ならイオン活動度積 \(a_{\mathrm{Ca^{2+}}}a_{\mathrm{CO_3^{2-}}}\) を溶解度積と比較する。
溶存総炭素とCO3の活量はpHによる種分布のため同じではない。
少量の固相が全部溶けた場合、観測濃度が飽和濃度とは限らない。
開放流動、別相生成、過飽和の維持もあり、「時間がたてば必ず飽和平衡」とはしない。
平衡は到達先を示すが、表面積、温度、反応経路、飽和からの距離が速度を変える\cite{gc_hem}。
令和4年度問20では、活量の理論・平衡・経験的抵抗性・速度の対象を区別する。

\section{地殻の存在度：対象と質量基準を固定する}
\label{theory:geochemistry:crust}
地殻は酸化物・ケイ酸塩を多く含み、質量基準ではO、Si、Alが上位である。
USGSの代表表はO46.60、Si27.72、Al8.13、K2.59、Ti0.44、Mn0.100、Cl0.020\%を示す\cite{gc_crust}。
平成25年度問10のCl>Kは逆である。
各小数は地殻モデルや試料集合に依存し、上部大陸地殻、全地殻、地球全体、海水では組成が違う。
化学式中の原子比と岩石中の質量存在度も別の比較である。

\section{酸化還元：電子供与体と受容体を先に置く}
\label{theory:geochemistry:redox}
電子を与える物質は自身が酸化されるので還元剤、受け取る物質は自身が還元されるので酸化剤である。
酸素供給の乏しい地下水では、利用可能な反応物と微生物活動に応じ、
酸素、硝酸、Mn・Fe酸化物、硫酸、CO2等が電子受容体として使われる\cite{gc_redox,gc_redox2009}。
これは各反応が常にきれいな順で完了するという保証ではない。
Fe(III)の固相が還元されるとFe(II)が溶存し得るため、還元地下水はFeMnに富みやすい。
湧出して酸素を得るとFe(III)水酸化物等が茶褐色に析出する。
沈殿の酸化状態を、その水が地下で持っていた状態と混同しない。

基準電極に対して測る電位を揃える。標準水素電極を0とする \(E_h\) と、
測定器の参照電極に対する表示は換算が必要である。
酸素の還元半反応
\[
 \mathrm{O_2+4H^++4e^-\rightleftharpoons2H_2O}
\]
にNernst式を用い、水活量1、気体の分圧を標準圧で無次元化すると、25℃では
\[
 E_h\simeq1.229-0.05916\,\mathrm{pH}
       +0.01479\log_{10}p_{\mathrm{O_2}}\quad(\mathrm V)
\]
である。pH7、空気 \(p_{\mathrm{O_2}}=0.21\) で約+0.805Vとなり、-0.250Vとは符号が異なる。
これは理想的な酸素/水の境界であり、白金電極の実測値が必ず一致するという説明ではない。
複数反応の混合電位、反応の遅さ、電極応答がある\cite{gc_hem}。
地下の長い滞留でも、酸化可能な物質との反応なしに酸素が消えるわけではない。
火山温泉の還元性と、湧出後の空気混合も区別する。

\section{硫化物の酸化と酸性鉱山排水}
\label{theory:geochemistry:mine}
代表例として黄鉄鉱の酸化は
\[
 \mathrm{FeS_2+\tfrac72O_2+H_2O\longrightarrow Fe^{2+}+2SO_4^{2-}+2H^+}
\]
と書ける。硫酸イオンと酸を生じ、さらに鉄の酸化・加水分解等が関わる。
露出面積、酸素供給、微生物、炭酸塩等の中和能が結果を変える。
硫化物があるだけで全ての排水が同じpHになるとは限らず、
中性や塩性で溶存金属の多い排水もある\cite{gc_mine,gc_mine2015}。
Ca-SO4型を生じるには、硫酸だけでなくCaの供給も必要である。

\section{化学的風化：一次鉱物から二次鉱物へ}
\label{theory:geochemistry:weathering}
地表の水・温度・酸素条件は、深部やマグマから形成された鉱物の条件と違う。
通常の造岩鉱物では、かんらん石、輝石、黒雲母、石英の順に化学的風化への抵抗が増す。
これは環境と鉱物の反応をまとめた経験系列で、物理的破砕の順位や全試料の反応速度を保証しない。
方解石は炭酸を含む水で比較的速く溶けるが、飽和に近い水では速度が低下する。

塩基性陽イオンと一部のSiが溶脱すると、残るAl・Feが酸化物・オキシ水酸化物や粘土を作る。
強い溶脱が進む湿潤熱帯ではカオリナイトが生じやすい。
スメクタイト等との違いは層状構造・組成と溶液条件の違いである。
岩種、水岩比、温度、生物、期間が重要で、気候名だけで粘土を決めない\cite{gc_clays,gc_hem}。
酸化的な通常土壌でのAl・Fe(III)沈殿と、酸性水・錯体・還元Fe(II)の溶存を区別する。

\section{炭素循環：速い生物過程と遅い岩石過程}
\label{theory:geochemistry:carbon}
光合成でCO2を有機炭素に取り込み、呼吸・分解で戻す速い循環と、
風化・炭酸塩堆積・有機物埋没・変成・火山による遅い循環を分ける\cite{gc_carbon}。
貯留が大きいことは、現在の吸収速度も大きいこととは別である。
海洋の溶存無機炭素等の貯留は大気よりはるかに多く、半分や同程度という関係ではない。

単純化したCaケイ酸塩の風化は
\[
 \mathrm{CaSiO_3+2CO_2+H_2O\longrightarrow Ca^{2+}+2HCO_3^-+SiO_2}
\]
である。海でCaCO3が沈殿してCO2の一部が戻る反応を加えると、
\[
 \mathrm{CaSiO_3+CO_2\longrightarrow CaCO_3+SiO_2}
\]
という正味の固定を表す。岩石を単一鉱物とみなした収支例で、実際の速度をこの式だけで求めない。
土壌の高いCO2は炭酸を供給し風化を促進し得る。
初期のCO2濃度の低下に風化と生物活動が関与したという地史的な説明は、
今日の大気が既に完全な定常状態にあるという主張ではない。

自然の温室効果に大きく寄与する水蒸気と、人為的な追加放射強制力のCO2を区別する。
CO2の主要な温室作用は赤外の吸収・放射であり、紫外の吸収との取り違えを避ける。
メタン生成は通常嫌気的な代謝で、光合成ではない。例として
\[
 \mathrm{CO_2+4H_2\longrightarrow CH_4+2H_2O}
\]
は還元条件での基質と電子収支を示す。原典の「メタン生成菌」という呼び方を、
生物の分類や全生成経路の説明にそのまま拡張しない。
\begin{tcolorbox}[title=炭素の収支を閉じるために]
平成23年度問19と平成28年度問27の正答肢は、
風化の除去作用が光合成・分解・呼吸で長期に均衡すると書く。
重要な過程を挙げる意図は理解できるが、遅い循環の火山・堆積・埋没等を省く。
教材では原語と公式を保存し、定常仮定の範囲を疑義として残す。
\end{tcolorbox}

\section{窒素循環：窒素の反応と生物の炭素源は別軸である}
\label{theory:geochemistry:nitrogen}
窒素固定はN2から反応性窒素を作り、硝化は還元態窒素を酸化し、脱窒は硝酸等を還元して気体へ戻す。
硝化の代表的な全反応は
\[
 \mathrm{NH_4^++2O_2\longrightarrow NO_3^-+2H^++H_2O}
\]
である。アンモニアから亜硝酸、さらに硝酸の段階を分ける。
固定ではまずN2の還元によりアンモニアを得てアミノ酸へ同化する。
雷や燃焼の熱的NOxは非生物固定だが、燃料中に既に存在する窒素の再排出と新規N2固定は違う。

有機物を電子供与体とする脱窒の収支例を \(\mathrm{CH_2O}\) で単純化すると
\[
 \mathrm{5CH_2O+4NO_3^-+4H^+\longrightarrow2N_2+5CO_2+7H_2O}
\]
となる。これを全ての脱窒生物の炭素源へ一般化しない。
無機硫黄を電子供与体とする例は
\[
 \mathrm{5S+6NO_3^-+2H_2O\longrightarrow3N_2+5SO_4^{2-}+4H^+}
\]
で、硫黄酸化と硝酸還元がエネルギーを供給する。
1989年の一次実験は、HCO3を無機炭素源、還元硫黄を電子供与体とする嫌気的な自養脱窒菌の集積を示す\cite{gc_sulfur_denit}。
電子供与体、受容体、同化する炭素源を別々に答えれば、硝化と脱窒の二語だけで栄養型を決める誤りを避けられる。
\begin{tcolorbox}[title=令和2年度問19の公式と科学的例外]
公式は肢5を不適切とする。独立初回も5であった。
ただし化学合成と脱窒を組み合わせて無機炭素を同化する過程は実在する。
硝化と従属栄養脱窒の対比という設問意図を説明し、肢5を全生物に一律で否定しない。
疑義と査読待ちを保持する。
\end{tcolorbox}

\section{溶質輸送：移流・拡散・分散を分ける}
\label{theory:geochemistry:transport}
同じ水とともに移動する移流、熱運動による分子拡散、流れの速度と経路の差による機械分散を分ける。
濃度 \(C\)、平均間隙流速 \(v\)、一定の拡散分散係数 \(D\) について、
一様な一次元、吸着と反応なしという模型では
\[
 J=vC-D\frac{\partial C}{\partial x},\qquad
 \frac{\partial C}{\partial t}+v\frac{\partial C}{\partial x}
 =D\frac{\partial^2 C}{\partial x^2}
\]
が流束と連続の式から得られる。
Fick則の比例する量は濃度勾配に対する流束であり、個々の分子の一定の流速ではない。
点状の初期分布を無限領域へ流せば中心は \(vt\)、分散は \(2Dt\) となり、移動と広がりを区別できる。

縦方向には \(D_L=\alpha_L|v|+D_m^\ast\)、横方向にも対応する式を使う。
\(\alpha_L,\alpha_T\) は長さ、\(D_m^\ast\) は媒体内の有効分子拡散係数で面積/時間である。
通常、横分散長は縦より小さい。野外では室内より大きい不均質性を含むため、
実験室の値をそのまま流域へ外挿できない。
59現場を比較した1992年研究は、観測尺度と分散長の傾向に加え、
同じ尺度でも大きな散らばりと信頼度差があるため普遍直線が不適切だとする\cite{gc_dispersion}。

\section{電気伝導率：イオンの量と動きやすさを測る}
\label{theory:geochemistry:conductivity}
水の電気伝導はイオンの移動による。希薄な独立寄与の近似なら
\[
 \kappa\simeq\sum_i\lambda_i c_i,\qquad
 \kappa=\frac1\rho=\frac{\ell}{RA}
\]
である。\(\lambda_i\) はイオン種に応じたモル伝導寄与、
\(\rho\) は抵抗率、\(R\) は長さ \(\ell\)・断面積 \(A\) の試料の電気抵抗である。
濃度の単位と \(\lambda_i\) の単位を合わせる。
1cm立方体では数値として \(1/R\) に対応するが、伝導率の単位はS/cmで、伝導度のSと異なる。
SIのS/mへの換算は \(1\,\mathrm{S/cm}=100\,\mathrm{S/m}\) である。

同じ天然水の通常温度域では温度上昇でイオン移動度が増し、伝導率は上がる。
比較には温度を揃え、25℃へ補正した値を用いる\cite{gc_hem,gc_conductance,gc_conductance1994}。
海水、河川水、雨の順位は塩類量の一般傾向で、例外はある。
同じ組成比なら濃度と蒸発残留物に関係するが、組成が違えば単一係数で総溶存量へ換算できない。

\section{水の同位体：正の相関と履歴の識別}
\label{theory:geochemistry:isotope}
標準に対する同位体比 \(R\) の偏差を
\[
 \delta=\left(\frac R{R_{\rm std}}-1\right)1000\quad(\text{千分率})
\]
で表す。重い同位体は蒸発・凝結の際に軽いものと同じ比率で移らない。
HとOの分別が連動するため、代表的な全球天水線は
\[
 \delta^2\mathrm H\simeq8\delta^{18}\mathrm O+10
\]
という正の相関である\cite{gc_meteoric}。
絶対同位体比が厳密な反比例であるという意味ではない。
局地天水線の傾き・切片は違い、蒸発や混合は天水線から逸脱させる。
逐次凝結で重い同位体が除かれると、高緯度へ向かう水蒸気と降水は重い同位体に乏しくなる一般傾向がある。
起源識別と単独の汚濁量評価を区別する。

\section{生物によるケイ酸の取り込みと、殻の溶解}
\label{theory:geochemistry:silica}
珪藻は溶存ケイ酸を取り込み、殻を形成する。
成長時の取り込みを、河川ケイ酸濃度を上げる直接作用として説明するのは逆である\cite{gc_hem,gc_diatom}。
しかし死後の殻の溶解はケイ酸を水へ戻す。
水中濃度の変化は、風化供給、取り込み、溶解、輸送、沈殿等の差で決まり、個々の過程を区別する。

海洋の粘土生成である逆風化は、アルカリ度とカチオンを消費してCO2を生じ、
pHを下げ得る。これは陸上の単純なケイ酸塩風化の酸消費とは違う。
2025年の原著研究はこの反応と生物起源シリカの役割を検討する\cite{gc_reverse}。
後年の科学説明であって、平成27年原典の意図を直接証明したものとはしない。
同問の「ケイ藻によってケイ素が放出されることがある」は生活史の範囲を省くため、
公式5を残しつつ初回4|5と条件疑義を保持する。

\section{自然ヒ素：化学形態・濃度・基準の時点}
\label{theory:geochemistry:arsenic}
ヒ素は鉱物中に存在し、酸化還元、溶解、吸着・脱着、二次鉱物生成により地下水等へ移り得る。
元素の全量、無機のIII・Vの形態、単体、溶存種、鉱物を分ける\cite{gc_hem,gc_redox}。
物理・化学的な取り込みの違いを、単体が一律に無害という説明にしない。
NIOSHの金属ヒ素を含む一次有害性資料は、原典の「毒性はない」という断言を支持しない\cite{gc_arsenic_hazard}。

\begin{tcolorbox}[title=平成30年度問28：単位の誤りと他肢の疑義を分ける]
原典肢2は0.01g/Lと書く。2004施行の水道基準表と現在の表はヒ素量0.01mg/Lである\cite{gc_arsenic2004,gc_arsenic_current}。
0.01g/Lは10mg/Lで1000倍違うため、公式2の根拠は確認できる。
しかし肢3の無毒という断言と肢4の岩種比較・揮発性の因果は別に検証を要する。
原典単位を正しいmg/Lへ校正して誤りを消したり、公式と一致しただけで全肢を確認済みにしたりしない。
\end{tcolorbox}
基準は用途・対象・施行時点に依存する。現在の表の項目数や他成分の変更を2018年へ遡及しない。

\section{CO2回収・貯留：分離と封じ込めの役割}
\label{theory:geochemistry:ccs}
CCSは排出等からCO2を回収し、移送して地中へ貯留する工程である。
アミン吸収や膜分離は回収、深部塩水帯水層等は貯留場所、
CO2-EORは利用と貯留へ接続する\cite{gc_ccs,gc_capture,gc_mea}。
帯水層という一候補だけで地下水流動の問題へ主分類を限定しない。
EORにはCO2以外の方式もあり、界面活性剤の油回収への利用が全てCCSそのものになるわけではない。
平成26年度問18は「関係性が最も薄い」という相対判断で、界面活性剤を無関係と断言する問いではない。
圧力、流動、鉱物反応、封じ込めと漏えいを別の評価課題として扱う。

\section{今回の出題形式と次の確認}
22問では、陸水の組成5問、炭素循環4問、水貯留3問が反復し、
指標、地殻存在度、CCS、窒素、輸送、伝導率、酸化還元、水鉱物、土壌風化、ヒ素が各1問である。
これは今回実際に読んだ問題の集計で、未処理の他分類や将来の出題確率を含む最終頻度ではない。
物理専攻者は、まず質量保存と流束、次に化学種を決める酸塩基・電子授受・活量、
最後に地質・生物・測定時点の条件を加えると、用語の暗記だけに依存せず判断できる。
正答位置が入れ替わった同型でも、方向・単位・分母・限定語を原典で再確認する。
疑義11問はAstra査読対象として保持し、以下の個別解説から条件を追跡できる。
\section{炭素貯留・生物循環と岩石風化}

\subsection{平成23年度 問19}
\label{q:h23-reg-applied-19}

\noindent\textbf{原典：} sources/H23/17\_applied\_science.pdf、PDF物理7頁。

\textbf{公式正答：} 3（single）。

\par\textbf{検討未確定：} 正答の風化・生物循環の重要性は支持されるが、「光合成、有機物分解、呼吸によって長期にバランス」という原典の収支範囲は不十分。火山・堆積・埋没も必要で、人為擾乱下の現在大気は定常でない。原典を修正せず部分校正。

\paragraph{問題の要旨} 炭素の貯留量、自然温室効果、生物過程と地史的風化を比較する。

\Needspace{4\baselineskip}\paragraph{確認した条件・記述（独自要約）}\mbox{}\par

\noindent\textbf{方向：} 正しいもの。\par

\noindent\textbf{対象：} 初期大気・現在の自然温室効果・炭素貯留と循環。\par

\paragraph{問われる核心} 炭素の速い生物循環と遅い岩石循環

\paragraph{必要知識} 炭素の速い生物循環と遅い岩石循環。設問が意図する3は風化と生物活動の重要性であり、初回・公式とも3である。他の候補は初期濃度、水蒸気の役割、メタン生成の酸化還元条件、海洋の大きな貯留量に反する。ただし風化の収支を生物3作用だけで閉じない。短い生物循環と火山・堆積を含む遅い循環を分け、現在の炭素増加を隠す定常仮定にも用いない。

\paragraph{考え方・正答を選ぶ理由（査読待ち初稿）} 設問が意図する3は風化と生物活動の重要性であり、初回・公式とも3である。他の候補は初期濃度、水蒸気の役割、メタン生成の酸化還元条件、海洋の大きな貯留量に反する。ただし風化の収支を生物3作用だけで閉じない。短い生物循環と火山・堆積を含む遅い循環を分け、現在の炭素増加を隠す定常仮定にも用いない。

\Needspace{4\baselineskip}\paragraph{選択肢の検討（査読待ち初稿）}\mbox{}\par

\noindent\textbf{1：不適切} 初期CO2は現在より低い。。初期大気のCO2は現在より高かったとする地史的な理解と逆である。初期という語の全時代の精密濃度を固定しない。\par

\noindent\textbf{2：不適切} 最大の温室効果はCO2、次いで水蒸気。。自然の温室効果では水蒸気の寄与が大きい。人為的な追加放射強制力のランキングを現在大気の全温室効果と混ぜている。\par

\noindent\textbf{3：意図上適切・収支範囲疑義} 風化によるCO2除去は重要で、光合成・分解・呼吸で長期に均衡。。風化による炭素固定と生物の吸収・呼吸・分解は重要である。ただし遅い循環の収支は火山・堆積・有機炭素埋没等も含み、3作用だけで任意時点の均衡が保証されるわけではない。\par

\noindent\textbf{4：不適切} 酸化環境でメタン生成。。メタン生成微生物は通常嫌気的な条件で活動する。酸化的な条件でのメタン酸化と生成を区別する。\par

\noindent\textbf{5：不適切} 大気炭素は海洋炭素の半分。。海洋の溶存無機炭素等は大気炭素よりはるかに大きい。大気が海洋の半分という量関係ではない。\par

\paragraph{再利用する考え方} 物質の貯留量・輸送速度・化学形態を分け、系の境界と反応の前提を確かめる。炭素の速い生物循環と遅い岩石循環を収支・平衡・反応速度のどのモデルで扱うか選ぶ。

\paragraph{共通理論} \hyperref[theory:geochemistry:carbon]{第\ref*{theory:geochemistry:carbon}節}, \hyperref[theory:geochemistry:weathering]{第\ref*{theory:geochemistry:weathering}節}

\par 参考：\cite{gc_carbon,atm_noaa_ghg}。

\paragraph{関連問題} \hyperref[q:h26-reg-applied-27]{H26-REG-APPLIED-27}, \hyperref[q:h28-reg-applied-27]{H28-REG-APPLIED-27}

\par 難易度 3（校準初稿）、優先度未評価。

\subsection{平成26年度 問27}
\label{q:h26-reg-applied-27}

\noindent\textbf{原典：} sources/H26/17\_applied\_science.pdf、PDF物理11頁。

\textbf{公式正答：} 4（single）。

\paragraph{問題の要旨} CO2吸収帯、風化、海洋貯留、地史的変化とメタン生成を比較する。

\Needspace{4\baselineskip}\paragraph{確認した条件・記述（独自要約）}\mbox{}\par

\noindent\textbf{方向：} 最も適切なもの。\par

\paragraph{問われる核心} 炭素循環の貯留・風化・微生物代謝

\paragraph{必要知識} 炭素循環の貯留・風化・微生物代謝。紫外を赤外、減速を促進、同程度を海洋の大きな貯留、光合成を嫌気代謝へそれぞれ区別する。地史的な風化と光合成の作用を述べる4が最も適切である。生物が吸ったCO2が全量恒久貯留するわけではない。

\paragraph{考え方・正答を選ぶ理由（査読待ち初稿）} 紫外を赤外、減速を促進、同程度を海洋の大きな貯留、光合成を嫌気代謝へそれぞれ区別する。地史的な風化と光合成の作用を述べる4が最も適切である。生物が吸ったCO2が全量恒久貯留するわけではない。

\Needspace{4\baselineskip}\paragraph{選択肢の検討（査読待ち初稿）}\mbox{}\par

\noindent\textbf{1：不適切} CO2は紫外に幅広い吸収帯。。主要な温室作用は赤外領域の吸収・放射である。紫外との取り違え。\par

\noindent\textbf{2：不適切} 土壌CO2が高いので風化のCO2除去が減速。。土壌呼吸によるCO2供給は炭酸を作り化学的風化を促す。実際の速度は水流、表面積、温度、飽和度等にも依存する。\par

\noindent\textbf{3：不適切} 大気炭素全量は海洋と同程度。。海洋の炭素貯留は大気よりはるかに多く、同程度ではない。\par

\noindent\textbf{4：適切} 初期CO2は高く、減少へ風化と光合成が作用。。地史的なCO2減少に風化と生物活動が関与する。光合成固定も呼吸で戻る部分と埋没する部分を区別する。\par

\noindent\textbf{5：不適切} 嫌気環境のメタン生成は光合成。。嫌気条件は整合するがメタン生成は光合成ではない。水素とCO2、酢酸等の基質を使う代謝である。\par

\paragraph{再利用する考え方} 物質の貯留量・輸送速度・化学形態を分け、系の境界と反応の前提を確かめる。炭素循環の貯留・風化・微生物代謝を収支・平衡・反応速度のどのモデルで扱うか選ぶ。

\paragraph{共通理論} \hyperref[theory:geochemistry:carbon]{第\ref*{theory:geochemistry:carbon}節}, \hyperref[theory:geochemistry:weathering]{第\ref*{theory:geochemistry:weathering}節}, \hyperref[theory:geochemistry:redox]{第\ref*{theory:geochemistry:redox}節}

\par 参考：\cite{gc_carbon,atm_noaa_ghg,gc_redox}。

\paragraph{関連問題} \hyperref[q:h23-reg-applied-19]{H23-REG-APPLIED-19}, \hyperref[q:r03-reg-applied-21]{R03-REG-APPLIED-21}

\par 難易度 3（校準初稿）、優先度未評価。

\subsection{平成28年度 問27}
\label{q:h28-reg-applied-27}

\noindent\textbf{原典：} sources/H28/17\_applied\_science.pdf、PDF物理11頁。

\textbf{公式正答：} 4（single）。

\par\textbf{検討未確定：} 正答の風化・生物循環の重要性は支持されるが、「光合成、有機物分解、呼吸によって長期にバランス」という原典の収支範囲は不十分。火山・堆積・埋没も必要で、人為擾乱下の現在大気は定常でない。原典を修正せず部分校正。

\paragraph{問題の要旨} 自然温室効果、初期大気、メタン条件、風化と炭素貯留を判定する。

\Needspace{4\baselineskip}\paragraph{確認した条件・記述（独自要約）}\mbox{}\par

\noindent\textbf{方向：} 最も適切なもの。\par

\paragraph{問われる核心} 速い生物循環と遅い炭素循環の収支範囲

\paragraph{必要知識} 速い生物循環と遅い炭素循環の収支範囲。初回・公式とも4であり、地史と生物過程を扱う意図で選ぶ。ただし遅い炭素循環は火山・堆積・埋没も含み、生物3作用だけで風化流束まで均衡するという説明にはしない。H23問19と同じ文言群だが正答位置と問いの方向を原典で再照合した。

\paragraph{考え方・正答を選ぶ理由（査読待ち初稿）} 初回・公式とも4であり、地史と生物過程を扱う意図で選ぶ。ただし遅い炭素循環は火山・堆積・埋没も含み、生物3作用だけで風化流束まで均衡するという説明にはしない。H23問19と同じ文言群だが正答位置と問いの方向を原典で再照合した。

\Needspace{4\baselineskip}\paragraph{選択肢の検討（査読待ち初稿）}\mbox{}\par

\noindent\textbf{1：不適切} 初期CO2は現在より低い。。初期大気のCO2は現在より高かったとする地史的な理解と逆である。初期という語の全時代の精密濃度を固定しない。\par

\noindent\textbf{2：不適切} 最大の温室効果はCO2で水蒸気が次。。自然の温室効果では水蒸気の寄与が大きい。人為的な追加放射強制力のランキングを現在大気の全温室効果と混ぜている。\par

\noindent\textbf{3：不適切} メタンは酸化環境で生成。。メタン生成微生物は通常嫌気的な条件で活動する。酸化的な条件でのメタン酸化と生成を区別する。\par

\noindent\textbf{4：意図上適切・収支範囲疑義} 風化によるCO2除去は重要で、光合成・分解・呼吸で長期に均衡。。風化による炭素固定と生物の吸収・呼吸・分解は重要である。ただし遅い循環の収支は火山・堆積・有機炭素埋没等も含み、3作用だけで任意時点の均衡が保証されるわけではない。\par

\noindent\textbf{5：不適切} 大気Cは海洋Cの半分。。海洋の溶存無機炭素等は大気炭素よりはるかに大きい。大気が海洋の半分という量関係ではない。\par

\paragraph{再利用する考え方} 物質の貯留量・輸送速度・化学形態を分け、系の境界と反応の前提を確かめる。速い生物循環と遅い炭素循環の収支範囲を収支・平衡・反応速度のどのモデルで扱うか選ぶ。

\paragraph{共通理論} \hyperref[theory:geochemistry:carbon]{第\ref*{theory:geochemistry:carbon}節}, \hyperref[theory:geochemistry:weathering]{第\ref*{theory:geochemistry:weathering}節}

\par 参考：\cite{gc_carbon,atm_noaa_ghg}。

\paragraph{関連問題} \hyperref[q:h23-reg-applied-19]{H23-REG-APPLIED-19}

\par 難易度 3（校準初稿）、優先度未評価。

\subsection{令和3年度 問21}
\label{q:r03-reg-applied-21}

\noindent\textbf{原典：} sources/R03/17\_applied\_science.pdf、PDF物理8頁。

\textbf{公式正答：} 2（single）。

\paragraph{問題の要旨} 炭素貯留、地史的風化、生物代謝とCO2の吸収帯を比較する。

\Needspace{4\baselineskip}\paragraph{確認した条件・記述（独自要約）}\mbox{}\par

\noindent\textbf{方向：} 最も適切。\par

\paragraph{問われる核心} 貯留・反応・放射を区別する炭素循環

\paragraph{必要知識} 貯留・反応・放射を区別する炭素循環。地史的な濃度減少への複数の関与を述べる2が適切。海洋と大気を同量にせず、メタン生成を光合成にせず、紫外と赤外を取り違えず、土壌CO2による炭酸供給の風化促進を理解する。風化速度・炭素埋没の時間尺度も意識する。

\paragraph{考え方・正答を選ぶ理由（査読待ち初稿）} 地史的な濃度減少への複数の関与を述べる2が適切。海洋と大気を同量にせず、メタン生成を光合成にせず、紫外と赤外を取り違えず、土壌CO2による炭酸供給の風化促進を理解する。風化速度・炭素埋没の時間尺度も意識する。

\Needspace{4\baselineskip}\paragraph{選択肢の検討（査読待ち初稿）}\mbox{}\par

\noindent\textbf{1：不適切} 大気炭素全量は海洋と同程度。。海洋の炭素貯留は大気よりはるかに多く、同程度ではない。\par

\noindent\textbf{2：適切} 初期CO2は高く、減少に風化と光合成生物が関与。。地史的なCO2減少に風化と生物活動が関与する。光合成固定も呼吸で戻る部分と埋没する部分を区別する。\par

\noindent\textbf{3：不適切} 嫌気環境のメタン生成は光合成。。嫌気条件は整合するがメタン生成は光合成ではない。水素とCO2、酢酸等の基質を使う代謝である。\par

\noindent\textbf{4：不適切} CO2は紫外に幅広い吸収帯。。主要な温室作用は赤外領域の吸収・放射である。紫外との取り違え。\par

\noindent\textbf{5：不適切} 土壌CO2が高いので風化のCO2除去が減速。。土壌呼吸によるCO2供給は炭酸を作り化学的風化を促す。実際の速度は水流、表面積、温度、飽和度等にも依存する。\par

\paragraph{再利用する考え方} 物質の貯留量・輸送速度・化学形態を分け、系の境界と反応の前提を確かめる。貯留・反応・放射を区別する炭素循環を収支・平衡・反応速度のどのモデルで扱うか選ぶ。

\paragraph{共通理論} \hyperref[theory:geochemistry:carbon]{第\ref*{theory:geochemistry:carbon}節}, \hyperref[theory:geochemistry:weathering]{第\ref*{theory:geochemistry:weathering}節}, \hyperref[theory:geochemistry:redox]{第\ref*{theory:geochemistry:redox}節}

\par 参考：\cite{gc_carbon,atm_noaa_ghg,gc_redox}。

\paragraph{関連問題} \hyperref[q:h26-reg-applied-27]{H26-REG-APPLIED-27}

\par 難易度 3（校準初稿）、優先度未評価。

\section{二酸化炭素回収・貯留}

\subsection{平成26年度 問18}
\label{q:h26-reg-applied-18}

\noindent\textbf{原典：} sources/H26/17\_applied\_science.pdf、PDF物理7頁。

\textbf{公式正答：} 4（single）。

\paragraph{問題の要旨} CCSの回収と地中貯留に直接つながる技術・場所を比較する。

\Needspace{4\baselineskip}\paragraph{確認した条件・記述（独自要約）}\mbox{}\par

\noindent\textbf{方向：} 関係性が最も薄いもの。\par

\noindent\textbf{対象：} CCS（Carbon dioxide Capture and Storage）。\par

\paragraph{問われる核心} CO2回収・利用・地中貯留の工程

\paragraph{必要知識} CO2回収・利用・地中貯留の工程。CCSを回収と貯留に分けると、アミン・膜は回収、帯水層は貯留、CO2-EORは利用と貯留に接続する。界面活性剤4が相対的に最も関係が薄い。原典の「薄い」を「無関係」に強めず、EORの種類も区別する。

\paragraph{考え方・正答を選ぶ理由（査読待ち初稿）} CCSを回収と貯留に分けると、アミン・膜は回収、帯水層は貯留、CO2-EORは利用と貯留に接続する。界面活性剤4が相対的に最も関係が薄い。原典の「薄い」を「無関係」に強めず、EORの種類も区別する。

\Needspace{4\baselineskip}\paragraph{選択肢の検討（査読待ち初稿）}\mbox{}\par

\noindent\textbf{1：直接関係} 帯水層。深部塩水帯水層は貯留候補である。圧力、封じ込め、漏えい等を評価する。\par

\noindent\textbf{2：CO2-EORで関係} EOR。CO2を用いるEORは石油増進回収とCO2貯留に結び付く。他のEORまで同一ではない。\par

\noindent\textbf{3：直接関係} モノエタノールアミン吸収法。MEAの反応吸収はCO2回収の代表方式である。\par

\noindent\textbf{4：提示中で最も薄い} 界面活性剤。界面活性剤は油回収等に使われるが、提示中ではCO2分離・貯留そのものとの直接関係が薄い。全く関係がないと断定しない。\par

\noindent\textbf{5：直接関係} 膜分離。選択透過する膜でCO2を分離回収できる。\par

\paragraph{再利用する考え方} 物質の貯留量・輸送速度・化学形態を分け、系の境界と反応の前提を確かめる。CO2回収・利用・地中貯留の工程を収支・平衡・反応速度のどのモデルで扱うか選ぶ。

\paragraph{共通理論} \hyperref[theory:geochemistry:ccs]{第\ref*{theory:geochemistry:ccs}節}

\par 参考：\cite{gc_ccs,gc_capture,gc_mea}。

\par 難易度 3（校準初稿）、優先度未評価。

\section{地殻元素存在度と比較基準}

\subsection{平成25年度 問10}
\label{q:h25-reg-applied-10}

\noindent\textbf{原典：} sources/H25/17\_applied\_science.pdf、PDF物理5頁。

\textbf{公式正答：} 4（single）。

\paragraph{問題の要旨} 地殻元素の存在度の5組の大小を比較する。

\Needspace{4\baselineskip}\paragraph{確認した条件・記述（独自要約）}\mbox{}\par

\noindent\textbf{方向：} 最も不適切なもの。\par

\noindent\textbf{対象：} 地球表層の自然水・岩石・元素循環。各肢の条件は独自要約として保持。\par

\paragraph{問われる核心} 地殻組成・質量存在度と微量元素

\paragraph{必要知識} 地殻組成・質量存在度と微量元素。地殻という対象と質量存在度の基準を置くと、Cl>Kという4だけが逆である。大気・海水や地球全体の組成を代用しない。代表表の精密小数はモデル依存であるが、この大小比較は頑健である。

\paragraph{考え方・正答を選ぶ理由（査読待ち初稿）} 地殻という対象と質量存在度の基準を置くと、Cl>Kという4だけが逆である。大気・海水や地球全体の組成を代用しない。代表表の精密小数はモデル依存であるが、この大小比較は頑健である。

\Needspace{4\baselineskip}\paragraph{選択肢の検討（査読待ち初稿）}\mbox{}\par

\noindent\textbf{1：適切（通常条件）} O > Si。代表質量存在度46.60\%対27.72\%で成立。\par

\noindent\textbf{2：適切（通常条件）} Si > Ti。27.72\%対0.44\%で成立。\par

\noindent\textbf{3：適切（通常条件）} Al > Cu。Al8.13\%は微量のCuより多い。\par

\noindent\textbf{4：不適切} Cl > K。Cl0.020\%はK2.59\%よりはるかに少なく逆である。\par

\noindent\textbf{5：適切（通常条件）} Ti > Mn。Ti0.44\%はMn0.100\%より多い。\par

\paragraph{再利用する考え方} 物質の貯留量・輸送速度・化学形態を分け、系の境界と反応の前提を確かめる。地殻組成・質量存在度と微量元素を収支・平衡・反応速度のどのモデルで扱うか選ぶ。

\paragraph{共通理論} \hyperref[theory:geochemistry:crust]{第\ref*{theory:geochemistry:crust}節}

\par 参考：\cite{gc_crust,gc_hem}。

\par 難易度 3（校準初稿）、優先度未評価。

\section{溶質の移流・拡散分散と観測尺度}

\subsection{令和3年度 問18}
\label{q:r03-reg-applied-18}

\noindent\textbf{原典：} sources/R03/17\_applied\_science.pdf、PDF物理7頁。

\textbf{公式正答：} 5（single）。

\paragraph{問題の要旨} 分子拡散、機械分散、縦横の分散長と野外の尺度を比較する。

\Needspace{4\baselineskip}\paragraph{確認した条件・記述（独自要約）}\mbox{}\par

\noindent\textbf{方向：} 最も不適切なもの。\par

\noindent\textbf{対象：} 地球表層の自然水・岩石・元素循環。各肢の条件は独自要約として保持。\par

\paragraph{問われる核心} 移流拡散・機械分散と尺度依存

\paragraph{必要知識} 移流拡散・機械分散と尺度依存。5の野外がはるかに小さいという一般化が逆である。\(D_L=\alpha_L|v|+D_m^\ast\)と置けば、分子拡散係数は面積/時間、分散長は長さであり別の量である。流れの平均速度に比例する機械分散と静止状態にもある分子拡散を分け、現場尺度への外挿を検証する。

\paragraph{考え方・正答を選ぶ理由（査読待ち初稿）} 5の野外がはるかに小さいという一般化が逆である。\(D_L=\alpha_L|v|+D_m^\ast\)と置けば、分子拡散係数は面積/時間、分散長は長さであり別の量である。流れの平均速度に比例する機械分散と静止状態にもある分子拡散を分け、現場尺度への外挿を検証する。

\Needspace{4\baselineskip}\paragraph{選択肢の検討（査読待ち初稿）}\mbox{}\par

\noindent\textbf{1：適切（通常条件）} 分子の不規則運動が分子拡散。。不規則な熱運動に基づく濃度の平滑化である。\par

\noindent\textbf{2：適切（通常条件）} 濃度勾配に比例する速度で分子拡散。。Fick則は拡散流束\(J=-D_m\partial C/\partial x\)を与える。「速度」は流束として解釈し、正味の分子流速との同一視を避ける。\par

\noindent\textbf{3：適切（通常条件）} 不均質な間隙構造で流動中の溶質集団が広がる。。流速・経路の差を平均化した見かけの広がりで、静止水の分子拡散とは区別する。\par

\noindent\textbf{4：適切（通常条件）} 横分散長は通常縦より小さい。。流れ方向の速度差が大きく、通常は\(\alpha_T<\alpha_L\)である。媒体構造に依存する経験傾向。\par

\noindent\textbf{5：不適切} 野外分散長は室内よりはるかに小さい。。野外は大きな不均質性を含み一般に大きい側の傾向がある。1992年59現場の研究は尺度傾向と大きな散らばりを示し、普遍直線は否定する。\par

\paragraph{再利用する考え方} 物質の貯留量・輸送速度・化学形態を分け、系の境界と反応の前提を確かめる。移流拡散・機械分散と尺度依存を収支・平衡・反応速度のどのモデルで扱うか選ぶ。

\paragraph{共通理論} \hyperref[theory:geochemistry:transport]{第\ref*{theory:geochemistry:transport}節}

\par 参考：\cite{gc_dispersion,gc_hem}。

\par 難易度 3（校準初稿）、優先度未評価。

\section{天然水の電子受容体・酸化還元電位}

\subsection{令和4年度 問19}
\label{q:r04-reg-applied-19}

\noindent\textbf{原典：} sources/R04/17\_applied\_science.pdf、PDF物理6頁。

\textbf{公式正答：} 3（single）。

\par\textbf{検討未確定：} 肢4の酸素減少は酸化可能な物質と反応する条件、5は還元ガス・空気混合の条件が必要。原典の省略を無条件の保証として教えず部分校正。

\paragraph{問題の要旨} 電子授受、基準電極、有酸素水のEhと地下・火山水の状態を比較する。

\Needspace{4\baselineskip}\paragraph{確認した条件・記述（独自要約）}\mbox{}\par

\noindent\textbf{方向：} 最も不適切なもの。\par

\noindent\textbf{対象：} 地球表層の自然水・岩石・元素循環。各肢の条件は独自要約として保持。\par

\paragraph{問われる核心} 酸化還元電位・電極基準と自然水の電子受容体

\paragraph{必要知識} 酸化還元電位・電極基準と自然水の電子受容体。有酸素水のEhを-250mVとする3が不適切。Nernst式の酸素/水境界はpH7・空気下で約+805mVを与え、還元側の符号との違いを確かめられる。ただし実測Ehは複数反応・電極応答を含む。4と5の一般傾向は物質の供給・反応・混合条件を付して説明する。

\paragraph{考え方・正答を選ぶ理由（査読待ち初稿）} 有酸素水のEhを-250mVとする3が不適切。Nernst式の酸素/水境界はpH7・空気下で約+805mVを与え、還元側の符号との違いを確かめられる。ただし実測Ehは複数反応・電極応答を含む。4と5の一般傾向は物質の供給・反応・混合条件を付して説明する。

\Needspace{4\baselineskip}\paragraph{選択肢の検討（査読待ち初稿）}\mbox{}\par

\noindent\textbf{1：適切} 還元剤が電子供与、酸化剤が受容。。還元剤自身は酸化され、酸化剤自身は還元される。\par

\noindent\textbf{2：適切（基準指定）} 標準水素電極0を基準とした電位をEhと表す。。比較には同じ基準電極へ換算する。ORP計の表示を無換算でEhと混ぜない。\par

\noindent\textbf{3：不適切} pH7・25℃の有酸素水Ehは約-250mV。。酸素還元境界は\(E_h=1.229-0.05916\,\mathrm{pH}+0.01479\log_{10}p_{O_2}\)V（25℃、水活量1）。空気0.21・pH7で約+0.805V、-0.250Vは還元側であり方向が違う。実測値は混合電位で理想境界と一致しなくてもよい。\par

\noindent\textbf{4：反応条件を付す} 深部をゆっくり通る水はO2を失い還元状態へ。。長い滞留中の酸化可能な有機物や鉱物との反応でO2を消費し得る。単に時間が長いだけで反応物なしに酸素が消えるという保存則違反の説明にしない。\par

\noindent\textbf{5：通常の還元熱水・混合注意} 火山温泉湧水はSO2・H2S等と反応し普通はO2なし。。熱水は還元性ガス等と反応して酸素が乏しい場合が多い。湧出後の空気混合などを含め全温泉が永遠に無酸素とはしない。\par

\paragraph{再利用する考え方} 物質の貯留量・輸送速度・化学形態を分け、系の境界と反応の前提を確かめる。酸化還元電位・電極基準と自然水の電子受容体を収支・平衡・反応速度のどのモデルで扱うか選ぶ。

\paragraph{共通理論} \hyperref[theory:geochemistry:redox]{第\ref*{theory:geochemistry:redox}節}

\par 参考：\cite{gc_hem,gc_redox,gc_redox2009}。

\par 難易度 3（校準初稿）、優先度未評価。

\section{窒素固定・硝化脱窒と炭素源}

\subsection{令和2年度 問19}
\label{q:r02-reg-applied-19}

\noindent\textbf{原典：} sources/R02/17\_applied\_science.pdf、PDF物理7頁。

\textbf{公式正答：} 5（single）。

\par\textbf{検討未確定：} 公式5を保持。自養脱窒の一次実験例があるため肢5を全化学合成生物で否定できない。肢3も燃料窒素由来NOxとN2の新規固定を区別する範囲疑義。意図と科学的適用範囲を分けて査読待ち。

\paragraph{問題の要旨} 固定・硝化・脱窒と、無機炭素を使う独立栄養を区別する。

\Needspace{4\baselineskip}\paragraph{確認した条件・記述（独自要約）}\mbox{}\par

\noindent\textbf{方向：} 最も不適切なもの。\par

\noindent\textbf{対象：} 地球表層の自然水・岩石・元素循環。各肢の条件は独自要約として保持。\par

\paragraph{問われる核心} 窒素反応の電子収支と炭素源の独立な分類

\paragraph{必要知識} 窒素反応の電子収支と炭素源の独立な分類。公式5と初回5は一致するが、硝化と従属栄養脱窒という教科書的な対照だけでは5を普遍的誤りとできない。硫黄自養脱窒では無機電子供与体の酸化と硝酸還元がエネルギーを供給し、無機炭素を同化する。電子供与体・受容体・炭素源を別軸にして、原典と科学的例外を併記する。

\paragraph{考え方・正答を選ぶ理由（査読待ち初稿）} 公式5と初回5は一致するが、硝化と従属栄養脱窒という教科書的な対照だけでは5を普遍的誤りとできない。硫黄自養脱窒では無機電子供与体の酸化と硝酸還元がエネルギーを供給し、無機炭素を同化する。電子供与体・受容体・炭素源を別軸にして、原典と科学的例外を併記する。

\Needspace{4\baselineskip}\paragraph{選択肢の検討（査読待ち初稿）}\mbox{}\par

\noindent\textbf{1：適切（脱窒過程）} 嫌気硝酸呼吸で硝酸を還元しN2O・N2生成。。硝酸を電子受容体とする脱窒でN2O等を経てN2となる。中間生成物の量は条件で変わる。\par

\noindent\textbf{2：適切} アンモニア→亜硝酸→硝酸は硝化。。酸化数を上げる硝化の二段階である。代表の独立栄養硝化生物は無機酸化のエネルギーで炭素を固定する。\par

\noindent\textbf{3：固定源を分ける条件付き} 雷・燃焼NOxは非生物的窒素固定。。雷や高温の熱的NOxはN2から反応性窒素を作る非生物固定である。ただし燃料に既に含まれた窒素の放出まで全て新規N2固定とすることはできない。\par

\noindent\textbf{4：適切（還元と同化）} 固定細菌・藍藻はN2からアミノ酸合成。。窒素固定ではN2をまず還元してアンモニアを得て、アミノ酸へ同化する。生成物と最初の酵素反応を区別する。\par

\noindent\textbf{5：意図上不適切・自養脱窒の例外} 化学合成生物は脱窒のエネルギーでCO2から有機物合成。。従属栄養脱窒を念頭にした設問意図では硝化との混同である。しかし1989年一次研究は嫌気条件で硫黄を電子供与体、HCO3-を炭素源とする自養脱窒菌を集積しているので、脱窒と無機炭素固定の結合は実在し一律に否定できない。\par

\paragraph{再利用する考え方} 物質の貯留量・輸送速度・化学形態を分け、系の境界と反応の前提を確かめる。窒素反応の電子収支と炭素源の独立な分類を収支・平衡・反応速度のどのモデルで扱うか選ぶ。

\paragraph{共通理論} \hyperref[theory:geochemistry:nitrogen]{第\ref*{theory:geochemistry:nitrogen}節}, \hyperref[theory:geochemistry:redox]{第\ref*{theory:geochemistry:redox}節}

\par 参考：\cite{gc_sulfur_denit,gc_redox,gc_hem}。

\par 難易度 3（校準初稿）、優先度未評価。

\section{陸水の主要イオン・地質と採水変化}

\subsection{平成24年度 問33}
\label{q:h24-reg-applied-33}

\noindent\textbf{原典：} sources/H24/17\_applied\_science.pdf、PDF物理10頁。

\textbf{公式正答：} 5（single）。

\paragraph{問題の要旨} 地下水の鉄沈殿、鉱山酸性水、日本のケイ酸、採水pHと世界河川のイオン源を判定する。

\Needspace{4\baselineskip}\paragraph{確認した条件・記述（独自要約）}\mbox{}\par

\noindent\textbf{方向：} 最も不適切なもの。\par

\noindent\textbf{対象：} 地球表層の自然水・岩石・元素循環。各肢の条件は独自要約として保持。\par

\paragraph{問われる核心} 陸水の組成を決める供給源と採水後の反応

\paragraph{必要知識} 陸水の組成を決める供給源と採水後の反応。風化、酸化還元、脱気、海塩を別の過程として読む。岩石との反応によるCa-HCO3系を落としてNaCl起源に限定する5が不適切である。日本と世界のケイ酸濃度は集計時期・SiとSiO2の表示基準を添える。

\paragraph{考え方・正答を選ぶ理由（査読待ち初稿）} 風化、酸化還元、脱気、海塩を別の過程として読む。岩石との反応によるCa-HCO3系を落としてNaCl起源に限定する5が不適切である。日本と世界のケイ酸濃度は集計時期・SiとSiO2の表示基準を添える。

\Needspace{4\baselineskip}\paragraph{選択肢の検討（査読待ち初稿）}\mbox{}\par

\noindent\textbf{1：適切（通常条件）} 低酸素地下水の湧出に伴い鉄が酸化して茶褐色沈殿。。還元的な地下でFe(II)が溶存し、湧出後に酸素を得るとFe(III)の水酸化物等が沈殿する。茶褐色沈殿を地下の酸化状態と同一視しない。\par

\noindent\textbf{2：適切（通常条件）} 硫化鉱物酸化で硫酸増加・pH低下。。硫化鉱物の酸化は硫酸と酸を生成する。炭酸塩等の中和能でpHの結果は変わるが、酸性鉱山排水の典型的原因である。\par

\noindent\textbf{3：適切（通常条件）} 日本の平均河川は世界より溶存ケイ酸が高い。。1970年一次論文の表4は、日本の蒸発残渣74.8mg/L・SiO2比26.84\%、世界100mg/L・11.67\%であり、濃度へ戻すと約20.1対11.7mg/Lである。古い集計の傾向であって全河川の普遍命題ではない。\par

\noindent\textbf{4：適切（通常条件）} 現場の非攪拌水は室内放置水よりpH低い傾向。。地下水は土壌CO2を多く溶かしている場合があり、採取・攪拌による脱気で炭酸平衡が動きpHが上がる。CO2過飽和という条件を付し、鉄酸化等で逆の変化もあり得るため現場測定を優先する。\par

\noindent\textbf{5：不適切} 世界河川の大半はNaCl主成分・海塩が主供給源。。河川の平均的な主要イオンはCa・HCO3等の岩石風化に依存する。海塩NaClは雨・沿岸・混合で重要だが世界河川の大半を一律に海塩支配とできない。\par

\paragraph{再利用する考え方} 物質の貯留量・輸送速度・化学形態を分け、系の境界と反応の前提を確かめる。陸水の組成を決める供給源と採水後の反応を収支・平衡・反応速度のどのモデルで扱うか選ぶ。

\paragraph{共通理論} \hyperref[theory:geochemistry:river]{第\ref*{theory:geochemistry:river}節}, \hyperref[theory:geochemistry:redox]{第\ref*{theory:geochemistry:redox}節}, \hyperref[theory:geochemistry:carbonate]{第\ref*{theory:geochemistry:carbonate}節}, \hyperref[theory:geochemistry:sampling]{第\ref*{theory:geochemistry:sampling}節}

\par 参考：\cite{gc_river1970,gc_hem,gc_redox,gc_mine}。

\paragraph{関連問題} \hyperref[q:h29-reg-applied-20]{H29-REG-APPLIED-20}

\par 難易度 3（校準初稿）、優先度未評価。

\subsection{平成25年度 問31}
\label{q:h25-reg-applied-31}

\noindent\textbf{原典：} sources/H25/17\_applied\_science.pdf、PDF物理11頁。

\textbf{公式正答：} 4（single）。

\par\textbf{検討未確定：} 原典は「炭酸カルシウム型」である。通常の主要溶存イオンによる名称はCa-HCO3（炭酸水素カルシウム）型で、CaCO3分子が主成分との理解は誤る。原典の名称を保持して部分校正。

\paragraph{問題の要旨} 希薄化、石灰岩反応、蒸発濃縮、水素酸素同位体と酸性河川の起源を比較する。

\Needspace{4\baselineskip}\paragraph{確認した条件・記述（独自要約）}\mbox{}\par

\noindent\textbf{方向：} 最も不適切なもの。\par

\noindent\textbf{対象：} 地球表層の自然水・岩石・元素循環。各肢の条件は独自要約として保持。\par

\paragraph{問われる核心} 陸水の反応・濃縮と同位体による水の起源

\paragraph{必要知識} 陸水の反応・濃縮と同位体による水の起源。重い同位体の分別がHとOで連動するため、天水線は正の相関を持つ。反比例とする4が不適切である。Ca-HCO3型の名称、蒸発濃縮と海水混合の違いを補足し、原典の全記述を無条件の一般法則へしない。

\paragraph{考え方・正答を選ぶ理由（査読待ち初稿）} 重い同位体の分別がHとOで連動するため、天水線は正の相関を持つ。反比例とする4が不適切である。Ca-HCO3型の名称、蒸発濃縮と海水混合の違いを補足し、原典の全記述を無条件の一般法則へしない。

\Needspace{4\baselineskip}\paragraph{選択肢の検討（査読待ち初稿）}\mbox{}\par

\noindent\textbf{1：適切（一般傾向）} 日本は多雨急流・低蒸発濃縮で成分が一般に希薄。。古い平均比較は支持する。流域地質・汚染・温泉混合等により例外がある。\par

\noindent\textbf{2：溶解は適切・名称疑義} 石灰岩溶解で弱アルカリ性・Ca増加・炭酸カルシウム型。。炭酸を含む水に石灰岩が溶けるとCaとHCO3が増す。「炭酸カルシウム型」という原本の呼び方と溶存種を区別する。\par

\noindent\textbf{3：適切（濃縮混合傾向）} 乾燥蒸発や河口で塩が濃くなり海水組成へ近づく傾向。。濃縮・海水混合は塩類濃度を上げる。ただし蒸発だけで全イオン比が海水へ必ず一致するわけではない。\par

\noindent\textbf{4：不適切} HとOの同位体比はおおむね反比例。。天水のδ2Hとδ18Oは正の相関で、代表線はδ2H=\allowbreak{}8δ18O+10。反比例ではない。海水起源の水蒸気が分別するという説明と区別する。\par

\noindent\textbf{5：適切} 火山・酸性降雨・鉱山排水が酸性水源。。これらは酸を供給する代表過程である。pHの結果は流域の中和反応にも依存する。\par

\paragraph{再利用する考え方} 物質の貯留量・輸送速度・化学形態を分け、系の境界と反応の前提を確かめる。陸水の反応・濃縮と同位体による水の起源を収支・平衡・反応速度のどのモデルで扱うか選ぶ。

\paragraph{共通理論} \hyperref[theory:geochemistry:river]{第\ref*{theory:geochemistry:river}節}, \hyperref[theory:geochemistry:carbonate]{第\ref*{theory:geochemistry:carbonate}節}, \hyperref[theory:geochemistry:isotope]{第\ref*{theory:geochemistry:isotope}節}

\par 参考：\cite{gc_river1970,gc_hem,gc_meteoric}。

\paragraph{関連問題} \hyperref[q:r01-reg-applied-20]{R01-REG-APPLIED-20}

\par 難易度 3（校準初稿）、優先度未評価。

\subsection{平成27年度 問27}
\label{q:h27-reg-applied-27}

\noindent\textbf{原典：} sources/H27/17\_applied\_science.pdf、PDF物理10頁。

\textbf{公式正答：} 5（single）。

\par\textbf{独立解答と公式正答に相違あり：} 独立解答 4|5。初回4|5と公式5を保持。肢2はCaと酸化条件、4は海洋逆風化条件、5は成長による取り込みと死後殻溶解の区別を省略。「ことがある」を原典通り保持し一意な科学的誤り肢の確定は査読待ち。

\paragraph{問題の要旨} ケイ酸供給、生物取り込み、硫化物反応、還元地下水と海洋反応を比較する。

\Needspace{4\baselineskip}\paragraph{確認した条件・記述（独自要約）}\mbox{}\par

\noindent\textbf{方向：} 最も不適切なもの。\par

\noindent\textbf{対象：} 地球表層の自然水・岩石・元素循環。各肢の条件は独自要約として保持。\par

\paragraph{問われる核心} 生物によるケイ酸取り込みと溶解・逆風化

\paragraph{必要知識} 生物によるケイ酸取り込みと溶解・逆風化。公式は5だが初回は4|5であった。珪藻成長では水からケイ酸を除くため意図上の誤答理由は説明できる。一方、死後殻の溶解放出と海洋逆風化の反応条件まで含めると、原典の省略された条件が重要である。公式から原典5を「成長時に」と改変せず、各肢の条件と未解決を保つ。

\paragraph{考え方・正答を選ぶ理由（査読待ち初稿）} 公式は5だが初回は4|5であった。珪藻成長では水からケイ酸を除くため意図上の誤答理由は説明できる。一方、死後殻の溶解放出と海洋逆風化の反応条件まで含めると、原典の省略された条件が重要である。公式から原典5を「成長時に」と改変せず、各肢の条件と未解決を保つ。

\Needspace{4\baselineskip}\paragraph{選択肢の検討（査読待ち初稿）}\mbox{}\par

\noindent\textbf{1：適切（集計傾向）} 日本の平均河川は世界より溶存ケイ酸が高い。。1970年一次論文の表4は、日本の蒸発残渣74.8mg/L・SiO2比26.84\%、世界100mg/L・11.67\%であり、濃度へ戻すと約20.1対11.7mg/Lである。古い集計の傾向であって全河川の普遍命題ではない。\par

\noindent\textbf{2：条件付き・省略疑義} 硫化物に富む底質なら硫酸カルシウム型河川水。。硫化物が酸化されれば硫酸が生じ、酸によりCa鉱物が溶ければCa-SO4型を形成し得る。ただし硫化物の存在だけでは酸素・Ca供給を保証しない。\par

\noindent\textbf{3：適切（反応条件）} 低酸素地下水は酸化反応でO2を消耗し還元的。。供給O2が少なく還元性成分と反応すると地下水の還元状態が進み、FeMn等の形態が変わる。\par

\noindent\textbf{4：逆風化条件で成立・省略疑義} アルミノケイ酸塩の海洋供給で海水pH低下。。海洋の粘土生成（逆風化）はアルカリ度を消費してCO2を生成し、pHを下げ得る。陸上のケイ酸塩風化による酸消費とは別過程である。原典が反応条件を省いており供給だけで必ず低下とは断定できない。\par

\noindent\textbf{5：意図上不適切・生活史範囲疑義} ケイ藻でケイ素が放出され河川濃度が上がる場合。。生きた珪藻はケイ酸を取り込み殻を作るので、成長を濃度上昇の機構とするなら逆である。ただし死後殻の溶解はケイ酸を再供給し得るため、「ことがある」という原典の広い表現は一律に誤りと確定できない。\par

\paragraph{再利用する考え方} 物質の貯留量・輸送速度・化学形態を分け、系の境界と反応の前提を確かめる。生物によるケイ酸取り込みと溶解・逆風化を収支・平衡・反応速度のどのモデルで扱うか選ぶ。

\paragraph{共通理論} \hyperref[theory:geochemistry:silica]{第\ref*{theory:geochemistry:silica}節}, \hyperref[theory:geochemistry:river]{第\ref*{theory:geochemistry:river}節}, \hyperref[theory:geochemistry:redox]{第\ref*{theory:geochemistry:redox}節}

\par 参考：\cite{gc_river1970,gc_hem,gc_diatom,gc_reverse}。

\paragraph{関連問題} \hyperref[q:r01-reg-applied-20]{R01-REG-APPLIED-20}

\par 難易度 3（校準初稿）、優先度未評価。

\subsection{平成29年度 問20}
\label{q:h29-reg-applied-20}

\noindent\textbf{原典：} sources/H29/17\_applied\_science.pdf、PDF物理6頁。

\textbf{公式正答：} 1（single）。

\paragraph{問題の要旨} 海塩支配の一般化、硫化物酸化、ケイ酸濃度、脱気pH、鉄沈殿を比較する。

\Needspace{4\baselineskip}\paragraph{確認した条件・記述（独自要約）}\mbox{}\par

\noindent\textbf{方向：} 最も不適切なもの。\par

\noindent\textbf{対象：} 地球表層の自然水・岩石・元素循環。各肢の条件は独自要約として保持。\par

\paragraph{問われる核心} 主要イオンの供給と採水後の化学反応

\paragraph{必要知識} 主要イオンの供給と採水後の化学反応。風化由来Ca-HCO3等を無視して世界の河川を海塩NaClに限定する1が不適切である。採水後のpH変化を地質の違いと混ぜず、CO2脱気の条件として考える。鉄は地下の溶存Fe(II)と地表沈殿Fe(III)を区別する。

\paragraph{考え方・正答を選ぶ理由（査読待ち初稿）} 風化由来Ca-HCO3等を無視して世界の河川を海塩NaClに限定する1が不適切である。採水後のpH変化を地質の違いと混ぜず、CO2脱気の条件として考える。鉄は地下の溶存Fe(II)と地表沈殿Fe(III)を区別する。

\Needspace{4\baselineskip}\paragraph{選択肢の検討（査読待ち初稿）}\mbox{}\par

\noindent\textbf{1：不適切} 世界河川の大半はNaCl主成分・海塩起源。。河川の平均的な主要イオンはCa・HCO3等の岩石風化に依存する。海塩NaClは雨・沿岸・混合で重要だが世界河川の大半を一律に海塩支配とできない。\par

\noindent\textbf{2：適切（通常条件）} 硫化鉱物酸化で硫酸増・pH低下。。硫化鉱物の酸化は硫酸と酸を生成する。炭酸塩等の中和能でpHの結果は変わるが、酸性鉱山排水の典型的原因である。\par

\noindent\textbf{3：適切（通常条件）} 日本平均河川ケイ酸は世界より高い。。1970年一次論文の表4は、日本の蒸発残渣74.8mg/L・SiO2比26.84\%、世界100mg/L・11.67\%であり、濃度へ戻すと約20.1対11.7mg/Lである。古い集計の傾向であって全河川の普遍命題ではない。\par

\noindent\textbf{4：適切（通常条件）} 非攪拌地下水pHは攪拌水より低い傾向。。地下水は土壌CO2を多く溶かしている場合があり、採取・攪拌による脱気で炭酸平衡が動きpHが上がる。CO2過飽和という条件を付し、鉄酸化等で逆の変化もあり得るため現場測定を優先する。\par

\noindent\textbf{5：適切（通常条件）} 低酸素地下水湧出で鉄の酸化沈殿。。還元的な地下でFe(II)が溶存し、湧出後に酸素を得るとFe(III)の水酸化物等が沈殿する。茶褐色沈殿を地下の酸化状態と同一視しない。\par

\paragraph{再利用する考え方} 物質の貯留量・輸送速度・化学形態を分け、系の境界と反応の前提を確かめる。主要イオンの供給と採水後の化学反応を収支・平衡・反応速度のどのモデルで扱うか選ぶ。

\paragraph{共通理論} \hyperref[theory:geochemistry:river]{第\ref*{theory:geochemistry:river}節}, \hyperref[theory:geochemistry:carbonate]{第\ref*{theory:geochemistry:carbonate}節}, \hyperref[theory:geochemistry:sampling]{第\ref*{theory:geochemistry:sampling}節}, \hyperref[theory:geochemistry:redox]{第\ref*{theory:geochemistry:redox}節}

\par 参考：\cite{gc_river1970,gc_hem,gc_redox,gc_mine}。

\paragraph{関連問題} \hyperref[q:h24-reg-applied-33]{H24-REG-APPLIED-33}

\par 難易度 3（校準初稿）、優先度未評価。

\subsection{令和元年度 問20}
\label{q:r01-reg-applied-20}

\noindent\textbf{原典：} sources/R01/17\_applied\_science.pdf、PDF物理6頁。

\textbf{公式正答：} 1（single）。

\par\textbf{検討未確定：} 原典は「炭酸カルシウム型」である。通常の主要溶存イオンによる名称はCa-HCO3（炭酸水素カルシウム）型で、CaCO3分子が主成分との理解は誤る。原典の名称を保持して部分校正。

\paragraph{問題の要旨} FeMnの溶存、火山地質ケイ酸、ORP、石灰岩反応と同位体緯度差を比較する。

\Needspace{4\baselineskip}\paragraph{確認した条件・記述（独自要約）}\mbox{}\par

\noindent\textbf{方向：} 最も不適切なもの。\par

\noindent\textbf{対象：} 地球表層の自然水・岩石・元素循環。各肢の条件は独自要約として保持。\par

\paragraph{問われる核心} 自然水の酸化還元・地質と溶存組成

\paragraph{必要知識} 自然水の酸化還元・地質と溶存組成。溶存FeMnを考えると河川と還元地下水の大小を逆にした1が不適切である。粒子を含む全濃度と溶存濃度も区別する。4の水質名称は原典を保持して疑義とし、炭酸平衡の溶存種はCa2+・HCO3-と説明する。

\paragraph{考え方・正答を選ぶ理由（査読待ち初稿）} 溶存FeMnを考えると河川と還元地下水の大小を逆にした1が不適切である。粒子を含む全濃度と溶存濃度も区別する。4の水質名称は原典を保持して疑義とし、炭酸平衡の溶存種はCa2+・HCO3-と説明する。

\Needspace{4\baselineskip}\paragraph{選択肢の検討（査読待ち初稿）}\mbox{}\par

\noindent\textbf{1：不適切} 河川水は地下水よりFeMnに富むことが多い。。酸化的河川ではFeMn酸化物が析出し、還元地下水はFe2+・Mn2+が溶存し得る。一般の大小方向が逆である。\par

\noindent\textbf{2：適切（地質傾向）} 新生代火山系を流れる河川でケイ酸が高い傾向。。火山性の反応しやすい物質からのケイ酸供給は高濃度へ寄与し得る。地質だけで全流域の濃度を固定しない。\par

\noindent\textbf{3：適切（基準と条件）} 大気共存水のORPは一般に正。。酸素が供給される通常の天然水は酸化的で、電位は正側の傾向。電極基準・温度pH・混合電位を区別する。\par

\noindent\textbf{4：反応は適切・名称疑義} 石灰岩溶解でCa増・弱アルカリ・炭酸カルシウム型。。石灰岩と炭酸の反応でCa・HCO3が増加する。原典の炭酸カルシウム型と通常のCa-HCO3名称を区別する。\par

\noindent\textbf{5：適切（一般傾向）} 降水酸素同位体値は赤道側で大きく極へ小さい。。水蒸気の逐次凝結で重い同位体が除かれ、高緯度降水は相対的に低いδ18Oとなる一般傾向。局地気候等の例外もある。\par

\paragraph{再利用する考え方} 物質の貯留量・輸送速度・化学形態を分け、系の境界と反応の前提を確かめる。自然水の酸化還元・地質と溶存組成を収支・平衡・反応速度のどのモデルで扱うか選ぶ。

\paragraph{共通理論} \hyperref[theory:geochemistry:river]{第\ref*{theory:geochemistry:river}節}, \hyperref[theory:geochemistry:redox]{第\ref*{theory:geochemistry:redox}節}, \hyperref[theory:geochemistry:carbonate]{第\ref*{theory:geochemistry:carbonate}節}, \hyperref[theory:geochemistry:isotope]{第\ref*{theory:geochemistry:isotope}節}

\par 参考：\cite{gc_hem,gc_river1970,gc_redox,gc_meteoric}。

\paragraph{関連問題} \hyperref[q:h25-reg-applied-31]{H25-REG-APPLIED-31}, \hyperref[q:h27-reg-applied-27]{H27-REG-APPLIED-27}

\par 難易度 3（校準初稿）、優先度未評価。

\section{水収支・貯留・平均滞留時間}

\subsection{平成23年度 問22}
\label{q:h23-reg-applied-22}

\noindent\textbf{原典：} sources/H23/17\_applied\_science.pdf、PDF物理7頁。

\textbf{公式正答：} 4（single）。

\paragraph{問題の要旨} 提示された陸上の水貯留の中で最大のものを選ぶ。

\Needspace{4\baselineskip}\paragraph{確認した条件・記述（独自要約）}\mbox{}\par

\noindent\textbf{方向：} 最も貯留量が大きいもの。\par

\noindent\textbf{対象：} 湖沼・河川・湿地・地下・土壌の5貯留。\par

\paragraph{問われる核心} 水の貯留量と対象集合

\paragraph{必要知識} 水の貯留量と対象集合。示された5形態では地下水4が最大である。氷河や海洋は候補にない。水の量をストックとして比較し、流量や取り出せる資源量のランキングへ移さない。

\paragraph{考え方・正答を選ぶ理由（査読待ち初稿）} 示された5形態では地下水4が最大である。氷河や海洋は候補にない。水の量をストックとして比較し、流量や取り出せる資源量のランキングへ移さない。

\Needspace{4\baselineskip}\paragraph{選択肢の検討（査読待ち初稿）}\mbox{}\par

\noindent\textbf{1：最大でない} 湖沼水。USGS表の全湖沼は約17.6万km3で、地下水より小さい。\par

\noindent\textbf{2：最大でない} 河川水。河川の瞬時貯留は約2120km3である。大きい年間輸送量と区別する。\par

\noindent\textbf{3：最大でない} 湿地の水。湿地の貯留は約11470km3で地下水より小さい。\par

\noindent\textbf{4：最大} 地下水。全地下水約2340万km3、淡水地下水だけでも約1053万km3で提示中最大。\par

\noindent\textbf{5：最大でない} 土壌水。土壌水約16500km3は地下水より小さい。\par

\paragraph{再利用する考え方} 物質の貯留量・輸送速度・化学形態を分け、系の境界と反応の前提を確かめる。水の貯留量と対象集合を収支・平衡・反応速度のどのモデルで扱うか選ぶ。

\paragraph{共通理論} \hyperref[theory:geochemistry:reservoir]{第\ref*{theory:geochemistry:reservoir}節}

\par 参考：\cite{gc_water}。

\paragraph{関連問題} \hyperref[q:h26-reg-applied-25]{H26-REG-APPLIED-25}

\par 難易度 3（校準初稿）、優先度未評価。

\subsection{平成26年度 問25}
\label{q:h26-reg-applied-25}

\noindent\textbf{原典：} sources/H26/17\_applied\_science.pdf、PDF物理11頁。

\textbf{公式正答：} 5（single）。

\par\textbf{検討未確定：} 原典の陸水に対する雪氷約70\%は、淡水貯留を分母にすれば68.7\%という代表表と合う。海洋以外の全水（塩性地下水等を含む）を陸水とすると約50\%に変わるため分母疑義を保持。

\paragraph{問題の要旨} 全球水量の概数・分母と、河川水と大気水の貯留量を比較する。

\Needspace{4\baselineskip}\paragraph{確認した条件・記述（独自要約）}\mbox{}\par

\noindent\textbf{方向：} 最も不適切なもの。\par

\noindent\textbf{対象：} 地球表層の自然水・岩石・元素循環。各肢の条件は独自要約として保持。\par

\paragraph{問われる核心} 全球水貯留・概数と割合の分母

\paragraph{必要知識} 全球水貯留・概数と割合の分母。同じ貯留量で比較すると河川2120と大気12900km3であり、多いとする5が不適切である。他の割合は丸めと分母を確認する。雪氷約70\%は淡水分母の説明として示し、原典の陸水という広い語は疑義を残す。

\paragraph{考え方・正答を選ぶ理由（査読待ち初稿）} 同じ貯留量で比較すると河川2120と大気12900km3であり、多いとする5が不適切である。他の割合は丸めと分母を確認する。雪氷約70\%は淡水分母の説明として示し、原典の陸水という広い語は疑義を残す。

\Needspace{4\baselineskip}\paragraph{選択肢の検討（査読待ち初稿）}\mbox{}\par

\noindent\textbf{1：適切（概数）} 総量約14億km3。1386000000km3という代表表を約14億km3へ丸めており妥当。\par

\noindent\textbf{2：適切（概数）} 海水約97\%。海水96.54\%を約97\%へ丸めている。\par

\noindent\textbf{3：分母疑義} 陸水の雪氷約70\%。淡水分母の氷は68.7\%で約70\%。陸水に塩性地下水まで含めるなら同じ割合でない。\par

\noindent\textbf{4：適切} 湖沼は地下水より少ない。。湖沼全体176400km3は全地下水23400000km3より小さい。\par

\noindent\textbf{5：不適切} 河川水は大気水蒸気より多い。。河川2120km3は大気12900km3より小さい。年間流出量をストックへ代入しない。\par

\paragraph{再利用する考え方} 物質の貯留量・輸送速度・化学形態を分け、系の境界と反応の前提を確かめる。全球水貯留・概数と割合の分母を収支・平衡・反応速度のどのモデルで扱うか選ぶ。

\paragraph{共通理論} \hyperref[theory:geochemistry:reservoir]{第\ref*{theory:geochemistry:reservoir}節}

\par 参考：\cite{gc_water}。

\paragraph{関連問題} \hyperref[q:h23-reg-applied-22]{H23-REG-APPLIED-22}, \hyperref[q:r01-ret-applied-18]{R01-RET-APPLIED-18}

\par 難易度 3（校準初稿）、優先度未評価。

\subsection{令和元年度再試験 問18}
\label{q:r01-ret-applied-18}

\noindent\textbf{原典：} sources/R01\_retest/17\_applied\_science.pdf、PDF物理6頁。

\textbf{公式正答：} 4（single）。

\par\textbf{検討未確定：} 原典の陸水に対する雪氷約70\%は、淡水貯留を分母にすれば68.7\%という代表表と合う。海洋以外の全水（塩性地下水等を含む）を陸水とすると約50\%に変わるため分母疑義を保持。

\paragraph{問題の要旨} 全球水の割合と地下・湖・河・大気の貯留を比較する。

\Needspace{4\baselineskip}\paragraph{確認した条件・記述（独自要約）}\mbox{}\par

\noindent\textbf{方向：} 最も不適切なもの。\par

\noindent\textbf{対象：} 地球表層の自然水・岩石・元素循環。各肢の条件は独自要約として保持。\par

\paragraph{問われる核心} 水貯留の分母と湖沼の淡水限定

\paragraph{必要知識} 水貯留の分母と湖沼の淡水限定。河川2120km3は大気12900km3より小さいため4が不適切。H26問25から候補順が変わり、湖沼に淡水という限定もあるので番号を転用しない。雪氷70\%の分母は淡水なら整合するが陸水という原語の範囲疑義を残す。

\paragraph{考え方・正答を選ぶ理由（査読待ち初稿）} 河川2120km3は大気12900km3より小さいため4が不適切。H26問25から候補順が変わり、湖沼に淡水という限定もあるので番号を転用しない。雪氷70\%の分母は淡水なら整合するが陸水という原語の範囲疑義を残す。

\Needspace{4\baselineskip}\paragraph{選択肢の検討（査読待ち初稿）}\mbox{}\par

\noindent\textbf{1：適切（概数）} 総量約14億km3。1386000000km3という代表表を約14億km3へ丸めており妥当。\par

\noindent\textbf{2：適切（概数）} 海水約97\%。海水96.54\%を約97\%へ丸めている。\par

\noindent\textbf{3：分母疑義} 陸水中雪氷約70\%。淡水分母の氷は68.7\%で約70\%。陸水に塩性地下水まで含めるなら同じ割合でない。\par

\noindent\textbf{4：不適切} 河川水は大気の水蒸気より多い。。河川2120km3は大気12900km3より小さい。年間流出量をストックへ代入しない。\par

\noindent\textbf{5：適切} 淡水湖は地下水より少ない。。淡水湖91000km3は淡水地下水10530000km3より小さく、全地下水と比較しても小さい。\par

\paragraph{再利用する考え方} 物質の貯留量・輸送速度・化学形態を分け、系の境界と反応の前提を確かめる。水貯留の分母と湖沼の淡水限定を収支・平衡・反応速度のどのモデルで扱うか選ぶ。

\paragraph{共通理論} \hyperref[theory:geochemistry:reservoir]{第\ref*{theory:geochemistry:reservoir}節}

\par 参考：\cite{gc_water}。

\paragraph{関連問題} \hyperref[q:h26-reg-applied-25]{H26-REG-APPLIED-25}

\par 難易度 3（校準初稿）、優先度未評価。

\section{天然水の電気伝導率・温度と組成}

\subsection{令和3年度 問19}
\label{q:r03-reg-applied-19}

\noindent\textbf{原典：} sources/R03/17\_applied\_science.pdf、PDF物理7頁。

\textbf{公式正答：} 2（single）。

\paragraph{問題の要旨} 天然水のイオン量・温度と電気伝導率、抵抗率と残留物を比較する。

\Needspace{4\baselineskip}\paragraph{確認した条件・記述（独自要約）}\mbox{}\par

\noindent\textbf{方向：} 最も不適切なもの。\par

\noindent\textbf{対象：} 地球表層の自然水・岩石・元素循環。各肢の条件は独自要約として保持。\par

\paragraph{問われる核心} イオンの伝導・抵抗率と温度補正

\paragraph{必要知識} イオンの伝導・抵抗率と温度補正。通常の天然水の電気伝導率は温度上昇で増すので2が不適切。\(\kappa=\sum_i \lambda_i c_i\)という希薄な加算近似では、イオン種別の伝導寄与と濃度の双方が必要である。比較時は温度を揃え、25℃換算値か実測温度の値かを明記する。

\paragraph{考え方・正答を選ぶ理由（査読待ち初稿）} 通常の天然水の電気伝導率は温度上昇で増すので2が不適切。\(\kappa=\sum_i \lambda_i c_i\)という希薄な加算近似では、イオン種別の伝導寄与と濃度の双方が必要である。比較時は温度を揃え、25℃換算値か実測温度の値かを明記する。

\Needspace{4\baselineskip}\paragraph{選択肢の検討（査読待ち初稿）}\mbox{}\par

\noindent\textbf{1：適切（通常条件）} 伝導率で溶存イオン量を大まかに比較。。イオン濃度と移動度に依存するので概略指標になる。中性有機物等の全溶存量の直接測定ではない。\par

\noindent\textbf{2：不適切} 温度が上がると伝導率低下。。同一試料・通常温度では温度上昇でイオン移動度が増し伝導率は上がる。25℃への温度補正値と生の温度依存を分ける。\par

\noindent\textbf{3：適切（通常条件）} 通常海水→河川→雨の順で低下。。塩類濃度の一般傾向に沿う。汚染・温泉・酸性水等の例外があり試料名だけで固定しない。\par

\noindent\textbf{4：適切（通常条件）} 1cm立方体の対面間に対応する電気抵抗率の逆数。。伝導率\(\kappa=1/\rho=\ell/(RA)\)。\(\ell=1\,\mathrm{cm},A=1\,\mathrm{cm}^2\)で測定抵抗Rの逆数の数値に対応するが単位はS/cm。抵抗と抵抗率の量を混ぜない。\par

\noindent\textbf{5：適切（通常条件）} 同じ組成比で濃度が違えば蒸発残留物と密接に関係。。同じ組成・温度・希薄条件なら比例的に変化する。ただし異なるイオン組成の試料を一つの係数で換算しない。\par

\paragraph{再利用する考え方} 物質の貯留量・輸送速度・化学形態を分け、系の境界と反応の前提を確かめる。イオンの伝導・抵抗率と温度補正を収支・平衡・反応速度のどのモデルで扱うか選ぶ。

\paragraph{共通理論} \hyperref[theory:geochemistry:conductivity]{第\ref*{theory:geochemistry:conductivity}節}

\par 参考：\cite{gc_hem,gc_conductance,gc_conductance1994}。

\par 難易度 3（校準初稿）、優先度未評価。

\section{鉱物水反応・活量と風化生成物}

\subsection{令和4年度 問20}
\label{q:r04-reg-applied-20}

\noindent\textbf{原典：} sources/R04/17\_applied\_science.pdf、PDF物理7頁。

\textbf{公式正答：} 1（single）。

\par\textbf{検討未確定：} 肢3は残存固相・固定条件での飽和平衡を必要とする。少量全溶解・開放流動を含めれば原典の無条件な最終平衡説明は成立しないため条件疑義を保持。

\paragraph{問題の要旨} 活量係数の理論、イオン強度、溶解平衡、風化系列と速度を区別する。

\Needspace{4\baselineskip}\paragraph{確認した条件・記述（独自要約）}\mbox{}\par

\noindent\textbf{方向：} 最も不適切なもの。\par

\noindent\textbf{対象：} 地球表層の自然水・岩石・元素循環。各肢の条件は独自要約として保持。\par

\paragraph{問われる核心} 溶解の平衡・活量と反応速度の分離

\paragraph{必要知識} 溶解の平衡・活量と反応速度の分離。1は理論の対象を取り違えている。Debye-Huckelはイオン強度から活量係数を近似する理論で、活量による平衡計算へ接続する。平衡でどこまで溶けるかと、そこへ至る速度、風化抵抗という経験系列を別の概念として整理する。

\paragraph{考え方・正答を選ぶ理由（査読待ち初稿）} 1は理論の対象を取り違えている。Debye-Huckelはイオン強度から活量係数を近似する理論で、活量による平衡計算へ接続する。平衡でどこまで溶けるかと、そこへ至る速度、風化抵抗という経験系列を別の概念として整理する。

\Needspace{4\baselineskip}\paragraph{選択肢の検討（査読待ち初稿）}\mbox{}\par

\noindent\textbf{1：不適切} Debye-Huckelは造岩鉱物の風化抵抗推定モデル。。Debye-Huckelは希薄溶液の静電相互作用と活量係数を扱う。風化抵抗系列そのもののモデルではない。\par

\noindent\textbf{2：適切} 各イオン濃度からイオン強度を計算。。価数zと濃度cを用いて\(I=\frac12\sum_i c_i z_i^2\)。価数の二乗を落とさず濃度基準を揃える。\par

\noindent\textbf{3：成立条件に疑義} 溶解の最終平衡濃度が溶解度。。固定温度・圧力・化学条件で残存固相と平衡に達した飽和濃度である。固相が全部溶けた場合や流動開放系はその濃度が溶解度とは限らない。\par

\noindent\textbf{4：適切（経験系列）} 風化抵抗はかんらん石→輝石→黒雲母→石英の順。。地表で高温形成鉱物が相対的に不安定な通常の風化系列の一部である。粒径や欠陥、流体条件で実際の速度は変わる。\par

\noindent\textbf{5：適切（通常反応条件）} 方解石はケイ酸塩鉱物より溶解が速い。。通常の炭酸を含む地表水では方解石は比較的速く溶ける。平衡に近いと速度は小さく、全温度・pH・鉱物に普遍とはしない。\par

\paragraph{再利用する考え方} 物質の貯留量・輸送速度・化学形態を分け、系の境界と反応の前提を確かめる。溶解の平衡・活量と反応速度の分離を収支・平衡・反応速度のどのモデルで扱うか選ぶ。

\paragraph{共通理論} \hyperref[theory:geochemistry:activity]{第\ref*{theory:geochemistry:activity}節}, \hyperref[theory:geochemistry:equilibrium]{第\ref*{theory:geochemistry:equilibrium}節}, \hyperref[theory:geochemistry:weathering]{第\ref*{theory:geochemistry:weathering}節}

\par 参考：\cite{gc_debye,gc_hem,gc_clays}。

\paragraph{関連問題} \hyperref[q:r06-reg-applied-22]{R06-REG-APPLIED-22}

\par 難易度 3（校準初稿）、優先度未評価。

\subsection{令和6年度 問22}
\label{q:r06-reg-applied-22}

\noindent\textbf{原典：} sources/R06/17\_applied\_science.pdf、PDF物理7頁。

\textbf{公式正答：} 2（single）。

\paragraph{問題の要旨} 土壌の酸化還元、炭酸塩析出、乾燥塩類、金属沈殿と粘土生成を比較する。

\Needspace{4\baselineskip}\paragraph{確認した条件・記述（独自要約）}\mbox{}\par

\noindent\textbf{方向：} 最も不適切なもの。\par

\noindent\textbf{対象：} 地球表層の自然水・岩石・元素循環。各肢の条件は独自要約として保持。\par

\paragraph{問われる核心} 化学的風化・炭酸平衡と土壌の二次鉱物

\paragraph{必要知識} 炭酸と炭酸水素イオンの平衡、CO2脱気による方解石析出、毛管上昇と蒸発濃縮、Al・Feの酸化状態と溶解度、溶脱と二次粘土鉱物。

\paragraph{考え方・正答を選ぶ理由（査読待ち初稿）} 2のCO2添加が逆である。\[\mathrm{CaCO_3+CO_2+H_2O\rightleftharpoons Ca^{2+}+2HCO_3^-}\]においてCO2を失うと左へ進み方解石を析出する。溶解、沈殿、蒸発濃縮、酸化還元、強溶脱の作用をそれぞれ区別して他肢を通常条件で判断する。

\Needspace{4\baselineskip}\paragraph{選択肢の検討（査読待ち初稿）}\mbox{}\par

\noindent\textbf{1：適切（通常条件）} 土壌有機物は還元型。。CO2に比べ有機炭素は還元されており、酸化で電子を供給できる。\par

\noindent\textbf{2：不適切} 石灰岩溶解地下水へ大気CO2が加わって鍾乳石・石筍形成。。普通の石灰岩洞窟ではCO2の脱気でCaCO3が析出する。添加では炭酸が増して溶解側へ動くので方向が逆。\par

\noindent\textbf{3：適切（通常条件）} 乾燥地で塩類が上昇・蒸発沈積。。毛管上昇と蒸発で溶存塩が濃縮される。原材料と飽和条件に応じ岩塩・石膏・炭酸塩等が析出する。\par

\noindent\textbf{4：適切（通常条件）} 通常土壌pHでAl3+やFeは酸化物・オキシ水酸化物沈殿。。通常の酸化的・弱酸性〜中性土壌ではAl・Fe(III)の水酸化物等が溶解を制限する。強酸、錯体、還元Fe(II)などの例外を付す。\par

\noindent\textbf{5：適切（通常条件）} 湿潤熱帯の強風化ではスメクタイトよりカオリナイト。。強い溶脱は塩基性陽イオンと一部Siを除き、カオリナイト等の比較的単純な相を残しやすい。気候だけで鉱物を一意に決めない。\par

\paragraph{再利用する考え方} 物質の貯留量・輸送速度・化学形態を分け、系の境界と反応の前提を確かめる。化学的風化・炭酸平衡と土壌の二次鉱物を収支・平衡・反応速度のどのモデルで扱うか選ぶ。

\paragraph{共通理論} \hyperref[theory:geochemistry:carbonate]{第\ref*{theory:geochemistry:carbonate}節}, \hyperref[theory:geochemistry:weathering]{第\ref*{theory:geochemistry:weathering}節}, \hyperref[theory:geochemistry:redox]{第\ref*{theory:geochemistry:redox}節}

\par 参考：\cite{gc_cave,gc_clays,gc_hem,gc_redox}。

\paragraph{関連問題} \hyperref[q:r04-reg-applied-20]{R04-REG-APPLIED-20}

\par 難易度 3（校準初稿）、優先度未評価。

\section{水質汚濁指標と起源トレーサー}

\subsection{平成23年度 問18}
\label{q:h23-reg-applied-18}

\noindent\textbf{原典：} sources/H23/17\_applied\_science.pdf、PDF物理6頁。

\textbf{公式正答：} 5（single）。

\paragraph{問題の要旨} 汚濁量を表す水質指標と、水の起源を追う同位体比を区別する。

\Needspace{4\baselineskip}\paragraph{確認した条件・記述（独自要約）}\mbox{}\par

\noindent\textbf{方向：} 最も不適切なもの。\par

\noindent\textbf{対象：} 地球表層の自然水・岩石・元素循環。各肢の条件は独自要約として保持。\par

\paragraph{問われる核心} 水質指標と起源トレーサーの用途

\paragraph{必要知識} 水質指標と起源トレーサーの用途。汚濁を評価する量と水の由来を同定する量は用途が異なる。1〜4は衛生・有機負荷・栄養塩・粒子の水質指標で、5が最も不適切である。同位体が汚染研究に利用できないという意味ではなく、単独の比を汚濁負荷へ変換できないという区別である。

\paragraph{考え方・正答を選ぶ理由（査読待ち初稿）} 汚濁を評価する量と水の由来を同定する量は用途が異なる。1〜4は衛生・有機負荷・栄養塩・粒子の水質指標で、5が最も不適切である。同位体が汚染研究に利用できないという意味ではなく、単独の比を汚濁負荷へ変換できないという区別である。

\Needspace{4\baselineskip}\paragraph{選択肢の検討（査読待ち初稿）}\mbox{}\par

\noindent\textbf{1：適切（通常条件）} 大腸菌群数。当時の衛生汚濁の指標である。菌数を物質濃度や発生源の同位体指標と混ぜない。\par

\noindent\textbf{2：適切（通常条件）} 生物化学的酸素要求量BOD。微生物分解に必要な酸素量で、分解可能な有機負荷を反映する。全有機物量そのものではない。\par

\noindent\textbf{3：適切（通常条件）} 全窒素。富栄養化に関与する窒素成分の総量を追う。\par

\noindent\textbf{4：適切（通常条件）} 浮遊物質SS。水に浮遊する粒子量を追う。溶存量とは区別する。\par

\noindent\textbf{5：不適切} 水素同位体比。蒸発・降水・混合等による水の起源と履歴を追うトレーサーであり、一般的な汚濁量の直接尺度ではない。\par

\paragraph{再利用する考え方} 物質の貯留量・輸送速度・化学形態を分け、系の境界と反応の前提を確かめる。水質指標と起源トレーサーの用途を収支・平衡・反応速度のどのモデルで扱うか選ぶ。

\paragraph{共通理論} \hyperref[theory:geochemistry:indicators]{第\ref*{theory:geochemistry:indicators}節}, \hyperref[theory:geochemistry:isotope]{第\ref*{theory:geochemistry:isotope}節}

\par 参考：\cite{gc_hem,gc_meteoric}。

\par 難易度 3（校準初稿）、優先度未評価。

\subsection{平成30年度 問28}
\label{q:h30-reg-applied-28}

\noindent\textbf{原典：} sources/H30/17\_applied\_science.pdf、PDF物理9頁。

\textbf{公式正答：} 2（single）。

\par\textbf{検討未確定：} 肢2の単位誤りは確認。肢3「金属ヒ素に毒性はない」は無条件に支持できず、肢4の岩種量比較と揮発性の因果は一次定量未確認。2018原典を科学的に正しい文へ修正せず疑義保持。

\paragraph{問題の要旨} 自然ヒ素、水道基準の単位、単体の有害性と岩石・鉱物での存在を比較する。

\Needspace{4\baselineskip}\paragraph{確認した条件・記述（独自要約）}\mbox{}\par

\noindent\textbf{方向：} 最も不適切。\par

\noindent\textbf{時点：} 2018年出題。2004施行表・2026年閲覧現行表を区別。\par

\noindent\textbf{単位：} 原典肢2はg/Lでありmg/Lへ校正しない。\par

\paragraph{問われる核心} 自然ヒ素の形態と濃度単位・基準の時点

\paragraph{必要知識} 自然ヒ素の形態と濃度単位・基準の時点。公式2・初回2を保持する。水道基準0.01mg/Lに対し原典は0.01g/Lであり1000倍違うため、2の不適切さは独立に確認できる。ただし3の無毒という断言は現在の一次有害性資料で支持できず、4の岩種比較も未確定である。正答だけが合ったことを全肢の科学的妥当性へ昇格させない。

\paragraph{考え方・正答を選ぶ理由（査読待ち初稿）} 公式2・初回2を保持する。水道基準0.01mg/Lに対し原典は0.01g/Lであり1000倍違うため、2の不適切さは独立に確認できる。ただし3の無毒という断言は現在の一次有害性資料で支持できず、4の岩種比較も未確定である。正答だけが合ったことを全肢の科学的妥当性へ昇格させない。

\Needspace{4\baselineskip}\paragraph{選択肢の検討（査読待ち初稿）}\mbox{}\par

\noindent\textbf{1：適切} 鉱物由来無機ヒ素が温泉・地下水へ溶出。。鉱物反応や酸化還元・吸着の変化によって自然由来の無機ヒ素が溶出する場合がある。\par

\noindent\textbf{2：不適切（単位1000倍）} 水道基準は0.01g/L以下。。2004施行表と現行表ともヒ素量0.01mg/Lで、0.01g/Lは10mg/L、1000倍の単位誤りである。\par

\noindent\textbf{3：無毒の断言に疑義} 半導体用金属ヒ素に毒性はない。。単体は溶解・取り込みが化合物と異なるが、NIOSHは金属ヒ素も有害性と曝露制限の対象として示す。無毒という絶対命題を教材の安全根拠にしない。\par

\noindent\textbf{4：比較と因果は要確認} 揮発性のため火成・変成岩にはあまり含まれない。。ヒ素は岩石中に微量含まれ火山ガスにも存在する。岩種ごとの含有量、鉱化、揮発移動を分ける必要があり、原典の因果・比較範囲は追加確認が必要である。\par

\noindent\textbf{5：適切（形態による）} 地表の水環境で亜ヒ酸塩・ヒ酸塩鉱物が形成。。酸化数III・Vに応じた亜ヒ酸塩・ヒ酸塩等の形態と二次鉱物があり、水と鉱物の反応に関わる。溶存種と固相鉱物を混ぜない。\par

\paragraph{再利用する考え方} 物質の貯留量・輸送速度・化学形態を分け、系の境界と反応の前提を確かめる。自然ヒ素の形態と濃度単位・基準の時点を収支・平衡・反応速度のどのモデルで扱うか選ぶ。

\paragraph{共通理論} \hyperref[theory:geochemistry:arsenic]{第\ref*{theory:geochemistry:arsenic}節}, \hyperref[theory:geochemistry:redox]{第\ref*{theory:geochemistry:redox}節}

\par 参考：\cite{gc_hem,gc_arsenic2004,gc_arsenic_current,gc_arsenic_hazard}。

\par 難易度 3（校準初稿）、優先度未評価。


\chapter{同位体地球化学}
\noindent 分野解説は執筆中。以下は登録問題と校正済み例の作業用一覧。
\subsection{平成23年度 問21}
\label{q:h23-reg-applied-21}

\noindent\textbf{原典：} sources/H23/17\_applied\_science.pdf、PDF物理7頁。

\textbf{公式正答：} 2（single）。

\par\textit{原典本文の校正と独自解説は未作成。分類のみの登録項目。}

\subsection{平成23年度 問24}
\label{q:h23-reg-applied-24}

\noindent\textbf{原典：} sources/H23/17\_applied\_science.pdf、PDF物理8頁。

\textbf{公式正答：} 4（single）。

\par\textit{原典本文の校正と独自解説は未作成。分類のみの登録項目。}

\subsection{平成24年度 問32}
\label{q:h24-reg-applied-32}

\noindent\textbf{原典：} sources/H24/17\_applied\_science.pdf、PDF物理9頁。

\textbf{公式正答：} 3（single）。

\par\textit{原典本文の校正と独自解説は未作成。分類のみの登録項目。}

\subsection{平成25年度 問27}
\label{q:h25-reg-applied-27}

\noindent\textbf{原典：} sources/H25/17\_applied\_science.pdf、PDF物理10頁。

\textbf{公式正答：} 2（single）。

\par\textit{原典本文の校正と独自解説は未作成。分類のみの登録項目。}

\subsection{平成26年度 問31}
\label{q:h26-reg-applied-31}

\noindent\textbf{原典：} sources/H26/17\_applied\_science.pdf、PDF物理13頁。

\textbf{公式正答：} 4（single）。

\par\textit{原典本文の校正と独自解説は未作成。分類のみの登録項目。}

\subsection{平成27年度 問25}
\label{q:h27-reg-applied-25}

\noindent\textbf{原典：} sources/H27/17\_applied\_science.pdf、PDF物理10頁。

\textbf{公式正答：} 3（single）。

\par\textit{原典本文の校正と独自解説は未作成。分類のみの登録項目。}

\subsection{平成28年度 問32}
\label{q:h28-reg-applied-32}

\noindent\textbf{原典：} sources/H28/17\_applied\_science.pdf、PDF物理13頁。

\textbf{公式正答：} 1（single）。

\par\textit{原典本文の校正と独自解説は未作成。分類のみの登録項目。}

\subsection{平成29年度 問30}
\label{q:h29-reg-applied-30}

\noindent\textbf{原典：} sources/H29/17\_applied\_science.pdf、PDF物理10頁。

\textbf{公式正答：} 3（single）。

\par\textit{原典本文の校正と独自解説は未作成。分類のみの登録項目。}

\subsection{平成30年度 問21}
\label{q:h30-reg-applied-21}

\noindent\textbf{原典：} sources/H30/17\_applied\_science.pdf、PDF物理7頁。

\textbf{公式正答：} 3（single）。

\par\textit{原典本文の校正と独自解説は未作成。分類のみの登録項目。}

\subsection{令和元年度再試験 問35}
\label{q:r01-ret-applied-35}

\noindent\textbf{原典：} sources/R01\_retest/17\_applied\_science.pdf、PDF物理11頁。

\textbf{公式正答：} 5（single）。

\par\textit{原典本文の校正と独自解説は未作成。分類のみの登録項目。}

\subsection{令和2年度 問20}
\label{q:r02-reg-applied-20}

\noindent\textbf{原典：} sources/R02/17\_applied\_science.pdf、PDF物理7頁。

\textbf{公式正答：} 3（single）。

\par\textit{原典本文の校正と独自解説は未作成。分類のみの登録項目。}

\subsection{令和4年度 問25}
\label{q:r04-reg-applied-25}

\noindent\textbf{原典：} sources/R04/17\_applied\_science.pdf、PDF物理9頁。

\textbf{公式正答：} 3（single）。

\par\textit{原典本文の校正と独自解説は未作成。分類のみの登録項目。}

\subsection{令和5年度 問21}
\label{q:r05-reg-applied-21}

\noindent\textbf{原典：} sources/R05/17\_applied\_science.pdf、PDF物理7頁。

\textbf{公式正答：} 1（single）。

\par\textit{原典本文の校正と独自解説は未作成。分類のみの登録項目。}

\subsection{令和6年度 問21}
\label{q:r06-reg-applied-21}

\noindent\textbf{原典：} sources/R06/17\_applied\_science.pdf、PDF物理7頁。

\textbf{公式正答：} 1（single）。

\par\textit{原典本文の校正と独自解説は未作成。分類のみの登録項目。}

\subsection{令和7年度 問22}
\label{q:r07-reg-applied-22}

\noindent\textbf{原典：} sources/R07/17\_applied\_science.pdf、PDF物理8頁。

\textbf{公式正答：} 4（single）。

\par\textit{原典本文の校正と独自解説は未作成。分類のみの登録項目。}


\chapter{地球物理・地球化学の補足}
\noindent 補足の章骨格。現在のDBにこの章への主分類割当はなく、追加の出題実績を意味しない。本文は未執筆。



\chapter{鉱物学}
\noindent 分野解説は執筆中。以下は登録問題と校正済み例の作業用一覧。
\subsection{平成24年度 問24}
\label{q:h24-reg-applied-24}

\noindent\textbf{原典：} sources/H24/17\_applied\_science.pdf、PDF物理8頁。

\textbf{公式正答：} 1（single）。

\par\textit{原典本文の校正と独自解説は未作成。分類のみの登録項目。}

\subsection{平成25年度 問18}
\label{q:h25-reg-applied-18}

\noindent\textbf{原典：} sources/H25/17\_applied\_science.pdf、PDF物理8頁。

\textbf{公式正答：} 3（single）。

\par\textit{原典本文の校正と独自解説は未作成。分類のみの登録項目。}

\subsection{平成25年度 問22}
\label{q:h25-reg-applied-22}

\noindent\textbf{原典：} sources/H25/17\_applied\_science.pdf、PDF物理9頁。

\textbf{公式正答：} 3（single）。

\par\textit{原典本文の校正と独自解説は未作成。分類のみの登録項目。}

\subsection{平成25年度 問23}
\label{q:h25-reg-applied-23}

\noindent\textbf{原典：} sources/H25/17\_applied\_science.pdf、PDF物理9頁。

\textbf{公式正答：} 4（single）。

\par\textit{原典本文の校正と独自解説は未作成。分類のみの登録項目。}

\subsection{平成26年度 問28}
\label{q:h26-reg-applied-28}

\noindent\textbf{原典：} sources/H26/17\_applied\_science.pdf、PDF物理12頁。

\textbf{公式正答：} 4（single）。

\par\textit{原典本文の校正と独自解説は未作成。分類のみの登録項目。}

\subsection{平成27年度 問30}
\label{q:h27-reg-applied-30}

\noindent\textbf{原典：} sources/H27/17\_applied\_science.pdf、PDF物理11頁。

\textbf{公式正答：} 2（single）。

\par\textit{原典本文の校正と独自解説は未作成。分類のみの登録項目。}

\subsection{平成29年度 問27}
\label{q:h29-reg-applied-27}

\noindent\textbf{原典：} sources/H29/17\_applied\_science.pdf、PDF物理9頁。

\textbf{公式正答：} 4（single）。

\par\textit{原典本文の校正と独自解説は未作成。分類のみの登録項目。}

\subsection{平成30年度 問31}
\label{q:h30-reg-applied-31}

\noindent\textbf{原典：} sources/H30/17\_applied\_science.pdf、PDF物理10頁。

\textbf{公式正答：} 4（single）。

\par\textit{原典本文の校正と独自解説は未作成。分類のみの登録項目。}

\subsection{令和元年度 問26}
\label{q:r01-reg-applied-26}

\noindent\textbf{原典：} sources/R01/17\_applied\_science.pdf、PDF物理9頁。

\textbf{公式正答：} 4（single）。

\par\textit{原典本文の校正と独自解説は未作成。分類のみの登録項目。}

\subsection{令和元年度再試験 問25}
\label{q:r01-ret-applied-25}

\noindent\textbf{原典：} sources/R01\_retest/17\_applied\_science.pdf、PDF物理8頁。

\textbf{公式正答：} 3（single）。

\par\textit{原典本文の校正と独自解説は未作成。分類のみの登録項目。}

\subsection{令和3年度 問22}
\label{q:r03-reg-applied-22}

\noindent\textbf{原典：} sources/R03/17\_applied\_science.pdf、PDF物理8頁。

\textbf{公式正答：} 2（single）。

\par\textit{原典本文の校正と独自解説は未作成。分類のみの登録項目。}

\subsection{令和3年度 問25}
\label{q:r03-reg-applied-25}

\noindent\textbf{原典：} sources/R03/17\_applied\_science.pdf、PDF物理9頁。

\textbf{公式正答：} 4（single）。

\par\textit{原典本文の校正と独自解説は未作成。分類のみの登録項目。}

\subsection{令和4年度 問23}
\label{q:r04-reg-applied-23}

\noindent\textbf{原典：} sources/R04/17\_applied\_science.pdf、PDF物理8頁。

\textbf{公式正答：} 4（single）。

\par\textit{原典本文の校正と独自解説は未作成。分類のみの登録項目。}

\subsection{令和6年度 問23}
\label{q:r06-reg-applied-23}

\noindent\textbf{原典：} sources/R06/17\_applied\_science.pdf、PDF物理8頁。

\textbf{公式正答：} 2（single）。

\par\textit{原典本文の校正と独自解説は未作成。分類のみの登録項目。}


\chapter{岩石学}
\noindent 分野解説は執筆中。以下は登録問題と校正済み例の作業用一覧。
\section{付加体の構成岩・形成過程}

\subsection{令和元年度再試験 問31}
\label{q:r01-ret-applied-31}

\noindent\textbf{原典：} sources/R01\_retest/17\_applied\_science.pdf、PDF物理10頁。

\textbf{公式正答：} 4（single）。

\paragraph{問題の要旨} 付加体へ取り込まれる典型的な岩石群から、最も関係の薄いものを選ぶ。

\Needspace{4\baselineskip}\paragraph{確認した条件・記述（独自要約）}\mbox{}\par

\noindent\textbf{方向：} 付加体の構成岩として最も関係のないもの。\par

\noindent\textbf{範囲：} 典型的な付加構成岩と後の貫入岩を区別。\par

\paragraph{問われる核心} 付加体の海洋プレート起源の岩石群

\paragraph{必要知識} 付加体の海洋プレート起源の岩石群。岩石名と生成過程、鉱物組合せ、温度・圧力・全岩組成、一般則と例外の区別。

\paragraph{考え方・正答を選ぶ理由（査読待ち初稿）} 海洋プレート本体・海山から玄武岩、遠洋の珪質堆積からチャート、海山上の炭酸塩堆積から石灰岩、海溝等の砕屑堆積から泥岩を対応させる。付加により取り込まれる典型的な岩石ではない花崗岩の4を選ぶ。現在の付加体地域に後で貫入した花崗岩が存在することとは両立する。

\Needspace{4\baselineskip}\paragraph{選択肢の検討（査読待ち初稿）}\mbox{}\par

\noindent\textbf{1：構成岩となり得る} 玄武岩。海洋地殻・海山などの火山岩が沈み込み時に付加され得る。\par

\noindent\textbf{2：構成岩となり得る} 泥岩。海溝などに運ばれる細粒砕屑堆積物が構成物となる。\par

\noindent\textbf{3：構成岩となり得る} チャート。海洋プレート上の遠洋性珪質堆積物として運搬・付加され得る。\par

\noindent\textbf{4：最も関係が薄い} 花崗岩。一般に地下の珪長質マグマの結晶化で形成され、海洋プレートの層序を剥ぎ取る付加物の典型ではない。付加後の花崗岩貫入まで否定しない。\par

\noindent\textbf{5：構成岩となり得る} 石灰岩。海山上の浅海生物起源の炭酸塩岩などが付加体へ取り込まれ得る。\par

\paragraph{再利用する考え方} 付加体の海洋プレート起源の岩石群を成因と観察される性質に結び付け、例示と全称命題、主要鉱物と副成分、公式解答と未解決の解釈を区別する。

\paragraph{共通理論} \hyperref[theory:petrology:accretion]{第\ref*{theory:petrology:accretion}節}

\par 参考：\cite{aist_accretion2026,gsj_saganoseki1993}。

\paragraph{関連問題} \hyperref[q:r05-reg-applied-26]{R05-REG-APPLIED-26}

\par 難易度 3（校準初稿）、優先度未評価。

\subsection{令和5年度 問26}
\label{q:r05-reg-applied-26}

\noindent\textbf{原典：} sources/R05/17\_applied\_science.pdf、PDF物理9頁。

\textbf{公式正答：} 2（single）。

\paragraph{問題の要旨} 付加体へ取り込まれる典型的な岩石群から、最も関係の薄いものを選ぶ。

\Needspace{4\baselineskip}\paragraph{確認した条件・記述（独自要約）}\mbox{}\par

\noindent\textbf{方向：} 付加体の構成岩として最も関係のないもの。\par

\noindent\textbf{範囲：} 典型的な付加構成岩と後の貫入岩を区別。\par

\paragraph{問われる核心} 付加体の海洋プレート起源の岩石群

\paragraph{必要知識} 付加体の海洋プレート起源の岩石群。岩石名と生成過程、鉱物組合せ、温度・圧力・全岩組成、一般則と例外の区別。

\paragraph{考え方・正答を選ぶ理由（査読待ち初稿）} 海洋プレート本体・海山から玄武岩、遠洋の珪質堆積からチャート、海山上の炭酸塩堆積から石灰岩、海溝等の砕屑堆積から泥岩を対応させる。付加により取り込まれる典型的な岩石ではない花崗岩の2を選ぶ。現在の付加体地域に後で貫入した花崗岩が存在することとは両立する。

\Needspace{4\baselineskip}\paragraph{選択肢の検討（査読待ち初稿）}\mbox{}\par

\noindent\textbf{1：構成岩となり得る} 玄武岩。海洋地殻・海山などの火山岩が沈み込み時に付加され得る。\par

\noindent\textbf{2：最も関係が薄い} 花崗岩。一般に地下の珪長質マグマの結晶化で形成され、海洋プレートの層序を剥ぎ取る付加物の典型ではない。付加後の花崗岩貫入まで否定しない。\par

\noindent\textbf{3：構成岩となり得る} 石灰岩。海山上の浅海生物起源の炭酸塩岩などが付加体へ取り込まれ得る。\par

\noindent\textbf{4：構成岩となり得る} チャート。海洋プレート上の遠洋性珪質堆積物として運搬・付加され得る。\par

\noindent\textbf{5：構成岩となり得る} 泥岩。海溝などに運ばれる細粒砕屑堆積物が構成物となる。\par

\paragraph{再利用する考え方} 付加体の海洋プレート起源の岩石群を成因と観察される性質に結び付け、例示と全称命題、主要鉱物と副成分、公式解答と未解決の解釈を区別する。

\paragraph{共通理論} \hyperref[theory:petrology:accretion]{第\ref*{theory:petrology:accretion}節}

\par 参考：\cite{aist_accretion2026,gsj_saganoseki1993}。

\paragraph{関連問題} \hyperref[q:r01-ret-applied-31]{R01-RET-APPLIED-31}

\par 難易度 3（校準初稿）、優先度未評価。

\section{鉱物・岩石}

\subsection{平成24年度 問29}
\label{q:h24-reg-applied-29}

\noindent\textbf{原典：} sources/H24/17\_applied\_science.pdf、PDF物理9頁。

\textbf{公式正答：} 3（single）。

\par\textbf{検討未確定：} 粘板岩は通常は低度変成岩に分類されるが、堆積原岩との関係がある。マイロナイトにも堆積原岩の例を排除できず「最も関係ない」という相対判断の範囲をAstraへ残す。

\paragraph{問題の要旨} 堆積岩名と、堆積原岩を変成した岩石、剪断変形で分類する岩石の関係を比較する。

\Needspace{4\baselineskip}\paragraph{確認した条件・記述（独自要約）}\mbox{}\par

\noindent\textbf{判断方向：} 堆積岩に最も関係のない名称を選ぶ\par

\noindent\textbf{原典照合：} 15原典画像で18問の問題幹、全5肢、否定、専門語、化学式、候補順を確認。\par

\noindent\textbf{適用条件：} 流れと構造の典型対応は一意成因を保証しない。用語、時間間隙、保存、試料処理の条件を個別に示す。\par

\paragraph{問われる核心} 堆積岩名と、堆積原岩を変成した岩石、剪断変形で分類する岩石の関係を比較する。

\paragraph{必要知識} 鉱物名と岩石名、炭酸塩と珪質堆積、泥質原岩と変成、剪断変形による分類。

\paragraph{考え方・正答を選ぶ理由（査読待ち初稿）} 本問の方向は「堆積岩でないもの」ではなく「最も関係のないもの」である。ドロマイトは苦灰岩と明示され、石灰岩とチャートも堆積に関連する。粘板岩は泥質堆積原岩の変成として関係を持つ。一方、マイロナイトは原岩の種類より剪断変形に基づく名称なので3を選ぶ。ただし堆積原岩からマイロナイトができる場合もあり、関係が絶対にゼロという解釈にはしない。

\Needspace{4\baselineskip}\paragraph{選択肢の検討（査読待ち初稿）}\mbox{}\par

\noindent\textbf{1：適切（条件を含めて検討）} ドロマイト（苦灰岩）。本問は鉱物名でなく括弧で苦灰岩という岩石を指定する。炭酸塩堆積岩に関係する。\par

\noindent\textbf{2：適切（条件を含めて検討）} 石灰岩。炭酸カルシウム主体の岩石で、生物・化学的堆積やその後の変化を考える。\par

\noindent\textbf{3：誤り・最も不適切} マイロナイト。強い剪断変形と動的再結晶などで生じる変形岩の分類であり、特定の堆積成因を表す名ではない。\par

\noindent\textbf{4：適切（条件を含めて検討）} チャート。シリカに富む堆積岩で、珪質生物殻の集積などが代表となる。\par

\noindent\textbf{5：適切（条件を含めて検討）} 粘板岩。泥質堆積原岩の低度変成と劈開に関わる。GSJにも堆積岩区分の登米粘板岩の標本があり、堆積岩との関係を切れない。\par

\paragraph{再利用する考え方} 観察できる構造と成因・環境の解釈を区別する。代表モデルに適用条件を添え、単一の形や測定値から唯一の環境を断定しない。

\paragraph{共通理論} \hyperref[theory:petrology:cycle]{第\ref*{theory:petrology:cycle}節}, \hyperref[theory:petrology:metamorphism]{第\ref*{theory:petrology:metamorphism}節}, \hyperref[theory:petrology:deposition]{第\ref*{theory:petrology:deposition}節}

\par 参考：\cite{sed_gsj_slate}。

\par 難易度 3（校準初稿）、優先度未評価。

\section{火成岩組成・鉱物組合せ}

\subsection{平成27年度 問28}
\label{q:h27-reg-applied-28}

\noindent\textbf{原典：} sources/H27/17\_applied\_science.pdf、PDF物理11頁。

\textbf{公式正答：} 3（single）。

\paragraph{問題の要旨} 玄武岩の組成・起源・粘性について不適切な記述を選ぶ。

\Needspace{4\baselineskip}\paragraph{確認した条件・記述（独自要約）}\mbox{}\par

\noindent\textbf{方向：} 最も不適切なもの。\par

\noindent\textbf{対象：} 一般的な岩石分類・形成過程。\par

\paragraph{問われる核心} 玄武岩質マグマの組成・部分溶融・流動性

\paragraph{必要知識} 玄武岩質マグマの組成・部分溶融・流動性。岩石名と生成過程、鉱物組合せ、温度・圧力・全岩組成、一般則と例外の区別。

\paragraph{考え方・正答を選ぶ理由（査読待ち初稿）} 苦鉄質という分類がFe・Mgに富む組成と対応する。流紋岩はSiO2に富む珪長質岩であり、比較を逆にした3が誤りである。粘性と形状は組成だけでなく温度や結晶量にも依存するため、一般的傾向として読む。

\Needspace{4\baselineskip}\paragraph{選択肢の検討（査読待ち初稿）}\mbox{}\par

\noindent\textbf{1：適切（一般的分類）} 富士火山や伊豆大島を構成する苦鉄質火山岩である。。代表的な玄武岩火山の例である。ただし火山全体の全噴出物を単一岩種と断定しない。\par

\noindent\textbf{2：適切（一般的分類）} SiO2はおよそ50質量\%である。。玄武岩の代表的値で、全ての試料を正確に50\%とする意味ではない。\par

\noindent\textbf{3：不適切} 鉄・マグネシウムの含有量は流紋岩より少ない。。玄武岩は流紋岩より一般にFe・Mgに富むので、比較方向が逆。\par

\noindent\textbf{4：適切（一般的分類）} 上部マントル岩石の部分溶融が起源の一つである。。マントルのかんらん岩が部分溶融して玄武岩質メルトを生じ得る。起源を全溶融と混同しない。\par

\noindent\textbf{5：適切（一般的分類）} 流紋岩より粘性が低く、平らな火山体を作りやすい。。一般に高温でSiO2に乏しい玄武岩質溶岩は流紋岩質より流れやすい。噴火様式・供給量・地形も火山形状へ影響する。\par

\paragraph{再利用する考え方} 玄武岩質マグマの組成・部分溶融・流動性を成因と観察される性質に結び付け、例示と全称命題、主要鉱物と副成分、公式解答と未解決の解釈を区別する。

\paragraph{共通理論} \hyperref[theory:petrology:igneous]{第\ref*{theory:petrology:igneous}節}, \hyperref[theory:petrology:magma]{第\ref*{theory:petrology:magma}節}

\par 参考：\cite{jma_magma1999,gsj_igneous2021}。

\par 難易度 3（校準初稿）、優先度未評価。

\subsection{平成30年度 問30}
\label{q:h30-reg-applied-30}

\noindent\textbf{原典：} sources/H30/17\_applied\_science.pdf、PDF物理10頁。

\textbf{公式正答：} 5（single）。

\par\textbf{独立解答と公式正答に相違あり：} 独立解答 2|5。原典は花崗岩の主な有色鉱物に白雲母を含める。通常の無色分類と矛盾し、初回独立解答2|5、公式5。白雲母を改変しない。石英はんれい岩の例外と「主な」の範囲を含めAstraで査読。

\paragraph{問題の要旨} 火成岩の主要鉱物・斑晶と有色／無色の区別について不適切な記述を選ぶ。

\Needspace{4\baselineskip}\paragraph{確認した条件・記述（独自要約）}\mbox{}\par

\noindent\textbf{方向：} 最も不適切なもの。\par

\noindent\textbf{対象：} 一般的な岩石分類・形成過程。\par

\paragraph{問われる核心} 火成岩の鉱物組合せと有色分類の原典疑義

\paragraph{必要知識} 火成岩の鉱物組合せと有色分類の原典疑義。岩石名と生成過程、鉱物組合せ、温度・圧力・全岩組成、一般則と例外の区別。

\paragraph{考え方・正答を選ぶ理由（査読待ち初稿）} 一般的なはんれい岩へ石英を主要鉱物として入れた5が公式正答である。一方、原典に白雲母を「有色鉱物」とする記述もある。産総研の結晶構造ギャラリーは白雲母の色を無色と記載しており、初回独立解答は2|5であった。原典を黒雲母へ直したり白雲母も通常有色と教えたりせず、疑義を保存する。石英はんれい岩の例外も含め、公式の一般的主要鉱物の比較と厳密な分類を分ける。

\Needspace{4\baselineskip}\paragraph{選択肢の検討（査読待ち初稿）}\mbox{}\par

\noindent\textbf{1：適切（鉱物群の例示）} 花崗岩質岩石は石英・アルカリ長石・斜長石を無色鉱物として含む。。珪長質岩の主要な明色鉱物である。ただし鉱物量比が分類を左右し、あらゆる花崗岩類が同量の3種を含むとはしない。\par

\noindent\textbf{2：原典分類に疑義（通常は白雲母は無色）} 花崗岩質岩石の主な有色鉱物は白雲母・黒雲母・角閃石である。。黒雲母・角閃石は代表的な有色鉱物だが、白雲母は通常無色として扱う。白雲母が花崗岩に存在することと、有色分類が正しいことは別であり、原典のまま疑義を残す。\par

\noindent\textbf{3：適切（鉱物群の例示）} 玄武岩の斑晶には普通輝石・紫蘇輝石・斜長石・かんらん石がある。。玄武岩の種類により斑晶に現れ得る鉱物群である。全試料に全て同時に存在するという意味には読まない。紫蘇輝石は古い名称で斜方輝石の一種を指す。\par

\noindent\textbf{4：適切（鉱物群の例示）} 安山岩の斑晶には斜長石・普通輝石・紫蘇輝石・角閃石・磁鉄鉱・黒雲母・かんらん石がある。。安山岩の組成・水分・酸化状態等により現れ得る鉱物群である。かんらん石や黒雲母の存在の可否を、単一の岩種名だけで禁止しない。\par

\noindent\textbf{5：不適切（公式の選択対象）} はんれい岩の主な構成鉱物は普通輝石・斜長石・石英である。。一般的なはんれい岩は斜長石・輝石等を主成分とし、石英を典型的主要鉱物に並べる記述が公式の選択対象である。ただし石英はんれい岩という区分はあり、全はんれい岩で石英が絶対にないとはしない。\par

\paragraph{再利用する考え方} 火成岩の鉱物組合せと有色分類の原典疑義を成因と観察される性質に結び付け、例示と全称命題、主要鉱物と副成分、公式解答と未解決の解釈を区別する。

\paragraph{共通理論} \hyperref[theory:petrology:igneous]{第\ref*{theory:petrology:igneous}節}, \hyperref[theory:petrology:minerals]{第\ref*{theory:petrology:minerals}節}, \hyperref[theory:petrology:exceptions]{第\ref*{theory:petrology:exceptions}節}

\par 参考：\cite{aist_crystal_colors,gsj_igneous2021,usgs_rock_thesaurus}。

\paragraph{関連問題} \hyperref[q:r02-reg-applied-25]{R02-REG-APPLIED-25}

\par 難易度 3（校準初稿）、優先度未評価。

\subsection{令和2年度 問25}
\label{q:r02-reg-applied-25}

\noindent\textbf{原典：} sources/R02/17\_applied\_science.pdf、PDF物理9頁。

\textbf{公式正答：} 2（single）。

\par\textbf{独立解答と公式正答に相違あり：} 独立解答 1|2。原典は花崗岩の主な有色鉱物に白雲母を含める。通常の無色分類と矛盾し、初回独立解答1|2、公式2。白雲母を改変しない。石英はんれい岩の例外と「主な」の範囲を含めAstraで査読。

\paragraph{問題の要旨} 火成岩の主要鉱物・斑晶と有色／無色の区別について不適切な記述を選ぶ。

\Needspace{4\baselineskip}\paragraph{確認した条件・記述（独自要約）}\mbox{}\par

\noindent\textbf{方向：} 最も不適切なもの。\par

\noindent\textbf{対象：} 一般的な岩石分類・形成過程。\par

\paragraph{問われる核心} 火成岩の鉱物組合せと有色分類の原典疑義

\paragraph{必要知識} 火成岩の鉱物組合せと有色分類の原典疑義。岩石名と生成過程、鉱物組合せ、温度・圧力・全岩組成、一般則と例外の区別。

\paragraph{考え方・正答を選ぶ理由（査読待ち初稿）} 一般的なはんれい岩へ石英を主要鉱物として入れた2が公式正答である。一方、原典に白雲母を「有色鉱物」とする記述もある。産総研の結晶構造ギャラリーは白雲母の色を無色と記載しており、初回独立解答は1|2であった。原典を黒雲母へ直したり白雲母も通常有色と教えたりせず、疑義を保存する。石英はんれい岩の例外も含め、公式の一般的主要鉱物の比較と厳密な分類を分ける。

\Needspace{4\baselineskip}\paragraph{選択肢の検討（査読待ち初稿）}\mbox{}\par

\noindent\textbf{1：原典分類に疑義（通常は白雲母は無色）} 花崗岩質岩石の主な有色鉱物は白雲母・黒雲母・角閃石である。。黒雲母・角閃石は代表的な有色鉱物だが、白雲母は通常無色として扱う。白雲母が花崗岩に存在することと、有色分類が正しいことは別であり、原典のまま疑義を残す。\par

\noindent\textbf{2：不適切（公式の選択対象）} はんれい岩の主な構成鉱物は普通輝石・斜長石・石英である。。一般的なはんれい岩は斜長石・輝石等を主成分とし、石英を典型的主要鉱物に並べる記述が公式の選択対象である。ただし石英はんれい岩という区分はあり、全はんれい岩で石英が絶対にないとはしない。\par

\noindent\textbf{3：適切（鉱物群の例示）} 花崗岩質岩石は石英・アルカリ長石・斜長石を無色鉱物として含む。。珪長質岩の主要な明色鉱物である。ただし鉱物量比が分類を左右し、あらゆる花崗岩類が同量の3種を含むとはしない。\par

\noindent\textbf{4：適切（鉱物群の例示）} 玄武岩の斑晶には普通輝石・紫蘇輝石・斜長石・かんらん石がある。。玄武岩の種類により斑晶に現れ得る鉱物群である。全試料に全て同時に存在するという意味には読まない。紫蘇輝石は古い名称で斜方輝石の一種を指す。\par

\noindent\textbf{5：適切（鉱物群の例示）} 安山岩の斑晶には斜長石・普通輝石・紫蘇輝石・角閃石・磁鉄鉱・黒雲母・かんらん石がある。。安山岩の組成・水分・酸化状態等により現れ得る鉱物群である。かんらん石や黒雲母の存在の可否を、単一の岩種名だけで禁止しない。\par

\paragraph{再利用する考え方} 火成岩の鉱物組合せと有色分類の原典疑義を成因と観察される性質に結び付け、例示と全称命題、主要鉱物と副成分、公式解答と未解決の解釈を区別する。

\paragraph{共通理論} \hyperref[theory:petrology:igneous]{第\ref*{theory:petrology:igneous}節}, \hyperref[theory:petrology:minerals]{第\ref*{theory:petrology:minerals}節}, \hyperref[theory:petrology:exceptions]{第\ref*{theory:petrology:exceptions}節}

\par 参考：\cite{aist_crystal_colors,gsj_igneous2021,usgs_rock_thesaurus}。

\paragraph{関連問題} \hyperref[q:h30-reg-applied-30]{H30-REG-APPLIED-30}

\par 難易度 3（校準初稿）、優先度未評価。

\section{変成作用・鉱物組合せ・組織}

\subsection{令和元年度 問24}
\label{q:r01-reg-applied-24}

\noindent\textbf{原典：} sources/R01/17\_applied\_science.pdf、PDF物理8頁。

\textbf{公式正答：} 3（single）。

\paragraph{問題の要旨} 変成岩の形成要因、変成帯、高圧多形について不適切な記述を選ぶ。

\Needspace{4\baselineskip}\paragraph{確認した条件・記述（独自要約）}\mbox{}\par

\noindent\textbf{方向：} 最も不適切なもの。\par

\noindent\textbf{対象：} 一般的な岩石分類・形成過程。\par

\paragraph{問われる核心} 変成鉱物組合せの温度・圧力・組成依存と多形

\paragraph{必要知識} 変成鉱物組合せの温度・圧力・組成依存と多形。岩石名と生成過程、鉱物組合せ、温度・圧力・全岩組成、一般則と例外の区別。

\paragraph{考え方・正答を選ぶ理由（査読待ち初稿）} 鉱物組合せの平衡は温度・圧力の下で、供給される元素量を制約として自由エネルギーを最小化して決まる。したがって化学組成を無視する3が誤りである。石灰岩と泥岩を同じ温度圧力で変成しても同じ鉱物群になるとは限らない。

\Needspace{4\baselineskip}\paragraph{選択肢の検討（査読待ち初稿）}\mbox{}\par

\noindent\textbf{1：適切（一般的分類）} 火成岩・堆積岩・変成岩のいずれも原岩となる。。既存の変成岩が再び違う条件を受ける複変成もある。\par

\noindent\textbf{2：適切（一般的分類）} 固体中の新鉱物形成と外力による変形が変成作用の要素である。。再結晶・反応と変形を区別する。ただし接触変成などで顕著な変形を必須とはしない。\par

\noindent\textbf{3：不適切} 鉱物組合せは岩石の化学組成によらず温度・圧力のみで決まる。。温度・圧力に加えて全岩組成・流体組成等が必要である。元素が存在しなければその元素を要する鉱物はできない。\par

\noindent\textbf{4：適切（一般的分類）} 西南日本の高圧型変成帯は北から蓮華・周防・三波川の順に分布する。。原画像では蓮華・周防・三波川でありOCRの進華等を採用しない。高圧帯の典型的配置をいう。\par

\noindent\textbf{5：適切（一般的分類）} コース石・ダイヤモンドは石英・石墨のそれぞれ高圧多形である。。同じ化学組成でも結晶構造が違い、高圧側ではより稠密な構造が安定となる。コース石はSiO2、両炭素鉱物はCである。\par

\paragraph{再利用する考え方} 変成鉱物組合せの温度・圧力・組成依存と多形を成因と観察される性質に結び付け、例示と全称命題、主要鉱物と副成分、公式解答と未解決の解釈を区別する。

\paragraph{共通理論} \hyperref[theory:petrology:metamorphism]{第\ref*{theory:petrology:metamorphism}節}, \hyperref[theory:petrology:equilibrium]{第\ref*{theory:petrology:equilibrium}節}, \hyperref[theory:petrology:paired]{第\ref*{theory:petrology:paired}節}

\par 参考：\cite{nelson_assemblages2011,gsj_metamorphic_words,gsj_highpressure2016,hokkaido_polymorph2012}。

\paragraph{関連問題} \hyperref[q:r03-reg-applied-23]{R03-REG-APPLIED-23}

\par 難易度 3（校準初稿）、優先度未評価。

\subsection{令和3年度 問23}
\label{q:r03-reg-applied-23}

\noindent\textbf{原典：} sources/R03/17\_applied\_science.pdf、PDF物理8頁。

\textbf{公式正答：} 3（single）。

\paragraph{問題の要旨} 泥質変成岩の組織と変成度による岩型名を判定する。

\Needspace{4\baselineskip}\paragraph{確認した条件・記述（独自要約）}\mbox{}\par

\noindent\textbf{方向：} 最も不適切なもの。\par

\noindent\textbf{対象：} 一般的な岩石分類・形成過程。\par

\paragraph{問われる核心} 泥質岩の変成度と片理・片麻状組織

\paragraph{必要知識} 泥質岩の変成度と片理・片麻状組織。岩石名と生成過程、鉱物組合せ、温度・圧力・全岩組成、一般則と例外の区別。

\paragraph{考え方・正答を選ぶ理由（査読待ち初稿）} 粒の大きさと面・縞の構造を別々に確認する。片麻岩は粗粒という点だけで識別せず、鉱物の偏在による縞状組織が典型なので、それを否定する3が誤りである。

\Needspace{4\baselineskip}\paragraph{選択肢の検討（査読待ち初稿）}\mbox{}\par

\noindent\textbf{1：適切（一般的分類）} スレートは剥離性の発達した細粒岩で石材にも使われる。。細粒でもよく割れる劈開を持つ点が特徴で、屋根材などへ加工できる。\par

\noindent\textbf{2：適切（一般的分類）} 泥質岩の変成温度上昇に伴いスレート・千枚岩・片岩・片麻岩の順の変化が見られる。。泥質の代表的な進行変成系列である。岩型名は組成・変形・組織にも依存し、温度だけの万能温度計ではない。\par

\noindent\textbf{3：不適切} 片麻岩は粗粒で片状構造や縞状構造が見られない。。片麻岩は明暗の鉱物群が作る縞状組織などを持つ。見られないとの否定が誤り。\par

\noindent\textbf{4：適切（一般的分類）} 片岩は片理の発達した広域変成岩の総称である。。配向した板状鉱物等が面構造を作る。片理は単純な堆積時の層理とは異なる。\par

\noindent\textbf{5：適切（一般的分類）} ホルンフェルスは一般に片理・縞を持たず熱変成で形成される。。方向性に乏しい緻密な再結晶組織が典型である。原岩の構造が残ることまで絶対に排除しない。\par

\paragraph{再利用する考え方} 泥質岩の変成度と片理・片麻状組織を成因と観察される性質に結び付け、例示と全称命題、主要鉱物と副成分、公式解答と未解決の解釈を区別する。

\paragraph{共通理論} \hyperref[theory:petrology:textures]{第\ref*{theory:petrology:textures}節}, \hyperref[theory:petrology:metamorphism]{第\ref*{theory:petrology:metamorphism}節}

\par 参考：\cite{gsj_metamorphic_words,nelson_regional2011}。

\paragraph{関連問題} \hyperref[q:r01-reg-applied-24]{R01-REG-APPLIED-24}

\par 難易度 3（校準初稿）、優先度未評価。

\subsection{令和5年度 問24}
\label{q:r05-reg-applied-24}

\noindent\textbf{原典：} sources/R05/17\_applied\_science.pdf、PDF物理9頁。

\textbf{公式正答：} 5（single）。

\paragraph{問題の要旨} 固体の変成作用、接触・広域変成、弧–海溝系の温度圧力配置を判定する。

\Needspace{4\baselineskip}\paragraph{確認した条件・記述（独自要約）}\mbox{}\par

\noindent\textbf{方向：} 最も不適切なもの。\par

\noindent\textbf{対象：} 一般的な岩石分類・形成過程。\par

\paragraph{問われる核心} 変成過程と対の変成帯の熱・圧力の成因

\paragraph{必要知識} 変成過程と対の変成帯の熱・圧力の成因。岩石名と生成過程、鉱物組合せ、温度・圧力・全岩組成、一般則と例外の区別。

\paragraph{考え方・正答を選ぶ理由（査読待ち初稿）} 深く沈む比較的冷たいプレートは圧力の上昇に比べ温度上昇が小さい。火山弧下のマグマ近傍は熱源による高温が支配的になる。両側の低温高圧／高温低圧を逆にした5が誤りである。大理石は接触でも広域でもできるが、石灰岩の接触変成で形成されるという例示自体は正しい。

\Needspace{4\baselineskip}\paragraph{選択肢の検討（査読待ち初稿）}\mbox{}\par

\noindent\textbf{1：適切（一般的分類）} 岩石が高い温度・圧力の下で固体のまま鉱物や組織を変えると変成岩になる。。変成は固体状態の反応・再結晶・変形を中心とする。鉱物組成の変化と全岩化学組成の変化を区別し、全岩の外部物質交換を必須とはしない。\par

\noindent\textbf{2：適切（一般的分類）} マグマの貫入による周囲の加熱で接触変成岩を生じる。。局所的な熱源により貫入岩周辺に変成の帯を生じる。\par

\noindent\textbf{3：適切（一般的分類）} プレート運動に伴う広範囲の温度・圧力条件で広域変成岩を生じる。。地殻変動や沈み込みにより広い範囲で変成条件が成立する。高温と高圧が常に同じ程度である必要はない。\par

\noindent\textbf{4：適切（一般的分類）} 石灰岩が接触変成を受けると大理石となる。。炭酸塩鉱物の再結晶で大理石となる。広域変成によっても形成され、接触のみという排他的分類ではない。\par

\noindent\textbf{5：不適切} 弧–海溝系では海溝側に高温低圧型、陸側の火山前線下に低温高圧型ができる。。沈み込む冷たいプレート側は相対的な低温高圧、マグマによる熱源側は高温低圧なので逆である。\par

\paragraph{再利用する考え方} 変成過程と対の変成帯の熱・圧力の成因を成因と観察される性質に結び付け、例示と全称命題、主要鉱物と副成分、公式解答と未解決の解釈を区別する。

\paragraph{共通理論} \hyperref[theory:petrology:metamorphism]{第\ref*{theory:petrology:metamorphism}節}, \hyperref[theory:petrology:paired]{第\ref*{theory:petrology:paired}節}

\par 参考：\cite{aist_awa2021,nelson_contact2011,gsj_metamorphic_words}。

\paragraph{関連問題} \hyperref[q:r06-reg-applied-27]{R06-REG-APPLIED-27}, \hyperref[q:r07-reg-applied-23]{R07-REG-APPLIED-23}

\par 難易度 3（校準初稿）、優先度未評価。

\subsection{令和6年度 問27}
\label{q:r06-reg-applied-27}

\noindent\textbf{原典：} sources/R06/17\_applied\_science.pdf、PDF物理9頁。

\textbf{公式正答：} 4（single）。

\par\textbf{検討未確定：} 初回独立解答4と公式4は一致。5の「大理石は接触変成岩に属する」は代表例としては通用するが広域変成でも大理石が形成される。排他的な全称分類か例示かの読みをAstraで確認。

\paragraph{問題の要旨} 変成作用の定義と領家・三波川の変成型を判定する。

\Needspace{4\baselineskip}\paragraph{確認した条件・記述（独自要約）}\mbox{}\par

\noindent\textbf{方向：} 最も不適切なもの。\par

\noindent\textbf{対象：} 一般的な岩石分類・形成過程。\par

\paragraph{問われる核心} 領家・三波川の変成型と大理石分類の適用範囲

\paragraph{必要知識} 領家・三波川の変成型と大理石分類の適用範囲。岩石名と生成過程、鉱物組合せ、温度・圧力・全岩組成、一般則と例外の区別。

\paragraph{考え方・正答を選ぶ理由（査読待ち初稿）} 領家はマグマに関連する高温低圧型、三波川は沈み込みに関連する低温高圧型である。逆にした4が公式正答かつ初回独立解答である。ただし5は大理石の広域変成での形成を省いており、全称的な分類として読む場合には追加の疑義がある。原典の断定をそのまま普遍則として教えず、解釈確認を残す。

\Needspace{4\baselineskip}\paragraph{選択肢の検討（査読待ち初稿）}\mbox{}\par

\noindent\textbf{1：適切} 岩石が高い温度・圧力の下で固体のまま鉱物や組織を変えると変成岩になる。。変成は固体状態の反応・再結晶・変形を中心とする。鉱物組成の変化と全岩化学組成の変化を区別し、全岩の外部物質交換を必須とはしない。\par

\noindent\textbf{2：適切} プレート運動に伴う広範囲の温度・圧力条件で広域変成岩を生じる。。地殻変動や沈み込みにより広い範囲で変成条件が成立する。高温と高圧が常に同じ程度である必要はない。\par

\noindent\textbf{3：適切} マグマの貫入による周囲の加熱で接触変成岩を生じる。。局所的な熱源により貫入岩周辺に変成の帯を生じる。\par

\noindent\textbf{4：不適切} 領家帯は低温高圧、三波川帯は高温低圧である。。領家は高温低圧、三波川は低温高圧の代表であり、対応が逆。\par

\noindent\textbf{5：代表例として適切／全称解釈に疑義} 結晶片岩・片麻岩は広域変成岩、ホルンフェルス・大理石は接触変成岩に属する。。試験では代表的な形成例の対応としているが、大理石は石灰質原岩の広域変成でもできる。全ての大理石を接触起源と断定するなら誤りとなるので、例示と排他的区分を区別する。\par

\paragraph{再利用する考え方} 領家・三波川の変成型と大理石分類の適用範囲を成因と観察される性質に結び付け、例示と全称命題、主要鉱物と副成分、公式解答と未解決の解釈を区別する。

\paragraph{共通理論} \hyperref[theory:petrology:paired]{第\ref*{theory:petrology:paired}節}, \hyperref[theory:petrology:exceptions]{第\ref*{theory:petrology:exceptions}節}

\par 参考：\cite{aist_awa2021,nelson_regional2011,nelson_contact2011}。

\paragraph{関連問題} \hyperref[q:r05-reg-applied-24]{R05-REG-APPLIED-24}

\par 難易度 3（校準初稿）、優先度未評価。

\subsection{令和7年度 問23}
\label{q:r07-reg-applied-23}

\noindent\textbf{原典：} sources/R07/17\_applied\_science.pdf、PDF物理8頁。

\textbf{公式正答：} 4（single）。

\paragraph{問題の要旨} 固体の変成作用、接触・広域変成、弧–海溝系の温度圧力配置を判定する。

\Needspace{4\baselineskip}\paragraph{確認した条件・記述（独自要約）}\mbox{}\par

\noindent\textbf{方向：} 最も不適切なもの。\par

\noindent\textbf{対象：} 一般的な岩石分類・形成過程。\par

\paragraph{問われる核心} 変成過程と対の変成帯の熱・圧力の成因

\paragraph{必要知識} 変成過程と対の変成帯の熱・圧力の成因。岩石名と生成過程、鉱物組合せ、温度・圧力・全岩組成、一般則と例外の区別。

\paragraph{考え方・正答を選ぶ理由（査読待ち初稿）} 深く沈む比較的冷たいプレートは圧力の上昇に比べ温度上昇が小さい。火山弧下のマグマ近傍は熱源による高温が支配的になる。両側の低温高圧／高温低圧を逆にした4が誤りである。大理石は接触でも広域でもできるが、石灰岩の接触変成で形成されるという例示自体は正しい。

\Needspace{4\baselineskip}\paragraph{選択肢の検討（査読待ち初稿）}\mbox{}\par

\noindent\textbf{1：適切（一般的分類）} 岩石が高い温度・圧力の下で固体のまま鉱物や組織を変えると変成岩になる。。変成は固体状態の反応・再結晶・変形を中心とする。鉱物組成の変化と全岩化学組成の変化を区別し、全岩の外部物質交換を必須とはしない。\par

\noindent\textbf{2：適切（一般的分類）} マグマの貫入による周囲の加熱で接触変成岩を生じる。。局所的な熱源により貫入岩周辺に変成の帯を生じる。\par

\noindent\textbf{3：適切（一般的分類）} プレート運動に伴う広範囲の温度・圧力条件で広域変成岩を生じる。。地殻変動や沈み込みにより広い範囲で変成条件が成立する。高温と高圧が常に同じ程度である必要はない。\par

\noindent\textbf{4：不適切} 弧–海溝系では海溝側に高温低圧型、陸側の火山前線下に低温高圧型ができる。。沈み込む冷たいプレート側は相対的な低温高圧、マグマによる熱源側は高温低圧なので逆である。\par

\noindent\textbf{5：適切（一般的分類）} 石灰岩が接触変成を受けると大理石となる。。炭酸塩鉱物の再結晶で大理石となる。広域変成によっても形成され、接触のみという排他的分類ではない。\par

\paragraph{再利用する考え方} 変成過程と対の変成帯の熱・圧力の成因を成因と観察される性質に結び付け、例示と全称命題、主要鉱物と副成分、公式解答と未解決の解釈を区別する。

\paragraph{共通理論} \hyperref[theory:petrology:metamorphism]{第\ref*{theory:petrology:metamorphism}節}, \hyperref[theory:petrology:paired]{第\ref*{theory:petrology:paired}節}

\par 参考：\cite{aist_awa2021,nelson_contact2011,gsj_metamorphic_words}。

\paragraph{関連問題} \hyperref[q:r05-reg-applied-24]{R05-REG-APPLIED-24}

\par 難易度 3（校準初稿）、優先度未評価。

\section{火山砕屑物・堆積構造}

\subsection{令和3年度 問24}
\label{q:r03-reg-applied-24}

\noindent\textbf{原典：} sources/R03/17\_applied\_science.pdf、PDF物理9頁。

\textbf{公式正答：} 4（single）。

\paragraph{問題の要旨} 火山砕屑物の名称・粒径・堆積様式・溶結について不適切な記述を選ぶ。

\Needspace{4\baselineskip}\paragraph{確認した条件・記述（独自要約）}\mbox{}\par

\noindent\textbf{方向：} 最も不適切なもの。\par

\noindent\textbf{対象：} 一般的な岩石分類・形成過程。\par

\paragraph{問われる核心} 火山砕屑物の粒径・支持構造・溶結

\paragraph{必要知識} 火山砕屑物の粒径・支持構造・溶結。岩石名と生成過程、鉱物組合せ、温度・圧力・全岩組成、一般則と例外の区別。

\paragraph{考え方・正答を選ぶ理由（査読待ち初稿）} 名称、粒径、支持構造、堆積後の結合を別軸で比較する。降下軽石の比較的よくそろった粒子支持と、火砕流の灰を含む基質支持を入れ替えた4が不適切である。火山弾を火山岩塊より必ず大粒とする粒径順序へは読み替えない。

\Needspace{4\baselineskip}\paragraph{選択肢の検討（査読待ち初稿）}\mbox{}\par

\noindent\textbf{1：適切（一般的分類）} 多孔質で低い見かけ密度の白〜灰色砕屑物は軽石、暗色のものはスコリアという。。発泡したマグマ由来の砕屑物で、色・組成に差がある。低い見かけ密度と鉱物そのものの密度を混同しない。\par

\noindent\textbf{2：適切（一般的分類）} 粒径区分は細かい方から火山灰・火山礫・火山岩塊および火山弾である。。灰は2 mm未満、火山礫は2〜64 mm、岩塊・火山弾は64 mm超が基準である。岩塊と火山弾は同じ粗粒区分で、主に固結状態・形状で分ける。\par

\noindent\textbf{3：適切（一般的分類）} アグロメレートは火山灰基質に火山弾・本質岩塊を多く含み火口近くに分布する。。火口付近に堆積しやすい大きな本質砕屑物を多く含む火砕岩の説明である。\par

\noindent\textbf{4：不適切} 降下軽石は基質支持、火砕流堆積物は岩片支持である。。典型的な降下軽石は比較的淘汰され粒子同士で支持し、火砕流は細粒灰を含む不淘汰な基質支持が一般的であり、比較方向が逆。流れの相や堆積位置による例外はある。\par

\noindent\textbf{5：適切（一般的分類）} アグルチネートは放出された高温の岩片が溶結したものである。。まだ高温で粘着性のある放出物が集積し結合した岩石である。冷たい砕屑物の機械的固結だけとは異なる。\par

\paragraph{再利用する考え方} 火山砕屑物の粒径・支持構造・溶結を成因と観察される性質に結び付け、例示と全称命題、主要鉱物と副成分、公式解答と未解決の解釈を区別する。

\paragraph{共通理論} \hyperref[theory:petrology:pyroclast]{第\ref*{theory:petrology:pyroclast}節}, \hyperref[theory:petrology:deposition]{第\ref*{theory:petrology:deposition}節}

\par 参考：\cite{usgs_tephra_sizes,usgs_pinatubo1996,usgs_ash_density,usgs_bishop_tuff}。

\paragraph{関連問題} \hyperref[q:r06-reg-applied-26]{R06-REG-APPLIED-26}

\par 難易度 3（校準初稿）、優先度未評価。

\subsection{令和6年度 問26}
\label{q:r06-reg-applied-26}

\noindent\textbf{原典：} sources/R06/17\_applied\_science.pdf、PDF物理9頁。

\textbf{公式正答：} 5（single）。

\paragraph{問題の要旨} 火山砕屑物の名称・粒径・堆積様式・溶結について不適切な記述を選ぶ。

\Needspace{4\baselineskip}\paragraph{確認した条件・記述（独自要約）}\mbox{}\par

\noindent\textbf{方向：} 最も不適切なもの。\par

\noindent\textbf{対象：} 一般的な岩石分類・形成過程。\par

\paragraph{問われる核心} 火山砕屑物の粒径・支持構造・溶結

\paragraph{必要知識} 火山砕屑物の粒径・支持構造・溶結。岩石名と生成過程、鉱物組合せ、温度・圧力・全岩組成、一般則と例外の区別。

\paragraph{考え方・正答を選ぶ理由（査読待ち初稿）} 名称、粒径、支持構造、堆積後の結合を別軸で比較する。降下軽石の比較的よくそろった粒子支持と、火砕流の灰を含む基質支持を入れ替えた5が不適切である。火山弾を火山岩塊より必ず大粒とする粒径順序へは読み替えない。

\Needspace{4\baselineskip}\paragraph{選択肢の検討（査読待ち初稿）}\mbox{}\par

\noindent\textbf{1：適切（一般的分類）} 多孔質で低い見かけ密度の白〜灰色砕屑物は軽石、暗色のものはスコリアという。。発泡したマグマ由来の砕屑物で、色・組成に差がある。低い見かけ密度と鉱物そのものの密度を混同しない。\par

\noindent\textbf{2：適切（一般的分類）} 粒径区分は細かい方から火山灰・火山礫・火山岩塊および火山弾である。。灰は2 mm未満、火山礫は2〜64 mm、岩塊・火山弾は64 mm超が基準である。岩塊と火山弾は同じ粗粒区分で、主に固結状態・形状で分ける。\par

\noindent\textbf{3：適切（一般的分類）} アグロメレートは火山灰基質に火山弾・本質岩塊を多く含み火口近くに分布する。。火口付近に堆積しやすい大きな本質砕屑物を多く含む火砕岩の説明である。\par

\noindent\textbf{4：適切（一般的分類）} アグルチネートは放出された高温の岩片が溶結したものである。。まだ高温で粘着性のある放出物が集積し結合した岩石である。冷たい砕屑物の機械的固結だけとは異なる。\par

\noindent\textbf{5：不適切} 降下軽石は基質支持、火砕流堆積物は岩片支持である。。典型的な降下軽石は比較的淘汰され粒子同士で支持し、火砕流は細粒灰を含む不淘汰な基質支持が一般的であり、比較方向が逆。流れの相や堆積位置による例外はある。\par

\paragraph{再利用する考え方} 火山砕屑物の粒径・支持構造・溶結を成因と観察される性質に結び付け、例示と全称命題、主要鉱物と副成分、公式解答と未解決の解釈を区別する。

\paragraph{共通理論} \hyperref[theory:petrology:pyroclast]{第\ref*{theory:petrology:pyroclast}節}, \hyperref[theory:petrology:deposition]{第\ref*{theory:petrology:deposition}節}

\par 参考：\cite{usgs_tephra_sizes,usgs_pinatubo1996,usgs_ash_density,usgs_bishop_tuff}。

\paragraph{関連問題} \hyperref[q:r03-reg-applied-24]{R03-REG-APPLIED-24}

\par 難易度 3（校準初稿）、優先度未評価。


\chapter{堆積学}
\section{背景：失われた流れを岩石から推定する}
\label{theory:sediment:background}
地層には、形成時の水流そのものは残らない。残るのは粒径、形、鉱物組成、層の幾何、生物の痕跡である。堆積学は、粒子の運搬と沈積という順問題を使い、現在観察できる地層から過去の環境を推定する逆問題である。同じ砂岩でも河川、砂丘、浅海など複数の候補がある。まず観測した特徴を書き、次に成因の仮説を立てる。観測と解釈を一語に押し込めない。\cite{sed_mit_environment}

\section{堆積相、組合せ、累重}
\label{theory:sediment:facies}
堆積相とは、岩相、粒径、堆積構造、化石などの組合せによって周囲と区別する地層の部分である。「細粒砂岩、斜交葉理、植物片」という記載と「河川堆積物」という解釈を分ける。相の組合せは堆積相アソシエーション、上下の配列は堆積相サクセションである。堆積シーケンスは単に相を集めたものではなく、不整合又はそれに対比される整合面を境界とする成因的に関連した地層群をいう。尺度と境界条件が異なる。\cite{sed_mit_environment,sed_walther_oregon}

\section{ワルターの法則：空間的隣接が時間的累重になる}
\label{theory:sediment:walther}
浅い所の砂と深い所の泥が側方に隣接して堆積している状況を考える。環境帯が移動すると、一地点では泥の上へ砂が重なる。連続した整合的・漸移的な累重において、上下の相はもともと側方で隣接する環境に対応するというのがワルターの法則の骨子である。これは全ての上下接触へ使える規則ではない。不整合や大きな侵食面を挟めば、長い時間後の別の環境が重なるため、上下関係だけから同時の側方関係を復元できない。\cite{sed_walther_oregon}

\section{粒子を動かす力と沈降}
\label{theory:sediment:transport}
水流は底面せん断によって粒子を動かし、乱流は沈降する粒子を再び上へ運ぶ。粒子は底面に沿う掃流と、水中の浮遊に分けて考えられる。球形の孤立粒子が低レイノルズ数で沈む理想条件では、重力と浮力の差を粘性抵抗と釣り合わせると
\[
 w_s=\frac{(\rho_s-\rho_f)gd^2}{18\eta}
\]
となる。\(w_s\) は終端沈降速度、\(d\) は直径、\(\rho_s,\rho_f\) は粒子と流体の密度、\(\eta\) は粘性係数である。粒径が大きいほど速く沈むという見通しを与えるが、粗い砂礫、非球形、高濃度の群れへこの式を直接代入しない。堆積構造は沈降だけでなく、底面の侵食と粒子供給の時間変化にも依存する。

\section{河川の二つの端成分と洪水}
\label{theory:sediment:river}
蛇行河川では湾曲した流路が側方へ移動し、内岸のポイントバーに砂がたまり、流路外の氾濫原には泥がたまりやすい。網状河川では複数の流路が砂礫州を挟み、粗粒負荷や流量変動の大きい条件が典型となる。しかし二つは排他的な分類ではなく、中間形態もある。勾配だけで流路を決めず、河岸の強さ、供給量、粒度を合わせて見る。\cite{sed_mit_environment}

\section{正級化と逆級化を区別する}
\label{theory:sediment:grading}
一枚の層の中で上へ細粒化するのが正級化、上へ粗粒化するのが逆級化である。流れの減衰と粗粒子の先行沈降は正級化の典型機構である。洪水の立上がりでは運搬可能粒径が時間とともに増し、氾濫原へ届く砂が上方粗粒化することがある。だが逆級化は土石流などにも見られ、河川だけの指標ではない。一枚のイベント層内の級化と、複数の相が重なる上方粗粒化の累重を区別する。

\section{デルタ、海進、海退と河口内湾}
\label{theory:sediment:delta}
デルタは河川が海や湖へ土砂を供給してできる河口の堆積体である。前進すると海側の細粒層の上へ浅い粗粒層が重なる。海退は海岸線が海側へ移ることであり、海水準の低下だけを意味しない。相対海水準が上がっても、土砂供給が新しく生じる堆積空間を埋める速さを上回れば前進する通常海退となる。相対海水準低下による強制海退と分ける。海進で河口や谷が沈水して形成されるエスチュアリーでは、河川と潮汐の双方が働く。\cite{sed_walther_oregon}

\section{移動するリップルと斜交層理}
\label{theory:sediment:cross}
流れが砂の床形を移動させると、粒子は緩い上流側を運ばれ、下流側の斜面へたまる。その斜面が保存されると斜交葉理又は斜交層理となる。小さな床形のリップルも大きなデューンも同じ幾何の考え方を使う。平板型とトラフ型の違いは床形の三次元性や切込みの形に関わり、別々の流体力学法則が必要という意味ではない。古流向は層の傾斜を補正し、複数断面から求める。\cite{sed_houston_structures}

\section{ヘリンボーン：往復流と断面の錯覚}
\label{theory:sediment:herringbone}
逆向きに傾く斜交層の組は、潮流のように向きが反転する流れを考える手掛かりとなる。ヘリンボーンは蛇行河川の一方向流の代表構造ではない。ただし一断面で逆向きに見えることだけでは、真の180度反転を保証しない。三次元の床形を斜めに切った断面でも似た形が見える。再活動面や泥の薄い覆い、古流向の分布を併用する。\cite{sed_mit_environment}

\section{ハンモック状斜交層理と暴浪}
\label{theory:sediment:hcs}
ゆるい凸部と凹部へ沿う低角度の層理がハンモック状斜交層理の特徴である。典型的にはストーム時の波浪による振動流、又は一方向成分を含む複合流を考える。堆積構造の形の名前と成因の名前を同一視しない。MITの講義では広い意味の同構造が多成因であることと、反砂堆に関わる例も扱う。したがって試験で「一方向流の代表として最も不適切」と選ぶことは、「一方向成分のある全ての条件で絶対に形成されない」という断定ではない。\cite{sed_mit_cross}

\section{混濁流とタービダイト}
\label{theory:sediment:turbidite}
混濁流は浮遊粒子によって周囲の水より密度が高くなった流れが、重力で斜面に沿って進む現象である。乱流が粒子の沈降に対抗する。沈積した層をタービダイトと呼び、流れが弱まると正級化が生じやすい。バウマシーケンスは級化砂、平行葉理、リップル葉理、細粒部などの一イベント内の典型配列である。全区分が毎回そろうわけではなく、侵食や堆積しない区分の省略がある。海底地すべりは発生源の一つであって、全ての混濁流の必要条件ではない。\cite{sed_mit_environment}

\section{未固結変形、荷重痕、火炎、脱水}
\label{theory:sediment:soft}
まだ固まらない泥の上へ砂が急速にたまると、重い砂が下へ突出する荷重痕と、泥が上へ伸びる火炎構造ができることがある。堆積後の未固結段階の変形であるため、粒子をたまらせる過程の葉理・級化・リップル・流痕と区別する。脱水による皿状構造も、層の記載だけでなく間隙水の逃げ方を考える対象である。液状化を伴う場合、有効応力 \(\sigma'=\sigma-u\) の低下が変形しやすさを説明する。\(u\) は間隙水圧である。火炎があるだけで地震起源を断定せず、過荷重などの候補と比較する。\cite{sed_soft_deformation,sed_houston_structures}

\section{地層の上下判定}
\label{theory:sediment:wayup}
現在の鉛直方向と形成時の上は、褶曲や転倒で一致しない。荷重痕は砂が下の泥へ突き出る向き、皿状構造は一般的な上に凹の形、流痕は侵食された面とその上の層が埋めた凸型の底面を手掛かりとする。生痕も巣穴の開口や分岐・埋積の幾何が分かる場合に上下面を示すが、名称だけでは全ての生痕に向きを付けられない。節理面の羽毛状構造は割れ目が広がった方向を示し、堆積時の上下を示すものではない。\cite{sed_houston_structures,sed_plumose_michigan}

\section{環境指標と分析装置の観測量}
\label{theory:sediment:proxy}
黄鉄鉱 \(\mathrm{FeS_2}\) の形成には鉄・硫黄の供給と還元状態が関わる。海成粘土を酸化・乾燥して懸濁すると硫酸塩の生成などが電気伝導度へ影響するため、古環境推定に利用される。だが作成手順、濃度、温度、後の溶脱も値を変えるため、海成と淡水成の万能な閾値を設けない。\cite{sed_gsj_ec2011}
蛍光X線分析は特性X線のエネルギーと強度から元素を調べる方法であり、内部の堆積構造を直接撮影する方法ではない。SEMは粒子表面の微細形状を観察する。運搬で生じた摩耗や衝突の痕跡は風成作用を検討する材料となるが、溶食や再運搬でも変わる。装置の名前、観測量、環境解釈という三段を区別する。\cite{sed_xrf_jaea,sed_sem_shimane}

\section{津波の地質記録：形成と保存は別である}
\label{theory:sediment:tsunami}
津波は海底や海岸の粒子を陸側へ運び、戻り流れは再び海側へ運ぶ。複数波、再侵食、滞留時間の差によって、塊状、葉理、級化、複数ユニットなどが生じる。一方、近年のUSGSデータ集では構造が稀な例も多い。したがって「単純が常に多い」「複雑が常に多い」のどちらにも一般化しない。\cite{sed_tsunami_fujino,sed_tsunami_usgs}
保存は形成後の侵食と再攪乱に依存する。静かな湿地や内湾の細粒背景層に異質な砂層が挟まれば調べやすいが、砂層だけで津波と認定しない。微化石、供給域、広がり、複数の構造、年代を合わせる。生物遺骸は供給源の水深を検討する手掛かりであり、混合・再堆積の影響を受ける。岩礁や礫の多い場所も記録を持ち得るが、識別と継続保存が難しい。\cite{sed_tsunami_2002,sed_tsunami_gsj_survey}

\section{典型的な問われ方}
過去問では、構造と流れの対応、初生と二次の区別、上下判定、分析手法の観測量、相・累重・シーケンスの名称の入替えが問われる。正答番号は同じテーマでも年度で変わる。全肢を条件付きで読むことに加え、「最も不適切」と科学的な全称命題を分ける。頻度は全問の分類確認後に集計する。

\section{過去問と個別の解法}
\section{流れ・波浪と斜交層理}

\subsection{平成30年度 問25}
\label{q:h30-reg-applied-25}

\noindent\textbf{原典：} sources/H30/17\_applied\_science.pdf、PDF物理8頁。

\textbf{公式正答：} 2（single）。

\par\textbf{検討未確定：} 海底地すべりを全タービダイトの起源、全Bouma区分を必須とする読みは過度である。HCSにも多成因がある。公式/独立2を保持して典型対応の範囲を部分校正に残す。

\paragraph{問題の要旨} 流れ・波浪と堆積構造を対応させ、蛇行河川と往復する潮流の構造を区別する。

\Needspace{4\baselineskip}\paragraph{確認した条件・記述（独自要約）}\mbox{}\par

\noindent\textbf{判断方向：} 最も不適切・誤りを選ぶ\par

\noindent\textbf{原典照合：} 15原典画像で18問の問題幹、全5肢、否定、専門語、化学式、候補順を確認。\par

\noindent\textbf{適用条件：} 流れと構造の典型対応は一意成因を保証しない。用語、時間間隙、保存、試料処理の条件を個別に示す。\par

\paragraph{問われる核心} 流れ・波浪と堆積構造を対応させ、蛇行河川と往復する潮流の構造を区別する。

\paragraph{必要知識} 床形の移動と斜交、往復流、振動・複合流、密度流、流痕、典型成因の適用条件。

\paragraph{考え方・正答を選ぶ理由（査読待ち初稿）} 床形の前進は斜交層理を保存し、流れの侵食は流痕を作る。混濁流の減衰は一般に級化したタービダイトを残す。ハンモック状斜交層理では暴浪時の振動流又は複合流が代表となる。ヘリンボーンを蛇行河川の主な構造に割り当てる対応が最も不適切であり、本問の番号は2である。潮流の反転を典型とするが、一断面の逆傾斜だけで真の180度反転を断定しない。原典のタービダイト説明も全例の起源を限定せず読む。

\Needspace{4\baselineskip}\paragraph{選択肢の検討（査読待ち初稿）}\mbox{}\par

\noindent\textbf{1：適切（条件を含めて検討）} 床形の移動で前面へ層が付き斜交層理となる。移動する床形の下流側斜面への堆積が内部の斜交を保存するという説明に対応する。\par

\noindent\textbf{2：誤り・最も不適切} 蛇行河川の代表としてヘリンボーンを挙げる。相反する流向で傾斜の異なる層が重なる構造は潮流を典型とし、蛇行の一方向流の主な構造という対応が違う。\par

\noindent\textbf{3：適切（条件を含めて検討）} 暴浪の波浪でハンモック状斜交層理が形成する。ストームの振動流・複合流が代表的な形成条件である。形の名称から全例の単一成因を断定しない。\par

\noindent\textbf{4：適切（条件を含めて検討）} 陸棚斜面の海底地すべりに起源を持つ混濁流の堆積物がタービダイトで級化・Boumaが特徴。海底地すべりは混濁流の発生源の一つである。級化とBoumaは典型だが全ての混濁流の起源や全区分の存在を限定する定義ではない。\par

\noindent\textbf{5：適切（条件を含めて検討）} 流れの小さな削込みを流痕とする。流れが底面を侵食した痕跡であり、充填された型が上位層底面に保存される場合もある。\par

\paragraph{再利用する考え方} 観察できる構造と成因・環境の解釈を区別する。代表モデルに適用条件を添え、単一の形や測定値から唯一の環境を断定しない。

\paragraph{共通理論} \hyperref[theory:sediment:cross]{第\ref*{theory:sediment:cross}節}, \hyperref[theory:sediment:herringbone]{第\ref*{theory:sediment:herringbone}節}, \hyperref[theory:sediment:hcs]{第\ref*{theory:sediment:hcs}節}, \hyperref[theory:sediment:turbidite]{第\ref*{theory:sediment:turbidite}節}

\par 参考：\cite{sed_mit_environment,sed_mit_cross,sed_houston_structures}。

\paragraph{関連問題} \hyperref[q:r04-reg-applied-28]{R04-REG-APPLIED-28}

\par 難易度 3（校準初稿）、優先度未評価。

\subsection{令和4年度 問28}
\label{q:r04-reg-applied-28}

\noindent\textbf{原典：} sources/R04/17\_applied\_science.pdf、PDF物理10頁。

\textbf{公式正答：} 1（single）。

\paragraph{問題の要旨} 流れ・波浪と堆積構造を対応させ、蛇行河川と往復する潮流の構造を区別する。

\Needspace{4\baselineskip}\paragraph{確認した条件・記述（独自要約）}\mbox{}\par

\noindent\textbf{判断方向：} 最も不適切・誤りを選ぶ\par

\noindent\textbf{原典照合：} 15原典画像で18問の問題幹、全5肢、否定、専門語、化学式、候補順を確認。\par

\noindent\textbf{適用条件：} 流れと構造の典型対応は一意成因を保証しない。用語、時間間隙、保存、試料処理の条件を個別に示す。\par

\paragraph{問われる核心} 流れ・波浪と堆積構造を対応させ、蛇行河川と往復する潮流の構造を区別する。

\paragraph{必要知識} 床形の移動と斜交、往復流、振動・複合流、密度流、流痕、典型成因の適用条件。

\paragraph{考え方・正答を選ぶ理由（査読待ち初稿）} 床形の前進は斜交層理を保存し、流れの侵食は流痕を作る。混濁流の減衰は一般に級化したタービダイトを残す。ハンモック状斜交層理では暴浪時の振動流又は複合流が代表となる。ヘリンボーンを蛇行河川の主な構造に割り当てる対応が最も不適切であり、本問の番号は1である。潮流の反転を典型とするが、一断面の逆傾斜だけで真の180度反転を断定しない。原典のタービダイト説明も全例の起源を限定せず読む。

\Needspace{4\baselineskip}\paragraph{選択肢の検討（査読待ち初稿）}\mbox{}\par

\noindent\textbf{1：誤り・最も不適切} 蛇行河川の代表としてヘリンボーンを挙げる。相反する流向で傾斜の異なる層が重なる構造は潮流を典型とし、蛇行の一方向流の主な構造という対応が違う。\par

\noindent\textbf{2：適切（条件を含めて検討）} 暴浪の波浪でハンモック状斜交層理が形成する。ストームの振動流・複合流が代表的な形成条件である。形の名称から全例の単一成因を断定しない。\par

\noindent\textbf{3：適切（条件を含めて検討）} 床形の移動で前面へ層が付き斜交層理となる。移動する床形の下流側斜面への堆積が内部の斜交を保存するという説明に対応する。\par

\noindent\textbf{4：適切（条件を含めて検討）} 流れの小さな削込みを流痕とする。流れが底面を侵食した痕跡であり、充填された型が上位層底面に保存される場合もある。\par

\noindent\textbf{5：適切（条件を含めて検討）} 混濁流がタービダイトを作り一般に級化する。密度差で流れ、乱流が粒子を支持し、流れの減衰と沈降で級化しやすい。全例で全Bouma区分がそろうとは限らない。\par

\paragraph{再利用する考え方} 観察できる構造と成因・環境の解釈を区別する。代表モデルに適用条件を添え、単一の形や測定値から唯一の環境を断定しない。

\paragraph{共通理論} \hyperref[theory:sediment:cross]{第\ref*{theory:sediment:cross}節}, \hyperref[theory:sediment:herringbone]{第\ref*{theory:sediment:herringbone}節}, \hyperref[theory:sediment:hcs]{第\ref*{theory:sediment:hcs}節}, \hyperref[theory:sediment:turbidite]{第\ref*{theory:sediment:turbidite}節}

\par 参考：\cite{sed_mit_environment,sed_mit_cross,sed_houston_structures}。

\paragraph{関連問題} \hyperref[q:h30-reg-applied-25]{H30-REG-APPLIED-25}, \hyperref[q:r07-reg-applied-27]{R07-REG-APPLIED-27}

\par 難易度 3（校準初稿）、優先度未評価。

\subsection{令和5年度 問31}
\label{q:r05-reg-applied-31}

\noindent\textbf{原典：} sources/R05/17\_applied\_science.pdf、PDF物理11頁。

\textbf{公式正答：} 2（single）。

\par\textbf{検討未確定：} 広義のHCSには多成因があり、一方向流で絶対に不可能という全称は未支持である。初回/公式2は典型モデルの比較として保持し、用語範囲をAstraへ残す。

\paragraph{問題の要旨} 一方向流で形成し得る層理と、波浪の振動・複合流が典型となる構造を比較する。

\Needspace{4\baselineskip}\paragraph{確認した条件・記述（独自要約）}\mbox{}\par

\noindent\textbf{判断方向：} 最も不適切・誤りを選ぶ\par

\noindent\textbf{原典照合：} 15原典画像で18問の問題幹、全5肢、否定、専門語、化学式、候補順を確認。\par

\noindent\textbf{適用条件：} 流れと構造の典型対応は一意成因を保証しない。用語、時間間隙、保存、試料処理の条件を個別に示す。\par

\paragraph{問われる核心} 一方向流で形成し得る層理と、波浪の振動・複合流が典型となる構造を比較する。

\paragraph{必要知識} 一方向の床形と層理、振動流と複合流、形の分類と成因の非一意性。

\paragraph{考え方・正答を選ぶ理由（査読待ち初稿）} 四つは一方向の掃流や床形の移動で形成し得るのに対し、ハンモック状斜交層理は波浪の振動流・複合流の代表なので2を選ぶ。ただしMITは広義のHCSに反砂堆を伴う例を含む多成因を挙げる。「最も不適切」という試験の代表比較を、全ての一方向流条件で絶対に形成されないという証明へ置き換えない。

\Needspace{4\baselineskip}\paragraph{選択肢の検討（査読待ち初稿）}\mbox{}\par

\noindent\textbf{1：適切（条件を含めて検討）} トラフ型斜交層理。三次元の床形の移動・侵食と堆積により一方向流で形成し得る。\par

\noindent\textbf{2：誤り・最も不適切} ハンモック状斜交層理。典型は暴浪の振動・複合流であり、本問で一方向流の代表として最も不適切である。ただし広義のHCSには複数成因が報告される。\par

\noindent\textbf{3：適切（条件を含めて検討）} フォーセット層理。床形の下流側斜面への堆積などにより一方向流で形成し得る。\par

\noindent\textbf{4：適切（条件を含めて検討）} 平行層理。床面に沿うシート状の運搬・堆積などで形成し得る。一方向流だけの固有指標という意味ではない。\par

\noindent\textbf{5：適切（条件を含めて検討）} 平板型斜交層理。二次元に近い床形の移動により一方向流で形成し得る。\par

\paragraph{再利用する考え方} 観察できる構造と成因・環境の解釈を区別する。代表モデルに適用条件を添え、単一の形や測定値から唯一の環境を断定しない。

\paragraph{共通理論} \hyperref[theory:sediment:cross]{第\ref*{theory:sediment:cross}節}, \hyperref[theory:sediment:hcs]{第\ref*{theory:sediment:hcs}節}

\par 参考：\cite{sed_mit_cross,sed_houston_structures}。

\par 難易度 3（校準初稿）、優先度未評価。

\subsection{令和7年度 問27}
\label{q:r07-reg-applied-27}

\noindent\textbf{原典：} sources/R07/17\_applied\_science.pdf、PDF物理10頁。

\textbf{公式正答：} 3（single）。

\paragraph{問題の要旨} 流れ・波浪と堆積構造を対応させ、蛇行河川と往復する潮流の構造を区別する。

\Needspace{4\baselineskip}\paragraph{確認した条件・記述（独自要約）}\mbox{}\par

\noindent\textbf{判断方向：} 最も不適切・誤りを選ぶ\par

\noindent\textbf{原典照合：} 15原典画像で18問の問題幹、全5肢、否定、専門語、化学式、候補順を確認。\par

\noindent\textbf{適用条件：} 流れと構造の典型対応は一意成因を保証しない。用語、時間間隙、保存、試料処理の条件を個別に示す。\par

\paragraph{問われる核心} 流れ・波浪と堆積構造を対応させ、蛇行河川と往復する潮流の構造を区別する。

\paragraph{必要知識} 床形の移動と斜交、往復流、振動・複合流、密度流、流痕、典型成因の適用条件。

\paragraph{考え方・正答を選ぶ理由（査読待ち初稿）} 床形の前進は斜交層理を保存し、流れの侵食は流痕を作る。混濁流の減衰は一般に級化したタービダイトを残す。ハンモック状斜交層理では暴浪時の振動流又は複合流が代表となる。ヘリンボーンを蛇行河川の主な構造に割り当てる対応が最も不適切であり、本問の番号は3である。潮流の反転を典型とするが、一断面の逆傾斜だけで真の180度反転を断定しない。原典のタービダイト説明も全例の起源を限定せず読む。

\Needspace{4\baselineskip}\paragraph{選択肢の検討（査読待ち初稿）}\mbox{}\par

\noindent\textbf{1：適切（条件を含めて検討）} 流れの小さな削込みを流痕とする。流れが底面を侵食した痕跡であり、充填された型が上位層底面に保存される場合もある。\par

\noindent\textbf{2：適切（条件を含めて検討）} 床形の移動で前面へ層が付き斜交層理となる。移動する床形の下流側斜面への堆積が内部の斜交を保存するという説明に対応する。\par

\noindent\textbf{3：誤り・最も不適切} 蛇行河川の代表としてヘリンボーンを挙げる。相反する流向で傾斜の異なる層が重なる構造は潮流を典型とし、蛇行の一方向流の主な構造という対応が違う。\par

\noindent\textbf{4：適切（条件を含めて検討）} 混濁流がタービダイトを作り一般に級化する。密度差で流れ、乱流が粒子を支持し、流れの減衰と沈降で級化しやすい。全例で全Bouma区分がそろうとは限らない。\par

\noindent\textbf{5：適切（条件を含めて検討）} 暴浪の波浪でハンモック状斜交層理が形成する。ストームの振動流・複合流が代表的な形成条件である。形の名称から全例の単一成因を断定しない。\par

\paragraph{再利用する考え方} 観察できる構造と成因・環境の解釈を区別する。代表モデルに適用条件を添え、単一の形や測定値から唯一の環境を断定しない。

\paragraph{共通理論} \hyperref[theory:sediment:cross]{第\ref*{theory:sediment:cross}節}, \hyperref[theory:sediment:herringbone]{第\ref*{theory:sediment:herringbone}節}, \hyperref[theory:sediment:hcs]{第\ref*{theory:sediment:hcs}節}, \hyperref[theory:sediment:turbidite]{第\ref*{theory:sediment:turbidite}節}

\par 参考：\cite{sed_mit_environment,sed_mit_cross,sed_houston_structures}。

\paragraph{関連問題} \hyperref[q:r04-reg-applied-28]{R04-REG-APPLIED-28}

\par 難易度 3（校準初稿）、優先度未評価。

\section{堆積相と河川・デルタ・河口システム}

\subsection{平成23年度 問25}
\label{q:h23-reg-applied-25}

\noindent\textbf{原典：} sources/H23/17\_applied\_science.pdf、PDF物理8頁。

\textbf{公式正答：} 5（single）。

\par\textbf{検討未確定：} デルタを海退期と呼ぶ範囲、氾濫原砂の逆級化の代表性は要限定である。逆級化は他成因にもあり河川に一意ではない。全河川・全デルタの排他的な診断とはしない。独立初回と公式の蛇行/網状混同肢を保持する。

\paragraph{問題の要旨} 河川と河口の堆積システムを比較し、蛇行と網状の典型条件を区別する。

\Needspace{4\baselineskip}\paragraph{確認した条件・記述（独自要約）}\mbox{}\par

\noindent\textbf{判断方向：} 最も不適切・誤りを選ぶ\par

\noindent\textbf{原典照合：} 15原典画像で18問の問題幹、全5肢、否定、専門語、化学式、候補順を確認。\par

\noindent\textbf{適用条件：} 流れと構造の典型対応は一意成因を保証しない。用語、時間間隙、保存、試料処理の条件を個別に示す。\par

\paragraph{問われる核心} 河川と河口の堆積システムを比較し、蛇行と網状の典型条件を区別する。

\paragraph{必要知識} 河川端成分、氾濫原とポイントバー、洪水時の級化、デルタ前進、相対海水準と供給、河口沈水。

\paragraph{考え方・正答を選ぶ理由（査読待ち初稿）} 五つは河川の端成分と河口堆積システムの対応を問う。蛇行の典型を、急勾配・大きな流量変動・粗粒子の大きな負荷で説明する肢を除く。この組合せは網状の典型に対応する。ただし実際の流路型には河岸の強度や供給も関わる。デルタは前進する堆積体で、海退と相対海水準低下は同義ではない。逆級化も洪水立上がりで形成し得るが、非一意の指標である。原典の典型モデルと、全称命題としての限界を分けて解く。本問の該当番号は5である。

\Needspace{4\baselineskip}\paragraph{選択肢の検討（査読待ち初稿）}\mbox{}\par

\noindent\textbf{1：適切（条件を含めて検討）} 河川を蛇行・網状の二端成分で考える。端成分モデルとして妥当である。中間形態や両特徴を持つ河川を除外する絶対二分ではない。\par

\noindent\textbf{2：適切（条件を含めて検討）} 海湖側へ突出するデルタを海退期のシステムとする。デルタの前進が海退を示す場合の説明である。海退を海水準低下と同義にせず、デルタの存在だけで全時期を決めない。\par

\noindent\textbf{3：適切（条件を含めて検討）} 海進で河口が沈水する内湾をエスチュアリーとする。沈水谷を河川と潮汐が埋める河口内湾というモデルに対応する。海進する全河口の唯一の形ではない。\par

\noindent\textbf{4：適切（条件を含めて検討）} 氾濫原の砂岩に見られる逆級化を河川の特徴とする。洪水の増水過程では砂が上方粗粒化する場合がある。ただし逆級化は河川に固有ではなく、全氾濫原砂がそうなるという意味ではない。\par

\noindent\textbf{5：誤り・最も不適切} 蛇行河川を流量変動・勾配・粗粒負荷の大きい河川とする。これは網状河川の典型条件との混同である。蛇行は流路の曲率と側方移動、ポイントバーと泥質氾濫原で捉える。\par

\paragraph{再利用する考え方} 観察できる構造と成因・環境の解釈を区別する。代表モデルに適用条件を添え、単一の形や測定値から唯一の環境を断定しない。

\paragraph{共通理論} \hyperref[theory:sediment:river]{第\ref*{theory:sediment:river}節}, \hyperref[theory:sediment:grading]{第\ref*{theory:sediment:grading}節}, \hyperref[theory:sediment:delta]{第\ref*{theory:sediment:delta}節}

\par 参考：\cite{sed_mit_environment,sed_walther_oregon}。

\paragraph{関連問題} \hyperref[q:h29-reg-applied-25]{H29-REG-APPLIED-25}, \hyperref[q:r06-reg-applied-31]{R06-REG-APPLIED-31}

\par 難易度 3（校準初稿）、優先度未評価。

\subsection{平成29年度 問25}
\label{q:h29-reg-applied-25}

\noindent\textbf{原典：} sources/H29/17\_applied\_science.pdf、PDF物理8頁。

\textbf{公式正答：} 5（single）。

\par\textbf{検討未確定：} デルタを海退期と呼ぶ範囲、氾濫原砂の逆級化の代表性は要限定である。逆級化は他成因にもあり河川に一意ではない。全河川・全デルタの排他的な診断とはしない。独立初回と公式の蛇行/網状混同肢を保持する。

\paragraph{問題の要旨} 河川と河口の堆積システムを比較し、蛇行と網状の典型条件を区別する。

\Needspace{4\baselineskip}\paragraph{確認した条件・記述（独自要約）}\mbox{}\par

\noindent\textbf{判断方向：} 最も不適切・誤りを選ぶ\par

\noindent\textbf{原典照合：} 15原典画像で18問の問題幹、全5肢、否定、専門語、化学式、候補順を確認。\par

\noindent\textbf{適用条件：} 流れと構造の典型対応は一意成因を保証しない。用語、時間間隙、保存、試料処理の条件を個別に示す。\par

\paragraph{問われる核心} 河川と河口の堆積システムを比較し、蛇行と網状の典型条件を区別する。

\paragraph{必要知識} 河川端成分、氾濫原とポイントバー、洪水時の級化、デルタ前進、相対海水準と供給、河口沈水。

\paragraph{考え方・正答を選ぶ理由（査読待ち初稿）} 五つは河川の端成分と河口堆積システムの対応を問う。蛇行の典型を、急勾配・大きな流量変動・粗粒子の大きな負荷で説明する肢を除く。この組合せは網状の典型に対応する。ただし実際の流路型には河岸の強度や供給も関わる。デルタは前進する堆積体で、海退と相対海水準低下は同義ではない。逆級化も洪水立上がりで形成し得るが、非一意の指標である。原典の典型モデルと、全称命題としての限界を分けて解く。本問の該当番号は5である。

\Needspace{4\baselineskip}\paragraph{選択肢の検討（査読待ち初稿）}\mbox{}\par

\noindent\textbf{1：適切（条件を含めて検討）} 河川を蛇行・網状の二端成分で考える。端成分モデルとして妥当である。中間形態や両特徴を持つ河川を除外する絶対二分ではない。\par

\noindent\textbf{2：適切（条件を含めて検討）} 海湖側へ突出するデルタを海退期のシステムとする。デルタの前進が海退を示す場合の説明である。海退を海水準低下と同義にせず、デルタの存在だけで全時期を決めない。\par

\noindent\textbf{3：適切（条件を含めて検討）} 海進で河口が沈水する内湾をエスチュアリーとする。沈水谷を河川と潮汐が埋める河口内湾というモデルに対応する。海進する全河口の唯一の形ではない。\par

\noindent\textbf{4：適切（条件を含めて検討）} 氾濫原の砂岩に見られる逆級化を河川の特徴とする。洪水の増水過程では砂が上方粗粒化する場合がある。ただし逆級化は河川に固有ではなく、全氾濫原砂がそうなるという意味ではない。\par

\noindent\textbf{5：誤り・最も不適切} 蛇行河川を流量変動・勾配・粗粒負荷の大きい河川とする。これは網状河川の典型条件との混同である。蛇行は流路の曲率と側方移動、ポイントバーと泥質氾濫原で捉える。\par

\paragraph{再利用する考え方} 観察できる構造と成因・環境の解釈を区別する。代表モデルに適用条件を添え、単一の形や測定値から唯一の環境を断定しない。

\paragraph{共通理論} \hyperref[theory:sediment:river]{第\ref*{theory:sediment:river}節}, \hyperref[theory:sediment:grading]{第\ref*{theory:sediment:grading}節}, \hyperref[theory:sediment:delta]{第\ref*{theory:sediment:delta}節}

\par 参考：\cite{sed_mit_environment,sed_walther_oregon}。

\paragraph{関連問題} \hyperref[q:h23-reg-applied-25]{H23-REG-APPLIED-25}, \hyperref[q:r05-reg-applied-30]{R05-REG-APPLIED-30}

\par 難易度 3（校準初稿）、優先度未評価。

\subsection{令和5年度 問30}
\label{q:r05-reg-applied-30}

\noindent\textbf{原典：} sources/R05/17\_applied\_science.pdf、PDF物理11頁。

\textbf{公式正答：} 3（single）。

\par\textbf{検討未確定：} デルタを海退期と呼ぶ範囲、氾濫原砂の逆級化の代表性は要限定である。逆級化は他成因にもあり河川に一意ではない。全河川・全デルタの排他的な診断とはしない。独立初回と公式の蛇行/網状混同肢を保持する。

\paragraph{問題の要旨} 河川と河口の堆積システムを比較し、蛇行と網状の典型条件を区別する。

\Needspace{4\baselineskip}\paragraph{確認した条件・記述（独自要約）}\mbox{}\par

\noindent\textbf{判断方向：} 最も不適切・誤りを選ぶ\par

\noindent\textbf{原典照合：} 15原典画像で18問の問題幹、全5肢、否定、専門語、化学式、候補順を確認。\par

\noindent\textbf{適用条件：} 流れと構造の典型対応は一意成因を保証しない。用語、時間間隙、保存、試料処理の条件を個別に示す。\par

\paragraph{問われる核心} 河川と河口の堆積システムを比較し、蛇行と網状の典型条件を区別する。

\paragraph{必要知識} 河川端成分、氾濫原とポイントバー、洪水時の級化、デルタ前進、相対海水準と供給、河口沈水。

\paragraph{考え方・正答を選ぶ理由（査読待ち初稿）} 五つは河川の端成分と河口堆積システムの対応を問う。蛇行の典型を、急勾配・大きな流量変動・粗粒子の大きな負荷で説明する肢を除く。この組合せは網状の典型に対応する。ただし実際の流路型には河岸の強度や供給も関わる。デルタは前進する堆積体で、海退と相対海水準低下は同義ではない。逆級化も洪水立上がりで形成し得るが、非一意の指標である。原典の典型モデルと、全称命題としての限界を分けて解く。本問の該当番号は3である。

\Needspace{4\baselineskip}\paragraph{選択肢の検討（査読待ち初稿）}\mbox{}\par

\noindent\textbf{1：適切（条件を含めて検討）} 河川を蛇行・網状の二端成分で考える。端成分モデルとして妥当である。中間形態や両特徴を持つ河川を除外する絶対二分ではない。\par

\noindent\textbf{2：適切（条件を含めて検討）} 氾濫原の砂岩に見られる逆級化を河川の特徴とする。洪水の増水過程では砂が上方粗粒化する場合がある。ただし逆級化は河川に固有ではなく、全氾濫原砂がそうなるという意味ではない。\par

\noindent\textbf{3：誤り・最も不適切} 蛇行河川を流量変動・勾配・粗粒負荷の大きい河川とする。これは網状河川の典型条件との混同である。蛇行は流路の曲率と側方移動、ポイントバーと泥質氾濫原で捉える。\par

\noindent\textbf{4：適切（条件を含めて検討）} 海湖側へ突出するデルタを海退期のシステムとする。デルタの前進が海退を示す場合の説明である。海退を海水準低下と同義にせず、デルタの存在だけで全時期を決めない。\par

\noindent\textbf{5：適切（条件を含めて検討）} 海進で河口が沈水する内湾をエスチュアリーとする。沈水谷を河川と潮汐が埋める河口内湾というモデルに対応する。海進する全河口の唯一の形ではない。\par

\paragraph{再利用する考え方} 観察できる構造と成因・環境の解釈を区別する。代表モデルに適用条件を添え、単一の形や測定値から唯一の環境を断定しない。

\paragraph{共通理論} \hyperref[theory:sediment:river]{第\ref*{theory:sediment:river}節}, \hyperref[theory:sediment:grading]{第\ref*{theory:sediment:grading}節}, \hyperref[theory:sediment:delta]{第\ref*{theory:sediment:delta}節}

\par 参考：\cite{sed_mit_environment,sed_walther_oregon}。

\paragraph{関連問題} \hyperref[q:h29-reg-applied-25]{H29-REG-APPLIED-25}

\par 難易度 3（校準初稿）、優先度未評価。

\section{堆積相の組合せ・累重とワルターの法則}

\subsection{令和6年度 問31}
\label{q:r06-reg-applied-31}

\noindent\textbf{原典：} sources/R06/17\_applied\_science.pdf、PDF物理10頁。

\textbf{公式正答：} 4（single）。

\paragraph{問題の要旨} 堆積相の記載、組合せ、累重の名前と、漸移的な上下関係の法則を区別する。

\Needspace{4\baselineskip}\paragraph{確認した条件・記述（独自要約）}\mbox{}\par

\noindent\textbf{判断方向：} 最も不適切・誤りを選ぶ\par

\noindent\textbf{原典照合：} 15原典画像で18問の問題幹、全5肢、否定、専門語、化学式、候補順を確認。\par

\noindent\textbf{適用条件：} 流れと構造の典型対応は一意成因を保証しない。用語、時間間隙、保存、試料処理の条件を個別に示す。\par

\paragraph{問われる核心} 堆積相の記載、組合せ、累重の名前と、漸移的な上下関係の法則を区別する。

\paragraph{必要知識} 観察属性と堆積相、association/succession、シーケンス境界、Waltherの連続・漸移条件。

\paragraph{考え方・正答を選ぶ理由（査読待ち初稿）} 堆積相を観察上の特徴で分け、組合せをアソシエーション、累重の順序をサクセションと呼ぶ。4は組合せをシーケンスとする用語の混同である。シーケンスは不整合又は対応する整合面で境界を定める地層群であり、単なる相の集合より条件が多い。ワルターの法則は漸移的・整合的な累重を扱うため、原典の条件を消して任意の上下接触へ使わない。

\Needspace{4\baselineskip}\paragraph{選択肢の検討（査読待ち初稿）}\mbox{}\par

\noindent\textbf{1：適切（条件を含めて検討）} 物理・生物・化学の特徴から堆積環境を復元する。観察記載から成因・環境を推定する解析の説明である。単一指標からの唯一解とはしない。\par

\noindent\textbf{2：適切（条件を含めて検討）} 側方隣接する環境の堆積物だけが漸移的に重なる。ワルターの法則の骨子である。原典の漸移的という条件を保持し、不整合を挟む全上下接触へ広げない。\par

\noindent\textbf{3：適切（条件を含めて検討）} 岩相・粒径・構造・化石で堆積相を区別する。相を区別する記載属性の組合せに対応する。環境の解釈と区別する。\par

\noindent\textbf{4：誤り・最も不適切} 堆積相の組合せを堆積シーケンスと呼ぶ。組合せの名称は堆積相アソシエーションである。シーケンスは境界で区分する成因的な地層群であり同じではない。\par

\noindent\textbf{5：適切（条件を含めて検討）} 累重・配列を堆積相サクセションと呼ぶ。相の上下の累重・配列というサクセションの説明に対応する。\par

\paragraph{再利用する考え方} 観察できる構造と成因・環境の解釈を区別する。代表モデルに適用条件を添え、単一の形や測定値から唯一の環境を断定しない。

\paragraph{共通理論} \hyperref[theory:sediment:facies]{第\ref*{theory:sediment:facies}節}, \hyperref[theory:sediment:walther]{第\ref*{theory:sediment:walther}節}

\par 参考：\cite{sed_mit_environment,sed_walther_oregon}。

\paragraph{関連問題} \hyperref[q:h23-reg-applied-25]{H23-REG-APPLIED-25}

\par 難易度 3（校準初稿）、優先度未評価。

\section{堆積環境指標と測定法の限界}

\subsection{平成23年度 問35}
\label{q:h23-reg-applied-35}

\noindent\textbf{原典：} sources/H23/17\_applied\_science.pdf、PDF物理11頁。

\textbf{公式正答：} 3（single）。

\par\textbf{検討未確定：} 陰イオン量の典型比較とECの古環境対応は濃度・試料処理・後の履歴に依存する。SEMの原典用語に対応する手法全体の追試は未確認。一次資料で確認した範囲と限界を明記して部分校正に残す。

\paragraph{問題の要旨} 堆積環境を反映する地球化学的指標と、元素分析・粒子表面観察の測定対象を分ける。

\Needspace{4\baselineskip}\paragraph{確認した条件・記述（独自要約）}\mbox{}\par

\noindent\textbf{判断方向：} 最も不適切・誤りを選ぶ\par

\noindent\textbf{原典照合：} 15原典画像で18問の問題幹、全5肢、否定、専門語、化学式、候補順を確認。\par

\noindent\textbf{適用条件：} 流れと構造の典型対応は一意成因を保証しない。用語、時間間隙、保存、試料処理の条件を個別に示す。\par

\paragraph{問われる核心} 堆積環境を反映する地球化学的指標と、元素分析・粒子表面観察の測定対象を分ける。

\paragraph{必要知識} XRFの特性X線と元素量、SEMの表面観察、硫化物と酸化還元、ECの試料処理依存。

\paragraph{考え方・正答を選ぶ理由（査読待ち初稿）} 蛍光X線分析は試料の原子を励起し、その特性X線を使う元素分析である。内部の葉理や粒子配列を観察する方法という3が誤りである。黄鉄鉱 \(\mathrm{FeS_2}\) は環境に応じて生成・酸化し、粘土の懸濁液のECもその反応に関わる。GSJの2011年資料は試料処理によるECの変化と環境推定の限界を示す。SEMの表面観察では風による運搬を推定する証拠を得られるが、後の化学的改変も検討する。

\Needspace{4\baselineskip}\paragraph{選択肢の検討（査読待ち初稿）}\mbox{}\par

\noindent\textbf{1：適切（条件を含めて検討）} FeS2含有量は堆積環境に依存する。硫化物形成は硫黄・鉄の供給、還元状態などに依存する。酸化・溶脱後の値を堆積時と無条件に同一視しない。\par

\noindent\textbf{2：適切（条件を含めて検討）} 自然水の陰イオン量は淡水より海水が多い。通常の海水は塩化物等の溶存イオン濃度が淡水より大きい。塩湖等の例外を含む全ての水への全称ではない。\par

\noindent\textbf{3：誤り・最も不適切} XRFで内部の堆積構造を観察する。蛍光X線のエネルギーと強度から元素を分析するので、内部堆積構造の画像を直接得る方法ではない。\par

\noindent\textbf{4：適切（条件を含めて検討）} 石英表面のSEMは風成作用の推定に有効。粒子表面の微細摩耗・衝突形態を環境推定に使える。ただし水成・風成・溶食や再運搬の履歴を単独で一意決定しない。\par

\noindent\textbf{5：適切（条件を含めて検討）} 粘土混濁水の電気伝導度変化は古環境に対応する。懸濁液のECを古環境指標とする研究があるが、乾燥と酸化、混合比や温度、後の溶脱にも依存し、単独判定の限界がある。\par

\paragraph{再利用する考え方} 観察できる構造と成因・環境の解釈を区別する。代表モデルに適用条件を添え、単一の形や測定値から唯一の環境を断定しない。

\paragraph{共通理論} \hyperref[theory:sediment:proxy]{第\ref*{theory:sediment:proxy}節}

\par 参考：\cite{sed_xrf_jaea,sed_gsj_ec2011,sed_sem_shimane}。

\par 難易度 3（校準初稿）、優先度未評価。

\section{未固結堆積物の二次変形}

\subsection{令和元年度 問28}
\label{q:r01-reg-applied-28}

\noindent\textbf{原典：} sources/R01/17\_applied\_science.pdf、PDF物理10頁。

\textbf{公式正答：} 2（single）。

\paragraph{問題の要旨} 堆積中にできる構造と、その後の未固結段階に変形してできる構造を区別する。

\Needspace{4\baselineskip}\paragraph{確認した条件・記述（独自要約）}\mbox{}\par

\noindent\textbf{判断方向：} 堆積後の未固結段階の二次構造を選ぶ\par

\noindent\textbf{原典照合：} 15原典画像で18問の問題幹、全5肢、否定、専門語、化学式、候補順を確認。\par

\noindent\textbf{適用条件：} 流れと構造の典型対応は一意成因を保証しない。用語、時間間隙、保存、試料処理の条件を個別に示す。\par

\paragraph{問われる核心} 堆積中にできる構造と、その後の未固結段階に変形してできる構造を区別する。

\paragraph{必要知識} 初生/二次の分類軸、未固結変形、荷重差・間隙水、葉理と後の変形の区別。

\paragraph{考え方・正答を選ぶ理由（査読待ち初稿）} 分類の軸は形の大きさではなく形成時期と過程である。流痕・リップル・葉理・級化は流れと沈積による初生構造に対応する。火炎構造は未固結の泥が上の砂へ突き出す荷重・流動化に関わるため二次構造であり、本問では2である。後から曲げられた葉理を見た場合は、葉理の形成と変形を別の段階として読む。火炎があるだけで地震と断定しない。

\Needspace{4\baselineskip}\paragraph{選択肢の検討（査読待ち初稿）}\mbox{}\par

\noindent\textbf{1：非該当・不適切} 流痕。流れが堆積面を侵食する初生構造である。\par

\noindent\textbf{2：該当・適切} 火炎構造。砂泥の荷重差や流動化による、堆積後の未固結変形の二次構造である。\par

\noindent\textbf{3：非該当・不適切} リップル。堆積しながら流れが床形を作る初生構造である。\par

\noindent\textbf{4：非該当・不適切} 葉理。堆積粒子の層状配列を記録する初生構造である。後から変形した葉理という別概念と区別する。\par

\noindent\textbf{5：非該当・不適切} 級化層理。堆積中の粒径の時間変化を記録する初生構造である。\par

\paragraph{再利用する考え方} 観察できる構造と成因・環境の解釈を区別する。代表モデルに適用条件を添え、単一の形や測定値から唯一の環境を断定しない。

\paragraph{共通理論} \hyperref[theory:sediment:soft]{第\ref*{theory:sediment:soft}節}, \hyperref[theory:sediment:cross]{第\ref*{theory:sediment:cross}節}, \hyperref[theory:sediment:grading]{第\ref*{theory:sediment:grading}節}

\par 参考：\cite{sed_soft_deformation,sed_houston_structures}。

\paragraph{関連問題} \hyperref[q:r02-reg-applied-29]{R02-REG-APPLIED-29}, \hyperref[q:r06-reg-applied-30]{R06-REG-APPLIED-30}

\par 難易度 3（校準初稿）、優先度未評価。

\subsection{令和6年度 問30}
\label{q:r06-reg-applied-30}

\noindent\textbf{原典：} sources/R06/17\_applied\_science.pdf、PDF物理10頁。

\textbf{公式正答：} 3（single）。

\paragraph{問題の要旨} 堆積中にできる構造と、その後の未固結段階に変形してできる構造を区別する。

\Needspace{4\baselineskip}\paragraph{確認した条件・記述（独自要約）}\mbox{}\par

\noindent\textbf{判断方向：} 堆積後の未固結段階の二次構造を選ぶ\par

\noindent\textbf{原典照合：} 15原典画像で18問の問題幹、全5肢、否定、専門語、化学式、候補順を確認。\par

\noindent\textbf{適用条件：} 流れと構造の典型対応は一意成因を保証しない。用語、時間間隙、保存、試料処理の条件を個別に示す。\par

\paragraph{問われる核心} 堆積中にできる構造と、その後の未固結段階に変形してできる構造を区別する。

\paragraph{必要知識} 初生/二次の分類軸、未固結変形、荷重差・間隙水、葉理と後の変形の区別。

\paragraph{考え方・正答を選ぶ理由（査読待ち初稿）} 分類の軸は形の大きさではなく形成時期と過程である。流痕・リップル・葉理・級化は流れと沈積による初生構造に対応する。火炎構造は未固結の泥が上の砂へ突き出す荷重・流動化に関わるため二次構造であり、本問では3である。後から曲げられた葉理を見た場合は、葉理の形成と変形を別の段階として読む。火炎があるだけで地震と断定しない。

\Needspace{4\baselineskip}\paragraph{選択肢の検討（査読待ち初稿）}\mbox{}\par

\noindent\textbf{1：非該当・不適切} 葉理。堆積粒子の層状配列を記録する初生構造である。後から変形した葉理という別概念と区別する。\par

\noindent\textbf{2：非該当・不適切} 級化層理。堆積中の粒径の時間変化を記録する初生構造である。\par

\noindent\textbf{3：該当・適切} 火炎構造。砂泥の荷重差や流動化による、堆積後の未固結変形の二次構造である。\par

\noindent\textbf{4：非該当・不適切} リップル。堆積しながら流れが床形を作る初生構造である。\par

\noindent\textbf{5：非該当・不適切} 流痕。流れが堆積面を侵食する初生構造である。\par

\paragraph{再利用する考え方} 観察できる構造と成因・環境の解釈を区別する。代表モデルに適用条件を添え、単一の形や測定値から唯一の環境を断定しない。

\paragraph{共通理論} \hyperref[theory:sediment:soft]{第\ref*{theory:sediment:soft}節}, \hyperref[theory:sediment:cross]{第\ref*{theory:sediment:cross}節}, \hyperref[theory:sediment:grading]{第\ref*{theory:sediment:grading}節}

\par 参考：\cite{sed_soft_deformation,sed_houston_structures}。

\paragraph{関連問題} \hyperref[q:r01-reg-applied-28]{R01-REG-APPLIED-28}

\par 難易度 3（校準初稿）、優先度未評価。

\section{津波堆積物の形成・認定・保存}

\subsection{令和3年度 問30}
\label{q:r03-reg-applied-30}

\noindent\textbf{原典：} sources/R03/17\_applied\_science.pdf、PDF物理11頁。

\textbf{公式正答：} 2（single）。

\par\textbf{検討未確定：} 公式/独立2を保持するが、単純/複雑の頻度は環境と尺度に依存し、USGSには構造が稀な例も多い。②の代表頻度を厳密には未確定。④の火山起源津波堆積物の個別報告も未確認。①③⑤の一般化の限界を含めAstra重点査読へ残す。

\paragraph{問題の要旨} 津波堆積物を、供給源、反復する流れ、非地震起源と形成後の保存条件から検討する。

\Needspace{4\baselineskip}\paragraph{確認した条件・記述（独自要約）}\mbox{}\par

\noindent\textbf{判断方向：} 最も不適切・誤りを選ぶ\par

\noindent\textbf{原典照合：} 15原典画像で18問の問題幹、全5肢、否定、専門語、化学式、候補順を確認。\par

\noindent\textbf{適用条件：} 流れと構造の典型対応は一意成因を保証しない。用語、時間間隙、保存、試料処理の条件を個別に示す。\par

\paragraph{問われる核心} 津波堆積物を、供給源、反復する流れ、非地震起源と形成後の保存条件から検討する。

\paragraph{必要知識} 反復波と遡上/戻り流れ、供給源、微化石、背景堆積と保存、非地震起源、事例頻度の母集団。

\paragraph{考え方・正答を選ぶ理由（査読待ち初稿）} 初回は波の反復・戻り流れによる多様性から2を最も不適切と判定し、公式2に一致した。藤野の報告は塊状から複数級化ユニットまでの多様性を示す。一方、USGSの近年の事例集では堆積構造が稀な例も多い。したがって「津波なら必ず複雑だから2が誤り」とは書けず、対象環境・記録の尺度による頻度の違いを疑義に残す。生物遺骸、分布、構造、年代などの複数証拠で認定し、形成と保存と観察可能性を分ける。

\Needspace{4\baselineskip}\paragraph{選択肢の検討（査読待ち初稿）}\mbox{}\par

\noindent\textbf{1：適切（条件を含めて検討）} 生物遺骸から侵食した海底の深度を検討する。生物の生息深度は供給源を制約する証拠となるが、混合と再堆積の影響を考慮する。\par

\noindent\textbf{2：誤り・最も不適切} 堆積構造は比較的単純な場合が多い。多様な層理・正逆級化・複数波のユニットがあるため一律な単純化が問題となる。ただしUSGS近年資料では構造が稀な例も多く、単純/複雑の一般頻度は資料依存である。\par

\noindent\textbf{3：適切（条件を含めて検討）} 礫床河川・扇状地・岩礁海岸は調査に不向き。通常の粗粒供給や侵食でイベント層の識別・保存が難しい。岩礁に津波記録が絶対にないという意味ではない。\par

\noindent\textbf{4：適切（条件を含めて検討）} 火山・隕石・海底地すべり起源の津波にも報告がある。隕石衝突とストレッガーの報告は確認した。火山起源堆積物の個別報告は今回未確認であり断定的な全確認扱いにしない。\par

\noindent\textbf{5：適切（条件を含めて検討）} 内湾は安定した環境で保存されやすい。静かな内湾・湿地の背景堆積と後の埋積は保存に有利な場合がある。侵食や擾乱を受けない保証ではない。\par

\paragraph{再利用する考え方} 観察できる構造と成因・環境の解釈を区別する。代表モデルに適用条件を添え、単一の形や測定値から唯一の環境を断定しない。

\paragraph{共通理論} \hyperref[theory:sediment:tsunami]{第\ref*{theory:sediment:tsunami}節}

\par 参考：\cite{sed_tsunami_fujino,sed_tsunami_usgs,sed_tsunami_2002,sed_tsunami_gsj_survey}。

\par 難易度 4（校準初稿）、優先度未評価。

\section{形成時の上下と堆積構造}

\subsection{令和2年度 問29}
\label{q:r02-reg-applied-29}

\noindent\textbf{原典：} sources/R02/17\_applied\_science.pdf、PDF物理10頁。

\textbf{公式正答：} 1（single）。

\paragraph{問題の要旨} 転倒した地層の形成時の上下を決める構造と、節理の伝播を示す構造を区別する。

\Needspace{4\baselineskip}\paragraph{確認した条件・記述（独自要約）}\mbox{}\par

\noindent\textbf{判断方向：} 最も不適切・誤りを選ぶ\par

\noindent\textbf{原典照合：} 15原典画像で18問の問題幹、全5肢、否定、専門語、化学式、候補順を確認。\par

\noindent\textbf{適用条件：} 流れと構造の典型対応は一意成因を保証しない。用語、時間間隙、保存、試料処理の条件を個別に示す。\par

\paragraph{問われる核心} 転倒した地層の形成時の上下を決める構造と、節理の伝播を示す構造を区別する。

\paragraph{必要知識} 節理と堆積構造、脱水皿の幾何、荷重方向、侵食型と底面型、生痕の開口。

\paragraph{考え方・正答を選ぶ理由（査読待ち初稿）} 羽毛状構造は固結岩の割れ目が広がった過程を記録する。堆積時の上を示す幾何ではないため1が最も不適切である。皿状・荷重・流痕は形成時の重力や層の界面との対応から上下を考え、生痕は開口や埋積を含めて読む。現在の鉛直と堆積時の鉛直は転倒で違うため、一つの曖昧な形だけでなく複数構造を合わせる。

\Needspace{4\baselineskip}\paragraph{選択肢の検討（査読待ち初稿）}\mbox{}\par

\noindent\textbf{1：誤り・最も不適切} 羽毛状構造。節理の破断面の起点から広がる伝播模様である。堆積時の上下面を直接示さない。\par

\noindent\textbf{2：適切（条件を含めて検討）} 皿状構造。一般に上に凹の脱水構造の形と周囲の層を使い上下を検討できる。\par

\noindent\textbf{3：適切（条件を含めて検討）} 荷重痕。上の砂が下の泥へ突出するため、荷重の向きを使える。\par

\noindent\textbf{4：適切（条件を含めて検討）} 流痕。侵食した凹面と上位層底面の凸型の充填を区別すると上下を検討できる。\par

\noindent\textbf{5：適切（条件を含めて検討）} 生痕化石。巣穴の開口・分岐や地表からの貫入などが明瞭なら上下を使える。全ての名称だけで上下が分かるとは限らない。\par

\paragraph{再利用する考え方} 観察できる構造と成因・環境の解釈を区別する。代表モデルに適用条件を添え、単一の形や測定値から唯一の環境を断定しない。

\paragraph{共通理論} \hyperref[theory:sediment:wayup]{第\ref*{theory:sediment:wayup}節}, \hyperref[theory:sediment:soft]{第\ref*{theory:sediment:soft}節}

\par 参考：\cite{sed_plumose_michigan,sed_mit_environment,sed_houston_structures}。

\paragraph{関連問題} \hyperref[q:r01-reg-applied-28]{R01-REG-APPLIED-28}

\par 難易度 3（校準初稿）、優先度未評価。


\chapter{層序・古生物}
\section{背景：離れた地層の時間関係を決める}
\label{theory:stratigraphy:background}
堆積学が形成環境を推定するのに対し、層序学は地層の順序、時間間隙、離れた地点間の対応を考える。似た岩相は異なる時期にもできるため、砂岩同士であることだけでは同時代を意味しない。露頭の上下関係、化石、火山灰、古地磁気などの独立した手掛かりを重ねる。

\section{整合と記録されなかった時間}
\label{theory:stratigraphy:hiatus}
下位の堆積が終わり、非堆積や侵食を経てから上位がたまると、境界に時間間隙を持つ不整合となる。不整合面は上位堆積の直前に存在した地表又は海底の一部を保存する。下位層の堆積直後の地形がそのまま残るとは限らない。埋没、固結、隆起、傾動、侵食を挟む場合もある。整合は有意な時間間隙がほとんどない関係であるが、実際の短い休止を全て識別できるという意味ではない。\cite{strat_iu_unconformity}

\section{平行、傾斜、無整合}
\label{theory:stratigraphy:unconformity}
平行不整合では上下の層理がほぼ平行でも、侵食や年代の抜けがある。傾斜不整合では下位層の傾動と侵食の後に上位がたまり、走向・傾斜に差ができる。無整合は、火成岩・変成岩などの基盤を堆積層が不整合に覆う関係である。層理の幾何と基盤の岩種を使って三つを分ける。平行に見える境界を自動的に整合としない。\cite{strat_iu_unconformity}

\section{鍵層と同時面}
\label{theory:stratigraphy:marker}
追跡しやすい特徴的な層を鍵層として対比に使う。広い海成層は役立つ場合があるが、海成というだけで一枚の同時面になるわけではない。環境帯の移動で同じ岩相が場所ごとに異なる時刻に形成されることがある。一回の噴火に対応する初生火山灰なら、堆積に要する時間が地質年代の尺度より短い場合、ほぼ同時の面として有効である。識別を誤った別噴火や再堆積層ではこの利点を失う。

\section{テフラとは何か}
\label{theory:stratigraphy:tephra}
テフラは火山から放出された破片状物質を扱う語であり、灰、軽石、スコリア、鉱物片などを含む。連続した溶岩流と区別する。文献によって降下物へ限る狭義と、火砕流物質を含む広義があるため、用語の範囲を確認する。テフロクロノロジーは噴火に対応する層を同定して、地域間の対比や年代の制約を得る手法である。\cite{volc_usgs_tephra}

\section{野外記載とガラス・鉱物分析の併用}
\label{theory:stratigraphy:identify}
テフラを同定するときは、層序位置、色、粒径、層厚、構成粒子と分布をまず記載する。火山ガラスや鉱物の屈折率・化学組成は重要な識別材料であるが、外観を使わないという規則ではない。GSJの対比訂正には、層厚・粒径が候補テフラと合わないこと、スコリアと軽石の岩種が異なることを検討した例がある。同じ色だけで同定しないのと同様に、同じ組成だけで上下順序や再堆積の可能性を無視しない。\cite{strat_gsj_tephra2013}

\section{適用域と再堆積の限界}
\label{theory:stratigraphy:scope}
広域テフラは遠隔地域の対比に有利であるが、火山近傍だけの局地的な層も、その地域の噴火順序を制約する。陸上に限らず、海や湖の泥層内でも使える。GSJの渥美層群の記載は中期更新世の海進堆積物中のテフラを対比する例であり、後期更新世以降だけという制限はない。初生堆積なら噴火時刻をほぼ示すが、侵食されて後に再堆積した粒子なら、噴火時刻とその層の堆積時刻は一致しない。\cite{strat_gsj_atsumi}

\section{花粉・微化石・古地磁気の役割}
\label{theory:stratigraphy:pollen}
花粉は陸上植生を反映するが、河川や風で海へ運ばれるため、淡水成層だけにあるとは限らない。徳島平野の海成層では花粉化石組成を層序対比に利用した報告がある。植生・気候の地域差と運搬を考慮し、同じ花粉群を無条件に一時刻へ結び付けない。珪藻、有孔虫、放散虫、貝形虫なども生息環境や年代を検討する材料となる。古地磁気の方向は定方位試料で測り、地層の回転や再磁化も検討する。\cite{strat_pollen_tokushima}

\section{典型的な問われ方と確認順序}
「岩相が似る」「海成である」「花粉を含む」「火山灰を含む」という観察と、時間的な対応の強さを分ける。対比では識別、初生性、保存、年代の順に確認する。不整合では幾何と時間間隙を分ける。個別問題は確認範囲と残る疑義を併記する。
\section{過去問と個別の解法}
\section{テフラ同定・年代対比と再堆積}

\subsection{平成28年度 問34}
\label{q:h28-reg-applied-34}

\noindent\textbf{原典：} sources/H28/17\_applied\_science.pdf、PDF物理13頁。

\textbf{公式正答：} 3（single）。

\par\textbf{検討未確定：} 海成層が常に鍵層、火山灰が常に一番確かという強い一般化は成立範囲を要する。花粉の淡水限定の誤りは確認済みだが、他肢の代表性を部分校正で追跡する。

\paragraph{問題の要旨} 地層対比で使う微化石・花粉・火山灰・古地磁気の情報と適用域を比較する。

\Needspace{4\baselineskip}\paragraph{確認した条件・記述（独自要約）}\mbox{}\par

\noindent\textbf{判断方向：} 最も不適切・誤りを選ぶ\par

\noindent\textbf{原典照合：} 15原典画像で18問の問題幹、全5肢、否定、専門語、化学式、候補順を確認。\par

\noindent\textbf{適用条件：} 流れと構造の典型対応は一意成因を保証しない。用語、時間間隙、保存、試料処理の条件を個別に示す。\par

\paragraph{問われる核心} 地層対比で使う微化石・花粉・火山灰・古地磁気の情報と適用域を比較する。

\paragraph{必要知識} 花粉の運搬・海成保存、微化石群集、鍵層と時間面、初生テフラ、定方位試料。

\paragraph{考え方・正答を選ぶ理由（査読待ち初稿）} 花粉は陸上で作られても風や河川によって海へ運ばれる。海成層の花粉分析による層序対比が実際に報告されているため、淡水成層だけとする3が最も不適切である。微化石と古地磁気は別の観測量を持ち、火山灰は一回の噴火を同定できる場合の対比に強い。広い海成層や火山灰の利点を、全条件での同時性・最確性へ広げない。

\Needspace{4\baselineskip}\paragraph{選択肢の検討（査読待ち初稿）}\mbox{}\par

\noindent\textbf{1：適切（条件を含めて検討）} 海成層の広い分布を鍵層として使う。広く追跡できる特徴的な海成層は鍵となることがある。ただし海成であるだけで同時面とはならない。\par

\noindent\textbf{2：適切（条件を含めて検討）} 微化石を堆積環境調査に使う。珪藻、有孔虫、放散虫、貝形虫などの組成が生息域や堆積環境の材料となる。再堆積も考慮する。\par

\noindent\textbf{3：誤り・最も不適切} 花粉は淡水成層のみに含まれる。海成層でも花粉化石が記録され、陸上植生の変遷や対比に使われる。「のみ」という限定が誤りである。\par

\noindent\textbf{4：適切（条件を含めて検討）} 同一火山灰層を最も確かな対比の決め手とする。同一噴火の初生テフラを正しく同定すれば強い時間対比となる。一律にどの手法より確かとはせず、再堆積と誤対比を除く。\par

\noindent\textbf{5：適切（条件を含めて検討）} 古地磁気は定方位試料で調べる。磁化方向を地理方向へ対応させる必要があるため方位を保持する。傾動や再磁化の補正も別に必要である。\par

\paragraph{再利用する考え方} 観察できる構造と成因・環境の解釈を区別する。代表モデルに適用条件を添え、単一の形や測定値から唯一の環境を断定しない。

\paragraph{共通理論} \hyperref[theory:stratigraphy:marker]{第\ref*{theory:stratigraphy:marker}節}, \hyperref[theory:stratigraphy:pollen]{第\ref*{theory:stratigraphy:pollen}節}, \hyperref[theory:stratigraphy:scope]{第\ref*{theory:stratigraphy:scope}節}

\par 参考：\cite{strat_pollen_tokushima,strat_gsj_atsumi}。

\paragraph{関連問題} \hyperref[q:r02-reg-applied-31]{R02-REG-APPLIED-31}

\par 難易度 3（校準初稿）、優先度未評価。

\subsection{令和2年度 問31}
\label{q:r02-reg-applied-31}

\noindent\textbf{原典：} sources/R02/17\_applied\_science.pdf、PDF物理11頁。

\textbf{公式正答：} 4（single）。

\paragraph{問題の要旨} 火山砕屑物を使う地層対比の適用範囲と、野外記載・鉱物分析の役割を比較する。

\Needspace{4\baselineskip}\paragraph{確認した条件・記述（独自要約）}\mbox{}\par

\noindent\textbf{判断方向：} 最も適切な記述を選ぶ\par

\noindent\textbf{原典照合：} 15原典画像で18問の問題幹、全5肢、否定、専門語、化学式、候補順を確認。\par

\noindent\textbf{適用条件：} 流れと構造の典型対応は一意成因を保証しない。用語、時間間隙、保存、試料処理の条件を個別に示す。\par

\paragraph{問われる核心} 火山砕屑物を使う地層対比の適用範囲と、野外記載・鉱物分析の役割を比較する。

\paragraph{必要知識} 破片状火山物質、局地/広域の用途、海湖成保存、中期更新世対比例、野外・分析の併用。

\paragraph{考え方・正答を選ぶ理由（査読待ち初稿）} テフラは火山の破片状噴出物を扱うので溶岩と区別する4が適切である。他は有効な材料を不当に除外する。局地層も局地対比に使え、海成・湖成層も対象となり、中期更新世の対比例がある。同定は野外の形質とガラス/鉱物分析を併用する。テフラの狭義・広義が文献で違うこと、初生と再堆積が年代解釈を変えることは別の確認事項である。

\Needspace{4\baselineskip}\paragraph{選択肢の検討（査読待ち初稿）}\mbox{}\par

\noindent\textbf{1：非該当・不適切} 局地的テフラを避け広域テフラだけを使う。局地的噴火層もその地域の時間対比に使える。広域性は用途の利点であり必須条件ではない。\par

\noindent\textbf{2：非該当・不適切} 陸上地層だけに適用する。海や湖の堆積物中のテフラも時間対比に使える。\par

\noindent\textbf{3：非該当・不適切} 色・粒度・層厚を使わず屈折率と組成で同定する。分析値は重要だが野外の層序位置・色・粒径・層厚も使う。GSJの対比訂正には粒径/層厚による候補排除がある。\par

\noindent\textbf{4：該当・適切} テフラは一般に火山砕屑物で溶岩と区別する。破片状の噴出物を連続した溶岩流と分けるという本問の一般的な用語に対応する。\par

\noindent\textbf{5：非該当・不適切} 後期更新世以降だけに適用する。GSJの中期更新世海成層の対比例だけでもこの限定への反例になる。古い層では保存と変質を別途検討する。\par

\paragraph{再利用する考え方} 観察できる構造と成因・環境の解釈を区別する。代表モデルに適用条件を添え、単一の形や測定値から唯一の環境を断定しない。

\paragraph{共通理論} \hyperref[theory:stratigraphy:tephra]{第\ref*{theory:stratigraphy:tephra}節}, \hyperref[theory:stratigraphy:identify]{第\ref*{theory:stratigraphy:identify}節}, \hyperref[theory:stratigraphy:scope]{第\ref*{theory:stratigraphy:scope}節}

\par 参考：\cite{volc_usgs_tephra,strat_gsj_tephra2013,strat_gsj_atsumi}。

\paragraph{関連問題} \hyperref[q:h28-reg-applied-34]{H28-REG-APPLIED-34}

\par 難易度 3（校準初稿）、優先度未評価。

\section{不整合の時間間隙と三分類}

\subsection{令和7年度 問28}
\label{q:r07-reg-applied-28}

\noindent\textbf{原典：} sources/R07/17\_applied\_science.pdf、PDF物理10頁。

\textbf{公式正答：} 1（single）。

\paragraph{問題の要旨} 地層境界の時間間隙と、基盤岩種・層理の幾何による不整合の分類を比較する。

\Needspace{4\baselineskip}\paragraph{確認した条件・記述（独自要約）}\mbox{}\par

\noindent\textbf{判断方向：} 最も不適切・誤りを選ぶ\par

\noindent\textbf{原典照合：} 15原典画像で18問の問題幹、全5肢、否定、専門語、化学式、候補順を確認。\par

\noindent\textbf{適用条件：} 流れと構造の典型対応は一意成因を保証しない。用語、時間間隙、保存、試料処理の条件を個別に示す。\par

\paragraph{問われる核心} 地層境界の時間間隙と、基盤岩種・層理の幾何による不整合の分類を比較する。

\paragraph{必要知識} 層序の時間間隙、侵食と非堆積、傾動、基盤岩種、層理の走向・傾斜。

\paragraph{考え方・正答を選ぶ理由（査読待ち初稿）} 不整合面を挟む間に記録されない時間と侵食がある。下位層の堆積後、固結・傾動・侵食が起きてから上位層が堆積することもあるため、下位の堆積直後の地形とする1が最も不適切である。無整合は基盤の火成・変成岩との関係、平行と傾斜不整合は上下層理の幾何で分ける。平行であるだけで整合とする判断を避ける。

\Needspace{4\baselineskip}\paragraph{選択肢の検討（査読待ち初稿）}\mbox{}\par

\noindent\textbf{1：誤り・最も不適切} 不整合面は多くの場合下位層の堆積直後の地形面。下位堆積後の非堆積、傾動や侵食を経て上位がたまるため、下位の堆積直後の面とは限らない。\par

\noindent\textbf{2：適切（条件を含めて検討）} 層状岩が非層状火成・変成岩を覆う不整合を無整合とする。堆積層が火成・変成基盤を不整合に覆う関係に対応する。\par

\noindent\textbf{3：適切（条件を含めて検討）} 平行な上下層で下位が顕著に削剥される境界を平行不整合とする。層理の平行性は時間的連続を保証しない。侵食と時間間隙がある平行不整合である。\par

\noindent\textbf{4：適切（条件を含めて検討）} 上下層の走向・傾斜が違う不整合を傾斜不整合とする。下位が傾動・侵食された後に上位が堆積する幾何に対応する。\par

\noindent\textbf{5：適切（条件を含めて検討）} 境界に時間間隙がほぼない関係を整合とする。本問の層序関係の定義である。検出できない短い休止までゼロと断定する意味ではない。\par

\paragraph{再利用する考え方} 観察できる構造と成因・環境の解釈を区別する。代表モデルに適用条件を添え、単一の形や測定値から唯一の環境を断定しない。

\paragraph{共通理論} \hyperref[theory:stratigraphy:hiatus]{第\ref*{theory:stratigraphy:hiatus}節}, \hyperref[theory:stratigraphy:unconformity]{第\ref*{theory:stratigraphy:unconformity}節}

\par 参考：\cite{strat_iu_unconformity}。

\par 難易度 3（校準初稿）、優先度未評価。


\chapter{構造地質}
\section{何を復元する学問か}
岩石がいつ、どこで、どのように変形したかを、現在の露頭から復元する分野である。
層理や鉱物の配列は変形前後を結ぶ標識となる。
一枚の露頭は三次元構造の一断面であり、姿勢測定・層序・微小組織・地下探査を合わせて
幾何、運動、形成条件、時間順序を推定する。
地質学的な形態名は、物理学の応力やひずみそのものではない。
形を認識した後に、どの力学モデルでその形を説明できるかを問う\cite{str_orientation}。

\section{面と線の姿勢：走向・傾斜・沈下角}
\label{theory:structure:orientation}
地層の境界や断層面は面として、条線や褶曲のヒンジは線として記載する。
面と水平面の交線の方位が走向である。最大の下降を示す方向は走向に直交し、
その鉛直断面内で水平から測る角が真の傾斜である。
走向に平行な鉛直断面では、平面との交線は水平となり、真の傾斜を測れない。
面の法線を \(\boldsymbol n\)、鉛直単位ベクトルを \(\boldsymbol e_z\) とすると、
走向の軸は \(\boldsymbol e_z\times\boldsymbol n\) で表せる。
水平面ではこの外積が零となり走向は一意でない。鉛直面では下降側の選び方にも注意する。

線の姿勢は、下向きの水平投影方位と水平からの沈下角（プランジ）で表す。
線の単位方向 \(\boldsymbol l\) では、沈下角は
\(\arcsin(|\boldsymbol l\cdot\boldsymbol e_z|)\) である。
絶対値は軸としての傾きのみを表し、変位センスの正負は別に定める。
クリノメータは水平を作る部分と鉛直面内の角を測る部分を使い分ける。
磁針を使う方位には偏角補正・近傍磁性体・計器の校正も必要である。
原典の「線構造の走向」は線の水平投影方位の呼称として保持し、面の走向とは区別する
\cite{str_orientation}。

\section{褶曲の幾何と年代}
\label{theory:structure:folds}
曲がった地層のうち曲率の大きい部分がヒンジ、比較的曲率の小さい両側が翼である。
一つの層上のヒンジを結ぶヒンジ線と、複数層のヒンジを含む軸面を分ける。
上に凸の形がアンチフォーム、下に凸の形がシンフォームである。
背斜は中心に古い層、向斜は中心に新しい層を持つという層序の概念である。
地層の逆転や再褶曲があると、上に凸という観察だけから背斜とは決められない
\cite{str_fold_geometry}。
カルミネーションは軸が延長方向に高まる部分であり、沈下角の変化と関係する。
単一断面での山形と、軸方向の高まりは別の情報である\cite{str_culmination}。

\begin{center}
\setlength{\unitlength}{1pt}
\begin{picture}(240,82)
\qbezier(20,15)(120,105)(220,15)
\qbezier(20,5)(120,95)(220,5)
\qbezier(20,-5)(120,85)(220,-5)
\put(112,63){\makebox(0,0){ヒンジ}}
\put(40,45){\makebox(0,0){翼}}
\put(200,45){\makebox(0,0){翼}}
\put(120,-5){\line(0,1){67}}
\end{picture}
\par\small 独自模式断面：上に凸の形態である。地層の年代が未指定なので背斜とは確定しない。
\end{center}

\section{座屈と曲げ：力の向きから読む}
\label{theory:structure:buckling}
層に沿う短縮により、初期の小さな曲がりが増幅するのが座屈褶曲である。
一方、層に大きな角度で荷重が作用して曲げる形式もある。
粘性の異なる層と周囲の物質は変形を分担し、層の厚さ・物性差・層間距離が形へ影響する。
地層を硬い弾性棒と同じものとせず、粘性流動、塑性変形、圧力溶解や層間すべりを考える
\cite{str_fold}。

物理学との接続として、曲げ剛性 \(EI\) の細い弾性棒に軸圧縮力 \(P\) が加わる小変位モデルでは、
\[
 EI\,\frac{d^4y}{dx^4}+P\,\frac{d^2y}{dx^2}=0
\]
となる。\(E\) はヤング率、\(I\) は断面二次モーメント、\(y\) は横変位である。
両端ピンの長さ \(L\) の棒なら最初の固有モードで \(P_{\rm cr}=\pi^2EI/L^2\) となる。
これは層に沿う圧縮が曲げを生むことを示す教材上の類比であり、
自然の粘性地層の臨界荷重や波長をこの式で確定するものではない。

\section{層面すべり・キンク・スランプ}
曲げた複数枚の紙では、外側と内側の弧長差を紙間のすべりが調整する。
地層でも弱い層や境界へ変形が集中する層面すべり褶曲がある。
層間が完全接着した場合と、すべり得る場合は異なるモデルである。
キンクは比較的直線的な翼と局所的な折れ曲がりで姿勢が変わる形式であり、
物性の異方性と関係する\cite{str_fold}。
スランプ褶曲は堆積物の斜面移動に伴う変形である。
移動物の前方停止と後方圧縮は形成の一例であり、すべての過程を一つの停止機構へ固定しない。
令和6年度の強度・剛性率の量的記述とスランプの全称的な説明は追加査読へ残す。

\section{断層の基本：上盤と下盤、相対変位}
\label{theory:structure:faults}
断層は両側の岩体に相対変位を生じた面又は帯である。
傾斜した面の上側が上盤、下側が下盤であり、上盤の相対下降が正断層、上昇が逆断層となる。
横ずれは走向に沿う成分である。向かい側の地塊が右へ動いて見えれば右横ずれ、
左へ動いて見えれば左横ずれという。斜めすべりは複数成分を持つ。
絶対標高や、一つの岩体が東へ動くという記述だけでは運動の分類を確定できない。
条線はすべり軸を記録するが、単一の直線には向きの正負の情報がない。
センスの決定には段差・非対称組織・変位した標識などの追加観察が必要となる。

\section{地塁・地溝と撓曲}
\label{theory:structure:horst}
両側に断層をもつ細長い地塊で、周囲より相対的に高いものが地塁、低いものが地溝である。
同じ正断層群でも、どの地塊を比較するかによって高低が異なる。
原典の「落ち込む地帯を地塁とする」は上下を取り違えた記述である。
基盤の変位が被覆層で広がって現れる場合、地層は線状の切断だけでなく
単斜又は撓曲という折れ曲がりを示す。地下断層と地表の形は常に同じではない。

\section{分岐・低角面・湾曲面}
\label{theory:structure:geometry}
衝上断層で大きく水平移動した岩体がナップ、その下の剥離面がデコルマン又はデタッチメントである。
底の平坦部から傾斜部（ランプ）へ移る上盤の運動は、断層折れ曲がり褶曲をつくり得る。
横ずれ断層帯には、浅部へ枝分かれし深部で主断層へ収束する花束構造が見られる。
これらは移動岩体、すべり面、断層群の断面形態という別々の記載である。
リストリック正断層は、典型的には浅部で急、深部で緩くなる湾曲面である。
「浅部が緩く深部が急」は逆の説明となる。
地すべりの地層境界にも湾曲して深部へ収束する正断層例があり、形態と形成場を分けて使う
\cite{str_listric}。

\section{横ずれの開く場所：プルアパート}
\label{theory:structure:pullapart}
横ずれ断層が一直線でなく、屈曲又は互い違いの区間を持つと、
主断層に沿う運動が局所的な伸張又は圧縮を生む。
開く配置に生じる伸張性の盆地がプルアパート構造である。
反対の配置では圧縮・隆起が生じ得るので、横ずれだからすべて盆地とするのは誤りである。
左右のずれと、段差が開くか閉じるかを別に図示して判定する。

\section{雁行と共役、配列名と変位センス}
\label{theory:structure:conjugate}
雁行は、複数の構造がほぼ平行に斜交しながら階段状に並ぶ幾何学的な配置である。
共役断層は同一の応力場等に対応する異なる方向の群であり、平行な段状配列の別名ではない。
交差する二方向が見えても、形成時期・応力場が違う二群を誤認する可能性がある。
切断関係はイベントの順序の手掛かりとなるが、既存断層の再活動があれば
最後の活動順序と初生形成の順序は一致しない場合がある\cite{str_orientation}。
ミ型・杉型という配列名を、そのまま右横ずれ・左横ずれという運動名へ置き換えない。
熊本地震の大学調査報告は、右横ずれに伴う杉型の地表配列を記録する
\cite{str_en2016}。亀裂の種類、観察方向、図の回転、引きずりの記載を固定して用いる。
原典3問の配列からセンスへの適用条件は追加査読へ残す。

\section{断層岩：形成機構と分類軸}
\label{theory:structure:rocks}
断層岩は、すべり帯で原岩が破砕・変形・融解した記録である。
未固結で細粒基質に富む断層ガウジと、固結した破砕岩のカタクレーサイトを区別する。
後者のプロト・通常・ウルトラ区分には、基質と破砕片の割合を用いる。
一つの分類案では細粒基質が10〜50、50〜90、90〜100\%という区分となる。
分類案や境界の記載方式を確認し、固結度をこの割合の代用にしない。
マイロナイトは結晶塑性等で粒径や配列が変化した岩石であり、
同じ細粒というだけで破砕岩と同一機構にしない\cite{str_faultrocks}。

\section{摩擦溶融と非対称微小構造}
\label{theory:structure:shear}
シュードタキライトは摩擦等の熱で局所的に溶融して急冷した断層岩である。
ガラス質基質、融食された鉱物片等は形成過程の手掛かりであり、
すべての試料で同じ特徴が全部残るとは限らない。
剛体粒子・層状鉱物等に見られる紡錘形のフィッシュや、
斑状粒子の非対称な尾は剪断センスの指標となる。
その判定は線構造に平行で面構造に垂直な断面を観察するなど、方位条件が必要である。
ガウジの粒子回転や粘土による可塑的な挙動は、
鉱物結晶の塑性変形や動的再結晶と区別する\cite{str_shear}。

\section{地下の破壊と地表の表れ}
\label{theory:structure:source}
地震を発生させた地下の破壊面が震源断層、地震時にずれが地表まで達して表れるものが
地表地震断層である。地下の破壊の全範囲が地表に露出するとは限らない。
活断層は最近の地質時代に活動し将来も活動すると考えられる断層である。
地質時代の採用期間には数十万年前以降を重視するものや第四紀中の証拠を含めるものがある。
過去の証拠から将来の可能性を評価することと、特定年月の予知は異なる
\cite{str_source}。
平成30年度の「地震断層」の史料条件は歴史例に限定した呼称の可能性があり、
地表地震断層一般に必要な条件かを査読対象として残す。

\section{活動度：累積した変位を時間で割る}
\label{theory:structure:activity}
長期間の変位量 \(\Delta u\) と期間 \(\Delta t\) から、
\(\bar v=\Delta u/\Delta t\) という平均変位速度を求める。
これは地震時の瞬間すべり速度ではない。
\[
1\,{\rm m}/1000\,{\rm yr}=1\,{\rm mm}/{\rm yr}.
\]
活動度Aは1以上10未満、Bは0.1以上1未満、Cは0.01以上0.1未満m/千年である
\cite{str_activity}。断層運動の成分、累積変位の測定対象、年代誤差を確認する。
この活動度A/B/Cと、一定期間の地震発生確率に付けたランクは別の指標である。
規模と断層長・すべり量の関係は統計的傾向であり、単一の長さから規模を厳密に決めない。

\section{変位地形と、別の過程で作られた地形}
\label{theory:structure:landforms}
閉塞丘は横ずれで移動した尾根等が谷を遮る地形、撓曲崖は層の曲がりとして現れる崖である。
三角末端面は山脚の切断とその後の侵食等を反映し、バルジは圧縮・隆起部の形態となる。
どれも形だけで断層を確定せず、地層対比や累積変位などを合わせる。
モレーンは氷河が運ぶ岩屑の堆積地形である\cite{str_moraine}。
その堆積物が後に断層で変位する場合も、堆積地形としての成因は氷河である。

\section{調査の役割：観察する量を先に決める}
\label{theory:structure:survey}
空中写真は変位地形と広域の連続性、踏査は露頭の地層・岩相・姿勢・断層岩、
トレンチは断層と地層の切断・被覆を観察する。
年代資料を組み合わせて活動履歴を制約する。
磁気探査は岩石の磁化のコントラストを利用して地下の境界を推定するものであり、
磁気異常を断層そのものと同一視しない。
GSJの出題前年の資料は沿岸活断層の連続性・活動性調査に空中磁気探査を含める
\cite{str_coastal2017}。
花粉分析は古植生・古環境や地層の対比を補う。
野坂断層の研究では、花粉・珪藻分析を堆積環境の変遷のために実施している
\cite{str_pollen2014}。直接変位と間接環境資料を分けると役割を理解できる。
平成30年度問23の初回候補5と公式4は、調査目的を省いた相対順位の疑義として保持する。

\section{貫入・浮力と地質構造}
\label{theory:structure:intrusions}
シルは層に沿う板状貫入岩体、ラコリスは屋根を膨らませた貫入岩体である
\cite{str_intrusion}。周囲の地層を曲げることはあっても、ヒンジや背斜などの褶曲用語とは
記載対象が違う。
ダイアピールは低密度等の塑性物質の上昇・貫入による構造である
\cite{str_diapir}。密度差のある流体系では重力による位置エネルギー低下が
上昇の駆動の一つとなる。蓋層の強度、粘性、差応力も形へ影響する。
断層との共存・相互制御を否定せず、「最も関係が薄い」は中心的な形成機構の比較として扱う。

\section{日本列島への接続}
\label{theory:structure:japan}
変成帯は温度・圧力の履歴を、断層は変位の履歴を表す。
領家の高温低圧と三波川の高圧型の対比は岩石学へ接続する\cite{str_ryoke}。
島弧の形成に関与する沈み込みプレートと、島弧を載せるプレートは同じ対象ではない。
方位と「年何cm」という速度には相手又は固定した参照系が必要である。
北海道の弧の衝突、九州の伸張など地域の形成史を、単一の局所断層の分類へ置換しない。

\section{典型的な問われ方と学習の順序}
今回照合した20問には、用語の関連を比べる5問、構造・活動・断層岩等の記述を比べる14問、
姿勢測定の手順1問がある。これは確認した群の内訳であり、全期間の頻度の確定値ではない。
最初に面と線の姿勢、上盤と下盤、褶曲の形態と年代を学ぶ。
次に断層岩の機構と分類軸、活動度と調査目的を結ぶ。
繰り返し問題でも、固結度と割合、地下と地表、平行と直交、地塁と地溝、
配列名と運動名、史料条件の追加という違いを比較する。
全問の分類・解説が揃ってから確認済み統計による頻度分析へ接続する。

\section{過去問}
\subsection{平成25年度 問20}
\label{q:h25-reg-applied-20}

\noindent\textbf{原典：} sources/H25/17\_applied\_science.pdf、PDF物理8頁。

\textbf{公式正答：} 5（single）。

\par\textit{原典本文の校正と独自解説は未作成。分類のみの登録項目。}

\subsection{平成25年度 問21}
\label{q:h25-reg-applied-21}

\noindent\textbf{原典：} sources/H25/17\_applied\_science.pdf、PDF物理8頁。

\textbf{公式正答：} 5（single）。

\par\textit{原典本文の校正と独自解説は未作成。分類のみの登録項目。}

\subsection{平成26年度 問35}
\label{q:h26-reg-applied-35}

\noindent\textbf{原典：} sources/H26/17\_applied\_science.pdf、PDF物理14頁。

\textbf{公式正答：} 2（single）。

\par\textit{原典本文の校正と独自解説は未作成。分類のみの登録項目。}

\subsection{平成30年度 問23}
\label{q:h30-reg-applied-23}

\noindent\textbf{原典：} sources/H30/17\_applied\_science.pdf、PDF物理8頁。

\textbf{公式正答：} 4（single）。

\par\textit{原典本文の校正と独自解説は未作成。分類のみの登録項目。}

\subsection{平成30年度 問26}
\label{q:h30-reg-applied-26}

\noindent\textbf{原典：} sources/H30/17\_applied\_science.pdf、PDF物理9頁。

\textbf{公式正答：} 5（single）。

\par\textit{原典本文の校正と独自解説は未作成。分類のみの登録項目。}

\subsection{令和元年度 問31}
\label{q:r01-reg-applied-31}

\noindent\textbf{原典：} sources/R01/17\_applied\_science.pdf、PDF物理11頁。

\textbf{公式正答：} 2（single）。

\par\textit{原典本文の校正と独自解説は未作成。分類のみの登録項目。}

\subsection{令和元年度再試験 問30}
\label{q:r01-ret-applied-30}

\noindent\textbf{原典：} sources/R01\_retest/17\_applied\_science.pdf、PDF物理10頁。

\textbf{公式正答：} 2（single）。

\par\textit{原典本文の校正と独自解説は未作成。分類のみの登録項目。}

\subsection{令和2年度 問27}
\label{q:r02-reg-applied-27}

\noindent\textbf{原典：} sources/R02/17\_applied\_science.pdf、PDF物理10頁。

\textbf{公式正答：} 4（single）。

\par\textit{原典本文の校正と独自解説は未作成。分類のみの登録項目。}

\subsection{令和2年度 問30}
\label{q:r02-reg-applied-30}

\noindent\textbf{原典：} sources/R02/17\_applied\_science.pdf、PDF物理11頁。

\textbf{公式正答：} 3（single）。

\par\textit{原典本文の校正と独自解説は未作成。分類のみの登録項目。}

\subsection{令和3年度 問27}
\label{q:r03-reg-applied-27}

\noindent\textbf{原典：} sources/R03/17\_applied\_science.pdf、PDF物理10頁。

\textbf{公式正答：} 5（single）。

\par\textit{原典本文の校正と独自解説は未作成。分類のみの登録項目。}

\subsection{令和4年度 問29}
\label{q:r04-reg-applied-29}

\noindent\textbf{原典：} sources/R04/17\_applied\_science.pdf、PDF物理10頁。

\textbf{公式正答：} 5（single）。

\par\textit{原典本文の校正と独自解説は未作成。分類のみの登録項目。}

\subsection{令和4年度 問30}
\label{q:r04-reg-applied-30}

\noindent\textbf{原典：} sources/R04/17\_applied\_science.pdf、PDF物理10頁。

\textbf{公式正答：} 2（single）。

\par\textit{原典本文の校正と独自解説は未作成。分類のみの登録項目。}

\subsection{令和4年度 問31}
\label{q:r04-reg-applied-31}

\noindent\textbf{原典：} sources/R04/17\_applied\_science.pdf、PDF物理11頁。

\textbf{公式正答：} 4（single）。

\par\textit{原典本文の校正と独自解説は未作成。分類のみの登録項目。}

\subsection{令和5年度 問27}
\label{q:r05-reg-applied-27}

\noindent\textbf{原典：} sources/R05/17\_applied\_science.pdf、PDF物理10頁。

\textbf{公式正答：} 4（single）。

\par\textit{原典本文の校正と独自解説は未作成。分類のみの登録項目。}

\subsection{令和5年度 問28}
\label{q:r05-reg-applied-28}

\noindent\textbf{原典：} sources/R05/17\_applied\_science.pdf、PDF物理10頁。

\textbf{公式正答：} 2（single）。

\par\textit{原典本文の校正と独自解説は未作成。分類のみの登録項目。}

\subsection{令和5年度 問29}
\label{q:r05-reg-applied-29}

\noindent\textbf{原典：} sources/R05/17\_applied\_science.pdf、PDF物理10頁。

\textbf{公式正答：} 1（single）。

\par\textit{原典本文の校正と独自解説は未作成。分類のみの登録項目。}

\subsection{令和6年度 問28}
\label{q:r06-reg-applied-28}

\noindent\textbf{原典：} sources/R06/17\_applied\_science.pdf、PDF物理9頁。

\textbf{公式正答：} 5（single）。

\par\textit{原典本文の校正と独自解説は未作成。分類のみの登録項目。}

\subsection{令和6年度 問29}
\label{q:r06-reg-applied-29}

\noindent\textbf{原典：} sources/R06/17\_applied\_science.pdf、PDF物理10頁。

\textbf{公式正答：} 3（single）。

\par\textit{原典本文の校正と独自解説は未作成。分類のみの登録項目。}

\subsection{令和7年度 問29}
\label{q:r07-reg-applied-29}

\noindent\textbf{原典：} sources/R07/17\_applied\_science.pdf、PDF物理11頁。

\textbf{公式正答：} 4（single）。

\par\textit{原典本文の校正と独自解説は未作成。分類のみの登録項目。}

\subsection{令和7年度 問30}
\label{q:r07-reg-applied-30}

\noindent\textbf{原典：} sources/R07/17\_applied\_science.pdf、PDF物理11頁。

\textbf{公式正答：} 5（single）。

\par\textit{原典本文の校正と独自解説は未作成。分類のみの登録項目。}


\chapter{地質年代の補足}
\noindent 補足の章骨格。現在のDBにこの章への主分類割当はなく、追加の出題実績を意味しない。本文は未執筆。



\chapter{地下水}
\section{水頭と含水量の定義}\label{theory:hydro:definitions}
飽和し地下水を伝える地層である帯水層、不飽和帯、自由水面を持つ不圧地下水、圧力下の被圧地下水を区別する。低速流れの全水頭は
\[h=z+\frac{p}{\rho g}\]
であり位置水頭と圧力水頭の和である。動水勾配は水頭差を距離で割る無次元量で、拡散率とは異なる。Darcy流束は\(q=-K\,dh/dl\)で、水頭の下がる方向へ流れる。\cite{inorganic_hydro_ref}
試料全体積\(V\)、間隙体積\(V_v\)、固体体積\(V_s\)、水体積\(V_w\)を使うと
\[\theta=V_w/V,\quad n=V_v/V,\quad e=V_v/V_s=\frac{n}{1-n},\quad S_r=V_w/V_v,\quad\theta=nS_r.\]
飽和では\(S_r=1\)なので\(\theta=n\)。間隙比eとは等しくない。土質の含水比は水質量\(M_w\)を乾燥固体質量\(M_s\)で割る\(w=M_w/M_s\)で、湿潤全質量を分母とした\(M_w/(M_s+M_w)\)とは異なる。百分率表示の100倍を定義と混同しない。
\section{河川交換と水質}
失水河川は河川水が地層へ浸透する区間、得水河川は地下水等の流入を受ける区間である。基底流は低水期にも流量を支える。降雨強度が浸透能を超えると浸透超過のHorton型地表流が生じる。蒸発と植物の蒸散を合わせて蒸発散という。
地下水水質は起源の水に加え、水岩反応、酸化還元、生物作用、人為影響に左右される。平均滞留時間は領域・推計・期間の定義に依存し、原典の大気約8日とUSGSの約10日の例を同じ厳密定数として統一しない。\cite{inorganic_waterchem,inorganic_waterbudget}
\section{流れを生む水頭差と流体の物性}\label{theory:hydro:flow}
一様な媒体の一次元層流でDarcy流束は \(q=-K\,dh/dl\)。流束は全断面積当たりの流量で、間隙内の平均移動速度とは違う。一定の有効間隙率 \(n_e\) で単純化すると実流速は \(q/n_e\) となる。被圧・不圧の区別だけで法則の成立を排除せず、層流・媒体の連続体近似等を確認する。
固有透過度 \(k\) の次元は面積、透水係数は \(K=k\rho g/\mu\) で長さ／時間である。媒体の構造が同じでも流体の密度 \(\rho\) と粘性係数 \(\mu\) が変わればKも変わる\cite{analysis49_darcy_ref}。
堆積層の配列等は水平・鉛直方向の透水性を変え得る。一様Kでの透水量係数 \(T=Kb\) は層厚bを含み、一定量Qの定常被圧井戸で単純化した低下量は \(s=Q\ln(R/r)/(2\pi T)\) となる。井戸半径・外側境界等を定めるモデル式で、全時間の非定常応答に無条件適用しない。
帯水層間の漏水は水頭差が決め、重力だけで常に下向きとはならない。難透水層の下の水頭が地表より高ければ井戸は自噴し得る。扇状地での例を地質と水頭を問わない普遍則へ広げない。
地表流では降雨が浸透能を超えるHorton型と、地表近くが飽和する飽和超過型を区別する\cite{analysis49_runoff_ref}。

\section{貯留量を体積と面積へ換算する}\label{theory:hydro:storage}
水頭を単位だけ下げたとき、単位体積から出る水量を比貯留率 \(S_s\) とすると次元は長さの逆数である。均一な被圧帯水層の厚さ \(b\) を積分すると、単位面積当たりの貯留係数
\[
 S=S_sb
\]
となり無次元である\cite{analysis49_storage_ref}。不圧帯水層の排水は比産出率の寄与も大きく、被圧の弾性貯留式と区別する。体積と面積の定義を反転すると厚さの掛け方も逆になる。
飽和した土の圧力水頭を下げ、最初に空気が入る閾値が空気侵入値である。負の圧力水頭で表す規約と、その絶対値を正の吸引量として表す規約を分ける。含水比は水の質量／乾燥固体質量であり、体積含水率は水体積／全体積である。

\section{全地球水と利用可能な淡水}\label{theory:hydro:global}
地球上の全水量には海水、氷、地下水、湖沼・河川、大気中の水等を含む。USGSの2010年資料の代表推計で淡水は約2.5\%である\cite{analysis49_fresh_ref}。この量は氷河や地下水を含むため、容易に取水できる地表水の割合ではない。提示値3\%を選ぶ問題では近い代表値として扱い、厳密な定数にしない。情報を読む際は分母・対象範囲・推計時点を保存する。

\section{原典と全五肢の解説初稿}
\section{Horton型地表流と飽和地表流の定義を分ける。}

\subsection{平成25年度 問29}
\label{q:h25-reg-applied-29}

\noindent\textbf{原典：} sources/H25/17\_applied\_science.pdf、PDF物理10頁。

\textbf{公式正答：} 1（single）。

\paragraph{問題の要旨} Horton型地表流と飽和地表流の定義を分ける。

\Needspace{4\baselineskip}\paragraph{確認した条件・記述（独自要約）}\mbox{}\par

\noindent\textbf{原典確認：} S7\_049の26原典画像で全五肢・否定・添字を確認。\par

\noindent\textbf{対象：} Horton型地表流と飽和地表流の定義を分ける。\par

\paragraph{問われる核心} Horton型地表流と飽和地表流の定義を分ける。

\paragraph{必要知識} 浸透能超過と飽和超過、Darcy則、TとKと固有透過度の区別。

\paragraph{考え方・正答を選ぶ理由（査読待ち初稿）} 地表流の発生を浸透能と飽和状態の二つで整理する。降雨強度が浸透能を上回るものはHorton型であり1が誤り。Darcy流束q=−K dh/dlと河川への流入を区別する。K=\allowbreak{}kρg/μ、一定KでT=\allowbreak{}Kbである。

\Needspace{4\baselineskip}\paragraph{選択肢の検討（査読待ち初稿）}\mbox{}\par

\noindent\textbf{1：不適切} 降雨が浸透能を超える流れを飽和地表流とする。これは浸透能超過型、すなわちHorton型。飽和型は土壌が飽和して追加の水を受け入れられなくなる場合である。\par

\noindent\textbf{2：通常の条件で適切} 透水量係数T=\allowbreak{}Kb。水平流を扱いKを厚さ方向に一定とする場合の式。一般には厚さ方向のKの積分となる。\par

\noindent\textbf{3：通常の条件で適切} 透水係数は媒体固有だけではない。Kは媒体の固有透過度に加え流体の密度と粘性に依存する。\par

\noindent\textbf{4：通常の条件で適切} 得水河川は地下水から水を受ける。地下水の流入により河川流量が増す区間をいう。\par

\noindent\textbf{5：通常の条件で適切} 絶対間隙率と有効間隙率の区別。総間隙体積の割合と、流れに有効な連結間隙の割合を分ける。\par

\paragraph{再利用する考え方} 測定対象と分離対象、定義の分母、絶対量と相対校正、全称と成立条件を分けて判断する。

\paragraph{共通理論} \hyperref[theory:hydro:flow]{第\ref*{theory:hydro:flow}節}

\par 参考：\cite{analysis49_runoff_ref,analysis49_darcy_ref}。

\par 難易度 3（校準初稿）、優先度未評価。

\subsection{令和元年度再試験 問29}
\label{q:r01-ret-applied-29}

\noindent\textbf{原典：} sources/R01\_retest/17\_applied\_science.pdf、PDF物理10頁。

\textbf{公式正答：} 5（single）。

\par\textbf{独立解答と公式正答に相違あり：} 独立解答 1|5。公式5、独立初回1|5を保持。原典1にのみ／だけの限定を追加して公式へ合わせない。Darcy流速と実流速の意味及び扇状地の地質・水頭条件も確認が必要。

\paragraph{問題の要旨} Darcy則・異方性・井戸水位と自噴条件を比較する。

\Needspace{4\baselineskip}\paragraph{確認した条件・記述（独自要約）}\mbox{}\par

\noindent\textbf{原典確認：} S7\_049の26原典画像で全五肢・否定・添字を確認。\par

\noindent\textbf{対象：} Darcy則・異方性・井戸水位と自噴条件を比較する。\par

\paragraph{問われる核心} Darcy則・異方性・井戸水位と自噴条件を比較する。

\paragraph{必要知識} 公式5は扇状地の地下構造と被圧水頭の代表的な説明である。1の原典は被圧の場合に成立するという文で、被圧の場合にのみという限定はない。Darcy則の一般性を指摘するだけでは1が誤りとは言えない。2～4は異方性、透水量係数の逆比例、水頭による流向から排除できるが、1と5の単一選択の整合は査読事項として残す。

\paragraph{考え方・正答を選ぶ理由（査読待ち初稿）} 公式5は扇状地の地下構造と被圧水頭の代表的な説明である。1の原典は被圧の場合に成立するという文で、被圧の場合にのみという限定はない。Darcy則の一般性を指摘するだけでは1が誤りとは言えない。2～4は異方性、透水量係数の逆比例、水頭による流向から排除できるが、1と5の単一選択の整合は査読事項として残す。

\Needspace{4\baselineskip}\paragraph{選択肢の検討（査読待ち初稿）}\mbox{}\par

\noindent\textbf{1：条件付き・解釈留保} Darcy則は被圧帯水層で成立する。適切な層流・連続体近似のもとでは被圧帯水層にも成立する。原典に被圧だけ、のみという語はない。非被圧でも使えることは、この肢を直接否定しない。\par

\noindent\textbf{2：不適切} 水平・鉛直の透水性の違いは稀。堆積層の配列や亀裂方向による異方性は一般的であり、等方性は近似条件である。\par

\noindent\textbf{3：不適切} 一定揚水の水位低下は厚さに比例。均一な被圧帯水層の定常井戸モデルではs=\allowbreak{}Q ln(R/r)/(2πKb)。厚さbへの依存は反比例である。\par

\noindent\textbf{4：不適切} 漏水は常に下向き。流れの向きは上下帯水層の水頭差で決まり、上向きも生じ得る。\par

\noindent\textbf{5：条件付き・解釈留保} 扇状地下流の被圧層と自噴。透水層を難透水層が覆い、水頭が地表より高い場合に自噴する。扇端等でその構造が成立する代表例であり、すべての扇状地に必ず起こるわけではない。\par

\paragraph{再利用する考え方} 測定対象と分離対象、定義の分母、絶対量と相対校正、全称と成立条件を分けて判断する。

\paragraph{共通理論} \hyperref[theory:hydro:flow]{第\ref*{theory:hydro:flow}節}

\par 参考：\cite{analysis49_darcy_ref,inorganic_hydro_ref}。

\par 難易度 4（校準初稿）、優先度未評価。

\section{全地球水に占める氷・地下水を含む淡水の割合を選ぶ。}

\subsection{平成24年度 問35}
\label{q:h24-reg-applied-35}

\noindent\textbf{原典：} sources/H24/17\_applied\_science.pdf、PDF物理10頁。

\textbf{公式正答：} 1（single）。

\par\textbf{検討未確定：} 水量は推計・対象範囲に依存する代表値。3\%を厳密な物理定数とはしない。

\paragraph{問題の要旨} 全地球水に占める氷・地下水を含む淡水の割合を選ぶ。

\Needspace{4\baselineskip}\paragraph{確認した条件・記述（独自要約）}\mbox{}\par

\noindent\textbf{原典確認：} S7\_049の26原典画像で全五肢・否定・添字を確認。\par

\noindent\textbf{対象：} 全地球水に占める氷・地下水を含む淡水の割合を選ぶ。\par

\paragraph{問われる核心} 全地球水に占める氷・地下水を含む淡水の割合を選ぶ。

\paragraph{必要知識} 分母は河川水だけでも利用可能な淡水だけでもなく、海水を含む地球上の水全体である。USGSの2010年資料では淡水約2.5\%なので、提示値で最も近い3\%を選ぶ。氷河・永久雪や地下水を淡水の内訳から除かない。

\paragraph{考え方・正答を選ぶ理由（査読待ち初稿）} 分母は河川水だけでも利用可能な淡水だけでもなく、海水を含む地球上の水全体である。USGSの2010年資料では淡水約2.5\%なので、提示値で最も近い3\%を選ぶ。氷河・永久雪や地下水を淡水の内訳から除かない。

\Needspace{4\baselineskip}\paragraph{選択肢の検討（査読待ち初稿）}\mbox{}\par

\noindent\textbf{1：通常の条件で適切} 3\%。全水量を分母とする代表推計約2.5\%に最も近い。\par

\noindent\textbf{2：不適切} 5\%。淡水総量の代表割合約2.5\%より大きい。\par

\noindent\textbf{3：不適切} 7\%。淡水総量の代表割合約2.5\%より大きい。\par

\noindent\textbf{4：不適切} 9\%。淡水総量の代表割合約2.5\%より大きい。\par

\noindent\textbf{5：不適切} 11\%。淡水総量の代表割合約2.5\%より大きい。\par

\paragraph{再利用する考え方} 測定対象と分離対象、定義の分母、絶対量と相対校正、全称と成立条件を分けて判断する。

\paragraph{共通理論} \hyperref[theory:hydro:global]{第\ref*{theory:hydro:global}節}

\par 参考：\cite{analysis49_fresh_ref}。

\par 難易度 4（校準初稿）、優先度未評価。

\section{比貯留率と貯留係数の次元と厚さ依存を確認する。}

\subsection{令和7年度 問19}
\label{q:r07-reg-applied-19}

\noindent\textbf{原典：} sources/R07/17\_applied\_science.pdf、PDF物理7頁。

\textbf{公式正答：} 4（single）。

\paragraph{問題の要旨} 比貯留率と貯留係数の次元と厚さ依存を確認する。

\Needspace{4\baselineskip}\paragraph{確認した条件・記述（独自要約）}\mbox{}\par

\noindent\textbf{原典確認：} S7\_049の26原典画像で全五肢・否定・添字を確認。\par

\noindent\textbf{対象：} 比貯留率と貯留係数の次元と厚さ依存を確認する。\par

\paragraph{問われる核心} 比貯留率と貯留係数の次元と厚さ依存を確認する。

\paragraph{必要知識} 比貯留率Ssは単位体積・単位水頭の放出量で次元は長さの逆数、被圧帯水層の貯留係数S=\allowbreak{}Ss bは単位面積・単位水頭で無次元である。4は体積と面積の定義を入れ替え、厚さの掛け方も逆にしている。T=\allowbreak{}Kbの次元は長さの二乗／時間。次元の比較で名称暗記の取り違えを防げる。

\paragraph{考え方・正答を選ぶ理由（査読待ち初稿）} 比貯留率Ssは単位体積・単位水頭の放出量で次元は長さの逆数、被圧帯水層の貯留係数S=\allowbreak{}Ss bは単位面積・単位水頭で無次元である。4は体積と面積の定義を入れ替え、厚さの掛け方も逆にしている。T=\allowbreak{}Kbの次元は長さの二乗／時間。次元の比較で名称暗記の取り違えを防げる。

\Needspace{4\baselineskip}\paragraph{選択肢の検討（査読待ち初稿）}\mbox{}\par

\noindent\textbf{1：通常の条件で適切} 空気侵入の負圧水頭。飽和から吸引を増して空気が入る閾値。正の吸引量として表す別の符号規約と区別する。\par

\noindent\textbf{2：通常の条件で適切} 透水係数は媒体固有だけではない。K=\allowbreak{}kρg/μで、流体の密度と粘性の影響も含む。\par

\noindent\textbf{3：通常の条件で適切} 含水比は水質量／乾燥固相質量。通常百分率で表す。体積含水率や湿潤総質量を分母とする割合とは異なる。\par

\noindent\textbf{4：不適切} 単位体積の貯留係数を厚さ倍して比貯留率。単位体積・単位水頭の量は比貯留率Ss。厚さを掛けて無次元貯留係数S=\allowbreak{}Ss bとなり、名前と関係が逆。\par

\noindent\textbf{5：通常の条件で適切} 透水量係数はKと厚さの積。一定Kの水平帯水層でT=\allowbreak{}Kb。SI次元は面積／時間である。\par

\paragraph{再利用する考え方} 測定対象と分離対象、定義の分母、絶対量と相対校正、全称と成立条件を分けて判断する。

\paragraph{共通理論} \hyperref[theory:hydro:storage]{第\ref*{theory:hydro:storage}節}

\par 参考：\cite{analysis49_storage_ref,analysis49_darcy_ref,inorganic_hydro_ref}。

\par 難易度 3（校準初稿）、優先度未評価。

\section{地下水の水質・水頭・含水量の定義を比較する。}

\subsection{平成26年度 問24}
\label{q:h26-reg-applied-24}

\noindent\textbf{原典：} sources/H26/17\_applied\_science.pdf、PDF物理10頁。

\textbf{公式正答：} 2（single）。

\paragraph{問題の要旨} 地下水の水質・水頭・含水量の定義を比較する。

\Needspace{4\baselineskip}\paragraph{確認した条件・記述（独自要約）}\mbox{}\par

\noindent\textbf{判断対象：} 地下水の水質・水頭・含水量の定義を比較する。\par

\noindent\textbf{原典確認：} 22原典頁画像で五候補・否定・化学式を照合。\par

\paragraph{問われる核心} 地下水の水質・水頭・含水量の定義を比較する。

\paragraph{必要知識} 地下水水質は起源の水だけでなく岩石との溶解・沈殿、酸化還元、生物・人為作用等を反映するので2。

\paragraph{考え方・正答を選ぶ理由（査読待ち初稿）} 地下水水質は起源の水だけでなく岩石との溶解・沈殿、酸化還元、生物・人為作用等を反映するので2。全水頭は位置と圧力の和、動水勾配は無次元。θは全体積を分母にし、飽和時θ=n、e=\allowbreak{}n/(1−n)。定義の名前だけでなく分母・次元を点検する。

\Needspace{4\baselineskip}\paragraph{選択肢の検討（査読待ち初稿）}\mbox{}\par

\noindent\textbf{1：選ばない候補} 帯水層に不飽和を含み被圧に自由水面。飽和し水を伝える帯水層と不飽和帯、自由水面の不圧を取り違え。\par

\noindent\textbf{2：選ぶ候補} 水質は降水/地表水/地層/環境に依存。起源と水岩反応・環境条件が水質を変える。\par

\noindent\textbf{3：選ばない候補} 全水頭は位置−圧力、勾配を拡散率。低速地下水のh=\allowbreak{}z+p/(ρg)は和で、差/距離は動水勾配。\par

\noindent\textbf{4：選ばない候補} 体積含水率は水/間隙、飽和で間隙比。θ=Vw/V全体、飽和では間隙率n。間隙比eではない。\par

\noindent\textbf{5：選ばない候補} 基底流は低水期除外、都市で直接流出減。基底流は低水期にも地下水等が支える流量。都市流出の断定も逆。\par

\paragraph{再利用する考え方} 原子番号と周期・族、物性と反応性、酸強度と危険性、生成と除去、定義の分母を別々に確かめる。

\paragraph{共通理論} \hyperref[theory:hydro:definitions]{第\ref*{theory:hydro:definitions}節}

\par 参考：\cite{inorganic_waterchem,inorganic_hydro_ref}。

\paragraph{関連問題} \hyperref[q:h27-reg-applied-23]{H27-REG-APPLIED-23}

\par 難易度 2（校準初稿）、優先度未評価。

\subsection{平成27年度 問23}
\label{q:h27-reg-applied-23}

\noindent\textbf{原典：} sources/H27/17\_applied\_science.pdf、PDF物理9頁。

\textbf{公式正答：} 3（single）。

\par\textbf{検討未確定：} 大気/河川平均滞留時間は資料・期間・推計の定義に依存。原典8日とUSGS約10日の違いを保持。

\paragraph{問題の要旨} 含水比の分母と陸水の移動過程を比較する。

\Needspace{4\baselineskip}\paragraph{確認した条件・記述（独自要約）}\mbox{}\par

\noindent\textbf{判断対象：} 含水比の分母と陸水の移動過程を比較する。\par

\noindent\textbf{原典確認：} 22原典頁画像で五候補・否定・化学式を照合。\par

\paragraph{問われる核心} 含水比の分母と陸水の移動過程を比較する。

\paragraph{必要知識} 通常の土質含水比は水質量/乾燥固体質量で、湿潤土全体で割った3は別の質量分率。

\paragraph{考え方・正答を選ぶ理由（査読待ち初稿）} 通常の土質含水比は水質量/乾燥固体質量で、湿潤土全体で割った3は別の質量分率。同じ試料でも値が異なる。失水/得水河川は交換の向き、Horton流は浸透能力を越える雨、蒸発散は気化の合算。大気滞留約8日の値を全推計の厳密定数とせず、比較関係と数値の留保を分ける。

\Needspace{4\baselineskip}\paragraph{選択肢の検討（査読待ち初稿）}\mbox{}\par

\noindent\textbf{1：選ばない候補} 大気平均滞留8日で河川より短い。約8日は推計条件依存。USGSには大気約10日・河川16–26日例もある。\par

\noindent\textbf{2：選ばない候補} 失水河川は地下へ浸透。河川から帯水層へ水が失われる区間。\par

\noindent\textbf{3：選ぶ候補} 含水比は水/湿潤土全体。通常土質含水比w=\allowbreak{}Mw/M乾燥であり湿潤全体ではない。\par

\noindent\textbf{4：選ばない候補} 浸透能より降雨強度大ならHorton流。浸透能力を超え余剰が地表へ流れる。\par

\noindent\textbf{5：選ばない候補} 蒸散と蒸発を合わせ蒸発散。植物経由と非植物面の気化を合算。\par

\paragraph{再利用する考え方} 原子番号と周期・族、物性と反応性、酸強度と危険性、生成と除去、定義の分母を別々に確かめる。

\paragraph{共通理論} \hyperref[theory:hydro:definitions]{第\ref*{theory:hydro:definitions}節}

\par 参考：\cite{inorganic_hydro_ref,inorganic_waterbudget}。

\paragraph{関連問題} \hyperref[q:h26-reg-applied-24]{H26-REG-APPLIED-24}

\par 難易度 2（校準初稿）、優先度未評価。


\chapter{応用地質・地盤}
\section{岩石の成因を地盤の挙動へつなぐ}
応用地質では、材料の強さだけでなく割れ目、弱層、水、応力の向きなどが構造物や斜面の安定性にどう影響するかを調べる。硬い鉱物からできた岩石でも、多数の割れ目を持つ岩盤では弱面に沿う破壊が支配的になり得る。地質名をそのまま一つの強度値と同一視しない。

\section{特殊土の名称を母材と堆積過程に戻す}
\label{theory:engineering:soil}
まさ土は花崗岩類の風化で粒状化した土である\cite{mlit_masa}。しらすは主に未溶結の火砕流堆積物などの火山性物質であり、花崗岩風化とは違う。しらすの再移動・再堆積もある\cite{maff_shiras2017,gsj_shiras1957}。白い砂状だから同じ材料という判断を避ける。
黄土（レス）は風成の細粒堆積物である\cite{gsj_loess2002,loess2003}。火山灰由来の粘性土と同義ではない。泥炭（ピート）は植物遺体が十分に分解されず蓄積した有機質土であり、植物繊維を残すことがある\cite{mlit_peat}。分解程度の幅もある。残留土は生成場所に残り、崩積土は斜面を移動して集積した土であり、名称が示す母材と後の移動を別に整理する。

\section{風化・熱水変質と強度}
\label{theory:engineering:weathering}
風化は地表近くの水・酸素・温度変化などの環境に岩石が応答する過程であり、熱水変質は温かい流体との反応を中心とする。どちらも鉱物の変化や結合の変化を起こす。粘土化・割れ目の増加などで弱化しやすい一方、珪化などで硬化する場合もある\cite{jseg_engineering_faq}。
有色鉱物の多い岩石が風化に不安定という一般的比較は、鉱物の形成条件と地表環境の違いを反映する。実際の耐風化性には粒子の結合、割れ目、流体経路が関係するので色だけで寿命を決めない。

\section{スレーキング：空気と固体の応力を分ける}
\label{theory:engineering:slaking}
乾燥した岩塊を水に浸すと崩れるスレーキングには複数の機構がある。封入空気の圧縮はその一つであるが、膨潤などの説明もあり、全ての岩石について一つのモデルへ還元しない\cite{slaking1979,slaking2006}。
閉じ込められた空気の体積が減れば空気の圧力が上がる。しかし固体の孔壁は内部から押し広げられるため、骨格には引張破壊を促す応力が生じる。例えば単純化した薄い球殻では内外圧差\(\Delta p\)、半径\(r\)、厚さ\(t\)に対して膜の引張応力は\(\sigma\simeq\Delta p r/(2t)\)となる。この式は岩石の実形状への予測式ではなく、空気の圧縮と骨格の圧縮を同じにしないための説明用モデルである。

\section{不均質・異方性と弱面の向き}
\label{theory:engineering:anisotropy}
メランジュは細粒の基質にさまざまな岩塊を含み、層としての連続性に乏しい。硬い岩塊と弱い基質の分布、割れ目、水の経路を確認する\cite{jseg_engineering_faq}。結晶片岩の片理面では力学応答が向きで変わる。面が斜面へ抜けるか、抵抗する向きかという幾何も重要である。
せん断抵抗の説明用モデルとして\(\tau_f=c'+(\sigma_n-u)\tan\phi'\)を考えると、同じ全垂直応力\(\sigma_n\)でも間隙水圧\(u\)が増すと有効応力が減る。\(c'\)は有効粘着力、\(\phi'\)は有効摩擦角である。弱面の強度と岩片の強度は同じとは限らず、この単純モデルだけで複雑な斜面の安全率を決められない。

\section{付加体の地盤リスクと解答不一致}
\label{theory:engineering:accretion}
付加・変形によって複雑な構造を持つ地域では、地下水や弱面の連続性が問題になる。メランジュ地域の深層崩壊は実際に報告されている\cite{kyoto_deep2018,sabo_deep2014}。しかし深層崩壊の事例があることから表層崩壊が起こり得ないとは導けない。
蛇紋岩を含む地山の問題は、変質や破砕・粘土化の程度に依存する。粘土質蛇紋岩の膨張性地圧等の事例は、全ての健全な蛇紋岩が吸水膨張するという主張ではない\cite{jsce_swelling2010,serpentinite_mechanics2003}。吸水膨潤と掘削後の塑性流動も区別する。
平成27年の付加体問題では独立判断5と公式1が一致しない。原典5の海山を斜面等へ改変せず、原典1の表層崩壊を否定する限定根拠も未確定である。問題幹・全肢・独立記録・公式表画像の照合を保存し、Astra査読事項として残す。

\section{過去問で確認すること}
土の名称なら母材・生成・運搬、岩盤なら材料強度・弱面・水・応力を分ける。一般的な傾向を全試料に通じる規則にしない。地盤調査・各種試験・対策の測定量と適用条件を次の共通理論から確認する。
\section{登録問題と解説初稿}
\section{岩片強度と岩盤の代理指標}
\label{theory:engineering:strengthindices}
一軸圧縮強度は側方を拘束しない岩石試験片の最大圧縮応力である。岩盤では割れ目、弱面、水、寸法効果が加わり、試験片の値をそのまま現地全体へ使わない。RQDは所定長以上の健全コアの長さ割合に基づく岩盤状態の指標で、打診音は硬さ・風化・割れ目等の現地判別である。弾性波速度は弾性率と密度・接触状態を反映する。いずれも強度を一意に測る式ではなく、調査・試験を組み合わせて判断する\cite{jseg_engineering_faq,geebec_ps}。
圧密降伏応力は通常、土の圧縮履歴や降伏を表す量である。令和元年度再試験問27では岩盤強度との関連が最も薄いものとして選ぶ。ただし弱い泥岩等の圧縮降伏との関連まで否定しない。物理探査のデータを含む設問でも、中心は岩盤の強度評価であるため本章へ置く。
\section{三相と分母：何に対する割合か}\label{theory:engineering:phases}
土粒子の体積を\(V_s\)、水を\(V_w\)、空気を\(V_a\)とし、
全体積を\(V=V_s+V_w+V_a\)、間隙体積を\(V_v=V_w+V_a\)とする。
乾燥土粒子の質量を\(M_s\)、水質量を\(M_w\)とすれば
\[
 w=\frac{M_w}{M_s},\quad n=\frac{V_v}{V},\quad
 e=\frac{V_v}{V_s},\quad \theta=\frac{V_w}{V},\quad
 S_r=\frac{V_w}{V_v}.
\]
順に含水比、間隙率、間隙比、体積含水率、飽和度である。
百分率で表す場合は該当する比に100を掛ける。\(n=e/(1+e)\)、\(\theta=nS_r\)は
同じ三相の体積から導ける。含水比は水が乾燥土より重ければ1を超えるが、
体積の定義に基づく\(n,\theta,S_r\)は通常0から1である。\cite{eng52_soil_phases}

湿潤質量\(M_s+M_w\)で割る含水率と、乾燥質量で割る土質の含水比を取り違えない。
質量を体積で割れば密度の次元が残るので、百分率の候補を除外できる。
乾燥質量は測定条件に依存する。平成25年度の指定\(110\pm5\,^\circ\mathrm C\)を保持し、
別の試験規約の温度へ勝手に置き換えない。

\section{有効応力から液状化を説明する}\label{theory:engineering:liquefaction}
飽和土の骨格へ作用する応力を、説明用に\(\sigma'=\sigma-u\)と表す。
全応力\(\sigma\)と間隙水圧\(u\)を分ける。緩い砂質土が短時間の繰返しせん断を受け、
水が十分に排出できない場合、収縮しようとする骨格の応答が過剰間隙水圧へ現れる。
uが上がると有効応力が下がり、抵抗・剛性が低下して大きな変形が生じ得る。\cite{eng52_liquefaction}

高い地下水位、若く緩い砂質堆積物、よく淘汰された粒度等は代表的な条件である。
旧河道・埋立地・浜堤列間の低地という地形履歴を地下水・粒度・密度へつなぐ。
N値は標準貫入の抵抗の指標で、低いNは緩い砂の傾向を示すが、上載圧や細粒分・地震外力も考える。
原典のN20以上やN10以下を全土層の安全・危険を分ける一つの絶対線とはしない。
淘汰の悪い土も組成・密度によって挙動が異なる。\cite{eng52_liq_sorting}

試料を採取して動的な変形・液状化抵抗を調べ、PS検層で応答に関係する波速を得る。
試料の乱れと現場の密度を確認する。通常の平板載荷は静的な支持・沈下が主目的なので、
発生の評価の相対比較では直接関係が薄い。しかし液状化後の地盤反力を平板載荷で測る研究もある。
「最も関係が薄い」を「絶対に利用されない」へ変えない。\cite{eng52_plate_liq,geebec_ps}

\section{試験名を測定量と条件へ戻す}\label{theory:engineering:tests}
三軸試験は拘束圧と排水・圧密の条件を指定し、破壊と変形を測る。
一軸試験は側方拘束を与えずに圧縮する。ベーン試験は回転のせん断抵抗、
標準貫入やサウンディングは貫入抵抗を調べる。
前二者の破壊試験と、貫入から経験的に推定する強度・支持性を同じ直接測定とはしない。\cite{eng52_geotech_tests}
圧密試験は荷重の増加に対する時間依存の体積変化・圧縮性・排水を調べる。
状態が変われば強度も変わり得るが、主目的はせん断破壊強度の直接測定ではない。\cite{eng52_consolidation}

\section{Mohr--Coulombの切片・傾きと測定丸め}\label{theory:engineering:mohr}
原典が直線の包絡線を指定したときは
\[
 \tau_f=c+\sigma_n\tan\phi
\]
として切片cと傾き\(\tan\phi\)を分ける。
有効応力の強度定数を使う場合は
\(\tau_f=c'+(\sigma_n-u)\tan\phi'\)となるが、問題の記号を無断で有効応力へ変更しない。
摩擦角の値をそのまま傾きへ入れず、正接を使う。\cite{eng52_mohr_ref}

平成29年度の測定は共通単位\(10^4\,\mathrm{Pa}\)で
\((1,3.1),(2,3.7),(4,4.8),(8,7.1)\)。
端点の傾きは\(4/7\simeq0.5714\)で指定\(\tan30^\circ=0.58\)に近い。
各点の\(\tau-0.58\sigma\)は2.52、2.54、2.48、2.46で、切片2.5の候補と整合する。
丸めた測定点は厳密には一直線ではない。通常の最小二乗でも
傾き0.5696、切片2.5391となり同じ選択となる。数値単位はcにも掛かる。
残差をゼロにする架空の測定値へ修正しない。

\section{揚水と圧密：水が抜けて骨格へ荷重が移る}\label{theory:engineering:consolidation}
地下水揚水で被圧帯水層の水頭が下がると、隣接する粘土層の水圧にも影響する。
同じ全応力でuが低下すれば有効応力は増加する。透水性が低い粘土では、
排水に時間を要して収縮し、地表の沈下として現れる。\cite{eng52_consolidation,eng52_subsidence_ref}
砂の即時的な変形、粘土の圧密、地震による液状化後の体積変化を区別する。
水頭低下と沈下が同時同量でないことにも注意する。

杭等に支えられた構造物の相対的な浮き上がり、道路・埋設物の変形に加え、
低平地では排水・浸水リスクも変わる。深度ごとに帯水層の水頭を観測し、
複数層をつないだ混合水位だけで原因を判断しない。
不圧地下水位の低下だけを広域沈下の唯一の機構としない。

\section{黄鉄鉱の酸化と、溶解・析出の時間順}\label{theory:engineering:pyrite}
還元的な環境で保存された黄鉄鉱は、掘削後に水と酸素へさらされると酸化する。
例えば二価鉄を生成する段階を収支式で書くと
\[
 \mathrm{FeS_2}+\frac72\mathrm{O_2}+\mathrm{H_2O}
 \longrightarrow \mathrm{Fe^{2+}}+2\mathrm{SO_4^{2-}}+2\mathrm{H^+}.
\]
さらに鉄の酸化・加水分解、炭酸塩による中和等が加わるので、実際のpHは材料と水の条件に依存する。
酸化そのものをアルカリ性排水の起源とはしない。\cite{eng52_weathering}

酸性水が膠結物等を溶かして岩石を弱める過程と、溶けた成分が石膏等の硫酸塩で再析出する過程を分ける。
蒸発・過飽和・水の経路によって析出し、結晶成長圧が盤膨れへつながる例がある。
青黒い未酸化岩の性質を、掘削後も酸素に触れないまま硫黄の酸化が進むという説明へ変えない。
酸化帯の充填、下の溶解帯の弱化などがあるため、酸化の存在だけで強化又は抑制を一方向へ決めない。\cite{eng52_weathering}

\section{鉄の色・還元溶出と交換サイト}\label{theory:engineering:soilredox}
酸化的な鉄の酸化物等は赤褐色・黄褐色、還元的なグライ化は青灰色等という代表的な対応がある。
土壌の色は腐植等にも依存するので一つの鉄化合物だけから全部を決めない。
湛水して酸素が消費されるとFe・Mnが還元され、二価の形で溶出しやすくなる場合がある。\cite{eng52_soil_colour}
水圧による力学的な不安定化と、酸化還元による化学的な物性変化を分ける。

粘土・腐植の交換サイトではカチオンが電荷を保ちながら入れ替わる。
同じ固定サイト系での保存条件と、pH等に伴う材料の交換容量の変化を区別する。
土壌有機炭素は重要だが、炭酸塩に富む土壌を含む全炭素を有機物だけへ限定しない。

\section{斜面：弱面と地下水を幾何へ重ねる}\label{theory:engineering:landslide}
重力性の移動を物質・移動様式・速度の履歴で記述する。
層理面・変質凝灰岩等の弱層が斜面へ抜ける流れ盤ではすべりへ注意する。
背斜・ドーム・キャップロックという名前だけで全斜面の安定を決めず、割れ目・水・地形を重ねる。\cite{eng52_slip_weaklayer}
すべり面の抵抗を\(c'+(\sigma_n-u)\tan\phi'\)で考えると、
水圧の上昇で抵抗が下がり得ることを説明できる。

滑落崖、亀裂、隆起・陥没は移動を示す観察対象である。長期緩慢・反復する地すべりも、
急激な滑動もある。「地すべり＝全例低速、崩壊＝全例脆性」の二分法は使わない。
原典の砂質/粘性の逆転は典型比較の中で説明し、全岩種を排他的に分類する根拠にはしない。\cite{eng52_landslide_ref}
浸透を減らす被覆・地下水排水・押え盛土等は機構に応じた対策候補であり、
現地の安全と施工条件を確認せず実行する一般指示ではない。

\section{トンネル：切羽・緩み・応力集中}\label{theory:engineering:tunnel}
掘削は空洞周囲の応力を変え、不連続面に沿う落下・緩み、水による切羽の弱化等を生じ得る。
切羽の天端・鏡面を観察し、地質構造と湧水を支保・設計へ反映する。\cite{eng52_tunnel_ref}
山はねは代表的には高応力・大きな土かぶり・硬く節理の少ない岩盤で問題になり、
小土かぶり・多節理を典型条件とする原典の候補は逆である。\cite{eng52_rockburst}
斜面地形・傾斜地層等の非対称な偏圧、浅部の地表沈下、地下水低下による沈下を分ける。
深さだけで沈下量を一つの単調則として決めない。

\section{骨材：岩種は品質試験の入口}\label{theory:engineering:aggregate}
岩石名から鉱物組成・強度・耐久性・吸水・反応性を考えるが、名称だけで品質を保証しない。
蛇紋岩骨材のポップアウト研究はある一方、骨材用採石の記録もある。
公式の一般的な不適性の比較と、全試料の使用不可能という全称を区別する。\cite{eng52_serpent_aggregate,eng52_popout}
チャートのシリカの反応性、砂岩の結合、片麻岩の方向性等も考える。
一つの健全試料の結果を全変質・風化状態へ広げない。

\section{浅部地中熱：伝導と移流を分離する}\label{theory:engineering:groundheat}
浅部の安定した温度と外気との差を利用して冷暖房する地中熱を、深部の高温地熱発電と区別する。
閉じた地中配管で媒体を循環するクローズドループと、地下水を利用するオープンループを分ける。
地盤・帯水層を熱の貯蔵へ使う場合は蓄熱の履歴と周辺環境への影響も考える。\cite{eng52_groundheat_ref}

熱伝導は温度勾配による輸送、地下水の移流は水が熱を運ぶ過程である。
熱応答試験を伝導モデルで整理すると移流の効果が見掛けの有効熱伝導率へ現れ、
材料自体の伝導率より大きく推定される例がある。\cite{eng52_heat_advection}
原典の一般的な大幅減少という向きを選ぶ誤りとして区別する。
その寄与は流速・方向・幾何・試験時間に依存し、全試験の一定倍率ではない。

\section{付加体の構成岩・形成過程}

\subsection{平成27年度 問34}
\label{q:h27-reg-applied-34}

\noindent\textbf{原典：} sources/H27/17\_applied\_science.pdf、PDF物理13頁。

\textbf{公式正答：} 1（single）。

\par\textbf{独立解答と公式正答に相違あり：} 独立解答 5。初回独立解答5、公式1。公式表107頁を実画像で再確認。1の表層崩壊を否定する排他的根拠は未確定。5の大陸側斜面と海山の関係も原典のまま保持。メランジュの深層崩壊事例を表層崩壊の不可能性へ読み替えない。Astraで原典と地質機構を重点査読。

\paragraph{問題の要旨} 付加体の構造・弱面と斜面安定性に関する記述を判定する。

\Needspace{4\baselineskip}\paragraph{確認した条件・記述（独自要約）}\mbox{}\par

\noindent\textbf{方向：} 最も不適切なもの。\par

\noindent\textbf{対象：} 一般的な岩石分類・形成過程。\par

\paragraph{問われる核心} 付加体の不均質・異方性・斜面崩壊と未解決の肢解釈

\paragraph{必要知識} 付加体の不均質・異方性・斜面崩壊と未解決の肢解釈。母材・運搬堆積・有機物起源、風化、弱面、水の作用、岩盤と土の物性。

\paragraph{考え方・正答を選ぶ理由（査読待ち初稿）} 公式正答は1であり、公式正答表107頁のIII–34の1を画像でも確認した。初回独立解答は5である。原典1の「表層崩壊」と5の「大陸側斜面」「海山」の限定を含め、公式1を支持する排他的理由は現在の一次資料照合では確定していない。メランジュに深層崩壊事例があることだけから1の表層崩壊を誤りとは導かない。相違を未解決として保持し、公式を独立の推論結果と偽らない。

\Needspace{4\baselineskip}\paragraph{選択肢の検討（査読待ち初稿）}\mbox{}\par

\noindent\textbf{1：公式では不適切／根拠未確定} 砂岩の岩塊を剥離性の泥質基質が囲むメランジュは表層崩壊等を生じやすい。。公式はこの肢を不適切とする。しかしメランジュの硬軟混在と割れ目が切土等で問題となることは学会資料で確認できる。深層崩壊との関連もあるが、深層の事例だけで表層崩壊を不可能とは言えない。誤りとする限定根拠は未確定。\par

\noindent\textbf{2：概ね適切} 片理による結晶片岩の力学的異方性は地盤安定性に影響する。。面に沿うせん断抵抗などが方向で変わり、弱面と斜面の位置関係を評価する必要がある。等方な単一強度では表せない。\par

\noindent\textbf{3：条件付きで適切} 蛇紋岩には膨張性地山や光沢ある鱗片状剥離面を持ち、地すべりを起こしやすいものがある。。蛇紋岩地域の地すべりや粘土化した蛇紋岩の膨張性地山は一次資料に事例がある。健全な全蛇紋岩が膨潤するとの主張にはしない。\par

\noindent\textbf{4：概ね適切} 海溝の堆積物は圧縮で形成された衝上断層に沿って大陸側へ付加される。。付加は海洋プレート上の堆積物等が剥ぎ取られ、陸側に取り込まれる過程で、衝上断層が発達する。侵食型の沈み込み域もある。\par

\noindent\textbf{5：初回独立判断で疑義／未確定} 海溝の大陸側斜面で海底地すべりが起こり、海山が断裂して海溝へ崩壊する。。海底斜面崩壊自体は起こり得る。一方、海山は一般に沈み込む海洋プレート上の地形で、この一文の大陸側斜面との関係を疑って初回独立解答5とした。沈み込み中の海山破壊も起こり得るため、単に海山という語だけで絶対的な誤りとは確定しない。\par

\paragraph{再利用する考え方} 付加体の不均質・異方性・斜面崩壊と未解決の肢解釈を成因と観察される性質に結び付け、例示と全称命題、主要鉱物と副成分、公式解答と未解決の解釈を区別する。

\paragraph{共通理論} \hyperref[theory:engineering:accretion]{第\ref*{theory:engineering:accretion}節}, \hyperref[theory:engineering:anisotropy]{第\ref*{theory:engineering:anisotropy}節}

\par 参考：\cite{jseg_engineering_faq,kyoto_deep2018,sabo_deep2014,jsce_swelling2010,aist_accretion2026}。

\par 難易度 3（校準初稿）、優先度未評価。

\section{骨材の岩種・変質と品質}

\subsection{平成24年度 問28}
\label{q:h24-reg-applied-28}

\noindent\textbf{原典：} sources/H24/17\_applied\_science.pdf、PDF物理8頁。

\textbf{公式正答：} 2（single）。

\par\textbf{検討未確定：} 蛇紋岩のポップアウト研究は土研報告の参考文献で確認し、研究原報告の全文は未取得。各岩種の試験当時の総合順位と全試料の適性は未確定。

\paragraph{問題の要旨} 岩石名だけでコンクリート骨材の適性を固定しない。

\Needspace{4\baselineskip}\paragraph{確認した条件・記述（独自要約）}\mbox{}\par

\noindent\textbf{原典照合：} 22原典画像で25問の全五肢・数値・単位・否定を確認。\par

\noindent\textbf{対象：} 岩石名だけでコンクリート骨材の適性を固定しない。\par

\paragraph{問われる核心} 岩石名だけでコンクリート骨材の適性を固定しない。

\paragraph{必要知識} 選択2は一般的な骨材の相対比較として保持する。蛇紋岩によるポップアウトの報告や品質問題はあるが、厚労省の一次資料には骨材用採石も記載される。したがって全蛇紋岩の使用が絶対不可能という解説にはしない。

\paragraph{考え方・正答を選ぶ理由（査読待ち初稿）} 選択2は一般的な骨材の相対比較として保持する。蛇紋岩によるポップアウトの報告や品質問題はあるが、厚労省の一次資料には骨材用採石も記載される。したがって全蛇紋岩の使用が絶対不可能という解説にはしない。

\Needspace{4\baselineskip}\paragraph{選択肢の検討（査読待ち初稿）}\mbox{}\par

\noindent\textbf{1：条件・全称を留保} チャート。硬い珪質岩だが、微細なシリカを含む骨材ではアルカリシリカ反応等に注意。全チャートが無条件に安全とはしない。\par

\noindent\textbf{2：不適切} 蛇紋岩。公式・独立とも2。蛇紋岩骨材に起因するポップアウトの研究例がある。風化・変質・鉱物組成により強度・耐久性の適性が変わる。\par

\noindent\textbf{3：代表的条件で適切} 硬質砂岩。硬く耐久性を持つ適切な試料は骨材候補。砂岩という名称だけで粒子結合や吸水率等を省略しない。\par

\noindent\textbf{4：代表的条件で適切} 斑れい岩。健全で緻密な試料は骨材候補。鉱物組成・変質・有害成分等の試験も必要。\par

\noindent\textbf{5：条件・全称を留保} 片麻岩。片理や異方性を含めて品質評価する。すべての方向・試料で同じ破砕形状・強度とはしない。\par

\paragraph{再利用する考え方} 直接測る量と代理指標を区別し、分母・有効応力・水の経路・材料の条件から判断する。

\paragraph{共通理論} \hyperref[theory:engineering:aggregate]{第\ref*{theory:engineering:aggregate}節}

\par 参考：\cite{eng52_serpent_aggregate,eng52_popout}。

\par 難易度 3（校準初稿）、優先度未評価。

\section{浅部地中熱の方式・移流}

\subsection{令和6年度 問18}
\label{q:r06-reg-applied-18}

\noindent\textbf{原典：} sources/R06/17\_applied\_science.pdf、PDF物理6頁。

\textbf{公式正答：} 5（single）。

\par\textbf{検討未確定：} 浅度200m・低温・騒音は代表的な説明であり全機器・地点の厳密境界ではない。移流の寄与は流速・方向・試験期間等に依存し全例の同じ倍率ではない。

\paragraph{問題の要旨} 浅部地中熱の循環方式と移流による熱輸送を区別する。

\Needspace{4\baselineskip}\paragraph{確認した条件・記述（独自要約）}\mbox{}\par

\noindent\textbf{原典照合：} 22原典画像で25問の全五肢・数値・単位・否定を確認。\par

\noindent\textbf{対象：} 浅部地中熱の循環方式と移流による熱輸送を区別する。\par

\paragraph{問われる核心} 浅部地中熱の循環方式と移流による熱輸送を区別する。

\paragraph{必要知識} 誤り5は地下水の移流による熱輸送の向き。実際の熱伝導率は材料物性、流動を含む見掛け値はモデルと観測条件を反映する。熱応答試験で推定する有効値を岩石自体の不変な熱伝導率と同一視しない。

\paragraph{考え方・正答を選ぶ理由（査読待ち初稿）} 誤り5は地下水の移流による熱輸送の向き。実際の熱伝導率は材料物性、流動を含む見掛け値はモデルと観測条件を反映する。熱応答試験で推定する有効値を岩石自体の不変な熱伝導率と同一視しない。

\Needspace{4\baselineskip}\paragraph{選択肢の検討（査読待ち初稿）}\mbox{}\par

\noindent\textbf{1：代表的条件で適切} 概ね200mより浅い低温地中熱。通常の地中熱利用の目安。深部の高温地熱発電と区別するが、200mを物理的な熱源の断絶境界とはしない。\par

\noindent\textbf{2：代表的条件で適切} 外部騒音が一般に小さい。空気熱交換部等の条件を考える。機器の音が絶対にゼロではない。\par

\noindent\textbf{3：代表的条件で適切} クローズド・オープンループ。地中の閉じた配管で媒体を循環する方式と、地下水を利用する方式を分ける。\par

\noindent\textbf{4：代表的条件で適切} 地盤を蓄熱体にする地下蓄熱。季節等の熱を地盤又は帯水層へ蓄えて利用する。すべての採熱装置が蓄熱運用という意味ではない。\par

\noindent\textbf{5：不適切} 移流で見掛け熱伝導率が大きく減少。地下水移流は熱を運び、熱応答試験の伝導モデルで見掛けの有効熱伝導率が増える場合がある。一般的な減少効果という向きが不適切。\par

\paragraph{再利用する考え方} 直接測る量と代理指標を区別し、分母・有効応力・水の経路・材料の条件から判断する。

\paragraph{共通理論} \hyperref[theory:engineering:groundheat]{第\ref*{theory:engineering:groundheat}節}

\par 参考：\cite{eng52_groundheat_ref,eng52_heat_advection}。

\par 難易度 3（校準初稿）、優先度未評価。

\section{弱面・水圧・地すべり滑動}

\subsection{平成26年度 問33}
\label{q:h26-reg-applied-33}

\noindent\textbf{原典：} sources/H26/17\_applied\_science.pdf、PDF物理13頁。

\textbf{公式正答：} 2（single）。

\par\textbf{検討未確定：} 列挙された全構造の発生頻度や全応急工法の適性は追加照合対象。岩種・構造名だけで安定性又は対策の安全性を保証しない。

\paragraph{問題の要旨} 地すべりの弱面・水・滑動と対策を対応させる。

\Needspace{4\baselineskip}\paragraph{確認した条件・記述（独自要約）}\mbox{}\par

\noindent\textbf{原典照合：} 22原典画像で25問の全五肢・数値・単位・否定を確認。\par

\noindent\textbf{対象：} 地すべりの弱面・水・滑動と対策を対応させる。\par

\paragraph{問われる核心} 地すべりの弱面・水・滑動と対策を対応させる。

\paragraph{必要知識} 公式2は各回で誤りの内容が変わる。H26は構造との関係を否定、R01再試験は泥岩すべりの長期緩慢・反復を否定した点が不適切。地質構造、間隙水圧、滑落崖、速度の履歴を組み合わせて読む。

\paragraph{考え方・正答を選ぶ理由（査読待ち初稿）} 公式2は各回で誤りの内容が変わる。H26は構造との関係を否定、R01再試験は泥岩すべりの長期緩慢・反復を否定した点が不適切。地質構造、間隙水圧、滑落崖、速度の履歴を組み合わせて読む。

\Needspace{4\baselineskip}\paragraph{選択肢の検討（査読待ち初稿）}\mbox{}\par

\noindent\textbf{1：代表的条件で適切} 重力によるマスムーブメント。地すべり・崩壊・土石流・落石は斜面下方への重力性移動に関係する。移動様式は区別する。\par

\noindent\textbf{2：不適切} 流れ盤・背斜・ドーム・キャップで発生しにくい。斜面へ抜ける層理等の弱面、構造による割れ目、水の経路は発生条件になり得る。列挙した構造をまとめて発生しにくいとはしない。\par

\noindent\textbf{3：代表的条件で適切} 最上部の滑落崖・馬蹄形。移動体上端の崖を滑落崖と呼ぶ。馬蹄形はよくある形で、全崖の唯一形ではない。\par

\noindent\textbf{4：代表的条件で適切} 泥岩で緩慢・長期、雨季・融雪で増加。地下水変動や弱いすべり面に対応した反復・継続の例がある。全地すべりの速度を一定にしない。\par

\noindent\textbf{5：条件・全称を留保} 浸透防止・横孔排水・押え盛土。水圧を下げたり滑動を抑えたりする対策。応急対策の施工位置・安全性とすべりの機構の確認が必要で、すべてを無条件に現場実施する指示ではない。\par

\paragraph{再利用する考え方} 直接測る量と代理指標を区別し、分母・有効応力・水の経路・材料の条件から判断する。

\paragraph{共通理論} \hyperref[theory:engineering:landslide]{第\ref*{theory:engineering:landslide}節}

\par 参考：\cite{eng52_slip_weaklayer,eng52_landslide_ref}。

\paragraph{関連問題} \hyperref[q:r01-ret-applied-26]{R01-RET-APPLIED-26}

\par 難易度 3（校準初稿）、優先度未評価。

\subsection{令和元年度再試験 問26}
\label{q:r01-ret-applied-26}

\noindent\textbf{原典：} sources/R01\_retest/17\_applied\_science.pdf、PDF物理9頁。

\textbf{公式正答：} 2（single）。

\par\textbf{検討未確定：} 列挙された全構造の発生頻度や全応急工法の適性は追加照合対象。岩種・構造名だけで安定性又は対策の安全性を保証しない。

\paragraph{問題の要旨} 地すべりの弱面・水・滑動と対策を対応させる。

\Needspace{4\baselineskip}\paragraph{確認した条件・記述（独自要約）}\mbox{}\par

\noindent\textbf{原典照合：} 22原典画像で25問の全五肢・数値・単位・否定を確認。\par

\noindent\textbf{対象：} 地すべりの弱面・水・滑動と対策を対応させる。\par

\paragraph{問われる核心} 地すべりの弱面・水・滑動と対策を対応させる。

\paragraph{必要知識} 公式2は各回で誤りの内容が変わる。H26は構造との関係を否定、R01再試験は泥岩すべりの長期緩慢・反復を否定した点が不適切。地質構造、間隙水圧、滑落崖、速度の履歴を組み合わせて読む。

\paragraph{考え方・正答を選ぶ理由（査読待ち初稿）} 公式2は各回で誤りの内容が変わる。H26は構造との関係を否定、R01再試験は泥岩すべりの長期緩慢・反復を否定した点が不適切。地質構造、間隙水圧、滑落崖、速度の履歴を組み合わせて読む。

\Needspace{4\baselineskip}\paragraph{選択肢の検討（査読待ち初稿）}\mbox{}\par

\noindent\textbf{1：代表的条件で適切} 重力によるマスムーブメント。地すべり・崩壊・土石流・落石は斜面下方への重力性移動に関係する。移動様式は区別する。\par

\noindent\textbf{2：不適切} 泥岩で長期反復が少なく短期急激。典型的な泥岩地域の地すべりには長期緩慢・反復的な滑動があり、その特徴を逆にした記述。急激な移動の例まで否定するものではない。\par

\noindent\textbf{3：条件・全称を留保} 流れ盤等で発生しやすい。原典の代表傾向として読む。弱面が斜面へ抜ける角度、強度、水、地形との組合せが必要で、背斜名だけから全地点を決めない。\par

\noindent\textbf{4：代表的条件で適切} 最上部の滑落崖・馬蹄形。移動体上端の崖を滑落崖と呼ぶ。馬蹄形はよくある形で、全崖の唯一形ではない。\par

\noindent\textbf{5：条件・全称を留保} 浸透防止・横孔排水・押え盛土。水圧を下げたり滑動を抑えたりする対策。応急対策の施工位置・安全性とすべりの機構の確認が必要で、すべてを無条件に現場実施する指示ではない。\par

\paragraph{再利用する考え方} 直接測る量と代理指標を区別し、分母・有効応力・水の経路・材料の条件から判断する。

\paragraph{共通理論} \hyperref[theory:engineering:landslide]{第\ref*{theory:engineering:landslide}節}

\par 参考：\cite{eng52_slip_weaklayer,eng52_landslide_ref}。

\paragraph{関連問題} \hyperref[q:h26-reg-applied-33]{H26-REG-APPLIED-33}

\par 難易度 3（校準初稿）、優先度未評価。

\section{飽和・粒度・密度の液状化条件}

\subsection{平成24年度 問31}
\label{q:h24-reg-applied-31}

\noindent\textbf{原典：} sources/H24/17\_applied\_science.pdf、PDF物理9頁。

\textbf{公式正答：} 5（single）。

\par\textbf{検討未確定：} N20を全地震動・全土層の不液状化保証へ使わない。R03の淘汰の悪さも唯一の排他的条件ではなく、原典は代表的な相対比較として保持。

\paragraph{問題の要旨} 液状化しやすい飽和・粒度・密度の条件を比べる。

\Needspace{4\baselineskip}\paragraph{確認した条件・記述（独自要約）}\mbox{}\par

\noindent\textbf{原典照合：} 22原典画像で25問の全五肢・数値・単位・否定を確認。\par

\noindent\textbf{対象：} 液状化しやすい飽和・粒度・密度の条件を比べる。\par

\paragraph{問われる核心} 液状化しやすい飽和・粒度・密度の条件を比べる。

\paragraph{必要知識} 緩い粒子骨格が飽和し、繰返しせん断で水圧が上がる条件を考える。公式5はこの代表比較で最も不適切な候補。土地名・N値・淘汰の一条件だけで判定を完結せず、密度・粒度・地下水と地震外力を組み合わせる。

\paragraph{考え方・正答を選ぶ理由（査読待ち初稿）} 緩い粒子骨格が飽和し、繰返しせん断で水圧が上がる条件を考える。公式5はこの代表比較で最も不適切な候補。土地名・N値・淘汰の一条件だけで判定を完結せず、密度・粒度・地下水と地震外力を組み合わせる。

\Needspace{4\baselineskip}\paragraph{選択肢の検討（査読待ち初稿）}\mbox{}\par

\noindent\textbf{1：代表的条件で適切} 埋立地・旧河道。若く緩い砂質堆積物や高い地下水位がある場所ではリスクが高まり得る。名称だけで全地点を判定しない。\par

\noindent\textbf{2：代表的条件で適切} 砂丘・浜堤列の間の低地。地下水が浅く飽和した砂質地盤が分布し得る。高い砂丘そのものと列間低地を区別する。\par

\noindent\textbf{3：代表的条件で適切} 高い地下水位。浅い飽和砂層が揺れを受ける条件を作る。水位だけで粒度・密度・揺れを省略しない。\par

\noindent\textbf{4：代表的条件で適切} 粒度が一定の堆積物。よく淘汰された緩い飽和砂は代表的に液状化しやすい。粒度だけで全密度の判定をしない。\par

\noindent\textbf{5：条件・全称を留保} N値20以上。密度・締まりが相対的に高い傾向で、緩い砂の典型条件との関係は小さい。地震動・有効上載圧・粒度等の補正があり、20を絶対的な安全線とはしない。\par

\paragraph{再利用する考え方} 直接測る量と代理指標を区別し、分母・有効応力・水の経路・材料の条件から判断する。

\paragraph{共通理論} \hyperref[theory:engineering:liquefaction]{第\ref*{theory:engineering:liquefaction}節}

\par 参考：\cite{eng52_liquefaction,eng52_liq_sorting}。

\paragraph{関連問題} \hyperref[q:h27-reg-applied-35]{H27-REG-APPLIED-35}

\par 難易度 3（校準初稿）、優先度未評価。

\subsection{平成27年度 問35}
\label{q:h27-reg-applied-35}

\noindent\textbf{原典：} sources/H27/17\_applied\_science.pdf、PDF物理14頁。

\textbf{公式正答：} 5（single）。

\par\textbf{検討未確定：} N20を全地震動・全土層の不液状化保証へ使わない。R03の淘汰の悪さも唯一の排他的条件ではなく、原典は代表的な相対比較として保持。

\paragraph{問題の要旨} 液状化しやすい飽和・粒度・密度の条件を比べる。

\Needspace{4\baselineskip}\paragraph{確認した条件・記述（独自要約）}\mbox{}\par

\noindent\textbf{原典照合：} 22原典画像で25問の全五肢・数値・単位・否定を確認。\par

\noindent\textbf{対象：} 液状化しやすい飽和・粒度・密度の条件を比べる。\par

\paragraph{問われる核心} 液状化しやすい飽和・粒度・密度の条件を比べる。

\paragraph{必要知識} 緩い粒子骨格が飽和し、繰返しせん断で水圧が上がる条件を考える。公式5はこの代表比較で最も不適切な候補。土地名・N値・淘汰の一条件だけで判定を完結せず、密度・粒度・地下水と地震外力を組み合わせる。

\paragraph{考え方・正答を選ぶ理由（査読待ち初稿）} 緩い粒子骨格が飽和し、繰返しせん断で水圧が上がる条件を考える。公式5はこの代表比較で最も不適切な候補。土地名・N値・淘汰の一条件だけで判定を完結せず、密度・粒度・地下水と地震外力を組み合わせる。

\Needspace{4\baselineskip}\paragraph{選択肢の検討（査読待ち初稿）}\mbox{}\par

\noindent\textbf{1：代表的条件で適切} 埋立地・旧河道。若く緩い砂質堆積物や高い地下水位がある場所ではリスクが高まり得る。名称だけで全地点を判定しない。\par

\noindent\textbf{2：代表的条件で適切} 砂丘・浜堤列の間の低地。地下水が浅く飽和した砂質地盤が分布し得る。高い砂丘そのものと列間低地を区別する。\par

\noindent\textbf{3：代表的条件で適切} 高い地下水位。浅い飽和砂層が揺れを受ける条件を作る。水位だけで粒度・密度・揺れを省略しない。\par

\noindent\textbf{4：代表的条件で適切} 粒度が一定の堆積物。よく淘汰された緩い飽和砂は代表的に液状化しやすい。粒度だけで全密度の判定をしない。\par

\noindent\textbf{5：条件・全称を留保} N値20以上。密度・締まりが相対的に高い傾向で、緩い砂の典型条件との関係は小さい。地震動・有効上載圧・粒度等の補正があり、20を絶対的な安全線とはしない。\par

\paragraph{再利用する考え方} 直接測る量と代理指標を区別し、分母・有効応力・水の経路・材料の条件から判断する。

\paragraph{共通理論} \hyperref[theory:engineering:liquefaction]{第\ref*{theory:engineering:liquefaction}節}

\par 参考：\cite{eng52_liquefaction,eng52_liq_sorting}。

\paragraph{関連問題} \hyperref[q:h24-reg-applied-31]{H24-REG-APPLIED-31}, \hyperref[q:r03-reg-applied-29]{R03-REG-APPLIED-29}

\par 難易度 3（校準初稿）、優先度未評価。

\subsection{令和3年度 問29}
\label{q:r03-reg-applied-29}

\noindent\textbf{原典：} sources/R03/17\_applied\_science.pdf、PDF物理10頁。

\textbf{公式正答：} 2（single）。

\par\textbf{検討未確定：} N20を全地震動・全土層の不液状化保証へ使わない。R03の淘汰の悪さも唯一の排他的条件ではなく、原典は代表的な相対比較として保持。

\paragraph{問題の要旨} 液状化しやすい飽和・粒度・密度の条件を比べる。

\Needspace{4\baselineskip}\paragraph{確認した条件・記述（独自要約）}\mbox{}\par

\noindent\textbf{原典照合：} 22原典画像で25問の全五肢・数値・単位・否定を確認。\par

\noindent\textbf{対象：} 液状化しやすい飽和・粒度・密度の条件を比べる。\par

\paragraph{問われる核心} 液状化しやすい飽和・粒度・密度の条件を比べる。

\paragraph{必要知識} 緩い粒子骨格が飽和し、繰返しせん断で水圧が上がる条件を考える。公式2はこの代表比較で最も不適切な候補。土地名・N値・淘汰の一条件だけで判定を完結せず、密度・粒度・地下水と地震外力を組み合わせる。

\paragraph{考え方・正答を選ぶ理由（査読待ち初稿）} 緩い粒子骨格が飽和し、繰返しせん断で水圧が上がる条件を考える。公式2はこの代表比較で最も不適切な候補。土地名・N値・淘汰の一条件だけで判定を完結せず、密度・粒度・地下水と地震外力を組み合わせる。

\Needspace{4\baselineskip}\paragraph{選択肢の検討（査読待ち初稿）}\mbox{}\par

\noindent\textbf{1：代表的条件で適切} 高い地下水位。浅い飽和砂層が揺れを受ける条件を作る。水位だけで粒度・密度・揺れを省略しない。\par

\noindent\textbf{2：条件・全称を留保} 淘汰の悪い堆積物。均一粒径の緩い飽和砂と比べれば典型条件から離れる。ただし細粒分の種類・密度・組成によって液状化し得るため、淘汰だけで全例の発生を否定しない。\par

\noindent\textbf{3：代表的条件で適切} N値10以下。低い貫入抵抗は緩い砂の代表的な指標。粘性土の低N値と飽和砂の液状化機構を同一視しない。\par

\noindent\textbf{4：代表的条件で適切} 埋立地・旧河道。若く緩い砂質堆積物や高い地下水位がある場所ではリスクが高まり得る。名称だけで全地点を判定しない。\par

\noindent\textbf{5：代表的条件で適切} 砂丘・浜堤列の間の低地。地下水が浅く飽和した砂質地盤が分布し得る。高い砂丘そのものと列間低地を区別する。\par

\paragraph{再利用する考え方} 直接測る量と代理指標を区別し、分母・有効応力・水の経路・材料の条件から判断する。

\paragraph{共通理論} \hyperref[theory:engineering:liquefaction]{第\ref*{theory:engineering:liquefaction}節}

\par 参考：\cite{eng52_liquefaction,eng52_liq_sorting}。

\paragraph{関連問題} \hyperref[q:h27-reg-applied-35]{H27-REG-APPLIED-35}

\par 難易度 3（校準初稿）、優先度未評価。

\section{液状化の動的評価と静的載荷}

\subsection{平成23年度 問33}
\label{q:h23-reg-applied-33}

\noindent\textbf{原典：} sources/H23/17\_applied\_science.pdf、PDF物理11頁。

\textbf{公式正答：} 5（single）。

\par\textbf{検討未確定：} 公式選択は維持するが平板載荷と液状化を絶対的に無関係とする定義ではない。土研・研究論文には液状化後の反力を載荷で検討する例がある。

\paragraph{問題の要旨} 液状化の動的評価と静的な平板載荷を相対比較する。

\Needspace{4\baselineskip}\paragraph{確認した条件・記述（独自要約）}\mbox{}\par

\noindent\textbf{原典照合：} 22原典画像で25問の全五肢・数値・単位・否定を確認。\par

\noindent\textbf{対象：} 液状化の動的評価と静的な平板載荷を相対比較する。\par

\paragraph{問われる核心} 液状化の動的評価と静的な平板載荷を相対比較する。

\paragraph{必要知識} 選択肢5は通常の静的な平板載荷試験。地震時の繰返しせん断、密度を保った試料、地盤応答を調べる動的試験・PS検層と主目的が違う。液状化後の支持力等の研究では載荷試験も使われるため、原典の相対比較として答える。

\paragraph{考え方・正答を選ぶ理由（査読待ち初稿）} 選択肢5は通常の静的な平板載荷試験。地震時の繰返しせん断、密度を保った試料、地盤応答を調べる動的試験・PS検層と主目的が違う。液状化後の支持力等の研究では載荷試験も使われるため、原典の相対比較として答える。

\Needspace{4\baselineskip}\paragraph{選択肢の検討（査読待ち初稿）}\mbox{}\par

\noindent\textbf{1：代表的条件で適切} 乱さない試料採取。密度・粒子骨格等の状態をなるべく保持し、動的試験や液状化抵抗の評価へつなぐ。採取だけで発生の有無を直接決めるものではない。\par

\noindent\textbf{2：代表的条件で適切} 動的変形特性試験。繰返し載荷に対する剛性・減衰等を測り、地震応答を検討する。特性試験と液状化強度試験の目的は区別する。\par

\noindent\textbf{3：代表的条件で適切} PS検層。P・S波速度から地盤の動的性質・応答を検討する。波速だけで液状化を一意に判定しない。\par

\noindent\textbf{4：代表的条件で適切} 繰返しせん断力。排水が追いつかない飽和土で過剰間隙水圧が上がり、有効応力とせん断抵抗が低下する機構に関係する。\par

\noindent\textbf{5：条件・全称を留保} 平板載荷試験。通常は静的な支持力・荷重沈下関係を調べるため、候補中では液状化発生の評価との直接関係が最も薄い。ただし液状化後の地盤反力を載荷で調べた実験もあり、全く無関係とはしない。\par

\paragraph{再利用する考え方} 直接測る量と代理指標を区別し、分母・有効応力・水の経路・材料の条件から判断する。

\paragraph{共通理論} \hyperref[theory:engineering:liquefaction]{第\ref*{theory:engineering:liquefaction}節}

\par 参考：\cite{eng52_liquefaction,eng52_plate_liq,geebec_ps}。

\paragraph{関連問題} \hyperref[q:h28-reg-applied-33]{H28-REG-APPLIED-33}

\par 難易度 3（校準初稿）、優先度未評価。

\subsection{平成28年度 問33}
\label{q:h28-reg-applied-33}

\noindent\textbf{原典：} sources/H28/17\_applied\_science.pdf、PDF物理13頁。

\textbf{公式正答：} 1（single）。

\par\textbf{検討未確定：} 公式選択は維持するが平板載荷と液状化を絶対的に無関係とする定義ではない。土研・研究論文には液状化後の反力を載荷で検討する例がある。

\paragraph{問題の要旨} 液状化の動的評価と静的な平板載荷を相対比較する。

\Needspace{4\baselineskip}\paragraph{確認した条件・記述（独自要約）}\mbox{}\par

\noindent\textbf{原典照合：} 22原典画像で25問の全五肢・数値・単位・否定を確認。\par

\noindent\textbf{対象：} 液状化の動的評価と静的な平板載荷を相対比較する。\par

\paragraph{問われる核心} 液状化の動的評価と静的な平板載荷を相対比較する。

\paragraph{必要知識} 選択肢1は通常の静的な平板載荷試験。地震時の繰返しせん断、密度を保った試料、地盤応答を調べる動的試験・PS検層と主目的が違う。液状化後の支持力等の研究では載荷試験も使われるため、原典の相対比較として答える。

\paragraph{考え方・正答を選ぶ理由（査読待ち初稿）} 選択肢1は通常の静的な平板載荷試験。地震時の繰返しせん断、密度を保った試料、地盤応答を調べる動的試験・PS検層と主目的が違う。液状化後の支持力等の研究では載荷試験も使われるため、原典の相対比較として答える。

\Needspace{4\baselineskip}\paragraph{選択肢の検討（査読待ち初稿）}\mbox{}\par

\noindent\textbf{1：条件・全称を留保} 平板載荷試験。通常は静的な支持力・荷重沈下関係を調べるため、候補中では液状化発生の評価との直接関係が最も薄い。ただし液状化後の地盤反力を載荷で調べた実験もあり、全く無関係とはしない。\par

\noindent\textbf{2：代表的条件で適切} 動的変形特性試験。繰返し載荷に対する剛性・減衰等を測り、地震応答を検討する。特性試験と液状化強度試験の目的は区別する。\par

\noindent\textbf{3：代表的条件で適切} PS検層。P・S波速度から地盤の動的性質・応答を検討する。波速だけで液状化を一意に判定しない。\par

\noindent\textbf{4：代表的条件で適切} 繰返しせん断力。排水が追いつかない飽和土で過剰間隙水圧が上がり、有効応力とせん断抵抗が低下する機構に関係する。\par

\noindent\textbf{5：代表的条件で適切} 乱さない試料採取。密度・粒子骨格等の状態をなるべく保持し、動的試験や液状化抵抗の評価へつなぐ。採取だけで発生の有無を直接決めるものではない。\par

\paragraph{再利用する考え方} 直接測る量と代理指標を区別し、分母・有効応力・水の経路・材料の条件から判断する。

\paragraph{共通理論} \hyperref[theory:engineering:liquefaction]{第\ref*{theory:engineering:liquefaction}節}

\par 参考：\cite{eng52_liquefaction,eng52_plate_liq,geebec_ps}。

\paragraph{関連問題} \hyperref[q:h23-reg-applied-33]{H23-REG-APPLIED-33}

\par 難易度 3（校準初稿）、優先度未評価。

\section{せん断包絡線の切片と傾き}

\subsection{平成29年度 問22}
\label{q:h29-reg-applied-22}

\noindent\textbf{原典：} sources/H29/17\_applied\_science.pdf、PDF物理7頁。

\textbf{公式正答：} 2（single）。

\paragraph{問題の要旨} 一面せん断の直線包絡線から切片と摩擦角を求める。

\Needspace{4\baselineskip}\paragraph{確認した条件・記述（独自要約）}\mbox{}\par

\noindent\textbf{原典照合：} H29応用原典7頁の四測定点・指数4・全五候補を画像確認。\par

\noindent\textbf{条件：} Mohr–Coulombの破壊包絡線を直線とする。指定tan25°=0.47、tan30°=0.58、tan35°=0.70。\par

\noindent\textbf{単位：} σ・τ・cは×10\textsuperscript{4}N/m\textsuperscript{2}。全応力又は有効応力という区別は原典に明示されず、勝手にuの値を追加しない。\par

\paragraph{問われる核心} 一面せん断の直線包絡線から切片と摩擦角を求める。

\paragraph{必要知識} 原典の測定点(σ,τ)は共通単位×10\textsuperscript{4}N/m\textsuperscript{2}で(1,3.1),(2,3.7),(4,4.8),(8,7.1)。端点の傾き(7.1−3.1)/(8−1)=\allowbreak{}0.5714で、指定0.58の30度に近い。φ=30度で各点からc=τ−0.58σを求めると約2.5×10\textsuperscript{4}Paなので2。丸め値は厳密に一つの直線上ではない。最小二乗でも傾き0.5696、切片2.5391となり同じ候補が最も近い。

\paragraph{考え方（査読待ち初稿）} 原典の測定点(σ,τ)は共通単位×10\textsuperscript{4}N/m\textsuperscript{2}で(1, 3.1)、(2, 3.7)、(4, 4.8)、(8, 7.1)。端点の傾き(7.1−3.1)/(8−1)=\allowbreak{}0.5714で、指定0.58の30度に近い。φ=30度で各点からc=τ−0.58σを求めると約2.5×10\textsuperscript{4}Paなので2。丸め値は厳密に一つの直線上ではない。最小二乗でも傾き0.5696、切片2.5391となり同じ候補が最も近い。

\paragraph{正答を選ぶ理由} 原典の測定点(σ,τ)は共通単位×10\textsuperscript{4}N/m\textsuperscript{2}で(1,3.1),(2,3.7),(4,4.8),(8,7.1)。端点の傾き(7.1−3.1)/(8−1)=\allowbreak{}0.5714で、指定0.58の30度に近い。φ=30度で各点からc=τ−0.58σを求めると約2.5×10\textsuperscript{4}Paなので2。丸め値は厳密に一つの直線上ではない。最小二乗でも傾き0.5696、切片2.5391となり同じ候補が最も近い。

\Needspace{4\baselineskip}\paragraph{選択肢の検討（査読待ち初稿）}\mbox{}\par

\noindent\textbf{1：不適切} c=\allowbreak{}2.5×10\textsuperscript{4}Pa、φ=25°。切片は近いが指定tan25°=0.47は測定の傾き約0.57〜0.60から離れる。\par

\noindent\textbf{2：代表的条件で適切} c=\allowbreak{}2.5×10\textsuperscript{4}Pa、φ=30°。指定tan30°=0.58、直線τ=c+σtanφから各点のcは2.52・2.54・2.48・2.46（×10\textsuperscript{4}Pa）で、丸めを含めて2.5と整合。\par

\noindent\textbf{3：不適切} c=\allowbreak{}2.0×10\textsuperscript{4}Pa、φ=25°。切片も傾きも測定と不一致。σ=1で予測τ=2.47に対し3.1。\par

\noindent\textbf{4：不適切} c=\allowbreak{}2.0×10\textsuperscript{4}Pa、φ=30°。傾きは近いがσ=1で予測τ=2.58に対し3.1。切片が不足。\par

\noindent\textbf{5：不適切} c=\allowbreak{}2.0×10\textsuperscript{4}Pa、φ=35°。傾き0.70は大きすぎ、低応力の切片も不足。全測定点で比較する。\par

\paragraph{再利用する考え方} 直接測る量と代理指標を区別し、分母・有効応力・水の経路・材料の条件から判断する。

\paragraph{共通理論} \hyperref[theory:engineering:mohr]{第\ref*{theory:engineering:mohr}節}

\par 参考：\cite{eng52_mohr_ref}。

\par 難易度 3（校準初稿）、優先度未評価。

\section{含水比・間隙率・飽和度の分母}

\subsection{平成25年度 問28}
\label{q:h25-reg-applied-28}

\noindent\textbf{原典：} sources/H25/17\_applied\_science.pdf、PDF物理10頁。

\textbf{公式正答：} 2（single）。

\paragraph{問題の要旨} 含水比・体積含水率・飽和度の分母を区別する。

\Needspace{4\baselineskip}\paragraph{確認した条件・記述（独自要約）}\mbox{}\par

\noindent\textbf{原典照合：} H25原典10頁で全五候補と記号を確認。\par

\noindent\textbf{指定：} 湿潤質量M、全体積V、水質量Mw、水体積Vw、110±5℃炉乾燥質量Ms、間隙体積Vφ。\par

\paragraph{問われる核心} 含水比・体積含水率・飽和度の分母を区別する。

\paragraph{必要知識} 原典指定の炉乾燥質量Msを分母にして水質量Mwを割るので2。水が乾燥土より重ければ含水比は100％を超え得る。分母が湿潤質量なら100％以内となる別量である。乾燥温度110±5℃も原典の定義として保持し、別資料の105℃測定へ置き換えない。

\paragraph{考え方・正答を選ぶ理由（査読待ち初稿）} 原典指定の炉乾燥質量Msを分母にして水質量Mwを割るので2。水が乾燥土より重ければ含水比は100％を超え得る。分母が湿潤質量なら100％以内となる別量である。乾燥温度110±5℃も原典の定義として保持し、別資料の105℃測定へ置き換えない。

\Needspace{4\baselineskip}\paragraph{選択肢の検討（査読待ち初稿）}\mbox{}\par

\noindent\textbf{1：不適切} Mw/M×100。湿潤試料全質量が分母で、乾燥土質量を分母とする土質の含水比とは違う。\par

\noindent\textbf{2：代表的条件で適切} Mw/Ms×100。水の質量を乾燥土粒子の質量で割る含水比。百分率で示すため100を掛ける。\par

\noindent\textbf{3：不適切} Mw/V×100。質量/体積で密度の次元がある。百分率の無次元比ではない。\par

\noindent\textbf{4：不適切} Vw/V×100。全体積に対する水体積の体積含水率。質量の含水比ではない。\par

\noindent\textbf{5：不適切} Vw/Vφ×100。間隙全体積に対する水体積で飽和度。試料全体積の体積含水率とも違う。\par

\paragraph{再利用する考え方} 直接測る量と代理指標を区別し、分母・有効応力・水の経路・材料の条件から判断する。

\paragraph{共通理論} \hyperref[theory:engineering:phases]{第\ref*{theory:engineering:phases}節}

\par 参考：\cite{eng52_soil_phases}。

\paragraph{関連問題} \hyperref[q:h28-reg-applied-24]{H28-REG-APPLIED-24}

\par 難易度 3（校準初稿）、優先度未評価。

\subsection{平成28年度 問24}
\label{q:h28-reg-applied-24}

\noindent\textbf{原典：} sources/H28/17\_applied\_science.pdf、PDF物理10頁。

\textbf{公式正答：} 3（single）。

\paragraph{問題の要旨} 三相の体積から間隙率と体積含水率を組み立てる。

\Needspace{4\baselineskip}\paragraph{確認した条件・記述（独自要約）}\mbox{}\par

\noindent\textbf{原典照合：} H28原典10頁で組合せ表全行を確認。\par

\noindent\textbf{記号：} Vs固相、Vw水、Va空気、V=\allowbreak{}Vs+Vw+Va。候補の百分率係数100は共通。\par

\paragraph{問われる核心} 三相の体積から間隙率と体積含水率を組み立てる。

\paragraph{必要知識} 全体積V=\allowbreak{}Vs+Vw+Vaと置く。間隙率n=\allowbreak{}(Vw+Va)/V、体積含水率θ=Vw/Vで、それぞれ百分率なら100を掛ける。両方の分母が全体積となる3を選ぶ。間隙比eと飽和度Srは別分母である。

\paragraph{考え方・正答を選ぶ理由（査読待ち初稿）} 全体積V=\allowbreak{}Vs+Vw+Vaと置く。間隙率n=\allowbreak{}(Vw+Va)/V、体積含水率θ=Vw/Vで、それぞれ百分率なら100を掛ける。両方の分母が全体積となる3を選ぶ。間隙比eと飽和度Srは別分母である。

\Needspace{4\baselineskip}\paragraph{選択肢の検討（査読待ち初稿）}\mbox{}\par

\noindent\textbf{1：不適切} A=\allowbreak{}(Vw+Va)/V、B=\allowbreak{}Vw/Vs。Aは間隙率だがBの分母が固相で、全体積の体積含水率ではない。\par

\noindent\textbf{2：不適切} A=\allowbreak{}(Vw+Va)/V、B=\allowbreak{}Vw/(Vw+Va)。Bは飽和度。\par

\noindent\textbf{3：代表的条件で適切} A=\allowbreak{}(Vw+Va)/V、B=\allowbreak{}Vw/V。間隙率も体積含水率も全体積V=\allowbreak{}Vs+Vw+Vaが分母。Aの分子は水と空気、Bの分子は水。\par

\noindent\textbf{4：不適切} A=\allowbreak{}(Vw+Va)/Vs、B=\allowbreak{}Vw/(Vw+Va)。Aは間隙比、Bは飽和度で両方違う。\par

\noindent\textbf{5：不適切} A=\allowbreak{}(Vw+Va)/Vs、B=\allowbreak{}Vw/V。Bは正しいがAは間隙比。\par

\paragraph{再利用する考え方} 直接測る量と代理指標を区別し、分母・有効応力・水の経路・材料の条件から判断する。

\paragraph{共通理論} \hyperref[theory:engineering:phases]{第\ref*{theory:engineering:phases}節}

\par 参考：\cite{eng52_soil_phases}。

\paragraph{関連問題} \hyperref[q:h25-reg-applied-28]{H25-REG-APPLIED-28}

\par 難易度 3（校準初稿）、優先度未評価。

\section{黄鉄鉱の酸化・溶解・塩析出}

\subsection{平成24年度 問30}
\label{q:h24-reg-applied-30}

\noindent\textbf{原典：} sources/H24/17\_applied\_science.pdf、PDF物理9頁。

\textbf{公式正答：} 5（single）。

\par\textbf{検討未確定：} 酸化フロントにすべりが多いという頻度、泥岩の全岩種・全層序への一般化は未確定。未酸化試料の掘削後酸化と硫酸塩析出を時間順で補足。

\paragraph{問題の要旨} 海成泥岩の黄鉄鉱酸化・酸性排水・析出を区別する。

\Needspace{4\baselineskip}\paragraph{確認した条件・記述（独自要約）}\mbox{}\par

\noindent\textbf{原典照合：} 22原典画像で25問の全五肢・数値・単位・否定を確認。\par

\noindent\textbf{対象：} 海成泥岩の黄鉄鉱酸化・酸性排水・析出を区別する。\par

\paragraph{問われる核心} 海成泥岩の黄鉄鉱酸化・酸性排水・析出を区別する。

\paragraph{必要知識} 誤り5を、黄鉄鉱の酸化が生む酸と弱化の過程から説明する。生成した硫酸イオンがCaと結び石膏等で析出すると盤膨れの別過程になり得る。層理・凝灰岩の弱層・地下水経路を合わせ、岩種と酸化だけで全地点のすべりを決めない。

\paragraph{考え方・正答を選ぶ理由（査読待ち初稿）} 誤り5を、黄鉄鉱の酸化が生む酸と弱化の過程から説明する。生成した硫酸イオンがCaと結び石膏等で析出すると盤膨れの別過程になり得る。層理・凝灰岩の弱層・地下水経路を合わせ、岩種と酸化だけで全地点のすべりを決めない。

\Needspace{4\baselineskip}\paragraph{選択肢の検討（査読待ち初稿）}\mbox{}\par

\noindent\textbf{1：代表的条件で適切} すべり面が層理に平行。弱い層・連続する面の向きがすべりの経路に関係する。全すべり面が層理という限定ではない。\par

\noindent\textbf{2：代表的条件で適切} 凝灰岩にすべり面。砂岩泥岩の間にある変質・軟化しやすい凝灰岩等が弱層となることがある。\par

\noindent\textbf{3：条件・全称を留保} 酸化フロント・下の溶解帯。黄鉄鉱の酸化と酸性水による膠結物の溶解は弱化へつながる。原典の多いという頻度と特定のすべり位置の全称は追加照合対象。\par

\noindent\textbf{4：代表的条件で適切} 未酸化青黒泥岩と硫酸塩・石膏の盤膨れ。新鮮な還元状態の岩石が掘削後に水・酸素と反応し、硫酸塩が析出して結晶成長圧が生じ得る。未酸化のままで酸化反応が進むという矛盾した理解は避ける。\par

\noindent\textbf{5：不適切} 黄鉄鉱でアルカリ性排水。黄鉄鉱の酸化は酸を生む。炭酸塩等で中和されることはあっても、その酸化自体をアルカリ化の起源とはしない。\par

\paragraph{再利用する考え方} 直接測る量と代理指標を区別し、分母・有効応力・水の経路・材料の条件から判断する。

\paragraph{共通理論} \hyperref[theory:engineering:pyrite]{第\ref*{theory:engineering:pyrite}節}

\par 参考：\cite{eng52_weathering,eng52_slip_weaklayer}。

\paragraph{関連問題} \hyperref[q:r01-reg-applied-29]{R01-REG-APPLIED-29}

\par 難易度 3（校準初稿）、優先度未評価。

\subsection{令和元年度 問29}
\label{q:r01-reg-applied-29}

\noindent\textbf{原典：} sources/R01/17\_applied\_science.pdf、PDF物理10頁。

\textbf{公式正答：} 5（single）。

\par\textbf{検討未確定：} 酸化フロントにすべりが多いという頻度、泥岩の全岩種・全層序への一般化は未確定。未酸化試料の掘削後酸化と硫酸塩析出を時間順で補足。

\paragraph{問題の要旨} 海成泥岩の黄鉄鉱酸化・酸性排水・析出を区別する。

\Needspace{4\baselineskip}\paragraph{確認した条件・記述（独自要約）}\mbox{}\par

\noindent\textbf{原典照合：} 22原典画像で25問の全五肢・数値・単位・否定を確認。\par

\noindent\textbf{対象：} 海成泥岩の黄鉄鉱酸化・酸性排水・析出を区別する。\par

\paragraph{問われる核心} 海成泥岩の黄鉄鉱酸化・酸性排水・析出を区別する。

\paragraph{必要知識} 誤り5を、黄鉄鉱の酸化が生む酸と弱化の過程から説明する。生成した硫酸イオンがCaと結び石膏等で析出すると盤膨れの別過程になり得る。層理・凝灰岩の弱層・地下水経路を合わせ、岩種と酸化だけで全地点のすべりを決めない。

\paragraph{考え方・正答を選ぶ理由（査読待ち初稿）} 誤り5を、黄鉄鉱の酸化が生む酸と弱化の過程から説明する。生成した硫酸イオンがCaと結び石膏等で析出すると盤膨れの別過程になり得る。層理・凝灰岩の弱層・地下水経路を合わせ、岩種と酸化だけで全地点のすべりを決めない。

\Needspace{4\baselineskip}\paragraph{選択肢の検討（査読待ち初稿）}\mbox{}\par

\noindent\textbf{1：代表的条件で適切} 未酸化青黒泥岩と硫酸塩・石膏の盤膨れ。新鮮な還元状態の岩石が掘削後に水・酸素と反応し、硫酸塩が析出して結晶成長圧が生じ得る。未酸化のままで酸化反応が進むという矛盾した理解は避ける。\par

\noindent\textbf{2：代表的条件で適切} 黄鉄鉱で酸性排水の可能性。酸素・水と反応すると硫酸イオンと酸を生む。反応量・緩衝能で実際の排水pHは変わる。\par

\noindent\textbf{3：代表的条件で適切} 凝灰岩にすべり面。砂岩泥岩の間にある変質・軟化しやすい凝灰岩等が弱層となることがある。\par

\noindent\textbf{4：代表的条件で適切} すべり面が層理に平行。弱い層・連続する面の向きがすべりの経路に関係する。全すべり面が層理という限定ではない。\par

\noindent\textbf{5：不適切} 酸化フロントが地すべりを抑制。酸化・溶解風化で結合が弱まり、すべりへつながり得るため、一方向の抑制とはしない。酸化生成物の間隙充填等がある場合も条件を分ける。\par

\paragraph{再利用する考え方} 直接測る量と代理指標を区別し、分母・有効応力・水の経路・材料の条件から判断する。

\paragraph{共通理論} \hyperref[theory:engineering:pyrite]{第\ref*{theory:engineering:pyrite}節}

\par 参考：\cite{eng52_weathering,eng52_slip_weaklayer}。

\paragraph{関連問題} \hyperref[q:h24-reg-applied-30]{H24-REG-APPLIED-30}

\par 難易度 3（校準初稿）、優先度未評価。

\section{地すべり・崩壊の代表比較}

\subsection{平成30年度 問24}
\label{q:h30-reg-applied-24}

\noindent\textbf{原典：} sources/H30/17\_applied\_science.pdf、PDF物理8頁。

\textbf{公式正答：} 3（single）。

\par\textbf{検討未確定：} 原典の塑性/脆性、緩慢/高速、砂質/粘性の全材料への排他性は未確認。公式選択と実際の多様な移動機構を分けて要査読とする。

\paragraph{問題の要旨} 地すべりと崩壊の代表比較を全称にしない。

\Needspace{4\baselineskip}\paragraph{確認した条件・記述（独自要約）}\mbox{}\par

\noindent\textbf{原典照合：} 22原典画像で25問の全五肢・数値・単位・否定を確認。\par

\noindent\textbf{対象：} 地すべりと崩壊の代表比較を全称にしない。\par

\paragraph{問われる核心} 地すべりと崩壊の代表比較を全称にしない。

\paragraph{必要知識} 公式3は砂質・粘性土の典型対応を逆にした候補。移動の継続性、弱面、水圧、降雨・地震、兆候の観察を合わせる。速度と破壊形式は代表比較であり、全事例を二種類へ固定する定義ではない。

\paragraph{考え方・正答を選ぶ理由（査読待ち初稿）} 公式3は砂質・粘性土の典型対応を逆にした候補。移動の継続性、弱面、水圧、降雨・地震、兆候の観察を合わせる。速度と破壊形式は代表比較であり、全事例を二種類へ固定する定義ではない。

\Needspace{4\baselineskip}\paragraph{選択肢の検討（査読待ち初稿）}\mbox{}\par

\noindent\textbf{1：条件・全称を留保} 地すべり塑性・崩壊脆性。原典の典型比較として保持するが、全地すべりを塑性・全崩壊を脆性という二分法で定義しない。せん断面の滑動と岩塊の移動には多様な機構がある。\par

\noindent\textbf{2：条件・全称を留保} 地すべり緩慢・崩壊高速。代表的な観察上の比較。地すべりにも急激な滑動があるため、名称から全速度を固定しない。\par

\noindent\textbf{3：不適切} 地すべり砂質・崩壊粘性。通常の教材の代表比較では、地すべりは粘性土等の弱いすべり面、崩壊は表層の砂質土等と関係し、この対応は逆。全材料を排他的に分類する規則ではない。\par

\noindent\textbf{4：代表的条件で適切} 地すべり地下水・崩壊豪雨地震。地下水圧や雨・地震が安定を低下させる。地すべりにも豪雨・地震が働き、唯一の引き金を分けるものではない。\par

\noindent\textbf{5：代表的条件で適切} 地すべりに亀裂・隆起・陥没の兆候。長期滑動等では変形の兆候を認める場合がある。崩壊に兆候が全くない、地すべりなら必ず発見できるとはしない。\par

\paragraph{再利用する考え方} 直接測る量と代理指標を区別し、分母・有効応力・水の経路・材料の条件から判断する。

\paragraph{共通理論} \hyperref[theory:engineering:landslide]{第\ref*{theory:engineering:landslide}節}

\par 参考：\cite{eng52_landslide_ref}。

\paragraph{関連問題} \hyperref[q:r04-reg-applied-27]{R04-REG-APPLIED-27}

\par 難易度 3（校準初稿）、優先度未評価。

\subsection{令和4年度 問27}
\label{q:r04-reg-applied-27}

\noindent\textbf{原典：} sources/R04/17\_applied\_science.pdf、PDF物理9頁。

\textbf{公式正答：} 4（single）。

\par\textbf{検討未確定：} 原典の塑性/脆性、緩慢/高速、砂質/粘性の全材料への排他性は未確認。公式選択と実際の多様な移動機構を分けて要査読とする。

\paragraph{問題の要旨} 地すべりと崩壊の代表比較を全称にしない。

\Needspace{4\baselineskip}\paragraph{確認した条件・記述（独自要約）}\mbox{}\par

\noindent\textbf{原典照合：} 22原典画像で25問の全五肢・数値・単位・否定を確認。\par

\noindent\textbf{対象：} 地すべりと崩壊の代表比較を全称にしない。\par

\paragraph{問われる核心} 地すべりと崩壊の代表比較を全称にしない。

\paragraph{必要知識} 公式4は砂質・粘性土の典型対応を逆にした候補。移動の継続性、弱面、水圧、降雨・地震、兆候の観察を合わせる。速度と破壊形式は代表比較であり、全事例を二種類へ固定する定義ではない。

\paragraph{考え方・正答を選ぶ理由（査読待ち初稿）} 公式4は砂質・粘性土の典型対応を逆にした候補。移動の継続性、弱面、水圧、降雨・地震、兆候の観察を合わせる。速度と破壊形式は代表比較であり、全事例を二種類へ固定する定義ではない。

\Needspace{4\baselineskip}\paragraph{選択肢の検討（査読待ち初稿）}\mbox{}\par

\noindent\textbf{1：条件・全称を留保} 地すべり塑性・崩壊脆性。原典の典型比較として保持するが、全地すべりを塑性・全崩壊を脆性という二分法で定義しない。せん断面の滑動と岩塊の移動には多様な機構がある。\par

\noindent\textbf{2：代表的条件で適切} 地すべりに亀裂・隆起・陥没の兆候。長期滑動等では変形の兆候を認める場合がある。崩壊に兆候が全くない、地すべりなら必ず発見できるとはしない。\par

\noindent\textbf{3：代表的条件で適切} 地すべり地下水・崩壊豪雨地震。地下水圧や雨・地震が安定を低下させる。地すべりにも豪雨・地震が働き、唯一の引き金を分けるものではない。\par

\noindent\textbf{4：不適切} 地すべり砂質・崩壊粘性。通常の教材の代表比較では、地すべりは粘性土等の弱いすべり面、崩壊は表層の砂質土等と関係し、この対応は逆。全材料を排他的に分類する規則ではない。\par

\noindent\textbf{5：条件・全称を留保} 地すべり緩慢・崩壊高速。代表的な観察上の比較。地すべりにも急激な滑動があるため、名称から全速度を固定しない。\par

\paragraph{再利用する考え方} 直接測る量と代理指標を区別し、分母・有効応力・水の経路・材料の条件から判断する。

\paragraph{共通理論} \hyperref[theory:engineering:landslide]{第\ref*{theory:engineering:landslide}節}

\par 参考：\cite{eng52_landslide_ref}。

\paragraph{関連問題} \hyperref[q:h30-reg-applied-24]{H30-REG-APPLIED-24}

\par 難易度 3（校準初稿）、優先度未評価。

\section{せん断強度・貫入と圧密の測定量}

\subsection{平成23年度 問34}
\label{q:h23-reg-applied-34}

\noindent\textbf{原典：} sources/H23/17\_applied\_science.pdf、PDF物理11頁。

\textbf{公式正答：} 3（single）。

\paragraph{問題の要旨} 圧密の体積変形とせん断強度試験の測定量を分ける。

\Needspace{4\baselineskip}\paragraph{確認した条件・記述（独自要約）}\mbox{}\par

\noindent\textbf{原典照合：} 22原典画像で25問の全五肢・数値・単位・否定を確認。\par

\noindent\textbf{対象：} 圧密の体積変形とせん断強度試験の測定量を分ける。\par

\paragraph{問われる核心} 圧密の体積変形とせん断強度試験の測定量を分ける。

\paragraph{必要知識} 最も関係の薄い候補3は圧密試験。せん断破壊、回転トルク、貫入抵抗を強度・支持性へ結び付ける他の試験と、圧密の圧縮・排水・時間の主目的を分ける。圧密と強度の状態依存まで無関係とはしない。

\paragraph{考え方・正答を選ぶ理由（査読待ち初稿）} 最も関係の薄い候補3は圧密試験。せん断破壊、回転トルク、貫入抵抗を強度・支持性へ結び付ける他の試験と、圧密の圧縮・排水・時間の主目的を分ける。圧密と強度の状態依存まで無関係とはしない。

\Needspace{4\baselineskip}\paragraph{選択肢の検討（査読待ち初稿）}\mbox{}\par

\noindent\textbf{1：代表的条件で適切} 三軸又は一軸の力学試験。拘束圧・排水条件を指定して試料を破壊へ載荷し、強度や応力ひずみ応答を調べる。\par

\noindent\textbf{2：代表的条件で適切} ベーンせん断試験。羽根を回すトルク等から軟らかい粘性土の非排水せん断抵抗を評価する。\par

\noindent\textbf{3：不適切} 圧密試験。主目的は荷重増加に対する時間依存の体積変化・圧縮性・排水の評価。せん断破壊強度を直接測る試験とは違う。\par

\noindent\textbf{4：代表的条件で適切} スウェーデン式サウンディング。貫入抵抗から地盤の硬軟や支持性・土質を経験的に推定する。推定量と直接破壊試験を区別する。\par

\noindent\textbf{5：代表的条件で適切} 標準貫入試験。N値は所定貫入の打撃数で、密度・硬軟等の経験的評価に使う。N値をせん断強度の単位としない。\par

\paragraph{再利用する考え方} 直接測る量と代理指標を区別し、分母・有効応力・水の経路・材料の条件から判断する。

\paragraph{共通理論} \hyperref[theory:engineering:tests]{第\ref*{theory:engineering:tests}節}

\par 参考：\cite{eng52_geotech_tests,eng52_consolidation}。

\paragraph{関連問題} \hyperref[q:h26-reg-applied-34]{H26-REG-APPLIED-34}

\par 難易度 2（校準初稿）、優先度未評価。

\subsection{平成26年度 問34}
\label{q:h26-reg-applied-34}

\noindent\textbf{原典：} sources/H26/17\_applied\_science.pdf、PDF物理14頁。

\textbf{公式正答：} 1（single）。

\paragraph{問題の要旨} 圧密の体積変形とせん断強度試験の測定量を分ける。

\Needspace{4\baselineskip}\paragraph{確認した条件・記述（独自要約）}\mbox{}\par

\noindent\textbf{原典照合：} 22原典画像で25問の全五肢・数値・単位・否定を確認。\par

\noindent\textbf{対象：} 圧密の体積変形とせん断強度試験の測定量を分ける。\par

\paragraph{問われる核心} 圧密の体積変形とせん断強度試験の測定量を分ける。

\paragraph{必要知識} 最も関係の薄い候補1は圧密試験。せん断破壊、回転トルク、貫入抵抗を強度・支持性へ結び付ける他の試験と、圧密の圧縮・排水・時間の主目的を分ける。圧密と強度の状態依存まで無関係とはしない。

\paragraph{考え方・正答を選ぶ理由（査読待ち初稿）} 最も関係の薄い候補1は圧密試験。せん断破壊、回転トルク、貫入抵抗を強度・支持性へ結び付ける他の試験と、圧密の圧縮・排水・時間の主目的を分ける。圧密と強度の状態依存まで無関係とはしない。

\Needspace{4\baselineskip}\paragraph{選択肢の検討（査読待ち初稿）}\mbox{}\par

\noindent\textbf{1：不適切} 圧密試験。主目的は荷重増加に対する時間依存の体積変化・圧縮性・排水の評価。せん断破壊強度を直接測る試験とは違う。\par

\noindent\textbf{2：代表的条件で適切} 三軸又は一軸の力学試験。拘束圧・排水条件を指定して試料を破壊へ載荷し、強度や応力ひずみ応答を調べる。\par

\noindent\textbf{3：代表的条件で適切} 標準貫入試験。N値は所定貫入の打撃数で、密度・硬軟等の経験的評価に使う。N値をせん断強度の単位としない。\par

\noindent\textbf{4：代表的条件で適切} ベーンせん断試験。羽根を回すトルク等から軟らかい粘性土の非排水せん断抵抗を評価する。\par

\noindent\textbf{5：代表的条件で適切} スウェーデン式サウンディング。貫入抵抗から地盤の硬軟や支持性・土質を経験的に推定する。推定量と直接破壊試験を区別する。\par

\paragraph{再利用する考え方} 直接測る量と代理指標を区別し、分母・有効応力・水の経路・材料の条件から判断する。

\paragraph{共通理論} \hyperref[theory:engineering:tests]{第\ref*{theory:engineering:tests}節}

\par 参考：\cite{eng52_geotech_tests,eng52_consolidation}。

\paragraph{関連問題} \hyperref[q:h23-reg-applied-34]{H23-REG-APPLIED-34}, \hyperref[q:h30-reg-applied-27]{H30-REG-APPLIED-27}

\par 難易度 2（校準初稿）、優先度未評価。

\subsection{平成30年度 問27}
\label{q:h30-reg-applied-27}

\noindent\textbf{原典：} sources/H30/17\_applied\_science.pdf、PDF物理9頁。

\textbf{公式正答：} 1（single）。

\paragraph{問題の要旨} 圧密の体積変形とせん断強度試験の測定量を分ける。

\Needspace{4\baselineskip}\paragraph{確認した条件・記述（独自要約）}\mbox{}\par

\noindent\textbf{原典照合：} 22原典画像で25問の全五肢・数値・単位・否定を確認。\par

\noindent\textbf{対象：} 圧密の体積変形とせん断強度試験の測定量を分ける。\par

\paragraph{問われる核心} 圧密の体積変形とせん断強度試験の測定量を分ける。

\paragraph{必要知識} 最も関係の薄い候補1は圧密試験。せん断破壊、回転トルク、貫入抵抗を強度・支持性へ結び付ける他の試験と、圧密の圧縮・排水・時間の主目的を分ける。圧密と強度の状態依存まで無関係とはしない。

\paragraph{考え方・正答を選ぶ理由（査読待ち初稿）} 最も関係の薄い候補1は圧密試験。せん断破壊、回転トルク、貫入抵抗を強度・支持性へ結び付ける他の試験と、圧密の圧縮・排水・時間の主目的を分ける。圧密と強度の状態依存まで無関係とはしない。

\Needspace{4\baselineskip}\paragraph{選択肢の検討（査読待ち初稿）}\mbox{}\par

\noindent\textbf{1：不適切} 圧密試験。主目的は荷重増加に対する時間依存の体積変化・圧縮性・排水の評価。せん断破壊強度を直接測る試験とは違う。\par

\noindent\textbf{2：代表的条件で適切} 三軸又は一軸の力学試験。拘束圧・排水条件を指定して試料を破壊へ載荷し、強度や応力ひずみ応答を調べる。\par

\noindent\textbf{3：代表的条件で適切} 標準貫入試験。N値は所定貫入の打撃数で、密度・硬軟等の経験的評価に使う。N値をせん断強度の単位としない。\par

\noindent\textbf{4：代表的条件で適切} ベーンせん断試験。羽根を回すトルク等から軟らかい粘性土の非排水せん断抵抗を評価する。\par

\noindent\textbf{5：代表的条件で適切} スウェーデン式サウンディング。貫入抵抗から地盤の硬軟や支持性・土質を経験的に推定する。推定量と直接破壊試験を区別する。\par

\paragraph{再利用する考え方} 直接測る量と代理指標を区別し、分母・有効応力・水の経路・材料の条件から判断する。

\paragraph{共通理論} \hyperref[theory:engineering:tests]{第\ref*{theory:engineering:tests}節}

\par 参考：\cite{eng52_geotech_tests,eng52_consolidation}。

\paragraph{関連問題} \hyperref[q:h26-reg-applied-34]{H26-REG-APPLIED-34}

\par 難易度 2（校準初稿）、優先度未評価。

\section{揚水水頭低下と粘土圧密沈下}

\subsection{平成25年度 問30}
\label{q:h25-reg-applied-30}

\noindent\textbf{原典：} sources/H25/17\_applied\_science.pdf、PDF物理11頁。

\textbf{公式正答：} 3（single）。

\paragraph{問題の要旨} 揚水による水圧低下と粘土層の圧密沈下をつなぐ。

\Needspace{4\baselineskip}\paragraph{確認した条件・記述（独自要約）}\mbox{}\par

\noindent\textbf{原典照合：} 22原典画像で25問の全五肢・数値・単位・否定を確認。\par

\noindent\textbf{対象：} 揚水による水圧低下と粘土層の圧密沈下をつなぐ。\par

\paragraph{問われる核心} 揚水による水圧低下と粘土層の圧密沈下をつなぐ。

\paragraph{必要知識} 3は不圧水位へ主因を限定した点が不適切。σ\ensuremath{'}=σ−uで、水圧uが下がると骨格の有効応力が増し、圧縮性の高い粘土層の排水・収縮が地表へ現れる。層ごとの水頭と地盤高を観測して時間差も追う。

\paragraph{考え方・正答を選ぶ理由（査読待ち初稿）} 3は不圧水位へ主因を限定した点が不適切。σ\ensuremath{'}=σ−uで、水圧uが下がると骨格の有効応力が増し、圧縮性の高い粘土層の排水・収縮が地表へ現れる。層ごとの水頭と地盤高を観測して時間差も追う。

\Needspace{4\baselineskip}\paragraph{選択肢の検討（査読待ち初稿）}\mbox{}\par

\noindent\textbf{1：代表的条件で適切} 構造物浮き上がり・道路変形・埋設破損。周囲の地盤が沈むと杭等に支えられた構造物が相対的に浮き上がった状態となり、道路・埋設物へ被害が生じる。\par

\noindent\textbf{2：代表的条件で適切} 帯水層に接する粘土も圧密収縮。揚水で水頭・間隙水圧が低下し、同じ全応力のもとで土粒子骨格の有効応力が増加する。\par

\noindent\textbf{3：不適切} 主に不圧地下水位の変動。広域の揚水沈下では被圧帯水層の水頭低下と隣接する粘土層の圧密が重要。不圧水位だけへ限定しない。\par

\noindent\textbf{4：代表的条件で適切} 低平地の浸水・冠水危険増大。標高が低下し排水勾配や治水の条件が悪化し、二次的被害のリスクを高める。\par

\noindent\textbf{5：代表的条件で適切} 帯水層深度ごとの観測井。水頭が異なる層を区別して観測する。複数帯水層をつなぐ井戸の混合水位だけでは各層を評価できない。\par

\paragraph{再利用する考え方} 直接測る量と代理指標を区別し、分母・有効応力・水の経路・材料の条件から判断する。

\paragraph{共通理論} \hyperref[theory:engineering:consolidation]{第\ref*{theory:engineering:consolidation}節}

\par 参考：\cite{eng52_subsidence_ref,eng52_consolidation}。

\par 難易度 2（校準初稿）、優先度未評価。

\section{切羽・偏圧・山はねと地表沈下}

\subsection{令和3年度 問28}
\label{q:r03-reg-applied-28}

\noindent\textbf{原典：} sources/R03/17\_applied\_science.pdf、PDF物理10頁。

\textbf{公式正答：} 2（single）。

\par\textbf{検討未確定：} 緩み土圧と応力損傷の全条件の境界、土かぶりと地表沈下の定量関係は未確定。代表傾向として保持し全工法へ一般化しない。

\paragraph{問題の要旨} 切羽崩壊・偏圧・山はねと土かぶりを区別する。

\Needspace{4\baselineskip}\paragraph{確認した条件・記述（独自要約）}\mbox{}\par

\noindent\textbf{原典照合：} 22原典画像で25問の全五肢・数値・単位・否定を確認。\par

\noindent\textbf{対象：} 切羽崩壊・偏圧・山はねと土かぶりを区別する。\par

\paragraph{問われる核心} 切羽崩壊・偏圧・山はねと土かぶりを区別する。

\paragraph{必要知識} 適切2は切羽の構造・水と不安定化の対応。応力の解放・集中、高応力の硬岩の山はね、割れ目に沿う緩み、非対称な偏圧、地表沈下の各機構を別々に比較する。

\paragraph{考え方・正答を選ぶ理由（査読待ち初稿）} 適切2は切羽の構造・水と不安定化の対応。応力の解放・集中、高応力の硬岩の山はね、割れ目に沿う緩み、非対称な偏圧、地表沈下の各機構を別々に比較する。

\Needspace{4\baselineskip}\paragraph{選択肢の検討（査読待ち初稿）}\mbox{}\par

\noindent\textbf{1：不適切} 割れ目の少ない硬岩の崩壊を緩み土圧。高応力の健全硬岩での損傷・山はねと、割れ目のある地山の重力的な緩みを混同している。現象名を応力・不連続面から分ける。\par

\noindent\textbf{2：代表的条件で適切} 切羽天端・鏡面の構造と湧水による崩壊。不連続面や弱い層、水による切羽の不安定化を観察する。施工中の観測を支保・設計へ反映する。\par

\noindent\textbf{3：不適切} 緩傾斜・水平地層で偏圧。典型的な偏圧は斜面地形や傾斜地層等の非対称な条件で問題になる。水平・対称をその原因へ固定しない。\par

\noindent\textbf{4：不適切} 小土かぶり・多節理で山はね。代表的な山はねは高応力・大土かぶり・硬く節理の少ない岩盤で起こりやすい。原典は条件を逆にした。\par

\noindent\textbf{5：条件・全称を留保} 土かぶりが大きいほど地表沈下大。地表沈下は浅い地山・緩み・水位低下・地層の条件等に依存する。大土かぶりだけで地表沈下が大きくなるという単純な単調関係にしない。\par

\paragraph{再利用する考え方} 直接測る量と代理指標を区別し、分母・有効応力・水の経路・材料の条件から判断する。

\paragraph{共通理論} \hyperref[theory:engineering:tunnel]{第\ref*{theory:engineering:tunnel}節}

\par 参考：\cite{eng52_tunnel_ref,eng52_rockburst}。

\par 難易度 3（校準初稿）、優先度未評価。

\section{岩石試験・岩盤強度の代理指標}

\subsection{令和元年度再試験 問27}
\label{q:r01-ret-applied-27}

\noindent\textbf{原典：} sources/R01\_retest/17\_applied\_science.pdf、PDF物理9頁。

\textbf{公式正答：} 3（single）。

\par\textbf{検討未確定：} 岩盤評価の相対比較である。RQD・打診・波速から強度を一意に決めず、弱岩の圧縮降伏の例外を保持。

\paragraph{問題の要旨} 岩盤強度の直接試験と間接指標、圧密降伏を比較する。

\Needspace{4\baselineskip}\paragraph{確認した条件・記述（独自要約）}\mbox{}\par

\noindent\textbf{原典照合：} 29問21原典画像で否定、全五候補、数式・上付き単位・観測対象を照合。\par

\noindent\textbf{前提：} 岩石試験片の破壊強度と、割れ目を含む岩盤の状態及び土の圧密履歴を区別する。\par

\paragraph{問われる核心} 岩盤強度の直接試験と間接指標、圧密降伏を比較する。

\paragraph{必要知識} 岩石試験片の破壊強度と、割れ目を含む岩盤の状態及び土の圧密履歴を区別する。

\paragraph{考え方・正答を選ぶ理由（査読待ち初稿）} 3の圧密降伏応力である。通常は土の一次元圧縮履歴と降伏を表す量であり、他の項目は岩片強度・割れ目状態・弾性応答など岩盤評価へ直接又は間接に使う。ただし弱い泥岩等の圧縮降伏もあり、関連性を完全ゼロとはしない。

\Needspace{4\baselineskip}\paragraph{選択肢の検討（査読待ち初稿）}\mbox{}\par

\noindent\textbf{1：条件付きで適切} ハンマー打診音。硬さ・風化・割れ目等の現地判別に用いるが、定量強度の直接測定ではなく観察者と条件に依存する。\par

\noindent\textbf{2：適切} 一軸圧縮強度。試験片を側方無拘束で圧縮破壊させた最大応力であり、岩片材料の強度を直接評価する。岩盤の弱面は別である。\par

\noindent\textbf{3：不適切} 圧密降伏応力。通常の土の圧密・応力履歴の指標であり、岩盤強度特性の比較では相対的に関係が薄い。弱岩の圧縮降伏との関連は否定しない。\par

\noindent\textbf{4：条件付きで適切} RQD。健全コアの所定長さ割合により割れ目による岩盤状態を表す。長さ品質指標そのものが一軸強度ではない。\par

\noindent\textbf{5：条件付きで適切} 弾性波速度。弾性率、密度、割れ目、含水と関係し岩盤評価に使う。破壊強度との一意換算ではない。\par

\paragraph{再利用する考え方} 岩片の材料試験と岩盤の割れ目・水・応力条件、代理指標を分ける。

\paragraph{共通理論} \hyperref[theory:engineering:strengthindices]{第\ref*{theory:engineering:strengthindices}節}

\par 参考：\cite{jseg_engineering_faq,geebec_ps}。

\par 難易度 3（校準初稿）、優先度未評価。

\section{岩石風化・変質・スレーキング}

\subsection{平成29年度 問24}
\label{q:h29-reg-applied-24}

\noindent\textbf{原典：} sources/H29/17\_applied\_science.pdf、PDF物理8頁。

\textbf{公式正答：} 3（single）。

\paragraph{問題の要旨} 風化・熱水変質・スレーキングが岩石強度へ及ぼす作用を判定する。

\Needspace{4\baselineskip}\paragraph{確認した条件・記述（独自要約）}\mbox{}\par

\noindent\textbf{方向：} 最も不適切なもの。\par

\noindent\textbf{対象：} 一般的な岩石分類・形成過程。\par

\paragraph{問われる核心} 風化・変質と間隙内圧による骨格破壊

\paragraph{必要知識} 風化・変質と間隙内圧による骨格破壊。母材・運搬堆積・有機物起源、風化、弱面、水の作用、岩盤と土の物性。

\paragraph{考え方・正答を選ぶ理由（査読待ち初稿）} スレーキングは水への浸漬などで岩塊が崩れる現象である。封入空気モデルでは、水の侵入で圧縮された空気の内圧が骨格を引っ張るので、岩塊の圧縮力とした3が不適切である。空気の状態変化と固体の応力を別々に記述する。膨潤・吸着・化学劣化なども機構となり得る。

\Needspace{4\baselineskip}\paragraph{選択肢の検討（査読待ち初稿）}\mbox{}\par

\noindent\textbf{1：適切（一般的分類）} 有色鉱物に富む黒い火成岩ほど一般に風化に不安定である。。高温形成のFe・Mgに富む鉱物は地表環境で変化しやすい。色そのものを原因とする法則ではなく、割れ目や組織も耐風化性を左右する。\par

\noindent\textbf{2：適切（一般的分類）} 磁鉄鉱から赤鉄鉱・褐鉄鉱への酸化が風化層の色に現れる。。Feの酸化と水和等により赤褐色系の鉱物を生じる。褐鉄鉱は単一の理想化学式を持つ一鉱物としてではなく含水酸化鉄の集合の名称として扱う。\par

\noindent\textbf{3：不適切} 浸水で間隙空気が圧縮され、岩塊内に圧縮力を生じてスレーキングする。。空気は圧縮されるが、内圧が孔壁・骨格を押し広げて引張破壊へ働く。空気の圧縮と岩塊への圧縮を混同している。封入空気だけが全試料の機構とも限らない。\par

\noindent\textbf{4：適切（一般的分類）} グリーンタフ地域の新第三紀層は熱水変質で軟弱化していることが多い。。火山性の岩石・堆積物が熱水反応で粘土化等を受け、低強度となる場合がある。全てのグリーンタフが同じ弱さという意味ではない。\par

\noindent\textbf{5：適切（一般的分類）} 一般に風化・変質は弱化させるが、珪化は強度を増すこともある。。シリカの充填・交代は緻密化や結合強化をもたらす場合があり、変質=常に強度低下とはしない。\par

\paragraph{再利用する考え方} 風化・変質と間隙内圧による骨格破壊を成因と観察される性質に結び付け、例示と全称命題、主要鉱物と副成分、公式解答と未解決の解釈を区別する。

\paragraph{共通理論} \hyperref[theory:engineering:weathering]{第\ref*{theory:engineering:weathering}節}, \hyperref[theory:engineering:slaking]{第\ref*{theory:engineering:slaking}節}

\par 参考：\cite{slaking1979,slaking2006,jseg_engineering_faq}。

\par 難易度 3（校準初稿）、優先度未評価。

\section{特殊土の母材・成因}

\subsection{平成25年度 問19}
\label{q:h25-reg-applied-19}

\noindent\textbf{原典：} sources/H25/17\_applied\_science.pdf、PDF物理8頁。

\textbf{公式正答：} 5（single）。

\paragraph{問題の要旨} まさ土・しらす・黄土・火山灰質粘性土・泥炭の成因について適切な記述を選ぶ。

\Needspace{4\baselineskip}\paragraph{確認した条件・記述（独自要約）}\mbox{}\par

\noindent\textbf{方向：} 最も適切なもの。\par

\noindent\textbf{対象：} 土の一般的な名称と成因。\par

\paragraph{問われる核心} 特殊土の母材・成因・残存植物組織

\paragraph{必要知識} 特殊土の母材・成因・残存植物組織。母材・運搬堆積・有機物起源、風化、弱面、水の作用、岩盤と土の物性。

\paragraph{考え方・正答を選ぶ理由（査読待ち初稿）} 色や粒状の見た目だけでなく母材と形成過程を対応させる。まさ土は花崗岩風化、しらすは火山噴出物、黄土は風成シルト、火山灰質粘性土は火山灰の風化等、泥炭は植物遺体の不完全分解である。適切なのは泥炭の5である。

\Needspace{4\baselineskip}\paragraph{選択肢の検討（査読待ち初稿）}\mbox{}\par

\noindent\textbf{1：不適切} まさ土の母材は流紋岩である。。まさ土は主に花崗岩類の風化により砂状となった土である。残留と崩積はその後の移動の区別であり、流紋岩を母材とする記述が誤り。\par

\noindent\textbf{2：不適切} しらすは花崗岩の風化物とその二次堆積物である。。しらすは主に珪長質火砕流の未溶結部などに由来する火山性堆積物で、二次堆積物もある。花崗岩の風化土とは起源が違う。\par

\noindent\textbf{3：不適切} 黄土は海洋堆積物である。。黄土（レス）は風で運ばれたシルトを主体とする風成堆積物で、海成という説明が誤り。\par

\noindent\textbf{4：不適切} 火山灰質粘性土をレスと呼ぶ。。火山灰由来の粘性土と風成シルトのレスを同一にしない。分布地域が一部重なるだけでは成因を一致させられない。\par

\noindent\textbf{5：適切} 泥炭はピートとも呼ばれ、未分解の植物組織を肉眼で判定できる。。水分の多い環境などで植物遺体の分解が遅れ、繊維・組織を残す有機質土が泥炭である。分解程度には幅があり、あらゆる泥炭の全組織が必ず肉眼で見えるという一般化は避ける。\par

\paragraph{再利用する考え方} 特殊土の母材・成因・残存植物組織を成因と観察される性質に結び付け、例示と全称命題、主要鉱物と副成分、公式解答と未解決の解釈を区別する。

\paragraph{共通理論} \hyperref[theory:engineering:soil]{第\ref*{theory:engineering:soil}節}

\par 参考：\cite{mlit_masa,maff_shiras2017,mlit_peat,gsj_loess2002}。

\paragraph{関連問題} \hyperref[q:h27-reg-applied-29]{H27-REG-APPLIED-29}

\par 難易度 3（校準初稿）、優先度未評価。

\subsection{平成27年度 問29}
\label{q:h27-reg-applied-29}

\noindent\textbf{原典：} sources/H27/17\_applied\_science.pdf、PDF物理11頁。

\textbf{公式正答：} 5（single）。

\paragraph{問題の要旨} まさ土・しらす・黄土・火山灰質粘性土・泥炭の成因について適切な記述を選ぶ。

\Needspace{4\baselineskip}\paragraph{確認した条件・記述（独自要約）}\mbox{}\par

\noindent\textbf{方向：} 最も適切なもの。\par

\noindent\textbf{対象：} 土の一般的な名称と成因。\par

\paragraph{問われる核心} 特殊土の母材・成因・残存植物組織

\paragraph{必要知識} 特殊土の母材・成因・残存植物組織。母材・運搬堆積・有機物起源、風化、弱面、水の作用、岩盤と土の物性。

\paragraph{考え方・正答を選ぶ理由（査読待ち初稿）} 色や粒状の見た目だけでなく母材と形成過程を対応させる。まさ土は花崗岩風化、しらすは火山噴出物、黄土は風成シルト、火山灰質粘性土は火山灰の風化等、泥炭は植物遺体の不完全分解である。適切なのは泥炭の5である。

\Needspace{4\baselineskip}\paragraph{選択肢の検討（査読待ち初稿）}\mbox{}\par

\noindent\textbf{1：不適切} まさ土の母材は流紋岩である。。まさ土は主に花崗岩類の風化により砂状となった土である。残留と崩積はその後の移動の区別であり、流紋岩を母材とする記述が誤り。\par

\noindent\textbf{2：不適切} しらすは花崗岩の風化物とその二次堆積物である。。しらすは主に珪長質火砕流の未溶結部などに由来する火山性堆積物で、二次堆積物もある。花崗岩の風化土とは起源が違う。\par

\noindent\textbf{3：不適切} 黄土は海洋堆積物である。。黄土（レス）は風で運ばれたシルトを主体とする風成堆積物で、海成という説明が誤り。\par

\noindent\textbf{4：不適切} 火山灰質粘性土をレスと呼ぶ。。火山灰由来の粘性土と風成シルトのレスを同一にしない。分布地域が一部重なるだけでは成因を一致させられない。\par

\noindent\textbf{5：適切} 泥炭はピートとも呼ばれ、未分解の植物組織を肉眼で判定できる。。水分の多い環境などで植物遺体の分解が遅れ、繊維・組織を残す有機質土が泥炭である。分解程度には幅があり、あらゆる泥炭の全組織が必ず肉眼で見えるという一般化は避ける。\par

\paragraph{再利用する考え方} 特殊土の母材・成因・残存植物組織を成因と観察される性質に結び付け、例示と全称命題、主要鉱物と副成分、公式解答と未解決の解釈を区別する。

\paragraph{共通理論} \hyperref[theory:engineering:soil]{第\ref*{theory:engineering:soil}節}

\par 参考：\cite{mlit_masa,maff_shiras2017,mlit_peat,gsj_loess2002}。

\paragraph{関連問題} \hyperref[q:h25-reg-applied-19]{H25-REG-APPLIED-19}, \hyperref[q:r01-ret-applied-28]{R01-RET-APPLIED-28}

\par 難易度 3（校準初稿）、優先度未評価。

\subsection{令和元年度再試験 問28}
\label{q:r01-ret-applied-28}

\noindent\textbf{原典：} sources/R01\_retest/17\_applied\_science.pdf、PDF物理9頁。

\textbf{公式正答：} 3（single）。

\paragraph{問題の要旨} まさ土・しらす・黄土・火山灰質粘性土・泥炭の成因について適切な記述を選ぶ。

\Needspace{4\baselineskip}\paragraph{確認した条件・記述（独自要約）}\mbox{}\par

\noindent\textbf{方向：} 最も適切なもの。\par

\noindent\textbf{対象：} 土の一般的な名称と成因。\par

\paragraph{問われる核心} 特殊土の母材・成因・残存植物組織

\paragraph{必要知識} 特殊土の母材・成因・残存植物組織。母材・運搬堆積・有機物起源、風化、弱面、水の作用、岩盤と土の物性。

\paragraph{考え方・正答を選ぶ理由（査読待ち初稿）} 色や粒状の見た目だけでなく母材と形成過程を対応させる。まさ土は花崗岩風化、しらすは火山噴出物、黄土は風成シルト、火山灰質粘性土は火山灰の風化等、泥炭は植物遺体の不完全分解である。適切なのは泥炭の3である。

\Needspace{4\baselineskip}\paragraph{選択肢の検討（査読待ち初稿）}\mbox{}\par

\noindent\textbf{1：不適切} 黄土は海洋堆積物である。。黄土（レス）は風で運ばれたシルトを主体とする風成堆積物で、海成という説明が誤り。\par

\noindent\textbf{2：不適切} しらすは花崗岩の風化物とその二次堆積物である。。しらすは主に珪長質火砕流の未溶結部などに由来する火山性堆積物で、二次堆積物もある。花崗岩の風化土とは起源が違う。\par

\noindent\textbf{3：適切} 泥炭はピートとも呼ばれ、未分解の植物組織を肉眼で判定できる。。水分の多い環境などで植物遺体の分解が遅れ、繊維・組織を残す有機質土が泥炭である。分解程度には幅があり、あらゆる泥炭の全組織が必ず肉眼で見えるという一般化は避ける。\par

\noindent\textbf{4：不適切} まさ土の母材は流紋岩である。。まさ土は主に花崗岩類の風化により砂状となった土である。残留と崩積はその後の移動の区別であり、流紋岩を母材とする記述が誤り。\par

\noindent\textbf{5：不適切} 火山灰質粘性土をレスと呼ぶ。。火山灰由来の粘性土と風成シルトのレスを同一にしない。分布地域が一部重なるだけでは成因を一致させられない。\par

\paragraph{再利用する考え方} 特殊土の母材・成因・残存植物組織を成因と観察される性質に結び付け、例示と全称命題、主要鉱物と副成分、公式解答と未解決の解釈を区別する。

\paragraph{共通理論} \hyperref[theory:engineering:soil]{第\ref*{theory:engineering:soil}節}

\par 参考：\cite{mlit_masa,maff_shiras2017,mlit_peat,gsj_loess2002}。

\paragraph{関連問題} \hyperref[q:h27-reg-applied-29]{H27-REG-APPLIED-29}

\par 難易度 3（校準初稿）、優先度未評価。


\chapter{資源・鉱床}
\section{普通の岩石から濃集を説明する}\label{theory:ore:origins}
鉱床は特定の元素・鉱物が周囲より濃集した地質体として考える。濃集の存在、品位、鉱量、採掘可能量、採算は別の評価段階である。品位が低くても大規模な鉱体は資源となり得る。
まず元素の供給源を考え、運搬媒体、通路、濃集条件、保存を順に追う。正マグマ性は結晶分化・集積や硫化物液の分離、熱水性は流体中の成分の沈殿、堆積性は化学沈殿や粒子運搬、風化残留は移動しにくい成分の残留を中心に整理する。石油は根源岩での生成・移動・貯留とトラップを分ける。
層状苦鉄質貫入体のクロムと白金族元素は代表的なマグマ性資源である\cite{ore50_magmatic_ref}。金銀の熱水成因が重要であることから、全マグマ性鉱床に金が絶対にないとは推論しない。自然金を含むPGE鉱床モデルもある\cite{ore50_gold_exception_ref}。
水中の化学種・pH・酸化還元状態で移動性は変わる。元素名の固定した溶けやすさ順位だけで全成因を説明しない。

\section{熱水の通路と沈殿条件}\label{theory:ore:hydrothermal}
高温の流体は金属等を運び、温度低下、圧力変化、流体混合、酸化還元やpHの変化で鉱物を沈殿させ得る。岩石の亀裂・断層は通路であり、通路を充填すれば鉱脈となる。流体の全てがマグマだけに由来するわけではない。
斑岩銅は浅所の貫入岩とその周辺に鉱染・細脈網を伴う大規模低品位の熱水系として考える。収束縁の大陸弧と海洋性島弧の両方のモデルがある\cite{ore50_porphyry_ref}。大陸弧に多いという傾向を島弧に絶対存在しないという全称に変えない。
炭酸塩岩との化学交代でCaケイ酸塩鉱物を作るスカルンは、熱だけによる再結晶と違い、成分の移動を伴う。高温から後期低温へ鉱物組合せが変化する例がある\cite{ore50_skarn_ref}。温度だけで金属種を全て決めず、母岩・流体組成・酸化還元条件も考える。

\section{黒鉱・グリーンタフ・別子型}\label{theory:ore:kuroko}
黒鉱は代表的に珪長質海底火山活動と関係する塊状硫化物鉱床で、閃亜鉛鉱、方鉛鉱、黄鉄鉱、黄銅鉱、四面銅鉱、重晶石等を含む\cite{ore50_kuroko_ref}。海底熱水からの化学的沈殿と、その下の流体通路を考える。
日本の代表的な産地は東北の中新世グリーンタフ地域にある。グリーンタフは変質した火山岩・火山砕屑岩を含む地域的呼称であり、緑色片岩という広域変成岩の岩種名とは違う。名称の「緑」だけから同じ母岩としない。
層状含銅硫化鉄鉱の別子型は、海底熱水性の硫化物が変形・変成を経たモデルで、三波川帯の例にも変成相の記録がある\cite{ore50_besshi_ref}。黒鉱の代表産状と、変成した苦鉄質・堆積岩系に産する硫化物を区別する。
年代区分の西黒沢期に「全て限られる」という原典の全称は当時の分類と例外を確認する。代表的な集中地域の証拠を、全世界の全型の排他的条件へ変換しない。

\section{プレート境界と鉱床の対応}\label{theory:ore:tectonic}
中央海嶺は発散、海溝での沈み込みは収束である\cite{usgs_plate_boundaries}。島弧・背弧でも海底火山と熱水系が発達し、黒鉱型の潜在鉱床が報告される\cite{ore50_backarc_ref}。斑岩銅も海洋性島弧・大陸弧等に関係するため、一つの境界呼称から一つの鉱床だけを決める対応表は不十分である。
平成24年度問25の「マリアナ型」に対する独立解答と公式番号の相違は問題別に保存した。母岩、伸張・圧縮、マグマ、熱水の観測を比べる必要があり、公式4を得るために黒鉱型の島弧・背弧での存在を否定しない。砂岩型ウランは堆積盆地内の透水性と酸化還元境界を考え、収束境界だけの資源としない。

\section{堆積資源と風化残留}\label{theory:ore:sedimentary}
縞状鉄鉱は珪質層と鉄に富む層の互層であり、代表モデルでもチャートと鉄鉱物層の組合せを扱う\cite{ore50_bif_ref}。鉄の酸化状態と海水での沈殿、堆積後の変成・変形を区別する。原典で銅へ置換した成因をそのまま覚えない。
漂砂では鉱物粒子を物理的に運搬・分級する。砂金の濃集と、溶存元素が還元帯で化学沈殿する砂岩型ウランを区別する。
ボーキサイトでは強い風化で他成分が溶脱し、Alを含む難溶物が相対的に残る\cite{ore50_bauxite_ref}。Al・FeとNa・K・Ca・Mgの残留・溶脱を逆転させた候補が出題される。濃度上昇は対象元素の新規供給だけでなく、分母となる他成分の減少でも起こる。
オイルシェールはケロジェンを含む細粒堆積岩である\cite{ore50_oilshale_ref}。粗粒・細粒の交換に注意する。メタンハイドレートは水の構造へガスが取り込まれる水和物で、温度圧力による安定条件をもつ。見た目のシャーベット状という説明だけで化学構造を定義しない。

\section{石炭化を年代だけで決めない}\label{theory:ore:coal}
植物由来の有機物が埋没・加熱を経て炭素に富み、水分・揮発分等が減る石炭化を考える。基本のランク順は褐炭、亜瀝青炭、瀝青炭、無煙炭である\cite{ore50_coal_ref}。石墨に近い状態までを扱う場合も、低ランクの褐炭を瀝青炭の後に置かない。
石炭化に伴うH/C・O/Cの変化は一般に減少方向となる。原子比、質量割合、揮発分、含水量は別の指標であり、分母を確認する。
日本には古第三紀の炭田が多く、主にジュラ系とする説明は適切でない\cite{ore50_japan_coal_ref}。石炭紀の代表的な湿地植生を針葉樹だけへ交換しない。
形成温度の概説に150--200℃の代表範囲があるが、ランクは最高到達温度と埋没熱史に依存する\cite{ore50_coal_temp_ref}。古生代という年代だけで全てが200℃以下と保証できない。年代と熟成度の二つを独立に扱う。

\section{地熱の熱・水・透水性}\label{theory:ore:geothermal}
地熱は地殻の放射性壊変、深部からの熱、火山近くの局所熱等を含む。熱いだけで流体を取り出せるとは限らず、流体と有効な透水経路を持つ貯留層が必要となる。亀裂・断層や封じるキャップロックが重要である\cite{ore50_reservoir_ref}。
発電では蒸気をタービンへ使う方式と、熱交換で低沸点二次媒体を蒸発させるバイナリー方式を分ける。中・低温の温水は暖房等へ直接利用でき、非火山地域にも温水系がある\cite{ore50_geothermal_ref}。
概略の取出し熱量率は \(\dot Q=\dot m c_p(T_{\rm in}-T_{\rm out})\)。流量・温度差・持続する涵養が関係し、熱の存在量と実際に取り出せる出力を区別する。バイナリーでも熱力学的効率は温度差に制約される。
地化学温度計は溶解平衡等の温度依存を利用する。平衡到達、上昇時の冷却、異なる水との混合、再平衡の有無を確認する。地温勾配の代表20--30、局所50--100℃/km以上等を全地点の固定値にしない。

\section{深海団塊の極めて遅い成長}\label{theory:ore:marine}
マンガン団塊は鉄・マンガン酸化物を主体に核の周囲へ成長し、Cu・Ni・Co等を伴う。海底の堆積が遅い場所に露出する団塊と、海山の岩盤を覆うクラストを区別する。
JOGMECの代表説明では成長速度は百万年で1mm程度である\cite{ore50_nodule_ref}。数mm／1万年への単位期間の置換は二桁以上の差になる。起源・続成作用・環境による幅を認めて代表率として読む。原典の副成分最大1--2\%という表現は全団塊の絶対上限として未確定であり、問題別に留保する。
成長の遅さ、資源量、採掘試験、商業生産、採算、環境影響の確認は別の事項である。海に鉱物があるという観測だけで商業生産中と結論しない。

\section{ウランの母岩・流体・年代・品位}\label{theory:ore:uranium}
代表的な分類に先カンブリア時代の石英礫岩、原生代の不整合関連、顕生代の砂岩型がある\cite{ore50_uranium_ref}。古い礫岩では酸素の少ない条件下で鉱物粒子の保存と漂砂濃集を考える。砂岩型では流体が溶存Uを運び、還元帯等で沈殿させる。どちらも堆積岩中だから同じ機構とはしない。
不整合関連では古い基盤と盆地堆積物の境界付近の流体鉱化、鉱脈型では地質体を切る亀裂充填、貫入岩型では岩体に伴う分散した鉱化等を考える。
IAEAの資源説明にも低・非常に低品位で大規模な貫入岩関連等の資源がある\cite{ore50_intrusive_ref}。貫入岩型の全てを高品位としない。品位が高いことと総量が大きいことは別である。
原典の4億年・23億年等の年代区切りは分類史・対象型の範囲と照合が残る。厳密な世界全例の境界として教えず、代表的地質時代と酸化環境を結び付けて理解する。

\section{学習順序}
まず濃集の機構を五種類に分け、次に母岩・鉱物・形態を対応させる。その後プレート環境と時代、温度・酸化還元条件を加える。最後に品位・鉱量・生産状況という資源評価を読む。似た名前や緑色という共通点だけで鉱床を統合せず、原典の存在命題と全称命題を区別する。

\section{原典と全五肢の解説初稿}
\section{石炭化の順序・原子比・時代と形成温度を比較する。}

\subsection{平成24年度 問26}
\label{q:h24-reg-applied-26}

\noindent\textbf{原典：} sources/H24/17\_applied\_science.pdf、PDF物理8頁。

\textbf{公式正答：} 4（single）。

\par\textbf{検討未確定：} 公式4を保持。200℃以下は典型的石炭化の代表説明であり、原典の古生代という年代だけの条件から全炭田の熱史を確定できない。無煙炭等の高温例も一次資料で確認。

\paragraph{問題の要旨} 石炭化の順序・原子比・時代と形成温度を比較する。

\Needspace{4\baselineskip}\paragraph{確認した条件・記述（独自要約）}\mbox{}\par

\noindent\textbf{原典確認：} S7\_050の29原典画像で34問の全五肢・否定・数値を照合。\par

\noindent\textbf{対象：} 石炭化の順序・原子比・時代と形成温度を比較する。\par

\paragraph{問われる核心} 石炭化の順序・原子比・時代と形成温度を比較する。

\paragraph{必要知識} 年代・ランク・植物・温度を別の軸で読む。日本の炭田年代を主にジュラ紀とする1、褐炭の位置を交換した2、H/C・O/Cを増加とする3、石炭紀の代表植生を針葉樹とする5は不適切。提示された相対比較では4を選ぶが、年代だけから石炭の最高到達温度を200℃以下と保証しない。

\paragraph{考え方・正答を選ぶ理由（査読待ち初稿）} 年代・ランク・植物・温度を別の軸で読む。日本の炭田年代を主にジュラ紀とする1、褐炭の位置を交換した2、H/C・O/Cを増加とする3、石炭紀の代表植生を針葉樹とする5は不適切。提示された相対比較では4を選ぶが、年代だけから石炭の最高到達温度を200℃以下と保証しない。

\Needspace{4\baselineskip}\paragraph{選択肢の検討（査読待ち初稿）}\mbox{}\par

\noindent\textbf{1：不適切} 日本の主な石炭はジュラ系。日本の代表的炭田には古第三系等が多い。少数の古い産出例を日本全体の主流としない。\par

\noindent\textbf{2：不適切} 亜瀝青・瀝青・褐炭の順で石炭化。基本の低ランクからの順は褐炭→亜瀝青炭→瀝青炭→無煙炭で、褐炭の位置が逆。\par

\noindent\textbf{3：不適切} 石炭化でH/C・O/Cが増加。脱水・脱炭酸・炭素濃集等で一般にはこれら比は減少する方向へ進む。\par

\noindent\textbf{4：条件・解釈を留保} 古生代石炭は一般に200℃以下。石炭化は埋没熱史による。USGSの概説には150–200℃を超えない典型の説明があるが、高ランク・熱変成では200–300℃の例もある。年代だけで上限は定めない。\par

\noindent\textbf{5：不適切} 石炭紀の主起源を針葉樹。石炭紀の代表的な湿地森林には大型のシダ植物等が多く、針葉樹だけを代表起源としない。\par

\paragraph{再利用する考え方} 地質環境・生成機構・物質の名称を対応させ、平均値と全称、典型例と例外、当時の分類と現在の知見を分ける。

\paragraph{共通理論} \hyperref[theory:ore:coal]{第\ref*{theory:ore:coal}節}

\par 参考：\cite{ore50_coal_ref,ore50_coal_temp_ref,ore50_japan_coal_ref}。

\par 難易度 4（校準初稿）、優先度未評価。

\section{地熱資源の利用域・貯留層・発電方式を比較する。}

\subsection{平成23年度 問26}
\label{q:h23-reg-applied-26}

\noindent\textbf{原典：} sources/H23/17\_applied\_science.pdf、PDF物理8頁。

\textbf{公式正答：} 2（single）。

\par\textbf{検討未確定：} 地温勾配と150℃区分は原典の代表値。全地域を固定した境界で分けず、具体地点の勾配・貯留温度と地化学平衡条件は追加照合対象。

\paragraph{問題の要旨} 地熱資源の利用域・貯留層・発電方式を比較する。

\Needspace{4\baselineskip}\paragraph{確認した条件・記述（独自要約）}\mbox{}\par

\noindent\textbf{原典確認：} S7\_050の29原典画像で34問の全五肢・否定・数値を照合。\par

\noindent\textbf{対象：} 地熱資源の利用域・貯留層・発電方式を比較する。\par

\paragraph{問われる核心} 地熱資源の利用域・貯留層・発電方式を比較する。

\paragraph{必要知識} 地熱利用には高温蒸気による発電だけでなく、中低温熱水の直接利用やバイナリー発電がある。非火山地域の温水利用もあるため火山域への限定2を選ぶ。流体貯留と循環には割れ目等の透水性が重要。地化学温度計は平衡・混合等の仮定を確認して地下温度を推定する。

\paragraph{考え方・正答を選ぶ理由（査読待ち初稿）} 地熱利用には高温蒸気による発電だけでなく、中低温熱水の直接利用やバイナリー発電がある。非火山地域の温水利用もあるため火山域への限定2を選ぶ。流体貯留と循環には割れ目等の透水性が重要。地化学温度計は平衡・混合等の仮定を確認して地下温度を推定する。

\Needspace{4\baselineskip}\paragraph{選択肢の検討（査読待ち初稿）}\mbox{}\par

\noindent\textbf{1：代表的条件で適切} 蒸気方式とバイナリー。蒸気でタービンを回す方式と、熱交換で低沸点の二次媒体を蒸発させる方式を区別する。\par

\noindent\textbf{2：不適切} 中・低温熱水の利用範囲。温室・暖房・温泉等への直接利用は火山周辺だけに限定されない。\par

\noindent\textbf{3：条件・解釈を留保} 広域と局所の地温勾配。放射性壊変・深部熱供給・火山等の熱源を考える。20–30、50–100℃/km以上は代表範囲で全地点の一定値ではない。\par

\noindent\textbf{4：代表的条件で適切} 割れ目に富む貯留層。流体が集まり循環するために有効な透水経路が必要で、亀裂や断層が重要となる。\par

\noindent\textbf{5：代表的条件で適切} 化学成分による温度推定。溶解平衡等の温度依存を利用する地化学温度計。冷却、混合、再平衡の影響は別に検討する。\par

\paragraph{再利用する考え方} 地質環境・生成機構・物質の名称を対応させ、平均値と全称、典型例と例外、当時の分類と現在の知見を分ける。

\paragraph{共通理論} \hyperref[theory:ore:geothermal]{第\ref*{theory:ore:geothermal}節}

\par 参考：\cite{ore50_geothermal_ref,ore50_reservoir_ref}。

\paragraph{関連問題} \hyperref[q:r03-reg-applied-20]{R03-REG-APPLIED-20}

\par 難易度 4（校準初稿）、優先度未評価。

\subsection{令和3年度 問20}
\label{q:r03-reg-applied-20}

\noindent\textbf{原典：} sources/R03/17\_applied\_science.pdf、PDF物理7頁。

\textbf{公式正答：} 4（single）。

\par\textbf{検討未確定：} 地温勾配の代表範囲と地化学温度計の平衡仮定は地域・化学種の条件を要する。原典の固定範囲を地球全域の厳密定数としない。

\paragraph{問題の要旨} 地熱貯留層の透水経路と温度利用を比較する。

\Needspace{4\baselineskip}\paragraph{確認した条件・記述（独自要約）}\mbox{}\par

\noindent\textbf{原典確認：} S7\_050の29原典画像で34問の全五肢・否定・数値を照合。\par

\noindent\textbf{対象：} 地熱貯留層の透水経路と温度利用を比較する。\par

\paragraph{問われる核心} 地熱貯留層の透水経路と温度利用を比較する。

\paragraph{必要知識} 発電・直接利用の方式と、地下の熱・流体・透水性を分ける。通常の熱水貯留層は流体が移動する割れ目等を持ち、乏しいとした4を選ぶ。高温だが透水性不足の岩体の利用には別の技術と条件が必要で、通常の貯留層の定義と混同しない。

\paragraph{考え方・正答を選ぶ理由（査読待ち初稿）} 発電・直接利用の方式と、地下の熱・流体・透水性を分ける。通常の熱水貯留層は流体が移動する割れ目等を持ち、乏しいとした4を選ぶ。高温だが透水性不足の岩体の利用には別の技術と条件が必要で、通常の貯留層の定義と混同しない。

\Needspace{4\baselineskip}\paragraph{選択肢の検討（査読待ち初稿）}\mbox{}\par

\noindent\textbf{1：条件・解釈を留保} 広域と局所の地温勾配。放射性壊変・深部熱供給・火山等の熱源を考える。20–30、50–100℃/km以上は代表範囲で全地点の一定値ではない。\par

\noindent\textbf{2：代表的条件で適切} 非火山地域も低温利用の対象。地殻の広域熱等による熱水も利用でき、火山周辺だけに限定されない。\par

\noindent\textbf{3：代表的条件で適切} 化学成分による温度推定。溶解平衡等の温度依存を利用する地化学温度計。冷却、混合、再平衡の影響は別に検討する。\par

\noindent\textbf{4：不適切} 割れ目に乏しい地熱貯留層。通常の開発可能な熱水循環には割れ目等の透水経路が重要。原典H23の富んだを乏しいへ反転したもの。\par

\noindent\textbf{5：代表的条件で適切} 蒸気方式とバイナリー。蒸気でタービンを回す方式と、熱交換で低沸点の二次媒体を蒸発させる方式を区別する。\par

\paragraph{再利用する考え方} 地質環境・生成機構・物質の名称を対応させ、平均値と全称、典型例と例外、当時の分類と現在の知見を分ける。

\paragraph{共通理論} \hyperref[theory:ore:geothermal]{第\ref*{theory:ore:geothermal}節}

\par 参考：\cite{ore50_geothermal_ref,ore50_reservoir_ref}。

\paragraph{関連問題} \hyperref[q:h23-reg-applied-26]{H23-REG-APPLIED-26}

\par 難易度 4（校準初稿）、優先度未評価。

\section{黒鉱の鉱物・海底熱水成因と母岩を比較する。}

\subsection{平成26年度 問29}
\label{q:h26-reg-applied-29}

\noindent\textbf{原典：} sources/H26/17\_applied\_science.pdf、PDF物理12頁。

\textbf{公式正答：} 4（single）。

\par\textbf{検討未確定：} 原典の生成年代を西黒沢期に限る全称と全鉱物組合せの例外は追加照合。代表地域の記録を全黒鉱の唯一の産状へ一般化しない。

\paragraph{問題の要旨} 黒鉱の鉱物・海底熱水成因と母岩を比較する。

\Needspace{4\baselineskip}\paragraph{確認した条件・記述（独自要約）}\mbox{}\par

\noindent\textbf{原典確認：} S7\_050の29原典画像で34問の全五肢・否定・数値を照合。\par

\noindent\textbf{対象：} 黒鉱の鉱物・海底熱水成因と母岩を比較する。\par

\paragraph{問われる核心} 黒鉱の鉱物・海底熱水成因と母岩を比較する。

\paragraph{必要知識} 黒鉱は金属を含む海底熱水と硫化物の沈殿に関係し、単なる熱による再結晶という4が不適切。グリーンタフは変質した火山岩・火山砕屑岩の地域的呼称で、緑色片岩という変成岩と取り違えない。

\paragraph{考え方・正答を選ぶ理由（査読待ち初稿）} 黒鉱は金属を含む海底熱水と硫化物の沈殿に関係し、単なる熱による再結晶という4が不適切。グリーンタフは変質した火山岩・火山砕屑岩の地域的呼称で、緑色片岩という変成岩と取り違えない。

\Needspace{4\baselineskip}\paragraph{選択肢の検討（査読待ち初稿）}\mbox{}\par

\noindent\textbf{1：代表的条件で適切} ZnS・PbS・重晶石等の黒色塊状鉱。黒鉱の代表組合せは閃亜鉛鉱・方鉛鉱・黄鉄鉱・重晶石等。四面銅鉱等の副成分もある。\par

\noindent\textbf{2：代表的条件で適切} 珪長質海底火山に関係。火山性塊状硫化物鉱床の一型で、海底熱水活動による鉱化を考える。\par

\noindent\textbf{3：条件・解釈を留保} 日本のグリーンタフ・西黒沢期。東北日本の中新世グリーンタフ地域と秋田県北部は代表産地。限られるという厳密な年代全称は当時の分類と例外の照合が残る。\par

\noindent\textbf{4：不適切} 単なる火成岩周辺の再結晶。これは接触変成の要点で、金属を運ぶ熱水の沈殿による黒鉱成因を説明しない。\par

\noindent\textbf{5：代表的条件で適切} キプロス型よりPb・Znに富む。珪長質に関係する黒鉱型と苦鉄質海洋地殻に関係するキプロス型の代表的組成差。\par

\paragraph{再利用する考え方} 地質環境・生成機構・物質の名称を対応させ、平均値と全称、典型例と例外、当時の分類と現在の知見を分ける。

\paragraph{共通理論} \hyperref[theory:ore:kuroko]{第\ref*{theory:ore:kuroko}節}

\par 参考：\cite{ore50_kuroko_ref,ore50_backarc_ref}。

\par 難易度 4（校準初稿）、優先度未評価。

\subsection{令和5年度 問23}
\label{q:r05-reg-applied-23}

\noindent\textbf{原典：} sources/R05/17\_applied\_science.pdf、PDF物理8頁。

\textbf{公式正答：} 3（single）。

\par\textbf{検討未確定：} スカルン金属種を高温と中低温だけで固定する原典4の一般化は追加照合。全形態分類の完全排他性も主張しない。

\paragraph{問題の要旨} 黒鉱のグリーンタフと緑色片岩を区別する。

\Needspace{4\baselineskip}\paragraph{確認した条件・記述（独自要約）}\mbox{}\par

\noindent\textbf{原典確認：} S7\_050の29原典画像で34問の全五肢・否定・数値を照合。\par

\noindent\textbf{対象：} 黒鉱のグリーンタフと緑色片岩を区別する。\par

\paragraph{問われる核心} 黒鉱のグリーンタフと緑色片岩を区別する。

\paragraph{必要知識} 黒鉱の母岩をグリーンタフと緑色片岩で取り違えた3を選ぶ。熱水成因が共通でも、斑岩銅の分散した細脈、鉱脈の通路、スカルンの化学交代、海底の黒鉱は位置と形態が違う。温度と金属種の対応は酸化還元・母岩・流体組成の条件も確認する。

\paragraph{考え方・正答を選ぶ理由（査読待ち初稿）} 黒鉱の母岩をグリーンタフと緑色片岩で取り違えた3を選ぶ。熱水成因が共通でも、斑岩銅の分散した細脈、鉱脈の通路、スカルンの化学交代、海底の黒鉱は位置と形態が違う。温度と金属種の対応は酸化還元・母岩・流体組成の条件も確認する。

\Needspace{4\baselineskip}\paragraph{選択肢の検討（査読待ち初稿）}\mbox{}\par

\noindent\textbf{1：代表的条件で適切} 熱水鉱床の形態分類。鉱脈、塊状、スカルン、斑岩銅等の代表タイプを区別する。全分類が完全に排他的とはしない。\par

\noindent\textbf{2：代表的条件で適切} 半深成岩と斑岩銅。浅所の斑岩質貫入岩とその熱水系に関係する。\par

\noindent\textbf{3：不適切} 黒鉱が東北の緑色片岩中。黒鉱の代表は中新世グリーンタフの海底火山・火山砕屑岩地域。広域変成の緑色片岩と名称を交換している。\par

\noindent\textbf{4：条件・解釈を留保} 温度によるスカルン金属の傾向。高温段階にW等、後期・遠方の低温段階にPb・Zn等の例がある。金属種を温度だけで一意決定する普遍則ではない。\par

\noindent\textbf{5：代表的条件で適切} 裂け目を通る熱水の鉱脈。流体が運んだ成分が亀裂・断層沿いで沈殿し脈状となる。\par

\paragraph{再利用する考え方} 地質環境・生成機構・物質の名称を対応させ、平均値と全称、典型例と例外、当時の分類と現在の知見を分ける。

\paragraph{共通理論} \hyperref[theory:ore:kuroko]{第\ref*{theory:ore:kuroko}節}

\par 参考：\cite{ore50_kuroko_ref,ore50_skarn_ref,ore50_porphyry_ref}。

\paragraph{関連問題} \hyperref[q:r07-reg-applied-25]{R07-REG-APPLIED-25}

\par 難易度 4（校準初稿）、優先度未評価。

\subsection{令和6年度 問24}
\label{q:r06-reg-applied-24}

\noindent\textbf{原典：} sources/R06/17\_applied\_science.pdf、PDF物理8頁。

\textbf{公式正答：} 3（single）。

\paragraph{問題の要旨} 黒鉱と炭酸塩岩を交代するスカルンを区別する。

\Needspace{4\baselineskip}\paragraph{確認した条件・記述（独自要約）}\mbox{}\par

\noindent\textbf{原典確認：} S7\_050の29原典画像で34問の全五肢・否定・数値を照合。\par

\noindent\textbf{対象：} 黒鉱と炭酸塩岩を交代するスカルンを区別する。\par

\paragraph{問われる核心} 黒鉱と炭酸塩岩を交代するスカルンを区別する。

\paragraph{必要知識} 3は黒鉱をスカルンの生成機構へ置換している。熱水が関与する点は共通でも、海底での硫化物沈殿と、炭酸塩岩を反応させCaケイ酸塩を作る交代は違う。鉱床名・母岩・鉱物・形態・成因を同時に対応させる。

\paragraph{考え方・正答を選ぶ理由（査読待ち初稿）} 3は黒鉱をスカルンの生成機構へ置換している。熱水が関与する点は共通でも、海底での硫化物沈殿と、炭酸塩岩を反応させCaケイ酸塩を作る交代は違う。鉱床名・母岩・鉱物・形態・成因を同時に対応させる。

\Needspace{4\baselineskip}\paragraph{選択肢の検討（査読待ち初稿）}\mbox{}\par

\noindent\textbf{1：代表的条件で適切} 縞状鉄鉱の珪質層と鉄鉱物層。珪質物と鉄鉱物に富む層が交互に現れる。銅の鉱床と交換しない。\par

\noindent\textbf{2：代表的条件で適切} 難溶鉱物の風化残留。風化に伴う移動しにくい成分の相対濃集をいう。\par

\noindent\textbf{3：不適切} 黒鉱が炭酸塩岩の交代で形成。炭酸塩岩との化学反応はスカルンの説明。黒鉱は代表的に海底火山熱水の塊状硫化物である。\par

\noindent\textbf{4：代表的条件で適切} 低品位大規模の斑岩銅。斑岩質貫入体とその周辺の鉱染・網状細脈を持つ代表タイプ。\par

\noindent\textbf{5：代表的条件で適切} 海底熱水からの沈殿。海底火山・熱水活動で溶けた金属が化学条件の変化で沈殿する。\par

\paragraph{再利用する考え方} 地質環境・生成機構・物質の名称を対応させ、平均値と全称、典型例と例外、当時の分類と現在の知見を分ける。

\paragraph{共通理論} \hyperref[theory:ore:kuroko]{第\ref*{theory:ore:kuroko}節}

\par 参考：\cite{ore50_kuroko_ref,ore50_porphyry_ref,ore50_bauxite_ref}。

\paragraph{関連問題} \hyperref[q:r07-reg-applied-25]{R07-REG-APPLIED-25}

\par 難易度 3（校準初稿）、優先度未評価。

\section{天然資源の鉱物・粒径・生成条件を対応させる。}

\subsection{平成25年度 問24}
\label{q:h25-reg-applied-24}

\noindent\textbf{原典：} sources/H25/17\_applied\_science.pdf、PDF物理9頁。

\textbf{公式正答：} 3（single）。

\paragraph{問題の要旨} 天然資源の鉱物・粒径・生成条件を対応させる。

\Needspace{4\baselineskip}\paragraph{確認した条件・記述（独自要約）}\mbox{}\par

\noindent\textbf{原典確認：} S7\_050の29原典画像で34問の全五肢・否定・数値を照合。\par

\noindent\textbf{対象：} 天然資源の鉱物・粒径・生成条件を対応させる。\par

\paragraph{問われる核心} 天然資源の鉱物・粒径・生成条件を対応させる。

\paragraph{必要知識} オイルシェールはケロジェン等を含む細粒堆積岩であり3の粗粒が誤り。縞状鉄鉱、斑岩銅、ガス水和物、風化で濃集したAl鉱石は生成機構が違う。メタンハイドレートのシャーベット状という見た目の表現を、水和物の結晶構造や相安定の定義と同一視しない。

\paragraph{考え方・正答を選ぶ理由（査読待ち初稿）} オイルシェールはケロジェン等を含む細粒堆積岩であり3の粗粒が誤り。縞状鉄鉱、斑岩銅、ガス水和物、風化で濃集したAl鉱石は生成機構が違う。メタンハイドレートのシャーベット状という見た目の表現を、水和物の結晶構造や相安定の定義と同一視しない。

\Needspace{4\baselineskip}\paragraph{選択肢の検討（査読待ち初稿）}\mbox{}\par

\noindent\textbf{1：代表的条件で適切} 縞状鉄鉱の珪質層と鉄鉱物層。珪質物と鉄鉱物に富む層が交互に現れる。銅の鉱床と交換しない。\par

\noindent\textbf{2：代表的条件で適切} 浅所斑岩と斑岩銅。花崗閃緑斑岩等に伴う熱水性の鉱染・網状細脈を作る。\par

\noindent\textbf{3：不適切} 粗粒堆積岩がオイルシェール。ケロジェンを含む細粒堆積岩として扱う。粗粒という説明が不適切。\par

\noindent\textbf{4：代表的条件で適切} メタンハイドレートの温度圧力条件。水分子の構造にガス分子が取り込まれた水和物で、低温・高圧等の安定条件がある。単に凍結したメタンではない。\par

\noindent\textbf{5：代表的条件で適切} Al水酸化物がボーキサイト。Alの水酸化物・酸化水酸化物等を含む主要なAl鉱石。原岩が風化してAlが相対濃集する。\par

\paragraph{再利用する考え方} 地質環境・生成機構・物質の名称を対応させ、平均値と全称、典型例と例外、当時の分類と現在の知見を分ける。

\paragraph{共通理論} \hyperref[theory:ore:sedimentary]{第\ref*{theory:ore:sedimentary}節}

\par 参考：\cite{ore50_oilshale_ref,ore50_bauxite_ref,ore50_bif_ref}。

\paragraph{関連問題} \hyperref[q:h28-reg-applied-31]{H28-REG-APPLIED-31}

\par 難易度 3（校準初稿）、優先度未評価。

\subsection{平成28年度 問31}
\label{q:h28-reg-applied-31}

\noindent\textbf{原典：} sources/H28/17\_applied\_science.pdf、PDF物理12頁。

\textbf{公式正答：} 3（single）。

\paragraph{問題の要旨} 天然資源の鉱物・粒径・生成条件を対応させる。

\Needspace{4\baselineskip}\paragraph{確認した条件・記述（独自要約）}\mbox{}\par

\noindent\textbf{原典確認：} S7\_050の29原典画像で34問の全五肢・否定・数値を照合。\par

\noindent\textbf{対象：} 天然資源の鉱物・粒径・生成条件を対応させる。\par

\paragraph{問われる核心} 天然資源の鉱物・粒径・生成条件を対応させる。

\paragraph{必要知識} オイルシェールはケロジェン等を含む細粒堆積岩であり3の粗粒が誤り。縞状鉄鉱、斑岩銅、ガス水和物、風化で濃集したAl鉱石は生成機構が違う。メタンハイドレートのシャーベット状という見た目の表現を、水和物の結晶構造や相安定の定義と同一視しない。

\paragraph{考え方・正答を選ぶ理由（査読待ち初稿）} オイルシェールはケロジェン等を含む細粒堆積岩であり3の粗粒が誤り。縞状鉄鉱、斑岩銅、ガス水和物、風化で濃集したAl鉱石は生成機構が違う。メタンハイドレートのシャーベット状という見た目の表現を、水和物の結晶構造や相安定の定義と同一視しない。

\Needspace{4\baselineskip}\paragraph{選択肢の検討（査読待ち初稿）}\mbox{}\par

\noindent\textbf{1：代表的条件で適切} 縞状鉄鉱の珪質層と鉄鉱物層。珪質物と鉄鉱物に富む層が交互に現れる。銅の鉱床と交換しない。\par

\noindent\textbf{2：代表的条件で適切} 浅所斑岩と斑岩銅。花崗閃緑斑岩等に伴う熱水性の鉱染・網状細脈を作る。\par

\noindent\textbf{3：不適切} 粗粒堆積岩がオイルシェール。ケロジェンを含む細粒堆積岩として扱う。粗粒という説明が不適切。\par

\noindent\textbf{4：代表的条件で適切} メタンハイドレートの温度圧力条件。水分子の構造にガス分子が取り込まれた水和物で、低温・高圧等の安定条件がある。単に凍結したメタンではない。\par

\noindent\textbf{5：代表的条件で適切} Al水酸化物がボーキサイト。Alの水酸化物・酸化水酸化物等を含む主要なAl鉱石。原岩が風化してAlが相対濃集する。\par

\paragraph{再利用する考え方} 地質環境・生成機構・物質の名称を対応させ、平均値と全称、典型例と例外、当時の分類と現在の知見を分ける。

\paragraph{共通理論} \hyperref[theory:ore:sedimentary]{第\ref*{theory:ore:sedimentary}節}

\par 参考：\cite{ore50_oilshale_ref,ore50_bauxite_ref,ore50_bif_ref}。

\paragraph{関連問題} \hyperref[q:h25-reg-applied-24]{H25-REG-APPLIED-24}

\par 難易度 3（校準初稿）、優先度未評価。

\section{マンガン団塊の組成・海底分布・成長時間を比較する。}

\subsection{平成23年度 問27}
\label{q:h23-reg-applied-27}

\noindent\textbf{原典：} sources/H23/17\_applied\_science.pdf、PDF物理9頁。

\textbf{公式正答：} 5（single）。

\par\textbf{検討未確定：} 副成分の最大1–2\%という上限と粒径の全分類範囲は未確定。成長も海水起源・続成作用等で異なり、一つの代表率を全試料へ固定しない。

\paragraph{問題の要旨} マンガン団塊の組成・海底分布・成長時間を比較する。

\Needspace{4\baselineskip}\paragraph{確認した条件・記述（独自要約）}\mbox{}\par

\noindent\textbf{原典確認：} S7\_050の29原典画像で34問の全五肢・否定・数値を照合。\par

\noindent\textbf{対象：} マンガン団塊の組成・海底分布・成長時間を比較する。\par

\paragraph{問われる核心} マンガン団塊の組成・海底分布・成長時間を比較する。

\paragraph{必要知識} 団塊は核の周囲へ鉄・マンガン酸化物等がゆっくり堆積する。JOGMECの例は百万年で1mm程度で、5の数mm／1万年は時間尺度が大きく違う。海底に露出した団塊と岩盤上のクラストを区別し、成長率は起源・環境で変わる代表値として読む。

\paragraph{考え方・正答を選ぶ理由（査読待ち初稿）} 団塊は核の周囲へ鉄・マンガン酸化物等がゆっくり堆積する。JOGMECの例は百万年で1mm程度で、5の数mm／1万年は時間尺度が大きく違う。海底に露出した団塊と岩盤上のクラストを区別し、成長率は起源・環境で変わる代表値として読む。

\Needspace{4\baselineskip}\paragraph{選択肢の検討（査読待ち初稿）}\mbox{}\par

\noindent\textbf{1：代表的条件で適切} Mn・Fe酸化物の黒色団塊。鉄・マンガン酸化物を主体に成長する団塊。原典の1mm以上は分類上の目安として保持する。\par

\noindent\textbf{2：条件・解釈を留保} Cu・Ni・Coを1–2\%含む場合。これら金属を副成分として含む。原典の最大1–2\%を全海域・全団塊の厳密上限とする根拠は未確認。\par

\noindent\textbf{3：代表的条件で適切} 堆積の遅い深海底の表層。団塊の埋没が遅い場所で海水・間隙水から成長する代表的な産状。露岩を覆うクラストと区別する。\par

\noindent\textbf{4：代表的条件で適切} 岩片や化石を核に成長。核を中心に層状に成長する場合がある。\par

\noindent\textbf{5：不適切} 数mm／1万年の成長。代表値は百万年で1mm程度と極めて遅い。原典の万年は代表的な海水沈殿成長の時間尺度と合わない。\par

\paragraph{再利用する考え方} 地質環境・生成機構・物質の名称を対応させ、平均値と全称、典型例と例外、当時の分類と現在の知見を分ける。

\paragraph{共通理論} \hyperref[theory:ore:marine]{第\ref*{theory:ore:marine}節}

\par 参考：\cite{ore50_nodule_ref}。

\par 難易度 4（校準初稿）、優先度未評価。

\section{風化・堆積・マグマ・熱水・石油の濃集機構を比較する。}

\subsection{平成27年度 問31}
\label{q:h27-reg-applied-31}

\noindent\textbf{原典：} sources/H27/17\_applied\_science.pdf、PDF物理12頁。

\textbf{公式正答：} 1（single）。

\paragraph{問題の要旨} 風化・堆積・マグマ・熱水・石油の濃集機構を比較する。

\Needspace{4\baselineskip}\paragraph{確認した条件・記述（独自要約）}\mbox{}\par

\noindent\textbf{原典確認：} S7\_050の29原典画像で34問の全五肢・否定・数値を照合。\par

\noindent\textbf{対象：} 風化・堆積・マグマ・熱水・石油の濃集機構を比較する。\par

\paragraph{問われる核心} 風化・堆積・マグマ・熱水・石油の濃集機構を比較する。

\paragraph{必要知識} 岩石から何が移動し、何が残るかを追う。風化残留でAl・Feが全て溶出し、Na・K・Ca・Mgが残るという交換が誤りである。他の成因は、粒子や溶存成分の堆積、硫化物液の分離と結晶集積、熱水の沈殿、炭化水素の移動とトラップとして区別する。pH・酸化還元条件で元素の移動性が変わる例まで固定した順位にしない。

\paragraph{考え方・正答を選ぶ理由（査読待ち初稿）} 岩石から何が移動し、何が残るかを追う。風化残留でAl・Feが全て溶出し、Na・K・Ca・Mgが残るという交換が誤りである。他の成因は、粒子や溶存成分の堆積、硫化物液の分離と結晶集積、熱水の沈殿、炭化水素の移動とトラップとして区別する。pH・酸化還元条件で元素の移動性が変わる例まで固定した順位にしない。

\Needspace{4\baselineskip}\paragraph{選択肢の検討（査読待ち初稿）}\mbox{}\par

\noindent\textbf{1：不適切} Al・Feを溶かしてNa・K・Ca・Mgを残留。通常の強い風化残留では移動しやすい成分が溶脱し、Al・Fe等の難溶性成分が相対濃集する。列挙した方向が逆。\par

\noindent\textbf{2：代表的条件で適切} 溶存成分沈殿又は粒子堆積。化学的な沈殿と物理的な運搬・分級による堆積鉱床の二つを含む。\par

\noindent\textbf{3：代表的条件で適切} 硫化物メルト分離・結晶集積。マグマ内の液相不混和や結晶分化・集積で金属が濃集する正マグマ成因。\par

\noindent\textbf{4：代表的条件で適切} 熱水中の元素の濃縮・沈殿。温度・圧力・酸化還元・pH等の変化で溶解度が変わり鉱化する。\par

\noindent\textbf{5：代表的条件で適切} 石油のトラップ濃集。移動する炭化水素が貯留層と封じる構造へ集まる。根源岩での生成とトラップは別段階。\par

\paragraph{再利用する考え方} 地質環境・生成機構・物質の名称を対応させ、平均値と全称、典型例と例外、当時の分類と現在の知見を分ける。

\paragraph{共通理論} \hyperref[theory:ore:origins]{第\ref*{theory:ore:origins}節}

\par 参考：\cite{ore50_bauxite_ref,ore50_porphyry_ref,ore50_magmatic_ref,ore50_besshi_ref}。

\paragraph{関連問題} \hyperref[q:r02-reg-applied-26]{R02-REG-APPLIED-26}

\par 難易度 3（校準初稿）、優先度未評価。

\subsection{平成29年度 問28}
\label{q:h29-reg-applied-28}

\noindent\textbf{原典：} sources/H29/17\_applied\_science.pdf、PDF物理9頁。

\textbf{公式正答：} 1（single）。

\par\textbf{検討未確定：} 公式1を保持。原典3のウランを一般の現代河川での漂砂と読むと成立しにくい。古代石英礫岩の還元環境と砂岩型の化学沈殿を区別する必要がある。 原典1は金銀の早期結晶がある場合と述べるため、代表例がCr/Ptであるだけでは全可能性を排除できない。マグマ性PGE鉱床に自然金を含むモデルもあり、公式1の排他理由を要査読。

\paragraph{問題の要旨} 正マグマ濃集・漂砂・熱水・風化・交代を比較する。

\Needspace{4\baselineskip}\paragraph{確認した条件・記述（独自要約）}\mbox{}\par

\noindent\textbf{原典確認：} S7\_050の29原典画像で34問の全五肢・否定・数値を照合。\par

\noindent\textbf{対象：} 正マグマ濃集・漂砂・熱水・風化・交代を比較する。\par

\paragraph{問われる核心} 正マグマ濃集・漂砂・熱水・風化・交代を比較する。

\paragraph{必要知識} 生成機構は正マグマの結晶・硫化物液集積、熱水の沈殿、粒子の漂砂、風化残留、母岩の化学的交代に分ける。早期に金銀を集積するという1を、クロム・白金等の代表例と区別して選ぶ。ウランの漂砂は酸化されにくい古代の環境等を条件とし、現代の河川で全ウラン鉱が砂金と同様に残ると説明しない。

\paragraph{考え方・正答を選ぶ理由（査読待ち初稿）} 生成機構は正マグマの結晶・硫化物液集積、熱水の沈殿、粒子の漂砂、風化残留、母岩の化学的交代に分ける。早期に金銀を集積するという1を、クロム・白金等の代表例と区別して選ぶ。ウランの漂砂は酸化されにくい古代の環境等を条件とし、現代の河川で全ウラン鉱が砂金と同様に残ると説明しない。

\Needspace{4\baselineskip}\paragraph{選択肢の検討（査読待ち初稿）}\mbox{}\par

\noindent\textbf{1：公式の選択・一般化を留保} 早期結晶する金・銀の正マグマ鉱。早期集積の代表はクロム・白金等。金銀鉱床は熱水等の成因が重要で、原典の例示が適切でない。 ただし原典は場合があるという存在命題で、全ての金銀マグマ成因を否定する論証にはならない。\par

\noindent\textbf{2：代表的条件で適切} 金属を運ぶ熱水。亀裂中又は海底等で金属成分を沈殿させる。全熱水鉱床の流体をマグマ起源だけに限定しない。\par

\noindent\textbf{3：条件・解釈を留保} ウラン鉱と砂金の漂砂濃集。機械的な粒子の運搬・堆積による砂金は代表例。ウラニナイトの漂砂は酸素の少ない古い時代等の条件が必要で、砂岩中の溶存Uの酸化還元沈殿とは違う。\par

\noindent\textbf{4：代表的条件で適切} 難溶物の風化残留。母岩の風化で移動しにくい成分が相対濃集し、ボーキサイト等となる。\par

\noindent\textbf{5：代表的条件で適切} 石灰岩との接触交代。熱水等との反応でCaを含むケイ酸塩鉱物を作るスカルン成因。単なる再結晶と分ける。\par

\paragraph{再利用する考え方} 地質環境・生成機構・物質の名称を対応させ、平均値と全称、典型例と例外、当時の分類と現在の知見を分ける。

\paragraph{共通理論} \hyperref[theory:ore:origins]{第\ref*{theory:ore:origins}節}

\par 参考：\cite{ore50_porphyry_ref,ore50_bauxite_ref,ore50_uranium_ref,ore50_magmatic_ref,ore50_besshi_ref,ore50_gold_exception_ref}。

\par 難易度 4（校準初稿）、優先度未評価。

\subsection{令和2年度 問26}
\label{q:r02-reg-applied-26}

\noindent\textbf{原典：} sources/R02/17\_applied\_science.pdf、PDF物理10頁。

\textbf{公式正答：} 3（single）。

\paragraph{問題の要旨} 風化・堆積・マグマ・熱水・石油の濃集機構を比較する。

\Needspace{4\baselineskip}\paragraph{確認した条件・記述（独自要約）}\mbox{}\par

\noindent\textbf{原典確認：} S7\_050の29原典画像で34問の全五肢・否定・数値を照合。\par

\noindent\textbf{対象：} 風化・堆積・マグマ・熱水・石油の濃集機構を比較する。\par

\paragraph{問われる核心} 風化・堆積・マグマ・熱水・石油の濃集機構を比較する。

\paragraph{必要知識} 岩石から何が移動し、何が残るかを追う。風化残留でAl・Feが全て溶出し、Na・K・Ca・Mgが残るという交換が誤りである。他の成因は、粒子や溶存成分の堆積、硫化物液の分離と結晶集積、熱水の沈殿、炭化水素の移動とトラップとして区別する。pH・酸化還元条件で元素の移動性が変わる例まで固定した順位にしない。

\paragraph{考え方・正答を選ぶ理由（査読待ち初稿）} 岩石から何が移動し、何が残るかを追う。風化残留でAl・Feが全て溶出し、Na・K・Ca・Mgが残るという交換が誤りである。他の成因は、粒子や溶存成分の堆積、硫化物液の分離と結晶集積、熱水の沈殿、炭化水素の移動とトラップとして区別する。pH・酸化還元条件で元素の移動性が変わる例まで固定した順位にしない。

\Needspace{4\baselineskip}\paragraph{選択肢の検討（査読待ち初稿）}\mbox{}\par

\noindent\textbf{1：代表的条件で適切} 熱水中の元素の濃縮・沈殿。温度・圧力・酸化還元・pH等の変化で溶解度が変わり鉱化する。\par

\noindent\textbf{2：代表的条件で適切} 硫化物メルト分離・結晶集積。マグマ内の液相不混和や結晶分化・集積で金属が濃集する正マグマ成因。\par

\noindent\textbf{3：不適切} Al・Feを溶かしてNa・K・Ca・Mgを残留。通常の強い風化残留では移動しやすい成分が溶脱し、Al・Fe等の難溶性成分が相対濃集する。列挙した方向が逆。\par

\noindent\textbf{4：代表的条件で適切} 石油のトラップ濃集。移動する炭化水素が貯留層と封じる構造へ集まる。根源岩での生成とトラップは別段階。\par

\noindent\textbf{5：代表的条件で適切} 溶存成分沈殿又は粒子堆積。化学的な沈殿と物理的な運搬・分級による堆積鉱床の二つを含む。\par

\paragraph{再利用する考え方} 地質環境・生成機構・物質の名称を対応させ、平均値と全称、典型例と例外、当時の分類と現在の知見を分ける。

\paragraph{共通理論} \hyperref[theory:ore:origins]{第\ref*{theory:ore:origins}節}

\par 参考：\cite{ore50_bauxite_ref,ore50_porphyry_ref,ore50_magmatic_ref,ore50_besshi_ref}。

\paragraph{関連問題} \hyperref[q:h27-reg-applied-31]{H27-REG-APPLIED-31}

\par 難易度 3（校準初稿）、優先度未評価。

\subsection{令和7年度 問25}
\label{q:r07-reg-applied-25}

\noindent\textbf{原典：} sources/R07/17\_applied\_science.pdf、PDF物理9頁。

\textbf{公式正答：} 3（single）。

\paragraph{問題の要旨} 正マグマ・黒鉱・別子型・スカルンの産状を比較する。

\Needspace{4\baselineskip}\paragraph{確認した条件・記述（独自要約）}\mbox{}\par

\noindent\textbf{原典確認：} S7\_050の29原典画像で34問の全五肢・否定・数値を照合。\par

\noindent\textbf{対象：} 正マグマ・黒鉱・別子型・スカルンの産状を比較する。\par

\paragraph{問われる核心} 正マグマ・黒鉱・別子型・スカルンの産状を比較する。

\paragraph{必要知識} 原典3の黒鉱の母岩が誤りで、緑色片岩ではなく代表的には中新世グリーンタフ地域である。2は正マグマ、4は変成を経た海底硫化物、5は炭酸塩岩の化学交代であり、同じ金属資源でも生成機構を別々に読む。

\paragraph{考え方・正答を選ぶ理由（査読待ち初稿）} 原典3の黒鉱の母岩が誤りで、緑色片岩ではなく代表的には中新世グリーンタフ地域である。2は正マグマ、4は変成を経た海底硫化物、5は炭酸塩岩の化学交代であり、同じ金属資源でも生成機構を別々に読む。

\Needspace{4\baselineskip}\paragraph{選択肢の検討（査読待ち初稿）}\mbox{}\par

\noindent\textbf{1：代表的条件で適切} 通常より鉱物・元素が濃集した鉱床。原典では地質学上の濃集をいう。経済的に採掘できるかは品位・鉱量・価格・技術等を別に評価する。\par

\noindent\textbf{2：代表的条件で適切} Cr・Ptの早期結晶と集積。苦鉄質～超苦鉄質のマグマ系でクロム・白金族等が濃集する代表的な正マグマ鉱床。全鉱体が厳密に下部だけという意味ではない。\par

\noindent\textbf{3：不適切} 黒鉱を東北の緑色片岩中とする。グリーンタフ地域の海底火山熱水に伴う黒鉱と、緑色片岩という変成岩を交換している。\par

\noindent\textbf{4：代表的条件で適切} 海底熱水鉱が変成した層状含銅硫化鉄鉱。海洋地殻上の熱水性硫化物が付加・沈み込みと変成を経る別子型等の代表モデル。\par

\noindent\textbf{5：代表的条件で適切} 炭酸塩との反応でCaケイ酸塩を作るスカルン。熱水による化学的交代で形成される。単なる熱だけの再結晶と分ける。\par

\paragraph{再利用する考え方} 地質環境・生成機構・物質の名称を対応させ、平均値と全称、典型例と例外、当時の分類と現在の知見を分ける。

\paragraph{共通理論} \hyperref[theory:ore:origins]{第\ref*{theory:ore:origins}節}

\par 参考：\cite{ore50_kuroko_ref,ore50_porphyry_ref,ore50_skarn_ref,ore50_magmatic_ref,ore50_besshi_ref,ore50_besshi_ref}。

\paragraph{関連問題} \hyperref[q:r05-reg-applied-23]{R05-REG-APPLIED-23}, \hyperref[q:r06-reg-applied-24]{R06-REG-APPLIED-24}

\par 難易度 3（校準初稿）、優先度未評価。

\section{斑岩銅・黒鉱・別子型・スカルンの生成場を比較する。}

\subsection{平成23年度 問23}
\label{q:h23-reg-applied-23}

\noindent\textbf{原典：} sources/H23/17\_applied\_science.pdf、PDF物理7頁。

\textbf{公式正答：} 1（single）。

\paragraph{問題の要旨} 斑岩銅・黒鉱・別子型・スカルンの生成場を比較する。

\Needspace{4\baselineskip}\paragraph{確認した条件・記述（独自要約）}\mbox{}\par

\noindent\textbf{原典確認：} S7\_050の29原典画像で34問の全五肢・否定・数値を照合。\par

\noindent\textbf{対象：} 斑岩銅・黒鉱・別子型・スカルンの生成場を比較する。\par

\paragraph{問われる核心} 斑岩銅・黒鉱・別子型・スカルンの生成場を比較する。

\paragraph{必要知識} 斑岩銅は浅所の斑岩質貫入と熱水系に関係し、岩体内外の鉱染・網状細脈として形成されるので1。硫化物を含むという共通点だけで、海底の塊状鉱、変成した層状鉱、炭酸塩岩を交代したスカルンを統合しない。

\paragraph{考え方・正答を選ぶ理由（査読待ち初稿）} 斑岩銅は浅所の斑岩質貫入と熱水系に関係し、岩体内外の鉱染・網状細脈として形成されるので1。硫化物を含むという共通点だけで、海底の塊状鉱、変成した層状鉱、炭酸塩岩を交代したスカルンを統合しない。

\Needspace{4\baselineskip}\paragraph{選択肢の検討（査読待ち初稿）}\mbox{}\par

\noindent\textbf{1：代表的条件で適切} 浅所貫入岩内外の鉱染・細脈。斑岩質岩体とその周辺に分散した硫化物と網状細脈を作り、大規模・低品位となる代表的な斑岩銅の産状。\par

\noindent\textbf{2：不適切} 珪長質海底火山の塊状硫化物。火山性塊状硫化物鉱床、特に黒鉱型の説明。斑岩銅と生成位置・形態が異なる。\par

\noindent\textbf{3：不適切} 苦鉄質起源の変成岩中の層状硫化物。別子型等の層状含銅硫化鉄鉱の説明に対応する。\par

\noindent\textbf{4：不適切} 炭酸塩岩の熱水交代。接触交代・スカルン鉱床の説明で、斑岩銅の鉱染・細脈とは区別する。\par

\noindent\textbf{5：不適切} 縞状鉱でCu2+を酸化沈殿。古い縞状鉄鉱の成因を銅へ置換している。鉄の酸化・沈殿を述べる場合と混同しない。\par

\paragraph{再利用する考え方} 地質環境・生成機構・物質の名称を対応させ、平均値と全称、典型例と例外、当時の分類と現在の知見を分ける。

\paragraph{共通理論} \hyperref[theory:ore:hydrothermal]{第\ref*{theory:ore:hydrothermal}節}

\par 参考：\cite{ore50_porphyry_ref,ore50_magmatic_ref,ore50_besshi_ref}。

\par 難易度 3（校準初稿）、優先度未評価。

\section{収束・発散境界と黒鉱・斑岩銅の関係を照合する。}

\subsection{平成24年度 問25}
\label{q:h24-reg-applied-25}

\noindent\textbf{原典：} sources/H24/17\_applied\_science.pdf、PDF物理8頁。

\textbf{公式正答：} 4（single）。

\par\textbf{独立解答と公式正答に相違あり：} 独立解答 2。初回2／公式4。正答原典p41画像も確認。海洋性島弧の斑岩銅と背弧海底の黒鉱型があり、肢2と4の排他性及びマリアナ型の当時定義は要査読。初回を公式へ合わせて置換しない。

\paragraph{問題の要旨} 収束・発散境界と黒鉱・斑岩銅の関係を照合する。

\Needspace{4\baselineskip}\paragraph{確認した条件・記述（独自要約）}\mbox{}\par

\noindent\textbf{原典確認：} S7\_050の29原典画像で34問の全五肢・否定・数値を照合。\par

\noindent\textbf{対象：} 収束・発散境界と黒鉱・斑岩銅の関係を照合する。\par

\paragraph{問われる核心} 収束・発散境界と黒鉱・斑岩銅の関係を照合する。

\paragraph{必要知識} 1と3は収束・発散を交換している。5は堆積盆地の砂岩型ウランをプレート境界へ一対一対応させている。公式原典の平成24年度正答表41頁で4を再確認した。一方、黒鉱型の海底熱水活動は島弧・背弧にも見られ、初回独立2を排除する根拠が不足している。斑岩銅をアンデス型だけと断定することも避け、マリアナ型という用語と鉱床タイプの対応を査読に残す。

\paragraph{考え方・正答を選ぶ理由（査読待ち初稿）} 1と3は収束・発散を交換している。5は堆積盆地の砂岩型ウランをプレート境界へ一対一対応させている。公式原典の平成24年度正答表41頁で4を再確認した。一方、黒鉱型の海底熱水活動は島弧・背弧にも見られ、初回独立2を排除する根拠が不足している。斑岩銅をアンデス型だけと断定することも避け、マリアナ型という用語と鉱床タイプの対応を査読に残す。

\Needspace{4\baselineskip}\paragraph{選択肢の検討（査読待ち初稿）}\mbox{}\par

\noindent\textbf{1：不適切} 発散境界の島弧海溝。島弧海溝系はプレート収束と沈み込みに関係する。発散は海嶺等の説明。\par

\noindent\textbf{2：条件・解釈を留保} マリアナ型と黒鉱。海洋性島弧や伸張する背弧の海底熱水活動に黒鉱型の潜在鉱床がある。マリアナ型の定義を固定しない限りこの肢を簡単に排除できない。\par

\noindent\textbf{3：不適切} 収束境界の中央海嶺。新しい海洋地殻を作る中央海嶺は発散境界である。\par

\noindent\textbf{4：条件・解釈を留保} マリアナ型と斑岩銅。斑岩銅は収束縁の島弧・陸弧等に産する。公式は4だが、マリアナ型への特定と2との差を一次資料から一意に説明する確認は未完了。\par

\noindent\textbf{5：不適切} 収束境界に砂岩型ウラン。砂岩中の流体移動と酸化還元境界等を使う分類であり、収束プレート境界だけを成立条件としない。\par

\paragraph{再利用する考え方} 地質環境・生成機構・物質の名称を対応させ、平均値と全称、典型例と例外、当時の分類と現在の知見を分ける。

\paragraph{共通理論} \hyperref[theory:ore:tectonic]{第\ref*{theory:ore:tectonic}節}

\par 参考：\cite{ore50_porphyry_ref,ore50_backarc_ref,usgs_plate_boundaries}。

\par 難易度 4（校準初稿）、優先度未評価。

\section{ウラン鉱床の母岩・年代・品位の分類を比較する。}

\subsection{令和4年度 問24}
\label{q:r04-reg-applied-24}

\noindent\textbf{原典：} sources/R04/17\_applied\_science.pdf、PDF物理8頁。

\textbf{公式正答：} 5（single）。

\par\textbf{検討未確定：} 公式5を保持。4億年以降と23億年前以前等の厳密な分類時代境界・例外は原典の当時分類とIAEA原分類表の追加照合が必要。低品位資源の存在のみで全鉱床の一律平均品位を作らない。

\paragraph{問題の要旨} ウラン鉱床の母岩・年代・品位の分類を比較する。

\Needspace{4\baselineskip}\paragraph{確認した条件・記述（独自要約）}\mbox{}\par

\noindent\textbf{原典確認：} S7\_050の29原典画像で34問の全五肢・否定・数値を照合。\par

\noindent\textbf{対象：} ウラン鉱床の母岩・年代・品位の分類を比較する。\par

\paragraph{問われる核心} ウラン鉱床の母岩・年代・品位の分類を比較する。

\paragraph{必要知識} ウラン濃集は母岩・流体・酸化還元と形態で分類する。貫入岩型には低～非常に低品位の大規模資源があり、全て高品位とした5を選ぶ。古い石英礫岩の漂砂、不整合近傍の流体、砂岩の還元帯、鉱脈の亀裂充填を区別する。品位と鉱量は別の評価量である。

\paragraph{考え方・正答を選ぶ理由（査読待ち初稿）} ウラン濃集は母岩・流体・酸化還元と形態で分類する。貫入岩型には低～非常に低品位の大規模資源があり、全て高品位とした5を選ぶ。古い石英礫岩の漂砂、不整合近傍の流体、砂岩の還元帯、鉱脈の亀裂充填を区別する。品位と鉱量は別の評価量である。

\Needspace{4\baselineskip}\paragraph{選択肢の検討（査読待ち初稿）}\mbox{}\par

\noindent\textbf{1：代表的条件で適切} 古い基盤と砂岩の不整合近傍。不整合関連型の代表は原生代砂岩盆地と古い変成基盤の不整合をまたぐ熱水・流体鉱化。\par

\noindent\textbf{2：条件・解釈を留保} 4億年前以降の河川砂岩・礫岩。顕生代の堆積盆地は砂岩型の代表だが、4億年を全砂岩型鉱床の唯一の年齢限界とする確認は未完了。\par

\noindent\textbf{3：条件・解釈を留保} 約23億年前以前の石英礫岩。古い酸素の少ない環境でU鉱物等が漂砂濃集する代表的な古代石英礫岩型。年代は酸化環境と地域史の条件として扱う。\par

\noindent\textbf{4：代表的条件で適切} 様々な時代の鉱脈型。亀裂等を切る流体の通路に脈状鉱体ができる形態で、特定時代だけに限定されない。\par

\noindent\textbf{5：不適切} 貫入岩型は他型より高品位。一般に低品位で大規模な貫入岩関連資源があり、高品位を共通特徴としない。\par

\paragraph{再利用する考え方} 地質環境・生成機構・物質の名称を対応させ、平均値と全称、典型例と例外、当時の分類と現在の知見を分ける。

\paragraph{共通理論} \hyperref[theory:ore:uranium]{第\ref*{theory:ore:uranium}節}

\par 参考：\cite{ore50_uranium_ref,ore50_intrusive_ref}。

\par 難易度 4（校準初稿）、優先度未評価。


\chapter{物理探査}
\section{観測量から地下の物性を推定する}
\label{theory:exploration:observables}
地下の岩石をすべて掘削して見ることはできない。物理探査では地表・空中・孔内などで測定した物理量を、地下の物性と幾何構造へ結び付ける。まず「どの信号を測るか」を定め、物性から観測値を計算する順問題を作り、その逆を推定する。観測量が異なれば感度のある対象も異なる。

\begin{longtable}{p{0.23\textwidth}p{0.29\textwidth}p{0.34\textwidth}}
\hline
手法 & 主な観測量 & 主な物性・構造\\\hline
重力 & 重力加速度の差 & 密度差と形状・深度\\
磁気 & 磁場・磁気異常 & 誘導磁化・残留磁化\\
地震・音波 & 弾性波・音波の走時や波形 & 波速、インピーダンス、境界\\
電気・電磁 & 電位差、電場・磁場応答 & 比抵抗・導電率\\
地中レーダ & 高周波電磁波の反射 & 誘電率・導電率の境界\\
放射能 & 主に天然ガンマ線の強度 & 表層付近の核種分布や移動\\
地温 & 地中・地表面の温度 & 熱の伝導・移流、熱水等\\
\hline
\end{longtable}

地震・音波探査では発信源と受振器を混同しない。地上で用いるジオフォンは一般に地盤の速度に応答する。一方、水中曳航のストリーマでは、管内に並べたハイドロフォンが音圧変動を測る。海水の水中音波受信を、地上の速度計を防水にしただけの装置と説明するのは不適切である\cite{usgs_marine_seismic}。

AE（アコースティック・エミッション）は、岩石などの微小破壊・変形に伴って放出される過渡的な弾性波である。岩石破壊実験や発生位置の推定に用いられる\cite{sato_lei2009_ae}。対象が岩石であることは重力・磁気探査との共通点だが、直接の観測量は異なる。令和6年度問34ではこの区別が問われる。

\section{MT・CSAMTとレーダを区別する}
\label{theory:exploration:electromagnetic}
MT（地磁気地電流法）は自然に変動する電場と磁場を測定し、地下の比抵抗構造を推定する\cite{usgs_mt_metadata}。電場と磁場の比を複素インピーダンスと呼ぶ。均質・等方な導電半空間で、変位電流を無視した平面波近似を考えると、導電率\(\sigma=1/\rho\)について
\[
 |Z|^2=\frac{\omega\mu}{\sigma},\qquad
 \rho_a=\frac{|E/H|^2}{\omega\mu},\qquad
 \delta=\sqrt{\frac{2\rho}{\omega\mu}}
\]
である。\(\delta\)は振幅が減衰する代表距離で、周波数が低いほど、また比抵抗が高いほど大きい。実際の層構造では\(\rho_a\)は各層の抵抗率そのものではなく、周波数ごとの応答を表す見掛比抵抗である。観測された\(B\)から\(H\)へ変換する場合には\(B=\mu H\)の単位と係数をそろえる。

CSAMTは人工送信源を用い、十分遠方の条件でMTに類似する電磁応答を解析する。送信源の近傍や遷移域では単純な平面波式が成り立たず、その効果の扱いが必要になる。CSAMTによる実地の比抵抗断面の例として西条平野の調査を参照する\cite{ochi_csamt_saijo}。磁場も測るからといって磁化率を主要な推定対象と考えるのが、平成27年度問19の誤りである。

地中レーダはこれとは異なり、電磁波を放射して境界からの反射を測る。通常の極浅部調査で用いる比較的高い周波数帯と、誘電率・導電率の違いを押さえる。高周波で常に深く見えるわけではなく、導電性の高い土壌や水による減衰が到達深度を制限する。

\section{測定と解釈の限界を分ける}
同じ重力異常を、浅い小岩体と深い大岩体が近似的に生むことがある。磁気には磁化方向、電磁探査には含水・塩分・粘土などの複数要因がある。したがって「異常を検出した」ことから特定岩種や製品を直ちに断定しない。物性、幾何形状、測定感度、ノイズの前提を列挙し、別の探査や露頭・ボーリングで制約を加える。

重力補正の具体的な導出は重力・測地章、誘導磁化と残留磁化の区別は地磁気章に置く。以下ではそれらを参照し、個別問題では指定条件と誤った対応だけを特定する。地震探査の走時・信号処理は次節で詳しく扱う。孔内計測の詳細は残る原典バッチとともに追加する。
\section{探査データ：測る、整える、地下へ戻す}
\label{theory:exploration:data}
受振器に入る波形は地下構造の写しそのものではない。震源波形、伝播経路、反射・屈折、計器の応答、必要でない波や環境雑音が重なった記録である。計測配置を設計し、時間と振幅をそろえ、対象成分を取り出して地下モデルへ結び付ける。フィルタ、重合、走時補正、画像化、逆解析はそれぞれ異なる役割を持つ。

\subsection{線形性と時不変性：応答を重ねる条件}
\label{theory:exploration:lti}
線形性は入力の和や定数倍に対する応答が同じ操作を保つことである。時不変性は入力を時間移動したとき、応答も同じだけ移動することである。両方が成立するLTI系では単位インパルス応答 \(h\) を使い、
\[
 y(t)=\int_{-\infty}^{\infty}x(\tau)h(t-\tau)\,d\tau,
 \qquad y[n]=\sum_k x[k]h[n-k]
\]
と表せる\cite{mit_signals2010}。任意入力を小さなインパルスの和へ分解し、その応答を加えるという意味である。

時変の線形系でも重ね合わせはできるが、一般には入力時刻にも依存する \(h(t,\tau)\) が必要である。「線形」だけから単一応答の時間移動による畳み込みまで保証しない。計器の応答を補償するときも、応答の零点付近では雑音を強く増幅する場合がある。

\subsection{フィルタ：残す周波数と抑える周波数}
\label{theory:exploration:filters}
低域通過は低周波数を残し高周波数を抑える。\(T=1/f\) なので、長周期を残す処理である。帯域通過は範囲内を残す。ノッチは狭い帯域を阻止する処理であり、商用電源等の周期的混入に用いる\cite{scipy_notch}。ノッチを「特定周波数だけを通す」と覚えると保持と除去が逆になる。

理想フィルタの通過利得1・阻止利得0はモデルの定義である。有限長の実装には遷移帯域、リップル、位相変化、端の処理がある。抑えたい雑音と対象信号の帯域が重なれば、雑音だけを無条件に取り除くことはできない。二次元の画像には空間方向の平滑化や微分を行い、エッジや構造を調べる。フィルタは二次元だけに限定された概念ではない。

\subsection{移動平均：分散の低減と周期消去の反例}
\label{theory:exploration:moving}
\(M\) 点の等重み平均 \(y[n]=M^{-1}\sum_{k=0}^{M-1}x[n-k]\) は平滑化である。各点に独立・平均ゼロ・分散 \(\sigma^2\) の雑音が含まれると、平均後の雑音分散は \(\sigma^2/M\)、標準偏差は \(\sigma/\sqrt{M}\) となる。有限点でランダム雑音が全て消えるわけではない。中央窓は未来側の値も使うため、即時の処理と後処理を区別する。

周波数応答は等比級数を足して独自に導くと、
\[
 H(e^{i\omega})=\frac1M\sum_{k=0}^{M-1}e^{-i\omega k}
 =e^{-i\omega(M-1)/2}\frac{\sin(M\omega/2)}{M\sin(\omega/2)}.
\]
\(\omega=2\pi j/M\) の非DC成分では零点を持つ。このため、周期がちょうど \(M\) 標本の正弦波は消せる。\(M=4\) なら一周期の四つの値の和がゼロとなる。一方、周期がずれれば同じ窓で完全には消えない。MITの移動平均の応答図も周波数によって利得と零点が変わることを示す\cite{mit_moving2012}。平成28年度問20の公式2を、移動平均では周期雑音を決して除去できないという数学的命題へ変更しない。

\subsection{スタックの平方根則を信号と分散から導く}
\label{theory:exploration:stack}
同じ対象信号 \(s\) と雑音 \(n_j\) の記録を \(N\) 本重ねる。
\[
 \sum_{j=1}^{N}(s+n_j)=Ns+\sum_jn_j,
 \qquad \operatorname{Var}\!\left(\sum_jn_j\right)=N\sigma^2
\]
は無相関・同分散雑音の条件で成り立つ。信号振幅は \(N\) 倍、雑音標準偏差は \(\sqrt N\) 倍なので、振幅S/Nの改善倍率は \(\sqrt N\) である。平均として \(N\) で割っても比は同じである。電力S/Nなら改善倍率は \(N\) となる。\(N^2\) はどちらの改善倍率でもない。

一般の雑音では分散に \(2\sum_{i<j}\operatorname{Cov}(n_i,n_j)\) が加わる。同じ電源雑音が各記録に残れば、回数だけ増やして平方根則を期待できない。対象波形も時刻がずれれば相殺されるため、走時差や極性をそろえる。USGSの実データ処理ではNMOとCMP重合、重合後の処理が別々に記録される\cite{usgs_processing2000}。

受振器群は独立雑音の平均化に加え、入射方向と空間波長によるフィルタとして働く。表面波は空間的に相関した波であり、単純な独立雑音とは異なる。群間隔と配列長によって対象反射も弱まる場合がある。同じ起振点で反復する垂直スタックと、共通中点の異なる経路をそろえるCMP重合を区別する。

\subsection{標本化：二倍という数と境界条件}
\label{theory:exploration:sampling}
標本化は \(x[n]=x(nT_s)\)、標本化周波数は \(f_s=1/T_s\) である。帯域制限信号を理想的に再構成するには、周期的に複写されるスペクトルが重ならないことが必要である。最高周波数成分を含む信号には \(f_s>2f_{\max}\) と余裕を取る\cite{mit_signals2010}。

等号では \(x(t)=\sin(2\pi f_{\max}t)\) を \(t=n/(2f_{\max})\) で測ると全てゼロになる。最高周波数を含まない帯域や境界成分の扱いを定義しないまま、「二倍以上だから全て完全復元」と断定しない。現実には量子化、有限記録長、時刻誤差、帯域外雑音もある。aliasingは異なる連続信号が同じ標本を与える現象であり、取得後にFFTをかけるだけでは区別できない。

\subsection{FFT・相互相関・コヒーレンス}
\label{theory:exploration:fft}
有限標本のDFTは \(X[k]=\sum_{n=0}^{N-1}x[n]e^{-2\pi i kn/N}\) である。FFTは同じ和を構造を利用して速く計算する算法である。radix-2では二のべき乗の標本数が扱いやすいが、それ以外が計算不能ではない\cite{scipy_fft}。ゼロを追加して計算点を増やしても、実測時間が延びたり近い二周波数を区別する情報が増えたりするわけではない。

\label{theory:exploration:correlation}
相互相関 \(R_{xy}[\ell]=\sum_n x[n]y[n+\ell]\) は一方をずらして類似を測る。採用する符号 convention と規格化によって遅れの符号や振幅を確認する。周期信号なら複数のピークがあり、最大値だけで唯一の到達遅れを保証しない。

代表的な二乗コヒーレンスは \(C_{xy}(f)=|S_{xy}(f)|^2/(S_{xx}(f)S_{yy}(f))\) であり、周波数ごとの関係を測る\cite{scipy_coherence}。相互相関と関連するが、時間遅れを横軸にする同じ量ではない。非零の自己スペクトルと適切な統計推定の条件を考え、二信号の因果関係をこの値だけで断定しない。

\subsection{走時・静補正・反射・CMPを結ぶ}
\label{theory:exploration:seismicterms}
単純な水平二層の屈折波では走時直線の切片を原点走時と呼ぶ。観測距離ゼロの実測初動と同一ではなく、直線の外挿である。中間層が薄く、初動として一度も現れない場合がある。また、低速度層は通常の臨界屈折条件を満たさず検出しにくい。USGSはこれらを別々に説明する\cite{usgs_refraction2013}。ブラインド層の日本語の使用範囲が広い場合もあるため、層が見えない仕組みを明示する。

単純な垂直入射の反射では \(Z=\rho V\)、\(R=(Z_2-Z_1)/(Z_2+Z_1)\) が反射振幅を支配する。\(Z\) はインピーダンス、\(R\) は反射係数であり別の量である\cite{epa_seismic_reflection}。静補正は地形や表層の低速度層等が生む走時ずれをそろえる\cite{usgs_static2004}。ブーゲー補正は質量による重力を扱い、走時補正ではない。

CMP重合では共通中点の記録を集め、通常NMO等の時間差を補正する。水平・均質の近似から外れると、共通中点が同じ反射点を意味するとは限らない。速度モデル、走時と深度の対応、データの範囲を確かめる。

\subsection{下方接続とmigrationは地震波にも用いる}
\label{theory:exploration:continuation}
migrationは傾斜面の反射や回折記録を地下の位置へ再配置する画像化である。Stanfordの説明では地表の波動場を深度方向へ下方接続し、適切な時刻の波動場から画像を作る\cite{stanford_downward2001}。重力・磁気の下方接続と名前を共有するが、伝播させる場と方程式が異なる。下方接続をポテンシャル場専用の処理と覚えるのは誤りである。

令和元年度問34の初回5はこの限定から生じた。初回記録は保存し、波動下方接続の一次説明と、送受信コイルを用いるrigid-boom EMIの電磁誘導の原理を別に確認して4へ訂正した\cite{rigid_boom2019}。後者は一次論文の検索収録を確認した範囲であり、PDF全文の取得・読了には至っていない。装置の日本語名の最終照合は査読待ちとする。

\subsection{順問題・逆問題・正則化：何を更新するのか}
\label{theory:exploration:inverse}
順問題はモデル \(m\) を与えてデータ \(F(m)\) を計算する。逆問題はデータ \(d\) からモデルを推定する。計測はデータの取得であり、計測そのものを順問題の数値計算と同一視しない。代表的な目的関数は
\[
 \Phi(m)=\|W_d(F(m)-d)\|_2^2+\lambda\|L(m-m_0)\|_2^2
\]
である。最初の項は誤差で重み付けした残差、第二項は滑らかさや参照モデルからの差への拘束である\cite{pygimli_inversion}。データを誤差より細かく再現しても、雑音をモデルへ取り込んでいる可能性がある。

Marquardt型の減衰付き更新、共役勾配、遺伝的探索はモデル更新の算法である。正則化・平滑化拘束は逆解の安定性を制御する。有限要素法や差分法は支配方程式の離散化を担い、\(F(m)\) や感度を計算するため逆解析の内部でも利用できる。平成30年度問35の公式5は直接の役割の比較として読むが、FEMがインバージョンに無関係という絶対命題にはしない。

\subsection{波動の基本則と等価層・検層}
\label{theory:exploration:principles}
線形波動では重ね合わせが成立する。相反性を持つ媒体では震源と受振点を交換した応答に対応がある。ホイヘンスは波面構成、スネルは境界の位相整合による角度と波速、フェルマは波線の走時の停留条件を表す。最短時間という単語だけで複数経路の可能性を排除しない。レンツの法則は磁束変化を妨げる電磁誘導の向きに関する法則であり、弾性波の屈折則とは異なる。

USGSの等価層は重力の密度層や磁気の双極子層によってポテンシャル場を表す説明である\cite{usgs_equivalent2007}。場の表現を波動伝播の基本則と混同しない。ラテラル法は孔内の電極間隔と配置による電気検層の方式である\cite{usgs_lateral_logging}。計測配置の名称と、取得後のフィルタ・差分・最適化の算法を区別する。

\section{屈折・反射・表面波：地下を拘束する情報の違い}
\label{theory:exploration:seismic}
地震探査の記録には直接波、屈折波、反射波、表面波が重なる。屈折法では通常初動走時、反射法では反射イベントの走時と波形、表面波法では周波数ごとの位相速度を主に使う。名前の違いは処理だけではなく、地下構造を拘束する観測情報の違いである。

\subsection{速度と地質：波種・状態・境界を分ける}
\label{theory:exploration:waveproperties}
安定な等方弾性体では、体積弾性率 \(K\)、せん断弾性率 \(G\)、密度 \(\rho\) を使い、
\[
 V_P=\sqrt{\frac{K+4G/3}{\rho}},\qquad
 V_S=\sqrt{\frac{G}{\rho}}
\]
となる。\(K>0,G>0\) なら \(V_P>V_S\) である。流体ではせん断を保てず通常のS波が伝わらないので、二つの伝播速度を常に観測して比較できるとは限らない。

岩種が同じでも、割れ目の開口・接触、風化、変質、空隙、飽和が有効弾性係数を変える。空隙の空気を水が置き換えると圧縮性等へ影響し、P波速度は変わり得る。USGSの測定報告では飽和帯の深さとP波速度を読み、表面波のS波逆解析を拘束する\cite{usgs_refraction2013}。地質学の境界は岩種・年代・堆積単位等の区分であり、速度差の境界と同一ではない。速度から風化や破砕を推定するとき、露頭やボーリングと併せて解釈する。

通常の浅部屈折ではP波が多く使われるが、S波屈折も存在する\cite{epa_seismic_refraction}。土質のS波調査を比較的平坦な地形で行うという運用の目安を、急峻地形では物理的に不可能という命題へ変更しない。強震動にはS波構造が重要であり、通常の地山P波屈折の主用途と、複数手法の統合による寄与を分ける\cite{segj_integrated2017,japansea_spac2017}。

\subsection{屈折初動：はぎとりと往復の整合}
\label{theory:exploration:refraction}
上層速度 \(V_1\)、下層速度 \(V_2>V_1\)、上層厚 \(h\) の単純な水平二層では、\(\sin i_c=V_1/V_2\) を満たす臨界屈折波の走時は
\[
 t_{\rm head}(x)=t_i+\frac{x}{V_2},\qquad
 t_i=\frac{2h\cos i_c}{V_1}
\]
となる。直接波の \(x/V_1\) と比べて最初に到達する時刻を読む。原点走時 \(t_i\) は走時直線の外挿切片であり、起振位置の実測初動ではない。臨界距離と直接波を追い越す交差距離も同じ概念ではない\cite{epa_seismic_refraction}。

萩原のはぎとり法は表層の走時寄与を分離して下層の速度と境界を求める屈折初動の方法である。JICA調査団の比較表は直接波＋屈折波を扱い、トモグラフィとの構造・速度の仮定を分けている\cite{jica_refraction2015}。反射波の走時をそのまま同じはぎとりへ入力する説明は誤りである。

端点AからBへの全走時とBからAへの全走時は、相反的な同じ経路・媒体で整合する。全走時一致はこの往復整合などを確かめる意味であり、全ての受振点の到達時刻が同じという意味ではない。原点走時の整合、曲線の平行性も観測モデルに沿って確認する。薄い中間層は屈折波があっても初動として現れない距離範囲しか持たない場合がある。低速度逆転、薄いブラインド層、連続速度勾配を分ける\cite{usgs_refraction2013}。

連続的に深度増速するミラージュ構造では波線が曲がり、走時は折れ線ではなく曲線となる。トモグラフィは波線が通る格子速度又は遅さを反復更新し、連続勾配も扱える。ただし通らない領域の速度を一意に決めることはできない。測線は地形と地質の変化を見ながら設計し、等高線に垂直で一定勾配という望ましい典型配置だけを絶対条件にしない。

\subsection{反射走時：双曲線と放物線の適用範囲}
\label{theory:exploration:reflectiongeometry}
単一水平反射面、上部の一定速度 \(V\)、深さ \(h\)、震源と受振器の距離 \(x\) なら、反射点は両者の中点にある。左右対称の経路長を足して独自に導くと、
\[
 t(x)=\frac{2}{V}\sqrt{h^2+(x/2)^2},\qquad
 t^2=t_0^2+\frac{x^2}{V^2},\quad t_0=\frac{2h}{V}.
\]
\(t\)-\(x\) 平面では双曲線である。多層では小〜中オフセットの範囲で有効な重合速度を使って双曲線近似する。大オフセット、強い不均質、異方性等には非双曲線成分が残る。

さらに \(|x|\ll Vt_0\) なら平方根を展開して、
\[
 t(x)=t_0+\frac{x^2}{2V^2t_0}
 -\frac{x^4}{8V^4t_0^3}+\cdots
\]
となる。二次項までの近似は放物線である。このため平成28年度問21と令和5年度問32の公式選択を、放物線ではどんな近似もできないという命題へ書き換えない。通常の走時モデルが双曲線であることと、原文に近似範囲の指定がない問題点を併記する。

NMOは各オフセットの走時をゼロオフセット走時へそろえる。重合速度はその目的に最適な速度であり、真の区間速度と無条件に同じではない\cite{usgs_processing2000}。反射とは異なるmoveoutの多重反射等は重合で弱まる場合があり、相関雑音も全て残るとは限らない。同じmoveoutの雑音は抑えにくい。

浅層反射の最大オフセットを目標深度程度とする説明は通常の展開目安である\cite{coastal_reflection2016}。目標反射角、近接雑音、速度解析、深度、開口と重合数で設計を変える。二倍を一律に必要とする説明を誤りと選ぶ場合も、二倍の観測が常に不適切という意味ではない。

\subsection{分解能と検出：薄層が存在することを知る}
\label{theory:exploration:resolution}
厚さ \(\Delta h\) の上下面反射の時間差は単純な垂直経路で \(\Delta t=2\Delta h/V\) である。二つのwaveletが強く重なると別々の反射として分離しにくくなる。代表的な垂直分解能の目安 \(\Delta h\sim\lambda/4\) は、上面と下面の分離に関するものである\cite{seg_vertical_resolution}。

分離できない薄層も、合成振幅や周波数応答の変化として存在を検出できる場合がある。検出にはS/N、インピーダンス差、位相、帯域が関わる。「四分の一波長より薄ければ必ず見えない」と断定しない。Rayleigh基準やRicker基準、waveletは反射イベントの分離に関係し、屈折直線のintercept timeとは役割が異なる。

横方向にはFresnel zoneの有限範囲から反射が寄与する\cite{seg_lateral_resolution}。単純なゼロオフセットの幾何では往復経路差を半波長とする一次zoneから \(r_F\approx\sqrt{h\lambda/2}\) の目安を独自に得る。migration、開口、帯域と速度モデルによって横方向の分解能は変わる。zone径を全処理の最終解像度と固定しない。

\subsection{表面波分散とS波モデル}
\label{theory:exploration:surfacewaves}
Rayleigh波は表面に沿って進み、深度へ振幅が減少する。不均質な成層では周波数により感度深度が変わるため位相速度 \(c(f)\) が分散する。長波長は深部、短波長は浅部を拘束する傾向があるが、一つの波長が一つの深度だけを測るわけではない\cite{epa_surfacewaves}。均質等方半無限体は非分散という退化例である。

観測位相速度から地下モデルを作り、理論分散曲線を計算して残差を小さくする。位相速度は \(V_P,V_S,\rho\)、層厚に依存するが、通常の地盤分散は \(V_S\) に強い感度を持つ。Xiaらの一次論文要約は高周波側でS波速度と層厚の影響が大きいことを述べる\cite{xia_rayleigh1999}。P波・密度の影響をゼロにせず、別の測定や拘束を使う。

水平多層の一次元逆解析は通常のモデルであり、地盤自体が常に横方向に均質と分かったわけではない。高次モード、方向による変化、モデル自由度とデータ数を確認する。能動表面波の通常の浅部対象を20m程度とする資料がある\cite{moe_surface2004}。しかし20mや100mは普遍的上限ではなく、震源の長波長、受振器の低周波応答、配置、雑音と解析範囲に依存する。微動で長波長側を補えれば深部拘束を増やせるが、併用だけで深度増加を保証しない。

\subsection{微動アレイの短波長と長波長の限界}
\label{theory:exploration:arraylimits}
空間標本化で間隔 \(d\) に対し波長が短すぎると、異なる空間波数が同じ標本を与える。等間隔・単方向の単純モデルでは \(\lambda>2d\) の条件が基本となる。実アレイの限界は配置と波の方向にも依存する。空間aliasingは主に最短波長側の制約であり、最長側と取り違えない\cite{japansea_spac2017}。

長波長ではアレイ内の位相差が小さく、有限開口、ノイズ、推定算法と観測時間の制約が効く。F-Kは周波数ごとの波数分布を求め、SPACは空間自己相関と分散の関係を使う。通常の条件でSPACがF-Kより長波長まで推定する例が一次論文にあるが、入射方向と多数の波、構造、配置、実装に依存する\cite{array_limits2002}。一定の半径倍率を全観測の保証値へ転記しない。複数サイズで重なる帯域を確認し、疑わしい分散曲線を機械的に一つのモデルへ合わせない。

\section{電気探査：地下へ流した電流から何が分かるか}
\label{theory:exploration:dcarrays}
地層を掘る前に、水や粘土、空隙、金属などによる電流の流れやすさの差を調べたい。比抵抗法は電流を注入して地表や孔内の電位差を測り、地下の抵抗率分布を推定する。抵抗\(R\)は試料の形状も含むが、均質な棒では\(R=\rho L/A\)である。長さ\(L\)、断面積\(A\)から\(\rho=RA/L\)となるため、抵抗率の単位は\(\Omega\,\mathrm m\)であり、\(\Omega\,\mathrm m^3\)ではない。局所形では\(\boldsymbol E=\rho\boldsymbol J\)で、導電率は\(\sigma=1/\rho\)である。

定常電流の保存と\(\boldsymbol E=-\nabla V\)を合わせると、源のない部分では
\[
 \nabla\cdot(\sigma\nabla V)=0
\]
となる。注入電極を源・回収電極を吸込みとして、地表の境界条件と地形を含めて電位を計算するのが順問題である。多数の電位差へ一致する\(\sigma\)を探すのが逆問題で、電位だけで地質名が一意に決まるわけではない\cite{epa_resistivity}。

\subsection*{半空間の電位とWennerの幾何係数}
教材の説明用に均質・等方な導電半空間、絶縁した平らな地表、小さい点電極を仮定する。地表の注入点から半球面を通る電流を\(I\)とすると、\(J_r=I/(2\pi r^2)\)である。したがって無限遠基準の電位は\(V(r)=\rho I/(2\pi r)\)となる。注入Aと回収Bの重ね合わせより、電位電極M,Nの差は
\[
 \Delta V=\frac{\rho I}{2\pi}
 \left(\frac1{AM}-\frac1{BM}-\frac1{AN}+\frac1{BN}\right).
\]
Wenner配置を地表直線上のA--M--N--B、全隣接間隔\(a\)とすると、括弧は\(1/a\)であり
\[
 \rho_a=2\pi a\frac{\Delta V}{I}
\]
を得る。この導出は教材独自である。符号は電極と電位差の定義をそろえる。不均質な地盤の\(\rho_a\)は「この配列で同じV/Iを出す均質半空間」の値であり、MとNの真下にある一層の真の抵抗率ではない。

Schlumberger、dipole--dipole等も幾何係数と地下感度が異なる。垂直探査では中心をおおむね固定し間隔を広げ、水平多層等のモデルを推定する。水平探査では配置を移動して横の変化を調べる。固定間隔でも感度は有限体積に分布するので「特定の深度だけの水平抵抗率」を直接測ると読まない。二次元・三次元の比抵抗トモグラフィは、多数配置のデータを順モデルへ合わせて分布を求める。偽断面の表示位置を実深度と同一視しない。

\subsection*{線形フィルタは配置でなく順応答の計算法}
層状大地モデルの比抵抗変換から見掛比抵抗曲線を計算する際、線形フィルタを用いる方法がある。USGSのTI-59プログラムはその具体例である\cite{usgs_linear_filter1981}。平成23年度問30では、配列名に混じるリニアフィルタを一般の雑音除去だけと解釈しない。IPは追加の分極応答を測るが、比抵抗と同じ電極を使うため、無関係という全称も避ける。

\subsection*{四電極で分ける接触電圧降下}
\label{theory:exploration:fourpoint}
二端子で電流を流しその端子電圧を測ると、地盤だけでなく両電極の接触抵抗による電圧降下も含む。別のM,Nを高入力抵抗\(R_{\rm in}\)の電圧計へつなぎ、測定電流\(I_v\)を十分小さくすれば、接触誤差の代表値\(I_v(R_{cM}+R_{cN})\)を小さくできる。これが四電極の要点である。送信側の接地抵抗は電流を確保する問題として残り、四電極なら接地・地形・電磁カップリングの影響が全てゼロになるわけではない。平成25年度問34はこの分離を問う。

\section{含水岩石の電流経路とArchie式}
\label{theory:exploration:archie}
多くの岩石骨格は絶縁性が高く、孔隙水に溶けたイオンが電流を運ぶ。孔隙率\(\phi\)、水飽和度\(S_w\)、水の抵抗率\(\rho_w\)について、代表的な実験モデルは
\[
 \rho_t=a\rho_w\phi^{-m}S_w^{-n}
\]
である。\(a,m,n\)は骨格・連結性・濡れ条件等に依存する係数で、全岩石に同じ値を置かない。飽和では\(\rho_t/\rho_w=a\phi^{-m}\)が地層係数となる。孔隙ネットワークの少ない導電断面と迂回路をまとめたモデルであり、普遍的な保存則から全係数が決まる式ではない。

粘土表面、導電鉱物、電子伝導に無視できない電流があれば、孔隙水のみという前提が破れる。単純な並列伝導モデル\(\sigma_{\rm bulk}=\sigma_{\rm pore}+\sigma_{\rm surface}\)を考えると、追加の表面伝導は抵抗率を下げる。低抵抗の水は相対的に孔隙水伝導を優勢にしやすい。逆に高抵抗水だけでは式の精度を保証できない。ただし清浄な絶縁骨格なら淡水でも利用し得る\cite{epa_borehole_resistivity}。

海水の塩分は通常導電率を上げ抵抗率を下げる。不飽和化は連続水経路を減らし高抵抗化しやすい。粘土化・含水割れ目は低抵抗化しやすいが、乾燥した割れ目や珪化まで同じ傾向にはしない。「礫層は高抵抗」という典型も周辺の粘土層との対比・塩分条件による。原典の一般傾向を、岩石名から値を一意に決める規則にしない。

\section{IPとSP：人工分極と自然電位の機構}
\label{theory:exploration:ipsp}
IP（誘導分極）は人工電流を流した際の地下界面の電荷蓄積を利用する。電流遮断後の電位減衰又は周波数応答の位相を測る。硫化鉱物の電子伝導と電解質のイオン伝導の界面、鉱物表面等が応答へ寄与する。比抵抗だけでは区別しにくい界面状態を調べるための追加情報であり、IPでも比抵抗を測る。ケーブル間の電磁結合が分極応答へ混ざり得るため、周波数・配置・測定時間を吟味する\cite{epa_ip}。

SP（自然電位）は人工的に電流を注入せず、基準点との電位差を非分極電極で測る。地下水流による電気二重層の電荷輸送、酸化還元差を結ぶ導電性鉱体、温度や化学勾配等が源となる。説明用に\(\boldsymbol J=-\sigma\nabla V+\boldsymbol J_s\)と置けば、電流保存は
\[
 \nabla\cdot(\sigma\nabla V)=\nabla\cdot\boldsymbol J_s
\]
となる。観測から源\(\boldsymbol J_s\)を推定するには抵抗率と源機構のモデルが必要である。電位の絶対零点は測らず、基準点の選び方も記録する。

典型的な硫化鉱体の鉱化電池では地表の負異常を利用する。一方、流動・熱・化学が重なる場の符号は一つの法則へ固定できない。正異常だから必ず熱水、負異常だから必ず鉱体と断定しない。SP異常の幅や幾何モデル、他観測との同時逆解析による深度推定は可能だが、SP分布だけで深度が一意に読めるわけではない\cite{epa_sp}。令和3年度問32の公式4は保持するが、「SPでは深度解析ができない」という誤った禁止を教材にしない。原典の否定範囲に残る疑義を問題ごとに示す。

\section{有限送信源・周波数領域・時間領域}
\label{theory:exploration:nearfield}
自然MTと人工CSAMTの抵抗率・表皮深度は前節\ref{theory:exploration:electromagnetic}の式へつながる。均質導電半空間で変位電流を無視すると、周期場の拡散による振幅減衰尺度は\(\delta=\sqrt{2\rho/(\omega\mu)}\)である。送受信距離\(r\)に対して\(r/\delta\)が十分大きい遠方では平面波近似を使いやすい。固定rで低周波化又は高抵抗化すればδが増し、近接場へ移る。USGS実地報告は三表皮深度超を目安として用いるが、その数を全装置・全構造の厳密閾値とはしない\cite{usgs_csamt2012}。

\(\rho_a=|E/H|^2/(\omega\mu)\)のµは透磁率である。真空ではµ0を用い、誘電率ε0ではない。HをBと取り違えると換算係数と単位も変わる。平成28年度問23は式の形だけでなく、この変数定義の誤りを問う。

\subsection*{周波数領域では振幅と位相を測る}
\label{theory:exploration:fdtd}
FDEMでは正弦電流が送信コイルの一次磁場を作り、地盤の渦電流が二次磁場を生む。定常応答の複素振幅比を同相・直交成分又は振幅・位相として測る。Slingramやloop--loopはこの仕組みを使う。電極を挿さなくてもよいが、コイル間隔、周波数、導電率、姿勢に感度が依存する\cite{epa_fdem}。

TDEM/TEMでは送信を遮断した後の過渡応答を測る。均質導電体の準静的近似では\(\partial\boldsymbol B/\partial t=(1/(\mu\sigma))\nabla^2\boldsymbol B\)となるため、代表拡散距離は\(\sqrt{t/(\mu\sigma)}\)の尺度で増す。遅い時刻には一般に深部への感度があるが、信号は弱くなり、装置寸法・電流・遮断時間・測定時間・ノイズが限界を決める。central loopで深部を測れることを「送受信間隔によらず無制限に深く測れる」としない\cite{epa_tdem_archive}。一次場の遮断後観測と、電場の局所歪みによるMTのstatic shiftの少なさは別の性質である。

VLFは既存通信局の3--30kHz帯の信号を使う代表的方式で、局の方位と対象構造の向きが重要である。調査者は既存局の送信条件を任意には変更できない。Schumann共振も電気電磁現象であり、MT観測記録で調べられている\cite{schumann_mt2009}。反射地震のNMO補正とは分野が異なる。

\section{GPR：走時断面から電気境界を推定する}
\label{theory:exploration:gpr}
地中レーダは高周波電磁パルスを放射し、電気特性の異なる境界からの反射を受ける。低損失・非磁性の均質媒質なら
\[
 v\simeq\frac{1}{\sqrt{\mu_0\epsilon}}
 =\frac{c}{\sqrt{\epsilon_r}},\qquad
 \epsilon_r=\frac{\epsilon}{\epsilon_0},\qquad
 z\simeq\frac{vt}{2}
\]
である。tは往復走時であり、直下水平面などの単純な幾何と速度が必要である。εはF/m、εrは無次元で、空気を分母にするのは通常の近似に過ぎない。高導電性・分散媒質ではこの単純速度式の条件が破れる。

波形をアンテナ位置へ並べるprofileはまず時間断面である。wide-angleや共通中点配置から速度を推定し、深度変換に用いる。埋設管等を横切る測線は位置を捉えやすい。点対象の回折双曲線をそのまま地層形状と思わず、速度・偏位・必要なmigrationを考える\cite{epa_gpr}。

\subsection*{分解能力と最小識別距離を区別する}
\label{theory:exploration:gprresolution}
二つの反射を時間差\(\Delta t\)で識別するにはパルス幅・帯域が制約になる。説明用には帯域Bで\(\Delta t\sim1/B\)、距離分解尺度\(\Delta z\sim v/(2B)\)である。帯域が中心周波数に比例する同系統装置なら、高周波ほど最小識別距離が小さく、能力としての「分解能」は良い。この語の二つの使い方で比例の向きが変わるので、式と変数を明示する。

実際の探査深度は導電減衰、散乱、幾何拡がり、アンテナ帯域、検出閾値等で決まる。高周波で分解を良くすると到達深度が制限されるという交換条件はあるが、深度が厳密に1/fへ比例する普遍式ではない。低損失非分散媒質の導電損失近似は\(\alpha\simeq(\sigma/2)\sqrt{\mu/\epsilon}\)で、fだけの反比例を導くものではない。

通常のGPRはMHz以上のパルスで、原典の10MHz程度は分類目安である。10kHzとの三桁差は誤りだが、10MHzに絶対的な物理境界があるわけではない。力学強度を直接観測しない一方、含水・空隙・劣化等の共通要因による間接相関はあり得る。「ほとんど相関なし」という原典はその限定を保持する。数mの空洞や埋設物の用途を、全媒質の到達上限へ変えない。

\section{孔内検層：深度目盛と観測物性を結び付ける}
\label{theory:exploration:borehole}
地表探査だけでは深度の曖昧さが残る。ボーリング孔に沿う連続計測は、その位置の深度と物性応答を結び付ける。ただし孔周辺の有限体積を見ており、孔壁・泥水侵入・ケーシング・装置結合の補正がある。検層記録の深度が分かることと、地盤全体の構造が完全に分かることは別である。

速度検層のdownholeは地表起振・孔内受振、upholeは孔内起振・地表受振である。走時から速度を求めるには距離と伝播経路、屈折・傾斜の条件を扱う。孔壁へ密着させるdownhole受振器なら水位のない部分でも測定できる\cite{geebec_ps}。水位以深だけという限定は誤りだが、孔内状態によらず精度が保証される意味ではない。

suspensionは孔内振源と複数受振器を持つゾンデで、受振点間の時間差から区間速度を得る\cite{oyo_suspension}。地盤のP/S応答を測り、水中をS波が伝わるという説明にはしない。VSPでは通常の地表振源・孔内受振と逆配置等を区別し、反射記録と孔内深度・地層を対比する。初動速度だけの検層と反射の解析目的は同じではない。

normal電気検層は電極の注入・電位差による比抵抗計測である。導電水・泥水を通じて地層へ電流を流す通常の裸孔条件を考える。孔内水の有無・抵抗率を無視しない。一方、コイルの誘導検層には別の接触条件があり、電極型の水位条件を全検層へ拡張しない。P波速度は弾性率と密度に関わり、含水・飽和で変わり得るが、含水比だけで全地盤の増減が決まるわけではない。

\subsection*{ガンマ密度：粒子密度でなくかさ密度}
\label{theory:exploration:gammadensity}
通常の密度検層はガンマ線を放射し、主にCompton散乱を介する応答から電子密度を測る。一般的岩石の組成では校正によりかさ密度へ結び付く。ベータ線やRayleigh散乱へ置き換えない。飽和二成分モデルなら
\[
 \rho_b=(1-\phi)\rho_m+\phi\rho_f,\qquad
 \phi=\frac{\rho_m-\rho_b}{\rho_m-\rho_f}
\]
である。\(\rho_m\)は粒子・骨格、\(\rho_f\)は孔隙流体の密度であり、これらが既知という条件から孔隙率を推定する。粒子密度そのものの直接測定ではない。孔壁接触や孔内流体、線源・検出器校正が必要である\cite{epa_gamma_density}。

\section{時間差・幾何・熱を使う補足調査}
\label{theory:exploration:multimethod}
time lapseは同じ地点・方式の反復測定を比較し、地下水移動や物性変化を推定する。差分には機器ドリフトや測線ずれ・温度差も入り得る。tomographyの覆いが不足すれば、特定方向の構造や薄い領域を分解しにくい。重力では同じ形状深度なら密度差の大きさが異常を増すが、密度差だけで異なる対象の検出難易度を順位付けしない。

1m深地温探査は、地表下1mの温度を測り流動水・熱源による平常温度との差を調べる。1mは測温深度であり、水脈の最大深度ではない。水温との温度差が小さい季節には検出しにくく、地況・境界条件を扱う熱モデルと観測を組み合わせる\cite{naro_thermal2008,landslide_thermal2007}。特定地の適期を全国共通の月にしない。

環境振動の振動レベルは人の感覚に応じた周波数補正を施す加速度レベルである。原典のJIS準拠測定を、単なる最大変位測定へ取り違えない\cite{rion_vibration2016}。出題時と現行の規格本文の版差は未照合であり、現在版の番号や性能要件を推測して補わない。

\subsection*{試験での確認順序}
励起源、受信器、直接観測量、推定物性、単位、有限距離・帯域・媒質の条件を順に照合する。今回の29問での論点を全16試験回の頻度の確定値とはしない。SPの深度推定、Archieの表面伝導、GPRの比例表現等の疑義は各問題へ残し、公式正答と科学的な説明の保証範囲を区別する。

\section{地震・弾性波探査}

\subsection{平成23年度 問29}
\label{q:h23-reg-applied-29}

\noindent\textbf{原典：} sources/H23/17\_applied\_science.pdf、PDF物理10頁。

\textbf{公式正答：} 5（single）。

\paragraph{問題の要旨} 物理探査の観測対象と装置を比較し、海上音波探査の受振器を誤って説明した記述を選ぶ。

\Needspace{4\baselineskip}\paragraph{確認した条件・記述（独自要約）}\mbox{}\par

\noindent\textbf{選択方向：} 誤り又は最も不適切・関係の薄いものを1つ\par

\paragraph{問われる核心} 観測量と受振器の対応

\paragraph{必要知識} ハイドロフォンの圧力応答、地上速度計、相対重力計、MT、電磁波の反射。

\paragraph{考え方・正答を選ぶ理由（査読待ち初稿）} 測定量を対照する。重力は加速度、MTは自然電磁応答、レーダは電磁反射である。水中の流体はせん断波を伝えず、海上の受振系では音圧を検出するハイドロフォンを用いる。速度計を防水加工するという第5肢の説明が誤りである。

\Needspace{4\baselineskip}\paragraph{選択肢の検討（査読待ち初稿）}\mbox{}\par

\noindent\textbf{1：適切} 磁気から鉄類の有無と概略位置を探るが、細かな形状や製品種類は判別できない。通常の磁気異常だけでは磁化・形状・深度が非一意であり、細かな製品識別までは確定できないという探査上の限界を述べる。現代の補助情報を用いる識別可能性全般を否定する命題とは解さない。\par

\noindent\textbf{2：適切} MTは自然の地磁気・地電流から深部比抵抗を探る。自然変動の水平電場・磁場の応答から電気抵抗率を推定する方法である。\par

\noindent\textbf{3：適切} 相対重力計は吊った錘のばね伸びを利用する。ばね定数と伸びの変化から重力の差を測る。絶対値を直接定める落下式とは区別する。\par

\noindent\textbf{4：適切} 地中レーダは概ね10 MHz以上の電磁波で極浅部を探る。原典が述べる通常の浅部レーダの周波数帯と原理に対応する。到達深度は周波数だけでなく導電率や減衰にも依存する。\par

\noindent\textbf{5：不適切} ストリーマケーブルには防水加工した速度計を多数入れる。水中音波を測る受振器は通常ハイドロフォンであり、圧力変動を検出する。地上用速度計の防水化と同一ではない。\par

\paragraph{再利用する考え方} 観測量と受振器の対応を、測定量・推定対象・基準面・近似条件に分けて理解する。名称だけで探査法を対応させない。

\paragraph{共通理論} \hyperref[theory:exploration:observables]{第\ref*{theory:exploration:observables}節}, \hyperref[theory:exploration:electromagnetic]{第\ref*{theory:exploration:electromagnetic}節}, \hyperref[theory:gravity:measurement]{第\ref*{theory:gravity:measurement}節}

\par 参考：\cite{usgs_marine_seismic,usgs_mt_metadata}。

\par 難易度 3（校準初稿）、優先度未評価。

\section{重力・磁気探査}

\subsection{平成27年度 問22}
\label{q:h27-reg-applied-22}

\noindent\textbf{原典：} sources/H27/17\_applied\_science.pdf、PDF物理8頁。

\textbf{公式正答：} 2（single）。

\paragraph{問題の要旨} 地形、無限水平層、潮汐、ばねドリフト、緯度の重力効果から誤りを選ぶ。

\Needspace{4\baselineskip}\paragraph{確認した条件・記述（独自要約）}\mbox{}\par

\noindent\textbf{選択方向：} 誤り又は最も不適切・関係の薄いものを1つ\par

\paragraph{問われる核心} 無限平板の重力と重力補正の基準

\paragraph{必要知識} ニュートン引力積分、無限平板と有限岩体の違い、地形補正の符号、潮汐、ドリフト、正規重力。

\paragraph{考え方・正答を選ぶ理由（査読待ち初稿）} 無限水平層の上に観測点がある場合、各薄層の寄与を積分すると\(g_z=2\pi G\rho t\)である。層と測点の距離は消えるので、浅いほど大きいという第2肢が誤りである。第1肢は無限平板に対する地形差、第4肢は短時間の線形ドリフト近似として条件を明示して読む。

\Needspace{4\baselineskip}\paragraph{選択肢の検討（査読待ち初稿）}\mbox{}\par

\noindent\textbf{1：適切} 山・谷の地形効果はともに重力を小さくする。測点高まで無限平板があるというブーゲー基準に対する地形効果として読む。測点より上の山の引力は上向き、谷は基準平板の下向き引力を欠くため、いずれも下向き重力を減らす。任意の山と谷の絶対重力を比較する命題ではない。\par

\noindent\textbf{2：不適切} 密度と厚さ一定の無限水平層は浅いほど重力が大きい。無限平板の外部重力は2πGρtで深度に依存しない。有限の岩体では距離依存があるため、無限という条件が重要である。\par

\noindent\textbf{3：適切} 月と太陽の引力による潮汐補正が必要である。月・太陽による時間変動を地下構造による異常と区別する必要がある。\par

\noindent\textbf{4：適切} ばねの時間変化を一定と近似し閉合誤差を各点へ配分する。同一点への戻り測定の閉合差から時間に比例するドリフトを近似して補正する操作を述べる。一定の伸びは測定区間のドリフトを一定率とする通常の近似として解釈し、あらゆる時間変化が必ず線形とはしない。\par

\noindent\textbf{5：適切} 絶対重力と緯度から求めた正規重力との差をとるのが緯度補正である。緯度に依存する理論基準を差し引き、基準地球の変化を除く操作である。\par

\paragraph{再利用する考え方} 無限平板の重力と重力補正の基準を、測定量・推定対象・基準面・近似条件に分けて理解する。名称だけで探査法を対応させない。

\paragraph{共通理論} \hyperref[theory:gravity:slab]{第\ref*{theory:gravity:slab}節}, \hyperref[theory:gravity:corrections]{第\ref*{theory:gravity:corrections}節}, \hyperref[theory:gravity:measurement]{第\ref*{theory:gravity:measurement}節}

\par 参考：\cite{gsi_gravity_definitions}。

\paragraph{関連問題} \hyperref[q:r07-reg-applied-33]{R07-REG-APPLIED-33}

\par 難易度 3（校準初稿）、優先度未評価。

\subsection{平成29年度 問31}
\label{q:h29-reg-applied-31}

\noindent\textbf{原典：} sources/H29/17\_applied\_science.pdf、PDF物理10頁。

\textbf{公式正答：} 5（single）。

\paragraph{問題の要旨} 重力探査で用いる補正名のうち、主に反射法地震探査に属するものを選ぶ。

\Needspace{4\baselineskip}\paragraph{確認した条件・記述（独自要約）}\mbox{}\par

\noindent\textbf{選択方向：} 誤り又は最も不適切・関係の薄いものを1つ\par

\paragraph{問われる核心} 重力補正と地震走時補正の区別

\paragraph{必要知識} 運動・質量・潮汐・標高の重力効果、浅部層による走時差。

\paragraph{考え方・正答を選ぶ理由（査読待ち初稿）} 補正の対象量を並べると、第1〜4肢は重力加速度の変化を扱う。静補正は主に地震探査の走時をそろえるもので、重力探査の標準的補正名としては該当しない。

\Needspace{4\baselineskip}\paragraph{選択肢の検討（査読待ち初稿）}\mbox{}\par

\noindent\textbf{1：適切} エトベス補正。船や航空機の運動に伴う回転系の見掛加速度の差を補正する。\par

\noindent\textbf{2：適切} ブーゲー補正。基準面と測点の間の岩石質量の引力を除く。\par

\noindent\textbf{3：適切} 潮汐補正。月・太陽による時間変動を除く。\par

\noindent\textbf{4：適切} フリーエア補正。基準面からの高さに伴う重力変化を補正する。\par

\noindent\textbf{5：不適切} 静補正。浅部低速度層や標高による地震波の走時差を、共通の基準面へ補正する地震探査の用語である。\par

\paragraph{再利用する考え方} 重力補正と地震走時補正の区別を、測定量・推定対象・基準面・近似条件に分けて理解する。名称だけで探査法を対応させない。

\paragraph{共通理論} \hyperref[theory:gravity:corrections]{第\ref*{theory:gravity:corrections}節}, \hyperref[theory:exploration:observables]{第\ref*{theory:exploration:observables}節}

\par 参考：\cite{gsi_gravity_definitions}。

\par 難易度 3（校準初稿）、優先度未評価。

\subsection{平成30年度 問32}
\label{q:h30-reg-applied-32}

\noindent\textbf{原典：} sources/H30/17\_applied\_science.pdf、PDF物理11頁。

\textbf{公式正答：} 3（single）。

\paragraph{問題の要旨} 重力・磁気探査の基準面、磁化特性、移動観測補正と、別の観測分野の語を区別する。

\Needspace{4\baselineskip}\paragraph{確認した条件・記述（独自要約）}\mbox{}\par

\noindent\textbf{選択方向：} 誤り又は最も不適切・関係の薄いものを1つ\par

\paragraph{問われる核心} 重力・磁気探査の物理量と用語

\paragraph{必要知識} 等ポテンシャル面、残留磁化と誘導磁化、磁気秩序の温度依存、回転座標系。

\paragraph{考え方・正答を選ぶ理由（査読待ち初稿）} ジオイドとエトベス補正は重力、ケーニヒスベルガー比とキュリー温度は磁気に直接関係する。プランク関数は熱平衡放射のスペクトルを記述する関数であり、この候補群では重力・磁気探査の基本量に当たらない。したがってこの肢を選ぶ。

\Needspace{4\baselineskip}\paragraph{選択肢の検討（査読待ち初稿）}\mbox{}\par

\noindent\textbf{1：適切} ジオイド。基準となる重力等ポテンシャル面であり、重力場と標高の基準に関係する。\par

\noindent\textbf{2：適切} ケーニヒスベルガー比。自然残留磁化の大きさと現在の場による誘導磁化の大きさの比であり、磁気異常の解釈に関係する。\par

\noindent\textbf{3：不適切} プランク関数。熱平衡放射のスペクトルを記述する関数であり、この候補群では重力・磁気探査の基本量に当たらない。\par

\noindent\textbf{4：適切} キュリー温度。強磁性・フェリ磁性の自発磁化が消失する温度に対応し、磁性鉱物や磁気探査の解釈に関係する。\par

\noindent\textbf{5：適切} エトベス補正。移動観測系の運動に伴う見掛重力の差を補正する。\par

\paragraph{再利用する考え方} 重力・磁気探査の物理量と用語を、測定量・推定対象・基準面・近似条件に分けて理解する。名称だけで探査法を対応させない。

\paragraph{共通理論} \hyperref[theory:gravity:corrections]{第\ref*{theory:gravity:corrections}節}, \hyperref[theory:gravity:geoid]{第\ref*{theory:gravity:geoid}節}, \hyperref[theory:geomagnetism:magnetization]{第\ref*{theory:geomagnetism:magnetization}節}, \hyperref[theory:exploration:observables]{第\ref*{theory:exploration:observables}節}

\par 参考：\cite{gsi_gravity_definitions,kanaya_okuma2011}。

\paragraph{関連問題} \hyperref[q:r04-reg-applied-35]{R04-REG-APPLIED-35}

\par 難易度 3（校準初稿）、優先度未評価。

\subsection{令和4年度 問35}
\label{q:r04-reg-applied-35}

\noindent\textbf{原典：} sources/R04/17\_applied\_science.pdf、PDF物理12頁。

\textbf{公式正答：} 4（single）。

\paragraph{問題の要旨} 重力・磁気探査の基準面、磁化特性、移動観測補正と、別の観測分野の語を区別する。

\Needspace{4\baselineskip}\paragraph{確認した条件・記述（独自要約）}\mbox{}\par

\noindent\textbf{選択方向：} 誤り又は最も不適切・関係の薄いものを1つ\par

\paragraph{問われる核心} 重力・磁気探査の物理量と用語

\paragraph{必要知識} 等ポテンシャル面、残留磁化と誘導磁化、磁気秩序の温度依存、回転座標系。

\paragraph{考え方・正答を選ぶ理由（査読待ち初稿）} ジオイドとエトベス補正は重力、ケーニヒスベルガー比とキュリー温度は磁気に直接関係する。プランク関数は熱平衡放射のスペクトルを記述する関数であり、この候補群では重力・磁気探査の基本量に当たらない。したがってこの肢を選ぶ。

\Needspace{4\baselineskip}\paragraph{選択肢の検討（査読待ち初稿）}\mbox{}\par

\noindent\textbf{1：適切} エトベス補正。移動観測系の運動に伴う見掛重力の差を補正する。\par

\noindent\textbf{2：適切} ジオイド。基準となる重力等ポテンシャル面であり、重力場と標高の基準に関係する。\par

\noindent\textbf{3：適切} キュリー温度。強磁性・フェリ磁性の自発磁化が消失する温度に対応し、磁性鉱物や磁気探査の解釈に関係する。\par

\noindent\textbf{4：不適切} プランク関数。熱平衡放射のスペクトルを記述する関数であり、この候補群では重力・磁気探査の基本量に当たらない。\par

\noindent\textbf{5：適切} ケーニヒスベルガー比。自然残留磁化の大きさと現在の場による誘導磁化の大きさの比であり、磁気異常の解釈に関係する。\par

\paragraph{再利用する考え方} 重力・磁気探査の物理量と用語を、測定量・推定対象・基準面・近似条件に分けて理解する。名称だけで探査法を対応させない。

\paragraph{共通理論} \hyperref[theory:gravity:corrections]{第\ref*{theory:gravity:corrections}節}, \hyperref[theory:gravity:geoid]{第\ref*{theory:gravity:geoid}節}, \hyperref[theory:geomagnetism:magnetization]{第\ref*{theory:geomagnetism:magnetization}節}, \hyperref[theory:exploration:observables]{第\ref*{theory:exploration:observables}節}

\par 参考：\cite{gsi_gravity_definitions,kanaya_okuma2011}。

\paragraph{関連問題} \hyperref[q:h30-reg-applied-32]{H30-REG-APPLIED-32}, \hyperref[q:r06-reg-applied-34]{R06-REG-APPLIED-34}

\par 難易度 3（校準初稿）、優先度未評価。

\subsection{令和6年度 問34}
\label{q:r06-reg-applied-34}

\noindent\textbf{原典：} sources/R06/17\_applied\_science.pdf、PDF物理12頁。

\textbf{公式正答：} 3（single）。

\paragraph{問題の要旨} 重力・磁気探査の基準面、磁化特性、移動観測補正と、別の観測分野の語を区別する。

\Needspace{4\baselineskip}\paragraph{確認した条件・記述（独自要約）}\mbox{}\par

\noindent\textbf{選択方向：} 誤り又は最も不適切・関係の薄いものを1つ\par

\paragraph{問われる核心} 重力・磁気探査の物理量と用語

\paragraph{必要知識} 等ポテンシャル面、残留磁化と誘導磁化、磁気秩序の温度依存、回転座標系。

\paragraph{考え方・正答を選ぶ理由（査読待ち初稿）} ジオイドとエトベス補正は重力、ケーニヒスベルガー比とキュリー温度は磁気に直接関係する。AE（アコースティック・エミッション）は微小破壊等に伴う過渡的な弾性波の放出・計測である。岩石研究にも用いるが、重力や磁化の測定とは異なる物理量を扱う。したがってこの肢を選ぶ。

\Needspace{4\baselineskip}\paragraph{選択肢の検討（査読待ち初稿）}\mbox{}\par

\noindent\textbf{1：適切} ケーニヒスベルガー比。自然残留磁化の大きさと現在の場による誘導磁化の大きさの比であり、磁気異常の解釈に関係する。\par

\noindent\textbf{2：適切} ジオイド。基準となる重力等ポテンシャル面であり、重力場と標高の基準に関係する。\par

\noindent\textbf{3：不適切} AE（アコースティック・エミッション）。微小破壊等に伴う過渡的な弾性波の放出・計測である。岩石研究にも用いるが、重力や磁化の測定とは異なる物理量を扱う。\par

\noindent\textbf{4：適切} キュリー温度。強磁性・フェリ磁性の自発磁化が消失する温度に対応し、磁性鉱物や磁気探査の解釈に関係する。\par

\noindent\textbf{5：適切} エトベス補正。移動観測系の運動に伴う見掛重力の差を補正する。\par

\paragraph{再利用する考え方} 重力・磁気探査の物理量と用語を、測定量・推定対象・基準面・近似条件に分けて理解する。名称だけで探査法を対応させない。

\paragraph{共通理論} \hyperref[theory:gravity:corrections]{第\ref*{theory:gravity:corrections}節}, \hyperref[theory:gravity:geoid]{第\ref*{theory:gravity:geoid}節}, \hyperref[theory:geomagnetism:magnetization]{第\ref*{theory:geomagnetism:magnetization}節}, \hyperref[theory:exploration:observables]{第\ref*{theory:exploration:observables}節}

\par 参考：\cite{gsi_gravity_definitions,kanaya_okuma2011,sato_lei2009_ae}。

\paragraph{関連問題} \hyperref[q:r04-reg-applied-35]{R04-REG-APPLIED-35}

\par 難易度 3（校準初稿）、優先度未評価。

\subsection{令和7年度 問33}
\label{q:r07-reg-applied-33}

\noindent\textbf{原典：} sources/R07/17\_applied\_science.pdf、PDF物理12頁。

\textbf{公式正答：} 5（single）。

\paragraph{問題の要旨} 正規重力と各補正を区別し、密度差によるブーゲー異常の符号を誤る記述を選ぶ。

\Needspace{4\baselineskip}\paragraph{確認した条件・記述（独自要約）}\mbox{}\par

\noindent\textbf{選択方向：} 誤り又は最も不適切・関係の薄いものを1つ\par

\paragraph{問われる核心} 基準密度と重力異常の符号

\paragraph{必要知識} 正規重力、フリーエア・ブーゲー・エトベス補正、密度差の引力積分。

\paragraph{考え方・正答を選ぶ理由（査読待ち初稿）} 測定重力から正規場・高さ・地表質量等の効果を取り除き、地下の密度差を評価する。同じ幾何条件なら重力の差は密度差\(\Delta\rho\)に線形であり、高密度は正の、低密度は負の異常を生む。第5肢はこの符号を逆に記述する。

\Needspace{4\baselineskip}\paragraph{選択肢の検討（査読待ち初稿）}\mbox{}\par

\noindent\textbf{1：適切} 正規重力は地球の質量・自転・形状から導く理論値。基準回転楕円体の外部重力ポテンシャルと自転による遠心ポテンシャルから定める。\par

\noindent\textbf{2：適切} ブーゲー補正は測点と海水準面の間の地殻物質の効果を除く。標高差だけでなく、その間の岩石質量の引力を平板等で近似して除く。\par

\noindent\textbf{3：適切} 船・航空機の移動は遠心効果を変える。地球自転と観測機の相対運動により見掛加速度が変わるためエトベス補正が必要である。\par

\noindent\textbf{4：適切} フリーエア補正は高さの違いの影響を補正する。高さによる重力の鉛直勾配の影響をそろえる操作である。\par

\noindent\textbf{5：不適切} 高密度地域のブーゲー異常は小さく、低密度地域では大きくなる。形状等を同条件にした基準密度との差では、密度の過剰は正、欠損は負の重力異常を生む。第5肢は基本的対応を逆にしている。\par

\paragraph{再利用する考え方} 基準密度と重力異常の符号を、測定量・推定対象・基準面・近似条件に分けて理解する。名称だけで探査法を対応させない。

\paragraph{共通理論} \hyperref[theory:gravity:slab]{第\ref*{theory:gravity:slab}節}, \hyperref[theory:gravity:corrections]{第\ref*{theory:gravity:corrections}節}

\par 参考：\cite{gsi_gravity_definitions}。

\paragraph{関連問題} \hyperref[q:h27-reg-applied-22]{H27-REG-APPLIED-22}

\par 難易度 3（校準初稿）、優先度未評価。

\section{電気・電磁探査}

\subsection{平成27年度 問19}
\label{q:h27-reg-applied-19}

\noindent\textbf{原典：} sources/H27/17\_applied\_science.pdf、PDF物理7頁。

\textbf{公式正答：} 3（single）。

\paragraph{問題の要旨} 放射能・地温・電磁・レーダ・重力探査の推定対象を比べ、CSAMTの物性を誤る記述を選ぶ。

\Needspace{4\baselineskip}\paragraph{確認した条件・記述（独自要約）}\mbox{}\par

\noindent\textbf{選択方向：} 誤り又は最も不適切・関係の薄いものを1つ\par

\paragraph{問われる核心} 電磁応答と比抵抗の対応

\paragraph{必要知識} 比抵抗と導電率、平面波インピーダンス、CSAMT遠方場、ガンマ線・温度・重力の観測。

\paragraph{考え方・正答を選ぶ理由（査読待ち初稿）} CSAMTは人工源の遠方場を用いる電磁探査であり、電磁応答から比抵抗構造を推定する。均質・平面波近似なら見掛比抵抗は\(\rho_a=|E/H|^2/(\omega\mu)\)である。磁化率を主対象とする第3肢は物性の対応を誤っている。

\Needspace{4\baselineskip}\paragraph{選択肢の検討（査読待ち初稿）}\mbox{}\par

\noindent\textbf{1：適切} 放射能探査は主にガンマ線強度から断層等の地質情報を得る。天然放射性核種の分布や移動の差を地質情報として利用する。ガンマ線の減衰により観測は主に表層からの寄与であり、深部核種をそのまま透視する意味ではない。\par

\noindent\textbf{2：適切} 地温探査は温泉・地熱・地下水等の調査に用いる。温度分布と熱の移流・伝導を手掛かりに熱水・地下水等を探る。\par

\noindent\textbf{3：不適切} CSAMTは遠方人工源の直交水平電場・磁場から磁化率を得る。主要な推定対象は比抵抗である。水平電場と磁場の比を磁化率そのものと取り違えている。\par

\noindent\textbf{4：適切} レーダは誘電率や導電率の違う境界の電磁反射を捉える。電気特性の差が波の伝播と反射を変えることを利用する。\par

\noindent\textbf{5：適切} 重力は密度構造を求め、標高精度が重要である。標高1 mの誤差は標準鉛直勾配で約0.3086 mGalに相当し、微小重力異常の解釈を妨げる。\par

\paragraph{再利用する考え方} 電磁応答と比抵抗の対応を、測定量・推定対象・基準面・近似条件に分けて理解する。名称だけで探査法を対応させない。

\paragraph{共通理論} \hyperref[theory:exploration:electromagnetic]{第\ref*{theory:exploration:electromagnetic}節}, \hyperref[theory:exploration:observables]{第\ref*{theory:exploration:observables}節}

\par 参考：\cite{ochi_csamt_saijo,gsi_gravity_definitions}。

\par 難易度 3（校準初稿）、優先度未評価。

\section{孔隙水導電とArchie式の適用条件}

\subsection{平成27年度 問18}
\label{q:h27-reg-applied-18}

\noindent\textbf{原典：} sources/H27/17\_applied\_science.pdf、PDF物理6頁。

\textbf{公式正答：} 4（single）。

\par\textbf{検討未確定：} 粘土と砂、割れ目と健岩の比較は含水・塩分・鉱物条件に依存する。風化の全称則とはしない。

\paragraph{問題の要旨} 粘土・不飽和・割れ目・海水が比抵抗へ与える傾向を比較する。

\Needspace{4\baselineskip}\paragraph{確認した条件・記述（独自要約）}\mbox{}\par

\noindent\textbf{原典照合：} 29問21原典画像で否定、全五候補、数式・上付き単位・観測対象を照合。\par

\noindent\textbf{前提：} 孔隙水のイオン伝導と粘土表面伝導、含水した割れ目を考える。\par

\paragraph{問われる核心} 粘土・不飽和・割れ目・海水が比抵抗へ与える傾向を比較する。

\paragraph{必要知識} 孔隙水のイオン伝導と粘土表面伝導、含水した割れ目を考える。

\paragraph{考え方・正答を選ぶ理由（査読待ち初稿）} 4が不適切である。海水のイオン濃度は高く導電率を増すため、孔内・地層水へ海水が入れば通常は抵抗率が低下する。孔内流体の導電率は電極検層の測定にも影響する。

\Needspace{4\baselineskip}\paragraph{選択肢の検討（査読待ち初稿）}\mbox{}\par

\noindent\textbf{1：条件付きで適切} 粘性土は砂質土より一般に低比抵抗。粘土表面伝導と含水が低抵抗へ寄与する。同じ流体・飽和等の比較が必要である。\par

\noindent\textbf{2：条件付きで適切} 不飽和部分は一般に高比抵抗。水を通る連続導電経路が減り通常は抵抗率が上がる。鉱物・流体条件も影響する。\par

\noindent\textbf{3：条件付きで適切} 割れ目と風化が進んだ岩盤は低比抵抗。含水割れ目や粘土化がある場合の傾向である。乾燥空隙や珪化を同じ低下としない。\par

\noindent\textbf{4：不適切} 海水混入で比抵抗が著しく高くなる。イオンを含む海水は通常導電性が高く、淡水状態より抵抗率を下げる方向である。\par

\noindent\textbf{5：適切} 電気検層は裸孔が基本。通常の電極法では地層への電流経路を確保しケーシングによる影響を避ける。\par

\paragraph{再利用する考え方} 低比抵抗を岩石名だけに対応させず、水・塩分・粘土を分ける。

\paragraph{共通理論} \hyperref[theory:exploration:archie]{第\ref*{theory:exploration:archie}節}, \hyperref[theory:exploration:borehole]{第\ref*{theory:exploration:borehole}節}

\par 参考：\cite{epa_resistivity,epa_borehole_resistivity}。

\par 難易度 3（校準初稿）、優先度未評価。

\subsection{平成29年度 問32}
\label{q:h29-reg-applied-32}

\noindent\textbf{原典：} sources/H29/17\_applied\_science.pdf、PDF物理11頁。

\textbf{公式正答：} 4（single）。

\par\textbf{検討未確定：} 高比抵抗水でも絶縁清浄骨格でArchieは成立し得る。礫層・割れ目・地熱SPの符号の全称は留保する。

\paragraph{問題の要旨} Archieの孔隙水伝導条件と地下水・SP解釈を比較する。

\Needspace{4\baselineskip}\paragraph{確認した条件・記述（独自要約）}\mbox{}\par

\noindent\textbf{原典照合：} 29問21原典画像で否定、全五候補、数式・上付き単位・観測対象を照合。\par

\noindent\textbf{前提：} 孔隙水だけで電流が流れ、骨格と表面伝導を無視できる場合にArchieを使う。\par

\paragraph{問われる核心} Archieの孔隙水伝導条件と地下水・SP解釈を比較する。

\paragraph{必要知識} 孔隙水だけで電流が流れ、骨格と表面伝導を無視できる場合にArchieを使う。

\paragraph{考え方・正答を選ぶ理由（査読待ち初稿）} 4の高比抵抗孔隙水なら精度よく適用という保証が不適切である。孔隙水導電が小さいと相対的に表面伝導等が重要となり、仮定が破れやすい。清浄な絶縁骨格なら淡水でも成立し得るため、高比抵抗では絶対に不可とはしない。

\Needspace{4\baselineskip}\paragraph{選択肢の検討（査読待ち初稿）}\mbox{}\par

\noindent\textbf{1：適切} 囲む電極の比抵抗tomographyで二次元・三次元分布。多数配置電位差を逆解析し、覆いとモデル制約の範囲で分布を求める。\par

\noindent\textbf{2：適切} CSAMT近接場は小間隔・高比抵抗で生じやすい。同一周波数でδ∝√ρとなるので、距離/δが小さくなる方向である。\par

\noindent\textbf{3：条件付きで適切} 平野礫帯水層は高抵抗、山地裂帯変質帯は低抵抗。粘土層との比較や含水割れ目・変質粘土の対比として典型的である。塩分、空隙率、乾燥、母岩で変わる。\par

\noindent\textbf{4：不適切} 孔隙水比抵抗が高い場合にArchieを精度よく適用。孔隙水以外の導電を無視できることが必要で、高抵抗水だけで精度は保証できない。表面伝導の相対的重要性が増す。\par

\noindent\textbf{5：条件付きで適切} 鉱体頂部は負異常、地熱対流は正異常。典型鉱化機構と上昇流の特定条件の傾向である。熱電・化学・流動電位と基準点に依存し、全地熱地点の正符号を保証しない。\par

\paragraph{再利用する考え方} 実験式の対象と地質解釈の典型を、全材料に通じる十分条件へ拡張しない。

\paragraph{共通理論} \hyperref[theory:exploration:archie]{第\ref*{theory:exploration:archie}節}, \hyperref[theory:exploration:nearfield]{第\ref*{theory:exploration:nearfield}節}, \hyperref[theory:exploration:ipsp]{第\ref*{theory:exploration:ipsp}節}

\par 参考：\cite{epa_borehole_resistivity,epa_resistivity,epa_sp,usgs_csamt2012}。

\paragraph{関連問題} \hyperref[q:h26-reg-applied-22]{H26-REG-APPLIED-22}

\par 難易度 3（校準初稿）、優先度未評価。

\section{微動アレイの長短波長限界}

\subsection{令和4年度 問34}
\label{q:r04-reg-applied-34}

\noindent\textbf{原典：} sources/R04/17\_applied\_science.pdf、PDF物理12頁。

\textbf{公式正答：} 1（single）。

\par\textbf{検討未確定：} F-K対SPACの大小は解析・入射条件と実装に依存する。水平多層モデルは通常の一次元逆解析の前提であり、手法全般の絶対定義とはしない。

\paragraph{問題の要旨} 微動アレイの短波長限界、水平層モデル、F-K/SPAC、S波逆解析を比較する。

\Needspace{4\baselineskip}\paragraph{確認した条件・記述（独自要約）}\mbox{}\par

\noindent\textbf{原典照合：} 五候補の順序、最適切と不適切の方向、波種、深度目安、曲線名、最長と最短を原典画像で照合。\par

\noindent\textbf{前提：} 短波長側の空間aliasingと長波長側の開口・位相差・統計精度を分ける。\par

\paragraph{問われる核心} 微動アレイの短波長限界、水平層モデル、F-K/SPAC、S波逆解析を比較する。

\paragraph{必要知識} 短波長側の空間aliasingと長波長側の開口・位相差・統計精度を分ける。

\paragraph{考え方・正答を選ぶ理由（査読待ち初稿）} 1が不適切である。空間aliasingは主に短波長・高波数側の限界であり最長波長ではない。長波長側はアレイ内の位相差が小さくなることや有限開口・解析手法の精度等に制約される。

\Needspace{4\baselineskip}\paragraph{選択肢の検討（査読待ち初稿）}\mbox{}\par

\noindent\textbf{1：不適切} 最長波長がアレイ固有の空間aliasingで決まる。最長と最短の役割が逆である。短波長を粗い測点で標本化すると空間波数が折り返して区別できなくなる。\par

\noindent\textbf{2：条件付きで適切} 水平多層構造を仮定して表面波から地下構造を推定する。通常の一次元分散逆解析のモデルである。横方向変化が強い場所や他の多次元算法まで常に水平層であるとはいえない。\par

\noindent\textbf{3：適切} 人工地表加振の表面波で一次元地下モデルを推定する。通常の浅層能動表面波探査の説明であり、周波数・波長ごとの速度を地盤の深度構造に結ぶ。\par

\noindent\textbf{4：条件付きで適切} 同じアレイサイズでF-Kの最大波長はSPACより短い。通常の解析性能の比較であり、実験論文の条件ではSPACがより長波長まで推定した。配置・波の入射・ノイズ・推定算法に依存するため全ての実装に無条件の大小関係を課さない。\par

\noindent\textbf{5：適切} 観測位相速度に近い理論値を持つS波速度構造を作る。順モデルによる分散曲線と観測を照合し主にS波速度を更新する。P波速度・密度・層厚の前提も必要であり唯一解の保証ではない。\par

\paragraph{再利用する考え方} 周波数ではなく波数の標本化として空間aliasingを理解する。

\paragraph{共通理論} \hyperref[theory:exploration:surfacewaves]{第\ref*{theory:exploration:surfacewaves}節}, \hyperref[theory:exploration:arraylimits]{第\ref*{theory:exploration:arraylimits}節}

\par 参考：\cite{epa_surfacewaves,xia_rayleigh1999,japansea_spac2017,array_limits2002}。

\paragraph{関連問題} \hyperref[q:r06-reg-applied-33]{R06-REG-APPLIED-33}

\par 難易度 3（校準初稿）、優先度未評価。

\subsection{令和6年度 問33}
\label{q:r06-reg-applied-33}

\noindent\textbf{原典：} sources/R06/17\_applied\_science.pdf、PDF物理11頁。

\textbf{公式正答：} 2（single）。

\par\textbf{検討未確定：} P波速度・密度・層厚の影響もゼロでない。F-K/SPACの最大波長比較と水平層モデルは通常の解析条件の範囲で読む。

\paragraph{問題の要旨} F-K/SPAC、S波逆解析、水平層、能動起振、最短波長を比較する。

\Needspace{4\baselineskip}\paragraph{確認した条件・記述（独自要約）}\mbox{}\par

\noindent\textbf{原典照合：} 五候補の順序、最適切と不適切の方向、波種、深度目安、曲線名、最長と最短を原典画像で照合。\par

\noindent\textbf{前提：} 微動のRayleigh位相速度は複数の物性へ依存するが、主にS波速度構造に強い感度を持つ。\par

\paragraph{問われる核心} F-K/SPAC、S波逆解析、水平層、能動起振、最短波長を比較する。

\paragraph{必要知識} 微動のRayleigh位相速度は複数の物性へ依存するが、主にS波速度構造に強い感度を持つ。

\paragraph{考え方・正答を選ぶ理由（査読待ち初稿）} 2が不適切である。通常の微動分散逆解析で作るのは主にS波速度モデルであり、P波モデルへ置換する説明が誤りである。P波速度が位相速度へ影響しないという意味ではない。

\Needspace{4\baselineskip}\paragraph{選択肢の検討（査読待ち初稿）}\mbox{}\par

\noindent\textbf{1：条件付きで適切} 同じアレイサイズでF-Kの最大波長はSPACより短い。通常の解析性能の比較であり、実験論文の条件ではSPACがより長波長まで推定した。配置・波の入射・ノイズ・推定算法に依存するため全ての実装に無条件の大小関係を課さない。\par

\noindent\textbf{2：不適切} 観測位相速度に近いP波速度構造モデルを作る。対象とする主な未知量をP波へ取り違える。分散はP波にも依存するが、通常はS波構造を推定し他の物性を別条件で拘束する。\par

\noindent\textbf{3：条件付きで適切} 水平多層構造を仮定して表面波から地下構造を推定する。通常の一次元分散逆解析のモデルである。横方向変化が強い場所や他の多次元算法まで常に水平層であるとはいえない。\par

\noindent\textbf{4：適切} 人工地表加振の表面波で一次元地下モデルを推定する。通常の浅層能動表面波探査の説明であり、周波数・波長ごとの速度を地盤の深度構造に結ぶ。\par

\noindent\textbf{5：適切} 推定位相速度の最短波長は空間aliasingに制約される。波長に対し測点が粗いと異なる空間波数が区別できなくなる。長波長側の位相差の小ささ・有限開口の制約とは区別する。\par

\paragraph{再利用する考え方} 位相速度の感度が強い物性と、影響がゼロではない他物性を区別する。

\paragraph{共通理論} \hyperref[theory:exploration:surfacewaves]{第\ref*{theory:exploration:surfacewaves}節}, \hyperref[theory:exploration:arraylimits]{第\ref*{theory:exploration:arraylimits}節}

\par 参考：\cite{epa_surfacewaves,xia_rayleigh1999,japansea_spac2017,array_limits2002}。

\paragraph{関連問題} \hyperref[q:r04-reg-applied-34]{R04-REG-APPLIED-34}, \hyperref[q:r07-reg-applied-35]{R07-REG-APPLIED-35}

\par 難易度 3（校準初稿）、優先度未評価。

\section{孔内検層の観測量・水位と接触条件}

\subsection{平成23年度 問31}
\label{q:h23-reg-applied-31}

\noindent\textbf{原典：} sources/H23/17\_applied\_science.pdf、PDF物理10頁。

\textbf{公式正答：} 2（single）。

\par\textbf{検討未確定：} VSPの振源又は受振器は通常型と逆配置を含む。裸孔条件は電極検層の基本条件であり全検層へ拡張しない。

\paragraph{問題の要旨} 孔内速度・電気・密度検層の水位条件を比較する。

\Needspace{4\baselineskip}\paragraph{確認した条件・記述（独自要約）}\mbox{}\par

\noindent\textbf{原典照合：} 29問21原典画像で否定、全五候補、数式・上付き単位・観測対象を照合。\par

\noindent\textbf{前提：} 速度検層は弾性波、電極検層は比抵抗、ガンマ密度検層は校正かさ密度である。\par

\paragraph{問われる核心} 孔内速度・電気・密度検層の水位条件を比較する。

\paragraph{必要知識} 速度検層は弾性波、電極検層は比抵抗、ガンマ密度検層は校正かさ密度である。

\paragraph{考え方・正答を選ぶ理由（査読待ち初稿）} 2が誤りである。孔壁に密着させるダウンホール受振器は孔内水位より上でも測定できる。水による振動結合を使う他方式や、導電流体を要する電極検層と区別する。

\Needspace{4\baselineskip}\paragraph{選択肢の検討（査読待ち初稿）}\mbox{}\par

\noindent\textbf{1：条件付きで適切} VSPで孔内に振源又は受振器を設置し反射波を測る。通常の地表振源・孔内受振と逆VSPなどの構成を区別する。孔内深度と反射記録を結び地表反射探査の解釈を助ける。\par

\noindent\textbf{2：不適切} ダウンホールは孔内水位以深だけに適用可能。孔壁へ固定・密着した受振器なら水位のない部分でも測定できる。孔内水位以深だけという限定が誤りである。\par

\noindent\textbf{3：適切} サスペンションで孔内振源と複数受振器を用いる。孔内ゾンデの起振と受振点間の走時差により、孔軸に沿う地盤の区間速度を求める。水中のS波伝播を直接測るという意味ではない。\par

\noindent\textbf{4：条件付きで適切} 電極検層は孔内と地表の電極で比抵抗を測り裸孔を基本とする。通常のnormal電極検層では導電孔内流体を介して地層へ電流を流す。金属ケーシングの影響を避ける。誘導検層等の別方式へこの条件をそのまま一般化しない。\par

\noindent\textbf{5：適切} ガンマ線散乱で孔壁周辺の密度を測り孔隙率を推定する。主にCompton散乱に関わる電子密度を校正してかさ密度を求める。粒子と流体の密度が既知なら孔隙率を推定できる。\par

\paragraph{再利用する考え方} 流体、孔壁結合、線種、センサ位置を観測量ごとに整理する。

\paragraph{共通理論} \hyperref[theory:exploration:borehole]{第\ref*{theory:exploration:borehole}節}, \hyperref[theory:exploration:gammadensity]{第\ref*{theory:exploration:gammadensity}節}

\par 参考：\cite{epa_borehole_resistivity,epa_gamma_density,geebec_ps,oyo_suspension}。

\paragraph{関連問題} \hyperref[q:r01-reg-applied-33]{R01-REG-APPLIED-33}

\par 難易度 3（校準初稿）、優先度未評価。

\subsection{平成24年度 問22}
\label{q:h24-reg-applied-22}

\noindent\textbf{原典：} sources/H24/17\_applied\_science.pdf、PDF物理7頁。

\textbf{公式正答：} 2（single）。

\paragraph{問題の要旨} 孔内弾性波速度測定とnormal電気検層を区別する。

\Needspace{4\baselineskip}\paragraph{確認した条件・記述（独自要約）}\mbox{}\par

\noindent\textbf{原典照合：} 29問21原典画像で否定、全五候補、数式・上付き単位・観測対象を照合。\par

\noindent\textbf{前提：} 振源・受振器の位置と直接観測する物性を対応させる。\par

\paragraph{問われる核心} 孔内弾性波速度測定とnormal電気検層を区別する。

\paragraph{必要知識} 振源・受振器の位置と直接観測する物性を対応させる。

\paragraph{考え方・正答を選ぶ理由（査読待ち初稿）} 2のノルマル法は電極配置による比抵抗検層である。downhole、uphole、VSP、suspensionは弾性波走時又は波形を測り速度構造に関連する。

\Needspace{4\baselineskip}\paragraph{選択肢の検討（査読待ち初稿）}\mbox{}\par

\noindent\textbf{1：適切} ダウンホール法。地表起振、孔内受振で到達走時から速度を求める。\par

\noindent\textbf{2：不適切} ノルマル法。電極の間隔を定めて孔周辺の見掛比抵抗を測り、弾性波速度を直接測る方法ではない。\par

\noindent\textbf{3：適切} アップホール法。孔内起振、地表受振で弾性波走時を測る。\par

\noindent\textbf{4：適切} VSP法。孔内受振等の波形と深度を用い、速度と反射構造を対応させる。\par

\noindent\textbf{5：適切} サスペンション法。孔内振源と複数受振器の到達時間差から区間速度を求める。\par

\paragraph{再利用する考え方} 名称だけでなくセンサと観測量を思い出す。

\paragraph{共通理論} \hyperref[theory:exploration:borehole]{第\ref*{theory:exploration:borehole}節}

\par 参考：\cite{epa_borehole_resistivity,geebec_ps,oyo_suspension}。

\par 難易度 2（校準初稿）、優先度未評価。

\subsection{令和元年度 問33}
\label{q:r01-reg-applied-33}

\noindent\textbf{原典：} sources/R01/17\_applied\_science.pdf、PDF物理12頁。

\textbf{公式正答：} 3（single）。

\par\textbf{検討未確定：} 水の有無だけでdownhole可否は決まらないが、全孔内条件の非依存性を保証しない。VSPと裸孔は方式を特定する。

\paragraph{問題の要旨} 密度検層の放射線と粒子密度・かさ密度を区別する。

\Needspace{4\baselineskip}\paragraph{確認した条件・記述（独自要約）}\mbox{}\par

\noindent\textbf{原典照合：} 29問21原典画像で否定、全五候補、数式・上付き単位・観測対象を照合。\par

\noindent\textbf{前提：} 速度検層は弾性波、電極検層は比抵抗、ガンマ密度検層は校正かさ密度である。\par

\paragraph{問われる核心} 密度検層の放射線と粒子密度・かさ密度を区別する。

\paragraph{必要知識} 速度検層は弾性波、電極検層は比抵抗、ガンマ密度検層は校正かさ密度である。

\paragraph{考え方・正答を選ぶ理由（査読待ち初稿）} 3が不適切である。密度検層はガンマ線の主にCompton散乱による電子密度から孔周辺のかさ密度を求める。ベータ線でも粒子密度の直接測定でもない。

\Needspace{4\baselineskip}\paragraph{選択肢の検討（査読待ち初稿）}\mbox{}\par

\noindent\textbf{1：条件付きで適切} VSPで孔内に振源又は受振器を設置し反射波を測る。通常の地表振源・孔内受振と逆VSPなどの構成を区別する。孔内深度と反射記録を結び地表反射探査の解釈を助ける。\par

\noindent\textbf{2：条件付きで適切} ダウンホールは孔内水の有無が問題とならない。孔壁結合が確保できれば水位より上も測定できる。ただし孔壁崩壊や受振器の結合条件は測定品質へ影響する。\par

\noindent\textbf{3：不適切} ベータ線の散乱で地層の粒子密度を測る。線種はガンマ線、主観測は電子密度から校正するかさ密度である。二つの観測対象を取り違える。\par

\noindent\textbf{4：条件付きで適切} 電極検層は孔内と地表の電極で比抵抗を測り裸孔を基本とする。通常のnormal電極検層では導電孔内流体を介して地層へ電流を流す。金属ケーシングの影響を避ける。誘導検層等の別方式へこの条件をそのまま一般化しない。\par

\noindent\textbf{5：適切} サスペンションで孔内振源と複数受振器を用いる。孔内ゾンデの起振と受振点間の走時差により、孔軸に沿う地盤の区間速度を求める。水中のS波伝播を直接測るという意味ではない。\par

\paragraph{再利用する考え方} 流体、孔壁結合、線種、センサ位置を観測量ごとに整理する。

\paragraph{共通理論} \hyperref[theory:exploration:borehole]{第\ref*{theory:exploration:borehole}節}, \hyperref[theory:exploration:gammadensity]{第\ref*{theory:exploration:gammadensity}節}

\par 参考：\cite{epa_borehole_resistivity,epa_gamma_density,geebec_ps,oyo_suspension}。

\paragraph{関連問題} \hyperref[q:h23-reg-applied-31]{H23-REG-APPLIED-31}, \hyperref[q:r05-reg-applied-34]{R05-REG-APPLIED-34}

\par 難易度 3（校準初稿）、優先度未評価。

\subsection{令和5年度 問34}
\label{q:r05-reg-applied-34}

\noindent\textbf{原典：} sources/R05/17\_applied\_science.pdf、PDF物理12頁。

\textbf{公式正答：} 4（single）。

\par\textbf{検討未確定：} VSPに通常型と逆配置があり、電極検層と誘導検層の条件を区別する。ガンマ密度の孔壁・孔内流体補正も保持。

\paragraph{問題の要旨} 電極検層の孔内流体と各検層の信号を対応させる。

\Needspace{4\baselineskip}\paragraph{確認した条件・記述（独自要約）}\mbox{}\par

\noindent\textbf{原典照合：} 29問21原典画像で否定、全五候補、数式・上付き単位・観測対象を照合。\par

\noindent\textbf{前提：} 速度検層は弾性波、電極検層は比抵抗、ガンマ密度検層は校正かさ密度である。\par

\paragraph{問われる核心} 電極検層の孔内流体と各検層の信号を対応させる。

\paragraph{必要知識} 速度検層は弾性波、電極検層は比抵抗、ガンマ密度検層は校正かさ密度である。

\paragraph{考え方・正答を選ぶ理由（査読待ち初稿）} 4が不適切である。通常の電極型電気検層は孔内の水又は導電泥水を通じて地層へ電流を流すため、その有無と電気特性が問題になる。ダウンホールの孔壁結合とは別である。

\Needspace{4\baselineskip}\paragraph{選択肢の検討（査読待ち初稿）}\mbox{}\par

\noindent\textbf{1：適切} ダウンホールで地表起振・孔内受振の初動走時を測る。孔壁に結合した受振器の位置と走時から区間速度を推定する。経路と屈折の補正が必要である。\par

\noindent\textbf{2：適切} サスペンションで孔内振源と複数受振器を用いる。孔内ゾンデの起振と受振点間の走時差により、孔軸に沿う地盤の区間速度を求める。水中のS波伝播を直接測るという意味ではない。\par

\noindent\textbf{3：適切} ガンマ線散乱で孔壁周辺の密度を測り孔隙率を推定する。主にCompton散乱に関わる電子密度を校正してかさ密度を求める。粒子と流体の密度が既知なら孔隙率を推定できる。\par

\noindent\textbf{4：不適切} 電極検層で孔内水の有無は問題とならない。電極から地層へ電流を流す導電接触が必要で、孔内流体は電流経路と測定値に影響する。電磁誘導検層との混同である。\par

\noindent\textbf{5：条件付きで適切} VSPで孔内に振源又は受振器を設置し反射波を測る。通常の地表振源・孔内受振と逆VSPなどの構成を区別する。孔内深度と反射記録を結び地表反射探査の解釈を助ける。\par

\paragraph{再利用する考え方} 流体、孔壁結合、線種、センサ位置を観測量ごとに整理する。

\paragraph{共通理論} \hyperref[theory:exploration:borehole]{第\ref*{theory:exploration:borehole}節}, \hyperref[theory:exploration:gammadensity]{第\ref*{theory:exploration:gammadensity}節}

\par 参考：\cite{epa_borehole_resistivity,epa_gamma_density,geebec_ps,oyo_suspension}。

\paragraph{関連問題} \hyperref[q:r01-reg-applied-33]{R01-REG-APPLIED-33}, \hyperref[q:r07-reg-applied-34]{R07-REG-APPLIED-34}

\par 難易度 3（校準初稿）、優先度未評価。

\subsection{令和7年度 問34}
\label{q:r07-reg-applied-34}

\noindent\textbf{原典：} sources/R07/17\_applied\_science.pdf、PDF物理13頁。

\textbf{公式正答：} 4（single）。

\par\textbf{検討未確定：} 水位条件は通常の電極検層についてである。VSPの逆配置や孔内流体・孔壁接触による密度補正を別に扱う。

\paragraph{問題の要旨} ガンマ密度検層のCompton散乱と誤った線種を比較する。

\Needspace{4\baselineskip}\paragraph{確認した条件・記述（独自要約）}\mbox{}\par

\noindent\textbf{原典照合：} 29問21原典画像で否定、全五候補、数式・上付き単位・観測対象を照合。\par

\noindent\textbf{前提：} 速度検層は弾性波、電極検層は比抵抗、ガンマ密度検層は校正かさ密度である。\par

\paragraph{問われる核心} ガンマ密度検層のCompton散乱と誤った線種を比較する。

\paragraph{必要知識} 速度検層は弾性波、電極検層は比抵抗、ガンマ密度検層は校正かさ密度である。

\paragraph{考え方・正答を選ぶ理由（査読待ち初稿）} 4が不適切である。通常の密度検層はガンマ線のCompton散乱を利用して校正したかさ密度を求める。ベータ線のRayleigh散乱とする原典記述を正しい文章へ改変しない。

\Needspace{4\baselineskip}\paragraph{選択肢の検討（査読待ち初稿）}\mbox{}\par

\noindent\textbf{1：条件付きで適切} 電極検層の測定区間は孔内水位以深。通常の裸孔電極検層では導電水・泥水を介する区間を測る。電極接触を持たない誘導検層などには一般化しない。\par

\noindent\textbf{2：適切} ダウンホールで地表起振・孔内受振の初動走時を測る。孔壁に結合した受振器の位置と走時から区間速度を推定する。経路と屈折の補正が必要である。\par

\noindent\textbf{3：条件付きで適切} VSPで孔内に振源又は受振器を設置し反射波を測る。通常の地表振源・孔内受振と逆VSPなどの構成を区別する。孔内深度と反射記録を結び地表反射探査の解釈を助ける。\par

\noindent\textbf{4：不適切} ベータ線のレイリー散乱で密度分布を測る。ガンマ線のCompton散乱が主要な原理である。電子密度と岩石かさ密度の校正を使う。\par

\noindent\textbf{5：適切} サスペンションで孔内振源と複数受振器を用いる。孔内ゾンデの起振と受振点間の走時差により、孔軸に沿う地盤の区間速度を求める。水中のS波伝播を直接測るという意味ではない。\par

\paragraph{再利用する考え方} 流体、孔壁結合、線種、センサ位置を観測量ごとに整理する。

\paragraph{共通理論} \hyperref[theory:exploration:borehole]{第\ref*{theory:exploration:borehole}節}, \hyperref[theory:exploration:gammadensity]{第\ref*{theory:exploration:gammadensity}節}

\par 参考：\cite{epa_borehole_resistivity,epa_gamma_density,geebec_ps,oyo_suspension}。

\paragraph{関連問題} \hyperref[q:r05-reg-applied-34]{R05-REG-APPLIED-34}

\par 難易度 3（校準初稿）、優先度未評価。

\section{直流比抵抗の電極配置・見掛比抵抗}

\subsection{平成23年度 問30}
\label{q:h23-reg-applied-30}

\noindent\textbf{原典：} sources/H23/17\_applied\_science.pdf、PDF物理10頁。

\textbf{公式正答：} 2（single）。

\par\textbf{検討未確定：} 最も関係ないは相対比較である。IPは比抵抗測定と密接に関連するため完全な非関連とは説明しない。

\paragraph{問題の要旨} 比抵抗法の電極配置・計算法とIP分極測定を比較する。

\Needspace{4\baselineskip}\paragraph{確認した条件・記述（独自要約）}\mbox{}\par

\noindent\textbf{原典照合：} 29問21原典画像で否定、全五候補、数式・上付き単位・観測対象を照合。\par

\noindent\textbf{前提：} 直流比抵抗の電極配列・順計算と、時間又は周波数に依存する分極応答を分ける。\par

\paragraph{問われる核心} 比抵抗法の電極配置・計算法とIP分極測定を比較する。

\paragraph{必要知識} 直流比抵抗の電極配列・順計算と、時間又は周波数に依存する分極応答を分ける。

\paragraph{考え方・正答を選ぶ理由（査読待ち初稿）} 2のIP法が、静的な比抵抗を求める三配列とlinear filter計算に対して相対的に別の観測を加える。IPも比抵抗と電極を共有するため、無関係とまでは言えない。4は層状構造の比抵抗変換から見掛比抵抗を計算する線形フィルタである。

\Needspace{4\baselineskip}\paragraph{選択肢の検討（査読待ち初稿）}\mbox{}\par

\noindent\textbf{1：適切} ウェンナー法。等間隔の四電極配置で電流と電位差から見掛比抵抗を求める。\par

\noindent\textbf{2：条件付きで適切} IP法。電流印加で生じる分極の減衰又は位相を測る。静的比抵抗法とは目的量が異なるが、同じ配置で比抵抗も測るので関係のなさは相対的である。\par

\noindent\textbf{3：適切} シュランベルジャー法。電流電極と短い電位電極対を使う比抵抗垂直探査の代表的配置である。\par

\noindent\textbf{4：適切} リニアフィルタ法。層状比抵抗モデルの順応答計算に用いられる。一般の信号ノイズ処理だけと解釈しない。\par

\noindent\textbf{5：適切} ダイポール・ダイポール法。電流双極子と電位双極子を離した比抵抗配置で、二次元調査等に用いる。\par

\paragraph{再利用する考え方} 測定配置、順計算法、追加の分極観測を同じ種類の名称と思わない。

\paragraph{共通理論} \hyperref[theory:exploration:dcarrays]{第\ref*{theory:exploration:dcarrays}節}, \hyperref[theory:exploration:ipsp]{第\ref*{theory:exploration:ipsp}節}

\par 参考：\cite{epa_resistivity,epa_ip,usgs_linear_filter1981}。

\paragraph{関連問題} \hyperref[q:h24-reg-applied-20]{H24-REG-APPLIED-20}

\par 難易度 3（校準初稿）、優先度未評価。

\subsection{平成24年度 問20}
\label{q:h24-reg-applied-20}

\noindent\textbf{原典：} sources/H24/17\_applied\_science.pdf、PDF物理7頁。

\textbf{公式正答：} 5（single）。

\paragraph{問題の要旨} 電極による比抵抗測定とコイルによる電磁誘導を区別する。

\Needspace{4\baselineskip}\paragraph{確認した条件・記述（独自要約）}\mbox{}\par

\noindent\textbf{原典照合：} 29問21原典画像で否定、全五候補、数式・上付き単位・観測対象を照合。\par

\noindent\textbf{前提：} 電気的接触を持つ電極測定と送受信ループの誘導測定を区別する。\par

\paragraph{問われる核心} 電極による比抵抗測定とコイルによる電磁誘導を区別する。

\paragraph{必要知識} 電気的接触を持つ電極測定と送受信ループの誘導測定を区別する。

\paragraph{考え方・正答を選ぶ理由（査読待ち初稿）} 5のループ・ループ法が電磁誘導による手法である。残りは電極の接触、配列、見掛比抵抗という電流注入型比抵抗測定に直接関係する。

\Needspace{4\baselineskip}\paragraph{選択肢の検討（査読待ち初稿）}\mbox{}\par

\noindent\textbf{1：適切} 接地抵抗。地盤へ電流を流す電極の接触抵抗が送信と測定条件へ影響する。\par

\noindent\textbf{2：適切} ウェンナー法。等間隔四電極の比抵抗配置である。\par

\noindent\textbf{3：適切} 見掛け比抵抗。測ったV/Iへ配置の幾何係数を掛けた等価半空間の抵抗率である。\par

\noindent\textbf{4：適切} 電極。電流印加と電位差測定に用いる。\par

\noindent\textbf{5：不適切} ループ・ループ法。送受信コイルから地盤の誘導電流応答を測る手法で、電極を使う直流比抵抗測定の用語ではない。比抵抗を推定する目的自体は共通する。\par

\paragraph{再利用する考え方} 推定する物性が同じでも励起・受信の仕組みが異なる。

\paragraph{共通理論} \hyperref[theory:exploration:dcarrays]{第\ref*{theory:exploration:dcarrays}節}, \hyperref[theory:exploration:fdtd]{第\ref*{theory:exploration:fdtd}節}

\par 参考：\cite{epa_resistivity,epa_fdem}。

\paragraph{関連問題} \hyperref[q:h23-reg-applied-30]{H23-REG-APPLIED-30}

\par 難易度 2（校準初稿）、優先度未評価。

\subsection{平成26年度 問23}
\label{q:h26-reg-applied-23}

\noindent\textbf{原典：} sources/H26/17\_applied\_science.pdf、PDF物理10頁。

\textbf{公式正答：} 4（single）。

\par\textbf{検討未確定：} 低比抵抗層の誘導検層精度と迅速広域性は装置条件による。Ω・m\textsuperscript{3}の原典上付き3を画像照合。

\paragraph{問題の要旨} 電気・電磁法の比抵抗単位と測定配置を比較する。

\Needspace{4\baselineskip}\paragraph{確認した条件・記述（独自要約）}\mbox{}\par

\noindent\textbf{原典照合：} 29問21原典画像で否定、全五候補、数式・上付き単位・観測対象を照合。\par

\noindent\textbf{前提：} 抵抗率ρ=RA/LなのでΩmである。電極電流と誘導電流を区別する。\par

\paragraph{問われる核心} 電気・電磁法の比抵抗単位と測定配置を比較する。

\paragraph{必要知識} 抵抗率ρ=RA/LなのでΩmである。電極電流と誘導電流を区別する。

\paragraph{考え方・正答を選ぶ理由（査読待ち初稿）} 4が不適切である。R=ρL/Aからρ=RA/Lとなり、Ω×m\textsuperscript{2}/m=Ωmである。原典のΩ・m\textsuperscript{3}は誤った単位である。

\Needspace{4\baselineskip}\paragraph{選択肢の検討（査読待ち初稿）}\mbox{}\par

\noindent\textbf{1：適切} 直流又は低周波交替直流の電場から比抵抗分布を推定。極性切替を含む低周波電流と電位差の測定を地盤モデルへ逆解析する。\par

\noindent\textbf{2：条件付きで適切} 電磁検層はコイル誘導、低比抵抗層に適する。導電率が高い地層は誘導二次場を生みやすい。ただし誘導数・装置校正・孔状態にも依存し全範囲で精度増大とはしない。\par

\noindent\textbf{3：条件付きで適切} 人工又は自然磁場変動から比抵抗分布を効率的に測る。電磁誘導応答は電気物性へ関係する。効率と面積は方式・機器に依存する。\par

\noindent\textbf{4：不適切} 比抵抗の単位はΩ・m\textsuperscript{3}。抵抗率は電流密度と電場の比例係数でΩmである。体積の単位を掛けると次元が違う。\par

\noindent\textbf{5：適切} 領域を囲む電極から二次元・三次元比抵抗tomography。多方向・多数電流配置の電位データを逆解析して比抵抗分布を推定する。覆い方により分解できる構造が変わる。\par

\paragraph{再利用する考え方} 物性単位を試料幾何から導出し、測定名と物性を一対一に混同しない。

\paragraph{共通理論} \hyperref[theory:exploration:dcarrays]{第\ref*{theory:exploration:dcarrays}節}, \hyperref[theory:exploration:fdtd]{第\ref*{theory:exploration:fdtd}節}

\par 参考：\cite{epa_resistivity,epa_fdem,epa_borehole_resistivity}。

\par 難易度 3（校準初稿）、優先度未評価。

\section{探査深度の目安と観測条件}

\subsection{平成27年度 問21}
\label{q:h27-reg-applied-21}

\noindent\textbf{原典：} sources/H27/17\_applied\_science.pdf、PDF物理8頁。

\textbf{公式正答：} 1（single）。

\par\textbf{検討未確定：} 公式1の100mと通常20mの区別は当時の一般的装置条件の読みである。長波長表面波で100mが不可能とはいえない。比抵抗300mと音響基盤も一律保証ではない。

\paragraph{問題の要旨} 表面波の深度目安、屈折ブラインド、反射展開、比抵抗、音響基盤を比較する。

\Needspace{4\baselineskip}\paragraph{確認した条件・記述（独自要約）}\mbox{}\par

\noindent\textbf{原典照合：} 五候補の順序、最適切と不適切の方向、波種、深度目安、曲線名、最長と最短を原典画像で照合。\par

\noindent\textbf{前提：} 能動の浅部表面波探査の通常目安と、長波長の利用により深度が変わる事実を分ける。\par

\paragraph{問われる核心} 表面波の深度目安、屈折ブラインド、反射展開、比抵抗、音響基盤を比較する。

\paragraph{必要知識} 能動の浅部表面波探査の通常目安と、長波長の利用により深度が変わる事実を分ける。

\paragraph{考え方・正答を選ぶ理由（査読待ち初稿）} 公式1は通常の浅層表面波の目安を100m程度とした説明を選ぶ。環境省資料では20m程度、数m〜数十mの通常対象を述べる。一方、長周期起振や微動で100m級以上を調べる方法もあるため、100mが原理的に不可能とは解説しない。

\Needspace{4\baselineskip}\paragraph{選択肢の検討（査読待ち初稿）}\mbox{}\par

\noindent\textbf{1：条件付きで適切} 表面波の分散でS波構造を調べ探査深度は100m程度までである。分散によるS波推定は正しい。通常の浅部能動測定では20m級の目安であり公式は100mを不適切とするが、装置と最大波長を指定しない普遍的上限は定められない。\par

\noindent\textbf{2：適切} 深部へ増速する層構造でも薄い中間層が初動に現れない。臨界屈折波が生じても、直接波や深い層を通る波より早くなる距離範囲がなければ初動に現れない。低速度逆転だけが見えない原因ではない。\par

\noindent\textbf{3：条件付きで適切} 浅層反射で最大オフセットを探査深度程度に設定する。浅層反射の一般的な展開目安として読む。速度解析・近接雑音・反射角・地形・重合数で調整し、全ての計測に必要な厳密等式ではない。\par

\noindent\textbf{4：条件付きで適切} 二次元比抵抗法は低比抵抗部を検出し実用深度が300m程度である。導電性領域へ電流が集まりやすい等の特性があるが、電極展開、コントラスト、ノイズで感度と深度が変わる。300mは原典の典型目安であり装置非依存の保証ではない。\par

\noindent\textbf{5：条件付きで適切} 海上音波で海底下反射を測り固結層等の上面が音響基盤となる。大きなインピーダンス差と減衰等によって深部への透過が弱まり、内部を捉えにくくなる。基盤より下は常に観測不能という絶対命題ではない。\par

\paragraph{再利用する考え方} 探査深度を手法名に固定せず、波長・配列・震源・物性・雑音条件を確認する。

\paragraph{共通理論} \hyperref[theory:exploration:surfacewaves]{第\ref*{theory:exploration:surfacewaves}節}, \hyperref[theory:exploration:reflectiongeometry]{第\ref*{theory:exploration:reflectiongeometry}節}

\par 参考：\cite{moe_surface2004,epa_surfacewaves,segj_integrated2017,epa_seismic_reflection}。

\par 難易度 3（校準初稿）、優先度未評価。

\section{周波数領域と時間領域電磁応答}

\subsection{平成28年度 問19}
\label{q:h28-reg-applied-19}

\noindent\textbf{原典：} sources/H28/17\_applied\_science.pdf、PDF物理7頁。

\textbf{公式正答：} 1（single）。

\paragraph{問題の要旨} 電気・電磁現象と地震反射のNMO補正を区別する。

\Needspace{4\baselineskip}\paragraph{確認した条件・記述（独自要約）}\mbox{}\par

\noindent\textbf{原典照合：} 29問21原典画像で否定、全五候補、数式・上付き単位・観測対象を照合。\par

\noindent\textbf{前提：} NMOは反射走時のoffset依存を補正する地震探査処理である。\par

\paragraph{問われる核心} 電気・電磁現象と地震反射のNMO補正を区別する。

\paragraph{必要知識} NMOは反射走時のoffset依存を補正する地震探査処理である。

\paragraph{考え方・正答を選ぶ理由（査読待ち初稿）} 1のNMO補正である。表皮・近接場・誘導分極・Schumann共振は電気・電磁応答に直接関係する。SchumannはMT電磁記録でも観測される。

\Needspace{4\baselineskip}\paragraph{選択肢の検討（査読待ち初稿）}\mbox{}\par

\noindent\textbf{1：不適切} NMO補正。normal moveoutは震源受振器の距離による反射走時差を補正する。電磁現象の名称ではない。\par

\noindent\textbf{2：適切} 表皮効果。導電媒質内の周期電磁場の減衰と深度感度へ関係する。\par

\noindent\textbf{3：適切} ニアフィールド効果。有限送信源からの近距離場が平面波近似と異なることをいう。\par

\noindent\textbf{4：適切} IP効果。電流による地下界面の分極で減衰・位相が変わる。\par

\noindent\textbf{5：適切} シューマン共振。地球と電離層の電磁共振であり、自然MTの電場磁場記録にも現れる。\par

\paragraph{再利用する考え方} 同じ探査で扱う補正でも信号の物理原理を識別する。

\paragraph{共通理論} \hyperref[theory:exploration:nearfield]{第\ref*{theory:exploration:nearfield}節}, \hyperref[theory:exploration:ipsp]{第\ref*{theory:exploration:ipsp}節}, \hyperref[theory:exploration:reflectiongeometry]{第\ref*{theory:exploration:reflectiongeometry}節}

\par 参考：\cite{schumann_mt2009,epa_ip,usgs_csamt2012}。

\par 難易度 2（校準初稿）、優先度未評価。

\subsection{令和2年度 問34}
\label{q:r02-reg-applied-34}

\noindent\textbf{原典：} sources/R02/17\_applied\_science.pdf、PDF物理12頁。

\textbf{公式正答：} 4（single）。

\par\textbf{検討未確定：} 精度0.01mGalは機器/現場の条件値、検層孔隙率は間接推定、AEの先行は全称ではない。

\paragraph{問題の要旨} 周波数領域の位相と時間領域の過渡応答を区別する。

\Needspace{4\baselineskip}\paragraph{確認した条件・記述（独自要約）}\mbox{}\par

\noindent\textbf{原典照合：} 29問21原典画像で否定、全五候補、数式・上付き単位・観測対象を照合。\par

\noindent\textbf{前提：} 正弦波定常応答は振幅と位相、遮断後応答は時間に対する減衰を測る。\par

\paragraph{問われる核心} 周波数領域の位相と時間領域の過渡応答を区別する。

\paragraph{必要知識} 正弦波定常応答は振幅と位相、遮断後応答は時間に対する減衰を測る。

\paragraph{考え方・正答を選ぶ理由（査読待ち初稿）} 4が不適切である。Slingram等の周波数領域電磁法は連続正弦送信に対する受信振幅と位相差を測る。過渡応答は時間領域の測定である。

\Needspace{4\baselineskip}\paragraph{選択肢の検討（査読待ち初稿）}\mbox{}\par

\noindent\textbf{1：条件付きで適切} 微小重力差の相対測定や鉛直重力勾配を利用。0.01mGal級以下の精度を持つ計器の利用と勾配観測がある。全装置と測定現場に共通の精度保証ではない。\par

\noindent\textbf{2：適切} 領域を囲む信号源センサで物性を逆解析。tomographyの観測覆いと感度に沿って物性分布を推定する。\par

\noindent\textbf{3：条件付きで適切} 検層で比抵抗密度速度孔隙率等を深度に記録。物性を連続記録し、孔隙率のように複数物性と校正から推定する量もある。すべてが直接測定ではない。\par

\noindent\textbf{4：不適切} 周波数領域で連続波の振幅と過渡応答を測る。定常の振幅と位相を測る方式と、遮断後の時間領域過渡を混同する。\par

\noindent\textbf{5：条件付きで適切} AEは微小破壊で主破壊に先行することが多い。微小破壊に伴う弾性波放出が前兆の研究対象になる。主破壊前だけの定義又は常に発生する保証ではない。\par

\paragraph{再利用する考え方} f領域の位相とt領域の減衰、直接物性と校正推定量を分ける。

\paragraph{共通理論} \hyperref[theory:exploration:fdtd]{第\ref*{theory:exploration:fdtd}節}, \hyperref[theory:exploration:multimethod]{第\ref*{theory:exploration:multimethod}節}

\par 参考：\cite{epa_fdem,epa_tdem_archive,sato_lei2009_ae,gsi_gravity_definitions}。

\par 難易度 3（校準初稿）、優先度未評価。

\section{探査信号のフィルタ選択}

\subsection{平成23年度 問28}
\label{q:h23-reg-applied-28}

\noindent\textbf{原典：} sources/H23/17\_applied\_science.pdf、PDF物理9頁。

\textbf{公式正答：} 4（single）。

\par\textbf{検討未確定：} 移動平均の除去は低減であり完全除去ではない。線形性だけで単一応答の畳み込みは保証しない。標本化の原文は一周期と間隔を混ぜた表現であるため、帯域と最高周波数を明示して条件付きで読む。

\paragraph{問題の要旨} 空間フィルタ・平均・インパルス応答・低域通過・標本化を区別する。

\Needspace{4\baselineskip}\paragraph{確認した条件・記述（独自要約）}\mbox{}\par

\noindent\textbf{原典照合：} 否定の対象、指数、周波数と周期、五つの候補の列順を原典画像で照合。\par

\noindent\textbf{前提：} 通常の二次元画像と時系列を対象とする。低域は低周波数すなわち長周期である。\par

\paragraph{問われる核心} 空間フィルタ・平均・インパルス応答・低域通過・標本化を区別する。

\paragraph{必要知識} 通常の二次元画像と時系列を対象とする。低域は低周波数すなわち長周期である。

\paragraph{考え方・正答を選ぶ理由（査読待ち初稿）} 4は低域通過の役割を逆にしている。長周期成分を残し、短周期成分を減衰させるのが低域通過である。他の説明は典型的処理として読むが、完全除去と畳み込みの省略条件を残す。

\Needspace{4\baselineskip}\paragraph{選択肢の検討（査読待ち初稿）}\mbox{}\par

\noindent\textbf{1：適切} 空間フィルタを二次元画像のエッジ・平滑化・差分に用いる。二次元データに対し空間方向の重み付けや差分を行う。一次元・三次元にも拡張でき、二次元しか扱えないという意味ではない。\par

\noindent\textbf{2：条件付きで適切} 注目点付近の移動平均でランダム成分を除去する。独立で平均ゼロの雑音をM点平均すれば分散は1/Mに減る。有限Mで雑音が厳密にゼロになる保証はない。中央窓以外に因果的な後方窓もある。\par

\noindent\textbf{3：条件付きで適切} 線形システムのインパルス応答を重ねて任意入力への応答を求める。線形性により各時刻のインパルスへの応答を重ねることはできる。単一の応答を時間移動する畳み込み表現には時不変性も必要であり、時変系には二時刻の応答を用いる。\par

\noindent\textbf{4：不適切} 低域通過は短周期を残し長周期を減衰させる。周期Tと周波数fは逆数の関係である。低周波数は長周期なので、記述の通過側と減衰側は逆である。\par

\noindent\textbf{5：条件付きで適切} 一周期二標本未満では元データを再構成できない。一般の帯域制限信号を一意に再構成するにはスペクトルの重なりを避ける標本化が必要である。特別な事前情報を持つ信号まで一律に不可能とする定理ではなく、最高周波数と標本化間隔を明示して読む。\par

\paragraph{再利用する考え方} フィルタの名称を周波数で定義してから周期へ変換する。

\paragraph{共通理論} \hyperref[theory:exploration:filters]{第\ref*{theory:exploration:filters}節}, \hyperref[theory:exploration:lti]{第\ref*{theory:exploration:lti}節}, \hyperref[theory:exploration:sampling]{第\ref*{theory:exploration:sampling}節}

\par 参考：\cite{mit_signals2010,mit_moving2012,scipy_notch,scipy_fft,usgs_processing2000}。

\paragraph{関連問題} \hyperref[q:h28-reg-applied-20]{H28-REG-APPLIED-20}

\par 難易度 3（校準初稿）、優先度未評価。

\subsection{平成24年度 問19}
\label{q:h24-reg-applied-19}

\noindent\textbf{原典：} sources/H24/17\_applied\_science.pdf、PDF物理7頁。

\textbf{公式正答：} 5（single）。

\par\textbf{検討未確定：} 移動平均でランダム成分を除去するとは統計的な低減の意味である。信号まで平均化で弱める場合と、相関雑音が残る場合を区別する。

\paragraph{問題の要旨} 相互相関・重合・移動平均・FFT・ノッチの目的を対応させる。

\Needspace{4\baselineskip}\paragraph{確認した条件・記述（独自要約）}\mbox{}\par

\noindent\textbf{原典照合：} 否定の対象、指数、周波数と周期、五つの候補の列順を原典画像で照合。\par

\noindent\textbf{前提：} 二つの時系列の相関、複数記録の平均化、DFT、狭帯域阻止を区別する。\par

\paragraph{問われる核心} 相互相関・重合・移動平均・FFT・ノッチの目的を対応させる。

\paragraph{必要知識} 二つの時系列の相関、複数記録の平均化、DFT、狭帯域阻止を区別する。

\paragraph{考え方・正答を選ぶ理由（査読待ち初稿）} 5が不適切である。ノッチは特定周波数の周囲を阻止するものであり、その周波数だけを通過させる説明は逆である。ランダム雑音の抑制は統計的効果であり有限平均での完全消失ではない。

\Needspace{4\baselineskip}\paragraph{選択肢の検討（査読待ち初稿）}\mbox{}\par

\noindent\textbf{1：適切} 相互相関から二波形の類似と時間遅れを検出する。一方の波形をずらして積の和又は平均を求め、相関の極大から遅れを推定する。極大が常に唯一とは限らず、周期性・雑音に注意する。\par

\noindent\textbf{2：条件付きで適切} 複数記録の重合でランダム雑音を抑える。信号の位相をそろえ、雑音が独立又は無相関なら振幅S/Nは回数の平方根だけ改善する。相関雑音や未補正の走時差がある場合は同じ効果を保証しない。\par

\noindent\textbf{3：条件付きで適切} 注目点付近の移動平均でランダム成分を除去する。独立で平均ゼロの雑音をM点平均すれば分散は1/Mに減る。有限Mで雑音が厳密にゼロになる保証はない。中央窓以外に因果的な後方窓もある。\par

\noindent\textbf{4：適切} FFTで離散時系列のフーリエ係数を高速に計算する。FFTはDFTと異なる変換ではなく、DFTを効率よく計算する算法である。標本化によるaliasingを自動的に取り除く算法ではない。\par

\noindent\textbf{5：不適切} ノッチで指定周波数成分だけを通過させる。ノッチは狭帯域阻止である。通過と遮断を逆にしており、帯域通過の説明と取り違えている。\par

\paragraph{再利用する考え方} 名称ではなく保持する成分と抑える成分を書く。

\paragraph{共通理論} \hyperref[theory:exploration:filters]{第\ref*{theory:exploration:filters}節}, \hyperref[theory:exploration:correlation]{第\ref*{theory:exploration:correlation}節}, \hyperref[theory:exploration:stack]{第\ref*{theory:exploration:stack}節}, \hyperref[theory:exploration:fft]{第\ref*{theory:exploration:fft}節}

\par 参考：\cite{mit_signals2010,mit_moving2012,scipy_notch,scipy_fft,usgs_processing2000}。

\paragraph{関連問題} \hyperref[q:h29-reg-applied-35]{H29-REG-APPLIED-35}

\par 難易度 3（校準初稿）、優先度未評価。

\subsection{平成29年度 問35}
\label{q:h29-reg-applied-35}

\noindent\textbf{原典：} sources/H29/17\_applied\_science.pdf、PDF物理12頁。

\textbf{公式正答：} 5（single）。

\par\textbf{検討未確定：} 移動平均でランダム成分を除去するとは統計的な低減の意味である。信号まで平均化で弱める場合と、相関雑音が残る場合を区別する。

\paragraph{問題の要旨} 相互相関・重合・移動平均・FFT・ノッチの目的を対応させる。

\Needspace{4\baselineskip}\paragraph{確認した条件・記述（独自要約）}\mbox{}\par

\noindent\textbf{原典照合：} 否定の対象、指数、周波数と周期、五つの候補の列順を原典画像で照合。\par

\noindent\textbf{前提：} 二つの時系列の相関、複数記録の平均化、DFT、狭帯域阻止を区別する。\par

\paragraph{問われる核心} 相互相関・重合・移動平均・FFT・ノッチの目的を対応させる。

\paragraph{必要知識} 二つの時系列の相関、複数記録の平均化、DFT、狭帯域阻止を区別する。

\paragraph{考え方・正答を選ぶ理由（査読待ち初稿）} 5が不適切である。ノッチは特定周波数の周囲を阻止するものであり、その周波数だけを通過させる説明は逆である。ランダム雑音の抑制は統計的効果であり有限平均での完全消失ではない。

\Needspace{4\baselineskip}\paragraph{選択肢の検討（査読待ち初稿）}\mbox{}\par

\noindent\textbf{1：適切} 相互相関から二波形の類似と時間遅れを検出する。一方の波形をずらして積の和又は平均を求め、相関の極大から遅れを推定する。極大が常に唯一とは限らず、周期性・雑音に注意する。\par

\noindent\textbf{2：条件付きで適切} 複数記録の重合でランダム雑音を抑える。信号の位相をそろえ、雑音が独立又は無相関なら振幅S/Nは回数の平方根だけ改善する。相関雑音や未補正の走時差がある場合は同じ効果を保証しない。\par

\noindent\textbf{3：条件付きで適切} 注目点付近の移動平均でランダム成分を除去する。独立で平均ゼロの雑音をM点平均すれば分散は1/Mに減る。有限Mで雑音が厳密にゼロになる保証はない。中央窓以外に因果的な後方窓もある。\par

\noindent\textbf{4：適切} FFTで離散時系列のフーリエ係数を高速に計算する。FFTはDFTと異なる変換ではなく、DFTを効率よく計算する算法である。標本化によるaliasingを自動的に取り除く算法ではない。\par

\noindent\textbf{5：不適切} ノッチで指定周波数成分だけを通過させる。ノッチは狭帯域阻止である。通過と遮断を逆にしており、帯域通過の説明と取り違えている。\par

\paragraph{再利用する考え方} 名称ではなく保持する成分と抑える成分を書く。

\paragraph{共通理論} \hyperref[theory:exploration:filters]{第\ref*{theory:exploration:filters}節}, \hyperref[theory:exploration:correlation]{第\ref*{theory:exploration:correlation}節}, \hyperref[theory:exploration:stack]{第\ref*{theory:exploration:stack}節}, \hyperref[theory:exploration:fft]{第\ref*{theory:exploration:fft}節}

\par 参考：\cite{mit_signals2010,mit_moving2012,scipy_notch,scipy_fft,usgs_processing2000}。

\paragraph{関連問題} \hyperref[q:h24-reg-applied-19]{H24-REG-APPLIED-19}

\par 難易度 3（校準初稿）、優先度未評価。

\section{探査の順問題・逆問題と対象雑音}

\subsection{令和7年度 問31}
\label{q:r07-reg-applied-31}

\noindent\textbf{原典：} sources/R07/17\_applied\_science.pdf、PDF物理11頁。

\textbf{公式正答：} 1（single）。

\par\textbf{検討未確定：} 標本化の等号端点と帯域制限、理想再構成の条件が省略される。S/Nの平方根は改善倍率であり、元のS/Nが常に1ではない。

\paragraph{問題の要旨} 順問題と逆問題、対象依存の雑音、標本化と重合を区別する。

\Needspace{4\baselineskip}\paragraph{確認した条件・記述（独自要約）}\mbox{}\par

\noindent\textbf{原典照合：} 否定の対象、指数、周波数と周期、五つの候補の列順を原典画像で照合。\par

\noindent\textbf{前提：} 地下モデルからデータを予測するのが順問題、観測データから地下を推定するのが逆問題である。\par

\paragraph{問われる核心} 順問題と逆問題、対象依存の雑音、標本化と重合を区別する。

\paragraph{必要知識} 地下モデルからデータを予測するのが順問題、観測データから地下を推定するのが逆問題である。

\paragraph{考え方・正答を選ぶ理由（査読待ち初稿）} 1が不適切である。地下構造を求める解析を順問題、計測を逆問題とする対応が逆である。計測そのものはデータ取得であり、順問題とはモデルからデータを計算する写像を指す。

\Needspace{4\baselineskip}\paragraph{選択肢の検討（査読待ち初稿）}\mbox{}\par

\noindent\textbf{1：不適切} 地下推定の解析が順問題、データを得る計測が逆問題である。F(m)でモデルからデータを計算するのが順問題である。与えたデータからmを推定するのが逆問題であり、原典の役割は逆転している。\par

\noindent\textbf{2：適切} 入力信号によって発生した表面波でも解析対象でなければ雑音である。ノイズは対象に対する不要成分の呼び方である。表面波探査では信号となる同じ波が、反射法では相関した妨害波となり得る。\par

\noindent\textbf{3：条件付きで適切} 最高周波数の二倍以上で標本化すれば完全復元できる。帯域制限・理想的標本と再構成を想定した説明である。最高周波数成分を含む場合の等号端点、量子化、有限長、雑音を無視できない。安全な条件はfs>2fmaxであり、境界では位相によって標本が全てゼロとなる例がある。\par

\noindent\textbf{4：適切} ノッチで特定の周波数成分を遮断する。狭い阻止帯域で商用電源などの周期的な混入を抑える。指定帯域だけを通過させるフィルタではない。実装には有限の帯域幅と位相応答がある。\par

\noindent\textbf{5：条件付きで適切} 重合によるランダム雑音に対するS/Nが回数の平方根になる。同一信号の位相と振幅がそろい、雑音が無相関・平均ゼロ・同分散なら成り立つ。一般には元のS/Nに平方根Nを掛ける改善倍率を意味する。\par

\paragraph{再利用する考え方} ノイズをランダムに限らず目的から定義し、写像の向きで順逆を区別する。

\paragraph{共通理論} \hyperref[theory:exploration:inverse]{第\ref*{theory:exploration:inverse}節}, \hyperref[theory:exploration:stack]{第\ref*{theory:exploration:stack}節}, \hyperref[theory:exploration:sampling]{第\ref*{theory:exploration:sampling}節}

\par 参考：\cite{mit_signals2010,mit_moving2012,scipy_notch,scipy_fft,usgs_processing2000,pygimli_inversion}。

\par 難易度 3（校準初稿）、優先度未評価。

\section{四電極と接地抵抗の分離}

\subsection{平成25年度 問34}
\label{q:h25-reg-applied-34}

\noindent\textbf{原典：} sources/H25/17\_applied\_science.pdf、PDF物理12頁。

\textbf{公式正答：} 3（single）。

\paragraph{問題の要旨} 四電極法で接触抵抗の電圧降下を分離する理由を選ぶ。

\Needspace{4\baselineskip}\paragraph{確認した条件・記述（独自要約）}\mbox{}\par

\noindent\textbf{原典照合：} 29問21原典画像で否定、全五候補、数式・上付き単位・観測対象を照合。\par

\noindent\textbf{前提：} 電流注入電極と高入力抵抗の電位測定電極を分ける。電位電極電流は十分小さい。\par

\paragraph{問われる核心} 四電極法で接触抵抗の電圧降下を分離する理由を選ぶ。

\paragraph{必要知識} 電流注入電極と高入力抵抗の電位測定電極を分ける。電位電極電流は十分小さい。

\paragraph{考え方・正答を選ぶ理由（査読待ち初稿）} 3である。二端子では地盤に加えて電流電極の接触抵抗の電圧降下も測る。別の電位電極へ高入力抵抗の電圧計を接続すれば、その電極を流れる電流が小さく、接触抵抗による誤差を軽減する。送信側の接地抵抗は依然として電流確保に影響する。

\Needspace{4\baselineskip}\paragraph{選択肢の検討（査読待ち初稿）}\mbox{}\par

\noindent\textbf{1：不適切} 電磁カップリングを小さくする。電磁誘導の結合はケーブル配置・周波数等にも依存し、四電極化の主要な理由ではない。\par

\noindent\textbf{2：不適切} 測定電圧を大きくする。Vは電極幾何と地盤応答で決まり、電極を分けることだけで増大しない。\par

\noindent\textbf{3：適切} 接地抵抗の影響を軽減する。電流端子の接触電圧降下を電位差観測から分け、電位端子の小電流によりその接触降下も抑える。\par

\noindent\textbf{4：不適切} 地形の影響を無視する。地形は電流分布を変えるため、別電位電極でも順モデルの地形補正が必要である。\par

\noindent\textbf{5：不適切} 探査深度を大きくする。深度感度は配列と間隔等で決まり、電極端子を分ける理由そのものではない。\par

\paragraph{再利用する考え方} 送信電流を確保する問題と、電位測定の接触誤差を減らす問題を分ける。

\paragraph{共通理論} \hyperref[theory:exploration:fourpoint]{第\ref*{theory:exploration:fourpoint}節}

\par 参考：\cite{epa_resistivity}。

\par 難易度 2（校準初稿）、優先度未評価。

\section{GPRの周波数・分解能・探査深度}

\subsection{平成24年度 問21}
\label{q:h24-reg-applied-21}

\noindent\textbf{原典：} sources/H24/17\_applied\_science.pdf、PDF物理7頁。

\textbf{公式正答：} 3（single）。

\par\textbf{検討未確定：} 固定波形の理想鉛直分解能は深度そのものだけで単調低下するわけではない。実地の交換条件と検出限界を区別。

\paragraph{問題の要旨} GPRの深度・周波数・分解能と調査配置を比較する。

\Needspace{4\baselineskip}\paragraph{確認した条件・記述（独自要約）}\mbox{}\par

\noindent\textbf{原典照合：} 29問21原典画像で否定、全五候補、数式・上付き単位・観測対象を照合。\par

\noindent\textbf{前提：} 低損失非磁性媒質でv≈c/√εr、波長λ=v/f、パルス帯域Bに対する距離分解はおよそv/(2B)である。導電率等が到達深度を制限する。\par

\paragraph{問われる核心} GPRの深度・周波数・分解能と調査配置を比較する。

\paragraph{必要知識} 低損失非磁性媒質でv≈c/√εr、波長λ=v/f、パルス帯域Bに対する距離分解はおよそv/(2B)である。導電率等が到達深度を制限する。

\paragraph{考え方・正答を選ぶ理由（査読待ち初稿）} 3が不適切である。深部まで測るための低周波化、波の減衰・拡がりやノイズにより一般に深部の分解・検出が難しくなる。深度が増しても分解能が低下しないという一般保証はできない。

\Needspace{4\baselineskip}\paragraph{選択肢の検討（査読待ち初稿）}\mbox{}\par

\noindent\textbf{1：適切} 地表型と孔内型のレーダがある。地表からの反射と孔内アンテナによる反射・透過などの構成がある。\par

\noindent\textbf{2：適切} 周波数は探査深度と分解能へ関係。波長・帯域と吸収が分解と到達深度を制限する。\par

\noindent\textbf{3：不適切} 一般に深度が増しても分解能は低下しない。深部の低周波化や減衰・ビーム拡がりが分解と検出を難しくする。非低下の一般保証は不適切である。\par

\noindent\textbf{4：適切} 直線埋設物へできる限り直角の測線。横切る測線で位置の変化と点状反射の走時を捉えやすくする。未知方位では複数方位を測る。\par

\noindent\textbf{5：適切} 高分解能なので路面下空洞に利用。浅部の空洞境界の電気特性差と十分な波長・信号品質を使う。\par

\paragraph{再利用する考え方} 反射走時を深度と混同せず、帯域・速度・損失と探査目的を対応させる。

\paragraph{共通理論} \hyperref[theory:exploration:gpr]{第\ref*{theory:exploration:gpr}節}, \hyperref[theory:exploration:gprresolution]{第\ref*{theory:exploration:gprresolution}節}

\par 参考：\cite{epa_gpr}。

\par 難易度 3（校準初稿）、優先度未評価。

\subsection{令和2年度 問33}
\label{q:r02-reg-applied-33}

\noindent\textbf{原典：} sources/R02/17\_applied\_science.pdf、PDF物理12頁。

\textbf{公式正答：} 3（single）。

\par\textbf{検討未確定：} 2の分解能は能力と最小距離で比例方向が変わる。深度の厳密反比例と5の間接相関全面否定は留保。

\paragraph{問題の要旨} GPRのMHz帯と測線・速度・強度の推定限界を照合する。

\Needspace{4\baselineskip}\paragraph{確認した条件・記述（独自要約）}\mbox{}\par

\noindent\textbf{原典照合：} 29問21原典画像で否定、全五候補、数式・上付き単位・観測対象を照合。\par

\noindent\textbf{前提：} 低損失非磁性媒質でv≈c/√εr、波長λ=v/f、パルス帯域Bに対する距離分解はおよそv/(2B)である。導電率等が到達深度を制限する。\par

\paragraph{問われる核心} GPRのMHz帯と測線・速度・強度の推定限界を照合する。

\paragraph{必要知識} 低損失非磁性媒質でv≈c/√εr、波長λ=v/f、パルス帯域Bに対する距離分解はおよそv/(2B)である。導電率等が到達深度を制限する。

\paragraph{考え方・正答を選ぶ理由（査読待ち初稿）} 3の10kHzが不適切である。通常の地中レーダはMHz帯以上の高周波パルスを使い、10kHzを分類境界とする説明は三桁違う。2の比例は定量則ではなく分解能力と深度の傾向として限定する。

\Needspace{4\baselineskip}\paragraph{選択肢の検討（査読待ち初稿）}\mbox{}\par

\noindent\textbf{1：適切} profileにより構造、wide angleで電磁波速度を推定。走査波形を位置へ並べるprofileと、送受信距離を変えて走時から速度を求める構成を使い分ける。\par

\noindent\textbf{2：条件付きで適切} 周波数と分解能は比例、深度は逆比例。周波数が高いほど分解能力が良く、深度を制限する通常の傾向である。最小識別距離は波長に比例し、実深度が厳密に1/fとなる保証はない。\par

\noindent\textbf{3：不適切} 10kHz程度以上をGPR、以下を電磁探査とする。通常の目安は10MHz程度であり、kHzの桁を取り違える。物理的境界としての厳密性もない。\par

\noindent\textbf{4：適切} 直線埋設物へできる限り直角の測線。横切る測線で位置の変化と点状反射の走時を捉えやすくする。未知方位では複数方位を測る。\par

\noindent\textbf{5：条件付きで適切} 電気境界を測るため力学特性との相関はほとんどない。強度や硬さを直接測らないことは正しい。ただし含水・空隙・劣化など共通要因を通じた間接相関まで全面否定はできない。\par

\paragraph{再利用する考え方} 反射走時を深度と混同せず、帯域・速度・損失と探査目的を対応させる。

\paragraph{共通理論} \hyperref[theory:exploration:gpr]{第\ref*{theory:exploration:gpr}節}, \hyperref[theory:exploration:gprresolution]{第\ref*{theory:exploration:gprresolution}節}

\par 参考：\cite{epa_gpr}。

\paragraph{関連問題} \hyperref[q:r05-reg-applied-33]{R05-REG-APPLIED-33}

\par 難易度 3（校準初稿）、優先度未評価。

\subsection{令和3年度 問35}
\label{q:r03-reg-applied-35}

\noindent\textbf{原典：} sources/R03/17\_applied\_science.pdf、PDF物理13頁。

\textbf{公式正答：} 4（single）。

\par\textbf{検討未確定：} 原典のF・mは画像照合済み。比誘電率の空気と真空、10MHz分類目安及び各用途の媒質条件を区別。

\paragraph{問題の要旨} GPRの比誘電率の基準と誘電率の単位を照合する。

\Needspace{4\baselineskip}\paragraph{確認した条件・記述（独自要約）}\mbox{}\par

\noindent\textbf{原典照合：} 29問21原典画像で否定、全五候補、数式・上付き単位・観測対象を照合。\par

\noindent\textbf{前提：} 低損失非磁性媒質でv≈c/√εr、波長λ=v/f、パルス帯域Bに対する距離分解はおよそv/(2B)である。導電率等が到達深度を制限する。\par

\paragraph{問われる核心} GPRの比誘電率の基準と誘電率の単位を照合する。

\paragraph{必要知識} 低損失非磁性媒質でv≈c/√εr、波長λ=v/f、パルス帯域Bに対する距離分解はおよそv/(2B)である。導電率等が到達深度を制限する。

\paragraph{考え方・正答を選ぶ理由（査読待ち初稿）} 4が不適切である。D=εEからεの単位は(C/m\textsuperscript{2})/(V/m)=\allowbreak{}F/mで、F・mではない。比誘電率は真空ε0に対する比ε/ε0である。空気は通常近似として真空に近いが厳密な定義ではない。

\Needspace{4\baselineskip}\paragraph{選択肢の検討（査読待ち初稿）}\mbox{}\par

\noindent\textbf{1：適切} 誘電率・導電率の異なる境界からの反射を測る。電磁インピーダンスと減衰の違いを波形に捉える。地質名や機械強度を直接測る方式ではない。\par

\noindent\textbf{2：適切} アンテナ移動の波形を順に並べ記録断面を作る。横軸は測線位置、縦軸は通常往復走時であり、速度モデルから深度へ変換する。\par

\noindent\textbf{3：条件付きで適切} 中心周波数10MHz程度以上をGPRの目安とする。通常のGPRの周波数帯を示す分類目安である。電磁探査との厳密な物理境界ではない。\par

\noindent\textbf{4：不適切} 誘電率単位F・m、比誘電率は空気に対する比。誘電率はF/mであり、比誘電率の基準は真空である。空気の基準は通常の近似に過ぎない。\par

\noindent\textbf{5：条件付きで適切} 埋設物・空洞・地層・鉄筋・遺跡・水等を調査。各対象が電磁特性の差を持ち、十分な到達深度と分解能を得られる場合の用途である。すべての現場で適用可能とはしない。\par

\paragraph{再利用する考え方} 反射走時を深度と混同せず、帯域・速度・損失と探査目的を対応させる。

\paragraph{共通理論} \hyperref[theory:exploration:gpr]{第\ref*{theory:exploration:gpr}節}, \hyperref[theory:exploration:gprresolution]{第\ref*{theory:exploration:gprresolution}節}

\par 参考：\cite{epa_gpr}。

\par 難易度 3（校準初稿）、優先度未評価。

\subsection{令和5年度 問33}
\label{q:r05-reg-applied-33}

\noindent\textbf{原典：} sources/R05/17\_applied\_science.pdf、PDF物理12頁。

\textbf{公式正答：} 3（single）。

\par\textbf{検討未確定：} 分解能の表現は最小距離と識別能力を区別する。実深度のf依存を厳密な比例則とせず、力学特性との共通要因による間接相関を留保。

\paragraph{問題の要旨} GPRの周波数と分解能力・到達深度の向きを判断する。

\Needspace{4\baselineskip}\paragraph{確認した条件・記述（独自要約）}\mbox{}\par

\noindent\textbf{原典照合：} 29問21原典画像で否定、全五候補、数式・上付き単位・観測対象を照合。\par

\noindent\textbf{前提：} 低損失非磁性媒質でv≈c/√εr、波長λ=v/f、パルス帯域Bに対する距離分解はおよそv/(2B)である。導電率等が到達深度を制限する。\par

\paragraph{問われる核心} GPRの周波数と分解能力・到達深度の向きを判断する。

\paragraph{必要知識} 低損失非磁性媒質でv≈c/√εr、波長λ=v/f、パルス帯域Bに対する距離分解はおよそv/(2B)である。導電率等が到達深度を制限する。

\paragraph{考え方・正答を選ぶ理由（査読待ち初稿）} 3が不適切である。高周波では通常波長が短く分解能力が高まる一方、実地では探査深度が制限される。原典の分解能逆比例・深度比例はこの通常の向きを逆にしている。最小識別距離と能力の語を混同しない。

\Needspace{4\baselineskip}\paragraph{選択肢の検討（査読待ち初稿）}\mbox{}\par

\noindent\textbf{1：条件付きで適切} 電気境界を測るため力学特性との相関はほとんどない。強度や硬さを直接測らないことは正しい。ただし含水・空隙・劣化など共通要因を通じた間接相関まで全面否定はできない。\par

\noindent\textbf{2：適切} 直線埋設物へできる限り直角の測線。横切る測線で位置の変化と点状反射の走時を捉えやすくする。未知方位では複数方位を測る。\par

\noindent\textbf{3：不適切} 中心周波数と分解能は逆比例、探査深度は比例。分解できる最小距離は通常高周波ほど小さく、能力としての分解能は良くなる。深度も高周波ほど大きいという説明は通常の交換条件に反する。\par

\noindent\textbf{4：適切} profileにより構造、wide angleで電磁波速度を推定。走査波形を位置へ並べるprofileと、送受信距離を変えて走時から速度を求める構成を使い分ける。\par

\noindent\textbf{5：条件付きで適切} 中心周波数10MHz程度以上をGPRの目安とする。通常のGPRの周波数帯を示す分類目安である。電磁探査との厳密な物理境界ではない。\par

\paragraph{再利用する考え方} 反射走時を深度と混同せず、帯域・速度・損失と探査目的を対応させる。

\paragraph{共通理論} \hyperref[theory:exploration:gpr]{第\ref*{theory:exploration:gpr}節}, \hyperref[theory:exploration:gprresolution]{第\ref*{theory:exploration:gprresolution}節}

\par 参考：\cite{epa_gpr}。

\paragraph{関連問題} \hyperref[q:r02-reg-applied-33]{R02-REG-APPLIED-33}, \hyperref[q:r06-reg-applied-35]{R06-REG-APPLIED-35}

\par 難易度 3（校準初稿）、優先度未評価。

\subsection{令和6年度 問35}
\label{q:r06-reg-applied-35}

\noindent\textbf{原典：} sources/R06/17\_applied\_science.pdf、PDF物理12頁。

\textbf{公式正答：} 2（single）。

\par\textbf{検討未確定：} 分解能の最小距離/能力と探査深度の交換条件を区別し、10MHzを全方式の厳密境界にしない。

\paragraph{問題の要旨} GPRの電気境界と周波数・断面・用途を比較する。

\Needspace{4\baselineskip}\paragraph{確認した条件・記述（独自要約）}\mbox{}\par

\noindent\textbf{原典照合：} 29問21原典画像で否定、全五候補、数式・上付き単位・観測対象を照合。\par

\noindent\textbf{前提：} 低損失非磁性媒質でv≈c/√εr、波長λ=v/f、パルス帯域Bに対する距離分解はおよそv/(2B)である。導電率等が到達深度を制限する。\par

\paragraph{問われる核心} GPRの電気境界と周波数・断面・用途を比較する。

\paragraph{必要知識} 低損失非磁性媒質でv≈c/√εr、波長λ=v/f、パルス帯域Bに対する距離分解はおよそv/(2B)である。導電率等が到達深度を制限する。

\paragraph{考え方・正答を選ぶ理由（査読待ち初稿）} 2が不適切である。高周波は通常分解能力を良くする一方、深度を大きくする保証はない。帯域と媒質導電性・ノイズを含む到達条件から判断する。

\Needspace{4\baselineskip}\paragraph{選択肢の検討（査読待ち初稿）}\mbox{}\par

\noindent\textbf{1：適切} 誘電率・導電率の異なる境界からの反射を測る。電磁インピーダンスと減衰の違いを波形に捉える。地質名や機械強度を直接測る方式ではない。\par

\noindent\textbf{2：不適切} 中心周波数と分解能は逆比例、探査深度は比例。分解できる最小距離は通常高周波ほど小さく、能力としての分解能は良くなる。深度も高周波ほど大きいという説明は通常の交換条件に反する。\par

\noindent\textbf{3：条件付きで適切} 埋設物・空洞・地層・鉄筋・遺跡・水等を調査。各対象が電磁特性の差を持ち、十分な到達深度と分解能を得られる場合の用途である。すべての現場で適用可能とはしない。\par

\noindent\textbf{4：適切} アンテナ移動の波形を順に並べ記録断面を作る。横軸は測線位置、縦軸は通常往復走時であり、速度モデルから深度へ変換する。\par

\noindent\textbf{5：条件付きで適切} 中心周波数10MHz程度以上をGPRの目安とする。通常のGPRの周波数帯を示す分類目安である。電磁探査との厳密な物理境界ではない。\par

\paragraph{再利用する考え方} 反射走時を深度と混同せず、帯域・速度・損失と探査目的を対応させる。

\paragraph{共通理論} \hyperref[theory:exploration:gpr]{第\ref*{theory:exploration:gpr}節}, \hyperref[theory:exploration:gprresolution]{第\ref*{theory:exploration:gprresolution}節}

\par 参考：\cite{epa_gpr}。

\paragraph{関連問題} \hyperref[q:r05-reg-applied-33]{R05-REG-APPLIED-33}

\par 難易度 3（校準初稿）、優先度未評価。

\section{逆問題の正則化・最適化と離散化}

\subsection{平成30年度 問35}
\label{q:h30-reg-applied-35}

\noindent\textbf{原典：} sources/H30/17\_applied\_science.pdf、PDF物理12頁。

\textbf{公式正答：} 5（single）。

\par\textbf{検討未確定：} 設問のインバージョンに関係するという広い語には有限要素順モデルも含められる。公式5を算法の直接的役割の比較として保持し、FEMを逆解析から排除しない。

\paragraph{問題の要旨} 逆問題の最適化・正則化と有限要素離散化を区別する。

\Needspace{4\baselineskip}\paragraph{確認した条件・記述（独自要約）}\mbox{}\par

\noindent\textbf{原典照合：} 否定の対象、指数、周波数と周期、五つの候補の列順を原典画像で照合。\par

\noindent\textbf{前提：} モデル更新算法と順モデル計算の離散化を区別する。ただし順モデルは逆問題を構成する内部計算でもある。\par

\paragraph{問われる核心} 逆問題の最適化・正則化と有限要素離散化を区別する。

\paragraph{必要知識} モデル更新算法と順モデル計算の離散化を区別する。ただし順モデルは逆問題を構成する内部計算でもある。

\paragraph{考え方・正答を選ぶ理由（査読待ち初稿）} 公式5は、他四項目が最適化又は逆解の安定化を直接担うのに対し、有限要素法が主に支配方程式の離散化を担うという比較による。ただしFEMで順問題を計算する逆解析は広くある。インバージョンにFEMが関係しないと断定しない。

\Needspace{4\baselineskip}\paragraph{選択肢の検討（査読待ち初稿）}\mbox{}\par

\noindent\textbf{1：適切} マルカート法。非線形最小二乗の更新を減衰項で安定化するLevenberg–Marquardt型の算法である。\par

\noindent\textbf{2：適切} 遺伝的アルゴリズム。モデル候補の集団を評価・選択・変化させて目的関数を探索する最適化法である。大域的最良解を常に保証するわけではない。\par

\noindent\textbf{3：適切} 正則化法。データだけで決まりにくいモデルへ滑らかさや大きさ等の制約を加え、逆問題の不安定性を抑える。\par

\noindent\textbf{4：適切} 共役勾配法。線形化された最小二乗等の連立系又は非線形目的関数の探索に用いる。対称正定値系のCGと非線形CGを区別する。\par

\noindent\textbf{5：条件付きで適切} 有限要素法。直接の最適化・正則化の名称ではなく支配方程式の離散化であるため公式はこれを選ぶ。しかし順モデルや感度計算として逆解析へ用いることができ、無関係という絶対判断は誤りである。\par

\paragraph{再利用する考え方} 目的関数、モデル更新、物理場の数値計算の三つの役割を別々に問う。

\paragraph{共通理論} \hyperref[theory:exploration:inverse]{第\ref*{theory:exploration:inverse}節}, \hyperref[theory:fem:background]{第\ref*{theory:fem:background}節}

\par 参考：\cite{pygimli_inversion,mit_roy_fea1999}。

\paragraph{関連問題} \hyperref[q:r01-reg-applied-35]{R01-REG-APPLIED-35}

\par 難易度 3（校準初稿）、優先度未評価。

\subsection{令和元年度 問35}
\label{q:r01-reg-applied-35}

\noindent\textbf{原典：} sources/R01/17\_applied\_science.pdf、PDF物理12頁。

\textbf{公式正答：} 2（single）。

\paragraph{問題の要旨} フィルタ・差分・共役勾配・平滑化とラテラル検層を区別する。

\Needspace{4\baselineskip}\paragraph{確認した条件・記述（独自要約）}\mbox{}\par

\noindent\textbf{原典照合：} 否定の対象、指数、周波数と周期、五つの候補の列順を原典画像で照合。\par

\noindent\textbf{前提：} 数値処理又はモデル拘束と、孔内で物理量を測る配置を区別する。\par

\paragraph{問われる核心} フィルタ・差分・共役勾配・平滑化とラテラル検層を区別する。

\paragraph{必要知識} 数値処理又はモデル拘束と、孔内で物理量を測る配置を区別する。

\paragraph{考え方・正答を選ぶ理由（査読待ち初稿）} 2が不適切である。ラテラル法は電極配置を用いる電気検層の測定方式である。他四項目は信号処理、方程式の離散化、最適化、解への拘束を担う。

\Needspace{4\baselineskip}\paragraph{選択肢の検討（査読待ち初稿）}\mbox{}\par

\noindent\textbf{1：適切} リニアフィルタ。入力の和と定数倍に対して同じ関係を保つ信号処理であり、畳み込み等の数値計算で実装できる。\par

\noindent\textbf{2：不適切} ラテラル法。USGSの電気検層資料では電極間隔と配置で定義されるlateral logとして確認できる。数値解法の算法を指す語ではない。\par

\noindent\textbf{3：適切} 差分法。微分演算を格子値の差へ置き換えて支配方程式の近似解を求める数値法である。\par

\noindent\textbf{4：適切} 共役勾配法。大型の連立系や最小化を反復して解く数値算法である。\par

\noindent\textbf{5：適切} 平滑化拘束。隣接するモデル値の差などを罰することで、逆解の不要な振動を抑える数値的な条件である。\par

\paragraph{再利用する考え方} 測定の幾何配置とデータ取得後の計算を分ける。

\paragraph{共通理論} \hyperref[theory:exploration:inverse]{第\ref*{theory:exploration:inverse}節}, \hyperref[theory:exploration:seismicterms]{第\ref*{theory:exploration:seismicterms}節}

\par 参考：\cite{usgs_lateral_logging,pygimli_inversion,mit_signals2010}。

\paragraph{関連問題} \hyperref[q:h30-reg-applied-35]{H30-REG-APPLIED-35}

\par 難易度 2（校準初稿）、優先度未評価。

\section{人工分極と自然電位の信号機構}

\subsection{平成28年度 問22}
\label{q:h28-reg-applied-22}

\noindent\textbf{原典：} sources/H28/17\_applied\_science.pdf、PDF物理9頁。

\textbf{公式正答：} 1（single）。

\par\textbf{検討未確定：} SPの負異常は典型的鉱化機構と基準点による。AMTの深部は相対表現であり一定深度を保証しない。

\paragraph{問題の要旨} SP鉱体異常の符号と電気・弾性波探査を比較する。

\Needspace{4\baselineskip}\paragraph{確認した条件・記述（独自要約）}\mbox{}\par

\noindent\textbf{原典照合：} 29問21原典画像で否定、全五候補、数式・上付き単位・観測対象を照合。\par

\noindent\textbf{前提：} SPは基準点に対する電位差であり、典型的な硫化鉱体上の負異常を考える。\par

\paragraph{問われる核心} SP鉱体異常の符号と電気・弾性波探査を比較する。

\paragraph{必要知識} SPは基準点に対する電位差であり、典型的な硫化鉱体上の負異常を考える。

\paragraph{考え方・正答を選ぶ理由（査読待ち初稿）} 1が不適切である。典型的な鉱化電池の鉱体上部では基準点に対し負の電位異常を利用する。全鉱体・全自然電位源の符号を固定する規則ではない。

\Needspace{4\baselineskip}\paragraph{選択肢の検討（査読待ち初稿）}\mbox{}\par

\noindent\textbf{1：不適切} 鉱体があると地表電位が周囲より高い。典型的な硫化鉱体探査は負の電位異常を利用し、周囲より高いという説明は逆である。酸化還元や流動の機構を特定する。\par

\noindent\textbf{2：適切} 比抵抗垂直探査は電極間隔を広げ深さ分布を推定。配置間隔を変えて深度感度を変え、水平多層等のモデルを推定する。\par

\noindent\textbf{3：条件付きで適切} AMTは自然磁場と地電流で深部比抵抗を推定。電場磁場の周波数応答で表層レーダより深い領域へ感度を持つ。実深度は周波数・抵抗率・信号品質による。\par

\noindent\textbf{4：適切} 弾性波tomographyは周囲の起振受振で伝播特性を画像化。多数経路走時などの逆解析で速度分布を求め、覆いの不足は分解能を制限する。\par

\noindent\textbf{5：適切} ダウンホールで地表起振・孔内受振の初動走時を測る。孔壁に結合した受振器の位置と走時から区間速度を推定する。経路と屈折の補正が必要である。\par

\paragraph{再利用する考え方} 自然電位の符号を単一の地下物体の存在へ直結させない。

\paragraph{共通理論} \hyperref[theory:exploration:ipsp]{第\ref*{theory:exploration:ipsp}節}, \hyperref[theory:exploration:dcarrays]{第\ref*{theory:exploration:dcarrays}節}, \hyperref[theory:exploration:borehole]{第\ref*{theory:exploration:borehole}節}

\par 参考：\cite{epa_sp,epa_resistivity,usgs_mt_metadata,geebec_ps}。

\par 難易度 3（校準初稿）、優先度未評価。

\subsection{令和3年度 問32}
\label{q:r03-reg-applied-32}

\noindent\textbf{原典：} sources/R03/17\_applied\_science.pdf、PDF物理12頁。

\textbf{公式正答：} 4（single）。

\par\textbf{検討未確定：} 4にはモデル依存の深度推定が可能という反例がある。公式4を保持するが「SPで深度解析はできない」という誤解を教えない。3の特定深さも有限感度体積の近似。

\paragraph{問題の要旨} 比抵抗の深度感度と自然電位法の逆解析限界を比較する。

\Needspace{4\baselineskip}\paragraph{確認した条件・記述（独自要約）}\mbox{}\par

\noindent\textbf{原典照合：} 29問21原典画像で否定、全五候補、数式・上付き単位・観測対象を照合。\par

\noindent\textbf{前提：} SPは自然電位差の分布を直接測り、源の深度には機構・形状・抵抗率のモデルが必要である。\par

\paragraph{問われる核心} 比抵抗の深度感度と自然電位法の逆解析限界を比較する。

\paragraph{必要知識} SPは自然電位差の分布を直接測り、源の深度には機構・形状・抵抗率のモデルが必要である。

\paragraph{考え方・正答を選ぶ理由（査読待ち初稿）} 公式の4を保持する。SP観測から地下水や鉱体の深さを簡単に直接求めるという読みなら不適切である。一方、半幅や源モデル、他観測との逆解析により深度推定の例はあり、深度を探査・解析するという広い原典表現自体が全面的に不可能とは言えない。疑義を査読対象にする。

\Needspace{4\baselineskip}\paragraph{選択肢の検討（査読待ち初稿）}\mbox{}\par

\noindent\textbf{1：適切} IPは印加電流の分極応答、硫化鉱物探査に用いる。金属・電解質界面等の分極で電位の時間減衰や周波数位相が生じ、鉱床や環境調査に使われる。\par

\noindent\textbf{2：適切} 垂直探査は水平多層の鉛直比抵抗を推定。電極間隔を変えて深度感度を変化させ、WennerやSchlumberger等の配置で層モデルを推定する。\par

\noindent\textbf{3：条件付きで適切} 水平探査は特定深さの水平比抵抗変化を調べる。固定間隔等で感度深度をおおむねそろえ横方向を調べる。ただし有限体積の重み付き応答で一点の厳密深度だけを測らない。\par

\noindent\textbf{4：条件付きで適切} SPから地下水や鉱床の深さを探査解析し簡単安全に測る。自然電位分布の簡単な受動測定はできるが、深度は非一意で源モデルや他データが必要である。深度解析が全く不可能とも言えず、この候補の否定範囲は未確定である。\par

\noindent\textbf{5：適切} 多電極の直流電位から比抵抗tomography。多数の注入・受信対の応答を順モデルへ合わせて比抵抗分布を推定する。\par

\paragraph{再利用する考え方} 公式の選択と、科学的に否定できる命題の範囲を分けて残す。

\paragraph{共通理論} \hyperref[theory:exploration:ipsp]{第\ref*{theory:exploration:ipsp}節}, \hyperref[theory:exploration:dcarrays]{第\ref*{theory:exploration:dcarrays}節}

\par 参考：\cite{epa_sp,epa_ip,epa_resistivity}。

\par 難易度 3（校準初稿）、優先度未評価。

\section{探査の観測量と物理手法の対応}

\subsection{令和元年度再試験 問32}
\label{q:r01-ret-applied-32}

\noindent\textbf{原典：} sources/R01\_retest/17\_applied\_science.pdf、PDF物理11頁。

\textbf{公式正答：} 4（single）。

\paragraph{問題の要旨} 地震探査を重力・磁気・電気・放射能探査から観測量で選ぶ。

\Needspace{4\baselineskip}\paragraph{確認した条件・記述（独自要約）}\mbox{}\par

\noindent\textbf{原典照合：} 五候補の順序、最適切と不適切の方向、波種、深度目安、曲線名、最長と最短を原典画像で照合。\par

\noindent\textbf{前提：} 最も適切な説明を選ぶ。弾性波の人工震源と地下からの応答を対象とする。\par

\paragraph{問われる核心} 地震探査を重力・磁気・電気・放射能探査から観測量で選ぶ。

\paragraph{必要知識} 最も適切な説明を選ぶ。弾性波の人工震源と地下からの応答を対象とする。

\paragraph{考え方・正答を選ぶ理由（査読待ち初稿）} 4が適切である。爆薬・バイブレータ・エアガンで弾性波を送り、反射や屈折から構造を推定する。地質を調べるという目的が同じでも、他の四つは測定する物理量が異なる。

\Needspace{4\baselineskip}\paragraph{選択肢の検討（査読待ち初稿）}\mbox{}\par

\noindent\textbf{1：不適切} ブーゲー異常で基盤・褶曲・断層を考察する。密度差による重力探査であり弾性波の走時・波形を直接測る地震探査ではない。\par

\noindent\textbf{2：不適切} 磁性鉱床・堆積盆地・地殻熱構造を調べる。磁気異常や磁化の性質を使う探査の説明である。\par

\noindent\textbf{3：不適切} 電気的構造と深度別の比抵抗を知る。電気・電磁応答から導電率又は比抵抗を推定する説明である。\par

\noindent\textbf{4：適切} 人工起振で弾性波を地下へ送り資源・土木調査に用いる。地震探査の震源と観測量に対応する。海上のエアガンなど場所ごとに装置は異なる。\par

\noindent\textbf{5：不適切} U・Th・K等の放射性元素を探査する。放射能探査の観測対象であり弾性波とは異なる。天然核種のガンマ線等を利用する。\par

\paragraph{再利用する考え方} 手法の目的が同じ場合も入力・出力の物理量で区別する。

\paragraph{共通理論} \hyperref[theory:exploration:observables]{第\ref*{theory:exploration:observables}節}

\par 参考：\cite{epa_seismic_reflection,segj_integrated2017}。

\par 難易度 2（校準初稿）、優先度未評価。

\section{移動平均の雑音低減と周波数零点}

\subsection{平成28年度 問20}
\label{q:h28-reg-applied-20}

\noindent\textbf{原典：} sources/H28/17\_applied\_science.pdf、PDF物理8頁。

\textbf{公式正答：} 2（single）。

\par\textbf{検討未確定：} 2には周期と窓長を合わせれば除去できる具体的反例がある。公式2を保持するが、数学的に周期雑音除去が不可能とは説明しない。3の時不変性と5の帯域・端点条件も省略される。

\paragraph{問題の要旨} 移動平均と周期雑音除去の関係を検討する。

\Needspace{4\baselineskip}\paragraph{確認した条件・記述（独自要約）}\mbox{}\par

\noindent\textbf{原典照合：} 否定の対象、指数、周波数と周期、五つの候補の列順を原典画像で照合。\par

\noindent\textbf{前提：} 通常の移動平均は平滑化を行うが、窓長に応じた周波数応答の零点を持つ。\par

\paragraph{問われる核心} 移動平均と周期雑音除去の関係を検討する。

\paragraph{必要知識} 通常の移動平均は平滑化を行うが、窓長に応じた周波数応答の零点を持つ。

\paragraph{考え方・正答を選ぶ理由（査読待ち初稿）} 公式は2であり、一般的な移動平均を周期雑音専用の選択除去と説明する意図を不適切と読んでいる。しかし周期と窓長が一致する正弦波は移動平均で厳密に消せる。したがって周期雑音を除去できないという一般則へ書き換えず、原典の曖昧さを残す。

\Needspace{4\baselineskip}\paragraph{選択肢の検討（査読待ち初稿）}\mbox{}\par

\noindent\textbf{1：適切} 空間フィルタを二次元画像のエッジ・平滑化・差分に用いる。二次元データに対し空間方向の重み付けや差分を行う。一次元・三次元にも拡張でき、二次元しか扱えないという意味ではない。\par

\noindent\textbf{2：条件付きで適切} 移動平均は時系列の周期的雑音を除去するために用いる。M点等重み平均の周波数応答はsin(Mω/2)/(M sin(ω/2))に位相因子を掛けたものである。ω=2πk/Mの非DC成分は消えるため、窓長を周期へ合わせる反例がある。一般に任意の周期雑音が消えるとは限らない。\par

\noindent\textbf{3：条件付きで適切} 線形システムのインパルス応答を重ねて任意入力への応答を求める。線形性により各時刻のインパルスへの応答を重ねることはできる。単一の応答を時間移動する畳み込み表現には時不変性も必要であり、時変系には二時刻の応答を用いる。\par

\noindent\textbf{4：条件付きで適切} 低域通過で高周波を減衰させ低周波を通す。低域通過の基本的な役割である。無減衰という文言は理想通過帯域のモデルであり、実装の位相・リップルは別に確認する。\par

\noindent\textbf{5：条件付きで適切} 一周期二標本未満では元データを再構成できない。一般の帯域制限信号を一意に再構成するにはスペクトルの重なりを避ける標本化が必要である。特別な事前情報を持つ信号まで一律に不可能とする定理ではなく、最高周波数と標本化間隔を明示して読む。\par

\paragraph{再利用する考え方} 移動平均の効果を目的名で決めず、入力周波数と窓長の応答を計算する。

\paragraph{共通理論} \hyperref[theory:exploration:moving]{第\ref*{theory:exploration:moving}節}, \hyperref[theory:exploration:lti]{第\ref*{theory:exploration:lti}節}, \hyperref[theory:exploration:sampling]{第\ref*{theory:exploration:sampling}節}

\par 参考：\cite{mit_signals2010,mit_moving2012,scipy_notch,scipy_fft,usgs_processing2000}。

\paragraph{関連問題} \hyperref[q:h23-reg-applied-28]{H23-REG-APPLIED-28}

\par 難易度 3（校準初稿）、優先度未評価。

\section{物理探査の複数手法と幾何・時間}

\subsection{平成25年度 問33}
\label{q:h25-reg-applied-33}

\noindent\textbf{原典：} sources/H25/17\_applied\_science.pdf、PDF物理12頁。

\textbf{公式正答：} 5（single）。

\par\textbf{検討未確定：} 地形起伏があれば反射法は絶対不可とはしない。速度比だけから層厚精度を保証せず、隠れ層と観測条件を保持。

\paragraph{問題の要旨} 重力・速度検層・屈折法と起伏地の反射法の制約を比較する。

\Needspace{4\baselineskip}\paragraph{確認した条件・記述（独自要約）}\mbox{}\par

\noindent\textbf{原典照合：} 29問21原典画像で否定、全五候補、数式・上付き単位・観測対象を照合。\par

\noindent\textbf{前提：} 地形起伏は走時・観測配置へ影響し、浅層反射法では静補正や起振受振条件を整える。\par

\paragraph{問われる核心} 重力・速度検層・屈折法と起伏地の反射法の制約を比較する。

\paragraph{必要知識} 地形起伏は走時・観測配置へ影響し、浅層反射法では静補正や起振受振条件を整える。

\paragraph{考え方・正答を選ぶ理由（査読待ち初稿）} 5が不適切である。急峻で凹凸の多い地形を浅層反射法の一般的な適用対象とする説明は、走時補正・起振結合・配置の難しさを落としている。ただし適切な補正と計測により適用可能な例を否定しない。

\Needspace{4\baselineskip}\paragraph{選択肢の検討（査読待ち初稿）}\mbox{}\par

\noindent\textbf{1：適切} 重力分布から地下密度構造を推定。重力異常に対して密度と形状のモデルを作り逆解析する。\par

\noindent\textbf{2：適切} 単孔で波速の深さ分布を求める速度検層。地表又は孔内起振と孔内受振の走時を用いる。\par

\noindent\textbf{3：条件付きで適切} 屈折波から地下速度構造を解析。通常は速い下層の境界を進み戻る臨界屈折波の初動等を使う。任意の低速度層まで検出できるわけではない。\par

\noindent\textbf{4：条件付きで適切} 下方増速で速度比大きいほど分割と厚さ精度が良い。層走時勾配の差が明瞭になり識別しやすい傾向である。隠れ層、薄層、配置、ノイズによる精度制限もある。\par

\noindent\textbf{5：不適切} 浅層反射法は急峻で凹凸の多い場所へ適用される。通常の適用上は地形起伏が観測と補正の障害となり、適性を積極的に示す一般記述は不適切である。実施そのものが絶対不可能とはしない。\par

\paragraph{再利用する考え方} 適用の一般的有利条件と、補正して実施できる可能性を分ける。

\paragraph{共通理論} \hyperref[theory:exploration:refraction]{第\ref*{theory:exploration:refraction}節}, \hyperref[theory:exploration:reflectiongeometry]{第\ref*{theory:exploration:reflectiongeometry}節}

\par 参考：\cite{epa_seismic_reflection,epa_seismic_refraction,usgs_static2004,geebec_ps}。

\par 難易度 3（校準初稿）、優先度未評価。

\subsection{平成26年度 問19}
\label{q:h26-reg-applied-19}

\noindent\textbf{原典：} sources/H26/17\_applied\_science.pdf、PDF物理8頁。

\textbf{公式正答：} 5（single）。

\par\textbf{検討未確定：} 含水比・飽和度・間隙流体を区別し、P波速度の全材料での単調変化を主張しない。P波初動の利点は通常の方式に限定する。

\paragraph{問題の要旨} 各物理探査の観測量とP波速度の含水依存を比較する。

\Needspace{4\baselineskip}\paragraph{確認した条件・記述（独自要約）}\mbox{}\par

\noindent\textbf{原典照合：} 29問21原典画像で否定、全五候補、数式・上付き単位・観測対象を照合。\par

\noindent\textbf{前提：} P波速度は弾性率と密度で決まり、飽和や孔隙流体は両者と接触状態を変え得る。\par

\paragraph{問われる核心} 各物理探査の観測量とP波速度の含水依存を比較する。

\paragraph{必要知識} P波速度は弾性率と密度で決まり、飽和や孔隙流体は両者と接触状態を変え得る。

\paragraph{考え方・正答を選ぶ理由（査読待ち初稿）} 5が不適切である。P波速度は飽和・含水状態で変わり得る。Vp=√((K+4µ/3)/ρ)の剛性と密度が含水で変化するため、一般に不変とすることはできない。単純な含水比だけから全地盤の単調変化は予測しない。

\Needspace{4\baselineskip}\paragraph{選択肢の検討（査読待ち初稿）}\mbox{}\par

\noindent\textbf{1：適切} 重力が密度・深度・形状を反映する。密度差の体積積分が重力へ寄与し、同じ密度でも深度と形状で観測異常が変わる。\par

\noindent\textbf{2：条件付きで適切} 通常の屈折法は地表でP波初動を測る。地表から広い対象を調査しP初動を使う構成が一般的である。S波方式や隠れ層の限界もある。\par

\noindent\textbf{3：適切} ダウンホールは孔内速度測定方式。地表起振・孔内受振の走時を使う。\par

\noindent\textbf{4：適切} ガンマ線で深度方向の密度検層。線源と受信器を降下させ孔壁周囲の校正かさ密度を測る。\par

\noindent\textbf{5：不適切} P波速度は一般に含水比によらず不変。孔隙流体による有効弾性率と密度・接触の変化があるため、含水依存を一般に除外できない。\par

\paragraph{再利用する考え方} 一つの観測物性を水分と切り離さず、波速と密度の感度を考える。

\paragraph{共通理論} \hyperref[theory:exploration:borehole]{第\ref*{theory:exploration:borehole}節}, \hyperref[theory:exploration:waveproperties]{第\ref*{theory:exploration:waveproperties}節}

\par 参考：\cite{epa_gamma_density,geebec_ps,epa_seismic_refraction}。

\par 難易度 3（校準初稿）、優先度未評価。

\subsection{平成29年度 問33}
\label{q:h29-reg-applied-33}

\noindent\textbf{原典：} sources/H29/17\_applied\_science.pdf、PDF物理11頁。

\textbf{公式正答：} 3（single）。

\par\textbf{検討未確定：} 磁気単独の非一意性を形状推定の絶対不可能性にしない。GPRの10MHzは典型的分類目安である。

\paragraph{問題の要旨} ばね式重力計・MT・GPR・海上音波の観測を比較する。

\Needspace{4\baselineskip}\paragraph{確認した条件・記述（独自要約）}\mbox{}\par

\noindent\textbf{原典照合：} 29問21原典画像で否定、全五候補、数式・上付き単位・観測対象を照合。\par

\noindent\textbf{前提：} ばね伸びの比較による重力差と、絶対加速度の時間長さ測定を分ける。\par

\paragraph{問われる核心} ばね式重力計・MT・GPR・海上音波の観測を比較する。

\paragraph{必要知識} ばね伸びの比較による重力差と、絶対加速度の時間長さ測定を分ける。

\paragraph{考え方・正答を選ぶ理由（査読待ち初稿）} 3が不適切である。ばねの伸びを使う多くの重力探査用装置は相対重力計である。絶対計は自由落下等の運動を長さと時間で測る。

\Needspace{4\baselineskip}\paragraph{選択肢の検討（査読待ち初稿）}\mbox{}\par

\noindent\textbf{1：条件付きで適切} 磁気探査で鉄の概略位置は測れるが細形状種類は判別できない。磁化と形状・深度の非一意性があるため磁気だけで製品種を一意に識別しない。追加観測で形状推定ができないという絶対禁止ではない。\par

\noindent\textbf{2：適切} MTは自然磁場と地電流から深部比抵抗を探査。電磁応答のインピーダンスと周波数から比抵抗モデルを求める。\par

\noindent\textbf{3：不適切} ばね伸びの装置は絶対重力計。標準点や基準との重力差を測る相対方式であり絶対方式と取り違える。\par

\noindent\textbf{4：条件付きで適切} GPRは透過反射屈折を使い中心周波数10MHz程度以上。通常のGPRの高周波帯の説明である。10MHzは厳密な全装置・全媒質の境界ではない。\par

\noindent\textbf{5：適切} 海上ストリーマ受振器はハイドロフォン。水中の音圧変動を受信する素子を曳航ケーブルへ多数配置する。\par

\paragraph{再利用する考え方} センサが測る量と、逆解析で推定する地下構造を分ける。

\paragraph{共通理論} \hyperref[theory:exploration:observables]{第\ref*{theory:exploration:observables}節}, \hyperref[theory:exploration:gpr]{第\ref*{theory:exploration:gpr}節}

\par 参考：\cite{gsi_gravity_definitions,usgs_mt_metadata,epa_gpr,usgs_marine_seismic}。

\par 難易度 3（校準初稿）、優先度未評価。

\subsection{令和元年度 問32}
\label{q:r01-reg-applied-32}

\noindent\textbf{原典：} sources/R01/17\_applied\_science.pdf、PDF物理11頁。

\textbf{公式正答：} 5（single）。

\par\textbf{検討未確定：} JISの規格本文と出題時/現行版差は未照合。重力ののみは処理目的と読み、残差誤差がないとはしない。

\paragraph{問題の要旨} 探査の時間差・環境振動と表面波のS波速度を比較する。

\Needspace{4\baselineskip}\paragraph{確認した条件・記述（独自要約）}\mbox{}\par

\noindent\textbf{原典照合：} 29問21原典画像で否定、全五候補、数式・上付き単位・観測対象を照合。\par

\noindent\textbf{前提：} Rayleigh波の分散は主にVs分布へ関係し、時間差調査は同じ測定条件を比較する。\par

\paragraph{問われる核心} 探査の時間差・環境振動と表面波のS波速度を比較する。

\paragraph{必要知識} Rayleigh波の分散は主にVs分布へ関係し、時間差調査は同じ測定条件を比較する。

\paragraph{考え方・正答を選ぶ理由（査読待ち初稿）} 5が不適切である。表面波探査は分散曲線から主にS波速度構造を推定する。Rayleigh波にP成分もあるが、主推定対象をP波速度とする説明は違う。

\Needspace{4\baselineskip}\paragraph{選択肢の検討（査読待ち初稿）}\mbox{}\par

\noindent\textbf{1：条件付きで適切} 重力補正で密度構造による異常のみ抽出。潮汐・標高・地形等の寄与を補正し地下密度異常を解釈する。補正誤差や地域背景が完全に消えるという保証ではない。\par

\noindent\textbf{2：適切} GPRは数m浅部の埋設物空洞遺跡に多用。典型的な浅部の利用であり、媒質・装置によってより深い例もある。\par

\noindent\textbf{3：適切} time lapseは繰返し探査差から地下変化を推定。同じ条件の反復観測で変化を見る。装置ドリフトや位置・気象差も区別する。\par

\noindent\textbf{4：条件付きで適切} 環境振動は感覚補正の振動レベルとJIS計器を使う。人の周波数感度を補正した加速度レベルで評価する。出題時と現在の規格版、計器性能の全文は追加照合対象である。\par

\noindent\textbf{5：不適切} 表面波分散から主にP波速度構造を推定。通常の表面波法の主推定物性はS波速度Vsである。P波速度は感度を持つ他パラメータとして扱う。\par

\paragraph{再利用する考え方} 観測波の粒子運動と、その逆解析が主に拘束する物性を区別する。

\paragraph{共通理論} \hyperref[theory:exploration:surfacewaves]{第\ref*{theory:exploration:surfacewaves}節}, \hyperref[theory:exploration:multimethod]{第\ref*{theory:exploration:multimethod}節}

\par 参考：\cite{epa_surfacewaves,epa_gpr,rion_vibration2016,gsi_gravity_definitions}。

\paragraph{関連問題} \hyperref[q:r05-reg-applied-35]{R05-REG-APPLIED-35}

\par 難易度 3（校準初稿）、優先度未評価。

\subsection{令和5年度 問35}
\label{q:r05-reg-applied-35}

\noindent\textbf{原典：} sources/R05/17\_applied\_science.pdf、PDF物理13頁。

\textbf{公式正答：} 2（single）。

\par\textbf{検討未確定：} 磁気数m・表面波20mは設定依存、熱探査は温度差と季節・地況による。1mを水脈の上限深度としない。

\paragraph{問題の要旨} 表面波のVsと磁気・tomography・地温の感度を比較する。

\Needspace{4\baselineskip}\paragraph{確認した条件・記述（独自要約）}\mbox{}\par

\noindent\textbf{原典照合：} 29問21原典画像で否定、全五候補、数式・上付き単位・観測対象を照合。\par

\noindent\textbf{前提：} 観測の空間覆い、密度差、温度差と、Rayleigh分散の主Vs感度を考える。\par

\paragraph{問われる核心} 表面波のVsと磁気・tomography・地温の感度を比較する。

\paragraph{必要知識} 観測の空間覆い、密度差、温度差と、Rayleigh分散の主Vs感度を考える。

\paragraph{考え方・正答を選ぶ理由（査読待ち初稿）} 2が不適切である。Rayleigh波分散から通常主にS波速度分布を推定する。20mは典型装置・設定の深度であり、全表面波法の深度限界ではない。

\Needspace{4\baselineskip}\paragraph{選択肢の検討（査読待ち初稿）}\mbox{}\par

\noindent\textbf{1：条件付きで適切} 地上磁気探査で水平分解能数mの浅部調査。細かい測点と近い対象で詳細調査が可能である。実分解能は深度、磁化、観測配置、ノイズによる。\par

\noindent\textbf{2：不適切} Rayleigh波から20m程度までP波速度分布。主推定対象はS波速度Vsである。深度の数値も測定波長・配置に依存する。\par

\noindent\textbf{3：適切} 地下tomographyの構造像は測定幾何に依存。全方向から測れず、覆いのない方向・領域の分解は制限される。\par

\noindent\textbf{4：条件付きで適切} 密度差が大きい対象は重力で検出しやすい。同じ形状・深度・ノイズ条件なら密度差が重力異常を増す。深い小対象との無条件比較ではない。\par

\noindent\textbf{5：条件付きで適切} 1m深地温で地熱温泉や地すべり地下水を調査。流動水と平常地温の温度差を利用する。測温1mと地下水源の深度は別で、適期や地表環境の補正が必要である。\par

\paragraph{再利用する考え方} 典型的数値を固定限界とせず、主観測量・感度・配置をそろえて比較する。

\paragraph{共通理論} \hyperref[theory:exploration:surfacewaves]{第\ref*{theory:exploration:surfacewaves}節}, \hyperref[theory:exploration:multimethod]{第\ref*{theory:exploration:multimethod}節}

\par 参考：\cite{epa_surfacewaves,gsi_gravity_definitions,naro_thermal2008,landslide_thermal2007}。

\paragraph{関連問題} \hyperref[q:r01-reg-applied-32]{R01-REG-APPLIED-32}

\par 難易度 3（校準初稿）、優先度未評価。

\section{CSAMTの表皮深度・近接場・平面波}

\subsection{平成25年度 問35}
\label{q:h25-reg-applied-35}

\noindent\textbf{原典：} sources/H25/17\_applied\_science.pdf、PDF物理13頁。

\textbf{公式正答：} 2（single）。

\par\textbf{検討未確定：} 3の距離によらずという全称にも疑義がある。CSAMT遠方条件の明確な欠落を2の根拠とするが、TDEMの無条件深度増大を承認しない。VLF局と周波数は方式の典型である。

\paragraph{問題の要旨} MT・CSAMT・TDEM・Slingram・VLFの深度条件を比較する。

\Needspace{4\baselineskip}\paragraph{確認した条件・記述（独自要約）}\mbox{}\par

\noindent\textbf{原典照合：} 29問21原典画像で否定、全五候補、数式・上付き単位・観測対象を照合。\par

\noindent\textbf{前提：} CSAMTの平面波近似には送受信距離が表皮深度より十分長い条件が必要である。\par

\paragraph{問われる核心} MT・CSAMT・TDEM・Slingram・VLFの深度条件を比較する。

\paragraph{必要知識} CSAMTの平面波近似には送受信距離が表皮深度より十分長い条件が必要である。

\paragraph{考え方・正答を選ぶ理由（査読待ち初稿）} 2が不適切である。δ=√(2ρ/(ωµ))は低周波ほど増す。固定送受信距離で周波数を下げ続けると遠方条件が破れ、近接場となるため距離によらず深く測れるとは言えない。3も装置規模・時間・ノイズから独立ではなく、その表現の疑義を保持する。

\Needspace{4\baselineskip}\paragraph{選択肢の検討（査読待ち初稿）}\mbox{}\par

\noindent\textbf{1：条件付きで適切} MTは自然電磁場を使い低周波で深部を探る。人工源を要せず低周波に深部感度があるが、信号強度と電気構造が推定を制限する。\par

\noindent\textbf{2：不適切} CSAMTは距離によらず低周波化で深度増大。有限送信源の近接場と平面波条件を無視する。固定距離では低周波ほどδが大きくなる。\par

\noindent\textbf{3：条件付きで適切} TDEMは遮断後応答を測り距離によらず深度増大。遮断後の時間応答を使いcentral loopでも深部に感度を持つ。ただし寸法・送信強度・測定時間・ノイズを無視した深度保証はできない。\par

\noindent\textbf{4：適切} Slingramは小型送受信コイルで浅部を迅速測定。周波数領域の二次磁場の振幅・位相を使い、浅部導電性分布を機動的に調べる。\par

\noindent\textbf{5：適切} VLFは20kHz付近、方向に注意して浅部探査。遠方送信局の電磁波と導電性構造の応答を使い、送信局に対する構造方位が感度へ影響する。\par

\paragraph{再利用する考え方} 深さを周波数だけから決めず、有限源と電磁拡散の条件を確認する。

\paragraph{共通理論} \hyperref[theory:exploration:nearfield]{第\ref*{theory:exploration:nearfield}節}, \hyperref[theory:exploration:fdtd]{第\ref*{theory:exploration:fdtd}節}

\par 参考：\cite{usgs_csamt2012,epa_tdem_archive,epa_fdem,usgs_mt_metadata}。

\par 難易度 3（校準初稿）、優先度未評価。

\subsection{平成26年度 問22}
\label{q:h26-reg-applied-22}

\noindent\textbf{原典：} sources/H26/17\_applied\_science.pdf、PDF物理9頁。

\textbf{公式正答：} 3（single）。

\par\textbf{検討未確定：} Archieの骨格・表面伝導条件、TEM一次場遮断の過渡とstatic shiftの比較を留保。生じにくいという原典限定を保持。

\paragraph{問題の要旨} Archie・電磁結合・CSAMT近接場・VLF・TEMを比較する。

\Needspace{4\baselineskip}\paragraph{確認した条件・記述（独自要約）}\mbox{}\par

\noindent\textbf{原典照合：} 29問21原典画像で否定、全五候補、数式・上付き単位・観測対象を照合。\par

\noindent\textbf{前提：} 遠方条件の指標は送受信距離/δであり、同一周波数では高比抵抗ほどδが増す。\par

\paragraph{問われる核心} Archie・電磁結合・CSAMT近接場・VLF・TEMを比較する。

\paragraph{必要知識} 遠方条件の指標は送受信距離/δであり、同一周波数では高比抵抗ほどδが増す。

\paragraph{考え方・正答を選ぶ理由（査読待ち初稿）} 3が不適切である。間隔が小さいほど近接場になり、同じ周波数・間隔では高比抵抗ほどδが大きく平面波条件が破れやすい。原典の低比抵抗を高比抵抗へ書き換えず、誤記述として保存する。

\Needspace{4\baselineskip}\paragraph{選択肢の検討（査読待ち初稿）}\mbox{}\par

\noindent\textbf{1：条件付きで適切} Archieは孔隙水だけの電流を仮定する。絶縁性骨格・表面伝導が無視できる孔隙ネットワークを想定する。低抵抗流体はその条件を満たしやすいが、有効孔隙率が大きければ十分とは限らない。\par

\noindent\textbf{2：適切} 電磁カップリングがIP誤差になる。送受信ケーブル間の誘導応答は分極による位相・減衰応答と重なり誤差へ寄与する。\par

\noindent\textbf{3：不適切} CSAMT近接場は低比抵抗で生じやすい。δは√ρに比例するため同一条件では高比抵抗の方が有限距離による近接場になりやすい。\par

\noindent\textbf{4：適切} VLFは3〜30kHz、VLF-EMとVLF-MT。磁場変化を測る方式と電場磁場比を用いる方式を区別する周波数帯・分類である。\par

\noindent\textbf{5：条件付きで適切} TEMは時間応答、一次磁場なしで測りstatic shiftが生じにくい。送信遮断後の二次応答を観測し、MT電場の局所歪みとは異なる。turn-off過程・送信器残留応答などは別途扱い、生じにくいを完全ゼロへ強めない。\par

\paragraph{再利用する考え方} 近接場は高/低抵抗を記憶せずδと送受信距離から判断する。

\paragraph{共通理論} \hyperref[theory:exploration:nearfield]{第\ref*{theory:exploration:nearfield}節}, \hyperref[theory:exploration:archie]{第\ref*{theory:exploration:archie}節}, \hyperref[theory:exploration:fdtd]{第\ref*{theory:exploration:fdtd}節}

\par 参考：\cite{usgs_csamt2012,epa_ip,epa_tdem_archive,epa_borehole_resistivity}。

\paragraph{関連問題} \hyperref[q:h29-reg-applied-32]{H29-REG-APPLIED-32}

\par 難易度 3（校準初稿）、優先度未評価。

\subsection{平成28年度 問23}
\label{q:h28-reg-applied-23}

\noindent\textbf{原典：} sources/H28/17\_applied\_science.pdf、PDF物理9頁。

\textbf{公式正答：} 3（single）。

\par\textbf{検討未確定：} 原典は真空の誘電率をµと定義していることを画像確認。粘土の表面伝導と既存VLF局の条件を保持。

\paragraph{問題の要旨} Wenner式・ArchieとMT式の透磁率・誘電率を照合する。

\Needspace{4\baselineskip}\paragraph{確認した条件・記述（独自要約）}\mbox{}\par

\noindent\textbf{原典照合：} 29問21原典画像で否定、全五候補、数式・上付き単位・観測対象を照合。\par

\noindent\textbf{前提：} 見掛比抵抗ρa=\allowbreak{}2πaV/Iとρa=|Ex/Hy|\textsuperscript{2}/(ωµ)を使い、µは透磁率とする。\par

\paragraph{問われる核心} Wenner式・ArchieとMT式の透磁率・誘電率を照合する。

\paragraph{必要知識} 見掛比抵抗ρa=\allowbreak{}2πaV/Iとρa=|Ex/Hy|\textsuperscript{2}/(ωµ)を使い、µは透磁率とする。

\paragraph{考え方・正答を選ぶ理由（査読待ち初稿）} 3が不適切である。式のµは真空の透磁率µ0であって誘電率ε0ではない。式の形が正しくても原典の変数定義が誤る。Wenner式とMT式の次元を独立に照合する。

\Needspace{4\baselineskip}\paragraph{選択肢の検討（査読待ち初稿）}\mbox{}\par

\noindent\textbf{1：適切} Wennerのρa=\allowbreak{}2πaV/I。等間隔aの四電極を均質半空間の点電極として配置した幾何係数2πaである。\par

\noindent\textbf{2：条件付きで適切} 粘土の含水岩はArchie予測より低比抵抗。粘土表面の追加導電がある場合、孔隙水のみのモデルより伝導が増える。全試料の固定差ではない。\par

\noindent\textbf{3：不適切} MT式のµを真空の誘電率と定義する。E/HからΩを持つインピーダンスを作りωµで割る。µはH/mの透磁率でありF/mの誘電率ではない。\par

\noindent\textbf{4：適切} TEMはループ又は接地両端電流源で深部探査。磁気ループ方式だけでなく接地電流源の時間領域応答を測る構成もある。\par

\noindent\textbf{5：条件付きで適切} VLFは3〜30kHzの既存局で送信条件を制御できない。既存通信送信局を利用する受信探査では局の位置・周波数・運用を調査者が任意には設定できない。専用人工源の別方式を含めない。\par

\paragraph{再利用する考え方} 式の記号の名前・単位と励起方式を照合し、形だけで正しいと判定しない。

\paragraph{共通理論} \hyperref[theory:exploration:dcarrays]{第\ref*{theory:exploration:dcarrays}節}, \hyperref[theory:exploration:archie]{第\ref*{theory:exploration:archie}節}, \hyperref[theory:exploration:nearfield]{第\ref*{theory:exploration:nearfield}節}

\par 参考：\cite{epa_resistivity,usgs_csamt2012,epa_borehole_resistivity,epa_tdem_archive}。

\par 難易度 3（校準初稿）、優先度未評価。

\section{反射走時の双曲線・局所近似・配置}

\subsection{平成28年度 問21}
\label{q:h28-reg-applied-21}

\noindent\textbf{原典：} sources/H28/17\_applied\_science.pdf、PDF物理8頁。

\textbf{公式正答：} 1（single）。

\par\textbf{検討未確定：} 放物線は小オフセットの正当な局所二次近似となる。原文の近似に範囲がなく、公式1を通常の双曲線近似の知識として保持し条件不足を査読対象にする。20mも装置条件付きの目安。

\paragraph{問題の要旨} 反射走時の双曲線と放物線近似、ミラージ、微動、表面波を比較する。

\Needspace{4\baselineskip}\paragraph{確認した条件・記述（独自要約）}\mbox{}\par

\noindent\textbf{原典照合：} 五候補の順序、最適切と不適切の方向、波種、深度目安、曲線名、最長と最短を原典画像で照合。\par

\noindent\textbf{前提：} 単一水平反射境界の走時t\textsuperscript{2}=t0\textsuperscript{2}+x\textsuperscript{2}/V\textsuperscript{2}を基本モデルとする。\par

\paragraph{問われる核心} 反射走時の双曲線と放物線近似、ミラージ、微動、表面波を比較する。

\paragraph{必要知識} 単一水平反射境界の走時t\textsuperscript{2}=t0\textsuperscript{2}+x\textsuperscript{2}/V\textsuperscript{2}を基本モデルとする。

\paragraph{考え方・正答を選ぶ理由（査読待ち初稿）} 公式1は通常の反射走時を双曲線とすることを問う。しかしt=\allowbreak{}t0 sqrt(1+x\textsuperscript{2}/(V\textsuperscript{2}t0\textsuperscript{2}))を小xで展開すればt≈t0+x\textsuperscript{2}/(2V\textsuperscript{2}t0)という放物線近似が成立する。近似されるという原文を絶対的に数学の誤りと断定しない。

\Needspace{4\baselineskip}\paragraph{選択肢の検討（査読待ち初稿）}\mbox{}\par

\noindent\textbf{1：条件付きで適切} 水平二層の反射走時は放物線で近似される。正確な単一反射の走時は双曲線である。一方、小オフセットの二次近似は放物線であり、適用範囲を指定すれば近似として成立する。\par

\noindent\textbf{2：適切} 連続増速のミラージ構造では初動走時が曲線になりトモグラフィが適する。離散境界の折れ線モデルでなく深度ごとに屈折しながら曲がる経路となる。格子速度を更新する走時トモグラフィは連続的な速度変化も扱える。解像度は波線分布に依存する。\par

\noindent\textbf{3：適切} 屈折法は通常初動を読み解析する。直接波・臨界屈折波のうち最も早く到達する時刻を距離ごとに読む。反射波を同じ初動としない。\par

\noindent\textbf{4：適切} 微動アレイで自然・人為的な微振動からS波構造を推定する。複数点の位相・空間相関から表面波の分散を求め、主にS波速度構造を推定する。\par

\noindent\textbf{5：条件付きで適切} 表面波から20m程度までのS波分布を求める。通常の浅部能動測定の典型的対象である。最大波長と配列・震源が変われば探査深度も変わり、20mを絶対上限とはしない。\par

\paragraph{再利用する考え方} 曲線名を暗記するだけでなく走時式を展開し、近似範囲を確認する。

\paragraph{共通理論} \hyperref[theory:exploration:reflectiongeometry]{第\ref*{theory:exploration:reflectiongeometry}節}, \hyperref[theory:exploration:surfacewaves]{第\ref*{theory:exploration:surfacewaves}節}

\par 参考：\cite{epa_seismic_reflection,jica_refraction2015,moe_surface2004,epa_surfacewaves}。

\paragraph{関連問題} \hyperref[q:r05-reg-applied-32]{R05-REG-APPLIED-32}

\par 難易度 3（校準初稿）、優先度未評価。

\subsection{令和5年度 問32}
\label{q:r05-reg-applied-32}

\noindent\textbf{原典：} sources/R05/17\_applied\_science.pdf、PDF物理11頁。

\textbf{公式正答：} 4（single）。

\par\textbf{検討未確定：} 放物線は小オフセットの局所近似として成立する。公式4を一般走時式の比較として保持し範囲不足を残す。CMPによる相関雑音抑制もmoveoutの差に依存する。

\paragraph{問題の要旨} ミラージ、測線方向、S波屈折、反射曲線、相関雑音を比較する。

\Needspace{4\baselineskip}\paragraph{確認した条件・記述（独自要約）}\mbox{}\par

\noindent\textbf{原典照合：} 五候補の順序、最適切と不適切の方向、波種、深度目安、曲線名、最長と最短を原典画像で照合。\par

\noindent\textbf{前提：} CMPの通常の双曲線近似と小オフセットの局所二次近似を分ける。\par

\paragraph{問われる核心} ミラージ、測線方向、S波屈折、反射曲線、相関雑音を比較する。

\paragraph{必要知識} CMPの通常の双曲線近似と小オフセットの局所二次近似を分ける。

\paragraph{考え方・正答を選ぶ理由（査読待ち初稿）} 公式4である。水平成層の通常近似は双曲線であり、放物線を一般的走時式とする説明は不十分である。一方、小オフセットのTaylor二次近似は放物線となるため、原文の近似的という語に対する範囲不足を疑義として保持する。

\Needspace{4\baselineskip}\paragraph{選択肢の検討（査読待ち初稿）}\mbox{}\par

\noindent\textbf{1：適切} 連続的な深度増速構造をミラージュ構造と呼ぶ。速度層境界を持つ折れ線モデルでなく、曲がる波線の走時曲線になる連続勾配モデルである。\par

\noindent\textbf{2：条件付きで適切} 屈折測線を等高線とほぼ垂直で一定勾配に配置するのが望ましい。地質構造と地形を評価しやすい代表的な測線設計である。全地形・全目的でその配置が唯一又は可能とはいえない。\par

\noindent\textbf{3：適切} 土質のS波屈折を比較的平坦な地形で多く行う。起振と受振及び地形補正の制約を考えた通常の調査条件である。\par

\noindent\textbf{4：条件付きで適切} 多層CMPの反射走時は近似的に放物線となる。通常は有効速度を用いる双曲線で近似する。小オフセットなら双曲線の二次展開は放物線となるため、その条件を除外せず公式判断と分ける。\par

\noindent\textbf{5：条件付きで適切} CMP重合がランダムとコヒーレント雑音の抑制に貢献する。適切なNMOで反射波をそろえ、異なるmoveoutの多重反射等は重合で弱められる。反射と同じmoveoutの相関雑音まで必ず消えるわけではない。\par

\paragraph{再利用する考え方} 近似の次数と範囲、雑音のmoveoutを示して判断する。

\paragraph{共通理論} \hyperref[theory:exploration:reflectiongeometry]{第\ref*{theory:exploration:reflectiongeometry}節}, \hyperref[theory:exploration:refraction]{第\ref*{theory:exploration:refraction}節}, \hyperref[theory:exploration:stack]{第\ref*{theory:exploration:stack}節}

\par 参考：\cite{epa_seismic_reflection,jica_refraction2015,usgs_processing2000}。

\paragraph{関連問題} \hyperref[q:h28-reg-applied-21]{H28-REG-APPLIED-21}

\par 難易度 3（校準初稿）、優先度未評価。

\section{反射分解能・薄層分離と存在検出}

\subsection{平成26年度 問20}
\label{q:h26-reg-applied-20}

\noindent\textbf{原典：} sources/H26/17\_applied\_science.pdf、PDF物理8頁。

\textbf{公式正答：} 1（single）。

\par\textbf{検討未確定：} 原典1の反射法は浅層の目安を想定した公式判断と読めるが、調査条件を一律固定できない。2は検出という語で分解能を説明する省略があり、薄層の振幅検出を否定しない。

\paragraph{問題の要旨} 反射法の配置・薄層分解能・NMOと屈折法のブラインドを比較する。

\Needspace{4\baselineskip}\paragraph{確認した条件・記述（独自要約）}\mbox{}\par

\noindent\textbf{原典照合：} 五候補の順序、最適切と不適切の方向、波種、深度目安、曲線名、最長と最短を原典画像で照合。\par

\noindent\textbf{前提：} 最大オフセットの目安と必要条件、薄層の検出と二境界の分離を区別する。\par

\paragraph{問われる核心} 反射法の配置・薄層分解能・NMOと屈折法のブラインドを比較する。

\paragraph{必要知識} 最大オフセットの目安と必要条件、薄層の検出と二境界の分離を区別する。

\paragraph{考え方・正答を選ぶ理由（査読待ち初稿）} 1が不適切である。浅層反射の展開長の通常目安を探査深度程度とする資料に対し、二倍を必須としている。必要距離は目的と速度構造次第であり、二倍が絶対禁止という意味でもない。2の四分の一波長は二境界の分離の目安であり、薄層の存在検出はそれ以下でも可能となる。

\Needspace{4\baselineskip}\paragraph{選択肢の検討（査読待ち初稿）}\mbox{}\par

\noindent\textbf{1：不適切} 反射法の最大オフセットを探査深度の二倍程度にする必要がある。浅層反射の通常目安の倍数を取り違える。一律に二倍を必要とする原理はなく、波群分離・速度解析・調査目的から設計する。\par

\noindent\textbf{2：条件付きで適切} 反射法で波長の四分の一未満の地層の検出が難しい。通常の分解能目安はλ/4程度で上面と下面を別々の反射として分離する話である。振幅変化で存在を検出できることと区別し、薄層は全て不可視とはしない。\par

\noindent\textbf{3：適切} CMP記録からNMO用の重合速度を速度解析する。異なるオフセットでの走時差をそろえる速度を推定する。重合速度を無条件に真の区間速度と同一視しない。\par

\noindent\textbf{4：適切} 深部へ増速する層構造でも薄い中間層が初動に現れない。臨界屈折波が生じても、直接波や深い層を通る波より早くなる距離範囲がなければ初動に現れない。低速度逆転だけが見えない原因ではない。\par

\noindent\textbf{5：条件付きで適切} 屈折法で全走時が一致することを確認する。相反的な同じ測線で起振点を往復した端点間走時の整合などを確認する。異なる受振点の全時刻が同じになるという意味ではなく、走時読み取り・経路の条件を要する。\par

\paragraph{再利用する考え方} 配置の数値、分離能力、検出の有無、速度の種類を別々に判断する。

\paragraph{共通理論} \hyperref[theory:exploration:reflectiongeometry]{第\ref*{theory:exploration:reflectiongeometry}節}, \hyperref[theory:exploration:resolution]{第\ref*{theory:exploration:resolution}節}, \hyperref[theory:exploration:refraction]{第\ref*{theory:exploration:refraction}節}

\par 参考：\cite{jica_refraction2015,usgs_refraction2013,epa_surfacewaves,xia_rayleigh1999,coastal_reflection2016,seg_vertical_resolution,usgs_processing2000}。

\paragraph{関連問題} \hyperref[q:r03-reg-applied-33]{R03-REG-APPLIED-33}

\par 難易度 3（校準初稿）、優先度未評価。

\subsection{令和3年度 問33}
\label{q:r03-reg-applied-33}

\noindent\textbf{原典：} sources/R03/17\_applied\_science.pdf、PDF物理12頁。

\textbf{公式正答：} 1（single）。

\paragraph{問題の要旨} 反射の分解能を走時切片と区別する。

\Needspace{4\baselineskip}\paragraph{確認した条件・記述（独自要約）}\mbox{}\par

\noindent\textbf{原典照合：} 五候補の順序、最適切と不適切の方向、波種、深度目安、曲線名、最長と最短を原典画像で照合。\par

\noindent\textbf{前提：} 薄層の上下面の分離、横方向の反射寄与範囲、入力waveletを区別する。\par

\paragraph{問われる核心} 反射の分解能を走時切片と区別する。

\paragraph{必要知識} 薄層の上下面の分離、横方向の反射寄与範囲、入力waveletを区別する。

\paragraph{考え方・正答を選ぶ理由（査読待ち初稿）} 1が不適切である。intercept timeは屈折走時の直線の切片であり、反射の二イベントの分離や横方向の寄与範囲を表す基準ではない。

\Needspace{4\baselineskip}\paragraph{選択肢の検討（査読待ち初稿）}\mbox{}\par

\noindent\textbf{1：不適切} インターセプトタイム。屈折波の走時直線を距離ゼロへ外挿した切片である。反射分解能の直接的な基準と異なる。\par

\noindent\textbf{2：適切} レイリー基準。近接する二つの応答を分離できるかという分解能の基準に関係する。表面波としてのRayleighと用語の役割を分ける。\par

\noindent\textbf{3：適切} ウェーブレット。反射の応答形状と有効帯域が近いイベントを分離できる程度に影響する。\par

\noindent\textbf{4：適切} リッカー基準。反射waveletの重なりを用いる分解能の基準に関係する。具体的閾値を基準名から一律数値へ補わない。\par

\noindent\textbf{5：適切} フレネルゾーン。同程度の位相で反射へ寄与する面上の範囲であり、横方向の分解能に関係する。migration前後と開口の条件を区別する。\par

\paragraph{再利用する考え方} 同じ人名の用語でも波の種類、wavelet、分離基準を混同しない。

\paragraph{共通理論} \hyperref[theory:exploration:resolution]{第\ref*{theory:exploration:resolution}節}

\par 参考：\cite{seg_vertical_resolution,seg_lateral_resolution,usgs_refraction2013}。

\paragraph{関連問題} \hyperref[q:h26-reg-applied-20]{H26-REG-APPLIED-20}

\par 難易度 2（校準初稿）、優先度未評価。

\section{屈折法の地山評価と主用途}

\subsection{令和2年度 問32}
\label{q:r02-reg-applied-32}

\noindent\textbf{原典：} sources/R02/17\_applied\_science.pdf、PDF物理11頁。

\textbf{公式正答：} 2（single）。

\par\textbf{検討未確定：} 屈折法にはS波もあり、強震動に寄与する地下情報を得る可能性はある。通常のP波地山調査の主用途という比較を明示し、用途の絶対排除をしない。

\paragraph{問題の要旨} 屈折法の主用途と強震動のS波モデルを区別する。

\Needspace{4\baselineskip}\paragraph{確認した条件・記述（独自要約）}\mbox{}\par

\noindent\textbf{原典照合：} 五候補の順序、最適切と不適切の方向、波種、深度目安、曲線名、最長と最短を原典画像で照合。\par

\noindent\textbf{前提：} 通常のP波屈折による地山の硬軟・割れ目・風化等の評価を主用途として比較する。\par

\paragraph{問われる核心} 屈折法の主用途と強震動のS波モデルを区別する。

\paragraph{必要知識} 通常のP波屈折による地山の硬軟・割れ目・風化等の評価を主用途として比較する。

\paragraph{考え方・正答を選ぶ理由（査読待ち初稿）} 公式2である。強震動・耐震評価では主にS波構造や減衰・震源条件が必要であり、通常の地山P波屈折をその主たる方法としない。他の四項目は岩盤・土構造物の地質条件の評価に結び付く。ただしS波屈折や複数手法の統合が耐震評価に役立つ可能性まで否定しない。

\Needspace{4\baselineskip}\paragraph{選択肢の検討（査読待ち初稿）}\mbox{}\par

\noindent\textbf{1：適切} ダムの基礎岩盤の評価。風化・破砕・硬軟の構造を把握する通常の地山評価に関係する。\par

\noindent\textbf{2：条件付きで適切} 強震動予測や耐震設計。P波屈折の代表的な主用途の比較では公式の選択対象である。S波屈折等の応用と地下モデルの補助情報としての役割はあり、全く利用できないという意味ではない。\par

\noindent\textbf{3：適切} のり面勾配等の土構造物の基礎調査。風化帯や基盤等の地質評価を補助する。速度だけで設計勾配を機械的に決定しない。\par

\noindent\textbf{4：適切} 不安定土塊の工学的判断。土塊・基盤の速度構造が地質状態の判断を助ける。力学強度や地下水等の他情報も要する。\par

\noindent\textbf{5：適切} トンネル沿いの地山評価。岩盤の風化・破砕帯等の分布を速度情報から把握する代表的用途である。\par

\paragraph{再利用する考え方} 主としてという比較を、利用可能性が一切ないという全称へ変えない。

\paragraph{共通理論} \hyperref[theory:exploration:waveproperties]{第\ref*{theory:exploration:waveproperties}節}, \hyperref[theory:exploration:surfacewaves]{第\ref*{theory:exploration:surfacewaves}節}

\par 参考：\cite{segj_integrated2017,jica_refraction2015,japansea_spac2017}。

\par 難易度 3（校準初稿）、優先度未評価。

\section{屈折初動のはぎとり・相反走時・連続勾配}

\subsection{平成23年度 問32}
\label{q:h23-reg-applied-32}

\noindent\textbf{原典：} sources/H23/17\_applied\_science.pdf、PDF物理11頁。

\textbf{公式正答：} 3（single）。

\par\textbf{検討未確定：} オフセットと深度の対応は目安であり一律の必要条件ではない。全走時の一致は相反的往復観測の整合であり、全測点の同時到達を意味しない。

\paragraph{問題の要旨} 表面波・浅層反射・はぎとり・全走時・ミラージ構造を比較する。

\Needspace{4\baselineskip}\paragraph{確認した条件・記述（独自要約）}\mbox{}\par

\noindent\textbf{原典照合：} 五候補の順序、最適切と不適切の方向、波種、深度目安、曲線名、最長と最短を原典画像で照合。\par

\noindent\textbf{前提：} 通常の屈折法の初動走時、浅層反射の配置目安、波長依存の表面波分散を考える。\par

\paragraph{問われる核心} 表面波・浅層反射・はぎとり・全走時・ミラージ構造を比較する。

\paragraph{必要知識} 通常の屈折法の初動走時、浅層反射の配置目安、波長依存の表面波分散を考える。

\paragraph{考え方・正答を選ぶ理由（査読待ち初稿）} 3が不適切である。はぎとり法は屈折波と直接波の初動走時の解析であり、反射波から同じ手法で境界を求める説明は対応が誤っている。全走時一致は往復の整合を意味し、各受振点の時刻が同じという命題ではない。

\Needspace{4\baselineskip}\paragraph{選択肢の検討（査読待ち初稿）}\mbox{}\par

\noindent\textbf{1：適切} レイリー波の波長別位相速度からS波速度分布を推定する。地盤の深度方向の不均質が分散曲線へ反映され、主にS波速度モデルを推定できる。P波速度・密度・層厚にも依存し、S波だけの関数ではない。\par

\noindent\textbf{2：条件付きで適切} 浅層反射で最大オフセットを探査深度程度に設定する。浅層反射の一般的な展開目安として読む。速度解析・近接雑音・反射角・地形・重合数で調整し、全ての計測に必要な厳密等式ではない。\par

\noindent\textbf{3：不適切} 反射波の走時をはぎとり法で解析して境界深度を求める。萩原のはぎとり法は直接波と屈折波の初動走時から層構造を求める手法である。反射波の走時を同じ手法の入力とする説明は誤りである。\par

\noindent\textbf{4：条件付きで適切} 屈折法で全走時が一致することを確認する。相反的な同じ測線で起振点を往復した端点間走時の整合などを確認する。異なる受振点の全時刻が同じになるという意味ではなく、走時読み取り・経路の条件を要する。\par

\noindent\textbf{5：適切} 連続増速のミラージ構造では初動走時が曲線になりトモグラフィが適する。離散境界の折れ線モデルでなく深度ごとに屈折しながら曲がる経路となる。格子速度を更新する走時トモグラフィは連続的な速度変化も扱える。解像度は波線分布に依存する。\par

\paragraph{再利用する考え方} 反射と屈折を、入力する波と走時解析の前提から分ける。

\paragraph{共通理論} \hyperref[theory:exploration:refraction]{第\ref*{theory:exploration:refraction}節}, \hyperref[theory:exploration:reflectiongeometry]{第\ref*{theory:exploration:reflectiongeometry}節}, \hyperref[theory:exploration:surfacewaves]{第\ref*{theory:exploration:surfacewaves}節}

\par 参考：\cite{jica_refraction2015,usgs_refraction2013,epa_surfacewaves,xia_rayleigh1999,coastal_reflection2016}。

\paragraph{関連問題} \hyperref[q:h29-reg-applied-34]{H29-REG-APPLIED-34}

\par 難易度 3（校準初稿）、優先度未評価。

\subsection{平成29年度 問34}
\label{q:h29-reg-applied-34}

\noindent\textbf{原典：} sources/H29/17\_applied\_science.pdf、PDF物理12頁。

\textbf{公式正答：} 3（single）。

\par\textbf{検討未確定：} オフセットと深度の対応は目安であり一律の必要条件ではない。全走時の一致は相反的往復観測の整合であり、全測点の同時到達を意味しない。

\paragraph{問題の要旨} 表面波・浅層反射・はぎとり・全走時・ミラージ構造を比較する。

\Needspace{4\baselineskip}\paragraph{確認した条件・記述（独自要約）}\mbox{}\par

\noindent\textbf{原典照合：} 五候補の順序、最適切と不適切の方向、波種、深度目安、曲線名、最長と最短を原典画像で照合。\par

\noindent\textbf{前提：} 通常の屈折法の初動走時、浅層反射の配置目安、波長依存の表面波分散を考える。\par

\paragraph{問われる核心} 表面波・浅層反射・はぎとり・全走時・ミラージ構造を比較する。

\paragraph{必要知識} 通常の屈折法の初動走時、浅層反射の配置目安、波長依存の表面波分散を考える。

\paragraph{考え方・正答を選ぶ理由（査読待ち初稿）} 3が不適切である。はぎとり法は屈折波と直接波の初動走時の解析であり、反射波から同じ手法で境界を求める説明は対応が誤っている。全走時一致は往復の整合を意味し、各受振点の時刻が同じという命題ではない。

\Needspace{4\baselineskip}\paragraph{選択肢の検討（査読待ち初稿）}\mbox{}\par

\noindent\textbf{1：適切} レイリー波の波長別位相速度からS波速度分布を推定する。地盤の深度方向の不均質が分散曲線へ反映され、主にS波速度モデルを推定できる。P波速度・密度・層厚にも依存し、S波だけの関数ではない。\par

\noindent\textbf{2：条件付きで適切} 浅層反射で最大オフセットを探査深度程度に設定する。浅層反射の一般的な展開目安として読む。速度解析・近接雑音・反射角・地形・重合数で調整し、全ての計測に必要な厳密等式ではない。\par

\noindent\textbf{3：不適切} 反射波の走時をはぎとり法で解析して境界深度を求める。萩原のはぎとり法は直接波と屈折波の初動走時から層構造を求める手法である。反射波の走時を同じ手法の入力とする説明は誤りである。\par

\noindent\textbf{4：条件付きで適切} 屈折法で全走時が一致することを確認する。相反的な同じ測線で起振点を往復した端点間走時の整合などを確認する。異なる受振点の全時刻が同じになるという意味ではなく、走時読み取り・経路の条件を要する。\par

\noindent\textbf{5：適切} 連続増速のミラージ構造では初動走時が曲線になりトモグラフィが適する。離散境界の折れ線モデルでなく深度ごとに屈折しながら曲がる経路となる。格子速度を更新する走時トモグラフィは連続的な速度変化も扱える。解像度は波線分布に依存する。\par

\paragraph{再利用する考え方} 反射と屈折を、入力する波と走時解析の前提から分ける。

\paragraph{共通理論} \hyperref[theory:exploration:refraction]{第\ref*{theory:exploration:refraction}節}, \hyperref[theory:exploration:reflectiongeometry]{第\ref*{theory:exploration:reflectiongeometry}節}, \hyperref[theory:exploration:surfacewaves]{第\ref*{theory:exploration:surfacewaves}節}

\par 参考：\cite{jica_refraction2015,usgs_refraction2013,epa_surfacewaves,xia_rayleigh1999,coastal_reflection2016}。

\paragraph{関連問題} \hyperref[q:h23-reg-applied-32]{H23-REG-APPLIED-32}

\par 難易度 3（校準初稿）、優先度未評価。

\section{地震速度境界・地質と地盤状態}

\subsection{平成25年度 問32}
\label{q:h25-reg-applied-32}

\noindent\textbf{原典：} sources/H25/17\_applied\_science.pdf、PDF物理12頁。

\textbf{公式正答：} 3（single）。

\par\textbf{検討未確定：} P>Sは安定な等方弾性体の条件、増速と初動のみは通常の屈折法の枠組みである。異方性や他の解析を無条件に排除しない。

\paragraph{問題の要旨} 屈折法の増速前提と地下水・風化・波種・初動を区別する。

\Needspace{4\baselineskip}\paragraph{確認した条件・記述（独自要約）}\mbox{}\par

\noindent\textbf{原典照合：} 五候補の順序、最適切と不適切の方向、波種、深度目安、曲線名、最長と最短を原典画像で照合。\par

\noindent\textbf{前提：} 通常の地盤の等方弾性を想定し、臨界屈折の層状解釈とP波の飽和依存を区別する。\par

\paragraph{問われる核心} 屈折法の増速前提と地下水・風化・波種・初動を区別する。

\paragraph{必要知識} 通常の地盤の等方弾性を想定し、臨界屈折の層状解釈とP波の飽和依存を区別する。

\paragraph{考え方・正答を選ぶ理由（査読待ち初稿）} 3が不適切である。空隙を空気又は水が満たす状態は圧縮性等へ影響し、P波速度は地下水の有無・飽和状態に左右され得る。速度だけで岩種を決めず、同じ岩種でも割れ目と風化の程度を考える。

\Needspace{4\baselineskip}\paragraph{選択肢の検討（査読待ち初稿）}\mbox{}\par

\noindent\textbf{1：条件付きで適切} 通常の屈折法で深部ほど速度が高いと仮定する。臨界屈折の層状モデルの前提である。あらゆる地下構造が必ず増速するという自然法則ではない。\par

\noindent\textbf{2：適切} 同じ岩種でも浅部速度の違いが風化・変質・割れ目の違いを反映する。弾性係数、接触状態、空隙等の変化が速度へ影響する。同じ岩種なら一定速度とは限らない。\par

\noindent\textbf{3：不適切} 土質のP波速度は地下水の有無に左右されない。空隙流体と飽和状態が圧縮性に影響するため無関係という断定は誤りである。\par

\noindent\textbf{4：条件付きで適切} P波速度は常にS波速度より大きい。安定な等方弾性体ではVP\textsuperscript{2}=(K+4G/3)/ρ、VS\textsuperscript{2}=G/ρよりVP>VSである。異方性のモードや流体のS波不存在まで常に二つの速度を比較できるわけではない。\par

\noindent\textbf{5：条件付きで適切} 解析に初動走時のみを用いる。通常の屈折初動解析の説明である。反射・全波形等を含める別の解析まで禁止する意味ではない。\par

\paragraph{再利用する考え方} 波種、弾性係数、空隙流体を結び、速度を岩種の識別子にしない。

\paragraph{共通理論} \hyperref[theory:exploration:refraction]{第\ref*{theory:exploration:refraction}節}, \hyperref[theory:exploration:waveproperties]{第\ref*{theory:exploration:waveproperties}節}

\par 参考：\cite{usgs_refraction2013,segj_integrated2017}。

\paragraph{関連問題} \hyperref[q:r03-reg-applied-34]{R03-REG-APPLIED-34}

\par 難易度 3（校準初稿）、優先度未評価。

\subsection{令和3年度 問34}
\label{q:r03-reg-applied-34}

\noindent\textbf{原典：} sources/R03/17\_applied\_science.pdf、PDF物理13頁。

\textbf{公式正答：} 2（single）。

\par\textbf{検討未確定：} S波屈折は平坦な土質調査という代表的運用の比較であり、急峻地形で絶対実施できないと断定しない。速度からの地質推定は複数解を持つ。

\paragraph{問題の要旨} 速度境界と地質境界、S波屈折の地形条件と用途を比較する。

\Needspace{4\baselineskip}\paragraph{確認した条件・記述（独自要約）}\mbox{}\par

\noindent\textbf{原典照合：} 五候補の順序、最適切と不適切の方向、波種、深度目安、曲線名、最長と最短を原典画像で照合。\par

\noindent\textbf{前提：} 通常の土質地盤のS波屈折は比較的平坦な配置で測ることが多い。地形により走時解釈や起振・受振方向が難しくなる。\par

\paragraph{問われる核心} 速度境界と地質境界、S波屈折の地形条件と用途を比較する。

\paragraph{必要知識} 通常の土質地盤のS波屈折は比較的平坦な配置で測ることが多い。地形により走時解釈や起振・受振方向が難しくなる。

\paragraph{考え方・正答を選ぶ理由（査読待ち初稿）} 2が不適切である。土質のS波屈折を急峻地形で多く行うという説明を、典型的な比較的平坦な条件と取り違えている。地形を克服する特別な測定が不可能という意味ではない。

\Needspace{4\baselineskip}\paragraph{選択肢の検討（査読待ち初稿）}\mbox{}\par

\noindent\textbf{1：適切} 屈折法の速度境界は地質境界と必ずしも一致しない。速度は岩種だけでなく風化・割れ目・飽和等に依存する。同じ岩種内でも速度が変わり、異なる地質が近い速度を持つ場合もある。\par

\noindent\textbf{2：不適切} 土質のS波屈折を比較的急峻な地形で多く行う。典型的なS波の土質調査の地形条件は比較的平坦である。急峻が通常条件という説明は逆であり、起振・受振・経路の制約が増す。\par

\noindent\textbf{3：適切} トンネル・ダム・のり面等の基礎調査へ屈折法を用いる。通常のP波初動を使って地山の速度構造を求め、工学的地質評価の一情報とする。\par

\noindent\textbf{4：条件付きで適切} 速度構造から硬軟・割れ目・風化・破砕帯等を推定する。物性と地質状態の相関を利用する推定である。速度だけから一意に種類や強度を決めず、露頭・ボーリング等で拘束する。\par

\noindent\textbf{5：適切} 速度層は岩種・膠結・風化・変質・亀裂の複合で決まる。速度への寄与が複数あるため、速度層を岩種ごとのラベルと同一視しない。\par

\paragraph{再利用する考え方} 地形の典型条件と技術的な実施可能性を分け、速度の多要因性を保つ。

\paragraph{共通理論} \hyperref[theory:exploration:waveproperties]{第\ref*{theory:exploration:waveproperties}節}, \hyperref[theory:exploration:refraction]{第\ref*{theory:exploration:refraction}節}

\par 参考：\cite{jica_refraction2015,segj_integrated2017,usgs_refraction2013}。

\paragraph{関連問題} \hyperref[q:h25-reg-applied-32]{H25-REG-APPLIED-32}, \hyperref[q:r06-reg-applied-32]{R06-REG-APPLIED-32}

\par 難易度 3（校準初稿）、優先度未評価。

\subsection{令和6年度 問32}
\label{q:r06-reg-applied-32}

\noindent\textbf{原典：} sources/R06/17\_applied\_science.pdf、PDF物理11頁。

\textbf{公式正答：} 1（single）。

\par\textbf{検討未確定：} オフセット深度の対応と双曲線は通常条件の目安・近似である。受振器群の効果は方向・間隔・波長に依存する。

\paragraph{問題の要旨} 速度境界、連続勾配、浅層展開、反射双曲線、受振器群を比較する。

\Needspace{4\baselineskip}\paragraph{確認した条件・記述（独自要約）}\mbox{}\par

\noindent\textbf{原典照合：} 五候補の順序、最適切と不適切の方向、波種、深度目安、曲線名、最長と最短を原典画像で照合。\par

\noindent\textbf{前提：} 地震波速度は岩種だけでなく、割れ目・風化・飽和などの状態にも依存する。\par

\paragraph{問われる核心} 速度境界、連続勾配、浅層展開、反射双曲線、受振器群を比較する。

\paragraph{必要知識} 地震波速度は岩種だけでなく、割れ目・風化・飽和などの状態にも依存する。

\paragraph{考え方・正答を選ぶ理由（査読待ち初稿）} 1が不適切である。屈折法の速度境界と地質学の境界がよく一致するという一般化はできない。同じ岩種内の風化境界が速度を変え、異なる岩種間で速度差が小さい場合もある。

\Needspace{4\baselineskip}\paragraph{選択肢の検討（査読待ち初稿）}\mbox{}\par

\noindent\textbf{1：不適切} 屈折法の層境界は地質境界とよく一致する。速度差で定まる境界を地質分類の境界と同一視している。岩種以外の要因が速度へ影響するため一般的な一致を保証しない。\par

\noindent\textbf{2：適切} 連続増速のミラージ構造では初動走時が曲線になりトモグラフィが適する。離散境界の折れ線モデルでなく深度ごとに屈折しながら曲がる経路となる。格子速度を更新する走時トモグラフィは連続的な速度変化も扱える。解像度は波線分布に依存する。\par

\noindent\textbf{3：条件付きで適切} 浅層反射で最大オフセットを探査深度程度に設定する。浅層反射の一般的な展開目安として読む。速度解析・近接雑音・反射角・地形・重合数で調整し、全ての計測に必要な厳密等式ではない。\par

\noindent\textbf{4：条件付きで適切} 水平多層CMPの走時は近似的に双曲線となる。通常のNMO近似で有効速度とゼロオフセット走時を用いる。大オフセットや強い不均質・異方性では非双曲線成分が重要になる。\par

\noindent\textbf{5：条件付きで適切} 受振器群を結合し表面波等を抑え反射を強調する。配列の空間応答により波群を選別し独立雑音も平均化する。方向・間隔が対象に適していなければ必要信号を弱める場合がある。\par

\paragraph{再利用する考え方} 速度構造と地質構造、近似モデルと実地の制約を別々に評価する。

\paragraph{共通理論} \hyperref[theory:exploration:waveproperties]{第\ref*{theory:exploration:waveproperties}節}, \hyperref[theory:exploration:reflectiongeometry]{第\ref*{theory:exploration:reflectiongeometry}節}, \hyperref[theory:exploration:stack]{第\ref*{theory:exploration:stack}節}

\par 参考：\cite{segj_integrated2017,jica_refraction2015,usgs_processing2000,epa_seismic_reflection}。

\paragraph{関連問題} \hyperref[q:r03-reg-applied-34]{R03-REG-APPLIED-34}

\par 難易度 3（校準初稿）、優先度未評価。

\section{地震探査の走時・画像化用語}

\subsection{平成24年度 問23}
\label{q:h24-reg-applied-23}

\noindent\textbf{原典：} sources/H24/17\_applied\_science.pdf、PDF物理7頁。

\textbf{公式正答：} 4（single）。

\paragraph{問題の要旨} 地震探査の走時・反射用語と重力補正を区別する。

\Needspace{4\baselineskip}\paragraph{確認した条件・記述（独自要約）}\mbox{}\par

\noindent\textbf{原典照合：} 否定の対象、指数、周波数と周期、五つの候補の列順を原典画像で照合。\par

\noindent\textbf{前提：} 地震探査は波の走時・波形、重力探査は加速度の異常を測る。\par

\paragraph{問われる核心} 地震探査の走時・反射用語と重力補正を区別する。

\paragraph{必要知識} 地震探査は波の走時・波形、重力探査は加速度の異常を測る。

\paragraph{考え方・正答を選ぶ理由（査読待ち初稿）} 4のブーゲー補正が不適切である。地震探査に関係する四つは、屈折波の走時切片、初動に現れにくい層、反射を支配するインピーダンス、地形や表層に由来する静的走時ずれに対応する。

\Needspace{4\baselineskip}\paragraph{選択肢の検討（査読待ち初稿）}\mbox{}\par

\noindent\textbf{1：適切} 原点走時。屈折波の走時直線を距離ゼロへ外挿した切片である。直線の傾きと合わせて層厚等を推定する。実測した距離ゼロの初動時刻とは同一視しない。\par

\noindent\textbf{2：適切} ブラインド層。屈折法の初動走時に現れにくい層である。薄い中間層が初動を与えない場合や低速度層を区別し、用語の範囲は資料により異なる。\par

\noindent\textbf{3：適切} 音響インピーダンス。垂直入射の簡単なモデルではZ=ρVであり、境界両側の差が反射振幅を支配する。重力で用いる密度そのものだけとは異なる。\par

\noindent\textbf{4：不適切} ブーゲー補正。測点と基準面の間の質量が重力へ及ぼす効果を補正する重力探査の処理である。地震波走時に対する静補正ではない。\par

\noindent\textbf{5：適切} 静補正。地形や表層低速度層などによる走時ずれを一定の時間シフトとしてそろえる処理である。重力の地形・Bouguer補正と観測量が異なる。\par

\paragraph{再利用する考え方} 処理名だけでなく何を補正し、単位は何かを問う。

\paragraph{共通理論} \hyperref[theory:exploration:seismicterms]{第\ref*{theory:exploration:seismicterms}節}, \hyperref[theory:gravity:corrections]{第\ref*{theory:gravity:corrections}節}

\par 参考：\cite{usgs_refraction2013,epa_seismic_reflection,usgs_static2004,usgs_processing2000}。

\paragraph{関連問題} \hyperref[q:h30-reg-applied-33]{H30-REG-APPLIED-33}

\par 難易度 2（校準初稿）、優先度未評価。

\subsection{平成30年度 問33}
\label{q:h30-reg-applied-33}

\noindent\textbf{原典：} sources/H30/17\_applied\_science.pdf、PDF物理11頁。

\textbf{公式正答：} 5（single）。

\paragraph{問題の要旨} 地震探査の走時・反射用語と重力補正を区別する。

\Needspace{4\baselineskip}\paragraph{確認した条件・記述（独自要約）}\mbox{}\par

\noindent\textbf{原典照合：} 否定の対象、指数、周波数と周期、五つの候補の列順を原典画像で照合。\par

\noindent\textbf{前提：} 地震探査は波の走時・波形、重力探査は加速度の異常を測る。\par

\paragraph{問われる核心} 地震探査の走時・反射用語と重力補正を区別する。

\paragraph{必要知識} 地震探査は波の走時・波形、重力探査は加速度の異常を測る。

\paragraph{考え方・正答を選ぶ理由（査読待ち初稿）} 5のブーゲー補正が不適切である。地震探査に関係する四つは、屈折波の走時切片、初動に現れにくい層、反射を支配するインピーダンス、地形や表層に由来する静的走時ずれに対応する。

\Needspace{4\baselineskip}\paragraph{選択肢の検討（査読待ち初稿）}\mbox{}\par

\noindent\textbf{1：適切} 原点走時。屈折波の走時直線を距離ゼロへ外挿した切片である。直線の傾きと合わせて層厚等を推定する。実測した距離ゼロの初動時刻とは同一視しない。\par

\noindent\textbf{2：適切} ブラインド層。屈折法の初動走時に現れにくい層である。薄い中間層が初動を与えない場合や低速度層を区別し、用語の範囲は資料により異なる。\par

\noindent\textbf{3：適切} 音響インピーダンス。垂直入射の簡単なモデルではZ=ρVであり、境界両側の差が反射振幅を支配する。重力で用いる密度そのものだけとは異なる。\par

\noindent\textbf{4：適切} 静補正。地形や表層低速度層などによる走時ずれを一定の時間シフトとしてそろえる処理である。重力の地形・Bouguer補正と観測量が異なる。\par

\noindent\textbf{5：不適切} ブーゲー補正。測点と基準面の間の質量が重力へ及ぼす効果を補正する重力探査の処理である。地震波走時に対する静補正ではない。\par

\paragraph{再利用する考え方} 処理名だけでなく何を補正し、単位は何かを問う。

\paragraph{共通理論} \hyperref[theory:exploration:seismicterms]{第\ref*{theory:exploration:seismicterms}節}, \hyperref[theory:gravity:corrections]{第\ref*{theory:gravity:corrections}節}

\par 参考：\cite{usgs_refraction2013,epa_seismic_reflection,usgs_static2004,usgs_processing2000}。

\paragraph{関連問題} \hyperref[q:h24-reg-applied-23]{H24-REG-APPLIED-23}

\par 難易度 2（校準初稿）、優先度未評価。

\subsection{令和元年度 問34}
\label{q:r01-reg-applied-34}

\noindent\textbf{原典：} sources/R01/17\_applied\_science.pdf、PDF物理12頁。

\textbf{公式正答：} 4（single）。

\par\textbf{独立解答と公式正答に相違あり：} 独立解答 5。初回独立5対公式4を保持し資料確認後の訂正4を別記する。rigid-boom論文は検索収録の当該原理を確認し、PDF取得は403及びHTML challengeで未達である。日本語表記と装置体系の最終照合を査読対象とする。

\paragraph{問題の要旨} 静補正・深度走時・CMP・リジッドブーム・下方接続を区別する。

\Needspace{4\baselineskip}\paragraph{確認した条件・記述（独自要約）}\mbox{}\par

\noindent\textbf{原典照合：} 否定の対象、指数、周波数と周期、五つの候補の列順を原典画像で照合。\par

\noindent\textbf{前提：} 原典の4はリジットブームシステム、5は下方接続と表記される。下方接続を重力・磁気にしか使わない手法と限定しない。\par

\paragraph{問われる核心} 静補正・深度走時・CMP・リジッドブーム・下方接続を区別する。

\paragraph{必要知識} 原典の4はリジットブームシステム、5は下方接続と表記される。下方接続を重力・磁気にしか使わない手法と限定しない。

\paragraph{考え方・正答を選ぶ理由（査読待ち初稿）} 公式4である。初回5は下方接続をポテンシャル場だけに限定した誤認であった。Stanfordの波動下方接続とmigrationの説明を確認すると、地下へ波動場を外挿する処理も地震探査に関係する。リジッドブームEMIは送受信コイルを用いる電磁誘導計測の配置であり、地震データ処理の他四項目と異なる。

\Needspace{4\baselineskip}\paragraph{選択肢の検討（査読待ち初稿）}\mbox{}\par

\noindent\textbf{1：適切} 静補正。地形や表層低速度層などによる走時ずれを一定の時間シフトとしてそろえる処理である。重力の地形・Bouguer補正と観測量が異なる。\par

\noindent\textbf{2：適切} 深度走時。速度モデルと対応する深度での伝播時間を結ぶ量であり、地震探査の時間・深度対応に用いる。\par

\noindent\textbf{3：適切} CMP重合。共通中点の反射記録を集め、走時差を補正して重ねる反射法の処理である。強い傾斜や不均質では中点が同じ反射点を意味するとは限らない。\par

\noindent\textbf{4：不適切} リジットブームシステム。原典表記を保存する。rigid-boom EMIの一次論文では一次磁場と導体の二次磁場を送受信コイルで扱う電磁計測として記述される。地震波の走時処理ではない。\par

\noindent\textbf{5：適切} 下方接続。地震波動場を深部へ外挿してmigration画像を得る方法に関係する。重力・磁気のポテンシャル場にも下方接続があるため、名前の共通性だけで地震分野から除外しない。\par

\paragraph{再利用する考え方} 同じ数学的処理が複数の物理分野で使われることを許し、計測装置の物理原理を確認する。

\paragraph{共通理論} \hyperref[theory:exploration:seismicterms]{第\ref*{theory:exploration:seismicterms}節}, \hyperref[theory:exploration:continuation]{第\ref*{theory:exploration:continuation}節}, \hyperref[theory:exploration:electromagnetic]{第\ref*{theory:exploration:electromagnetic}節}

\par 参考：\cite{usgs_refraction2013,epa_seismic_reflection,usgs_static2004,usgs_processing2000,stanford_downward2001,rigid_boom2019}。

\paragraph{関連問題} \hyperref[q:r04-reg-applied-32]{R04-REG-APPLIED-32}

\par 難易度 3（校準初稿）、優先度未評価。

\subsection{令和4年度 問32}
\label{q:r04-reg-applied-32}

\noindent\textbf{原典：} sources/R04/17\_applied\_science.pdf、PDF物理11頁。

\textbf{公式正答：} 3（single）。

\paragraph{問題の要旨} 深度走時・migration・静補正・CMPとブーゲー補正を区別する。

\Needspace{4\baselineskip}\paragraph{確認した条件・記述（独自要約）}\mbox{}\par

\noindent\textbf{原典照合：} 否定の対象、指数、周波数と周期、五つの候補の列順を原典画像で照合。\par

\noindent\textbf{前提：} 走時・反射波形の処理と重力の質量補正を分ける。\par

\paragraph{問われる核心} 深度走時・migration・静補正・CMPとブーゲー補正を区別する。

\paragraph{必要知識} 走時・反射波形の処理と重力の質量補正を分ける。

\paragraph{考え方・正答を選ぶ理由（査読待ち初稿）} 3のブーゲー補正が不適切である。migrationは反射・回折記録を地下の反射位置へ戻す処理であり、波動場を下方接続する方法もある。

\Needspace{4\baselineskip}\paragraph{選択肢の検討（査読待ち初稿）}\mbox{}\par

\noindent\textbf{1：適切} 深度走時。速度モデルと深度から伝播時間を対応させる地震探査の量である。\par

\noindent\textbf{2：適切} マイグレーション。傾斜面や回折の見掛け位置を地下の適切な位置へ再配置する画像化処理である。観測時間を一律定数で換算するだけの処理ではない。\par

\noindent\textbf{3：不適切} ブーゲー補正。測点と基準面の間の質量が重力へ及ぼす効果を補正する重力探査の処理である。地震波走時に対する静補正ではない。\par

\noindent\textbf{4：適切} 静補正。地形や表層低速度層などによる走時ずれを一定の時間シフトとしてそろえる処理である。重力の地形・Bouguer補正と観測量が異なる。\par

\noindent\textbf{5：適切} CMP重合。共通中点の記録を集め、NMO等の走時差を整えて重ねる。単に未補正の任意波形を足す方法ではない。\par

\paragraph{再利用する考え方} 同じ補正という語でも観測量と誤差の原因を分ける。

\paragraph{共通理論} \hyperref[theory:exploration:seismicterms]{第\ref*{theory:exploration:seismicterms}節}, \hyperref[theory:exploration:continuation]{第\ref*{theory:exploration:continuation}節}, \hyperref[theory:gravity:corrections]{第\ref*{theory:gravity:corrections}節}

\par 参考：\cite{usgs_refraction2013,epa_seismic_reflection,usgs_static2004,usgs_processing2000,stanford_downward2001}。

\paragraph{関連問題} \hyperref[q:r01-reg-applied-34]{R01-REG-APPLIED-34}

\par 難易度 2（校準初稿）、優先度未評価。

\section{探査の重合とS/N平方根則}

\subsection{平成26年度 問21}
\label{q:h26-reg-applied-21}

\noindent\textbf{原典：} sources/H26/17\_applied\_science.pdf、PDF物理9頁。

\textbf{公式正答：} 3（single）。

\par\textbf{検討未確定：} 平方根則は振幅S/Nと無相関雑音・同位相信号の条件付きであり、群設置も方向・間隔に依存する。

\paragraph{問題の要旨} S/Nの平方根則・群設置・ADC・ノッチ・逆解析を区別する。

\Needspace{4\baselineskip}\paragraph{確認した条件・記述（独自要約）}\mbox{}\par

\noindent\textbf{原典照合：} 否定の対象、指数、周波数と周期、五つの候補の列順を原典画像で照合。\par

\noindent\textbf{前提：} 同じ信号に独立同分散の雑音が重なる場合の振幅S/Nを考える。ADCの標本化周波数は信号帯域に対して設定する。\par

\paragraph{問われる核心} S/Nの平方根則・群設置・ADC・ノッチ・逆解析を区別する。

\paragraph{必要知識} 同じ信号に独立同分散の雑音が重なる場合の振幅S/Nを考える。ADCの標本化周波数は信号帯域に対して設定する。

\paragraph{考え方・正答を選ぶ理由（査読待ち初稿）} 3が不適切である。標本化周波数を最高周波数と同じにすると、一般には異なる連続信号が同じ標本を与える。帯域制限を前提に二倍を超える周波数を用いる。他の候補は代表的な処理を説明する。

\Needspace{4\baselineskip}\paragraph{選択肢の検討（査読待ち初稿）}\mbox{}\par

\noindent\textbf{1：条件付きで適切} 重合のS/N改善が回数の平方根に比例する。同一位相信号はN倍、独立雑音の標準偏差は平方根N倍なので比は平方根N倍となる。\par

\noindent\textbf{2：条件付きで適切} 複数センサーの出力を群として結合し雑音を低減する。独立雑音の平均化に加え、センサーの間隔と入射方向に応じた空間フィルタ効果を利用する。対象の反射波を強めるには位相・配置が適切でなければならない。\par

\noindent\textbf{3：不適切} ADCの標本化周波数を信号の最高周波数と同じにする。一般の信号に必要なNyquist条件を満たさずaliasingを起こす。同じ周波数にするのは正しい一般的設定ではない。\par

\noindent\textbf{4：適切} ノッチで特定の周波数成分を遮断する。狭い阻止帯域で商用電源などの周期的な混入を抑える。指定帯域だけを通過させるフィルタではない。実装には有限の帯域幅と位相応答がある。\par

\noindent\textbf{5：適切} モデルの理論値と観測値の差を非線形最小二乗等で小さくする。順モデルF(m)の残差を目的関数にして地下物性を更新する代表的な逆解析である。線形問題や他のノルムもあり、非線形最小二乗だけが唯一の方法ではない。\par

\paragraph{再利用する考え方} 雑音を振幅で測るか電力で測るかを先に固定する。

\paragraph{共通理論} \hyperref[theory:exploration:stack]{第\ref*{theory:exploration:stack}節}, \hyperref[theory:exploration:sampling]{第\ref*{theory:exploration:sampling}節}, \hyperref[theory:exploration:inverse]{第\ref*{theory:exploration:inverse}節}

\par 参考：\cite{mit_signals2010,mit_moving2012,scipy_notch,scipy_fft,usgs_processing2000,pygimli_inversion}。

\paragraph{関連問題} \hyperref[q:h30-reg-applied-34]{H30-REG-APPLIED-34}

\par 難易度 3（校準初稿）、優先度未評価。

\subsection{平成27年度 問20}
\label{q:h27-reg-applied-20}

\noindent\textbf{原典：} sources/H27/17\_applied\_science.pdf、PDF物理7頁。

\textbf{公式正答：} 1（single）。

\par\textbf{検討未確定：} S/Nの振幅・電力定義が省略される。標本化の等号端点と完全復元、理想帯域通過の実装上の限界を区別する。

\paragraph{問題の要旨} 重合回数・コヒーレンス・標本化・帯域通過・FFTを比較する。

\Needspace{4\baselineskip}\paragraph{確認した条件・記述（独自要約）}\mbox{}\par

\noindent\textbf{原典照合：} 否定の対象、指数、周波数と周期、五つの候補の列順を原典画像で照合。\par

\noindent\textbf{前提：} 振幅S/N、周波数ごとの二信号の関係、理想フィルタを区別する。\par

\paragraph{問われる核心} 重合回数・コヒーレンス・標本化・帯域通過・FFTを比較する。

\paragraph{必要知識} 振幅S/N、周波数ごとの二信号の関係、理想フィルタを区別する。

\paragraph{考え方・正答を選ぶ理由（査読待ち初稿）} 1が不適切である。振幅S/Nは独立雑音のとき平方根N倍であり、N倍ではない。コヒーレンスは相互スペクトルを両信号の自己スペクトルで規格化して比較する。

\Needspace{4\baselineskip}\paragraph{選択肢の検討（査読待ち初稿）}\mbox{}\par

\noindent\textbf{1：不適切} スタックのS/Nが回数Nに比例する。信号はN倍、雑音の振幅尺度は平方根N倍なので振幅S/Nは平方根N倍である。電力比ならN倍となるため比の定義を明示する。\par

\noindent\textbf{2：適切} コヒーレンスは二つの時系列の周波数ごとの類似を示す相関係数である。周波数別の結び付きであり、代表的な二乗コヒーレンスは相互スペクトルの絶対値二乗を自己スペクトルの積で割る。時間領域の遅れ別相関そのものとは区別する。\par

\noindent\textbf{3：条件付きで適切} 最高周波数の二倍以上で標本化すれば完全復元できる。帯域制限・理想的標本と再構成を想定した説明である。最高周波数成分を含む場合の等号端点、量子化、有限長、雑音を無視できない。安全な条件はfs>2fmaxであり、境界では位相によって標本が全てゼロとなる例がある。\par

\noindent\textbf{4：条件付きで適切} 帯域通過は設定帯域を減衰させず、それ以外を遮断する。理想帯域通過の定義として読む。実際の設計では通過帯域のリップル、遷移帯域、有限減衰を持つため完全遮断を一般保証しない。\par

\noindent\textbf{5：適切} FFTは二のべき乗の標本数で効率的である。radix-2では段数がlog2Nとなる。二のべき乗でなければ計算不能という意味ではなく、混合基数やBluestein法もある。\par

\paragraph{再利用する考え方} 振幅比と電力比、理想定義と実装の差を明示する。

\paragraph{共通理論} \hyperref[theory:exploration:stack]{第\ref*{theory:exploration:stack}節}, \hyperref[theory:exploration:correlation]{第\ref*{theory:exploration:correlation}節}, \hyperref[theory:exploration:sampling]{第\ref*{theory:exploration:sampling}節}, \hyperref[theory:exploration:fft]{第\ref*{theory:exploration:fft}節}

\par 参考：\cite{mit_signals2010,mit_moving2012,scipy_notch,scipy_fft,usgs_processing2000,scipy_coherence}。

\par 難易度 3（校準初稿）、優先度未評価。

\subsection{平成30年度 問34}
\label{q:h30-reg-applied-34}

\noindent\textbf{原典：} sources/H30/17\_applied\_science.pdf、PDF物理11頁。

\textbf{公式正答：} 1（single）。

\par\textbf{検討未確定：} 群設置と平方根則は雑音相関・信号整列に依存し、標本化の等号と帯域制限を省略して完全復元を保証できない。

\paragraph{問題の要旨} 重合のS/N指数と電源雑音・群設置・標本化・逆解析を照合する。

\Needspace{4\baselineskip}\paragraph{確認した条件・記述（独自要約）}\mbox{}\par

\noindent\textbf{原典照合：} 否定の対象、指数、周波数と周期、五つの候補の列順を原典画像で照合。\par

\noindent\textbf{前提：} 振幅S/Nは同一位相信号と独立雑音を想定する。\par

\paragraph{問われる核心} 重合のS/N指数と電源雑音・群設置・標本化・逆解析を照合する。

\paragraph{必要知識} 振幅S/Nは同一位相信号と独立雑音を想定する。

\paragraph{考え方・正答を選ぶ理由（査読待ち初稿）} 1が不適切である。平方根Nの向上をNの二乗とする誤りである。電力S/NでもN倍でありNの二乗にはならない。周期雑音には狭帯域阻止、物性推定には順モデルの残差最小化を対応させる。

\Needspace{4\baselineskip}\paragraph{選択肢の検討（査読待ち初稿）}\mbox{}\par

\noindent\textbf{1：不適切} スタックのS/N改善が回数Nの二乗に比例する。信号の振幅Nに対して雑音の標準偏差は平方根Nである。振幅比も電力比もNの二乗とはならない。\par

\noindent\textbf{2：条件付きで適切} 複数センサーの出力を群として結合し雑音を低減する。独立雑音の平均化に加え、センサーの間隔と入射方向に応じた空間フィルタ効果を利用する。対象の反射波を強めるには位相・配置が適切でなければならない。\par

\noindent\textbf{3：適切} ノッチで特定の周波数成分を遮断する。狭い阻止帯域で商用電源などの周期的な混入を抑える。指定帯域だけを通過させるフィルタではない。実装には有限の帯域幅と位相応答がある。\par

\noindent\textbf{4：条件付きで適切} 最高周波数の二倍以上で標本化すれば完全復元できる。帯域制限・理想的標本と再構成を想定した説明である。最高周波数成分を含む場合の等号端点、量子化、有限長、雑音を無視できない。安全な条件はfs>2fmaxであり、境界では位相によって標本が全てゼロとなる例がある。\par

\noindent\textbf{5：適切} モデルの理論値と観測値の差を非線形最小二乗等で小さくする。順モデルF(m)の残差を目的関数にして地下物性を更新する代表的な逆解析である。線形問題や他のノルムもあり、非線形最小二乗だけが唯一の方法ではない。\par

\paragraph{再利用する考え方} 加算した信号と雑音を別々に計算して指数の誤りを検出する。

\paragraph{共通理論} \hyperref[theory:exploration:stack]{第\ref*{theory:exploration:stack}節}, \hyperref[theory:exploration:filters]{第\ref*{theory:exploration:filters}節}, \hyperref[theory:exploration:inverse]{第\ref*{theory:exploration:inverse}節}

\par 参考：\cite{mit_signals2010,mit_moving2012,scipy_notch,scipy_fft,usgs_processing2000,pygimli_inversion}。

\paragraph{関連問題} \hyperref[q:h26-reg-applied-21]{H26-REG-APPLIED-21}, \hyperref[q:r04-reg-applied-33]{R04-REG-APPLIED-33}

\par 難易度 3（校準初稿）、優先度未評価。

\subsection{令和4年度 問33}
\label{q:r04-reg-applied-33}

\noindent\textbf{原典：} sources/R04/17\_applied\_science.pdf、PDF物理11頁。

\textbf{公式正答：} 5（single）。

\par\textbf{検討未確定：} 完全復元の等号端点・帯域制限、群配置の方向と間隔を要する。表面波抑制を独立雑音の平方根則だけで説明しない。

\paragraph{問題の要旨} モデル更新・ノッチ・受振器群・標本化・垂直重合を区別する。

\Needspace{4\baselineskip}\paragraph{確認した条件・記述（独自要約）}\mbox{}\par

\noindent\textbf{原典照合：} 否定の対象、指数、周波数と周期、五つの候補の列順を原典画像で照合。\par

\noindent\textbf{前提：} 同一起振点の反復で対象反射の位相をそろえる。群配置は対象波と妨害波の空間応答を考える。\par

\paragraph{問われる核心} モデル更新・ノッチ・受振器群・標本化・垂直重合を区別する。

\paragraph{必要知識} 同一起振点の反復で対象反射の位相をそろえる。群配置は対象波と妨害波の空間応答を考える。

\paragraph{考え方・正答を選ぶ理由（査読待ち初稿）} 5が不適切である。垂直重合の振幅S/Nは理論上平方根N倍でありNの二乗ではない。表面波は反射波に対する妨害波となるが、空間的に相関した波なので独立ランダム雑音の平均化と同じ説明だけでは足りない。

\Needspace{4\baselineskip}\paragraph{選択肢の検討（査読待ち初稿）}\mbox{}\par

\noindent\textbf{1：適切} モデルの理論値と観測値の差を非線形最小二乗等で小さくする。順モデルF(m)の残差を目的関数にして地下物性を更新する代表的な逆解析である。線形問題や他のノルムもあり、非線形最小二乗だけが唯一の方法ではない。\par

\noindent\textbf{2：適切} ノッチで特定の周波数成分を遮断する。狭い阻止帯域で商用電源などの周期的な混入を抑える。指定帯域だけを通過させるフィルタではない。実装には有限の帯域幅と位相応答がある。\par

\noindent\textbf{3：条件付きで適切} 受振器群で表面波を抑え反射波を強めランダム雑音を減らす。複数受振器の結合には空間応答があり、反射と表面波の見掛け速度・波長に応じた選別を行える。配列が悪ければ必要信号も弱まるため無条件ではない。\par

\noindent\textbf{4：条件付きで適切} 最高周波数の二倍以上で標本化すれば完全復元できる。帯域制限・理想的標本と再構成を想定した説明である。最高周波数成分を含む場合の等号端点、量子化、有限長、雑音を無視できない。安全な条件はfs>2fmaxであり、境界では位相によって標本が全てゼロとなる例がある。\par

\noindent\textbf{5：不適切} 垂直スタックのS/Nが回数Nの二乗に比例する。独立同分散雑音の振幅S/N改善は平方根N倍である。Nの二乗は信号電力だけの増加を雑音の増加と比較せず取り出した誤りである。\par

\paragraph{再利用する考え方} 信号の種類、統計的相関、配列の空間応答を分ける。

\paragraph{共通理論} \hyperref[theory:exploration:stack]{第\ref*{theory:exploration:stack}節}, \hyperref[theory:exploration:filters]{第\ref*{theory:exploration:filters}節}, \hyperref[theory:exploration:inverse]{第\ref*{theory:exploration:inverse}節}

\par 参考：\cite{mit_signals2010,mit_moving2012,scipy_notch,scipy_fft,usgs_processing2000,pygimli_inversion}。

\paragraph{関連問題} \hyperref[q:h30-reg-applied-34]{H30-REG-APPLIED-34}

\par 難易度 3（校準初稿）、優先度未評価。

\section{表面波分散・S波感度と統合逆解析}

\subsection{令和7年度 問35}
\label{q:r07-reg-applied-35}

\noindent\textbf{原典：} sources/R07/17\_applied\_science.pdf、PDF物理13頁。

\textbf{公式正答：} 2（single）。

\par\textbf{検討未確定：} 均質半無限体は非分散の退化例であり1を全称にしない。P波・密度の感度はゼロではない。微動併用で長波長が得られることとモード同定の確認が必要。

\paragraph{問題の要旨} 表面波の分散、物性感度、反復逆解析、波長別深度、能動と微動の併用を比較する。

\Needspace{4\baselineskip}\paragraph{確認した条件・記述（独自要約）}\mbox{}\par

\noindent\textbf{原典照合：} 五候補の順序、最適切と不適切の方向、波種、深度目安、曲線名、最長と最短を原典画像で照合。\par

\noindent\textbf{前提：} 通常の不均質水平成層モデルとRayleigh分散の主なS波速度感度を考える。\par

\paragraph{問われる核心} 表面波の分散、物性感度、反復逆解析、波長別深度、能動と微動の併用を比較する。

\paragraph{必要知識} 通常の不均質水平成層モデルとRayleigh分散の主なS波速度感度を考える。

\paragraph{考え方・正答を選ぶ理由（査読待ち初稿）} 2が不適切である。Rayleigh分散は一般にS波速度へ強い感度を持ち、P波の影響がS波より大きいとする説明は逆である。波長を広い範囲で得ると浅部と深部を拘束しやすくなるが、観測に合うモデルが唯一とは限らない。

\Needspace{4\baselineskip}\paragraph{選択肢の検討（査読待ち初稿）}\mbox{}\par

\noindent\textbf{1：条件付きで適切} 水平成層のRayleigh位相速度は周波数による分散を持つ。深度で物性が変わる成層の通常モデルでは分散を持つ。均質等方半無限体や全層同物性の退化例は非分散なので、全てのモデルへの全称としない。\par

\noindent\textbf{2：不適切} Rayleigh位相速度へのP波の影響がS波より大きい。通常の地盤分散はS波速度が支配的に影響し、P波だけを主たる未知量とする説明は誤りである。\par

\noindent\textbf{3：適切} 理論と観測の位相速度差を減らすようモデル修正と計算を反復する。分散曲線を順計算し地下構造パラメータを更新する逆問題である。初期値・誤差と拘束にも依存する。\par

\noindent\textbf{4：適切} 探査深度は最大波長、浅部解像度は最小波長と強く関係する。長波長ほど深部に感度を持ち、短波長は浅部の小さな構造を拘束する傾向がある。一波長が一深度だけを測るわけではない。\par

\noindent\textbf{5：条件付きで適切} 能動表面波と微動を併用して探査深度を増やす。微動で長波長側を補える観測条件なら深部の拘束が増す。低周波が得られない場合やモードが不整合の場合は併用だけで深度増加を保証しない。\par

\paragraph{再利用する考え方} 分散の主感度、波長の感度範囲、データ併用の条件を示し唯一解を仮定しない。

\paragraph{共通理論} \hyperref[theory:exploration:surfacewaves]{第\ref*{theory:exploration:surfacewaves}節}, \hyperref[theory:exploration:arraylimits]{第\ref*{theory:exploration:arraylimits}節}, \hyperref[theory:exploration:inverse]{第\ref*{theory:exploration:inverse}節}

\par 参考：\cite{epa_surfacewaves,xia_rayleigh1999,japansea_spac2017,array_limits2002,pygimli_inversion}。

\paragraph{関連問題} \hyperref[q:r06-reg-applied-33]{R06-REG-APPLIED-33}

\par 難易度 3（校準初稿）、優先度未評価。

\section{探査波動の伝播原理}

\subsection{令和2年度 問35}
\label{q:r02-reg-applied-35}

\noindent\textbf{原典：} sources/R02/17\_applied\_science.pdf、PDF物理12頁。

\textbf{公式正答：} 2（single）。

\par\textbf{検討未確定：} 相反性と重ね合わせは媒体・線形性の条件付きである。等価層は異なる分野で使われる語でもあるため、ポテンシャル場の基本的用法と波動伝播の原理を区別する。

\paragraph{問題の要旨} 弾性波の伝播原理とポテンシャル場の等価層を区別する。

\Needspace{4\baselineskip}\paragraph{確認した条件・記述（独自要約）}\mbox{}\par

\noindent\textbf{原典照合：} 否定の対象、指数、周波数と周期、五つの候補の列順を原典画像で照合。\par

\noindent\textbf{前提：} 線形で相反的な媒体における波動の重ね合わせ・波面構成・伝播経路を対象とする。\par

\paragraph{問われる核心} 弾性波の伝播原理とポテンシャル場の等価層を区別する。

\paragraph{必要知識} 線形で相反的な媒体における波動の重ね合わせ・波面構成・伝播経路を対象とする。

\paragraph{考え方・正答を選ぶ理由（査読待ち初稿）} 2の等価層の原理が不適切である。USGSの等価層は重力の密度層や磁気の双極子層でポテンシャル場を表現するものとして確認できる。これは地震波伝播を記述する基本原理の四項目とは役割が異なる。

\Needspace{4\baselineskip}\paragraph{選択肢の検討（査読待ち初稿）}\mbox{}\par

\noindent\textbf{1：適切} フェルマの原理。波線経路に対する走時の停留条件を述べる。あらゆる経路が絶対的最短時間になるという一般化は避ける。\par

\noindent\textbf{2：不適切} 等価層の原理。重力・磁気の場を等価な密度又は双極子の層で表す考え方である。弾性波の波面や伝播経路を定める基本則とは異なる。\par

\noindent\textbf{3：適切} 相反原理。相反性が成立する線形媒体では震源と受振点を交換した応答が対応する。媒体や成分の条件を保つ必要がある。\par

\noindent\textbf{4：適切} ホイヘンスの原理。波面上の点からの二次波の構成で次の波面を考える。地震波の伝播・画像化の理解に関係する。\par

\noindent\textbf{5：条件付きで適切} 重ね合わせの原理。線形の波動方程式では解の線形結合も解となる。破壊を伴う大きな非線形変形へ無条件に適用しない。\par

\paragraph{再利用する考え方} 原理の名前を、場の表現・波面・経路・線形性・相反性という役割へ対応させる。

\paragraph{共通理論} \hyperref[theory:exploration:principles]{第\ref*{theory:exploration:principles}節}

\par 参考：\cite{usgs_equivalent2007,stanford_downward2001,mit_signals2010}。

\paragraph{関連問題} \hyperref[q:r07-reg-applied-32]{R07-REG-APPLIED-32}

\par 難易度 3（校準初稿）、優先度未評価。

\subsection{令和7年度 問32}
\label{q:r07-reg-applied-32}

\noindent\textbf{原典：} sources/R07/17\_applied\_science.pdf、PDF物理12頁。

\textbf{公式正答：} 3（single）。

\par\textbf{検討未確定：} 重ね合わせは線形媒体に限る。スネルの法則は波種と境界条件を指定して用いる。

\paragraph{問題の要旨} ホイヘンス・重ね合わせ・レンツ・スネル・フェルマを区別する。

\Needspace{4\baselineskip}\paragraph{確認した条件・記述（独自要約）}\mbox{}\par

\noindent\textbf{原典照合：} 否定の対象、指数、周波数と周期、五つの候補の列順を原典画像で照合。\par

\noindent\textbf{前提：} 地震探査で扱う弾性波の波面、線形性、屈折、波線の走時を対象とする。\par

\paragraph{問われる核心} ホイヘンス・重ね合わせ・レンツ・スネル・フェルマを区別する。

\paragraph{必要知識} 地震探査で扱う弾性波の波面、線形性、屈折、波線の走時を対象とする。

\paragraph{考え方・正答を選ぶ理由（査読待ち初稿）} 3のレンツの法則は電磁誘導が磁束変化を妨げる向きに生じることを表す。電磁探査には関係するが、他四項目のような弾性波伝播の基本則ではない。

\Needspace{4\baselineskip}\paragraph{選択肢の検討（査読待ち初稿）}\mbox{}\par

\noindent\textbf{1：適切} ホイヘンスの原理。二次波で波面を構成する考え方であり、波動伝播の理解に関係する。\par

\noindent\textbf{2：条件付きで適切} 重ね合わせの原理。線形波動では複数波の和を解として扱うことができる。大変形や非線形な破壊過程は同じ前提ではない。\par

\noindent\textbf{3：不適切} レンツの法則。電磁誘導の向きを磁束変化との関係から決める法則であり、地震波の反射・屈折の基本則とは異なる。\par

\noindent\textbf{4：適切} スネルの法則。境界の位相整合により角度と波速を結ぶ。地震波ではP波・S波へのモード変換を含む各波速で考える。\par

\noindent\textbf{5：適切} フェルマの原理。波線の走時が経路変分に対して停留する条件を用いる。最小という単語だけで複数経路の可能性を捨てない。\par

\paragraph{再利用する考え方} 電磁気の誘導と弾性波の伝播を、扱う物理量から分ける。

\paragraph{共通理論} \hyperref[theory:exploration:principles]{第\ref*{theory:exploration:principles}節}, \hyperref[theory:optics:refraction]{第\ref*{theory:optics:refraction}節}

\par 参考：\cite{mit_signals2010,stanford_downward2001,epa_seismic_reflection}。

\paragraph{関連問題} \hyperref[q:r02-reg-applied-35]{R02-REG-APPLIED-35}

\par 難易度 3（校準初稿）、優先度未評価。


\chapter{地質の補足}
\noindent 補足の章骨格。現在のDBにこの章への主分類割当はなく、追加の出題実績を意味しない。本文は未執筆。




\chapter{出題頻度（暫定集計）}
分類未確認を含む全登録値と、確認済みの部分集合を区別する。16試験回、15年度が分母。難易度未知を0として平均しない。
\begin{longtable}{p{.16\linewidth}p{.42\linewidth}rrr}
科目 & 大分類 & 暫定問数 & 確認済 & 出題回数\\\hline
APPLIED & 地球物理及び地球化学 & 98 & 1 & 16\\
APPLIED & 地質 & 179 & 1 & 16\\
APPLIED & 物理及び化学 & 283 & 5 & 16\\
APT & コンプライアンス・説明責任 & 15 & 1 & 9\\
APT & 安全・リスク & 69 & 2 & 16\\
APT & 技術と社会・科学技術政策 & 1 & 0 & 1\\
APT & 技術士法・義務責務 & 26 & 0 & 16\\
APT & 技術者倫理・意思決定 & 26 & 3 & 12\\
APT & 環境・持続可能性 & 44 & 0 & 16\\
APT & 研究・情報倫理 & 59 & 0 & 16\\
BASIC & 情報・論理 & 94 & 2 & 16\\
BASIC & 材料・化学・バイオ & 94 & 2 & 16\\
BASIC & 環境・エネルギー・技術 & 94 & 12 & 16\\
BASIC & 解析 & 94 & 18 & 16\\
BASIC & 設計・計画 & 94 & 1 & 16\\
\end{longtable}

\chapter{頻度を読む}
試験回頻度と年度頻度を区別する。再試験を独立回として扱う一方、年度の出現数では令和元年度を一つとして数える。未確認の分類を含む暫定数を学習優先度へ自動変換しない。

\chapter{年度推移と比較の分母}
平成23・24年度は基礎25問、それ以後は基礎30問である。適性15問、応用理学35問は各試験回で同じ登録範囲を使う。出題数の変化を概念の流行と判断する前に分母を確認する。
\begin{longtable}{lrrrr}
試験回 & 基礎 & 適性 & 応用理学 & 合計\\\hline
H23 & 25 & 15 & 35 & 75\\
H24 & 25 & 15 & 35 & 75\\
H25 & 30 & 15 & 35 & 80\\
H26 & 30 & 15 & 35 & 80\\
H27 & 30 & 15 & 35 & 80\\
H28 & 30 & 15 & 35 & 80\\
H29 & 30 & 15 & 35 & 80\\
H30 & 30 & 15 & 35 & 80\\
R01 & 30 & 15 & 35 & 80\\
R01\_retest & 30 & 15 & 35 & 80\\
R02 & 30 & 15 & 35 & 80\\
R03 & 30 & 15 & 35 & 80\\
R04 & 30 & 15 & 35 & 80\\
R05 & 30 & 15 & 35 & 80\\
R06 & 30 & 15 & 35 & 80\\
R07 & 30 & 15 & 35 & 80\\
\end{longtable}
\section{令和3〜7年度の五回と全期間の比較}
後期五回は基礎150、適性75、応用理学175問を分母にする。全期間の構成比と後期の構成比は、同じ主分類で比較する。数ポイントの差だけで単調な増減や次年度の出題を断定しない。
\begin{longtable}{lp{.40\linewidth}rrr}
科目 & 中分類 & 全期間\% & 後期問数 & 後期\%\\\hline
基礎 & エネルギー & 6.2 & 9 & 6.0\\
基礎 & 生物・バイオ & 6.0 & 10 & 6.7\\
基礎 & 科学技術史・社会 & 5.5 & 9 & 6.0\\
基礎 & 材料・金属 & 5.5 & 9 & 6.0\\
基礎 & 環境・持続可能性 & 5.5 & 10 & 6.7\\
基礎 & 微積分・数値解析 & 5.5 & 8 & 5.3\\
適性 & 安全・リスク & 24.6 & 16 & 21.3\\
適性 & コンプライアンス・説明責任 & 22.1 & 17 & 22.7\\
適性 & 技術者倫理・意思決定 & 17.1 & 11 & 14.7\\
適性 & 技術士法・義務責務 & 12.1 & 8 & 10.7\\
適性 & 研究・情報倫理 & 12.1 & 8 & 10.7\\
適性 & 環境・持続可能性 & 7.9 & 9 & 12.0\\
応用理学 & 物理探査 & 12.1 & 21 & 12.0\\
応用理学 & 物理化学 & 8.6 & 16 & 9.1\\
応用理学 & 電磁気・物性 & 7.1 & 14 & 8.0\\
応用理学 & 有機・高分子化学 & 5.9 & 8 & 4.6\\
応用理学 & 応用地質・地盤 & 4.8 & 4 & 2.3\\
応用理学 & 力学 & 4.6 & 9 & 5.1\\
\end{longtable}
再試験を通常年の傾向へ二重に算入する効果、分類境界の変更、主題が同じでも問い方が変わる効果を区別する。ここでは長期構成と後期構成の記述的比較までを行い、有意差や予測確率は計算していない。

\chapter{類似問題と再出題}
同概念、同解法、条件違い、近似再出題、完全再出題を区別する。ばね系の比較例では同じ接続図と条件を確認して近似再出題を承認した。OCR類似度だけの候補は未承認として残す。

\chapter{学習順序を頻度と転用性から組み立てる}
大学院物理の背景がある読者は、既知の力学・熱・電磁気を原典の条件確認で回収し、情報・化学・地球科学・制度の新規概念へ時間を配分する。次の列挙は出題数による着手候補であり、習得時間を測ったS/A/B/Cの確定優先度ではない。難易度は校準初稿、学習時間と得点効率は未評価として残す。
\section{基礎の頻度上位に接続する基礎}
エネルギー（29問）、生物・バイオ（28問）、科学技術史・社会（26問）、材料・金属（26問）、環境・持続可能性（26問）を最初の点検対象とする。
\section{適性の頻度上位に接続する基礎}
安全・リスク（59問）、コンプライアンス・説明責任（53問）、技術者倫理・意思決定（41問）、技術士法・義務責務（29問）、研究・情報倫理（29問）を最初の点検対象とする。
\section{応用理学の頻度上位に接続する基礎}
物理探査（68問）、物理化学（48問）、電磁気・物性（40問）、有機・高分子化学（33問）、応用地質・地盤（27問）を最初の点検対象とする。
基礎は保存則、単位、論理、確率と待ち行列の成立条件を共通の道具にする。適性は技術士法の義務責務と、公益・安全・説明責任を別の根拠で整理する。応用理学は波の伝播・場・収支を物理探査へつなぎ、鉱物・岩石・地層・年代・同位体を用語と観測証拠から学ぶ。
頻度の低い問題も索引と解説から除外しない。高リスクの法令や歴史統計は解説の留保を先に読み、出題時点の原典と現在の参考説明を分ける。頻度だけで誤った初稿を反復学習しない。

\chapter{学習計画への接続}
既知の数学・物理を使える問題は原典条件を短く確認し、新規の化学・情報・制度概念に学習時間を配分する。未検証の得点効率から選択戦略を確定しない。

\appendix
\chapter{校正済み例に用いた式}
\section{線形ばね系}
ばねの微小変位に対し復元力が $-k_{\rm eff}x$ で、質量が$m$の一自由度系では
\[m\ddot{x}+k_{\rm eff}x=0,\qquad \omega_0=\sqrt{k_{\rm eff}/m},\qquad f_0=\omega_0/(2\pi).\]
直列接続では共通の力に対する伸びを加えるため
$1/k_{\rm eff}=\sum_i1/k_i$、並列接続では共通の伸びに対する力を加えるため$k_{\rm eff}=\sum_i k_i$となる。
重力などの定常力が平衡点を移動させても、その点の周囲の線形復元剛性が同じなら固有振動数は変わらない。

\chapter{初稿の用語}
\begin{description}
\item[LIFO] 後から入れた要素を先に取り出す規則。スタックの基本操作に用いる。
\item[固有振動数] 外部駆動のない線形振動系のモードに対応する振動数。材質だけでなく形状・質量・支持条件に依存する。
\item[系統分析] 複数成分を、条件を変えた分離操作によって段階的に識別する分析。
\item[主分類] 1問につき一つ割り当てる学習分野。横断概念のタグとは別に集計する。
\item[全員得点] 公式に受験者全員へ得点を与える取扱い。全選択肢が正しいという意味ではない。
\end{description}

\chapter{全問の収録方針}
年度別索引から全1,270問の登録項目へ移動できる。独自解説が未作成の問題は、原典情報と未完了状態のみを表示している。
全問の原文画像をこのPDFへ転載する方式は採らず、原典と独自要約・解説を識別する。

\chapter{年度別索引}
全1,270問の登録索引。リンク先の校正・執筆状態を確認する。
\section{平成23年度}
\paragraph{基礎（25問）} \hyperref[q:h23-reg-basic-1-1]{1-1}, \hyperref[q:h23-reg-basic-1-2]{1-2}, \hyperref[q:h23-reg-basic-1-3]{1-3}, \hyperref[q:h23-reg-basic-1-4]{1-4}, \hyperref[q:h23-reg-basic-1-5]{1-5}, \hyperref[q:h23-reg-basic-2-1]{2-1}, \hyperref[q:h23-reg-basic-2-2]{2-2}, \hyperref[q:h23-reg-basic-2-3]{2-3}, \hyperref[q:h23-reg-basic-2-4]{2-4}, \hyperref[q:h23-reg-basic-2-5]{2-5}, \hyperref[q:h23-reg-basic-3-1]{3-1}, \hyperref[q:h23-reg-basic-3-2]{3-2}, \hyperref[q:h23-reg-basic-3-3]{3-3}, \hyperref[q:h23-reg-basic-3-4]{3-4}, \hyperref[q:h23-reg-basic-3-5]{3-5}, \hyperref[q:h23-reg-basic-4-1]{4-1}, \hyperref[q:h23-reg-basic-4-2]{4-2}, \hyperref[q:h23-reg-basic-4-3]{4-3}, \hyperref[q:h23-reg-basic-4-4]{4-4}, \hyperref[q:h23-reg-basic-4-5]{4-5}, \hyperref[q:h23-reg-basic-5-1]{5-1}, \hyperref[q:h23-reg-basic-5-2]{5-2}, \hyperref[q:h23-reg-basic-5-3]{5-3}, \hyperref[q:h23-reg-basic-5-4]{5-4}, \hyperref[q:h23-reg-basic-5-5]{5-5}
\paragraph{適性（15問）} \hyperref[q:h23-reg-apt-01]{1}, \hyperref[q:h23-reg-apt-02]{2}, \hyperref[q:h23-reg-apt-03]{3}, \hyperref[q:h23-reg-apt-04]{4}, \hyperref[q:h23-reg-apt-05]{5}, \hyperref[q:h23-reg-apt-06]{6}, \hyperref[q:h23-reg-apt-07]{7}, \hyperref[q:h23-reg-apt-08]{8}, \hyperref[q:h23-reg-apt-09]{9}, \hyperref[q:h23-reg-apt-10]{10}, \hyperref[q:h23-reg-apt-11]{11}, \hyperref[q:h23-reg-apt-12]{12}, \hyperref[q:h23-reg-apt-13]{13}, \hyperref[q:h23-reg-apt-14]{14}, \hyperref[q:h23-reg-apt-15]{15}
\paragraph{応用理学（35問）} \hyperref[q:h23-reg-applied-01]{1}, \hyperref[q:h23-reg-applied-02]{2}, \hyperref[q:h23-reg-applied-03]{3}, \hyperref[q:h23-reg-applied-04]{4}, \hyperref[q:h23-reg-applied-05]{5}, \hyperref[q:h23-reg-applied-06]{6}, \hyperref[q:h23-reg-applied-07]{7}, \hyperref[q:h23-reg-applied-08]{8}, \hyperref[q:h23-reg-applied-09]{9}, \hyperref[q:h23-reg-applied-10]{10}, \hyperref[q:h23-reg-applied-11]{11}, \hyperref[q:h23-reg-applied-12]{12}, \hyperref[q:h23-reg-applied-13]{13}, \hyperref[q:h23-reg-applied-14]{14}, \hyperref[q:h23-reg-applied-15]{15}, \hyperref[q:h23-reg-applied-16]{16}, \hyperref[q:h23-reg-applied-17]{17}, \hyperref[q:h23-reg-applied-18]{18}, \hyperref[q:h23-reg-applied-19]{19}, \hyperref[q:h23-reg-applied-20]{20}, \hyperref[q:h23-reg-applied-21]{21}, \hyperref[q:h23-reg-applied-22]{22}, \hyperref[q:h23-reg-applied-23]{23}, \hyperref[q:h23-reg-applied-24]{24}, \hyperref[q:h23-reg-applied-25]{25}, \hyperref[q:h23-reg-applied-26]{26}, \hyperref[q:h23-reg-applied-27]{27}, \hyperref[q:h23-reg-applied-28]{28}, \hyperref[q:h23-reg-applied-29]{29}, \hyperref[q:h23-reg-applied-30]{30}, \hyperref[q:h23-reg-applied-31]{31}, \hyperref[q:h23-reg-applied-32]{32}, \hyperref[q:h23-reg-applied-33]{33}, \hyperref[q:h23-reg-applied-34]{34}, \hyperref[q:h23-reg-applied-35]{35}
\section{平成24年度}
\paragraph{基礎（25問）} \hyperref[q:h24-reg-basic-1-1]{1-1}, \hyperref[q:h24-reg-basic-1-2]{1-2}, \hyperref[q:h24-reg-basic-1-3]{1-3}, \hyperref[q:h24-reg-basic-1-4]{1-4}, \hyperref[q:h24-reg-basic-1-5]{1-5}, \hyperref[q:h24-reg-basic-2-1]{2-1}, \hyperref[q:h24-reg-basic-2-2]{2-2}, \hyperref[q:h24-reg-basic-2-3]{2-3}, \hyperref[q:h24-reg-basic-2-4]{2-4}, \hyperref[q:h24-reg-basic-2-5]{2-5}, \hyperref[q:h24-reg-basic-3-1]{3-1}, \hyperref[q:h24-reg-basic-3-2]{3-2}, \hyperref[q:h24-reg-basic-3-3]{3-3}, \hyperref[q:h24-reg-basic-3-4]{3-4}, \hyperref[q:h24-reg-basic-3-5]{3-5}, \hyperref[q:h24-reg-basic-4-1]{4-1}, \hyperref[q:h24-reg-basic-4-2]{4-2}, \hyperref[q:h24-reg-basic-4-3]{4-3}, \hyperref[q:h24-reg-basic-4-4]{4-4}, \hyperref[q:h24-reg-basic-4-5]{4-5}, \hyperref[q:h24-reg-basic-5-1]{5-1}, \hyperref[q:h24-reg-basic-5-2]{5-2}, \hyperref[q:h24-reg-basic-5-3]{5-3}, \hyperref[q:h24-reg-basic-5-4]{5-4}, \hyperref[q:h24-reg-basic-5-5]{5-5}
\paragraph{適性（15問）} \hyperref[q:h24-reg-apt-01]{1}, \hyperref[q:h24-reg-apt-02]{2}, \hyperref[q:h24-reg-apt-03]{3}, \hyperref[q:h24-reg-apt-04]{4}, \hyperref[q:h24-reg-apt-05]{5}, \hyperref[q:h24-reg-apt-06]{6}, \hyperref[q:h24-reg-apt-07]{7}, \hyperref[q:h24-reg-apt-08]{8}, \hyperref[q:h24-reg-apt-09]{9}, \hyperref[q:h24-reg-apt-10]{10}, \hyperref[q:h24-reg-apt-11]{11}, \hyperref[q:h24-reg-apt-12]{12}, \hyperref[q:h24-reg-apt-13]{13}, \hyperref[q:h24-reg-apt-14]{14}, \hyperref[q:h24-reg-apt-15]{15}
\paragraph{応用理学（35問）} \hyperref[q:h24-reg-applied-01]{1}, \hyperref[q:h24-reg-applied-02]{2}, \hyperref[q:h24-reg-applied-03]{3}, \hyperref[q:h24-reg-applied-04]{4}, \hyperref[q:h24-reg-applied-05]{5}, \hyperref[q:h24-reg-applied-06]{6}, \hyperref[q:h24-reg-applied-07]{7}, \hyperref[q:h24-reg-applied-08]{8}, \hyperref[q:h24-reg-applied-09]{9}, \hyperref[q:h24-reg-applied-10]{10}, \hyperref[q:h24-reg-applied-11]{11}, \hyperref[q:h24-reg-applied-12]{12}, \hyperref[q:h24-reg-applied-13]{13}, \hyperref[q:h24-reg-applied-14]{14}, \hyperref[q:h24-reg-applied-15]{15}, \hyperref[q:h24-reg-applied-16]{16}, \hyperref[q:h24-reg-applied-17]{17}, \hyperref[q:h24-reg-applied-18]{18}, \hyperref[q:h24-reg-applied-19]{19}, \hyperref[q:h24-reg-applied-20]{20}, \hyperref[q:h24-reg-applied-21]{21}, \hyperref[q:h24-reg-applied-22]{22}, \hyperref[q:h24-reg-applied-23]{23}, \hyperref[q:h24-reg-applied-24]{24}, \hyperref[q:h24-reg-applied-25]{25}, \hyperref[q:h24-reg-applied-26]{26}, \hyperref[q:h24-reg-applied-27]{27}, \hyperref[q:h24-reg-applied-28]{28}, \hyperref[q:h24-reg-applied-29]{29}, \hyperref[q:h24-reg-applied-30]{30}, \hyperref[q:h24-reg-applied-31]{31}, \hyperref[q:h24-reg-applied-32]{32}, \hyperref[q:h24-reg-applied-33]{33}, \hyperref[q:h24-reg-applied-34]{34}, \hyperref[q:h24-reg-applied-35]{35}
\section{平成25年度}
\paragraph{基礎（30問）} \hyperref[q:h25-reg-basic-1-1]{1-1}, \hyperref[q:h25-reg-basic-1-2]{1-2}, \hyperref[q:h25-reg-basic-1-3]{1-3}, \hyperref[q:h25-reg-basic-1-4]{1-4}, \hyperref[q:h25-reg-basic-1-5]{1-5}, \hyperref[q:h25-reg-basic-1-6]{1-6}, \hyperref[q:h25-reg-basic-2-1]{2-1}, \hyperref[q:h25-reg-basic-2-2]{2-2}, \hyperref[q:h25-reg-basic-2-3]{2-3}, \hyperref[q:h25-reg-basic-2-4]{2-4}, \hyperref[q:h25-reg-basic-2-5]{2-5}, \hyperref[q:h25-reg-basic-2-6]{2-6}, \hyperref[q:h25-reg-basic-3-1]{3-1}, \hyperref[q:h25-reg-basic-3-2]{3-2}, \hyperref[q:h25-reg-basic-3-3]{3-3}, \hyperref[q:h25-reg-basic-3-4]{3-4}, \hyperref[q:h25-reg-basic-3-5]{3-5}, \hyperref[q:h25-reg-basic-3-6]{3-6}, \hyperref[q:h25-reg-basic-4-1]{4-1}, \hyperref[q:h25-reg-basic-4-2]{4-2}, \hyperref[q:h25-reg-basic-4-3]{4-3}, \hyperref[q:h25-reg-basic-4-4]{4-4}, \hyperref[q:h25-reg-basic-4-5]{4-5}, \hyperref[q:h25-reg-basic-4-6]{4-6}, \hyperref[q:h25-reg-basic-5-1]{5-1}, \hyperref[q:h25-reg-basic-5-2]{5-2}, \hyperref[q:h25-reg-basic-5-3]{5-3}, \hyperref[q:h25-reg-basic-5-4]{5-4}, \hyperref[q:h25-reg-basic-5-5]{5-5}, \hyperref[q:h25-reg-basic-5-6]{5-6}
\paragraph{適性（15問）} \hyperref[q:h25-reg-apt-01]{1}, \hyperref[q:h25-reg-apt-02]{2}, \hyperref[q:h25-reg-apt-03]{3}, \hyperref[q:h25-reg-apt-04]{4}, \hyperref[q:h25-reg-apt-05]{5}, \hyperref[q:h25-reg-apt-06]{6}, \hyperref[q:h25-reg-apt-07]{7}, \hyperref[q:h25-reg-apt-08]{8}, \hyperref[q:h25-reg-apt-09]{9}, \hyperref[q:h25-reg-apt-10]{10}, \hyperref[q:h25-reg-apt-11]{11}, \hyperref[q:h25-reg-apt-12]{12}, \hyperref[q:h25-reg-apt-13]{13}, \hyperref[q:h25-reg-apt-14]{14}, \hyperref[q:h25-reg-apt-15]{15}
\paragraph{応用理学（35問）} \hyperref[q:h25-reg-applied-01]{1}, \hyperref[q:h25-reg-applied-02]{2}, \hyperref[q:h25-reg-applied-03]{3}, \hyperref[q:h25-reg-applied-04]{4}, \hyperref[q:h25-reg-applied-05]{5}, \hyperref[q:h25-reg-applied-06]{6}, \hyperref[q:h25-reg-applied-07]{7}, \hyperref[q:h25-reg-applied-08]{8}, \hyperref[q:h25-reg-applied-09]{9}, \hyperref[q:h25-reg-applied-10]{10}, \hyperref[q:h25-reg-applied-11]{11}, \hyperref[q:h25-reg-applied-12]{12}, \hyperref[q:h25-reg-applied-13]{13}, \hyperref[q:h25-reg-applied-14]{14}, \hyperref[q:h25-reg-applied-15]{15}, \hyperref[q:h25-reg-applied-16]{16}, \hyperref[q:h25-reg-applied-17]{17}, \hyperref[q:h25-reg-applied-18]{18}, \hyperref[q:h25-reg-applied-19]{19}, \hyperref[q:h25-reg-applied-20]{20}, \hyperref[q:h25-reg-applied-21]{21}, \hyperref[q:h25-reg-applied-22]{22}, \hyperref[q:h25-reg-applied-23]{23}, \hyperref[q:h25-reg-applied-24]{24}, \hyperref[q:h25-reg-applied-25]{25}, \hyperref[q:h25-reg-applied-26]{26}, \hyperref[q:h25-reg-applied-27]{27}, \hyperref[q:h25-reg-applied-28]{28}, \hyperref[q:h25-reg-applied-29]{29}, \hyperref[q:h25-reg-applied-30]{30}, \hyperref[q:h25-reg-applied-31]{31}, \hyperref[q:h25-reg-applied-32]{32}, \hyperref[q:h25-reg-applied-33]{33}, \hyperref[q:h25-reg-applied-34]{34}, \hyperref[q:h25-reg-applied-35]{35}
\section{平成26年度}
\paragraph{基礎（30問）} \hyperref[q:h26-reg-basic-1-1]{1-1}, \hyperref[q:h26-reg-basic-1-2]{1-2}, \hyperref[q:h26-reg-basic-1-3]{1-3}, \hyperref[q:h26-reg-basic-1-4]{1-4}, \hyperref[q:h26-reg-basic-1-5]{1-5}, \hyperref[q:h26-reg-basic-1-6]{1-6}, \hyperref[q:h26-reg-basic-2-1]{2-1}, \hyperref[q:h26-reg-basic-2-2]{2-2}, \hyperref[q:h26-reg-basic-2-3]{2-3}, \hyperref[q:h26-reg-basic-2-4]{2-4}, \hyperref[q:h26-reg-basic-2-5]{2-5}, \hyperref[q:h26-reg-basic-2-6]{2-6}, \hyperref[q:h26-reg-basic-3-1]{3-1}, \hyperref[q:h26-reg-basic-3-2]{3-2}, \hyperref[q:h26-reg-basic-3-3]{3-3}, \hyperref[q:h26-reg-basic-3-4]{3-4}, \hyperref[q:h26-reg-basic-3-5]{3-5}, \hyperref[q:h26-reg-basic-3-6]{3-6}, \hyperref[q:h26-reg-basic-4-1]{4-1}, \hyperref[q:h26-reg-basic-4-2]{4-2}, \hyperref[q:h26-reg-basic-4-3]{4-3}, \hyperref[q:h26-reg-basic-4-4]{4-4}, \hyperref[q:h26-reg-basic-4-5]{4-5}, \hyperref[q:h26-reg-basic-4-6]{4-6}, \hyperref[q:h26-reg-basic-5-1]{5-1}, \hyperref[q:h26-reg-basic-5-2]{5-2}, \hyperref[q:h26-reg-basic-5-3]{5-3}, \hyperref[q:h26-reg-basic-5-4]{5-4}, \hyperref[q:h26-reg-basic-5-5]{5-5}, \hyperref[q:h26-reg-basic-5-6]{5-6}
\paragraph{適性（15問）} \hyperref[q:h26-reg-apt-01]{1}, \hyperref[q:h26-reg-apt-02]{2}, \hyperref[q:h26-reg-apt-03]{3}, \hyperref[q:h26-reg-apt-04]{4}, \hyperref[q:h26-reg-apt-05]{5}, \hyperref[q:h26-reg-apt-06]{6}, \hyperref[q:h26-reg-apt-07]{7}, \hyperref[q:h26-reg-apt-08]{8}, \hyperref[q:h26-reg-apt-09]{9}, \hyperref[q:h26-reg-apt-10]{10}, \hyperref[q:h26-reg-apt-11]{11}, \hyperref[q:h26-reg-apt-12]{12}, \hyperref[q:h26-reg-apt-13]{13}, \hyperref[q:h26-reg-apt-14]{14}, \hyperref[q:h26-reg-apt-15]{15}
\paragraph{応用理学（35問）} \hyperref[q:h26-reg-applied-01]{1}, \hyperref[q:h26-reg-applied-02]{2}, \hyperref[q:h26-reg-applied-03]{3}, \hyperref[q:h26-reg-applied-04]{4}, \hyperref[q:h26-reg-applied-05]{5}, \hyperref[q:h26-reg-applied-06]{6}, \hyperref[q:h26-reg-applied-07]{7}, \hyperref[q:h26-reg-applied-08]{8}, \hyperref[q:h26-reg-applied-09]{9}, \hyperref[q:h26-reg-applied-10]{10}, \hyperref[q:h26-reg-applied-11]{11}, \hyperref[q:h26-reg-applied-12]{12}, \hyperref[q:h26-reg-applied-13]{13}, \hyperref[q:h26-reg-applied-14]{14}, \hyperref[q:h26-reg-applied-15]{15}, \hyperref[q:h26-reg-applied-16]{16}, \hyperref[q:h26-reg-applied-17]{17}, \hyperref[q:h26-reg-applied-18]{18}, \hyperref[q:h26-reg-applied-19]{19}, \hyperref[q:h26-reg-applied-20]{20}, \hyperref[q:h26-reg-applied-21]{21}, \hyperref[q:h26-reg-applied-22]{22}, \hyperref[q:h26-reg-applied-23]{23}, \hyperref[q:h26-reg-applied-24]{24}, \hyperref[q:h26-reg-applied-25]{25}, \hyperref[q:h26-reg-applied-26]{26}, \hyperref[q:h26-reg-applied-27]{27}, \hyperref[q:h26-reg-applied-28]{28}, \hyperref[q:h26-reg-applied-29]{29}, \hyperref[q:h26-reg-applied-30]{30}, \hyperref[q:h26-reg-applied-31]{31}, \hyperref[q:h26-reg-applied-32]{32}, \hyperref[q:h26-reg-applied-33]{33}, \hyperref[q:h26-reg-applied-34]{34}, \hyperref[q:h26-reg-applied-35]{35}
\section{平成27年度}
\paragraph{基礎（30問）} \hyperref[q:h27-reg-basic-1-1]{1-1}, \hyperref[q:h27-reg-basic-1-2]{1-2}, \hyperref[q:h27-reg-basic-1-3]{1-3}, \hyperref[q:h27-reg-basic-1-4]{1-4}, \hyperref[q:h27-reg-basic-1-5]{1-5}, \hyperref[q:h27-reg-basic-1-6]{1-6}, \hyperref[q:h27-reg-basic-2-1]{2-1}, \hyperref[q:h27-reg-basic-2-2]{2-2}, \hyperref[q:h27-reg-basic-2-3]{2-3}, \hyperref[q:h27-reg-basic-2-4]{2-4}, \hyperref[q:h27-reg-basic-2-5]{2-5}, \hyperref[q:h27-reg-basic-2-6]{2-6}, \hyperref[q:h27-reg-basic-3-1]{3-1}, \hyperref[q:h27-reg-basic-3-2]{3-2}, \hyperref[q:h27-reg-basic-3-3]{3-3}, \hyperref[q:h27-reg-basic-3-4]{3-4}, \hyperref[q:h27-reg-basic-3-5]{3-5}, \hyperref[q:h27-reg-basic-3-6]{3-6}, \hyperref[q:h27-reg-basic-4-1]{4-1}, \hyperref[q:h27-reg-basic-4-2]{4-2}, \hyperref[q:h27-reg-basic-4-3]{4-3}, \hyperref[q:h27-reg-basic-4-4]{4-4}, \hyperref[q:h27-reg-basic-4-5]{4-5}, \hyperref[q:h27-reg-basic-4-6]{4-6}, \hyperref[q:h27-reg-basic-5-1]{5-1}, \hyperref[q:h27-reg-basic-5-2]{5-2}, \hyperref[q:h27-reg-basic-5-3]{5-3}, \hyperref[q:h27-reg-basic-5-4]{5-4}, \hyperref[q:h27-reg-basic-5-5]{5-5}, \hyperref[q:h27-reg-basic-5-6]{5-6}
\paragraph{適性（15問）} \hyperref[q:h27-reg-apt-01]{1}, \hyperref[q:h27-reg-apt-02]{2}, \hyperref[q:h27-reg-apt-03]{3}, \hyperref[q:h27-reg-apt-04]{4}, \hyperref[q:h27-reg-apt-05]{5}, \hyperref[q:h27-reg-apt-06]{6}, \hyperref[q:h27-reg-apt-07]{7}, \hyperref[q:h27-reg-apt-08]{8}, \hyperref[q:h27-reg-apt-09]{9}, \hyperref[q:h27-reg-apt-10]{10}, \hyperref[q:h27-reg-apt-11]{11}, \hyperref[q:h27-reg-apt-12]{12}, \hyperref[q:h27-reg-apt-13]{13}, \hyperref[q:h27-reg-apt-14]{14}, \hyperref[q:h27-reg-apt-15]{15}
\paragraph{応用理学（35問）} \hyperref[q:h27-reg-applied-01]{1}, \hyperref[q:h27-reg-applied-02]{2}, \hyperref[q:h27-reg-applied-03]{3}, \hyperref[q:h27-reg-applied-04]{4}, \hyperref[q:h27-reg-applied-05]{5}, \hyperref[q:h27-reg-applied-06]{6}, \hyperref[q:h27-reg-applied-07]{7}, \hyperref[q:h27-reg-applied-08]{8}, \hyperref[q:h27-reg-applied-09]{9}, \hyperref[q:h27-reg-applied-10]{10}, \hyperref[q:h27-reg-applied-11]{11}, \hyperref[q:h27-reg-applied-12]{12}, \hyperref[q:h27-reg-applied-13]{13}, \hyperref[q:h27-reg-applied-14]{14}, \hyperref[q:h27-reg-applied-15]{15}, \hyperref[q:h27-reg-applied-16]{16}, \hyperref[q:h27-reg-applied-17]{17}, \hyperref[q:h27-reg-applied-18]{18}, \hyperref[q:h27-reg-applied-19]{19}, \hyperref[q:h27-reg-applied-20]{20}, \hyperref[q:h27-reg-applied-21]{21}, \hyperref[q:h27-reg-applied-22]{22}, \hyperref[q:h27-reg-applied-23]{23}, \hyperref[q:h27-reg-applied-24]{24}, \hyperref[q:h27-reg-applied-25]{25}, \hyperref[q:h27-reg-applied-26]{26}, \hyperref[q:h27-reg-applied-27]{27}, \hyperref[q:h27-reg-applied-28]{28}, \hyperref[q:h27-reg-applied-29]{29}, \hyperref[q:h27-reg-applied-30]{30}, \hyperref[q:h27-reg-applied-31]{31}, \hyperref[q:h27-reg-applied-32]{32}, \hyperref[q:h27-reg-applied-33]{33}, \hyperref[q:h27-reg-applied-34]{34}, \hyperref[q:h27-reg-applied-35]{35}
\section{平成28年度}
\paragraph{基礎（30問）} \hyperref[q:h28-reg-basic-1-1]{1-1}, \hyperref[q:h28-reg-basic-1-2]{1-2}, \hyperref[q:h28-reg-basic-1-3]{1-3}, \hyperref[q:h28-reg-basic-1-4]{1-4}, \hyperref[q:h28-reg-basic-1-5]{1-5}, \hyperref[q:h28-reg-basic-1-6]{1-6}, \hyperref[q:h28-reg-basic-2-1]{2-1}, \hyperref[q:h28-reg-basic-2-2]{2-2}, \hyperref[q:h28-reg-basic-2-3]{2-3}, \hyperref[q:h28-reg-basic-2-4]{2-4}, \hyperref[q:h28-reg-basic-2-5]{2-5}, \hyperref[q:h28-reg-basic-2-6]{2-6}, \hyperref[q:h28-reg-basic-3-1]{3-1}, \hyperref[q:h28-reg-basic-3-2]{3-2}, \hyperref[q:h28-reg-basic-3-3]{3-3}, \hyperref[q:h28-reg-basic-3-4]{3-4}, \hyperref[q:h28-reg-basic-3-5]{3-5}, \hyperref[q:h28-reg-basic-3-6]{3-6}, \hyperref[q:h28-reg-basic-4-1]{4-1}, \hyperref[q:h28-reg-basic-4-2]{4-2}, \hyperref[q:h28-reg-basic-4-3]{4-3}, \hyperref[q:h28-reg-basic-4-4]{4-4}, \hyperref[q:h28-reg-basic-4-5]{4-5}, \hyperref[q:h28-reg-basic-4-6]{4-6}, \hyperref[q:h28-reg-basic-5-1]{5-1}, \hyperref[q:h28-reg-basic-5-2]{5-2}, \hyperref[q:h28-reg-basic-5-3]{5-3}, \hyperref[q:h28-reg-basic-5-4]{5-4}, \hyperref[q:h28-reg-basic-5-5]{5-5}, \hyperref[q:h28-reg-basic-5-6]{5-6}
\paragraph{適性（15問）} \hyperref[q:h28-reg-apt-01]{1}, \hyperref[q:h28-reg-apt-02]{2}, \hyperref[q:h28-reg-apt-03]{3}, \hyperref[q:h28-reg-apt-04]{4}, \hyperref[q:h28-reg-apt-05]{5}, \hyperref[q:h28-reg-apt-06]{6}, \hyperref[q:h28-reg-apt-07]{7}, \hyperref[q:h28-reg-apt-08]{8}, \hyperref[q:h28-reg-apt-09]{9}, \hyperref[q:h28-reg-apt-10]{10}, \hyperref[q:h28-reg-apt-11]{11}, \hyperref[q:h28-reg-apt-12]{12}, \hyperref[q:h28-reg-apt-13]{13}, \hyperref[q:h28-reg-apt-14]{14}, \hyperref[q:h28-reg-apt-15]{15}
\paragraph{応用理学（35問）} \hyperref[q:h28-reg-applied-01]{1}, \hyperref[q:h28-reg-applied-02]{2}, \hyperref[q:h28-reg-applied-03]{3}, \hyperref[q:h28-reg-applied-04]{4}, \hyperref[q:h28-reg-applied-05]{5}, \hyperref[q:h28-reg-applied-06]{6}, \hyperref[q:h28-reg-applied-07]{7}, \hyperref[q:h28-reg-applied-08]{8}, \hyperref[q:h28-reg-applied-09]{9}, \hyperref[q:h28-reg-applied-10]{10}, \hyperref[q:h28-reg-applied-11]{11}, \hyperref[q:h28-reg-applied-12]{12}, \hyperref[q:h28-reg-applied-13]{13}, \hyperref[q:h28-reg-applied-14]{14}, \hyperref[q:h28-reg-applied-15]{15}, \hyperref[q:h28-reg-applied-16]{16}, \hyperref[q:h28-reg-applied-17]{17}, \hyperref[q:h28-reg-applied-18]{18}, \hyperref[q:h28-reg-applied-19]{19}, \hyperref[q:h28-reg-applied-20]{20}, \hyperref[q:h28-reg-applied-21]{21}, \hyperref[q:h28-reg-applied-22]{22}, \hyperref[q:h28-reg-applied-23]{23}, \hyperref[q:h28-reg-applied-24]{24}, \hyperref[q:h28-reg-applied-25]{25}, \hyperref[q:h28-reg-applied-26]{26}, \hyperref[q:h28-reg-applied-27]{27}, \hyperref[q:h28-reg-applied-28]{28}, \hyperref[q:h28-reg-applied-29]{29}, \hyperref[q:h28-reg-applied-30]{30}, \hyperref[q:h28-reg-applied-31]{31}, \hyperref[q:h28-reg-applied-32]{32}, \hyperref[q:h28-reg-applied-33]{33}, \hyperref[q:h28-reg-applied-34]{34}, \hyperref[q:h28-reg-applied-35]{35}
\section{平成29年度}
\paragraph{基礎（30問）} \hyperref[q:h29-reg-basic-1-1]{1-1}, \hyperref[q:h29-reg-basic-1-2]{1-2}, \hyperref[q:h29-reg-basic-1-3]{1-3}, \hyperref[q:h29-reg-basic-1-4]{1-4}, \hyperref[q:h29-reg-basic-1-5]{1-5}, \hyperref[q:h29-reg-basic-1-6]{1-6}, \hyperref[q:h29-reg-basic-2-1]{2-1}, \hyperref[q:h29-reg-basic-2-2]{2-2}, \hyperref[q:h29-reg-basic-2-3]{2-3}, \hyperref[q:h29-reg-basic-2-4]{2-4}, \hyperref[q:h29-reg-basic-2-5]{2-5}, \hyperref[q:h29-reg-basic-2-6]{2-6}, \hyperref[q:h29-reg-basic-3-1]{3-1}, \hyperref[q:h29-reg-basic-3-2]{3-2}, \hyperref[q:h29-reg-basic-3-3]{3-3}, \hyperref[q:h29-reg-basic-3-4]{3-4}, \hyperref[q:h29-reg-basic-3-5]{3-5}, \hyperref[q:h29-reg-basic-3-6]{3-6}, \hyperref[q:h29-reg-basic-4-1]{4-1}, \hyperref[q:h29-reg-basic-4-2]{4-2}, \hyperref[q:h29-reg-basic-4-3]{4-3}, \hyperref[q:h29-reg-basic-4-4]{4-4}, \hyperref[q:h29-reg-basic-4-5]{4-5}, \hyperref[q:h29-reg-basic-4-6]{4-6}, \hyperref[q:h29-reg-basic-5-1]{5-1}, \hyperref[q:h29-reg-basic-5-2]{5-2}, \hyperref[q:h29-reg-basic-5-3]{5-3}, \hyperref[q:h29-reg-basic-5-4]{5-4}, \hyperref[q:h29-reg-basic-5-5]{5-5}, \hyperref[q:h29-reg-basic-5-6]{5-6}
\paragraph{適性（15問）} \hyperref[q:h29-reg-apt-01]{1}, \hyperref[q:h29-reg-apt-02]{2}, \hyperref[q:h29-reg-apt-03]{3}, \hyperref[q:h29-reg-apt-04]{4}, \hyperref[q:h29-reg-apt-05]{5}, \hyperref[q:h29-reg-apt-06]{6}, \hyperref[q:h29-reg-apt-07]{7}, \hyperref[q:h29-reg-apt-08]{8}, \hyperref[q:h29-reg-apt-09]{9}, \hyperref[q:h29-reg-apt-10]{10}, \hyperref[q:h29-reg-apt-11]{11}, \hyperref[q:h29-reg-apt-12]{12}, \hyperref[q:h29-reg-apt-13]{13}, \hyperref[q:h29-reg-apt-14]{14}, \hyperref[q:h29-reg-apt-15]{15}
\paragraph{応用理学（35問）} \hyperref[q:h29-reg-applied-01]{1}, \hyperref[q:h29-reg-applied-02]{2}, \hyperref[q:h29-reg-applied-03]{3}, \hyperref[q:h29-reg-applied-04]{4}, \hyperref[q:h29-reg-applied-05]{5}, \hyperref[q:h29-reg-applied-06]{6}, \hyperref[q:h29-reg-applied-07]{7}, \hyperref[q:h29-reg-applied-08]{8}, \hyperref[q:h29-reg-applied-09]{9}, \hyperref[q:h29-reg-applied-10]{10}, \hyperref[q:h29-reg-applied-11]{11}, \hyperref[q:h29-reg-applied-12]{12}, \hyperref[q:h29-reg-applied-13]{13}, \hyperref[q:h29-reg-applied-14]{14}, \hyperref[q:h29-reg-applied-15]{15}, \hyperref[q:h29-reg-applied-16]{16}, \hyperref[q:h29-reg-applied-17]{17}, \hyperref[q:h29-reg-applied-18]{18}, \hyperref[q:h29-reg-applied-19]{19}, \hyperref[q:h29-reg-applied-20]{20}, \hyperref[q:h29-reg-applied-21]{21}, \hyperref[q:h29-reg-applied-22]{22}, \hyperref[q:h29-reg-applied-23]{23}, \hyperref[q:h29-reg-applied-24]{24}, \hyperref[q:h29-reg-applied-25]{25}, \hyperref[q:h29-reg-applied-26]{26}, \hyperref[q:h29-reg-applied-27]{27}, \hyperref[q:h29-reg-applied-28]{28}, \hyperref[q:h29-reg-applied-29]{29}, \hyperref[q:h29-reg-applied-30]{30}, \hyperref[q:h29-reg-applied-31]{31}, \hyperref[q:h29-reg-applied-32]{32}, \hyperref[q:h29-reg-applied-33]{33}, \hyperref[q:h29-reg-applied-34]{34}, \hyperref[q:h29-reg-applied-35]{35}
\section{平成30年度}
\paragraph{基礎（30問）} \hyperref[q:h30-reg-basic-1-1]{1-1}, \hyperref[q:h30-reg-basic-1-2]{1-2}, \hyperref[q:h30-reg-basic-1-3]{1-3}, \hyperref[q:h30-reg-basic-1-4]{1-4}, \hyperref[q:h30-reg-basic-1-5]{1-5}, \hyperref[q:h30-reg-basic-1-6]{1-6}, \hyperref[q:h30-reg-basic-2-1]{2-1}, \hyperref[q:h30-reg-basic-2-2]{2-2}, \hyperref[q:h30-reg-basic-2-3]{2-3}, \hyperref[q:h30-reg-basic-2-4]{2-4}, \hyperref[q:h30-reg-basic-2-5]{2-5}, \hyperref[q:h30-reg-basic-2-6]{2-6}, \hyperref[q:h30-reg-basic-3-1]{3-1}, \hyperref[q:h30-reg-basic-3-2]{3-2}, \hyperref[q:h30-reg-basic-3-3]{3-3}, \hyperref[q:h30-reg-basic-3-4]{3-4}, \hyperref[q:h30-reg-basic-3-5]{3-5}, \hyperref[q:h30-reg-basic-3-6]{3-6}, \hyperref[q:h30-reg-basic-4-1]{4-1}, \hyperref[q:h30-reg-basic-4-2]{4-2}, \hyperref[q:h30-reg-basic-4-3]{4-3}, \hyperref[q:h30-reg-basic-4-4]{4-4}, \hyperref[q:h30-reg-basic-4-5]{4-5}, \hyperref[q:h30-reg-basic-4-6]{4-6}, \hyperref[q:h30-reg-basic-5-1]{5-1}, \hyperref[q:h30-reg-basic-5-2]{5-2}, \hyperref[q:h30-reg-basic-5-3]{5-3}, \hyperref[q:h30-reg-basic-5-4]{5-4}, \hyperref[q:h30-reg-basic-5-5]{5-5}, \hyperref[q:h30-reg-basic-5-6]{5-6}
\paragraph{適性（15問）} \hyperref[q:h30-reg-apt-01]{1}, \hyperref[q:h30-reg-apt-02]{2}, \hyperref[q:h30-reg-apt-03]{3}, \hyperref[q:h30-reg-apt-04]{4}, \hyperref[q:h30-reg-apt-05]{5}, \hyperref[q:h30-reg-apt-06]{6}, \hyperref[q:h30-reg-apt-07]{7}, \hyperref[q:h30-reg-apt-08]{8}, \hyperref[q:h30-reg-apt-09]{9}, \hyperref[q:h30-reg-apt-10]{10}, \hyperref[q:h30-reg-apt-11]{11}, \hyperref[q:h30-reg-apt-12]{12}, \hyperref[q:h30-reg-apt-13]{13}, \hyperref[q:h30-reg-apt-14]{14}, \hyperref[q:h30-reg-apt-15]{15}
\paragraph{応用理学（35問）} \hyperref[q:h30-reg-applied-01]{1}, \hyperref[q:h30-reg-applied-02]{2}, \hyperref[q:h30-reg-applied-03]{3}, \hyperref[q:h30-reg-applied-04]{4}, \hyperref[q:h30-reg-applied-05]{5}, \hyperref[q:h30-reg-applied-06]{6}, \hyperref[q:h30-reg-applied-07]{7}, \hyperref[q:h30-reg-applied-08]{8}, \hyperref[q:h30-reg-applied-09]{9}, \hyperref[q:h30-reg-applied-10]{10}, \hyperref[q:h30-reg-applied-11]{11}, \hyperref[q:h30-reg-applied-12]{12}, \hyperref[q:h30-reg-applied-13]{13}, \hyperref[q:h30-reg-applied-14]{14}, \hyperref[q:h30-reg-applied-15]{15}, \hyperref[q:h30-reg-applied-16]{16}, \hyperref[q:h30-reg-applied-17]{17}, \hyperref[q:h30-reg-applied-18]{18}, \hyperref[q:h30-reg-applied-19]{19}, \hyperref[q:h30-reg-applied-20]{20}, \hyperref[q:h30-reg-applied-21]{21}, \hyperref[q:h30-reg-applied-22]{22}, \hyperref[q:h30-reg-applied-23]{23}, \hyperref[q:h30-reg-applied-24]{24}, \hyperref[q:h30-reg-applied-25]{25}, \hyperref[q:h30-reg-applied-26]{26}, \hyperref[q:h30-reg-applied-27]{27}, \hyperref[q:h30-reg-applied-28]{28}, \hyperref[q:h30-reg-applied-29]{29}, \hyperref[q:h30-reg-applied-30]{30}, \hyperref[q:h30-reg-applied-31]{31}, \hyperref[q:h30-reg-applied-32]{32}, \hyperref[q:h30-reg-applied-33]{33}, \hyperref[q:h30-reg-applied-34]{34}, \hyperref[q:h30-reg-applied-35]{35}
\section{令和元年度}
\paragraph{基礎（30問）} \hyperref[q:r01-reg-basic-1-1]{1-1}, \hyperref[q:r01-reg-basic-1-2]{1-2}, \hyperref[q:r01-reg-basic-1-3]{1-3}, \hyperref[q:r01-reg-basic-1-4]{1-4}, \hyperref[q:r01-reg-basic-1-5]{1-5}, \hyperref[q:r01-reg-basic-1-6]{1-6}, \hyperref[q:r01-reg-basic-2-1]{2-1}, \hyperref[q:r01-reg-basic-2-2]{2-2}, \hyperref[q:r01-reg-basic-2-3]{2-3}, \hyperref[q:r01-reg-basic-2-4]{2-4}, \hyperref[q:r01-reg-basic-2-5]{2-5}, \hyperref[q:r01-reg-basic-2-6]{2-6}, \hyperref[q:r01-reg-basic-3-1]{3-1}, \hyperref[q:r01-reg-basic-3-2]{3-2}, \hyperref[q:r01-reg-basic-3-3]{3-3}, \hyperref[q:r01-reg-basic-3-4]{3-4}, \hyperref[q:r01-reg-basic-3-5]{3-5}, \hyperref[q:r01-reg-basic-3-6]{3-6}, \hyperref[q:r01-reg-basic-4-1]{4-1}, \hyperref[q:r01-reg-basic-4-2]{4-2}, \hyperref[q:r01-reg-basic-4-3]{4-3}, \hyperref[q:r01-reg-basic-4-4]{4-4}, \hyperref[q:r01-reg-basic-4-5]{4-5}, \hyperref[q:r01-reg-basic-4-6]{4-6}, \hyperref[q:r01-reg-basic-5-1]{5-1}, \hyperref[q:r01-reg-basic-5-2]{5-2}, \hyperref[q:r01-reg-basic-5-3]{5-3}, \hyperref[q:r01-reg-basic-5-4]{5-4}, \hyperref[q:r01-reg-basic-5-5]{5-5}, \hyperref[q:r01-reg-basic-5-6]{5-6}
\paragraph{適性（15問）} \hyperref[q:r01-reg-apt-01]{1}, \hyperref[q:r01-reg-apt-02]{2}, \hyperref[q:r01-reg-apt-03]{3}, \hyperref[q:r01-reg-apt-04]{4}, \hyperref[q:r01-reg-apt-05]{5}, \hyperref[q:r01-reg-apt-06]{6}, \hyperref[q:r01-reg-apt-07]{7}, \hyperref[q:r01-reg-apt-08]{8}, \hyperref[q:r01-reg-apt-09]{9}, \hyperref[q:r01-reg-apt-10]{10}, \hyperref[q:r01-reg-apt-11]{11}, \hyperref[q:r01-reg-apt-12]{12}, \hyperref[q:r01-reg-apt-13]{13}, \hyperref[q:r01-reg-apt-14]{14}, \hyperref[q:r01-reg-apt-15]{15}
\paragraph{応用理学（35問）} \hyperref[q:r01-reg-applied-01]{1}, \hyperref[q:r01-reg-applied-02]{2}, \hyperref[q:r01-reg-applied-03]{3}, \hyperref[q:r01-reg-applied-04]{4}, \hyperref[q:r01-reg-applied-05]{5}, \hyperref[q:r01-reg-applied-06]{6}, \hyperref[q:r01-reg-applied-07]{7}, \hyperref[q:r01-reg-applied-08]{8}, \hyperref[q:r01-reg-applied-09]{9}, \hyperref[q:r01-reg-applied-10]{10}, \hyperref[q:r01-reg-applied-11]{11}, \hyperref[q:r01-reg-applied-12]{12}, \hyperref[q:r01-reg-applied-13]{13}, \hyperref[q:r01-reg-applied-14]{14}, \hyperref[q:r01-reg-applied-15]{15}, \hyperref[q:r01-reg-applied-16]{16}, \hyperref[q:r01-reg-applied-17]{17}, \hyperref[q:r01-reg-applied-18]{18}, \hyperref[q:r01-reg-applied-19]{19}, \hyperref[q:r01-reg-applied-20]{20}, \hyperref[q:r01-reg-applied-21]{21}, \hyperref[q:r01-reg-applied-22]{22}, \hyperref[q:r01-reg-applied-23]{23}, \hyperref[q:r01-reg-applied-24]{24}, \hyperref[q:r01-reg-applied-25]{25}, \hyperref[q:r01-reg-applied-26]{26}, \hyperref[q:r01-reg-applied-27]{27}, \hyperref[q:r01-reg-applied-28]{28}, \hyperref[q:r01-reg-applied-29]{29}, \hyperref[q:r01-reg-applied-30]{30}, \hyperref[q:r01-reg-applied-31]{31}, \hyperref[q:r01-reg-applied-32]{32}, \hyperref[q:r01-reg-applied-33]{33}, \hyperref[q:r01-reg-applied-34]{34}, \hyperref[q:r01-reg-applied-35]{35}
\section{令和元年度再試験}
\paragraph{基礎（30問）} \hyperref[q:r01-ret-basic-1-1]{1-1}, \hyperref[q:r01-ret-basic-1-2]{1-2}, \hyperref[q:r01-ret-basic-1-3]{1-3}, \hyperref[q:r01-ret-basic-1-4]{1-4}, \hyperref[q:r01-ret-basic-1-5]{1-5}, \hyperref[q:r01-ret-basic-1-6]{1-6}, \hyperref[q:r01-ret-basic-2-1]{2-1}, \hyperref[q:r01-ret-basic-2-2]{2-2}, \hyperref[q:r01-ret-basic-2-3]{2-3}, \hyperref[q:r01-ret-basic-2-4]{2-4}, \hyperref[q:r01-ret-basic-2-5]{2-5}, \hyperref[q:r01-ret-basic-2-6]{2-6}, \hyperref[q:r01-ret-basic-3-1]{3-1}, \hyperref[q:r01-ret-basic-3-2]{3-2}, \hyperref[q:r01-ret-basic-3-3]{3-3}, \hyperref[q:r01-ret-basic-3-4]{3-4}, \hyperref[q:r01-ret-basic-3-5]{3-5}, \hyperref[q:r01-ret-basic-3-6]{3-6}, \hyperref[q:r01-ret-basic-4-1]{4-1}, \hyperref[q:r01-ret-basic-4-2]{4-2}, \hyperref[q:r01-ret-basic-4-3]{4-3}, \hyperref[q:r01-ret-basic-4-4]{4-4}, \hyperref[q:r01-ret-basic-4-5]{4-5}, \hyperref[q:r01-ret-basic-4-6]{4-6}, \hyperref[q:r01-ret-basic-5-1]{5-1}, \hyperref[q:r01-ret-basic-5-2]{5-2}, \hyperref[q:r01-ret-basic-5-3]{5-3}, \hyperref[q:r01-ret-basic-5-4]{5-4}, \hyperref[q:r01-ret-basic-5-5]{5-5}, \hyperref[q:r01-ret-basic-5-6]{5-6}
\paragraph{適性（15問）} \hyperref[q:r01-ret-apt-01]{1}, \hyperref[q:r01-ret-apt-02]{2}, \hyperref[q:r01-ret-apt-03]{3}, \hyperref[q:r01-ret-apt-04]{4}, \hyperref[q:r01-ret-apt-05]{5}, \hyperref[q:r01-ret-apt-06]{6}, \hyperref[q:r01-ret-apt-07]{7}, \hyperref[q:r01-ret-apt-08]{8}, \hyperref[q:r01-ret-apt-09]{9}, \hyperref[q:r01-ret-apt-10]{10}, \hyperref[q:r01-ret-apt-11]{11}, \hyperref[q:r01-ret-apt-12]{12}, \hyperref[q:r01-ret-apt-13]{13}, \hyperref[q:r01-ret-apt-14]{14}, \hyperref[q:r01-ret-apt-15]{15}
\paragraph{応用理学（35問）} \hyperref[q:r01-ret-applied-01]{1}, \hyperref[q:r01-ret-applied-02]{2}, \hyperref[q:r01-ret-applied-03]{3}, \hyperref[q:r01-ret-applied-04]{4}, \hyperref[q:r01-ret-applied-05]{5}, \hyperref[q:r01-ret-applied-06]{6}, \hyperref[q:r01-ret-applied-07]{7}, \hyperref[q:r01-ret-applied-08]{8}, \hyperref[q:r01-ret-applied-09]{9}, \hyperref[q:r01-ret-applied-10]{10}, \hyperref[q:r01-ret-applied-11]{11}, \hyperref[q:r01-ret-applied-12]{12}, \hyperref[q:r01-ret-applied-13]{13}, \hyperref[q:r01-ret-applied-14]{14}, \hyperref[q:r01-ret-applied-15]{15}, \hyperref[q:r01-ret-applied-16]{16}, \hyperref[q:r01-ret-applied-17]{17}, \hyperref[q:r01-ret-applied-18]{18}, \hyperref[q:r01-ret-applied-19]{19}, \hyperref[q:r01-ret-applied-20]{20}, \hyperref[q:r01-ret-applied-21]{21}, \hyperref[q:r01-ret-applied-22]{22}, \hyperref[q:r01-ret-applied-23]{23}, \hyperref[q:r01-ret-applied-24]{24}, \hyperref[q:r01-ret-applied-25]{25}, \hyperref[q:r01-ret-applied-26]{26}, \hyperref[q:r01-ret-applied-27]{27}, \hyperref[q:r01-ret-applied-28]{28}, \hyperref[q:r01-ret-applied-29]{29}, \hyperref[q:r01-ret-applied-30]{30}, \hyperref[q:r01-ret-applied-31]{31}, \hyperref[q:r01-ret-applied-32]{32}, \hyperref[q:r01-ret-applied-33]{33}, \hyperref[q:r01-ret-applied-34]{34}, \hyperref[q:r01-ret-applied-35]{35}
\section{令和2年度}
\paragraph{基礎（30問）} \hyperref[q:r02-reg-basic-1-1]{1-1}, \hyperref[q:r02-reg-basic-1-2]{1-2}, \hyperref[q:r02-reg-basic-1-3]{1-3}, \hyperref[q:r02-reg-basic-1-4]{1-4}, \hyperref[q:r02-reg-basic-1-5]{1-5}, \hyperref[q:r02-reg-basic-1-6]{1-6}, \hyperref[q:r02-reg-basic-2-1]{2-1}, \hyperref[q:r02-reg-basic-2-2]{2-2}, \hyperref[q:r02-reg-basic-2-3]{2-3}, \hyperref[q:r02-reg-basic-2-4]{2-4}, \hyperref[q:r02-reg-basic-2-5]{2-5}, \hyperref[q:r02-reg-basic-2-6]{2-6}, \hyperref[q:r02-reg-basic-3-1]{3-1}, \hyperref[q:r02-reg-basic-3-2]{3-2}, \hyperref[q:r02-reg-basic-3-3]{3-3}, \hyperref[q:r02-reg-basic-3-4]{3-4}, \hyperref[q:r02-reg-basic-3-5]{3-5}, \hyperref[q:r02-reg-basic-3-6]{3-6}, \hyperref[q:r02-reg-basic-4-1]{4-1}, \hyperref[q:r02-reg-basic-4-2]{4-2}, \hyperref[q:r02-reg-basic-4-3]{4-3}, \hyperref[q:r02-reg-basic-4-4]{4-4}, \hyperref[q:r02-reg-basic-4-5]{4-5}, \hyperref[q:r02-reg-basic-4-6]{4-6}, \hyperref[q:r02-reg-basic-5-1]{5-1}, \hyperref[q:r02-reg-basic-5-2]{5-2}, \hyperref[q:r02-reg-basic-5-3]{5-3}, \hyperref[q:r02-reg-basic-5-4]{5-4}, \hyperref[q:r02-reg-basic-5-5]{5-5}, \hyperref[q:r02-reg-basic-5-6]{5-6}
\paragraph{適性（15問）} \hyperref[q:r02-reg-apt-01]{1}, \hyperref[q:r02-reg-apt-02]{2}, \hyperref[q:r02-reg-apt-03]{3}, \hyperref[q:r02-reg-apt-04]{4}, \hyperref[q:r02-reg-apt-05]{5}, \hyperref[q:r02-reg-apt-06]{6}, \hyperref[q:r02-reg-apt-07]{7}, \hyperref[q:r02-reg-apt-08]{8}, \hyperref[q:r02-reg-apt-09]{9}, \hyperref[q:r02-reg-apt-10]{10}, \hyperref[q:r02-reg-apt-11]{11}, \hyperref[q:r02-reg-apt-12]{12}, \hyperref[q:r02-reg-apt-13]{13}, \hyperref[q:r02-reg-apt-14]{14}, \hyperref[q:r02-reg-apt-15]{15}
\paragraph{応用理学（35問）} \hyperref[q:r02-reg-applied-01]{1}, \hyperref[q:r02-reg-applied-02]{2}, \hyperref[q:r02-reg-applied-03]{3}, \hyperref[q:r02-reg-applied-04]{4}, \hyperref[q:r02-reg-applied-05]{5}, \hyperref[q:r02-reg-applied-06]{6}, \hyperref[q:r02-reg-applied-07]{7}, \hyperref[q:r02-reg-applied-08]{8}, \hyperref[q:r02-reg-applied-09]{9}, \hyperref[q:r02-reg-applied-10]{10}, \hyperref[q:r02-reg-applied-11]{11}, \hyperref[q:r02-reg-applied-12]{12}, \hyperref[q:r02-reg-applied-13]{13}, \hyperref[q:r02-reg-applied-14]{14}, \hyperref[q:r02-reg-applied-15]{15}, \hyperref[q:r02-reg-applied-16]{16}, \hyperref[q:r02-reg-applied-17]{17}, \hyperref[q:r02-reg-applied-18]{18}, \hyperref[q:r02-reg-applied-19]{19}, \hyperref[q:r02-reg-applied-20]{20}, \hyperref[q:r02-reg-applied-21]{21}, \hyperref[q:r02-reg-applied-22]{22}, \hyperref[q:r02-reg-applied-23]{23}, \hyperref[q:r02-reg-applied-24]{24}, \hyperref[q:r02-reg-applied-25]{25}, \hyperref[q:r02-reg-applied-26]{26}, \hyperref[q:r02-reg-applied-27]{27}, \hyperref[q:r02-reg-applied-28]{28}, \hyperref[q:r02-reg-applied-29]{29}, \hyperref[q:r02-reg-applied-30]{30}, \hyperref[q:r02-reg-applied-31]{31}, \hyperref[q:r02-reg-applied-32]{32}, \hyperref[q:r02-reg-applied-33]{33}, \hyperref[q:r02-reg-applied-34]{34}, \hyperref[q:r02-reg-applied-35]{35}
\section{令和3年度}
\paragraph{基礎（30問）} \hyperref[q:r03-reg-basic-1-1]{1-1}, \hyperref[q:r03-reg-basic-1-2]{1-2}, \hyperref[q:r03-reg-basic-1-3]{1-3}, \hyperref[q:r03-reg-basic-1-4]{1-4}, \hyperref[q:r03-reg-basic-1-5]{1-5}, \hyperref[q:r03-reg-basic-1-6]{1-6}, \hyperref[q:r03-reg-basic-2-1]{2-1}, \hyperref[q:r03-reg-basic-2-2]{2-2}, \hyperref[q:r03-reg-basic-2-3]{2-3}, \hyperref[q:r03-reg-basic-2-4]{2-4}, \hyperref[q:r03-reg-basic-2-5]{2-5}, \hyperref[q:r03-reg-basic-2-6]{2-6}, \hyperref[q:r03-reg-basic-3-1]{3-1}, \hyperref[q:r03-reg-basic-3-2]{3-2}, \hyperref[q:r03-reg-basic-3-3]{3-3}, \hyperref[q:r03-reg-basic-3-4]{3-4}, \hyperref[q:r03-reg-basic-3-5]{3-5}, \hyperref[q:r03-reg-basic-3-6]{3-6}, \hyperref[q:r03-reg-basic-4-1]{4-1}, \hyperref[q:r03-reg-basic-4-2]{4-2}, \hyperref[q:r03-reg-basic-4-3]{4-3}, \hyperref[q:r03-reg-basic-4-4]{4-4}, \hyperref[q:r03-reg-basic-4-5]{4-5}, \hyperref[q:r03-reg-basic-4-6]{4-6}, \hyperref[q:r03-reg-basic-5-1]{5-1}, \hyperref[q:r03-reg-basic-5-2]{5-2}, \hyperref[q:r03-reg-basic-5-3]{5-3}, \hyperref[q:r03-reg-basic-5-4]{5-4}, \hyperref[q:r03-reg-basic-5-5]{5-5}, \hyperref[q:r03-reg-basic-5-6]{5-6}
\paragraph{適性（15問）} \hyperref[q:r03-reg-apt-01]{1}, \hyperref[q:r03-reg-apt-02]{2}, \hyperref[q:r03-reg-apt-03]{3}, \hyperref[q:r03-reg-apt-04]{4}, \hyperref[q:r03-reg-apt-05]{5}, \hyperref[q:r03-reg-apt-06]{6}, \hyperref[q:r03-reg-apt-07]{7}, \hyperref[q:r03-reg-apt-08]{8}, \hyperref[q:r03-reg-apt-09]{9}, \hyperref[q:r03-reg-apt-10]{10}, \hyperref[q:r03-reg-apt-11]{11}, \hyperref[q:r03-reg-apt-12]{12}, \hyperref[q:r03-reg-apt-13]{13}, \hyperref[q:r03-reg-apt-14]{14}, \hyperref[q:r03-reg-apt-15]{15}
\paragraph{応用理学（35問）} \hyperref[q:r03-reg-applied-01]{1}, \hyperref[q:r03-reg-applied-02]{2}, \hyperref[q:r03-reg-applied-03]{3}, \hyperref[q:r03-reg-applied-04]{4}, \hyperref[q:r03-reg-applied-05]{5}, \hyperref[q:r03-reg-applied-06]{6}, \hyperref[q:r03-reg-applied-07]{7}, \hyperref[q:r03-reg-applied-08]{8}, \hyperref[q:r03-reg-applied-09]{9}, \hyperref[q:r03-reg-applied-10]{10}, \hyperref[q:r03-reg-applied-11]{11}, \hyperref[q:r03-reg-applied-12]{12}, \hyperref[q:r03-reg-applied-13]{13}, \hyperref[q:r03-reg-applied-14]{14}, \hyperref[q:r03-reg-applied-15]{15}, \hyperref[q:r03-reg-applied-16]{16}, \hyperref[q:r03-reg-applied-17]{17}, \hyperref[q:r03-reg-applied-18]{18}, \hyperref[q:r03-reg-applied-19]{19}, \hyperref[q:r03-reg-applied-20]{20}, \hyperref[q:r03-reg-applied-21]{21}, \hyperref[q:r03-reg-applied-22]{22}, \hyperref[q:r03-reg-applied-23]{23}, \hyperref[q:r03-reg-applied-24]{24}, \hyperref[q:r03-reg-applied-25]{25}, \hyperref[q:r03-reg-applied-26]{26}, \hyperref[q:r03-reg-applied-27]{27}, \hyperref[q:r03-reg-applied-28]{28}, \hyperref[q:r03-reg-applied-29]{29}, \hyperref[q:r03-reg-applied-30]{30}, \hyperref[q:r03-reg-applied-31]{31}, \hyperref[q:r03-reg-applied-32]{32}, \hyperref[q:r03-reg-applied-33]{33}, \hyperref[q:r03-reg-applied-34]{34}, \hyperref[q:r03-reg-applied-35]{35}
\section{令和4年度}
\paragraph{基礎（30問）} \hyperref[q:r04-reg-basic-1-1]{1-1}, \hyperref[q:r04-reg-basic-1-2]{1-2}, \hyperref[q:r04-reg-basic-1-3]{1-3}, \hyperref[q:r04-reg-basic-1-4]{1-4}, \hyperref[q:r04-reg-basic-1-5]{1-5}, \hyperref[q:r04-reg-basic-1-6]{1-6}, \hyperref[q:r04-reg-basic-2-1]{2-1}, \hyperref[q:r04-reg-basic-2-2]{2-2}, \hyperref[q:r04-reg-basic-2-3]{2-3}, \hyperref[q:r04-reg-basic-2-4]{2-4}, \hyperref[q:r04-reg-basic-2-5]{2-5}, \hyperref[q:r04-reg-basic-2-6]{2-6}, \hyperref[q:r04-reg-basic-3-1]{3-1}, \hyperref[q:r04-reg-basic-3-2]{3-2}, \hyperref[q:r04-reg-basic-3-3]{3-3}, \hyperref[q:r04-reg-basic-3-4]{3-4}, \hyperref[q:r04-reg-basic-3-5]{3-5}, \hyperref[q:r04-reg-basic-3-6]{3-6}, \hyperref[q:r04-reg-basic-4-1]{4-1}, \hyperref[q:r04-reg-basic-4-2]{4-2}, \hyperref[q:r04-reg-basic-4-3]{4-3}, \hyperref[q:r04-reg-basic-4-4]{4-4}, \hyperref[q:r04-reg-basic-4-5]{4-5}, \hyperref[q:r04-reg-basic-4-6]{4-6}, \hyperref[q:r04-reg-basic-5-1]{5-1}, \hyperref[q:r04-reg-basic-5-2]{5-2}, \hyperref[q:r04-reg-basic-5-3]{5-3}, \hyperref[q:r04-reg-basic-5-4]{5-4}, \hyperref[q:r04-reg-basic-5-5]{5-5}, \hyperref[q:r04-reg-basic-5-6]{5-6}
\paragraph{適性（15問）} \hyperref[q:r04-reg-apt-01]{1}, \hyperref[q:r04-reg-apt-02]{2}, \hyperref[q:r04-reg-apt-03]{3}, \hyperref[q:r04-reg-apt-04]{4}, \hyperref[q:r04-reg-apt-05]{5}, \hyperref[q:r04-reg-apt-06]{6}, \hyperref[q:r04-reg-apt-07]{7}, \hyperref[q:r04-reg-apt-08]{8}, \hyperref[q:r04-reg-apt-09]{9}, \hyperref[q:r04-reg-apt-10]{10}, \hyperref[q:r04-reg-apt-11]{11}, \hyperref[q:r04-reg-apt-12]{12}, \hyperref[q:r04-reg-apt-13]{13}, \hyperref[q:r04-reg-apt-14]{14}, \hyperref[q:r04-reg-apt-15]{15}
\paragraph{応用理学（35問）} \hyperref[q:r04-reg-applied-01]{1}, \hyperref[q:r04-reg-applied-02]{2}, \hyperref[q:r04-reg-applied-03]{3}, \hyperref[q:r04-reg-applied-04]{4}, \hyperref[q:r04-reg-applied-05]{5}, \hyperref[q:r04-reg-applied-06]{6}, \hyperref[q:r04-reg-applied-07]{7}, \hyperref[q:r04-reg-applied-08]{8}, \hyperref[q:r04-reg-applied-09]{9}, \hyperref[q:r04-reg-applied-10]{10}, \hyperref[q:r04-reg-applied-11]{11}, \hyperref[q:r04-reg-applied-12]{12}, \hyperref[q:r04-reg-applied-13]{13}, \hyperref[q:r04-reg-applied-14]{14}, \hyperref[q:r04-reg-applied-15]{15}, \hyperref[q:r04-reg-applied-16]{16}, \hyperref[q:r04-reg-applied-17]{17}, \hyperref[q:r04-reg-applied-18]{18}, \hyperref[q:r04-reg-applied-19]{19}, \hyperref[q:r04-reg-applied-20]{20}, \hyperref[q:r04-reg-applied-21]{21}, \hyperref[q:r04-reg-applied-22]{22}, \hyperref[q:r04-reg-applied-23]{23}, \hyperref[q:r04-reg-applied-24]{24}, \hyperref[q:r04-reg-applied-25]{25}, \hyperref[q:r04-reg-applied-26]{26}, \hyperref[q:r04-reg-applied-27]{27}, \hyperref[q:r04-reg-applied-28]{28}, \hyperref[q:r04-reg-applied-29]{29}, \hyperref[q:r04-reg-applied-30]{30}, \hyperref[q:r04-reg-applied-31]{31}, \hyperref[q:r04-reg-applied-32]{32}, \hyperref[q:r04-reg-applied-33]{33}, \hyperref[q:r04-reg-applied-34]{34}, \hyperref[q:r04-reg-applied-35]{35}
\section{令和5年度}
\paragraph{基礎（30問）} \hyperref[q:r05-reg-basic-1-1]{1-1}, \hyperref[q:r05-reg-basic-1-2]{1-2}, \hyperref[q:r05-reg-basic-1-3]{1-3}, \hyperref[q:r05-reg-basic-1-4]{1-4}, \hyperref[q:r05-reg-basic-1-5]{1-5}, \hyperref[q:r05-reg-basic-1-6]{1-6}, \hyperref[q:r05-reg-basic-2-1]{2-1}, \hyperref[q:r05-reg-basic-2-2]{2-2}, \hyperref[q:r05-reg-basic-2-3]{2-3}, \hyperref[q:r05-reg-basic-2-4]{2-4}, \hyperref[q:r05-reg-basic-2-5]{2-5}, \hyperref[q:r05-reg-basic-2-6]{2-6}, \hyperref[q:r05-reg-basic-3-1]{3-1}, \hyperref[q:r05-reg-basic-3-2]{3-2}, \hyperref[q:r05-reg-basic-3-3]{3-3}, \hyperref[q:r05-reg-basic-3-4]{3-4}, \hyperref[q:r05-reg-basic-3-5]{3-5}, \hyperref[q:r05-reg-basic-3-6]{3-6}, \hyperref[q:r05-reg-basic-4-1]{4-1}, \hyperref[q:r05-reg-basic-4-2]{4-2}, \hyperref[q:r05-reg-basic-4-3]{4-3}, \hyperref[q:r05-reg-basic-4-4]{4-4}, \hyperref[q:r05-reg-basic-4-5]{4-5}, \hyperref[q:r05-reg-basic-4-6]{4-6}, \hyperref[q:r05-reg-basic-5-1]{5-1}, \hyperref[q:r05-reg-basic-5-2]{5-2}, \hyperref[q:r05-reg-basic-5-3]{5-3}, \hyperref[q:r05-reg-basic-5-4]{5-4}, \hyperref[q:r05-reg-basic-5-5]{5-5}, \hyperref[q:r05-reg-basic-5-6]{5-6}
\paragraph{適性（15問）} \hyperref[q:r05-reg-apt-01]{1}, \hyperref[q:r05-reg-apt-02]{2}, \hyperref[q:r05-reg-apt-03]{3}, \hyperref[q:r05-reg-apt-04]{4}, \hyperref[q:r05-reg-apt-05]{5}, \hyperref[q:r05-reg-apt-06]{6}, \hyperref[q:r05-reg-apt-07]{7}, \hyperref[q:r05-reg-apt-08]{8}, \hyperref[q:r05-reg-apt-09]{9}, \hyperref[q:r05-reg-apt-10]{10}, \hyperref[q:r05-reg-apt-11]{11}, \hyperref[q:r05-reg-apt-12]{12}, \hyperref[q:r05-reg-apt-13]{13}, \hyperref[q:r05-reg-apt-14]{14}, \hyperref[q:r05-reg-apt-15]{15}
\paragraph{応用理学（35問）} \hyperref[q:r05-reg-applied-01]{1}, \hyperref[q:r05-reg-applied-02]{2}, \hyperref[q:r05-reg-applied-03]{3}, \hyperref[q:r05-reg-applied-04]{4}, \hyperref[q:r05-reg-applied-05]{5}, \hyperref[q:r05-reg-applied-06]{6}, \hyperref[q:r05-reg-applied-07]{7}, \hyperref[q:r05-reg-applied-08]{8}, \hyperref[q:r05-reg-applied-09]{9}, \hyperref[q:r05-reg-applied-10]{10}, \hyperref[q:r05-reg-applied-11]{11}, \hyperref[q:r05-reg-applied-12]{12}, \hyperref[q:r05-reg-applied-13]{13}, \hyperref[q:r05-reg-applied-14]{14}, \hyperref[q:r05-reg-applied-15]{15}, \hyperref[q:r05-reg-applied-16]{16}, \hyperref[q:r05-reg-applied-17]{17}, \hyperref[q:r05-reg-applied-18]{18}, \hyperref[q:r05-reg-applied-19]{19}, \hyperref[q:r05-reg-applied-20]{20}, \hyperref[q:r05-reg-applied-21]{21}, \hyperref[q:r05-reg-applied-22]{22}, \hyperref[q:r05-reg-applied-23]{23}, \hyperref[q:r05-reg-applied-24]{24}, \hyperref[q:r05-reg-applied-25]{25}, \hyperref[q:r05-reg-applied-26]{26}, \hyperref[q:r05-reg-applied-27]{27}, \hyperref[q:r05-reg-applied-28]{28}, \hyperref[q:r05-reg-applied-29]{29}, \hyperref[q:r05-reg-applied-30]{30}, \hyperref[q:r05-reg-applied-31]{31}, \hyperref[q:r05-reg-applied-32]{32}, \hyperref[q:r05-reg-applied-33]{33}, \hyperref[q:r05-reg-applied-34]{34}, \hyperref[q:r05-reg-applied-35]{35}
\section{令和6年度}
\paragraph{基礎（30問）} \hyperref[q:r06-reg-basic-1-1]{1-1}, \hyperref[q:r06-reg-basic-1-2]{1-2}, \hyperref[q:r06-reg-basic-1-3]{1-3}, \hyperref[q:r06-reg-basic-1-4]{1-4}, \hyperref[q:r06-reg-basic-1-5]{1-5}, \hyperref[q:r06-reg-basic-1-6]{1-6}, \hyperref[q:r06-reg-basic-2-1]{2-1}, \hyperref[q:r06-reg-basic-2-2]{2-2}, \hyperref[q:r06-reg-basic-2-3]{2-3}, \hyperref[q:r06-reg-basic-2-4]{2-4}, \hyperref[q:r06-reg-basic-2-5]{2-5}, \hyperref[q:r06-reg-basic-2-6]{2-6}, \hyperref[q:r06-reg-basic-3-1]{3-1}, \hyperref[q:r06-reg-basic-3-2]{3-2}, \hyperref[q:r06-reg-basic-3-3]{3-3}, \hyperref[q:r06-reg-basic-3-4]{3-4}, \hyperref[q:r06-reg-basic-3-5]{3-5}, \hyperref[q:r06-reg-basic-3-6]{3-6}, \hyperref[q:r06-reg-basic-4-1]{4-1}, \hyperref[q:r06-reg-basic-4-2]{4-2}, \hyperref[q:r06-reg-basic-4-3]{4-3}, \hyperref[q:r06-reg-basic-4-4]{4-4}, \hyperref[q:r06-reg-basic-4-5]{4-5}, \hyperref[q:r06-reg-basic-4-6]{4-6}, \hyperref[q:r06-reg-basic-5-1]{5-1}, \hyperref[q:r06-reg-basic-5-2]{5-2}, \hyperref[q:r06-reg-basic-5-3]{5-3}, \hyperref[q:r06-reg-basic-5-4]{5-4}, \hyperref[q:r06-reg-basic-5-5]{5-5}, \hyperref[q:r06-reg-basic-5-6]{5-6}
\paragraph{適性（15問）} \hyperref[q:r06-reg-apt-01]{1}, \hyperref[q:r06-reg-apt-02]{2}, \hyperref[q:r06-reg-apt-03]{3}, \hyperref[q:r06-reg-apt-04]{4}, \hyperref[q:r06-reg-apt-05]{5}, \hyperref[q:r06-reg-apt-06]{6}, \hyperref[q:r06-reg-apt-07]{7}, \hyperref[q:r06-reg-apt-08]{8}, \hyperref[q:r06-reg-apt-09]{9}, \hyperref[q:r06-reg-apt-10]{10}, \hyperref[q:r06-reg-apt-11]{11}, \hyperref[q:r06-reg-apt-12]{12}, \hyperref[q:r06-reg-apt-13]{13}, \hyperref[q:r06-reg-apt-14]{14}, \hyperref[q:r06-reg-apt-15]{15}
\paragraph{応用理学（35問）} \hyperref[q:r06-reg-applied-01]{1}, \hyperref[q:r06-reg-applied-02]{2}, \hyperref[q:r06-reg-applied-03]{3}, \hyperref[q:r06-reg-applied-04]{4}, \hyperref[q:r06-reg-applied-05]{5}, \hyperref[q:r06-reg-applied-06]{6}, \hyperref[q:r06-reg-applied-07]{7}, \hyperref[q:r06-reg-applied-08]{8}, \hyperref[q:r06-reg-applied-09]{9}, \hyperref[q:r06-reg-applied-10]{10}, \hyperref[q:r06-reg-applied-11]{11}, \hyperref[q:r06-reg-applied-12]{12}, \hyperref[q:r06-reg-applied-13]{13}, \hyperref[q:r06-reg-applied-14]{14}, \hyperref[q:r06-reg-applied-15]{15}, \hyperref[q:r06-reg-applied-16]{16}, \hyperref[q:r06-reg-applied-17]{17}, \hyperref[q:r06-reg-applied-18]{18}, \hyperref[q:r06-reg-applied-19]{19}, \hyperref[q:r06-reg-applied-20]{20}, \hyperref[q:r06-reg-applied-21]{21}, \hyperref[q:r06-reg-applied-22]{22}, \hyperref[q:r06-reg-applied-23]{23}, \hyperref[q:r06-reg-applied-24]{24}, \hyperref[q:r06-reg-applied-25]{25}, \hyperref[q:r06-reg-applied-26]{26}, \hyperref[q:r06-reg-applied-27]{27}, \hyperref[q:r06-reg-applied-28]{28}, \hyperref[q:r06-reg-applied-29]{29}, \hyperref[q:r06-reg-applied-30]{30}, \hyperref[q:r06-reg-applied-31]{31}, \hyperref[q:r06-reg-applied-32]{32}, \hyperref[q:r06-reg-applied-33]{33}, \hyperref[q:r06-reg-applied-34]{34}, \hyperref[q:r06-reg-applied-35]{35}
\section{令和7年度}
\paragraph{基礎（30問）} \hyperref[q:r07-reg-basic-1-1]{1-1}, \hyperref[q:r07-reg-basic-1-2]{1-2}, \hyperref[q:r07-reg-basic-1-3]{1-3}, \hyperref[q:r07-reg-basic-1-4]{1-4}, \hyperref[q:r07-reg-basic-1-5]{1-5}, \hyperref[q:r07-reg-basic-1-6]{1-6}, \hyperref[q:r07-reg-basic-2-1]{2-1}, \hyperref[q:r07-reg-basic-2-2]{2-2}, \hyperref[q:r07-reg-basic-2-3]{2-3}, \hyperref[q:r07-reg-basic-2-4]{2-4}, \hyperref[q:r07-reg-basic-2-5]{2-5}, \hyperref[q:r07-reg-basic-2-6]{2-6}, \hyperref[q:r07-reg-basic-3-1]{3-1}, \hyperref[q:r07-reg-basic-3-2]{3-2}, \hyperref[q:r07-reg-basic-3-3]{3-3}, \hyperref[q:r07-reg-basic-3-4]{3-4}, \hyperref[q:r07-reg-basic-3-5]{3-5}, \hyperref[q:r07-reg-basic-3-6]{3-6}, \hyperref[q:r07-reg-basic-4-1]{4-1}, \hyperref[q:r07-reg-basic-4-2]{4-2}, \hyperref[q:r07-reg-basic-4-3]{4-3}, \hyperref[q:r07-reg-basic-4-4]{4-4}, \hyperref[q:r07-reg-basic-4-5]{4-5}, \hyperref[q:r07-reg-basic-4-6]{4-6}, \hyperref[q:r07-reg-basic-5-1]{5-1}, \hyperref[q:r07-reg-basic-5-2]{5-2}, \hyperref[q:r07-reg-basic-5-3]{5-3}, \hyperref[q:r07-reg-basic-5-4]{5-4}, \hyperref[q:r07-reg-basic-5-5]{5-5}, \hyperref[q:r07-reg-basic-5-6]{5-6}
\paragraph{適性（15問）} \hyperref[q:r07-reg-apt-01]{1}, \hyperref[q:r07-reg-apt-02]{2}, \hyperref[q:r07-reg-apt-03]{3}, \hyperref[q:r07-reg-apt-04]{4}, \hyperref[q:r07-reg-apt-05]{5}, \hyperref[q:r07-reg-apt-06]{6}, \hyperref[q:r07-reg-apt-07]{7}, \hyperref[q:r07-reg-apt-08]{8}, \hyperref[q:r07-reg-apt-09]{9}, \hyperref[q:r07-reg-apt-10]{10}, \hyperref[q:r07-reg-apt-11]{11}, \hyperref[q:r07-reg-apt-12]{12}, \hyperref[q:r07-reg-apt-13]{13}, \hyperref[q:r07-reg-apt-14]{14}, \hyperref[q:r07-reg-apt-15]{15}
\paragraph{応用理学（35問）} \hyperref[q:r07-reg-applied-01]{1}, \hyperref[q:r07-reg-applied-02]{2}, \hyperref[q:r07-reg-applied-03]{3}, \hyperref[q:r07-reg-applied-04]{4}, \hyperref[q:r07-reg-applied-05]{5}, \hyperref[q:r07-reg-applied-06]{6}, \hyperref[q:r07-reg-applied-07]{7}, \hyperref[q:r07-reg-applied-08]{8}, \hyperref[q:r07-reg-applied-09]{9}, \hyperref[q:r07-reg-applied-10]{10}, \hyperref[q:r07-reg-applied-11]{11}, \hyperref[q:r07-reg-applied-12]{12}, \hyperref[q:r07-reg-applied-13]{13}, \hyperref[q:r07-reg-applied-14]{14}, \hyperref[q:r07-reg-applied-15]{15}, \hyperref[q:r07-reg-applied-16]{16}, \hyperref[q:r07-reg-applied-17]{17}, \hyperref[q:r07-reg-applied-18]{18}, \hyperref[q:r07-reg-applied-19]{19}, \hyperref[q:r07-reg-applied-20]{20}, \hyperref[q:r07-reg-applied-21]{21}, \hyperref[q:r07-reg-applied-22]{22}, \hyperref[q:r07-reg-applied-23]{23}, \hyperref[q:r07-reg-applied-24]{24}, \hyperref[q:r07-reg-applied-25]{25}, \hyperref[q:r07-reg-applied-26]{26}, \hyperref[q:r07-reg-applied-27]{27}, \hyperref[q:r07-reg-applied-28]{28}, \hyperref[q:r07-reg-applied-29]{29}, \hyperref[q:r07-reg-applied-30]{30}, \hyperref[q:r07-reg-applied-31]{31}, \hyperref[q:r07-reg-applied-32]{32}, \hyperref[q:r07-reg-applied-33]{33}, \hyperref[q:r07-reg-applied-34]{34}, \hyperref[q:r07-reg-applied-35]{35}

\chapter{分野別索引}
分類は確認済みと暫定を含む。概念タグの横断頻度は主分類の問題数と別に扱う。
\section{APPLIED：光学}
\hyperref[q:h23-reg-applied-07]{H23:7}, \hyperref[q:h24-reg-applied-03]{H24:3}, \hyperref[q:h26-reg-applied-07]{H26:7}, \hyperref[q:h27-reg-applied-03]{H27:3}, \hyperref[q:h28-reg-applied-05]{H28:5}, \hyperref[q:h30-reg-applied-12]{H30:12}, \hyperref[q:h30-reg-applied-14]{H30:14}, \hyperref[q:r01-ret-applied-03]{R01-再:3}, \hyperref[q:r02-reg-applied-05]{R02:5}, \hyperref[q:r02-reg-applied-07]{R02:7}, \hyperref[q:r02-reg-applied-12]{R02:12}, \hyperref[q:r03-reg-applied-05]{R03:5}, \hyperref[q:r05-reg-applied-09]{R05:9}, \hyperref[q:r06-reg-applied-05]{R06:5}
\section{APPLIED：分析化学}
\hyperref[q:h24-reg-applied-16]{H24:16}, \hyperref[q:h24-reg-applied-27]{H24:27}, \hyperref[q:h25-reg-applied-15]{H25:15}, \hyperref[q:h25-reg-applied-16]{H25:16}, \hyperref[q:h26-reg-applied-16]{H26:16}, \hyperref[q:h28-reg-applied-16]{H28:16}, \hyperref[q:h29-reg-applied-12]{H29:12}, \hyperref[q:h30-reg-applied-13]{H30:13}, \hyperref[q:r01-reg-applied-12]{R01:12}, \hyperref[q:r02-reg-applied-11]{R02:11}, \hyperref[q:r03-reg-applied-16]{R03:16}, \hyperref[q:r05-reg-applied-15]{R05:15}, \hyperref[q:r05-reg-applied-22]{R05:22}, \hyperref[q:r06-reg-applied-14]{R06:14}, \hyperref[q:r06-reg-applied-19]{R06:19}, \hyperref[q:r07-reg-applied-14]{R07:14}
\section{APPLIED：力学}
\hyperref[q:h23-reg-applied-01]{H23:1}, \hyperref[q:h23-reg-applied-04]{H23:4}, \hyperref[q:h24-reg-applied-01]{H24:1}, \hyperref[q:h24-reg-applied-09]{H24:9}, \hyperref[q:h25-reg-applied-01]{H25:1}, \hyperref[q:h25-reg-applied-06]{H25:6}, \hyperref[q:h26-reg-applied-01]{H26:1}, \hyperref[q:h26-reg-applied-03]{H26:3}, \hyperref[q:h28-reg-applied-01]{H28:1}, \hyperref[q:h28-reg-applied-02]{H28:2}, \hyperref[q:h30-reg-applied-01]{H30:1}, \hyperref[q:h30-reg-applied-02]{H30:2}, \hyperref[q:r01-reg-applied-01]{R01:1}, \hyperref[q:r01-reg-applied-02]{R01:2}, \hyperref[q:r01-ret-applied-02]{R01-再:2}, \hyperref[q:r02-reg-applied-01]{R02:1}, \hyperref[q:r02-reg-applied-02]{R02:2}, \hyperref[q:r03-reg-applied-01]{R03:1}, \hyperref[q:r03-reg-applied-02]{R03:2}, \hyperref[q:r04-reg-applied-01]{R04:1}, \hyperref[q:r04-reg-applied-02]{R04:2}, \hyperref[q:r05-reg-applied-01]{R05:1}, \hyperref[q:r05-reg-applied-02]{R05:2}, \hyperref[q:r06-reg-applied-01]{R06:1}, \hyperref[q:r06-reg-applied-02]{R06:2}, \hyperref[q:r07-reg-applied-01]{R07:1}
\section{APPLIED：単位・計測}
\hyperref[q:h23-reg-applied-06]{H23:6}, \hyperref[q:h25-reg-applied-04]{H25:4}, \hyperref[q:h27-reg-applied-04]{H27:4}, \hyperref[q:h28-reg-applied-06]{H28:6}, \hyperref[q:h29-reg-applied-05]{H29:5}, \hyperref[q:r01-ret-applied-01]{R01-再:1}, \hyperref[q:r02-reg-applied-09]{R02:9}, \hyperref[q:r03-reg-applied-09]{R03:9}, \hyperref[q:r05-reg-applied-13]{R05:13}, \hyperref[q:r06-reg-applied-16]{R06:16}, \hyperref[q:r07-reg-applied-09]{R07:9}
\section{APPLIED：原子・分子・結晶}
\hyperref[q:h24-reg-applied-05]{H24:5}, \hyperref[q:h26-reg-applied-11]{H26:11}, \hyperref[q:h26-reg-applied-14]{H26:14}, \hyperref[q:h27-reg-applied-09]{H27:9}, \hyperref[q:h28-reg-applied-07]{H28:7}, \hyperref[q:h28-reg-applied-10]{H28:10}, \hyperref[q:h28-reg-applied-14]{H28:14}, \hyperref[q:h29-reg-applied-11]{H29:11}, \hyperref[q:h30-reg-applied-09]{H30:9}, \hyperref[q:r01-reg-applied-14]{R01:14}, \hyperref[q:r01-reg-applied-17]{R01:17}, \hyperref[q:r01-ret-applied-07]{R01-再:7}, \hyperref[q:r01-ret-applied-14]{R01-再:14}, \hyperref[q:r02-reg-applied-10]{R02:10}, \hyperref[q:r03-reg-applied-08]{R03:8}, \hyperref[q:r03-reg-applied-14]{R03:14}, \hyperref[q:r04-reg-applied-08]{R04:8}, \hyperref[q:r04-reg-applied-12]{R04:12}, \hyperref[q:r05-reg-applied-10]{R05:10}, \hyperref[q:r06-reg-applied-08]{R06:8}, \hyperref[q:r07-reg-applied-08]{R07:8}
\section{APPLIED：同位体地球化学}
\hyperref[q:h23-reg-applied-21]{H23:21}, \hyperref[q:h23-reg-applied-24]{H23:24}, \hyperref[q:h24-reg-applied-32]{H24:32}, \hyperref[q:h25-reg-applied-27]{H25:27}, \hyperref[q:h26-reg-applied-31]{H26:31}, \hyperref[q:h27-reg-applied-25]{H27:25}, \hyperref[q:h28-reg-applied-32]{H28:32}, \hyperref[q:h29-reg-applied-18]{H29:18}, \hyperref[q:h29-reg-applied-30]{H29:30}, \hyperref[q:h30-reg-applied-21]{H30:21}, \hyperref[q:r01-ret-applied-35]{R01-再:35}, \hyperref[q:r02-reg-applied-20]{R02:20}, \hyperref[q:r04-reg-applied-25]{R04:25}, \hyperref[q:r05-reg-applied-21]{R05:21}, \hyperref[q:r06-reg-applied-21]{R06:21}, \hyperref[q:r07-reg-applied-22]{R07:22}
\section{APPLIED：地下水}
\hyperref[q:h24-reg-applied-35]{H24:35}, \hyperref[q:h25-reg-applied-29]{H25:29}, \hyperref[q:h26-reg-applied-24]{H26:24}, \hyperref[q:h27-reg-applied-23]{H27:23}, \hyperref[q:r01-ret-applied-29]{R01-再:29}, \hyperref[q:r07-reg-applied-19]{R07:19}
\section{APPLIED：地球内部・プレート}
\hyperref[q:h26-reg-applied-30]{H26:30}, \hyperref[q:h27-reg-applied-32]{H27:32}, \hyperref[q:h28-reg-applied-30]{H28:30}, \hyperref[q:h29-reg-applied-26]{H29:26}, \hyperref[q:r01-reg-applied-30]{R01:30}, \hyperref[q:r01-ret-applied-33]{R01-再:33}, \hyperref[q:r01-ret-applied-34]{R01-再:34}, \hyperref[q:r02-reg-applied-28]{R02:28}, \hyperref[q:r03-reg-applied-31]{R03:31}, \hyperref[q:r04-reg-applied-26]{R04:26}, \hyperref[q:r07-reg-applied-24]{R07:24}
\section{APPLIED：地球化学}
\hyperref[q:h23-reg-applied-18]{H23:18}, \hyperref[q:h23-reg-applied-19]{H23:19}, \hyperref[q:h23-reg-applied-22]{H23:22}, \hyperref[q:h24-reg-applied-33]{H24:33}, \hyperref[q:h25-reg-applied-10]{H25:10}, \hyperref[q:h25-reg-applied-31]{H25:31}, \hyperref[q:h26-reg-applied-18]{H26:18}, \hyperref[q:h26-reg-applied-25]{H26:25}, \hyperref[q:h26-reg-applied-27]{H26:27}, \hyperref[q:h27-reg-applied-27]{H27:27}, \hyperref[q:h28-reg-applied-27]{H28:27}, \hyperref[q:h29-reg-applied-20]{H29:20}, \hyperref[q:h30-reg-applied-28]{H30:28}, \hyperref[q:r01-reg-applied-20]{R01:20}, \hyperref[q:r01-ret-applied-18]{R01-再:18}, \hyperref[q:r02-reg-applied-19]{R02:19}, \hyperref[q:r03-reg-applied-18]{R03:18}, \hyperref[q:r03-reg-applied-19]{R03:19}, \hyperref[q:r03-reg-applied-21]{R03:21}, \hyperref[q:r04-reg-applied-19]{R04:19}, \hyperref[q:r04-reg-applied-20]{R04:20}, \hyperref[q:r06-reg-applied-22]{R06:22}
\section{APPLIED：地磁気}
\hyperref[q:r01-reg-applied-27]{R01:27}
\section{APPLIED：地震学}
\hyperref[q:h26-reg-applied-32]{H26:32}
\section{APPLIED：堆積学}
\hyperref[q:h23-reg-applied-25]{H23:25}, \hyperref[q:h23-reg-applied-35]{H23:35}, \hyperref[q:h29-reg-applied-25]{H29:25}, \hyperref[q:h30-reg-applied-25]{H30:25}, \hyperref[q:r01-reg-applied-28]{R01:28}, \hyperref[q:r02-reg-applied-29]{R02:29}, \hyperref[q:r03-reg-applied-30]{R03:30}, \hyperref[q:r04-reg-applied-28]{R04:28}, \hyperref[q:r05-reg-applied-30]{R05:30}, \hyperref[q:r05-reg-applied-31]{R05:31}, \hyperref[q:r06-reg-applied-30]{R06:30}, \hyperref[q:r06-reg-applied-31]{R06:31}, \hyperref[q:r07-reg-applied-27]{R07:27}
\section{APPLIED：層序・古生物}
\hyperref[q:h28-reg-applied-34]{H28:34}, \hyperref[q:r02-reg-applied-31]{R02:31}, \hyperref[q:r07-reg-applied-28]{R07:28}
\section{APPLIED：岩石学}
\hyperref[q:h24-reg-applied-29]{H24:29}, \hyperref[q:h27-reg-applied-28]{H27:28}, \hyperref[q:h30-reg-applied-30]{H30:30}, \hyperref[q:r01-reg-applied-24]{R01:24}, \hyperref[q:r01-ret-applied-31]{R01-再:31}, \hyperref[q:r02-reg-applied-25]{R02:25}, \hyperref[q:r03-reg-applied-23]{R03:23}, \hyperref[q:r03-reg-applied-24]{R03:24}, \hyperref[q:r05-reg-applied-24]{R05:24}, \hyperref[q:r05-reg-applied-26]{R05:26}, \hyperref[q:r06-reg-applied-26]{R06:26}, \hyperref[q:r06-reg-applied-27]{R06:27}, \hyperref[q:r07-reg-applied-23]{R07:23}
\section{APPLIED：応用地質・地盤}
\hyperref[q:h23-reg-applied-33]{H23:33}, \hyperref[q:h23-reg-applied-34]{H23:34}, \hyperref[q:h24-reg-applied-28]{H24:28}, \hyperref[q:h24-reg-applied-30]{H24:30}, \hyperref[q:h24-reg-applied-31]{H24:31}, \hyperref[q:h25-reg-applied-19]{H25:19}, \hyperref[q:h25-reg-applied-28]{H25:28}, \hyperref[q:h25-reg-applied-30]{H25:30}, \hyperref[q:h26-reg-applied-33]{H26:33}, \hyperref[q:h26-reg-applied-34]{H26:34}, \hyperref[q:h27-reg-applied-29]{H27:29}, \hyperref[q:h27-reg-applied-34]{H27:34}, \hyperref[q:h27-reg-applied-35]{H27:35}, \hyperref[q:h28-reg-applied-24]{H28:24}, \hyperref[q:h28-reg-applied-33]{H28:33}, \hyperref[q:h29-reg-applied-22]{H29:22}, \hyperref[q:h29-reg-applied-24]{H29:24}, \hyperref[q:h30-reg-applied-24]{H30:24}, \hyperref[q:h30-reg-applied-27]{H30:27}, \hyperref[q:r01-reg-applied-29]{R01:29}, \hyperref[q:r01-ret-applied-26]{R01-再:26}, \hyperref[q:r01-ret-applied-27]{R01-再:27}, \hyperref[q:r01-ret-applied-28]{R01-再:28}, \hyperref[q:r03-reg-applied-28]{R03:28}, \hyperref[q:r03-reg-applied-29]{R03:29}, \hyperref[q:r04-reg-applied-27]{R04:27}, \hyperref[q:r06-reg-applied-18]{R06:18}
\section{APPLIED：振動・波動}
\hyperref[q:h23-reg-applied-03]{H23:3}, \hyperref[q:h24-reg-applied-02]{H24:2}, \hyperref[q:h25-reg-applied-02]{H25:2}, \hyperref[q:h25-reg-applied-05]{H25:5}, \hyperref[q:h27-reg-applied-01]{H27:1}, \hyperref[q:h29-reg-applied-01]{H29:1}, \hyperref[q:h29-reg-applied-04]{H29:4}, \hyperref[q:h30-reg-applied-05]{H30:5}, \hyperref[q:r01-reg-applied-05]{R01:5}, \hyperref[q:r04-reg-applied-05]{R04:5}, \hyperref[q:r05-reg-applied-05]{R05:5}, \hyperref[q:r07-reg-applied-02]{R07:2}, \hyperref[q:r07-reg-applied-05]{R07:5}
\section{APPLIED：有機・高分子化学}
\hyperref[q:h23-reg-applied-09]{H23:9}, \hyperref[q:h23-reg-applied-14]{H23:14}, \hyperref[q:h23-reg-applied-15]{H23:15}, \hyperref[q:h23-reg-applied-16]{H23:16}, \hyperref[q:h24-reg-applied-15]{H24:15}, \hyperref[q:h24-reg-applied-18]{H24:18}, \hyperref[q:h25-reg-applied-14]{H25:14}, \hyperref[q:h25-reg-applied-17]{H25:17}, \hyperref[q:h26-reg-applied-17]{H26:17}, \hyperref[q:h27-reg-applied-13]{H27:13}, \hyperref[q:h27-reg-applied-14]{H27:14}, \hyperref[q:h27-reg-applied-16]{H27:16}, \hyperref[q:h28-reg-applied-15]{H28:15}, \hyperref[q:h28-reg-applied-17]{H28:17}, \hyperref[q:h29-reg-applied-15]{H29:15}, \hyperref[q:h29-reg-applied-16]{H29:16}, \hyperref[q:h29-reg-applied-17]{H29:17}, \hyperref[q:h30-reg-applied-15]{H30:15}, \hyperref[q:r01-reg-applied-11]{R01:11}, \hyperref[q:r01-reg-applied-13]{R01:13}, \hyperref[q:r01-ret-applied-11]{R01-再:11}, \hyperref[q:r01-ret-applied-13]{R01-再:13}, \hyperref[q:r01-ret-applied-17]{R01-再:17}, \hyperref[q:r02-reg-applied-15]{R02:15}, \hyperref[q:r02-reg-applied-18]{R02:18}, \hyperref[q:r03-reg-applied-12]{R03:12}, \hyperref[q:r04-reg-applied-16]{R04:16}, \hyperref[q:r04-reg-applied-18]{R04:18}, \hyperref[q:r05-reg-applied-11]{R05:11}, \hyperref[q:r05-reg-applied-14]{R05:14}, \hyperref[q:r06-reg-applied-17]{R06:17}, \hyperref[q:r07-reg-applied-11]{R07:11}, \hyperref[q:r07-reg-applied-12]{R07:12}
\section{APPLIED：材料物性}
\hyperref[q:r02-reg-applied-08]{R02:8}
\section{APPLIED：構造地質}
\hyperref[q:h25-reg-applied-20]{H25:20}, \hyperref[q:h25-reg-applied-21]{H25:21}, \hyperref[q:h26-reg-applied-35]{H26:35}, \hyperref[q:h30-reg-applied-23]{H30:23}, \hyperref[q:h30-reg-applied-26]{H30:26}, \hyperref[q:r01-reg-applied-31]{R01:31}, \hyperref[q:r01-ret-applied-30]{R01-再:30}, \hyperref[q:r02-reg-applied-27]{R02:27}, \hyperref[q:r02-reg-applied-30]{R02:30}, \hyperref[q:r03-reg-applied-27]{R03:27}, \hyperref[q:r04-reg-applied-29]{R04:29}, \hyperref[q:r04-reg-applied-30]{R04:30}, \hyperref[q:r04-reg-applied-31]{R04:31}, \hyperref[q:r05-reg-applied-27]{R05:27}, \hyperref[q:r05-reg-applied-28]{R05:28}, \hyperref[q:r05-reg-applied-29]{R05:29}, \hyperref[q:r06-reg-applied-28]{R06:28}, \hyperref[q:r06-reg-applied-29]{R06:29}, \hyperref[q:r07-reg-applied-29]{R07:29}, \hyperref[q:r07-reg-applied-30]{R07:30}
\section{APPLIED：気象・大気}
\hyperref[q:h23-reg-applied-20]{H23:20}, \hyperref[q:h24-reg-applied-34]{H24:34}, \hyperref[q:h26-reg-applied-26]{H26:26}, \hyperref[q:h27-reg-applied-24]{H27:24}, \hyperref[q:h27-reg-applied-26]{H27:26}, \hyperref[q:h28-reg-applied-25]{H28:25}, \hyperref[q:h28-reg-applied-26]{H28:26}, \hyperref[q:h29-reg-applied-19]{H29:19}, \hyperref[q:h29-reg-applied-21]{H29:21}, \hyperref[q:h30-reg-applied-20]{H30:20}, \hyperref[q:h30-reg-applied-22]{H30:22}, \hyperref[q:r01-reg-applied-19]{R01:19}, \hyperref[q:r01-reg-applied-21]{R01:21}, \hyperref[q:r01-reg-applied-22]{R01:22}, \hyperref[q:r01-ret-applied-20]{R01-再:20}, \hyperref[q:r01-ret-applied-21]{R01-再:21}, \hyperref[q:r02-reg-applied-21]{R02:21}, \hyperref[q:r02-reg-applied-22]{R02:22}, \hyperref[q:r04-reg-applied-22]{R04:22}, \hyperref[q:r05-reg-applied-18]{R05:18}, \hyperref[q:r05-reg-applied-20]{R05:20}, \hyperref[q:r06-reg-applied-20]{R06:20}, \hyperref[q:r07-reg-applied-21]{R07:21}
\section{APPLIED：海洋}
\hyperref[q:h30-reg-applied-19]{H30:19}, \hyperref[q:r01-reg-applied-18]{R01:18}, \hyperref[q:r01-ret-applied-19]{R01-再:19}, \hyperref[q:r02-reg-applied-23]{R02:23}, \hyperref[q:r04-reg-applied-21]{R04:21}, \hyperref[q:r05-reg-applied-19]{R05:19}
\section{APPLIED：火山・マグマ}
\hyperref[q:h25-reg-applied-25]{H25:25}, \hyperref[q:h25-reg-applied-26]{H25:26}, \hyperref[q:h28-reg-applied-28]{H28:28}, \hyperref[q:h28-reg-applied-29]{H28:29}, \hyperref[q:h28-reg-applied-35]{H28:35}, \hyperref[q:h29-reg-applied-23]{H29:23}, \hyperref[q:h29-reg-applied-29]{H29:29}, \hyperref[q:h30-reg-applied-29]{H30:29}, \hyperref[q:r01-reg-applied-23]{R01:23}, \hyperref[q:r01-reg-applied-25]{R01:25}, \hyperref[q:r01-ret-applied-22]{R01-再:22}, \hyperref[q:r01-ret-applied-23]{R01-再:23}, \hyperref[q:r01-ret-applied-24]{R01-再:24}, \hyperref[q:r02-reg-applied-24]{R02:24}, \hyperref[q:r03-reg-applied-26]{R03:26}, \hyperref[q:r05-reg-applied-25]{R05:25}, \hyperref[q:r06-reg-applied-25]{R06:25}, \hyperref[q:r07-reg-applied-20]{R07:20}, \hyperref[q:r07-reg-applied-26]{R07:26}
\section{APPLIED：無機化学}
\hyperref[q:h23-reg-applied-11]{H23:11}, \hyperref[q:h23-reg-applied-13]{H23:13}, \hyperref[q:h23-reg-applied-17]{H23:17}, \hyperref[q:h24-reg-applied-12]{H24:12}, \hyperref[q:h25-reg-applied-13]{H25:13}, \hyperref[q:h26-reg-applied-12]{H26:12}, \hyperref[q:h27-reg-applied-10]{H27:10}, \hyperref[q:h27-reg-applied-12]{H27:12}, \hyperref[q:h28-reg-applied-11]{H28:11}, \hyperref[q:h28-reg-applied-13]{H28:13}, \hyperref[q:h29-reg-applied-08]{H29:8}, \hyperref[q:h29-reg-applied-10]{H29:10}, \hyperref[q:h30-reg-applied-10]{H30:10}, \hyperref[q:r01-reg-applied-09]{R01:9}, \hyperref[q:r01-ret-applied-09]{R01-再:9}, \hyperref[q:r01-ret-applied-16]{R01-再:16}, \hyperref[q:r02-reg-applied-14]{R02:14}, \hyperref[q:r03-reg-applied-13]{R03:13}, \hyperref[q:r04-reg-applied-13]{R04:13}, \hyperref[q:r05-reg-applied-12]{R05:12}, \hyperref[q:r06-reg-applied-09]{R06:9}, \hyperref[q:r06-reg-applied-13]{R06:13}, \hyperref[q:r07-reg-applied-15]{R07:15}, \hyperref[q:r07-reg-applied-16]{R07:16}, \hyperref[q:r07-reg-applied-17]{R07:17}
\section{APPLIED：熱力学}
\hyperref[q:h23-reg-applied-05]{H23:5}, \hyperref[q:h24-reg-applied-08]{H24:8}, \hyperref[q:h25-reg-applied-08]{H25:8}, \hyperref[q:h25-reg-applied-09]{H25:9}, \hyperref[q:h26-reg-applied-06]{H26:6}, \hyperref[q:h26-reg-applied-08]{H26:8}, \hyperref[q:h27-reg-applied-08]{H27:8}, \hyperref[q:h28-reg-applied-08]{H28:8}, \hyperref[q:h29-reg-applied-06]{H29:6}, \hyperref[q:h30-reg-applied-06]{H30:6}, \hyperref[q:h30-reg-applied-07]{H30:7}, \hyperref[q:r01-reg-applied-06]{R01:6}, \hyperref[q:r01-ret-applied-04]{R01-再:4}, \hyperref[q:r02-reg-applied-06]{R02:6}, \hyperref[q:r03-reg-applied-07]{R03:7}, \hyperref[q:r04-reg-applied-06]{R04:6}, \hyperref[q:r04-reg-applied-07]{R04:7}, \hyperref[q:r04-reg-applied-11]{R04:11}, \hyperref[q:r04-reg-applied-17]{R04:17}, \hyperref[q:r05-reg-applied-06]{R05:6}, \hyperref[q:r06-reg-applied-06]{R06:6}, \hyperref[q:r07-reg-applied-06]{R07:6}
\section{APPLIED：物理化学}
\hyperref[q:h23-reg-applied-12]{H23:12}, \hyperref[q:h24-reg-applied-10]{H24:10}, \hyperref[q:h24-reg-applied-11]{H24:11}, \hyperref[q:h24-reg-applied-13]{H24:13}, \hyperref[q:h24-reg-applied-14]{H24:14}, \hyperref[q:h24-reg-applied-17]{H24:17}, \hyperref[q:h25-reg-applied-11]{H25:11}, \hyperref[q:h25-reg-applied-12]{H25:12}, \hyperref[q:h26-reg-applied-10]{H26:10}, \hyperref[q:h26-reg-applied-13]{H26:13}, \hyperref[q:h26-reg-applied-15]{H26:15}, \hyperref[q:h27-reg-applied-11]{H27:11}, \hyperref[q:h27-reg-applied-15]{H27:15}, \hyperref[q:h27-reg-applied-17]{H27:17}, \hyperref[q:h28-reg-applied-12]{H28:12}, \hyperref[q:h28-reg-applied-18]{H28:18}, \hyperref[q:h29-reg-applied-13]{H29:13}, \hyperref[q:h29-reg-applied-14]{H29:14}, \hyperref[q:h30-reg-applied-08]{H30:8}, \hyperref[q:h30-reg-applied-11]{H30:11}, \hyperref[q:h30-reg-applied-16]{H30:16}, \hyperref[q:h30-reg-applied-17]{H30:17}, \hyperref[q:h30-reg-applied-18]{H30:18}, \hyperref[q:r01-reg-applied-10]{R01:10}, \hyperref[q:r01-reg-applied-15]{R01:15}, \hyperref[q:r01-reg-applied-16]{R01:16}, \hyperref[q:r01-ret-applied-10]{R01-再:10}, \hyperref[q:r01-ret-applied-12]{R01-再:12}, \hyperref[q:r01-ret-applied-15]{R01-再:15}, \hyperref[q:r02-reg-applied-13]{R02:13}, \hyperref[q:r02-reg-applied-16]{R02:16}, \hyperref[q:r02-reg-applied-17]{R02:17}, \hyperref[q:r03-reg-applied-10]{R03:10}, \hyperref[q:r03-reg-applied-11]{R03:11}, \hyperref[q:r03-reg-applied-15]{R03:15}, \hyperref[q:r03-reg-applied-17]{R03:17}, \hyperref[q:r04-reg-applied-10]{R04:10}, \hyperref[q:r04-reg-applied-14]{R04:14}, \hyperref[q:r04-reg-applied-15]{R04:15}, \hyperref[q:r05-reg-applied-16]{R05:16}, \hyperref[q:r05-reg-applied-17]{R05:17}, \hyperref[q:r06-reg-applied-10]{R06:10}, \hyperref[q:r06-reg-applied-11]{R06:11}, \hyperref[q:r06-reg-applied-12]{R06:12}, \hyperref[q:r06-reg-applied-15]{R06:15}, \hyperref[q:r07-reg-applied-10]{R07:10}, \hyperref[q:r07-reg-applied-13]{R07:13}, \hyperref[q:r07-reg-applied-18]{R07:18}
\section{APPLIED：物理探査}
\hyperref[q:h23-reg-applied-28]{H23:28}, \hyperref[q:h23-reg-applied-29]{H23:29}, \hyperref[q:h23-reg-applied-30]{H23:30}, \hyperref[q:h23-reg-applied-31]{H23:31}, \hyperref[q:h23-reg-applied-32]{H23:32}, \hyperref[q:h24-reg-applied-19]{H24:19}, \hyperref[q:h24-reg-applied-20]{H24:20}, \hyperref[q:h24-reg-applied-21]{H24:21}, \hyperref[q:h24-reg-applied-22]{H24:22}, \hyperref[q:h24-reg-applied-23]{H24:23}, \hyperref[q:h25-reg-applied-32]{H25:32}, \hyperref[q:h25-reg-applied-33]{H25:33}, \hyperref[q:h25-reg-applied-34]{H25:34}, \hyperref[q:h25-reg-applied-35]{H25:35}, \hyperref[q:h26-reg-applied-19]{H26:19}, \hyperref[q:h26-reg-applied-20]{H26:20}, \hyperref[q:h26-reg-applied-21]{H26:21}, \hyperref[q:h26-reg-applied-22]{H26:22}, \hyperref[q:h26-reg-applied-23]{H26:23}, \hyperref[q:h27-reg-applied-18]{H27:18}, \hyperref[q:h27-reg-applied-19]{H27:19}, \hyperref[q:h27-reg-applied-20]{H27:20}, \hyperref[q:h27-reg-applied-21]{H27:21}, \hyperref[q:h27-reg-applied-22]{H27:22}, \hyperref[q:h28-reg-applied-19]{H28:19}, \hyperref[q:h28-reg-applied-20]{H28:20}, \hyperref[q:h28-reg-applied-21]{H28:21}, \hyperref[q:h28-reg-applied-22]{H28:22}, \hyperref[q:h28-reg-applied-23]{H28:23}, \hyperref[q:h29-reg-applied-31]{H29:31}, \hyperref[q:h29-reg-applied-32]{H29:32}, \hyperref[q:h29-reg-applied-33]{H29:33}, \hyperref[q:h29-reg-applied-34]{H29:34}, \hyperref[q:h29-reg-applied-35]{H29:35}, \hyperref[q:h30-reg-applied-32]{H30:32}, \hyperref[q:h30-reg-applied-33]{H30:33}, \hyperref[q:h30-reg-applied-34]{H30:34}, \hyperref[q:h30-reg-applied-35]{H30:35}, \hyperref[q:r01-reg-applied-32]{R01:32}, \hyperref[q:r01-reg-applied-33]{R01:33}, \hyperref[q:r01-reg-applied-34]{R01:34}, \hyperref[q:r01-reg-applied-35]{R01:35}, \hyperref[q:r01-ret-applied-32]{R01-再:32}, \hyperref[q:r02-reg-applied-32]{R02:32}, \hyperref[q:r02-reg-applied-33]{R02:33}, \hyperref[q:r02-reg-applied-34]{R02:34}, \hyperref[q:r02-reg-applied-35]{R02:35}, \hyperref[q:r03-reg-applied-32]{R03:32}, \hyperref[q:r03-reg-applied-33]{R03:33}, \hyperref[q:r03-reg-applied-34]{R03:34}, \hyperref[q:r03-reg-applied-35]{R03:35}, \hyperref[q:r04-reg-applied-32]{R04:32}, \hyperref[q:r04-reg-applied-33]{R04:33}, \hyperref[q:r04-reg-applied-34]{R04:34}, \hyperref[q:r04-reg-applied-35]{R04:35}, \hyperref[q:r05-reg-applied-32]{R05:32}, \hyperref[q:r05-reg-applied-33]{R05:33}, \hyperref[q:r05-reg-applied-34]{R05:34}, \hyperref[q:r05-reg-applied-35]{R05:35}, \hyperref[q:r06-reg-applied-32]{R06:32}, \hyperref[q:r06-reg-applied-33]{R06:33}, \hyperref[q:r06-reg-applied-34]{R06:34}, \hyperref[q:r06-reg-applied-35]{R06:35}, \hyperref[q:r07-reg-applied-31]{R07:31}, \hyperref[q:r07-reg-applied-32]{R07:32}, \hyperref[q:r07-reg-applied-33]{R07:33}, \hyperref[q:r07-reg-applied-34]{R07:34}, \hyperref[q:r07-reg-applied-35]{R07:35}
\section{APPLIED：現代物理・原子核}
\hyperref[q:h23-reg-applied-10]{H23:10}, \hyperref[q:h24-reg-applied-07]{H24:7}, \hyperref[q:h25-reg-applied-07]{H25:7}, \hyperref[q:h26-reg-applied-09]{H26:9}, \hyperref[q:h27-reg-applied-05]{H27:5}, \hyperref[q:h27-reg-applied-06]{H27:6}, \hyperref[q:h28-reg-applied-09]{H28:9}, \hyperref[q:r01-ret-applied-05]{R01-再:5}, \hyperref[q:r04-reg-applied-09]{R04:9}, \hyperref[q:r05-reg-applied-07]{R05:7}
\section{APPLIED：資源・鉱床}
\hyperref[q:h23-reg-applied-23]{H23:23}, \hyperref[q:h23-reg-applied-26]{H23:26}, \hyperref[q:h23-reg-applied-27]{H23:27}, \hyperref[q:h24-reg-applied-25]{H24:25}, \hyperref[q:h24-reg-applied-26]{H24:26}, \hyperref[q:h25-reg-applied-24]{H25:24}, \hyperref[q:h26-reg-applied-29]{H26:29}, \hyperref[q:h27-reg-applied-31]{H27:31}, \hyperref[q:h28-reg-applied-31]{H28:31}, \hyperref[q:h29-reg-applied-28]{H29:28}, \hyperref[q:r02-reg-applied-26]{R02:26}, \hyperref[q:r03-reg-applied-20]{R03:20}, \hyperref[q:r04-reg-applied-24]{R04:24}, \hyperref[q:r05-reg-applied-23]{R05:23}, \hyperref[q:r06-reg-applied-24]{R06:24}, \hyperref[q:r07-reg-applied-25]{R07:25}
\section{APPLIED：重力・測地}
\hyperref[q:h27-reg-applied-33]{H27:33}
\section{APPLIED：鉱物学}
\hyperref[q:h24-reg-applied-24]{H24:24}, \hyperref[q:h25-reg-applied-18]{H25:18}, \hyperref[q:h25-reg-applied-22]{H25:22}, \hyperref[q:h25-reg-applied-23]{H25:23}, \hyperref[q:h26-reg-applied-28]{H26:28}, \hyperref[q:h27-reg-applied-30]{H27:30}, \hyperref[q:h29-reg-applied-27]{H29:27}, \hyperref[q:h30-reg-applied-31]{H30:31}, \hyperref[q:r01-reg-applied-26]{R01:26}, \hyperref[q:r01-ret-applied-25]{R01-再:25}, \hyperref[q:r03-reg-applied-22]{R03:22}, \hyperref[q:r03-reg-applied-25]{R03:25}, \hyperref[q:r04-reg-applied-23]{R04:23}, \hyperref[q:r06-reg-applied-23]{R06:23}
\section{APPLIED：電磁気・物性}
\hyperref[q:h23-reg-applied-02]{H23:2}, \hyperref[q:h23-reg-applied-08]{H23:8}, \hyperref[q:h24-reg-applied-04]{H24:4}, \hyperref[q:h24-reg-applied-06]{H24:6}, \hyperref[q:h25-reg-applied-03]{H25:3}, \hyperref[q:h26-reg-applied-02]{H26:2}, \hyperref[q:h26-reg-applied-04]{H26:4}, \hyperref[q:h26-reg-applied-05]{H26:5}, \hyperref[q:h27-reg-applied-02]{H27:2}, \hyperref[q:h27-reg-applied-07]{H27:7}, \hyperref[q:h28-reg-applied-03]{H28:3}, \hyperref[q:h28-reg-applied-04]{H28:4}, \hyperref[q:h29-reg-applied-02]{H29:2}, \hyperref[q:h29-reg-applied-03]{H29:3}, \hyperref[q:h29-reg-applied-07]{H29:7}, \hyperref[q:h29-reg-applied-09]{H29:9}, \hyperref[q:h30-reg-applied-03]{H30:3}, \hyperref[q:h30-reg-applied-04]{H30:4}, \hyperref[q:r01-reg-applied-03]{R01:3}, \hyperref[q:r01-reg-applied-04]{R01:4}, \hyperref[q:r01-reg-applied-07]{R01:7}, \hyperref[q:r01-reg-applied-08]{R01:8}, \hyperref[q:r01-ret-applied-06]{R01-再:6}, \hyperref[q:r01-ret-applied-08]{R01-再:8}, \hyperref[q:r02-reg-applied-03]{R02:3}, \hyperref[q:r02-reg-applied-04]{R02:4}, \hyperref[q:r03-reg-applied-03]{R03:3}, \hyperref[q:r03-reg-applied-04]{R03:4}, \hyperref[q:r03-reg-applied-06]{R03:6}, \hyperref[q:r04-reg-applied-03]{R04:3}, \hyperref[q:r04-reg-applied-04]{R04:4}, \hyperref[q:r05-reg-applied-03]{R05:3}, \hyperref[q:r05-reg-applied-04]{R05:4}, \hyperref[q:r05-reg-applied-08]{R05:8}, \hyperref[q:r06-reg-applied-03]{R06:3}, \hyperref[q:r06-reg-applied-04]{R06:4}, \hyperref[q:r06-reg-applied-07]{R06:7}, \hyperref[q:r07-reg-applied-03]{R07:3}, \hyperref[q:r07-reg-applied-04]{R07:4}, \hyperref[q:r07-reg-applied-07]{R07:7}
\section{APT：AIと社会変革}
\hyperref[q:r01-ret-apt-15]{R01-再:15}, \hyperref[q:r03-reg-apt-06]{R03:6}, \hyperref[q:r04-reg-apt-04]{R04:4}, \hyperref[q:r07-reg-apt-07]{R07:7}, \hyperref[q:r07-reg-apt-09]{R07:9}
\section{APT：インフラ維持管理}
\hyperref[q:h26-reg-apt-13]{H26:13}, \hyperref[q:r01-reg-apt-08]{R01:8}, \hyperref[q:r05-reg-apt-13]{R05:13}
\section{APT：コンプライアンス・説明責任}
\hyperref[q:h23-reg-apt-09]{H23:9}, \hyperref[q:h23-reg-apt-15]{H23:15}, \hyperref[q:h24-reg-apt-03]{H24:3}, \hyperref[q:h24-reg-apt-08]{H24:8}, \hyperref[q:h24-reg-apt-09]{H24:9}, \hyperref[q:h25-reg-apt-03]{H25:3}, \hyperref[q:h25-reg-apt-04]{H25:4}, \hyperref[q:h25-reg-apt-10]{H25:10}, \hyperref[q:h25-reg-apt-13]{H25:13}, \hyperref[q:h26-reg-apt-05]{H26:5}, \hyperref[q:h26-reg-apt-06]{H26:6}, \hyperref[q:h26-reg-apt-09]{H26:9}, \hyperref[q:h26-reg-apt-14]{H26:14}, \hyperref[q:h27-reg-apt-09]{H27:9}, \hyperref[q:h27-reg-apt-10]{H27:10}, \hyperref[q:h27-reg-apt-11]{H27:11}, \hyperref[q:h27-reg-apt-12]{H27:12}, \hyperref[q:h28-reg-apt-06]{H28:6}, \hyperref[q:h28-reg-apt-10]{H28:10}, \hyperref[q:h28-reg-apt-11]{H28:11}, \hyperref[q:h28-reg-apt-15]{H28:15}, \hyperref[q:h29-reg-apt-04]{H29:4}, \hyperref[q:h29-reg-apt-05]{H29:5}, \hyperref[q:h29-reg-apt-10]{H29:10}, \hyperref[q:h30-reg-apt-06]{H30:6}, \hyperref[q:h30-reg-apt-07]{H30:7}, \hyperref[q:h30-reg-apt-08]{H30:8}, \hyperref[q:h30-reg-apt-12]{H30:12}, \hyperref[q:r01-reg-apt-05]{R01:5}, \hyperref[q:r01-reg-apt-12]{R01:12}, \hyperref[q:r01-ret-apt-08]{R01-再:8}, \hyperref[q:r01-ret-apt-09]{R01-再:9}, \hyperref[q:r02-reg-apt-04]{R02:4}, \hyperref[q:r02-reg-apt-05]{R02:5}, \hyperref[q:r02-reg-apt-13]{R02:13}, \hyperref[q:r02-reg-apt-15]{R02:15}, \hyperref[q:r03-reg-apt-04]{R03:4}, \hyperref[q:r03-reg-apt-07]{R03:7}, \hyperref[q:r03-reg-apt-13]{R03:13}, \hyperref[q:r04-reg-apt-05]{R04:5}, \hyperref[q:r04-reg-apt-08]{R04:8}, \hyperref[q:r04-reg-apt-09]{R04:9}, \hyperref[q:r04-reg-apt-12]{R04:12}, \hyperref[q:r05-reg-apt-03]{R05:3}, \hyperref[q:r05-reg-apt-04]{R05:4}, \hyperref[q:r05-reg-apt-12]{R05:12}, \hyperref[q:r06-reg-apt-03]{R06:3}, \hyperref[q:r06-reg-apt-10]{R06:10}, \hyperref[q:r06-reg-apt-12]{R06:12}, \hyperref[q:r07-reg-apt-08]{R07:8}, \hyperref[q:r07-reg-apt-11]{R07:11}, \hyperref[q:r07-reg-apt-13]{R07:13}, \hyperref[q:r07-reg-apt-15]{R07:15}
\section{APT：安全・リスク}
\hyperref[q:h23-reg-apt-05]{H23:5}, \hyperref[q:h23-reg-apt-06]{H23:6}, \hyperref[q:h23-reg-apt-08]{H23:8}, \hyperref[q:h23-reg-apt-12]{H23:12}, \hyperref[q:h24-reg-apt-06]{H24:6}, \hyperref[q:h24-reg-apt-07]{H24:7}, \hyperref[q:h24-reg-apt-10]{H24:10}, \hyperref[q:h24-reg-apt-11]{H24:11}, \hyperref[q:h25-reg-apt-05]{H25:5}, \hyperref[q:h25-reg-apt-07]{H25:7}, \hyperref[q:h25-reg-apt-08]{H25:8}, \hyperref[q:h25-reg-apt-09]{H25:9}, \hyperref[q:h25-reg-apt-14]{H25:14}, \hyperref[q:h26-reg-apt-07]{H26:7}, \hyperref[q:h26-reg-apt-08]{H26:8}, \hyperref[q:h26-reg-apt-10]{H26:10}, \hyperref[q:h26-reg-apt-11]{H26:11}, \hyperref[q:h26-reg-apt-15]{H26:15}, \hyperref[q:h27-reg-apt-13]{H27:13}, \hyperref[q:h27-reg-apt-14]{H27:14}, \hyperref[q:h28-reg-apt-05]{H28:5}, \hyperref[q:h28-reg-apt-07]{H28:7}, \hyperref[q:h28-reg-apt-08]{H28:8}, \hyperref[q:h29-reg-apt-08]{H29:8}, \hyperref[q:h29-reg-apt-09]{H29:9}, \hyperref[q:h29-reg-apt-12]{H29:12}, \hyperref[q:h30-reg-apt-09]{H30:9}, \hyperref[q:h30-reg-apt-10]{H30:10}, \hyperref[q:h30-reg-apt-11]{H30:11}, \hyperref[q:h30-reg-apt-14]{H30:14}, \hyperref[q:r01-reg-apt-03]{R01:3}, \hyperref[q:r01-reg-apt-06]{R01:6}, \hyperref[q:r01-reg-apt-07]{R01:7}, \hyperref[q:r01-reg-apt-11]{R01:11}, \hyperref[q:r01-reg-apt-13]{R01:13}, \hyperref[q:r01-ret-apt-07]{R01-再:7}, \hyperref[q:r01-ret-apt-12]{R01-再:12}, \hyperref[q:r01-ret-apt-13]{R01-再:13}, \hyperref[q:r02-reg-apt-06]{R02:6}, \hyperref[q:r02-reg-apt-07]{R02:7}, \hyperref[q:r02-reg-apt-08]{R02:8}, \hyperref[q:r02-reg-apt-09]{R02:9}, \hyperref[q:r02-reg-apt-12]{R02:12}, \hyperref[q:r03-reg-apt-08]{R03:8}, \hyperref[q:r03-reg-apt-10]{R03:10}, \hyperref[q:r03-reg-apt-12]{R03:12}, \hyperref[q:r03-reg-apt-15]{R03:15}, \hyperref[q:r04-reg-apt-06]{R04:6}, \hyperref[q:r04-reg-apt-11]{R04:11}, \hyperref[q:r05-reg-apt-06]{R05:6}, \hyperref[q:r05-reg-apt-08]{R05:8}, \hyperref[q:r05-reg-apt-09]{R05:9}, \hyperref[q:r05-reg-apt-10]{R05:10}, \hyperref[q:r05-reg-apt-14]{R05:14}, \hyperref[q:r06-reg-apt-07]{R06:7}, \hyperref[q:r06-reg-apt-13]{R06:13}, \hyperref[q:r06-reg-apt-14]{R06:14}, \hyperref[q:r07-reg-apt-12]{R07:12}, \hyperref[q:r07-reg-apt-14]{R07:14}
\section{APT：専門職・経営管理・内部統制}
\hyperref[q:h23-reg-apt-13]{H23:13}, \hyperref[q:r04-reg-apt-02]{R04:2}
\section{APT：技術士法・義務責務}
\hyperref[q:h23-reg-apt-01]{H23:1}, \hyperref[q:h23-reg-apt-04]{H23:4}, \hyperref[q:h24-reg-apt-01]{H24:1}, \hyperref[q:h24-reg-apt-14]{H24:14}, \hyperref[q:h25-reg-apt-01]{H25:1}, \hyperref[q:h26-reg-apt-01]{H26:1}, \hyperref[q:h27-reg-apt-01]{H27:1}, \hyperref[q:h27-reg-apt-02]{H27:2}, \hyperref[q:h27-reg-apt-03]{H27:3}, \hyperref[q:h28-reg-apt-01]{H28:1}, \hyperref[q:h28-reg-apt-02]{H28:2}, \hyperref[q:h29-reg-apt-01]{H29:1}, \hyperref[q:h29-reg-apt-02]{H29:2}, \hyperref[q:h30-reg-apt-01]{H30:1}, \hyperref[q:h30-reg-apt-02]{H30:2}, \hyperref[q:h30-reg-apt-03]{H30:3}, \hyperref[q:r01-reg-apt-01]{R01:1}, \hyperref[q:r01-reg-apt-02]{R01:2}, \hyperref[q:r01-ret-apt-01]{R01-再:1}, \hyperref[q:r01-ret-apt-02]{R01-再:2}, \hyperref[q:r02-reg-apt-01]{R02:1}, \hyperref[q:r03-reg-apt-01]{R03:1}, \hyperref[q:r04-reg-apt-01]{R04:1}, \hyperref[q:r04-reg-apt-15]{R04:15}, \hyperref[q:r05-reg-apt-01]{R05:1}, \hyperref[q:r05-reg-apt-05]{R05:5}, \hyperref[q:r06-reg-apt-01]{R06:1}, \hyperref[q:r06-reg-apt-15]{R06:15}, \hyperref[q:r07-reg-apt-01]{R07:1}
\section{APT：技術者倫理・意思決定}
\hyperref[q:h23-reg-apt-11]{H23:11}, \hyperref[q:h24-reg-apt-02]{H24:2}, \hyperref[q:h24-reg-apt-12]{H24:12}, \hyperref[q:h24-reg-apt-13]{H24:13}, \hyperref[q:h25-reg-apt-02]{H25:2}, \hyperref[q:h25-reg-apt-06]{H25:6}, \hyperref[q:h25-reg-apt-12]{H25:12}, \hyperref[q:h25-reg-apt-15]{H25:15}, \hyperref[q:h26-reg-apt-12]{H26:12}, \hyperref[q:h27-reg-apt-04]{H27:4}, \hyperref[q:h27-reg-apt-05]{H27:5}, \hyperref[q:h27-reg-apt-06]{H27:6}, \hyperref[q:h27-reg-apt-15]{H27:15}, \hyperref[q:h28-reg-apt-03]{H28:3}, \hyperref[q:h28-reg-apt-09]{H28:9}, \hyperref[q:h28-reg-apt-12]{H28:12}, \hyperref[q:h29-reg-apt-03]{H29:3}, \hyperref[q:h29-reg-apt-07]{H29:7}, \hyperref[q:h29-reg-apt-11]{H29:11}, \hyperref[q:h29-reg-apt-13]{H29:13}, \hyperref[q:h29-reg-apt-15]{H29:15}, \hyperref[q:h30-reg-apt-04]{H30:4}, \hyperref[q:h30-reg-apt-05]{H30:5}, \hyperref[q:h30-reg-apt-15]{H30:15}, \hyperref[q:r01-reg-apt-14]{R01:14}, \hyperref[q:r01-ret-apt-03]{R01-再:3}, \hyperref[q:r01-ret-apt-04]{R01-再:4}, \hyperref[q:r01-ret-apt-05]{R01-再:5}, \hyperref[q:r02-reg-apt-02]{R02:2}, \hyperref[q:r02-reg-apt-11]{R02:11}, \hyperref[q:r03-reg-apt-02]{R03:2}, \hyperref[q:r03-reg-apt-09]{R03:9}, \hyperref[q:r04-reg-apt-03]{R04:3}, \hyperref[q:r04-reg-apt-07]{R04:7}, \hyperref[q:r05-reg-apt-11]{R05:11}, \hyperref[q:r06-reg-apt-04]{R06:4}, \hyperref[q:r06-reg-apt-11]{R06:11}, \hyperref[q:r07-reg-apt-02]{R07:2}, \hyperref[q:r07-reg-apt-03]{R07:3}, \hyperref[q:r07-reg-apt-04]{R07:4}, \hyperref[q:r07-reg-apt-06]{R07:6}
\section{APT：環境・持続可能性}
\hyperref[q:h23-reg-apt-14]{H23:14}, \hyperref[q:h24-reg-apt-05]{H24:5}, \hyperref[q:h25-reg-apt-11]{H25:11}, \hyperref[q:h28-reg-apt-14]{H28:14}, \hyperref[q:h30-reg-apt-13]{H30:13}, \hyperref[q:r01-reg-apt-15]{R01:15}, \hyperref[q:r01-ret-apt-11]{R01-再:11}, \hyperref[q:r01-ret-apt-14]{R01-再:14}, \hyperref[q:r02-reg-apt-10]{R02:10}, \hyperref[q:r02-reg-apt-14]{R02:14}, \hyperref[q:r03-reg-apt-05]{R03:5}, \hyperref[q:r03-reg-apt-11]{R03:11}, \hyperref[q:r04-reg-apt-10]{R04:10}, \hyperref[q:r04-reg-apt-14]{R04:14}, \hyperref[q:r05-reg-apt-15]{R05:15}, \hyperref[q:r06-reg-apt-02]{R06:2}, \hyperref[q:r06-reg-apt-05]{R06:5}, \hyperref[q:r06-reg-apt-06]{R06:6}, \hyperref[q:r07-reg-apt-10]{R07:10}
\section{APT：研究・情報倫理}
\hyperref[q:h23-reg-apt-02]{H23:2}, \hyperref[q:h23-reg-apt-03]{H23:3}, \hyperref[q:h23-reg-apt-07]{H23:7}, \hyperref[q:h23-reg-apt-10]{H23:10}, \hyperref[q:h24-reg-apt-04]{H24:4}, \hyperref[q:h24-reg-apt-15]{H24:15}, \hyperref[q:h26-reg-apt-02]{H26:2}, \hyperref[q:h26-reg-apt-03]{H26:3}, \hyperref[q:h26-reg-apt-04]{H26:4}, \hyperref[q:h27-reg-apt-07]{H27:7}, \hyperref[q:h27-reg-apt-08]{H27:8}, \hyperref[q:h28-reg-apt-04]{H28:4}, \hyperref[q:h28-reg-apt-13]{H28:13}, \hyperref[q:h29-reg-apt-06]{H29:6}, \hyperref[q:h29-reg-apt-14]{H29:14}, \hyperref[q:r01-reg-apt-04]{R01:4}, \hyperref[q:r01-reg-apt-09]{R01:9}, \hyperref[q:r01-reg-apt-10]{R01:10}, \hyperref[q:r01-ret-apt-06]{R01-再:6}, \hyperref[q:r01-ret-apt-10]{R01-再:10}, \hyperref[q:r02-reg-apt-03]{R02:3}, \hyperref[q:r03-reg-apt-03]{R03:3}, \hyperref[q:r03-reg-apt-14]{R03:14}, \hyperref[q:r04-reg-apt-13]{R04:13}, \hyperref[q:r05-reg-apt-02]{R05:2}, \hyperref[q:r05-reg-apt-07]{R05:7}, \hyperref[q:r06-reg-apt-08]{R06:8}, \hyperref[q:r06-reg-apt-09]{R06:9}, \hyperref[q:r07-reg-apt-05]{R07:5}
\section{BASIC：アルゴリズム・計算量}
\hyperref[q:h23-reg-basic-2-2]{H23:2-2}, \hyperref[q:h24-reg-basic-2-5]{H24:2-5}, \hyperref[q:h25-reg-basic-2-2]{H25:2-2}, \hyperref[q:h26-reg-basic-2-2]{H26:2-2}, \hyperref[q:h26-reg-basic-2-6]{H26:2-6}, \hyperref[q:h29-reg-basic-2-3]{H29:2-3}, \hyperref[q:r01-ret-basic-2-2]{R01-再:2-2}, \hyperref[q:r03-reg-basic-2-6]{R03:2-6}, \hyperref[q:r05-reg-basic-2-2]{R05:2-2}, \hyperref[q:r06-reg-basic-2-5]{R06:2-5}, \hyperref[q:r07-reg-basic-2-6]{R07:2-6}
\section{BASIC：エネルギー}
\hyperref[q:h23-reg-basic-5-3]{H23:5-3}, \hyperref[q:h24-reg-basic-5-2]{H24:5-2}, \hyperref[q:h24-reg-basic-5-3]{H24:5-3}, \hyperref[q:h25-reg-basic-5-1]{H25:5-1}, \hyperref[q:h25-reg-basic-5-2]{H25:5-2}, \hyperref[q:h26-reg-basic-5-4]{H26:5-4}, \hyperref[q:h27-reg-basic-5-3]{H27:5-3}, \hyperref[q:h27-reg-basic-5-4]{H27:5-4}, \hyperref[q:h28-reg-basic-5-3]{H28:5-3}, \hyperref[q:h28-reg-basic-5-4]{H28:5-4}, \hyperref[q:h29-reg-basic-5-3]{H29:5-3}, \hyperref[q:h29-reg-basic-5-4]{H29:5-4}, \hyperref[q:h30-reg-basic-5-3]{H30:5-3}, \hyperref[q:h30-reg-basic-5-4]{H30:5-4}, \hyperref[q:r01-reg-basic-5-3]{R01:5-3}, \hyperref[q:r01-reg-basic-5-4]{R01:5-4}, \hyperref[q:r01-ret-basic-5-3]{R01-再:5-3}, \hyperref[q:r01-ret-basic-5-4]{R01-再:5-4}, \hyperref[q:r02-reg-basic-5-3]{R02:5-3}, \hyperref[q:r02-reg-basic-5-4]{R02:5-4}, \hyperref[q:r03-reg-basic-5-3]{R03:5-3}, \hyperref[q:r03-reg-basic-5-4]{R03:5-4}, \hyperref[q:r04-reg-basic-5-3]{R04:5-3}, \hyperref[q:r04-reg-basic-5-4]{R04:5-4}, \hyperref[q:r05-reg-basic-5-3]{R05:5-3}, \hyperref[q:r05-reg-basic-5-4]{R05:5-4}, \hyperref[q:r06-reg-basic-5-4]{R06:5-4}, \hyperref[q:r07-reg-basic-5-3]{R07:5-3}, \hyperref[q:r07-reg-basic-5-4]{R07:5-4}
\section{BASIC：データベース・AI}
\hyperref[q:r01-reg-basic-2-3]{R01:2-3}
\section{BASIC：データ構造・形式言語}
\hyperref[q:h24-reg-basic-2-2]{H24:2-2}, \hyperref[q:h26-reg-basic-2-4]{H26:2-4}, \hyperref[q:h26-reg-basic-2-5]{H26:2-5}, \hyperref[q:h27-reg-basic-2-6]{H27:2-6}, \hyperref[q:h30-reg-basic-2-5]{H30:2-5}, \hyperref[q:r01-reg-basic-2-2]{R01:2-2}, \hyperref[q:r01-reg-basic-2-4]{R01:2-4}, \hyperref[q:r01-reg-basic-2-6]{R01:2-6}, \hyperref[q:r03-reg-basic-2-5]{R03:2-5}
\section{BASIC：ネットワーク}
\hyperref[q:r03-reg-basic-2-3]{R03:2-3}, \hyperref[q:r06-reg-basic-2-3]{R06:2-3}, \hyperref[q:r07-reg-basic-2-3]{R07:2-3}, \hyperref[q:r07-reg-basic-2-5]{R07:2-5}
\section{BASIC：プロジェクト・最適化}
\hyperref[q:h23-reg-basic-1-2]{H23:1-2}, \hyperref[q:h23-reg-basic-1-3]{H23:1-3}, \hyperref[q:h25-reg-basic-1-3]{H25:1-3}, \hyperref[q:h25-reg-basic-1-4]{H25:1-4}, \hyperref[q:h25-reg-basic-1-5]{H25:1-5}, \hyperref[q:h26-reg-basic-2-1]{H26:2-1}, \hyperref[q:h27-reg-basic-1-2]{H27:1-2}, \hyperref[q:h28-reg-basic-1-4]{H28:1-4}, \hyperref[q:h28-reg-basic-1-5]{H28:1-5}, \hyperref[q:h29-reg-basic-1-1]{H29:1-1}, \hyperref[q:h30-reg-basic-1-2]{H30:1-2}, \hyperref[q:h30-reg-basic-1-4]{H30:1-4}, \hyperref[q:r01-reg-basic-1-5]{R01:1-5}, \hyperref[q:r01-ret-basic-1-2]{R01-再:1-2}, \hyperref[q:r01-ret-basic-1-4]{R01-再:1-4}, \hyperref[q:r02-reg-basic-1-4]{R02:1-4}, \hyperref[q:r03-reg-basic-1-3]{R03:1-3}, \hyperref[q:r06-reg-basic-1-5]{R06:1-5}, \hyperref[q:r07-reg-basic-1-2]{R07:1-2}
\section{BASIC：ベクトル解析・幾何}
\hyperref[q:h23-reg-basic-3-5]{H23:3-5}, \hyperref[q:h25-reg-basic-3-6]{H25:3-6}, \hyperref[q:h26-reg-basic-3-1]{H26:3-1}, \hyperref[q:h27-reg-basic-3-2]{H27:3-2}, \hyperref[q:h29-reg-basic-3-2]{H29:3-2}, \hyperref[q:h30-reg-basic-3-2]{H30:3-2}, \hyperref[q:r01-reg-basic-3-1]{R01:3-1}, \hyperref[q:r01-ret-basic-3-2]{R01-再:3-2}, \hyperref[q:r02-reg-basic-3-1]{R02:3-1}, \hyperref[q:r02-reg-basic-3-2]{R02:3-2}, \hyperref[q:r03-reg-basic-3-1]{R03:3-1}, \hyperref[q:r03-reg-basic-3-6]{R03:3-6}, \hyperref[q:r04-reg-basic-3-2]{R04:3-2}, \hyperref[q:r06-reg-basic-3-2]{R06:3-2}, \hyperref[q:r07-reg-basic-3-2]{R07:3-2}
\section{BASIC：信頼性設計}
\hyperref[q:h23-reg-basic-1-1]{H23:1-1}, \hyperref[q:h24-reg-basic-1-1]{H24:1-1}, \hyperref[q:h26-reg-basic-1-4]{H26:1-4}, \hyperref[q:h26-reg-basic-1-5]{H26:1-5}, \hyperref[q:h27-reg-basic-1-1]{H27:1-1}, \hyperref[q:h28-reg-basic-1-1]{H28:1-1}, \hyperref[q:h30-reg-basic-1-1]{H30:1-1}, \hyperref[q:r01-ret-basic-1-3]{R01-再:1-3}, \hyperref[q:r01-ret-basic-1-6]{R01-再:1-6}, \hyperref[q:r02-reg-basic-1-6]{R02:1-6}, \hyperref[q:r03-reg-basic-1-2]{R03:1-2}, \hyperref[q:r03-reg-basic-1-4]{R03:1-4}, \hyperref[q:r05-reg-basic-1-4]{R05:1-4}, \hyperref[q:r05-reg-basic-1-5]{R05:1-5}
\section{BASIC：力学設計}
\hyperref[q:r06-reg-basic-1-4]{R06:1-4}, \hyperref[q:r07-reg-basic-1-4]{R07:1-4}
\section{BASIC：化学・物質}
\hyperref[q:r01-reg-basic-4-4]{R01:4-4}, \hyperref[q:r01-ret-basic-4-1]{R01-再:4-1}, \hyperref[q:r01-ret-basic-4-2]{R01-再:4-2}, \hyperref[q:r05-reg-basic-4-2]{R05:4-2}
\section{BASIC：原子・無機化学}
\hyperref[q:h23-reg-basic-4-2]{H23:4-2}, \hyperref[q:h23-reg-basic-4-3]{H23:4-3}, \hyperref[q:h24-reg-basic-4-2]{H24:4-2}, \hyperref[q:h26-reg-basic-4-1]{H26:4-1}, \hyperref[q:h26-reg-basic-4-2]{H26:4-2}, \hyperref[q:h27-reg-basic-4-1]{H27:4-1}, \hyperref[q:h27-reg-basic-4-2]{H27:4-2}, \hyperref[q:h28-reg-basic-4-2]{H28:4-2}, \hyperref[q:h29-reg-basic-4-1]{H29:4-1}, \hyperref[q:h30-reg-basic-4-1]{H30:4-1}, \hyperref[q:r01-reg-basic-4-1]{R01:4-1}, \hyperref[q:r01-reg-basic-4-2]{R01:4-2}, \hyperref[q:r02-reg-basic-4-1]{R02:4-1}, \hyperref[q:r03-reg-basic-4-1]{R03:4-1}, \hyperref[q:r04-reg-basic-4-2]{R04:4-2}, \hyperref[q:r05-reg-basic-4-1]{R05:4-1}, \hyperref[q:r06-reg-basic-4-1]{R06:4-1}, \hyperref[q:r06-reg-basic-4-2]{R06:4-2}, \hyperref[q:r07-reg-basic-4-1]{R07:4-1}, \hyperref[q:r07-reg-basic-4-4]{R07:4-4}
\section{BASIC：品質・統計・計測}
\hyperref[q:h23-reg-basic-5-4]{H23:5-4}, \hyperref[q:h24-reg-basic-5-4]{H24:5-4}
\section{BASIC：品質・統計設計}
\hyperref[q:h24-reg-basic-1-3]{H24:1-3}, \hyperref[q:h25-reg-basic-1-6]{H25:1-6}, \hyperref[q:h26-reg-basic-1-3]{H26:1-3}, \hyperref[q:h27-reg-basic-1-5]{H27:1-5}, \hyperref[q:h27-reg-basic-1-6]{H27:1-6}, \hyperref[q:h28-reg-basic-1-2]{H28:1-2}, \hyperref[q:h30-reg-basic-1-5]{H30:1-5}, \hyperref[q:r02-reg-basic-1-2]{R02:1-2}, \hyperref[q:r04-reg-basic-1-2]{R04:1-2}, \hyperref[q:r04-reg-basic-1-3]{R04:1-3}, \hyperref[q:r06-reg-basic-1-6]{R06:1-6}, \hyperref[q:r07-reg-basic-1-5]{R07:1-5}, \hyperref[q:r07-reg-basic-1-6]{R07:1-6}
\section{BASIC：安全・強度設計}
\hyperref[q:h23-reg-basic-1-5]{H23:1-5}, \hyperref[q:h24-reg-basic-1-2]{H24:1-2}, \hyperref[q:h25-reg-basic-1-2]{H25:1-2}, \hyperref[q:h27-reg-basic-1-3]{H27:1-3}, \hyperref[q:h28-reg-basic-1-3]{H28:1-3}, \hyperref[q:h29-reg-basic-1-2]{H29:1-2}, \hyperref[q:h29-reg-basic-1-4]{H29:1-4}, \hyperref[q:h29-reg-basic-1-6]{H29:1-6}, \hyperref[q:r01-reg-basic-1-4]{R01:1-4}, \hyperref[q:r01-ret-basic-1-5]{R01-再:1-5}, \hyperref[q:r02-reg-basic-1-3]{R02:1-3}, \hyperref[q:r03-reg-basic-1-5]{R03:1-5}, \hyperref[q:r04-reg-basic-1-1]{R04:1-1}, \hyperref[q:r04-reg-basic-1-4]{R04:1-4}, \hyperref[q:r04-reg-basic-1-5]{R04:1-5}, \hyperref[q:r05-reg-basic-1-1]{R05:1-1}, \hyperref[q:r05-reg-basic-1-2]{R05:1-2}, \hyperref[q:r05-reg-basic-1-3]{R05:1-3}, \hyperref[q:r06-reg-basic-1-2]{R06:1-2}, \hyperref[q:r06-reg-basic-1-3]{R06:1-3}, \hyperref[q:r07-reg-basic-1-3]{R07:1-3}
\section{BASIC：微積分・数値解析}
\hyperref[q:h24-reg-basic-3-3]{H24:3-3}, \hyperref[q:h24-reg-basic-3-4]{H24:3-4}, \hyperref[q:h26-reg-basic-3-2]{H26:3-2}, \hyperref[q:h26-reg-basic-3-3]{H26:3-3}, \hyperref[q:h27-reg-basic-3-1]{H27:3-1}, \hyperref[q:h27-reg-basic-3-3]{H27:3-3}, \hyperref[q:h28-reg-basic-3-1]{H28:3-1}, \hyperref[q:h28-reg-basic-3-2]{H28:3-2}, \hyperref[q:h28-reg-basic-3-3]{H28:3-3}, \hyperref[q:h28-reg-basic-3-4]{H28:3-4}, \hyperref[q:h29-reg-basic-3-1]{H29:3-1}, \hyperref[q:h30-reg-basic-3-1]{H30:3-1}, \hyperref[q:h30-reg-basic-3-4]{H30:3-4}, \hyperref[q:r01-reg-basic-3-2]{R01:3-2}, \hyperref[q:r01-ret-basic-3-1]{R01-再:3-1}, \hyperref[q:r01-ret-basic-3-3]{R01-再:3-3}, \hyperref[q:r01-ret-basic-3-4]{R01-再:3-4}, \hyperref[q:r02-reg-basic-3-3]{R02:3-3}, \hyperref[q:r03-reg-basic-3-2]{R03:3-2}, \hyperref[q:r04-reg-basic-3-1]{R04:3-1}, \hyperref[q:r04-reg-basic-3-3]{R04:3-3}, \hyperref[q:r05-reg-basic-3-2]{R05:3-2}, \hyperref[q:r05-reg-basic-3-3]{R05:3-3}, \hyperref[q:r06-reg-basic-3-1]{R06:3-1}, \hyperref[q:r06-reg-basic-3-3]{R06:3-3}, \hyperref[q:r07-reg-basic-3-3]{R07:3-3}
\section{BASIC：情報セキュリティ}
\hyperref[q:h23-reg-basic-2-5]{H23:2-5}, \hyperref[q:h27-reg-basic-2-5]{H27:2-5}, \hyperref[q:h29-reg-basic-2-1]{H29:2-1}, \hyperref[q:h30-reg-basic-2-1]{H30:2-1}, \hyperref[q:r01-ret-basic-2-1]{R01-再:2-1}, \hyperref[q:r02-reg-basic-2-3]{R02:2-3}, \hyperref[q:r03-reg-basic-2-1]{R03:2-1}, \hyperref[q:r04-reg-basic-2-1]{R04:2-1}, \hyperref[q:r05-reg-basic-2-1]{R05:2-1}, \hyperref[q:r06-reg-basic-2-6]{R06:2-6}, \hyperref[q:r07-reg-basic-2-4]{R07:2-4}
\section{BASIC：情報理論・符号・確率モデル}
\hyperref[q:h24-reg-basic-2-3]{H24:2-3}, \hyperref[q:h24-reg-basic-2-4]{H24:2-4}, \hyperref[q:h26-reg-basic-2-3]{H26:2-3}, \hyperref[q:h28-reg-basic-2-1]{H28:2-1}, \hyperref[q:h28-reg-basic-2-5]{H28:2-5}, \hyperref[q:r01-reg-basic-2-5]{R01:2-5}, \hyperref[q:r02-reg-basic-2-1]{R02:2-1}, \hyperref[q:r04-reg-basic-2-4]{R04:2-4}, \hyperref[q:r05-reg-basic-2-3]{R05:2-3}, \hyperref[q:r05-reg-basic-2-4]{R05:2-4}, \hyperref[q:r06-reg-basic-2-4]{R06:2-4}, \hyperref[q:r07-reg-basic-2-2]{R07:2-2}
\section{BASIC：意思決定・最適化}
\hyperref[q:h29-reg-basic-1-3]{H29:1-3}, \hyperref[q:r01-reg-basic-1-1]{R01:1-1}, \hyperref[q:r04-reg-basic-1-6]{R04:1-6}
\section{BASIC：振動・運動}
\hyperref[q:h24-reg-basic-3-2]{H24:3-2}, \hyperref[q:h26-reg-basic-3-5]{H26:3-5}, \hyperref[q:r01-reg-basic-3-6]{R01:3-6}, \hyperref[q:r01-ret-basic-3-5]{R01-再:3-5}, \hyperref[q:r02-reg-basic-3-5]{R02:3-5}, \hyperref[q:r04-reg-basic-3-6]{R04:3-6}, \hyperref[q:r06-reg-basic-3-5]{R06:3-5}
\section{BASIC：数値表現・誤差}
\hyperref[q:h23-reg-basic-2-1]{H23:2-1}, \hyperref[q:h24-reg-basic-2-1]{H24:2-1}, \hyperref[q:h25-reg-basic-2-1]{H25:2-1}, \hyperref[q:h25-reg-basic-2-5]{H25:2-5}, \hyperref[q:h27-reg-basic-2-1]{H27:2-1}, \hyperref[q:h27-reg-basic-2-3]{H27:2-3}, \hyperref[q:h27-reg-basic-2-4]{H27:2-4}, \hyperref[q:h28-reg-basic-2-3]{H28:2-3}, \hyperref[q:h28-reg-basic-2-6]{H28:2-6}, \hyperref[q:h29-reg-basic-2-2]{H29:2-2}, \hyperref[q:h30-reg-basic-2-3]{H30:2-3}, \hyperref[q:r01-reg-basic-2-1]{R01:2-1}, \hyperref[q:r01-ret-basic-2-3]{R01-再:2-3}, \hyperref[q:r01-ret-basic-2-4]{R01-再:2-4}, \hyperref[q:r02-reg-basic-2-4]{R02:2-4}, \hyperref[q:r02-reg-basic-2-5]{R02:2-5}, \hyperref[q:r04-reg-basic-2-5]{R04:2-5}, \hyperref[q:r04-reg-basic-2-6]{R04:2-6}, \hyperref[q:r06-reg-basic-2-1]{R06:2-1}
\section{BASIC：最適化・経済性}
\hyperref[q:h23-reg-basic-1-4]{H23:1-4}, \hyperref[q:h24-reg-basic-1-5]{H24:1-5}, \hyperref[q:h26-reg-basic-1-2]{H26:1-2}, \hyperref[q:r01-reg-basic-1-2]{R01:1-2}
\section{BASIC：材料・金属}
\hyperref[q:h23-reg-basic-4-4]{H23:4-4}, \hyperref[q:h24-reg-basic-4-3]{H24:4-3}, \hyperref[q:h24-reg-basic-4-4]{H24:4-4}, \hyperref[q:h25-reg-basic-4-3]{H25:4-3}, \hyperref[q:h25-reg-basic-4-4]{H25:4-4}, \hyperref[q:h26-reg-basic-4-3]{H26:4-3}, \hyperref[q:h26-reg-basic-4-4]{H26:4-4}, \hyperref[q:h27-reg-basic-4-4]{H27:4-4}, \hyperref[q:h28-reg-basic-4-3]{H28:4-3}, \hyperref[q:h28-reg-basic-4-4]{H28:4-4}, \hyperref[q:h29-reg-basic-4-3]{H29:4-3}, \hyperref[q:h29-reg-basic-4-4]{H29:4-4}, \hyperref[q:h30-reg-basic-4-3]{H30:4-3}, \hyperref[q:h30-reg-basic-4-4]{H30:4-4}, \hyperref[q:r01-ret-basic-4-4]{R01-再:4-4}, \hyperref[q:r02-reg-basic-4-3]{R02:4-3}, \hyperref[q:r02-reg-basic-4-4]{R02:4-4}, \hyperref[q:r03-reg-basic-4-3]{R03:4-3}, \hyperref[q:r03-reg-basic-4-4]{R03:4-4}, \hyperref[q:r04-reg-basic-4-3]{R04:4-3}, \hyperref[q:r04-reg-basic-4-4]{R04:4-4}, \hyperref[q:r05-reg-basic-4-3]{R05:4-3}, \hyperref[q:r05-reg-basic-4-4]{R05:4-4}, \hyperref[q:r06-reg-basic-4-3]{R06:4-3}, \hyperref[q:r06-reg-basic-4-4]{R06:4-4}, \hyperref[q:r07-reg-basic-4-3]{R07:4-3}
\section{BASIC：物理モデル・回路}
\hyperref[q:h29-reg-basic-3-4]{H29:3-4}, \hyperref[q:r02-reg-basic-3-6]{R02:3-6}, \hyperref[q:r05-reg-basic-3-6]{R05:3-6}, \hyperref[q:r06-reg-basic-3-6]{R06:3-6}, \hyperref[q:r07-reg-basic-3-5]{R07:3-5}, \hyperref[q:r07-reg-basic-3-6]{R07:3-6}
\section{BASIC：物理化学・反応}
\hyperref[q:h23-reg-basic-4-1]{H23:4-1}, \hyperref[q:h24-reg-basic-4-1]{H24:4-1}, \hyperref[q:h25-reg-basic-4-1]{H25:4-1}, \hyperref[q:h25-reg-basic-4-2]{H25:4-2}, \hyperref[q:h25-reg-basic-4-6]{H25:4-6}, \hyperref[q:h27-reg-basic-4-3]{H27:4-3}, \hyperref[q:h28-reg-basic-4-1]{H28:4-1}, \hyperref[q:h29-reg-basic-4-2]{H29:4-2}, \hyperref[q:h30-reg-basic-4-2]{H30:4-2}, \hyperref[q:r01-reg-basic-4-3]{R01:4-3}, \hyperref[q:r01-ret-basic-4-3]{R01-再:4-3}, \hyperref[q:r02-reg-basic-4-2]{R02:4-2}, \hyperref[q:r02-reg-basic-4-5]{R02:4-5}, \hyperref[q:r03-reg-basic-4-2]{R03:4-2}, \hyperref[q:r04-reg-basic-4-1]{R04:4-1}, \hyperref[q:r07-reg-basic-4-2]{R07:4-2}
\section{BASIC：状態遷移・形式モデル}
\hyperref[q:h30-reg-basic-2-2]{H30:2-2}
\section{BASIC：環境・公害}
\hyperref[q:h25-reg-basic-5-4]{H25:5-4}, \hyperref[q:h26-reg-basic-5-1]{H26:5-1}, \hyperref[q:r01-reg-basic-5-1]{R01:5-1}, \hyperref[q:r03-reg-basic-5-2]{R03:5-2}, \hyperref[q:r05-reg-basic-5-2]{R05:5-2}
\section{BASIC：環境・持続可能性}
\hyperref[q:h23-reg-basic-5-1]{H23:5-1}, \hyperref[q:h23-reg-basic-5-2]{H23:5-2}, \hyperref[q:h25-reg-basic-5-3]{H25:5-3}, \hyperref[q:h26-reg-basic-5-2]{H26:5-2}, \hyperref[q:h27-reg-basic-5-1]{H27:5-1}, \hyperref[q:h28-reg-basic-5-1]{H28:5-1}, \hyperref[q:h28-reg-basic-5-2]{H28:5-2}, \hyperref[q:h29-reg-basic-5-1]{H29:5-1}, \hyperref[q:h29-reg-basic-5-2]{H29:5-2}, \hyperref[q:h30-reg-basic-5-1]{H30:5-1}, \hyperref[q:h30-reg-basic-5-2]{H30:5-2}, \hyperref[q:r01-reg-basic-5-2]{R01:5-2}, \hyperref[q:r01-ret-basic-5-1]{R01-再:5-1}, \hyperref[q:r01-ret-basic-5-6]{R01-再:5-6}, \hyperref[q:r02-reg-basic-5-1]{R02:5-1}, \hyperref[q:r02-reg-basic-5-2]{R02:5-2}, \hyperref[q:r03-reg-basic-5-1]{R03:5-1}, \hyperref[q:r04-reg-basic-5-1]{R04:5-1}, \hyperref[q:r04-reg-basic-5-2]{R04:5-2}, \hyperref[q:r04-reg-basic-5-5]{R04:5-5}, \hyperref[q:r05-reg-basic-5-1]{R05:5-1}, \hyperref[q:r06-reg-basic-5-1]{R06:5-1}, \hyperref[q:r06-reg-basic-5-2]{R06:5-2}, \hyperref[q:r06-reg-basic-5-3]{R06:5-3}, \hyperref[q:r07-reg-basic-5-1]{R07:5-1}, \hyperref[q:r07-reg-basic-5-2]{R07:5-2}
\section{BASIC：生物・バイオ}
\hyperref[q:h23-reg-basic-4-5]{H23:4-5}, \hyperref[q:h24-reg-basic-4-5]{H24:4-5}, \hyperref[q:h25-reg-basic-4-5]{H25:4-5}, \hyperref[q:h26-reg-basic-4-5]{H26:4-5}, \hyperref[q:h26-reg-basic-4-6]{H26:4-6}, \hyperref[q:h27-reg-basic-4-5]{H27:4-5}, \hyperref[q:h27-reg-basic-4-6]{H27:4-6}, \hyperref[q:h28-reg-basic-4-5]{H28:4-5}, \hyperref[q:h28-reg-basic-4-6]{H28:4-6}, \hyperref[q:h29-reg-basic-4-5]{H29:4-5}, \hyperref[q:h29-reg-basic-4-6]{H29:4-6}, \hyperref[q:h30-reg-basic-4-5]{H30:4-5}, \hyperref[q:h30-reg-basic-4-6]{H30:4-6}, \hyperref[q:r01-reg-basic-4-5]{R01:4-5}, \hyperref[q:r01-reg-basic-4-6]{R01:4-6}, \hyperref[q:r01-ret-basic-4-5]{R01-再:4-5}, \hyperref[q:r01-ret-basic-4-6]{R01-再:4-6}, \hyperref[q:r02-reg-basic-4-6]{R02:4-6}, \hyperref[q:r03-reg-basic-4-5]{R03:4-5}, \hyperref[q:r03-reg-basic-4-6]{R03:4-6}, \hyperref[q:r04-reg-basic-4-5]{R04:4-5}, \hyperref[q:r04-reg-basic-4-6]{R04:4-6}, \hyperref[q:r05-reg-basic-4-5]{R05:4-5}, \hyperref[q:r05-reg-basic-4-6]{R05:4-6}, \hyperref[q:r06-reg-basic-4-5]{R06:4-5}, \hyperref[q:r06-reg-basic-4-6]{R06:4-6}, \hyperref[q:r07-reg-basic-4-5]{R07:4-5}, \hyperref[q:r07-reg-basic-4-6]{R07:4-6}
\section{BASIC：知的財産・制度}
\hyperref[q:h24-reg-basic-5-1]{H24:5-1}, \hyperref[q:h26-reg-basic-5-3]{H26:5-3}, \hyperref[q:h27-reg-basic-5-2]{H27:5-2}, \hyperref[q:h27-reg-basic-5-5]{H27:5-5}, \hyperref[q:r01-reg-basic-5-6]{R01:5-6}, \hyperref[q:r01-ret-basic-5-2]{R01-再:5-2}
\section{BASIC：科学技術史・社会}
\hyperref[q:h23-reg-basic-5-5]{H23:5-5}, \hyperref[q:h24-reg-basic-5-5]{H24:5-5}, \hyperref[q:h25-reg-basic-5-5]{H25:5-5}, \hyperref[q:h25-reg-basic-5-6]{H25:5-6}, \hyperref[q:h26-reg-basic-5-5]{H26:5-5}, \hyperref[q:h26-reg-basic-5-6]{H26:5-6}, \hyperref[q:h27-reg-basic-5-6]{H27:5-6}, \hyperref[q:h28-reg-basic-5-5]{H28:5-5}, \hyperref[q:h28-reg-basic-5-6]{H28:5-6}, \hyperref[q:h29-reg-basic-5-5]{H29:5-5}, \hyperref[q:h29-reg-basic-5-6]{H29:5-6}, \hyperref[q:h30-reg-basic-5-5]{H30:5-5}, \hyperref[q:h30-reg-basic-5-6]{H30:5-6}, \hyperref[q:r01-reg-basic-5-5]{R01:5-5}, \hyperref[q:r01-ret-basic-5-5]{R01-再:5-5}, \hyperref[q:r02-reg-basic-5-5]{R02:5-5}, \hyperref[q:r02-reg-basic-5-6]{R02:5-6}, \hyperref[q:r03-reg-basic-5-5]{R03:5-5}, \hyperref[q:r03-reg-basic-5-6]{R03:5-6}, \hyperref[q:r04-reg-basic-5-6]{R04:5-6}, \hyperref[q:r05-reg-basic-5-5]{R05:5-5}, \hyperref[q:r05-reg-basic-5-6]{R05:5-6}, \hyperref[q:r06-reg-basic-5-5]{R06:5-5}, \hyperref[q:r06-reg-basic-5-6]{R06:5-6}, \hyperref[q:r07-reg-basic-5-5]{R07:5-5}, \hyperref[q:r07-reg-basic-5-6]{R07:5-6}
\section{BASIC：線形代数・ベクトル}
\hyperref[q:h25-reg-basic-3-4]{H25:3-4}, \hyperref[q:h26-reg-basic-3-4]{H26:3-4}, \hyperref[q:h30-reg-basic-3-3]{H30:3-3}, \hyperref[q:r05-reg-basic-3-1]{R05:3-1}, \hyperref[q:r07-reg-basic-3-1]{R07:3-1}
\section{BASIC：製図}
\hyperref[q:h26-reg-basic-1-6]{H26:1-6}, \hyperref[q:h29-reg-basic-1-5]{H29:1-5}, \hyperref[q:r01-reg-basic-1-3]{R01:1-3}, \hyperref[q:r02-reg-basic-1-5]{R02:1-5}
\section{BASIC：言語・データ構造}
\hyperref[q:h29-reg-basic-2-5]{H29:2-5}
\section{BASIC：計算力学}
\hyperref[q:h23-reg-basic-3-1]{H23:3-1}, \hyperref[q:h23-reg-basic-3-3]{H23:3-3}, \hyperref[q:h24-reg-basic-3-1]{H24:3-1}, \hyperref[q:h25-reg-basic-3-3]{H25:3-3}, \hyperref[q:h29-reg-basic-3-3]{H29:3-3}, \hyperref[q:r02-reg-basic-3-4]{R02:3-4}, \hyperref[q:r03-reg-basic-3-3]{R03:3-3}
\section{BASIC：計算機アーキテクチャ}
\hyperref[q:h23-reg-basic-2-4]{H23:2-4}, \hyperref[q:h25-reg-basic-2-6]{H25:2-6}, \hyperref[q:h28-reg-basic-2-4]{H28:2-4}, \hyperref[q:h29-reg-basic-2-6]{H29:2-6}, \hyperref[q:r02-reg-basic-2-6]{R02:2-6}, \hyperref[q:r04-reg-basic-2-3]{R04:2-3}
\section{BASIC：設計の応用数学}
\hyperref[q:r01-reg-basic-1-6]{R01:1-6}, \hyperref[q:r01-ret-basic-1-1]{R01-再:1-1}, \hyperref[q:r05-reg-basic-1-6]{R05:1-6}
\section{BASIC：設計原理・管理}
\hyperref[q:h24-reg-basic-1-4]{H24:1-4}, \hyperref[q:h25-reg-basic-1-1]{H25:1-1}, \hyperref[q:h26-reg-basic-1-1]{H26:1-1}, \hyperref[q:h27-reg-basic-1-4]{H27:1-4}, \hyperref[q:h28-reg-basic-1-6]{H28:1-6}, \hyperref[q:h30-reg-basic-1-3]{H30:1-3}, \hyperref[q:h30-reg-basic-1-6]{H30:1-6}, \hyperref[q:r02-reg-basic-1-1]{R02:1-1}, \hyperref[q:r03-reg-basic-1-1]{R03:1-1}, \hyperref[q:r03-reg-basic-1-6]{R03:1-6}, \hyperref[q:r06-reg-basic-1-1]{R06:1-1}, \hyperref[q:r07-reg-basic-1-1]{R07:1-1}
\section{BASIC：論理・離散数学}
\hyperref[q:h23-reg-basic-2-3]{H23:2-3}, \hyperref[q:h25-reg-basic-2-3]{H25:2-3}, \hyperref[q:h25-reg-basic-2-4]{H25:2-4}, \hyperref[q:h27-reg-basic-2-2]{H27:2-2}, \hyperref[q:h28-reg-basic-2-2]{H28:2-2}, \hyperref[q:h29-reg-basic-2-4]{H29:2-4}, \hyperref[q:h30-reg-basic-2-4]{H30:2-4}, \hyperref[q:h30-reg-basic-2-6]{H30:2-6}, \hyperref[q:r01-ret-basic-2-5]{R01-再:2-5}, \hyperref[q:r01-ret-basic-2-6]{R01-再:2-6}, \hyperref[q:r02-reg-basic-2-2]{R02:2-2}, \hyperref[q:r03-reg-basic-2-2]{R03:2-2}, \hyperref[q:r03-reg-basic-2-4]{R03:2-4}, \hyperref[q:r04-reg-basic-2-2]{R04:2-2}, \hyperref[q:r05-reg-basic-2-5]{R05:2-5}, \hyperref[q:r05-reg-basic-2-6]{R05:2-6}, \hyperref[q:r06-reg-basic-2-2]{R06:2-2}, \hyperref[q:r07-reg-basic-2-1]{R07:2-1}
\section{BASIC：連続体・構造力学}
\hyperref[q:h23-reg-basic-3-2]{H23:3-2}, \hyperref[q:h23-reg-basic-3-4]{H23:3-4}, \hyperref[q:h24-reg-basic-3-5]{H24:3-5}, \hyperref[q:h25-reg-basic-3-1]{H25:3-1}, \hyperref[q:h25-reg-basic-3-2]{H25:3-2}, \hyperref[q:h25-reg-basic-3-5]{H25:3-5}, \hyperref[q:h26-reg-basic-3-6]{H26:3-6}, \hyperref[q:h27-reg-basic-3-5]{H27:3-5}, \hyperref[q:h27-reg-basic-3-6]{H27:3-6}, \hyperref[q:h28-reg-basic-3-5]{H28:3-5}, \hyperref[q:h28-reg-basic-3-6]{H28:3-6}, \hyperref[q:h29-reg-basic-3-5]{H29:3-5}, \hyperref[q:h29-reg-basic-3-6]{H29:3-6}, \hyperref[q:h30-reg-basic-3-6]{H30:3-6}, \hyperref[q:r01-reg-basic-3-4]{R01:3-4}, \hyperref[q:r01-reg-basic-3-5]{R01:3-5}, \hyperref[q:r01-ret-basic-3-6]{R01-再:3-6}, \hyperref[q:r03-reg-basic-3-4]{R03:3-4}, \hyperref[q:r04-reg-basic-3-4]{R04:3-4}, \hyperref[q:r05-reg-basic-3-4]{R05:3-4}, \hyperref[q:r06-reg-basic-3-4]{R06:3-4}, \hyperref[q:r07-reg-basic-3-4]{R07:3-4}
\section{BASIC：運動・エネルギー}
\hyperref[q:h27-reg-basic-3-4]{H27:3-4}, \hyperref[q:h30-reg-basic-3-5]{H30:3-5}, \hyperref[q:r01-reg-basic-3-3]{R01:3-3}, \hyperref[q:r03-reg-basic-3-5]{R03:3-5}, \hyperref[q:r04-reg-basic-3-5]{R04:3-5}, \hyperref[q:r05-reg-basic-3-5]{R05:3-5}

\chapter{方法と確認状態}
CSVを正本とし、索引・頻度表・問題ブロックを派生出力する。本文の手書き部分と自動生成部分は保存先を分ける。
生OCRは未校正の抽出値として保持し、校正済みの条件・選択肢対応・図参照を別レコードに置く。
確認証跡には原典ハッシュと対象内容のハッシュを記録する。原典が変われば旧確認を引き継がない。
\section{統計の限界}
暫定分類による数と、確認済みの部分集合の数を並べる。確認済みの部分集合は全出題頻度の推定値ではない。
難易度・学習時間が空欄なら未評価であり、0や易問として処理しない。
\section{完成の条件}
実際の校正・解説、他肢の検証、引用・参照、TeXビルド、専門査読が必要である。
ラベルの存在とコンパイルの成功だけでは、教材の内容品質や完成を認定しない。

\chapter{作業状態と監査範囲}
登録 1270 問、主分類確認済み 1270 問、要旨確認済み 1270 問、解説初稿 1270 問。
\par 初稿には未確定事項を明記し、専門査読を待つ。初期監査と最終監査を分け、最終完成は未判定。

\printbibliography[heading=bibintoc,title={参照資料}]
\end{document}
